% paper12/matrix_spencer_lifts.tex
% Clean revision of the owner-supplied draft, with theorem labels, display
% delimiters, references, and the multilevel Toeplitz statement repaired.
%
\documentclass[11pt]{amsart}

\usepackage[T1]{fontenc}
\usepackage[utf8]{inputenc}
\usepackage{lmodern}
\usepackage{microtype}
\usepackage{amsmath,amssymb,amsthm,mathtools}
\usepackage{mathrsfs}
\usepackage{enumitem}
\usepackage[colorlinks=true,linkcolor=blue,citecolor=blue,urlcolor=blue]{hyperref}

\numberwithin{equation}{section}

\newtheorem{theorem}{Theorem}[section]
\newtheorem{proposition}[theorem]{Proposition}
\newtheorem{corollary}[theorem]{Corollary}
\newtheorem{lemma}[theorem]{Lemma}
\newtheorem{question}[theorem]{Question}
\theoremstyle{definition}
\newtheorem{definition}[theorem]{Definition}
\theoremstyle{remark}
\newtheorem{remark}[theorem]{Remark}

\newcommand{\C}{\mathbb C}
\newcommand{\R}{\mathbb R}
\newcommand{\Z}{\mathbb Z}
\newcommand{\T}{\mathbb T}
\newcommand{\M}{\mathrm M}
\newcommand{\Id}{I}
\newcommand{\op}{\mathrm{op}}
\newcommand{\Span}{\operatorname{span}}
\newcommand{\disc}{\operatorname{disc}}
\newcommand{\discop}{\operatorname{disc}_{\op}}
\newcommand{\Ent}{\mathsf E}
\newcommand{\LiftEnt}{\mathsf L}
\newcommand{\supp}{\operatorname{supp}}
\newcommand{\diag}{\operatorname{diag}}
\newcommand{\wh}{\widehat}
\newcommand{\norm}[1]{\left\lVert #1\right\rVert}
\newcommand{\abs}[1]{\left\lvert #1\right\rvert}
\newcommand{\eps}{\varepsilon}
\newcommand{\mcA}{\mathcal A}
\newcommand{\mcB}{\mathcal B}
\newcommand{\mcF}{\mathcal F}
\newcommand{\mcG}{\mathcal G}
\newcommand{\lesssimC}{\lesssim}

\DeclareMathOperator*{\argmin}{arg\,min}

\title[Matrix Spencer bounds from finite-entropy lifts]{Matrix Spencer Bounds from Finite-Entropy Lifts}
\author{Yitzchak Shmalo}
\address{Einstein Institute of Mathematics, The Hebrew University of Jerusalem, Jerusalem, Israel}
\date{June 2026}

\subjclass[2020]{Primary 15A60, 11K38; Secondary 15B05, 42A05, 46L07, 47B35, 68Q25}
\keywords{Matrix discrepancy, Matrix Spencer conjecture, finite-dimensional C*-algebras, Toeplitz matrices, Hankel matrices, group convolution, block diagonal matrices}

\begin{document}

\begin{abstract}
The Matrix Spencer conjecture asks whether every collection of self-adjoint
$n\times n$ contractions $A_1,\ldots,A_n$ admits signs
$\eps_i\in\{-1,1\}$ for which
$\|\sum_i \eps_i A_i\|_{\op}=O(\sqrt n)$. We formulate a reusable
finite-lift principle for structured instances. Let
\[
\mcB=\bigoplus_{\alpha=1}^m \M_{r_\alpha}(\C),\qquad
h(\mcB)=\max_\alpha r_\alpha,\qquad
R(\mcB)=\sum_\alpha r_\alpha,
\]
and set $\Ent(\mcB)=h(\mcB)R(\mcB)$. If a tuple of matrices has a
$K$-bounded lift to $\mcB$ through a contraction on the lifted span, then
\[
\discop(A_1,\ldots,A_n)
\le C K \sqrt{n\log\left(e+\frac{\Ent(\mcB)}{n}\right)} .
\]
This is a direct consequence of the block-diagonal Matrix Spencer theorem of
Dadush, Jiang and Reis. The same argument also covers matrix-valued
multiplier models after sampling on a finite norming set.

The framework gives deterministic Matrix Spencer bounds for scalar Toeplitz
and Hankel contractions, fixed block-size multilevel Toeplitz and Hankel classes
at any fixed level (full bandwidth, with level-dependent constants) together
with dimension-free constants and a matching lower bound in the banded
case, matrices in a common finite-dimensional C*-algebra, finite-group
convolution algebras with controlled irreducible dimensions, and the usual
compression and bounded-image operations. The paper also records the correct entropy accounting
for finite sums: for direct sums one has
$\Ent(\bigoplus_\ell \mcB_\ell)= (\max_\ell h(\mcB_\ell))\sum_\ell R(\mcB_\ell)$,
not generally $\sum_\ell \Ent(\mcB_\ell)$.
The same direct-sum bookkeeping gives a finite-menu common-sign corollary:
many finite-lift tuples have one signing whenever their combined lift entropy is
\(O(n)\), and this supplies deterministic common core signings for the later
active-prefix random-amplitude perturbation menus. A packaged corollary therefore
covers finite-lift core libraries plus localized random-amplitude matrix noise
under the same union-count and energy schedules. Finally, an elementary dimension
argument shows that any exact, or even fixed-error approximate faithful, finite
block lift whose image covers the full unit ball of $\M_D(\C)$ has entropy at least $D^2$;
the associated lift-entropy invariant records the exact complexity seen by this
finite-lift engine and explains why the full-unit-ball finite-lift route cannot
by itself settle the square case of the full Matrix Spencer conjecture. We also record the June 2026
algebraic theorem of Akbas--Sra, which improves the exact finite-dimensional
C*-algebra input from the elementary $hR$ parameter to intrinsic dimension
$\dim_\C\mcA$; this gives a sharper AF-stage corollary when the approximating
stages have $\dim_\C A_k=O(n)$.  Their low-Frobenius perturbation theorem also
gives a matrix-representation AF corollary: an AF-stage algebraic core of
intrinsic dimension \(O(n)\), perturbed by self-adjoint tails with
\(\|L_i\|_{\mathrm F}^2\le n/\log^3(en)\) -- in particular by
operator-norm-one tails of rank at most \(n/\log^3(en)\) -- still has
discrepancy \(O(\sqrt n)\) up to the AF approximation error. The same
perturbative input applies to the square group-convolution regime: if
\(|G|=O(n)\), arbitrary self-adjoint convolution contractions remain at
discrepancy \(O(\sqrt n)\) after adding self-adjoint low-Frobenius, hence
low-rank, perturbations.  The same deterministic perturbative input is packaged
as a distribution-free random-matrix corollary: with the same probability as the
structural event, arbitrary random algebraic cores of intrinsic dimension
\(O(n)\), including full \(D\times D\) random self-adjoint models for
\(D\lesssim\sqrt n\), remain Spencer-scale after arbitrary dependent random
low-Frobenius, low-rank, spiked, compressed, or coordinate-supported
self-adjoint tails.  A separate average-case corollary covers full centered,
possibly dependent, random matrix ensembles in every ambient dimension.  First,
a distribution-free all-plus gauge says that arbitrary self-adjoint or
rectangular random matrices become a Matrix-Spencer-scale instance after
pathwise scaling by the larger of the largest individual operator norm and the
all-plus norm divided by \(\sqrt n\); deterministic scale choices only have to
dominate this gauge eventually, and moment/quantile schedules make that check
explicit from \(L^p\) gauge bounds with no independence assumptions, including
closed-form all-integer and sparse-dimension thresholds when the active menu is
polynomial in the ambient dimension; direct sub-Weibull gauge tails give the
same almost-sure conclusion with logarithmic normalizations.  The same input
also gives finite-confidence, finite-menu certificates with explicit
\(\delta\)-dependent scales, including conditional versions whose scales are
chosen from side information via conditional moment or tail envelopes; random
active finite-confidence menus pay only their selected conditional error
budget.  The same conditional budgets also give adaptive almost-sure
triangular-array schedules by Borel--Cantelli, and random active subcatalogues
pay only their selected conditional budget.  Second,
when their scalar \(\ell_2\)-stable
sums have the usual Gaussian operator-norm tail, deterministic scaling by
\(K(\sqrt D+\sqrt{\log(en)})\) makes the all-plus signing itself give
discrepancy at most \(\sqrt n\).  A direct active-menu
form makes this simultaneous for finite or countable libraries under the
summability of \(J_m(n_m+1)e^{-\tau_m^2}\), without any sign-symmetry
assumption; random-active self-adjoint and rectangular versions pay only the
expected selected operator-tail cost.  A quadratic-form net criterion verifies the same self-adjoint
active-menu input from scalar tails for
\(\langle u,\sum_i a_iY_i u\rangle\), and Hermitian sub-Gaussian atom series
are controlled by their atom square functions.  For independent upper-triangular
Hermitian entries this proxy is an explicit quadratic-form entry square
function, bounded by \(\sqrt2\,\sigma\) in the homogeneous case and by row
variance budgets in sparse or banded profiles.  A Bernstein quadratic-form
variant handles centered rank-one covariance and averaged Wishart blocks at scale
\(\kappa^2((\sqrt D+\tau)/\sqrt r+(\sqrt D+\tau)^2/r)\), and Hermitian
independent entries with Bernstein variance and jump proxies at scale
\(Q(\sqrt D+\tau)+B(\sqrt D+\tau)^2\).  The same all-plus active-menu
argument applies to rectangular or
non-self-adjoint full operator-tail matrices with \(D=d_1+d_2\).  A bilinear-form
net lemma verifies the rectangular hypothesis from scalar sub-Gaussian bilinear
forms, giving the same active schedules and random-active selected-cost
schedules at scale \(q(\sqrt{d_1+d_2}+\tau)\).  Haar unitary and orthogonal matrices have
\(q=O(D^{-1/2})\), hence need only a universal normalization when
\(\tau=O(\sqrt D)\).  For independent-entry rectangular arrays the
bilinear proxy is the worst weighted entry \(\ell_2\) scale, bounded in
particular by the maximum entry sub-Gaussian proxy.  More generally, rectangular
sub-Gaussian coefficient series are covered with proxy given by the worst
fixed-bilinear atom energy, bounded by the left/right atom variance operators.
Bernstein bilinear-form tails give the same rectangular all-plus route, with
the same selected-cost random-active form, at scale
\(q(\sqrt{d_1+d_2}+\tau)+b(\sqrt{d_1+d_2}+\tau)^2\), covering rectangular
bounded, Poissonized, exponential-entry, and gamma-entry libraries, as well as
balanced Haar rank-one dyads after only constant normalization.
The selected-cost random-active Bernstein form is also made entrywise
checkable for independent rectangular arrays, and ordinary conditionally
independent entrywise Bernstein mgf proxies imply that checkable form.
A one-row finite-confidence version gives the same explicit normalizations for
single finite random-matrix libraries, and a bounded-entry specialization
reduces the hypotheses to conditional second-moment and uniform entry bounds.
A bounded matrix-Bernstein/Freedman all-plus criterion also removes sign
symmetry for independent or adapted centered self-adjoint and rectangular blocks
whenever their variance budgets make the all-plus tail summable; the independent
bounded matrix-Bernstein and adapted bounded matrix-Freedman sides also have
one-row finite-confidence forms.  A quantum-effect specialization applies this
directly to centered random POVM effects or density matrices, with variance
proxy \(\|\sum_i\mathbb E(E_i\mid\mathcal F)\|_{\rm op}\).
A proxy-scale form separates the typical variance scale from the largest jump,
so heterogeneous bounded streams normalize by \(\max\{R,K\}q\) once their jumps
are at most \(Rq\).
A tail-capped Bernstein/Freedman form then admits unbounded independent or
adapted centered self-adjoint and rectangular streams after summable norm-cap
failures, with only a truncated-tail drift budget and a closed-form cap margin;
its finite-confidence selected-library form pays the selected cap failures and
the selected matrix Bernstein/Freedman tails in one row.
An explicit finite-confidence sub-Weibull operator-norm selector chooses caps
from the requested confidence level, while the asymptotic selector uses the
same logarithmic cap and variance scale to give all-integer and sparse
active-menu schedules for independent or adapted matrix streams with
conditional norm tails.  A finite-moment finite-confidence companion and its
asymptotic schedule replace exponential cap tails by Markov cap tails and give
the corresponding polynomial normalization for arbitrary conditional
operator-norm moment tails.  A heterogeneous cap
profile form lets each increment use its own deterministic cap, so the jump
scale is a maximum cap but the variance budget still sees the aggregate
truncated second-moment profile; its finite-confidence version pays only the
selected per-increment cap budgets and the selected matrix tail.
Its finite-moment selector, including its finite-confidence form, lets the
moment order, scale, and Markov cap parameter vary by increment, with cap cost
\(\sum_iU_i^{-1}\) and variance energy
\(\sum_i s_i^2U_i^{(2-p_i)_+/p_i}\).  A tail-profile selector, again with a
finite-confidence form, then accepts arbitrary deterministic survival profiles
\(\pi_i\), using cap cost \(\sum_i\pi_i(u_i)\) and truncated variance energy
\(2\sum_i s_i^2\int_0^{u_i}t\pi_i(t)\,dt\).  A quantile-budget selector
allocates a prescribed cap-failure budget across the increments and chooses
the corresponding profile quantiles automatically.  Its mixed sub-Weibull and
finite-moment specializations give closed-form caps for increment-dependent
tail constants, exponents, and moment orders while keeping the variance scale
at the aggregate tail-energy level.  When the sub-Weibull exponent is common,
a single inverse tail-budget equation gives the optimal common cap.  When the
scales are not meant to be optimized, a closed-form envelope using
\(\max_i s_i\) and \(\sum_iA_i\) gives an immediately checkable schedule, and
is exact when the scales are common; uniformly dominated sub-Weibull streams
therefore have a direct \(A^\Sigma=nA\), \(s^{\max}=s\) plug-in formula.  When
the moment order is common, the
budget shares have an explicit optimizer in terms of the aggregate
\(p\)-moment mass; deterministic envelopes for that mass give closed-form
heavy-tail schedules under polynomial dimension growth, and uniform
increment envelopes reduce further to explicit \(n\)- and \(D\)-power
normalizations.  A mixed active-library form lets sub-Weibull and finite-moment
dominated streams coexist in the same row under one cap budget, and the same
selector transfers to already-signed core libraries after per-stream scaling.
The raw noncentered form instead keeps matrices \(W=B+Y\) together and
normalizes by \(\max\{1,\lambda\}\), so structured conditional-mean cores absorb
mixed dominated centered fluctuations directly.  Its scaled variant lets the
structured mean carry an environment-measurable size \(\rho\) and normalizes by
\(\max\{\rho,\lambda\}\).  A two-scale variant further decouples the structured
row size from the common-signing size and normalizes by
\(\max\{\rho,\sigma,\lambda\}\), with a structured specialization for finite
lifts and intrinsic algebraic cores.
For coordinatewise
sign-symmetric random tuples, a Fubini transfer shows that any all-plus
\(O(\sqrt n)\) tail is itself a discrepancy tail for the same
realization-normalized ensemble.  Its Borel--Cantelli finite-menu form gives
simultaneous eventual signings for active libraries of such sign-symmetric
models whenever their all-plus failure series is summable; for independent or
conditionally independent bounded sign-symmetric tuples this yields an explicit
matrix-Bernstein active-menu criterion in terms of the all-plus variance proxy,
with summable environmental envelope exceptions, all-integer and sparse
dimension schedules under logarithmic-square sample growth, and closed-form
discrepancy margins; symmetric scalar-multiplier environments inherit the same
explicit thresholds from their conditional second-moment profile, and
sub-Gaussian symmetric multipliers admit cap-normalized schedules with no
Bernstein jump term.  Arbitrary symmetric multipliers with summable cap
failures also admit tail-capped schedules with the same bounded-Bernstein
margins, with sub-Weibull tails and polynomial moment tails as explicit
specializations; the cap may even be chosen from the conditioning environment,
so random variance profiles can use adaptive normalizations.  More generally, any
random core tuple
carrying an environment-measurable \(K\sqrt n\) signing remains
Spencer-scale after adding full operator-tail noise whose scalar linear
combinations obey the same tail conditionally on the core; dimension-controlled
random algebraic cores are one special case.  The finite-lift common-menu
principle gives an active-prefix version for finite-lift core libraries, with
union cost \(J_m(n_m+1)e^{-\tau_m^2}\) after scaling the full noise by
\(\sqrt D+\tau_m\); the explicit schedule permits arbitrary dimension growth,
and if \(J_m(n_m+1)\le(eD_m)^\Gamma\), the all-integer margin is
\(\tau_D^2\ge R\log(eD)\) with \(R>\Gamma+1\), while sparse schedules need only
\(R>\Gamma\).  The same finite-lift signs now have a mixed dominated-tail
version: each active perturbation stream may be tagged sub-Weibull or
finite-moment and scaled by its own \(q_{m,h}\), with the same cap-budget and
matrix-tail summability thresholds as the general mixed selector.
Dimension-controlled algebraic core menus now have the same mixed dominated-tail
version: applying Akbas--Sra to the direct sum of the active core algebras gives
one common signing whenever the active-prefix intrinsic C*-dimension is
\(O(n)\).  A catalogue corollary spells this out for Toeplitz/Hankel,
banded multilevel, finite C*-algebra, and finite group-convolution core menus
under their corresponding active-prefix entropy or intrinsic-dimension budgets.
A concrete size-budget corollary makes those budgets directly checkable through
active-prefix sums of ambient dimensions, bandwidth products, C*-dimensions, or
group sizes.
An orbit-closure version then keeps the same mixed-tail conclusion under random
unitary or permutation conjugacy, corners, fixed-size tensor amplification, block
direct sums, and bounded self-adjoint images with the corresponding summed size
charge.
A selected-submenu form allows the active structured library itself to be chosen
randomly or adaptively from the row environment, provided its realized
cardinality and total charge are summably controlled.
A Bernstein-selected specialization verifies those selection failures from
independent candidate indicators with deterministic probability and charge
envelopes.
A Laplace-dominated specialization gives the same checks for negatively
dependent selections whenever the selected-cardinality and charge transforms are
dominated by the corresponding independent Bernoulli envelopes.
A random-envelope Laplace version allows those Bernoulli envelopes to be
environment-measurable, provided their aggregate cardinality, charge, and
quadratic charge profiles have deterministic summable bounds.
A random-envelope Poisson-count version handles proposal-count and occupancy
catalogues by passing from dominated compound-Poisson count profiles to the
occupied selected support.
A Cox-Poisson occupied-support refinement charges each label by
\(1-e^{-\Lambda}\), so duplicate proposals do not consume extra catalogue
budget.
A product-categorical specialization covers randomized choice of one structured
template per slot, with deterministic selected-cardinality and only a scalar
Bernstein charge tail.
An adaptive-categorical version keeps the same deterministic cardinality while
allowing sequential slot choices controlled only by predictable probability
envelopes, using a scalar Freedman charge tail.
A random-envelope adaptive-categorical version permits those predictable
envelopes to be state-dependent, once their aggregate charge drift and quadratic
charge are summably certified.
A Freedman-selected version removes conditional independence and covers
sequential adaptive indicators controlled by predictable probability envelopes.
A random-envelope Freedman-selected version further allows those predictable
envelopes to be random, after summable certification of the aggregate
cardinality, charge, and quadratic charge profiles.
A Serfling-selected specialization does the same for fixed-size and stratified
sampling without replacement, using finite-population charge averages.
A hybrid composition corollary combines Bernoulli, Laplace, Poisson/Cox,
categorical, adaptive, and without-replacement submenus in one active library,
with no independence required between mechanisms once their component
cardinality, charge, and failure budgets are summably certified.
A random-budget hybrid version allows those component budgets themselves to be
environment-random, provided their aggregate cardinality and charge are
summably certified.
A hierarchical random-budget version flattens recursive random trees, branching
catalogues, and adaptive search frontiers to their selected terminal antichains,
so only the realized frontier cardinality and charge need summable
certification.
For Gaussian,
Rademacher, or scalar sub-Gaussian matrix-series
noise with predictable coefficient variance, the same finite-lift core signs
use the sharper matrix-Laplace scale \(q(\sqrt{\log(e\kappa)}+\tau)\), where
\(\kappa\) is a rank envelope for the coefficient square function
(\(\kappa=D\) in the full-rank fallback case), instead of the generic
\(q(\sqrt D+\tau)\).  A common low-rank projection specialization makes this
checkable for spiked, block-localized, and random-subspace matrix-series noise.
A norm-budget atomic specialization reduces the remaining square-function
checks to scalar sums of squared atom amplitudes, covering rank-one spike
fields, self-adjoint edge atoms, and finite frame dictionaries in a low-rank
support.
A block-load variant improves this for orthogonal block libraries: the scale is
set by the largest block incident energy, not by the total atom energy.
A bounded weighted block-design specialization removes the load-concentration
step entirely for regular, matching, permutation, and balanced
without-replacement block designs whose weighted block degrees are deterministic
or fail only summably.  A bounded-overlap block-hypergraph corollary gives the
concrete \(q/\sqrt d\) amplitude normalization for block designs, regular block
hypergraphs, and balanced block dictionaries with overlap degree \(d\).  A
heterogeneous degree-profile version uses local row and total block degrees,
avoiding a worst-block reduction for irregular block hypergraphs.  A row-wise
block-permutation corollary gives the same normalization for \(d\)-round
permutation, matching, and balanced without-replacement block assignments.
A conditional Bernstein specialization verifies those incident-load bounds for
independent random block choices, random edge supports, and clustered
dictionaries.  A weighted conditional Bernstein variant also allows the atom
amplitudes or intensities to be random, as long as the realized incident
weights have summable conditional Bernstein envelopes.
A uniform block-support specialization gives explicit formulas when atoms
choose bounded-size random block sets, Bernoulli block thinnings, or random
block hyperedges.  A profile schedule form converts the resulting Bernstein
exponents into arbitrary-growth, all-integer, and sparse Borel--Cantelli
criteria.  In the homogeneous case these exponents reduce to the expected
number of active atom incidences per block.
A predictable block-support profile form removes atom independence for
sequentially exposed supports: one-step block inclusion probabilities and
predictable amplitude envelopes give the same scalar-Freedman budget, covering
adaptive and without-replacement block-sampled atoms.
A profile schedule for these predictable supports requires only the resulting
Freedman exponents to dominate the active block-count logarithm.
With a common predictable amplitude cap, deterministic row-block and total-block
incidence-mass bounds give a concrete inhomogeneous \(8/3\)-threshold schedule
for biased adaptive block choices.
In the homogeneous predictable uniform case, the same all-integer
\(8(\Gamma+1)/3\) and sparse \(8\Gamma/3\) incidence thresholds hold without
atom independence.
A predictable weighted variant replaces conditional independence by scalar
Freedman budgets, covering adaptive atom sampling, online random weighted
dictionaries, and without-replacement block exposure.  The weighted block-load
profile schedule turns both the Bernstein and Freedman exponents into
all-integer and sparse Borel--Cantelli criteria.
A bounded weighted incidence specialization replaces those exponents by the
normalized mean incident load \(\rho/\omega\), retaining the same \(8/3\)
threshold constants for bounded random intensities and adaptive weighted
dictionaries.
A quadratic-form criterion makes this operator-tail input
checkable from scalar conditional sub-Gaussian bounds on
\(\langle u,\sum_i a_iY_i u\rangle\), so conditionally Gaussian, correlated
Gaussian, and Wigner-type noise libraries enter the same finite-lift menu
theorem at scale \(q(\sqrt D+\tau)\).  An entrywise specialization bounds this
proxy by the upper-triangular
sub-Gaussian square function; in particular uniform entry scale \(\sigma\)
gives the usual \(\sigma(\sqrt D+\tau)\) spectral normalization for
GOE/GUE/Wigner-type active-prefix libraries.  A row-variance envelope further
covers sparse or banded entry profiles at scale
\((\alpha^2+2\Delta\beta^2)^{1/2}(\sqrt D+\tau)\), hence normalized sparse
profiles with \(\beta\lesssim\Delta^{-1/2}\) stay at the same spectral scale;
the all-integer and sparse-dimension finite-lift schedules then apply with the
same thresholds \(R>\Gamma+1\) and \(R>\Gamma\).  A random graph-mask extension
adds inhomogeneous Erdos--Renyi support profiles after a scalar Bernstein
max-degree exception; when the expected degree dominates the active-prefix union
logarithm, expected-degree normalization has no additional loss.  Dependent
bounded-degree support laws such as random regular graphs and random matchings
need no degree exception.  A weighted random-mask version replaces degree by a
Bernstein bound for the actual row variance
\(2\sum_q\eta_{pq}c_{pq}^2\), covering highly inhomogeneous variance-profile
entry models; its schedule form gives the same arbitrary-growth,
all-integer, and sparse-dimension thresholds once the support Bernstein exponent
dominates the active-prefix union logarithm.  A bounded-incidence specialization
reduces this to the normalized mean row variance \(\rho/\omega\), keeping the
same \(8/3\) threshold constants for bounded inhomogeneous entry masks.  The
same argument is recorded in
a predictable-profile form: the inclusion probabilities, proxy caps, and
variance budgets may be chosen from the core environment, as long as their
conditional Bernstein envelopes satisfy the same summable schedule.  This
covers adaptive and randomly sampled variance profiles without conditioning
back to a fixed deterministic mask model.  A dependent-mask budget version then
removes support independence altogether: arbitrary random support laws are
allowed once their realized weighted row budgets are controlled, so random
regular, matching, without-replacement, and other correlated masks enter through
the same finite-lift menu.  A bounded weighted-degree specialization gives the
usual \(d^{-1/2}\) normalization for any dependent mask law supported on
weighted-degree-\(d\) graphs, with no support-probability payment.
A separate variance-profile
certificate, obtained from
noncommutative Khintchine, gives \(O(\sqrt n)\) on every realization satisfying
\(\|\sum_i A_i^2\|_{\op}\lesssim n/\log(eD)\); in fact the logarithm may use the
realized rank of \(\sum_iA_i^2\) whenever the tuple lives in a common low-rank
corner.  Matrix Bernstein verifies this profile for independent random
contractions whenever
\(\|\sum_i\mathbb E A_i^2\|_{\op}=O(n/\log(eD))\), and in particular for
Haar-random rotations of arbitrary fixed or random spectra with average
Frobenius-square \(O(D/\log(eD))\).  The same argument now applies to
conditional square-twirl and exact finite-design rotations, including finite
irreducible unitary-group rotations and signed-permutation/monomial twirls,
under the identical average Frobenius-square profile.  A block-design twirl
variant replaces the scalar Frobenius profile by the operator-valued block
energy
\(\|\sum_{i,\alpha}d_{i,\alpha}^{-1}\operatorname{tr}(P_{i,\alpha}T_i^2P_{i,\alpha})P_{i,\alpha}\|_{\op}\),
covering independent block-Haar, block finite-design, and block
signed-permutation randomizations.  A dominated approximate-twirl variant
allows exact identities to be replaced by positive-semidefinite upper bounds,
so approximate scalar or block designs pay only the operator norm of the
aggregate twirl error.  A matrix-Freedman version gives the same conclusion
for adapted dependent random contractions whose predictable square profile is
\(O(n/\log(eD))\), and its predictable dominated-twirl specialization covers
online adaptive scalar, block, and approximate random-basis samplers with only
one-step twirl domination; a rank-sensitive common-corner variant replaces
\(D\) by a deterministic or random joint support rank \(\kappa\), and a finite
active-menu schedule pays \(\sum_h\kappa_{m,h}\) instead of \(J_mD_m\).  The
same square-function method has a two-scale form: a \(\kappa\)-dimensional core
may be perturbed by a full ambient residual under the same signing, paying
\(\log(e\kappa)\) for the core profile and \(\log(eD)\) only for the residual;
the corresponding predictable finite-menu schedule combines the rank cost
\(\sum_h\kappa_{m,h}\) with the residual cost \(J_mD_m\).  A finite-depth
multiscale form handles several low-rank layers plus one ambient residual under
one signing; active menus pay
\(\sum_{h,q}\kappa_{m,h,q}\) for all low-rank layers and \(J_mD_m\) only for
the residual.  A summable-depth version allows countably many low-rank layers
whenever the layer square-profile costs are summable, with active menus paying
the corresponding triple series \(\sum_{m,h,q}\) plus the same ambient residual
series.  A flag-compressible version generates those layers automatically from
nested initial projections: the \(q\)-th compression shell pays
\(\log(e\kappa_q)\), the part outside the terminal flag pays the ambient
logarithm, and active menus use the same summable triple structural tail.  A
compression-error flag form further replaces shell square-profile estimates by
predictable approximation budgets
\(\sum_i(\tau_{i,q-1}+\tau_{i,q})^2\), making common-flag spectral truncation,
wavelet tail, and low-rank approximation estimates directly usable.  A
tail-envelope specialization gives explicit common-flag criteria: polynomial
tails \((1+q)^{-p}\) beat rank growth \(\kappa_q\le\exp(\gamma q^s)\) when
\(p>1+s/2\), while any positive power tail works along geometrically growing
rank windows, and active menus pay only the summable envelope-exception and
sampled-search series.  The
same rank-sensitive replacement now has a direct dominated square-twirl form,
so random low-rank Haar, finite-design, signed-permutation, and approximate
twirl ensembles use \(O(n/\log(e\kappa))\) in their common corner.  Its
finite-menu form controls active libraries and countable exhaustions at cost
\(\sum_h\kappa_{m,h}\).  A rank-sensitive predictable dominated-twirl version
does the same for online adaptive random-basis streams chosen from the past.  A
Borel--Cantelli array form gives almost-sure
eventual signings, and almost-sure sampled certificates, whenever the row
twirl-budget and search-failure series are summable; explicit schedule rules
give the all-dimensions threshold \(n(D)\gtrsim\log^2(eD)\), and the sparse
stretched-exponential-dimension threshold with the sharp one-row exponent
constant.  Under the same square-function hypothesis, \(m\)
independent random sign trials fail to find a Spencer-scale signing with
probability at most \(2^{-m}\) beyond the profile failure.  A conditional
scalar-multiplier version covers arbitrary random self-adjoint contraction
environments after independent bounded thinning or amplitudes with conditional
second moments \(O(1/\log(eD))\), and a predictable version covers adaptive
online directions without conditional independence.  The same Borel--Cantelli
and explicit schedule rules give almost-sure scalar-multiplier signings along
all integer dimensions and sparse stretched-exponential dimension schedules.
Finite menus of such adaptive scalar-multiplier streams are simultaneous after
replacing the row factor \(D_m\) by \(|\mathcal H_m|D_m\), with common sampled
sign batches paying only \(|\mathcal H_m|2^{-M_m}\).  Exhaustive countable
menus are controlled on prefixes \(h\le J_m\) at cost \(J_mD_m\), so every fixed
adaptive scalar-multiplier stream is eventually covered.
Hybrid versions show
that any core tuple with a known \(K\sqrt n\) signing remains balanced after
adding centered independent scalar noise, bounded scalar martingale noise, or
bounded matrix-valued martingale-difference noise whose predictable quadratic
variation is scalar- or matrix-profile-dominated by independent random variance
budgets, including trace-budget and trace-envelope/low-rank/Schatten variants
with no dimension factor in the budget exception.  The signed-core predictable
matrix martingale-noise theorem also has an almost-sure triangular-array form
and explicit all-integer/sparse dimension schedules under the
\(O(n/\log(eD))\) predictable quadratic-variation budget, with scalar
martingale noise as a special case.  Countable libraries of such signed-core
matrix-noise arrays are controlled on active prefixes at cost \(J_mD_m\), so
every fixed already-signed core-plus-noise model is eventually covered.  Its
sampled localized specialization gives
the same schedules for coordinate-support and conditional Haar-subspace matrix
martingale arrays under the predictable weighted-overlap or average-rank budget.
The localized branch also has an active-prefix countable-library form with the
same \(J_mD_m\) cost, using those weighted-overlap or Haar average-rank budgets
directly.
A quantile/finite-moment
version for heavy-tailed total trace envelopes, an optimized grouped
tail-profile form covering exact survival functions, log-Weibull/lognormal,
polynomial, and sub-Weibull group tails, with optimized profile-menu and
common-profile aggregate criteria when groups share tail shapes, finite-menu
and countable weighted heterogeneous polynomial-index criteria, finite-menu and
countable weighted heterogeneous stretched-exponential and log-Weibull profile
criteria, a
closed-form common-index polynomial menu optimizer and growth schedule, including
finite class-count and power-decaying countable-class criteria, and explicit stretched-exponential and
log-Weibull growth schedules for polynomially many groups, plus a regularly
varying log-corrected schedule with an endpoint
logarithmic summability rule, and a
sub-gamma version for unbounded trace envelopes with termwise,
atomized/countable bounded,
gamma/chi-square,
gamma-atom, countable gamma-atom, weighted-chi-square/real-Wishart,
weighted-exponential/complex-Wishart, noncentral/correlated
Gaussian/Wishart, compound-Poisson, and mixed-atom
specializations,
or in predictable or conditionally sampled localized
directions, with the same signing and explicit matrix-Bernstein/Freedman tails.
Thus random localized, projection, frame, and
sparse-support ensembles are reduced to their weighted empirical overlap bound
\(\|\sum_i \rho_i^2P_i\|_{\op}\lesssim n/\log(eD)\).  An active-menu dependent
version needs only summable certification of the realized overlap budgets, and a
tail-capped dependent version adds sub-Weibull or finite-moment localized blocks
by paying cap-tail summability plus the capped overlap certificate.  A
common-sign version direct-sums the active menu and gives one signing for all
dependent localized streams after replacing \(\log(eD)\) by
\(\log(eJ_mD_m)\).  A support-rank variant replaces \(J_mD_m\), and in the
rectangular case \(J_m(d_{m,1}+d_{m,2})\), by the total active
support-envelope rank, with the same replacement available after tail-capping
unbounded localized streams under summable cap-tail failures.  Coordinate-window
Bernstein corollaries verify these support-envelope hypotheses from scalar
occupancy tails inside the active windows; predictable coordinate-window
Freedman corollaries remove independence for adapted support exposure, with
explicit arbitrary-growth, all-support-rank, and sparse-support-rank schedules
under the half-budget load condition.
Subspace-envelope Bernstein corollaries do the same for random projection
supports inside deterministic low-rank envelopes.  Predictable
subspace-envelope Freedman corollaries remove independence by checking adaptive
projection supports from conditional overlap budgets, with the same explicit
support-rank schedule consequences.  Rectangular coordinate-support and
row/column random-subspace Bernstein menus now also have explicit
active-menu schedules: arbitrary dimension growth pays
\(J_m(d_{m,1}+d_{m,2})\), all integer rectangular dimensions need
\(c_0a>\Gamma+2\), sparse schedules need \(c_0a>\Gamma+1\), and sampled
certificates add only the corresponding \(J_m2^{-M_m}\) search series.  A
tail-capped rectangular coordinate-support variant pays only deterministic
operator-norm cap tails plus the same capped row/column occupancy series, and
its cap-selection schedule gives explicit sub-Weibull and finite-moment sparse
rectangular block criteria.  For independent inhomogeneous random coordinate supports, scalar Bernstein controls
the weighted occupancies \(\max_j\sum_i\rho_i^2{\bf 1}_{\{j\in S_i\}}\); the
homogeneous support-density \(O(1/\log(eD))\) case is the unweighted specialization.  A
matrix Bernstein variant gives the same conclusion for independent random
subspace projections with
\(\|\sum_i\rho_i^2\mathbb E P_i\|_{\op}=O(n/\log(eD))\), with Haar ranks entering
through \(D^{-1}\sum_i\rho_i^2r_i\).  The deterministic amplitude weights may
also be independent random amplitudes; the coordinate budget becomes
\(\max_j\sum_i p_{ij}\mathbb E\rho_i^2\), and the random-subspace budget becomes
\(\|\sum_i\mathbb E(\rho_i^2P_i)\|_{\op}\), with the same Bernstein tails.  An
explicit random-amplitude schedule corollary gives the same all-integer,
sparse, arbitrary-growth, and sampled-certificate thresholds.  For a fixed
orthogonal block decomposition the same argument pays the block count \(B\)
instead of the ambient coordinate count \(D\), with block budget
\(\max_\alpha\sum_iq_{i\alpha}\mathbb E\rho_i^2\).  Random conditional
inclusion profiles \(Q\) are also allowed, replacing these deterministic
budgets by \(\max_j\sum_i\mathbb E(\rho_i^2Q_{ij})\) or
\(\max_\alpha\sum_i\mathbb E(\rho_i^2Q_{i\alpha})\), which covers random
support sizes without independence from the amplitudes.  The random-profile
tail also has a variance-sensitive Borel--Cantelli form: with \(N_m=D_m\) or
\(B_m\), summability of
\[
  N_m\exp\{-t_m^2/(2(\sigma_m^2+t_m/3))\}
\]
gives eventual almost-sure signings and sampled certificates.  A predictable
profile/Freedman variant removes independence across rows: coordinate or block
supports may be sampled online from \(\mcF_{i-1}\)-measurable profiles, paying
only the probability that the realized predictable overlap budget is too large.
Its schedule corollary gives the same arbitrary-growth rule with
\(\log(eN_m)\), and the same all-integer and sparse thresholds whenever the
predictable-budget exceptions are summable.  A finite-menu predictable-profile
version gives simultaneous eventual signings for many online coordinate/block
support models in the same dimension, paying only
\(L_m=\sum_hN_{m,h}\) plus the summed budget exceptions.  An exhaustive
countable-menu version controls a growing prefix \(h\le J_m\) of a countable
library at each row, paying
\(L_m^{(J)}=\sum_{h\le J_m}N_{m,h}\), so every fixed online support model is
eventually covered.  A deterministic-envelope
specialization gives zero budget exception for fully adaptive support processes
whose conditional coordinate or block inclusion probabilities are bounded by
fixed profiles \(p_{ij}\) or \(q_{i\alpha}\), with random amplitudes dominated
by deterministic caps.  Its deterministic-envelope schedule corollary removes
the exception entirely under the fixed half-budget and gives the same
\(\log(eN_m)\) arbitrary-growth rule, all-integer threshold
\(a>16(\Gamma+1)/(3C_0)\), sparse threshold \(a>16\Gamma/(3C_0)\), and
logarithmic sampled-sign counts.  A finite-menu version gives simultaneous
eventual signings for many coordinate/block envelopes in the same dimension,
paying only \(L_m=\sum_hN_{m,h}\) in place of a single structural count.  Its
exhaustive countable-menu version pays
\(L_m^{(J)}=\sum_{h\le J_m}N_{m,h}\) for the active prefix and eventually covers
every fixed deterministic envelope model with no budget-exception term.  A
Borel--Cantelli array form gives the same conclusion eventually almost surely
along any dimension/sample schedule whose coordinate or subspace Bernstein
failure series is summable, and the same
tail is automatically summable along all integer dimensions under
\(n_D\ge a\log^2(eD)\) once \(a>32/(3C_0)\), or along sparse
stretched-exponential dimension schedules once \(a>16/(3C_0)\).  The same
weighted support families have sampled-sign certificates with only an additional
\(2^{-M}\) finite-search failure; under the preceding all-dimension hypotheses,
\(M_D\ge(1+\delta)\log_2(eD)\) sign trials suffice along all integer dimensions,
and \(M_m\ge(1+\delta)\log_2(e+m)\) suffices along sparse schedules.  More
generally, the same certificate is eventual whenever
\[
n_m\gtrsim \log(eD_m)\bigl(\log(eD_m)+\log(e+m)\bigr)
\]
with the explicit constant \(16/(3C_0)\) below.  In the
signed-core setting the
same martingale-noise tail is made explicit for adaptively sampled coordinate
supports with coordinate-dependent predictable inclusion bounds and for
conditional Haar subspaces under the corresponding random weighted conditional
inclusion-probability or average-rank budget, where predictable amplitude weights
replace \(P_i\) by \(\rho_i^2P_i\) in the conditional overlap.  For
coordinate-support and Haar-subspace matrix noise with independent random
amplitudes, scalar Bernstein gives explicit tails for that weighted budget.  For
coordinate supports the inclusion profile may be deterministic or independently
random across samples; in the homogeneous coordinate case the budget scale is
\(p\sum_i\mathbb E\rho_i^2\).  The same argument applies to arbitrary fixed
orthogonal block decompositions, with the budget union over the number of
blocks.  In the fixed-rank Haar rank-\(r_i\) case the scale is
\(\sum_i(r_i/D)\mathbb E\rho_i^2\).  Random rank fractions obey the same
Bernstein profile with mean \(\sum_i\mathbb E[(R_i/D)\rho_i^2]\).  More
generally, conditional isotropy
\(\mathbb E(P_i\mid\mathcal F_{i-1})=\alpha_iI\) gives the identical profile
with random intensities \(\alpha_i\), and the same proof only needs spectral
domination by \(\alpha_iI\).  A Borel--Cantelli menu theorem converts all of
these random-amplitude localized tails into almost-sure all-integer and sparse
dimension schedules once the scalar Bernstein budget exceptions are summable.
An exhaustive countable-menu version controls active prefixes after replacing
the one-row budget union count by \(\sum_{h\le J_m}N_{m,h}\), and therefore
eventually covers every fixed random-amplitude localized model.
A common-sign variant shows that if one core signing balances every active-prefix
core, then that same signing balances all active-prefix random perturbations
after the identical union-count payment.
A finite-lift common-menu corollary supplies such common core signings
deterministically whenever the active-prefix direct-sum lift entropy is \(O(n)\).
Combining these statements gives finite-lift core libraries plus localized
random-amplitude noise with the same active-prefix union-count and homogeneous
energy schedules.
In homogeneous coordinate, block, Haar, dominated projection, and mixed
coordinate/Haar models this reduces to the concrete energy rule
\(\beta\sum_i\mathbb E\rho_i^2=O(n/\log(eD))\), with
\(\sum_i\mathbb E\rho_i^2\gtrsim\log(eD)\) only needed to make the scalar
budget exception summable.
The homogeneous specialization also has an active-prefix countable-menu form
paying \(\Lambda_m^{(J)}=\sum_{h\le J_m}N_{m,h}\); if
\(\Lambda_D^{(J)}\le(eD)^\Gamma\), the all-integer energy margin is
\(R>8(\Gamma+1)/3\), while sparse stretched-exponential schedules need only
\(R>8\Gamma/3\).
The same homogeneous energy margins hold in common-sign form when the
active-prefix cores share one signing.
The
same scalar-Bernstein profile also handles independent samples which choose
between a coordinate support and a Haar subspace by replacing the coordinate
profile with \(M_iQ_{ij}+(1-M_i)R_i/D\).  In the commutative
AF case a clopen partition with
\(M\) atoms and oscillation \(\eta\) gives
\(C\sqrt{n\log(e+M/n)}+2n\eta\), while interval-range examples force
\(M=\Omega(\eta^{-1})\) at scale \(\eta\).  For H\"older tuples, ordinary
covering growth with exponent \(d\) has Spencer threshold \(d\le2\alpha\). This
includes clopen/metric Minkowski endpoints and, via zero-dimensional
resolutions, arbitrary compact metric targets, H\"older quotients, snowflaked
targets, finite unions/products, finite and exhaustive finite-scale H\"older
atlases, and finite product-parametrized strata. H\"older quotient and stratum
maps divide the polynomial exponent by their exponent; polynomial exhaustive
atlas schedules add the stage-complexity exponent divided by the approximation
exponent, while stretched-exponential exhaustive schedules pay only logarithmic
stage complexity; signed logarithmic cover powers follow the same
critical-window rules.  At the purely Banach-space level, totally
bounded classes have hereditary \(o(n)\) discrepancy by a finite-net fiber
signing, with explicit polynomial, stretched-exponential, and signed-log
entropy rates; in
particular compact operator images of unit balls have sublinear signed sums.
Direct dyadic entropy-number decay gives sharper logarithmic rates, for example
\(e_k\lesssim k^{-s}\) implies \(\Delta_n\lesssim n(\log n)^{-s}\), and the
same package is stable under passage from standard external entropy numbers to
internal nets, finite unions, finite and weighted summable Minkowski sums, and H\"older
images of compact Banach classes.  Weighted series are also recorded in a
profile form, balancing the component entropy profile against the weight-tail
profile before the final dyadic discrepancy step; polynomial accumulated
amplitude and tail profiles give a closed optimization rule, and for polynomial
weights with logarithmic corrections, including polynomially growing component
entropy amplitudes, this yields optimized log-balanced exponents.
The same principle gives a finite-rank approximation envelope: if a compact
class lies within \(\tau_m\) of an \(m\)-dimensional subspace, then
\(\Delta_n\) is bounded by a rank-\(m\) net term plus \(n(\eta+\tau_m)\), with
concrete polynomial, log-refined polynomial, and stretched-exponential rates
from approximation numbers and, in Hilbert spaces, from singular value and
Schatten-class decay.
\end{abstract}

\maketitle

\tableofcontents

\section{Introduction}

Spencer's theorem that ``six standard deviations suffice'' is a central result
of discrepancy theory \cite{Spencer,Matousek,Chazelle}. In one standard form,
$n$ vectors in the unit cube of $\R^n$ can be signed so that their signed sum has
$\ell_\infty$ norm $O(\sqrt n)$. The Matrix Spencer conjecture asks for the
analogous statement with vectors replaced by self-adjoint matrices and
$\ell_\infty$ norm replaced by operator norm.

\begin{definition}[Matrix discrepancy]
For self-adjoint matrices $A_1,\ldots,A_n$, set
\[
\disc(A_1,\ldots,A_n)
= \min_{\eps\in\{-1,1\}^n}
\norm{\sum_{i=1}^n \eps_i A_i}_{\op} .
\]
For arbitrary complex matrices we use the same formula and write
\[
\discop(A_1,\ldots,A_n)
=\min_{\eps\in\{-1,1\}^n}
\norm{\sum_{i=1}^n\eps_i A_i}_{\op}.
\]
\end{definition}

\begin{question}[Matrix Spencer]
Is there a universal constant $C$ such that every $n$ self-adjoint matrices
$A_1,\ldots,A_n\in \M_n(\C)$ satisfying $\norm{A_i}_{\op}\le 1$ obey
\[
\disc(A_1,\ldots,A_n)\le C\sqrt n?
\]
More generally, one asks for the same bound for $d\times d$ matrices with
$d\le n$, with a constant independent of $d$.
\end{question}

The general problem remains open. Noncommutative Khintchine inequalities give
logarithmic losses \cite{LustPisier,Pisier}. Important recent advances include
the intrinsic-freeness concentration theory of Bandeira, Boedihardjo and van
Handel \cite{BBvH}; the quantum-communication route of Hopkins, Raghavendra and
Shetty \cite{HRS}; the polylogarithmic-rank theorem of Bansal, Jiang and Meka
\cite{BJM}; and the mirror-descent matrix-discrepancy framework of Dadush,
Jiang and Reis \cite{DJR}. The engine we use is the block-diagonal case: a bound
of $O(\sqrt{n}\,\sqrt{\log h})$ for block size $h$ was first proved by Levy,
Ramadas and Rothvoss \cite{LRR} via the multiplicative weight update method, and
Dadush, Jiang and Reis \cite{DJR} later removed the restriction $h\le n$ through
their mirror-descent partial-coloring framework. These results do not directly
cover many full-rank structured families, such as normalized Toeplitz, Hankel,
or convolution matrices in the square regime.

The point of this paper is to isolate a simple sufficient condition that does
cover these structured classes. The condition is the existence of a bounded
finite block lift. If each input matrix is a compression, quotient, or bounded
linear image of a block-diagonal element whose total block complexity is small,
one signs the lifted block-diagonal matrices and then maps back.

\subsection{The finite-lift principle}

Let
\[
\mcB=\bigoplus_{\alpha=1}^m \M_{r_\alpha}(\C),
\qquad
\norm{Y}_{\mcB}=\max_\alpha \norm{Y_\alpha}_{\op}.
\]
We write
\[
h(\mcB)=\max_\alpha r_\alpha,\qquad
R(\mcB)=\sum_{\alpha=1}^m r_\alpha,
\qquad
\Ent(\mcB)=h(\mcB)R(\mcB).
\]
The number $\Ent(\mcB)$ is the entropy dimension seen by block-diagonal Matrix
Spencer. It is $D$ for a diagonal algebra of dimension $D$, $D^2$ for the full
matrix algebra $\M_D$, and $bD$ for a block-diagonal algebra on $\C^D$ with
blocks of size at most $b$.

\begin{theorem}[Finite block-lift Matrix Spencer]\label{thm:main}
There is a universal constant $C$ with the following property. Let
$A_1,\ldots,A_n\in \M_D(\C)$ be arbitrary complex matrices with
$\norm{A_i}_{\op}\le 1$. Suppose there are elements $Y_i\in\mcB$ and a fixed
linear map
\[
\Pi: \Span\{Y_1,\ldots,Y_n\}\longrightarrow \M_D(\C)
\]
such that
\[
\Pi(Y_i)=A_i,
\qquad
\norm{\Pi(Y)}_{\op}\le \Lambda\norm{Y}_{\mcB}
\quad (Y\in\Span\{Y_1,\ldots,Y_n\}),
\]
and
\[
\norm{Y_i}_{\mcB}\le K \qquad (1\le i\le n).
\]
Then
\[
\discop(A_1,\ldots,A_n)
\le C\Lambda K
\sqrt{n\log\left(e+\frac{\Ent(\mcB)}{n}\right)} .
\]
If the $A_i$ are self-adjoint and the $Y_i$ may be chosen self-adjoint, the
same estimate holds for $\disc(A_1,\ldots,A_n)$.
\end{theorem}

Thus every class with uniformly bounded lift loss and $\Ent(\mcB)=O(n)$
satisfies the conjectured $O(\sqrt n)$ Matrix Spencer bound. Scalar Toeplitz
matrices are the commutative case: after a bounded finite Toeplitz extension and
Fourier sampling, the lift algebra is diagonal of dimension $O(D)$. Finite group
convolution and block Toeplitz matrices are noncommutative variants in which the
finite lift has bounded block size.

The theorem has a built-in limitation. If the full unit ball of $\M_D(\C)$ is lifted
exactly through a finite block algebra, or even approximated to fixed
operator-norm error \(<1\) by such a lift, dimension forces
$\sum_\alpha r_\alpha^2\ge D^2$, hence $\Ent(\mcB)\ge D^2$. Therefore this
full-unit-ball route cannot, by itself, turn arbitrary $D\times D$ matrices in
the square regime $D=n$ into an $O(\sqrt n)$ theorem. It yields instead the usual
$\sqrt{n\log(e+D^2/n)}$ scale. The point is not that finite lifts are the only
possible method, but that this reusable class-level sufficient condition applies
when the structured class has small finite-lift entropy.

Since this note is organized around finite lifts, it keeps the elementary $hR$
theorem and the corresponding obstruction bookkeeping explicit.  The current
strongest general input for exact finite-dimensional C*-algebra containment is
stronger: Akbas and Sra \cite{AkbasSra} prove that contractions contained in a
finite-dimensional C*-algebra of intrinsic dimension $O(n)$ already satisfy the
$O(\sqrt n)$ Matrix Spencer bound.  We use that result below as an external
upgrade for finite stages and for group convolution algebras.  We also use
their perturbative theorem for algebraic families plus low-Frobenius
self-adjoint tails to record a matrix-representation AF corollary and a
distribution-free random algebraic-core ensemble theorem.  The
finite-lift theorem remains the mechanism for the Toeplitz, Hankel, multiplier,
and obstruction statements proved here.

\subsection{Main deterministic consequences}

The abstract theorem gives the following consequences. Constants are universal
unless a dependence is displayed in the subscript.

\begin{enumerate}[leftmargin=2em]
\item \textbf{Scalar Toeplitz and Hankel.} Every collection of arbitrary complex
$d\times d$ Toeplitz or Hankel contractions satisfies
\[
\discop(A_1,\ldots,A_n)
\le C\sqrt{n\log\left(e+\frac{d}{n}\right)} .
\]
In particular, self-adjoint Toeplitz and Hankel matrices satisfy the Matrix
Spencer bound when $d\le n$.

\item \textbf{Fixed-block and fixed-level multilevel Toeplitz/Hankel.} For
fixed matrix block size $b$, all one-level $b\times b$ block Toeplitz and block
Hankel contractions of total dimension $D$ satisfy
\[
\discop(A_1,\ldots,A_n)
\le C_b\sqrt{n\log\left(e+\frac{D}{n}\right)} .
\]
More generally, for any fixed multilevel order $r\ge 1$ and block size $b$, all
$r$-level $b\times b$ block Toeplitz and block Hankel contractions of full
bandwidth and ambient dimension $D=b\prod_j d_j$ satisfy
\[
\discop(A_1,\ldots,A_n)
\le C_{r,b}\sqrt{n\log\left(e+\frac{D}{n}\right)}
\qquad(\text{Theorem \ref{thm:blocktoepfull}}),
\]
so this is $O_{r,b}(\sqrt n)$ whenever $D\le n$. The constant $C_{r,b}$ here is
non-constructive and grows with $r$, since it passes through a multilevel
Toeplitz symbol extension (Theorem \ref{thm:mltext}) whose constant cannot be
taken uniform in $r$. A complementary, fully explicit form is available for
banded symbols: if the matrix symbols have per-coordinate degree at most $w$
(with each box dimension at least $4w$), then
\[
\discop(A_1,\ldots,A_n)
\le C_{r,b}\sqrt{n\log\left(e+\frac{w^r}{n}\right)}
\]
with a dimension-free, explicit constant and no appeal to the symbol extension.
This banded bound is sample-optimal: the factor $w^r$ is forced by a norming-set
dimension count, and a matching lower bound of the same form is realized by
multilevel circulants, which are exactly complex vector Spencer in disguise. The
formal anisotropic version replaces \(w^r\) by \(\prod_j(2w_j+1)\), and the
periodic circulant lower bound has the same product form.  In the self-adjoint
product-bandwidth regime \(\prod_j(2w_j+1)\le n\) it is
\(O_{r,b}(\sqrt n)\) with \(D\) arbitrary, so the class escapes the \(D^2\)
entropy barrier.

\item \textbf{Common finite-dimensional C*-algebras.} Suppose the $A_i$ lie in a
common finite-dimensional C*-algebra represented as
\[
\mcA \simeq \bigoplus_{\alpha=1}^m
\left(\M_{r_\alpha}(\C)\otimes I_{m_\alpha}\right)
\subseteq \M_D(\C).
\]
Then
\[
\discop(A_1,\ldots,A_n)
\le C\sqrt{n\log\left(e+
\frac{(\max_\alpha r_\alpha)\sum_\alpha r_\alpha}{n}
\right)} .
\]
This includes simultaneous diagonalization, common bounded-block
diagonalization, and commutative association-scheme algebras.

\item \textbf{Finite group convolution.} For left convolution operators on
$\ell^2(G)$, where $G$ is finite and $q(G)=\max_{\pi\in\widehat G}\dim\pi$,
\[
\discop(L_{f_1},\ldots,L_{f_n})
\le C\sqrt{n\log\left(e+
\frac{q(G)\sum_{\pi\in\widehat G}\dim\pi}{n}
\right)} .
\]
For finite abelian groups this becomes
$C\sqrt{n\log(e+|G|/n)}$.
Using the Akbas--Sra intrinsic-dimension theorem, the square regime
\(|G|=O(n)\) has the sharper \(O(\sqrt n)\) bound for arbitrary finite groups,
with no bounded-irrep hypothesis; the low-Frobenius perturbative version also
allows self-adjoint rank-\(n/\log^3(en)\) noise on top of those convolution
cores.

\item \textbf{Closure operations.} The framework is stable under compression,
bounded linear images, completely bounded fixed-size tensor amplification, and
finite sums when the combined direct-sum entropy remains controlled. The last
condition is important: the entropy of a direct sum is not the sum of the
individual entropies in general.

\item \textbf{Structured random inputs.} All of the preceding statements are
pathwise deterministic. They therefore apply to random Toeplitz, Hankel,
fixed-block multilevel Toeplitz/Hankel, circulant, and group-convolution samples
after normalization by their realized operator norms; a complementary parity
obstruction to the stability of random banded Toeplitz matrices is established in
the companion paper \cite{Shmalo-paper5}.  The Akbas--Sra perturbative input
also gives a distribution-free random theorem for arbitrary random algebraic
cores of intrinsic dimension \(O(n)\), including \(D\times D\) full random
self-adjoint models for \(D^2=O(n)\), after adding arbitrary dependent
low-Frobenius, low-rank, spiked, compressed, or coordinate-supported tails.
At the completely distribution-free level, the all-plus gauge
\[
  \max\left\{\max_i\|X_i\|_{\op},
  n^{-1/2}\left\|\sum_iX_i\right\|_{\op}\right\}
\]
normalizes arbitrary self-adjoint or rectangular random matrices into a
Spencer-scale instance, and \(L^p\) moment or quantile budgets for this gauge
give checkable deterministic almost-sure schedules, with explicit
all-integer and sparse-dimension thresholds for polynomial active menus.
Direct tail-profile bounds, including sub-Weibull gauge tails, give the
corresponding logarithmic deterministic normalizations without independence.
The finite-confidence version gives one-row or finite-library certificates with
explicit dependence on the requested failure probability \(\delta\), and its
conditional form allows the normalizing scales to be chosen from side
information through conditional moment or tail envelopes, and a random-active
finite-confidence form charges only the selected menu.  The corresponding
conditional Borel--Cantelli form gives almost-sure adaptive schedules when the
conditional failure budgets are summable, including random active submenus
chosen by the environment.
For full centered random matrix ensembles in arbitrary dimension, even with
dependence among the matrices, an operator-tail hypothesis stable under scalar
\(\ell_2\) sums gives another average-case route: after scaling by
\(K(\sqrt D+\sqrt{\log(en)})\), all matrices are contractions and the all-plus
signing has norm at most \(\sqrt n\) with high probability.  The corresponding
active-menu version is simultaneous over finite or countable libraries under
the operator-tail union cost \(J_m(n_m+1)e^{-\tau_m^2}\), with no
coordinatewise sign-symmetry requirement.  If the active operator-tail streams
are selected randomly by the environment, the same proof pays only the selected
conditional tail cost.  The rectangular version has the same selected-cost
random-active form and uses the same
all-plus certificate and the left/right dimension sum \(d_1+d_2\).  Bilinear
scalar sub-Gaussian tails for
\(\langle u,\sum_i a_iX_i v\rangle\) verify that rectangular operator-tail
hypothesis by a two-net argument, and their random-active form pays only the
selected bilinear tail cost; rectangular Gaussian and feature-map families only
need a uniform bilinear proxy \(q\).
The Bernstein bilinear analogue also has a random-active selected-cost form,
with the additional jump term
\(b(\sqrt{d_1+d_2}+\tau)^2\), and its independent-entry specialization is
checkable from selected variance and jump proxies; a further entrywise
Bernstein-mgf corollary verifies those proxies from conditional independence of
the individual matrix entries, and its finite-confidence form gives explicit
one-row failure budgets.  For bounded conditionally independent entries the
same finite-confidence certificate needs only conditional entry variances and
uniform entry caps.
The bounded matrix-Bernstein all-plus route also has a finite-confidence
selected-library version, using only matrix norm caps and all-plus variance
budgets, with a rectangular form obtained by self-adjoint dilation.  The
adapted bounded matrix-Freedman route has the same finite-confidence version
with predictable row and column variance budgets.
The tail-capped Bernstein/Freedman route likewise has a finite-confidence
selected-library version for unbounded streams after paying the selected
cap-failure budget.
Combining the core-signing and operator-tail inputs, the paper also allows any
random core class with an environment-measurable \(K\sqrt n\) signing, including
dimension-controlled algebraic cores, finite-lift cores, Toeplitz cores, and
Hankel cores, to absorb full random operator-tail noise whose scalar linear
combinations obey the same Gaussian tail conditionally on the core, in arbitrary
ambient dimension after the same natural spectral scaling of the noise.
For active-prefix finite-lift core libraries, the same mechanism is simultaneous
after paying \(J_m(n_m+1)e^{-\tau_m^2}\) and scaling the full noise by
\(\sqrt{D_m}+\tau_m\); the concrete schedules are arbitrary-growth
\(\tau_m^2\gtrsim\log(J_m(n_m+1))+\log m\), all-integer
\(\tau_D^2=R\log(eD)\) with \(R>\Gamma+1\), and sparse
\(R>\Gamma\) when \(J_m(n_m+1)\le(eD_m)^\Gamma\).
For Gaussian, Rademacher, or scalar sub-Gaussian matrix-series noise with
predictable self-adjoint coefficient directions, a matrix-Laplace tail improves
the finite-lift perturbation scale to
\(q_{m,h}(\sqrt{\log(e\kappa_{m,h})}+\tau_m)\), assuming the realized
individual and signed-sum coefficient variances are bounded by \(q_{m,h}^2\)
and \(n_mq_{m,h}^2\), with coefficient square-function ranks at most
\(\kappa_{m,h}\).  Taking \(\kappa_{m,h}=D_m\) recovers the full-rank
fallback.
A common-projection specialization packages this as low-rank matrix-series
noise: if all coefficient directions live in a rank-\(\kappa_{m,h}\)
subspace, the rank hypotheses are automatic.
A norm-budget specialization then replaces the square-function hypotheses by
the scalar checks \(\max_i\sum_r b_{m,h,i,r}^2\le q_{m,h}^2\) and
\(\sum_{i,r}b_{m,h,i,r}^2\le n_mq_{m,h}^2\), whenever
\(\|G_{m,h,i,r}\|_{\op}\le b_{m,h,i,r}\) inside the common support.
A block-load specialization lets these atoms range over an orthogonal block
family, replacing total atom energy by the maximum incident block energy.
A bounded weighted block-design specialization covers regular, matching,
permutation, and balanced without-replacement designs with no load-tail cost
once their weighted block degrees are deterministic or fail only summably.
The bounded-overlap block-hypergraph specialization makes this a direct
\(q/\sqrt d\) rule for block designs, regular block hypergraphs, and balanced
block dictionaries with overlap degree \(d\).
The heterogeneous bounded-overlap version replaces \(d\) by local row-block and
total-block degree profiles, so irregular block hypergraphs need not be
normalized by their worst active block.
The row-wise block-permutation specialization makes this explicit for
\(d\)-round permutation, matching, and balanced without-replacement assignments
with atom amplitudes at scale \(q/\sqrt d\).
A random block-load specialization verifies the incident-energy event by
conditional scalar Bernstein, with union cost over active blocks and atoms.
A uniform random block-support specialization makes the same test explicit in
terms of each atom's expected support size divided by the number of blocks.
A profile schedule version requires only the resulting Bernstein exponents to
dominate the active block-count logarithm.
A homogeneous uniform version further reduces those exponents to
\(N_{\rm ind}s/B\) and \(N_{\rm tot}s/B\), yielding all-integer threshold
\(8(\Gamma+1)/3\) and sparse threshold \(8\Gamma/3\) for the incidence counts.
A weighted version applies the same argument directly to random incident
weights, covering random amplitudes, bounded compound-thinned atom fields, and
random weighted dictionaries.
A predictable block-support profile specialization removes independence for
sequential support exposure by using one-step inclusion probabilities and
predictable amplitude envelopes.
Its profile schedule has the same active block-count logarithm as the weighted
Freedman schedule, with the predictable-support exponents in place of the
weighted incident-load exponents.
An inhomogeneous incidence-mass specialization replaces the raw exponents by
row-block and total-block predictable probability masses, retaining the same
\(8/3\) threshold constants for biased adaptive block samplers.
Its homogeneous predictable uniform schedule has the same
\(8(\Gamma+1)/3\) all-integer and \(8\Gamma/3\) sparse incidence thresholds as
the independent homogeneous block-support schedule.
A predictable weighted version replaces independence by scalar Freedman
control of the incident weights, allowing adaptive atom sampling and
without-replacement block exposure.
A profile schedule form covers both weighted variants once the corresponding
Bernstein or Freedman exponents dominate the active block-count logarithm.
A bounded weighted incidence specialization needs only the normalized
row-block and total-block mean loads \(\rho/\omega\), with the same \(8/3\)
threshold constants.
A quadratic-form-to-operator-tail lemma then verifies the same hypothesis from
conditional scalar sub-Gaussian bounds on
\(\langle u,\sum_i a_iY_i u\rangle\), so Gaussian, correlated Gaussian, and
Wigner-type active-prefix noise libraries only need their one-dimensional
quadratic-form proxy \(q_{m,h}\), and the natural scale becomes
\(q_{m,h}(\sqrt{D_m}+\tau_m)\).
For conditionally independent upper-triangular sub-Gaussian entries, that proxy
is controlled by the entry square-function envelope \(Q_{m,h}\); uniform entry
scale \(\sigma_{m,h}\) gives \(Q_{m,h}\le\sqrt2\,\sigma_{m,h}\), hence the
usual \(\sigma_{m,h}(\sqrt{D_m}+\tau_m)\) Wigner/GOE/GUE normalization.
A row-variance specialization replaces the vector supremum by
\(\max_p(\alpha_p^2+2\sum_{q\ne p}\beta_{pq}^2)^{1/2}\), giving sparse and
banded entry-profile libraries the explicit scale
\((\alpha^2+2\Delta\beta^2)^{1/2}(\sqrt{D_m}+\tau_m)\).
At normalized graph scale \(\beta\lesssim\Delta^{-1/2}\), the same
all-integer and sparse-dimension active-prefix schedules apply with no degree
loss.  The same row-variance reduction now covers independent random graph masks
after adding the scalar Bernstein max-degree failure probability; for
inhomogeneous Erdos--Renyi masks this uses expected-degree normalization whenever
the expected degree is larger than the menu union logarithm.  Random regular,
matching, and other dependent bounded-degree masks skip this support exception.
A weighted random-mask refinement controls the actual row-variance profile
\(2\sum_q\eta_{pq}c_{pq}^2\) by scalar Bernstein, so highly inhomogeneous
variance-profile masks need not be reduced to their maximum degree; the same
arbitrary-growth, all-integer, and sparse-dimension schedules follow when the
support Bernstein exponent dominates the active-prefix union logarithm.
A bounded-incidence specialization reduces this support condition to the
normalized mean row variance \(\rho/\omega\), keeping the same \(8/3\)
threshold constants for bounded inhomogeneous entry masks.
A predictable-profile variant lets the whole intensity and proxy field be
chosen from the core environment, paying only the probability that the
conditional Bernstein envelopes exceed their deterministic schedule budgets.
A dependent weighted-budget variant removes support independence entirely once
the realized row budgets are controlled outside a summable event.
A bounded weighted-degree specialization recovers the normalized
\(d^{-1/2}\) edge scale for arbitrary dependent weighted-degree-bounded masks
with no support exception.
For coordinatewise sign-symmetric random tuples, the paper also records a
direct transfer from all-plus operator-norm tails to discrepancy tails for the
same realization-normalized ensemble, and a Borel--Cantelli active-menu version
turns summable all-plus tails into simultaneous eventual signings for finite or
countably exhausted sign-symmetric libraries.  A matrix-Bernstein specialization
turns independent or conditionally independent bounded sign-symmetric arrays
with summable variance-envelope failures into simultaneous eventual normalized
signings, with explicit all-integer and sparse schedules once
\(n_m\gtrsim\log^2(eD_m)\), including closed-form constants for the admissible
discrepancy margin.  Symmetric scalar-multiplier menus therefore have the same
closed-form margins under the conditional second-moment profile
\(\|\sum_iq_iB_i^2\|_{\op}=O(n_m/\log(eD_m))\).  Unbounded symmetric
sub-Gaussian multipliers are also covered after deterministic cap
normalization; the matrix-Laplace tail gives all-integer and sparse margins
\(K>\sqrt{2\beta(\Gamma+1)}\) and \(K>\sqrt{2\beta\Gamma}\), without the
bounded-jump correction.  A tail-capped companion covers arbitrary symmetric
multiplier tails after summable cap failures; in the sub-Weibull case the cap
cost is \(2\sum_i\exp\{-(L/a_i)^\alpha\}\), while the truncated second-moment
budget inherits the same bounded-Bernstein margins \(\Theta(\Gamma+1,a,\beta)\)
and \(\Theta(\Gamma,a,\beta)\).  A polynomial moment-proxy specialization
covers every tail index \(p>0\): writing \(r=p\wedge2\), the variance budget is
controlled by \(L^{-r}\|\sum_i a_i^rB_i^2\|_{\op}\) and the cap cost by
\(\sum_i(a_i/L)^p\).  The same conclusions hold when the caps and moment
proxies are chosen from the conditioning environment, provided the conditional
cap-failure series is summable.
A noncommutative-Khintchine square-function certificate gives a variance-profile
random route as well: any realization with
\(\|\sum_i A_i^2\|_{\op}=O(n/\log(eD))\) is Spencer-scale; more precisely, the
logarithm can be taken at the realized rank of \(\sum_iA_i^2\).  Hermitian
dilation gives the same certificate for rectangular or non-self-adjoint
operators from the two natural profiles
\(\|\sum_iX_iX_i^*\|_{\op}\) and \(\|\sum_iX_i^*X_i\|_{\op}\), with the
logarithm at \(d_1+d_2\) or at the realized left-plus-right profile rank.
Independent
random sign search finds such a rank-sensitive signing after \(m\) trials except
with an additional \(2^{-m}\) failure probability.  A common-sign self-adjoint
menu version direct-sums the active square-profile streams and gives one
signing for all active models at the scale \(\log(eJ_mD_m)\), with a shared
sampled sign list losing only \(2^{-M_m}\).  A rank-sensitive common-sign
version replaces \(J_mD_m\) by the certified total active square-function rank,
so low-rank active libraries pay \(\log(eR_m)\) instead of the ambient menu
dimension.  The rectangular common-sign version has the same rank-sensitive
upgrade with the total active left-plus-right square-profile rank in place of
the active ambient rectangular dimension.  Matrix Bernstein turns this into a
probabilistic theorem for independent random contractions with
\(\|\sum_i\mathbb E A_i^2\|_{\op}=O(n/\log(eD))\); Haar-random rotations of fixed
or conditionally random spectra are included whenever the average
Frobenius-square is \(O(D/\log(eD))\).  More generally, conditional
square-twirl rotations and exact finite-design rotations, including finite
irreducible unitary-group rotations and signed-permutation/monomial twirls,
obey the same theorem with the same average Frobenius-square profile.  A
block-design twirl version covers partially randomized bases and fixed
orthogonal decompositions by replacing the scalar Frobenius energy with the
corresponding block-energy operator norm; a dominated approximate-twirl version
also covers approximate scalar or block designs with the aggregate PSD error
added to the same profile.  A matrix-Freedman variant covers adapted dependent
random contractions with predictable square profile
\(\|\sum_i\mathbb E(A_i^2\mid\mcF_{i-1})\|_{\op}=O(n/\log(eD))\), with the
same randomized sign-search certificate.  For Hermitian matrices with
conditionally independent centered upper-triangular entries, this profile is
checked by the maximum cumulative row variance, and whole-matrix cap tails give
sub-Weibull or finite-moment Wigner/GOE/GUE-type variants with one sign per
whole matrix.  The rectangular variant only asks for the corresponding
predictable left and right square profiles to be
\(O(n/\log(e(d_1+d_2)))\) and has the same active-menu Borel--Cantelli schedule.
For centered rectangular matrices with conditionally independent entries, this
reduces to the maximum cumulative row and column variance, giving one
discrepancy sign per whole random matrix after deterministic operator-norm
normalization.  A tail-capped version replaces the almost-sure operator-norm
bound by summable whole-matrix cap tails, with explicit sub-Weibull and
finite-moment cap schedules.
Its low-overlap specialization reduces sparse rectangular and bipartite support
models to the two weighted occupancy profiles on rows and columns.  A
realized-profile common-sign menu version direct-sums the active rectangular
library, pays only the scale \(\log(eJ_m(d_{m,1}+d_{m,2}))\), and removes the
active-menu factor from a sampled shared sign search.  The dependent
overlap-budget active-menu version needs only summable certification of the
realized left and right overlap budgets, and a tail-capped dependent version
adds sub-Weibull or finite-moment coefficients by paying cap-tail summability
plus the capped realized overlap certificate.  Scalar Bernstein controls
inhomogeneous coordinate supports and matrix Bernstein controls independent
random row/column subspaces.
Another tail-capped version allows unbounded random
low-rank or block coefficients by paying their operator-norm cap exceptions.  A
cap-selection schedule makes this explicit for sub-Weibull and finite-moment
operator tails, including the all-integer \(R>\Gamma+1\) and sparse
\(R>\Gamma\) cap-tail thresholds.  The rank-one
specialization covers sampled rectangular dyads from
arbitrary frame laws by the two expected weighted frame operators, including
Haar-isotropic dyads, randomized features, and nonuniform sampled atom
dictionaries, with the same active-menu schedule for countable dyad libraries.
A tail-capped version adds Gaussian, sub-Gaussian, sub-exponential, and
heavy-tailed scalar weights by paying only their cap-exception probabilities
plus the capped frame-Bernstein term.  Its cap-selection schedule gives
explicit sub-Weibull and finite-moment scalar-weight criteria for dyad/frame
menus with the same all-integer and sparse cap-tail thresholds; in the
homogeneous isotropic case this reduces to the scalar frame cap
\(q^2\lesssim d_*/\log(eD)\) and the explicit growth thresholds
\(R>\Gamma+1,\ c_0a>\Gamma+2\) on all integer dimensions, or
\(R>\Gamma,\ c_0a>\Gamma+1\) on sparse stretched-exponential schedules.  A
coordinate-sampled specialization replaces the Haar frame budget by the largest
row or column marginal load, giving the same schedules for iid sampled-entry,
random bipartite-edge, and sparse observation matrices.  Its predictable
version removes independence between samples: conditional marginal-load
envelopes cover adaptive sampling, without-replacement exposure, and permutation
or matching windows after the same cap-tail accounting.  A profile-capped form
keeps the full row and column load arrays, so nonuniform coordinate laws and
varying amplitude caps pay only their maximum cumulative capped marginal mass.
Its finite-menu schedule gives general, all-integer, sparse-dimension, and
sampled-certificate thresholds for active nonuniform coordinate libraries.
An entry-intensity specialization lets adaptive inhomogeneous sparse matrices
be checked directly from their predictable weighted row and column degrees.
A closed-form cap schedule then reduces sub-Weibull or finite-moment entry
weights to this maximum predictable degree budget.  A position-dependent cap
refinement charges rare heavy coordinates by
\(\sum q_{i,r,s}^2p_{i,r,s}\), rather than by a worst-case scalar cap times
the whole expected degree.
If the stream is supported in a common
random \(\kappa\)-dimensional corner, the same theorem uses
\(\log(e\kappa)\) and a \(\kappa\)-dimensional Freedman tail, with active menus
paying \(\sum_h\kappa_{m,h}\).  A two-scale square-function certificate keeps
one signing for a rank-\(\kappa\) core plus a full ambient residual, and its
predictable menu schedule pays the low-rank structural cost
\(\sum_h\kappa_{m,h}\) together with the residual structural cost \(J_mD_m\).
A finite-depth multiscale extension handles several low-rank layers plus one
ambient residual, with active structural cost
\(\sum_{h,q}\kappa_{m,h,q}+J_mD_m\).  A summable-depth extension permits
countably many low-rank layers under a summable \(\sum_q\sqrt{a_q}\) budget and
uses the corresponding \(\sum_{m,h,q}\) structural tail in active menus.  A
flag-compressible specialization turns any nested initial projection flag into
summable compression shells, so multiresolution, wavelet, spectral-window, and
hierarchical basis models only need summable shell square profiles.  A
compression-error version bounds those shell profiles by predictable
approximation errors, so common-flag spectral truncation or low-rank
approximation tails can be checked without computing the shell squares.  An
explicit tail-envelope version covers polynomial approximation tails against
stretched-exponential flag ranks and power tails against geometrically growing
rank windows, with active menus paying only summable envelope exceptions and
sampled-search losses.  The
corresponding dominated square-twirl
corollary makes Haar, finite-design, signed-permutation, and approximate
random-basis rotations inside a common random low-rank corner pay only the
\(\kappa\)-dimensional profile \(O(n/\log(e\kappa))\), and the matching
finite-menu dominated-twirl schedule gives simultaneous eventual signings for
active libraries at structural cost \(\sum_h\kappa_{m,h}\).  The predictable
dominated-twirl theorem also has a rank-sensitive common-corner version, so
online adaptive low-rank random-basis streams and their active-menu schedules
pay the same \(\sum_h\kappa_{m,h}\) structural count.  A finite-menu predictable
square-profile schedule gives simultaneous almost-sure signings for active
libraries of arbitrary adapted random-matrix streams after replacing \(D_m\) by
\(J_mD_m\), with common sampled-sign batches paying \(J_m2^{-M_m}\).  Its
predictable dominated-twirl
specialization lets the spectrum, block decomposition, and approximate
random-basis sampler be chosen from the past at each step; only the one-step
PSD twirl-error profile enters the same square-function budget.  The
corresponding Borel--Cantelli array version yields almost-sure eventual
Matrix-Spencer signings, and sampled sign certificates, when the row
twirl-budget tails and \(2^{-M_m}\) search losses are summable.  Explicit
schedule consequences include all integer dimensions once
\(n(D)\gtrsim\log^2(eD)\), and sparse stretched-exponential dimension schedules
under the one-row logarithmic-square threshold.  A finite-menu predictable
twirl schedule gives simultaneous eventual signings for active families of
online random-basis samplers after replacing the one-row structural count
\(D_m\) by \(J_mD_m\); common sampled-sign batches add only
\(J_m2^{-M_m}\).  Its
subspace-localized specialization covers adaptive coordinate supports, online
random subspaces, and conditionally sampled localized Wigner/GOE/GUE blocks
under the same predictable overlap scale.  A finite-menu subspace-localized
schedule makes these adaptive support/subspace models simultaneous on active
menus, with the same \(J_mD_m\) structural count and common sampled-sign
penalty \(J_m2^{-M_m}\).  A conditional scalar-multiplier
specialization also covers arbitrary
adversarial matrix directions after independent bounded thinning or amplitudes
with conditional second moments \(O(1/\log(eD))\), while the predictable
specialization allows the directions and amplitudes to be generated adaptively.
Its triangular-array form gives the same all-integer and sparse-dimension
almost-sure schedules under the half-budget conditional second-moment bound.
Finite menus of such adaptive scalar-multiplier streams are simultaneous under
the same schedules after replacing \(D_m\) by \(|\mathcal H_m|D_m\); common
sampled-sign batches add only \(|\mathcal H_m|2^{-M_m}\).  An exhaustive
countable-menu version controls the active prefix \(h\le J_m\) after replacing
\(D_m\) by \(J_mD_m\), hence every fixed adaptive scalar-multiplier stream is
eventually signed.
The same centered scalar-multiplier noise can be placed on
top of any core class whose signs are already controlled, including algebraic or
finite-lift cores, at the cost of an explicit matrix-Bernstein tail under the
core signing.  A matrix-valued martingale-difference version replaces
conditional independence by one-step centering; a scalar variance-domination
corollary only assumes
\(\mathbb E(Y_i^2\mid\mcF_{i-1})\preceq\theta_i\rho_i^2I\) and gives a
scalar-Bernstein tail for the random variance budget, while a matrix-profile
version allows anisotropic random PSD variance profiles and pays a
matrix-Bernstein budget tail.  A trace-budget variant gives a dimension-free
scalar budget exception when only the total trace of the random variance profile
is controlled; the predictable matrix martingale-noise branch now has an
almost-sure triangular-array schedule form for all integer dimensions and sparse
stretched-exponential dimension subsequences, with the scalar martingale-noise
case included.  Its active-prefix countable-menu version pays \(J_mD_m\) and
therefore eventually covers every fixed already-signed core-plus-noise model in
a countable library.  The sampled localized specialization gives the same
Borel--Cantelli schedules for adaptive coordinate-support and conditional
Haar-subspace matrix-noise arrays when the predictable weighted occupancy, or
the weighted average-rank budget, is \(O(n/\log(eD))\).  The same argument accepts
an active-prefix countable localized-noise library, replacing the variance
budget by the corresponding weighted-overlap or Haar average-rank budget for
each active model, and also accepts
independent trace envelopes, giving
explicit low-rank and effective-rank variance-profile tails from
\(\operatorname{rank}(V_i)\|V_i\|_{\op}\), and a Schatten-envelope corollary
interpolates between trace, Frobenius-plus-rank, and operator-plus-rank control.
A quantile trace-envelope variant only needs a high-probability bound, or a
finite moment, for the total weighted trace budget; it therefore covers
heavy-tailed trace-class random covariance and spike profiles outside the
Bernstein regime.  Its summable version gives almost-sure eventual
Spencer-scale signings for triangular arrays whose trace-budget quantiles are
at the \(n/\log(eD)\) scale; a normalized-moment specialization makes this
checkable by summing moments of \(S_m/(n_m/\log(eD_m))\), and a termwise
moment version reduces this to individual heavy-tail atom moments without
dependence assumptions.  A termwise-tail version only asks for summable
individual truncation failures, and a scale-mixture specialization converts
this into an aggregate tail check for deterministic atom weights.  A
common-scale version removes the atom-count union factor when all atoms share
one heavy-tailed multiplier, and finite or countable grouped-scale versions pay
one tail per volatility group.  An optimized tail-profile form allows exact
survival functions and log-Weibull/lognormal tails, while the heterogeneous
polynomial and sub-Weibull specializations have optimized share-cost criteria.
Profile-menu aggregation groups volatility components by tail class and optimizes
only over class-level trace shares.  A finite heterogeneous polynomial-index menu
is controlled by its worst tail index once the budget dominates the largest
class-count/trace product.  A countable weighted heterogeneous version gives the
same worst-index budget rule for power-decaying coefficient and trace tails
whenever a summable class-share distribution dominates the trace tail.  For a
finite heterogeneous stretched-exponential profile menu, the corresponding
logarithmic budget rule has exponent \(\delta=c_0a^{\beta_0}\) and pays the
active class-count exponent \(w\), giving \(\delta>p+w+1\) with the same
log-log endpoint convention.  A countable weighted stretched-exponential
version removes this count penalty under class-index trace decay faster than the
summable share weights, giving \(\delta>p+1\).  The analogous finite
log-Weibull menu has the same \(\delta>p+w+1\) rule with a subpolynomial
budget margin
\(\exp\{a(\Lambda_m+bM_m)^{1/\beta_0}\}\), while a countable weighted
log-Weibull version removes the count penalty once the residual class-index
log-Weibull series is summable.  For a
common polynomial tail index this infimum has a closed-form effective class
mass, and the resulting polynomial/logarithmic schedule has the rule
\(\alpha b>(\alpha+1)p+1\).  If the number of active classes grows like
\(m^w\), the finite-menu rule becomes
\(\alpha b>p+\alpha u+(\alpha+1)w+1\).  For countably many classes with
power-decaying coefficients and trace masses, the rule becomes
\(\alpha b>p+\alpha u+1\) once the class-index decay is summable.  When all
groups share one tail profile, proportional mass shares give a single aggregate
criterion.  A
stretched-exponential schedule gives the concrete rule \(ca^\beta>p+1\), with a
log-log endpoint \(q-ca^\beta b<-1\), while a log-Weibull schedule gives the
subpolynomial rule \(ca^\beta>p+1\).  A regularly varying schedule with log
corrections gives the matching exponent rule \(\alpha a>p+1\), with endpoint
condition \(q+\gamma-\alpha b<-1\); the common-index polynomial-tail form also
has an explicit effective group-tail mass.
A sub-gamma trace-envelope variant removes the bounded-envelope hypothesis and
only asks for a Bernstein-type tail for the total weighted envelope; a termwise
version verifies this from independent sub-gamma envelopes with summable
variance proxies.
A gamma-dominated specialization gives explicit chi-square and Wishart-type
trace-envelope tails.
Gamma-atom specializations allow heterogeneous finite or countable gamma and
exponential atom scales in the same weighted trace envelope.
A compound-Poisson dominated specialization covers random counts of bounded
trace-envelope atoms and bounded spike-count variance profiles.
Atomized bounded-envelope specializations give finite or countable
Bernoulli/binomial spike and low-rank atom sums an atom-level jump parameter.
A mixed-atom calculus combines bounded, gamma, and compound-Poisson trace
atoms in the same envelope with one Bernstein parameter.
A weighted chi-square specialization covers unequal Gaussian directions and
real Wishart trace envelopes.
A weighted exponential specialization covers complex Gaussian directions and
complex Wishart trace envelopes.
The same calculation with deterministic shifts and residual constants covers
noncentral and correlated real or complex Gaussian/Wishart trace envelopes.
A scalar specialization
allows predictable noise directions with the same matrix-Freedman tail.  Localized
versions control online coordinate-supported, subspace-localized, or
frame-localized martingale noise by realized predictable overlap, conditional
predictable overlap, or weighted-overlap operators.  Thus random localized,
projection, frame, and sparse-support
ensembles are covered whenever their weighted empirical support overlap, or
weighted frame operator, has norm \(O(n/\log(eD))\).
For independent inhomogeneous random coordinate supports, scalar Bernstein
controls the weighted coordinate occupancies; in the homogeneous case support
density \(O(1/\log(eD))\) is enough for a high-probability \(O(\sqrt n)\)
theorem.
The same variance route also has a noncommutative concentration form: independent
random subspace projections with
\(\|\sum_i\rho_i^2\mathbb E P_i\|_{\op}=O(n/\log(eD))\) have the required
weighted empirical frame bound with matrix-Bernstein probability; Haar ranks
enter through \(D^{-1}\sum_i\rho_i^2r_i\).  The same weighted-support theorem
allows independent random amplitudes: scalar Bernstein controls
\(\max_j\sum_i\rho_i^2{\bf 1}_{\{j\in S_i\}}\) from
\(\max_j\sum_i p_{ij}\mathbb E\rho_i^2\), while matrix Bernstein controls
\(\|\sum_i\rho_i^2P_i\|_{\op}\) from
\(\|\sum_i\mathbb E(\rho_i^2P_i)\|_{\op}\).  The coordinate and subspace
Bernstein tails also give an almost-sure triangular-array version whenever the
row failure probabilities are summable.  Under the half-budget condition these
tails are summable along all integer dimensions once
\(n_D\ge a\log^2(eD)\) with \(a>32/(3C_0)\), and along sparse
stretched-exponential dimension schedules once \(a>16/(3C_0)\); the explicit
random-amplitude schedule form also covers arbitrary growth by paying the
\(\log(e+m)\) summability margin.  A fixed orthogonal block decomposition has
the same random-amplitude support theorem with \(D\) in the union bound replaced
by the block count \(B\), and its schedules depend on \(\log(eB_m)\) in the
summability margin.  A random-profile version permits conditional inclusion
profiles \(Q_{ij}\) or \(Q_{i\alpha}\), replacing deterministic budgets by
their \(\rho_i^2\)-weighted expectations and justifying random support-size
budgets even when sizes and amplitudes are correlated.  Its triangular-array
form is variance-sensitive: summability of
\[
  N_m\exp\{-t_m^2/(2(\sigma_m^2+t_m/3))\},\qquad N_m=D_m\text{ or }B_m,
\]
gives eventual almost-sure profile signings, and the half-budget case is
summable under
\[
  n_m\ge \frac{16}{3C_0}\log(eD_m)
  \bigl(\log(eN_m)+(1+\gamma)\log(e+m)\bigr).
\]
A predictable profile/Freedman version further permits sequentially sampled
coordinate or block supports whose conditional inclusion profiles are
\(\mcF_{i-1}\)-measurable; the same conclusion holds outside the realized
predictable-overlap budget exception, and sampled signs still add only
\(2^{-M}\).  Its explicit schedule form uses \(N_m=D_m\) or \(B_m\): arbitrary
growth needs
\[
  n_m\ge \frac{16}{3C_0}\log(eD_m)
  \bigl(\log(eN_m)+(1+\gamma)\log(e+m)\bigr),
\]
while \(N_m\le(eD_m)^\Gamma\) gives the all-integer threshold
\(a>16(\Gamma+1)/(3C_0)\) and the sparse stretched-exponential threshold
\(a>16\Gamma/(3C_0)\), up to the summable predictable-budget exceptions.
A finite-menu predictable-profile schedule gives simultaneous eventual signings
for many online coordinate/block support models after replacing \(N_m\) by
\(L_m=\sum_hN_{m,h}\) and summing the predictable-budget exceptions.
An exhaustive countable-menu schedule controls the active prefix \(h\le J_m\)
after replacing \(N_m\) by \(L_m^{(J)}=\sum_{h\le J_m}N_{m,h}\), hence every
fixed member of a countable online support library is eventually signed.
A deterministic-envelope specialization has no such exception: if the adaptive
conditional inclusion probabilities are bounded by fixed \(p_{ij}\) or
\(q_{i\alpha}\), and the predictable amplitudes are bounded by deterministic
caps \(a_i\), the budget is simply
\(\max_j\sum_i a_i^2p_{ij}\) or
\(\max_\alpha\sum_i a_i^2q_{i\alpha}\).
The deterministic-envelope schedule form imposes only this fixed half-budget,
so the arbitrary-growth, all-integer, sparse, and sampled-certificate thresholds
hold with no budget-exception summability term.
For finite menus of deterministic-envelope coordinate/block models, the same
simultaneous conclusion holds after replacing the structural count by
\(L_m=\sum_{h\in\mathcal H_m}N_{m,h}\), and common sampled-sign batches pay only
\(|\mathcal H_m|2^{-M_m}\).
For a countable deterministic-envelope library, the exhaustive-prefix version
controls \(h\le J_m\) at cost \(L_m^{(J)}=\sum_{h\le J_m}N_{m,h}\), so every
fixed envelope model is eventually signed without adding any budget-exception
series.
Also,
\(M\) independent sign trials add only \(2^{-M}\) to the finite failure
probability, so logarithmically many trials give eventual sampled certificates:
\((1+\delta)\log_2(eD)\) trials along all integer dimensions and
\((1+\delta)\log_2(e+m)\) trials along sparse schedules.
For arbitrary triangular arrays, the same proof gives a direct growth rule:
\[
  n_m\ge \frac{16}{3C_0}\log(eD_m)
  \bigl(\log(eD_m)+(1+\gamma)\log(e+m)\bigr)
\]
makes both the matrix failure and, with \(M_m\ge(1+\delta)\log_2(e+m)\), the
sampled-search failure summable.
The signed-core martingale-noise version has the analogous conditional forms:
coordinate-dependent adaptive inclusion probabilities and conditional Haar
ranks control the same core signing through a random weighted budget \(W\), with
failure term \(\mathbb P\{W>v\}\), and through the overlap
\(\|\sum_i\rho_i^2\mathbb E(P_i\mid\mcF_{i-1})\|_{\op}\), recovering the
unweighted deterministic form when \(\rho_i=1\) and \(W\) is bounded.  Its
triangular-array form gives all-integer and sparse stretched-exponential
almost-sure schedules under the same \(n/\log(eD)\) weighted-overlap or
average-rank scale.  For
coordinate supports with deterministic, or independent random, inclusion
profiles and independent random amplitudes, a scalar Bernstein union bound makes
this budget tail explicit; in the homogeneous support-density case the effective
overlap is \(pA\), where \(A=\sum_i\mathbb E\rho_i^2\).  The same coordinate
argument works for fixed orthogonal block decompositions, replacing the
coordinate union by a union over the block count.  For fixed-rank Haar subspaces with independent
random amplitudes the analogous scalar Bernstein bound has mean
\(\sum_i(r_i/D)\mathbb E\rho_i^2\), with no coordinate-union loss in the budget
exception.  The same budget-tail statement holds for independent random ranks
\(R_i\), replacing \(r_i/D\) by \(R_i/D\) inside the scalar Bernstein mean and
variance proxy.  More generally, any conditionally isotropic projection law has
the same bound after replacing \(R_i/D\) by the random intensity \(\alpha_i\).
The equality can be weakened to the spectral domination
\(\mathbb E(P_i\mid\mcF_{i-1})\preceq\alpha_iI\), which covers
leverage-bounded and approximate-design projection laws.  The resulting
random-amplitude localized menu has a Borel--Cantelli array theorem: if the
scalar Bernstein budget exceptions and the matrix-Freedman row tails are
summable, the same core signing is eventually \(O(\sqrt n)\), with explicit
all-integer and sparse dimension schedules under half-budget \(n/\log(eD)\)
profiles.  Its countable-menu form controls active prefixes by replacing
\(N_m\) with \(\sum_{h\le J_m}N_{m,h}\), hence every fixed localized
random-amplitude model is eventually covered.  In the homogeneous subcases this becomes an effective-energy
criterion: the budget scale is \(\beta H\) and the scalar budget exception is
summable once \(H=\sum_i\mathbb E\rho_i^2\) dominates \(\log(eD)\) by the
appropriate union-count margin.
If the signed cores share one active-prefix signing, a common-sign countable
form preserves that same signing through all random-amplitude localized
perturbations with the same \(\sum_{h\le J_m}N_{m,h}\) payment.
A finite-lift common-menu corollary supplies this shared core signing
deterministically whenever the active-prefix direct-sum lift entropy is \(O(n)\).
Thus finite-lift core libraries absorb localized random-amplitude noise under
the same active-prefix union-count and homogeneous energy schedules.
The homogeneous criterion also has a countable active-prefix version with
\(\Lambda_m^{(J)}=\sum_{h\le J_m}N_{m,h}\), giving all-integer coverage for
\(R>8(\Gamma+1)/3\) when \(\Lambda_D^{(J)}\le(eD)^\Gamma\), and sparse
coverage for \(R>8\Gamma/3\).
The same energy rule has a common-sign version, so one active-prefix core
signing remains valid for every homogeneous random-amplitude perturbation in
the prefix.
The coordinate and Haar branches can also be mixed
sample-by-sample: if \(M_i\in\{0,1\}\) selects the coordinate model and
\(\beta_{ij}=M_iQ_{ij}+(1-M_i)R_i/D\), then the same tail holds with
\(\max_j\sum_i\mathbb E(\beta_{ij}\rho_i^2)\) and
\(\max_j\sum_i\mathbb E(\beta_{ij}^2\rho_i^2)\).
\end{enumerate}

\section{Discrepancy inputs}

We use two discrepancy inputs. The first is the rectangular vector Spencer
theorem. All logarithms are natural.

\begin{theorem}[Vector Spencer]\label{thm:vspencer}
There is a universal constant $C_S$ such that, for every
$x_1,\ldots,x_n\in\C^m$ with $\norm{x_i}_\infty\le 1$, there are signs
$\eps\in\{-1,1\}^n$ satisfying
\[
\norm{\sum_{i=1}^n \eps_i x_i}_\infty
\le C_S \sqrt{n\log\left(e+\frac{m}{n}\right)} .
\]
\end{theorem}

The complex case follows from the real theorem by applying it to real and
imaginary parts and absorbing the harmless factor of two into the constant.

The second input is the block-diagonal Matrix Spencer theorem. For block size
$h$ it was first established by Levy, Ramadas and Rothvoss \cite{LRR}, with the
restriction $h\le n$, and the restriction was later removed by Dadush, Jiang and
Reis \cite{DJR}. We state it in the standard form with total matrix dimension
$R$ and maximum block size $h$.

\begin{theorem}[Block-diagonal Matrix Spencer]\label{thm:djr}
There is a universal constant $C_{BD}$ such that the following holds. Let
$\C^R=\bigoplus_\alpha \C^{r_\alpha}$ be a fixed block decomposition with
$h=\max_\alpha r_\alpha$. Let $Y_1,\ldots,Y_n$ be self-adjoint matrices that
are block diagonal with respect to this decomposition and satisfy
$\norm{Y_i}_{\op}\le 1$. Then there are signs $\eps\in\{-1,1\}^n$ such that
\[
\norm{\sum_{i=1}^n \eps_i Y_i}_{\op}
\le C_{BD}\sqrt{n\log\left(e+\frac{hR}{n}\right)} .
\]
Equivalently, for $\mcB=\bigoplus_\alpha \M_{r_\alpha}(\C)$,
\[
\norm{\sum_{i=1}^n \eps_i Y_i}_{\mcB}
\le C_{BD}\sqrt{n\log\left(e+\frac{\Ent(\mcB)}{n}\right)} .
\]
\end{theorem}

\begin{remark}
Dadush, Jiang and Reis prove this as an application of their partial-coloring
framework; in their notation the logarithmic factor is of the form
$\log(hm/n)$ for $m\times m$ block-diagonal matrices with block size at most
$h$ \cite{DJR}. Their $m$ is the total matrix dimension, which is our $R$, so
$\log(hm/n)=\log(hR/n)$. The displayed version inserts $e+$ to cover all
regimes.
\end{remark}

\begin{remark}[Realification]\label{rem:realification}
Theorem \ref{thm:djr} is stated for real symmetric matrices. One passes to
$\C$ through the realification
\[
\iota(A)=\begin{pmatrix}\operatorname{Re}A & -\operatorname{Im}A\\
\operatorname{Im}A & \operatorname{Re}A\end{pmatrix},
\]
which is $\R$-linear, sends self-adjoint matrices to real symmetric ones,
preserves the operator norm, and respects the block decomposition. It doubles
the block size and the total dimension, so the entropy is multiplied by four;
this is absorbed into the universal constant.
\end{remark}

For non-self-adjoint block-diagonal matrices we use Hermitian dilation. For a
matrix $Y$, set
\[
\mathscr H(Y)=\begin{pmatrix}0&Y\\Y^*&0\end{pmatrix} .
\]
Then $\mathscr H(Y)$ is self-adjoint and $\norm{\mathscr H(Y)}_{\op}=\norm{Y}_{\op}$.
Applying Theorem \ref{thm:djr} to the dilations doubles all block sizes
and changes only the universal constant.

\begin{corollary}[Non-self-adjoint block diagonal signing]\label{cor:nsa}
Let $Y_1,\ldots,Y_n\in\mcB$ be arbitrary complex block-diagonal matrices with
$\norm{Y_i}_{\mcB}\le 1$. Then
\[
\min_{\eps\in\{-1,1\}^n}
\norm{\sum_{i=1}^n \eps_iY_i}_{\mcB}
\le C\sqrt{n\log\left(e+\frac{\Ent(\mcB)}{n}\right)} .
\]
\end{corollary}

\section{Finite block and sampled multiplier lifts}

The lift formulation is tuple-level: for a given input collection it is enough
to choose bounded lifts of the individual matrices and a contraction on their
span.

\begin{definition}[Finite block lift]
Let $A_1,\ldots,A_n\in\M_D(\C)$ and
$\mcB=\bigoplus_{\alpha=1}^m \M_{r_\alpha}(\C)$. We say that the tuple admits a
$(K,\Lambda)$-lift to $\mcB$ if there exist elements $Y_i\in\mcB$ and a fixed
linear map
\[
\Pi:\Span\{Y_1,\ldots,Y_n\}\to\M_D(\C)
\]
such that
\[
\Pi(Y_i)=A_i,
\qquad
\norm{\Pi(Y)}_{\op}\le \Lambda\norm{Y}_{\mcB}
\quad (Y\in\Span\{Y_1,\ldots,Y_n\}),
\]
and $\norm{Y_i}_{\mcB}\le K$ for all $i$.
\end{definition}

The map $\Pi$ is usually a compression, a quotient map, or a fixed bounded
linear image. Contractive lifts correspond to $\Lambda=1$.

\begin{proof}[Proof of Theorem \ref{thm:main}]
If $K=0$ then all $Y_i=0$, so all $A_i=0$ and the bound is trivial; assume
$K>0$. First assume that the $Y_i$ are self-adjoint. Apply Theorem
\ref{thm:djr} to $Y_i/K$. There are signs for which
\[
\norm{\sum_i \eps_i Y_i}_{\mcB}
\le C K\sqrt{n\log\left(e+\frac{\Ent(\mcB)}{n}\right)} .
\]
Since $\Pi$ is linear and has norm at most $\Lambda$ on the relevant span,
\[
\norm{\sum_i \eps_i A_i}_{\op}
= \norm{\Pi\left(\sum_i \eps_iY_i\right)}_{\op}
\le \Lambda\norm{\sum_i \eps_iY_i}_{\mcB} .
\]
This proves the self-adjoint case. In the arbitrary complex case, replace each
block $Y_{i,\alpha}$ by its Hermitian dilation. The block sizes and total
reduced dimension are both doubled, so the entropy changes by a factor of four,
and the identity
$\norm{\sum_i\eps_i\mathscr H(Y_i)}=\norm{\sum_i\eps_iY_i}$ transfers the bound.
Since Theorem \ref{thm:djr} is over $\R$, one further applies the realification
$\iota$ of Remark \ref{rem:realification} before invoking it, which again
doubles the block size and the total dimension. The general complex case thus
uses both Hermitian dilation and realification; each doubles the dimension, so
the entropy grows by a factor of at most $16$, absorbed into the universal
constant $C$.
\end{proof}

\begin{corollary}[Common signs for finite menus of finite-lift tuples]
\label{cor:finite-lift-common-menu}
Let \(\mathcal H\) be finite and nonempty.  For each \(h\in\mathcal H\), let
\(A_1^{(h)},\ldots,A_n^{(h)}\in \M_{D_h}(\C)\) be contractions which admit a
\((K_h,\Lambda_h)\)-lift to
\[
\mcB_h=\bigoplus_{\alpha=1}^{m_h}\M_{r_{h,\alpha}}(\C).
\]
Put
\[
\begin{aligned}
M&=\max_{h\in\mathcal H}\Lambda_hK_h,
&
h_*&=\max_{h,\alpha}r_{h,\alpha},\\
R_*&=\sum_{h\in\mathcal H}\sum_{\alpha=1}^{m_h}r_{h,\alpha},
&
E_*&=h_*R_*.
\end{aligned}
\]
Then there is one sign vector \(\eps\in\{-1,1\}^n\) such that
\[
\max_{h\in\mathcal H}
\norm{\sum_{i=1}^n\eps_iA_i^{(h)}}_{\op}
\le
CM\sqrt{n\log\left(e+\frac{E_*}{n}\right)} .
\]
If all tuples and their chosen lifts are self-adjoint, the same bound holds for
self-adjoint discrepancy.  In particular, if \(E_*\le C_E n\) and \(M\le M_0\),
then the whole finite menu has one \(O_{C_E,M_0}(\sqrt n)\) signing.

Consequently, in a triangular-array countable active-prefix library, if each row
\(m\) and active prefix \(h\le J_m\) has contraction cores
\(B_{m,1}^{(h)},\ldots,B_{m,n_m}^{(h)}\) admitting such finite lifts with
\(M_m^{(J)}\le M_0\) and \(E_m^{(J)}\le C_E n_m\), then the common-core signing
hypotheses in Corollaries
\ref{cor:core-plus-common-sign-countable-localized-random-amplitude-menu-borel}
and
\ref{cor:core-plus-common-sign-homogeneous-localized-random-amplitude-schedules}
hold deterministically, with no exceptional core event.
\end{corollary}

\begin{proof}
For \(K_h>0\), choose lifts \(Y_i^{(h)}\in\mcB_h\) and set
\(Z_i^{(h)}=Y_i^{(h)}/K_h\).  If \(K_h=0\), then
\(A_i^{(h)}=0\) for all \(i\), and we set \(Z_i^{(h)}=0\).  Form the direct-sum
algebra
\[
\mcB_*=\bigoplus_{h\in\mathcal H}\mcB_h
\]
and \(Z_i=(Z_i^{(h)})_{h\in\mathcal H}\in\mcB_*\).  Then
\(\norm{Z_i}_{\mcB_*}\le1\), and
\(\Ent(\mcB_*)=h_*R_*=E_*\).  Define, on the span of the \(Z_i\), the map whose
\(h\)-component is
\[
\Theta(W)^{(h)}
=
\begin{cases}
K_h\Pi_h(W^{(h)}), & K_h>0,\\
0, & K_h=0,
\end{cases}
\qquad
\Theta(W)\in\bigoplus_{h\in\mathcal H}\M_{D_h}(\C).
\]
For \(W\in\Span\{Z_1,\ldots,Z_n\}\),
\[
\norm{\Theta(W)}_{\op}
\le \max_{h\in\mathcal H}\Lambda_hK_h\,\norm{W}_{\mcB_*}
=M\norm{W}_{\mcB_*},
\]
and \(\Theta(Z_i)=\bigoplus_{h\in\mathcal H}A_i^{(h)}\).  Applying Theorem
\ref{thm:main} to the block-diagonal target tuple
\(\bigoplus_h A_i^{(h)}\) gives the displayed estimate, because the operator
norm in the target direct sum is the maximum over \(h\).  The self-adjoint case
uses self-adjoint direct sums and the self-adjoint part of Theorem
\ref{thm:main}.  The final assertion applies the finite-menu statement row by
row to the deterministic active prefix.
\end{proof}

\subsection{Sampled multiplier lifts}

Toeplitz arguments naturally produce continuous multiplier lifts rather than
finite block lifts at the first step. The following sampling principle converts
them into finite block lifts.

\begin{definition}[Norming set]
Let $\Omega$ be a compact space and let
$\mcF\subset C(\Omega;\M_b(\C))$ be a finite-dimensional linear space. A finite
set $Q\subset\Omega$ is $L_Q$-norming for $\mcF$ if
\[
\norm{F}_{L^\infty(\Omega;\M_b)}
\le L_Q\max_{\omega\in Q}\norm{F(\omega)}_{\op}
\qquad(F\in\mcF).
\]
\end{definition}

\begin{definition}[Sampled multiplier lift]
A tuple $A_1,\ldots,A_n\in\M_D(\C)$ admits a
$(K,\Lambda,L_Q,Q,b)$ sampled multiplier lift if there are functions
$F_i\in\mcF\subset C(\Omega;\M_b(\C))$ and a fixed linear map
\[
\Pi:\Span\{F_1,\ldots,F_n\}\to\M_D(\C)
\]
such that
\[
\Pi(F_i)=A_i,
\qquad
\norm{\Pi(F)}_{\op}\le \Lambda\norm{F}_{L^\infty(\Omega;\M_b)},
\]
$\norm{F_i}_{L^\infty}\le K$, and $Q$ is $L_Q$-norming for
$\Span\{F_1,\ldots,F_n\}$.
\end{definition}

\begin{theorem}[Sampled multiplier Matrix Spencer]\label{thm:sampled}
If $A_1,\ldots,A_n$ admit a $(K,\Lambda,L_Q,Q,b)$ sampled multiplier lift, then
\[
\discop(A_1,\ldots,A_n)
\le C K\Lambda L_Q
\sqrt{n\log\left(e+\frac{b^2|Q|}{n}\right)} .
\]
For $b=1$ this is the vector-Spencer scale
$CK\Lambda L_Q\sqrt{n\log(e+|Q|/n)}$.
\end{theorem}

\begin{proof}
Evaluate each $F_i$ on $Q$ and form
\[
Y_i=\bigoplus_{\omega\in Q}F_i(\omega)
\in \bigoplus_{\omega\in Q}\M_b(\C).
\]
This finite block algebra has $h=b$, $R=b|Q|$, and entropy $b^2|Q|$. By
Corollary \ref{cor:nsa}, there are signs such that
\[
\max_{\omega\in Q}
\norm{\sum_i\eps_iF_i(\omega)}_{\op}
\le C K\sqrt{n\log\left(e+\frac{b^2|Q|}{n}\right)} .
\]
Since $Q$ is $L_Q$-norming,
\[
\norm{\sum_i\eps_iF_i}_{L^\infty}
\le C K L_Q\sqrt{n\log\left(e+\frac{b^2|Q|}{n}\right)} .
\]
Applying $\Pi$ gives the result.
\end{proof}

\subsection{A grid lemma for trigonometric polynomial lifts}

We shall use the following elementary norming lemma. It is stated for
matrix-valued trigonometric polynomials because this is the form needed for
block Toeplitz matrices.

\begin{lemma}[Grid norming]\label{lem:grid}
Fix $r\ge 1$. Let $P:\T^r\to\M_b(\C)$ be a matrix-valued trigonometric
polynomial with coordinate degrees at most $N_1,\ldots,N_r$. For each $j$, let
$G_j$ be a uniform grid in $\T$ of size at least $8rN_j$ when $N_j\ge 1$, and of
size $1$ when $N_j=0$. Put $G=G_1\times\cdots\times G_r$. Then
\[
\norm{P}_{L^\infty(\T^r;\M_b)}
\le 2\max_{\theta\in G}\norm{P(\theta)}_{\op} .
\]
In particular, $|G|=O_r(\prod_{j=1}^r(1+N_j))$.
\end{lemma}

\begin{proof}
The Banach-valued Bernstein inequality gives
$\norm{\partial_j P}_\infty\le N_j\norm{P}_\infty$ for each coordinate. One way
to see this is to apply the scalar Bernstein inequality to $\varphi(P(\cdot))$
for norm-one linear functionals $\varphi$ on $\M_b$. Choose $\theta^*$ with
$\norm{P(\theta^*)}\ge\norm{P}_\infty-\delta$, and choose $\theta\in G$ with
$|\theta_j-\theta^*_j|\le \pi/|G_j|$ for every $j$. Then
\[
\norm{P(\theta^*)-P(\theta)}
\le \sum_{j=1}^r \frac{\pi}{|G_j|}N_j\norm{P}_\infty
\le \frac{\pi}{8}\norm{P}_\infty .
\]
Letting $\delta\downarrow 0$ gives
$\norm{P}_\infty\le (1-\pi/8)^{-1}\max_G\norm{P}\le 2\max_G\norm{P}$.
\end{proof}

\section{Finite-dimensional C*-algebras and group convolution}

The cleanest applications are matrices that already lie in a common
finite-dimensional C*-algebra.

\begin{theorem}[Common finite-dimensional C*-algebra]\label{thm:cstar}
Let $\mcA\subset\M_D(\C)$ be a finite-dimensional C*-algebra. On its support,
and omitting any zero summand, write its representation after unitary conjugacy
as
\[
\mcA\simeq \bigoplus_{\alpha=1}^m
\left(\M_{r_\alpha}(\C)\otimes I_{m_\alpha}\right).
\]
Let $h=\max_\alpha r_\alpha$ and $R=\sum_\alpha r_\alpha$. If
$A_1,\ldots,A_n\in\mcA$ satisfy $\norm{A_i}_{\op}\le 1$, then
\[
\discop(A_1,\ldots,A_n)
\le C\sqrt{n\log\left(e+\frac{hR}{n}\right)} .
\]
If the matrices are self-adjoint, this is a self-adjoint discrepancy bound.
\end{theorem}

\begin{proof}
An element of $\mcA$ is unitarily equivalent to
$0\oplus\bigoplus_\alpha(B_\alpha\otimes I_{m_\alpha})$, with the initial zero
summand absent in the unital-on-$\C^D$ case. Delete the zero summand and the
multiplicity identities to obtain the reduced algebra
$\mcB=\bigoplus_\alpha\M_{r_\alpha}(\C)$. The deletion map and the inverse map
$\bigoplus B_\alpha\mapsto\bigoplus(B_\alpha\otimes I_{m_\alpha})$ are isometric
for the operator norm. Theorem \ref{thm:main} applies with
$K=\Lambda=1$.
\end{proof}

\begin{theorem}[Akbas--Sra algebraic Matrix Spencer]\label{thm:akbassra}
For every \(C_0\ge1\) there is a constant \(C(C_0)\) with the following
property. Let \(\mcA\) be a finite-dimensional C*-algebra with
\(\dim_\C\mcA\le C_0n\). If
\(A_1,\ldots,A_n\in\mcA\) are contractions, then
\[
\discop(A_1,\ldots,A_n)\le C(C_0)\sqrt n .
\]
The self-adjoint version holds for self-adjoint tuples.
\end{theorem}

\begin{proof}
This is the algebraic Matrix Spencer theorem of Akbas and Sra
\cite[Thm.~1.3]{AkbasSra}.  Their proof reduces a finite-dimensional
C*-algebra by Wedderburn decomposition to common block-diagonal matrices and
uses a multiscale partial-coloring argument whose complexity is the intrinsic
square-sum \(\sum_\alpha r_\alpha^2=\dim_\C\mcA\), rather than the cruder
\((\max_\alpha r_\alpha)\sum_\alpha r_\alpha\) used in
Theorem~\ref{thm:cstar}.
\end{proof}

\begin{theorem}[Akbas--Sra low-Frobenius perturbation]\label{thm:akbassra-pert}
For every \(C_0\ge1\) there is \(C(C_0)\) with the following property.  Let
\[
X_i=B_i+L_i,\qquad i=1,\ldots,n,
\]
be self-adjoint matrices.  Suppose \(B_i\in\mcA\), where \(\mcA\) is a
finite-dimensional C*-algebra with \(\dim_\C\mcA\le C_0n\), and suppose
\[
\|B_i\|_{\op}\le1,\qquad
\|L_i\|_{\op}\le1,\qquad
\|L_i\|_{\mathrm F}^2\le \frac{n}{\log^3(en)}
\]
for all \(i\).  Then there are signs \(\eps_i\in\{-1,1\}\) such that
\[
\left\|\sum_{i=1}^n\eps_iX_i\right\|_{\op}\le C(C_0)\sqrt n .
\]
\end{theorem}

\begin{proof}
This is \cite[Thm.~1.5]{AkbasSra}.  It combines their algebraic partial-coloring
body with the Bansal--Jiang--Meka low-Frobenius/rank-constrained Matrix Spencer
input.
\end{proof}

\begin{corollary}[Low-rank perturbations of algebraic cores]
\label{cor:akbassra-lowrank}
For every \(C_0\ge1\) there is \(C(C_0)\) with the following property.  Let
\[
X_i=B_i+L_i,\qquad i=1,\ldots,n,
\]
be self-adjoint matrices.  Suppose \(B_i\in\mcA\), where \(\mcA\) is a
finite-dimensional C*-algebra with \(\dim_\C\mcA\le C_0n\), and suppose
\[
\|B_i\|_{\op}\le1,\qquad
\|L_i\|_{\op}\le1,\qquad
\operatorname{rank}(L_i)\le \frac{n}{\log^3(en)}
\]
for all \(i\).  Then there are signs \(\eps_i\in\{-1,1\}\) such that
\[
\left\|\sum_{i=1}^n\eps_iX_i\right\|_{\op}\le C(C_0)\sqrt n .
\]
\end{corollary}

\begin{proof}
Since \(L_i\) is self-adjoint,
\[
\|L_i\|_{\mathrm F}^2
=\sum_j s_j(L_i)^2
\le \operatorname{rank}(L_i)\,\|L_i\|_{\op}^2
\le \frac{n}{\log^3(en)}.
\]
The result is therefore Theorem~\ref{thm:akbassra-pert}.
\end{proof}

\begin{corollary}[Simultaneous diagonalization]
If $A_1,\ldots,A_n$ are normal matrices diagonalized by a common unitary and
$\norm{A_i}_{\op}\le 1$, then
\[
\discop(A_1,\ldots,A_n)
\le C\sqrt{n\log\left(e+\frac{D}{n}\right)} .
\]
Thus the conjectured $O(\sqrt n)$ bound holds when $D\le n$.
\end{corollary}

\begin{corollary}[Common bounded-block diagonalization]
Suppose all $A_i$ are simultaneously unitarily block diagonal with common block
sizes at most $b$. Let $D$ be the ambient dimension. Then
\[
\discop(A_1,\ldots,A_n)
\le C\sqrt{n\log\left(e+\frac{bD}{n}\right)} .
\]
For fixed $b$ and $D\le n$, this is $O_b(\sqrt n)$.
\end{corollary}

\begin{corollary}[Commutative association schemes]
Let $\mcA\subset\M_D(\C)$ be a commutative self-adjoint matrix algebra with $s$
minimal spectral projections. If $A_1,\ldots,A_n\in\mcA$ are contractions, then
\[
\discop(A_1,\ldots,A_n)
\le C\sqrt{n\log\left(e+\frac{s}{n}\right)} .
\]
In particular, any family contained in the Bose--Mesner algebra of a symmetric
association scheme with $s$ primitive idempotents satisfies the Matrix Spencer
bound whenever $s=O(n)$, independently of the number of vertices; a $d$-class
symmetric scheme contributes $s=d+1$ primitive idempotents.
\end{corollary}

\subsection{Finite group convolution}

Let $G$ be a finite group. A left convolution operator on $\ell^2(G)$ has the
form
\[
L_f\xi(x)=\sum_{y\in G} f(xy^{-1})\xi(y).
\]
By the finite Fourier transform,
\[
L_f \simeq \bigoplus_{\pi\in\widehat G}
\left(\widehat f(\pi)\otimes I_{d_\pi}\right),
\qquad d_\pi=\dim\pi,
\]
and
$\norm{L_f}_{\op}=\max_{\pi\in\widehat G}\norm{\widehat f(\pi)}_{\op}$.
Thus the reduced convolution algebra is
$\bigoplus_{\pi\in\widehat G}\M_{d_\pi}(\C)$.

\begin{theorem}[Finite group convolution]\label{thm:group}
Let $G$ be finite, let $q(G)=\max_{\pi\in\widehat G}d_\pi$, and put
$S(G)=\sum_{\pi\in\widehat G}d_\pi$. If
$L_{f_1},\ldots,L_{f_n}$ are left convolution operators on $\ell^2(G)$ with
$\norm{L_{f_i}}_{\op}\le 1$, then
\[
\discop(L_{f_1},\ldots,L_{f_n})
\le C\sqrt{n\log\left(e+\frac{q(G)S(G)}{n}\right)} .
\]
Consequently, for finite abelian groups,
\[
\discop(L_{f_1},\ldots,L_{f_n})
\le C\sqrt{n\log\left(e+\frac{|G|}{n}\right)} .
\]
For any family of groups with $q(G)\le q$ fixed, arbitrary convolution kernels
satisfy $O_q(\sqrt n)$ whenever $|G|\le n$.
\end{theorem}

\begin{proof}
Apply Theorem \ref{thm:cstar} to the group convolution algebra. In the
abelian case all irreducible representations have dimension one and
$S(G)=|G|$. In general $S(G)\le |G|$, and therefore $q(G)S(G)\le q|G|$ on
bounded-irrep families.
\end{proof}

\begin{corollary}[Group convolution with no irrep bound]\label{cor:group-akbassra}
For every \(C_0\ge1\) there is \(C(C_0)\) such that the following holds.  Let
\(G\) be a finite group with \(|G|\le C_0n\).  If
\(L_{f_1},\ldots,L_{f_n}\) are arbitrary left convolution contractions on
\(\ell^2(G)\), then
\[
\discop(L_{f_1},\ldots,L_{f_n})\le C(C_0)\sqrt n .
\]
\end{corollary}

\begin{proof}
The left convolution algebra is finite-dimensional and, by the Fourier
decomposition above, is isomorphic to
\(\bigoplus_{\pi\in\widehat G}\M_{d_\pi}(\C)\).  Its intrinsic dimension is
\[
\sum_{\pi\in\widehat G}d_\pi^2=|G|.
\]
Apply Theorem~\ref{thm:akbassra}.
\end{proof}

\begin{corollary}[Group convolution with low-rank noise]\label{cor:group-akbassra-pert}
For every \(C_0\ge1\) there is \(C(C_0)\) such that the following holds.  Let
\(G\) be a finite group with \(|G|\le C_0n\), and let
\[
X_i=L_{f_i}+L_i,\qquad i=1,\ldots,n,
\]
be self-adjoint operators on \(\ell^2(G)\).  Suppose that each \(L_{f_i}\) is a
self-adjoint left convolution contraction and that
\[
\|L_i\|_{\op}\le1,\qquad
\|L_i\|_{\mathrm F}^2\le \frac{n}{\log^3(en)}
\qquad(1\le i\le n).
\]
Then
\[
\discop(X_1,\ldots,X_n)\le C(C_0)\sqrt n .
\]
In particular the same conclusion holds if
\(\operatorname{rank}(L_i)\le n/\log^3(en)\) and \(\|L_i\|_{\op}\le1\).
\end{corollary}

\begin{proof}
The self-adjoint left convolution operators form the self-adjoint part of the
finite-dimensional convolution C*-algebra
\(\bigoplus_{\pi\in\widehat G}\M_{d_\pi}(\C)\), whose intrinsic dimension is
\(\sum_{\pi\in\widehat G}d_\pi^2=|G|\le C_0n\).  Apply
Theorem~\ref{thm:akbassra-pert} with \(B_i=L_{f_i}\).  The rank specialization
follows from
\[
\|L_i\|_{\mathrm F}^2
\le \operatorname{rank}(L_i)\|L_i\|_{\op}^2
\le \frac{n}{\log^3(en)} .
\]
\end{proof}

\begin{remark}
Theorem \ref{thm:group} signs arbitrary convolution kernels. This is
different from the finite-group regular-representation theorem of Bandeira and
B\"olcskei \cite{BandeiraBolcskei}, which signs the specific list of left-regular
representation matrices indexed by group elements with a universal
$C\sqrt{|G|}$ bound for every finite group. The two results are complementary:
Theorem \ref{thm:group} is the elementary finite-lift estimate and handles
arbitrary kernels with an explicit \(q(G)S(G)\) entropy parameter; the
Akbas--Sra corollary above is stronger in the square group-algebra regime
\(|G|=O(n)\), because it uses the intrinsic dimension \(|G|\) and does not
require the irreducible dimensions to be controlled.  Corollary~\ref{cor:group-akbassra-pert}
adds the perturbative version in the same square regime: arbitrary
self-adjoint low-Frobenius, in particular low-rank, noise can be placed on top
of the convolution cores.
\end{remark}

\section{Toeplitz, Hankel, and multilevel Toeplitz classes}

We now recover and extend the Toeplitz and Hankel endpoint. The full-bandwidth
results, at one level and at any fixed level, use bounded finite Toeplitz
extension; at higher level the extension constant depends on the level. The
banded result uses an elementary finite-section estimate, is dimension-free in
its constant, and is included with full details.

\subsection{Scalar Toeplitz matrices}

For a finite sequence $(a_k)_{-(d-1)\le k\le d-1}$, write
\[
T_d(a)= (a_{r-s})_{0\le r,s<d} .
\]
If $\varphi\in L^\infty(\T)$ has Fourier coefficients
$\widehat\varphi(k)=a_k$ for $|k|<d$, then $\varphi$ induces $T_d(a)$.

\begin{theorem}[Finite Toeplitz extension]\label{thm:text}
There is an absolute constant $K_T$ such that every complex $d\times d$ Toeplitz
matrix $A=T_d(a)$ has an inducer $\varphi_A\in L^\infty(\T)$ satisfying
\[
\widehat\varphi_A(k)=a_k \quad (|k|<d),
\qquad
\norm{\varphi_A}_{L^\infty}\le K_T\norm{A}_{\op} .
\]
One may take $K_T\le 3$.
\end{theorem}

This is the finite Toeplitz extension theorem of Nikolskaya and Farforovskaya
\cite{NikolskayaFarforovskaya}; Courtney and Sarason record the bound and note
that the same reasoning applies in the non-self-adjoint case \cite{CourtneySarason}.

We also use a standard de la Vall\'ee Poussin kernel $V_N$ satisfying
\[
\widehat V_N(k)=1\quad(|k|\le N),
\qquad
\widehat V_N(k)=0\quad(|k|>2N),
\qquad
\norm{V_N}_{L^1}\le C_V .
\]

\begin{theorem}[Arbitrary scalar Toeplitz signing]\label{thm:toep}
There is a universal constant $C$ such that, for every collection of arbitrary
complex Toeplitz matrices $A_1,\ldots,A_n\in\M_d(\C)$ with
$\norm{A_i}_{\op}\le 1$,
\[
\discop(A_1,\ldots,A_n)
\le C\sqrt{n\log\left(e+\frac{d}{n}\right)} .
\]
In particular, self-adjoint Toeplitz matrices satisfy Matrix Spencer in the
regime $d\le n$.
\end{theorem}

\begin{proof}
For each $i$, choose an inducer $\varphi_i$ from Theorem
\ref{thm:text}. Put
$p_i=\varphi_i*V_{d-1}$. Then $p_i$ is a trigonometric polynomial of degree at
most $2d-2$, its Fourier coefficients agree with those of $\varphi_i$ for
$|k|<d$, and
\[
\norm{p_i}_{\infty}\le C_VK_T .
\]
Hence $T_d(p_i)=A_i$.

Choose a uniform grid $Q\subset\T$ of size $O(d)$ that is $2$-norming for
trigonometric polynomials of degree at most $2d$ by Lemma \ref{lem:grid}.
Apply Theorem \ref{thm:sampled} with $b=1$, $K=C_VK_T$, and $\Lambda=1$. There
are signs for which the signed polynomial $P=\sum_i\eps_i p_i$ satisfies
\[
\norm{P}_{\infty}
\le C\sqrt{n\log\left(e+\frac{d}{n}\right)} .
\]
Toeplitz compression is contractive from $L^\infty(\T)$ to operator norm, so
\[
\norm{\sum_i\eps_i A_i}_{\op}
= \norm{T_d(P)}_{\op}
\le \norm{P}_{\infty} .
\]
\end{proof}

\begin{remark}[Circulant-lift interpretation]
If the grid size is larger than twice the degree of the $p_i$, then the vector
$(p_i(\theta))_{\theta\in Q}$ defines a circulant matrix whose upper-left
$d\times d$ block is $A_i$. Thus the preceding proof is a bounded
circulant-lift proof.
\end{remark}

\subsection{Scalar Hankel matrices}

Let $J_d$ be the reversal matrix, $(J_dx)_r=x_{d-1-r}$. If $H$ is Hankel, then
$HJ_d$ is Toeplitz and $\norm{HJ_d}=\norm H$.

\begin{corollary}[Arbitrary scalar Hankel signing]\label{cor:hankel}
For arbitrary complex Hankel matrices $H_1,\ldots,H_n\in\M_d(\C)$ with
$\norm{H_i}_{\op}\le 1$,
\[
\discop(H_1,\ldots,H_n)
\le C\sqrt{n\log\left(e+\frac{d}{n}\right)} .
\]
Thus self-adjoint Hankel matrices satisfy Matrix Spencer when $d\le n$.
\end{corollary}

\begin{proof}
Apply Theorem \ref{thm:toep} to the Toeplitz matrices $H_iJ_d$ and
observe that
\[
\norm{\sum_i\eps_iH_i}
=\norm{\left(\sum_i\eps_iH_iJ_d\right)J_d}
=\norm{\sum_i\eps_iH_iJ_d} .
\]
\end{proof}

\subsection{Banded Toeplitz matrices}

For bandwidth $w\ll d$, one needs only $O(w)$ Fourier samples. The following
finite-section lemma replaces the bounded extension theorem in this case; for
the spectral theory of banded Toeplitz matrices in general see B\"ottcher and
Grudsky \cite{BottcherGrudsky}.

\begin{lemma}[Finite-section lower bound for banded symbols]\label{lem:band}
Let $p(e^{it})=\sum_{|k|\le w}a_ke^{ikt}$ and let
$T_d(p)=(a_{r-s})_{0\le r,s<d}$. If $d>2w$, then
\[
\norm{p}_{\infty}
\le \left(1-\frac{2w}{d}\right)^{-1/2}\norm{T_d(p)}_{\op} .
\]
In particular, if $d\ge 4w$, then $\norm{p}_{\infty}\le \sqrt2\norm{T_d(p)}$.
\end{lemma}

\begin{proof}
Choose $t_0$ such that $|p(e^{-it_0})|=\norm p_\infty$, and set
$x_s=d^{-1/2}e^{is t_0}$ for $0\le s<d$. Then $\norm x_2=1$. For every
interior row $w\le r\le d-1-w$,
\[
(T_d(p)x)_r
= \frac{1}{\sqrt d}\sum_{|k|\le w} a_k e^{i(r-k)t_0}
= \frac{e^{ir t_0}}{\sqrt d}p(e^{-it_0}).
\]
Therefore
\[
\norm{T_d(p)}_{\op}^2
\ge \norm{T_d(p)x}_2^2
\ge \frac{d-2w}{d}\norm p_\infty^2 .
\]
\end{proof}

\begin{theorem}[Banded scalar Toeplitz signing]\label{thm:banded}
There is a universal constant $C$ such that the following holds. Let
$A_i=T_d(p_i)$ be arbitrary complex Toeplitz matrices whose symbols have degree
at most $w$, and assume $\norm{A_i}_{\op}\le 1$. If $d\ge 4w$ and $w\ge 1$,
then
\[
\discop(A_1,\ldots,A_n)
\le C\sqrt{n\log\left(e+\frac{w}{n}\right)} .
\]
For $w=0$ the same statement holds with $w$ replaced by $1$ inside the logarithm.
\end{theorem}

\begin{proof}
By Lemma \ref{lem:band}, $\norm{p_i}_\infty\le\sqrt2$ for all $i$.
Choose an $O(w)$-point uniform grid that is $2$-norming for degree-$w$
trigonometric polynomials by Lemma \ref{lem:grid}. Apply Theorem
\ref{thm:sampled} with $b=1$ to obtain signs with
\[
\norm{\sum_i\eps_i p_i}_\infty
\le C\sqrt{n\log\left(e+\frac{w}{n}\right)} .
\]
Contractivity of Toeplitz compression gives the same bound for
$\sum_i\eps_iA_i$.
\end{proof}

\subsection{Block and multilevel Toeplitz matrices}

An $r$-level $b\times b$ block Toeplitz matrix is indexed by a box
\[
\Lambda=\prod_{j=1}^r\{0,\ldots,d_j-1\}
\]
and has $b\times b$ blocks depending only on coordinate differences. Its total
dimension is
\[
D=b\prod_{j=1}^r d_j .
\]

At one level the bounded finite Toeplitz extension of Theorem \ref{thm:text}
supplies a symbol of comparable norm, and this drives the sampling argument. We
record one consequence of this fact for emphasis.

\begin{theorem}[Finite Toeplitz extension; one level]\label{thm:text2}
There is an absolute constant $K_T$ (one may take $K_T\le 3$) such that every
$d\times d$ scalar Toeplitz matrix $A=T_d(a)$ has an inducer
$\varphi_A\in L^\infty(\T)$ with $\widehat\varphi_A(k)=a_k$ for $|k|<d$ and
$\norm{\varphi_A}_{L^\infty}\le K_T\norm A_{\op}$.
\end{theorem}

This is Theorem \ref{thm:text}, restated for reference; we use it below at one
level. At higher level the extension still holds for each fixed $r$, by Theorem
\ref{thm:mltext} below, but the constant is no longer absolute: it depends on
$r$ and cannot be taken uniform in $r$.

\begin{remark}\label{rem:nomultilevel}
The one-level constant $K_T$ has no analogue that is uniform in the number of
levels. For each fixed $r$ a bounded symbol $\varphi\in L^\infty(\T^r)$ of
comparable norm does exist (Theorem \ref{thm:mltext}), but the best constant
$K_r$ grows with $r$: the underlying polydisc Nehari constant is at least
$(\pi^2/8)^{r/4}$ \cite{OrtegaCerdaSeip}, so there is no level-free bound and no
infinite-dimensional analogue. For the purposes of the present method this means
the symbol route is available at every fixed level, but with an $r$-dependent
loss in the constant rather than the absolute constant of the one-level case.
\end{remark}

\begin{theorem}[Fixed-block, one-level block Toeplitz/Hankel]\label{thm:blocktoep1}
Fix $b\ge 1$. There is a constant $C_b$ such that every collection of
$b\times b$ (one-level) block Toeplitz, or block Hankel, contractions
$A_1,\ldots,A_n\in\M_D(\C)$ satisfies
\[
\discop(A_1,\ldots,A_n)
\le C_b\sqrt{n\log\left(e+\frac{D}{n}\right)} .
\]
\end{theorem}

\begin{proof}
First consider block Toeplitz matrices, so $r=1$. For a fixed pair of block
coordinates $p,q\in\{1,\ldots,b\}$, the scalar Toeplitz matrix formed by the
$(p,q)$ entries of all blocks is a compression of $A_i$ between coordinate
subspaces, hence has norm at most $\norm{A_i}_{\op}$. By Theorem
\ref{thm:text}, each scalar entry has an $L^\infty(\T)$ inducer with norm at
most $K_T\norm{A_i}_{\op}$. Combining the $b^2$ scalar inducers gives a
matrix-valued inducer $\Phi_i\in L^\infty(\T;\M_b)$ with
\[
\norm{\Phi_i}_{L^\infty(\T;\M_b)}\le bK_T\norm{A_i}_{\op} .
\]
Convolve $\Phi_i$ with the de la Vall\'ee Poussin kernel $V_{D/b-1}$ to obtain a
matrix trigonometric polynomial $P_i$ of degree $O(D/b)$ with the same finite
Toeplitz coefficients as $A_i$ and $\norm{P_i}_{L^\infty(\T;\M_b)}\le
C_b\norm{A_i}_{\op}$. By Lemma \ref{lem:grid} with $r=1$ there is a $2$-norming
grid $Q\subset\T$ of size $O(D/b)$. Theorem \ref{thm:sampled} with block size
$b$ gives
\[
\discop(A_1,\ldots,A_n)
\le C_b\sqrt{n\log\left(e+\frac{b^2|Q|}{n}\right)}
\le C_b\sqrt{n\log\left(e+\frac{D}{n}\right)},
\]
since $b$ is fixed and $b^2|Q|=O_b(D)$, using contractivity of Toeplitz
compression to transfer the bound to $\sum_i\eps_iA_i$. The block Hankel
statement follows by multiplying by the reversal matrix, exactly as in
Corollary \ref{cor:hankel}.
\end{proof}

\subsection{The multilevel symbol extension theorem}

The finite Toeplitz extension of Theorem \ref{thm:text} is multilevel. For each
fixed number of levels there is a bounded $L^\infty(\T^r)$ symbol carrying the
matrix entries, with a constant that depends on $r$ but not on the box.

\begin{theorem}[Finite multilevel Toeplitz symbol extension]\label{thm:mltext}
For each fixed $r\ge 1$ there is a finite constant $K_r$ such that every finite
scalar $r$-level Toeplitz matrix $A$ on a box $\prod_j\{0,\ldots,d_j-1\}$ admits
a symbol $\varphi\in L^\infty(\T^r)$ with
\[
\wh\varphi(k)=a_k
\quad\text{on the difference indices }
k\in\prod_j\{-(d_j-1),\ldots,d_j-1\},
\qquad
\norm{\varphi}_{L^\infty(\T^r)}\le K_r\norm{A}_{\op}.
\]
\end{theorem}

This is the multivariable finite Toeplitz extension theorem of Carlsson and
Perfekt \cite[Thm.~7.1 and Cor.~5.4]{CarlssonPerfekt}, whose proof realizes the
matrix as a bounded convolution operator with an $L^\infty$ symbol. For $r=1$ it
is Theorem \ref{thm:text}. For $r\ge 2$ it rests on the polydisc Nehari theorem,
proved for $r=2$ by Ferguson and Lacey \cite{FergusonLacey} and for $r\ge 3$ by
Lacey and Terwilleger \cite{LaceyTerwilleger}.

\begin{remark}\label{rem:notuniform}
The constant $K_r$ cannot be taken uniform in $r$. The corresponding polydisc
Nehari constant is at least $(\pi^2/8)^{r/4}$ by the lower bound of
Ortega-Cerd\`a and Seip \cite{OrtegaCerdaSeip}, so it grows at least
exponentially in $r$, and there is no infinite-dimensional analogue of the
extension with a level-free constant.
\end{remark}

For higher level the symbol of Theorem \ref{thm:mltext} is bounded only by the
$r$-dependent constant $K_r$ of Remark \ref{rem:notuniform}. If the symbols are
in addition banded, an elementary finite-section estimate replaces the extension
theorem with no loss of constant and no dependence on $r$. The next lemma is the
multilevel analogue of Lemma \ref{lem:band}.

\begin{lemma}[Multilevel banded finite-section bound]\label{lem:bandml}
Let $\mathbf w=(w_1,\ldots,w_r)$, and let
\[
p(\theta_1,\ldots,\theta_r)
=\sum_{|k_j|\le w_j\ \forall j}a_k e^{i\langle k,\theta\rangle}
\]
be a trigonometric polynomial of per-coordinate half-bandwidth at most $w_j$,
and let $T_{\mathbf d}(p)$ be
its $r$-level Toeplitz matrix on $\prod_j\{0,\ldots,d_j-1\}$. If $d_j>2w_j$ for
all $j$, then
\[
\norm{p}_{L^\infty(\T^r)}
\le \left(\prod_{j=1}^r\Big(1-\frac{2w_j}{d_j}\Big)\right)^{-1/2}
\norm{T_{\mathbf d}(p)}_{\op} .
\]
In particular, if $d_j\ge 4w_j$ for all $j$, then
$\norm p_\infty\le 2^{r/2}\norm{T_{\mathbf d}(p)}_{\op}$.
\end{lemma}

\begin{proof}
Pick $(t_1^0,\ldots,t_r^0)$ with
$|p(e^{-it_1^0},\ldots,e^{-it_r^0})|=\norm p_\infty$, and take the tensor unit
vector $x_{\mathbf s}=\prod_j d_j^{-1/2}e^{is_jt_j^0}$, so $\norm x_2=1$. For
every multi-index $\mathbf m$ interior in every coordinate
($w_j\le m_j\le d_j-1-w_j$) the multilevel-Toeplitz action factors
coordinatewise, giving
\[
(T_{\mathbf d}(p)x)_{\mathbf m}
=\Big(\prod_j d_j^{-1/2}\Big)
e^{i\langle\mathbf m,\mathbf t^0\rangle}
p(e^{-it_1^0},\ldots,e^{-it_r^0}).
\]
The number of fully-interior indices is $\prod_j(d_j-2w_j)$, so
\[
\norm{T_{\mathbf d}(p)}_{\op}^2
\ge \norm{T_{\mathbf d}(p)x}_2^2
\ge \prod_{j=1}^r\frac{d_j-2w_j}{d_j}\,\norm p_\infty^2 .
\]
This is Lemma \ref{lem:band} run on a product test vector; no extension theorem
is used.
\end{proof}

\begin{theorem}[Banded multilevel block Toeplitz/Hankel]\label{thm:blocktoepml}
Fix $r,b\ge 1$ and bandwidth $w\ge 1$. There is a constant $C_{r,b}$ such that
for arbitrary complex $r$-level $b\times b$ block Toeplitz contractions
$A_1,\ldots,A_n\in\M_D(\C)$ whose matrix symbols have per-coordinate degree at
most $w$, with $d_j\ge 4w$ for all $j$,
\[
\discop(A_1,\ldots,A_n)
\le C_{r,b}\sqrt{n\log\left(e+\frac{w^r}{n}\right)} .
\]
The same holds for block Hankel matrices for which $H_iJ_{\mathbf d}$ is a
block Toeplitz matrix satisfying the same bandwidth hypothesis, where
$J_{\mathbf d}$ is the tensor-product reversal. For fixed $r,b,w$ with the box
large ($d_j\ge 4w$) this is $O_{r,b,w}(\sqrt n)$.
\end{theorem}

\begin{proof}
The matrix symbols $P_i\in L^\infty(\T^r;\M_b)$ of the $A_i$ are already
trigonometric polynomials of per-axis degree at most $w$, so no extension step
is needed. For each block-coordinate pair, the corresponding scalar finite
section is a compression of $A_i$ between coordinate subspaces, hence has
operator norm at most $\norm{A_i}_{\op}$. Applying Lemma \ref{lem:bandml}
entrywise to the $b^2$ scalar symbols gives
\[
\norm{P_i}_{L^\infty(\T^r;\M_b)}\le 2^{r/2}b\,\norm{A_i}_{\op} .
\]
By Lemma \ref{lem:grid} there is a $2$-norming grid $Q\subset\T^r$ of size
$|Q|=O_r(w^r)$. Theorem \ref{thm:sampled} with block size $b$ gives
\[
\discop(A_1,\ldots,A_n)
\le C_{r,b}\sqrt{n\log\left(e+\frac{b^2|Q|}{n}\right)}
\le C_{r,b}\sqrt{n\log\left(e+\frac{w^r}{n}\right)},
\]
and multilevel Toeplitz-compression contractivity transfers the bound to
$\sum_i\eps_iA_i$. For the Hankel case, apply the Toeplitz statement to
$H_iJ_{\mathbf d}$ and then multiply on the right by the unitary
$J_{\mathbf d}$, exactly as in Corollary \ref{cor:hankel}.
\end{proof}

\subsection{Banded multilevel block Hankel matrices}

The Hankel half of Theorem \ref{thm:blocktoepml} was phrased as a hypothesis on
$H_iJ_{\mathbf d}$. We now make the reduction unconditional by recording the
exact properties of the tensor-product reversal $J_{\mathbf d}$, so that an
arbitrary multilevel block Hankel matrix is turned into a multilevel block
Toeplitz matrix of the same operator norm, bandwidth, and box.

\begin{lemma}[Tensor reversal]\label{lem:tensorrev}
For $d\ge 1$ let $J_d$ be the reversal matrix $(J_d)_{rs}=\delta_{s,\,d-1-r}$ on
$\C^d$, and set
\[
\mathbf J_\Lambda
=J_{d_1}\otimes\cdots\otimes J_{d_r}\otimes I_b
\]
on the box $\Lambda=\prod_{j=1}^r\{0,\ldots,d_j-1\}$ with $b\times b$ blocks.
Then $\mathbf J_\Lambda$ is real, symmetric, orthogonal, and an involution. If
$H$ is an $r$-level $b\times b$ block Hankel matrix, that is, its block at
$(\mathbf m,\mathbf m')$ depends only on $\mathbf m+\mathbf m'$, then
$H\mathbf J_\Lambda$ is $r$-level $b\times b$ block Toeplitz on the same box,
\[
\norm{H\mathbf J_\Lambda}_{\op}=\norm{H}_{\op},
\]
and per-coordinate bandedness with half-bandwidth $w$ is preserved.
\end{lemma}

\begin{proof}
Each $J_{d_j}$ is a permutation matrix, hence orthogonal, and it is symmetric
with $J_{d_j}^2=I$, so it is a real symmetric orthogonal involution; the same is
true of $I_b$. A tensor product of real symmetric orthogonal involutions is
again one, so $\mathbf J_\Lambda$ has all four properties. Write the block of
$H$ at $(\mathbf m,\mathbf m')$ as $C_{\mathbf m+\mathbf m'}$. Since
$\mathbf J_\Lambda$ acts on the index $\mathbf m'$ by the coordinatewise
reflection $\rho(\mathbf m')_j=d_j-1-m'_j$ and is the identity on the block
index,
\[
(H\mathbf J_\Lambda)_{\mathbf m,\mathbf m'}
=C_{\mathbf m+\rho(\mathbf m')},
\qquad
\big(\mathbf m+\rho(\mathbf m')\big)_j=(d_j-1)+(m_j-m'_j).
\]
The right-hand side depends only on the differences $m_j-m'_j$, so
$H\mathbf J_\Lambda$ is $r$-level block Toeplitz; its symbol coefficient at lag
$\mathbf k$ is $C_{(\mathbf d-\mathbf 1)+\mathbf k}$, which is nonzero only when
$|k_j|\le w$ for every $j$ exactly when $H$ has per-coordinate half-bandwidth
$w$. Finally $\mathbf J_\Lambda$ is orthogonal, so right multiplication by it is
an isometry for the operator norm.
\end{proof}

\begin{theorem}[Banded multilevel block Hankel signing]\label{thm:blockhankelml}
Fix $r,b\ge 1$ and bandwidth $w\ge 1$. There is a constant $C_{r,b}$ such that
every collection of arbitrary complex $r$-level $b\times b$ block Hankel
contractions $H_1,\ldots,H_n\in\M_D(\C)$ with per-coordinate half-bandwidth at
most $w$, on a box with $d_j\ge 4w$ for all $j$, satisfies
\[
\discop(H_1,\ldots,H_n)
\le C_{r,b}\sqrt{n\log\left(e+\frac{w^r}{n}\right)} .
\]
\end{theorem}

\begin{proof}
Put $A_i=H_i\mathbf J_\Lambda$. By Lemma \ref{lem:tensorrev} each $A_i$ is an
$r$-level $b\times b$ block Toeplitz contraction with the same per-coordinate
half-bandwidth $w$ on the same box, so the hypotheses of Theorem
\ref{thm:blocktoepml} hold. Because $\mathbf J_\Lambda$ is one fixed orthogonal
matrix, independent of $i$,
\[
\sum_i\eps_iH_i=\Big(\sum_i\eps_iA_i\Big)\mathbf J_\Lambda,
\qquad
\norm{\sum_i\eps_iH_i}_{\op}=\norm{\sum_i\eps_iA_i}_{\op},
\]
and the signs furnished by Theorem \ref{thm:blocktoepml} give the bound.
\end{proof}

\subsection{Sample-optimality of the grid}

The grids used above have size \(O_r(w^r)\) in the isotropic case, and
\(O_r\prod_j(2w_j+1)\) in the anisotropic case. The following dimension count
shows that no norming set can be smaller, so the bandwidth factor inside the
logarithm is intrinsic to the sampling method.

\begin{proposition}[Norming-set lower bound]\label{prop:samplelb}
Let $V\subset C(\T^r;\M_b(\C))$ be the space of $\M_b$-valued trigonometric
polynomials of degree at most \(w_j\) in the \(j\)-th coordinate, so that
\[
\dim_\C V=b^2\prod_{j=1}^r(2w_j+1).
\]
If $Q\subset\T^r$ is an $L_Q$-norming set for $V$ in the sense of the norming
definition of \S3, then the evaluation map
$V\to\bigoplus_{\omega\in Q}\M_b(\C)$, $F\mapsto(F(\omega))_{\omega\in Q}$, is
injective, and therefore
\[
b^2\abs{Q}\ge\dim_\C V=b^2\prod_{j=1}^r(2w_j+1),
\qquad\text{that is}\qquad
\abs{Q}\ge\prod_{j=1}^r(2w_j+1).
\]
In particular, when \(w_1=\cdots=w_r=w\), the grid size \(O_r(w^r)\) of
Lemma \ref{lem:grid} is optimal up to the constant, and the appearance of
\(w^r\) inside the logarithm in Theorems \ref{thm:blocktoepml} and
\ref{thm:blockhankelml} is inherent to any norming-set argument for this class.
\end{proposition}

\begin{proof}
If $F\in V$ vanishes on $Q$ then $\max_{\omega\in Q}\norm{F(\omega)}_{\op}=0$,
so the norming inequality forces $\norm{F}_{L^\infty}=0$ and $F=0$; thus
evaluation is injective. An injective linear map cannot increase dimension, and
$\dim_\C\bigoplus_{\omega\in Q}\M_b(\C)=b^2\abs{Q}$, which gives the stated
inequality.
\end{proof}

\begin{corollary}[Anisotropic banded multilevel block Toeplitz/Hankel]
\label{cor:anisotropic-bandml}
With per-axis half-bandwidths $w_1,\ldots,w_r$ the same arguments, using the
anisotropic grid of Lemma \ref{lem:grid} and the count
\(\dim_\C V=b^2\prod_j(2w_j+1)\), give the sharp form
\[
\discop(A_1,\ldots,A_n)
\le C_{r,b}\sqrt{n\log\left(e+\frac{\prod_{j=1}^r(2w_j+1)}{n}\right)} ,
\]
for arbitrary complex \(r\)-level \(b\times b\) block Toeplitz contractions
whose matrix symbols have per-coordinate degree at most \(w_j\), on boxes with
\(d_j\ge4w_j\).  The same bound holds for arbitrary complex \(r\)-level
\(b\times b\) block Hankel contractions with the same anisotropic
half-bandwidths.  Moreover any norming-set proof for this anisotropic symbol
class must use at least \(\prod_j(2w_j+1)\) sample points.
\end{corollary}

\begin{proof}
For the Toeplitz case, repeat the proof of Theorem~\ref{thm:blocktoepml} with
the anisotropic grid from Lemma~\ref{lem:grid}.  Lemma~\ref{lem:bandml} gives
the same entrywise symbol norm control, while the grid has
\(|Q|=O_r\prod_j(2w_j+1)\).  Theorem~\ref{thm:sampled} then gives the displayed
bound.  The Hankel case follows by applying the Toeplitz statement to
\(H_i\mathbf J_\Lambda\) and using Lemma~\ref{lem:tensorrev}.  The lower bound
on the number of sample points is Proposition~\ref{prop:samplelb} with
\(\dim_\C V=b^2\prod_j(2w_j+1)\).
\end{proof}

\subsection{The self-adjoint regime and the entropy obstruction}

For fixed $r,b$ and total anisotropic bandwidth small relative to $n$, the
banded multilevel results give the conjectured square-root bound with no
restriction on the ambient dimension $D=b\prod_j d_j$.

\begin{corollary}[Self-adjoint anisotropic banded multilevel regime]
\label{cor:saregime}
Fix $r,b\ge1$ and half-bandwidths $w_1,\ldots,w_r$ with
\[
  \prod_{j=1}^r(2w_j+1)\le n.
\]
Let $A_1,\ldots,A_n$ be self-adjoint banded $r$-level $b\times b$ block
Toeplitz or block Hankel contractions with per-coordinate half-bandwidth at
most \(w_j\), on any box with $d_j\ge 4w_j$ for all $j$. Then
\[
\disc(A_1,\ldots,A_n)=O_{r,b}(\sqrt n),
\]
with the ambient dimension $D=b\prod_{j=1}^r d_j$ unrestricted.
\end{corollary}

\begin{proof}
The logarithm in Corollary~\ref{cor:anisotropic-bandml} is at most
\(\log(e+1)\), a constant, so the bound reads \(O_{r,b}(\sqrt n)\).  The
hypotheses constrain only the symbol and the bandwidths, never
\(\prod_jd_j\), so \(D\) may be arbitrarily large.
\end{proof}

This avoids the full-unit-ball obstruction of Proposition~\ref{prop:obstruct}:
the banded symbol family has only \(b^2\prod_j(2w_j+1)=O_{r,b}(n)\) effective
scalar degrees of freedom, whereas faithful finite-block reductions of the
whole unit ball/full algebra \(\M_D(\C)\) must pay \(D^2\) entropy.

\subsection{A matching lower bound from circulants}

The upper bounds above are sharp in the sampled regime. We show that for the
periodic model the discrepancy problem is, on the nose, complex vector Spencer,
and we read off a matching lower bound from the tightness of Spencer's theorem.

Fix \(N_1,\ldots,N_r\ge1\) and work on the product torus
\(\Lambda=\prod_{j=1}^r\Z_{N_j}\), with \(M=\prod_jN_j\) points. For a lag
function \(a:\Lambda\to\C\) the \(r\)-level circulant
\[
A=\sum_{\mathbf k\in\Lambda}a_{\mathbf k}\,S^{\mathbf k},
\qquad
S^{\mathbf k}=S_1^{k_1}\cdots S_r^{k_r},
\]
is built from the coordinate cyclic shifts \(S_j\) on \(\Z_{N_j}\). All
\(r\)-level circulants are simultaneously diagonalized by the product discrete
Fourier transform, and the eigenvalue at the character indexed by
\(\xi\in\Lambda\) is
\[
v(A)_\xi=\sum_{\mathbf k\in\Lambda}a_{\mathbf k}
\prod_{j=1}^r\omega_j^{\xi_jk_j},
\qquad \omega_j=e^{2\pi i/N_j}.
\]
Thus
\[
\norm{A}_{\op}=\max_{\xi\in\Lambda}\abs{v(A)_\xi}=\norm{v(A)}_\infty,
\]
and $a\mapsto v(A)$ is a (scaled) Fourier bijection $\C^M\to\C^M$.

\begin{theorem}[Anisotropic circulant lower bound]\label{thm:circulantlb}
Let \(w_1,\ldots,w_r\ge0\) be integers, let \(N_j=2w_j+1\), let
\(\Lambda=\prod_j\Z_{N_j}\), and let
\(M=\prod_jN_j=\prod_j(2w_j+1)\). Since every residue in \(\Z_{N_j}\) has a
representative in \([-w_j,w_j]\), these are the periodic analogues of
anisotropic half-bandwidths \(w_j\). Signing \(r\)-level circulants on
\(\Lambda\) is exactly complex vector Spencer on \(M\) coordinates: for any
circulants \(A_1,\ldots,A_n\),
\[
\norm{\sum_{i=1}^n\eps_iA_i}_{\op}
=\norm{\sum_{i=1}^n\eps_i v(A_i)}_\infty
\qquad(\eps\in\{-1,1\}^n),
\]
where $v(A_i)\in\C^M$ are the Fourier symbols. Consequently, for $n\le M$ there
exist circulant contractions $A_1,\ldots,A_n$ with
\[
\disc(A_1,\ldots,A_n)
\ge c\sqrt{n\log\left(e+\frac{\prod_{j=1}^r(2w_j+1)}{n}\right)}
\]
for an absolute constant $c>0$. The same lower bound already holds among
self-adjoint scalar circulants, and hence among self-adjoint \(b\times b\) block
circulants for every \(b\ge1\). In particular this matches the upper bound of
Corollary~\ref{cor:anisotropic-bandml} up to constants.
\end{theorem}

\begin{proof}
The displayed identity is immediate from
$\norm{A}_{\op}=\norm{v(A)}_\infty$ and the linearity of $a\mapsto v(A)$ over
the common signs. Because $a\mapsto v(A)$ is a bijection $\C^M\to\C^M$, choosing
the symbols $v(A_i)$ freely realizes any tuple of vectors in $\C^M$ as circulant
symbols, the contraction normalization $\norm{A_i}_{\op}\le1$ being exactly
$\norm{v(A_i)}_\infty\le1$. The lower bound is then the tightness of vector
Spencer: for $n\le M$ there are vectors $x_1,\ldots,x_n\in\C^M$ with
$\norm{x_i}_\infty\le1$ for which every signing has
$\norm{\sum_i\eps_ix_i}_\infty\ge c\sqrt{n\log(e+M/n)}$. One may take the
$x_i$ to be independent uniform signs on a suitable subset of coordinates; the
matching \(\Omega(\sqrt{n\log(e+M/n)})\) lower bound for the rectangular
$\ell_\infty$ problem is classical \cite[Ch.~4]{Matousek}, and is also a case of
Gluskin's lower bounds for \(\ell_\infty\)-balls \cite{Gluskin}. Substituting
\(M=\prod_j(2w_j+1)\) and transporting through the bijection \(v\) gives the
claim.

For the self-adjoint version, choose the lower-bound vectors \(x_i\) real. The
Fourier transform maps self-adjoint scalar circulants bijectively onto real
functions on the dual product torus: if \(v(A)\) is real, then the inverse
Fourier coefficients satisfy \(a_{-\mathbf k}=\overline{a_{\mathbf k}}\), and
the circulant is self-adjoint. Thus the same real \(\ell_\infty^M\) lower bound
is realized by self-adjoint scalar circulants. Finally, for block size \(b\), fix
a rank-one projection \(P\in\M_b(\C)\) and replace \(A_i\) by \(A_i\otimes P\).
This is a self-adjoint \(b\times b\) block circulant contraction and
\[
\norm{\sum_i\eps_i(A_i\otimes P)}_{\op}
=\norm{\sum_i\eps_iA_i}_{\op},
\]
so the scalar lower bound is inherited unchanged.
\end{proof}

\begin{corollary}[The clean one-level case]\label{cor:circ1}
For $r=1$, $N=2w+1$, and $n\le N$, there are scalar circulant contractions
$A_1,\ldots,A_n\in\M_N(\C)$ with
\[
\disc(A_1,\ldots,A_n)\ge c\sqrt{n\log\left(e+\frac{2w+1}{n}\right)} .
\]
Since scalar circulants are the periodic analogue of banded scalar Toeplitz
matrices, this gives the exact lower-bound companion for the sampled periodic
model behind Theorem \ref{thm:banded}.
\end{corollary}

\begin{proof}
This is Theorem \ref{thm:circulantlb} at $r=1$, where $M=N=2w+1$ and a circulant
is a Toeplitz matrix with cyclically wrapped, hence $w$-banded, diagonals.
\end{proof}

\begin{question}[Non-periodic matching bound]\label{q:nonperiodic}
The lower bound of Theorem \ref{thm:circulantlb} is for the periodic
(circulant) model, where the symbol-to-operator passage is an exact isometry.
Does the same
\[
\Omega\!\left(\sqrt{n\log\left(e+\frac{\prod_j(2w_j+1)}{n}\right)}\right)
\]
lower bound hold for genuine finite-section banded Toeplitz matrices
\(T_{\mathbf d}(p)\) on a box with \(d_j\) comparable to \(w_j\)? The
finite-section symbol map is only a near-isometry: Lemma \ref{lem:bandml}
controls it from below by a constant, but transporting a worst-case symbol to
the matrix incurs the Lebesgue-type loss of the grid, which for the natural
isotropic construction degrades the bound by a factor \((\log w)^{-r}\). Closing
this residual gap, so that the finite-section banded class matches the
anisotropic circulant lower bound exactly, is open.
\end{question}

\subsection{The full-bandwidth multilevel case}

For fixed level the full-bandwidth class is also covered, at the Spencer rate,
once the multilevel symbol extension of Theorem \ref{thm:mltext} is fed into the
sampling argument. No bandwidth restriction is needed; only $r$ and $b$ are
fixed.

\begin{theorem}[Fixed-level multilevel block Toeplitz/Hankel]\label{thm:blocktoepfull}
Fix $r,b\ge 1$. There is a constant $C_{r,b}$ such that every collection of
$r$-level $b\times b$ block Toeplitz, or block Hankel, contractions
$A_1,\ldots,A_n\in\M_D(\C)$ of full bandwidth and ambient dimension
$D=b\prod_j d_j$ satisfies
\[
\discop(A_1,\ldots,A_n)
\le C_{r,b}\sqrt{n\log\left(e+\frac{D}{n}\right)} .
\]
In particular, for $D\le n$ this is $O_{r,b}(\sqrt n)$, with the box
$\prod_j d_j$ otherwise unrestricted.
\end{theorem}

\begin{proof}
Treat the block Toeplitz case first. For each ordered pair of block coordinates
$p,q\in\{1,\ldots,b\}$, the scalar $r$-level Toeplitz matrix formed by the
$(p,q)$ entries of the blocks of $A_i$ is a compression of $A_i$ between
coordinate subspaces, hence has operator norm at most $\norm{A_i}_{\op}$. By
Theorem \ref{thm:mltext} each of these $b^2$ scalar entries has an
$L^\infty(\T^r)$ symbol agreeing with the matrix on the difference indices and
of norm at most $K_r\norm{A_i}_{\op}$. Assembling the $b^2$ scalar symbols gives
a matrix-valued inducer $\Phi_i\in L^\infty(\T^r;\M_b)$ with
\[
\norm{\Phi_i}_{L^\infty(\T^r;\M_b)}\le bK_r\norm{A_i}_{\op}.
\]
Convolve $\Phi_i$ with the product de la Vall\'ee Poussin kernel
$\prod_j V_{d_j-1}(\theta_j)$ to obtain a matrix trigonometric polynomial $P_i$
of per-axis degree $O(d_j)$ that retains the finite Toeplitz coefficients of
$A_i$ and satisfies $\norm{P_i}_{L^\infty(\T^r;\M_b)}\le C_{r,b}\norm{A_i}_{\op}$.
By Lemma \ref{lem:grid} there is a $2$-norming grid $Q\subset\T^r$ of size
$|Q|=O_r(\prod_j d_j)=O(D/b)$. Theorem \ref{thm:sampled} with block size $b$
then has entropy $b^2|Q|=O_{r,b}(D)$ and gives
\[
\discop(A_1,\ldots,A_n)
\le C_{r,b}\sqrt{n\log\left(e+\frac{b^2|Q|}{n}\right)}
\le C_{r,b}\sqrt{n\log\left(e+\frac{D}{n}\right)},
\]
where multilevel Toeplitz-compression contractivity transfers the bound to
$\sum_i\eps_iA_i$. The block Hankel case follows by applying the Toeplitz
statement to $H_i\mathbf J_\Lambda$ and multiplying on the right by the unitary
tensor reversal $\mathbf J_\Lambda$ of Lemma \ref{lem:tensorrev}, exactly as in
Corollary \ref{cor:hankel}.
\end{proof}

\begin{remark}[An alternative, weaker operator route]\label{rem:fullbandwidth}
There is a second, purely operator-theoretic route to a bound that does not pass
through a symbol. Carlsson and Perfekt \cite{CarlssonPerfekt} also extend a
finite-section multilevel Toeplitz matrix to a bounded $\ell^2$ multilevel
Toeplitz operator with a level-only constant; taking that operator as a single
block of size $D$ and applying Theorem \ref{thm:main} gives the unconditional
but weaker rate $\sqrt{n\log(e+D^2/n)}$, with the generic $D^2$ entropy of
Proposition \ref{prop:obstruct}. The symbol route of Theorem
\ref{thm:blocktoepfull} feeds entropy $O(D)$ rather than $D^2$, so it gives the
better dimension dependence and upgrades the fixed-level full-bandwidth case
from this trivial operator rate to a theorem at the Spencer rate.
\end{remark}

\begin{remark}[Comparison with the banded results]\label{rem:bandedcompare}
Theorem \ref{thm:blocktoepfull} is non-constructive and its constant $K_r$
depends on $r$ through the polydisc Nehari theorem (Remark
\ref{rem:notuniform}). The banded results of Theorem \ref{thm:blocktoepml}
onward are complementary: there the symbol is the matrix symbol itself, so the
finite
section bound of Lemma \ref{lem:bandml} gives a dimension-free, explicit
constant with no appeal to Nehari's theorem, and Theorem \ref{thm:circulantlb}
supplies a matching lower bound. The full-bandwidth theorem trades these
explicit, matched constants for coverage of the entire fixed-level class.
\end{remark}

\section{Closure operations}

The finite-lift framework is stable under the operations normally used to build
structured matrices. The entropy accounting in this section is explicit because
this is a common source of mistakes.

\begin{proposition}[Compressions and bounded images]\label{prop:comp}
Suppose $B_i$ satisfy a discrepancy bound
\[
\discop(B_1,\ldots,B_n)\le \Delta .
\]
If $A_i=V^*B_iV$ for a fixed isometry $V$, then
$\discop(A_1,\ldots,A_n)\le\Delta$. More generally, if
$A_i=\Phi(B_i)$ for a fixed linear map $\Phi$ with $\norm\Phi\le L$ on the
relevant span, then
\[
\discop(A_1,\ldots,A_n)
\le L\Delta .
\]
\end{proposition}

\begin{proof}
Use the same signs and the inequality
$\norm{\sum_i\eps_i\Phi(B_i)}\le L\norm{\sum_i\eps_iB_i}$.
\end{proof}

\begin{proposition}[Finite sums and direct-sum entropy]\label{prop:sums}
Suppose
\[
A_i=A_i^{(1)}+\cdots+A_i^{(s)} .
\]
For each $\ell$, assume the tuple $(A_i^{(\ell)})_{i=1}^n$ has a
$(K_\ell,\Lambda_\ell)$-lift to
\[
\mcB_\ell=\bigoplus_\alpha \M_{r_{\ell,\alpha}}(\C).
\]
Put
\[
h=\max_{\ell,\alpha} r_{\ell,\alpha},
\qquad
R=\sum_{\ell=1}^s\sum_\alpha r_{\ell,\alpha} .
\]
Then $(A_i)_{i=1}^n$ has a lift to $\bigoplus_{\ell=1}^s\mcB_\ell$ and
\[
\discop(A_1,\ldots,A_n)
\le C\left(\sum_{\ell=1}^s\Lambda_\ell K_\ell\right)
\sqrt{n\log\left(e+\frac{hR}{n}\right)} .
\]
\end{proposition}

\begin{proof}
For each component choose lifts $Y_i^{(\ell)}$ in $\mcB_\ell$. If $K_\ell>0$,
set $Z_i^{(\ell)}=Y_i^{(\ell)}/K_\ell$; zero components may be omitted. Then
$Z_i=(Z_i^{(1)},\ldots,Z_i^{(s)})$ has norm at most one in
$\bigoplus_\ell\mcB_\ell$. Define
\[
\Theta(Z^{(1)},\ldots,Z^{(s)})
=\sum_{\ell=1}^s K_\ell\Pi_\ell(Z^{(\ell)}).
\]
Then $\Theta(Z_i)=A_i$ and
$\norm{\Theta}\le\sum_\ell\Lambda_\ell K_\ell$ on the relevant span. Apply
Theorem \ref{thm:main} to the direct-sum algebra, whose maximum block
size is $h$ and whose reduced dimension is $R$.
\end{proof}

\begin{remark}[Direct sums are not additive in entropy]\label{rem:nonadd}
For $\mcB=\bigoplus_{\ell=1}^s\mcB_\ell$ one has
\[
\Ent(\mcB)=\left(\max_\ell h(\mcB_\ell)\right)
\sum_{\ell=1}^s R(\mcB_\ell),
\]
not generally $\sum_\ell\Ent(\mcB_\ell)$. Consequently, a finite sum of classes
that are individually in the Spencer regime remains in the Spencer regime only
when this combined entropy is $O(n)$. This condition is automatic, for example,
for a fixed number of scalar sampled multiplier classes of total sample size
$O(n)$, and for a fixed number of uniformly bounded-block classes with total
reduced dimension $O(n)$.
\end{remark}

\begin{corollary}[Toeplitz plus Hankel with bounded decomposition]
If $A_i=T_i+H_i$, where $T_i$ is Toeplitz, $H_i$ is Hankel, and
\[
\norm{T_i}_{\op}+\norm{H_i}_{\op}\le K
\qquad(1\le i\le n),
\]
then
\[
\discop(A_1,\ldots,A_n)
\le CK\sqrt{n\log\left(e+\frac{d}{n}\right)} .
\]
\end{corollary}

\begin{proof}
Use the scalar Toeplitz and Hankel multiplier lifts from Theorem
\ref{thm:toep} and Corollary \ref{cor:hankel}. Their sampled lift
algebras are diagonal with $O(d)$ total samples. The direct sum of the two
sampled algebras is still diagonal and has $O(d)$ total samples. Proposition
\ref{prop:sums} gives the result.
\end{proof}

\begin{proposition}[Completely bounded tensor amplification]\label{prop:cb}
Let $S\subset\M_D(\C)$ be an operator space and suppose there are linear maps
\[
\sigma:S\to\mcB=\bigoplus_\alpha\M_{r_\alpha}(\C),
\qquad
\Pi:\sigma(S)\to S,
\]
with $\Pi\sigma=\operatorname{id}_S$, $\norm{\sigma}_{cb}\le K$, and
$\norm{\Pi}_{cb}\le\Lambda$. Then, for each $c\ge1$, the amplified space
$S\otimes\M_c$ has a $(K,\Lambda)$ lift to
\[
\mcB\otimes\M_c\simeq\bigoplus_\alpha\M_{cr_\alpha}(\C).
\]
The entropy is multiplied by $c^2$. Hence fixed-size completely bounded tensor
amplification preserves the Spencer-regime conclusion.
\end{proposition}

\begin{proof}
Use the maps $\sigma\otimes\operatorname{id}_{\M_c}$ and
$\Pi\otimes\operatorname{id}_{\M_c}$. The cb-norm assumptions give the same
lift constants at matrix level $c$. The new block sizes are $cr_\alpha$, so both
$h$ and $R$ are multiplied by $c$.
\end{proof}

\begin{remark}
The completely bounded hypothesis in Proposition \ref{prop:cb} is necessary
for this general formulation. For concrete lifts coming from $*$-homomorphisms,
compressions, or C*-quotient maps it is automatically satisfied with the
expected constants.
\end{remark}

\section{Random structured inputs}

The preceding results are deterministic. They therefore apply pathwise to
random samples from the covered structured classes after normalization. For
random banded Toeplitz matrices a complementary structural obstruction, a parity
obstruction to stability, is established in the companion paper
\cite{Shmalo-paper5}.

\begin{corollary}[Pathwise random structured matrices]
Let $A_1,\ldots,A_n$ be random matrices which, almost surely, lie in one of the
classes covered above: scalar Toeplitz, scalar Hankel, one-level block
Toeplitz/Hankel or its banded multilevel form, a fixed finite-dimensional
C*-algebra of entropy $\Ent$, or a finite group convolution algebra. On each realization define
\[
\widetilde A_i=
\begin{cases}
A_i/\norm{A_i}_{\op},& A_i\ne0,\\
0,& A_i=0.
\end{cases}
\]
Then the corresponding deterministic bound holds for
$\widetilde A_1,\ldots,\widetilde A_n$ on that realization.
\end{corollary}

\begin{corollary}[Distribution-free random algebraic-core ensembles]
\label{cor:random-algebraic-lowfrobenius}
For every \(C_0\ge1\) there is \(C(C_0)\) with the following property.  Let
\(X_1,\ldots,X_n\) be random self-adjoint matrices in a common ambient
\(\M_D(\C)\).  Suppose that, outside an exceptional event \(E\), there is a
finite-dimensional C*-subalgebra \(\mcA_\omega\subset\M_D(\C)\) and a
decomposition
\[
  X_i(\omega)=B_i(\omega)+L_i(\omega),\qquad i=1,\ldots,n,
\]
such that
\[
  \dim_\C\mcA_\omega\le C_0n,\qquad
  B_i(\omega)\in\mcA_\omega,\qquad
  \|B_i(\omega)\|_{\op}\le1,
\]
and
\[
  L_i(\omega)=L_i(\omega)^*,\qquad
  \|L_i(\omega)\|_{\op}\le1,\qquad
  \|L_i(\omega)\|_{\mathrm F}^2\le \frac{n}{\log^3(en)}
\]
for every \(i\).  Then
\[
  \mathbb P\left\{\disc(X_1,\ldots,X_n)>C(C_0)\sqrt n\right\}
  \le \mathbb P(E).
\]
No independence, identical distribution, or deterministic choice of
\(\mcA_\omega\) is required.

Consequently the Matrix Spencer bound holds almost surely, or with the same
high probability as the structural event above, for all of the following random
ensembles:
\begin{enumerate}[label=\textup{(\roman*)}]
\item arbitrary random self-adjoint contractions in a full matrix algebra
\(\M_D(\C)\) whenever \(D^2\le C_0n\); this includes normalized
\(\mathrm{GOE}\), \(\mathrm{GUE}\), Wigner, sample-covariance, and Wishart-type
models in the \(D\lesssim\sqrt n\) square-block range;
\item arbitrary random self-adjoint contractions in a random common block
algebra \(U_\omega(\bigoplus_\alpha\M_{r_\alpha}(\C))U_\omega^*\) whenever
\(\sum_\alpha r_\alpha^2\le C_0n\);
\item any of the preceding algebraic-core ensembles after adding arbitrary
dependent random self-adjoint tails of rank at most \(n/\log^3(en)\) and
operator norm at most \(1\);
\item in particular, random spiked tails \(U_iD_iU_i^*\), random compressions
\(P_iY_iP_i\), and random coordinate-supported tails
\(P_{S_i}Y_iP_{S_i}\), provided the projections or supports have rank/cardinality
at most \(n/\log^3(en)\) and the inserted self-adjoint operators have operator
norm at most \(1\).
\end{enumerate}
\end{corollary}

\begin{proof}
On each realization outside \(E\), apply the Akbas--Sra low-Frobenius
perturbation theorem, Theorem~\ref{thm:akbassra-pert}, to the finite-dimensional
C*-algebra \(\mcA_\omega\) and to the decomposition \(X_i=B_i+L_i\).  This gives
the displayed discrepancy bound on that realization.  The probability bound is
therefore immediate.

The listed examples are just specializations of the structural event.  In
\(\M_D(\C)\) the intrinsic algebra dimension is \(D^2\).  A common block algebra
has intrinsic dimension \(\sum_\alpha r_\alpha^2\), unitary conjugation does not
change it, and no distributional assumption enters the deterministic theorem.
Finally,
\[
  \|L_i\|_{\mathrm F}^2
  \le \operatorname{rank}(L_i)\|L_i\|_{\op}^2
  \le \frac{n}{\log^3(en)}
\]
for every rank-controlled tail with \(\|L_i\|_{\op}\le1\).  Spiked,
compressed, and coordinate-supported tails have rank bounded by the displayed
projection rank or support size.
\end{proof}

\begin{corollary}[Distribution-free all-plus gauge for arbitrary random matrices]
\label{cor:distribution-free-all-plus-gauge-random-matrices}
Let \(X_1,\ldots,X_n\) be arbitrary random self-adjoint matrices in a common
\(\M_D(\C)\).  No independence, centering, identical distribution, tail, or
moment assumption is imposed.  Define the pathwise gauge
\[
  G(X):=\max\left\{
  \max_{i\le n}\|X_i\|_{\op},
  \frac1{\sqrt n}\left\|\sum_{i=1}^nX_i\right\|_{\op}
  \right\}.
\]
Set \(A_i=0\) on \(\{G(X)=0\}\), and \(A_i=X_i/G(X)\) otherwise.  Then,
pathwise,
\[
  \max_{i\le n}\|A_i\|_{\op}\le1,
  \qquad
  \disc(A_1,\ldots,A_n)\le\sqrt n .
\]
The same assertion holds for arbitrary rectangular matrices
\(X_i\in\M_{d_1\times d_2}(\C)\), with \(\discop\) in place of \(\disc\).

More generally, for an active triangular menu
\((X_{m,h,i})_{i\le n_m}\), \(1\le h\le J_m\), of arbitrary self-adjoint or
rectangular random matrices, let \(G_{m,h}\) be the corresponding all-plus
gauge.  If \(s_{m,h}>0\) are deterministic or environment-measurable scales and
\[
  \sum_m\mathbb P\left\{\max_{h\le J_m}\frac{G_{m,h}}{s_{m,h}}>1\right\}<\infty,
\]
then almost surely, for all sufficiently large \(m\), every active scaled tuple
\[
  X_{m,h,1}/s_{m,h},\ldots,X_{m,h,n_m}/s_{m,h}
\]
is contraction-normalized and has the all-plus certificate
\[
  \disc\left(X_{m,h,1}/s_{m,h},\ldots,X_{m,h,n_m}/s_{m,h}\right)
  \le\sqrt{n_m}
\]
in the self-adjoint case, and the analogous \(\discop\) bound in the
rectangular case.
\end{corollary}

\begin{proof}
The pathwise statement is immediate from the definition of \(G(X)\).  If
\(G(X)=0\), then all normalized matrices are zero.  Otherwise,
\[
  \|A_i\|_{\op}\le1
  \qquad\text{and}\qquad
  \left\|\sum_{i=1}^nA_i\right\|_{\op}
  \le\sqrt n .
\]
The all-plus sign vector is an admissible signing, so it gives the displayed
discrepancy bound.  The rectangular assertion is the same argument with
\(\discop\), since the all-plus sum is still a valid signed sum.  For the
active triangular statement, Borel--Cantelli gives
\(G_{m,h}\le s_{m,h}\) for every active \(h\le J_m\) and all large \(m\), and
the preceding pathwise argument applies to each active scaled tuple.
\end{proof}

\begin{corollary}[Moment schedules for the distribution-free gauge]
\label{cor:distribution-free-gauge-moment-schedules}
In the active triangular menu of
Corollary~\ref{cor:distribution-free-all-plus-gauge-random-matrices}, let
\(G_{m,h}\) be the all-plus gauge of the \(h\)-th active tuple.  Fix
\(p>0\), and let \(s_{m,h}>0\) be deterministic scales.  If
\[
  \sum_m\sum_{h=1}^{J_m}
  \mathbb E\left(\frac{G_{m,h}}{s_{m,h}}\right)^p<\infty,
\]
then almost surely, for all sufficiently large \(m\), every active scaled tuple
\[
  X_{m,h,1}/s_{m,h},\ldots,X_{m,h,n_m}/s_{m,h}
\]
is contraction-normalized and has the all-plus Matrix-Spencer certificate
\[
  \disc\left(X_{m,h,1}/s_{m,h},\ldots,X_{m,h,n_m}/s_{m,h}\right)
  \le\sqrt{n_m}
\]
in the self-adjoint case, with the analogous \(\discop\) bound in the
rectangular case.

Equivalently, if \(\eta_{m,h}>0\) satisfy
\(\sum_m\sum_{h\le J_m}\eta_{m,h}<\infty\), then it is enough to choose
\[
  s_{m,h}\ge \eta_{m,h}^{-1/p}
  \left(\mathbb E G_{m,h}^p\right)^{1/p}.
\]
In particular, a common row scale \(s_m\) works whenever
\[
  s_m\ge \eta_m^{-1/p}J_m^{1/p}
  \max_{h\le J_m}\left(\mathbb E G_{m,h}^p\right)^{1/p}
  \qquad\text{and}\qquad
  \sum_m\eta_m<\infty .
\]
No independence between the matrices, between menu elements, or across
dimensions is required.
\end{corollary}

\begin{proof}
Markov's inequality gives
\[
  \mathbb P\left\{\max_{h\le J_m}\frac{G_{m,h}}{s_{m,h}}>1\right\}
  \le
  \sum_{h=1}^{J_m}
  \mathbb E\left(\frac{G_{m,h}}{s_{m,h}}\right)^p .
\]
The first displayed summability therefore implies the tail summability
hypothesis of
Corollary~\ref{cor:distribution-free-all-plus-gauge-random-matrices}, and that
corollary gives the almost-sure conclusion.  The quantile-budget and common
row-scale forms are the same estimate with
\[
  \mathbb E(G_{m,h}/s_{m,h})^p\le\eta_{m,h}
  \quad\text{or}\quad
  \sum_{h\le J_m}\mathbb E(G_{m,h}/s_m)^p\le\eta_m .
\]
\end{proof}

\begin{corollary}[Closed-form distribution-free gauge schedules]
\label{cor:distribution-free-gauge-closed-schedules}
In the active triangular menu of
Corollary~\ref{cor:distribution-free-all-plus-gauge-random-matrices}, let
\(G_{m,h}\) be the all-plus gauge of the \(h\)-th active tuple.  Fix
\(p>0\), and suppose that deterministic envelopes \(r_{m,h}>0\) satisfy
\[
  \mathbb E G_{m,h}^p\le r_{m,h}^p
  \qquad(m\ge1,\ 1\le h\le J_m).
\]
Then the following deterministic normalizations are valid for arbitrary
self-adjoint or rectangular random matrices, with no independence assumptions.

\begin{enumerate}[label=\textup{(\roman*)}]
\item For arbitrary growth of \(J_m\), if \(\gamma>0\) and
\[
  s_{m,h}\ge r_{m,h}\bigl(J_m(e+m)^{1+\gamma}\bigr)^{1/p},
\]
then almost surely all active scaled tuples are eventually
contraction-normalized and have the all-plus
\(\sqrt{n_m}\) discrepancy certificate.

\item If \(D_m\ge2\), \(J_m\le(eD_m)^\Gamma\), and
\[
  \max_{h\le J_m}r_{m,h}\le r_m,
\]
then the common row scale
\[
  s_m=r_m(eD_m)^B
\]
is valid along all integer dimensions \(D_m=m\) whenever \(pB>\Gamma+1\).
Along every sparse stretched-exponential dimension schedule
\[
  D_m\ge \exp(cm^\kappa)
  \qquad(c,\kappa>0),
\]
the same scale is valid whenever \(pB>\Gamma\).

\item More generally, if \(L_m:=\sum_{h=1}^{J_m}r_{m,h}^p\), then any common
row scale \(s_m\) satisfying
\[
  \sum_m\frac{L_m}{s_m^p}<\infty
\]
is valid.
\end{enumerate}
\end{corollary}

\begin{proof}
Corollary~\ref{cor:distribution-free-gauge-moment-schedules} shows that it is
enough to verify
\[
  \sum_m\sum_{h=1}^{J_m}\mathbb E(G_{m,h}/s_{m,h})^p<\infty .
\]
For (i), the \(m\)-th row contribution is at most
\[
  \sum_{h=1}^{J_m}
  \frac{r_{m,h}^p}
       {r_{m,h}^pJ_m(e+m)^{1+\gamma}}
  =(e+m)^{-1-\gamma},
\]
which is summable.  For (ii), the row contribution is bounded by
\[
  J_m(eD_m)^{-pB}\le(eD_m)^{\Gamma-pB}.
\]
This is summable over all integer dimensions when \(pB>\Gamma+1\), and is
stretched-exponentially summable under \(D_m\ge\exp(cm^\kappa)\) when
\(pB>\Gamma\).  The last assertion is exactly the same moment-summability
criterion with \(s_{m,h}=s_m\).
\end{proof}

\begin{corollary}[Tail-profile distribution-free gauge schedules]
\label{cor:distribution-free-gauge-tail-schedules}
In the active triangular menu of
Corollary~\ref{cor:distribution-free-all-plus-gauge-random-matrices}, let
\(G_{m,h}\) be the all-plus gauge of the \(h\)-th active tuple.  Suppose that
there are deterministic scales \(r_{m,h}>0\) and tail controls
\(\Phi_{m,h}:[1,\infty)\to[0,1]\) such that
\[
  \mathbb P\{G_{m,h}>r_{m,h}u\}\le \Phi_{m,h}(u)
  \qquad(u\ge1).
\]
If \(u_{m,h}\ge1\) satisfy
\[
  \sum_m\sum_{h=1}^{J_m}\Phi_{m,h}(u_{m,h})<\infty,
\]
then the deterministic scales \(s_{m,h}=r_{m,h}u_{m,h}\) are valid: almost
surely, for all sufficiently large \(m\), every active scaled self-adjoint or
rectangular tuple is contraction-normalized and has the all-plus
\(\sqrt{n_m}\) discrepancy certificate.

In particular, assume the uniform sub-Weibull gauge tail
\[
  \mathbb P\{G_{m,h}>r_{m,h}u\}
  \le C_0\exp(-c_0u^q)
  \qquad(u\ge1)
\]
with \(C_0,c_0,q>0\).  Then, for every \(\gamma>0\), it is enough to take
\[
  u_m=\max\left\{1,
  \left(\frac{\max\{0,\log(C_0J_m)+(1+\gamma)\log(e+m)\}}{c_0}\right)^{1/q}
  \right\}
\]
and \(s_{m,h}=r_{m,h}u_m\).  If \(D_m\ge2\),
\(J_m\le(eD_m)^\Gamma\), and \(\max_{h\le J_m}r_{m,h}\le r_m\), then the common
scale
\[
  s_m=r_m A(\log(eD_m))^{1/q}
\]
is valid along all integer dimensions \(D_m=m\) whenever \(c_0A^q>\Gamma+1\).
Along every sparse stretched-exponential dimension schedule
\[
  D_m\ge\exp(cm^\kappa)
  \qquad(c,\kappa>0),
\]
the same common scale is valid whenever \(c_0A^q>\Gamma\).
\end{corollary}

\begin{proof}
The first assertion is exactly the active-menu hypothesis in
Corollary~\ref{cor:distribution-free-all-plus-gauge-random-matrices}, since
\[
  \mathbb P\left\{\max_{h\le J_m}\frac{G_{m,h}}{r_{m,h}u_{m,h}}>1\right\}
  \le \sum_{h=1}^{J_m}\Phi_{m,h}(u_{m,h}).
\]
For the sub-Weibull choice of \(u_m\), the \(m\)-th row contribution is at most
\[
  J_mC_0\exp(-c_0u_m^q)\le(e+m)^{-1-\gamma},
\]
which is summable.  For the polynomial-menu common scale, the row contribution
is bounded by
\[
  C_0J_m(eD_m)^{-c_0A^q}
  \le C_0(eD_m)^{\Gamma-c_0A^q}.
\]
This is summable over all integer dimensions if \(c_0A^q>\Gamma+1\), and is
stretched-exponentially summable under \(D_m\ge\exp(cm^\kappa)\) if
\(c_0A^q>\Gamma\).
\end{proof}

\begin{corollary}[Finite-confidence distribution-free gauge certificates]
\label{cor:distribution-free-gauge-finite-confidence}
Let \((X_{h,i})_{i\le n_h}\), \(1\le h\le J\), be a finite menu of arbitrary
self-adjoint or rectangular random matrix tuples, and let \(G_h\) be the
all-plus gauge of the \(h\)-th tuple.  Fix \(0<\delta<1\).

\begin{enumerate}[label=\textup{(\roman*)}]
\item If \(p>0\) and deterministic envelopes \(r_h>0\) satisfy
\[
  \mathbb E G_h^p\le r_h^p,
\]
then the deterministic scales
\[
  s_h=r_h\left(\frac{J}{\delta}\right)^{1/p}
\]
make all \(J\) scaled tuples contraction-normalized and all-plus certified with
probability at least \(1-\delta\).

\item More generally, if
\[
  \mathbb P\{G_h>r_hu\}\le \Phi_h(u)\qquad(u\ge1)
\]
and \(u_h\ge1\) are chosen so that
\[
  \sum_{h=1}^J\Phi_h(u_h)\le\delta,
\]
then the scales \(s_h=r_hu_h\) have the same \(1-\delta\) conclusion.

\item In particular, under the uniform sub-Weibull gauge tail
\[
  \mathbb P\{G_h>r_hu\}\le C_0\exp(-c_0u^q)\qquad(u\ge1),
\]
the common multiplier
\[
  u=\max\left\{1,
  \left(\frac{\max\{0,\log(C_0J/\delta)\}}{c_0}\right)^{1/q}
  \right\}
\]
and scales \(s_h=r_hu\) have the same \(1-\delta\) conclusion.
\end{enumerate}

The conclusion means, simultaneously for all \(h\le J\), that the scaled
self-adjoint tuples obey
\[
  \disc(X_{h,1}/s_h,\ldots,X_{h,n_h}/s_h)\le \sqrt{n_h},
\]
with the analogous \(\discop\) bound in the rectangular case.  No independence,
centering, or identical-distribution assumption is imposed.
\end{corollary}

\begin{proof}
For the moment form, Markov's inequality and the union bound give
\[
  \mathbb P\left\{\max_{h\le J}\frac{G_h}{s_h}>1\right\}
  \le \sum_{h=1}^J \left(\frac{r_h}{s_h}\right)^p
  \le\delta .
\]
The tail-profile form is the same union bound:
\[
  \mathbb P\left\{\max_{h\le J}\frac{G_h}{r_hu_h}>1\right\}
  \le \sum_{h=1}^J\Phi_h(u_h)\le\delta .
\]
The sub-Weibull scale makes this last sum at most
\(JC_0\exp(-c_0u^q)\le\delta\).  On the complementary event, each scaled tuple
has gauge at most one, and
Corollary~\ref{cor:distribution-free-all-plus-gauge-random-matrices} gives the
pathwise all-plus discrepancy certificate for every menu element.
\end{proof}

\begin{corollary}[Conditional gauge certificates]
\label{cor:distribution-free-gauge-conditional-confidence}
Let \((X_{h,i})_{i\le n_h}\), \(1\le h\le J\), be a finite menu of arbitrary
self-adjoint or rectangular random matrix tuples, let \(G_h\) be the all-plus
gauge of the \(h\)-th tuple, and let \(\mcF\) be a sub-\(\sigma\)-algebra.
Fix \(0<\delta<1\).

\begin{enumerate}[label=\textup{(\roman*)}]
\item If \(p>0\) and positive \(\mcF\)-measurable random variables \(R_h\)
satisfy
\[
  \mathbb E(G_h^p\mid\mcF)\le R_h^p
  \qquad\text{a.s.},
\]
then the \(\mcF\)-measurable scales
\[
  S_h:=R_h\left(\frac{J}{\delta}\right)^{1/p}
\]
make all \(J\) scaled tuples contraction-normalized and all-plus certified with
probability at least \(1-\delta\).

\item More generally, suppose positive \(\mcF\)-measurable \(R_h\),
\(\mcF\)-measurable \(U_h\ge1\), and \(\mcF\)-measurable error budgets
\(\eta_h\ge0\) satisfy the chosen-threshold conditional tail bounds
\[
  \mathbb P\{G_h>R_hU_h\mid\mcF\}\le \eta_h
\]
and
\[
  \sum_{h=1}^J\eta_h\le\delta
  \qquad\text{a.s.}
\]
Then \(S_h:=R_hU_h\) has the same \(1-\delta\) conclusion.

\item In particular, if the conditional sub-Weibull gauge tails
\[
  \mathbb P\{G_h>R_hu\mid\mcF\}
  \le C_0\exp(-c_0u^q)
  \qquad(u\ge1)
\]
hold a.s. for positive \(\mcF\)-measurable \(R_h\), then
\[
  u=\max\left\{1,
  \left(\frac{\max\{0,\log(C_0J/\delta)\}}{c_0}\right)^{1/q}
  \right\},
  \qquad S_h:=R_hu,
\]
has the same \(1-\delta\) conclusion.
\end{enumerate}

The conclusion is simultaneous over all \(h\le J\): after division by the
random scales \(S_h\), the self-adjoint tuples have discrepancy at most
\(\sqrt{n_h}\), with the analogous \(\discop\) bound for rectangular tuples.
The side information \(\mcF\) may contain arbitrary dependence between menu
elements; no conditional independence, centering, or identical-distribution
assumption is imposed.
\end{corollary}

\begin{proof}
For the moment form, conditional Markov and the conditional union bound give
\[
  \mathbb P\left\{\max_{h\le J}\frac{G_h}{S_h}>1\mid\mcF\right\}
  \le
  \sum_{h=1}^J\frac{\mathbb E(G_h^p\mid\mcF)}{S_h^p}
  \le\delta
  \qquad\text{a.s.}
\]
Taking expectations gives the unconditional failure probability bound.  The
tail-profile form is identical:
\[
  \mathbb P\left\{\max_{h\le J}\frac{G_h}{R_hU_h}>1\mid\mcF\right\}
  \le \sum_{h=1}^J\eta_h\le\delta
  \qquad\text{a.s.}
\]
The sub-Weibull choice makes this last sum at most
\(JC_0\exp(-c_0u^q)\le\delta\).  On the complementary event, every scaled tuple
has all-plus gauge at most one, so
Corollary~\ref{cor:distribution-free-all-plus-gauge-random-matrices} gives the
pathwise discrepancy certificate.
\end{proof}

\begin{corollary}[Random active finite-confidence gauge certificates]
\label{cor:distribution-free-gauge-random-active-confidence}
Let \(\mcF\) be side information and let
\(\mathcal H\subset\mathbb N\) be an almost surely finite
\(\mcF\)-measurable active set.  For each \(h\in\mathbb N\), let
\((X_{h,i})_{i\le n_h}\) be an arbitrary self-adjoint or rectangular random
matrix tuple, and let \(G_h\) be its all-plus gauge.  Fix \(0<\delta<1\).

\begin{enumerate}[label=\textup{(\roman*)}]
\item Fix \(p>0\).  Suppose positive \(\mcF\)-measurable random variables
\(R_h\) and positive \(\mcF\)-measurable budgets \(\eta_h\) satisfy
\[
  \mathbb E(G_h^p\mid\mcF)\le R_h^p
  \qquad\text{a.s.}
\]
and
\[
  \mathbb E\sum_{h\in\mathcal H}\eta_h\le\delta .
\]
Then the random scales
\[
  S_h:=R_h\eta_h^{-1/p}
\]
make every active scaled tuple contraction-normalized and all-plus certified
with probability at least \(1-\delta\).

\item More generally, suppose positive \(\mcF\)-measurable \(R_h\),
\(\mcF\)-measurable \(U_h\ge1\), and nonnegative \(\mcF\)-measurable budgets
\(\eta_h\) satisfy
\[
  \mathbb P\{G_h>R_hU_h\mid\mcF\}\le\eta_h
  \qquad\text{a.s.}
\]
and
\[
  \mathbb E\sum_{h\in\mathcal H}\eta_h\le\delta .
\]
Then \(S_h:=R_hU_h\) has the same \(1-\delta\) conclusion on the random active
menu.

\item In particular, suppose the conditional sub-Weibull gauge tails
\[
  \mathbb P\{G_h>R_hu\mid\mcF\}
  \le C_0\exp(-c_0u^q)
  \qquad(u\ge1)
\]
hold a.s. for positive \(\mcF\)-measurable \(R_h\), with fixed
\(C_0,c_0,q>0\).  If \(M\ge1\) is \(\mcF\)-measurable, \(\abs{\mathcal H}\le M\)
a.s., and \(u\ge1\) is \(\mcF\)-measurable with
\[
  \mathbb E\left[MC_0e^{-c_0u^q}\right]\le\delta,
\]
then \(S_h:=R_hu\) has the same \(1-\delta\) conclusion.  In particular, this
holds for the random multiplier
\[
  u=\max\left\{1,
  \left(\frac{\max\{0,\log(C_0M/\delta)\}}{c_0}\right)^{1/q}
  \right\}.
\]
\end{enumerate}

Thus a finite-confidence certificate may be requested only for the active
random catalogue selected by the environment.  The confidence cost is the
selected conditional budget, not the size of a deterministic ambient library.
No independence, conditional independence, centering, or identical-distribution
assumption is imposed.
\end{corollary}

\begin{proof}
For the moment form, conditional Markov gives
\[
  \mathbb P\{G_h>S_h\mid\mcF\}\le\eta_h
  \qquad\text{a.s.}
\]
Since \(\mathcal H\) is \(\mcF\)-measurable and finite,
\[
  \mathbb P\left\{\exists h\in\mathcal H:\,G_h>S_h\right\}
  \le
  \mathbb E\sum_{h\in\mathcal H}\eta_h
  \le\delta .
\]
On the complementary event, every selected scaled tuple has all-plus gauge at
most one, so
Corollary~\ref{cor:distribution-free-all-plus-gauge-random-matrices} gives the
simultaneous all-plus discrepancy certificate.

The chosen-threshold form is identical with
\[
  \mathbb P\{G_h>R_hU_h\mid\mcF\}\le\eta_h .
\]
For the conditional sub-Weibull form, the selected failure probability is at
most
\[
  \mathbb E\sum_{h\in\mathcal H}C_0e^{-c_0u^q}
  \le \mathbb E\left[MC_0e^{-c_0u^q}\right]\le\delta .
\]
The displayed random multiplier makes \(MC_0e^{-c_0u^q}\le\delta\) almost
surely.
\end{proof}

\begin{corollary}[Conditional gauge schedules]
\label{cor:distribution-free-gauge-conditional-schedules}
In the active triangular menu of
Corollary~\ref{cor:distribution-free-all-plus-gauge-random-matrices}, let
\(\mcF_m\) be side information for the \(m\)-th row and let \(G_{m,h}\) be the
all-plus gauge of the \(h\)-th active tuple.

\begin{enumerate}[label=\textup{(\roman*)}]
\item Fix \(p>0\).  Suppose positive \(\mcF_m\)-measurable random variables
\(R_{m,h}\) and positive deterministic budgets \(\eta_{m,h}\) satisfy
\[
  \mathbb E(G_{m,h}^p\mid\mcF_m)\le R_{m,h}^p
  \qquad\text{a.s.}
\]
and
\[
  \sum_m\sum_{h=1}^{J_m}\eta_{m,h}<\infty .
\]
Then the random scales
\[
  S_{m,h}:=R_{m,h}\eta_{m,h}^{-1/p}
\]
are valid: almost surely, for all sufficiently large \(m\), every active scaled
self-adjoint or rectangular tuple is contraction-normalized and has the
all-plus \(\sqrt{n_m}\) discrepancy certificate.

\item More generally, suppose positive \(\mcF_m\)-measurable \(R_{m,h}\),
\(\mcF_m\)-measurable \(U_{m,h}\ge1\), and nonnegative
\(\mcF_m\)-measurable budgets \(\eta_{m,h}\) satisfy
\[
  \mathbb P\{G_{m,h}>R_{m,h}U_{m,h}\mid\mcF_m\}
  \le \eta_{m,h}
  \qquad\text{a.s.}
\]
and
\[
  \sum_m\sum_{h=1}^{J_m}\mathbb E\eta_{m,h}<\infty .
\]
Then \(S_{m,h}:=R_{m,h}U_{m,h}\) is valid in the same almost-sure sense.

\item In particular, if the conditional sub-Weibull gauge tails
\[
  \mathbb P\{G_{m,h}>R_{m,h}u\mid\mcF_m\}
  \le C_0\exp(-c_0u^q)
  \qquad(u\ge1)
\]
hold a.s. for positive \(\mcF_m\)-measurable \(R_{m,h}\), with fixed
\(C_0,c_0,q>0\), then for every \(\gamma>0\) the row multiplier
\[
  u_m=\max\left\{1,
  \left(\frac{\max\{0,\log(C_0J_m)+(1+\gamma)\log(e+m)\}}{c_0}\right)^{1/q}
  \right\}
\]
and scales \(S_{m,h}:=R_{m,h}u_m\) are valid.
\end{enumerate}

The scales may depend on the side information and on the active environment.
No independence, conditional independence, centering, or identical-distribution
assumption is imposed.
\end{corollary}

\begin{proof}
For the moment form, conditional Markov gives
\[
  \mathbb P\{G_{m,h}>S_{m,h}\mid\mcF_m\}
  \le \frac{\mathbb E(G_{m,h}^p\mid\mcF_m)}{S_{m,h}^p}
  \le\eta_{m,h}
  \qquad\text{a.s.}
\]
After taking expectations and summing over \(m,h\), the bad events
\(\{G_{m,h}>S_{m,h}\}\) have summable probabilities.  Borel--Cantelli therefore
implies that almost surely only finitely many active scale failures occur.
On the complementary eventual event, the pathwise all-plus gauge certificate of
Corollary~\ref{cor:distribution-free-all-plus-gauge-random-matrices} applies to
each scaled tuple.

The chosen-threshold tail form is identical, since
\[
  \mathbb P\{G_{m,h}>R_{m,h}U_{m,h}\}
  =
  \mathbb E\,\mathbb P\{G_{m,h}>R_{m,h}U_{m,h}\mid\mcF_m\}
  \le \mathbb E\eta_{m,h}.
\]
For the conditional sub-Weibull choice, the \(m\)-th row failure probability is
bounded by
\[
  \sum_{h=1}^{J_m}C_0\exp(-c_0u_m^q)\le(e+m)^{-1-\gamma},
\]
which is summable.  The preceding Borel--Cantelli argument completes the proof.
\end{proof}

\begin{corollary}[Random active gauge schedules]
\label{cor:distribution-free-gauge-random-active-schedules}
For each \(m\), let \(\mcF_m\) be side information and let
\(\mathcal H_m\subset\mathbb N\) be an almost surely finite
\(\mcF_m\)-measurable active set.  For each \(h\in\mathbb N\), let
\((X_{m,h,i})_{i\le n_m}\) be an arbitrary self-adjoint or rectangular random
matrix tuple, and let \(G_{m,h}\) be its all-plus gauge.

\begin{enumerate}[label=\textup{(\roman*)}]
\item Fix \(p>0\).  Suppose positive \(\mcF_m\)-measurable random variables
\(R_{m,h}\) and positive \(\mcF_m\)-measurable budgets \(\eta_{m,h}\) satisfy
\[
  \mathbb E(G_{m,h}^p\mid\mcF_m)\le R_{m,h}^p
  \qquad\text{a.s.}
\]
and
\[
  \sum_m
  \mathbb E\sum_{h\in\mathcal H_m}\eta_{m,h}<\infty .
\]
Then the random scales
\[
  S_{m,h}:=R_{m,h}\eta_{m,h}^{-1/p}
\]
are valid on the random active menu: almost surely, for all sufficiently large
\(m\) and every \(h\in\mathcal H_m\), the scaled tuple is
contraction-normalized and has the all-plus \(\sqrt{n_m}\) discrepancy
certificate.

\item More generally, suppose positive \(\mcF_m\)-measurable \(R_{m,h}\),
\(\mcF_m\)-measurable \(U_{m,h}\ge1\), and nonnegative
\(\mcF_m\)-measurable budgets \(\eta_{m,h}\) satisfy
\[
  \mathbb P\{G_{m,h}>R_{m,h}U_{m,h}\mid\mcF_m\}
  \le\eta_{m,h}
  \qquad\text{a.s.}
\]
and
\[
  \sum_m
  \mathbb E\sum_{h\in\mathcal H_m}\eta_{m,h}<\infty .
\]
Then \(S_{m,h}:=R_{m,h}U_{m,h}\) is valid on the random active menu in the same
almost-sure sense.

\item In particular, suppose the conditional sub-Weibull gauge tails
\[
  \mathbb P\{G_{m,h}>R_{m,h}u\mid\mcF_m\}
  \le C_0\exp(-c_0u^q)
  \qquad(u\ge1)
\]
hold a.s. for positive \(\mcF_m\)-measurable \(R_{m,h}\), with fixed
\(C_0,c_0,q>0\).  If \(M_m\ge1\) is an \(\mcF_m\)-measurable random variable such
that \(\abs{\mathcal H_m}\le M_m\) a.s. and
\[
  \sum_m \mathbb E\left[M_mC_0e^{-c_0u_m^q}\right]<\infty,
\]
then \(S_{m,h}:=R_{m,h}u_m\) is valid on the random active menu.  In particular,
this holds for
\[
  u_m=\max\left\{1,
  \left(\frac{\max\{0,\log(C_0M_m)+(1+\gamma)\log(e+m)\}}{c_0}\right)^{1/q}
  \right\}
\]
whenever \(M_m\) is deterministic, or more generally whenever the displayed
series is summable.
\end{enumerate}

Thus the active catalogue itself may be random and environment-selected.  The
cost is only the selected conditional budget, not a deterministic ambient
library size.  No independence, conditional independence, centering, or
identical-distribution assumption is imposed.
\end{corollary}

\begin{proof}
For the moment form, conditional Markov gives
\[
  \mathbb P\{G_{m,h}>S_{m,h}\mid\mcF_m\}\le\eta_{m,h}
  \qquad\text{a.s.}
\]
Since \(\mathcal H_m\) is \(\mcF_m\)-measurable and finite,
\[
  \mathbb P\left\{\exists h\in\mathcal H_m:\,G_{m,h}>S_{m,h}\right\}
  \le
  \mathbb E\sum_{h\in\mathcal H_m}\eta_{m,h}.
\]
The assumed summability and Borel--Cantelli give eventual simultaneous control
on every selected active set.  On that event,
Corollary~\ref{cor:distribution-free-all-plus-gauge-random-matrices} gives the
pathwise all-plus discrepancy certificate for each selected scaled tuple.

The chosen-threshold tail form is the same calculation with
\[
  \mathbb P\{G_{m,h}>R_{m,h}U_{m,h}\mid\mcF_m\}\le\eta_{m,h}.
\]
For the conditional sub-Weibull form, the selected-row failure probability is at
most
\[
  \mathbb E\sum_{h\in\mathcal H_m}C_0e^{-c_0u_m^q}
  \le \mathbb E\left[M_mC_0e^{-c_0u_m^q}\right],
\]
which is summable by assumption.  If \(M_m\) is deterministic, the displayed
choice gives \(M_mC_0e^{-c_0u_m^q}\le(e+m)^{-1-\gamma}\), and hence
summability.
\end{proof}

\begin{corollary}[Full random matrices with Gaussian operator tails]
\label{cor:random-gaussian-tail-full}
Let \(Y_1,\ldots,Y_n\) be centered self-adjoint random matrices in
\(\M_D(\C)\).  Suppose that there is \(K_0\ge1\) such that for every
\(a=(a_1,\ldots,a_n)\in\R^n\) and every \(t\ge0\),
\[
  \mathbb P\left\{
  \left\|\sum_{i=1}^n a_iY_i\right\|_{\op}
  >K_0\|a\|_2(\sqrt D+t)
  \right\}\le 2e^{-t^2}.
\]
Put
\[
  s_{D,n}:=\sqrt D+\sqrt{3\log(en)},\qquad
  A_i:=\frac{Y_i}{K_0s_{D,n}} .
\]
Then, with probability at least \(1-2(n+1)(en)^{-3}\), all \(A_i\) are
contractions and
\[
  \disc(A_1,\ldots,A_n)\le \sqrt n .
\]
Thus, after the deterministic spectral scaling \(K_0(\sqrt D+\sqrt{\log(en)})\),
full random matrix ensembles satisfying the displayed operator-tail stability
obey the Matrix Spencer bound in every ambient dimension \(D\), with no
independence among the \(Y_i\) needed beyond the displayed tail.  This includes
the usual independent \(\mathrm{GOE}\), \(\mathrm{GUE}\), and centered
sub-Gaussian Wigner-type ensembles, as well as correlated centered ensembles
for which the same linear-combination operator-norm tail holds.
\end{corollary}

\begin{proof}
Take \(t=\sqrt{3\log(en)}\).  Applying the assumed tail with \(a=e_i\) gives
\[
  \mathbb P\{\|Y_i\|_{\op}>K_0s_{D,n}\}\le 2(en)^{-3}.
\]
A union bound over \(i=1,\ldots,n\) shows that all \(A_i\) are contractions
except with probability at most \(2n(en)^{-3}\).  Applying the same tail with
\(a=(1,\ldots,1)\) gives
\[
  \mathbb P\left\{
  \left\|\sum_{i=1}^nY_i\right\|_{\op}>K_0\sqrt n\,s_{D,n}
  \right\}\le 2(en)^{-3}.
\]
On the intersection of these two events, the all-plus signing certifies
\[
  \disc(A_1,\ldots,A_n)
  \le \left\|\sum_{i=1}^nA_i\right\|_{\op}
  \le \sqrt n .
\]
\end{proof}

\begin{corollary}[Active menus of full operator-tail random matrices]
\label{cor:active-menu-full-operator-tail-random-matrices}
For each \(m\ge1\), let \(D_m\ge2\), \(n_m\ge1\), and \(J_m\ge1\).  Let
\(\mathcal H_m\) be an environment, and for every active \(1\le h\le J_m\)
let
\[
  Y_{m,h,1},\ldots,Y_{m,h,n_m}\in\M_{D_m}(\C)
\]
be self-adjoint random matrices.  Suppose that \(K_{m,h}>0\) is
\(\mathcal H_m\)-measurable and that, conditionally on \(\mathcal H_m\), for
every \(a=(a_1,\ldots,a_{n_m})\in\R^{n_m}\) and every \(t\ge0\),
\[
  \mathbb P\left\{
  \left\|\sum_{i=1}^{n_m}a_iY_{m,h,i}\right\|_{\op}
  >K_{m,h}\|a\|_2(\sqrt{D_m}+t)
  \,\middle|\,\mathcal H_m
  \right\}\le 2e^{-t^2}.
\]
Let \(\tau_m>0\) be deterministic and define the normalized tuple
\[
  A_{m,h,i}:=
  \frac{Y_{m,h,i}}{K_{m,h}(\sqrt{D_m}+\tau_m)} .
\]
If
\[
  \sum_m J_m(n_m+1)e^{-\tau_m^2}<\infty,
\]
then, almost surely, for all sufficiently large \(m\), every active full
operator-tail stream is contraction-normalized and, for \(1\le h\le J_m\), has
the all-plus certificate
\[
  \max_{i\le n_m}\|A_{m,h,i}\|_{\op}\le1,
  \qquad
  \disc(A_{m,h,1},\ldots,A_{m,h,n_m})\le\sqrt{n_m}.
\]
No coordinatewise sign symmetry is required.  If \(J_m\to\infty\) and the
active prefixes exhaust a countable library, then every fixed library element
is eventually covered.

The summability condition has the following explicit schedules.  If, for some
\(\gamma>0\),
\[
  \tau_m^2\ge
  \log\bigl(2J_m(n_m+1)\bigr)+(1+\gamma)\log(e+m)
\]
eventually, then the active-menu conclusion holds.  If
\[
  J_m(n_m+1)\le(eD_m)^\Gamma
\]
eventually, then along all integer dimensions \(D_m=m\) it is enough to take
\[
  \tau_m^2=R\log(eD_m),\qquad R>\Gamma+1,
\]
while along every stretched-exponential sparse dimension schedule
\(D_m\ge\exp(cm^\kappa)\) it is enough that \(R>\Gamma\).

Consequently, Gaussian, Wigner/GOE/GUE, correlated Gaussian, and other full
random matrix libraries whose scalar linear combinations satisfy the displayed
operator-tail bound are simultaneously Matrix-Spencer scale after their natural
spectral normalization \(K_{m,h}(\sqrt{D_m}+\tau_m)\), even when the
normalizing tail scale is chosen from the environment.
\end{corollary}

\begin{proof}
Fix \(m\) and condition on \(\mathcal H_m\).  For a fixed active stream \(h\),
apply the conditional operator-tail hypothesis with \(a=e_i\),
\(1\le i\le n_m\), and with \(a=(1,\ldots,1)\), all at \(t=\tau_m\).  A union
bound gives conditional failure probability at most
\[
  2(n_m+1)e^{-\tau_m^2}
\]
for the event that either some \(\|A_{m,h,i}\|_{\op}>1\), or
\[
  \left\|\sum_{i=1}^{n_m}A_{m,h,i}\right\|_{\op}>\sqrt{n_m}.
\]
Union over \(h\le J_m\) and then average over \(\mathcal H_m\).  The displayed
summability and Borel--Cantelli imply that, almost surely, all these events
fail for every active \(h\) and all sufficiently large \(m\).  On that event
the all-plus sign vector gives the discrepancy bound.

The arbitrary-growth schedule gives
\[
  2J_m(n_m+1)e^{-\tau_m^2}\le(e+m)^{-1-\gamma},
\]
which is summable.  Under \(J_m(n_m+1)\le(eD_m)^\Gamma\) and
\(\tau_m^2=R\log(eD_m)\), the failure term is at most a constant multiple of
\((eD_m)^{\Gamma-R}\).  This is summable over integer dimensions when
\(R>\Gamma+1\), and is a stretched-exponential summable sequence along the
sparse dimension schedules when \(R>\Gamma\).
\end{proof}

\begin{corollary}[Random active full operator-tail random matrices]
\label{cor:random-active-full-operator-tail-random-matrices}
For each \(m\ge1\), let \(D_m\ge2\), \(n_m\ge1\), let \(\mcF_m\) be side
information, and let \(\mathcal A_m\subset\mathbb N\) be an almost surely finite
\(\mcF_m\)-measurable active set.  For every \(h\in\mathbb N\), let
\[
  Y_{m,h,1},\ldots,Y_{m,h,n_m}\in\M_{D_m}(\C)
\]
be self-adjoint random matrices.  Suppose that \(K_{m,h}>0\) and
\(\tau_{m,h}>0\) are \(\mcF_m\)-measurable and that, conditionally on
\(\mcF_m\), for every \(h\), every \(a\in\R^{n_m}\), and every \(t\ge0\),
\[
  \mathbb P\left\{
  \left\|\sum_{i=1}^{n_m}a_iY_{m,h,i}\right\|_{\op}
  >K_{m,h}\|a\|_2(\sqrt{D_m}+t)
  \,\middle|\,\mcF_m
  \right\}\le 2e^{-t^2}.
\]
Define
\[
  A_{m,h,i}:=
  \frac{Y_{m,h,i}}{K_{m,h}(\sqrt{D_m}+\tau_{m,h})}.
\]
If
\[
  \sum_m
  \mathbb E\sum_{h\in\mathcal A_m}
  2(n_m+1)e^{-\tau_{m,h}^2}<\infty,
\]
then, almost surely, for all sufficiently large \(m\), every selected active
full operator-tail stream is contraction-normalized and has the all-plus
certificate
\[
  \max_{i\le n_m}\|A_{m,h,i}\|_{\op}\le1,
  \qquad
  \disc(A_{m,h,1},\ldots,A_{m,h,n_m})\le\sqrt{n_m}
  \qquad(h\in\mathcal A_m).
\]
Thus random active operator-tail libraries pay only their selected tail cost,
not a deterministic ambient library size, and no coordinatewise sign symmetry
is required.

In particular, if \(M_m\ge1\) is \(\mcF_m\)-measurable,
\(\abs{\mathcal A_m}\le M_m\) a.s., and a common \(\tau_m>0\) satisfies
\[
  \sum_m \mathbb E\left[2M_m(n_m+1)e^{-\tau_m^2}\right]<\infty,
\]
then the conclusion holds with \(\tau_{m,h}=\tau_m\).  If \(M_m\) is
deterministic, it is enough to take, for any \(\gamma>0\),
\[
  \tau_m^2\ge
  \log\bigl(2M_m(n_m+1)\bigr)+(1+\gamma)\log(e+m)
\]
eventually.  More generally, if
\[
  \mathbb E[M_m](n_m+1)\le(eD_m)^\Gamma
\]
eventually, then \(\tau_m^2=R\log(eD_m)\) is enough along all integer dimensions
\(D_m=m\) whenever \(R>\Gamma+1\), and along every stretched-exponential sparse
dimension schedule \(D_m\ge\exp(cm^\kappa)\) whenever \(R>\Gamma\).
\end{corollary}

\begin{proof}
Fix \(m\) and condition on \(\mcF_m\).  For a selected stream \(h\), apply the
conditional operator-tail hypothesis with \(a=e_i\), \(1\le i\le n_m\), and
with \(a=(1,\ldots,1)\), all at \(t=\tau_{m,h}\).  The conditional probability
that either some \(\|A_{m,h,i}\|_{\op}>1\) or
\[
  \left\|\sum_{i=1}^{n_m}A_{m,h,i}\right\|_{\op}>\sqrt{n_m}
\]
is at most \(2(n_m+1)e^{-\tau_{m,h}^2}\).  Since \(\mathcal A_m\) is
\(\mcF_m\)-measurable and finite, the conditional selected-row failure
probability is at most
\[
  \sum_{h\in\mathcal A_m}2(n_m+1)e^{-\tau_{m,h}^2}.
\]
After taking expectations, the assumed summability and Borel--Cantelli give
eventual simultaneous control on every selected active set.  On that event the
all-plus sign vector gives the discrepancy bound.

The common-\(\tau_m\) assertion follows by the bound
\[
  \sum_{h\in\mathcal A_m}2(n_m+1)e^{-\tau_m^2}
  \le 2M_m(n_m+1)e^{-\tau_m^2}.
\]
The deterministic-size schedule makes this upper bound at most
\((e+m)^{-1-\gamma}\).  Under
\(\mathbb E[M_m](n_m+1)\le(eD_m)^\Gamma\) and
\(\tau_m^2=R\log(eD_m)\), the expected failure contribution is at most a
constant multiple of \((eD_m)^{\Gamma-R}\), which is summable in the two stated
regimes.
\end{proof}

\begin{corollary}[Active menus from quadratic-form tails]
\label{cor:active-menu-quadratic-form-full-random-matrices}
In the setting of
Corollary~\ref{cor:active-menu-full-operator-tail-random-matrices}, replace the
operator-tail hypothesis by the following scalar quadratic-form hypothesis.
For every active \(1\le h\le J_m\), suppose that \(q_{m,h}>0\) is
\(\mathcal H_m\)-measurable and that, conditionally on \(\mathcal H_m\), the
self-adjoint matrices \(Y_{m,h,i}\) are centered and, for every
\(a\in\R^{n_m}\), every deterministic unit vector \(u\in\C^{D_m}\), and every
\(s\ge0\),
\[
  \mathbb P\left\{
  \left|\left\langle u,
  \sum_{i=1}^{n_m}a_iY_{m,h,i}u
  \right\rangle\right|
  >q_{m,h}\|a\|_2s
  \,\middle|\,\mathcal H_m
  \right\}\le2e^{-s^2}.
\]
Set
\[
  \widehat A_{m,h,i}:=
  \frac{Y_{m,h,i}}{5q_{m,h}(\sqrt{D_m}+\tau_m)} .
\]
If
\[
  \sum_m J_m(n_m+1)e^{-\tau_m^2}<\infty,
\]
then, almost surely, for all sufficiently large \(m\), all active
\(\widehat A_{m,h,i}\) are contractions and the all-plus signing gives
\[
  \disc(\widehat A_{m,h,1},\ldots,\widehat A_{m,h,n_m})
  \le\sqrt{n_m}
  \qquad(1\le h\le J_m).
\]
The arbitrary-growth, all-integer, and sparse schedules of
Corollary~\ref{cor:active-menu-full-operator-tail-random-matrices} apply with
quadratic-form scale \(5q_{m,h}(\sqrt{D_m}+\tau_m)\).

Consequently, Gaussian, correlated Gaussian, Wigner-type, and other
self-adjoint random matrix libraries with uniform conditional scalar
quadratic-form sub-Gaussian proxies are simultaneously Matrix-Spencer scale
without requiring sign symmetry or entry independence.
\end{corollary}

\begin{proof}
Apply Lemma~\ref{lem:quadratic-form-to-operator-tail} conditionally on
\(\mathcal H_m\), for each active \(h\).  This gives the operator-tail
hypothesis of
Corollary~\ref{cor:active-menu-full-operator-tail-random-matrices} with
\(K_{m,h}=5q_{m,h}\).  The conclusions and schedules are therefore those of
that corollary.
\end{proof}

\begin{corollary}[Hermitian coefficient-series quadratic-form proxies]
\label{cor:active-menu-hermitian-coefficient-series-quadratic}
For each \(m\ge1\), let \(D_m\ge2\), \(n_m\ge1\), and \(J_m\ge1\).  Let
\(\mathcal H_m\) be an environment.  For each active \(1\le h\le J_m\), suppose
that
\[
  Y_{m,h,i}=\sum_{\alpha\in\mathcal A_{m,h,i}}
  \xi_{m,h,i,\alpha}B_{m,h,i,\alpha}\in\M_{D_m}(\C),
  \qquad 1\le i\le n_m,
\]
where the atom sets \(\mathcal A_{m,h,i}\) are finite, the atoms
\(B_{m,h,i,\alpha}\) are self-adjoint, and the atoms and coefficient proxies
\(\sigma_{m,h,i,\alpha}\ge0\) are \(\mathcal H_m\)-measurable.  Assume that,
conditionally on \(\mathcal H_m\), the real scalar coefficients are centered
and have the sub-Gaussian linear-combination estimate
\[
  \mathbb P\left\{
  \left|\sum_{i,\alpha}c_{i,\alpha}\xi_{m,h,i,\alpha}\right|
  >C_{\rm hser}
  \left(\sum_{i,\alpha}c_{i,\alpha}^2\sigma_{m,h,i,\alpha}^2\right)^{1/2}t
  \,\middle|\,\mathcal H_m
  \right\}\le2e^{-t^2}
\]
for every real coefficient array \(c=(c_{i,\alpha})\) and every \(t\ge0\).
This is implied, after changing \(C_{\rm hser}\), by conditional independence
and the usual real scalar sub-Gaussian mgf bounds.  Define
\[
  Q^{\rm h}_{m,h}:=
  \max_{1\le i\le n_m}
  \sup_{\|u\|_2=1}
  \left(
  \sum_{\alpha\in\mathcal A_{m,h,i}}
  \sigma_{m,h,i,\alpha}^2
  |\langle u,B_{m,h,i,\alpha}u\rangle|^2
  \right)^{1/2}.
\]
Assume \(Q^{\rm h}_{m,h}>0\) for the nonzero active streams, and set
\[
  H_{m,h,i}:=
  \frac{Y_{m,h,i}}
       {5C_{\rm hser}Q^{\rm h}_{m,h}(\sqrt{D_m}+\tau_m)} .
\]
If
\[
  \sum_m J_m(n_m+1)e^{-\tau_m^2}<\infty,
\]
then, almost surely, for all sufficiently large \(m\), all active
\(H_{m,h,i}\) are contractions and the all-plus signing gives
\[
  \disc(H_{m,h,1},\ldots,H_{m,h,n_m})\le\sqrt{n_m}
  \qquad(1\le h\le J_m).
\]
The arbitrary-growth, all-integer, and sparse schedules of
Corollary~\ref{cor:active-menu-full-operator-tail-random-matrices} apply with
scale
\[
  5C_{\rm hser}Q^{\rm h}_{m,h}(\sqrt{D_m}+\tau_m).
\]
Moreover
\[
  (Q^{\rm h}_{m,h})^2\le
  \max_{1\le i\le n_m}
  \left\|\sum_{\alpha\in\mathcal A_{m,h,i}}
  \sigma_{m,h,i,\alpha}^2B_{m,h,i,\alpha}^2\right\|_{\op}.
\]
Thus Hermitian sub-Gaussian atom expansions, correlated Wigner-type patterns,
random Hermitian feature blocks, and structured self-adjoint measurement
libraries are covered once this atom square function is controlled.
\end{corollary}

\begin{proof}
Fix \(m,h\) and condition on \(\mathcal H_m\).  For a deterministic unit vector
\(u\in\C^{D_m}\) and \(a\in\R^{n_m}\), apply the scalar coefficient tail with
\[
  c_{i,\alpha}=a_i\langle u,B_{m,h,i,\alpha}u\rangle .
\]
Since the atoms are self-adjoint, these coefficients are real.  The variance
proxy is bounded by
\[
  \begin{aligned}
  \sum_{i,\alpha}c_{i,\alpha}^2\sigma_{m,h,i,\alpha}^2
  &=
  \sum_{i=1}^{n_m}a_i^2
  \sum_{\alpha\in\mathcal A_{m,h,i}}
  \sigma_{m,h,i,\alpha}^2
  |\langle u,B_{m,h,i,\alpha}u\rangle|^2  \\
  &\le (Q^{\rm h}_{m,h})^2\|a\|_2^2 .
  \end{aligned}
\]
Therefore the hypothesis of
Corollary~\ref{cor:active-menu-quadratic-form-full-random-matrices} holds with
\(q_{m,h}=C_{\rm hser}Q^{\rm h}_{m,h}\), and that corollary gives the
simultaneous all-plus conclusion and schedules.

Finally, for each self-adjoint atom \(B\) and unit vector \(u\),
\(|\langle u,Bu\rangle|^2\le\langle u,B^2u\rangle\).  Summing this inequality
over \(\alpha\), taking the supremum over \(u\), and then the maximum over
\(i\) gives the displayed square-function bound.
\end{proof}

\begin{corollary}[Independent-entry Hermitian quadratic-form proxies]
\label{cor:active-menu-hermitian-independent-entry-quadratic}
For each \(m\ge1\), let \(D_m\ge2\), \(n_m\ge1\), and \(J_m\ge1\).  Let
\(\mathcal H_m\) be an environment.  For each active \(1\le h\le J_m\), let
\[
  \begin{aligned}
  Y_{m,h,i}
  &=
  \sum_{p=1}^{D_m}\xi_{m,h,i,p}E_{pp}  \\
  &\quad+
  \sum_{1\le p<q\le D_m}
  \left(\zeta_{m,h,i,pq}E_{pq}
  +\overline{\zeta_{m,h,i,pq}}E_{qp}\right),
  \qquad 1\le i\le n_m,
  \end{aligned}
\]
be centered self-adjoint random matrices.  Suppose that, conditionally on
\(\mathcal H_m\), all upper-triangular entries in each active stream are
independent.  Assume that \(\xi_{m,h,i,p}\) is real sub-Gaussian with proxy
\(\alpha_{m,h,i,p}\), and that every real projection
\(\operatorname{Re}(e^{\sqrt{-1}\theta}\zeta_{m,h,i,pq})\) is sub-Gaussian
with proxy \(\beta_{m,h,i,pq}\), uniformly in \(\theta\in\R\).  Define
\[
  Q^{\rm ent}_{m,h}:=
  \max_{1\le i\le n_m}\sup_{\|u\|_2=1}
  \left[
  \sum_p\alpha_{m,h,i,p}^2|u_p|^4
  +4\sum_{p<q}\beta_{m,h,i,pq}^2|u_p|^2|u_q|^2
  \right]^{1/2}.
\]
Assume \(Q^{\rm ent}_{m,h}>0\) for the nonzero active streams, and set
\[
  G_{m,h,i}:=
  \frac{Y_{m,h,i}}
       {5C_{\rm ent}Q^{\rm ent}_{m,h}(\sqrt{D_m}+\tau_m)} .
\]
If
\[
  \sum_m J_m(n_m+1)e^{-\tau_m^2}<\infty,
\]
then, almost surely, for all sufficiently large \(m\), all active
\(G_{m,h,i}\) are contractions and the all-plus signing gives
\[
  \disc(G_{m,h,1},\ldots,G_{m,h,n_m})\le\sqrt{n_m}
  \qquad(1\le h\le J_m).
\]
The arbitrary-growth, all-integer, and sparse schedules of
Corollary~\ref{cor:active-menu-full-operator-tail-random-matrices} apply with
scale
\[
  5C_{\rm ent}Q^{\rm ent}_{m,h}(\sqrt{D_m}+\tau_m).
\]
If all displayed entry proxies are bounded by \(\sigma_{m,h}\), then
\[
  Q^{\rm ent}_{m,h}\le\sqrt2\,\sigma_{m,h}.
\]
More generally, if
\[
  R_{m,h}^2\ge
  \max_{i,p}\left(
  \alpha_{m,h,i,p}^2
  +2\sum_{q\ne p}\beta_{m,h,i,\{p,q\}}^2
  \right),
\]
where \(\beta_{m,h,i,\{p,q\}}\) denotes the corresponding off-diagonal proxy,
then \(Q^{\rm ent}_{m,h}\le R_{m,h}\).  Thus homogeneous Wigner/GOE/GUE
libraries, variance-profile Hermitian ensembles, and sparse or banded
upper-triangular models are covered at the natural
\(Q(\sqrt D+\tau)\) spectral scale by the direct all-plus route.
\end{corollary}

\begin{proof}
Fix \(m,h\) and condition on \(\mathcal H_m\).  For a deterministic unit vector
\(u\in\C^{D_m}\), the quadratic form is
\[
  \langle u,Y_{m,h,i}u\rangle
  =\sum_p\xi_{m,h,i,p}|u_p|^2
  +2\sum_{p<q}\operatorname{Re}
  \left(\zeta_{m,h,i,pq}\overline{u_p}u_q\right).
\]
By the standard linear-combination rule for independent sub-Gaussian
variables, this scalar is conditionally sub-Gaussian with proxy at most
\[
  C_{\rm ent}
  \left[
  \sum_p\alpha_{m,h,i,p}^2|u_p|^4
  +4\sum_{p<q}\beta_{m,h,i,pq}^2|u_p|^2|u_q|^2
  \right]^{1/2}.
\]
Since the entries are independent across the whole active stream, for every
\(a\in\R^{n_m}\) the quadratic form
\(\langle u,\sum_i a_iY_{m,h,i}u\rangle\) has proxy at most
\(C_{\rm ent}Q^{\rm ent}_{m,h}\|a\|_2\), after increasing the universal
constant to match the \(2e^{-s^2}\) tail normalization.  Therefore
Corollary~\ref{cor:active-menu-quadratic-form-full-random-matrices} applies
with \(q_{m,h}=C_{\rm ent}Q^{\rm ent}_{m,h}\), giving the active all-plus
conclusion and schedules.

If all proxies are at most \(\sigma_{m,h}\), then, with
\(x_p=|u_p|^2\),
\[
  \sum_p x_p^2+4\sum_{p<q}x_px_q
  =2-\sum_p x_p^2\le2,
\]
which gives \(Q^{\rm ent}_{m,h}\le\sqrt2\,\sigma_{m,h}\).  For the row-variance
bound, use \(2x_px_q\le x_p^2+x_q^2\) to get
\[
  4\sum_{p<q}\beta_{m,h,i,pq}^2x_px_q
  \le
  2\sum_p x_p^2\sum_{q\ne p}\beta_{m,h,i,\{p,q\}}^2.
\]
The displayed row-budget hypothesis then bounds the whole square function by
\(R_{m,h}^2\sum_p x_p^2\le R_{m,h}^2\).
\end{proof}

\begin{corollary}[Bernstein quadratic-form active menus]
\label{cor:active-menu-bernstein-quadratic-form-random-matrices}
For each \(m\ge1\), let \(D_m\ge2\), \(n_m\ge1\), and \(J_m\ge1\).  Let
\(\mathcal H_m\) be an environment.  For each active \(1\le h\le J_m\), let
\[
  Y_{m,h,1},\ldots,Y_{m,h,n_m}\in\M_{D_m}(\C)
\]
be conditionally centered and self-adjoint.  Suppose that \(q_{m,h}\ge0\) and
\(b_{m,h}\ge0\) are \(\mathcal H_m\)-measurable.  Assume that, conditionally on
\(\mathcal H_m\), for every \(a\in\R^{n_m}\), every deterministic unit vector
\(u\in\C^{D_m}\), and every \(s\ge0\),
\[
  \mathbb P\left\{
  \left|\left\langle u,
  \sum_{i=1}^{n_m}a_iY_{m,h,i}u
  \right\rangle\right|
  >q_{m,h}\|a\|_2s+b_{m,h}\|a\|_\infty s^2
  \,\middle|\,\mathcal H_m
  \right\}\le2e^{-s^2}.
\]
Put
\[
  S_{m,h}:=C_{\rm Bq}\left[
  q_{m,h}(\sqrt{D_m}+\tau_m)
  +b_{m,h}(\sqrt{D_m}+\tau_m)^2
  \right],
\]
where \(C_{\rm Bq}\) is a sufficiently large numerical constant, and assume
\(S_{m,h}>0\) for nonzero active streams.  Define
\[
  B_{m,h,i}:=\frac{Y_{m,h,i}}{S_{m,h}} .
\]
If
\[
  \sum_m J_m(n_m+1)e^{-\tau_m^2}<\infty,
\]
then, almost surely, for all sufficiently large \(m\), all active
\(B_{m,h,i}\) are contractions and the all-plus signing gives
\[
  \disc(B_{m,h,1},\ldots,B_{m,h,n_m})\le\sqrt{n_m}
  \qquad(1\le h\le J_m).
\]
The arbitrary-growth, all-integer, and sparse schedules of
Corollary~\ref{cor:active-menu-full-operator-tail-random-matrices} apply with
the Bernstein scale \(S_{m,h}\).  Thus self-adjoint random matrix libraries whose
fixed quadratic forms have a Bernstein rather than sub-Gaussian linear-tail are
still covered by the direct all-plus route.
\end{corollary}

\begin{proof}
Fix \(m,h\) and condition on \(\mathcal H_m\).  The usual \(1/4\)-net argument
from Lemma~\ref{lem:quadratic-form-to-operator-tail}, applied at scalar level
\(c_0(\sqrt{D_m}+\tau_m)\) with \(c_0\) large enough, gives the following bound.
Put \(A_m=\sqrt{D_m}+\tau_m\) and
\[
  R_{m,h}(a):=
  C_{\rm Bq}\left[
  q_{m,h}\|a\|_2A_m+b_{m,h}\|a\|_\infty A_m^2
  \right].
\]
Then for every deterministic \(a\in\R^{n_m}\),
\[
  \mathbb P\left\{
  \left\|\sum_i a_iY_{m,h,i}\right\|_{\op}>R_{m,h}(a)
  \,\middle|\,\mathcal H_m\right\}
  \le 2e^{-\tau_m^2}.
\]
Indeed, the net has cardinality at most \(9^{2D_m}\), and the constant \(c_0\)
absorbs this entropy exactly as in Lemma~\ref{lem:quadratic-form-to-operator-tail}.
Apply the displayed bound to \(a=e_i\), \(1\le i\le n_m\), and to
\(a=(1,\ldots,1)\).  The first choices give \(\|B_{m,h,i}\|_{\op}\le1\) outside
conditional failure probability \(2n_me^{-\tau_m^2}\).  The all-plus choice gives
\[
\begin{aligned}
  \left\|\sum_iY_{m,h,i}\right\|_{\op}
  &\le
  C_{\rm Bq}\left[
  q_{m,h}\sqrt{n_m}A_m+b_{m,h}A_m^2
  \right]  \\
  &\le \sqrt{n_m}\,S_{m,h},
\end{aligned}
\]
because \(n_m\ge1\).  Thus the all-plus normalized sum has norm at most
\(\sqrt{n_m}\).  Union over \(h\le J_m\), average over the environment, and use
Borel--Cantelli.  The schedule statements are the same summability estimates as
in Corollary~\ref{cor:active-menu-full-operator-tail-random-matrices}.
\end{proof}

\begin{corollary}[Centered covariance and Wishart active menus]
\label{cor:active-menu-centered-covariance-wishart}
For each \(m\ge1\), let \(D_m\ge2\), \(n_m\ge1\), and \(J_m\ge1\).  Let
\(\mathcal H_m\) be an environment.  For each active \(1\le h\le J_m\), let
\(r_{m,h,i}\ge1\) and let
\[
  x_{m,h,i,\ell}\in\C^{D_m},
  \qquad
  1\le i\le n_m,\quad 1\le \ell\le r_{m,h,i},
\]
be conditionally independent random vectors.  Assume that, conditionally on
\(\mathcal H_m\), every one-dimensional real projection is centered
sub-Gaussian with proxy \(\kappa_{m,h,i,\ell}\): for every deterministic unit
vector \(u\in\C^{D_m}\), every \(\theta\in\R\), and every \(t\ge0\),
\[
  \mathbb P\left\{
  \left|\operatorname{Re}\left(e^{\sqrt{-1}\theta}
  \langle x_{m,h,i,\ell},u\rangle\right)\right|
  >\kappa_{m,h,i,\ell}t
  \,\middle|\,\mathcal H_m
  \right\}\le 2e^{-t^2}.
\]
Set
\[
  Y_{m,h,i}:=
  \frac1{r_{m,h,i}}\sum_{\ell=1}^{r_{m,h,i}}
  \left(
  x_{m,h,i,\ell}x_{m,h,i,\ell}^*
  -\mathbb E(x_{m,h,i,\ell}x_{m,h,i,\ell}^*\mid\mathcal H_m)
  \right),
\]
and define
\[
\begin{aligned}
  q^{\rm cov}_{m,h}
  &:=
  C_{\rm cov}\max_{1\le i\le n_m}
  \left(
  \frac1{r_{m,h,i}^2}
  \sum_{\ell=1}^{r_{m,h,i}}\kappa_{m,h,i,\ell}^4
  \right)^{1/2},\\
  b^{\rm cov}_{m,h}
  &:=
  C_{\rm cov}\max_{i,\ell}
  \frac{\kappa_{m,h,i,\ell}^2}{r_{m,h,i}} .
\end{aligned}
\]
With
\[
  S^{\rm cov}_{m,h}:=C_{\rm Bq}\left[
  q^{\rm cov}_{m,h}(\sqrt{D_m}+\tau_m)
  +b^{\rm cov}_{m,h}(\sqrt{D_m}+\tau_m)^2
  \right],
\]
the normalized matrices \(Y_{m,h,i}/S^{\rm cov}_{m,h}\) satisfy the conclusion
of Corollary~\ref{cor:active-menu-bernstein-quadratic-form-random-matrices}.
In particular, if all \(\kappa_{m,h,i,\ell}\le\kappa_{m,h}\) and all
\(r_{m,h,i}\ge r_{m,h}\), then the active covariance/Wishart scale is
\[
  C\kappa_{m,h}^2\left[
  \frac{\sqrt{D_m}+\tau_m}{\sqrt{r_{m,h}}}
  +\frac{(\sqrt{D_m}+\tau_m)^2}{r_{m,h}}
  \right].
\]
Rank-one centered covariance matrices correspond to \(r_{m,h}=1\), while
averaged real or complex Wishart increments and random feature covariance blocks
gain the displayed \(r^{-1/2}\) and \(r^{-1}\) improvements.
\end{corollary}

\begin{proof}
Fix \(m,h\) and condition on \(\mathcal H_m\).  For a unit vector \(u\), put
\[
  Z_{i,\ell}(u):=
  |\langle x_{m,h,i,\ell},u\rangle|^2
  -\mathbb E\left(|\langle x_{m,h,i,\ell},u\rangle|^2
  \,\middle|\,\mathcal H_m\right).
\]
The one-dimensional sub-Gaussian hypothesis implies the standard scalar
Bernstein estimate for centered squares: for every real coefficient array
\(c_{i,\ell}\) and every \(s\ge0\),
\[
  V_c:=
  \left(\sum_{i,\ell}c_{i,\ell}^2\kappa_{m,h,i,\ell}^4\right)^{1/2},
  \qquad
  M_c:=\max_{i,\ell}|c_{i,\ell}|\kappa_{m,h,i,\ell}^2,
\]
\[
  \mathbb P\left\{
  \left|\sum_{i,\ell}c_{i,\ell}Z_{i,\ell}(u)\right|
  >C_{\rm cov}\left(V_cs+M_cs^2\right)
  \,\middle|\,\mathcal H_m
  \right\}\le2e^{-s^2}.
\]
Apply this with \(c_{i,\ell}=a_i/r_{m,h,i}\).  Since
\[
  \left\langle u,\sum_i a_iY_{m,h,i}u\right\rangle
  =
  \sum_{i,\ell}\frac{a_i}{r_{m,h,i}}Z_{i,\ell}(u),
\]
the variance and jump terms are bounded by
\[
  q^{\rm cov}_{m,h}\|a\|_2,
  \qquad
  b^{\rm cov}_{m,h}\|a\|_\infty,
\]
respectively.  Hence the hypothesis of
Corollary~\ref{cor:active-menu-bernstein-quadratic-form-random-matrices} holds
with \(q_{m,h}=q^{\rm cov}_{m,h}\) and \(b_{m,h}=b^{\rm cov}_{m,h}\), and that
corollary gives the active-menu conclusion and schedules.  The homogeneous
display follows by the definitions of \(q^{\rm cov}_{m,h}\) and
\(b^{\rm cov}_{m,h}\).
\end{proof}

\begin{corollary}[Independent-entry Hermitian Bernstein-form proxies]
\label{cor:active-menu-hermitian-independent-entry-bernstein}
For each \(m\ge1\), let \(D_m\ge2\), \(n_m\ge1\), and \(J_m\ge1\).  Let
\(\mathcal H_m\) be an environment.  For each active \(1\le h\le J_m\), let
\[
  \begin{aligned}
  Y_{m,h,i}
  &=
  \sum_{p=1}^{D_m}\xi_{m,h,i,p}E_{pp}  \\
  &\quad+
  \sum_{1\le p<q\le D_m}
  \left(\zeta_{m,h,i,pq}E_{pq}
  +\overline{\zeta_{m,h,i,pq}}E_{qp}\right),
  \qquad 1\le i\le n_m,
  \end{aligned}
\]
be centered self-adjoint random matrices.  Suppose that, conditionally on
\(\mathcal H_m\), the upper-triangular variables in each active stream satisfy
the following scalar Bernstein linear-combination estimate.  There are
\(\mathcal H_m\)-measurable proxies
\(\alpha_{m,h,i,p},\lambda_{m,h,i,p}\ge0\) and
\(\beta_{m,h,i,pq},\mu_{m,h,i,pq}\ge0\) such that, for every real coefficient
arrays \(c_{i,p}\), \(d_{i,pq}\), every phases \(\theta_{i,pq}\in\R\), and
every \(s\ge0\),
\[
  Z(c,d,\theta):=
  \sum_{i,p}c_{i,p}\xi_{m,h,i,p}
  +\sum_{i,p<q}d_{i,pq}
  \operatorname{Re}\left(e^{\sqrt{-1}\theta_{i,pq}}
  \zeta_{m,h,i,pq}\right).
\]
Then
\[
  \mathbb P\left\{
  |Z(c,d,\theta)|>C_{\rm eB}\left(Vs+Ms^2\right)
  \,\middle|\,\mathcal H_m
  \right\}\le2e^{-s^2},
\]
where
\[
\begin{aligned}
  V^2&:=
  \sum_{i,p}c_{i,p}^2\alpha_{m,h,i,p}^2
  +\sum_{i,p<q}d_{i,pq}^2\beta_{m,h,i,pq}^2,\\
  M&:=
  \max\left\{
  \max_{i,p}|c_{i,p}|\lambda_{m,h,i,p},
  \max_{i,p<q}|d_{i,pq}|\mu_{m,h,i,pq}
  \right\}.
\end{aligned}
\]
This estimate is implied, after changing \(C_{\rm eB}\), by conditional
independence and the usual scalar Bernstein mgf bounds for the displayed real
variables.  Define
\[
  Q^{\rm eB}_{m,h}:=
  \max_{1\le i\le n_m}\sup_{\|u\|_2=1}
  \left[
  \sum_p\alpha_{m,h,i,p}^2|u_p|^4
  +4\sum_{p<q}\beta_{m,h,i,pq}^2|u_p|^2|u_q|^2
  \right]^{1/2}
\]
and
\[
  B^{\rm eB}_{m,h}:=
  \max_{1\le i\le n_m}
  \max\left\{
  \max_p\lambda_{m,h,i,p},
  \max_{p<q}\mu_{m,h,i,pq}
  \right\}.
\]
Set
\[
  S^{\rm eB}_{m,h}:=
  C_{\rm Bq}C_{\rm eB}\left[
  Q^{\rm eB}_{m,h}(\sqrt{D_m}+\tau_m)
  +B^{\rm eB}_{m,h}(\sqrt{D_m}+\tau_m)^2
  \right],
\]
and assume this scale is positive for nonzero active streams.  If
\[
  \sum_mJ_m(n_m+1)e^{-\tau_m^2}<\infty,
\]
then, almost surely, for all sufficiently large \(m\), all active matrices
\(Y_{m,h,i}/S^{\rm eB}_{m,h}\) are contractions and the all-plus signing gives
\[
  \disc\left(
  Y_{m,h,1}/S^{\rm eB}_{m,h},\ldots,
  Y_{m,h,n_m}/S^{\rm eB}_{m,h}
  \right)\le\sqrt{n_m}.
\]
The same arbitrary-growth, all-integer, and sparse schedules apply.  If
\(\alpha,\beta\le\sigma_{m,h}\) and \(\lambda,\mu\le L_{m,h}\) throughout the
active stream, then
\[
  Q^{\rm eB}_{m,h}\le\sqrt2\,\sigma_{m,h},
  \qquad
  B^{\rm eB}_{m,h}\le L_{m,h},
\]
so centered bounded, Poissonized, exponential-entry, gamma-entry, and other
Bernstein-entry Hermitian libraries enter at scale
\[
  C\left[
  \sigma_{m,h}(\sqrt{D_m}+\tau_m)
  +L_{m,h}(\sqrt{D_m}+\tau_m)^2
  \right].
\]
\end{corollary}

\begin{proof}
Fix \(m,h\) and condition on \(\mathcal H_m\).  For a deterministic unit vector
\(u\), write \(x_p=|u_p|^2\).  The quadratic form is
\[
  \langle u,Y_{m,h,i}u\rangle
  =
  \sum_p\xi_{m,h,i,p}x_p
  +2\sum_{p<q}\operatorname{Re}
  \left(\zeta_{m,h,i,pq}\overline{u_p}u_q\right).
\]
For \(\sum_i a_i\langle u,Y_{m,h,i}u\rangle\), choose
\[
  c_{i,p}=a_ix_p,\qquad
  d_{i,pq}=2a_i|u_p||u_q|,
\]
and choose the phases so that the second term agrees with \(Z(c,d,\theta)\)'s
rotated off-diagonal real parts.  The variance term is bounded by
\((Q^{\rm eB}_{m,h})^2\|a\|_2^2\).  Also
\(|c_{i,p}|\le|a_i|\) and
\[
  |d_{i,pq}|=2|a_i||u_p||u_q|\le |a_i|,
\]
so the jump term is bounded by \(B^{\rm eB}_{m,h}\|a\|_\infty\).  Therefore the
hypothesis of
Corollary~\ref{cor:active-menu-bernstein-quadratic-form-random-matrices} holds
with \(q_{m,h}=C_{\rm eB}Q^{\rm eB}_{m,h}\) and
\(b_{m,h}=C_{\rm eB}B^{\rm eB}_{m,h}\), giving the active all-plus conclusion.
The homogeneous bound for \(Q^{\rm eB}\) is the same computation as in
Corollary~\ref{cor:active-menu-hermitian-independent-entry-quadratic}; the
bound for \(B^{\rm eB}\) is immediate.
\end{proof}

\begin{corollary}[Active menus of rectangular full operator-tail matrices]
\label{cor:active-menu-rectangular-full-operator-tail-random-matrices}
For each \(m\ge1\), let \(d_{m,1},d_{m,2}\ge1\),
\[
  D_m^{\rm rec}:=d_{m,1}+d_{m,2},
\]
let \(n_m\ge1\), and let \(J_m\ge1\).  Let \(\mathcal H_m\) be an
environment, and for every active \(1\le h\le J_m\) let
\[
  X_{m,h,1},\ldots,X_{m,h,n_m}
  \in\M_{d_{m,1}\times d_{m,2}}(\C)
\]
be random rectangular matrices.  Suppose that \(K_{m,h}>0\) is
\(\mathcal H_m\)-measurable and that, conditionally on \(\mathcal H_m\), for
every \(a=(a_1,\ldots,a_{n_m})\in\R^{n_m}\) and every \(t\ge0\),
\[
  \mathbb P\left\{
  \left\|\sum_{i=1}^{n_m}a_iX_{m,h,i}\right\|_{\op}
  >K_{m,h}\|a\|_2(\sqrt{D_m^{\rm rec}}+t)
  \,\middle|\,\mathcal H_m
  \right\}\le 2e^{-t^2}.
\]
Let \(\tau_m>0\) be deterministic and define
\[
  Z_{m,h,i}:=
  \frac{X_{m,h,i}}{K_{m,h}(\sqrt{D_m^{\rm rec}}+\tau_m)} .
\]
If
\[
  \sum_m J_m(n_m+1)e^{-\tau_m^2}<\infty,
\]
then, almost surely, for all sufficiently large \(m\), every active
rectangular full operator-tail stream is contraction-normalized and, for
\(1\le h\le J_m\), has the all-plus certificate
\[
  \max_{i\le n_m}\|Z_{m,h,i}\|_{\op}\le1,
\]
\[
  \begin{aligned}
  \disc(Z_{m,h,1},\ldots,Z_{m,h,n_m})
  &\le \left\|\sum_{i=1}^{n_m}Z_{m,h,i}\right\|_{\op} \\
  &\le \sqrt{n_m}.
  \end{aligned}
\]
No coordinatewise sign symmetry, entry independence, or self-adjointness is
required beyond the displayed full operator-tail hypothesis.  If \(J_m\to\infty\)
and the active prefixes exhaust a countable library, then every fixed library
element is eventually covered.

The same schedules as in
Corollary~\ref{cor:active-menu-full-operator-tail-random-matrices} hold with
\(D_m\) replaced by \(D_m^{\rm rec}\).  In particular, if, for some
\(\gamma>0\),
\[
  \tau_m^2\ge
  \log\bigl(2J_m(n_m+1)\bigr)+(1+\gamma)\log(e+m)
\]
eventually, then the active-menu conclusion holds.  If
\[
  J_m(n_m+1)\le(eD_m^{\rm rec})^\Gamma
\]
eventually, then along all integer left-plus-right dimensions
\(D_m^{\rm rec}=m\) it is enough to take
\[
  \tau_m^2=R\log(eD_m^{\rm rec}),\qquad R>\Gamma+1,
\]
while along every stretched-exponential sparse schedule
\(D_m^{\rm rec}\ge\exp(cm^\kappa)\) it is enough that \(R>\Gamma\).

Consequently, rectangular Gaussian, rectangular sub-Gaussian, correlated
rectangular matrix-series, randomized feature-block, measurement-matrix, and
other non-self-adjoint full operator-tail libraries are simultaneously
Matrix-Spencer scale after their natural spectral normalization
\(K_{m,h}(\sqrt{D_m^{\rm rec}}+\tau_m)\).
\end{corollary}

\begin{proof}
Fix \(m\) and condition on \(\mathcal H_m\).  For a fixed active stream \(h\),
apply the conditional operator-tail hypothesis with \(a=e_i\),
\(1\le i\le n_m\), and with \(a=(1,\ldots,1)\), all at \(t=\tau_m\).  The
conditional probability that some \(\|Z_{m,h,i}\|_{\op}>1\), or that
\[
  \left\|\sum_{i=1}^{n_m}Z_{m,h,i}\right\|_{\op}>\sqrt{n_m},
\]
is at most \(2(n_m+1)e^{-\tau_m^2}\).  Union over \(h\le J_m\), average over
\(\mathcal H_m\), and apply Borel--Cantelli to the displayed summability
assumption.  On the resulting eventual event, the all-plus sign vector is an
admissible rectangular signing and gives the discrepancy bound.

The schedule estimates are identical to those in
Corollary~\ref{cor:active-menu-full-operator-tail-random-matrices}: the
arbitrary-growth hypothesis gives
\[
  2J_m(n_m+1)e^{-\tau_m^2}\le(e+m)^{-1-\gamma},
\]
and the polynomial active-menu condition with
\(\tau_m^2=R\log(eD_m^{\rm rec})\) gives a failure term bounded by a constant
multiple of \((eD_m^{\rm rec})^{\Gamma-R}\).  This is summable over all integer
left-plus-right dimensions for \(R>\Gamma+1\), and along the sparse schedules
for \(R>\Gamma\).
\end{proof}

\begin{corollary}[Random active rectangular full operator-tail matrices]
\label{cor:random-active-rectangular-full-operator-tail-random-matrices}
For each \(m\ge1\), let \(d_{m,1},d_{m,2}\ge1\),
\[
  D_m^{\rm rec}:=d_{m,1}+d_{m,2},
\]
let \(n_m\ge1\), let \(\mcF_m\) be side information, and let
\(\mathcal A_m\subset\mathbb N\) be an almost surely finite
\(\mcF_m\)-measurable active set.  For every \(h\in\mathbb N\), let
\[
  X_{m,h,1},\ldots,X_{m,h,n_m}
  \in\M_{d_{m,1}\times d_{m,2}}(\C)
\]
be random rectangular matrices.  Suppose that \(K_{m,h}>0\) and
\(\tau_{m,h}>0\) are \(\mcF_m\)-measurable and that, conditionally on
\(\mcF_m\), for every \(h\), every \(a\in\R^{n_m}\), and every \(t\ge0\),
\[
  \mathbb P\left\{
  \left\|\sum_{i=1}^{n_m}a_iX_{m,h,i}\right\|_{\op}
  >K_{m,h}\|a\|_2(\sqrt{D_m^{\rm rec}}+t)
  \,\middle|\,\mcF_m
  \right\}\le 2e^{-t^2}.
\]
Define
\[
  Z_{m,h,i}:=
  \frac{X_{m,h,i}}{K_{m,h}(\sqrt{D_m^{\rm rec}}+\tau_{m,h})}.
\]
If
\[
  \sum_m
  \mathbb E\sum_{h\in\mathcal A_m}
  2(n_m+1)e^{-\tau_{m,h}^2}<\infty,
\]
then, almost surely, for all sufficiently large \(m\), every selected active
rectangular full operator-tail stream is contraction-normalized and has the
all-plus certificate
\[
  \max_{i\le n_m}\|Z_{m,h,i}\|_{\op}\le1,
\]
\[
  \disc(Z_{m,h,1},\ldots,Z_{m,h,n_m})
  \le\left\|\sum_{i=1}^{n_m}Z_{m,h,i}\right\|_{\op}
  \le\sqrt{n_m}
  \qquad(h\in\mathcal A_m).
\]
Thus random active rectangular or non-self-adjoint operator-tail libraries pay
only their selected tail cost, not a deterministic ambient library size, and no
coordinatewise sign symmetry, entry independence, or self-adjointness is
required beyond the displayed full operator-tail hypothesis.

In particular, if \(M_m\ge1\) is \(\mcF_m\)-measurable,
\(\abs{\mathcal A_m}\le M_m\) a.s., and a common \(\tau_m>0\) satisfies
\[
  \sum_m \mathbb E\left[2M_m(n_m+1)e^{-\tau_m^2}\right]<\infty,
\]
then the conclusion holds with \(\tau_{m,h}=\tau_m\).  If \(M_m\) is
deterministic, it is enough to take, for any \(\gamma>0\),
\[
  \tau_m^2\ge
  \log\bigl(2M_m(n_m+1)\bigr)+(1+\gamma)\log(e+m)
\]
eventually.  More generally, if
\[
  \mathbb E[M_m](n_m+1)\le(eD_m^{\rm rec})^\Gamma
\]
eventually, then \(\tau_m^2=R\log(eD_m^{\rm rec})\) is enough along all integer
left-plus-right dimensions \(D_m^{\rm rec}=m\) whenever \(R>\Gamma+1\), and
along every stretched-exponential sparse schedule
\(D_m^{\rm rec}\ge\exp(cm^\kappa)\) whenever \(R>\Gamma\).
\end{corollary}

\begin{proof}
Fix \(m\) and condition on \(\mcF_m\).  For a selected stream \(h\), apply the
conditional operator-tail hypothesis with \(a=e_i\), \(1\le i\le n_m\), and
with \(a=(1,\ldots,1)\), all at \(t=\tau_{m,h}\).  The conditional probability
that either some \(\|Z_{m,h,i}\|_{\op}>1\) or
\[
  \left\|\sum_{i=1}^{n_m}Z_{m,h,i}\right\|_{\op}>\sqrt{n_m}
\]
is at most \(2(n_m+1)e^{-\tau_{m,h}^2}\).  Since \(\mathcal A_m\) is
\(\mcF_m\)-measurable and finite, the conditional selected-row failure
probability is at most
\[
  \sum_{h\in\mathcal A_m}2(n_m+1)e^{-\tau_{m,h}^2}.
\]
Taking expectations and using Borel--Cantelli gives eventual simultaneous
control on every selected active set.  On that event the all-plus sign vector
is an admissible rectangular signing and gives the discrepancy bound.

The common-\(\tau_m\) assertion follows from
\[
  \sum_{h\in\mathcal A_m}2(n_m+1)e^{-\tau_m^2}
  \le 2M_m(n_m+1)e^{-\tau_m^2}.
\]
The deterministic-size and polynomial expected-size schedules are the same
calculations as in Corollary~\ref{cor:random-active-full-operator-tail-random-matrices},
with \(D_m\) replaced by \(D_m^{\rm rec}\).
\end{proof}

\begin{lemma}[Bilinear-form tails imply rectangular operator tails]
\label{lem:bilinear-form-to-rectangular-operator-tail}
Let \(X_1,\ldots,X_n\in\M_{d_1\times d_2}(\C)\) be random rectangular matrices
relative to a sigma-field \(\mathcal H\), and put \(D_{\rm rec}=d_1+d_2\).
Suppose that there is \(q>0\) such that, conditionally on \(\mathcal H\), for
every \(a\in\R^n\), every pair of deterministic unit vectors
\(u\in\C^{d_1}\), \(v\in\C^{d_2}\), and every \(s\ge0\),
\[
  \mathbb P\left\{
  \left|\left\langle u,\sum_{i=1}^n a_iX_i v\right\rangle\right|
  >q\|a\|_2s
  \,\middle|\,\mathcal H
  \right\}\le 2e^{-s^2}.
\]
Then, conditionally on \(\mathcal H\), for every \(a\in\R^n\) and every
\(t\ge0\),
\[
  \mathbb P\left\{
  \left\|\sum_{i=1}^n a_iX_i\right\|_{\op}
  >5q\|a\|_2(\sqrt{D_{\rm rec}}+t)
  \,\middle|\,\mathcal H
  \right\}\le 2e^{-t^2}.
\]
\end{lemma}

\begin{proof}
Fix \(a\) and \(t\), and put \(S=\sum_i a_iX_i\).  Let
\(\mathcal N_1\) and \(\mathcal N_2\) be deterministic \(1/4\)-nets of the
unit spheres of \(\C^{d_1}\) and \(\C^{d_2}\), viewed as real spheres, with
\(|\mathcal N_j|\le9^{2d_j}\).  The standard two-net estimate gives
\[
  \|S\|_{\op}
  \le2\max_{u\in\mathcal N_1,\ v\in\mathcal N_2}|\langle u,Sv\rangle|.
\]
Therefore the conditional failure probability is at most
\[
  \begin{aligned}
  &2|\mathcal N_1||\mathcal N_2|
  \exp\left(-\frac{25}{4}(\sqrt{D_{\rm rec}}+t)^2\right) \\
  &\qquad\le
  2\exp\left(
    2D_{\rm rec}\log9-\frac{25}{4}D_{\rm rec}-\frac{25}{4}t^2
  \right)
  \le2e^{-t^2},
  \end{aligned}
\]
since \(25/4>2\log9\).
\end{proof}

\begin{corollary}[Active menus of bilinear-form rectangular matrices]
\label{cor:active-menu-bilinear-form-rectangular-random-matrices}
In the setting of
Corollary~\ref{cor:active-menu-rectangular-full-operator-tail-random-matrices},
replace the rectangular operator-tail hypothesis by the following scalar
bilinear-form hypothesis.  For each active \(1\le h\le J_m\), suppose that
\(q_{m,h}>0\) is \(\mathcal H_m\)-measurable and that, conditionally on
\(\mathcal H_m\), for every \(a\in\R^{n_m}\), every deterministic unit
\(u\in\C^{d_{m,1}}\), every deterministic unit \(v\in\C^{d_{m,2}}\), and every
\(s\ge0\),
\[
  \mathbb P\left\{
  \left|\left\langle u,
  \sum_{i=1}^{n_m}a_iX_{m,h,i}v
  \right\rangle\right|
  >q_{m,h}\|a\|_2s
  \,\middle|\,\mathcal H_m
  \right\}\le2e^{-s^2}.
\]
Let \(\tau_m>0\) be deterministic and set
\[
  \widetilde Z_{m,h,i}:=
  \frac{X_{m,h,i}}{5q_{m,h}(\sqrt{D_m^{\rm rec}}+\tau_m)} .
\]
If
\[
  \sum_m J_m(n_m+1)e^{-\tau_m^2}<\infty,
\]
then, almost surely, for all sufficiently large \(m\), all active
\(\widetilde Z_{m,h,i}\) are contractions and the all-plus signing gives
\[
  \disc(\widetilde Z_{m,h,1},\ldots,\widetilde Z_{m,h,n_m})
  \le\sqrt{n_m}
  \qquad(1\le h\le J_m).
\]
The arbitrary-growth, all-integer, and sparse schedules of
Corollary~\ref{cor:active-menu-rectangular-full-operator-tail-random-matrices}
apply with the bilinear scale \(5q_{m,h}(\sqrt{D_m^{\rm rec}}+\tau_m)\).

Consequently, conditionally Gaussian rectangular matrix series, correlated
rectangular Gaussian blocks, independent-entry rectangular sub-Gaussian
matrices, randomized feature blocks, and measurement-matrix libraries are
covered whenever all fixed bilinear forms have conditional sub-Gaussian proxy at
most \(q_{m,h}\|a\|_2\).
\end{corollary}

\begin{proof}
Apply Lemma~\ref{lem:bilinear-form-to-rectangular-operator-tail} conditionally
on \(\mathcal H_m\), for each active \(h\).  This gives the operator-tail
hypothesis of
Corollary~\ref{cor:active-menu-rectangular-full-operator-tail-random-matrices}
with \(K_{m,h}=5q_{m,h}\), and all conclusions follow from that corollary.
The listed examples are exactly the standard one-dimensional Gaussian or
sub-Gaussian fixed-bilinear-form estimates plugged into the displayed
hypothesis.
\end{proof}

\begin{corollary}[Random active bilinear-form rectangular matrices]
\label{cor:random-active-bilinear-form-rectangular-random-matrices}
In the setting of
Corollary~\ref{cor:random-active-rectangular-full-operator-tail-random-matrices},
replace the rectangular operator-tail hypothesis by the following scalar
bilinear-form hypothesis.  Suppose that \(q_{m,h}>0\) and \(\tau_{m,h}>0\)
are \(\mcF_m\)-measurable and that, conditionally on \(\mcF_m\), for every
\(h\), every \(a\in\R^{n_m}\), every deterministic unit
\(u\in\C^{d_{m,1}}\), every deterministic unit \(v\in\C^{d_{m,2}}\), and every
\(s\ge0\),
\[
  \mathbb P\left\{
  \left|\left\langle u,
  \sum_{i=1}^{n_m}a_iX_{m,h,i}v
  \right\rangle\right|
  >q_{m,h}\|a\|_2s
  \,\middle|\,\mcF_m
  \right\}\le2e^{-s^2}.
\]
Define
\[
  \widetilde Z_{m,h,i}:=
  \frac{X_{m,h,i}}{5q_{m,h}(\sqrt{D_m^{\rm rec}}+\tau_{m,h})}.
\]
If
\[
  \sum_m
  \mathbb E\sum_{h\in\mathcal A_m}
  2(n_m+1)e^{-\tau_{m,h}^2}<\infty,
\]
then, almost surely, for all sufficiently large \(m\), every selected active
bilinear-form rectangular stream is contraction-normalized and has the
all-plus certificate
\[
  \disc(\widetilde Z_{m,h,1},\ldots,\widetilde Z_{m,h,n_m})
  \le \sqrt{n_m}
  \qquad(h\in\mathcal A_m).
\]
The common-\(\tau_m\), deterministic-size, and dimension-schedule conclusions
from the random-active rectangular operator-tail corollary apply with the
bilinear scale
\[
  5q_{m,h}(\sqrt{D_m^{\rm rec}}+\tau_{m,h}).
\]

Thus random active conditionally Gaussian rectangular series, correlated
rectangular Gaussian blocks, independent-entry rectangular sub-Gaussian
matrices, randomized feature blocks, and measurement-matrix libraries pay only
their selected expected bilinear tail count.
\end{corollary}

\begin{proof}
Apply Lemma~\ref{lem:bilinear-form-to-rectangular-operator-tail} conditionally
on \(\mcF_m\), for each selected \(h\).  This gives the operator-tail
hypothesis of
Corollary~\ref{cor:random-active-rectangular-full-operator-tail-random-matrices}
with \(K_{m,h}=5q_{m,h}\).  The conclusion and schedules are those of that
corollary.
\end{proof}

\begin{corollary}[Haar unitary and orthogonal active menus]
\label{cor:active-menu-haar-unitary-orthogonal}
For each \(m\ge1\), let \(D_m\ge2\), \(n_m\ge1\), and \(J_m\ge1\).  Let
\(\mathcal H_m\) be an environment.  For each active \(1\le h\le J_m\), let
\[
  U_{m,h,1},\ldots,U_{m,h,n_m}\in\M_{D_m}(\C)
\]
be conditionally independent random matrices, each of which is either Haar
unitary on \(\C^{D_m}\) or Haar orthogonal on \(\R^{D_m}\), embedded in
\(\M_{D_m}(\C)\).  The choices of unitary versus orthogonal may depend on
\(\mathcal H_m\).  Put
\[
  q^{\rm Haar}_{m}:=C_{\rm Haar}D_m^{-1/2}
\]
and define
\[
  \widehat U_{m,h,i}:=
  \frac{U_{m,h,i}}
       {5q^{\rm Haar}_{m}(\sqrt{2D_m}+\tau_m)} .
\]
If
\[
  \sum_mJ_m(n_m+1)e^{-\tau_m^2}<\infty,
\]
then, almost surely, for all sufficiently large \(m\), all active
\(\widehat U_{m,h,i}\) are contractions and the all-plus signing gives
\[
  \disc(\widehat U_{m,h,1},\ldots,\widehat U_{m,h,n_m})
  \le\sqrt{n_m}
  \qquad(1\le h\le J_m).
\]
The arbitrary-growth, all-integer, and sparse schedules of
Corollary~\ref{cor:active-menu-rectangular-full-operator-tail-random-matrices}
apply with \(D_m^{\rm rec}=2D_m\).  In particular, if
\(\tau_m\le c\sqrt{D_m}\) eventually, then the normalizing denominator is bounded
above and below by positive constants depending only on \(c\).  Thus active
libraries of independent Haar unitary, Haar orthogonal, and mixed Haar
unitary/orthogonal matrices are Matrix-Spencer scale after only a universal
constant normalization.
\end{corollary}

\begin{proof}
Fix \(m,h\) and condition on \(\mathcal H_m\).  For deterministic unit vectors
\(u,v\in\C^{D_m}\), if \(U\) is Haar unitary, then \(Uv\) is uniform on the
complex unit sphere, and the standard spherical-coordinate tail gives
\[
  \mathbb P\left\{
  \left|\operatorname{Re}\left(e^{\sqrt{-1}\theta}\langle u,Uv\rangle\right)\right|
  >C D_m^{-1/2}s
  \,\middle|\,\mathcal H_m
  \right\}\le2e^{-s^2}
\]
for every \(\theta\in\R\).  For Haar orthogonal matrices embedded in
\(\M_{D_m}(\C)\), decompose \(u\) and \(v\) into real and imaginary parts and
apply the same real spherical-coordinate tail to the resulting four real
bilinear forms, increasing the numerical constant.  Conditional independence
therefore gives, for every \(a\in\R^{n_m}\),
\[
  \mathbb P\left\{
  \left|\left\langle u,\sum_i a_iU_{m,h,i}v\right\rangle\right|
  >q^{\rm Haar}_m\|a\|_2s
  \,\middle|\,\mathcal H_m
  \right\}\le2e^{-s^2}.
\]
Apply
Corollary~\ref{cor:active-menu-bilinear-form-rectangular-random-matrices} with
\(d_{m,1}=d_{m,2}=D_m\) and \(q_{m,h}=q^{\rm Haar}_m\).  The final constant-scale
claim follows from
\[
  5C_{\rm Haar}D_m^{-1/2}(\sqrt{2D_m}+\tau_m)
  =5C_{\rm Haar}(\sqrt2+\tau_m/\sqrt{D_m}).
\]
\end{proof}

\begin{corollary}[Independent-entry rectangular bilinear proxies]
\label{cor:active-menu-independent-entry-bilinear-rectangular}
For each \(m\ge1\), let \(d_{m,1},d_{m,2}\ge1\),
\(D_m^{\rm rec}=d_{m,1}+d_{m,2}\), \(n_m\ge1\), and \(J_m\ge1\).  Let
\(\mathcal H_m\) be an environment.  For each active \(1\le h\le J_m\), let
\[
  X_{m,h,i}=(x_{m,h,i,r,s})_{1\le r\le d_{m,1},\,1\le s\le d_{m,2}}
  \in\M_{d_{m,1}\times d_{m,2}}(\C),
  \qquad 1\le i\le n_m,
\]
be centered rectangular random matrices.  Suppose that, conditionally on
\(\mathcal H_m\), the entries in the active stream are independent and have
\(\mathcal H_m\)-measurable sub-Gaussian proxies
\(\sigma_{m,h,i,r,s}\ge0\) in the following linear-combination sense: for every
coefficient array \(c=(c_{i,r,s})\) and every \(t\ge0\),
\[
  \mathbb P\left\{
  \left|\sum_{i,r,s}c_{i,r,s}x_{m,h,i,r,s}\right|
  >
  C_{\rm ent}
  \left(\sum_{i,r,s}|c_{i,r,s}|^2\sigma_{m,h,i,r,s}^2\right)^{1/2}t
  \,\middle|\,\mathcal H_m
  \right\}\le2e^{-t^2},
\]
where \(C_{\rm ent}\) is an absolute constant.  This displayed hypothesis is
implied, after changing \(C_{\rm ent}\), by the usual conditional real or
complex sub-Gaussian entry mgf bounds.  Define
\[
  Q_{m,h}:=
  \max_{1\le i\le n_m}
  \sup_{\substack{\|u\|_2=1\\ \|v\|_2=1}}
  \left(
  \sum_{r=1}^{d_{m,1}}\sum_{s=1}^{d_{m,2}}
  |u_r|^2|v_s|^2\sigma_{m,h,i,r,s}^2
  \right)^{1/2}.
\]
Assume \(Q_{m,h}>0\) for the nonzero active streams, and define
\[
  W_{m,h,i}:=
  \frac{X_{m,h,i}}
       {5C_{\rm ent}Q_{m,h}(\sqrt{D_m^{\rm rec}}+\tau_m)} .
\]
If
\[
  \sum_m J_m(n_m+1)e^{-\tau_m^2}<\infty,
\]
then, almost surely, for all sufficiently large \(m\), all active
\(W_{m,h,i}\) are contractions and the all-plus signing gives
\[
  \disc(W_{m,h,1},\ldots,W_{m,h,n_m})\le\sqrt{n_m}
  \qquad(1\le h\le J_m).
\]
The arbitrary-growth, all-integer, and sparse schedules of
Corollary~\ref{cor:active-menu-rectangular-full-operator-tail-random-matrices}
apply with scale
\[
  5C_{\rm ent}Q_{m,h}(\sqrt{D_m^{\rm rec}}+\tau_m).
\]
In particular,
\[
  Q_{m,h}\le \max_{i,r,s}\sigma_{m,h,i,r,s},
\]
so homogeneous independent-entry rectangular sub-Gaussian matrices with entry
proxy \(\sigma_{m,h}\) are covered at the natural spectral scale
\(\sigma_{m,h}(\sqrt{D_m^{\rm rec}}+\tau_m)\), up to the universal constant.
\end{corollary}

\begin{proof}
Fix an active stream and condition on \(\mathcal H_m\).  For deterministic unit
vectors \(u\in\C^{d_{m,1}}\), \(v\in\C^{d_{m,2}}\), and
\(a\in\R^{n_m}\), apply the displayed scalar tail with
\[
  c_{i,r,s}=a_i\overline{u_r}v_s .
\]
The variance proxy is bounded by
\[
  \begin{aligned}
  \sum_{i,r,s}|c_{i,r,s}|^2\sigma_{m,h,i,r,s}^2
  &=
  \sum_{i=1}^{n_m}a_i^2
  \sum_{r,s}|u_r|^2|v_s|^2\sigma_{m,h,i,r,s}^2  \\
  &\le Q_{m,h}^2\|a\|_2^2 .
  \end{aligned}
\]
Thus the bilinear-form hypothesis of
Corollary~\ref{cor:active-menu-bilinear-form-rectangular-random-matrices}
holds with \(q_{m,h}=C_{\rm ent}Q_{m,h}\).  That corollary gives the
simultaneous all-plus conclusion and schedules.  Finally,
\(Q_{m,h}\le\max_{i,r,s}\sigma_{m,h,i,r,s}\), since
\(\sum_{r,s}|u_r|^2|v_s|^2=1\).
\end{proof}

\begin{corollary}[Sub-Gaussian coefficient-series rectangular bilinear proxies]
\label{cor:active-menu-coefficient-series-bilinear-rectangular}
For each \(m\ge1\), let \(d_{m,1},d_{m,2}\ge1\),
\(D_m^{\rm rec}=d_{m,1}+d_{m,2}\), \(n_m\ge1\), and \(J_m\ge1\).  Let
\(\mathcal H_m\) be an environment.  For each active \(1\le h\le J_m\), suppose
that
\[
  X_{m,h,i}=\sum_{\alpha\in\mathcal A_{m,h,i}}
  \xi_{m,h,i,\alpha}B_{m,h,i,\alpha}
  \in\M_{d_{m,1}\times d_{m,2}}(\C),
  \qquad 1\le i\le n_m,
\]
where the atom sets \(\mathcal A_{m,h,i}\) are finite and the atom matrices
\(B_{m,h,i,\alpha}\) and coefficient proxies
\(\sigma_{m,h,i,\alpha}\ge0\) are \(\mathcal H_m\)-measurable.  Assume that,
conditionally on \(\mathcal H_m\), the scalar coefficients are centered and have
the following sub-Gaussian linear-combination estimate: for every coefficient
array \(c=(c_{i,\alpha})\) and every \(t\ge0\),
\[
  V_c:=
  \left(\sum_{i=1}^{n_m}\sum_{\alpha\in\mathcal A_{m,h,i}}
  |c_{i,\alpha}|^2\sigma_{m,h,i,\alpha}^2\right)^{1/2}t
\]
satisfies
\[
  \mathbb P\left\{
  \left|\sum_{i=1}^{n_m}\sum_{\alpha\in\mathcal A_{m,h,i}}
  c_{i,\alpha}\xi_{m,h,i,\alpha}\right|
  >C_{\rm ser}V_c
  \,\middle|\,\mathcal H_m
  \right\}\le2e^{-t^2}.
\]
This is implied, after changing the absolute constant \(C_{\rm ser}\), by
conditional independence and the usual real or complex scalar sub-Gaussian mgf
bounds.  Define the bilinear atom-energy proxy
\[
  R_{m,h}:=\max_{1\le i\le n_m}
  \sup_{\substack{\|u\|_2=1\\ \|v\|_2=1}}
  \left(
  \sum_{\alpha\in\mathcal A_{m,h,i}}
  \sigma_{m,h,i,\alpha}^2
  |\langle u,B_{m,h,i,\alpha}v\rangle|^2
  \right)^{1/2}.
\]
Assume \(R_{m,h}>0\) for the nonzero active streams, and set
\[
  U_{m,h,i}:=
  \frac{X_{m,h,i}}
       {5C_{\rm ser}R_{m,h}(\sqrt{D_m^{\rm rec}}+\tau_m)} .
\]
If
\[
  \sum_m J_m(n_m+1)e^{-\tau_m^2}<\infty,
\]
then, almost surely, for all sufficiently large \(m\), all active
\(U_{m,h,i}\) are contractions and the all-plus signing gives
\[
  \disc(U_{m,h,1},\ldots,U_{m,h,n_m})\le\sqrt{n_m}
  \qquad(1\le h\le J_m).
\]
The arbitrary-growth, all-integer, and sparse schedules of
Corollary~\ref{cor:active-menu-rectangular-full-operator-tail-random-matrices}
apply with scale
\[
  5C_{\rm ser}R_{m,h}(\sqrt{D_m^{\rm rec}}+\tau_m).
\]
Moreover
\[
  R_{m,h}^2\le
  \max_{1\le i\le n_m}
  \min\{L^{\rm at}_{m,h,i},R^{\rm at}_{m,h,i}\},
\]
where
\[
  \begin{aligned}
  L^{\rm at}_{m,h,i}
  &:=
  \left\|\sum_{\alpha\in\mathcal A_{m,h,i}}
  \sigma_{m,h,i,\alpha}^2
  B_{m,h,i,\alpha}B_{m,h,i,\alpha}^*\right\|_{\op},\\
  R^{\rm at}_{m,h,i}
  &:=
  \left\|\sum_{\alpha\in\mathcal A_{m,h,i}}
  \sigma_{m,h,i,\alpha}^2
  B_{m,h,i,\alpha}^*B_{m,h,i,\alpha}\right\|_{\op}.
  \end{aligned}
\]
Consequently this covers atomized rectangular sub-Gaussian series,
common-feature correlated entries, random feature blocks, and structured
measurement libraries once the fixed-bilinear atom energy is controlled.
\end{corollary}

\begin{proof}
Fix \(m,h\) and condition on \(\mathcal H_m\).  For deterministic unit vectors
\(u\in\C^{d_{m,1}}\), \(v\in\C^{d_{m,2}}\), and \(a\in\R^{n_m}\), apply the
scalar coefficient tail with
\[
  c_{i,\alpha}=a_i\langle u,B_{m,h,i,\alpha}v\rangle .
\]
The variance proxy is bounded by
\[
  \begin{aligned}
  \sum_{i,\alpha}|c_{i,\alpha}|^2\sigma_{m,h,i,\alpha}^2
  &=
  \sum_{i=1}^{n_m}a_i^2
  \sum_{\alpha\in\mathcal A_{m,h,i}}
  \sigma_{m,h,i,\alpha}^2
  |\langle u,B_{m,h,i,\alpha}v\rangle|^2  \\
  &\le R_{m,h}^2\|a\|_2^2 .
  \end{aligned}
\]
Therefore the hypothesis of
Corollary~\ref{cor:active-menu-bilinear-form-rectangular-random-matrices}
holds with \(q_{m,h}=C_{\rm ser}R_{m,h}\), and that corollary gives the
simultaneous all-plus conclusion and schedules.

For the displayed computable bound, for each \(i,u,v\),
\[
  \sum_{\alpha}\sigma_{m,h,i,\alpha}^2
  |\langle u,B_{m,h,i,\alpha}v\rangle|^2
  \le
  \left\|\sum_{\alpha}
  \sigma_{m,h,i,\alpha}^2B_{m,h,i,\alpha}B_{m,h,i,\alpha}^*
  \right\|_{\op},
\]
because \(vv^*\le I\).  The same argument with \(uu^*\le I\) gives the
right-variance bound with \(B_{m,h,i,\alpha}^*B_{m,h,i,\alpha}\).  Taking the
minimum of the two bounds and then the maximum over \(i\) proves the claim.
\end{proof}

\begin{corollary}[Bernstein bilinear-form rectangular active menus]
\label{cor:active-menu-bernstein-bilinear-rectangular}
In the setting of
Corollary~\ref{cor:active-menu-rectangular-full-operator-tail-random-matrices},
replace the rectangular operator-tail hypothesis by the following scalar
Bernstein bilinear-form hypothesis.  For each active \(1\le h\le J_m\), suppose
that \(q_{m,h}\ge0\) and \(b_{m,h}\ge0\) are \(\mathcal H_m\)-measurable and
that, conditionally on \(\mathcal H_m\), the rectangular matrices
\(X_{m,h,i}\) are centered and, for every \(a\in\R^{n_m}\), every deterministic
unit \(u\in\C^{d_{m,1}}\), every deterministic unit \(v\in\C^{d_{m,2}}\), and
every \(s\ge0\),
\[
  \mathbb P\left\{
  \left|\left\langle u,
  \sum_{i=1}^{n_m}a_iX_{m,h,i}v
  \right\rangle\right|
  >q_{m,h}\|a\|_2s+b_{m,h}\|a\|_\infty s^2
  \,\middle|\,\mathcal H_m
  \right\}\le2e^{-s^2}.
\]
Put
\[
  T_{m,h}:=C_{\rm Bb}\left[
  q_{m,h}(\sqrt{D_m^{\rm rec}}+\tau_m)
  +b_{m,h}(\sqrt{D_m^{\rm rec}}+\tau_m)^2
  \right],
\]
where \(C_{\rm Bb}\) is a sufficiently large numerical constant, and assume
\(T_{m,h}>0\) for nonzero active streams.  Define
\[
  V_{m,h,i}:=\frac{X_{m,h,i}}{T_{m,h}} .
\]
If
\[
  \sum_mJ_m(n_m+1)e^{-\tau_m^2}<\infty,
\]
then, almost surely, for all sufficiently large \(m\), all active
\(V_{m,h,i}\) are contractions and the all-plus signing gives
\[
  \disc(V_{m,h,1},\ldots,V_{m,h,n_m})\le\sqrt{n_m}
  \qquad(1\le h\le J_m).
\]
The arbitrary-growth, all-integer, and sparse schedules of
Corollary~\ref{cor:active-menu-rectangular-full-operator-tail-random-matrices}
apply with the Bernstein bilinear scale \(T_{m,h}\).
\end{corollary}

\begin{proof}
Fix \(m,h\) and condition on \(\mathcal H_m\).  Let
\(A_m=\sqrt{D_m^{\rm rec}}+\tau_m\).  The same two-net argument as in
Lemma~\ref{lem:bilinear-form-to-rectangular-operator-tail}, using the scalar
Bernstein hypothesis at level \(cA_m\) with \(c\) large enough, gives for every
deterministic \(a\in\R^{n_m}\)
\[
  \mathbb P\left\{
  \left\|\sum_i a_iX_{m,h,i}\right\|_{\op}
  >C_{\rm Bb}\left[
  q_{m,h}\|a\|_2A_m+b_{m,h}\|a\|_\infty A_m^2
  \right]
  \,\middle|\,\mathcal H_m
  \right\}\le2e^{-\tau_m^2}.
\]
Apply this bound to \(a=e_i\), \(1\le i\le n_m\), and to
\(a=(1,\ldots,1)\).  The \(e_i\) choices give \(\|V_{m,h,i}\|_{\op}\le1\)
outside conditional failure probability \(2n_me^{-\tau_m^2}\).  The all-plus
choice gives
\[
  \left\|\sum_iX_{m,h,i}\right\|_{\op}
  \le C_{\rm Bb}\left[
  q_{m,h}\sqrt{n_m}A_m+b_{m,h}A_m^2
  \right]
  \le\sqrt{n_m}\,T_{m,h},
\]
since \(n_m\ge1\).  Union over \(h\le J_m\), average over the environment, and
apply Borel--Cantelli.  The schedule statements are the same summability
estimates as in the rectangular full operator-tail corollary.
\end{proof}

\begin{corollary}[Random active Bernstein bilinear-form rectangular matrices]
\label{cor:random-active-bernstein-bilinear-rectangular}
In the setting of
Corollary~\ref{cor:random-active-rectangular-full-operator-tail-random-matrices},
replace the rectangular operator-tail hypothesis by the following scalar
Bernstein bilinear-form hypothesis.  Suppose that \(q_{m,h}\ge0\),
\(b_{m,h}\ge0\), and \(\tau_{m,h}>0\) are \(\mcF_m\)-measurable and that,
conditionally on \(\mcF_m\), the rectangular matrices \(X_{m,h,i}\) are
centered and, for every \(h\), every \(a\in\R^{n_m}\), every deterministic unit
\(u\in\C^{d_{m,1}}\), every deterministic unit \(v\in\C^{d_{m,2}}\), and every
\(s\ge0\),
\[
  \mathbb P\left\{
  \left|\left\langle u,
  \sum_{i=1}^{n_m}a_iX_{m,h,i}v
  \right\rangle\right|
  >q_{m,h}\|a\|_2s+b_{m,h}\|a\|_\infty s^2
  \,\middle|\,\mcF_m
  \right\}\le2e^{-s^2}.
\]
Put
\[
  A_{m,h}:=\sqrt{D_m^{\rm rec}}+\tau_{m,h},
  \qquad
  T_{m,h}:=C_{\rm Bb}\bigl(q_{m,h}A_{m,h}+b_{m,h}A_{m,h}^2\bigr),
\]
where \(C_{\rm Bb}\) is a sufficiently large numerical constant, and define
\[
  \widehat Z_{m,h,i}:=
  \begin{cases}
  X_{m,h,i}/T_{m,h},& T_{m,h}>0,\\
  0,& T_{m,h}=0.
  \end{cases}
\]
If
\[
  \sum_m
  \mathbb E\sum_{h\in\mathcal A_m}
  2(n_m+1)e^{-\tau_{m,h}^2}<\infty,
\]
then, almost surely, for all sufficiently large \(m\), every selected active
Bernstein bilinear-form rectangular stream is contraction-normalized and has
the all-plus certificate
\[
  \disc(\widehat Z_{m,h,1},\ldots,\widehat Z_{m,h,n_m})
  \le \sqrt{n_m}
  \qquad(h\in\mathcal A_m).
\]
The common-\(\tau_m\), deterministic-size, and dimension-schedule conclusions
from the random-active rectangular operator-tail corollary apply with the
Bernstein bilinear scale \(T_{m,h}\).  Thus selected bounded-entry,
Poissonized, exponential-entry, gamma-entry, and balanced random-dyad
rectangular libraries pay only their selected expected Bernstein bilinear tail
count.
\end{corollary}

\begin{proof}
Condition on \(\mcF_m\) and fix a selected \(h\).  Let
\(A_{m,h}=\sqrt{D_m^{\rm rec}}+\tau_{m,h}\).  The two-net argument used in
Corollary~\ref{cor:active-menu-bernstein-bilinear-rectangular} gives, for every
deterministic \(a\in\R^{n_m}\),
\[
  \begin{aligned}
  &\mathbb P\left\{
  \left\|\sum_i a_iX_{m,h,i}\right\|_{\op}
  >C_{\rm Bb}\left[
  q_{m,h}\|a\|_2A_{m,h}+b_{m,h}\|a\|_\infty A_{m,h}^2
  \right]
  \,\middle|\,\mcF_m
  \right\} \\
  &\qquad\le2e^{-\tau_{m,h}^2}.
  \end{aligned}
\]
Apply this estimate to \(a=e_i\), \(1\le i\le n_m\), and to
\(a=(1,\ldots,1)\).  These \(n_m+1\) estimates show that, outside conditional
failure probability \(2(n_m+1)e^{-\tau_{m,h}^2}\), every
\(\widehat Z_{m,h,i}\) is a contraction and
\[
  \left\|\sum_i\widehat Z_{m,h,i}\right\|_{\op}\le\sqrt{n_m}.
\]
Union over \(h\in\mathcal A_m\), average over the side information, and apply
Borel--Cantelli.  The common-\(\tau_m\) and dimension-schedule statements are
the same summability estimates as in
Corollary~\ref{cor:random-active-rectangular-full-operator-tail-random-matrices}.
\end{proof}

\begin{corollary}[Independent-entry rectangular Bernstein-form proxies]
\label{cor:active-menu-independent-entry-bernstein-rectangular}
For each \(m\ge1\), let \(d_{m,1},d_{m,2}\ge1\),
\(D_m^{\rm rec}=d_{m,1}+d_{m,2}\), \(n_m\ge1\), and \(J_m\ge1\).  Let
\(\mathcal H_m\) be an environment.  For each active \(1\le h\le J_m\), let
\[
  X_{m,h,i}=(x_{m,h,i,r,s})_{1\le r\le d_{m,1},\,1\le s\le d_{m,2}}
  \in\M_{d_{m,1}\times d_{m,2}}(\C),
  \qquad 1\le i\le n_m,
\]
be centered rectangular random matrices.  Suppose that, conditionally on
\(\mathcal H_m\), the entries in the active stream satisfy the following
complex scalar Bernstein linear-combination estimate.  There are
\(\mathcal H_m\)-measurable proxies
\(\sigma_{m,h,i,r,s},L_{m,h,i,r,s}\ge0\) such that, for every coefficient
array \(c=(c_{i,r,s})\) and every \(t\ge0\),
\[
  \mathbb P\left\{
  \left|\sum_{i,r,s}c_{i,r,s}x_{m,h,i,r,s}\right|
  >C_{\rm rB}\left(V_ct+M_ct^2\right)
  \,\middle|\,\mathcal H_m
  \right\}\le2e^{-t^2},
\]
where
\[
  V_c^2:=\sum_{i,r,s}|c_{i,r,s}|^2\sigma_{m,h,i,r,s}^2,
  \qquad
  M_c:=\max_{i,r,s}|c_{i,r,s}|L_{m,h,i,r,s}.
\]
This estimate is implied, after changing \(C_{\rm rB}\), by conditional
independence and the usual real and imaginary scalar Bernstein mgf bounds.
Define
\[
  Q^{\rm rB}_{m,h}:=
  \max_{1\le i\le n_m}
  \sup_{\substack{\|u\|_2=1\\ \|v\|_2=1}}
  \left(
  \sum_{r=1}^{d_{m,1}}\sum_{s=1}^{d_{m,2}}
  |u_r|^2|v_s|^2\sigma_{m,h,i,r,s}^2
  \right)^{1/2}
\]
and
\[
  B^{\rm rB}_{m,h}:=\max_{i,r,s}L_{m,h,i,r,s}.
\]
Set
\[
  T^{\rm rB}_{m,h}:=C_{\rm Bb}C_{\rm rB}\left[
  Q^{\rm rB}_{m,h}(\sqrt{D_m^{\rm rec}}+\tau_m)
  +B^{\rm rB}_{m,h}(\sqrt{D_m^{\rm rec}}+\tau_m)^2
  \right],
\]
and assume this scale is positive for nonzero active streams.  Then, under
\[
  \sum_mJ_m(n_m+1)e^{-\tau_m^2}<\infty,
\]
the normalized matrices \(X_{m,h,i}/T^{\rm rB}_{m,h}\) are eventually
contractions and the all-plus signing gives discrepancy at most \(\sqrt{n_m}\)
simultaneously for every active \(h\).  If
\(\sigma_{m,h,i,r,s}\le\sigma_{m,h}\) and
\(L_{m,h,i,r,s}\le L_{m,h}\), then
\[
  Q^{\rm rB}_{m,h}\le\sigma_{m,h},\qquad
  B^{\rm rB}_{m,h}\le L_{m,h},
\]
so bounded, Poissonized, exponential-entry, gamma-entry, and other
light-tailed non-sub-Gaussian rectangular libraries are covered at scale
\[
  C\left[
  \sigma_{m,h}(\sqrt{D_m^{\rm rec}}+\tau_m)
  +L_{m,h}(\sqrt{D_m^{\rm rec}}+\tau_m)^2
  \right].
\]
\end{corollary}

\begin{proof}
Fix \(m,h\) and condition on \(\mathcal H_m\).  For deterministic unit vectors
\(u\in\C^{d_{m,1}}\), \(v\in\C^{d_{m,2}}\), and \(a\in\R^{n_m}\), apply the
displayed scalar Bernstein estimate with
\[
  c_{i,r,s}=a_i\overline{u_r}v_s .
\]
Then \(V_c\le Q^{\rm rB}_{m,h}\|a\|_2\), while
\[
  M_c\le B^{\rm rB}_{m,h}\|a\|_\infty
\]
because \(|u_r|\le1\) and \(|v_s|\le1\).  Hence the hypothesis of
Corollary~\ref{cor:active-menu-bernstein-bilinear-rectangular} holds with
\(q_{m,h}=C_{\rm rB}Q^{\rm rB}_{m,h}\) and
\(b_{m,h}=C_{\rm rB}B^{\rm rB}_{m,h}\).  The homogeneous bound follows from
\(\sum_{r,s}|u_r|^2|v_s|^2=1\).
\end{proof}

\begin{corollary}[Random active independent-entry rectangular Bernstein proxies]
\label{cor:random-active-independent-entry-bernstein-rectangular}
For each \(m\ge1\), let \(d_{m,1},d_{m,2}\ge1\),
\(D_m^{\rm rec}=d_{m,1}+d_{m,2}\), \(n_m\ge1\), let \(\mcF_m\) be side
information, and let \(\mathcal A_m\subset\mathbb N\) be an almost surely
finite \(\mcF_m\)-measurable active set.  For every \(h\in\mathbb N\), let
\[
  X_{m,h,i}=(x_{m,h,i,r,s})_{1\le r\le d_{m,1},\,1\le s\le d_{m,2}}
  \in\M_{d_{m,1}\times d_{m,2}}(\C),
  \qquad 1\le i\le n_m,
\]
be centered rectangular random matrices.  Suppose that, conditionally on
\(\mcF_m\), the entries in every selected stream satisfy the following complex
scalar Bernstein linear-combination estimate.  There are
\(\mcF_m\)-measurable proxies \(\sigma_{m,h,i,r,s},L_{m,h,i,r,s}\ge0\) such
that, for every \(h\), every coefficient array \(c=(c_{i,r,s})\), and every
\(t\ge0\),
\[
  \mathbb P\left\{
  \left|\sum_{i,r,s}c_{i,r,s}x_{m,h,i,r,s}\right|
  >C_{\rm rB}\left(V_{m,h,c}t+M_{m,h,c}t^2\right)
  \,\middle|\,\mcF_m
  \right\}\le2e^{-t^2},
\]
where
\[
  V_{m,h,c}^2:=\sum_{i,r,s}|c_{i,r,s}|^2\sigma_{m,h,i,r,s}^2,
  \qquad
  M_{m,h,c}:=\max_{i,r,s}|c_{i,r,s}|L_{m,h,i,r,s}.
\]
Define
\[
  Q^{\rm rB}_{m,h}:=
  \max_{1\le i\le n_m}
  \sup_{\substack{\|u\|_2=1\\ \|v\|_2=1}}
  \left(
  \sum_{r=1}^{d_{m,1}}\sum_{s=1}^{d_{m,2}}
  |u_r|^2|v_s|^2\sigma_{m,h,i,r,s}^2
  \right)^{1/2},
\]
and
\[
  B^{\rm rB}_{m,h}:=\max_{i,r,s}L_{m,h,i,r,s}.
\]
Let \(\tau_{m,h}>0\) be \(\mcF_m\)-measurable and put
\[
  T^{\rm rB}_{m,h}:=C_{\rm Bb}C_{\rm rB}\left[
  Q^{\rm rB}_{m,h}(\sqrt{D_m^{\rm rec}}+\tau_{m,h})
  +B^{\rm rB}_{m,h}(\sqrt{D_m^{\rm rec}}+\tau_{m,h})^2
  \right].
\]
Define
\[
  Z^{\rm rB}_{m,h,i}:=
  \begin{cases}
  X_{m,h,i}/T^{\rm rB}_{m,h},& T^{\rm rB}_{m,h}>0,\\
  0,& T^{\rm rB}_{m,h}=0.
  \end{cases}
\]
If
\[
  \sum_m
  \mathbb E\sum_{h\in\mathcal A_m}
  2(n_m+1)e^{-\tau_{m,h}^2}<\infty,
\]
then, almost surely, for all sufficiently large \(m\), every selected
independent-entry Bernstein stream is contraction-normalized after the
displayed normalization and has the all-plus certificate
\[
  \disc\left(
  Z^{\rm rB}_{m,h,1},\ldots,
  Z^{\rm rB}_{m,h,n_m}
  \right)\le\sqrt{n_m}
  \qquad(h\in\mathcal A_m).
\]
If \(\sigma_{m,h,i,r,s}\le\sigma_{m,h}\) and
\(L_{m,h,i,r,s}\le L_{m,h}\), then one may use
\[
  C\left[
  \sigma_{m,h}(\sqrt{D_m^{\rm rec}}+\tau_{m,h})
  +L_{m,h}(\sqrt{D_m^{\rm rec}}+\tau_{m,h})^2
  \right]
\]
instead.  Thus random active bounded-entry, Poissonized, exponential-entry,
gamma-entry, and other light-tailed independent-entry rectangular libraries pay
only their selected expected Bernstein entry-tail count.
\end{corollary}

\begin{proof}
For deterministic unit vectors \(u\in\C^{d_{m,1}}\), \(v\in\C^{d_{m,2}}\),
and \(a\in\R^{n_m}\), use the scalar Bernstein estimate with
\[
  c_{i,r,s}=a_i\overline{u_r}v_s .
\]
Then \(V_{m,h,c}\le Q^{\rm rB}_{m,h}\|a\|_2\) and
\(M_{m,h,c}\le B^{\rm rB}_{m,h}\|a\|_\infty\).  Hence the hypothesis of
Corollary~\ref{cor:random-active-bernstein-bilinear-rectangular} holds with
\(q_{m,h}=C_{\rm rB}Q^{\rm rB}_{m,h}\) and
\(b_{m,h}=C_{\rm rB}B^{\rm rB}_{m,h}\).  That corollary gives the selected
active conclusion.  The homogeneous scale follows from
\(\sum_{r,s}|u_r|^2|v_s|^2=1\).
\end{proof}

\begin{corollary}[Random active conditionally independent Bernstein entries]
\label{cor:random-active-entrywise-bernstein-rectangular}
In Corollary~\ref{cor:random-active-independent-entry-bernstein-rectangular},
replace the scalar linear-combination hypothesis by the following entrywise
condition.  Conditionally on \(\mcF_m\), for every selected \(h\) the variables
\((x_{m,h,i,r,s})_{i,r,s}\) are independent, centered, and satisfy the uniform
rotated Bernstein mgf bound: for every phase \(\theta\in\R\) and every real
\(\lambda\) with \(|\lambda|L_{m,h,i,r,s}<1\),
\[
  \mathbb E\left[
  \exp\left\{
  \lambda\operatorname{Re}(e^{-i\theta}x_{m,h,i,r,s})
  \right\}
  \,\middle|\,\mcF_m
  \right]
  \le
  \exp\left\{
  \frac{\lambda^2\sigma_{m,h,i,r,s}^2}
       {2(1-|\lambda|L_{m,h,i,r,s})}
  \right\}.
\]
Then the conclusions of
Corollary~\ref{cor:random-active-independent-entry-bernstein-rectangular} hold,
after changing only the absolute numerical constant in the normalization.
In particular, the same selected summability condition
\[
  \sum_m
  \mathbb E\sum_{h\in\mathcal A_m}
  2(n_m+1)e^{-\tau_{m,h}^2}<\infty
\]
gives eventual all-plus Matrix-Spencer certificates for bounded, Poissonized,
centered exponential or gamma, and other conditionally independent rectangular
entry models with the displayed variance and jump proxies.
\end{corollary}

\begin{proof}
Fix \(m,h\) and condition on \(\mcF_m\).  For a deterministic coefficient array
\(c=(c_{i,r,s})\), put
\[
\begin{aligned}
  Z_c&:=\sum_{i,r,s}c_{i,r,s}x_{m,h,i,r,s},\\
  V_c^2&:=\sum_{i,r,s}|c_{i,r,s}|^2\sigma_{m,h,i,r,s}^2,\\
  M_c&:=\max_{i,r,s}|c_{i,r,s}|L_{m,h,i,r,s}.
\end{aligned}
\]
For every fixed phase \(\phi\), independence and the entrywise mgf bound give
\[
  \mathbb E\left[
  \exp\{\lambda\operatorname{Re}(e^{-i\phi}Z_c)\}
  \,\middle|\,\mcF_m
  \right]
  \le
  \exp\left\{\frac{\lambda^2V_c^2}{2(1-|\lambda|M_c)}\right\},
  \qquad |\lambda|M_c<1.
\]
The scalar Bernstein inequality therefore gives
\[
  \mathbb P\left\{
  |\operatorname{Re}(e^{-i\phi}Z_c)|
  > C(V_ct+M_ct^2)
  \,\middle|\,\mcF_m
  \right\}\le 2e^{-t^2}.
\]
Taking \(\phi=0\) and \(\phi=\pi/2\), and using
\(|Z_c|\le\sqrt2\max\{|\operatorname{Re}Z_c|,|\operatorname{Im}Z_c|\}\),
gives the scalar complex Bernstein linear-combination estimate required in
Corollary~\ref{cor:random-active-independent-entry-bernstein-rectangular},
with a changed absolute constant.  Applying that corollary proves the claim.
\end{proof}

\begin{corollary}[Finite-confidence entrywise Bernstein random matrices]
\label{cor:random-active-entrywise-bernstein-finite-confidence}
Fix one dimension pair \(d_1,d_2\ge1\), put \(D^{\rm rec}=d_1+d_2\), and let
\(\mcF\) be side information.  Let \(\mathcal A\subset\mathbb N\) be an almost
surely finite \(\mcF\)-measurable active set.  For every \(h\), let
\[
  X_{h,i}=(x_{h,i,r,s})_{1\le r\le d_1,\,1\le s\le d_2},
  \qquad 1\le i\le n,
\]
be centered rectangular matrices whose selected entries satisfy the conditional
independence and rotated Bernstein mgf hypothesis of
Corollary~\ref{cor:random-active-entrywise-bernstein-rectangular}, with
proxies \(\sigma_{h,i,r,s}\) and \(L_{h,i,r,s}\).  Define \(Q_h^{\rm rB}\),
\(B_h^{\rm rB}\), \(T_h^{\rm rB}\), and \(Z_{h,i}^{\rm rB}\) as in
Corollary~\ref{cor:random-active-independent-entry-bernstein-rectangular}, with
an \(\mcF\)-measurable choice \(\tau_h>0\).  If
\[
  \mathbb E\sum_{h\in\mathcal A}2(n+1)e^{-\tau_h^2}\le\delta,
\]
then, with probability at least \(1-\delta\), every selected entrywise
Bernstein stream is contraction-normalized and has the all-plus certificate
\[
  \disc\left(
  Z_{h,1}^{\rm rB},\ldots,Z_{h,n}^{\rm rB}
  \right)\le\sqrt n
  \qquad(h\in\mathcal A).
\]
In particular, if \(\abs{\mathcal A}\le M\) almost surely, the common choice
\[
  \tau^2\ge \log\frac{2M(n+1)}{\delta}
\]
gives a simultaneous \(1-\delta\) certificate after the homogeneous scale
\[
  C\left[
  \sigma_h(\sqrt{D^{\rm rec}}+\tau)
  +L_h(\sqrt{D^{\rm rec}}+\tau)^2
  \right]
\]
whenever \(\sigma_{h,i,r,s}\le\sigma_h\) and \(L_{h,i,r,s}\le L_h\).
\end{corollary}

\begin{proof}
Condition on \(\mcF\).  By the proof of
Corollary~\ref{cor:random-active-entrywise-bernstein-rectangular}, each selected
stream satisfies the scalar Bernstein linear-combination hypothesis of
Corollary~\ref{cor:random-active-independent-entry-bernstein-rectangular}.  The
conditional failure probability for a fixed selected \(h\) is therefore at most
\(2(n+1)e^{-\tau_h^2}\).  Summing over the \(\mcF\)-measurable active set and
then taking expectation gives total failure at most the displayed budget
\(\delta\).  The common-\(\tau\) and homogeneous scale statements are the
corresponding specializations.
\end{proof}

\begin{corollary}[Finite-confidence bounded-entry random matrices]
\label{cor:bounded-entry-bernstein-finite-confidence}
Fix one dimension pair \(d_1,d_2\ge1\), put \(D^{\rm rec}=d_1+d_2\), and let
\(\mcF\) be side information.  Let \(\mathcal A\subset\mathbb N\) be an almost
surely finite \(\mcF\)-measurable active set.  For every \(h\), let
\[
  X_{h,i}=(x_{h,i,r,s})_{1\le r\le d_1,\,1\le s\le d_2},
  \qquad 1\le i\le n,
\]
be centered rectangular matrices.  Assume that, conditionally on \(\mcF\), for
every selected \(h\) the entries \((x_{h,i,r,s})_{i,r,s}\) are independent and
there are \(\mcF\)-measurable proxies
\(\nu_{h,i,r,s},R_{h,i,r,s}\ge0\) such that
\[
  \mathbb E\left(|x_{h,i,r,s}|^2\mid\mcF\right)\le\nu_{h,i,r,s}^2,
  \qquad
  |x_{h,i,r,s}|\le R_{h,i,r,s}
  \quad\text{a.s.}
\]
Define
\[
  Q^{\rm bd}_h:=
  \max_{1\le i\le n}
  \sup_{\substack{\|u\|_2=1\\ \|v\|_2=1}}
  \left(
  \sum_{r=1}^{d_1}\sum_{s=1}^{d_2}
  |u_r|^2|v_s|^2\nu_{h,i,r,s}^2
  \right)^{1/2},
\]
and
\[
  R^{\rm bd}_h:=\max_{i,r,s}R_{h,i,r,s}.
\]
Let \(\tau_h>0\) be \(\mcF\)-measurable and put
\[
  T^{\rm bd}_h:=
  C_{\rm bd}\left[
  Q^{\rm bd}_h(\sqrt{D^{\rm rec}}+\tau_h)
  +R^{\rm bd}_h(\sqrt{D^{\rm rec}}+\tau_h)^2
  \right],
\]
with
\[
  Z^{\rm bd}_{h,i}:=
  \begin{cases}
  X_{h,i}/T^{\rm bd}_h,& T^{\rm bd}_h>0,\\
  0,& T^{\rm bd}_h=0.
  \end{cases}
\]
If
\[
  \mathbb E\sum_{h\in\mathcal A}2(n+1)e^{-\tau_h^2}\le\delta,
\]
then, with probability at least \(1-\delta\), every selected bounded-entry
stream is contraction-normalized and has the all-plus certificate
\[
  \disc\left(
  Z_{h,1}^{\rm bd},\ldots,Z_{h,n}^{\rm bd}
  \right)\le\sqrt n
  \qquad(h\in\mathcal A).
\]
In particular, if \(\abs{\mathcal A}\le M\) almost surely, the common choice
\[
  \tau^2\ge\log\frac{2M(n+1)}{\delta}
\]
gives a simultaneous \(1-\delta\) certificate after the homogeneous scale
\[
  C\left[
  \nu_h(\sqrt{D^{\rm rec}}+\tau)
  +R_h(\sqrt{D^{\rm rec}}+\tau)^2
  \right]
\]
whenever \(\nu_{h,i,r,s}\le\nu_h\) and \(R_{h,i,r,s}\le R_h\).
\end{corollary}

\begin{proof}
Condition on \(\mcF\) and fix a selected stream.  For every phase
\(\theta\in\R\), the scalar variable
\[
  Y_{h,i,r,s,\theta}:=
  \operatorname{Re}(e^{-i\theta}x_{h,i,r,s})
\]
is centered, has conditional variance at most \(\nu_{h,i,r,s}^2\), and satisfies
\(|Y_{h,i,r,s,\theta}|\le R_{h,i,r,s}\).  The standard bounded-variable
Bernstein mgf bound gives, for \(|\lambda|R_{h,i,r,s}<1\),
\[
  \mathbb E\left[
  \exp\{\lambda Y_{h,i,r,s,\theta}\}\mid\mcF
  \right]
  \le
  \exp\left\{
  \frac{C\lambda^2\nu_{h,i,r,s}^2}
       {2(1-|\lambda|R_{h,i,r,s})}
  \right\}.
\]
Absorbing the absolute constant into the final normalization, this is precisely
the rotated Bernstein mgf hypothesis of
Corollary~\ref{cor:random-active-entrywise-bernstein-finite-confidence}, with
\(\sigma_{h,i,r,s}=C\nu_{h,i,r,s}\) and
\(L_{h,i,r,s}=R_{h,i,r,s}\).  That corollary gives the finite-confidence
certificate, and the homogeneous bound follows from the displayed definitions of
\(Q_h^{\rm bd}\) and \(R_h^{\rm bd}\).
\end{proof}

\begin{corollary}[Haar rank-one dyad active menus]
\label{cor:active-menu-haar-rank-one-dyads}
For each \(m\ge1\), let \(d_{m,1},d_{m,2}\ge2\),
\(D_m^{\rm rec}=d_{m,1}+d_{m,2}\), \(n_m\ge1\), and \(J_m\ge1\).  Let
\(\mathcal H_m\) be an environment.  For each active \(1\le h\le J_m\), let
\[
  x_{m,h,i}\in\C^{d_{m,1}},
  \qquad
  y_{m,h,i}\in\C^{d_{m,2}},
  \qquad 1\le i\le n_m,
\]
be conditionally independent unit vectors, each Haar-distributed on the
corresponding complex unit sphere, and let
\(\rho_{m,h,i}\) be \(\mathcal H_m\)-measurable amplitudes with
\(|\rho_{m,h,i}|\le\rho_{m,h}\).  Real Haar unit vectors, embedded in the
complex spaces, are also allowed after changing the absolute constant.  Put
\[
  X_{m,h,i}:=\rho_{m,h,i}\,x_{m,h,i}y_{m,h,i}^*
\]
and
\[
  T^{\rm dyad}_{m,h}:=
  C_{\rm dyad}\rho_{m,h}
  \left[
  \frac{\sqrt{D_m^{\rm rec}}+\tau_m}{\sqrt{d_{m,1}d_{m,2}}}
  +\frac{(\sqrt{D_m^{\rm rec}}+\tau_m)^2}{\sqrt{d_{m,1}d_{m,2}}}
  \right].
\]
Set \(Z^{\rm dyad}_{m,h,i}:=X_{m,h,i}/T^{\rm dyad}_{m,h}\).  If
\[
  \sum_mJ_m(n_m+1)e^{-\tau_m^2}<\infty,
\]
then, almost surely, for all sufficiently large \(m\), all active
\(Z^{\rm dyad}_{m,h,i}\) are contractions and the all-plus signing gives
\[
  \disc\left(
  Z^{\rm dyad}_{m,h,1},\ldots,Z^{\rm dyad}_{m,h,n_m}
  \right)\le\sqrt{n_m}.
\]
In particular, if \(d_{m,1}\) and \(d_{m,2}\) are comparable and
\(\tau_m=O(\sqrt{D_m^{\rm rec}})\), then independent Haar rank-one dyads,
random feature outer products, and mixed real/complex Haar dyad libraries are
Matrix-Spencer scale after only a universal constant multiple of the amplitude
normalization.
\end{corollary}

\begin{proof}
Fix \(m,h\) and condition on \(\mathcal H_m\).  For deterministic unit vectors
\(u\in\C^{d_{m,1}}\) and \(v\in\C^{d_{m,2}}\),
\[
  \langle u,X_{m,h,i}v\rangle
  =
  \rho_{m,h,i}\langle u,x_{m,h,i}\rangle
  \langle y_{m,h,i},v\rangle .
\]
The spherical-coordinate tail gives sub-Gaussian proxies
\(Cd_{m,1}^{-1/2}\) and \(Cd_{m,2}^{-1/2}\) for the two scalar factors,
uniformly over real projections.  The product of two independent sub-Gaussian
variables is sub-exponential with proxy the product of their proxies.  Therefore
the standard scalar Bernstein estimate for independent sub-exponential
variables gives, for every \(a\in\R^{n_m}\),
\[
\begin{aligned}
  &\mathbb P\left\{
  \left|\left\langle u,\sum_i a_iX_{m,h,i}v\right\rangle\right|
  >C\rho_{m,h}
  \left(
  \frac{\|a\|_2s}{\sqrt{d_{m,1}d_{m,2}}}
  +\frac{\|a\|_\infty s^2}{\sqrt{d_{m,1}d_{m,2}}}
  \right)
  \,\middle|\,\mathcal H_m
  \right\}  \\
  &\qquad\le2e^{-s^2}.
\end{aligned}
\]
The same proof applies to real Haar vectors after decomposing complex
bilinear forms into real and imaginary parts and increasing the numerical
constant.  Thus
Corollary~\ref{cor:active-menu-bernstein-bilinear-rectangular} applies with
\[
  q_{m,h}=C\rho_{m,h}(d_{m,1}d_{m,2})^{-1/2},
  \qquad
  b_{m,h}=C\rho_{m,h}(d_{m,1}d_{m,2})^{-1/2}.
\]
The displayed dyad scale is the resulting Bernstein bilinear scale.  If
\(d_{m,1}\asymp d_{m,2}\) and \(\tau_m=O(\sqrt{D_m^{\rm rec}})\), then
\((D_m^{\rm rec})/\sqrt{d_{m,1}d_{m,2}}\) is bounded, giving the final
constant-normalization claim.
\end{proof}

\begin{corollary}[Bounded matrix-Bernstein all-plus active menus]
\label{cor:active-menu-bounded-bernstein-all-plus}
For each \(m\ge1\), let \(D_m\ge2\), \(n_m\ge1\), and \(J_m\ge1\).  Let
\(\mathcal H_m\) be an environment.  For each active \(1\le h\le J_m\), let
\[
  X_{m,h,1},\ldots,X_{m,h,n_m}\in\M_{D_m}(\C)
\]
be conditionally independent centered self-adjoint random matrices.  Let
\(E_{m,h}\in\mathcal H_m\) be an exceptional event and let
\(s_{m,h},v_{m,h},u_{m,h}>0\) be deterministic.  Assume that, on
\(E_{m,h}^c\),
\[
  \|X_{m,h,i}\|_{\op}\le s_{m,h}
  \qquad(1\le i\le n_m)
\]
almost surely and
\[
  \left\|\sum_{i=1}^{n_m}
  \mathbb E(X_{m,h,i}^2\mid\mathcal H_m)\right\|_{\op}
  \le v_{m,h}.
\]
Set
\[
  S_{m,h}:=\max\{s_{m,h},\,u_{m,h}/\sqrt{n_m}\}.
\]
If
\[
  \sum_m\sum_{h=1}^{J_m}
  \left[
  \mathbb P(E_{m,h})
  +2D_m\exp\left\{
  -\frac{u_{m,h}^2}{2(v_{m,h}+s_{m,h}u_{m,h}/3)}
  \right\}
  \right]<\infty,
\]
then, almost surely, for all sufficiently large \(m\), all active matrices
\(X_{m,h,i}/S_{m,h}\) are contractions and the all-plus signing gives
\[
  \disc(X_{m,h,1}/S_{m,h},\ldots,X_{m,h,n_m}/S_{m,h})
  \le\sqrt{n_m}
  \qquad(1\le h\le J_m).
\]
No coordinatewise sign symmetry is required.  In particular, if
\(v_{m,h}\le \beta s_{m,h}^2n_m/\log(eD_m)\) and
\(u_{m,h}=Ks_{m,h}\sqrt{n_m}\), this gives a constant multiple of the
\(s_{m,h}\)-normalization whenever the displayed Bernstein failure series is
summable.  This covers centered bounded random projections, bounded sparse
adjacency increments, bounded random block matrices, and centered bounded
self-adjoint measurement libraries under their all-plus variance budgets.
\end{corollary}

\begin{proof}
Fix \(m,h\).  On \(E_{m,h}^c\), the two-sided matrix Bernstein inequality
\cite{Tropp} gives
\[
  \mathbb P\left\{
  \left\|\sum_iX_{m,h,i}\right\|_{\op}>u_{m,h}
  \,\middle|\,\mathcal H_m
  \right\}
  \le
  2D_m\exp\left\{
  -\frac{u_{m,h}^2}{2(v_{m,h}+s_{m,h}u_{m,h}/3)}
  \right\}.
\]
Therefore the unconditional probability that either \(E_{m,h}\) occurs or the
all-plus sum exceeds \(u_{m,h}\) is bounded by the displayed summand.  The first
Borel--Cantelli lemma gives, almost surely, eventual simultaneous control for
all active \(h\).  On that event, \(\|X_{m,h,i}/S_{m,h}\|_{\op}\le1\) and
\[
  \left\|\sum_iX_{m,h,i}/S_{m,h}\right\|_{\op}
  \le u_{m,h}/S_{m,h}\le\sqrt{n_m},
\]
so the all-plus signing is an admissible Matrix-Spencer certificate.
\end{proof}

\begin{corollary}[Rectangular bounded matrix-Bernstein all-plus active menus]
\label{cor:active-menu-rectangular-bounded-bernstein-all-plus}
For each \(m\ge1\), let \(d_{m,1},d_{m,2}\ge1\),
\(D_m^{\rm rec}=d_{m,1}+d_{m,2}\), \(n_m\ge1\), and \(J_m\ge1\).  Let
\(\mathcal H_m\) be an environment.  For each active \(1\le h\le J_m\), let
\[
  X_{m,h,1},\ldots,X_{m,h,n_m}\in
  \M_{d_{m,1}\times d_{m,2}}(\C)
\]
be conditionally independent centered rectangular random matrices.  Let
\(E_{m,h}\in\mathcal H_m\) and let
\(s_{m,h},v_{m,h},u_{m,h}>0\) be deterministic.  Assume that, on
\(E_{m,h}^c\),
\[
  \|X_{m,h,i}\|_{\op}\le s_{m,h}\qquad(1\le i\le n_m)
\]
almost surely and
\[
\begin{aligned}
  \left\|\sum_i\mathbb E(X_{m,h,i}X_{m,h,i}^*\mid\mathcal H_m)\right\|_{\op}
  &\le v_{m,h},\\
  \left\|\sum_i\mathbb E(X_{m,h,i}^*X_{m,h,i}\mid\mathcal H_m)\right\|_{\op}
  &\le v_{m,h}.
\end{aligned}
\]
Set \(S_{m,h}:=\max\{s_{m,h},u_{m,h}/\sqrt{n_m}\}\).  If
\[
  \sum_m\sum_{h=1}^{J_m}
  \left[
  \mathbb P(E_{m,h})
  +2D_m^{\rm rec}\exp\left\{
  -\frac{u_{m,h}^2}{2(v_{m,h}+s_{m,h}u_{m,h}/3)}
  \right\}
  \right]<\infty,
\]
then, almost surely, for all sufficiently large \(m\), all active rectangular
matrices \(X_{m,h,i}/S_{m,h}\) are contractions and the all-plus signing gives
\[
  \disc(X_{m,h,1}/S_{m,h},\ldots,X_{m,h,n_m}/S_{m,h})
  \le\sqrt{n_m}
  \qquad(1\le h\le J_m).
\]
Thus centered bounded rectangular blocks, sparse rectangular increments,
centered bounded random feature blocks, and bounded measurement matrices are
Matrix-Spencer scale whenever their row and column all-plus variance budgets
make the displayed Bernstein series summable.
\end{corollary}

\begin{proof}
Apply Corollary~\ref{cor:active-menu-bounded-bernstein-all-plus} to the
self-adjoint dilations
\[
  \mathscr D(X_{m,h,i})=
  \begin{pmatrix}
  0&X_{m,h,i}\\
  X_{m,h,i}^*&0
  \end{pmatrix}
  \in\M_{D_m^{\rm rec}}(\C).
\]
The dilation has operator norm \(\|X_{m,h,i}\|_{\op}\), and its variance is
block diagonal with diagonal blocks
\(\sum_i\mathbb E(X_{m,h,i}X_{m,h,i}^*\mid\mathcal H_m)\) and
\(\sum_i\mathbb E(X_{m,h,i}^*X_{m,h,i}\mid\mathcal H_m)\).  Hence the
self-adjoint corollary applies with dimension \(D_m^{\rm rec}\), jump
\(s_{m,h}\), and variance bound \(v_{m,h}\).  Finally,
\(\|\sum_i\mathscr D(X_{m,h,i})\|_{\op}=\|\sum_iX_{m,h,i}\|_{\op}\), so the
dilation certificate is exactly the rectangular all-plus certificate.
\end{proof}

\begin{corollary}[Finite-confidence bounded matrix-Bernstein all-plus certificates]
\label{cor:bounded-matrix-bernstein-finite-confidence}
Let \(\mcF\) be side information, let \(\mathcal A\subset\mathbb N\) be an
almost surely finite \(\mcF\)-measurable active set, and fix
\(n\ge1\) and \(0<\delta<1\).
For each \(h\), let
\[
  X_{h,1},\ldots,X_{h,n}\in\M_D(\C)
\]
be centered self-adjoint random matrices.  Assume that, conditionally on
\(\mcF\), every selected stream is independent and has \(\mcF\)-measurable
proxies \(s_h,v_h,u_h>0\) satisfying
\[
  \|X_{h,i}\|_{\op}\le s_h\qquad(1\le i\le n)
\]
almost surely and
\[
  \left\|\sum_{i=1}^n\mathbb E(X_{h,i}^2\mid\mcF)\right\|_{\op}\le v_h .
\]
Set
\[
  S_h:=\max\{s_h,u_h/\sqrt n\}.
\]
If
\[
  \mathbb E\sum_{h\in\mathcal A}
  2D\exp\left\{-\frac{u_h^2}{2(v_h+s_hu_h/3)}\right\}\le\delta,
\]
then, with probability at least \(1-\delta\), every selected stream is
contraction-normalized and all-plus certified:
\[
  \disc(X_{h,1}/S_h,\ldots,X_{h,n}/S_h)\le\sqrt n
  \qquad(h\in\mathcal A).
\]

The same statement holds for rectangular
\(X_{h,i}\in\M_{d_1\times d_2}(\C)\), with \(D^{\rm rec}=d_1+d_2\) replacing
\(D\), provided
\[
  \|X_{h,i}\|_{\op}\le s_h
\]
and
\[
\begin{aligned}
  \left\|\sum_i\mathbb E(X_{h,i}X_{h,i}^*\mid\mcF)\right\|_{\op}
  &\le v_h,\\
  \left\|\sum_i\mathbb E(X_{h,i}^*X_{h,i}\mid\mcF)\right\|_{\op}
  &\le v_h .
\end{aligned}
\]
In particular, if \(\abs{\mathcal A}\le M\) almost surely and
\[
  L=\log\frac{2M\Delta}{\delta},\qquad
  \Delta=D\ \text{or}\ D^{\rm rec},
\]
then the homogeneous choice
\[
  u=C\bigl(\sqrt{vL}+sL\bigr)
\]
gives a simultaneous \(1-\delta\) certificate after normalizing by
\[
  C\max\left\{s,\sqrt{\frac{vL}{n}}+\frac{sL}{\sqrt n}\right\},
\]
whenever \(s_h\le s\) and \(v_h\le v\) for all selected \(h\).
\end{corollary}

\begin{proof}
Condition on \(\mcF\).  For a fixed selected self-adjoint stream, the two-sided
matrix Bernstein inequality \cite{Tropp} gives
\[
  \mathbb P\left\{
  \left\|\sum_iX_{h,i}\right\|_{\op}>u_h
  \,\middle|\,\mcF
  \right\}
  \le
  2D\exp\left\{-\frac{u_h^2}{2(v_h+s_hu_h/3)}\right\}.
\]
Summing over the \(\mcF\)-measurable active set and taking expectation gives
total failure at most \(\delta\).  On the complementary event,
\(\|X_{h,i}/S_h\|_{\op}\le1\) and
\[
  \left\|\sum_iX_{h,i}/S_h\right\|_{\op}\le u_h/S_h\le\sqrt n,
\]
so the all-plus signing is a Matrix-Spencer certificate.

The rectangular assertion follows by applying the self-adjoint statement to the
dilations \(\mathscr D(X_{h,i})\).  The dilation has dimension \(D^{\rm rec}\),
operator norm \(\|X_{h,i}\|_{\op}\), and variance bounded by the two displayed
row and column variance budgets.  Finally,
\(\|\sum_i\mathscr D(X_{h,i})\|_{\op}=\|\sum_iX_{h,i}\|_{\op}\).

For the homogeneous finite-menu form, the usual Bernstein algebra gives
\[
  \frac{u^2}{2(v+su/3)}\ge L
\]
after increasing the numerical constant in
\(u=C(\sqrt{vL}+sL)\).  The displayed scale is
\(\max\{s,u/\sqrt n\}\), up to the same absolute constant.
\end{proof}

\begin{corollary}[Finite-confidence quantum effect fluctuations]
\label{cor:quantum-effect-bernstein-finite-confidence}
Let \(\mcF\) be side information, let \(\mathcal A\subset\mathbb N\) be an
almost surely finite \(\mcF\)-measurable active set, and fix
\(n,D\ge1\) and \(0<\delta<1\).  For each \(h\), let
\[
  E_{h,1},\ldots,E_{h,n}\in\M_D(\C)
\]
be conditionally independent quantum effects:
\[
  0\preceq E_{h,i}\preceq I_D\qquad(1\le i\le n).
\]
Put
\[
  \mu_{h,i}:=\mathbb E(E_{h,i}\mid\mcF),
  \qquad
  X_{h,i}:=E_{h,i}-\mu_{h,i},
\]
and suppose that \(v_h,u_h>0\) are \(\mcF\)-measurable proxies satisfying
\[
  \left\|\sum_{i=1}^n\mu_{h,i}\right\|_{\op}\le v_h
  \qquad(h\in\mathcal A).
\]
Set
\[
  S_h:=\max\{1,u_h/\sqrt n\}.
\]
If
\[
  \mathbb E\sum_{h\in\mathcal A}
  2D\exp\left\{-\frac{u_h^2}{2(v_h+u_h/3)}\right\}\le\delta,
\]
then, with probability at least \(1-\delta\), every selected centered
quantum-effect stream is contraction-normalized and all-plus certified:
\[
  \disc(X_{h,1}/S_h,\ldots,X_{h,n}/S_h)\le\sqrt n
  \qquad(h\in\mathcal A).
\]
In particular, if \(\abs{\mathcal A}\le M\) almost surely and
\[
  L=\log\frac{2MD}{\delta},
  \qquad
  v_h\le v\quad(h\in\mathcal A),
\]
then the homogeneous normalization
\[
  S=C\max\left\{1,\sqrt{\frac{vL}{n}}+\frac{L}{\sqrt n}\right\}
\]
gives the same simultaneous \(1-\delta\) certificate.  If
\(C(\sqrt{vL}+L)\le\sqrt n\), the unscaled centered effects \(X_{h,i}\) are
already all-plus Matrix-Spencer certified.

For random density matrices with conditionally maximally mixed means
\(\mu_{h,i}=I_D/D\), one may take \(v=n/D\).  Hence, for a finite active menu,
the unscaled centered density-matrix fluctuations are all-plus certified as
soon as \(D\gtrsim L\) and \(n\gtrsim L^2\).
\end{corollary}

\begin{proof}
Condition on \(\mcF\).  Since \(0\preceq E_{h,i}\preceq I_D\) and
\(0\preceq\mu_{h,i}\preceq I_D\), we have
\[
  -I_D\preceq -\mu_{h,i}\preceq X_{h,i}\preceq E_{h,i}\preceq I_D,
\]
so \(\|X_{h,i}\|_{\op}\le1\).  Also,
\[
  \mathbb E(X_{h,i}^2\mid\mcF)
  =
  \mathbb E(E_{h,i}^2\mid\mcF)-\mu_{h,i}^2
  \preceq
  \mathbb E(E_{h,i}^2\mid\mcF)
  \preceq
  \mu_{h,i},
\]
because \(E_{h,i}^2\preceq E_{h,i}\).  Thus
\[
  \left\|\sum_i\mathbb E(X_{h,i}^2\mid\mcF)\right\|_{\op}
  \le
  \left\|\sum_i\mu_{h,i}\right\|_{\op}
  \le v_h.
\]
Corollary~\ref{cor:bounded-matrix-bernstein-finite-confidence} applies with
\(s_h=1\).  The homogeneous scale and the unscaled criterion are the displayed
specializations of \(u=C(\sqrt{vL}+L)\).  If \(\mu_{h,i}=I_D/D\), then
\(\|\sum_i\mu_{h,i}\|_{\op}=n/D\).
\end{proof}

\begin{corollary}[Bounded matrix-Freedman all-plus active menus]
\label{cor:active-menu-bounded-freedman-all-plus}
For each \(m\ge1\), let \(D_m,n_m,J_m\ge1\).  For each active
\(1\le h\le J_m\), let
\[
  (X_{m,h,i},\mcF_{m,h,i})_{i=1}^{n_m}
\]
be an adapted self-adjoint matrix martingale-difference sequence in
\(\M_{D_m}(\C)\).  Let \(E_{m,h}\) be an event, and let
\(s_{m,h},v_{m,h},u_{m,h}>0\) be deterministic.  Assume that
\[
  \|X_{m,h,i}\|_{\op}\le s_{m,h}\qquad(1\le i\le n_m)
\]
almost surely and that, on \(E_{m,h}^c\),
\[
  \left\|\sum_{i=1}^{n_m}
  \mathbb E(X_{m,h,i}^2\mid\mcF_{m,h,i-1})\right\|_{\op}
  \le v_{m,h}.
\]
Set \(S_{m,h}:=\max\{s_{m,h},u_{m,h}/\sqrt{n_m}\}\).  If
\[
  \sum_m\sum_{h=1}^{J_m}
  \left[
  \mathbb P(E_{m,h})
  +2D_m\exp\left\{
  -\frac{u_{m,h}^2}{2(v_{m,h}+s_{m,h}u_{m,h}/3)}
  \right\}
  \right]<\infty,
\]
then, almost surely, for all sufficiently large \(m\), all active matrices
\(X_{m,h,i}/S_{m,h}\) are contractions and the all-plus signing gives
\[
  \disc(X_{m,h,1}/S_{m,h},\ldots,X_{m,h,n_m}/S_{m,h})
  \le\sqrt{n_m}
  \qquad(1\le h\le J_m).
\]
Thus the bounded Bernstein all-plus route remains valid for online or adapted
centered self-adjoint streams after replacing the ordinary variance by the
predictable quadratic variation.
\end{corollary}

\begin{proof}
Fix \(m,h\) and put
\[
  \nu_{m,h}:=
  \left\|\sum_i
  \mathbb E(X_{m,h,i}^2\mid\mcF_{m,h,i-1})\right\|_{\op}.
\]
The two-sided matrix Freedman inequality \cite{Tropp}, applied to the
martingale differences \(X_{m,h,i}\) and \(-X_{m,h,i}\), gives
\[
\begin{aligned}
  &\mathbb P\left\{
  \left\|\sum_iX_{m,h,i}\right\|_{\op}>u_{m,h},
  \ \nu_{m,h}\le v_{m,h}
  \right\} \\
  &\qquad\le
  2D_m\exp\left\{
  -\frac{u_{m,h}^2}{2(v_{m,h}+s_{m,h}u_{m,h}/3)}
  \right\}.
\end{aligned}
\]
On \(E_{m,h}^c\) the event \(\nu_{m,h}\le v_{m,h}\) holds, so the displayed
series controls all active all-plus failures together with the exceptional
events.  Borel--Cantelli and the same normalization calculation as in
Corollary~\ref{cor:active-menu-bounded-bernstein-all-plus} finish the proof.
\end{proof}

\begin{corollary}[Rectangular bounded matrix-Freedman all-plus active menus]
\label{cor:active-menu-rectangular-bounded-freedman-all-plus}
For each \(m\ge1\), let \(d_{m,1},d_{m,2}\ge1\),
\(D_m^{\rm rec}=d_{m,1}+d_{m,2}\), \(n_m\ge1\), and \(J_m\ge1\).  For each
active \(1\le h\le J_m\), let
\[
  (X_{m,h,i},\mcF_{m,h,i})_{i=1}^{n_m}
\]
be adapted rectangular martingale differences with values in
\[
  \M_{d_{m,1}\times d_{m,2}}(\C).
\]
Let \(E_{m,h}\) be an event, and let \(s_{m,h},v_{m,h},u_{m,h}>0\) be
deterministic.  Assume that
\[
  \|X_{m,h,i}\|_{\op}\le s_{m,h}\qquad(1\le i\le n_m)
\]
almost surely and that, on \(E_{m,h}^c\),
\[
\begin{aligned}
  \left\|\sum_i
  \mathbb E(X_{m,h,i}X_{m,h,i}^*\mid\mcF_{m,h,i-1})\right\|_{\op}
  &\le v_{m,h},\\
  \left\|\sum_i
  \mathbb E(X_{m,h,i}^*X_{m,h,i}\mid\mcF_{m,h,i-1})\right\|_{\op}
  &\le v_{m,h}.
\end{aligned}
\]
Set \(S_{m,h}:=\max\{s_{m,h},u_{m,h}/\sqrt{n_m}\}\).  If
\[
  \sum_m\sum_{h=1}^{J_m}
  \left[
  \mathbb P(E_{m,h})
  +2D_m^{\rm rec}\exp\left\{
  -\frac{u_{m,h}^2}{2(v_{m,h}+s_{m,h}u_{m,h}/3)}
  \right\}
  \right]<\infty,
\]
then, almost surely, for all sufficiently large \(m\), all active rectangular
matrices \(X_{m,h,i}/S_{m,h}\) are contractions and the all-plus signing gives
\[
  \disc(X_{m,h,1}/S_{m,h},\ldots,X_{m,h,n_m}/S_{m,h})
  \le\sqrt{n_m}
  \qquad(1\le h\le J_m).
\]
This covers online rectangular feature blocks, adaptive measurement increments,
and other non-self-adjoint martingale-difference streams under row and column
predictable variance budgets.
\end{corollary}

\begin{proof}
Apply Corollary~\ref{cor:active-menu-bounded-freedman-all-plus} to the
self-adjoint dilations \(\mathscr D(X_{m,h,i})\).  The dilation is a
self-adjoint martingale difference, has norm \(\|X_{m,h,i}\|_{\op}\), and has
predictable quadratic variation block diagonal with the two displayed row and
column variance sums.  The operator norm of the dilated all-plus sum is the
operator norm of the rectangular all-plus sum, so the self-adjoint conclusion is
exactly the desired rectangular certificate.
\end{proof}

\begin{corollary}[Finite-confidence bounded matrix-Freedman all-plus certificates]
\label{cor:bounded-matrix-freedman-finite-confidence}
Let \(\mcG\) be side information, let \(\mathcal A\subset\mathbb N\) be an
almost surely finite \(\mcG\)-measurable active set, and fix
\(n\ge1\) and \(0<\delta<1\).
For each \(h\), let
\[
  (\mcF_{h,i})_{i=0}^n
\]
be a filtration with \(\mcG\subset\mcF_{h,0}\), and let
\(X_{h,1},\ldots,X_{h,n}\in\M_D(\C)\) be a self-adjoint matrix
martingale-difference sequence adapted to it.  Assume that every selected
stream has \(\mcG\)-measurable proxies \(s_h,v_h,u_h>0\) satisfying
\[
  \|X_{h,i}\|_{\op}\le s_h\qquad(1\le i\le n)
\]
almost surely and
\[
  \left\|\sum_{i=1}^n
  \mathbb E(X_{h,i}^2\mid\mcF_{h,i-1})\right\|_{\op}\le v_h
\]
almost surely.  Set
\[
  S_h:=\max\{s_h,u_h/\sqrt n\}.
\]
If
\[
  \mathbb E\sum_{h\in\mathcal A}
  2D\exp\left\{-\frac{u_h^2}{2(v_h+s_hu_h/3)}\right\}\le\delta,
\]
then, with probability at least \(1-\delta\), every selected stream is
contraction-normalized and all-plus certified:
\[
  \disc(X_{h,1}/S_h,\ldots,X_{h,n}/S_h)\le\sqrt n
  \qquad(h\in\mathcal A).
\]

The same statement holds for rectangular martingale differences
\(X_{h,i}\in\M_{d_1\times d_2}(\C)\), with
\(D^{\rm rec}=d_1+d_2\) replacing \(D\), provided
\[
  \|X_{h,i}\|_{\op}\le s_h
\]
and
\[
\begin{aligned}
  \left\|\sum_i
  \mathbb E(X_{h,i}X_{h,i}^*\mid\mcF_{h,i-1})\right\|_{\op}
  &\le v_h,\\
  \left\|\sum_i
  \mathbb E(X_{h,i}^*X_{h,i}\mid\mcF_{h,i-1})\right\|_{\op}
  &\le v_h
\end{aligned}
\]
almost surely.
In particular, if \(\abs{\mathcal A}\le M\) almost surely and
\[
  L=\log\frac{2M\Delta}{\delta},\qquad
  \Delta=D\ \text{or}\ D^{\rm rec},
\]
then the homogeneous choice
\[
  u=C\bigl(\sqrt{vL}+sL\bigr)
\]
gives a simultaneous \(1-\delta\) certificate after normalizing by
\[
  C\max\left\{s,\sqrt{\frac{vL}{n}}+\frac{sL}{\sqrt n}\right\},
\]
whenever \(s_h\le s\) and \(v_h\le v\) for all selected \(h\).
\end{corollary}

\begin{proof}
Condition on \(\mcG\).  For a fixed selected self-adjoint stream, put
\[
  \nu_h:=
  \left\|\sum_i
  \mathbb E(X_{h,i}^2\mid\mcF_{h,i-1})\right\|_{\op}.
\]
The two-sided matrix Freedman inequality \cite{Tropp}, applied conditionally
on the initial side information, gives
\[
\begin{aligned}
  &\mathbb P\left\{
  \left\|\sum_iX_{h,i}\right\|_{\op}>u_h,\ \nu_h\le v_h
  \,\middle|\,\mcG
  \right\} \\
  &\qquad\le
  2D\exp\left\{-\frac{u_h^2}{2(v_h+s_hu_h/3)}\right\}.
\end{aligned}
\]
The predictable variance bound removes the event \(\nu_h\le v_h\).  Summing
over the \(\mcG\)-measurable active set and taking expectation gives total
failure at most \(\delta\).  On the complementary event,
\(\|X_{h,i}/S_h\|_{\op}\le1\) and
\[
  \left\|\sum_iX_{h,i}/S_h\right\|_{\op}\le u_h/S_h\le\sqrt n,
\]
so the all-plus signing is a Matrix-Spencer certificate.

The rectangular assertion follows by applying the self-adjoint statement to
the dilations \(\mathscr D(X_{h,i})\).  The dilation is a self-adjoint
martingale difference, has dimension \(D^{\rm rec}\), has operator norm
\(\|X_{h,i}\|_{\op}\), and has predictable quadratic variation bounded by the
two displayed row and column variance budgets.  The all-plus norm is preserved
by dilation.

The homogeneous finite-menu scale is the same Bernstein algebra used in
Corollary~\ref{cor:bounded-matrix-bernstein-finite-confidence}.
\end{proof}

\begin{corollary}[Explicit bounded Bernstein/Freedman all-plus schedules]
\label{cor:active-menu-bounded-bernstein-all-plus-schedules}
Consider any one of the four bounded all-plus settings in
Corollaries~\ref{cor:active-menu-bounded-bernstein-all-plus},
\ref{cor:active-menu-rectangular-bounded-bernstein-all-plus},
\ref{cor:active-menu-bounded-freedman-all-plus}, and
\ref{cor:active-menu-rectangular-bounded-freedman-all-plus}.  Let
\(\Delta_m=D_m\) in the self-adjoint case and
\(\Delta_m=D_m^{\rm rec}\) in the rectangular case, and put
\(\ell_m=\log(e\Delta_m)\).  Fix \(K\ge1\), \(\beta\ge0\), and \(B>0\).
Assume that, outside the exceptional events \(E_{m,h}\),
\[
  v_{m,h}\le \beta s_{m,h}^2\,\frac{n_m}{\ell_m}
  \qquad(1\le h\le J_m),
\]
that
\[
  \sum_m\sum_{h=1}^{J_m}\mathbb P(E_{m,h})<\infty,
\]
and that \(\ell_m\le B\sqrt{n_m}\) eventually.  Define
\[
  \gamma(K,\beta,B):=
  \frac{K^2}{2(\beta+KB/3)} .
\]
If
\[
  \sum_m J_m\Delta_m(e\Delta_m)^{-\gamma(K,\beta,B)}<\infty,
\]
then, almost surely, for all sufficiently large \(m\), every active
self-adjoint or rectangular bounded Bernstein/Freedman stream is
contraction-normalized by \(Ks_{m,h}\), and the all-plus signing gives
\[
  \disc(X_{m,h,1}/(Ks_{m,h}),\ldots,X_{m,h,n_m}/(Ks_{m,h}))
  \le\sqrt{n_m}
  \qquad(1\le h\le J_m).
\]
In particular, if \(J_m\Delta_m\le(e\Delta_m)^\Gamma\) eventually and
\(n_m\ge a\ell_m^2\), then the conclusion holds along all integer dimensions
\(\Delta_m=m\) whenever
\[
  \gamma(K,\beta,a^{-1/2})>\Gamma+1,
\]
and along every stretched-exponential sparse dimension schedule whenever
\[
  \gamma(K,\beta,a^{-1/2})>\Gamma.
\]
\end{corollary}

\begin{proof}
Use the appropriate preceding self-adjoint or rectangular bounded all-plus
corollary with
\(u_{m,h}=Ks_{m,h}\sqrt{n_m}\).  Since \(K\ge1\), its normalizing scale is
\[
  S_{m,h}=\max\{s_{m,h},u_{m,h}/\sqrt{n_m}\}=Ks_{m,h}.
\]
On \(E_{m,h}^c\), the common Bernstein/Freedman exponent is at least
\[
\begin{aligned}
  \frac{u_{m,h}^2}{2(v_{m,h}+s_{m,h}u_{m,h}/3)}
  &\ge
  \frac{K^2n_m}{2(\beta n_m/\ell_m+K\sqrt{n_m}/3)}  \\
  &\ge \gamma(K,\beta,B)\ell_m,
\end{aligned}
\]
because \(\ell_m\le B\sqrt{n_m}\).  Hence the active-row failure probability is
bounded by a constant multiple of
\[
  J_m\Delta_m(e\Delta_m)^{-\gamma(K,\beta,B)}
\]
plus the exceptional-event probabilities.  The displayed summability and
Borel--Cantelli give the general schedule.  If
\(J_m\Delta_m\le(e\Delta_m)^\Gamma\), the remaining summand is bounded by a
constant multiple of \((e\Delta_m)^{\Gamma-\gamma}\).  This is summable over all
integer dimensions for \(\gamma>\Gamma+1\), and is summable along every
stretched-exponential sparse schedule for \(\gamma>\Gamma\).
\end{proof}

\begin{corollary}[Proxy-scale bounded all-plus schedules]
\label{cor:active-menu-bounded-proxy-scale-all-plus-schedules}
Consider any one of the four bounded all-plus settings in
Corollaries~\ref{cor:active-menu-bounded-bernstein-all-plus},
\ref{cor:active-menu-rectangular-bounded-bernstein-all-plus},
\ref{cor:active-menu-bounded-freedman-all-plus}, and
\ref{cor:active-menu-rectangular-bounded-freedman-all-plus}.  Let
\(\Delta_m=D_m\) in the self-adjoint case and
\(\Delta_m=D_m^{\rm rec}\) in the rectangular case, and put
\(\ell_m=\log(e\Delta_m)\).  Fix \(K>0\), \(R>0\), \(\beta\ge0\), and
\(B>0\).  For each active stream let \(q_{m,h}>0\) be a deterministic proxy
scale and let \(b_{m,h}\) be a deterministic jump bound, so that all increments
have operator norm at most \(b_{m,h}\).  Assume that, outside \(E_{m,h}\),
\[
  b_{m,h}\le Rq_{m,h},
  \qquad
  v_{m,h}\le \beta q_{m,h}^2\,\frac{n_m}{\ell_m}
  \qquad(1\le h\le J_m),
\]
that
\[
  \sum_m\sum_{h=1}^{J_m}\mathbb P(E_{m,h})<\infty,
\]
and that \(\ell_m\le B\sqrt{n_m}\) eventually.  Put
\[
  \gamma_{\rm pr}(K,R,\beta,B):=
  \frac{K^2}{2(\beta+KRB/3)} .
\]
If
\[
  \sum_m J_m\Delta_m(e\Delta_m)^{-\gamma_{\rm pr}(K,R,\beta,B)}<\infty,
\]
then, almost surely, for all sufficiently large \(m\), every active
self-adjoint or rectangular bounded Bernstein/Freedman stream is
contraction-normalized by
\[
  Mq_{m,h},\qquad M:=\max\{R,K\},
\]
and the all-plus signing gives
\[
  \disc(X_{m,h,1}/(Mq_{m,h}),\ldots,X_{m,h,n_m}/(Mq_{m,h}))
  \le\sqrt{n_m}
  \qquad(1\le h\le J_m).
\]
In particular, if \(J_m\Delta_m\le(e\Delta_m)^\Gamma\) eventually and
\(n_m\ge a\ell_m^2\), then the conclusion holds along all integer dimensions
whenever
\[
  \gamma_{\rm pr}(K,R,\beta,a^{-1/2})>\Gamma+1,
\]
and along every stretched-exponential sparse dimension schedule whenever
\[
  \gamma_{\rm pr}(K,R,\beta,a^{-1/2})>\Gamma.
\]
This lets heterogeneous bounded libraries use a typical variance scale
\(q_{m,h}\) rather than the worst increment scale, provided the worst increment
is within a fixed factor \(R\) of \(q_{m,h}\).
\end{corollary}

\begin{proof}
Apply the appropriate bounded Bernstein or bounded Freedman all-plus corollary
with common jump \(s_{m,h}=b_{m,h}\) and
\[
  u_{m,h}=Kq_{m,h}\sqrt{n_m}.
\]
The normalizing scale in that corollary is at most
\[
  \max\{b_{m,h},u_{m,h}/\sqrt{n_m}\}
  \le \max\{R,K\}q_{m,h}=Mq_{m,h}.
\]
Using the larger scale \(Mq_{m,h}\) preserves contraction and only reduces the
all-plus norm.  On \(E_{m,h}^c\), the common Bernstein/Freedman exponent is at
least
\[
\begin{aligned}
  \frac{u_{m,h}^2}{2(v_{m,h}+b_{m,h}u_{m,h}/3)}
  &\ge
  \frac{K^2n_m}{2(\beta n_m/\ell_m+KR\sqrt{n_m}/3)}\\
  &\ge \gamma_{\rm pr}(K,R,\beta,B)\ell_m .
\end{aligned}
\]
The same Borel--Cantelli and summability argument as in
Corollary~\ref{cor:active-menu-bounded-bernstein-all-plus-schedules} gives the
active-menu, all-integer, and sparse conclusions.
\end{proof}

\begin{corollary}[Tail-capped centered Bernstein all-plus schedules]
\label{cor:tail-capped-centered-bernstein-all-plus-schedules}
For each \(m\ge1\), consider either self-adjoint dimensions \(D_m\ge2\) or
rectangular dimensions \(d_{m,1},d_{m,2}\ge1\).  Put
\(\Delta_m=D_m\) in the self-adjoint case, put
\(\Delta_m=d_{m,1}+d_{m,2}\) in the rectangular case, and let
\(\ell_m=\log(e\Delta_m)\).  Let \(\mathcal H_m\) be an environment.  For each
active \(1\le h\le J_m\), let
\[
  X_{m,h,1},\ldots,X_{m,h,n_m}
\]
be conditionally independent random matrices, self-adjoint in
\(\M_{D_m}(\C)\) in the first case and rectangular in
\(\M_{d_{m,1}\times d_{m,2}}(\C)\) in the second.  Let
\(E_{m,h}\in\mathcal H_m\) be exceptional events.  Let \(q_{m,h}>0\) and
\(b_{m,h}>0\) be deterministic scales, and define the capped variables
\[
  T_{m,h,i}:=
  X_{m,h,i}\mathbf 1_{\{\|X_{m,h,i}\|_{\op}\le b_{m,h}\}},
  \qquad
  \mu_{m,h,i}:=\mathbb E(T_{m,h,i}\mid\mathcal H_m),
\]
\[
  Y_{m,h,i}:=T_{m,h,i}-\mu_{m,h,i},
  \qquad
  F^{\rm cap}_{m,h}:=
  \left\{\max_{1\le i\le n_m}\|X_{m,h,i}\|_{\op}>b_{m,h}\right\}.
\]
Fix constants \(K>0\), \(R>0\), \(\beta\ge0\), \(\eta\ge0\), and \(B>0\).
Assume that \(b_{m,h}\le Rq_{m,h}\), that
\[
  \sum_m\sum_{h=1}^{J_m}
  \bigl[\mathbb P(E_{m,h})+\mathbb P(F^{\rm cap}_{m,h})\bigr]<\infty,
\]
and that, outside \(E_{m,h}\),
\[
  \left\|\sum_{i=1}^{n_m}\mu_{m,h,i}\right\|_{\op}
  \le \eta q_{m,h}\sqrt{n_m}.
\]
In the self-adjoint case assume also, outside \(E_{m,h}\), that
\[
  \left\|\sum_{i=1}^{n_m}
  \mathbb E(Y_{m,h,i}^2\mid\mathcal H_m)\right\|_{\op}
  \le \beta q_{m,h}^2\,\frac{n_m}{\ell_m}.
\]
In the rectangular case replace this by the two row and column bounds
\[
\begin{aligned}
  \left\|\sum_i
  \mathbb E(Y_{m,h,i}Y_{m,h,i}^*\mid\mathcal H_m)\right\|_{\op}
  &\le \beta q_{m,h}^2\,\frac{n_m}{\ell_m},\\
  \left\|\sum_i
  \mathbb E(Y_{m,h,i}^*Y_{m,h,i}\mid\mathcal H_m)\right\|_{\op}
  &\le \beta q_{m,h}^2\,\frac{n_m}{\ell_m}.
\end{aligned}
\]
Suppose that \(\ell_m\le B\sqrt{n_m}\) eventually and put
\[
  \gamma_{\rm tc}(K,R,\beta,B):=
  \frac{K^2}{2(\beta+2KRB/3)} .
\]
If
\[
  \sum_mJ_m\Delta_m(e\Delta_m)^{-\gamma_{\rm tc}(K,R,\beta,B)}<\infty,
\]
then, almost surely, for all sufficiently large \(m\), every active original
stream is contraction-normalized by
\[
  Mq_{m,h},
  \qquad
  M:=\max\{R,K+\eta\},
\]
and the all-plus signing gives
\[
  \disc(X_{m,h,1}/(Mq_{m,h}),\ldots,X_{m,h,n_m}/(Mq_{m,h}))
  \le\sqrt{n_m}
  \qquad(1\le h\le J_m).
\]

For \(A\ge0\) define the closed-form margin
\[
  \Theta_{\rm tc}(A,a,\beta,R):=
  \frac{2AR}{3\sqrt a}
  +
  \sqrt{\frac{4A^2R^2}{9a}+2A\beta}.
\]
If \(J_m\Delta_m\le(e\Delta_m)^\Gamma\) eventually and
\(n_m\ge a\ell_m^2\), then the conclusion holds along all integer dimensions
\(\Delta_m=m\) whenever
\[
  K>\Theta_{\rm tc}(\Gamma+1,a,\beta,R),
\]
and along every stretched-exponential sparse dimension schedule whenever
\[
  K>\Theta_{\rm tc}(\Gamma,a,\beta,R).
\]
When the original variables are conditionally centered, the drift condition is
only a tail-mean budget, since
\(\mu_{m,h,i}=-\mathbb E(X_{m,h,i}
\mathbf 1_{\{\|X_{m,h,i}\|_{\op}>b_{m,h}\}}\mid\mathcal H_m)\).
\end{corollary}

\begin{proof}
Fix \(m,h\) and condition on \(\mathcal H_m\).  Outside \(E_{m,h}\), the
matrices \(Y_{m,h,i}\) are conditionally independent and centered.  Moreover,
\[
  \|Y_{m,h,i}\|_{\op}
  \le \|T_{m,h,i}\|_{\op}
  +\|\mathbb E(T_{m,h,i}\mid\mathcal H_m)\|_{\op}
  \le 2b_{m,h}
  \le 2Rq_{m,h}.
\]
In the self-adjoint case, matrix Bernstein gives
\[
\begin{aligned}
  &\mathbb P\left\{
  \left\|\sum_iY_{m,h,i}\right\|_{\op}>Kq_{m,h}\sqrt{n_m}
  \,\middle|\,\mathcal H_m
  \right\}  \\
  &\qquad\le
  2\Delta_m
  \exp\left\{
  -\frac{K^2n_m}
          {2(\beta n_m/\ell_m+2RK\sqrt{n_m}/3)}
  \right\}.
\end{aligned}
\]
The rectangular case is identical after applying the same inequality to the
self-adjoint dilations; the two displayed row and column variance bounds are
exactly the dilation variance bound.  Since \(\ell_m\le B\sqrt{n_m}\), the
exponent is at least \(\gamma_{\rm tc}(K,R,\beta,B)\ell_m\).  Hence the
active centered-sum failure probabilities are bounded by a constant multiple of
\[
  J_m\Delta_m(e\Delta_m)^{-\gamma_{\rm tc}(K,R,\beta,B)}.
\]
Together with the summable \(E_{m,h}\) and cap failures,
Borel--Cantelli gives, almost surely, eventual simultaneous control over all
active \(h\).

On that eventual event, \(F^{\rm cap}_{m,h}\) does not occur, so
\(T_{m,h,i}=X_{m,h,i}\) for every \(i\).  Therefore
\[
  \left\|\sum_iX_{m,h,i}\right\|_{\op}
  \le
  \left\|\sum_iY_{m,h,i}\right\|_{\op}
  +
  \left\|\sum_i\mu_{m,h,i}\right\|_{\op}
  \le (K+\eta)q_{m,h}\sqrt{n_m},
\]
while \(\max_i\|X_{m,h,i}\|_{\op}\le b_{m,h}\le Rq_{m,h}\).  Dividing by
\(Mq_{m,h}\), where \(M=\max\{R,K+\eta\}\), gives contractions and an all-plus
sum of norm at most \(\sqrt{n_m}\).

For the schedule statements, \(n_m\ge a\ell_m^2\) lets us run the preceding
argument with \(B=a^{-1/2}\).  If
\(J_m\Delta_m\le(e\Delta_m)^\Gamma\), the remaining row summand is
\(O((e\Delta_m)^{\Gamma-\gamma_{\rm tc}})\), summable over all integer
dimensions when \(\gamma_{\rm tc}>\Gamma+1\) and along stretched-exponential
schedules when \(\gamma_{\rm tc}>\Gamma\).  Finally,
\[
  \gamma_{\rm tc}(K,R,\beta,a^{-1/2})>A
\]
is equivalent to
\[
  K^2-\frac{4AR}{3\sqrt a}K-2A\beta>0,
\]
whose positive root is \(\Theta_{\rm tc}(A,a,\beta,R)\).
\end{proof}

\begin{corollary}[Tail-capped matrix-Freedman all-plus schedules]
\label{cor:tail-capped-freedman-all-plus-schedules}
For each \(m\ge1\), consider either self-adjoint dimensions \(D_m\ge2\) or
rectangular dimensions \(d_{m,1},d_{m,2}\ge1\).  Put
\(\Delta_m=D_m\) in the self-adjoint case, put
\(\Delta_m=d_{m,1}+d_{m,2}\) in the rectangular case, and let
\(\ell_m=\log(e\Delta_m)\).  For each active \(1\le h\le J_m\), let
\[
  (X_{m,h,i},\mcF_{m,h,i})_{i=1}^{n_m}
\]
be an adapted martingale-difference sequence, self-adjoint in
\(\M_{D_m}(\C)\) in the first case and rectangular in
\(\M_{d_{m,1}\times d_{m,2}}(\C)\) in the second.  Let \(E_{m,h}\) be an event,
let \(q_{m,h}>0\) and \(b_{m,h}>0\) be deterministic scales, and define
\[
  T_{m,h,i}:=
  X_{m,h,i}\mathbf 1_{\{\|X_{m,h,i}\|_{\op}\le b_{m,h}\}},
  \qquad
  \mu_{m,h,i}:=
  \mathbb E(T_{m,h,i}\mid\mcF_{m,h,i-1}),
\]
\[
  Y_{m,h,i}:=T_{m,h,i}-\mu_{m,h,i},
  \qquad
  F^{\rm cap}_{m,h}:=
  \left\{\max_{1\le i\le n_m}\|X_{m,h,i}\|_{\op}>b_{m,h}\right\}.
\]
Fix constants \(K>0\), \(R>0\), \(\beta\ge0\), \(\eta\ge0\), and \(B>0\).
Assume that \(b_{m,h}\le Rq_{m,h}\), that
\[
  \sum_m\sum_{h=1}^{J_m}
  \bigl[\mathbb P(E_{m,h})+\mathbb P(F^{\rm cap}_{m,h})\bigr]<\infty,
\]
and that, outside \(E_{m,h}\),
\[
  \left\|\sum_{i=1}^{n_m}\mu_{m,h,i}\right\|_{\op}
  \le \eta q_{m,h}\sqrt{n_m}.
\]
In the self-adjoint case assume also, outside \(E_{m,h}\), that
\[
  \left\|\sum_{i=1}^{n_m}
  \mathbb E(Y_{m,h,i}^2\mid\mcF_{m,h,i-1})\right\|_{\op}
  \le \beta q_{m,h}^2\,\frac{n_m}{\ell_m}.
\]
In the rectangular case replace this by the two predictable row and column
bounds
\[
\begin{aligned}
  \left\|\sum_i
  \mathbb E(Y_{m,h,i}Y_{m,h,i}^*\mid\mcF_{m,h,i-1})\right\|_{\op}
  &\le \beta q_{m,h}^2\,\frac{n_m}{\ell_m},\\
  \left\|\sum_i
  \mathbb E(Y_{m,h,i}^*Y_{m,h,i}\mid\mcF_{m,h,i-1})\right\|_{\op}
  &\le \beta q_{m,h}^2\,\frac{n_m}{\ell_m}.
\end{aligned}
\]
Suppose that \(\ell_m\le B\sqrt{n_m}\) eventually.  If
\[
  \sum_mJ_m\Delta_m(e\Delta_m)^{-\gamma_{\rm tc}(K,R,\beta,B)}<\infty,
\]
where \(\gamma_{\rm tc}\) is the exponent from
Corollary~\ref{cor:tail-capped-centered-bernstein-all-plus-schedules}, then,
almost surely, for all sufficiently large \(m\), every active original stream
is contraction-normalized by
\[
  Mq_{m,h},
  \qquad
  M:=\max\{R,K+\eta\},
\]
and the all-plus signing gives
\[
  \disc(X_{m,h,1}/(Mq_{m,h}),\ldots,X_{m,h,n_m}/(Mq_{m,h}))
  \le\sqrt{n_m}
  \qquad(1\le h\le J_m).
\]

If \(J_m\Delta_m\le(e\Delta_m)^\Gamma\) eventually and
\(n_m\ge a\ell_m^2\), then the same closed-form margin
\(\Theta_{\rm tc}\) gives the all-integer threshold
\[
  K>\Theta_{\rm tc}(\Gamma+1,a,\beta,R)
\]
and the stretched-exponential sparse threshold
\[
  K>\Theta_{\rm tc}(\Gamma,a,\beta,R).
\]
Since the original stream is a martingale difference,
\[
  \mu_{m,h,i}
  =
  -\mathbb E\left(
  X_{m,h,i}\mathbf 1_{\{\|X_{m,h,i}\|_{\op}>b_{m,h}\}}
  \mid\mcF_{m,h,i-1}\right),
\]
so the drift hypothesis is again only a predictable tail-mean budget.
\end{corollary}

\begin{proof}
Fix \(m,h\).  The variables \(Y_{m,h,i}\) are adapted martingale differences,
and
\[
  \|Y_{m,h,i}\|_{\op}\le2b_{m,h}\le2Rq_{m,h}.
\]
In the self-adjoint case, matrix Freedman applied to \(Y_{m,h,i}\) and
\(-Y_{m,h,i}\) gives
\[
\begin{aligned}
  &\mathbb P\left\{
  \left\|\sum_iY_{m,h,i}\right\|_{\op}>Kq_{m,h}\sqrt{n_m},
  \ \nu_{m,h}\le \beta q_{m,h}^2n_m/\ell_m
  \right\} \\
  &\qquad\le
  2\Delta_m
  \exp\left\{
  -\frac{K^2n_m}
          {2(\beta n_m/\ell_m+2RK\sqrt{n_m}/3)}
  \right\},
\end{aligned}
\]
where
\[
  \nu_{m,h}:=
  \left\|\sum_i
  \mathbb E(Y_{m,h,i}^2\mid\mcF_{m,h,i-1})\right\|_{\op}.
\]
In the rectangular case, apply the same argument to the self-adjoint dilations;
the two displayed predictable row and column bounds are exactly the dilation
variance bound.  The exponent is at least
\(\gamma_{\rm tc}(K,R,\beta,B)\ell_m\) whenever
\(\ell_m\le B\sqrt{n_m}\).  The same Borel--Cantelli argument as in
Corollary~\ref{cor:tail-capped-centered-bernstein-all-plus-schedules} gives
eventual simultaneous control of the centered capped sums, together with the
summable exceptional and cap-failure events.

On the eventual no-cap event,
\[
  X_{m,h,i}=T_{m,h,i}=Y_{m,h,i}+\mu_{m,h,i},
\]
so the drift estimate gives
\[
  \left\|\sum_iX_{m,h,i}\right\|_{\op}
  \le (K+\eta)q_{m,h}\sqrt{n_m}.
\]
The same event gives
\(\max_i\|X_{m,h,i}\|_{\op}\le b_{m,h}\le Rq_{m,h}\).  Dividing by
\(Mq_{m,h}\) proves the all-plus certificate.  The all-integer and sparse
thresholds are the identical summability calculation and quadratic-margin
calculation from
Corollary~\ref{cor:tail-capped-centered-bernstein-all-plus-schedules}.
\end{proof}

\begin{corollary}[Finite-confidence tail-capped all-plus certificates]
\label{cor:tail-capped-finite-confidence}
Let \(\mcG\) be side information, let \(\mathcal A\subset\mathbb N\) be an
almost surely finite \(\mcG\)-measurable active set, and fix
\(n\ge1\) and \(0<\delta<1\).  Work in one of the following two settings.

\emph{Independent setting.}  Conditionally on \(\mcG\), the variables
\(X_{h,1},\ldots,X_{h,n}\) in each selected stream are independent.  They are
either self-adjoint in \(\M_D(\C)\) or rectangular in
\(\M_{d_1\times d_2}(\C)\).  Put
\(\mathcal P_{h,i}=\mcG\).

\emph{Adapted setting.}  For each selected \(h\), let
\((\mcF_{h,i})_{i=0}^n\) be a filtration with \(\mcG\subset\mcF_{h,0}\), and
let \(X_{h,1},\ldots,X_{h,n}\) be a self-adjoint or rectangular
martingale-difference sequence adapted to it.  Put
\(\mathcal P_{h,i}=\mcF_{h,i-1}\).

Let \(\Delta=D\) in the self-adjoint case and
\(\Delta=d_1+d_2\) in the rectangular case.  Let
\(b_h,v_h,u_h,r_h>0\) be \(\mcG\)-measurable proxies.  Define
\[
  T_{h,i}:=X_{h,i}\mathbf 1_{\{\norm{X_{h,i}}_{\op}\le b_h\}},
  \qquad
  \mu_{h,i}:=\mathbb E(T_{h,i}\mid\mathcal P_{h,i}),
\]
\[
  Y_{h,i}:=T_{h,i}-\mu_{h,i},
  \qquad
  F_h^{\rm cap}:=
  \left\{\max_{1\le i\le n}\norm{X_{h,i}}_{\op}>b_h\right\}.
\]
Assume that the selected streams satisfy the drift bound
\[
  \left\|\sum_{i=1}^n\mu_{h,i}\right\|_{\op}\le r_h
\]
almost surely.  In the self-adjoint case assume
\[
  \left\|\sum_{i=1}^n
  \mathbb E(Y_{h,i}^2\mid\mathcal P_{h,i})\right\|_{\op}\le v_h
\]
almost surely.  In the rectangular case replace this by
\[
\begin{aligned}
  \left\|\sum_i
  \mathbb E(Y_{h,i}Y_{h,i}^*\mid\mathcal P_{h,i})\right\|_{\op}
  &\le v_h,\\
  \left\|\sum_i
  \mathbb E(Y_{h,i}^*Y_{h,i}\mid\mathcal P_{h,i})\right\|_{\op}
  &\le v_h
\end{aligned}
\]
almost surely.  If
\[
  \mathbb E\sum_{h\in\mathcal A}
  \left[
  \mathbb P(F_h^{\rm cap}\mid\mcG)
  +2\Delta\exp\left\{-\frac{u_h^2}{2(v_h+2b_hu_h/3)}\right\}
  \right]\le\delta,
\]
then, with probability at least \(1-\delta\), every selected original stream is
all-plus certified after the contraction-normalizing scale
\[
  S_h:=\max\left\{b_h,\frac{u_h+r_h}{\sqrt n}\right\}:
\]
\[
  \disc(X_{h,1}/S_h,\ldots,X_{h,n}/S_h)\le\sqrt n
  \qquad(h\in\mathcal A).
\]

In particular, if \(\abs{\mathcal A}\le M\) almost surely,
\(\mathbb E\sum_{h\in\mathcal A}\mathbb P(F_h^{\rm cap}\mid\mcG)\le\delta/2\),
and \(b_h\le b\), \(v_h\le v\), \(r_h\le r\) on the selected set, then with
\[
  L=\log\frac{4M\Delta}{\delta},
  \qquad
  u=C(\sqrt{vL}+bL),
\]
the same conclusion holds with probability at least \(1-\delta\) after the
homogeneous normalization
\[
  C\max\left\{b,\frac{r+\sqrt{vL}+bL}{\sqrt n}\right\}.
\]
\end{corollary}

\begin{proof}
Condition on \(\mcG\) and fix a selected stream.  In the independent setting
the matrices \(Y_{h,i}\) are conditionally independent and centered.  In the
adapted setting they are martingale differences.  In either case,
\[
  \norm{Y_{h,i}}_{\op}
  \le \norm{T_{h,i}}_{\op}
  +\norm{\mathbb E(T_{h,i}\mid\mathcal P_{h,i})}_{\op}
  \le 2b_h .
\]
The two-sided matrix Bernstein inequality in the independent self-adjoint
case, and the two-sided matrix Freedman inequality in the adapted
self-adjoint case, give
\[
  \mathbb P\left\{
  \left\|\sum_iY_{h,i}\right\|_{\op}>u_h
  \,\middle|\,\mcG
  \right\}
  \le
  2D\exp\left\{-\frac{u_h^2}{2(v_h+2b_hu_h/3)}\right\}.
\]
For rectangular streams apply the same argument to the self-adjoint dilations;
the row and column variance hypotheses are exactly the dilation variance
bound, and the dimension is \(\Delta=d_1+d_2\).

Taking the union bound over the \(\mcG\)-measurable active set and then
expectation gives total failure at most \(\delta\).  On the complementary
event no selected cap failure occurs, so \(X_{h,i}=T_{h,i}\) and
\(\norm{X_{h,i}}_{\op}\le b_h\).  Also,
\[
  \left\|\sum_iX_{h,i}\right\|_{\op}
  \le
  \left\|\sum_iY_{h,i}\right\|_{\op}
  +
  \left\|\sum_i\mu_{h,i}\right\|_{\op}
  \le u_h+r_h .
\]
Dividing by \(S_h\) gives contractions and all-plus sum norm at most
\(\sqrt n\).

For the homogeneous form, the cap failures cost at most \(\delta/2\).  The
usual Bernstein algebra gives
\[
  \frac{u^2}{2(v+2bu/3)}\ge L
\]
after increasing the numerical constant in \(u=C(\sqrt{vL}+bL)\), so the
centered capped-sum failures also cost at most \(\delta/2\).
\end{proof}

\begin{corollary}[Finite-confidence sub-Weibull norm-tail certificates]
\label{cor:subweibull-finite-confidence}
Work in the independent or adapted setting of
Corollary~\ref{cor:tail-capped-finite-confidence}.  Fix
\(\alpha>0\) and \(A_*\ge1\).  For each selected stream let
\(s_h>0\), \(L_h\ge1\), \(v_h>0\), \(u_h>0\), and \(r_h\ge0\) be
\(\mcG\)-measurable, and put
\[
  b_h:=s_hL_h^{1/\alpha}.
\]
Assume the conditional operator-norm cap tails
\[
  \mathbb P\{\norm{X_{h,i}}_{\op}>b_h\mid\mathcal P_{h,i}\}
  \le A_*e^{-L_h}
  \qquad(1\le i\le n).
\]
Define
\[
  T_{h,i}:=X_{h,i}\mathbf 1_{\{\norm{X_{h,i}}_{\op}\le b_h\}},
  \qquad
  \mu_{h,i}:=\mathbb E(T_{h,i}\mid\mathcal P_{h,i}).
\]
Let
\[
  Y_{h,i}:=T_{h,i}-\mu_{h,i},
  \qquad
  F_h^{\rm cap}:=
  \left\{\max_{1\le i\le n}\norm{X_{h,i}}_{\op}>b_h\right\}.
\]
Assume the truncated drift bound
\[
  \left\|\sum_i\mu_{h,i}\right\|_{\op}\le r_h .
\]
In the self-adjoint case assume the uncentered capped variance bound
\[
  \left\|\sum_i
  \mathbb E(T_{h,i}^2\mid\mathcal P_{h,i})\right\|_{\op}\le v_h .
\]
In the rectangular case assume
\[
\begin{aligned}
  \left\|\sum_i
  \mathbb E(T_{h,i}T_{h,i}^*\mid\mathcal P_{h,i})\right\|_{\op}
  &\le v_h,\\
  \left\|\sum_i
  \mathbb E(T_{h,i}^*T_{h,i}\mid\mathcal P_{h,i})\right\|_{\op}
  &\le v_h .
\end{aligned}
\]
If
\[
  \mathbb E\sum_{h\in\mathcal A}
  \left[
  A_*ne^{-L_h}
  +2\Delta\exp\left\{-\frac{u_h^2}{2(v_h+2b_hu_h/3)}\right\}
  \right]\le\delta,
\]
then, with probability at least \(1-\delta\), every selected original stream is
all-plus certified after the contraction-normalizing scale
\[
  S_h:=\max\left\{b_h,\frac{u_h+r_h}{\sqrt n}\right\}.
\]
Equivalently,
\[
  \disc(X_{h,1}/S_h,\ldots,X_{h,n}/S_h)\le\sqrt n
  \qquad(h\in\mathcal A).
\]

In particular, suppose that \(\abs{\mathcal A}\le M\) almost surely, with
\(M\ge1\) deterministic, and that on the selected set
\(s_h\le s\), \(v_h\le v\), and \(r_h\le r\).  Suppose also that the tails hold
at all levels,
\[
  \mathbb P\{\norm{X_{h,i}}_{\op}>s_ht\mid\mathcal P_{h,i}\}
  \le A_*e^{-t^\alpha}
  \qquad(t\ge1).
\]
Set
\[
  L_{\rm cap}:=\max\left\{1,\log\frac{2A_*Mn}{\delta}\right\},
  \qquad
  L_{\rm mat}:=\log\frac{4M\Delta}{\delta},
  \qquad
  b:=sL_{\rm cap}^{1/\alpha}.
\]
Then the same conclusion holds with probability at least \(1-\delta\) after
the homogeneous normalization
\[
  C\max\left\{b,\frac{r+\sqrt{vL_{\rm mat}}+bL_{\rm mat}}{\sqrt n}\right\}.
\]
\end{corollary}

\begin{proof}
For each selected \(h\), the conditional union bound and the tower property give
\[
\begin{aligned}
  \mathbb P(F_h^{\rm cap}\mid\mcG)
  &\le
  \sum_{i=1}^n
  \mathbb E\left[
  \mathbb P\{\norm{X_{h,i}}_{\op}>b_h\mid\mathcal P_{h,i}\}
  \,\middle|\,\mcG\right]  \\
  &\le A_*ne^{-L_h}.
\end{aligned}
\]
It remains to compare the variance hypothesis with the centered capped
variance required in
Corollary~\ref{cor:tail-capped-finite-confidence}.  In the self-adjoint case,
\[
  \mathbb E(Y_{h,i}^2\mid\mathcal P_{h,i})
  =
  \mathbb E(T_{h,i}^2\mid\mathcal P_{h,i})-\mu_{h,i}^2
  \preceq
  \mathbb E(T_{h,i}^2\mid\mathcal P_{h,i}).
\]
In the rectangular case,
\[
  \mathbb E(Y_{h,i}Y_{h,i}^*\mid\mathcal P_{h,i})
  =
  \mathbb E(T_{h,i}T_{h,i}^*\mid\mathcal P_{h,i})-\mu_{h,i}\mu_{h,i}^*
\]
and the analogous identity holds for \(Y_{h,i}^*Y_{h,i}\).  Thus the displayed
uncentered capped variance bounds imply the centered variance bounds of
Corollary~\ref{cor:tail-capped-finite-confidence}.  Applying that corollary
proves the first assertion.

For the homogeneous assertion take \(L_h=L_{\rm cap}\) and
\(b_h=s_hL_{\rm cap}^{1/\alpha}\).  The selected cap cost is at most
\(\delta/2\).  With
\[
  u=C(\sqrt{vL_{\rm mat}}+bL_{\rm mat}),
\]
the usual Bernstein algebra gives
\[
  \frac{u^2}{2(v+2bu/3)}\ge L_{\rm mat}
\]
after increasing the numerical constant \(C\).  Hence the selected centered
matrix-tail cost is at most \(\delta/2\), and the first assertion gives the
homogeneous scale.
\end{proof}

\begin{corollary}[Finite-confidence finite-moment norm-tail certificates]
\label{cor:moment-finite-confidence}
Work in the independent or adapted setting of
Corollary~\ref{cor:tail-capped-finite-confidence}.  Fix \(p>0\) and
\(A_*\ge1\).  For each selected stream let
\(s_h>0\), \(U_h\ge1\), \(v_h>0\), \(u_h>0\), and \(r_h\ge0\) be
\(\mcG\)-measurable, and put
\[
  b_h:=s_hU_h^{1/p}.
\]
Assume the conditional operator-norm moment bound
\[
  \mathbb E(\norm{X_{h,i}}_{\op}^p\mid\mathcal P_{h,i})
  \le A_*s_h^p
  \qquad(1\le i\le n).
\]
Define
\[
  T_{h,i}:=X_{h,i}\mathbf 1_{\{\norm{X_{h,i}}_{\op}\le b_h\}},
  \qquad
  \mu_{h,i}:=\mathbb E(T_{h,i}\mid\mathcal P_{h,i}).
\]
Let
\[
  F_h^{\rm cap}:=
  \left\{\max_{1\le i\le n}\norm{X_{h,i}}_{\op}>b_h\right\}.
\]
Assume the truncated drift bound
\[
  \left\|\sum_i\mu_{h,i}\right\|_{\op}\le r_h .
\]
In the self-adjoint case assume the uncentered capped variance bound
\[
  \left\|\sum_i
  \mathbb E(T_{h,i}^2\mid\mathcal P_{h,i})\right\|_{\op}\le v_h .
\]
In the rectangular case assume
\[
\begin{aligned}
  \left\|\sum_i
  \mathbb E(T_{h,i}T_{h,i}^*\mid\mathcal P_{h,i})\right\|_{\op}
  &\le v_h,\\
  \left\|\sum_i
  \mathbb E(T_{h,i}^*T_{h,i}\mid\mathcal P_{h,i})\right\|_{\op}
  &\le v_h .
\end{aligned}
\]
If
\[
  \mathbb E\sum_{h\in\mathcal A}
  \left[
  A_*nU_h^{-1}
  +2\Delta\exp\left\{-\frac{u_h^2}{2(v_h+2b_hu_h/3)}\right\}
  \right]\le\delta,
\]
then, with probability at least \(1-\delta\), every selected original stream
satisfies
\[
  \disc(X_{h,1}/S_h,\ldots,X_{h,n}/S_h)\le\sqrt n,
  \qquad
  S_h:=\max\left\{b_h,\frac{u_h+r_h}{\sqrt n}\right\}.
\]

The uncentered capped variance hypotheses are automatic with
\[
  v_h:=A_*s_h^2n\,U_h^{(2-p)_+/p}.
\]
If, in addition, the independent streams are conditionally centered given
\(\mcG\), or the adapted streams are martingale differences, and \(p\ge1\),
then the truncated drift bound is automatic with
\[
  r_h:=A_*ns_h\,U_h^{(1-p)/p}.
\]

In particular, suppose that \(\abs{\mathcal A}\le M\) almost surely, with
\(M\ge1\) deterministic, that \(s_h\le s\) on the selected set, and that the
displayed moment bound holds.  Set
\[
  U:=\max\left\{1,\frac{2A_*Mn}{\delta}\right\},
  \qquad
  L:=\log\frac{4M\Delta}{\delta},
  \qquad
  b:=sU^{1/p}.
\]
If \(v_h\le v\) and \(r_h\le r\) on the selected set, then the same conclusion
holds with probability at least \(1-\delta\) after the homogeneous
normalization
\[
  C\max\left\{b,\frac{r+\sqrt{vL}+bL}{\sqrt n}\right\}.
\]
One may take
\[
  v=A_*s^2n\,U^{(2-p)_+/p},
\]
and, in the centered or martingale-difference case with \(p\ge1\),
\[
  r=A_*ns\,U^{(1-p)/p}.
\]
\end{corollary}

\begin{proof}
Markov's inequality gives, for every selected \(h\) and \(i\),
\[
  \mathbb P\{\norm{X_{h,i}}_{\op}>b_h\mid\mathcal P_{h,i}\}
  \le A_*U_h^{-1}.
\]
Thus the same conditional union-bound calculation as in
Corollary~\ref{cor:subweibull-finite-confidence} gives
\[
  \mathbb P(F_h^{\rm cap}\mid\mcG)\le A_*nU_h^{-1}.
\]
The uncentered capped variance bounds imply the centered capped variance
bounds by the same subtraction-of-squares argument used in
Corollary~\ref{cor:subweibull-finite-confidence}.  Hence
Corollary~\ref{cor:tail-capped-finite-confidence} proves the first assertion.

It remains to justify the automatic proxies.  In the self-adjoint case,
\(T_{h,i}^2\preceq\norm{T_{h,i}}_{\op}^2I\); in the rectangular case,
\(T_{h,i}T_{h,i}^*\) and \(T_{h,i}^*T_{h,i}\) are bounded by the same scalar
multiple of the identity.  Therefore
\[
  \left\|\mathbb E(T_{h,i}^2\mid\mathcal P_{h,i})\right\|_{\op}
  \le \mathbb E(\norm{T_{h,i}}_{\op}^2\mid\mathcal P_{h,i})
\]
and similarly for the rectangular row and column terms.  If \(p\ge2\), Jensen's
inequality and \(A_*\ge1\) give
\[
  \mathbb E(\norm{T_{h,i}}_{\op}^2\mid\mathcal P_{h,i})
  \le A_*s_h^2.
\]
If \(0<p<2\), then on the cap event
\(\norm{T_{h,i}}_{\op}^2\le b_h^{2-p}\norm{X_{h,i}}_{\op}^p\), so the same
conditional moment bound gives
\[
  \mathbb E(\norm{T_{h,i}}_{\op}^2\mid\mathcal P_{h,i})
  \le A_*s_h^2U_h^{(2-p)/p}.
\]
Summing over \(i\) gives the displayed \(v_h\).

If the original stream is conditionally centered or a martingale difference,
then
\[
  \mu_{h,i}
  =
  -\mathbb E\left(
  X_{h,i}\mathbf 1_{\{\norm{X_{h,i}}_{\op}>b_h\}}
  \,\middle|\,\mathcal P_{h,i}\right).
\]
For \(p\ge1\),
\[
  \norm{\mu_{h,i}}_{\op}
  \le b_h^{1-p}
  \mathbb E(\norm{X_{h,i}}_{\op}^p\mid\mathcal P_{h,i})
  \le A_*s_hU_h^{(1-p)/p}.
\]
Summing over \(i\) gives the displayed \(r_h\).

For the homogeneous form take \(U_h=U\).  The selected cap cost is at most
\(\delta/2\).  With \(u=C(\sqrt{vL}+bL)\), the usual Bernstein algebra gives
\[
  \frac{u^2}{2(v+2bu/3)}\ge L
\]
after increasing \(C\), so the selected matrix-tail cost is also at most
\(\delta/2\).
\end{proof}

\begin{corollary}[Finite-confidence heterogeneous cap-profile certificates]
\label{cor:heterogeneous-cap-profile-finite-confidence}
Work in the independent or adapted setting of
Corollary~\ref{cor:tail-capped-finite-confidence}.  For each selected stream
let \(b_{h,i}>0\) and \(\eta_{h,i}\ge0\) be \(\mcG\)-measurable, and assume
\[
  \mathbb P\{\norm{X_{h,i}}_{\op}>b_{h,i}\mid\mathcal P_{h,i}\}
  \le \eta_{h,i}
  \qquad(1\le i\le n).
\]
Let \(v_h>0\), \(u_h>0\), and \(r_h\ge0\) be \(\mcG\)-measurable.  Put
\[
  b_h^*:=\max_{1\le i\le n}b_{h,i}.
\]
Define the heterogeneously capped variables
\[
  T_{h,i}:=X_{h,i}\mathbf 1_{\{\norm{X_{h,i}}_{\op}\le b_{h,i}\}},
  \qquad
  \mu_{h,i}:=\mathbb E(T_{h,i}\mid\mathcal P_{h,i}),
\]
and assume the drift bound
\[
  \left\|\sum_i\mu_{h,i}\right\|_{\op}\le r_h .
\]
In the self-adjoint case assume the uncentered capped variance bound
\[
  \left\|\sum_i
  \mathbb E(T_{h,i}^2\mid\mathcal P_{h,i})\right\|_{\op}\le v_h .
\]
In the rectangular case assume
\[
\begin{aligned}
  \left\|\sum_i
  \mathbb E(T_{h,i}T_{h,i}^*\mid\mathcal P_{h,i})\right\|_{\op}
  &\le v_h,\\
  \left\|\sum_i
  \mathbb E(T_{h,i}^*T_{h,i}\mid\mathcal P_{h,i})\right\|_{\op}
  &\le v_h .
\end{aligned}
\]
If
\[
  \mathbb E\sum_{h\in\mathcal A}
  \left[
  \sum_{i=1}^n\eta_{h,i}
  +2\Delta\exp\left\{-\frac{u_h^2}{2(v_h+2b_h^*u_h/3)}\right\}
  \right]\le\delta,
\]
then, with probability at least \(1-\delta\), every selected original stream
satisfies
\[
  \disc(X_{h,1}/S_h,\ldots,X_{h,n}/S_h)\le\sqrt n,
  \qquad
  S_h:=\max\left\{b_h^*,\frac{u_h+r_h}{\sqrt n}\right\}.
\]

In particular, if \(\abs{\mathcal A}\le M\) almost surely,
\[
  \mathbb E\sum_{h\in\mathcal A}\sum_i\eta_{h,i}\le\delta/2,
\]
and \(b_h^*\le b\), \(v_h\le v\), \(r_h\le r\) on the selected set, then with
\[
  L=\log\frac{4M\Delta}{\delta},
  \qquad
  u=C(\sqrt{vL}+bL),
\]
the same conclusion holds with probability at least \(1-\delta\) after the
homogeneous normalization
\[
  C\max\left\{b,\frac{r+\sqrt{vL}+bL}{\sqrt n}\right\}.
\]
\end{corollary}

\begin{proof}
Let
\[
  F_h^{\rm cap}:=
  \left\{\max_i\norm{X_{h,i}}_{\op}>b_{h,i}\right\}.
\]
The conditional union bound and the tower property give
\[
  \mathbb P(F_h^{\rm cap}\mid\mcG)\le\sum_i\eta_{h,i}.
\]
Set \(Y_{h,i}:=T_{h,i}-\mu_{h,i}\).  In both the independent and adapted
settings these are centered independent variables, conditionally on \(\mcG\),
or martingale differences.  Moreover
\[
  \norm{Y_{h,i}}_{\op}\le2b_h^* .
\]
The uncentered capped variance bounds imply the centered capped variance
bounds by the same subtraction-of-squares argument used in
Corollary~\ref{cor:subweibull-finite-confidence}.

Matrix Bernstein in the independent case and matrix Freedman in the adapted
case, with self-adjoint dilation for rectangular streams, give
\[
  \mathbb P\left\{
  \left\|\sum_iY_{h,i}\right\|_{\op}>u_h
  \,\middle|\,\mcG
  \right\}
  \le
  2\Delta\exp\left\{-\frac{u_h^2}{2(v_h+2b_h^*u_h/3)}\right\}.
\]
Taking the union bound over the active set and then expectation gives total
failure at most \(\delta\).  On the complementary event
\(X_{h,i}=T_{h,i}\), \(\norm{X_{h,i}}_{\op}\le b_{h,i}\le b_h^*\), and
\[
  \left\|\sum_iX_{h,i}\right\|_{\op}
  \le u_h+r_h .
\]
Dividing by \(S_h\) proves the first assertion.  The homogeneous form is the
same Bernstein algebra used in
Corollary~\ref{cor:tail-capped-finite-confidence}.
\end{proof}

\begin{corollary}[Finite-confidence heterogeneous finite-moment certificates]
\label{cor:heterogeneous-moment-finite-confidence}
Work in the independent or adapted setting of
Corollary~\ref{cor:tail-capped-finite-confidence}.  Fix \(A_*\ge1\).  For
each selected stream and each \(1\le i\le n\), let
\(p_{h,i}>0\), \(s_{h,i}>0\), and \(U_{h,i}\ge1\) be
\(\mcG\)-measurable, and assume the conditional moment envelope
\[
  \mathbb E(\norm{X_{h,i}}_{\op}^{p_{h,i}}\mid\mathcal P_{h,i})
  \le A_*s_{h,i}^{p_{h,i}} .
\]
Put
\[
  b_{h,i}:=s_{h,i}U_{h,i}^{1/p_{h,i}},
  \qquad
  b_h^*:=\max_{1\le i\le n}b_{h,i},
\]
and
\[
  V_h:=
  A_*\sum_{i=1}^n
  s_{h,i}^2U_{h,i}^{(2-p_{h,i})_+/p_{h,i}} .
\]
For \(\mcG\)-measurable \(u_h>0\), define the heterogeneously capped variables
\[
  T_{h,i}:=X_{h,i}\mathbf 1_{\{\norm{X_{h,i}}_{\op}\le b_{h,i}\}},
  \qquad
  \mu_{h,i}:=\mathbb E(T_{h,i}\mid\mathcal P_{h,i}).
\]
Assume the truncated drift bound
\[
  \left\|\sum_i\mu_{h,i}\right\|_{\op}\le r_h
\]
for some \(\mcG\)-measurable \(r_h\ge0\).  If
\[
  \mathbb E\sum_{h\in\mathcal A}
  \left[
  A_*\sum_{i=1}^nU_{h,i}^{-1}
  +2\Delta\exp\left\{-\frac{u_h^2}{2(V_h+2b_h^*u_h/3)}\right\}
  \right]\le\delta,
\]
then, with probability at least \(1-\delta\), every selected original stream
satisfies
\[
  \disc(X_{h,1}/S_h,\ldots,X_{h,n}/S_h)\le\sqrt n,
  \qquad
  S_h:=\max\left\{b_h^*,\frac{u_h+r_h}{\sqrt n}\right\}.
\]
The same conclusion holds in the rectangular case, with the same
\(\Delta\)-convention as in
Corollary~\ref{cor:heterogeneous-cap-profile-finite-confidence}.

If the independent streams are conditionally centered given \(\mcG\), or the
adapted streams are martingale differences, and \(p_{h,i}\ge1\) for every
selected \(h,i\), then the drift hypothesis is automatic with
\[
  r_h:=
  A_*\sum_{i=1}^n
  s_{h,i}U_{h,i}^{(1-p_{h,i})/p_{h,i}} .
\]

In particular, if \(\abs{\mathcal A}\le M\) almost surely,
\[
  \mathbb E\sum_{h\in\mathcal A}
  A_*\sum_{i=1}^nU_{h,i}^{-1}\le\delta/2,
\]
and \(b_h^*\le b\), \(V_h\le v\), \(r_h\le r\) on the selected set, then with
\[
  L=\log\frac{4M\Delta}{\delta},
  \qquad
  u=C(\sqrt{vL}+bL),
\]
the same conclusion holds with probability at least \(1-\delta\) after the
homogeneous normalization
\[
  C\max\left\{b,\frac{r+\sqrt{vL}+bL}{\sqrt n}\right\}.
\]
\end{corollary}

\begin{proof}
Markov's inequality gives
\[
  \mathbb P\{\norm{X_{h,i}}_{\op}>b_{h,i}\mid\mathcal P_{h,i}\}
  \le A_*U_{h,i}^{-1}.
\]
It remains only to verify the capped second-moment input in
Corollary~\ref{cor:heterogeneous-cap-profile-finite-confidence}.  In the
self-adjoint case,
\(T_{h,i}^2\preceq\norm{T_{h,i}}_{\op}^2I\); in the rectangular case, the row
and column squares have the same scalar upper bound.  If \(p_{h,i}\ge2\),
Jensen's inequality and \(A_*\ge1\) give
\[
  \mathbb E(\norm{T_{h,i}}_{\op}^2\mid\mathcal P_{h,i})
  \le A_*s_{h,i}^2 .
\]
If \(0<p_{h,i}<2\), then on the cap event
\[
  \norm{T_{h,i}}_{\op}^2
  \le b_{h,i}^{2-p_{h,i}}\norm{X_{h,i}}_{\op}^{p_{h,i}},
\]
so
\[
  \mathbb E(\norm{T_{h,i}}_{\op}^2\mid\mathcal P_{h,i})
  \le A_*s_{h,i}^2U_{h,i}^{(2-p_{h,i})/p_{h,i}} .
\]
Summing these bounds gives the self-adjoint, and rectangular row and column,
uncentered capped variance bound with \(V_h\).  Applying
Corollary~\ref{cor:heterogeneous-cap-profile-finite-confidence} proves the
first assertion.

In the centered or martingale-difference case,
\[
  \mu_{h,i}
  =
  -\mathbb E\left(
  X_{h,i}\mathbf 1_{\{\norm{X_{h,i}}_{\op}>b_{h,i}\}}
  \,\middle|\,\mathcal P_{h,i}\right).
\]
For \(p_{h,i}\ge1\),
\[
  \norm{\mu_{h,i}}_{\op}
  \le b_{h,i}^{1-p_{h,i}}
  \mathbb E(\norm{X_{h,i}}_{\op}^{p_{h,i}}\mid\mathcal P_{h,i})
  \le A_*s_{h,i}U_{h,i}^{(1-p_{h,i})/p_{h,i}} .
\]
Summing over \(i\) gives the displayed automatic drift proxy.

For the homogeneous form, the selected cap cost is at most \(\delta/2\).  With
\(u=C(\sqrt{vL}+bL)\), the usual Bernstein algebra gives
\[
  \frac{u^2}{2(v+2bu/3)}\ge L
\]
after increasing \(C\), so the selected centered matrix-tail cost is also at
most \(\delta/2\).
\end{proof}

\begin{corollary}[Finite-confidence tail-profile certificates]
\label{cor:tail-profile-finite-confidence}
Work in the independent or adapted setting of
Corollary~\ref{cor:tail-capped-finite-confidence}.  For each selected stream
and each \(1\le i\le n\), let \(s_{h,i}>0\), \(\lambda_{h,i}>0\), and
\(u_h>0\) be \(\mcG\)-measurable, and let
\(\pi_{h,i}:[0,\infty)\to[0,1]\) be deterministic measurable profiles
satisfying
\[
  \mathbb P\{\norm{X_{h,i}}_{\op}>s_{h,i}t\mid\mathcal P_{h,i}\}
  \le \pi_{h,i}(t)
  \qquad(t\ge0).
\]
Put
\[
  b_{h,i}:=s_{h,i}\lambda_{h,i},
  \qquad
  b_h^*:=\max_{1\le i\le n}b_{h,i},
\]
and
\[
  V_h:=
  2\sum_{i=1}^n
  s_{h,i}^2\int_0^{\lambda_{h,i}}t\,\pi_{h,i}(t)\,dt .
\]
Define
\[
  T_{h,i}:=X_{h,i}\mathbf 1_{\{\norm{X_{h,i}}_{\op}\le b_{h,i}\}},
  \qquad
  \mu_{h,i}:=\mathbb E(T_{h,i}\mid\mathcal P_{h,i}).
\]
Assume
\[
  \left\|\sum_i\mu_{h,i}\right\|_{\op}\le r_h
\]
for some \(\mcG\)-measurable \(r_h\ge0\).  If
\[
  \mathbb E\sum_{h\in\mathcal A}
  \left[
  \sum_{i=1}^n\pi_{h,i}(\lambda_{h,i})
  +2\Delta\exp\left\{-\frac{u_h^2}{2(V_h+2b_h^*u_h/3)}\right\}
  \right]\le\delta,
\]
then, with probability at least \(1-\delta\), every selected original stream
satisfies
\[
  \disc(X_{h,1}/S_h,\ldots,X_{h,n}/S_h)\le\sqrt n,
  \qquad
  S_h:=\max\left\{b_h^*,\frac{u_h+r_h}{\sqrt n}\right\}.
\]
The same conclusion holds in the rectangular case, with the usual
\(\Delta=d_1+d_2\) dilation convention.

If the independent streams are conditionally centered given \(\mcG\), or the
adapted streams are martingale differences, then the drift hypothesis is
automatic whenever
\[
  \sum_i s_{h,i}
  \left[
    \lambda_{h,i}\pi_{h,i}(\lambda_{h,i})
    +\int_{\lambda_{h,i}}^\infty\pi_{h,i}(t)\,dt
  \right]\le r_h .
\]

In particular, if \(\abs{\mathcal A}\le M\) almost surely,
\[
  \mathbb E\sum_{h\in\mathcal A}
  \sum_i\pi_{h,i}(\lambda_{h,i})\le\delta/2,
\]
and \(b_h^*\le b\), \(V_h\le v\), \(r_h\le r\) on the selected set, then with
\[
  L=\log\frac{4M\Delta}{\delta},
  \qquad
  u=C(\sqrt{vL}+bL),
\]
the same conclusion holds with probability at least \(1-\delta\) after the
homogeneous normalization
\[
  C\max\left\{b,\frac{r+\sqrt{vL}+bL}{\sqrt n}\right\}.
\]
\end{corollary}

\begin{proof}
The profile tail bound gives
\[
  \mathbb P\{\norm{X_{h,i}}_{\op}>b_{h,i}\mid\mathcal P_{h,i}\}
  \le \pi_{h,i}(\lambda_{h,i}).
\]
For the capped second moment, in the self-adjoint case
\(T_{h,i}^2\preceq\norm{T_{h,i}}_{\op}^2I\), and similarly for rectangular row
and column squares.  By layer cake and the profile tail bound,
\[
\begin{aligned}
  \mathbb E(\norm{T_{h,i}}_{\op}^2\mid\mathcal P_{h,i})
  &\le s_{h,i}^2
  \int_0^{\lambda_{h,i}^2}
  \mathbb P\{\norm{X_{h,i}}_{\op}^2>s_{h,i}^2 r
  \mid\mathcal P_{h,i}\}\,dr\\
  &\le
  2s_{h,i}^2\int_0^{\lambda_{h,i}}t\,\pi_{h,i}(t)\,dt .
\end{aligned}
\]
Summing gives the self-adjoint and rectangular uncentered capped variance
bound with \(V_h\).  Corollary~\ref{cor:heterogeneous-cap-profile-finite-confidence}
therefore gives the first assertion.

For centered or martingale-difference streams,
\[
  \mu_{h,i}
  =
  -\mathbb E\left(
  X_{h,i}\mathbf 1_{\{\norm{X_{h,i}}_{\op}>b_{h,i}\}}
  \,\middle|\,\mathcal P_{h,i}\right),
\]
and another layer-cake estimate gives
\[
  \norm{\mu_{h,i}}_{\op}
  \le s_{h,i}
  \left[
    \lambda_{h,i}\pi_{h,i}(\lambda_{h,i})
    +\int_{\lambda_{h,i}}^\infty\pi_{h,i}(t)\,dt
  \right].
\]
Summing proves the drift claim.  The homogeneous form is the same Bernstein
algebra used in
Corollary~\ref{cor:heterogeneous-cap-profile-finite-confidence}.
\end{proof}

\begin{corollary}[Sub-Weibull norm-tail cap selection]
\label{cor:subweibull-norm-tail-capped-all-plus-schedules}
Work in either the independent setting of
Corollary~\ref{cor:tail-capped-centered-bernstein-all-plus-schedules} or the
adapted martingale-difference setting of
Corollary~\ref{cor:tail-capped-freedman-all-plus-schedules}.  Let
\(\mathcal P_{m,h,i}\) denote the conditioning field
\(\mathcal H_m\) in the independent case and
\(\mcF_{m,h,i-1}\) in the adapted case.  Fix
\(\alpha>0\), \(A_*\ge1\), \(K>0\), \(\beta\ge0\),
\(\eta\ge0\), and \(B>0\).  For each active stream let
\(s_{m,h}>0\) be a deterministic norm-tail scale, and let \(L_m\ge1\)
be deterministic.  Put
\[
  b_{m,h}:=s_{m,h}L_m^{1/\alpha},
  \qquad
  q_{m,h}:=
  s_{m,h}\max\{L_m^{1/\alpha},\sqrt{\ell_m}\}.
\]
Assume the conditional operator-norm tails
\[
  \mathbb P\{\|X_{m,h,i}\|_{\op}>b_{m,h}\mid\mathcal P_{m,h,i}\}
  \le A_*e^{-L_m}
  \qquad(1\le i\le n_m)
\]
and let
\[
  T_{m,h,i}:=
  X_{m,h,i}\mathbf 1_{\{\|X_{m,h,i}\|_{\op}\le b_{m,h}\}},
  \qquad
  \mu_{m,h,i}:=\mathbb E(T_{m,h,i}\mid\mathcal P_{m,h,i}).
\]
Outside events \(E_{m,h}\), assume the truncated-tail drift bound
\[
  \left\|\sum_{i=1}^{n_m}\mu_{m,h,i}\right\|_{\op}
  \le \eta q_{m,h}\sqrt{n_m}.
\]
In the self-adjoint case assume also, outside \(E_{m,h}\), that
\[
  \left\|\sum_{i=1}^{n_m}
  \mathbb E(T_{m,h,i}^2\mid\mathcal P_{m,h,i})\right\|_{\op}
  \le \beta s_{m,h}^2n_m.
\]
In the rectangular case replace this by the two truncated row and column
profile bounds
\[
\begin{aligned}
  \left\|\sum_i
  \mathbb E(T_{m,h,i}T_{m,h,i}^*\mid\mathcal P_{m,h,i})\right\|_{\op}
  &\le \beta s_{m,h}^2n_m,\\
  \left\|\sum_i
  \mathbb E(T_{m,h,i}^*T_{m,h,i}\mid\mathcal P_{m,h,i})\right\|_{\op}
  &\le \beta s_{m,h}^2n_m.
\end{aligned}
\]
Suppose that \(\ell_m\le B\sqrt{n_m}\) eventually and that
\[
  \sum_m\sum_{h=1}^{J_m}
  \left[\mathbb P(E_{m,h})+A_*n_m e^{-L_m}\right]<\infty
\]
and
\[
  \sum_mJ_m\Delta_m
  (e\Delta_m)^{-\gamma_{\rm tc}(K,1,\beta,B)}<\infty .
\]
Then, almost surely, for all sufficiently large \(m\), every active original
stream is contraction-normalized by
\[
  Mq_{m,h},
  \qquad
  M:=\max\{1,K+\eta\},
\]
and the all-plus signing gives
\[
  \disc(X_{m,h,1}/(Mq_{m,h}),\ldots,X_{m,h,n_m}/(Mq_{m,h}))
  \le\sqrt{n_m}
  \qquad(1\le h\le J_m).
\]

Consequently, if
\[
  J_m\Delta_m\le(e\Delta_m)^\Gamma,\qquad
  J_mn_m\le(e\Delta_m)^\Lambda,\qquad
  n_m\ge a\ell_m^2
\]
eventually and \(L_m=c\ell_m\), then the conclusion holds along all integer
dimensions whenever
\[
  c>\Lambda+1,
  \qquad
  K>\Theta_{\rm tc}(\Gamma+1,a,\beta,1),
\]
and along every stretched-exponential sparse dimension schedule whenever
\[
  c>\Lambda,
  \qquad
  K>\Theta_{\rm tc}(\Gamma,a,\beta,1).
\]
In these closed-form schedules the normalizing scale is
\[
  q_{m,h}
  =
  s_{m,h}\max\{(c\ell_m)^{1/\alpha},\sqrt{\ell_m}\}.
\]
For conditionally centered independent streams, and for the martingale
streams above, the displayed sub-Weibull tail bound makes the drift hypothesis
automatic for every fixed \(\eta>0\) under the same closed-form cap choices.
\end{corollary}

\begin{proof}
The conditional tail bound and a union bound give
\[
  \mathbb P(F^{\rm cap}_{m,h})
  \le A_*n_m e^{-L_m}.
\]
Moreover \(b_{m,h}\le q_{m,h}\), so the tail-capped corollaries apply with
\(R=1\) once their capped variance hypotheses are verified.  In the
self-adjoint case,
\[
  \mathbb E\bigl((T_{m,h,i}-\mu_{m,h,i})^2\mid\mathcal P_{m,h,i}\bigr)
  =
  \mathbb E(T_{m,h,i}^2\mid\mathcal P_{m,h,i})
  -\mu_{m,h,i}^2
  \preceq
  \mathbb E(T_{m,h,i}^2\mid\mathcal P_{m,h,i}).
\]
The rectangular row and column identities are the same with
\(TT^*\) and \(T^*T\).  Since
\[
  q_{m,h}^2/\ell_m
  =
  s_{m,h}^2
  \max\{L_m^{2/\alpha}/\ell_m,1\}
  \ge s_{m,h}^2,
\]
the displayed truncated profiles imply the required ordinary or predictable
variance bounds with the same \(\beta\).  Corollaries
\ref{cor:tail-capped-centered-bernstein-all-plus-schedules} and
\ref{cor:tail-capped-freedman-all-plus-schedules}, with \(R=1\), now give the
general conclusion.

For the closed-form schedules, the cap-tail summand is at most
\((e\Delta_m)^{\Lambda-c}\), up to the fixed factor \(A_*\).  This is
summable over all integer dimensions when \(c>\Lambda+1\), and along every
stretched-exponential sparse schedule when \(c>\Lambda\).  The all-plus
matrix-tail summability is exactly the corresponding
\(\Theta_{\rm tc}\) calculation from the two tail-capped corollaries, again
with \(R=1\).

Finally suppose the original variables are centered relative to
\(\mathcal P_{m,h,i}\).  For
\[
  \psi_\alpha(u):=
  u e^{-u^\alpha}+\int_u^\infty e^{-t^\alpha}\,dt
\]
the conditional tail bound gives
\[
  \|\mu_{m,h,i}\|_{\op}
  \le A_*s_{m,h}\psi_\alpha(L_m^{1/\alpha}).
\]
With \(L_m=c\ell_m\) and \(J_mn_m\le(e\Delta_m)^\Lambda\),
\[
  \frac{\sum_i\|\mu_{m,h,i}\|_{\op}}
       {q_{m,h}\sqrt{n_m}}
  \le
  A_*\frac{\sqrt{n_m}\,\psi_\alpha((c\ell_m)^{1/\alpha})}
            {\max\{(c\ell_m)^{1/\alpha},\sqrt{\ell_m}\}}
  \longrightarrow0
\]
under the displayed closed-form cap choices: indeed
\(\sqrt{n_m}\le(e\Delta_m)^{\Lambda/2}\), while
\(\psi_\alpha((c\ell_m)^{1/\alpha})\) is a polynomial factor in \(\ell_m\)
times \((e\Delta_m)^{-c}\).  Hence the drift bound holds eventually for every
fixed \(\eta>0\).
\end{proof}

\begin{corollary}[Finite-moment norm-tail cap selection]
\label{cor:moment-norm-tail-capped-all-plus-schedules}
Work in the same independent or adapted setting, with the same conditioning
fields \(\mathcal P_{m,h,i}\), as in
Corollary~\ref{cor:subweibull-norm-tail-capped-all-plus-schedules}.  Fix
\(p>0\), \(A_*\ge1\), \(K>0\), \(\eta\ge0\), and \(B>0\).  For each active
stream let \(s_{m,h}>0\) be deterministic and assume the conditional moment
envelope
\[
  \mathbb E(\|X_{m,h,i}\|_{\op}^p\mid\mathcal P_{m,h,i})
  \le A_*s_{m,h}^p
  \qquad(1\le i\le n_m).
\]
Let \(U_m\ge1\) be deterministic and put
\[
  \theta_p:=\frac{(2-p)_+}{2p},
  \qquad
  b_{m,h}:=s_{m,h}U_m^{1/p},
\]
\[
  q_{m,h}:=
  s_{m,h}\max\{U_m^{1/p},\sqrt{\ell_m}\,U_m^{\theta_p}\}.
\]
Define \(T_{m,h,i}\) and \(\mu_{m,h,i}\) as in
Corollary~\ref{cor:subweibull-norm-tail-capped-all-plus-schedules}.  Outside
events \(E_{m,h}\), assume the truncated-tail drift bound
\[
  \left\|\sum_{i=1}^{n_m}\mu_{m,h,i}\right\|_{\op}
  \le \eta q_{m,h}\sqrt{n_m}.
\]
Suppose that \(\ell_m\le B\sqrt{n_m}\) eventually and that
\[
  \sum_m\sum_{h=1}^{J_m}
  \left[\mathbb P(E_{m,h})+A_*n_mU_m^{-1}\right]<\infty
\]
and
\[
  \sum_mJ_m\Delta_m
  (e\Delta_m)^{-\gamma_{\rm tc}(K,1,A_*,B)}<\infty .
\]
Then, almost surely, for all sufficiently large \(m\), every active original
stream is contraction-normalized by
\[
  Mq_{m,h},
  \qquad
  M:=\max\{1,K+\eta\},
\]
and the all-plus signing gives
\[
  \disc(X_{m,h,1}/(Mq_{m,h}),\ldots,X_{m,h,n_m}/(Mq_{m,h}))
  \le\sqrt{n_m}
  \qquad(1\le h\le J_m).
\]

Consequently, if
\[
  J_m\Delta_m\le(e\Delta_m)^\Gamma,\qquad
  J_mn_m\le(e\Delta_m)^\Lambda,\qquad
  n_m\ge a\ell_m^2
\]
eventually and \(U_m=(e\Delta_m)^c\), then the conclusion holds along all
integer dimensions whenever
\[
  c>\Lambda+1,
  \qquad
  K>\Theta_{\rm tc}(\Gamma+1,a,A_*,1),
\]
and along every stretched-exponential sparse dimension schedule whenever
\[
  c>\Lambda,
  \qquad
  K>\Theta_{\rm tc}(\Gamma,a,A_*,1).
\]
In these closed-form schedules the normalizing scale is
\[
  q_{m,h}
  =
  s_{m,h}\max\{(e\Delta_m)^{c/p},
  \sqrt{\ell_m}\,(e\Delta_m)^{c\theta_p}\}.
\]
If \(p\ge1\), and if the original streams are conditionally centered in the
independent case or martingale differences in the adapted case, then the moment
envelope makes the drift hypothesis automatic for every fixed \(\eta>0\) under
the same closed-form cap choices.
\end{corollary}

\begin{proof}
Markov's inequality gives
\[
  \mathbb P(F^{\rm cap}_{m,h})
  \le A_*n_mU_m^{-1}.
\]
Also \(b_{m,h}\le q_{m,h}\), so the tail-capped corollaries may be used with
\(R=1\) once the capped variance budget is verified.  In the self-adjoint
case,
\[
  T_{m,h,i}^2\preceq \|T_{m,h,i}\|_{\op}^2I.
\]
If \(0<p<2\), then on the cap event
\[
  \|T_{m,h,i}\|_{\op}^2
  \le b_{m,h}^{2-p}\|X_{m,h,i}\|_{\op}^p,
\]
while if \(p\ge2\), conditional Lyapunov monotonicity gives
\[
  \mathbb E(\|T_{m,h,i}\|_{\op}^2\mid\mathcal P_{m,h,i})
  \le A_*s_{m,h}^2 .
\]
Thus, in all cases,
\[
  \mathbb E(T_{m,h,i}^2\mid\mathcal P_{m,h,i})
  \preceq
  A_*s_{m,h}^2U_m^{(2-p)_+/p}I.
\]
The same argument with \(TT^*\) and \(T^*T\) gives the rectangular row and
column bounds.  Since
\[
  q_{m,h}^2/\ell_m
  \ge s_{m,h}^2U_m^{(2-p)_+/p},
\]
the centered capped ordinary or predictable variance profile is at most
\(A_*q_{m,h}^2n_m/\ell_m\).  Corollaries
\ref{cor:tail-capped-centered-bernstein-all-plus-schedules} and
\ref{cor:tail-capped-freedman-all-plus-schedules}, with \(R=1\) and
\(\beta=A_*\), prove the general conclusion.

For the closed-form schedules, the cap-tail summand is bounded by
\((e\Delta_m)^{\Lambda-c}\), up to \(A_*\), so it is summable over all integer
dimensions when \(c>\Lambda+1\) and along every stretched-exponential sparse
schedule when \(c>\Lambda\).  The matrix-tail summability is exactly the same
\(\Theta_{\rm tc}\) calculation with \(\beta=A_*\).

If \(p\ge1\) and the original variables are centered relative to
\(\mathcal P_{m,h,i}\), then
\[
  \|\mu_{m,h,i}\|_{\op}
  \le A_*s_{m,h}U_m^{-(p-1)/p}.
\]
Because \(q_{m,h}\ge s_{m,h}U_m^{1/p}\), we get
\[
  \frac{\sum_i\|\mu_{m,h,i}\|_{\op}}
       {q_{m,h}\sqrt{n_m}}
  \le A_*\sqrt{n_m}\,U_m^{-1}.
\]
With \(U_m=(e\Delta_m)^c\) and \(J_mn_m\le(e\Delta_m)^\Lambda\), this tends
to \(0\) under the displayed closed-form cap choices.  Hence the drift
hypothesis holds eventually for every fixed \(\eta>0\).
\end{proof}

\begin{corollary}[Heterogeneous norm-tail cap profiles]
\label{cor:heterogeneous-norm-tail-cap-profiles}
Work in the same independent or adapted setting, with the same conditioning
fields \(\mathcal P_{m,h,i}\), as in
Corollary~\ref{cor:subweibull-norm-tail-capped-all-plus-schedules}.  For each
active stream let \(q_{m,h}>0\) be deterministic, and let
\(b_{m,h,i}>0\) be deterministic caps satisfying
\[
  b_{m,h,i}\le Rq_{m,h}\qquad(1\le i\le n_m)
\]
for a fixed \(R>0\).  Define
\[
  T_{m,h,i}:=
  X_{m,h,i}\mathbf 1_{\{\|X_{m,h,i}\|_{\op}\le b_{m,h,i}\}},
  \qquad
  \mu_{m,h,i}:=\mathbb E(T_{m,h,i}\mid\mathcal P_{m,h,i}),
\]
and
\[
  F^{\rm cap}_{m,h}:=
  \left\{\exists i:\|X_{m,h,i}\|_{\op}>b_{m,h,i}\right\}.
\]
Fix \(K>0\), \(\beta\ge0\), \(\eta\ge0\), and \(B>0\).  Assume
\[
  \sum_m\sum_{h=1}^{J_m}
  \bigl[\mathbb P(E_{m,h})+\mathbb P(F^{\rm cap}_{m,h})\bigr]<\infty
\]
and, outside \(E_{m,h}\),
\[
  \left\|\sum_{i=1}^{n_m}\mu_{m,h,i}\right\|_{\op}
  \le \eta q_{m,h}\sqrt{n_m}.
\]
In the self-adjoint case assume also, outside \(E_{m,h}\), that
\[
  \left\|\sum_i
  \mathbb E(T_{m,h,i}^2\mid\mathcal P_{m,h,i})\right\|_{\op}
  \le \beta q_{m,h}^2\,\frac{n_m}{\ell_m}.
\]
In the rectangular case replace this by the two row and column bounds
\[
\begin{aligned}
  \left\|\sum_i
  \mathbb E(T_{m,h,i}T_{m,h,i}^*\mid\mathcal P_{m,h,i})\right\|_{\op}
  &\le \beta q_{m,h}^2\,\frac{n_m}{\ell_m},\\
  \left\|\sum_i
  \mathbb E(T_{m,h,i}^*T_{m,h,i}\mid\mathcal P_{m,h,i})\right\|_{\op}
  &\le \beta q_{m,h}^2\,\frac{n_m}{\ell_m}.
\end{aligned}
\]
If \(\ell_m\le B\sqrt{n_m}\) eventually and
\[
  \sum_mJ_m\Delta_m
  (e\Delta_m)^{-\gamma_{\rm tc}(K,R,\beta,B)}<\infty,
\]
then, almost surely, for all sufficiently large \(m\), every active original
stream is contraction-normalized by
\[
  Mq_{m,h},
  \qquad
  M:=\max\{R,K+\eta\},
\]
and the all-plus signing gives
\[
  \disc(X_{m,h,1}/(Mq_{m,h}),\ldots,X_{m,h,n_m}/(Mq_{m,h}))
  \le\sqrt{n_m}
  \qquad(1\le h\le J_m).
\]

If \(J_m\Delta_m\le(e\Delta_m)^\Gamma\) eventually and
\(n_m\ge a\ell_m^2\), then the same closed-form margins as in the
tail-capped corollaries give the all-integer threshold
\[
  K>\Theta_{\rm tc}(\Gamma+1,a,\beta,R)
\]
and the stretched-exponential sparse threshold
\[
  K>\Theta_{\rm tc}(\Gamma,a,\beta,R).
\]

In particular, suppose the conditional sub-Weibull norm tails
\[
  \mathbb P\{\|X_{m,h,i}\|_{\op}>s_{m,h,i}u
  \mid\mathcal P_{m,h,i}\}
  \le A_*e^{-u^\alpha}
  \qquad(u>0)
\]
hold.  With \(L_m\ge1\), choose
\[
  b_{m,h,i}:=s_{m,h,i}L_m^{1/\alpha}
\]
and
\[
  q_{m,h}:=
  \max\left\{
  \max_i b_{m,h,i},
  \left(\frac{\ell_m}{n_m}\sum_i s_{m,h,i}^2\right)^{1/2}
  \right\}.
\]
Then the preceding theorem applies with \(\beta=C_\alpha A_*\), where
\(C_\alpha\) is an absolute constant depending only on \(\alpha\), whenever
the displayed drift bound holds and
\[
  \sum_m\sum_{h=1}^{J_m}
  \left[\mathbb P(E_{m,h})+A_*n_m e^{-L_m}\right]<\infty .
\]
\end{corollary}

\begin{proof}
Let \(Y_{m,h,i}=T_{m,h,i}-\mu_{m,h,i}\).  In both the independent and adapted
settings these are centered, respectively martingale-difference, capped
increments relative to \(\mathcal P_{m,h,i}\), and
\[
  \|Y_{m,h,i}\|_{\op}\le2b_{m,h,i}\le2Rq_{m,h}.
\]
The displayed uncentered second-moment profiles dominate the corresponding
centered variance profiles:
\[
  \mathbb E(Y_{m,h,i}^2\mid\mathcal P_{m,h,i})
  \preceq
  \mathbb E(T_{m,h,i}^2\mid\mathcal P_{m,h,i})
\]
in the self-adjoint case, and similarly with \(YY^*,TT^*\) and
\(Y^*Y,T^*T\) in the rectangular case.  Matrix Bernstein in the independent
case, and matrix Freedman in the adapted case, therefore give the same
all-plus tail used in
Corollaries~\ref{cor:tail-capped-centered-bernstein-all-plus-schedules} and
\ref{cor:tail-capped-freedman-all-plus-schedules}, with the common jump
\(2Rq_{m,h}\) and variance parameter \(\beta q_{m,h}^2n_m/\ell_m\).
Borel--Cantelli, the drift estimate, and the cap event then give the stated
contraction normalization and all-plus discrepancy bound.  The all-integer and
sparse thresholds are the same summability calculation with
\(B=a^{-1/2}\).

For the sub-Weibull specialization, the cap-failure probability is bounded by
\[
  \mathbb P(F^{\rm cap}_{m,h})
  \le A_*n_m e^{-L_m}.
\]
The same tail integration used in
Corollary~\ref{cor:subweibull-norm-tail-capped-all-plus-schedules} gives
\[
  \mathbb E(\|T_{m,h,i}\|_{\op}^2\mid\mathcal P_{m,h,i})
  \le C_\alpha A_*s_{m,h,i}^2.
\]
Hence the self-adjoint and rectangular uncentered row/column variance profiles
are at most \(C_\alpha A_*\sum_i s_{m,h,i}^2\), which is bounded by
\(C_\alpha A_*q_{m,h}^2n_m/\ell_m\) by the definition of \(q_{m,h}\).
\end{proof}

\begin{corollary}[Heterogeneous finite-moment norm-tail selection]
\label{cor:heterogeneous-moment-norm-tail-selection}
Work in the same independent or adapted setting, with the same conditioning
fields \(\mathcal P_{m,h,i}\), as in
Corollary~\ref{cor:heterogeneous-norm-tail-cap-profiles}.  Fix
\(A_*\ge1\), \(K>0\), \(\eta\ge0\), and \(B>0\).  For each active stream let
\(p_{m,h,i}>0\), \(s_{m,h,i}>0\), and \(U_{m,h,i}\ge1\) be deterministic, and
assume the conditional moment envelopes
\[
  \mathbb E(\|X_{m,h,i}\|_{\op}^{p_{m,h,i}}\mid\mathcal P_{m,h,i})
  \le A_*s_{m,h,i}^{p_{m,h,i}}
  \qquad(1\le i\le n_m).
\]
Put
\[
  b_{m,h,i}:=
  s_{m,h,i}U_{m,h,i}^{1/p_{m,h,i}},
  \qquad
  \rho_{m,h,i}:=
  s_{m,h,i}^2
  U_{m,h,i}^{(2-p_{m,h,i})_+/p_{m,h,i}},
\]
and choose
\[
  q_{m,h}:=
  \max\left\{
  \max_i b_{m,h,i},
  \left(\frac{\ell_m}{n_m}\sum_i\rho_{m,h,i}\right)^{1/2}
  \right\}.
\]
Define \(T_{m,h,i}\) and \(\mu_{m,h,i}\) using these individual caps as in
Corollary~\ref{cor:heterogeneous-norm-tail-cap-profiles}.  Outside events
\(E_{m,h}\), assume the truncated-tail drift bound
\[
  \left\|\sum_{i=1}^{n_m}\mu_{m,h,i}\right\|_{\op}
  \le \eta q_{m,h}\sqrt{n_m}.
\]
Suppose that \(\ell_m\le B\sqrt{n_m}\) eventually and that
\[
  \sum_m\sum_{h=1}^{J_m}
  \left[\mathbb P(E_{m,h})
  + A_*\sum_{i=1}^{n_m}U_{m,h,i}^{-1}\right]<\infty
\]
and
\[
  \sum_mJ_m\Delta_m
  (e\Delta_m)^{-\gamma_{\rm tc}(K,1,A_*,B)}<\infty .
\]
Then, almost surely, for all sufficiently large \(m\), every active original
stream is contraction-normalized by
\[
  Mq_{m,h},
  \qquad
  M:=\max\{1,K+\eta\},
\]
and the all-plus signing gives
\[
  \disc(X_{m,h,1}/(Mq_{m,h}),\ldots,X_{m,h,n_m}/(Mq_{m,h}))
  \le\sqrt{n_m}
  \qquad(1\le h\le J_m).
\]

In particular, assume \(U_{m,h,i}=U_m=(e\Delta_m)^c\) for every active
\((h,i)\), and suppose
\[
  J_m\Delta_m\le(e\Delta_m)^\Gamma,\qquad
  J_mn_m\le(e\Delta_m)^\Lambda,\qquad
  n_m\ge a\ell_m^2
\]
eventually.  Then the conclusion holds along all integer dimensions whenever
\[
  c>\Lambda+1,
  \qquad
  K>\Theta_{\rm tc}(\Gamma+1,a,A_*,1),
\]
and along every stretched-exponential sparse dimension schedule whenever
\[
  c>\Lambda,
  \qquad
  K>\Theta_{\rm tc}(\Gamma,a,A_*,1).
\]
Under these common-\(U_m\) schedules, write
\[
\begin{aligned}
  Q^{\rm cap}_{m,h}
  &:=
  \max_i s_{m,h,i}(e\Delta_m)^{c/p_{m,h,i}},\\
  Q^{\rm var}_{m,h}
  &:=
  \left(
    \frac{\ell_m}{n_m}
    \sum_i s_{m,h,i}^2
    (e\Delta_m)^{c(2-p_{m,h,i})_+/p_{m,h,i}}
  \right)^{1/2}.
\end{aligned}
\]
Then \(q_{m,h}=\max\{Q^{\rm cap}_{m,h},Q^{\rm var}_{m,h}\}\).
If \(p_{m,h,i}\ge1\) for all \(i\), and if the original streams are
conditionally centered in the independent case or martingale differences in
the adapted case, then the displayed drift hypothesis is automatic for every
fixed \(\eta>0\) under the same common-\(U_m\) closed-form cap choices.
\end{corollary}

\begin{proof}
Markov's inequality gives the cap-failure bound
\[
  \mathbb P(F^{\rm cap}_{m,h})
  \le A_*\sum_{i=1}^{n_m}U_{m,h,i}^{-1}.
\]
Also \(b_{m,h,i}\le q_{m,h}\), so
Corollary~\ref{cor:heterogeneous-norm-tail-cap-profiles} can be used with
\(R=1\) once the uncentered second-moment profiles are checked.  In the
self-adjoint case,
\[
  T_{m,h,i}^2\preceq\|T_{m,h,i}\|_{\op}^2I.
\]
If \(p_{m,h,i}<2\), then on the cap event
\[
  \|T_{m,h,i}\|_{\op}^2
  \le b_{m,h,i}^{2-p_{m,h,i}}
     \|X_{m,h,i}\|_{\op}^{p_{m,h,i}}.
\]
If \(p_{m,h,i}\ge2\), conditional Lyapunov monotonicity and \(A_*\ge1\) give
\[
  \mathbb E(\|T_{m,h,i}\|_{\op}^2\mid\mathcal P_{m,h,i})
  \le A_*s_{m,h,i}^2 .
\]
Thus in all cases
\[
  \mathbb E(T_{m,h,i}^2\mid\mathcal P_{m,h,i})
  \preceq A_*\rho_{m,h,i}I.
\]
The same scalar domination with \(TT^*\) and \(T^*T\) gives the rectangular
row and column profiles.  Hence each uncentered profile is bounded by
\[
  A_*\sum_i\rho_{m,h,i}
  \le A_*q_{m,h}^2\,\frac{n_m}{\ell_m},
\]
so the preceding heterogeneous cap-profile corollary applies with
\(\beta=A_*\).

For the common-\(U_m\) schedules, the cap-tail summand is bounded by
\((e\Delta_m)^{\Lambda-c}\), up to the fixed factor \(A_*\).  This is summable
over all integer dimensions when \(c>\Lambda+1\), and along every
stretched-exponential sparse schedule when \(c>\Lambda\).  The matrix-tail
summability is the same \(\Theta_{\rm tc}\) calculation with \(\beta=A_*\) and
\(R=1\).

Finally suppose \(p_{m,h,i}\ge1\) and the original variables are centered
relative to \(\mathcal P_{m,h,i}\).  Then, under the common-\(U_m\) cap choice,
\[
  \|\mu_{m,h,i}\|_{\op}
  \le
  A_*s_{m,h,i}U_m^{-1+1/p_{m,h,i}}
  =
  A_*U_m^{-1}b_{m,h,i}.
\]
Therefore
\[
  \frac{\sum_i\|\mu_{m,h,i}\|_{\op}}
       {q_{m,h}\sqrt{n_m}}
  \le
  A_*\sqrt{n_m}\,U_m^{-1}.
\]
Since \(n_m\le(e\Delta_m)^\Lambda\) and \(U_m=(e\Delta_m)^c\), this tends to
zero under either displayed cap-tail margin.  Hence the drift bound holds
eventually for every fixed \(\eta>0\).
\end{proof}

\begin{corollary}[Heterogeneous tail-profile norm-tail selection]
\label{cor:heterogeneous-tail-profile-norm-tail-selection}
Work in the same independent or adapted setting, with the same conditioning
fields \(\mathcal P_{m,h,i}\), as in
Corollary~\ref{cor:heterogeneous-norm-tail-cap-profiles}.  Fix
\(K>0\), \(\eta\ge0\), and \(B>0\).  For each active stream let
\(s_{m,h,i}>0\), \(u_{m,h,i}>0\), and deterministic measurable profiles
\(\pi_{m,h,i}:[0,\infty)\to[0,1]\) satisfy the conditional norm-tail bounds
\[
  \mathbb P\{\|X_{m,h,i}\|_{\op}>s_{m,h,i}t
  \mid\mathcal P_{m,h,i}\}
  \le \pi_{m,h,i}(t)
  \qquad(t\ge0).
\]
Set
\[
  b_{m,h,i}:=s_{m,h,i}u_{m,h,i},
  \qquad
  \rho_{m,h,i}:=
  2s_{m,h,i}^2\int_0^{u_{m,h,i}}t\,\pi_{m,h,i}(t)\,dt,
\]
and choose
\[
  q_{m,h}:=
  \max\left\{
  \max_i b_{m,h,i},
  \left(\frac{\ell_m}{n_m}\sum_i\rho_{m,h,i}\right)^{1/2}
  \right\}.
\]
Define \(T_{m,h,i}\) and \(\mu_{m,h,i}\) using these individual caps as in
Corollary~\ref{cor:heterogeneous-norm-tail-cap-profiles}.  Outside events
\(E_{m,h}\), assume the truncated-tail drift bound
\[
  \left\|\sum_{i=1}^{n_m}\mu_{m,h,i}\right\|_{\op}
  \le \eta q_{m,h}\sqrt{n_m}.
\]
Suppose that \(\ell_m\le B\sqrt{n_m}\) eventually and that
\[
  \sum_m\sum_{h=1}^{J_m}
  \left[\mathbb P(E_{m,h})
  +\sum_{i=1}^{n_m}\pi_{m,h,i}(u_{m,h,i})\right]<\infty
\]
and
\[
  \sum_mJ_m\Delta_m
  (e\Delta_m)^{-\gamma_{\rm tc}(K,1,1,B)}<\infty .
\]
Then, almost surely, for all sufficiently large \(m\), every active original
stream is contraction-normalized by
\[
  Mq_{m,h},
  \qquad
  M:=\max\{1,K+\eta\},
\]
and the all-plus signing gives
\[
  \disc(X_{m,h,1}/(Mq_{m,h}),\ldots,X_{m,h,n_m}/(Mq_{m,h}))
  \le\sqrt{n_m}
  \qquad(1\le h\le J_m).
\]

If \(J_m\Delta_m\le(e\Delta_m)^\Gamma\) eventually and
\(n_m\ge a\ell_m^2\), then the matrix-tail part is summable along all integer
dimensions whenever
\[
  K>\Theta_{\rm tc}(\Gamma+1,a,1,1),
\]
and along every stretched-exponential sparse dimension schedule whenever
\[
  K>\Theta_{\rm tc}(\Gamma,a,1,1),
\]
provided the displayed profile cap-tail series is summable.

Moreover, if the original streams are conditionally centered in the
independent case or martingale differences in the adapted case, then the drift
hypothesis is implied by the deterministic tail-drift profile bound
\[
  \sum_i s_{m,h,i}
  \left[
    u_{m,h,i}\pi_{m,h,i}(u_{m,h,i})
    +\int_{u_{m,h,i}}^\infty\pi_{m,h,i}(t)\,dt
  \right]
  \le \eta q_{m,h}\sqrt{n_m}
\]
outside \(E_{m,h}\).
\end{corollary}

\begin{proof}
The cap-failure event satisfies
\[
  \mathbb P(F^{\rm cap}_{m,h})
  \le\sum_{i=1}^{n_m}\pi_{m,h,i}(u_{m,h,i}).
\]
Also \(b_{m,h,i}\le q_{m,h}\), so it remains to verify the uncentered
second-moment profiles in
Corollary~\ref{cor:heterogeneous-norm-tail-cap-profiles} with \(R=1\) and
\(\beta=1\).  In the self-adjoint case,
\[
  T_{m,h,i}^2\preceq\|T_{m,h,i}\|_{\op}^2I.
\]
The conditional layer-cake formula and the displayed tail profile give
\[
\begin{aligned}
  \mathbb E(\|T_{m,h,i}\|_{\op}^2\mid\mathcal P_{m,h,i})
  &\le
  s_{m,h,i}^2
  \int_0^{u_{m,h,i}^2}
  \mathbb P\{\|X_{m,h,i}\|_{\op}^2>s_{m,h,i}^2r
  \mid\mathcal P_{m,h,i}\}\,dr\\
  &\le
  2s_{m,h,i}^2
  \int_0^{u_{m,h,i}}t\,\pi_{m,h,i}(t)\,dt
  =\rho_{m,h,i}.
\end{aligned}
\]
The same scalar domination gives the rectangular \(TT^*\) and \(T^*T\)
profiles.  Therefore every uncentered profile is bounded by
\[
  \sum_i\rho_{m,h,i}
  \le q_{m,h}^2\,\frac{n_m}{\ell_m},
\]
and the heterogeneous cap-profile corollary proves the general conclusion.
The two closed-form margins are exactly its all-integer and sparse
\(\Theta_{\rm tc}\) calculation with \(\beta=R=1\).

For the final assertion, conditional centering gives
\[
  \mu_{m,h,i}
  =
  -\mathbb E\!\left(
    X_{m,h,i}\mathbf 1_{\{\|X_{m,h,i}\|_{\op}>b_{m,h,i}\}}
    \mid\mathcal P_{m,h,i}
  \right).
\]
Another layer-cake estimate bounds
\[
  \|\mu_{m,h,i}\|_{\op}
  \le
  s_{m,h,i}
  \left[
    u_{m,h,i}\pi_{m,h,i}(u_{m,h,i})
    +\int_{u_{m,h,i}}^\infty\pi_{m,h,i}(t)\,dt
  \right].
\]
Summing in \(i\) gives the stated sufficient drift condition.
\end{proof}

\begin{corollary}[Quantile-budget norm-tail profile selection]
\label{cor:quantile-budget-norm-tail-profile-selection}
Work in the setting of
Corollary~\ref{cor:heterogeneous-tail-profile-norm-tail-selection}, with the
same deterministic profiles \(\pi_{m,h,i}\) and scales \(s_{m,h,i}\).  For
\(\delta\in(0,1]\), let \(Q_{m,h,i}(\delta)\) be any deterministic upper
\(\delta\)-tail level for the profile, meaning
\[
  \pi_{m,h,i}(Q_{m,h,i}(\delta))\le\delta .
\]
For each active stream choose a deterministic cap-failure budget
\(\varepsilon_{m,h}\in(0,1]\) and deterministic shares
\(\omega_{m,h,i}>0\) with \(\sum_i\omega_{m,h,i}\le1\).  Put
\[
  u_{m,h,i}:=
  Q_{m,h,i}(\varepsilon_{m,h}\omega_{m,h,i})
\]
and define
\[
\begin{aligned}
  b_{m,h,i}&:=s_{m,h,i}u_{m,h,i},\\
  \rho_{m,h,i}&:=
  2s_{m,h,i}^2\int_0^{u_{m,h,i}}t\,\pi_{m,h,i}(t)\,dt,\\
  q_{m,h}&:=
  \max\left\{
  \max_i b_{m,h,i},
  \left(\frac{\ell_m}{n_m}\sum_i\rho_{m,h,i}\right)^{1/2}
  \right\}.
\end{aligned}
\]
Assume the truncated-tail drift bound from
Corollary~\ref{cor:heterogeneous-tail-profile-norm-tail-selection}, with
these caps, holds outside \(E_{m,h}\).  If \(\ell_m\le B\sqrt{n_m}\)
eventually,
\[
  \sum_m\sum_{h=1}^{J_m}\bigl[\mathbb P(E_{m,h})+\varepsilon_{m,h}\bigr]
  <\infty,
\]
and
\[
  \sum_mJ_m\Delta_m
  (e\Delta_m)^{-\gamma_{\rm tc}(K,1,1,B)}<\infty,
\]
then the original streams are eventually contraction-normalized by
\(\max\{1,K+\eta\}q_{m,h}\), and the all-plus signing gives discrepancy at
most \(\sqrt{n_m}\), simultaneously for all active \(h\).

Consequently, if \(J_m\Delta_m\le(e\Delta_m)^\Gamma\),
\(J_m\le(e\Delta_m)^\Lambda\), \(n_m\ge a\ell_m^2\) eventually, and
\(\varepsilon_{m,h}=(e\Delta_m)^{-c}\), then the cap-budget part is summable
for all integer dimensions when \(c>\Lambda+1\), and along every
stretched-exponential sparse dimension schedule when \(c>\Lambda\).  The
corresponding matrix-tail margins are
\[
  K>\Theta_{\rm tc}(\Gamma+1,a,1,1)
\]
in the all-integer case and
\[
  K>\Theta_{\rm tc}(\Gamma,a,1,1)
\]
in the sparse case.  The same conclusion holds for any deterministic choice of
shares, including a near-minimizer of the displayed scale \(q_{m,h}\) over the
simplex \(\sum_i\omega_{m,h,i}\le1\).
\end{corollary}

\begin{proof}
The selected upper quantiles give
\[
  \sum_i\pi_{m,h,i}(u_{m,h,i})
  \le
  \varepsilon_{m,h}\sum_i\omega_{m,h,i}
  \le \varepsilon_{m,h}.
\]
Thus all hypotheses of
Corollary~\ref{cor:heterogeneous-tail-profile-norm-tail-selection} are exactly
the displayed assumptions here, and that corollary gives the general
conclusion.  If \(\varepsilon_{m,h}=(e\Delta_m)^{-c}\) and
\(J_m\le(e\Delta_m)^\Lambda\), then
\[
  \sum_{h=1}^{J_m}\varepsilon_{m,h}
  \le (e\Delta_m)^{\Lambda-c}.
\]
This series is summable over all integer dimensions when \(c>\Lambda+1\), and
over every stretched-exponential sparse schedule when \(c>\Lambda\).  The
matrix-tail part is precisely the \(\Theta_{\rm tc}\) summability calculation
from the tail-profile corollary.
\end{proof}

\begin{corollary}[Mixed sub-Weibull quantile-budget norm-tail selection]
\label{cor:mixed-subweibull-quantile-budget-norm-tail-selection}
Work in the setting of
Corollary~\ref{cor:quantile-budget-norm-tail-profile-selection}.  For each
active increment suppose that, for deterministic \(A_{m,h,i}\ge1\),
\(\alpha_{m,h,i}>0\), and \(s_{m,h,i}>0\),
\[
  \mathbb P\{\|X_{m,h,i}\|_{\op}>s_{m,h,i}t
  \mid\mathcal P_{m,h,i}\}
  \le
  \min\{1,A_{m,h,i}e^{-t^{\alpha_{m,h,i}}}\}
  \qquad(t\ge0).
\]
Choose a stream cap budget \(\varepsilon_{m,h}\in(0,1]\) and positive shares
\(\omega_{m,h,i}\) with \(\sum_i\omega_{m,h,i}\le1\).  Put
\[
  \lambda_{m,h,i}:=
  \log\frac{eA_{m,h,i}}{\varepsilon_{m,h}\omega_{m,h,i}},
  \qquad
  b_{m,h,i}:=
  s_{m,h,i}\lambda_{m,h,i}^{1/\alpha_{m,h,i}},
\]
and
\[
  \widehat\rho_{m,h,i}:=
  \frac{2A_{m,h,i}}{\alpha_{m,h,i}}
  \Gamma\!\left(\frac2{\alpha_{m,h,i}}\right)s_{m,h,i}^2 .
\]
Define the closed-form mixed sub-Weibull scale
\[
  q_{m,h}:=
  \max\left\{
  \max_i b_{m,h,i},
  \left(\frac{\ell_m}{n_m}
  \sum_i\widehat\rho_{m,h,i}\right)^{1/2}
  \right\}.
\]
Assume the truncated-tail drift bound from
Corollary~\ref{cor:heterogeneous-tail-profile-norm-tail-selection}, with
these caps, holds outside \(E_{m,h}\).  If \(\ell_m\le B\sqrt{n_m}\)
eventually,
\[
  \sum_m\sum_{h=1}^{J_m}\bigl[\mathbb P(E_{m,h})+\varepsilon_{m,h}\bigr]
  <\infty,
\]
and
\[
  \sum_mJ_m\Delta_m
  (e\Delta_m)^{-\gamma_{\rm tc}(K,1,1,B)}<\infty,
\]
then the original streams are eventually contraction-normalized by
\(\max\{1,K+\eta\}q_{m,h}\), and the all-plus signing gives discrepancy at
most \(\sqrt{n_m}\), simultaneously for all active \(h\).

If \(J_m\Delta_m\le(e\Delta_m)^\Gamma\),
\(J_m\le(e\Delta_m)^\Lambda\), \(n_m\ge a\ell_m^2\) eventually, and
\(\varepsilon_{m,h}=(e\Delta_m)^{-c}\), then the cap-budget part is summable
for all integer dimensions when \(c>\Lambda+1\), and along every
stretched-exponential sparse dimension schedule when \(c>\Lambda\).  The
matrix-tail margins remain
\[
  K>\Theta_{\rm tc}(\Gamma+1,a,1,1)
\]
and
\[
  K>\Theta_{\rm tc}(\Gamma,a,1,1),
\]
respectively.  In the conditionally centered independent case, or in the
adapted martingale-difference case, the drift hypothesis is implied by
\[
  \sum_i s_{m,h,i}
  \left[
    \lambda_{m,h,i}^{1/\alpha_{m,h,i}}
    \varepsilon_{m,h}\omega_{m,h,i}
    +
    A_{m,h,i}\int_{\lambda_{m,h,i}^{1/\alpha_{m,h,i}}}^\infty
    e^{-t^{\alpha_{m,h,i}}}\,dt
  \right]
  \le \eta q_{m,h}\sqrt{n_m}
\]
outside \(E_{m,h}\).  The shares may be optimized to reduce the maximum cap
term across mixed tail constants, scales, and exponents.
\end{corollary}

\begin{proof}
Let
\[
  \pi_{m,h,i}(t):=
  \min\{1,A_{m,h,i}e^{-t^{\alpha_{m,h,i}}}\}.
\]
With \(u_{m,h,i}:=\lambda_{m,h,i}^{1/\alpha_{m,h,i}}\),
\[
  \pi_{m,h,i}(u_{m,h,i})
  \le
  A_{m,h,i}e^{-\lambda_{m,h,i}}
  =
  e^{-1}\varepsilon_{m,h}\omega_{m,h,i}
  \le \varepsilon_{m,h}\omega_{m,h,i}.
\]
Thus these \(u_{m,h,i}\) are admissible upper quantiles for the
quantile-budget corollary.  Its exact variance profile satisfies
\[
  2s_{m,h,i}^2
  \int_0^{u_{m,h,i}}t\,\pi_{m,h,i}(t)\,dt
  \le
  2A_{m,h,i}s_{m,h,i}^2
  \int_0^\infty t e^{-t^{\alpha_{m,h,i}}}\,dt
  =
  \widehat\rho_{m,h,i}.
\]
Consequently the displayed \(q_{m,h}\) dominates the exact
tail-profile scale furnished by
Corollary~\ref{cor:quantile-budget-norm-tail-profile-selection}.  Applying
that corollary gives the general conclusion, and enlarging the normalization
from the exact scale to \(q_{m,h}\) preserves both contraction normalization
and the \(\sqrt{n_m}\) discrepancy bound.

The summability thresholds are the same cap-budget and matrix-tail calculations
as in the quantile-budget corollary.  Finally, in the centered or martingale
case, the tail-drift sufficient condition from
Corollary~\ref{cor:heterogeneous-tail-profile-norm-tail-selection} is bounded
above by the displayed mixed sub-Weibull drift sum because
\(\pi_{m,h,i}(u_{m,h,i})\le\varepsilon_{m,h}\omega_{m,h,i}\) and
\(\pi_{m,h,i}(t)\le A_{m,h,i}e^{-t^{\alpha_{m,h,i}}}\).
\end{proof}

\begin{corollary}[Optimized common-exponent sub-Weibull budget selection]
\label{cor:optimized-common-exponent-subweibull-budget-norm-tail-selection}
Work in the setting of
Corollary~\ref{cor:mixed-subweibull-quantile-budget-norm-tail-selection}, but
assume that the tail exponent is a fixed common \(\alpha>0\).  For each active
stream define the common-cap tail budget
\[
  \Phi_{m,h}(r):=
  \sum_{i=1}^{n_m}
  A_{m,h,i}\exp\left\{-\left(\frac{r}{s_{m,h,i}}\right)^\alpha\right\}
  \qquad(r\ge0).
\]
Choose a stream cap budget \(\varepsilon_{m,h}\in(0,1]\), and let
\(r_{m,h}\) be any deterministic cap satisfying
\[
  \Phi_{m,h}(r_{m,h})\le e^{-1}\varepsilon_{m,h}.
\]
In the optimized form, take the smallest such cap
\[
  r_{m,h}^*:=
  \inf\{r\ge0:\Phi_{m,h}(r)\le e^{-1}\varepsilon_{m,h}\}.
\]
Put \(u_{m,h,i}:=r_{m,h}/s_{m,h,i}\) and define the exact capped
tail-energy
\[
  V_{m,h}:=
  2\sum_i s_{m,h,i}^2
  \int_0^{u_{m,h,i}}
  t\,\min\{1,A_{m,h,i}e^{-t^\alpha}\}\,dt .
\]
It is always bounded by the closed-form full-tail energy
\[
  V_{m,h}\le
  \frac{2}{\alpha}\Gamma\!\left(\frac2\alpha\right)
  \sum_i A_{m,h,i}s_{m,h,i}^2 .
\]
Set
\[
  q_{m,h}:=
  \max\left\{
  r_{m,h},
  \left(\frac{\ell_m}{n_m}V_{m,h}\right)^{1/2}
  \right\}.
\]
Assume the truncated-tail drift bound from
Corollary~\ref{cor:heterogeneous-tail-profile-norm-tail-selection}, with the
common caps \(b_{m,h,i}=r_{m,h}\), holds outside \(E_{m,h}\).  If
\(\ell_m\le B\sqrt{n_m}\) eventually,
\[
  \sum_m\sum_{h=1}^{J_m}\bigl[\mathbb P(E_{m,h})+\varepsilon_{m,h}\bigr]
  <\infty,
\]
and
\[
  \sum_mJ_m\Delta_m
  (e\Delta_m)^{-\gamma_{\rm tc}(K,1,1,B)}<\infty,
\]
then the original streams are eventually contraction-normalized by
\(\max\{1,K+\eta\}q_{m,h}\), and the all-plus signing gives discrepancy at
most \(\sqrt{n_m}\), simultaneously for all active \(h\).

If \(J_m\Delta_m\le(e\Delta_m)^\Gamma\),
\(J_m\le(e\Delta_m)^\Lambda\), \(n_m\ge a\ell_m^2\) eventually, and
\(\varepsilon_{m,h}=(e\Delta_m)^{-c}\), then the cap-budget part is summable
for all integer dimensions when \(c>\Lambda+1\), and along every
stretched-exponential sparse dimension schedule when \(c>\Lambda\).  The
matrix-tail margins remain
\[
  K>\Theta_{\rm tc}(\Gamma+1,a,1,1)
\]
and
\[
  K>\Theta_{\rm tc}(\Gamma,a,1,1),
\]
respectively.

If the stream is conditionally centered independent, or is an adapted
martingale difference, the drift hypothesis is implied by
\[
  \sum_i\left[
    r_{m,h}A_{m,h,i}e^{-u_{m,h,i}^{\alpha}}
    +A_{m,h,i}s_{m,h,i}
    \int_{u_{m,h,i}}^\infty e^{-t^\alpha}\,dt
  \right]
  \le \eta q_{m,h}\sqrt{n_m}
\]
outside \(E_{m,h}\).

Moreover, \(r_{m,h}^*\) is the optimal common cap among all deterministic
budget-share choices in
Corollary~\ref{cor:mixed-subweibull-quantile-budget-norm-tail-selection}: a
share choice with every mixed sub-Weibull cap at most \(r\) exists if and only
if \(\Phi_{m,h}(r)\le e^{-1}\varepsilon_{m,h}\).
\end{corollary}

\begin{proof}
For fixed \(m,h\), set
\[
  \pi_i(t):=\min\{1,A_{m,h,i}e^{-t^\alpha}\},
  \qquad
  u_i:=r_{m,h}/s_{m,h,i}.
\]
Then
\[
  \sum_i\pi_i(u_i)
  \le
  \sum_i A_{m,h,i}e^{-u_i^\alpha}
  =
  \Phi_{m,h}(r_{m,h})
  \le e^{-1}\varepsilon_{m,h}
  \le \varepsilon_{m,h}.
\]
The displayed \(V_{m,h}\) is exactly the sum of the capped tail-energy terms
from Corollary~\ref{cor:heterogeneous-tail-profile-norm-tail-selection}.  That
corollary therefore gives the general conclusion with \(\beta=R=1\).  The
closed-form bound on \(V_{m,h}\) follows from
\[
  2s_{m,h,i}^2\int_0^{u_i}t\,\pi_i(t)\,dt
  \le
  2A_{m,h,i}s_{m,h,i}^2\int_0^\infty t e^{-t^\alpha}\,dt
  =
  \frac{2A_{m,h,i}}{\alpha}
  \Gamma\!\left(\frac2\alpha\right)s_{m,h,i}^2 .
\]
The summability thresholds are the same cap-budget and matrix-tail
calculations as in the quantile-budget corollary.  In the centered or
martingale case, the stated drift condition dominates the tail-drift bound in
Corollary~\ref{cor:heterogeneous-tail-profile-norm-tail-selection}, since
\(\pi_i(t)\le A_{m,h,i}e^{-t^\alpha}\) and \(s_{m,h,i}u_i=r_{m,h}\).

It remains only to record optimality of the common cap.  If shares
\(\omega_{m,h,i}>0\) with \(\sum_i\omega_{m,h,i}\le1\) produce mixed
sub-Weibull caps no larger than \(r\), then
\[
  \log\frac{eA_{m,h,i}}{\varepsilon_{m,h}\omega_{m,h,i}}
  \le
  \left(\frac{r}{s_{m,h,i}}\right)^\alpha,
\]
hence
\[
  \omega_{m,h,i}\ge
  \frac{eA_{m,h,i}}{\varepsilon_{m,h}}
  \exp\left\{-\left(\frac{r}{s_{m,h,i}}\right)^\alpha\right\}.
\]
Summing in \(i\) gives \(\Phi_{m,h}(r)\le e^{-1}\varepsilon_{m,h}\).
Conversely, if \(\Phi_{m,h}(r)\le e^{-1}\varepsilon_{m,h}\), choose
\[
  \omega_{m,h,i}:=
  \frac{eA_{m,h,i}}{\varepsilon_{m,h}}
  \exp\left\{-\left(\frac{r}{s_{m,h,i}}\right)^\alpha\right\}.
\]
Then \(\sum_i\omega_{m,h,i}\le1\), and the mixed sub-Weibull cap attached to
each \(i\) is exactly \(r\).  Thus the infimum \(r_{m,h}^*\) is the smallest
common cap available to the budget-share rule.
\end{proof}

\begin{corollary}[Closed-form common-exponent sub-Weibull schedules]
\label{cor:closed-form-common-exponent-subweibull-schedules}
Keep the hypotheses and notation of
Corollary~\ref{cor:optimized-common-exponent-subweibull-budget-norm-tail-selection}.
Set
\[
  A^\Sigma_{m,h}:=\sum_i A_{m,h,i},
  \qquad
  s^{\max}_{m,h}:=\max_i s_{m,h,i},
  \qquad
  C_\alpha:=\frac2\alpha\Gamma\!\left(\frac2\alpha\right).
\]
For a stream budget \(\varepsilon_{m,h}\in(0,1]\), define the explicit
envelope cap
\[
  r^{\rm env}_{m,h}:=
  s^{\max}_{m,h}
  \left(\log\frac{eA^\Sigma_{m,h}}{\varepsilon_{m,h}}\right)^{1/\alpha}
\]
and the closed variance envelope
\[
  V^{\rm env}_{m,h}:=
  C_\alpha\sum_iA_{m,h,i}s_{m,h,i}^2 .
\]
Then the scale
\[
  q^{\rm env}_{m,h}:=
  \max\left\{
  r^{\rm env}_{m,h},
  \left(\frac{\ell_m}{n_m}V^{\rm env}_{m,h}\right)^{1/2}
  \right\}
\]
dominates the optimized common-exponent scale.  Consequently, under the
summability and matrix-tail hypotheses of
Corollary~\ref{cor:optimized-common-exponent-subweibull-budget-norm-tail-selection},
with the truncated-tail drift checked at the common cap \(r^{\rm env}_{m,h}\),
the original streams are eventually contraction-normalized by
\(\max\{1,K+\eta\}q^{\rm env}_{m,h}\), and the all-plus signing gives
discrepancy at most \(\sqrt{n_m}\), simultaneously for all active \(h\).

If all scales in a stream are equal, \(s_{m,h,i}=s_{m,h}\), then
\[
  r^{\rm env}_{m,h}=
  s_{m,h}
  \left(\log\frac{eA^\Sigma_{m,h}}{\varepsilon_{m,h}}\right)^{1/\alpha}
\]
is the optimized common cap \(r^*_{m,h}\).  In particular, if
\(A^\Sigma_{m,h}\le(e\Delta_m)^\chi\) and
\(\varepsilon_{m,h}=(e\Delta_m)^{-c}\), then
\[
  r^{\rm env}_{m,h}
  \le
  s^{\max}_{m,h}\bigl((\chi+c+1)\ell_m\bigr)^{1/\alpha}.
\]
\end{corollary}

\begin{proof}
For every \(i\),
\[
  \left(\frac{r^{\rm env}_{m,h}}{s_{m,h,i}}\right)^\alpha
  \ge
  \left(\frac{r^{\rm env}_{m,h}}{s^{\max}_{m,h}}\right)^\alpha
  =
  \log\frac{eA^\Sigma_{m,h}}{\varepsilon_{m,h}}.
\]
Hence the common-cap budget from
Corollary~\ref{cor:optimized-common-exponent-subweibull-budget-norm-tail-selection}
satisfies
\[
  \Phi_{m,h}(r^{\rm env}_{m,h})
  \le
  A^\Sigma_{m,h}
  \exp\left\{-\log\frac{eA^\Sigma_{m,h}}{\varepsilon_{m,h}}\right\}
  =
  e^{-1}\varepsilon_{m,h}.
\]
The full-tail estimate in that corollary gives
\[
  V_{m,h}\le V^{\rm env}_{m,h}.
\]
Thus \(q^{\rm env}_{m,h}\) dominates the admissible common-exponent
tail-profile scale, and the stated conclusion is exactly the previous
corollary applied with this explicit cap and enlarged variance envelope.

When all \(s_{m,h,i}\) are equal to \(s_{m,h}\),
\[
  \Phi_{m,h}(r)=A^\Sigma_{m,h}e^{-(r/s_{m,h})^\alpha},
\]
so the equality
\(\Phi_{m,h}(r)=e^{-1}\varepsilon_{m,h}\) is solved by the displayed
\(r^{\rm env}_{m,h}\).  This is the minimal admissible common cap by the
optimality assertion in the previous corollary.  The final growth bound follows
from
\[
  \log\frac{eA^\Sigma_{m,h}}{\varepsilon_{m,h}}
  \le
  (\chi+c+1)\log(e\Delta_m)
  =
  (\chi+c+1)\ell_m .
\]
\end{proof}

\begin{corollary}[Uniform common-exponent sub-Weibull schedules]
\label{cor:uniform-common-exponent-subweibull-schedules}
Work in the setting of
Corollary~\ref{cor:optimized-common-exponent-subweibull-budget-norm-tail-selection}.
Assume that the sub-Weibull exponent is a fixed common \(\alpha>0\), and that
for each active stream there are deterministic \(A_{m,h}\ge1\) and
\(s_{m,h}>0\) such that, for every increment and every \(t\ge0\),
\[
  \mathbb P\{\|X_{m,h,i}\|_{\op}>t\mid\mathcal P_{m,h,i}\}
  \le
  A_{m,h}\exp\{-(t/s_{m,h})^\alpha\}.
\]
Put \(C_\alpha=(2/\alpha)\Gamma(2/\alpha)\).  For a stream budget
\(\varepsilon_{m,h}\in(0,1]\), define
\[
  r^{\psi,{\rm unif}}_{m,h}:=
  s_{m,h}
  \left(\log\frac{e n_mA_{m,h}}{\varepsilon_{m,h}}\right)^{1/\alpha}
\]
and
\[
  V^{\psi,{\rm unif}}_{m,h}:=
  C_\alpha n_mA_{m,h}s_{m,h}^2 .
\]
Set
\[
  q^{\psi,{\rm unif}}_{m,h}:=
  \max\left\{
  r^{\psi,{\rm unif}}_{m,h},
  \left(\frac{\ell_m}{n_m}V^{\psi,{\rm unif}}_{m,h}\right)^{1/2}
  \right\}.
\]
Under the summability and matrix-tail hypotheses of
Corollary~\ref{cor:optimized-common-exponent-subweibull-budget-norm-tail-selection},
with the truncated-tail drift checked at the common cap
\(r^{\psi,{\rm unif}}_{m,h}\), the original streams are eventually
contraction-normalized by
\(\max\{1,K+\eta\}q^{\psi,{\rm unif}}_{m,h}\), and the all-plus signing gives
discrepancy at most \(\sqrt{n_m}\), simultaneously for all active \(h\).

In particular, if \(\varepsilon_{m,h}=(e\Delta_m)^{-c}\),
\[
  A_{m,h}\le(e\Delta_m)^\chi,
  \qquad
  n_m\le(e\Delta_m)^\nu,
\]
then
\[
  r^{\psi,{\rm unif}}_{m,h}
  \le
  s_{m,h}\bigl((\nu+\chi+c+1)\ell_m\bigr)^{1/\alpha}
\]
and
\[
  V^{\psi,{\rm unif}}_{m,h}
  \le
  C_\alpha s_{m,h}^2(e\Delta_m)^{\nu+\chi}.
\]
Equivalently,
\[
  q^{\psi,{\rm unif}}_{m,h}
  \le
  s_{m,h}\max\left\{
  \bigl((\nu+\chi+c+1)\ell_m\bigr)^{1/\alpha},
  \left(C_\alpha\ell_m(e\Delta_m)^{\nu+\chi}/n_m\right)^{1/2}
  \right\}.
\]
\end{corollary}

\begin{proof}
Apply
Corollary~\ref{cor:closed-form-common-exponent-subweibull-schedules} after
choosing the admissible tail parameters
\(A_{m,h,i}=A_{m,h}\) and \(s_{m,h,i}=s_{m,h}\) for every \(i\).  Then
\[
  A^\Sigma_{m,h}=n_mA_{m,h},
  \qquad
  s^{\max}_{m,h}=s_{m,h},
\]
and the closed-form envelope in that corollary is exactly the displayed
\(r^{\psi,{\rm unif}}_{m,h}\), \(V^{\psi,{\rm unif}}_{m,h}\), and
\(q^{\psi,{\rm unif}}_{m,h}\).  The eventual normalization and discrepancy
statement follow.

The growth estimates follow from
\[
  \log\frac{e n_mA_{m,h}}{\varepsilon_{m,h}}
  \le
  (\nu+\chi+c+1)\ell_m
\]
and \(n_mA_{m,h}\le(e\Delta_m)^{\nu+\chi}\).  Substituting these bounds into
the definition of \(q^{\psi,{\rm unif}}_{m,h}\) gives the final display.
\end{proof}

\begin{corollary}[Mixed finite-moment quantile-budget norm-tail selection]
\label{cor:mixed-moment-quantile-budget-norm-tail-selection}
Work in the same independent or adapted setting as
Corollary~\ref{cor:quantile-budget-norm-tail-profile-selection}.  For each
active increment suppose that, for deterministic \(A_{m,h,i}\ge1\),
\(p_{m,h,i}>0\), and \(s_{m,h,i}>0\),
\[
  \mathbb E\left(\|X_{m,h,i}\|_{\op}^{p_{m,h,i}}
  \mid\mathcal P_{m,h,i}\right)
  \le A_{m,h,i}s_{m,h,i}^{p_{m,h,i}} .
\]
Choose a stream cap budget \(\varepsilon_{m,h}\in(0,1]\) and positive shares
\(\omega_{m,h,i}\) with \(\sum_i\omega_{m,h,i}\le1\).  Put
\[
  \lambda_{m,h,i}:=
  \frac{eA_{m,h,i}}{\varepsilon_{m,h}\omega_{m,h,i}},
  \qquad
  b_{m,h,i}:=
  s_{m,h,i}\lambda_{m,h,i}^{1/p_{m,h,i}},
\]
and define the moment-energy proxy
\[
  \widehat\rho_{m,h,i}:=
  \begin{cases}
  A_{m,h,i}s_{m,h,i}^2
  \lambda_{m,h,i}^{(2-p_{m,h,i})/p_{m,h,i}},
  &0<p_{m,h,i}<2,\\[3pt]
  A_{m,h,i}^{2/p_{m,h,i}}s_{m,h,i}^2,
  &p_{m,h,i}\ge2.
  \end{cases}
\]
Set
\[
  q_{m,h}:=
  \max\left\{
  \max_i b_{m,h,i},
  \left(\frac{\ell_m}{n_m}
  \sum_i\widehat\rho_{m,h,i}\right)^{1/2}
  \right\}.
\]
Assume the truncated-tail drift bound from
Corollary~\ref{cor:heterogeneous-norm-tail-cap-profiles}, with these caps,
holds outside \(E_{m,h}\).  If \(\ell_m\le B\sqrt{n_m}\) eventually,
\[
  \sum_m\sum_{h=1}^{J_m}\bigl[\mathbb P(E_{m,h})+\varepsilon_{m,h}\bigr]
  <\infty,
\]
and
\[
  \sum_mJ_m\Delta_m
  (e\Delta_m)^{-\gamma_{\rm tc}(K,1,1,B)}<\infty,
\]
then the original streams are eventually contraction-normalized by
\(\max\{1,K+\eta\}q_{m,h}\), and the all-plus signing gives discrepancy at
most \(\sqrt{n_m}\), simultaneously for all active \(h\).

If \(J_m\Delta_m\le(e\Delta_m)^\Gamma\),
\(J_m\le(e\Delta_m)^\Lambda\), \(n_m\ge a\ell_m^2\) eventually, and
\(\varepsilon_{m,h}=(e\Delta_m)^{-c}\), then the cap-budget part is summable
for all integer dimensions when \(c>\Lambda+1\), and along every
stretched-exponential sparse dimension schedule when \(c>\Lambda\).  The
matrix-tail margins remain
\[
  K>\Theta_{\rm tc}(\Gamma+1,a,1,1)
\]
and
\[
  K>\Theta_{\rm tc}(\Gamma,a,1,1),
\]
respectively.

If every \(p_{m,h,i}>1\), then in the conditionally centered independent case,
or in the adapted martingale-difference case, the drift hypothesis is implied
by
\[
  \sum_iA_{m,h,i}s_{m,h,i}
  \lambda_{m,h,i}^{-(p_{m,h,i}-1)/p_{m,h,i}}
  \le \eta q_{m,h}\sqrt{n_m}
\]
outside \(E_{m,h}\).  The shares may be optimized to trade off the largest cap
against the aggregate truncated second-moment energy for mixed moment orders.
\end{corollary}

\begin{proof}
Markov's inequality gives
\[
  \mathbb P\{\|X_{m,h,i}\|_{\op}>b_{m,h,i}\mid\mathcal P_{m,h,i}\}
  \le
  A_{m,h,i}\lambda_{m,h,i}^{-1}
  =
  e^{-1}\varepsilon_{m,h}\omega_{m,h,i}
  \le \varepsilon_{m,h}\omega_{m,h,i}.
\]
Thus the total cap-failure probability is at most \(\varepsilon_{m,h}\).
Also \(b_{m,h,i}\le q_{m,h}\), so the heterogeneous cap-profile theorem will
apply with \(R=1\) once the uncentered second-moment profile is bounded.

For \(0<p_{m,h,i}<2\), on the capped event,
\[
  \|T_{m,h,i}\|_{\op}^2
  \le b_{m,h,i}^{2-p_{m,h,i}}
  \|X_{m,h,i}\|_{\op}^{p_{m,h,i}},
\]
and hence
\[
  \mathbb E(\|T_{m,h,i}\|_{\op}^2\mid\mathcal P_{m,h,i})
  \le \widehat\rho_{m,h,i}.
\]
For \(p_{m,h,i}\ge2\), conditional Lyapunov monotonicity gives the same bound:
\[
  \mathbb E(\|T_{m,h,i}\|_{\op}^2\mid\mathcal P_{m,h,i})
  \le
  \left[
  \mathbb E(\|X_{m,h,i}\|_{\op}^{p_{m,h,i}}
  \mid\mathcal P_{m,h,i})
  \right]^{2/p_{m,h,i}}
  \le \widehat\rho_{m,h,i}.
\]
In the self-adjoint case \(T_{m,h,i}^2\preceq\|T_{m,h,i}\|_{\op}^2I\), and in
the rectangular case the same scalar bound applies to \(T_{m,h,i}T_{m,h,i}^*\)
and \(T_{m,h,i}^*T_{m,h,i}\).  Therefore every uncentered second-moment
profile is bounded by
\[
  \sum_i\widehat\rho_{m,h,i}
  \le q_{m,h}^2\,\frac{n_m}{\ell_m}.
\]
Corollary~\ref{cor:heterogeneous-norm-tail-cap-profiles}, with
\(\beta=R=1\), gives the general conclusion.  The cap-budget and matrix-tail
summability thresholds are the same as in the quantile-budget corollary.

Finally, if \(p_{m,h,i}>1\) and the original variables are centered relative
to \(\mathcal P_{m,h,i}\), then
\[
\begin{aligned}
  \|\mu_{m,h,i}\|_{\op}
  &\le
  \mathbb E\Bigl(
  \|X_{m,h,i}\|_{\op}
  \mathbf 1_{\{\|X_{m,h,i}\|_{\op}>b_{m,h,i}\}}
  \,\Bigm|\,\mathcal P_{m,h,i}\Bigr)\\
  &\le
  A_{m,h,i}s_{m,h,i}
  \lambda_{m,h,i}^{-(p_{m,h,i}-1)/p_{m,h,i}}.
\end{aligned}
\]
Summing gives the displayed sufficient drift bound.
\end{proof}

\begin{corollary}[Optimized common-moment quantile-budget selection]
\label{cor:optimized-common-moment-quantile-budget-norm-tail-selection}
Work in the setting of
Corollary~\ref{cor:mixed-moment-quantile-budget-norm-tail-selection}, but
assume that the moment order is a fixed common \(p>0\).  For each active stream
put
\[
  S_{m,h}:=\sum_{i=1}^{n_m}A_{m,h,i}s_{m,h,i}^p
\]
and choose a stream cap budget \(\varepsilon_{m,h}\in(0,1]\).  Define the
equalized Markov cap
\[
  r_{m,h}:=\left(\frac{eS_{m,h}}{\varepsilon_{m,h}}\right)^{1/p}
\]
and the aggregate moment-energy proxy
\[
  V_{m,h}:=
  \begin{cases}
  \left(\dfrac{e}{\varepsilon_{m,h}}\right)^{(2-p)/p}
  S_{m,h}^{2/p},&0<p<2,\\[6pt]
  \displaystyle\sum_iA_{m,h,i}^{2/p}s_{m,h,i}^2,&p\ge2 .
  \end{cases}
\]
Set
\[
  q_{m,h}:=
  \max\left\{
  r_{m,h},
  \left(\frac{\ell_m}{n_m}V_{m,h}\right)^{1/2}
  \right\}.
\]
Assume the truncated-tail drift bound from
Corollary~\ref{cor:heterogeneous-norm-tail-cap-profiles}, with the common caps
\(b_{m,h,i}=r_{m,h}\), holds outside \(E_{m,h}\).  If
\(\ell_m\le B\sqrt{n_m}\) eventually,
\[
  \sum_m\sum_{h=1}^{J_m}\bigl[\mathbb P(E_{m,h})+\varepsilon_{m,h}\bigr]
  <\infty,
\]
and
\[
  \sum_mJ_m\Delta_m
  (e\Delta_m)^{-\gamma_{\rm tc}(K,1,1,B)}<\infty,
\]
then the original streams are eventually contraction-normalized by
\(\max\{1,K+\eta\}q_{m,h}\), and the all-plus signing gives discrepancy at
most \(\sqrt{n_m}\), simultaneously for all active \(h\).

If \(J_m\Delta_m\le(e\Delta_m)^\Gamma\),
\(J_m\le(e\Delta_m)^\Lambda\), \(n_m\ge a\ell_m^2\) eventually, and
\(\varepsilon_{m,h}=(e\Delta_m)^{-c}\), then the cap-budget part is summable
for all integer dimensions when \(c>\Lambda+1\), and along every
stretched-exponential sparse dimension schedule when \(c>\Lambda\).  The
matrix-tail margins remain
\[
  K>\Theta_{\rm tc}(\Gamma+1,a,1,1)
\]
and
\[
  K>\Theta_{\rm tc}(\Gamma,a,1,1),
\]
respectively.

If \(p>1\), then in the conditionally centered independent case, or in the
adapted martingale-difference case, the drift hypothesis is implied by
\[
  S_{m,h}r_{m,h}^{1-p}\le \eta q_{m,h}\sqrt{n_m}
\]
outside \(E_{m,h}\).  Equivalently, the optimal budget shares are
\[
  \omega_{m,h,i}:=\frac{A_{m,h,i}s_{m,h,i}^p}{S_{m,h}},
\]
for which every cap \(b_{m,h,i}\) in
Corollary~\ref{cor:mixed-moment-quantile-budget-norm-tail-selection} is equal
to \(r_{m,h}\).
\end{corollary}

\begin{proof}
Apply Corollary~\ref{cor:mixed-moment-quantile-budget-norm-tail-selection} with
\[
  \omega_{m,h,i}=
  \frac{A_{m,h,i}s_{m,h,i}^p}{S_{m,h}}.
\]
Then \(\sum_i\omega_{m,h,i}=1\) and
\[
  \lambda_{m,h,i}
  =
  \frac{eA_{m,h,i}}{\varepsilon_{m,h}\omega_{m,h,i}}
  =
  \frac{eS_{m,h}}{\varepsilon_{m,h}s_{m,h,i}^p},
\]
so
\[
  b_{m,h,i}=s_{m,h,i}\lambda_{m,h,i}^{1/p}=r_{m,h}
  \qquad(1\le i\le n_m).
\]
If \(0<p<2\), the mixed moment-energy terms sum to
\[
\begin{aligned}
  \sum_iA_{m,h,i}s_{m,h,i}^2
  \lambda_{m,h,i}^{(2-p)/p}
  &=
  \left(\frac{eS_{m,h}}{\varepsilon_{m,h}}\right)^{(2-p)/p}
  \sum_iA_{m,h,i}s_{m,h,i}^{p}\\
  &=
  \left(\frac{e}{\varepsilon_{m,h}}\right)^{(2-p)/p}
  S_{m,h}^{2/p}.
\end{aligned}
\]
For \(p\ge2\), the mixed theorem gives the displayed sum
\(\sum_iA_{m,h,i}^{2/p}s_{m,h,i}^2\).  Thus its scale is exactly the
displayed \(q_{m,h}\), and all conclusions and summability thresholds follow
from that corollary.  The drift bound is the mixed finite-moment drift
specialized to the same shares:
\[
  \sum_iA_{m,h,i}s_{m,h,i}\lambda_{m,h,i}^{-(p-1)/p}
  =
  S_{m,h}r_{m,h}^{1-p}.
\]
\end{proof}

\begin{corollary}[Closed-form common-moment growth schedules]
\label{cor:closed-form-common-moment-growth-schedules}
Keep the hypotheses and notation of
Corollary~\ref{cor:optimized-common-moment-quantile-budget-norm-tail-selection}.
Let \(\bar S_{m,h}\) be any deterministic envelope with
\[
  S_{m,h}\le \bar S_{m,h}.
\]
When \(p\ge2\), also let \(\bar W_{m,h}\) satisfy
\[
  \sum_iA_{m,h,i}^{2/p}s_{m,h,i}^2\le \bar W_{m,h}.
\]
For a stream budget \(\varepsilon_{m,h}\in(0,1]\), set
\[
  r^{\rm env}_{m,h}:=
  \left(\frac{e\bar S_{m,h}}{\varepsilon_{m,h}}\right)^{1/p}
\]
and
\[
  V^{\rm env}_{m,h}:=
  \begin{cases}
  \left(\dfrac{e}{\varepsilon_{m,h}}\right)^{(2-p)/p}
  \bar S_{m,h}^{2/p},&0<p<2,\\[6pt]
  \bar W_{m,h},&p\ge2 .
  \end{cases}
\]
Define
\[
  q^{\rm env}_{m,h}:=
  \max\left\{
  r^{\rm env}_{m,h},
  \left(\frac{\ell_m}{n_m}V^{\rm env}_{m,h}\right)^{1/2}
  \right\}.
\]
Under the summability and matrix-tail hypotheses of
Corollary~\ref{cor:optimized-common-moment-quantile-budget-norm-tail-selection},
with the truncated-tail drift checked for the exact cap \(r_{m,h}\), the
original streams are eventually contraction-normalized by
\(\max\{1,K+\eta\}q^{\rm env}_{m,h}\), and the all-plus signing gives
discrepancy at most \(\sqrt{n_m}\), simultaneously for all active \(h\).
If \(p>1\), the centered independent or martingale-difference drift condition
is implied by the envelope bound
\[
  \bar S_{m,h}\bigl(r^{\rm env}_{m,h}\bigr)^{1-p}
  \le \eta q^{\rm env}_{m,h}\sqrt{n_m}.
\]

In particular, if \(\varepsilon_{m,h}=(e\Delta_m)^{-c}\) and
\(\bar S_{m,h}\le(e\Delta_m)^\chi\), then
\[
  r^{\rm env}_{m,h}\le
  (e\Delta_m)^{(\chi+c+1)/p}.
\]
For \(0<p<2\),
\[
  V^{\rm env}_{m,h}
  \le
  (e\Delta_m)^{((2-p)(c+1)+2\chi)/p},
\]
whereas for \(p\ge2\), if \(\bar W_{m,h}\le(e\Delta_m)^\upsilon\), then
\[
  V^{\rm env}_{m,h}\le(e\Delta_m)^\upsilon .
\]
\end{corollary}

\begin{proof}
The exact common-moment cap from
Corollary~\ref{cor:optimized-common-moment-quantile-budget-norm-tail-selection}
satisfies
\[
  r_{m,h}=
  \left(\frac{eS_{m,h}}{\varepsilon_{m,h}}\right)^{1/p}
  \le r^{\rm env}_{m,h}.
\]
For \(0<p<2\), the exact energy is monotone in \(S_{m,h}\), hence is at most
the displayed \(V^{\rm env}_{m,h}\).  For \(p\ge2\), this is exactly the
assumption defining \(\bar W_{m,h}\).  Thus \(q^{\rm env}_{m,h}\) dominates
the exact optimized common-moment scale, so the previous corollary gives the
same eventual normalization after enlarging the scale.

If \(p>1\), then
\[
  S_{m,h}r_{m,h}^{1-p}
  =
  S_{m,h}^{1/p}
  \left(\frac{e}{\varepsilon_{m,h}}\right)^{-(p-1)/p}
  \le
  \bar S_{m,h}^{1/p}
  \left(\frac{e}{\varepsilon_{m,h}}\right)^{-(p-1)/p}
  =
  \bar S_{m,h}(r^{\rm env}_{m,h})^{1-p},
\]
so the envelope drift condition implies the exact drift condition.  The final
growth displays follow from
\[
  \frac{e}{\varepsilon_{m,h}}\le(e\Delta_m)^{c+1}
  \qquad\text{and}\qquad
  \bar S_{m,h}\le(e\Delta_m)^\chi .
\]
\end{proof}

\begin{corollary}[Uniform common-moment operator-tail schedules]
\label{cor:uniform-common-moment-operator-tail-schedules}
Work in the setting of
Corollary~\ref{cor:optimized-common-moment-quantile-budget-norm-tail-selection}.
Assume that the moment order is a fixed common \(p>0\), and that for each
active stream there are deterministic \(A_{m,h}\ge1\) and \(s_{m,h}>0\) such
that, for every increment,
\[
  \mathbb E\left(\|X_{m,h,i}\|_{\op}^p
  \mid\mathcal P_{m,h,i}\right)
  \le A_{m,h}s_{m,h}^p .
\]
For a stream budget \(\varepsilon_{m,h}\in(0,1]\), define
\[
  r^{\rm unif}_{m,h}:=
  s_{m,h}\left(\frac{e n_mA_{m,h}}{\varepsilon_{m,h}}\right)^{1/p}
\]
and
\[
  V^{\rm unif}_{m,h}:=
  \begin{cases}
  s_{m,h}^2
  \left(\dfrac{e}{\varepsilon_{m,h}}\right)^{(2-p)/p}
  (n_mA_{m,h})^{2/p},&0<p<2,\\[6pt]
  n_mA_{m,h}^{2/p}s_{m,h}^2,&p\ge2 .
  \end{cases}
\]
Set
\[
  q^{\rm unif}_{m,h}:=
  \max\left\{
  r^{\rm unif}_{m,h},
  \left(\frac{\ell_m}{n_m}V^{\rm unif}_{m,h}\right)^{1/2}
  \right\}.
\]
Under the summability and matrix-tail hypotheses of
Corollary~\ref{cor:optimized-common-moment-quantile-budget-norm-tail-selection},
with the truncated-tail drift checked at the common cap
\(r^{\rm unif}_{m,h}\), the original streams are eventually
contraction-normalized by \(\max\{1,K+\eta\}q^{\rm unif}_{m,h}\), and the
all-plus signing gives discrepancy at most \(\sqrt{n_m}\), simultaneously for
all active \(h\).

If \(p>1\), then in the centered independent or martingale-difference case the
drift condition is implied by
\[
  n_mA_{m,h}s_{m,h}^p
  \bigl(r^{\rm unif}_{m,h}\bigr)^{1-p}
  =
  s_{m,h}(n_mA_{m,h})^{1/p}
  \left(\frac{e}{\varepsilon_{m,h}}\right)^{-(p-1)/p}
  \le \eta q^{\rm unif}_{m,h}\sqrt{n_m}.
\]
In particular, if \(\varepsilon_{m,h}=(e\Delta_m)^{-c}\),
\[
  A_{m,h}\le(e\Delta_m)^\chi,
  \qquad
  n_m\le(e\Delta_m)^\nu,
\]
then
\[
  r^{\rm unif}_{m,h}
  \le
  s_{m,h}(e\Delta_m)^{(\nu+\chi+c+1)/p}.
\]
For \(0<p<2\),
\[
  V^{\rm unif}_{m,h}
  \le
  s_{m,h}^2
  (e\Delta_m)^{((2-p)(c+1)+2(\nu+\chi))/p},
\]
whereas for \(p\ge2\),
\[
  V^{\rm unif}_{m,h}
  \le
  s_{m,h}^2(e\Delta_m)^{\nu+2\chi/p}.
\]
\end{corollary}

\begin{proof}
Apply
Corollary~\ref{cor:closed-form-common-moment-growth-schedules} after choosing
the admissible moment parameters
\(A_{m,h,i}=A_{m,h}\) and \(s_{m,h,i}=s_{m,h}\) for every \(i\).  Then
\[
  S_{m,h}=n_mA_{m,h}s_{m,h}^p .
\]
When \(p\ge2\), also
\[
  \sum_iA_{m,h,i}^{2/p}s_{m,h,i}^2
  =
  n_mA_{m,h}^{2/p}s_{m,h}^2 .
\]
Use these exact aggregates as \(\bar S_{m,h}\) and \(\bar W_{m,h}\) in that
corollary.  The definitions of
\(r^{\rm unif}_{m,h}\), \(V^{\rm unif}_{m,h}\), and \(q^{\rm unif}_{m,h}\) are
exactly the envelope definitions from that corollary with these choices, so
the eventual normalization and discrepancy conclusion follow.

For \(p>1\), the envelope drift condition from the previous corollary becomes
\[
  \bar S_{m,h}
  \bigl(r^{\rm unif}_{m,h}\bigr)^{1-p}
  \le \eta q^{\rm unif}_{m,h}\sqrt{n_m}.
\]
Substituting \(\bar S_{m,h}=n_mA_{m,h}s_{m,h}^p\) gives the displayed identity.
The three growth estimates follow from
\[
  \frac{e}{\varepsilon_{m,h}}\le(e\Delta_m)^{c+1},
  \qquad
  n_mA_{m,h}\le(e\Delta_m)^{\nu+\chi},
  \qquad
  A_{m,h}^{2/p}\le(e\Delta_m)^{2\chi/p}.
\]
\end{proof}

\begin{corollary}[Mixed dominated operator-tail active libraries]
\label{cor:mixed-dominated-operator-tail-active-libraries}
For each \(m\ge1\), let \(D_m,n_m,J_m\ge1\), put
\(\Delta_m:=D_m+n_m+1\) and \(\ell_m:=\log(e\Delta_m)\), and let each active
stream \(1\le h\le J_m\) be tagged either \(\psi\) or \(\mu\).  Choose
\(\varepsilon_{m,h}\in(0,1]\), \(A_{m,h}\ge1\), and \(s_{m,h}>0\).  If
\(h\) is tagged \(\psi\), assume that for some \(\alpha_{m,h}>0\),
\[
  \mathbb P\{\|X_{m,h,i}\|_{\op}>t\mid\mathcal P_{m,h,i}\}
  \le
  A_{m,h}\exp\{-(t/s_{m,h})^{\alpha_{m,h}}\}
  \qquad(t\ge0)
\]
for every increment.  Set
\[
  C_{\alpha_{m,h}}:=
  \frac2{\alpha_{m,h}}\Gamma\!\left(\frac2{\alpha_{m,h}}\right),
  \qquad
  r_{m,h}:=
  s_{m,h}
  \left(\log\frac{e n_mA_{m,h}}{\varepsilon_{m,h}}\right)^{1/\alpha_{m,h}},
\]
and
\[
  V_{m,h}:=C_{\alpha_{m,h}}n_mA_{m,h}s_{m,h}^2 .
\]
If \(h\) is tagged \(\mu\), assume that for some \(p_{m,h}>0\),
\[
  \mathbb E\left(\|X_{m,h,i}\|_{\op}^{p_{m,h}}
  \mid\mathcal P_{m,h,i}\right)
  \le A_{m,h}s_{m,h}^{p_{m,h}}
\]
for every increment.  Set
\[
  r_{m,h}:=
  s_{m,h}\left(\frac{e n_mA_{m,h}}{\varepsilon_{m,h}}\right)^{1/p_{m,h}}
\]
and
\[
  V_{m,h}:=
  \begin{cases}
  s_{m,h}^2
  \left(\dfrac{e}{\varepsilon_{m,h}}\right)^{(2-p_{m,h})/p_{m,h}}
  (n_mA_{m,h})^{2/p_{m,h}},&0<p_{m,h}<2,\\[6pt]
  n_mA_{m,h}^{2/p_{m,h}}s_{m,h}^2,&p_{m,h}\ge2 .
  \end{cases}
\]
In both cases define
\[
  q_{m,h}:=
  \max\left\{
  r_{m,h},
  \left(\frac{\ell_m}{n_m}V_{m,h}\right)^{1/2}
  \right\}.
\]
Assume that the corresponding truncated-tail drift hypothesis from
Corollary~\ref{cor:mixed-subweibull-quantile-budget-norm-tail-selection} or
Corollary~\ref{cor:mixed-moment-quantile-budget-norm-tail-selection}, with
the common cap \(r_{m,h}\), holds outside \(E_{m,h}\).  If
\(\ell_m\le B\sqrt{n_m}\) eventually,
\[
  \sum_m\sum_{h=1}^{J_m}\bigl[\mathbb P(E_{m,h})+\varepsilon_{m,h}\bigr]
  <\infty,
\]
and
\[
  \sum_mJ_m\Delta_m
  (e\Delta_m)^{-\gamma_{\rm tc}(K,1,1,B)}<\infty,
\]
then the original streams are eventually contraction-normalized by
\(\max\{1,K+\eta\}q_{m,h}\), and the all-plus signing gives discrepancy at
most \(\sqrt{n_m}\), simultaneously for all active \(h\).

If \(J_m\Delta_m\le(e\Delta_m)^\Gamma\),
\(J_m\le(e\Delta_m)^\Lambda\), \(n_m\ge a\ell_m^2\) eventually, and
\(\varepsilon_{m,h}=(e\Delta_m)^{-c}\), then the cap-budget part is summable
for all integer dimensions when \(c>\Lambda+1\), and along every
stretched-exponential sparse dimension schedule when \(c>\Lambda\).  The
matrix-tail margins remain
\[
  K>\Theta_{\rm tc}(\Gamma+1,a,1,1)
\]
and
\[
  K>\Theta_{\rm tc}(\Gamma,a,1,1),
\]
respectively.

In the centered independent or martingale-difference case, the drift
hypothesis for a \(\psi\)-tagged stream is implied by
\[
  \begin{gathered}
  s_{m,h}\left[
    \varepsilon_{m,h}u_{m,h}
    +
    n_mA_{m,h}\int_{u_{m,h}}^\infty
    e^{-t^{\alpha_{m,h}}}\,dt
  \right]
  \le \eta q_{m,h}\sqrt{n_m},\\
  u_{m,h}:=
  \left(\log\frac{e n_mA_{m,h}}{\varepsilon_{m,h}}\right)^{1/\alpha_{m,h}}.
  \end{gathered}
\]
For a \(\mu\)-tagged stream with \(p_{m,h}>1\), it is implied by
\[
  n_mA_{m,h}s_{m,h}^{p_{m,h}}r_{m,h}^{1-p_{m,h}}
  \le \eta q_{m,h}\sqrt{n_m}.
\]
\end{corollary}

\begin{proof}
For a \(\psi\)-tagged stream, apply
Corollary~\ref{cor:mixed-subweibull-quantile-budget-norm-tail-selection} with
\[
  A_{m,h,i}=A_{m,h},\qquad
  s_{m,h,i}=s_{m,h},\qquad
  \alpha_{m,h,i}=\alpha_{m,h},\qquad
  \omega_{m,h,i}=1/n_m .
\]
Then every cap \(b_{m,h,i}\) equals \(r_{m,h}\), and the mixed sub-Weibull
energy is \(V_{m,h}\).  The displayed centered drift condition is exactly the
drift sum in that corollary with these equal shares.

For a \(\mu\)-tagged stream, apply
Corollary~\ref{cor:mixed-moment-quantile-budget-norm-tail-selection} with
\[
  A_{m,h,i}=A_{m,h},\qquad
  s_{m,h,i}=s_{m,h},\qquad
  p_{m,h,i}=p_{m,h},\qquad
  \omega_{m,h,i}=1/n_m .
\]
Again every cap equals \(r_{m,h}\).  If \(0<p_{m,h}<2\), the aggregate
moment-energy is
\[
  \begin{aligned}
  &n_mA_{m,h}s_{m,h}^2
  \left(\frac{e n_mA_{m,h}}{\varepsilon_{m,h}}\right)^{(2-p_{m,h})/p_{m,h}}\\
  &\hspace{2em}=
  s_{m,h}^2
  \left(\frac{e}{\varepsilon_{m,h}}\right)^{(2-p_{m,h})/p_{m,h}}
  (n_mA_{m,h})^{2/p_{m,h}},
  \end{aligned}
\]
and if \(p_{m,h}\ge2\) it is
\(n_mA_{m,h}^{2/p_{m,h}}s_{m,h}^2\).  The displayed moment drift condition is
the same specialization of the mixed finite-moment drift bound.

Apply the appropriate selector to every active stream.  The two selectors have
the same cap-budget summability and matrix-tail requirement, so summing over
the active menu gives the simultaneous eventual conclusion by Borel--Cantelli.
The all-integer and sparse dimension thresholds are the common thresholds from
those two selector corollaries.
\end{proof}

\begin{corollary}[Signed cores plus mixed dominated operator-tail libraries]
\label{cor:signed-core-mixed-dominated-operator-tail-libraries}
For each \(m\ge1\), let \(D_m,n_m,J_m\ge1\), put
\(\Delta_m:=D_m+n_m+1\) and \(\ell_m:=\log(e\Delta_m)\), and let
\(\mathcal H_m\) be a row environment.  For \(h\ge1\), let
\[
  B_{m,h,1},\ldots,B_{m,h,n_m},
  \qquad
  Y_{m,h,1},\ldots,Y_{m,h,n_m}
\]
be self-adjoint random matrices in \(\M_{D_m}(\C)\).  Assume that there are
events \(F_m^{\rm core}\in\mathcal H_m\) with
\(\sum_m\mathbb P(F_m^{\rm core})<\infty\) such that, outside
\(F_m^{\rm core}\), the active core tuples are contractions and admit one
\(\mathcal H_m\)-measurable active-prefix signing:
\[
  \|B_{m,h,i}\|_{\op}\le1,
  \qquad
  \left\|\sum_{i=1}^{n_m}\eps_{m,i}^\star B_{m,h,i}\right\|_{\op}
  \le K_{\rm core}\sqrt{n_m}
\]
for every \(h\le J_m\) and \(i\le n_m\).

For each active \(h\le J_m\), tag the perturbation stream either \(\psi\) or
\(\mu\), choose the parameters
\(\varepsilon_{m,h},A_{m,h},s_{m,h}\), and define \(r_{m,h}\), \(V_{m,h}\),
and \(q_{m,h}\) exactly as in
Corollary~\ref{cor:mixed-dominated-operator-tail-active-libraries}, with
\(Y_{m,h,i}\) in place of \(X_{m,h,i}\).  Assume that the corresponding
conditional tail or moment hypotheses, and the matching truncated-tail drift
hypotheses outside events \(E_{m,h}\), hold after adjoining
\(\mathcal H_m\) to the predictable sigma-fields.  Let
\(K_{\rm tail},\eta>0\), and put
\[
  \begin{aligned}
  \lambda_{m,h}&:=\max\{1,K_{\rm tail}+\eta\}\,q_{m,h},\\
  Z_{m,h,i}&:=\frac{Y_{m,h,i}}{\lambda_{m,h}},\\
  X_{m,h,i}&:=B_{m,h,i}+Z_{m,h,i}.
  \end{aligned}
\]
If \(\ell_m\le B\sqrt{n_m}\) eventually,
\[
  \sum_m\left[
    \mathbb P(F_m^{\rm core})
    +\sum_{h=1}^{J_m}\bigl(\mathbb P(E_{m,h})+\varepsilon_{m,h}\bigr)
  \right]<\infty,
\]
and
\[
  \sum_mJ_m\Delta_m
  (e\Delta_m)^{-\gamma_{\rm tc}(K_{\rm tail},1,1,B)}<\infty,
\]
then, with probability one, for all sufficiently large \(m\) and all
\(h\le J_m\),
\[
  \max_{i\le n_m}\|Z_{m,h,i}\|_{\op}\le1,\qquad
  \max_{i\le n_m}\|X_{m,h,i}\|_{\op}\le2,
\]
and the same active-prefix core signs satisfy
\[
  \left\|\sum_{i=1}^{n_m}\eps_{m,i}^\star X_{m,h,i}\right\|_{\op}
  \le (K_{\rm core}+1)\sqrt{n_m}.
\]
Consequently the normalized contractions
\(\widehat X_{m,h,i}:=X_{m,h,i}/2\) have common active-prefix discrepancy at
most \((K_{\rm core}+1)\sqrt{n_m}/2\), and every fixed menu element is
eventually covered whenever \(J_m\to\infty\).

The all-integer and sparse active-menu thresholds are the same as in
Corollary~\ref{cor:mixed-dominated-operator-tail-active-libraries}: after
adding \(\sum_m\mathbb P(F_m^{\rm core})<\infty\), the cap-budget part is
summable under \(c>\Lambda+1\) for all integer dimensions and under
\(c>\Lambda\) on stretched-exponential sparse schedules, with the same
\(\Theta_{\rm tc}\) margins for \(K_{\rm tail}\).
\end{corollary}

\begin{proof}
Work outside \(F_m^{\rm core}\).  The signing \(\eps_m^\star\) is
\(\mathcal H_m\)-measurable, so after conditioning on the row environment the
signed perturbation streams
\[
  \widetilde Y_{m,h,i}:=\eps_{m,i}^\star Y_{m,h,i}
\]
have the same operator-norm tail and moment bounds as \(Y_{m,h,i}\).  The
centered independent and martingale-difference drift conditions are also
unchanged after multiplication by these predictable signs, because all
tail-drift estimates use only operator norms.

Apply
Corollary~\ref{cor:mixed-dominated-operator-tail-active-libraries} to the
active library \((\widetilde Y_{m,h,i})\).  The displayed summability
assumptions, together with Borel--Cantelli, imply that eventually
\[
  \max_i\|Z_{m,h,i}\|_{\op}\le1
  \quad\text{and}\quad
  \left\|\sum_i\eps_{m,i}^\star Z_{m,h,i}\right\|_{\op}\le\sqrt{n_m}
\]
simultaneously for all \(h\le J_m\).  Combining this with the core signing
gives
\[
  \begin{aligned}
  \left\|\sum_i\eps_{m,i}^\star X_{m,h,i}\right\|_{\op}
  &\le
  \left\|\sum_i\eps_{m,i}^\star B_{m,h,i}\right\|_{\op}
  +
  \left\|\sum_i\eps_{m,i}^\star Z_{m,h,i}\right\|_{\op}\\
  &\le (K_{\rm core}+1)\sqrt{n_m},
  \end{aligned}
\]
and the contraction bound follows from
\[
  \|X_{m,h,i}\|_{\op}
  \le
  \|B_{m,h,i}\|_{\op}+\|Z_{m,h,i}\|_{\op}
  \le2.
\]
The final threshold statement is inherited verbatim from the mixed dominated
operator-tail corollary.
\end{proof}

\begin{corollary}[Raw noncentered core-plus-tail libraries]
\label{cor:raw-mean-core-mixed-dominated-operator-tail}
Keep the hypotheses and notation of
Corollary~\ref{cor:signed-core-mixed-dominated-operator-tail-libraries}, but
regard
\[
  W_{m,h,i}:=B_{m,h,i}+Y_{m,h,i}
\]
as the raw active random matrices.  Define
\[
  \lambda_{m,h}:=\max\{1,K_{\rm tail}+\eta\}\,q_{m,h},
  \qquad
  \Lambda_{m,h}:=\max\{1,\lambda_{m,h}\}.
\]
Then, under the same summability and matrix-tail hypotheses, with probability
one, for all sufficiently large \(m\), one \(\mathcal H_m\)-measurable sign
vector \(\eps_m^\star\) satisfies, simultaneously for every active
\(h\le J_m\),
\[
  \max_{i\le n_m}
  \left\|\frac{W_{m,h,i}}{\Lambda_{m,h}}\right\|_{\op}\le2
\]
and
\[
  \left\|
  \sum_{i=1}^{n_m}\eps_{m,i}^\star
  \frac{W_{m,h,i}}{\Lambda_{m,h}}
  \right\|_{\op}
  \le (K_{\rm core}+1)\sqrt{n_m}.
\]
Consequently the raw normalized contractions
\[
  \widehat W_{m,h,i}:=\frac{W_{m,h,i}}{2\Lambda_{m,h}}
\]
have common active-prefix discrepancy at most
\((K_{\rm core}+1)\sqrt{n_m}/2\), and every fixed menu element is eventually
covered whenever \(J_m\to\infty\).

In particular, one may take
\[
  B_{m,h,i}=\mathbb E(W_{m,h,i}\mid\mathcal H_m),
  \qquad
  Y_{m,h,i}=W_{m,h,i}-\mathbb E(W_{m,h,i}\mid\mathcal H_m),
\]
whenever the conditional-mean core tuple satisfies the common-signing
hypothesis and the centered fluctuations satisfy the same mixed dominated
tail, moment, and drift assumptions.  Thus noncentered random matrix libraries
with structured conditional means reduce to the same per-stream
\(\Lambda_{m,h}\) normalizations and the same all-integer and sparse
thresholds.
\end{corollary}

\begin{proof}
The proof of
Corollary~\ref{cor:signed-core-mixed-dominated-operator-tail-libraries}, applied
to the signed perturbation streams
\(\widetilde Y_{m,h,i}=\eps_{m,i}^\star Y_{m,h,i}\), gives eventually,
simultaneously for every active \(h\le J_m\),
\[
  \max_i\left\|\frac{Y_{m,h,i}}{\lambda_{m,h}}\right\|_{\op}\le1,
  \qquad
  \left\|\sum_i\eps_{m,i}^\star
  \frac{Y_{m,h,i}}{\lambda_{m,h}}\right\|_{\op}\le\sqrt{n_m}.
\]
Equivalently,
\[
  \max_i\|Y_{m,h,i}\|_{\op}\le\lambda_{m,h},
  \qquad
  \left\|\sum_i\eps_{m,i}^\star Y_{m,h,i}\right\|_{\op}
  \le \lambda_{m,h}\sqrt{n_m}.
\]
Since \(\Lambda_{m,h}\ge1\) and
\(\Lambda_{m,h}\ge\lambda_{m,h}\), the signed-core hypothesis gives
\[
  \left\|\sum_i\eps_{m,i}^\star
  \frac{W_{m,h,i}}{\Lambda_{m,h}}\right\|_{\op}
  \le
  \frac{K_{\rm core}\sqrt{n_m}}{\Lambda_{m,h}}
  +
  \frac{\lambda_{m,h}\sqrt{n_m}}{\Lambda_{m,h}}
  \le (K_{\rm core}+1)\sqrt{n_m}.
\]
Similarly,
\[
  \left\|\frac{W_{m,h,i}}{\Lambda_{m,h}}\right\|_{\op}
  \le
  \frac{\|B_{m,h,i}\|_{\op}}{\Lambda_{m,h}}
  +
  \frac{\|Y_{m,h,i}\|_{\op}}{\Lambda_{m,h}}
  \le2 .
\]
The conditional-mean specialization is only the decomposition
\(W=B+Y\) with \(B\) chosen to be \(\mathcal H_m\)-measurable; the hypotheses
are exactly those of the signed-core corollary for that decomposition.
\end{proof}

\begin{corollary}[Scaled raw noncentered core-plus-tail libraries]
\label{cor:scaled-raw-mean-core-mixed-dominated-operator-tail}
In the setting of
Corollary~\ref{cor:raw-mean-core-mixed-dominated-operator-tail}, replace the
unit core bounds by the following scale-aware bounds.  There are positive
\(\mathcal H_m\)-measurable numbers \(\rho_{m,h}\) such that, outside
\(F_m^{\rm core}\), the same active-prefix sign vector \(\eps_m^\star\)
satisfies
\[
  \|B_{m,h,i}\|_{\op}\le \rho_{m,h},
  \qquad
  \left\|\sum_{i=1}^{n_m}\eps_{m,i}^\star B_{m,h,i}\right\|_{\op}
  \le K_{\rm core}\rho_{m,h}\sqrt{n_m}
\]
for every \(h\le J_m\) and \(i\le n_m\).  Keep the same mixed dominated
operator-tail hypotheses for \(Y_{m,h,i}\), and define
\[
  \lambda_{m,h}:=\max\{1,K_{\rm tail}+\eta\}\,q_{m,h},
  \qquad
  \Lambda^{\rm sc}_{m,h}:=\max\{\rho_{m,h},\lambda_{m,h}\}.
\]
Then, under the same summability and matrix-tail assumptions, with probability
one, for all sufficiently large \(m\), the same \(\mathcal H_m\)-measurable sign
vector \(\eps_m^\star\) satisfies, simultaneously for every active
\(h\le J_m\),
\[
  \max_{i\le n_m}
  \left\|\frac{W_{m,h,i}}{\Lambda^{\rm sc}_{m,h}}\right\|_{\op}\le2
\]
and
\[
  \left\|
  \sum_{i=1}^{n_m}\eps_{m,i}^\star
  \frac{W_{m,h,i}}{\Lambda^{\rm sc}_{m,h}}
  \right\|_{\op}
  \le (K_{\rm core}+1)\sqrt{n_m}.
\]
Thus the raw matrices
\[
  \widehat W^{\rm sc}_{m,h,i}:=
  \frac{W_{m,h,i}}{2\Lambda^{\rm sc}_{m,h}}
\]
are contractions with common active-prefix discrepancy at most
\((K_{\rm core}+1)\sqrt{n_m}/2\), and every fixed menu element is eventually
covered whenever \(J_m\to\infty\).

In particular, for noncentered ensembles one may take
\(B_{m,h,i}=\mathbb E(W_{m,h,i}\mid\mathcal H_m)\) and
\(Y_{m,h,i}=W_{m,h,i}-\mathbb E(W_{m,h,i}\mid\mathcal H_m)\) whenever the
conditional-mean tuple has a common signing at scale \(\rho_{m,h}\).  The
centered fluctuation scale \(\lambda_{m,h}\) and the structured-mean scale
\(\rho_{m,h}\) then compete only through the single realized normalizer
\(\Lambda^{\rm sc}_{m,h}\).
\end{corollary}

\begin{proof}
The mixed dominated perturbation proof used in
Corollary~\ref{cor:raw-mean-core-mixed-dominated-operator-tail} gives
eventually, simultaneously for every active \(h\le J_m\),
\[
  \max_i\|Y_{m,h,i}\|_{\op}\le\lambda_{m,h},
  \qquad
  \left\|\sum_i\eps_{m,i}^\star Y_{m,h,i}\right\|_{\op}
  \le \lambda_{m,h}\sqrt{n_m}.
\]
Since \(\Lambda^{\rm sc}_{m,h}\ge\rho_{m,h}\) and
\(\Lambda^{\rm sc}_{m,h}\ge\lambda_{m,h}\), the scaled core hypothesis gives
\[
  \left\|\sum_i\eps_{m,i}^\star
  \frac{W_{m,h,i}}{\Lambda^{\rm sc}_{m,h}}\right\|_{\op}
  \le
  \frac{K_{\rm core}\rho_{m,h}\sqrt{n_m}}{\Lambda^{\rm sc}_{m,h}}
  +
  \frac{\lambda_{m,h}\sqrt{n_m}}{\Lambda^{\rm sc}_{m,h}}
  \le (K_{\rm core}+1)\sqrt{n_m}.
\]
The contraction estimate is the same pathwise comparison:
\[
  \left\|\frac{W_{m,h,i}}{\Lambda^{\rm sc}_{m,h}}\right\|_{\op}
  \le
  \frac{\rho_{m,h}}{\Lambda^{\rm sc}_{m,h}}
  +
  \frac{\lambda_{m,h}}{\Lambda^{\rm sc}_{m,h}}
  \le2 .
\]
No summability threshold changes, because the tail event for \(Y\) is exactly
the event used in the raw noncentered corollary and the core event is still
summable.
\end{proof}

\begin{corollary}[Two-scale raw noncentered core-plus-tail libraries]
\label{cor:two-scale-raw-mean-core-mixed-dominated-operator-tail}
In the setting of
Corollary~\ref{cor:raw-mean-core-mixed-dominated-operator-tail}, suppose instead
that the active cores have separate row and signing scales.  Namely, outside
\(F_m^{\rm core}\), the same active-prefix sign vector \(\eps_m^\star\)
satisfies
\[
  \|B_{m,h,i}\|_{\op}\le \rho_{m,h},
  \qquad
  \left\|\sum_{i=1}^{n_m}\eps_{m,i}^\star B_{m,h,i}\right\|_{\op}
  \le K_{\rm core}\sigma_{m,h}\sqrt{n_m}
\]
for every \(h\le J_m\) and \(i\le n_m\), where
\(\rho_{m,h}>0\) and \(\sigma_{m,h}>0\) are
\(\mathcal H_m\)-measurable.  Keep the same mixed dominated operator-tail
hypotheses for \(Y_{m,h,i}\), and define
\[
  \lambda_{m,h}:=\max\{1,K_{\rm tail}+\eta\}\,q_{m,h},
  \qquad
  \Lambda^{\rm rs}_{m,h}:=
  \max\{\rho_{m,h},\sigma_{m,h},\lambda_{m,h}\}.
\]
Then, under the same summability and matrix-tail assumptions, with probability
one, for all sufficiently large \(m\), the same signs satisfy, simultaneously
for every active \(h\le J_m\),
\[
  \max_{i\le n_m}
  \left\|\frac{W_{m,h,i}}{\Lambda^{\rm rs}_{m,h}}\right\|_{\op}\le2
\]
and
\[
  \left\|
  \sum_{i=1}^{n_m}\eps_{m,i}^\star
  \frac{W_{m,h,i}}{\Lambda^{\rm rs}_{m,h}}
  \right\|_{\op}
  \le (K_{\rm core}+1)\sqrt{n_m}.
\]
Consequently
\[
  \widehat W^{\rm rs}_{m,h,i}:=
  \frac{W_{m,h,i}}{2\Lambda^{\rm rs}_{m,h}}
\]
is a contraction tuple with common active-prefix discrepancy at most
\((K_{\rm core}+1)\sqrt{n_m}/2\), and every fixed menu element is eventually
covered whenever \(J_m\to\infty\).

Thus the row-size estimate for a structured conditional mean and the scale of
its available common signing need not be the same.  Finite-lift, algebraic,
Toeplitz/Hankel, group-convolution, or selected structured cores may be used
with whichever realized row bound \(\rho_{m,h}\) and signed-core certificate
\(\sigma_{m,h}\) their construction supplies; the centered mixed dominated
fluctuation competes with both only through the maximum normalizer
\(\Lambda^{\rm rs}_{m,h}\).
\end{corollary}

\begin{proof}
As in
Corollary~\ref{cor:raw-mean-core-mixed-dominated-operator-tail}, the mixed
dominated perturbation event gives eventually, simultaneously for every active
\(h\le J_m\),
\[
  \max_i\|Y_{m,h,i}\|_{\op}\le\lambda_{m,h},
  \qquad
  \left\|\sum_i\eps_{m,i}^\star Y_{m,h,i}\right\|_{\op}
  \le \lambda_{m,h}\sqrt{n_m}.
\]
Because \(\Lambda^{\rm rs}_{m,h}\) dominates
\(\rho_{m,h}\), \(\sigma_{m,h}\), and \(\lambda_{m,h}\),
\[
  \left\|\sum_i\eps_{m,i}^\star
  \frac{W_{m,h,i}}{\Lambda^{\rm rs}_{m,h}}\right\|_{\op}
  \le
  \frac{K_{\rm core}\sigma_{m,h}\sqrt{n_m}}{\Lambda^{\rm rs}_{m,h}}
  +
  \frac{\lambda_{m,h}\sqrt{n_m}}{\Lambda^{\rm rs}_{m,h}}
  \le (K_{\rm core}+1)\sqrt{n_m},
\]
while
\[
  \left\|\frac{W_{m,h,i}}{\Lambda^{\rm rs}_{m,h}}\right\|_{\op}
  \le
  \frac{\rho_{m,h}}{\Lambda^{\rm rs}_{m,h}}
  +
  \frac{\lambda_{m,h}}{\Lambda^{\rm rs}_{m,h}}
  \le2 .
\]
The probabilistic input is unchanged from the raw noncentered theorem, so the
same Borel--Cantelli and active-menu threshold conclusions apply.
\end{proof}

\begin{corollary}[Two-scale raw structured cores plus mixed dominated tails]
\label{cor:two-scale-raw-structured-core-mixed-dominated-operator-tail}
For each \(m\ge1\), let \(D_m,n_m,J_m\ge1\), put
\(\Delta_m:=D_m+n_m+1\) and \(\ell_m:=\log(e\Delta_m)\), and let
\(\mathcal H_m\) be a row environment.  Let
\[
  W_{m,h,i}:=B_{m,h,i}+Y_{m,h,i}
\]
be active raw self-adjoint matrices in \(\M_{D_m}(\C)\).  Assume that there are
events \(F_m^{\rm str}\in\mathcal H_m\), with
\(\sum_m\mathbb P(F_m^{\rm str})<\infty\), and positive
\(\mathcal H_m\)-measurable scales \(\rho_{m,h}\), \(\sigma_{m,h}\) such that,
outside \(F_m^{\rm str}\),
\[
  \|B_{m,h,i}\|_{\op}\le\rho_{m,h}\qquad(h\le J_m,\ i\le n_m).
\]
Suppose moreover that one of the following active-prefix structural
certificates holds for
\[
  C_{m,h,i}:=\frac{B_{m,h,i}}{\sigma_{m,h}} .
\]

\begin{enumerate}[label=\textup{(\alph*)}]
\item The tuples \((C_{m,h,i})_{i\le n_m}\) are contractions and admit finite
block lifts whose active-prefix parameters in the sense of
Corollary~\ref{cor:finite-lift-common-menu} obey
\[
  M_m^{(J)}\le M_0,\qquad E_m^{(J)}\le C_E n_m .
\]
Set \(K_{\rm str}:=C M_0\sqrt{\log(e+C_E)}\).

\item Each \((C_{m,h,i})_{i\le n_m}\) is a contraction tuple in a
finite-dimensional C*-algebra \(\mcA_{m,h}\subset\M_{D_m}(\C)\), and
\[
  \sum_{h=1}^{J_m}\dim_\C\mcA_{m,h}\le C_0n_m .
\]
Set \(K_{\rm str}:=C(C_0)\), the Akbas--Sra constant from
Theorem~\ref{thm:akbassra}.
\end{enumerate}

For each active \(h\le J_m\), tag the centered fluctuation stream either
\(\psi\) or \(\mu\), choose the parameters
\(\varepsilon_{m,h},A_{m,h},s_{m,h}\), and define \(r_{m,h}\), \(V_{m,h}\), and
\(q_{m,h}\) as in
Corollary~\ref{cor:mixed-dominated-operator-tail-active-libraries}, with
\(Y_{m,h,i}\) in place of \(X_{m,h,i}\).  Assume that the corresponding mixed
dominated tail or moment hypotheses, and the matching truncated-tail drift
hypotheses outside events \(E_{m,h}\), hold after adjoining \(\mathcal H_m\) to
the predictable sigma-fields.  Let \(K_{\rm tail},\eta>0\), and put
\[
  \lambda_{m,h}:=\max\{1,K_{\rm tail}+\eta\}\,q_{m,h},
  \qquad
  \Lambda^{\rm str}_{m,h}:=
  \max\{\rho_{m,h},\sigma_{m,h},\lambda_{m,h}\}.
\]
If \(\ell_m\le B\sqrt{n_m}\) eventually,
\[
  \sum_m\left[
    \mathbb P(F_m^{\rm str})
    +\sum_{h=1}^{J_m}\bigl(\mathbb P(E_{m,h})+\varepsilon_{m,h}\bigr)
  \right]<\infty,
\]
and
\[
  \sum_mJ_m\Delta_m
  (e\Delta_m)^{-\gamma_{\rm tc}(K_{\rm tail},1,1,B)}<\infty,
\]
then, with probability one, for all sufficiently large \(m\), one
\(\mathcal H_m\)-measurable sign vector \(\eps_m^\star\) satisfies,
simultaneously for every active \(h\le J_m\),
\[
  \max_{i\le n_m}
  \left\|\frac{W_{m,h,i}}{\Lambda^{\rm str}_{m,h}}\right\|_{\op}\le2
\]
and
\[
  \left\|
  \sum_{i=1}^{n_m}\eps_{m,i}^\star
  \frac{W_{m,h,i}}{\Lambda^{\rm str}_{m,h}}
  \right\|_{\op}
  \le (K_{\rm str}+1)\sqrt{n_m}.
\]
Thus the raw matrices
\[
  \widehat W^{\rm str}_{m,h,i}:=
  \frac{W_{m,h,i}}{2\Lambda^{\rm str}_{m,h}}
\]
are contractions with common active-prefix discrepancy at most
\((K_{\rm str}+1)\sqrt{n_m}/2\), and every fixed menu element is eventually
covered whenever \(J_m\to\infty\).

The all-integer and sparse active-menu thresholds are unchanged from
Corollary~\ref{cor:mixed-dominated-operator-tail-active-libraries}: if
\(J_m\Delta_m\le(e\Delta_m)^\Gamma\),
\(J_m\le(e\Delta_m)^\Lambda\), \(n_m\ge a\ell_m^2\) eventually, and
\(\varepsilon_{m,h}=(e\Delta_m)^{-c}\), then the cap-budget part is summable
under \(c>\Lambda+1\) for all integer dimensions and under \(c>\Lambda\) on
stretched-exponential sparse schedules, with the same
\(\Theta_{\rm tc}\) margins for \(K_{\rm tail}\).
\end{corollary}

\begin{proof}
Work outside \(F_m^{\rm str}\).  In case (a),
Corollary~\ref{cor:finite-lift-common-menu} gives signs satisfying
\[
  \left\|\sum_i\eps_{m,i}^\star C_{m,h,i}\right\|_{\op}
  \le K_{\rm str}\sqrt{n_m}\qquad(h\le J_m).
\]
In case (b), apply Theorem~\ref{thm:akbassra} to the direct-sum algebra
\(\bigoplus_{h\le J_m}\mcA_{m,h}\) and the direct-sum tuple
\(\bigoplus_{h\le J_m}C_{m,h,i}\) to get the same estimate.  Choose the
lexicographically first successful sign vector so that it is
\(\mathcal H_m\)-measurable.  Multiplying by \(\sigma_{m,h}\) gives the
two-scale core signing estimate
\[
  \left\|\sum_i\eps_{m,i}^\star B_{m,h,i}\right\|_{\op}
  \le K_{\rm str}\sigma_{m,h}\sqrt{n_m}.
\]
Together with the row bound \(\|B_{m,h,i}\|_{\op}\le\rho_{m,h}\), this verifies
the core hypotheses of
Corollary~\ref{cor:two-scale-raw-mean-core-mixed-dominated-operator-tail} with
\(K_{\rm core}=K_{\rm str}\) and \(F_m^{\rm core}=F_m^{\rm str}\).  The mixed
dominated assumptions and the displayed summability hypotheses are exactly the
perturbation and Borel--Cantelli hypotheses required there, so the raw
normalization, fixed-menu coverage, and all-integer/sparse threshold conclusions
follow.
\end{proof}

\begin{proposition}[Sign-symmetric random-tuple transfer]
\label{prop:sign-symmetric-transfer}
Let \(A=(A_1,\ldots,A_n)\) be a random self-adjoint tuple in
\(\M_D(\C)\).  Assume that its law is invariant under coordinatewise sign
changes: for every \(\delta\in\{\pm1\}^n\),
\[
  (\delta_1A_1,\ldots,\delta_nA_n)
  \quad\text{has the same law as}\quad
  (A_1,\ldots,A_n).
\]
If \(H\) is any event invariant under the same coordinatewise sign changes,
then, for every \(t\ge0\),
\[
  \mathbb P\left(H\cap\{\disc(A_1,\ldots,A_n)>t\}\right)
  \le
  \mathbb P\left(
  H\cap\left\{\left\|\sum_{i=1}^nA_i\right\|_{\op}>t\right\}
  \right).
\]
In particular, if \(H\subset\{\max_i\|A_i\|_{\op}\le1\}\), then any
high-probability all-plus estimate
\[
  \mathbb P\left(
  H^c\cup\left\{\left\|\sum_{i=1}^nA_i\right\|_{\op}>K\sqrt n\right\}
  \right)\le p
\]
implies
\[
  \mathbb P\left\{
  \max_i\|A_i\|_{\op}\le1
  \ \text{and}\
  \disc(A_1,\ldots,A_n)\le K\sqrt n
  \right\}\ge1-p .
\]
Thus, for coordinatewise sign-symmetric random matrix models, the Matrix
Spencer problem is reduced to the all-plus operator-norm tail for the same
realization-normalized ensemble.
\end{proposition}

\begin{proof}
Let \(\eta=(\eta_1,\ldots,\eta_n)\) be uniform on \(\{\pm1\}^n\), independent
of \(A\).  On the event \(H\cap\{\disc(A)>t\}\), every signing has norm larger
than \(t\).  Hence
\[
  \mathbb P\left(H\cap\{\disc(A)>t\}\right)
  \le
  \mathbb P\left(
  H\cap\left\{\left\|\sum_i\eta_iA_i\right\|_{\op}>t\right\}
  \right).
\]
By the sign-invariance of the law of \(A\), and by the sign-invariance of
\(H\), the last probability is exactly
\[
  \mathbb P\left(
  H\cap\left\{\left\|\sum_iA_i\right\|_{\op}>t\right\}
  \right).
\]
The final statement applies this inequality with \(t=K\sqrt n\) and with
\(H\) chosen to enforce the contraction normalization.
\end{proof}

\begin{corollary}[Borel--Cantelli sign-symmetric all-plus menus]
\label{cor:sign-symmetric-all-plus-menu-borel}
For each \(m\ge1\), let \(J_m\ge1\) be a deterministic active-menu size.  For
every \(1\le h\le J_m\), let
\[
  A_{m,h}=(A_{m,h,1},\ldots,A_{m,h,n_m})
\]
be a random self-adjoint tuple in \(\M_{D_m}(\C)\) whose law is invariant under
coordinatewise sign changes, and let \(H_{m,h}\) be a sign-invariant event such
that
\[
  H_{m,h}\subset
  \left\{\max_{1\le i\le n_m}\|A_{m,h,i}\|_{\op}\le1\right\}.
\]
Let \(K_m\ge0\) be deterministic.  If
\[
  \sum_m\sum_{h=1}^{J_m}
  \mathbb P\left(
  H_{m,h}^c
  \cup
  \left\{
  \left\|\sum_{i=1}^{n_m}A_{m,h,i}\right\|_{\op}
  >K_m\sqrt{n_m}
  \right\}
  \right)<\infty,
\]
then, with probability one, for all sufficiently large \(m\), every active menu
element is a normalized Spencer-scale instance:
\[
  \max_i\|A_{m,h,i}\|_{\op}\le1,
  \qquad
  \disc(A_{m,h,1},\ldots,A_{m,h,n_m})
  \le K_m\sqrt{n_m}
  \qquad(1\le h\le J_m).
\]

Equivalently, let \(X_{m,h,i}\) be coordinatewise sign-symmetric random
self-adjoint matrices, let \(H_{m,h}\) be sign-invariant events, and let
\(s_{m,h}>0\) be deterministic scales.  If
\[
  H_{m,h}\subset
  \left\{\max_i\|X_{m,h,i}\|_{\op}\le s_{m,h}\right\}
\]
and
\[
  \sum_m\sum_{h=1}^{J_m}
  \mathbb P\left(
  H_{m,h}^c
  \cup
  \left\{
  \left\|\sum_{i=1}^{n_m}X_{m,h,i}\right\|_{\op}
  >K_m s_{m,h}\sqrt{n_m}
  \right\}
  \right)<\infty,
\]
then, almost surely for all sufficiently large \(m\) and all \(1\le h\le J_m\),
\[
  \disc(X_{m,h,1},\ldots,X_{m,h,n_m})
  \le K_ms_{m,h}\sqrt{n_m}.
\]
If the streams are defined for all \(h\ge1\), \(J_m\to\infty\), and the active
menu is a deterministic exhaustion of a countable library, then every fixed
sign-symmetric model in the library is covered for all sufficiently large
\(m\).
\end{corollary}

\begin{proof}
By Proposition~\ref{prop:sign-symmetric-transfer}, for each \(m,h\) the
probability that \(H_{m,h}\) holds but the discrepancy exceeds
\(K_m\sqrt{n_m}\) is at most the displayed all-plus tail.  Hence the total
probability of active-menu failures is summable, and the first
Borel--Cantelli lemma gives the simultaneous eventual conclusion.  The scaled
version applies the first statement to \(A_{m,h,i}=s_{m,h}^{-1}X_{m,h,i}\).  If
\(J_m\to\infty\), each fixed library element belongs to the active menu
eventually.
\end{proof}

\begin{corollary}[Matrix-Bernstein sign-symmetric all-plus menus]
\label{cor:sign-symmetric-bernstein-menu-borel}
For each \(m\ge1\), let \(J_m\ge1\) be a deterministic active-menu size.  For
every \(1\le h\le J_m\), let
\[
  X_{m,h,1},\ldots,X_{m,h,n_m}\in\M_{D_m}(\C)
\]
be independent self-adjoint random matrices.  Assume that each
\(X_{m,h,i}\) has symmetric law,
\[
  X_{m,h,i}\stackrel{d}{=}-X_{m,h,i},
\]
and that, for a deterministic scale \(s_{m,h}>0\),
\[
  \|X_{m,h,i}\|_{\op}\le s_{m,h}
  \qquad\text{almost surely for all } i .
\]
Put
\[
  \sigma_{m,h}^2
  :=
  \left\|\sum_{i=1}^{n_m}\mathbb E X_{m,h,i}^2\right\|_{\op}.
\]
For deterministic thresholds \(u_{m,h}>0\), suppose that
\[
  \sum_m\sum_{h=1}^{J_m}
  2D_m
  \exp\left\{
  -\frac{u_{m,h}^2}
          {2(\sigma_{m,h}^2+s_{m,h}u_{m,h}/3)}
  \right\}
  <\infty .
\]
Then, with probability one, for all sufficiently large \(m\) and all active
\(1\le h\le J_m\),
\[
  \disc(X_{m,h,1},\ldots,X_{m,h,n_m})\le u_{m,h}.
\]
In particular, if \(u_{m,h}=K_m s_{m,h}\sqrt{n_m}\), then the normalized tuples
\(s_{m,h}^{-1}X_{m,h,i}\) are contractions and have discrepancy at most
\(K_m\sqrt{n_m}\) simultaneously on the active menu, eventually almost surely.
If the streams are defined for all \(h\ge1\), \(J_m\to\infty\), and the active
menu exhausts a countable library, then every fixed independent bounded
sign-symmetric model in that library is eventually covered.
\end{corollary}

\begin{proof}
Fix \(m\) and \(h\).  Independence together with the individual symmetry
\(X_{m,h,i}\stackrel{d}{=}-X_{m,h,i}\) makes the whole tuple coordinatewise
sign-symmetric, and also gives \(\mathbb E X_{m,h,i}=0\).  Applying
Proposition~\ref{prop:sign-symmetric-transfer} with \(H=\Omega\) gives
\[
  \mathbb P\left\{
  \disc(X_{m,h,1},\ldots,X_{m,h,n_m})>u_{m,h}
  \right\}
  \le
  \mathbb P\left\{
  \left\|\sum_iX_{m,h,i}\right\|_{\op}>u_{m,h}
  \right\}.
\]
The two-sided matrix Bernstein inequality \cite{Tropp}, with jump bound
\(s_{m,h}\) and variance proxy \(\sigma_{m,h}^2\), bounds the last probability
by the displayed summand.  The summability assumption and the first
Borel--Cantelli lemma give the simultaneous active-menu conclusion.  The
normalized and countable-exhaustion statements are immediate from the
deterministic scale bound and from \(J_m\to\infty\).
\end{proof}

\begin{corollary}[Explicit sign-symmetric Bernstein schedules]
\label{cor:sign-symmetric-bernstein-menu-schedules}
In the setting of Corollary~\ref{cor:sign-symmetric-bernstein-menu-borel}, put
\(\ell_m=\log(eD_m)\).  Fix constants \(\beta\ge0\) and \(K>0\).  Suppose that
outside events \(E_m\),
\[
  \sigma_{m,h}^2
  \le
  \beta\,s_{m,h}^2\,\frac{n_m}{\ell_m}
  \qquad(1\le h\le J_m),
\]
and that \(\sum_m\mathbb P(E_m)<\infty\).  Then the following deterministic
schedule conditions imply that, almost surely, for all sufficiently large \(m\)
and every active \(1\le h\le J_m\),
\[
  \disc(X_{m,h,1},\ldots,X_{m,h,n_m})
  \le K\,s_{m,h}\sqrt{n_m}.
\]

\begin{enumerate}[label=\textup{(\alph*)}]
\item General dimension schedules.  If there is \(B>0\) such that
\[
  \ell_m\le B\sqrt{n_m}
\]
eventually, and
\[
  \sum_m J_mD_m(eD_m)^{-\gamma(K,\beta,B)}<\infty,
  \qquad
  \gamma(K,\beta,B):=
  \frac{K^2}{2(\beta+KB/3)},
\]
then the conclusion holds.

\item All integer dimensions.  Suppose the rows are indexed by \(D=2,3,\ldots\),
with ambient dimension \(D\), and
\[
  J_DD\le(eD)^\Gamma
\]
eventually.  If
\[
  n_D\ge a\,\log^2(eD)
\]
eventually and
\[
  \frac{K^2}{2(\beta+K/(3\sqrt a))}>\Gamma+1,
\]
then the conclusion holds along all integer dimensions.

\item Sparse stretched-exponential dimensions.  If
\[
  D_m\ge \exp(c m^\kappa)
\]
eventually for some \(c,\kappa>0\), if
\[
  J_mD_m\le(eD_m)^\Gamma
\]
eventually, if
\[
  n_m\ge a\,\log^2(eD_m)
\]
eventually, and if
\[
  \frac{K^2}{2(\beta+K/(3\sqrt a))}>\Gamma,
\]
then the conclusion holds along the sparse schedule.
\end{enumerate}

In the two active-prefix cases, if \(J_m\to\infty\), the same schedules cover
every fixed member of the countable independent bounded sign-symmetric library
eventually.
\end{corollary}

\begin{proof}
Apply Corollary~\ref{cor:sign-symmetric-bernstein-menu-borel} with
\(u_{m,h}=Ks_{m,h}\sqrt{n_m}\).  On \(E_m^c\), the \(h\)-th active summand in
that corollary is bounded by
\[
  2D_m
  \exp\left\{
  -\frac{K^2n_m}
          {2(\beta n_m/\ell_m+K\sqrt{n_m}/3)}
  \right\}.
\]
If \(\ell_m\le B\sqrt{n_m}\), then
\[
  \frac{K^2n_m}
       {2(\beta n_m/\ell_m+K\sqrt{n_m}/3)}
  \ge
  \gamma(K,\beta,B)\ell_m.
\]
Thus the total active-row tail is at most
\[
  2J_mD_m(eD_m)^{-\gamma(K,\beta,B)}.
\]
The first schedule is therefore summable, and Borel--Cantelli gives the
eventual conclusion together with the summable variance-budget exceptions.

For all integer dimensions, \(n_D\ge a\ell_D^2\) gives \(B=1/\sqrt a\), so the
active-row tail is bounded by
\[
  2(eD)^{\Gamma-\gamma(K,\beta,1/\sqrt a)}.
\]
This is summable over \(D=2,3,\ldots\) when
\(\gamma(K,\beta,1/\sqrt a)>\Gamma+1\).  Along a sparse
stretched-exponential schedule the same bound is summable as soon as the
exponent is strictly negative, which is exactly the displayed sparse
condition.  If \(J_m\to\infty\), each fixed countable-library element is in the
active prefix eventually.
\end{proof}

\begin{corollary}[Closed-form sign-symmetric Bernstein margins]
\label{cor:sign-symmetric-bernstein-menu-closed-margins}
In the setting of Corollary~\ref{cor:sign-symmetric-bernstein-menu-schedules},
for \(A\ge0\) define
\[
  \Theta(A,a,\beta)
  :=
  \frac{A}{3\sqrt a}
  +
  \sqrt{\frac{A^2}{9a}+2A\beta}.
\]
Assume the same variance-budget hypothesis
\[
  \sigma_{m,h}^2
  \le
  \beta\,s_{m,h}^2\,\frac{n_m}{\log(eD_m)}
\]
outside summable exceptions.

If the rows are indexed by all integer dimensions \(D=2,3,\ldots\), if
\[
  J_DD\le(eD)^\Gamma,\qquad
  n_D\ge a\log^2(eD)
\]
eventually, and if
\[
  K>\Theta(\Gamma+1,a,\beta),
\]
then almost surely, for all sufficiently large \(D\) and all active
\(1\le h\le J_D\),
\[
  \disc(X_{D,h,1},\ldots,X_{D,h,n_D})
  \le K\,s_{D,h}\sqrt{n_D}.
\]

Along any sparse stretched-exponential dimension schedule
\(D_m\ge\exp(c m^\kappa)\), the same conclusion holds under
\[
  J_mD_m\le(eD_m)^\Gamma,\qquad
  n_m\ge a\log^2(eD_m),
  \qquad
  K>\Theta(\Gamma,a,\beta).
\]
If \(J_m\to\infty\), both statements exhaust countable independent bounded
sign-symmetric libraries.

Equivalently, for a prescribed \(K\) and target exponent \(A>0\), the condition
\[
  \frac{K^2}{2(\beta+K/(3\sqrt a))}>A
\]
is guaranteed whenever \(K^2>2A\beta\) and
\[
  a>
  \left(\frac{2AK}{3(K^2-2A\beta)}\right)^2 .
\]
\end{corollary}

\begin{proof}
The inequality
\[
  \frac{K^2}{2(\beta+K/(3\sqrt a))}>A
\]
is equivalent, for \(A>0\), to
\[
  K^2-\frac{2A}{3\sqrt a}K-2A\beta>0.
\]
The positive root of the quadratic on the left is exactly
\(\Theta(A,a,\beta)\).  Hence \(K>\Theta(\Gamma+1,a,\beta)\) implies the
all-integer hypothesis in
Corollary~\ref{cor:sign-symmetric-bernstein-menu-schedules}, and
\(K>\Theta(\Gamma,a,\beta)\) implies the sparse hypothesis.  The case \(A=0\)
is immediate for every \(K>0\).  Solving the same inequality for \(a\) gives the
final equivalent margin condition.
\end{proof}

\begin{corollary}[Conditional sign-symmetric Bernstein menus]
\label{cor:conditional-sign-symmetric-bernstein-menu}
For each \(m\ge1\), let \(J_m\ge1\) be a deterministic active-menu size.  For
every \(1\le h\le J_m\), let
\[
  X_{m,h,1},\ldots,X_{m,h,n_m}\in\M_{D_m}(\C)
\]
be self-adjoint random matrices and let \(\mcG_{m,h}\) be a conditioning
sigma-field.  Assume that, conditionally on \(\mcG_{m,h}\), the matrices are
independent and each \(X_{m,h,i}\) has symmetric conditional law.  Let
\(E_{m,h}\in\mcG_{m,h}\), and let \(s_{m,h}>0\), \(v_{m,h}\ge0\), and
\(u_{m,h}>0\) be deterministic.  Suppose that on \(E_{m,h}^c\),
\[
  \mathbb P\left\{
  \max_i\|X_{m,h,i}\|_{\op}\le s_{m,h}
  \,\middle|\,\mcG_{m,h}
  \right\}=1
\]
and
\[
  \left\|
  \sum_{i=1}^{n_m}\mathbb E\left(X_{m,h,i}^2\mid\mcG_{m,h}\right)
  \right\|_{\op}
  \le v_{m,h}.
\]
If
\[
  \sum_m\sum_{h=1}^{J_m}
  \left[
  \mathbb P(E_{m,h})
  +2D_m
  \exp\left\{
  -\frac{u_{m,h}^2}
          {2(v_{m,h}+s_{m,h}u_{m,h}/3)}
  \right\}
  \right]
  <\infty,
\]
then, with probability one, for all sufficiently large \(m\) and all active
\(1\le h\le J_m\),
\[
  \max_i\|X_{m,h,i}\|_{\op}\le s_{m,h},
  \qquad
  \disc(X_{m,h,1},\ldots,X_{m,h,n_m})\le u_{m,h}.
\]

Consequently, put \(\ell_m=\log(eD_m)\) and take
\(u_{m,h}=K s_{m,h}\sqrt{n_m}\).  If, outside the same summable envelope
failures,
\[
  v_{m,h}\le
  \beta\,s_{m,h}^2\,\frac{n_m}{\ell_m},
\]
then the three schedule conclusions of
Corollary~\ref{cor:sign-symmetric-bernstein-menu-schedules} hold with
\(\sigma_{m,h}^2\) replaced by \(v_{m,h}\).  In particular, under
\(J_DD\le(eD)^\Gamma\) and \(n_D\ge a\log^2(eD)\), all integer dimensions are
covered for every
\[
  K>\Theta(\Gamma+1,a,\beta),
\]
while sparse stretched-exponential dimension schedules with
\(J_mD_m\le(eD_m)^\Gamma\) are covered for every
\[
  K>\Theta(\Gamma,a,\beta).
\]
If \(J_m\to\infty\), these conditional schedules exhaust countable
conditionally independent bounded sign-symmetric libraries.

In particular, the hypotheses apply to random-direction scalar-multiplier
models \(X_{m,h,i}=\xi_{m,h,i}B_{m,h,i}\), where the directions \(B_{m,h,i}\)
are \(\mcG_{m,h}\)-measurable self-adjoint matrices and the real multipliers
\(\xi_{m,h,i}\) are conditionally independent, conditionally symmetric, and
bounded by \(1\) in absolute value, with
\[
  v_{m,h}=
  \left\|
  \sum_i\mathbb E(\xi_{m,h,i}^2\mid\mcG_{m,h})B_{m,h,i}^2
  \right\|_{\op}.
\]
\end{corollary}

\begin{proof}
Fix \(m,h\) and condition on \(\mcG_{m,h}\).  On \(E_{m,h}^c\), the conditional
law is coordinatewise sign-symmetric and is supported on the operator-norm ball
of radius \(s_{m,h}\).  Applying the proof of
Proposition~\ref{prop:sign-symmetric-transfer} to the regular conditional law
gives
\[
  \mathbb P\left\{
  \disc(X_{m,h,1},\ldots,X_{m,h,n_m})>u_{m,h}
  \,\middle|\,\mcG_{m,h}
  \right\}
  \le
  \mathbb P\left\{
  \left\|\sum_iX_{m,h,i}\right\|_{\op}>u_{m,h}
  \,\middle|\,\mcG_{m,h}
  \right\}
\]
on \(E_{m,h}^c\).  Conditional matrix Bernstein, with jump bound \(s_{m,h}\)
and conditional variance proxy bounded by \(v_{m,h}\), bounds the right-hand
side by the displayed exponential term.  Integrating and adding the bad event
\(E_{m,h}\) gives a summable active-menu failure series, so the first
Borel--Cantelli lemma gives the simultaneous eventual conclusion.  Substituting
\(u_{m,h}=Ks_{m,h}\sqrt{n_m}\) and the variance envelope gives exactly the
estimates used in
Corollary~\ref{cor:sign-symmetric-bernstein-menu-schedules}; the closed-form
thresholds are Corollary~\ref{cor:sign-symmetric-bernstein-menu-closed-margins}.
The scalar-multiplier specialization follows by conditioning on the directions.
\end{proof}

\begin{corollary}[Symmetric scalar-multiplier Bernstein schedules]
\label{cor:symmetric-scalar-multiplier-bernstein-schedules}
For each \(m\ge1\), let \(D_m\ge2\), let \(J_m\ge1\) be a deterministic
active-menu size, and put \(\ell_m=\log(eD_m)\).  For every
\(1\le h\le J_m\), let
\(\mcG_{m,h}\) be a conditioning sigma-field, let
\[
  B_{m,h,1},\ldots,B_{m,h,n_m}\in\M_{D_m}(\C)
\]
be \(\mcG_{m,h}\)-measurable self-adjoint contractions, and, conditionally on
\(\mcG_{m,h}\), let
\(\xi_{m,h,1},\ldots,\xi_{m,h,n_m}\) be independent real random variables with
symmetric conditional laws and \(|\xi_{m,h,i}|\le1\).  Put
\[
  X_{m,h,i}:=\xi_{m,h,i}B_{m,h,i},\qquad
  q_{m,h,i}:=\mathbb E(\xi_{m,h,i}^2\mid\mcG_{m,h}),
\]
and
\[
  \nu_{m,h}:=
  \left\|
  \sum_{i=1}^{n_m}q_{m,h,i}B_{m,h,i}^2
  \right\|_{\op}.
\]
Fix \(\beta\ge0\) and \(K>0\).  Suppose that outside events \(E_{m,h}\),
\[
  \nu_{m,h}\le \beta\,\frac{n_m}{\ell_m}
  \qquad(1\le h\le J_m),
\]
and that \(\sum_m\sum_{h\le J_m}\mathbb P(E_{m,h})<\infty\).  Then the
following deterministic schedule conditions imply that, almost surely, for all
sufficiently large \(m\) and every active \(1\le h\le J_m\),
\[
  \disc(X_{m,h,1},\ldots,X_{m,h,n_m})\le K\sqrt{n_m}.
\]

\begin{enumerate}[label=\textup{(\alph*)}]
\item General dimension schedules.  If there is \(B>0\) such that
\[
  \ell_m\le B\sqrt{n_m}
\]
eventually, and
\[
  \sum_m J_mD_m(eD_m)^{-\gamma(K,\beta,B)}<\infty,
  \qquad
  \gamma(K,\beta,B):=
  \frac{K^2}{2(\beta+KB/3)},
\]
then the conclusion holds.

\item All integer dimensions.  Suppose the rows are indexed by
\(D=2,3,\ldots\), with ambient dimension \(D\), and
\[
  J_DD\le(eD)^\Gamma
\]
eventually.  If
\[
  n_D\ge a\log^2(eD)
\]
eventually and
\[
  K>\Theta(\Gamma+1,a,\beta),
\]
then the conclusion holds along all integer dimensions.

\item Sparse stretched-exponential dimensions.  If
\[
  D_m\ge\exp(c m^\kappa)
\]
eventually for some \(c,\kappa>0\), if
\[
  J_mD_m\le(eD_m)^\Gamma
\]
eventually, if
\[
  n_m\ge a\log^2(eD_m)
\]
eventually, and if
\[
  K>\Theta(\Gamma,a,\beta),
\]
then the conclusion holds along the sparse schedule.
\end{enumerate}

The variance-envelope hypothesis is automatic, with no \(E_{m,h}\), whenever
\[
  q_{m,h,i}\le p_m\le \frac{\beta}{\ell_m}
  \qquad(1\le h\le J_m,\ 1\le i\le n_m),
\]
because \(0\preceq B_{m,h,i}^2\preceq I\).  If \(J_m\to\infty\), the same
schedules exhaust countable libraries of symmetric scalar-multiplier streams.
\end{corollary}

\begin{proof}
The conditional hypotheses imply that the \(X_{m,h,i}\)'s are conditionally
independent, conditionally sign-symmetric, and satisfy
\(\|X_{m,h,i}\|_{\op}\le1\).  Their conditional variance proxy is exactly
\(\nu_{m,h}\).  Corollary~\ref{cor:conditional-sign-symmetric-bernstein-menu}
with \(s_{m,h}=1\), \(v_{m,h}=\beta n_m/\ell_m\), and
\(u_{m,h}=K\sqrt{n_m}\) gives the general summability criterion.  The three
displayed schedules and the closed-form margins are precisely
Corollaries~\ref{cor:sign-symmetric-bernstein-menu-schedules} and
\ref{cor:sign-symmetric-bernstein-menu-closed-margins} in this scalar-multiplier
normalization.

If \(q_{m,h,i}\le p_m\), then
\[
  \nu_{m,h}
  \le p_m\left\|\sum_iB_{m,h,i}^2\right\|_{\op}
  \le p_m n_m
  \le \beta\,\frac{n_m}{\ell_m}.
\]
The countable-library assertion follows from the active-prefix exhaustion.
\end{proof}

\begin{corollary}[Sub-Gaussian symmetric scalar-multiplier normalizations]
\label{cor:subgaussian-symmetric-scalar-multiplier-normalizations}
For each \(m\ge1\), let \(D_m\ge2\), let \(J_m\ge1\) be a deterministic
active-menu size, and put \(\ell_m=\log(eD_m)\).  For every
\(1\le h\le J_m\), let \(\mcG_{m,h}\) be a conditioning sigma-field and let
\[
  B_{m,h,1},\ldots,B_{m,h,n_m}\in\M_{D_m}(\C)
\]
be \(\mcG_{m,h}\)-measurable self-adjoint contractions.  Conditionally on
\(\mcG_{m,h}\), let
\(\xi_{m,h,1},\ldots,\xi_{m,h,n_m}\) be independent real random variables with
symmetric conditional laws and deterministic sub-Gaussian proxies
\(a_{m,h,i}\ge0\):
\[
  \mathbb E\{\exp(\lambda\xi_{m,h,i})\mid\mcG_{m,h}\}
  \le
  \exp(\lambda^2a_{m,h,i}^2/2)
  \qquad(\lambda\in\R).
\]
Let \(L_{m,h}>0\) be deterministic normalization caps and set
\[
  A_{m,h,i}:=L_{m,h}^{-1}\xi_{m,h,i}B_{m,h,i},
  \qquad
  \nu_{m,h}:=
  L_{m,h}^{-2}
  \left\|\sum_{i=1}^{n_m}a_{m,h,i}^2B_{m,h,i}^2\right\|_{\op}.
\]
Fix \(\beta>0\) and \(K>0\).  Suppose that outside events
\(E_{m,h}\in\mcG_{m,h}\),
\[
  \nu_{m,h}\le \beta\,\frac{n_m}{\ell_m}
  \qquad(1\le h\le J_m),
\]
and that the variance-budget and cap-failure series
\[
  \sum_m\sum_{h=1}^{J_m}
  \left[
  \mathbb P(E_{m,h})
  +
  2\sum_{i=1}^{n_m}
  \exp\left\{-\frac{L_{m,h}^2}{2a_{m,h,i}^2}\right\}
  \right]
  <\infty
\]
is finite, with the exponential term interpreted as \(0\) when
\(a_{m,h,i}=0\).  Put
\[
  \gamma_{\rm sg}(K,\beta):=\frac{K^2}{2\beta}.
\]
Then each of the following deterministic schedule conditions implies that,
almost surely, for all sufficiently large \(m\) and every active
\(1\le h\le J_m\),
\[
  \max_i\|A_{m,h,i}\|_{\op}\le1,
  \qquad
  \disc(A_{m,h,1},\ldots,A_{m,h,n_m})\le K\sqrt{n_m}.
\]

\begin{enumerate}[label=\textup{(\alph*)}]
\item General dimension schedules.  If
\[
  \sum_mJ_mD_m(eD_m)^{-\gamma_{\rm sg}(K,\beta)}<\infty,
\]
then the conclusion holds.

\item All integer dimensions.  Suppose the rows are indexed by
\(D=2,3,\ldots\), with ambient dimension \(D\), and
\[
  J_DD\le(eD)^\Gamma
\]
eventually.  If
\[
  K>\sqrt{2\beta(\Gamma+1)},
\]
then the conclusion holds along all integer dimensions.

\item Sparse stretched-exponential dimensions.  If
\[
  D_m\ge\exp(c m^\kappa)
\]
eventually for some \(c,\kappa>0\), if
\[
  J_mD_m\le(eD_m)^\Gamma
\]
eventually, and if
\[
  K>\sqrt{2\beta\Gamma},
\]
then the conclusion holds along the sparse schedule.
\end{enumerate}

If \(J_m\to\infty\), the same schedules exhaust countable libraries of
sub-Gaussian symmetric scalar-multiplier streams.  In particular, for uniformly
proxied multipliers \(a_{m,h,i}\le a_{m,h}\), the cap-failure contribution is
bounded by
\[
  2n_m\exp\left\{-\frac{L_{m,h}^2}{2a_{m,h}^2}\right\};
\]
choosing \(L_{m,h}\) above the active logarithmic union scale makes the
unbounded-multiplier normalization exception summable.
\end{corollary}

\begin{proof}
Let
\[
  H_{m,h}:=\left\{\max_i|\xi_{m,h,i}|\le L_{m,h}\right\}.
\]
On \(H_{m,h}\), the normalized matrices \(A_{m,h,i}\) are contractions.  The
event \(H_{m,h}\) is invariant under coordinatewise sign changes.  Conditional
on \(\mcG_{m,h}\), the tuple \(A_{m,h}\) is coordinatewise sign-symmetric, so
Proposition~\ref{prop:sign-symmetric-transfer}, applied to the regular
conditional law and to \(H_{m,h}\), gives on \(E_{m,h}^c\)
\[
  \mathbb P\left(
  H_{m,h}\cap
  \{\disc(A_{m,h,1},\ldots,A_{m,h,n_m})>K\sqrt{n_m}\}
  \,\middle|\,\mcG_{m,h}
  \right)
  \le
  \mathbb P\left\{
  \left\|\sum_iA_{m,h,i}\right\|_{\op}>K\sqrt{n_m}
  \,\middle|\,\mcG_{m,h}
  \right\}.
\]
If \(\nu_{m,h}=0\), then \(a_{m,h,i}B_{m,h,i}=0\) for every \(i\), hence
\(\xi_{m,h,i}B_{m,h,i}=0\) almost surely under the conditional law.  Otherwise
the matrix Laplace bound of Tropp \cite{Tropp}, as in
Lemma~\ref{lem:matrix-series-subgaussian-tail}, bounds the last probability by
\[
  2D_m\exp\left\{-\frac{K^2n_m}{2\nu_{m,h}}\right\}
  \le
  2D_m(eD_m)^{-\gamma_{\rm sg}(K,\beta)}
\]
on \(E_{m,h}^c\).  Also, by the scalar sub-Gaussian tail bound and a union
bound,
\[
  \mathbb P(H_{m,h}^c\mid\mcG_{m,h})
  \le
  2\sum_{i=1}^{n_m}
  \exp\left\{-\frac{L_{m,h}^2}{2a_{m,h,i}^2}\right\}.
\]
After adding \(\mathbb P(E_{m,h})\), summing over the active menu, and using
the displayed summability hypotheses, Borel--Cantelli gives the simultaneous
eventual contraction and discrepancy conclusion.

The all-integer and sparse assertions are immediate from the general schedule:
under \(J_DD\le(eD)^\Gamma\), the row tail is
\(O((eD)^{\Gamma-\gamma_{\rm sg}(K,\beta)})\), summable over all integer
dimensions when \(\gamma_{\rm sg}(K,\beta)>\Gamma+1\), and summable along every
stretched-exponential schedule when \(\gamma_{\rm sg}(K,\beta)>\Gamma\).
The countable-library statement follows from active-prefix exhaustion.
\end{proof}

\begin{corollary}[Tail-capped symmetric scalar-multiplier schedules]
\label{cor:tail-capped-symmetric-scalar-multiplier-schedules}
For each \(m\ge1\), let \(D_m\ge2\), let \(J_m\ge1\) be a deterministic
active-menu size, and put \(\ell_m=\log(eD_m)\).  For every
\(1\le h\le J_m\), let \(\mcG_{m,h}\) be a conditioning sigma-field and let
\[
  B_{m,h,1},\ldots,B_{m,h,n_m}\in\M_{D_m}(\C)
\]
be \(\mcG_{m,h}\)-measurable self-adjoint contractions.  Conditionally on
\(\mcG_{m,h}\), let
\(\xi_{m,h,1},\ldots,\xi_{m,h,n_m}\) be independent real random variables with
symmetric conditional laws.  Let \(L_{m,h}>0\) be deterministic caps and set
\[
  A_{m,h,i}:=L_{m,h}^{-1}\xi_{m,h,i}B_{m,h,i},
\]
\[
  \widehat q_{m,h,i}:=
  L_{m,h}^{-2}
  \mathbb E\left(\xi_{m,h,i}^2
  \mathbf 1_{\{|\xi_{m,h,i}|\le L_{m,h}\}}\mid\mcG_{m,h}\right),
\]
and
\[
  \widehat\nu_{m,h}:=
  \left\|\sum_{i=1}^{n_m}\widehat q_{m,h,i}B_{m,h,i}^2\right\|_{\op}.
\]
Fix \(\beta\ge0\) and \(K>0\).  Suppose that outside events
\(E_{m,h}\in\mcG_{m,h}\),
\[
  \widehat\nu_{m,h}\le \beta\,\frac{n_m}{\ell_m}
  \qquad(1\le h\le J_m),
\]
and that the cap and variance-exception series
\[
  \sum_m\sum_{h=1}^{J_m}
  \left[
  \mathbb P(E_{m,h})
  +
  \sum_{i=1}^{n_m}
  \mathbb P\{|\xi_{m,h,i}|>L_{m,h}\}
  \right]
  <\infty
\]
is finite.  Then each of the following deterministic schedule conditions
implies that, almost surely, for all sufficiently large \(m\) and every active
\(1\le h\le J_m\),
\[
  \max_i\|A_{m,h,i}\|_{\op}\le1,
  \qquad
  \disc(A_{m,h,1},\ldots,A_{m,h,n_m})\le K\sqrt{n_m}.
\]

\begin{enumerate}[label=\textup{(\alph*)}]
\item General dimension schedules.  If there is \(B>0\) such that
\[
  \ell_m\le B\sqrt{n_m}
\]
eventually, and
\[
  \sum_mJ_mD_m(eD_m)^{-\gamma(K,\beta,B)}<\infty,
  \qquad
  \gamma(K,\beta,B):=
  \frac{K^2}{2(\beta+KB/3)},
\]
then the conclusion holds.

\item All integer dimensions.  Suppose the rows are indexed by
\(D=2,3,\ldots\), with ambient dimension \(D\), and
\[
  J_DD\le(eD)^\Gamma
\]
eventually.  If
\[
  n_D\ge a\log^2(eD)
\]
eventually and
\[
  K>\Theta(\Gamma+1,a,\beta),
\]
then the conclusion holds along all integer dimensions.

\item Sparse stretched-exponential dimensions.  If
\[
  D_m\ge\exp(c m^\kappa)
\]
eventually for some \(c,\kappa>0\), if
\[
  J_mD_m\le(eD_m)^\Gamma
\]
eventually, if
\[
  n_m\ge a\log^2(eD_m)
\]
eventually, and if
\[
  K>\Theta(\Gamma,a,\beta),
\]
then the conclusion holds along the sparse schedule.
\end{enumerate}

If \(J_m\to\infty\), the same schedules exhaust countable libraries of
tail-capped symmetric scalar-multiplier streams.  Moreover, suppose that for
some \(\alpha>0\) the conditional tails obey the deterministic sub-Weibull
bounds
\[
  \mathbb P\{|\xi_{m,h,i}|>t\mid\mcG_{m,h}\}
  \le
  2\exp\left\{-(t/a_{m,h,i})^\alpha\right\}
  \qquad(t>0),
\]
with the right-hand side interpreted as \(0\) for \(a_{m,h,i}=0\).  Then the
cap contribution is bounded by
\[
  2\sum_{i=1}^{n_m}
  \exp\left\{-(L_{m,h}/a_{m,h,i})^\alpha\right\},
\]
and it suffices for the variance budget to have
\[
  C_\alpha L_{m,h}^{-2}
  \left\|\sum_{i=1}^{n_m}a_{m,h,i}^2B_{m,h,i}^2\right\|_{\op}
  \le \beta\,\frac{n_m}{\ell_m},
  \qquad
  C_\alpha:=\frac4\alpha\Gamma(2/\alpha).
\]
For uniformly proxied multipliers \(a_{m,h,i}\le a_{m,h}\), the cap
contribution is at most
\[
  2n_m\exp\left\{-(L_{m,h}/a_{m,h})^\alpha\right\}.
\]
\end{corollary}

\begin{proof}
Define the truncated multipliers and matrices
\[
  \widetilde\xi_{m,h,i}:=
  \xi_{m,h,i}\mathbf 1_{\{|\xi_{m,h,i}|\le L_{m,h}\}},
  \qquad
  \widetilde A_{m,h,i}:=
  L_{m,h}^{-1}\widetilde\xi_{m,h,i}B_{m,h,i}.
\]
Conditionally on \(\mcG_{m,h}\), the variables
\(\widetilde\xi_{m,h,i}/L_{m,h}\) are independent, symmetric, and bounded by
\(1\) in absolute value.  The conditional variance proxy of the bounded
scalar-multiplier tuple \((\widetilde A_{m,h,i})_i\) is exactly
\(\widehat\nu_{m,h}\).  Hence
Corollary~\ref{cor:conditional-sign-symmetric-bernstein-menu}, or equivalently
Corollary~\ref{cor:symmetric-scalar-multiplier-bernstein-schedules} applied to
the truncated tuple, gives the stated Bernstein active-menu failure bound under
the three displayed schedule hypotheses.

Let
\[
  H_{m,h}:=\{\max_i|\xi_{m,h,i}|\le L_{m,h}\}.
\]
The cap-failure series and the first Borel--Cantelli lemma imply that
\(H_{m,h}\) holds for all sufficiently large active pairs \((m,h)\).  On this
event \(\widetilde A_{m,h,i}=A_{m,h,i}\) for every \(i\), and the \(A_{m,h,i}\)
are contractions.  Combining this eventual cap event with the eventual
discrepancy conclusion for the truncated tuple proves the assertion for
\((A_{m,h,i})_i\).

It remains only to justify the advertised sub-Weibull reduction.  Integrating
the tail estimate gives, conditionally on \(\mcG_{m,h}\),
\[
  \mathbb E(\xi_{m,h,i}^2\mid\mcG_{m,h})
  =
  \int_0^\infty 2t\,
  \mathbb P\{|\xi_{m,h,i}|>t\mid\mcG_{m,h}\}\,dt
  \le
  \frac4\alpha\Gamma(2/\alpha)a_{m,h,i}^2.
\]
Thus \(\widehat q_{m,h,i}\le C_\alpha L_{m,h}^{-2}a_{m,h,i}^2\), while the
same tail display at \(t=L_{m,h}\) gives the cap contribution.  The uniform
proxy bound is immediate.
\end{proof}

\begin{corollary}[Polynomial-moment tail-capped scalar-multiplier schedules]
\label{cor:polynomial-moment-tail-capped-scalar-multiplier-schedules}
In the setting of
Corollary~\ref{cor:tail-capped-symmetric-scalar-multiplier-schedules}, suppose
that for some \(p>0\) there are deterministic proxies \(a_{m,h,i}\ge0\) such
that
\[
  \mathbb E\left(|\xi_{m,h,i}|^p\mid\mcG_{m,h}\right)
  \le a_{m,h,i}^p
  \qquad(1\le h\le J_m,\ 1\le i\le n_m).
\]
Put \(r:=p\wedge2\).  If, outside events \(E_{m,h}\in\mcG_{m,h}\),
\[
  L_{m,h}^{-r}
  \left\|\sum_{i=1}^{n_m}a_{m,h,i}^rB_{m,h,i}^2\right\|_{\op}
  \le \beta\,\frac{n_m}{\ell_m},
\]
and
\[
  \sum_m\sum_{h=1}^{J_m}
  \left[
  \mathbb P(E_{m,h})
  +
  \sum_{i=1}^{n_m}\left(\frac{a_{m,h,i}}{L_{m,h}}\right)^p
  \right]
  <\infty,
\]
then the three schedule conclusions of
Corollary~\ref{cor:tail-capped-symmetric-scalar-multiplier-schedules} hold.
Thus the normalized matrices
\[
  A_{m,h,i}=L_{m,h}^{-1}\xi_{m,h,i}B_{m,h,i}
\]
are eventually contractions and satisfy
\[
  \disc(A_{m,h,1},\ldots,A_{m,h,n_m})\le K\sqrt{n_m}
\]
simultaneously over the active menus under the same general, all-integer, and
sparse stretched-exponential dimension schedules.

In particular, assume \(a_{m,h,i}\le a_{m,h}\) for all \(i\), take
\(\beta>0\), and choose some \(\delta>0\).  If
\[
  L_{m,h}\ge
  a_{m,h}
  \max\left\{
  \left(\frac{\ell_m}{\beta}\right)^{1/r},
  \left(J_mn_m(e+m)^{1+\delta}\right)^{1/p}
  \right\}
\]
for every active \(h\), then the variance and cap hypotheses are automatic
with no \(E_{m,h}\).  Hence polynomial-moment symmetric multipliers of every
tail index \(p>0\) have active-menu Matrix Spencer signings after this explicit
deterministic cap normalization.
\end{corollary}

\begin{proof}
For the cap failures, Markov's inequality gives
\[
  \mathbb P\{|\xi_{m,h,i}|>L_{m,h}\mid\mcG_{m,h}\}
  \le
  \left(\frac{a_{m,h,i}}{L_{m,h}}\right)^p .
\]
It remains to compare the truncated second-moment profile with the displayed
moment profile.  If \(0<p\le2\), then on \(\{|\xi_{m,h,i}|\le L_{m,h}\}\),
\[
  |\xi_{m,h,i}|^2\le
  L_{m,h}^{2-p}|\xi_{m,h,i}|^p,
\]
so
\[
  \widehat q_{m,h,i}
  \le
  L_{m,h}^{-p}a_{m,h,i}^p
  =
  L_{m,h}^{-r}a_{m,h,i}^r.
\]
If \(p>2\), monotonicity of conditional \(L^q\)-norms gives
\[
  \mathbb E(\xi_{m,h,i}^2\mid\mcG_{m,h})
  \le a_{m,h,i}^2,
\]
hence again
\[
  \widehat q_{m,h,i}
  \le
  L_{m,h}^{-2}a_{m,h,i}^2
  =
  L_{m,h}^{-r}a_{m,h,i}^r.
\]
Thus the hypotheses imply the variance and cap hypotheses of
Corollary~\ref{cor:tail-capped-symmetric-scalar-multiplier-schedules}, which
gives the first assertion.

For the uniform proxy statement, use \(0\preceq B_{m,h,i}^2\preceq I\) to get
\[
  L_{m,h}^{-r}
  \left\|\sum_i a_{m,h,i}^rB_{m,h,i}^2\right\|_{\op}
  \le
  n_m\left(\frac{a_{m,h}}{L_{m,h}}\right)^r
  \le
  \beta\,\frac{n_m}{\ell_m}.
\]
The same lower bound on \(L_{m,h}\) gives
\[
  \sum_{h=1}^{J_m}\sum_{i=1}^{n_m}
  \left(\frac{a_{m,h,i}}{L_{m,h}}\right)^p
  \le
  (e+m)^{-1-\delta},
\]
which is summable in \(m\).  Applying the first part proves the explicit
normalization claim.
\end{proof}

\begin{corollary}[Environment-adaptive tail-capped scalar-multiplier schedules]
\label{cor:environment-adaptive-tail-capped-scalar-multiplier-schedules}
In the setting of
Corollary~\ref{cor:tail-capped-symmetric-scalar-multiplier-schedules}, replace
the deterministic caps by positive \(\mcG_{m,h}\)-measurable caps
\(L_{m,h}\).  Define
\[
  A_{m,h,i}:=L_{m,h}^{-1}\xi_{m,h,i}B_{m,h,i},
\]
\[
  \widehat q_{m,h,i}:=
  L_{m,h}^{-2}
  \mathbb E\left(\xi_{m,h,i}^2
  \mathbf 1_{\{|\xi_{m,h,i}|\le L_{m,h}\}}\mid\mcG_{m,h}\right),
\]
and
\[
  \widehat\nu_{m,h}:=
  \left\|\sum_{i=1}^{n_m}\widehat q_{m,h,i}B_{m,h,i}^2\right\|_{\op}.
\]
Fix \(\beta\ge0\) and \(K>0\).  Suppose that, outside events
\(E_{m,h}\in\mcG_{m,h}\),
\[
  \widehat\nu_{m,h}\le \beta\,\frac{n_m}{\ell_m},
\]
and that
\[
  \sum_m\sum_{h=1}^{J_m}
  \left[
  \mathbb P(E_{m,h})
  +
  \mathbb E\sum_{i=1}^{n_m}
  \mathbb P\{|\xi_{m,h,i}|>L_{m,h}\mid\mcG_{m,h}\}
  \right]
  <\infty .
\]
Then the three schedule conclusions of
Corollary~\ref{cor:tail-capped-symmetric-scalar-multiplier-schedules} remain
valid.  Thus, almost surely, all sufficiently large active normalized matrices
are contractions and satisfy
\[
  \disc(A_{m,h,1},\ldots,A_{m,h,n_m})\le K\sqrt{n_m}.
\]

More concretely, let \(p>0\), put \(r=p\wedge2\), and let
\(a_{m,h,i}\) be nonnegative \(\mcG_{m,h}\)-measurable moment proxies with
\[
  \mathbb E(|\xi_{m,h,i}|^p\mid\mcG_{m,h})\le a_{m,h,i}^p .
\]
If, outside \(E_{m,h}\),
\[
  L_{m,h}^{-r}
  \left\|\sum_i a_{m,h,i}^rB_{m,h,i}^2\right\|_{\op}
  \le \beta\,\frac{n_m}{\ell_m},
\]
and
\[
  \sum_m\sum_{h=1}^{J_m}
  \left[
  \mathbb P(E_{m,h})
  +
  \mathbb E\sum_i\left(\frac{a_{m,h,i}}{L_{m,h}}\right)^p
  \right]
  <\infty ,
\]
then the same conclusion holds.

In particular, if \(\beta>0\) and \(a_{m,h,i}\le a_{m,h}\) for an
\(\mcG_{m,h}\)-measurable envelope \(a_{m,h}\), then for any \(\delta>0\) the
conditions are automatic, with no \(E_{m,h}\), whenever
\[
  L_{m,h}\ge
  a_{m,h}
  \max\left\{
  \left(\frac{\ell_m}{\beta}\right)^{1/r},
  \left(J_mn_m(e+m)^{1+\delta}\right)^{1/p}
  \right\}
\]
almost surely for each active \(h\).  When \(a_{m,h}=0\), any positive
\(\mcG_{m,h}\)-measurable cap is allowed.
\end{corollary}

\begin{proof}
Fix \(m,h\) and condition on \(\mcG_{m,h}\).  The cap \(L_{m,h}\), the
directions \(B_{m,h,i}\), and the event \(E_{m,h}\) are then fixed.  Truncate
\[
  \widetilde\xi_{m,h,i}:=
  \xi_{m,h,i}\mathbf 1_{\{|\xi_{m,h,i}|\le L_{m,h}\}},
  \qquad
  \widetilde A_{m,h,i}:=
  L_{m,h}^{-1}\widetilde\xi_{m,h,i}B_{m,h,i}.
\]
Conditionally on \(\mcG_{m,h}\), the tuple
\((\widetilde A_{m,h,i})_i\) is independent, sign-symmetric, and made of
contractions, with conditional variance proxy \(\widehat\nu_{m,h}\).  On
\(E_{m,h}^c\), Corollary~\ref{cor:conditional-sign-symmetric-bernstein-menu}
therefore gives the same active-menu Bernstein failure estimate as in the
deterministic-cap theorem.  The displayed conditional cap-failure summability
and Borel--Cantelli imply that eventually the truncated and untruncated
normalized tuples agree over all active menus.  This proves the adaptive-cap
schedule conclusions.

For the moment-proxy specialization, conditional Markov gives
\[
  \mathbb P\{|\xi_{m,h,i}|>L_{m,h}\mid\mcG_{m,h}\}
  \le \left(\frac{a_{m,h,i}}{L_{m,h}}\right)^p.
\]
As in the proof of
Corollary~\ref{cor:polynomial-moment-tail-capped-scalar-multiplier-schedules},
the truncated second moment obeys
\[
  \widehat q_{m,h,i}
  \le L_{m,h}^{-r}a_{m,h,i}^r.
\]
This gives the second assertion.  The uniform-envelope cap makes
\[
  L_{m,h}^{-r}\left\|\sum_i a_{m,h,i}^rB_{m,h,i}^2\right\|_{\op}
  \le \beta\,\frac{n_m}{\ell_m}
\]
and
\[
  \sum_{h=1}^{J_m}\sum_{i=1}^{n_m}
  \left(\frac{a_{m,h,i}}{L_{m,h}}\right)^p
  \le (e+m)^{-1-\delta},
\]
which is summable.
\end{proof}

\begin{corollary}[Signed cores plus full random operator-tail noise]
\label{cor:signed-core-plus-operator-tail-noise}
Let \(B_1,\ldots,B_n\) be random self-adjoint matrices in \(\M_D(\C)\), let
\(E\) be an event, let \(K\ge0\), and let \(K_0\ge1\).  Suppose that, outside
\(E\),
\[
  \|B_i\|_{\op}\le1\quad (1\le i\le n)
\]
and there are signs \(\eps_1^\star,\ldots,\eps_n^\star\), measurable with
respect to \(B_1,\ldots,B_n\), such that
\[
  \left\|\sum_{i=1}^n\eps_i^\star B_i\right\|_{\op}\le K\sqrt n .
\]
Let \(Y_1,\ldots,Y_n\) be centered self-adjoint random matrices such that,
conditionally on \(B_1,\ldots,B_n\), for every
\(a=(a_1,\ldots,a_n)\in\R^n\) and every \(t\ge0\),
\[
  \mathbb P\left\{
  \left\|\sum_{i=1}^n a_iY_i\right\|_{\op}
  >K_0\|a\|_2(\sqrt D+t)
  \,\middle|\,B_1,\ldots,B_n
  \right\}\le 2e^{-t^2}.
\]
Put
\[
  s_{D,n}:=\sqrt D+\sqrt{3\log(en)},\qquad
  Z_i:=\frac{Y_i}{K_0s_{D,n}},\qquad
  X_i:=B_i+Z_i .
\]
Then
\[
  \begin{aligned}
  &\mathbb P\left\{
  \begin{gathered}
  \disc(X_1,\ldots,X_n)>(K+1)\sqrt n,\\
  \text{or }\max_i\|X_i\|_{\op}>2
  \end{gathered}
  \right\} \\
  &\qquad\le
  \mathbb P(E)+2(n+1)(en)^{-3}.
  \end{aligned}
\]
Consequently, for \(\widehat X_i=X_i/2\),
\[
  \begin{aligned}
  &\mathbb P\left\{
  \begin{gathered}
  \disc(\widehat X_1,\ldots,\widehat X_n)>\frac{K+1}{2}\sqrt n,\\
  \text{or }\max_i\|\widehat X_i\|_{\op}>1
  \end{gathered}
  \right\} \\
  &\qquad\le
  \mathbb P(E)+2(n+1)(en)^{-3}.
  \end{aligned}
\]
Thus any deterministic or random core class whose existing theorem supplies a
\(K\sqrt n\) signing--including finite-lift cores, scalar Toeplitz or Hankel
cores, algebraic cores, and their random versions--remains Spencer-scale after
adding full centered GOE/GUE, Wigner-type, correlated Gaussian, or other
Gaussian-operator-tail noise at its natural spectral scale.
\end{corollary}

\begin{proof}
Condition on a core realization outside \(E\).  Then the signs
\(\eps^\star\) are fixed from the point of view of the conditional tail.  The
conditional operator-tail hypothesis applied with \(a=\eps^\star\) gives
\[
  \mathbb P\left\{
  \left\|\sum_i\eps_i^\star Y_i\right\|_{\op}>
  K_0\sqrt n\,s_{D,n}
  \,\middle|\,B_1,\ldots,B_n
  \right\}
  \le 2(en)^{-3}.
\]
The same conditional tail applied with \(a=e_i\), followed by a union bound over \(i\),
shows that all \(\|Z_i\|_{\op}\le1\) except with conditional probability at most
\(2n(en)^{-3}\).  On the complementary noise event,
\[
  \left\|\sum_i\eps_i^\star X_i\right\|_{\op}
  \le
  \left\|\sum_i\eps_i^\star B_i\right\|_{\op}
  +\left\|\sum_i\eps_i^\star Z_i\right\|_{\op}
  \le (K+1)\sqrt n,
\]
and \(\|X_i\|_{\op}\le\|B_i\|_{\op}+\|Z_i\|_{\op}\le2\) for every \(i\).
Averaging over the core realization gives the displayed probability bound.
\end{proof}

\begin{corollary}[Finite-lift core menus plus full operator-tail noise]
\label{cor:finite-lift-core-operator-tail-menu-borel}
For each \(m\ge1\), let \(D_m\ge2\), let
\[
  J_1\le J_2\le\cdots,\qquad J_m\to\infty,
\]
be a deterministic active-prefix exhaustion, and let \(\mathcal G_m\) be a row
environment.  For each \(h\ge1\), let
\[
  B_{m,1}^{(h)},\ldots,B_{m,n_m}^{(h)},\qquad
  Y_{m,1}^{(h)},\ldots,Y_{m,n_m}^{(h)}
\]
be self-adjoint random matrices in \(\M_{D_m}(\C)\).  Put
\[
  \mathcal H_m=\sigma\left(\mathcal G_m,\,
  B_{m,i}^{(h)}:1\le i\le n_m,\ 1\le h\le J_m\right).
\]
Assume that there are events \(F_m^{\rm lift}\in\mathcal H_m\) with
\(\sum_m\mathbb P(F_m^{\rm lift})<\infty\) such that, outside
\(F_m^{\rm lift}\), the active-prefix core tuples are contractions and admit
finite block lifts with active-prefix loss and entropy
\[
  M_m^{(J)}\le M_0,\qquad E_m^{(J)}\le C_E n_m
\]
in the sense of Corollary~\ref{cor:finite-lift-common-menu}.  Suppose also that
for each \(h\le J_m\) there is \(K_{0,m,h}\ge1\) such that, conditionally on
\(\mathcal H_m\), for every \(a\in\R^{n_m}\) and every \(t\ge0\),
\[
  \mathbb P\left\{
  \left\|\sum_{i=1}^{n_m}a_iY_{m,i}^{(h)}\right\|_{\op}
  >K_{0,m,h}\|a\|_2(\sqrt{D_m}+t)
  \,\middle|\,\mathcal H_m
  \right\}
  \le 2e^{-t^2}.
\]
Let \(\tau_m>0\) be deterministic and define
\[
  Z_{m,i}^{(h)}
  =\frac{Y_{m,i}^{(h)}}{K_{0,m,h}(\sqrt{D_m}+\tau_m)},
  \qquad
  X_{m,i}^{(h)}=B_{m,i}^{(h)}+Z_{m,i}^{(h)} .
\]
If
\[
  \sum_m\left[
  \mathbb P(F_m^{\rm lift})+2J_m(n_m+1)e^{-\tau_m^2}
  \right]<\infty,
\]
then, with probability one, for all sufficiently large \(m\), one
\(\mathcal H_m\)-measurable sign vector \(\eps_m^\star\) satisfies
\[
  \left\|\sum_{i=1}^{n_m}\eps_{m,i}^\star X_{m,i}^{(h)}\right\|_{\op}
  \le (K_{\rm fl}+1)\sqrt{n_m}
  \qquad(1\le h\le J_m),
\]
where \(K_{\rm fl}=C M_0\sqrt{\log(e+C_E)}\).  On the same event,
\[
  \max_{1\le h\le J_m}\max_{1\le i\le n_m}
  \|X_{m,i}^{(h)}\|_{\op}\le2.
\]
Consequently the normalized contractions
\(\widehat X_{m,i}^{(h)}=X_{m,i}^{(h)}/2\) have a common active-prefix signing
of size \((K_{\rm fl}+1)\sqrt{n_m}/2\), and every fixed menu element is
eventually covered.
\end{corollary}

\begin{proof}
On \((F_m^{\rm lift})^c\), Corollary~\ref{cor:finite-lift-common-menu} gives
one sign vector \(\eps_m^\star\) satisfying
\[
  \left\|\sum_i\eps_{m,i}^\star B_{m,i}^{(h)}\right\|_{\op}
  \le K_{\rm fl}\sqrt{n_m}\qquad(1\le h\le J_m).
\]
As before, choose the lexicographically first successful sign vector to make it
\(\mathcal H_m\)-measurable.  Conditional on \(\mathcal H_m\), apply the
operator-tail hypothesis for each \(h\le J_m\) with \(a=\eps_m^\star\) and
with \(a=e_i\), \(1\le i\le n_m\), at \(t=\tau_m\).  A union bound gives
conditional failure probability at most \(2J_m(n_m+1)e^{-\tau_m^2}\).  Outside
this noise event,
\[
  \left\|\sum_i\eps_{m,i}^\star Z_{m,i}^{(h)}\right\|_{\op}\le\sqrt{n_m},
  \qquad
  \max_i\|Z_{m,i}^{(h)}\|_{\op}\le1
  \qquad(1\le h\le J_m).
\]
Combining the core and noise estimates gives the displayed active-prefix
signing and norm bounds.  The summability assumption and Borel--Cantelli give
the almost-sure eventual statement; \(J_m\to\infty\) gives the fixed-menu
coverage.
\end{proof}

\begin{corollary}[Operator-tail schedules for finite-lift core menus]
\label{cor:finite-lift-core-operator-tail-menu-schedules}
In the setting of
Corollary~\ref{cor:finite-lift-core-operator-tail-menu-borel}, put
\[
  L_m^{\rm ot}:=J_m(n_m+1).
\]
If \(\gamma>0\) and
\[
  \tau_m^2\ge \log\bigl(2L_m^{\rm ot}\bigr)+(1+\gamma)\log(e+m)
\]
eventually, while \(\sum_m\mathbb P(F_m^{\rm lift})<\infty\), then the
active-prefix and fixed-menu conclusions of
Corollary~\ref{cor:finite-lift-core-operator-tail-menu-borel} hold almost
surely.  Thus arbitrary dimension growth is allowed after scaling the \(h\)-th
full operator-tail noise by
\[
  K_{0,m,h}\left(\sqrt{D_m}+\tau_m\right).
\]

In particular, suppose \(L_m^{\rm ot}\le(eD_m)^\Gamma\) eventually.  Along all
integer dimensions \(D=2,3,\ldots\), if
\[
  \tau_D=\sqrt{R\log(eD)},\qquad R>\Gamma+1,
\]
then the same active-prefix and fixed-menu conclusions hold.  Along any sparse
schedule \(D_m\ge\exp(c m^\kappa)\), the same conclusions hold whenever
\[
  \tau_m=\sqrt{R\log(eD_m)},\qquad R>\Gamma .
\]
\end{corollary}

\begin{proof}
The arbitrary-growth hypothesis gives
\[
  2L_m^{\rm ot}e^{-\tau_m^2}\le(e+m)^{-(1+\gamma)},
\]
which is summable.  Therefore the summability condition in
Corollary~\ref{cor:finite-lift-core-operator-tail-menu-borel} follows from the
lift-exception summability.

For the all-integer specialization,
\[
  2J_D(n_D+1)e^{-\tau_D^2}
  \le 2(eD)^{\Gamma-R},
\]
which is summable over \(D=2,3,\ldots\) when \(R-\Gamma>1\).  For the sparse
schedule the same bound is a stretched-exponential summable sequence whenever
\(R-\Gamma>0\).
\end{proof}

\begin{corollary}[Finite-lift core menus plus mixed dominated operator-tail noise]
\label{cor:finite-lift-core-mixed-dominated-operator-tail-menu-borel}
Keep the finite-lift core assumptions and notation of
Corollary~\ref{cor:finite-lift-core-operator-tail-menu-borel}, but replace the
displayed full operator-tail hypothesis and its scaling as follows.  Put
\(\Delta_m:=D_m+n_m+1\) and \(\ell_m:=\log(e\Delta_m)\).  For each active
\(h\le J_m\), tag the perturbation stream either \(\psi\) or \(\mu\), choose
\(\varepsilon_{m,h}\in(0,1]\), \(A_{m,h}\ge1\), and \(s_{m,h}>0\), and define
\(r_{m,h}\), \(V_{m,h}\), and \(q_{m,h}\) as in
Corollary~\ref{cor:mixed-dominated-operator-tail-active-libraries}, with
\(Y_{m,i}^{(h)}\) in place of \(X_{m,h,i}\).  Assume that, after adjoining
\(\mathcal H_m\) to the predictable sigma-fields, the corresponding
sub-Weibull or finite-moment dominated hypotheses and the matching truncated
tail-drift hypotheses hold outside events \(E_{m,h}\).

Let \(K_{\rm tail},\eta>0\) and set
\[
  \begin{aligned}
  \lambda_{m,h}&:=\max\{1,K_{\rm tail}+\eta\}\,q_{m,h},\\
  Z_{m,i}^{(h)}&:=\frac{Y_{m,i}^{(h)}}{\lambda_{m,h}},\\
  X_{m,i}^{(h)}&:=B_{m,i}^{(h)}+Z_{m,i}^{(h)} .
  \end{aligned}
\]
If \(\ell_m\le B\sqrt{n_m}\) eventually,
\[
  \sum_m\left[
    \mathbb P(F_m^{\rm lift})
    +\sum_{h=1}^{J_m}\bigl(\mathbb P(E_{m,h})+\varepsilon_{m,h}\bigr)
  \right]<\infty,
\]
and
\[
  \sum_mJ_m\Delta_m
  (e\Delta_m)^{-\gamma_{\rm tc}(K_{\rm tail},1,1,B)}
  <\infty,
\]
then, with probability one, for all sufficiently large \(m\), one
\(\mathcal H_m\)-measurable sign vector \(\eps_m^\star\) satisfies
\[
  \left\|\sum_{i=1}^{n_m}\eps_{m,i}^\star X_{m,i}^{(h)}\right\|_{\op}
  \le (K_{\rm fl}+1)\sqrt{n_m}
  \qquad(1\le h\le J_m),
\]
where \(K_{\rm fl}=C M_0\sqrt{\log(e+C_E)}\).  On the same event,
\[
  \max_{1\le h\le J_m}\max_{1\le i\le n_m}
  \|X_{m,i}^{(h)}\|_{\op}\le2 .
\]
Consequently the normalized contractions
\(\widehat X_{m,i}^{(h)}=X_{m,i}^{(h)}/2\) have a common active-prefix signing
of size \((K_{\rm fl}+1)\sqrt{n_m}/2\), and every fixed menu element is
eventually covered.

In particular, if \(J_m\Delta_m\le(e\Delta_m)^\Gamma\),
\(J_m\le(e\Delta_m)^\Lambda\), \(n_m\ge a\ell_m^2\) eventually, and
\(\varepsilon_{m,h}=(e\Delta_m)^{-c}\), then after adding
\(\sum_m\mathbb P(F_m^{\rm lift})<\infty\) the cap-budget part is summable for
all integer dimensions when \(c>\Lambda+1\), and along every
stretched-exponential sparse dimension schedule when \(c>\Lambda\).  The
matrix-tail margins are
\[
  K_{\rm tail}>\Theta_{\rm tc}(\Gamma+1,a,1,1)
\]
and
\[
  K_{\rm tail}>\Theta_{\rm tc}(\Gamma,a,1,1),
\]
respectively.
\end{corollary}

\begin{proof}
Outside \(F_m^{\rm lift}\), Corollary~\ref{cor:finite-lift-common-menu} gives
one active-prefix signing for the finite-lift core menu:
\[
  \left\|\sum_i\eps_{m,i}^\star B_{m,i}^{(h)}\right\|_{\op}
  \le K_{\rm fl}\sqrt{n_m}
  \qquad(1\le h\le J_m).
\]
As in Corollary~\ref{cor:finite-lift-core-operator-tail-menu-borel}, choose the
lexicographically first successful sign vector to make it
\(\mathcal H_m\)-measurable.  Now apply
Corollary~\ref{cor:signed-core-mixed-dominated-operator-tail-libraries} with
\[
  \begin{gathered}
  B_{m,h,i}=B_{m,i}^{(h)},\qquad
  Y_{m,h,i}=Y_{m,i}^{(h)},\\
  K_{\rm core}=K_{\rm fl},\qquad
  F_m^{\rm core}=F_m^{\rm lift}.
  \end{gathered}
\]
The finite-lift conclusion above supplies the signed-core hypothesis, while the
displayed mixed dominated assumptions are exactly the perturbation hypotheses.
The almost-sure active-prefix statement, fixed-menu coverage, and the
all-integer and sparse schedule thresholds are therefore the corresponding
conclusions of the signed-core mixed dominated corollary.
\end{proof}

\begin{lemma}[Matrix-series sub-Gaussian tail]
\label{lem:matrix-series-subgaussian-tail}
Let \(\mathcal K\) be a sigma-field.  Conditionally on \(\mathcal K\), let
\(\theta_1,\ldots,\theta_N\) be independent centered real variables satisfying
\[
  \mathbb E\{\exp(\lambda\theta_r)\mid\mathcal K\}
  \le \exp(\lambda^2/2)
  \qquad(\lambda\in\R).
\]
Let \(A_1,\ldots,A_N\in\M_D(\C)\) be \(\mathcal K\)-measurable self-adjoint
matrices, and put
\[
  v=\left\|\sum_{r=1}^NA_r^2\right\|_{\op}.
\]
Let
\[
  \rho=\max\left\{1,\operatorname{rank}\left(\sum_{r=1}^NA_r^2\right)\right\}.
\]
There is a universal constant \(C_{\rm ms}\) such that, conditionally on
\(\mathcal K\), for every \(\tau\ge0\),
\[
  \mathbb P\left\{
  \left\|\sum_{r=1}^N\theta_rA_r\right\|_{\op}
  >
  C_{\rm ms}\sqrt v\,(\sqrt{\log(e\rho)}+\tau)
  \,\middle|\,\mathcal K
  \right\}
  \le 2e^{-\tau^2}.
\]
The same conclusion applies to conditional Gaussian, Rademacher, or any
uniformly scalar-sub-Gaussian coefficient family after changing the universal
constant by the corresponding proxy.
\end{lemma}

\begin{proof}
Condition on \(\mathcal K\).  The matrix Laplace bound of Tropp \cite{Tropp},
applied with
\(\mathbb E e^{\lambda\theta_rA_r}\preceq e^{\lambda^2A_r^2/2}\), gives
\[
  \mathbb P\left\{
  \left\|\sum_r\theta_rA_r\right\|_{\op}>s
  \right\}
  \le
  2D\exp\left\{-\frac{s^2}{2v}\right\}
\]
when \(v>0\).  If \(v=0\), then all \(A_r=0\) and the claim is trivial.  When
\(v>0\), let \(V=\sum_rA_r^2\) and \(E=\operatorname{ran}V\).  Since
\[
  \ker V=\bigcap_r\ker A_r,
\]
each \(A_r\) kills \(E^\perp\), and self-adjointness makes \(E\) invariant.
Thus the random series is supported on the \(\rho\)-dimensional space \(E\).
Applying the preceding Laplace bound to the restricted matrices on \(E\)
replaces the factor \(D\) by \(\rho\).  Taking
\(s=C_{\rm ms}\sqrt v(\sqrt{\log(e\rho)}+\tau)\) with \(C_{\rm ms}\) large
enough absorbs this dimension factor and leaves \(2e^{-\tau^2}\).
\end{proof}

\begin{corollary}[Finite-lift cores plus matrix-series noise]
\label{cor:finite-lift-core-matrix-series-noise}
In the setting of
Corollary~\ref{cor:finite-lift-core-operator-tail-menu-borel}, keep the same
finite-lift core assumptions and active-prefix exhaustion.  Let
\(\mathcal K_m\supseteq\mathcal H_m\) be a sigma-field.  For every active
\((m,h,i)\), let \(\mathcal R_{m,h,i}\) be a finite index set and let
\[
  G_{m,h,i,r}\in\M_{D_m}(\C)\qquad(r\in\mathcal R_{m,h,i})
\]
be \(\mathcal K_m\)-measurable self-adjoint matrices.  Conditionally on
\(\mathcal K_m\), let the scalar coefficients
\(\theta_{m,h,i,r}\) be independent, centered, and satisfy the scalar
sub-Gaussian bound in Lemma~\ref{lem:matrix-series-subgaussian-tail}.  Set
\[
  Y_{m,i}^{(h)}
  =
  \sum_{r\in\mathcal R_{m,h,i}}
  \theta_{m,h,i,r}G_{m,h,i,r}.
\]
Let \(q_{m,h}>0\), \(\tau_m>0\), and \(1\le\kappa_{m,h}\le D_m\) be
deterministic.  Let \(F_m^{\rm ms}\in\mathcal K_m\) be the event that for some
\(h\le J_m\) at least one of the following four bounds fails:
\[
  \max_{i\le n_m}
  \left\|\sum_{r\in\mathcal R_{m,h,i}}G_{m,h,i,r}^2\right\|_{\op}
  \le q_{m,h}^2,
\]
\[
  \left\|
  \sum_{i=1}^{n_m}\sum_{r\in\mathcal R_{m,h,i}}
  G_{m,h,i,r}^2
  \right\|_{\op}
  \le n_mq_{m,h}^2,
\]
\[
  \max_{i\le n_m}
  \operatorname{rank}\left(
  \sum_{r\in\mathcal R_{m,h,i}}G_{m,h,i,r}^2
  \right)\le\kappa_{m,h},
\]
\[
  \operatorname{rank}\left(
  \sum_{i=1}^{n_m}\sum_{r\in\mathcal R_{m,h,i}}
  G_{m,h,i,r}^2
  \right)\le\kappa_{m,h}.
\]
Define
\[
  Z_{m,i}^{(h)}
  =
  \frac{Y_{m,i}^{(h)}}
       {C_{\rm ms}q_{m,h}(\sqrt{\log(e\kappa_{m,h})}+\tau_m)},
  \qquad
  X_{m,i}^{(h)}=B_{m,i}^{(h)}+Z_{m,i}^{(h)} .
\]
If
\[
  \sum_m\left[
  \mathbb P(F_m^{\rm lift})
  +\mathbb P(F_m^{\rm ms})
  +2J_m(n_m+1)e^{-\tau_m^2}
  \right]<\infty,
\]
then, with probability one, for all sufficiently large \(m\), the
active-prefix and fixed-menu conclusions of
Corollary~\ref{cor:finite-lift-core-operator-tail-menu-borel} hold for these
matrix-series perturbations.  Thus finite-lift core libraries absorb
conditionally Gaussian, Rademacher, or scalar-sub-Gaussian matrix-series noise
at the sharper scale
\[
  q_{m,h}(\sqrt{\log(e\kappa_{m,h})}+\tau_m),
\]
rather than the generic full-operator-tail scale
\(q_{m,h}(\sqrt{D_m}+\tau_m)\), whenever the realized coefficient square
function satisfies the displayed variance and rank bounds outside a summable
event.  The full-rank choice \(\kappa_{m,h}=D_m\) recovers the preceding
log-dimensional matrix-series scale.

The same arbitrary-growth, all-integer, and sparse-dimension schedules as in
Corollary~\ref{cor:finite-lift-core-operator-tail-menu-schedules} apply with
this log-dimensional matrix-series scale: if \(J_m(n_m+1)\le(eD_m)^\Gamma\),
then \(\tau_D^2=S\log(eD)\) with \(S>\Gamma+1\) suffices along all integer
dimensions, while sparse stretched-exponential dimension schedules need only
\(S>\Gamma\).
\end{corollary}

\begin{proof}
Work outside \(F_m^{\rm lift}\cup F_m^{\rm ms}\).  The finite-lift common-menu
corollary gives one active-prefix core signing \(\eps_m^\star\), which is
\(\mathcal H_m\)-measurable and hence \(\mathcal K_m\)-measurable.  Conditional
on \(\mathcal K_m\), Lemma~\ref{lem:matrix-series-subgaussian-tail} applied to
the signed series
\[
  \sum_{i=1}^{n_m}\eps_{m,i}^\star Y_{m,i}^{(h)}
  =
  \sum_{i,r}\theta_{m,h,i,r}\eps_{m,i}^\star G_{m,h,i,r}
\]
has variance parameter at most \(n_mq_{m,h}^2\) and rank parameter at most
\(\kappa_{m,h}\).  Therefore, except with conditional probability at most
\(2e^{-\tau_m^2}\),
\[
  \left\|\sum_i\eps_{m,i}^\star Z_{m,i}^{(h)}\right\|_{\op}
  \le \sqrt{n_m}.
\]
The same lemma applied to each individual \(Y_{m,i}^{(h)}\) has variance
parameter at most \(q_{m,h}^2\) and rank parameter at most
\(\kappa_{m,h}\), and gives
\(\|Z_{m,i}^{(h)}\|_{\op}\le1\) except with conditional probability at most
\(2e^{-\tau_m^2}\).  A union bound over \(h\le J_m\) and
\(i=1,\ldots,n_m\), together with the signed-sum events, gives the failure
probability \(2J_m(n_m+1)e^{-\tau_m^2}\).  On the complementary event, the
same estimates as in
Corollary~\ref{cor:finite-lift-core-operator-tail-menu-borel} give the
active-prefix discrepancy and contraction bounds.  Borel--Cantelli gives the
almost-sure eventual conclusion, and the schedule paragraph is the same
summability calculation as in
Corollary~\ref{cor:finite-lift-core-operator-tail-menu-schedules}.
\end{proof}

\begin{corollary}[Low-rank finite-lift matrix-series noise]
\label{cor:finite-lift-core-low-rank-matrix-series-noise}
In the setting of
Corollary~\ref{cor:finite-lift-core-matrix-series-noise}, suppose that for
each active \(h\le J_m\) there is a \(\mathcal K_m\)-measurable orthogonal
projection \(P_{m,h}\) with
\[
  \operatorname{rank}P_{m,h}\le\kappa_{m,h}
\]
such that every active coefficient direction satisfies
\[
  G_{m,h,i,r}=P_{m,h}G_{m,h,i,r}P_{m,h}.
\]
Let \(F_m^{\rm lr}\in\mathcal K_m\) be the event that for some \(h\le J_m\)
one of the two square-function bounds
\[
  \max_{i\le n_m}
  \left\|\sum_{r\in\mathcal R_{m,h,i}}G_{m,h,i,r}^2\right\|_{\op}
  \le q_{m,h}^2
\]
or
\[
  \left\|
  \sum_{i=1}^{n_m}\sum_{r\in\mathcal R_{m,h,i}}
  G_{m,h,i,r}^2
  \right\|_{\op}
  \le n_mq_{m,h}^2
\]
fails.  If
\[
  \sum_m\left[
  \mathbb P(F_m^{\rm lift})
  +\mathbb P(F_m^{\rm lr})
  +2J_m(n_m+1)e^{-\tau_m^2}
  \right]<\infty,
\]
then, with probability one, for all sufficiently large \(m\), the
active-prefix and fixed-menu conclusions of
Corollary~\ref{cor:finite-lift-core-operator-tail-menu-borel} hold after
scaling the \(h\)-th low-rank matrix-series noise stream by
\[
  C_{\rm ms}q_{m,h}(\sqrt{\log(e\kappa_{m,h})}+\tau_m).
\]
Thus finite-lift core libraries absorb conditionally Gaussian, Rademacher, or
scalar-sub-Gaussian matrix-series noise supported on a common
\(\kappa_{m,h}\)-dimensional subspace at the rank-sensitive spectral scale.
This covers spiked matrix-series perturbations, block-localized matrix-series
noise, and random-subspace matrix-series perturbations after conditioning on
the support projection.  If the two square-function bounds hold
deterministically, then \(F_m^{\rm lr}=\emptyset\).

The arbitrary-growth, all-integer, and sparse-dimension schedules are the same
as in Corollary~\ref{cor:finite-lift-core-operator-tail-menu-schedules}; only
the scale changes from \(\sqrt{\log(eD_m)}\) to
\(\sqrt{\log(e\kappa_{m,h})}\).
\end{corollary}

\begin{proof}
The common support projection gives
\[
  \operatorname{ran}\left(\sum_rG_{m,h,i,r}^2\right)
  \subseteq\operatorname{ran}P_{m,h}
  \qquad\text{and}\qquad
  \operatorname{ran}\left(\sum_{i,r}G_{m,h,i,r}^2\right)
  \subseteq\operatorname{ran}P_{m,h}.
\]
Hence both rank bounds required in
Corollary~\ref{cor:finite-lift-core-matrix-series-noise} are automatic with
rank envelope \(\kappa_{m,h}\).  On \((F_m^{\rm lr})^c\), the two variance
bounds required there also hold, so the corollary applies with
\(F_m^{\rm ms}=F_m^{\rm lr}\).  The examples are exactly the displayed
support-projection choices, and the schedule statement is unchanged.
\end{proof}

\begin{corollary}[Atomic norm-budget low-rank matrix-series noise]
\label{cor:finite-lift-core-atomic-low-rank-matrix-series-noise}
In the setting of
Corollary~\ref{cor:finite-lift-core-low-rank-matrix-series-noise}, suppose the
common support projections \(P_{m,h}\) are given as there.  Suppose also that
there are \(\mathcal K_m\)-measurable nonnegative amplitudes
\(b_{m,h,i,r}\) such that every active coefficient direction satisfies
\[
  G_{m,h,i,r}=P_{m,h}G_{m,h,i,r}P_{m,h},
  \qquad
  \|G_{m,h,i,r}\|_{\op}\le b_{m,h,i,r}.
\]
Let \(F_m^{\rm at}\in\mathcal K_m\) be the event that for some \(h\le J_m\)
one of the two scalar atom-budget bounds
\[
  \max_{i\le n_m}\sum_{r\in\mathcal R_{m,h,i}}b_{m,h,i,r}^2
  \le q_{m,h}^2
\]
or
\[
  \sum_{i=1}^{n_m}\sum_{r\in\mathcal R_{m,h,i}}b_{m,h,i,r}^2
  \le n_mq_{m,h}^2
\]
fails.  If
\[
  \sum_m\left[
  \mathbb P(F_m^{\rm lift})
  +\mathbb P(F_m^{\rm at})
  +2J_m(n_m+1)e^{-\tau_m^2}
  \right]<\infty,
\]
then, with probability one, for all sufficiently large \(m\), the
active-prefix and fixed-menu conclusions of
Corollary~\ref{cor:finite-lift-core-operator-tail-menu-borel} hold after
scaling the \(h\)-th atomic matrix-series noise stream by
\[
  C_{\rm ms}q_{m,h}(\sqrt{\log(e\kappa_{m,h})}+\tau_m).
\]
Thus finite-lift core libraries absorb rank-one spike fields, self-adjoint
edge-atom series, and finite frame/dictionary matrix-series noise inside a
common low-rank support whenever the scalar squared-amplitude budgets are
summable outside \(F_m^{\rm at}\).  If the scalar budgets hold deterministically,
then \(F_m^{\rm at}=\emptyset\).
\end{corollary}

\begin{proof}
For a self-adjoint matrix supported on \(P_{m,h}\),
\[
  G_{m,h,i,r}^2\preceq \|G_{m,h,i,r}\|_{\op}^2P_{m,h}
  \preceq b_{m,h,i,r}^2P_{m,h}.
\]
Therefore, outside \(F_m^{\rm at}\),
\[
  \left\|\sum_rG_{m,h,i,r}^2\right\|_{\op}
  \le \sum_rb_{m,h,i,r}^2\le q_{m,h}^2
\]
for every active \(i,h\), and
\[
  \left\|\sum_{i,r}G_{m,h,i,r}^2\right\|_{\op}
  \le \sum_{i,r}b_{m,h,i,r}^2\le n_mq_{m,h}^2.
\]
The common support projection gives the rank bounds exactly as in
Corollary~\ref{cor:finite-lift-core-low-rank-matrix-series-noise}.  Applying
that corollary with \(F_m^{\rm lr}=F_m^{\rm at}\) proves the conclusion.
\end{proof}

\begin{corollary}[Block-load atomic matrix-series noise]
\label{cor:finite-lift-core-block-load-atomic-matrix-series-noise}
In the setting of
Corollary~\ref{cor:finite-lift-core-matrix-series-noise}, suppose that for
each active \(h\le J_m\) there is a finite nonempty block set
\(\mathcal B_{m,h}\) and pairwise orthogonal
\(\mathcal K_m\)-measurable projections
\[
  P_{m,h,\beta}\qquad(\beta\in\mathcal B_{m,h})
\]
such that
\[
  \operatorname{rank}\left(\sum_{\beta\in\mathcal B_{m,h}}P_{m,h,\beta}\right)
  \le\kappa_{m,h}.
\]
For every active coefficient direction let
\(S_{m,h,i,r}\subseteq\mathcal B_{m,h}\) be a
\(\mathcal K_m\)-measurable support set, put
\[
  P_{m,h,S_{m,h,i,r}}
  =\sum_{\beta\in S_{m,h,i,r}}P_{m,h,\beta},
\]
and assume that for some \(\mathcal K_m\)-measurable nonnegative amplitudes
\(b_{m,h,i,r}\),
\[
  G_{m,h,i,r}
  =
  P_{m,h,S_{m,h,i,r}}G_{m,h,i,r}P_{m,h,S_{m,h,i,r}},
  \qquad
  \|G_{m,h,i,r}\|_{\op}\le b_{m,h,i,r}.
\]
Let \(F_m^{\rm bl}\in\mathcal K_m\) be the event that for some \(h\le J_m\)
one of the two block-load bounds
\[
  \max_{i\le n_m}\max_{\beta\in\mathcal B_{m,h}}
  \sum_{\substack{r\in\mathcal R_{m,h,i}\\
                  \beta\in S_{m,h,i,r}}}
  b_{m,h,i,r}^2
  \le q_{m,h}^2
\]
or
\[
  \max_{\beta\in\mathcal B_{m,h}}
  \sum_{i=1}^{n_m}
  \sum_{\substack{r\in\mathcal R_{m,h,i}\\
                  \beta\in S_{m,h,i,r}}}
  b_{m,h,i,r}^2
  \le n_mq_{m,h}^2
\]
fails.  If
\[
  \sum_m\left[
  \mathbb P(F_m^{\rm lift})
  +\mathbb P(F_m^{\rm bl})
  +2J_m(n_m+1)e^{-\tau_m^2}
  \right]<\infty,
\]
then, with probability one, for all sufficiently large \(m\), the
active-prefix and fixed-menu conclusions of
Corollary~\ref{cor:finite-lift-core-operator-tail-menu-borel} hold after
scaling the \(h\)-th block-atomic matrix-series noise stream by
\[
  C_{\rm ms}q_{m,h}(\sqrt{\log(e\kappa_{m,h})}+\tau_m).
\]
Thus block-diagonal atom libraries, self-adjoint off-diagonal edge atoms, and
cluster/frame dictionaries supported on a low-total-rank orthogonal block
family are controlled by their maximum incident squared-amplitude load, rather
than by the total squared-amplitude load.
\end{corollary}

\begin{proof}
For each active atom, self-adjoint support on \(P_{m,h,S_{m,h,i,r}}\) gives
\[
  G_{m,h,i,r}^2
  \preceq b_{m,h,i,r}^2P_{m,h,S_{m,h,i,r}}.
\]
Hence
\[
  \sum_{r\in\mathcal R_{m,h,i}}G_{m,h,i,r}^2
  \preceq
  \sum_{\beta\in\mathcal B_{m,h}}
  \left(
  \sum_{\substack{r\in\mathcal R_{m,h,i}\\
                  \beta\in S_{m,h,i,r}}}
  b_{m,h,i,r}^2
  \right)P_{m,h,\beta}.
\]
The projections are orthogonal, so the operator norm of the right side is the
largest displayed block load, at most \(q_{m,h}^2\) outside \(F_m^{\rm bl}\).
The same calculation after summing over \(i\) gives
\[
  \left\|\sum_{i,r}G_{m,h,i,r}^2\right\|_{\op}\le n_mq_{m,h}^2.
\]
Both square functions are supported on
\(\sum_{\beta\in\mathcal B_{m,h}}P_{m,h,\beta}\), whose rank is at most
\(\kappa_{m,h}\).  Therefore all four hypotheses defining
\((F_m^{\rm ms})^c\) in
Corollary~\ref{cor:finite-lift-core-matrix-series-noise} hold on
\((F_m^{\rm bl})^c\).  Apply that corollary with
\(F_m^{\rm ms}=F_m^{\rm bl}\).
\end{proof}

\begin{corollary}[Bounded weighted block-design matrix-series noise]
\label{cor:finite-lift-core-bounded-weighted-block-design-matrix-series-noise}
In the setting of
Corollary~\ref{cor:finite-lift-core-block-load-atomic-matrix-series-noise},
suppose that for every active \(h\le J_m\) there are deterministic
\(q_{m,h}\ge0\), \(s_{m,h}>0\), and weights
\(\omega_{m,h,i,r,\beta}\ge0\).  Let \(F_m^{\rm wbd}\in\mathcal K_m\) be the
event that, for some \(h\le J_m\), at least one of the following weighted
block-design bounds fails:
\[
  b_{m,h,i,r}^2
  \le
  \frac{q_{m,h}^2}{s_{m,h}}\,
  \omega_{m,h,i,r,\beta}
  \qquad
  \text{whenever }\beta\in S_{m,h,i,r},
\]
\[
  \max_{i\le n_m}\max_{\beta\in\mathcal B_{m,h}}
  \sum_{\substack{r\in\mathcal R_{m,h,i}\\
                  \beta\in S_{m,h,i,r}}}
  \omega_{m,h,i,r,\beta}
  \le s_{m,h},
\]
and
\[
  \max_{\beta\in\mathcal B_{m,h}}
  \sum_{i=1}^{n_m}
  \sum_{\substack{r\in\mathcal R_{m,h,i}\\
                  \beta\in S_{m,h,i,r}}}
  \omega_{m,h,i,r,\beta}
  \le n_ms_{m,h}.
\]
If
\[
  \sum_m\left[
  \mathbb P(F_m^{\rm lift})
  +\mathbb P(F_m^{\rm wbd})
  +2J_m(n_m+1)e^{-\tau_m^2}
  \right]<\infty,
\]
then, with probability one, for all sufficiently large \(m\), the
active-prefix and fixed-menu conclusions of
Corollary~\ref{cor:finite-lift-core-operator-tail-menu-borel} hold after
scaling the \(h\)-th weighted block-design matrix-series noise stream by
\[
  C_{\rm ms}q_{m,h}(\sqrt{\log(e\kappa_{m,h})}+\tau_m).
\]

In particular, if the support law is almost surely supported on the displayed
weighted block-design constraints, then \(F_m^{\rm wbd}=\emptyset\).  In the
unweighted case, if every active block is incident to at most \(d_{m,h}\) atoms
from each row and at most \(n_md_{m,h}\) atoms in total, and if
\(b_{m,h,i,r}\le q_{m,h}/\sqrt{d_{m,h}}\) for \(d_{m,h}\ge1\), then the same
conclusion holds with no support-concentration payment.  Thus random regular
block hypergraphs, matching and permutation block designs, balanced
without-replacement assignments, and deterministic or randomly chosen
weighted-degree block samplers enter at the same rank-sensitive
matrix-series scale whenever their weighted block degrees are deterministic or
fail only summably.

The explicit operator-tail schedules are exactly those of
Corollary~\ref{cor:finite-lift-core-operator-tail-menu-schedules}: if
\(J_m(n_m+1)\le(eD_m)^\Gamma\), then all integer dimensions are covered by
\(\tau_D^2=S\log(eD)\) with \(S>\Gamma+1\), while sparse
stretched-exponential dimension schedules need only \(S>\Gamma\).
\end{corollary}

\begin{proof}
On \((F_m^{\rm wbd})^c\), for every active \(h,i,\beta\),
\[
\begin{aligned}
  \sum_{\substack{r\in\mathcal R_{m,h,i}\\
                  \beta\in S_{m,h,i,r}}}
  b_{m,h,i,r}^2
  &\le
  \frac{q_{m,h}^2}{s_{m,h}}
  \sum_{\substack{r\in\mathcal R_{m,h,i}\\
                  \beta\in S_{m,h,i,r}}}
  \omega_{m,h,i,r,\beta}  \\
  &\le q_{m,h}^2 .
\end{aligned}
\]
The same calculation after summing over all \(i\) gives, for every
\(\beta\in\mathcal B_{m,h}\),
\[
  \sum_{i=1}^{n_m}
  \sum_{\substack{r\in\mathcal R_{m,h,i}\\
                  \beta\in S_{m,h,i,r}}}
  b_{m,h,i,r}^2
  \le n_mq_{m,h}^2 .
\]
Therefore the block-load failure event \(F_m^{\rm bl}\) from
Corollary~\ref{cor:finite-lift-core-block-load-atomic-matrix-series-noise}
is contained in \(F_m^{\rm wbd}\).  Applying that corollary gives the
almost-sure active-prefix and fixed-menu conclusions.  The examples are the
corresponding choices of deterministic weights \(\omega\), and the schedule
statement is
Corollary~\ref{cor:finite-lift-core-operator-tail-menu-schedules}.
\end{proof}

\begin{corollary}[Bounded-overlap block-hypergraph matrix-series noise]
\label{cor:finite-lift-core-bounded-overlap-block-hypergraph-matrix-series-noise}
In the setting of
Corollary~\ref{cor:finite-lift-core-block-load-atomic-matrix-series-noise},
suppose that for every active \(h\le J_m\) there is a deterministic
degree parameter \(d_{m,h}\ge1\).  Let \(F_m^{\rm bhg}\in\mathcal K_m\) be the
event that, for some \(h\le J_m\), at least one of the two block-overlap
bounds
\[
  \max_{i\le n_m}\max_{\beta\in\mathcal B_{m,h}}
  \#\{r\in\mathcal R_{m,h,i}:\beta\in S_{m,h,i,r}\}
  \le d_{m,h}
\]
or
\[
  \max_{\beta\in\mathcal B_{m,h}}
  \sum_{i=1}^{n_m}
  \#\{r\in\mathcal R_{m,h,i}:\beta\in S_{m,h,i,r}\}
  \le n_md_{m,h}
\]
fails.  Assume also that, on \((F_m^{\rm bhg})^c\),
\[
  b_{m,h,i,r}\le \frac{q_{m,h}}{\sqrt{d_{m,h}}}
  \qquad(r\in\mathcal R_{m,h,i})
\]
for deterministic \(q_{m,h}\ge0\).  If
\[
  \sum_m\left[
  \mathbb P(F_m^{\rm lift})
  +\mathbb P(F_m^{\rm bhg})
  +2J_m(n_m+1)e^{-\tau_m^2}
  \right]<\infty,
\]
then, with probability one, for all sufficiently large \(m\), the
active-prefix and fixed-menu conclusions of
Corollary~\ref{cor:finite-lift-core-operator-tail-menu-borel} hold after
scaling the \(h\)-th bounded-overlap block-hypergraph matrix-series noise
stream by
\[
  C_{\rm ms}q_{m,h}(\sqrt{\log(e\kappa_{m,h})}+\tau_m).
\]

In particular, if the support law is supported on block hypergraphs with the
displayed row and total block-overlap degrees, then \(F_m^{\rm bhg}=\emptyset\).
This includes deterministic block designs, random regular block hypergraphs,
balanced block dictionaries, and without-replacement block-hyperedge samplers
conditioned on their degree sequence.  Their all-integer and sparse schedules
are exactly the usual operator-tail schedules, with no block-load concentration
payment.
\end{corollary}

\begin{proof}
Apply
Corollary~\ref{cor:finite-lift-core-bounded-weighted-block-design-matrix-series-noise}
with \(s_{m,h}=d_{m,h}\) and
\[
  \omega_{m,h,i,r,\beta}
  =
  {\bf 1}_{\{\beta\in S_{m,h,i,r}\}}.
\]
On \((F_m^{\rm bhg})^c\), the two displayed overlap bounds are exactly the row
and total weighted-degree constraints.  The amplitude condition gives
\[
  b_{m,h,i,r}^2\le q_{m,h}^2/d_{m,h}
  =
  q_{m,h}^2\omega_{m,h,i,r,\beta}/d_{m,h}
\]
whenever \(\beta\in S_{m,h,i,r}\).  Hence the weighted block-design failure
event is contained in \(F_m^{\rm bhg}\), and the cited corollary proves the
claim.
\end{proof}

\begin{corollary}[Heterogeneous bounded-overlap block-hypergraph matrix-series noise]
\label{cor:finite-lift-core-heterogeneous-bounded-overlap-block-hypergraph-matrix-series-noise}
In the setting of
Corollary~\ref{cor:finite-lift-core-block-load-atomic-matrix-series-noise},
suppose that for every active \(h\le J_m\) there are deterministic local degree
profiles
\[
  \delta_{m,h,i,\beta}^{\rm ind}\ge1,
  \qquad
  \Delta_{m,h,\beta}^{\rm tot}\ge1
  \qquad(i\le n_m,\ \beta\in\mathcal B_{m,h}).
\]
Let \(F_m^{\rm hbhg}\in\mathcal K_m\) be the event that, for some \(h\le J_m\),
at least one of the heterogeneous block-overlap bounds
\[
  \#\{r\in\mathcal R_{m,h,i}:\beta\in S_{m,h,i,r}\}
  \le \delta_{m,h,i,\beta}^{\rm ind}
  \qquad(i\le n_m,\ \beta\in\mathcal B_{m,h})
\]
or
\[
  \sum_{i=1}^{n_m}
  \#\{r\in\mathcal R_{m,h,i}:\beta\in S_{m,h,i,r}\}
  \le \Delta_{m,h,\beta}^{\rm tot}
  \qquad(\beta\in\mathcal B_{m,h})
\]
fails.  Assume also that, on \((F_m^{\rm hbhg})^c\), every active incidence
\(\beta\in S_{m,h,i,r}\) satisfies
\[
  b_{m,h,i,r}^2
  \le
  \min\left\{
    \frac{q_{m,h}^2}{\delta_{m,h,i,\beta}^{\rm ind}},
    \frac{n_mq_{m,h}^2}{\Delta_{m,h,\beta}^{\rm tot}}
  \right\}
\]
for deterministic \(q_{m,h}\ge0\).  If
\[
  \sum_m\left[
  \mathbb P(F_m^{\rm lift})
  +\mathbb P(F_m^{\rm hbhg})
  +2J_m(n_m+1)e^{-\tau_m^2}
  \right]<\infty,
\]
then, with probability one, for all sufficiently large \(m\), the
active-prefix and fixed-menu conclusions of
Corollary~\ref{cor:finite-lift-core-operator-tail-menu-borel} hold after
scaling the \(h\)-th heterogeneous bounded-overlap block-hypergraph
matrix-series noise stream by
\[
  C_{\rm ms}q_{m,h}(\sqrt{\log(e\kappa_{m,h})}+\tau_m).
\]

Thus irregular degree-certified block hypergraphs, stratified block
dictionaries, and conditioned configuration-model block samplers are covered at
the local-degree normalization scale, with no reduction to the worst active
block degree and no support-concentration payment when the displayed degree
profiles hold deterministically.
\end{corollary}

\begin{proof}
On \((F_m^{\rm hbhg})^c\), for every active \(h,i,\beta\),
\[
\begin{aligned}
  \sum_{\substack{r\in\mathcal R_{m,h,i}\\
                  \beta\in S_{m,h,i,r}}}
  b_{m,h,i,r}^2
  &\le
  \#\{r\in\mathcal R_{m,h,i}:\beta\in S_{m,h,i,r}\}
  \frac{q_{m,h}^2}{\delta_{m,h,i,\beta}^{\rm ind}}\\
  &\le q_{m,h}^2 .
\end{aligned}
\]
Similarly, for every active \(h,\beta\),
\[
\begin{aligned}
  \sum_{i=1}^{n_m}
  \sum_{\substack{r\in\mathcal R_{m,h,i}\\
                  \beta\in S_{m,h,i,r}}}
  b_{m,h,i,r}^2
  &\le
  \left(\sum_{i=1}^{n_m}
  \#\{r\in\mathcal R_{m,h,i}:\beta\in S_{m,h,i,r}\}\right)
  \frac{n_mq_{m,h}^2}{\Delta_{m,h,\beta}^{\rm tot}}\\
  &\le n_mq_{m,h}^2 .
\end{aligned}
\]
Therefore the block-load failure event \(F_m^{\rm bl}\) from
Corollary~\ref{cor:finite-lift-core-block-load-atomic-matrix-series-noise}
is contained in \(F_m^{\rm hbhg}\), and that corollary gives the asserted
finite-lift active-prefix and fixed-menu conclusions.
\end{proof}

\begin{corollary}[Row-wise block-permutation matrix-series noise]
\label{cor:finite-lift-core-row-wise-block-permutation-matrix-series-noise}
In the setting of
Corollary~\ref{cor:finite-lift-core-block-load-atomic-matrix-series-noise},
suppose that every active atom has singleton block support
\[
  S_{m,h,i,r}=\{\beta_{m,h,i,r}\}
  \subseteq\mathcal B_{m,h}.
\]
For each active \(h\le J_m\), let \(d_{m,h}\ge1\) be deterministic and suppose
that each row index set \(\mathcal R_{m,h,i}\) is decomposed into
\[
  \mathcal R_{m,h,i}
  =
  \mathcal R_{m,h,i}^{(1)}\cup\cdots\cup
  \mathcal R_{m,h,i}^{(d_{m,h})},
\]
so that, outside an event \(F_m^{\rm perm}\in\mathcal K_m\), each layer is
block-injective:
\[
  \#\{r\in\mathcal R_{m,h,i}^{(\ell)}:
       \beta_{m,h,i,r}=\beta\}
  \le1
  \qquad
  (i\le n_m,\ \beta\in\mathcal B_{m,h},\ 1\le\ell\le d_{m,h}).
\]
Assume also that, on \((F_m^{\rm perm})^c\),
\[
  b_{m,h,i,r}\le \frac{q_{m,h}}{\sqrt{d_{m,h}}}
  \qquad(r\in\mathcal R_{m,h,i})
\]
for deterministic \(q_{m,h}\ge0\).  If
\[
  \sum_m\left[
  \mathbb P(F_m^{\rm lift})
  +\mathbb P(F_m^{\rm perm})
  +2J_m(n_m+1)e^{-\tau_m^2}
  \right]<\infty,
\]
then, with probability one, for all sufficiently large \(m\), the
active-prefix and fixed-menu conclusions of
Corollary~\ref{cor:finite-lift-core-operator-tail-menu-borel} hold after
scaling the \(h\)-th row-wise block-permutation matrix-series noise stream by
\[
  C_{\rm ms}q_{m,h}(\sqrt{\log(e\kappa_{m,h})}+\tau_m).
\]

In particular, \(F_m^{\rm perm}=\emptyset\) for laws supported on row-wise
block permutations, row-wise partial matchings, \(d_{m,h}\)-round balanced
without-replacement block assignments, or mixtures of such laws.  Thus these
random block assignment models have the same active-prefix finite-lift schedule
as deterministic block designs: if \(J_m(n_m+1)\le(eD_m)^\Gamma\), all integer
dimensions are covered by \(\tau_D^2=S\log(eD)\) with \(S>\Gamma+1\), while
sparse stretched-exponential dimension schedules need only \(S>\Gamma\).
\end{corollary}

\begin{proof}
Apply
Corollary~\ref{cor:finite-lift-core-bounded-weighted-block-design-matrix-series-noise}
with \(s_{m,h}=d_{m,h}\) and
\[
  \omega_{m,h,i,r,\beta}
  =
  {\bf 1}_{\{\beta=\beta_{m,h,i,r}\}}.
\]
On \((F_m^{\rm perm})^c\), the block-injectivity of each layer gives row-block
weighted degree at most \(d_{m,h}\).  Summing this row bound over
\(i\le n_m\) gives total block weighted degree at most \(n_md_{m,h}\).  The
amplitude cap is exactly
\[
  b_{m,h,i,r}^2\le q_{m,h}^2/d_{m,h}
  =
  q_{m,h}^2\omega_{m,h,i,r,\beta_{m,h,i,r}}/d_{m,h}.
\]
Thus the weighted block-design failure event is contained in
\(F_m^{\rm perm}\), and the cited corollary proves the claim.
\end{proof}

\begin{corollary}[Random block-load atomic matrix-series noise]
\label{cor:finite-lift-core-random-block-load-atomic-matrix-series-noise}
In the setting of
Corollary~\ref{cor:finite-lift-core-block-load-atomic-matrix-series-noise},
suppose that the block projections and amplitudes are measurable with respect
to a sigma-field \(\mathcal L_m\subseteq\mathcal K_m\).  Write
\[
  \eta_{m,h,i,r,\beta}
  ={\bf 1}_{\{\beta\in S_{m,h,i,r}\}}.
\]
Assume that, for each fixed active \(h\le J_m\) and
\(\beta\in\mathcal B_{m,h}\), the variables
\(\eta_{m,h,i,r,\beta}\) are conditionally independent over the pairs
\((i,r)\), given \(\mathcal L_m\), and that
\[
  \mathbb P\{\eta_{m,h,i,r,\beta}=1\mid\mathcal L_m\}
  \le p_{m,h,i,r,\beta}.
\]
Here the \(p_{m,h,i,r,\beta}\) are \(\mathcal L_m\)-measurable.  Let
deterministic envelopes \(\mu_{m,h}^{\rm ind}\),
\(\sigma_{m,h}^{2,{\rm ind}}\), \(\mu_{m,h}^{\rm tot}\),
\(\sigma_{m,h}^{2,{\rm tot}}\), and \(\omega_{m,h}\) satisfy, almost surely,
\[
  \max_{i,\beta}
  \sum_{r\in\mathcal R_{m,h,i}}
  b_{m,h,i,r}^2p_{m,h,i,r,\beta}
  \le \mu_{m,h}^{\rm ind},
\]
\[
  \max_{i,\beta}
  \sum_{r\in\mathcal R_{m,h,i}}
  b_{m,h,i,r}^4p_{m,h,i,r,\beta}
  \le \sigma_{m,h}^{2,{\rm ind}},
\]
\[
  \max_{\beta}
  \sum_{i=1}^{n_m}\sum_{r\in\mathcal R_{m,h,i}}
  b_{m,h,i,r}^2p_{m,h,i,r,\beta}
  \le \mu_{m,h}^{\rm tot},
\]
\[
  \max_{\beta}
  \sum_{i=1}^{n_m}\sum_{r\in\mathcal R_{m,h,i}}
  b_{m,h,i,r}^4p_{m,h,i,r,\beta}
  \le \sigma_{m,h}^{2,{\rm tot}},
  \qquad
  \max_{i,r} b_{m,h,i,r}^2\le\omega_{m,h}.
\]
Let \(u_{m,h}^{\rm ind},u_{m,h}^{\rm tot}\ge0\) be deterministic and suppose
\[
  \mu_{m,h}^{\rm ind}+u_{m,h}^{\rm ind}\le q_{m,h}^2,
  \qquad
  \mu_{m,h}^{\rm tot}+u_{m,h}^{\rm tot}\le n_mq_{m,h}^2.
\]
With the convention that a Bernstein term with zero denominator is omitted,
put
\[
\begin{aligned}
  \mathfrak B_m
  :=
  \sum_{h=1}^{J_m}|\mathcal B_{m,h}|
  \Bigg[
  &n_m\exp\left\{
  -\frac{(u_{m,h}^{\rm ind})^2}
         {2(\sigma_{m,h}^{2,{\rm ind}}
            +\omega_{m,h}u_{m,h}^{\rm ind}/3)}
  \right\}\\
  &+
  \exp\left\{
  -\frac{(u_{m,h}^{\rm tot})^2}
         {2(\sigma_{m,h}^{2,{\rm tot}}
            +\omega_{m,h}u_{m,h}^{\rm tot}/3)}
  \right\}
  \Bigg].
\end{aligned}
\]
If
\[
  \sum_m\left[
  \mathbb P(F_m^{\rm lift})
  +\mathfrak B_m
  +2J_m(n_m+1)e^{-\tau_m^2}
  \right]<\infty,
\]
then, with probability one, for all sufficiently large \(m\), the
active-prefix and fixed-menu conclusions of
Corollary~\ref{cor:finite-lift-core-operator-tail-menu-borel} hold after
scaling the \(h\)-th random block-atomic matrix-series noise stream by
\[
  C_{\rm ms}q_{m,h}(\sqrt{\log(e\kappa_{m,h})}+\tau_m).
\]
In particular, independent random block choices, random edge-atom supports,
and random clustered frame dictionaries are covered whenever their conditional
incident block loads have summable scalar-Bernstein tails.
\end{corollary}

\begin{proof}
For fixed \(m,h,i,\beta\), condition on \(\mathcal L_m\) and apply scalar
Bernstein to the independent variables
\[
  b_{m,h,i,r}^2
  \left(\eta_{m,h,i,r,\beta}
  -\mathbb E[\eta_{m,h,i,r,\beta}\mid\mathcal L_m]\right).
\]
Their conditional mean budget, variance proxy, and jump are bounded by
\(\mu_{m,h}^{\rm ind}\), \(\sigma_{m,h}^{2,{\rm ind}}\), and
\(\omega_{m,h}\).  Therefore the corresponding individual block load exceeds
\(q_{m,h}^2\) with conditional probability at most the displayed individual
Bernstein exponential.  A union bound over \(i\le n_m\) and
\(\beta\in\mathcal B_{m,h}\) gives the first part of \(\mathfrak B_m\).

The same conditional Bernstein argument, now summing over all pairs \((i,r)\)
for fixed \(h,\beta\), shows that the total block load exceeds
\(n_mq_{m,h}^2\) with conditional probability at most the displayed total-load
exponential.  Union over \(\beta\) and \(h\le J_m\) gives
\(\mathbb P(F_m^{\rm bl})\le\mathfrak B_m\).  The result follows from
Corollary~\ref{cor:finite-lift-core-block-load-atomic-matrix-series-noise}.
\end{proof}

\begin{corollary}[Uniform random block-support matrix-series noise]
\label{cor:finite-lift-core-uniform-random-block-support-matrix-series-noise}
In the setting of
Corollary~\ref{cor:finite-lift-core-block-load-atomic-matrix-series-noise},
suppose that the block projections and amplitudes are measurable with respect
to a sigma-field \(\mathcal L_m\subseteq\mathcal K_m\).  Put
\[
  B_{m,h}^{\#}=|\mathcal B_{m,h}|.
\]
For each active atom let \(s_{m,h,i,r}\in[0,B_{m,h}^{\#}]\) and
\(a_{m,h,i,r}\ge0\) be \(\mathcal L_m\)-measurable quantities such that
\[
  b_{m,h,i,r}\le a_{m,h,i,r}.
\]
Assume that, conditionally on \(\mathcal L_m\), the support sets
\(S_{m,h,i,r}\) are independent over the pairs \((i,r)\), and for every
\(\beta\in\mathcal B_{m,h}\),
\[
  \mathbb P\{\beta\in S_{m,h,i,r}\mid\mathcal L_m\}
  \le \frac{s_{m,h,i,r}}{B_{m,h}^{\#}}.
\]
Let deterministic envelopes
\(\mu_{m,h}^{\rm ind}\), \(\sigma_{m,h}^{2,{\rm ind}}\),
\(\mu_{m,h}^{\rm tot}\), \(\sigma_{m,h}^{2,{\rm tot}}\), and
\(\omega_{m,h}\) satisfy, almost surely,
\[
  \max_{i\le n_m}
  \sum_{r\in\mathcal R_{m,h,i}}
  a_{m,h,i,r}^2\frac{s_{m,h,i,r}}{B_{m,h}^{\#}}
  \le \mu_{m,h}^{\rm ind},
\]
\[
  \max_{i\le n_m}
  \sum_{r\in\mathcal R_{m,h,i}}
  a_{m,h,i,r}^4\frac{s_{m,h,i,r}}{B_{m,h}^{\#}}
  \le \sigma_{m,h}^{2,{\rm ind}},
\]
\[
  \sum_{i=1}^{n_m}\sum_{r\in\mathcal R_{m,h,i}}
  a_{m,h,i,r}^2\frac{s_{m,h,i,r}}{B_{m,h}^{\#}}
  \le \mu_{m,h}^{\rm tot},
\]
\[
  \sum_{i=1}^{n_m}\sum_{r\in\mathcal R_{m,h,i}}
  a_{m,h,i,r}^4\frac{s_{m,h,i,r}}{B_{m,h}^{\#}}
  \le \sigma_{m,h}^{2,{\rm tot}},
  \qquad
  \max_{i,r}a_{m,h,i,r}^2\le\omega_{m,h}.
\]
Let \(u_{m,h}^{\rm ind},u_{m,h}^{\rm tot}\ge0\) be deterministic and suppose
\[
  \mu_{m,h}^{\rm ind}+u_{m,h}^{\rm ind}\le q_{m,h}^2,
  \qquad
  \mu_{m,h}^{\rm tot}+u_{m,h}^{\rm tot}\le n_mq_{m,h}^2.
\]
With the convention that a Bernstein term with zero denominator is omitted,
put
\[
\begin{aligned}
  \mathfrak U_m
  :=
  \sum_{h=1}^{J_m}B_{m,h}^{\#}
  \Bigg[
  &n_m\exp\left\{
  -\frac{(u_{m,h}^{\rm ind})^2}
         {2(\sigma_{m,h}^{2,{\rm ind}}
            +\omega_{m,h}u_{m,h}^{\rm ind}/3)}
  \right\}\\
  &+
  \exp\left\{
  -\frac{(u_{m,h}^{\rm tot})^2}
         {2(\sigma_{m,h}^{2,{\rm tot}}
            +\omega_{m,h}u_{m,h}^{\rm tot}/3)}
  \right\}
  \Bigg].
\end{aligned}
\]
If
\[
  \sum_m\left[
  \mathbb P(F_m^{\rm lift})
  +\mathfrak U_m
  +2J_m(n_m+1)e^{-\tau_m^2}
  \right]<\infty,
\]
then, with probability one, for all sufficiently large \(m\), the
active-prefix and fixed-menu conclusions of
Corollary~\ref{cor:finite-lift-core-operator-tail-menu-borel} hold after
scaling the \(h\)-th uniformly block-sampled matrix-series noise stream by
\[
  C_{\rm ms}q_{m,h}(\sqrt{\log(e\kappa_{m,h})}+\tau_m).
\]
This covers atoms choosing a fixed number of blocks uniformly without
replacement, independent Bernoulli block thinning, random block hyperedges,
and randomly clustered dictionaries, with \(s_{m,h,i,r}\) equal to the
conditional expected support size or any upper bound for it.
\end{corollary}

\begin{proof}
Set
\[
  p_{m,h,i,r,\beta}=\frac{s_{m,h,i,r}}{B_{m,h}^{\#}}.
\]
For each fixed \(h,\beta\), the indicators
\({\bf 1}_{\{\beta\in S_{m,h,i,r}\}}\) are conditionally independent over
\((i,r)\), and their conditional probabilities are bounded by
\(p_{m,h,i,r,\beta}\).  Since \(b_{m,h,i,r}\le a_{m,h,i,r}\), the displayed
envelopes imply the four mean and variance hypotheses and the jump bound in
Corollary~\ref{cor:finite-lift-core-random-block-load-atomic-matrix-series-noise}.
The conclusion follows from that corollary, with
\(\mathfrak B_m\le\mathfrak U_m\).  The examples have exactly the displayed
uniform inclusion bound.
\end{proof}

\begin{corollary}[Uniform block-support profile schedules]
\label{cor:finite-lift-core-uniform-block-support-profile-schedules}
In the setting of
Corollary~\ref{cor:finite-lift-core-uniform-random-block-support-matrix-series-noise},
retain the deterministic envelopes and slacks
\[
  \mu_{m,h}^{\rm ind},\ \sigma_{m,h}^{2,{\rm ind}},\
  \mu_{m,h}^{\rm tot},\ \sigma_{m,h}^{2,{\rm tot}},\
  \omega_{m,h},\ u_{m,h}^{\rm ind},\ u_{m,h}^{\rm tot}.
\]
Define the two Bernstein exponents, with the zero-denominator convention
interpreted as \(+\infty\), by
\[
  \Phi_{m,h}^{\rm ind}
  =
  \begin{cases}
  \dfrac{(u_{m,h}^{\rm ind})^2}
        {2(\sigma_{m,h}^{2,{\rm ind}}
           +\omega_{m,h}u_{m,h}^{\rm ind}/3)},
  &\sigma_{m,h}^{2,{\rm ind}}
     +\omega_{m,h}u_{m,h}^{\rm ind}/3>0,\\[1.1em]
  +\infty,
  &\sigma_{m,h}^{2,{\rm ind}}
     +\omega_{m,h}u_{m,h}^{\rm ind}/3=0,
  \end{cases}
\]
\[
  \Phi_{m,h}^{\rm tot}
  =
  \begin{cases}
  \dfrac{(u_{m,h}^{\rm tot})^2}
        {2(\sigma_{m,h}^{2,{\rm tot}}
           +\omega_{m,h}u_{m,h}^{\rm tot}/3)},
  &\sigma_{m,h}^{2,{\rm tot}}
     +\omega_{m,h}u_{m,h}^{\rm tot}/3>0,\\[1.1em]
  +\infty,
  &\sigma_{m,h}^{2,{\rm tot}}
     +\omega_{m,h}u_{m,h}^{\rm tot}/3=0.
  \end{cases}
\]
Set
\[
  \Phi_m
  =
  \min_{1\le h\le J_m}
  \min\{\Phi_{m,h}^{\rm ind},\Phi_{m,h}^{\rm tot}\},
  \qquad
  L_m^{\rm ubs}
  =
  (n_m+1)\left(J_m+\sum_{h=1}^{J_m}B_{m,h}^{\#}\right).
\]
Assume that the half-budget inequalities in
Corollary~\ref{cor:finite-lift-core-uniform-random-block-support-matrix-series-noise}
hold and that \(\sum_m\mathbb P(F_m^{\rm lift})<\infty\).  If, for some
\(\gamma>0\),
\[
  \Phi_m
  \ge
  \log(eL_m^{\rm ubs})+(1+\gamma)\log(e+m)
\]
and
\[
  \tau_m^2
  \ge
  \log(2L_m^{\rm ubs})+(1+\gamma)\log(e+m)
\]
eventually, then the active-prefix and fixed-menu conclusions of
Corollary~\ref{cor:finite-lift-core-operator-tail-menu-borel} hold almost
surely after scaling the \(h\)-th uniformly block-sampled matrix-series noise
stream by
\[
  C_{\rm ms}q_{m,h}(\sqrt{\log(e\kappa_{m,h})}+\tau_m).
\]

In particular, suppose \(L_m^{\rm ubs}\le(eD_m)^\Gamma\) eventually.  Along all
integer dimensions \(D=2,3,\ldots\), the same conclusion holds if
\[
  \Phi_D\ge R\log(eD),\qquad
  \tau_D^2\ge S\log(eD),
\]
eventually for constants \(R>\Gamma+1\) and \(S>\Gamma+1\).  Along any sparse
schedule \(D_m\ge\exp(c m^\alpha)\), \(c,\alpha>0\), the same conclusion holds
if
\[
  \Phi_m\ge R\log(eD_m),\qquad
  \tau_m^2\ge S\log(eD_m),
\]
eventually for constants \(R>\Gamma\) and \(S>\Gamma\).  Thus heterogeneous
amplitudes and expected support sizes need only produce summable Bernstein
profile exponents; no homogeneity is required.
\end{corollary}

\begin{proof}
By the definition of \(\Phi_m\),
\[
  \mathfrak U_m
  \le
  (n_m+1)\sum_{h=1}^{J_m}B_{m,h}^{\#}e^{-\Phi_m}
  \le L_m^{\rm ubs}e^{-\Phi_m}.
\]
The arbitrary-growth lower bound on \(\Phi_m\) gives
\[
  \mathfrak U_m\le e^{-1}(e+m)^{-(1+\gamma)}.
\]
Also \(J_m(n_m+1)\le L_m^{\rm ubs}\), so the lower bound on \(\tau_m^2\) gives
\[
  2J_m(n_m+1)e^{-\tau_m^2}\le(e+m)^{-(1+\gamma)}.
\]
The summability hypothesis in
Corollary~\ref{cor:finite-lift-core-uniform-random-block-support-matrix-series-noise}
therefore follows.

For all integer dimensions,
\[
  \mathfrak U_D\le(eD)^{\Gamma-R},
  \qquad
  2J_D(n_D+1)e^{-\tau_D^2}\le2(eD)^{\Gamma-S},
\]
which is summable when \(R-\Gamma>1\) and \(S-\Gamma>1\).  For sparse
dimension schedules the same bounds are stretched-exponential summable as soon
as \(R>\Gamma\) and \(S>\Gamma\).
\end{proof}

\begin{corollary}[Homogeneous uniform block-support matrix-series noise]
\label{cor:finite-lift-core-homogeneous-uniform-block-support-matrix-series-noise}
In the setting of
Corollary~\ref{cor:finite-lift-core-uniform-random-block-support-matrix-series-noise},
suppose that for each active \(h\le J_m\) there are deterministic parameters
\[
  N_{m,h}^{\rm ind},\quad N_{m,h}^{\rm tot},\quad
  a_{m,h}>0,\quad s_{m,h}\in[0,B_{m,h}^{\#}]
\]
such that, almost surely,
\[
  |\mathcal R_{m,h,i}|\le N_{m,h}^{\rm ind}\quad(i\le n_m),
  \qquad
  \sum_{i=1}^{n_m}|\mathcal R_{m,h,i}|\le N_{m,h}^{\rm tot},
\]
and every active atom has \(a_{m,h,i,r}\le a_{m,h}\) and
\(s_{m,h,i,r}\le s_{m,h}\).  Put
\[
  \rho_{m,h}^{\rm ind}
  =
  N_{m,h}^{\rm ind}a_{m,h}^2\frac{s_{m,h}}{B_{m,h}^{\#}},
  \qquad
  \rho_{m,h}^{\rm tot}
  =
  N_{m,h}^{\rm tot}a_{m,h}^2\frac{s_{m,h}}{B_{m,h}^{\#}}.
\]
Assume the half-budget bounds
\[
  2\rho_{m,h}^{\rm ind}\le q_{m,h}^2,
  \qquad
  2\rho_{m,h}^{\rm tot}\le n_mq_{m,h}^2.
\]
With the convention that an exponential term with \(\rho=0\) is omitted, set
\[
  \mathfrak H_m
  =
  \sum_{h=1}^{J_m}B_{m,h}^{\#}
  \left[
  n_m\exp\left\{
    -\frac{3\rho_{m,h}^{\rm ind}}{8a_{m,h}^2}
  \right\}
  +
  \exp\left\{
    -\frac{3\rho_{m,h}^{\rm tot}}{8a_{m,h}^2}
  \right\}
  \right].
\]
If
\[
  \sum_m\left[
  \mathbb P(F_m^{\rm lift})
  +\mathfrak H_m
  +2J_m(n_m+1)e^{-\tau_m^2}
  \right]<\infty,
\]
then, with probability one, for all sufficiently large \(m\), the
active-prefix and fixed-menu conclusions of
Corollary~\ref{cor:finite-lift-core-operator-tail-menu-borel} hold after
scaling the \(h\)-th homogeneous uniformly block-sampled matrix-series noise
stream by
\[
  C_{\rm ms}q_{m,h}(\sqrt{\log(e\kappa_{m,h})}+\tau_m).
\]
Equivalently, in this homogeneous model the block-load concentration exponent
is controlled by the expected number of active atom incidences per block:
\[
  \frac{\rho_{m,h}^{\rm ind}}{a_{m,h}^2}
  =
  \frac{N_{m,h}^{\rm ind}s_{m,h}}{B_{m,h}^{\#}},
  \qquad
  \frac{\rho_{m,h}^{\rm tot}}{a_{m,h}^2}
  =
  \frac{N_{m,h}^{\rm tot}s_{m,h}}{B_{m,h}^{\#}}.
\]
\end{corollary}

\begin{proof}
Apply Corollary~\ref{cor:finite-lift-core-uniform-random-block-support-matrix-series-noise}
with
\[
  \mu_{m,h}^{\rm ind}=\rho_{m,h}^{\rm ind},
  \quad
  \sigma_{m,h}^{2,{\rm ind}}=a_{m,h}^2\rho_{m,h}^{\rm ind},
  \quad
  \mu_{m,h}^{\rm tot}=\rho_{m,h}^{\rm tot},
  \quad
  \sigma_{m,h}^{2,{\rm tot}}=a_{m,h}^2\rho_{m,h}^{\rm tot},
  \quad
  \omega_{m,h}=a_{m,h}^2,
\]
and choose
\(u_{m,h}^{\rm ind}=\rho_{m,h}^{\rm ind}\) and
\(u_{m,h}^{\rm tot}=\rho_{m,h}^{\rm tot}\).  The half-budget assumptions give
the two required inequalities.  When \(\rho>0\),
\[
  \frac{u^2}{2(\sigma^2+\omega u/3)}
  =
  \frac{\rho^2}{2(a_{m,h}^2\rho+a_{m,h}^2\rho/3)}
  =
  \frac{3\rho}{8a_{m,h}^2},
\]
which is the displayed \(\mathfrak H_m\).  The zero-load terms are
deterministic and need no exceptional probability.
\end{proof}

\begin{corollary}[Homogeneous uniform block-support schedules]
\label{cor:finite-lift-core-homogeneous-uniform-block-support-schedules}
In the setting of
Corollary~\ref{cor:finite-lift-core-homogeneous-uniform-block-support-matrix-series-noise},
put
\[
  \lambda_{m,h}^{\rm ind}
  =
  \frac{N_{m,h}^{\rm ind}s_{m,h}}{B_{m,h}^{\#}},
  \qquad
  \lambda_{m,h}^{\rm tot}
  =
  \frac{N_{m,h}^{\rm tot}s_{m,h}}{B_{m,h}^{\#}},
\]
\[
  \lambda_m
  =
  \min_{1\le h\le J_m}
  \min\{\lambda_{m,h}^{\rm ind},\lambda_{m,h}^{\rm tot}\},
  \qquad
  L_m^{\rm blk}
  =
  \sum_{h=1}^{J_m}B_{m,h}^{\#}(n_m+1).
\]
Assume the half-budget bounds of that corollary and
\(\sum_m\mathbb P(F_m^{\rm lift})<\infty\).
If, for some \(\gamma>0\),
\[
  \lambda_m
  \ge
  \frac83\left[
    \log(eL_m^{\rm blk})+(1+\gamma)\log(e+m)
  \right]
\]
and
\[
  \tau_m^2
  \ge
  \log\bigl(2J_m(n_m+1)\bigr)+(1+\gamma)\log(e+m)
\]
eventually, then the active-prefix and fixed-menu conclusions of
Corollary~\ref{cor:finite-lift-core-operator-tail-menu-borel} hold almost
surely after scaling the \(h\)-th homogeneous uniformly block-sampled
matrix-series noise stream by
\[
  C_{\rm ms}q_{m,h}(\sqrt{\log(e\kappa_{m,h})}+\tau_m).
\]

In particular, suppose \(L_m^{\rm blk}\le(eD_m)^\Gamma\) eventually.  Along
all integer dimensions \(D=2,3,\ldots\), the same conclusion holds if
\[
  \lambda_D\ge R\log(eD),\qquad
  \tau_D^2\ge S\log(eD),
\]
eventually for constants
\[
  R>\frac{8(\Gamma+1)}3,\qquad S>\Gamma+1.
\]
Along any sparse schedule \(D_m\ge\exp(c m^\alpha)\), \(c,\alpha>0\), the same
conclusion holds if
\[
  \lambda_m\ge R\log(eD_m),\qquad
  \tau_m^2\ge S\log(eD_m),
\]
eventually for constants
\[
  R>\frac{8\Gamma}3,\qquad S>\Gamma .
\]
Thus homogeneous block-sampled matrix-series perturbations are covered whenever
the expected individual and total atom incidences per active block dominate the
active block-count logarithm at these margins.
\end{corollary}

\begin{proof}
For every active \(h\),
\[
  \frac{\rho_{m,h}^{\rm ind}}{a_{m,h}^2}
  =
  \lambda_{m,h}^{\rm ind},
  \qquad
  \frac{\rho_{m,h}^{\rm tot}}{a_{m,h}^2}
  =
  \lambda_{m,h}^{\rm tot}.
\]
Hence the homogeneous Bernstein budget in
Corollary~\ref{cor:finite-lift-core-homogeneous-uniform-block-support-matrix-series-noise}
satisfies
\[
  \mathfrak H_m
  \le
  L_m^{\rm blk}\exp\left\{-\frac{3\lambda_m}{8}\right\}.
\]
The arbitrary-growth hypothesis gives
\[
  L_m^{\rm blk}\exp\left\{-\frac{3\lambda_m}{8}\right\}
  \le e^{-1}(e+m)^{-(1+\gamma)},
\]
and the displayed lower bound on \(\tau_m^2\) gives
\[
  2J_m(n_m+1)e^{-\tau_m^2}\le(e+m)^{-(1+\gamma)}.
\]
Together with the lift-exception summability, the summability hypothesis of
Corollary~\ref{cor:finite-lift-core-homogeneous-uniform-block-support-matrix-series-noise}
follows.

For all integer dimensions,
\[
  \mathfrak H_D\le(eD)^{\Gamma-3R/8},
  \qquad
  2J_D(n_D+1)e^{-\tau_D^2}\le 2(eD)^{\Gamma-S}.
\]
These two series are summable when \(3R/8-\Gamma>1\) and \(S-\Gamma>1\).
For a sparse schedule \(D_m\ge\exp(c m^\alpha)\), the same bounds are
stretched-exponential summable as soon as \(3R/8>\Gamma\) and \(S>\Gamma\).
\end{proof}

\begin{corollary}[Random weighted block-load atomic matrix-series noise]
\label{cor:finite-lift-core-random-weighted-block-load-atomic-matrix-series-noise}
In the setting of
Corollary~\ref{cor:finite-lift-core-block-load-atomic-matrix-series-noise},
let \(\mathcal L_m\subseteq\mathcal K_m\) be a sigma-field with respect to
which the block projections are measurable.  The support sets and amplitudes
may be \(\mathcal K_m\)-measurable random objects.  Put
\[
  W_{m,h,i,r,\beta}
  =
  b_{m,h,i,r}^2{\bf 1}_{\{\beta\in S_{m,h,i,r}\}}.
\]
Assume that, for each fixed active \(h\le J_m\) and
\(\beta\in\mathcal B_{m,h}\), the nonnegative variables
\(W_{m,h,i,r,\beta}\) are conditionally independent over the pairs
\((i,r)\), given \(\mathcal L_m\).  Let deterministic envelopes
\(\mu_{m,h}^{\rm ind}\), \(\sigma_{m,h}^{2,{\rm ind}}\),
\(\mu_{m,h}^{\rm tot}\), \(\sigma_{m,h}^{2,{\rm tot}}\), and
\(\omega_{m,h}\) satisfy, almost surely,
\[
  \max_{i,\beta}
  \sum_{r\in\mathcal R_{m,h,i}}
  \mathbb E[W_{m,h,i,r,\beta}\mid\mathcal L_m]
  \le \mu_{m,h}^{\rm ind},
\]
\[
  \max_{i,\beta}
  \sum_{r\in\mathcal R_{m,h,i}}
  \mathbb E[W_{m,h,i,r,\beta}^2\mid\mathcal L_m]
  \le \sigma_{m,h}^{2,{\rm ind}},
\]
\[
  \max_{\beta}
  \sum_{i=1}^{n_m}\sum_{r\in\mathcal R_{m,h,i}}
  \mathbb E[W_{m,h,i,r,\beta}\mid\mathcal L_m]
  \le \mu_{m,h}^{\rm tot},
\]
\[
  \max_{\beta}
  \sum_{i=1}^{n_m}\sum_{r\in\mathcal R_{m,h,i}}
  \mathbb E[W_{m,h,i,r,\beta}^2\mid\mathcal L_m]
  \le \sigma_{m,h}^{2,{\rm tot}},
  \qquad
  \max_{i,r,\beta}W_{m,h,i,r,\beta}\le\omega_{m,h}.
\]
Let \(u_{m,h}^{\rm ind},u_{m,h}^{\rm tot}\ge0\) be deterministic and suppose
\[
  \mu_{m,h}^{\rm ind}+u_{m,h}^{\rm ind}\le q_{m,h}^2,
  \qquad
  \mu_{m,h}^{\rm tot}+u_{m,h}^{\rm tot}\le n_mq_{m,h}^2.
\]
With the convention that a Bernstein term with zero denominator is omitted,
put
\[
\begin{aligned}
  \mathfrak W_m
  :=
  \sum_{h=1}^{J_m}|\mathcal B_{m,h}|
  \Bigg[
  &n_m\exp\left\{
  -\frac{(u_{m,h}^{\rm ind})^2}
         {2(\sigma_{m,h}^{2,{\rm ind}}
            +\omega_{m,h}u_{m,h}^{\rm ind}/3)}
  \right\}\\
  &+
  \exp\left\{
  -\frac{(u_{m,h}^{\rm tot})^2}
         {2(\sigma_{m,h}^{2,{\rm tot}}
            +\omega_{m,h}u_{m,h}^{\rm tot}/3)}
  \right\}
  \Bigg].
\end{aligned}
\]
If
\[
  \sum_m\left[
  \mathbb P(F_m^{\rm lift})
  +\mathfrak W_m
  +2J_m(n_m+1)e^{-\tau_m^2}
  \right]<\infty,
\]
then, with probability one, for all sufficiently large \(m\), the
active-prefix and fixed-menu conclusions of
Corollary~\ref{cor:finite-lift-core-operator-tail-menu-borel} hold after
scaling the \(h\)-th random weighted block-atomic matrix-series noise stream by
\[
  C_{\rm ms}q_{m,h}(\sqrt{\log(e\kappa_{m,h})}+\tau_m).
\]
This covers independent random amplitudes, random atom intensities, bounded
compound-thinned atom fields, and random weighted dictionaries whose realized
incident block weights have the displayed conditional Bernstein envelopes.
\end{corollary}

\begin{proof}
For fixed \(m,h,i,\beta\), condition on \(\mathcal L_m\) and apply scalar
Bernstein to
\[
  \sum_{r\in\mathcal R_{m,h,i}}
  \left(
  W_{m,h,i,r,\beta}
  -\mathbb E[W_{m,h,i,r,\beta}\mid\mathcal L_m]
  \right).
\]
The conditional variance is bounded by the sum of the conditional second
moments, and each centered summand has jump at most \(\omega_{m,h}\).  Hence
the individual block load exceeds \(q_{m,h}^2\) only with the displayed
individual Bernstein probability.  Union over \(i\le n_m\) and
\(\beta\in\mathcal B_{m,h}\) gives the first part of \(\mathfrak W_m\).

The same argument applied to the sum over all pairs \((i,r)\), for fixed
\(h,\beta\), gives the displayed total-load Bernstein probability, and union
over \(h\le J_m\) and \(\beta\in\mathcal B_{m,h}\) yields
\(\mathbb P(F_m^{\rm bl})\le\mathfrak W_m\).  The conclusion is then exactly
Corollary~\ref{cor:finite-lift-core-block-load-atomic-matrix-series-noise}.
\end{proof}

\begin{corollary}[Predictable weighted block-load atomic matrix-series noise]
\label{cor:finite-lift-core-predictable-weighted-block-load-atomic-matrix-series-noise}
In the setting of
Corollary~\ref{cor:finite-lift-core-block-load-atomic-matrix-series-noise},
fix for each active \(m,h\) a deterministic enumeration
\[
  \mathcal I_{m,h}
  =
  \{(i,r):1\le i\le n_m,\ r\in\mathcal R_{m,h,i}\}
  =
  \{a=1,\ldots,N_{m,h}\},
\]
and write \(i(a),r(a)\) for the corresponding pair.  Let
\((\mathcal L_{m,h,a})_{a=0}^{N_{m,h}}\) be a filtration contained in
\(\mathcal K_m\), with the block projections
\(P_{m,h,\beta}\) measurable with respect to \(\mathcal L_{m,h,0}\).  Put
\[
  W_{m,h,a,\beta}
  =
  b_{m,h,i(a),r(a)}^2
  {\bf 1}_{\{\beta\in S_{m,h,i(a),r(a)}\}}.
\]
Assume \(W_{m,h,a,\beta}\) is \(\mathcal L_{m,h,a}\)-measurable for every
\(a,\beta\).  Define the predictable first and second moment profiles
\[
  M_{m,h,a,\beta}
  =
  \mathbb E[W_{m,h,a,\beta}\mid\mathcal L_{m,h,a-1}],
  \qquad
  V_{m,h,a,\beta}
  =
  \mathbb E[W_{m,h,a,\beta}^2\mid\mathcal L_{m,h,a-1}].
\]
Let deterministic envelopes
\(\mu_{m,h}^{\rm ind}\), \(\sigma_{m,h}^{2,{\rm ind}}\),
\(\mu_{m,h}^{\rm tot}\), \(\sigma_{m,h}^{2,{\rm tot}}\), and
\(\omega_{m,h}\) satisfy, almost surely,
\[
  \max_{i,\beta}
  \sum_{\substack{a\le N_{m,h}\\ i(a)=i}}
  M_{m,h,a,\beta}
  \le \mu_{m,h}^{\rm ind},
\]
\[
  \max_{i,\beta}
  \sum_{\substack{a\le N_{m,h}\\ i(a)=i}}
  V_{m,h,a,\beta}
  \le \sigma_{m,h}^{2,{\rm ind}},
\]
\[
  \max_{\beta}
  \sum_{a=1}^{N_{m,h}}M_{m,h,a,\beta}
  \le \mu_{m,h}^{\rm tot},
\]
\[
  \max_{\beta}
  \sum_{a=1}^{N_{m,h}}V_{m,h,a,\beta}
  \le \sigma_{m,h}^{2,{\rm tot}},
  \qquad
  \max_{a,\beta}W_{m,h,a,\beta}\le\omega_{m,h}.
\]
Let \(u_{m,h}^{\rm ind},u_{m,h}^{\rm tot}\ge0\) be deterministic and suppose
\[
  \mu_{m,h}^{\rm ind}+u_{m,h}^{\rm ind}\le q_{m,h}^2,
  \qquad
  \mu_{m,h}^{\rm tot}+u_{m,h}^{\rm tot}\le n_mq_{m,h}^2.
\]
With the convention that a Freedman term with zero denominator is omitted, put
\[
\begin{aligned}
  \mathfrak P_m
  :=
  \sum_{h=1}^{J_m}|\mathcal B_{m,h}|
  \Bigg[
  &n_m\exp\left\{
  -\frac{(u_{m,h}^{\rm ind})^2}
         {2(\sigma_{m,h}^{2,{\rm ind}}
            +\omega_{m,h}u_{m,h}^{\rm ind}/3)}
  \right\}\\
  &+
  \exp\left\{
  -\frac{(u_{m,h}^{\rm tot})^2}
         {2(\sigma_{m,h}^{2,{\rm tot}}
            +\omega_{m,h}u_{m,h}^{\rm tot}/3)}
  \right\}
  \Bigg].
\end{aligned}
\]
If
\[
  \sum_m\left[
  \mathbb P(F_m^{\rm lift})
  +\mathfrak P_m
  +2J_m(n_m+1)e^{-\tau_m^2}
  \right]<\infty,
\]
then, with probability one, for all sufficiently large \(m\), the
active-prefix and fixed-menu conclusions of
Corollary~\ref{cor:finite-lift-core-operator-tail-menu-borel} hold after
scaling the \(h\)-th predictable weighted block-atomic matrix-series noise
stream by
\[
  C_{\rm ms}q_{m,h}(\sqrt{\log(e\kappa_{m,h})}+\tau_m).
\]
Thus the block-load matrix-series theorem also covers adaptive atom sampling,
online random weighted dictionaries, sequential block choices, and
without-replacement block exposure whenever the predictable incident-weight
budgets have summable scalar-Freedman tails.
\end{corollary}

\begin{proof}
Fix \(m,h,i,\beta\).  Along the enumeration define
\[
  X_a
  =
  {\bf 1}_{\{i(a)=i\}}
  \left(W_{m,h,a,\beta}-M_{m,h,a,\beta}\right).
\]
Then \((X_a,\mathcal L_{m,h,a})\) is a scalar martingale-difference sequence,
\(|X_a|\le\omega_{m,h}\), and its predictable quadratic variation is bounded
by
\[
  \sum_{\substack{a\le N_{m,h}\\ i(a)=i}}V_{m,h,a,\beta}
  \le\sigma_{m,h}^{2,{\rm ind}}.
\]
Since the corresponding predictable mean is at most
\(\mu_{m,h}^{\rm ind}\), the individual block load can exceed \(q_{m,h}^2\)
only if \(\sum_aX_a>u_{m,h}^{\rm ind}\).  The scalar Freedman inequality
therefore gives the displayed individual exponential.  Union over
\(i\le n_m\) and \(\beta\in\mathcal B_{m,h}\) gives the first part of
\(\mathfrak P_m\).

For the total load at fixed \(h,\beta\), use the martingale differences
\(W_{m,h,a,\beta}-M_{m,h,a,\beta}\) over all \(a\).  The predictable mean,
quadratic variation, and jump are bounded by
\(\mu_{m,h}^{\rm tot}\), \(\sigma_{m,h}^{2,{\rm tot}}\), and
\(\omega_{m,h}\), respectively.  Scalar Freedman gives the displayed total
exponential, and union over \(h\) and \(\beta\) yields
\(\mathbb P(F_m^{\rm bl})\le\mathfrak P_m\).  Apply
Corollary~\ref{cor:finite-lift-core-block-load-atomic-matrix-series-noise}.
\end{proof}

\begin{corollary}[Predictable block-support profile matrix-series noise]
\label{cor:finite-lift-core-predictable-block-support-profile-matrix-series-noise}
In the setting of
Corollary~\ref{cor:finite-lift-core-block-load-atomic-matrix-series-noise},
fix for each active \(m,h\) the deterministic enumeration and filtration from
Corollary~\ref{cor:finite-lift-core-predictable-weighted-block-load-atomic-matrix-series-noise}.
Assume that the block projections are
\(\mathcal L_{m,h,0}\)-measurable, that
\(S_{m,h,i(a),r(a)}\) is \(\mathcal L_{m,h,a}\)-measurable, and that there are
predictable profiles
\[
  c_{m,h,a}\ge0,\qquad
  0\le p_{m,h,a,\beta}\le1
  \qquad(a\le N_{m,h},\ \beta\in\mathcal B_{m,h})
\]
which are \(\mathcal L_{m,h,a-1}\)-measurable and satisfy
\[
  b_{m,h,i(a),r(a)}\le c_{m,h,a},
  \qquad
  \mathbb P\{\beta\in S_{m,h,i(a),r(a)}
        \mid \mathcal L_{m,h,a-1}\}
  \le p_{m,h,a,\beta}.
\]
Let deterministic envelopes
\(\mu_{m,h}^{\rm ind}\), \(\sigma_{m,h}^{2,{\rm ind}}\),
\(\mu_{m,h}^{\rm tot}\), \(\sigma_{m,h}^{2,{\rm tot}}\), and
\(\omega_{m,h}\) satisfy, almost surely,
\[
  \max_{i,\beta}
  \sum_{\substack{a\le N_{m,h}\\ i(a)=i}}
  c_{m,h,a}^2p_{m,h,a,\beta}
  \le \mu_{m,h}^{\rm ind},
\]
\[
  \max_{i,\beta}
  \sum_{\substack{a\le N_{m,h}\\ i(a)=i}}
  c_{m,h,a}^4p_{m,h,a,\beta}
  \le \sigma_{m,h}^{2,{\rm ind}},
\]
\[
  \max_{\beta}
  \sum_{a=1}^{N_{m,h}}c_{m,h,a}^2p_{m,h,a,\beta}
  \le \mu_{m,h}^{\rm tot},
\]
\[
  \max_{\beta}
  \sum_{a=1}^{N_{m,h}}c_{m,h,a}^4p_{m,h,a,\beta}
  \le \sigma_{m,h}^{2,{\rm tot}},
  \qquad
  \max_a c_{m,h,a}^2\le\omega_{m,h}.
\]
Let \(u_{m,h}^{\rm ind},u_{m,h}^{\rm tot}\ge0\) be deterministic and suppose
\[
  \mu_{m,h}^{\rm ind}+u_{m,h}^{\rm ind}\le q_{m,h}^2,
  \qquad
  \mu_{m,h}^{\rm tot}+u_{m,h}^{\rm tot}\le n_mq_{m,h}^2.
\]
With the convention that a Freedman term with zero denominator is omitted, put
\[
\begin{aligned}
  \mathfrak A_m
  :=
  \sum_{h=1}^{J_m}|\mathcal B_{m,h}|
  \Bigg[
  &n_m\exp\left\{
  -\frac{(u_{m,h}^{\rm ind})^2}
         {2(\sigma_{m,h}^{2,{\rm ind}}
            +\omega_{m,h}u_{m,h}^{\rm ind}/3)}
  \right\}\\
  &+
  \exp\left\{
  -\frac{(u_{m,h}^{\rm tot})^2}
         {2(\sigma_{m,h}^{2,{\rm tot}}
            +\omega_{m,h}u_{m,h}^{\rm tot}/3)}
  \right\}
  \Bigg].
\end{aligned}
\]
If
\[
  \sum_m\left[
  \mathbb P(F_m^{\rm lift})
  +\mathfrak A_m
  +2J_m(n_m+1)e^{-\tau_m^2}
  \right]<\infty,
\]
then, with probability one, for all sufficiently large \(m\), the
active-prefix and fixed-menu conclusions of
Corollary~\ref{cor:finite-lift-core-operator-tail-menu-borel} hold after
scaling the \(h\)-th predictable block-support matrix-series noise stream by
\[
  C_{\rm ms}q_{m,h}(\sqrt{\log(e\kappa_{m,h})}+\tau_m).
\]
Thus adaptive block samplers, online block dictionaries, urn-type sampling,
and without-replacement block exposure need no independence between different
atoms: it is enough to control the one-step conditional inclusion
probabilities and the predictable amplitude profile.  The uniform predictable
case is obtained by taking
\(p_{m,h,a,\beta}\le s_{m,h,a}/|\mathcal B_{m,h}|\).
\end{corollary}

\begin{proof}
Put
\[
  W_{m,h,a,\beta}
  =
  b_{m,h,i(a),r(a)}^2
  {\bf 1}_{\{\beta\in S_{m,h,i(a),r(a)}\}}.
\]
Then \(W_{m,h,a,\beta}\) is \(\mathcal L_{m,h,a}\)-measurable and
\[
  \mathbb E[W_{m,h,a,\beta}\mid\mathcal L_{m,h,a-1}]
  \le c_{m,h,a}^2p_{m,h,a,\beta},
\]
while
\[
  \mathbb E[W_{m,h,a,\beta}^2\mid\mathcal L_{m,h,a-1}]
  \le c_{m,h,a}^4p_{m,h,a,\beta},
  \qquad
  W_{m,h,a,\beta}\le c_{m,h,a}^2\le\omega_{m,h}.
\]
The displayed deterministic envelopes therefore imply all predictable first
moment, second moment, and jump hypotheses in
Corollary~\ref{cor:finite-lift-core-predictable-weighted-block-load-atomic-matrix-series-noise}.
The half-budget inequalities are the same, and its Freedman budget
\(\mathfrak P_m\) is bounded by \(\mathfrak A_m\).  Applying that corollary
gives the result.
\end{proof}

\begin{corollary}[Predictable block-support profile schedules]
\label{cor:finite-lift-core-predictable-block-support-profile-schedules}
In the setting of
Corollary~\ref{cor:finite-lift-core-predictable-block-support-profile-matrix-series-noise},
retain the deterministic envelopes and slacks
\[
  \mu_{m,h}^{\rm ind},\ \sigma_{m,h}^{2,{\rm ind}},\
  \mu_{m,h}^{\rm tot},\ \sigma_{m,h}^{2,{\rm tot}},\
  \omega_{m,h},\ u_{m,h}^{\rm ind},\ u_{m,h}^{\rm tot}.
\]
Define the two predictable-support Freedman exponents by
\[
  \Theta_{m,h}^{\rm ind}
  =
  \begin{cases}
  \dfrac{(u_{m,h}^{\rm ind})^2}
        {2(\sigma_{m,h}^{2,{\rm ind}}
           +\omega_{m,h}u_{m,h}^{\rm ind}/3)},
  &\sigma_{m,h}^{2,{\rm ind}}
     +\omega_{m,h}u_{m,h}^{\rm ind}/3>0,\\[1.1em]
  +\infty,
  &\sigma_{m,h}^{2,{\rm ind}}
     +\omega_{m,h}u_{m,h}^{\rm ind}/3=0,
  \end{cases}
\]
\[
  \Theta_{m,h}^{\rm tot}
  =
  \begin{cases}
  \dfrac{(u_{m,h}^{\rm tot})^2}
        {2(\sigma_{m,h}^{2,{\rm tot}}
           +\omega_{m,h}u_{m,h}^{\rm tot}/3)},
  &\sigma_{m,h}^{2,{\rm tot}}
     +\omega_{m,h}u_{m,h}^{\rm tot}/3>0,\\[1.1em]
  +\infty,
  &\sigma_{m,h}^{2,{\rm tot}}
     +\omega_{m,h}u_{m,h}^{\rm tot}/3=0.
  \end{cases}
\]
Set
\[
  \Theta_m
  =
  \min_{1\le h\le J_m}
  \min\{\Theta_{m,h}^{\rm ind},\Theta_{m,h}^{\rm tot}\},
  \qquad
  L_m^{\rm pbs}
  =
  (n_m+1)\left(J_m+\sum_{h=1}^{J_m}|\mathcal B_{m,h}|\right).
\]
Assume the half-budget inequalities from
Corollary~\ref{cor:finite-lift-core-predictable-block-support-profile-matrix-series-noise}
and \(\sum_m\mathbb P(F_m^{\rm lift})<\infty\).  If, for some \(\gamma>0\),
\[
  \Theta_m
  \ge
  \log(eL_m^{\rm pbs})+(1+\gamma)\log(e+m)
\]
and
\[
  \tau_m^2
  \ge
  \log(2L_m^{\rm pbs})+(1+\gamma)\log(e+m)
\]
eventually, then the active-prefix and fixed-menu conclusions of
Corollary~\ref{cor:finite-lift-core-operator-tail-menu-borel} hold almost
surely after scaling the \(h\)-th predictable block-support matrix-series
noise stream by
\[
  C_{\rm ms}q_{m,h}(\sqrt{\log(e\kappa_{m,h})}+\tau_m).
\]

In particular, suppose \(L_m^{\rm pbs}\le(eD_m)^\Gamma\) eventually.  Along
all integer dimensions \(D=2,3,\ldots\), the same conclusion holds if
\[
  \Theta_D\ge R\log(eD),\qquad
  \tau_D^2\ge S\log(eD),
\]
eventually for constants \(R>\Gamma+1\) and \(S>\Gamma+1\).  Along any sparse
schedule \(D_m\ge\exp(c m^\alpha)\), \(c,\alpha>0\), the same conclusion holds
if
\[
  \Theta_m\ge R\log(eD_m),\qquad
  \tau_m^2\ge S\log(eD_m),
\]
eventually for constants \(R>\Gamma\) and \(S>\Gamma\).  Thus heterogeneous
adaptive block samplers need only make their predictable support Freedman
exponents dominate the active block-count logarithm.
\end{corollary}

\begin{proof}
By the definition of \(\Theta_m\),
\[
  \mathfrak A_m
  \le
  (n_m+1)\sum_{h=1}^{J_m}|\mathcal B_{m,h}|e^{-\Theta_m}
  \le L_m^{\rm pbs}e^{-\Theta_m}.
\]
The arbitrary-growth lower bound on \(\Theta_m\) gives
\[
  \mathfrak A_m\le e^{-1}(e+m)^{-(1+\gamma)}.
\]
Also \(J_m(n_m+1)\le L_m^{\rm pbs}\), so the lower bound on \(\tau_m^2\) gives
\[
  2J_m(n_m+1)e^{-\tau_m^2}\le(e+m)^{-(1+\gamma)}.
\]
The summability hypothesis in
Corollary~\ref{cor:finite-lift-core-predictable-block-support-profile-matrix-series-noise}
therefore follows.

For all integer dimensions,
\[
  \mathfrak A_D\le(eD)^{\Gamma-R},
  \qquad
  2J_D(n_D+1)e^{-\tau_D^2}\le2(eD)^{\Gamma-S},
\]
which is summable when \(R-\Gamma>1\) and \(S-\Gamma>1\).  For sparse dimension
schedules the same bounds are stretched-exponential summable as soon as
\(R>\Gamma\) and \(S>\Gamma\).
\end{proof}

\begin{corollary}[Inhomogeneous predictable block-support incidence schedules]
\label{cor:finite-lift-core-inhomogeneous-predictable-block-support-incidence-schedules}
In the setting of
Corollary~\ref{cor:finite-lift-core-predictable-block-support-profile-matrix-series-noise},
suppose that for each active \(h\le J_m\) there are deterministic parameters
\[
  a_{m,h}>0,\qquad
  \lambda_{m,h}^{\rm ind}>0,\qquad
  \lambda_{m,h}^{\rm tot}>0
\]
such that, almost surely,
\[
  c_{m,h,a}\le a_{m,h}\qquad(a\le N_{m,h}),
\]
\[
  \max_{i\le n_m}\max_{\beta\in\mathcal B_{m,h}}
  \sum_{\substack{a\le N_{m,h}\\ i(a)=i}}
  p_{m,h,a,\beta}
  \le \lambda_{m,h}^{\rm ind},
\]
and
\[
  \max_{\beta\in\mathcal B_{m,h}}
  \sum_{a=1}^{N_{m,h}}p_{m,h,a,\beta}
  \le \lambda_{m,h}^{\rm tot}.
\]
Assume the half-budget bounds
\[
  2a_{m,h}^2\lambda_{m,h}^{\rm ind}\le q_{m,h}^2,
  \qquad
  2a_{m,h}^2\lambda_{m,h}^{\rm tot}\le n_mq_{m,h}^2.
\]
Put
\[
  \lambda_m
  =
  \min_{1\le h\le J_m}
  \min\{\lambda_{m,h}^{\rm ind},\lambda_{m,h}^{\rm tot}\},
  \qquad
  L_m^{\rm pinh}
  =
  \sum_{h=1}^{J_m}|\mathcal B_{m,h}|(n_m+1).
\]
Assume \(\sum_m\mathbb P(F_m^{\rm lift})<\infty\).  If, for some
\(\gamma>0\),
\[
  \lambda_m
  \ge
  \frac83\left[
    \log(eL_m^{\rm pinh})+(1+\gamma)\log(e+m)
  \right]
\]
and
\[
  \tau_m^2
  \ge
  \log\bigl(2J_m(n_m+1)\bigr)+(1+\gamma)\log(e+m)
\]
eventually, then the active-prefix and fixed-menu conclusions of
Corollary~\ref{cor:finite-lift-core-operator-tail-menu-borel} hold almost
surely after scaling the \(h\)-th inhomogeneous predictable block-support
matrix-series noise stream by
\[
  C_{\rm ms}q_{m,h}(\sqrt{\log(e\kappa_{m,h})}+\tau_m).
\]

In particular, suppose \(L_m^{\rm pinh}\le(eD_m)^\Gamma\) eventually.  Along
all integer dimensions \(D=2,3,\ldots\), the same conclusion holds if
\[
  \lambda_D\ge R\log(eD),\qquad
  \tau_D^2\ge S\log(eD),
\]
eventually for constants
\[
  R>\frac{8(\Gamma+1)}3,\qquad S>\Gamma+1.
\]
Along any sparse schedule \(D_m\ge\exp(c m^\alpha)\), \(c,\alpha>0\), the same
conclusion holds if
\[
  \lambda_m\ge R\log(eD_m),\qquad
  \tau_m^2\ge S\log(eD_m),
\]
eventually for constants
\[
  R>\frac{8\Gamma}3,\qquad S>\Gamma .
\]
Thus biased adaptive categorical block samplers, inhomogeneous urn exposure,
and nonuniform without-replacement block choices are covered once their
predictable row-block and total-block incidence masses dominate the active
block-count logarithm.
\end{corollary}

\begin{proof}
The displayed bounds imply the deterministic envelopes in
Corollary~\ref{cor:finite-lift-core-predictable-block-support-profile-matrix-series-noise}
with
\[
  \mu_{m,h}^{\rm ind}=a_{m,h}^2\lambda_{m,h}^{\rm ind},
  \quad
  \sigma_{m,h}^{2,{\rm ind}}=a_{m,h}^4\lambda_{m,h}^{\rm ind},
  \quad
  \mu_{m,h}^{\rm tot}=a_{m,h}^2\lambda_{m,h}^{\rm tot},
  \quad
  \sigma_{m,h}^{2,{\rm tot}}=a_{m,h}^4\lambda_{m,h}^{\rm tot},
  \quad
  \omega_{m,h}=a_{m,h}^2.
\]
Choose
\[
  u_{m,h}^{\rm ind}=a_{m,h}^2\lambda_{m,h}^{\rm ind},
  \qquad
  u_{m,h}^{\rm tot}=a_{m,h}^2\lambda_{m,h}^{\rm tot}.
\]
The half-budget assumptions give the two budget inequalities.  For either
load parameter \(\lambda\),
\[
  \frac{(a_{m,h}^2\lambda)^2}
       {2(a_{m,h}^4\lambda+a_{m,h}^2a_{m,h}^2\lambda/3)}
  =
  \frac{3\lambda}{8}.
\]
Hence the predictable support budget is bounded by
\[
  \mathfrak A_m
  \le
  L_m^{\rm pinh}\exp\left\{-\frac{3\lambda_m}{8}\right\}.
\]
The arbitrary-growth, all-integer, and sparse conclusions now follow exactly
as in
Corollary~\ref{cor:finite-lift-core-homogeneous-predictable-uniform-block-support-schedules}.
\end{proof}

\begin{corollary}[Homogeneous predictable uniform block-support schedules]
\label{cor:finite-lift-core-homogeneous-predictable-uniform-block-support-schedules}
In the setting of
Corollary~\ref{cor:finite-lift-core-predictable-block-support-profile-matrix-series-noise},
put \(B_{m,h}^{\#}=|\mathcal B_{m,h}|\).  Suppose that for each active
\(h\le J_m\) there are deterministic parameters
\[
  N_{m,h}^{\rm ind},\quad N_{m,h}^{\rm tot},\quad
  a_{m,h}>0,\quad s_{m,h}\in[0,B_{m,h}^{\#}]
\]
such that
\[
  \#\{a\le N_{m,h}:i(a)=i\}\le N_{m,h}^{\rm ind}\quad(i\le n_m),
  \qquad
  N_{m,h}\le N_{m,h}^{\rm tot},
\]
and the predictable profiles satisfy, almost surely,
\[
  c_{m,h,a}\le a_{m,h},
  \qquad
  p_{m,h,a,\beta}\le \frac{s_{m,h}}{B_{m,h}^{\#}}
  \qquad(a\le N_{m,h},\ \beta\in\mathcal B_{m,h}).
\]
Define
\[
  \lambda_{m,h}^{\rm ind}
  =
  \frac{N_{m,h}^{\rm ind}s_{m,h}}{B_{m,h}^{\#}},
  \qquad
  \lambda_{m,h}^{\rm tot}
  =
  \frac{N_{m,h}^{\rm tot}s_{m,h}}{B_{m,h}^{\#}},
\]
and assume the half-budget bounds
\[
  2a_{m,h}^2\lambda_{m,h}^{\rm ind}\le q_{m,h}^2,
  \qquad
  2a_{m,h}^2\lambda_{m,h}^{\rm tot}\le n_mq_{m,h}^2.
\]
Set
\[
  \lambda_m
  =
  \min_{1\le h\le J_m}
  \min\{\lambda_{m,h}^{\rm ind},\lambda_{m,h}^{\rm tot}\},
  \qquad
  L_m^{\rm pblk}
  =
  \sum_{h=1}^{J_m}B_{m,h}^{\#}(n_m+1).
\]
Assume \(\sum_m\mathbb P(F_m^{\rm lift})<\infty\).  If, for some
\(\gamma>0\),
\[
  \lambda_m
  \ge
  \frac83\left[
    \log(eL_m^{\rm pblk})+(1+\gamma)\log(e+m)
  \right]
\]
and
\[
  \tau_m^2
  \ge
  \log\bigl(2J_m(n_m+1)\bigr)+(1+\gamma)\log(e+m)
\]
eventually, then the active-prefix and fixed-menu conclusions of
Corollary~\ref{cor:finite-lift-core-operator-tail-menu-borel} hold almost
surely after scaling the \(h\)-th homogeneous predictable uniformly
block-sampled matrix-series noise stream by
\[
  C_{\rm ms}q_{m,h}(\sqrt{\log(e\kappa_{m,h})}+\tau_m).
\]

In particular, suppose \(L_m^{\rm pblk}\le(eD_m)^\Gamma\) eventually.  Along
all integer dimensions \(D=2,3,\ldots\), the same conclusion holds if
\[
  \lambda_D\ge R\log(eD),\qquad
  \tau_D^2\ge S\log(eD),
\]
eventually for constants
\[
  R>\frac{8(\Gamma+1)}3,\qquad S>\Gamma+1.
\]
Along any sparse schedule \(D_m\ge\exp(c m^\alpha)\), \(c,\alpha>0\), the same
conclusion holds if
\[
  \lambda_m\ge R\log(eD_m),\qquad
  \tau_m^2\ge S\log(eD_m),
\]
eventually for constants
\[
  R>\frac{8\Gamma}3,\qquad S>\Gamma .
\]
Thus adaptive uniform block samplers, online urn models, and
without-replacement block exposures have the same incidence thresholds as the
independent homogeneous block-support model; independence is replaced by
predictability of the one-step inclusion profile.
\end{corollary}

\begin{proof}
For every active \(h\), set
\[
  \rho_{m,h}^{\rm ind}=a_{m,h}^2\lambda_{m,h}^{\rm ind},
  \qquad
  \rho_{m,h}^{\rm tot}=a_{m,h}^2\lambda_{m,h}^{\rm tot}.
\]
The homogeneous bounds imply the deterministic envelopes in
Corollary~\ref{cor:finite-lift-core-predictable-block-support-profile-matrix-series-noise}
with
\[
  \mu_{m,h}^{\rm ind}=\rho_{m,h}^{\rm ind},
  \quad
  \sigma_{m,h}^{2,{\rm ind}}=a_{m,h}^2\rho_{m,h}^{\rm ind},
  \quad
  \mu_{m,h}^{\rm tot}=\rho_{m,h}^{\rm tot},
  \quad
  \sigma_{m,h}^{2,{\rm tot}}=a_{m,h}^2\rho_{m,h}^{\rm tot},
  \quad
  \omega_{m,h}=a_{m,h}^2.
\]
Choose \(u_{m,h}^{\rm ind}=\rho_{m,h}^{\rm ind}\) and
\(u_{m,h}^{\rm tot}=\rho_{m,h}^{\rm tot}\).  The half-budget assumptions are
exactly the two required budget inequalities, and the Freedman exponent for a
nonzero load \(\rho\) is
\[
  \frac{\rho^2}{2(a_{m,h}^2\rho+a_{m,h}^2\rho/3)}
  =
  \frac{3\rho}{8a_{m,h}^2}.
\]
Zero-load terms are deterministic and omitted.  Hence the predictable
block-support budget satisfies
\[
  \mathfrak A_m
  \le
  L_m^{\rm pblk}\exp\left\{-\frac{3\lambda_m}{8}\right\}.
\]
The arbitrary-growth lower bound on \(\lambda_m\) makes this budget at most
\(e^{-1}(e+m)^{-(1+\gamma)}\), while the lower bound on \(\tau_m^2\) gives
\[
  2J_m(n_m+1)e^{-\tau_m^2}\le(e+m)^{-(1+\gamma)}.
\]
Thus the summability hypothesis of
Corollary~\ref{cor:finite-lift-core-predictable-block-support-profile-matrix-series-noise}
holds.

For all integer dimensions,
\[
  \mathfrak A_D\le(eD)^{\Gamma-3R/8},
  \qquad
  2J_D(n_D+1)e^{-\tau_D^2}\le2(eD)^{\Gamma-S},
\]
which is summable when \(3R/8-\Gamma>1\) and \(S-\Gamma>1\).  Along a sparse
dimension schedule the same estimates are stretched-exponential summable once
\(3R/8>\Gamma\) and \(S>\Gamma\).
\end{proof}

\begin{corollary}[Weighted block-load profile schedules]
\label{cor:finite-lift-core-weighted-block-load-profile-schedules}
Work in the triangular-array setting of either
Corollary~\ref{cor:finite-lift-core-random-weighted-block-load-atomic-matrix-series-noise}
or
Corollary~\ref{cor:finite-lift-core-predictable-weighted-block-load-atomic-matrix-series-noise}.
Retain the deterministic envelopes and slacks
\[
  \mu_{m,h}^{\rm ind},\ \sigma_{m,h}^{2,{\rm ind}},\
  \mu_{m,h}^{\rm tot},\ \sigma_{m,h}^{2,{\rm tot}},\
  \omega_{m,h},\ u_{m,h}^{\rm ind},\ u_{m,h}^{\rm tot}.
\]
Define
\[
  \Psi_{m,h}^{\rm ind}
  =
  \begin{cases}
  \dfrac{(u_{m,h}^{\rm ind})^2}
        {2(\sigma_{m,h}^{2,{\rm ind}}
           +\omega_{m,h}u_{m,h}^{\rm ind}/3)},
  &\sigma_{m,h}^{2,{\rm ind}}
     +\omega_{m,h}u_{m,h}^{\rm ind}/3>0,\\[1.1em]
  +\infty,
  &\sigma_{m,h}^{2,{\rm ind}}
     +\omega_{m,h}u_{m,h}^{\rm ind}/3=0,
  \end{cases}
\]
\[
  \Psi_{m,h}^{\rm tot}
  =
  \begin{cases}
  \dfrac{(u_{m,h}^{\rm tot})^2}
        {2(\sigma_{m,h}^{2,{\rm tot}}
           +\omega_{m,h}u_{m,h}^{\rm tot}/3)},
  &\sigma_{m,h}^{2,{\rm tot}}
     +\omega_{m,h}u_{m,h}^{\rm tot}/3>0,\\[1.1em]
  +\infty,
  &\sigma_{m,h}^{2,{\rm tot}}
     +\omega_{m,h}u_{m,h}^{\rm tot}/3=0.
  \end{cases}
\]
Put
\[
  \Psi_m
  =
  \min_{1\le h\le J_m}
  \min\{\Psi_{m,h}^{\rm ind},\Psi_{m,h}^{\rm tot}\},
  \qquad
  L_m^{\rm wbl}
  =
  (n_m+1)\left(J_m+\sum_{h=1}^{J_m}|\mathcal B_{m,h}|\right).
\]
Assume the relevant half-budget inequalities and
\(\sum_m\mathbb P(F_m^{\rm lift})<\infty\).  If, for some \(\gamma>0\),
\[
  \Psi_m
  \ge
  \log(eL_m^{\rm wbl})+(1+\gamma)\log(e+m)
\]
and
\[
  \tau_m^2
  \ge
  \log(2L_m^{\rm wbl})+(1+\gamma)\log(e+m)
\]
eventually, then the active-prefix and fixed-menu conclusions of
Corollary~\ref{cor:finite-lift-core-operator-tail-menu-borel} hold almost
surely for the corresponding random weighted or predictable weighted
block-atomic matrix-series perturbations, after scaling the \(h\)-th noise
stream by
\[
  C_{\rm ms}q_{m,h}(\sqrt{\log(e\kappa_{m,h})}+\tau_m).
\]

In particular, suppose \(L_m^{\rm wbl}\le(eD_m)^\Gamma\) eventually.  Along all
integer dimensions \(D=2,3,\ldots\), the same conclusion holds if
\[
  \Psi_D\ge R\log(eD),\qquad
  \tau_D^2\ge S\log(eD),
\]
eventually for constants \(R>\Gamma+1\) and \(S>\Gamma+1\).  Along any sparse
schedule \(D_m\ge\exp(c m^\alpha)\), \(c,\alpha>0\), the same conclusion holds
if
\[
  \Psi_m\ge R\log(eD_m),\qquad
  \tau_m^2\ge S\log(eD_m),
\]
eventually for constants \(R>\Gamma\) and \(S>\Gamma\).  Thus random
amplitudes, random atom intensities, adaptive atom sampling, and
without-replacement block exposure all obey the same active-prefix schedule once
their scalar Bernstein or Freedman exponents dominate the active block count.
\end{corollary}

\begin{proof}
Let \(\mathfrak Q_m\) denote the budget term \(\mathfrak W_m\) in the independent
weighted case and \(\mathfrak P_m\) in the predictable case.  By the definition
of \(\Psi_m\),
\[
  \mathfrak Q_m
  \le
  (n_m+1)\sum_{h=1}^{J_m}|\mathcal B_{m,h}|e^{-\Psi_m}
  \le L_m^{\rm wbl}e^{-\Psi_m}.
\]
The arbitrary-growth lower bound on \(\Psi_m\) gives
\[
  \mathfrak Q_m\le e^{-1}(e+m)^{-(1+\gamma)}.
\]
Also \(J_m(n_m+1)\le L_m^{\rm wbl}\), so the lower bound on \(\tau_m^2\) gives
\[
  2J_m(n_m+1)e^{-\tau_m^2}\le(e+m)^{-(1+\gamma)}.
\]
The summability hypothesis in the relevant weighted block-load corollary
follows.

For all integer dimensions,
\[
  \mathfrak Q_D\le(eD)^{\Gamma-R},
  \qquad
  2J_D(n_D+1)e^{-\tau_D^2}\le2(eD)^{\Gamma-S},
\]
which is summable when \(R-\Gamma>1\) and \(S-\Gamma>1\).  For sparse dimension
schedules, the same bounds are stretched-exponential summable as soon as
\(R>\Gamma\) and \(S>\Gamma\).
\end{proof}

\begin{corollary}[Bounded weighted block-load incidence schedules]
\label{cor:finite-lift-core-bounded-weighted-block-load-incidence-schedules}
Work in the triangular-array setting of
Corollary~\ref{cor:finite-lift-core-weighted-block-load-profile-schedules}.
Suppose that, for each active \(h\le J_m\), there are deterministic parameters
\[
  \omega_{m,h}^{\rm bd}>0,\qquad
  \rho_{m,h}^{\rm ind}>0,\qquad
  \rho_{m,h}^{\rm tot}>0
\]
such that the deterministic envelopes may be chosen to satisfy
\[
  \mu_{m,h}^{\rm ind}\le\rho_{m,h}^{\rm ind},
  \qquad
  \sigma_{m,h}^{2,{\rm ind}}
  \le \omega_{m,h}^{\rm bd}\rho_{m,h}^{\rm ind},
  \qquad
  \omega_{m,h}\le\omega_{m,h}^{\rm bd},
\]
and
\[
  \mu_{m,h}^{\rm tot}\le\rho_{m,h}^{\rm tot},
  \qquad
  \sigma_{m,h}^{2,{\rm tot}}
  \le \omega_{m,h}^{\rm bd}\rho_{m,h}^{\rm tot}.
\]
Assume the half-budget bounds
\[
  2\rho_{m,h}^{\rm ind}\le q_{m,h}^2,
  \qquad
  2\rho_{m,h}^{\rm tot}\le n_mq_{m,h}^2.
\]
Put
\[
  \Lambda_m^{\rm bw}
  =
  \min_{1\le h\le J_m}
  \min\left\{
    \frac{\rho_{m,h}^{\rm ind}}{\omega_{m,h}^{\rm bd}},
    \frac{\rho_{m,h}^{\rm tot}}{\omega_{m,h}^{\rm bd}}
  \right\},
  \qquad
  L_m^{\rm bw}
  =
  (n_m+1)\left(J_m+\sum_{h=1}^{J_m}|\mathcal B_{m,h}|\right).
\]
Assume \(\sum_m\mathbb P(F_m^{\rm lift})<\infty\).  If, for some
\(\gamma>0\),
\[
  \Lambda_m^{\rm bw}
  \ge
  \frac83\left[
    \log(eL_m^{\rm bw})+(1+\gamma)\log(e+m)
  \right]
\]
and
\[
  \tau_m^2
  \ge
  \log(2L_m^{\rm bw})+(1+\gamma)\log(e+m)
\]
eventually, then the active-prefix and fixed-menu conclusions of
Corollary~\ref{cor:finite-lift-core-operator-tail-menu-borel} hold almost
surely for the corresponding random weighted or predictable weighted
block-atomic matrix-series perturbations, after scaling the \(h\)-th noise
stream by
\[
  C_{\rm ms}q_{m,h}(\sqrt{\log(e\kappa_{m,h})}+\tau_m).
\]

In particular, suppose \(L_m^{\rm bw}\le(eD_m)^\Gamma\) eventually.  Along all
integer dimensions \(D=2,3,\ldots\), the same conclusion holds if
\[
  \Lambda_D^{\rm bw}\ge R\log(eD),\qquad
  \tau_D^2\ge S\log(eD),
\]
eventually for constants
\[
  R>\frac{8(\Gamma+1)}3,\qquad S>\Gamma+1.
\]
Along any sparse schedule \(D_m\ge\exp(c m^\alpha)\), \(c,\alpha>0\), the same
conclusion holds if
\[
  \Lambda_m^{\rm bw}\ge R\log(eD_m),\qquad
  \tau_m^2\ge S\log(eD_m),
\]
eventually for constants
\[
  R>\frac{8\Gamma}3,\qquad S>\Gamma .
\]
Thus bounded random intensities, bounded compound-thinned atom fields,
adaptive weighted dictionaries, and without-replacement weighted exposure are
covered once their normalized row-block and total-block mean loads dominate the
active block-count logarithm.
\end{corollary}

\begin{proof}
Use the weighted block-load profile schedule with the deterministic slacks
\[
  u_{m,h}^{\rm ind}=\rho_{m,h}^{\rm ind},
  \qquad
  u_{m,h}^{\rm tot}=\rho_{m,h}^{\rm tot}.
\]
The displayed half-budget assumptions give
\(\mu_{m,h}^{\rm ind}+u_{m,h}^{\rm ind}\le q_{m,h}^2\) and
\(\mu_{m,h}^{\rm tot}+u_{m,h}^{\rm tot}\le n_mq_{m,h}^2\).  Moreover
\[
  \Psi_{m,h}^{\rm ind}
  \ge
  \frac{(\rho_{m,h}^{\rm ind})^2}
       {2(\omega_{m,h}^{\rm bd}\rho_{m,h}^{\rm ind}
          +\omega_{m,h}^{\rm bd}\rho_{m,h}^{\rm ind}/3)}
  =
  \frac{3\rho_{m,h}^{\rm ind}}{8\omega_{m,h}^{\rm bd}},
\]
and the same estimate gives
\[
  \Psi_{m,h}^{\rm tot}
  \ge
  \frac{3\rho_{m,h}^{\rm tot}}{8\omega_{m,h}^{\rm bd}}.
\]
Hence \(\Psi_m\ge3\Lambda_m^{\rm bw}/8\).  The arbitrary-growth, all-integer,
and sparse conclusions are exactly those of
Corollary~\ref{cor:finite-lift-core-weighted-block-load-profile-schedules}
with this lower bound for \(\Psi_m\).
\end{proof}

\begin{lemma}[Quadratic-form tails imply operator tails]
\label{lem:quadratic-form-to-operator-tail}
Let \(Y_1,\ldots,Y_n\) be conditionally centered self-adjoint random matrices
in \(\M_D(\C)\), relative to a sigma-field \(\mathcal H\).  Suppose that
there is \(q>0\) such that, conditionally on \(\mathcal H\), for every
\(a\in\R^n\), every deterministic unit vector \(u\in\C^D\), and every
\(s\ge0\),
\[
  \mathbb P\left\{
  \left|\left\langle u,\sum_{i=1}^n a_iY_i u\right\rangle\right|
  >q\|a\|_2s
  \,\middle|\,\mathcal H
  \right\}\le 2e^{-s^2}.
\]
Then, conditionally on \(\mathcal H\), for every \(a\in\R^n\) and every
\(t\ge0\),
\[
  \mathbb P\left\{
  \left\|\sum_{i=1}^n a_iY_i\right\|_{\op}
  >5q\|a\|_2(\sqrt D+t)
  \,\middle|\,\mathcal H
  \right\}\le 2e^{-t^2}.
\]
\end{lemma}

\begin{proof}
Fix \(a\) and \(t\), and put \(S=\sum_i a_iY_i\).  Let \(\mathcal N\) be a
deterministic \(1/4\)-net of the unit sphere of \(\C^D\), viewed as a real
\(2D\)-dimensional sphere, with \(|\mathcal N|\le 9^{2D}\).  For self-adjoint
\(S\),
\[
  \|S\|_{\op}\le 2\max_{v\in\mathcal N}|\langle v,Sv\rangle|.
\]
Therefore the conditional failure probability is at most
\[
  2|\mathcal N|
  \exp\left(-\frac{25}{4}(\sqrt D+t)^2\right)
  \le
  2\exp\left(2D\log 9-\frac{25}{4}D-\frac{25}{4}t^2\right)
  \le 2e^{-t^2},
\]
since \(25/4>2\log 9\).
\end{proof}

\begin{corollary}[Quadratic-form noise for finite-lift core menus]
\label{cor:finite-lift-core-quadratic-form-menu}
In the setting of
Corollary~\ref{cor:finite-lift-core-operator-tail-menu-borel}, keep the same
finite-lift core assumptions and the same active-prefix exhaustion.  Replace
the operator-tail hypothesis by the following quadratic-form hypothesis: for
each \(h\le J_m\) there is \(q_{m,h}>0\) such that, conditionally on
\(\mathcal H_m\), for every \(a\in\R^{n_m}\), every deterministic unit vector
\(u\in\C^{D_m}\), and every \(s\ge0\),
\[
  \mathbb P\left\{
  \left|\left\langle u,
  \sum_{i=1}^{n_m}a_iY_{m,i}^{(h)}u
  \right\rangle\right|
  >q_{m,h}\|a\|_2s
  \,\middle|\,\mathcal H_m
  \right\}\le 2e^{-s^2}.
\]
For deterministic \(\tau_m>0\), define
\[
  Z_{m,i}^{(h)}
  =\frac{Y_{m,i}^{(h)}}{5q_{m,h}(\sqrt{D_m}+\tau_m)},
  \qquad
  X_{m,i}^{(h)}=B_{m,i}^{(h)}+Z_{m,i}^{(h)} .
\]
If
\[
  \sum_m\left[
  \mathbb P(F_m^{\rm lift})+2J_m(n_m+1)e^{-\tau_m^2}
  \right]<\infty,
\]
then the active-prefix and fixed-menu conclusions of
Corollary~\ref{cor:finite-lift-core-operator-tail-menu-borel} hold almost
surely for these \(X_{m,i}^{(h)}\).  In particular, the explicit schedules of
Corollary~\ref{cor:finite-lift-core-operator-tail-menu-schedules} apply with
noise scale \(5q_{m,h}(\sqrt{D_m}+\tau_m)\).

Consequently the finite-lift core menu theorem applies directly to
conditionally Gaussian matrix series whose scalar quadratic forms have
conditional variance at most \(q_{m,h}^2\|a\|_2^2\), after enlarging
\(q_{m,h}\) by an absolute constant.  The same conclusion applies to
conditionally sub-Gaussian Wigner-type or correlated entry models whenever
their conditional quadratic-form sub-Gaussian proxy is uniformly bounded by
\(q_{m,h}\|a\|_2\).
\end{corollary}

\begin{proof}
Lemma~\ref{lem:quadratic-form-to-operator-tail}, applied conditionally on
\(\mathcal H_m\) for each active-prefix \(h\), gives the operator-tail
hypothesis of Corollary~\ref{cor:finite-lift-core-operator-tail-menu-borel}
with \(K_{0,m,h}=5q_{m,h}\).  The displayed conclusion is therefore exactly
that corollary, and the schedule statement is
Corollary~\ref{cor:finite-lift-core-operator-tail-menu-schedules}.

For the Gaussian specialization, the displayed variance bound gives the
quadratic-form tail by the one-dimensional Gaussian estimate, with only an
absolute change of constants.  The sub-Gaussian version is the same
one-dimensional linear-form estimate applied to each fixed quadratic form.
\end{proof}

\begin{corollary}[Entrywise sub-Gaussian noise menus]
\label{cor:finite-lift-core-entrywise-subgaussian-menu}
In the setting of
Corollary~\ref{cor:finite-lift-core-operator-tail-menu-borel}, keep the same
finite-lift core assumptions and the same active-prefix exhaustion.  Suppose
that, conditionally on \(\mathcal H_m\), for each \(h\le J_m\) the
upper-triangular entries of
\(\{Y_{m,i}^{(h)}:1\le i\le n_m\}\) are independent and centered, with
\[
  Y_{m,i}^{(h)}
  =\sum_{p=1}^{D_m}\xi_{m,i,p}^{(h)}E_{pp}
  +\sum_{1\le p<q\le D_m}
  \left(\zeta_{m,i,pq}^{(h)}E_{pq}
  +\overline{\zeta_{m,i,pq}^{(h)}}E_{qp}\right).
\]
Assume that, for every \(h\le J_m\), there is a positive deterministic
envelope \(Q_{m,h}\) such that for every \(i\le n_m\) and every unit
\(u\in\C^{D_m}\),
\[
  \left[
  \sum_p\alpha_{m,i,p}^{(h)2}|u_p|^4
  +4\sum_{p<q}\beta_{m,i,pq}^{(h)2}|u_p|^2|u_q|^2
  \right]^{1/2}
  \le Q_{m,h},
\]
where conditionally on \(\mathcal H_m\),
\(\xi_{m,i,p}^{(h)}\) is real sub-Gaussian with proxy
\(\alpha_{m,i,p}^{(h)}\), and every real projection
\(\operatorname{Re}(e^{\sqrt{-1}\theta}\zeta_{m,i,pq}^{(h)})\) is
sub-Gaussian with proxy \(\beta_{m,i,pq}^{(h)}\), uniformly in
\(\theta\in\R\).

Then there is a universal constant \(C_{\rm ent}\) such that, for every
deterministic \(\tau_m>0\), the matrices
\[
  Z_{m,i}^{(h)}
  =\frac{Y_{m,i}^{(h)}}{C_{\rm ent}Q_{m,h}(\sqrt{D_m}+\tau_m)},
  \qquad
  X_{m,i}^{(h)}=B_{m,i}^{(h)}+Z_{m,i}^{(h)}
\]
have the active-prefix and fixed-menu conclusions of
Corollary~\ref{cor:finite-lift-core-operator-tail-menu-borel} whenever
\[
  \sum_m\left[
  \mathbb P(F_m^{\rm lift})+2J_m(n_m+1)e^{-\tau_m^2}
  \right]<\infty .
\]
The explicit schedules of
Corollary~\ref{cor:finite-lift-core-operator-tail-menu-schedules} therefore
apply with entrywise noise scale
\[
  C_{\rm ent}Q_{m,h}(\sqrt{D_m}+\tau_m).
\]
In particular, if all displayed entry proxies are bounded by
\(\sigma_{m,h}\), then one may take
\(Q_{m,h}=\sqrt2\,\sigma_{m,h}\).  Thus normalized Wigner, GOE, GUE, and
variance-profile Gaussian or sub-Gaussian active-prefix noise libraries with
entry scale \(\sigma_{m,h}\) are covered by the same finite-lift core signs at
the natural spectral scale \(\sigma_{m,h}(\sqrt{D_m}+\tau_m)\).
\end{corollary}

\begin{proof}
Fix \(m,h,i\), a unit vector \(u\), and condition on \(\mathcal H_m\).  The
quadratic form is
\[
  \left\langle u,Y_{m,i}^{(h)}u\right\rangle
  =\sum_p\xi_{m,i,p}^{(h)}|u_p|^2
  +2\sum_{p<q}\operatorname{Re}
  \left(\zeta_{m,i,pq}^{(h)}\overline{u_p}u_q\right).
\]
By independence of the upper-triangular entries and the standard linear
combination rule for sub-Gaussian variables, this scalar is conditionally
sub-Gaussian with proxy at most a universal constant times the displayed
square-function, hence at most \(C Q_{m,h}\).  Independence across the
matrices \(i\), conditionally on \(\mathcal H_m\), gives for every
\(a\in\R^{n_m}\) the proxy \(C Q_{m,h}\|a\|_2\) for
\(\langle u,\sum_i a_iY_{m,i}^{(h)}u\rangle\).  After increasing the universal
constant so that the tail is in the \(2e^{-s^2}\) normalization, the
hypothesis of Corollary~\ref{cor:finite-lift-core-quadratic-form-menu} holds.
That corollary gives the asserted finite-lift menu conclusion and schedules.

If all entry proxies are at most \(\sigma_{m,h}\), then
\[
  \sum_p |u_p|^4+4\sum_{p<q}|u_p|^2|u_q|^2
  =2-\sum_p|u_p|^4\le2,
\]
so \(Q_{m,h}=\sqrt2\,\sigma_{m,h}\) is admissible.
\end{proof}

\begin{corollary}[Row-variance entry envelopes]
\label{cor:finite-lift-core-entrywise-row-variance-menu}
In the setting and notation of
Corollary~\ref{cor:finite-lift-core-entrywise-subgaussian-menu}, suppose that
for every active-prefix \(h\le J_m\) there is \(R_{m,h}>0\) such that, for all
\(i\le n_m\) and all \(1\le p\le D_m\),
\[
  \alpha_{m,i,p}^{(h)2}
  +2\sum_{q<p}\beta_{m,i,qp}^{(h)2}
  +2\sum_{p<q}\beta_{m,i,pq}^{(h)2}
  \le R_{m,h}^2 .
\]
Then \(Q_{m,h}=R_{m,h}\) is admissible in
Corollary~\ref{cor:finite-lift-core-entrywise-subgaussian-menu}.  Hence the
active-prefix finite-lift core menu theorem, and all schedules of
Corollary~\ref{cor:finite-lift-core-operator-tail-menu-schedules}, hold after
scaling the \(h\)-th entrywise noise stream by
\[
  C_{\rm ent}R_{m,h}(\sqrt{D_m}+\tau_m).
\]

In particular, if the off-diagonal proxies in the \(h\)-th stream are
supported on graphs of maximum degree \(\Delta_{m,h}\), while
\(\alpha_{m,i,p}^{(h)}\le\alpha_{m,h}\) and
\(\beta_{m,i,pq}^{(h)}\le\beta_{m,h}\), then one may take
\[
  R_{m,h}=\left(\alpha_{m,h}^2+2\Delta_{m,h}\beta_{m,h}^2\right)^{1/2}.
\]
Thus sparse or banded Wigner/GOE/GUE-type active-prefix noise with degree
\(\Delta_{m,h}\) and off-diagonal proxy \(\beta_{m,h}\) is covered at scale
\[
  \left(\alpha_{m,h}^2+2\Delta_{m,h}\beta_{m,h}^2\right)^{1/2}
  (\sqrt{D_m}+\tau_m),
\]
up to the universal constant \(C_{\rm ent}\).  In the normalized sparse case
\(\beta_{m,h}\lesssim\Delta_{m,h}^{-1/2}\) and
\(\alpha_{m,h}=O(1)\), this is the same \(O(\sqrt{D_m}+\tau_m)\) operator-tail
scale as in the dense normalized models.
\end{corollary}

\begin{proof}
Fix \(m,h,i\) and put \(x_p=|u_p|^2\), so \(x_p\ge0\) and \(\sum_p x_p=1\).
The square-function in
Corollary~\ref{cor:finite-lift-core-entrywise-subgaussian-menu} is
\[
  \sum_p\alpha_{m,i,p}^{(h)2}x_p^2
  +4\sum_{p<q}\beta_{m,i,pq}^{(h)2}x_px_q .
\]
This is \(x^TVx\), where \(V_{pp}=\alpha_{m,i,p}^{(h)2}\) and
\(V_{pq}=2\beta_{m,i,pq}^{(h)2}\) for \(p<q\), symmetrically extended.  Since
all entries of \(V\) are nonnegative and every row sum is at most \(R_{m,h}^2\),
\[
  x^TVx
  \le \sum_p x_p R_{m,h}^2
  =R_{m,h}^2 .
\]
Therefore \(Q_{m,h}=R_{m,h}\) satisfies the entrywise square-function
hypothesis, and
Corollary~\ref{cor:finite-lift-core-entrywise-subgaussian-menu} gives the
conclusion.

The graph-degree specialization is the displayed row-sum bound with at most
\(\Delta_{m,h}\) nonzero off-diagonal proxies touching each row.
\end{proof}

\begin{corollary}[Sparse and banded entry-profile schedules]
\label{cor:finite-lift-core-sparse-banded-entry-schedules}
In the setting of
Corollary~\ref{cor:finite-lift-core-entrywise-row-variance-menu}, suppose that
for each active-prefix \(h\le J_m\) the off-diagonal support of every
\(Y_{m,i}^{(h)}\) is contained in a deterministic graph
\(G_{m,h}\) on \(\{1,\ldots,D_m\}\) with maximum degree
\(\Delta_{m,h}\).  Put \(d_{m,h}=\max(1,\Delta_{m,h})\).  Suppose also that
for deterministic \(a_{m,h},b_{m,h}\ge0\), with
\(a_{m,h}^2+b_{m,h}^2>0\),
\[
  \alpha_{m,i,p}^{(h)}\le a_{m,h},\qquad
  \beta_{m,i,pq}^{(h)}
  \le \frac{b_{m,h}}{\sqrt{d_{m,h}}}
\]
for every active \(i,p,q,h\).  Then all active-prefix and fixed-menu
conclusions of
Corollary~\ref{cor:finite-lift-core-operator-tail-menu-borel} hold, with the
same Borel--Cantelli schedules as
Corollary~\ref{cor:finite-lift-core-operator-tail-menu-schedules}, after
scaling the \(h\)-th entrywise noise stream by
\[
  C_{\rm ent}(a_{m,h}^2+2b_{m,h}^2)^{1/2}
  (\sqrt{D_m}+\tau_m).
\]
Equivalently, if \(a_{m,h}\) and \(b_{m,h}\) are uniformly bounded on the
active prefixes, sparse, banded, and variance-profile Wigner/GOE/GUE-type
noise libraries at the normalized graph scale \(d_{m,h}^{-1/2}\) are absorbed
at the same \(O(\sqrt{D_m}+\tau_m)\) spectral scale as dense normalized
ensembles.

More explicitly, if
\[
  L_m^{\rm ot}=J_m(n_m+1)\le(eD_m)^\Gamma
\]
eventually, then along all integer dimensions \(D=2,3,\ldots\) it is enough to
take \(\tau_D=\sqrt{R\log(eD)}\) with \(R>\Gamma+1\).  Along sparse schedules
\(D_m\ge\exp(c m^\kappa)\), the same conclusion holds for
\(\tau_m=\sqrt{R\log(eD_m)}\) whenever \(R>\Gamma\).

For a one-dimensional bandwidth-\(w_{m,h}\) profile
\(|p-q|\le w_{m,h}\), one may take \(\Delta_{m,h}\le2w_{m,h}\).  Thus
off-diagonal proxy \(b_{m,h}/\sqrt{\max(1,2w_{m,h})}\) gives the displayed
scale with no additional dependence on the bandwidth.
\end{corollary}

\begin{proof}
The row-variance envelope in
Corollary~\ref{cor:finite-lift-core-entrywise-row-variance-menu} is bounded by
\[
  a_{m,h}^2+2\Delta_{m,h}\frac{b_{m,h}^2}{d_{m,h}}
  \le a_{m,h}^2+2b_{m,h}^2.
\]
That corollary gives the active-prefix finite-lift conclusion at the displayed
entry-profile scale, and
Corollary~\ref{cor:finite-lift-core-operator-tail-menu-schedules} gives the
explicit all-integer and sparse-dimension schedules.  The banded specialization
uses the elementary degree bound \(\Delta_{m,h}\le2w_{m,h}\).
\end{proof}

\begin{corollary}[Random graph-mask entry schedules]
\label{cor:finite-lift-core-random-graph-mask-entry-schedules}
In the setting of
Corollary~\ref{cor:finite-lift-core-entrywise-subgaussian-menu}, suppose that
the off-diagonal support is random.  More precisely, for each active
\((m,h,i)\) and each \(p<q\), let
\(\eta_{m,h,i,pq}\in\{0,1\}\) be a support indicator.  Assume that,
conditionally on the core environment, these indicators are independent within
each fixed \((m,h,i)\) and satisfy
\[
  \mathbb P\{\eta_{m,h,i,pq}=1\mid\mathcal H_m\}
  \le p_{m,h,i,pq}.
\]
Put
\[
  \lambda_{m,h}
  =\max_{i,p}
  \left(
  \sum_{q<p}p_{m,h,i,qp}+\sum_{p<q}p_{m,h,i,pq}
  \right),
  \qquad
  s_{m,h}=\max(1,\lambda_{m,h}+u_{m,h}),
\]
where \(u_{m,h}>0\) is deterministic.  Conditional on the core environment and
on all support indicators, assume that the upper-triangular entries are
independent and centered, the diagonal proxies satisfy
\(\alpha_{m,i,p}^{(h)}\le a_{m,h}\), and the off-diagonal entries vanish when
\(\eta_{m,h,i,pq}=0\) while, when \(\eta_{m,h,i,pq}=1\), their sub-Gaussian
proxies satisfy
\[
  \beta_{m,i,pq}^{(h)}\le \frac{b_{m,h}}{\sqrt{s_{m,h}}}.
\]
If
\[
  \sum_m\left[
  \mathbb P(F_m^{\rm lift})
  +\sum_{h\le J_m}n_mD_m
    \exp\left\{
    -\frac{u_{m,h}^2}{2(\lambda_{m,h}+u_{m,h}/3)}
    \right\}
  +2J_m(n_m+1)e^{-\tau_m^2}
  \right]<\infty,
\]
then all active-prefix and fixed-menu conclusions of
Corollary~\ref{cor:finite-lift-core-operator-tail-menu-borel} hold, after
scaling the \(h\)-th random-mask entrywise noise stream by
\[
  C_{\rm ent}(a_{m,h}^2+2b_{m,h}^2)^{1/2}
  (\sqrt{D_m}+\tau_m).
\]

The same conclusion holds without the Bernstein degree-exception term for any
dependent random mask law which is almost surely supported on graphs of maximum
degree at most \(\Delta_{m,h}\), after setting
\(s_{m,h}=\max(1,\Delta_{m,h})\).  Thus random \(d\)-regular masks, random
matchings, random bounded-degree subgraphs of a band, and other dependent
bounded-degree support laws are covered with no support-probability loss.

In particular, for inhomogeneous Erdos--Renyi masks with
\(p_{m,h,i,pq}\le p_{m,h}\), one has
\(\lambda_{m,h}\le(D_m-1)p_{m,h}\).  If
\[
  L_m^{\rm deg}=\log(eJ_mn_mD_m)+(1+\gamma)\log(e+m),
  \qquad
  u_{m,h}=\max\{\lambda_{m,h},A L_m^{\rm deg}\},
\]
with \(A>8/3\), then the degree-exception series is summable.  Thus random
graph-supported Wigner/GOE/GUE-type active-prefix libraries are covered at the
same displayed scale with off-diagonal normalization
\(s_{m,h}^{-1/2}\).  When the expected degree dominates the union logarithm,
this is the usual expected-degree normalization, up to an absolute constant.
\end{corollary}

\begin{proof}
Let \(F_m^{\rm deg}\) be the event that for some \(h\le J_m\), \(i\le n_m\),
and \(1\le p\le D_m\), the random support degree
\[
  N_{m,h,i,p}
  =\sum_{q<p}\eta_{m,h,i,qp}+\sum_{p<q}\eta_{m,h,i,pq}
\]
exceeds \(\lambda_{m,h}+u_{m,h}\).  For fixed \(m,h,i,p\), this is a sum of
independent \([0,1]\)-valued variables with conditional mean at most
\(\lambda_{m,h}\).  Bernstein's inequality gives
\[
  \mathbb P\{N_{m,h,i,p}>\lambda_{m,h}+u_{m,h}\mid\mathcal H_m\}
  \le
  \exp\left\{
    -\frac{u_{m,h}^2}{2(\lambda_{m,h}+u_{m,h}/3)}
  \right\}.
\]
A union bound over \(h,i,p\) gives the displayed degree-exception term.

Adjoin the support indicators to the conditioning sigma-field.  On
\((F_m^{\rm deg})^c\), every active off-diagonal support graph has maximum
degree at most \(s_{m,h}\), and the finite-lift core hypotheses are unchanged.
Hence the row-variance envelope is bounded by
\[
  a_{m,h}^2+2s_{m,h}\frac{b_{m,h}^2}{s_{m,h}}
  =a_{m,h}^2+2b_{m,h}^2.
\]
Apply Corollary~\ref{cor:finite-lift-core-entrywise-row-variance-menu} on this
event, with the exceptional event enlarged from \(F_m^{\rm lift}\) to
\(F_m^{\rm lift}\cup F_m^{\rm deg}\), and then use the displayed summability
condition and Borel--Cantelli.

For the homogeneous schedule, \(u_{m,h}\ge\lambda_{m,h}\), so
\[
  \frac{u_{m,h}^2}{2(\lambda_{m,h}+u_{m,h}/3)}
  \ge \frac38 u_{m,h}\ge\frac{3A}{8}L_m^{\rm deg}.
\]
Since \(A>8/3\), the union-count factor is summable in \(m\).

For the deterministic-degree variant, skip the Bernstein step and take
\(F_m^{\rm deg}=\emptyset\).  The same row-variance bound holds with
\(s_{m,h}=\max(1,\Delta_{m,h})\).
\end{proof}

\begin{corollary}[Weighted random-mask variance profiles]
\label{cor:finite-lift-core-weighted-random-mask-entry-schedules}
In the setting of
Corollary~\ref{cor:finite-lift-core-entrywise-subgaussian-menu}, suppose that
for each active \((m,h,i)\) and each \(p<q\) there is a random support
indicator \(\eta_{m,h,i,pq}\in\{0,1\}\).  Assume that, conditionally on
\(\mathcal H_m\), these indicators are independent within each fixed
\((m,h,i)\) and satisfy
\[
  \mathbb P\{\eta_{m,h,i,pq}=1\mid\mathcal H_m\}
  \le p_{m,h,i,pq}.
\]
Let \(a_{m,h}\ge0\), \(c_{m,h,i,pq}\ge0\), and \(t_{m,h}>0\) be deterministic.
Conditional on the core environment and on the support indicators, assume that
the upper-triangular entries are independent and centered, the diagonal proxies
satisfy \(\alpha_{m,i,p}^{(h)}\le a_{m,h}\), and the off-diagonal entries vanish
when \(\eta_{m,h,i,pq}=0\), while when \(\eta_{m,h,i,pq}=1\) their proxies
satisfy
\[
  \beta_{m,i,pq}^{(h)}\le c_{m,h,i,pq}.
\]
Put
\[
\begin{aligned}
  \mu_{m,h}
  &=
  \max_{i,p}
  2\left(
  \sum_{q<p}p_{m,h,i,qp}c_{m,h,i,qp}^2
  +\sum_{p<q}p_{m,h,i,pq}c_{m,h,i,pq}^2
  \right),\\
  \sigma_{m,h}^2
  &=
  \max_{i,p}
  4\left(
  \sum_{q<p}p_{m,h,i,qp}c_{m,h,i,qp}^4
  +\sum_{p<q}p_{m,h,i,pq}c_{m,h,i,pq}^4
  \right),\\
  b_{m,h}
  &=2\max_{i,p<q}c_{m,h,i,pq}^2,
  \qquad
  R_{m,h}^2=a_{m,h}^2+\mu_{m,h}+t_{m,h}.
\end{aligned}
\]
If
\[
  \sum_m\left[
  \mathbb P(F_m^{\rm lift})
  +\sum_{h\le J_m}n_mD_m
    \exp\left\{
    -\frac{t_{m,h}^2}{2(\sigma_{m,h}^2+b_{m,h}t_{m,h}/3)}
    \right\}
  +2J_m(n_m+1)e^{-\tau_m^2}
  \right]<\infty,
\]
with the Bernstein term omitted when \(\sigma_{m,h}^2+b_{m,h}t_{m,h}/3=0\),
then all active-prefix and fixed-menu conclusions of
Corollary~\ref{cor:finite-lift-core-operator-tail-menu-borel} hold, after
scaling the \(h\)-th weighted random-mask entrywise noise stream by
\[
  C_{\rm ent}R_{m,h}(\sqrt{D_m}+\tau_m).
\]
Thus the finite-lift core signs absorb inhomogeneous random variance-profile
Wigner/GOE/GUE-type masks whenever the random weighted row-variance
\(2\sum_q\eta_{pq}c_{pq}^2\) is eventually bounded by a summable Bernstein
budget.  This criterion can be strictly sharper than a maximum-degree bound
when the potential entries have highly nonuniform variance proxies.
\end{corollary}

\begin{proof}
For fixed \(m,h,i,p\), set
\[
  W_{m,h,i,p}
  =
  2\sum_{q<p}\eta_{m,h,i,qp}c_{m,h,i,qp}^2
  +2\sum_{p<q}\eta_{m,h,i,pq}c_{m,h,i,pq}^2 .
\]
Conditionally on \(\mathcal H_m\), this is a sum of independent
\([0,b_{m,h}]\)-valued variables, with conditional mean at most \(\mu_{m,h}\)
and total variance at most \(\sigma_{m,h}^2\).  Bernstein's inequality gives
\[
  \mathbb P\{W_{m,h,i,p}>\mu_{m,h}+t_{m,h}\mid\mathcal H_m\}
  \le
  \exp\left\{
    -\frac{t_{m,h}^2}{2(\sigma_{m,h}^2+b_{m,h}t_{m,h}/3)}
  \right\},
\]
with the stated zero-variance convention.  A union bound over \(h,i,p\) gives
the displayed support-profile exception.

Adjoin the support indicators to the conditioning sigma-field.  Outside this
exception, every active row satisfies
\[
  \alpha_{m,i,p}^{(h)2}
  +2\sum_{q<p}\beta_{m,i,qp}^{(h)2}
  +2\sum_{p<q}\beta_{m,i,pq}^{(h)2}
  \le a_{m,h}^2+\mu_{m,h}+t_{m,h}
  =R_{m,h}^2.
\]
Corollary~\ref{cor:finite-lift-core-entrywise-row-variance-menu} therefore
applies on the complement of the enlarged exceptional event, and
Borel--Cantelli gives the asserted active-prefix and fixed-menu conclusions.
\end{proof}

\begin{corollary}[Weighted random-mask Bernstein schedules]
\label{cor:finite-lift-core-weighted-random-mask-bernstein-schedules}
In the setting of
Corollary~\ref{cor:finite-lift-core-weighted-random-mask-entry-schedules}, put
\[
  L_m^{\rm var}=J_mn_mD_m,\qquad
  L_m^{\rm ot}=J_m(n_m+1),
\]
and, when \(\sigma_{m,h}^2+b_{m,h}t_{m,h}/3>0\),
\[
  \Phi_{m,h}
  =
  \frac{t_{m,h}^2}{2(\sigma_{m,h}^2+b_{m,h}t_{m,h}/3)}.
\]
If \(\sum_m\mathbb P(F_m^{\rm lift})<\infty\), \(\gamma>0\), and eventually
\[
  \Phi_{m,h}\ge \log L_m^{\rm var}+(1+\gamma)\log(e+m)
  \qquad(1\le h\le J_m),
\]
with the support term omitted in the zero-variance cases, while
\[
  \tau_m^2\ge \log(2L_m^{\rm ot})+(1+\gamma)\log(e+m),
\]
then the active-prefix and fixed-menu conclusions of
Corollary~\ref{cor:finite-lift-core-weighted-random-mask-entry-schedules} hold
almost surely.

In particular, suppose
\[
  \max(L_m^{\rm var},L_m^{\rm ot})\le(eD_m)^\Gamma
\]
eventually.  Along all integer dimensions \(D=2,3,\ldots\), the same conclusion
holds if
\[
  \Phi_{D,h}\ge R\log(eD),\qquad
  \tau_D^2\ge R\log(eD)
  \qquad(1\le h\le J_D)
\]
eventually for some \(R>\Gamma+1\).  Along sparse schedules
\(D_m\ge\exp(c m^\kappa)\), the same conclusion holds under the same
inequalities whenever \(R>\Gamma\).  The resulting entrywise noise scale is
\[
  C_{\rm ent}(a_{m,h}^2+\mu_{m,h}+t_{m,h})^{1/2}
  (\sqrt{D_m}+\tau_m).
\]
\end{corollary}

\begin{proof}
The arbitrary-growth inequalities give
\[
  \sum_{h\le J_m}n_mD_me^{-\Phi_{m,h}}
  \le(e+m)^{-(1+\gamma)}
\]
and
\[
  2J_m(n_m+1)e^{-\tau_m^2}
  \le(e+m)^{-(1+\gamma)}.
\]
Together with the lift-exception summability, this is the summability condition
in Corollary~\ref{cor:finite-lift-core-weighted-random-mask-entry-schedules}.

For all integer dimensions, both exceptional row factors are bounded by
constant multiples of \((eD)^{\Gamma-R}\), which is summable when
\(R>\Gamma+1\).  For sparse stretched-exponential dimension schedules, the same
bound is summable as soon as \(R>\Gamma\).
\end{proof}

\begin{corollary}[Bounded weighted random-mask incidence schedules]
\label{cor:finite-lift-core-bounded-weighted-random-mask-incidence-schedules}
In the setting of
Corollary~\ref{cor:finite-lift-core-weighted-random-mask-bernstein-schedules},
suppose that, for each active \(h\le J_m\), deterministic parameters
\[
  A_{m,h}\ge0,\qquad \rho_{m,h}>0,\qquad \omega_{m,h}>0
\]
satisfy
\[
  a_{m,h}^2\le A_{m,h}^2,\qquad
  \mu_{m,h}\le\rho_{m,h},\qquad
  \sigma_{m,h}^2\le\omega_{m,h}\rho_{m,h},\qquad
  b_{m,h}\le\omega_{m,h}.
\]
Put
\[
  \Lambda_m^{\rm mask}
  =
  \min_{1\le h\le J_m}\frac{\rho_{m,h}}{\omega_{m,h}},
  \qquad
  L_m^{\rm var}=J_mn_mD_m,\qquad
  L_m^{\rm ot}=J_m(n_m+1).
\]
Assume \(\sum_m\mathbb P(F_m^{\rm lift})<\infty\).  If, for some
\(\gamma>0\),
\[
  \Lambda_m^{\rm mask}
  \ge
  \frac83\left[\log(eL_m^{\rm var})+(1+\gamma)\log(e+m)\right]
\]
and
\[
  \tau_m^2
  \ge
  \log(2L_m^{\rm ot})+(1+\gamma)\log(e+m)
\]
eventually, then the active-prefix and fixed-menu conclusions of
Corollary~\ref{cor:finite-lift-core-operator-tail-menu-borel} hold after
scaling the \(h\)-th bounded weighted random-mask entrywise noise stream by
\[
  C_{\rm ent}(A_{m,h}^2+2\rho_{m,h})^{1/2}(\sqrt{D_m}+\tau_m).
\]

In particular, suppose
\(\max\{L_m^{\rm var},L_m^{\rm ot}\}\le(eD_m)^\Gamma\) eventually.  Along all
integer dimensions \(D=2,3,\ldots\), the same conclusion holds if
\[
  \Lambda_D^{\rm mask}\ge R\log(eD),
  \qquad
  \tau_D^2\ge S\log(eD),
\]
eventually for constants
\[
  R>\frac{8(\Gamma+1)}3,\qquad S>\Gamma+1.
\]
Along any sparse schedule \(D_m\ge\exp(c m^\alpha)\), \(c,\alpha>0\), the same
conclusion holds if
\[
  \Lambda_m^{\rm mask}\ge R\log(eD_m),
  \qquad
  \tau_m^2\ge S\log(eD_m),
\]
eventually for constants
\[
  R>\frac{8\Gamma}3,\qquad S>\Gamma .
\]

The same statement holds in the predictable-profile setting of
Corollary~\ref{cor:finite-lift-core-predictable-weighted-random-mask-schedules}
after replacing \(a_{m,h},\mu_{m,h},\sigma_{m,h},b_{m,h}\) by the deterministic
envelopes \(\bar a_{m,h},\bar\mu_{m,h},\bar\sigma_{m,h},\bar b_{m,h}\) and
adding \(\sum_m\mathbb P(F_m^{\rm pr})<\infty\).  Thus independent or
environment-predictable inhomogeneous Wigner/GOE/GUE-type masks are covered by
the normalized mean row-variance ratio \(\rho/\omega\), rather than by maximum
degree.
\end{corollary}

\begin{proof}
Apply
Corollary~\ref{cor:finite-lift-core-weighted-random-mask-bernstein-schedules}
with \(t_{m,h}=\rho_{m,h}\).  The resulting row-variance scale satisfies
\[
  a_{m,h}^2+\mu_{m,h}+t_{m,h}
  \le A_{m,h}^2+2\rho_{m,h}.
\]
The support exponent obeys
\[
  \Phi_{m,h}
  =
  \frac{\rho_{m,h}^2}
       {2(\sigma_{m,h}^2+b_{m,h}\rho_{m,h}/3)}
  \ge
  \frac{\rho_{m,h}^2}
       {2(\omega_{m,h}\rho_{m,h}
          +\omega_{m,h}\rho_{m,h}/3)}
  =
  \frac{3\rho_{m,h}}{8\omega_{m,h}}.
\]
Hence \(\Phi_{m,h}\ge3\Lambda_m^{\rm mask}/8\) for every active \(h\).  The
arbitrary-growth, all-integer, and sparse estimates are the corresponding
estimates in
Corollary~\ref{cor:finite-lift-core-weighted-random-mask-bernstein-schedules}
with this lower bound.  The predictable-profile variant is identical after
working on \((F_m^{\rm pr})^c\) and using the deterministic barred envelopes.
\end{proof}

\begin{corollary}[Predictable weighted random-mask profiles]
\label{cor:finite-lift-core-predictable-weighted-random-mask-schedules}
In the setting of
Corollary~\ref{cor:finite-lift-core-entrywise-subgaussian-menu}, keep the same
finite-lift core assumptions and active-prefix exhaustion.  Allow the
random-mask intensity and proxy profile to be chosen from the core environment:
for each active \(h\le J_m\), let
\[
  a_{m,h},\qquad
  p_{m,h,i,pq}\in[0,1],\qquad
  c_{m,h,i,pq}\ge0
\]
be \(\mathcal H_m\)-measurable.  Conditionally on \(\mathcal H_m\), assume
that the support indicators \(\eta_{m,h,i,pq}\) are independent within each
fixed \((m,h,i)\) and satisfy
\[
  \mathbb P\{\eta_{m,h,i,pq}=1\mid\mathcal H_m\}
  \le p_{m,h,i,pq}.
\]
Conditional on \(\mathcal H_m\) and on the supports, assume that the
upper-triangular entries are independent and centered, the diagonal proxies are
bounded by \(a_{m,h}\), and the off-diagonal entries vanish when
\(\eta_{m,h,i,pq}=0\), while their proxies are bounded by
\(c_{m,h,i,pq}\) when \(\eta_{m,h,i,pq}=1\).

Define the conditional profile budgets
\[
\begin{aligned}
  \mu_{m,h}^{\rm pr}
  &=
  \max_{i,p}
  2\left(
  \sum_{q<p}p_{m,h,i,qp}c_{m,h,i,qp}^2
  +\sum_{p<q}p_{m,h,i,pq}c_{m,h,i,pq}^2
  \right),\\
  (\sigma_{m,h}^{\rm pr})^2
  &=
  \max_{i,p}
  4\left(
  \sum_{q<p}p_{m,h,i,qp}c_{m,h,i,qp}^4
  +\sum_{p<q}p_{m,h,i,pq}c_{m,h,i,pq}^4
  \right),\\
  b_{m,h}^{\rm pr}
  &=2\max_{i,p<q}c_{m,h,i,pq}^2 .
\end{aligned}
\]
Let \(\bar a_{m,h},\bar\mu_{m,h},\bar\sigma_{m,h},\bar b_{m,h}\ge0\),
\(t_{m,h}>0\), and \(\tau_m>0\) be deterministic, and let
\(F_m^{\rm pr}\in\mathcal H_m\) be the event that for some \(h\le J_m\),
\[
  a_{m,h}^2>\bar a_{m,h}^2,\qquad
  \mu_{m,h}^{\rm pr}>\bar\mu_{m,h},\qquad
  (\sigma_{m,h}^{\rm pr})^2>\bar\sigma_{m,h}^2,
  \quad\text{or}\quad
  b_{m,h}^{\rm pr}>\bar b_{m,h}.
\]
If
\[
  \begin{aligned}
  \sum_m\Bigg[
  &\mathbb P(F_m^{\rm lift})
  +\mathbb P(F_m^{\rm pr})\\
  &+\sum_{h\le J_m}n_mD_m
    \exp\left\{
    -\frac{t_{m,h}^2}
    {2(\bar\sigma_{m,h}^2+\bar b_{m,h}t_{m,h}/3)}
    \right\}\\
  &+2J_m(n_m+1)e^{-\tau_m^2}
  \Bigg]<\infty,
  \end{aligned}
\]
with the Bernstein term omitted when
\(\bar\sigma_{m,h}^2+\bar b_{m,h}t_{m,h}/3=0\), then all active-prefix and
fixed-menu conclusions of
Corollary~\ref{cor:finite-lift-core-operator-tail-menu-borel} hold, after
scaling the \(h\)-th predictable random-mask entrywise noise stream by
\[
  C_{\rm ent}(\bar a_{m,h}^2+\bar\mu_{m,h}+t_{m,h})^{1/2}
  (\sqrt{D_m}+\tau_m).
\]

Equivalently, the deterministic-profile theorem remains valid after replacing
the fixed variance profile by an \(\mathcal H_m\)-measurable random or adaptive
profile, provided the conditional Bernstein envelopes exceed their prescribed
deterministic budgets only summably.  If
\[
  \bar\Phi_{m,h}
  =
  \frac{t_{m,h}^2}
  {2(\bar\sigma_{m,h}^2+\bar b_{m,h}t_{m,h}/3)}
\]
and \(L_m^{\rm var}=J_mn_mD_m\), \(L_m^{\rm ot}=J_m(n_m+1)\), then the
arbitrary-growth, all-integer, and sparse-dimension schedule conclusions of
Corollary~\ref{cor:finite-lift-core-weighted-random-mask-bernstein-schedules}
hold with \(\Phi_{m,h}\) replaced by \(\bar\Phi_{m,h}\), as long as
\(\sum_m\mathbb P(F_m^{\rm pr})<\infty\) is added to the lift-exception
summability.
\end{corollary}

\begin{proof}
Work on \((F_m^{\rm pr})^c\) and condition on \(\mathcal H_m\).  For fixed
\(m,h,i,p\), the weighted row profile
\[
  W_{m,h,i,p}
  =
  2\sum_{q<p}\eta_{m,h,i,qp}c_{m,h,i,qp}^2
  +2\sum_{p<q}\eta_{m,h,i,pq}c_{m,h,i,pq}^2
\]
is a sum of independent variables bounded by \(\bar b_{m,h}\), with conditional
mean at most \(\bar\mu_{m,h}\) and total conditional variance at most
\(\bar\sigma_{m,h}^2\).  Bernstein's inequality gives
\[
  \mathbb P\{W_{m,h,i,p}>\bar\mu_{m,h}+t_{m,h}\mid\mathcal H_m\}
  \le
  \exp\left\{
    -\frac{t_{m,h}^2}
    {2(\bar\sigma_{m,h}^2+\bar b_{m,h}t_{m,h}/3)}
  \right\},
\]
with the stated zero-variance convention.  A union bound over \(h,i,p\) gives
the displayed predictable-profile exception.

Adjoin the supports to the conditioning sigma-field.  Outside the predictable
profile exception, every active row satisfies
\[
  \alpha_{m,i,p}^{(h)2}
  +2\sum_{q<p}\beta_{m,i,qp}^{(h)2}
  +2\sum_{p<q}\beta_{m,i,pq}^{(h)2}
  \le \bar a_{m,h}^2+\bar\mu_{m,h}+t_{m,h}.
\]
Corollary~\ref{cor:finite-lift-core-entrywise-row-variance-menu} therefore
applies on the complement of
\(F_m^{\rm lift}\cup F_m^{\rm pr}\) and the support-profile exception, while
Corollary~\ref{cor:finite-lift-core-operator-tail-menu-borel} supplies the
same active-prefix conclusion after the operator-tail union term is included.
Borel--Cantelli gives the almost-sure eventual statement.

The schedule paragraph is exactly the proof of
Corollary~\ref{cor:finite-lift-core-weighted-random-mask-bernstein-schedules},
with \(\Phi_{m,h}\) replaced by \(\bar\Phi_{m,h}\) and the summable events
\(F_m^{\rm pr}\) added to \(F_m^{\rm lift}\).
\end{proof}

\begin{corollary}[Dependent weighted-budget mask laws]
\label{cor:finite-lift-core-dependent-weighted-budget-masks}
In the setting of
Corollary~\ref{cor:finite-lift-core-entrywise-subgaussian-menu}, keep the same
finite-lift core assumptions and active-prefix exhaustion.  For each active
\((m,h,i)\) and each \(p<q\), let
\(\eta_{m,h,i,pq}\in\{0,1\}\) be an arbitrary random support indicator; no
independence, product structure, or identical distribution is assumed for these
supports.  Let \(a_{m,h}\ge0\), \(c_{m,h,i,pq}\ge0\), \(R_{m,h}>0\), and
\(\tau_m>0\) be deterministic.  Conditional on the core environment and on all
active supports, assume that the upper-triangular entries are independent and
centered, the diagonal proxies satisfy
\(\alpha_{m,i,p}^{(h)}\le a_{m,h}\), and the off-diagonal entries vanish when
\(\eta_{m,h,i,pq}=0\), while when \(\eta_{m,h,i,pq}=1\) their proxies satisfy
\[
  \beta_{m,i,pq}^{(h)}\le c_{m,h,i,pq}.
\]
Let \(F_m^{\rm wb}\) be the weighted-budget failure event that for some
\(h\le J_m\), \(i\le n_m\), and \(1\le p\le D_m\),
\[
  a_{m,h}^2
  +2\sum_{q<p}\eta_{m,h,i,qp}c_{m,h,i,qp}^2
  +2\sum_{p<q}\eta_{m,h,i,pq}c_{m,h,i,pq}^2
  >R_{m,h}^2 .
\]
If
\[
  \sum_m\left[
  \mathbb P(F_m^{\rm lift})
  +\mathbb P(F_m^{\rm wb})
  +2J_m(n_m+1)e^{-\tau_m^2}
  \right]<\infty,
\]
then, with probability one, for all sufficiently large \(m\), the active-prefix
and fixed-menu conclusions of
Corollary~\ref{cor:finite-lift-core-operator-tail-menu-borel} hold after
scaling the \(h\)-th dependent-mask entrywise noise stream by
\[
  C_{\rm ent}R_{m,h}(\sqrt{D_m}+\tau_m).
\]
The common sign vector may be chosen from the finite-lift core environment; it
does not need to observe the random support masks.

Consequently every support law whose realized weighted row budgets are bounded
deterministically has no support-probability cost.  This includes random regular
or biregular masks, matchings, sampling without replacement, random
bounded-degree subgraphs of a band, determinantal or negatively dependent
support samplers, and mixtures of such laws, whenever the displayed weighted
row budget is bounded by \(R_{m,h}^2\) outside a summable event.  The explicit
operator-tail schedules of
Corollary~\ref{cor:finite-lift-core-operator-tail-menu-schedules} apply with no
additional change: if \(J_m(n_m+1)\le(eD_m)^\Gamma\), then
\(\tau_D^2=S\log(eD)\) with \(S>\Gamma+1\) suffices along all integer
dimensions, while sparse stretched-exponential dimension schedules need only
\(S>\Gamma\).
\end{corollary}

\begin{proof}
Enlarge the conditioning sigma-field by all active support indicators.  On
\((F_m^{\rm wb})^c\), every active row satisfies the row-variance hypothesis
\[
  \alpha_{m,i,p}^{(h)2}
  +2\sum_{q<p}\beta_{m,i,qp}^{(h)2}
  +2\sum_{p<q}\beta_{m,i,pq}^{(h)2}
  \le R_{m,h}^2 .
\]
Thus Corollary~\ref{cor:finite-lift-core-entrywise-row-variance-menu} applies
conditionally on the supports, with \(Q_{m,h}=R_{m,h}\).  The finite-lift core
signs are obtained before revealing the supports, so the same
\(\mathcal H_m\)-measurable sign vector is admissible after the conditioning is
enlarged.  The only remaining random failures are
\(F_m^{\rm lift}\), \(F_m^{\rm wb}\), and the usual operator-tail union event
of size \(2J_m(n_m+1)e^{-\tau_m^2}\).  Borel--Cantelli gives the active-prefix
and fixed-menu almost-sure conclusion, and the schedule statement is
Corollary~\ref{cor:finite-lift-core-operator-tail-menu-schedules}.
\end{proof}

\begin{corollary}[Bounded weighted-degree dependent masks]
\label{cor:finite-lift-core-bounded-weighted-degree-dependent-masks}
In the setting of
Corollary~\ref{cor:finite-lift-core-dependent-weighted-budget-masks}, suppose
that for every active \(h\le J_m\) there are deterministic
\(a_{m,h},b_{m,h}\ge0\), \(s_{m,h}>0\), and weights
\(\omega_{m,h,i,pq}\ge0\) such that the diagonal proxies are at most
\(a_{m,h}\) and, whenever \(\eta_{m,h,i,pq}=1\),
\[
  \beta_{m,i,pq}^{(h)}
  \le b_{m,h}\left(\frac{\omega_{m,h,i,pq}}{s_{m,h}}\right)^{1/2}.
\]
Let \(F_m^{\rm wdeg}\) be the event that for some \(h\le J_m\), \(i\le n_m\),
and \(1\le p\le D_m\),
\[
  \sum_{q<p}\eta_{m,h,i,qp}\omega_{m,h,i,qp}
  +\sum_{p<q}\eta_{m,h,i,pq}\omega_{m,h,i,pq}
  >s_{m,h}.
\]
If
\[
  \sum_m\left[
  \mathbb P(F_m^{\rm lift})
  +\mathbb P(F_m^{\rm wdeg})
  +2J_m(n_m+1)e^{-\tau_m^2}
  \right]<\infty,
\]
then, with probability one, for all sufficiently large \(m\), all active-prefix
and fixed-menu conclusions of
Corollary~\ref{cor:finite-lift-core-operator-tail-menu-borel} hold after
scaling the \(h\)-th dependent weighted-degree noise stream by
\[
  C_{\rm ent}(a_{m,h}^2+2b_{m,h}^2)^{1/2}
  (\sqrt{D_m}+\tau_m).
\]

In particular, when the support law is almost surely supported on the displayed
weighted-degree constraint, \(F_m^{\rm wdeg}=\emptyset\).  For unweighted
degree-bounded masks this says: if every active support graph has maximum
degree at most \(\Delta_{m,h}\) and the off-diagonal proxies are bounded by
\(b_{m,h}/\sqrt{\max(1,\Delta_{m,h})}\), then the finite-lift core signs absorb
the noise at scale
\[
  C_{\rm ent}(a_{m,h}^2+2b_{m,h}^2)^{1/2}(\sqrt{D_m}+\tau_m)
\]
with no support exception.  Thus random \(d\)-regular and biregular masks,
matchings, \(k\)-nearest-neighbor masks sampled without replacement, bounded
band subgraphs, and weighted-degree graph samplers all have the same normalized
finite-lift scale whenever their weighted degree is deterministically bounded
or fails only summably.

The explicit operator-tail schedules are exactly those of
Corollary~\ref{cor:finite-lift-core-operator-tail-menu-schedules}: if
\(J_m(n_m+1)\le(eD_m)^\Gamma\), then all integer dimensions are covered by
\(\tau_D^2=S\log(eD)\) with \(S>\Gamma+1\), and sparse stretched-exponential
dimension schedules need only \(S>\Gamma\).
\end{corollary}

\begin{proof}
On \((F_m^{\rm wdeg})^c\), put
\[
  W_{m,h,i,p}^{\omega}
  =
  \sum_{q<p}\eta_{m,h,i,qp}\omega_{m,h,i,qp}
  +\sum_{p<q}\eta_{m,h,i,pq}\omega_{m,h,i,pq}
  \le s_{m,h}.
\]
Then every active row satisfies
\[
\begin{aligned}
  &a_{m,h}^2
  +2\sum_{q<p}\eta_{m,h,i,qp}\beta_{m,i,qp}^{(h)2}
  +2\sum_{p<q}\eta_{m,h,i,pq}\beta_{m,i,pq}^{(h)2}\\
  &\qquad\le
  a_{m,h}^2
  +\frac{2b_{m,h}^2}{s_{m,h}}W_{m,h,i,p}^{\omega}
  \le a_{m,h}^2+2b_{m,h}^2 .
\end{aligned}
\]
Therefore
Corollary~\ref{cor:finite-lift-core-dependent-weighted-budget-masks} applies
with \(R_{m,h}=(a_{m,h}^2+2b_{m,h}^2)^{1/2}\) and
\(F_m^{\rm wb}=F_m^{\rm wdeg}\).  The examples are the corresponding choices
of the deterministic weights \(\omega\).  The final schedule statement again
comes from Corollary~\ref{cor:finite-lift-core-operator-tail-menu-schedules}.
\end{proof}

\begin{corollary}[Algebraic cores plus full random operator-tail noise]
\label{cor:random-algebraic-plus-gaussian-noise}
For every \(C_0\ge1\) there is \(C(C_0)\) with the following property.  Let
\(B_1,\ldots,B_n\) be random self-adjoint matrices in \(\M_D(\C)\).  Suppose
that, outside an event \(E\), there is a finite-dimensional C*-subalgebra
\(\mcA_\omega\subset\M_D(\C)\) such that
\[
  \dim_\C\mcA_\omega\le C_0n,\qquad
  B_i(\omega)\in\mcA_\omega,\qquad
  \|B_i(\omega)\|_{\op}\le1
\]
for every \(i\).  Let \(Y_1,\ldots,Y_n\) be centered self-adjoint random
matrices, independent of \(B_1,\ldots,B_n\), satisfying the operator-tail
hypothesis of Corollary~\ref{cor:random-gaussian-tail-full} with constant
\(K_0\).  Put
\[
  s_{D,n}:=\sqrt D+\sqrt{3\log(en)},\qquad
  Z_i:=\frac{Y_i}{K_0s_{D,n}},\qquad
  X_i:=B_i+Z_i .
\]
Then, with probability at least \(1-\mathbb P(E)-2(n+1)(en)^{-3}\),
\[
  \disc(X_1,\ldots,X_n)\le C(C_0)\sqrt n .
\]
On the same event, all \(X_i\) have operator norm at most \(2\).  Hence the
normalized contractions
\[
  \widehat X_i:=\frac12X_i
\]
satisfy
\[
  \disc(\widehat X_1,\ldots,\widehat X_n)\le C(C_0)\sqrt n .
\]
Thus dimension-controlled random algebraic cores remain in the Matrix Spencer
regime after adding full centered GOE/GUE, Wigner-type, or other
Gaussian-operator-tail noise in arbitrary ambient dimension, once that noise is
put at its natural spectral scale.  The noise matrices themselves may be
dependent, provided the displayed scalar-combination tail holds.
\end{corollary}

\begin{proof}
On a realization outside \(E\), apply Theorem~\ref{thm:akbassra} to
\(B_1,\ldots,B_n\) and choose signs \(\eps_i^\star\) such that
\[
  \left\|\sum_i\eps_i^\star B_i\right\|_{\op}\le C(C_0)\sqrt n .
\]
The signs are measurable with respect to the algebraic-core data and are
independent of the noise.  Conditional on the core realization, the operator-tail
hypothesis applied with \(a=\eps^\star\) gives
\[
  \mathbb P\left\{
  \left\|\sum_i\eps_i^\star Y_i\right\|_{\op}>
  K_0\sqrt n\,s_{D,n}
  \,\middle|\,B_1,\ldots,B_n
  \right\}
  \le 2(en)^{-3}.
\]
Also, applying the same tail to \(a=e_i\) and union-bounding over \(i\) shows
that all \(\|Z_i\|_{\op}\le1\) except with conditional probability at most
\(2n(en)^{-3}\).  Therefore, outside a further event of probability at most
\(2(n+1)(en)^{-3}\),
\[
  \left\|\sum_i\eps_i^\star X_i\right\|_{\op}
  \le C(C_0)\sqrt n+\sqrt n
  \le C'(C_0)\sqrt n,
\]
and all \(\|X_i\|_{\op}\le2\).  Absorb the harmless additive \(1\) into
\(C(C_0)\) and divide by \(2\) for the normalized contraction statement.
\end{proof}

\begin{corollary}[Algebraic core menus plus mixed dominated operator-tail noise]
\label{cor:algebraic-core-mixed-dominated-operator-tail-menu-borel}
For each \(m\ge1\), let \(D_m,n_m,J_m\ge1\), put
\(\Delta_m:=D_m+n_m+1\) and \(\ell_m:=\log(e\Delta_m)\), and let
\(\mathcal H_m\) be a row environment.  For \(h\ge1\), let
\[
  B_{m,h,1},\ldots,B_{m,h,n_m},
  \qquad
  Y_{m,h,1},\ldots,Y_{m,h,n_m}
\]
be self-adjoint random matrices in \(\M_{D_m}(\C)\).  Fix \(C_0\ge1\), and let
\(K_{\rm alg}=C(C_0)\) be the Akbas--Sra constant from
Theorem~\ref{thm:akbassra}.  Assume that there are events
\(F_m^{\rm alg}\in\mathcal H_m\) with
\(\sum_m\mathbb P(F_m^{\rm alg})<\infty\) such that, outside
\(F_m^{\rm alg}\), for every active \(h\le J_m\) there is a
finite-dimensional C*-subalgebra
\(\mcA_{m,h}\subset\M_{D_m}(\C)\) with
\[
  B_{m,h,i}\in\mcA_{m,h},\qquad
  \|B_{m,h,i}\|_{\op}\le1
  \quad(1\le i\le n_m),
\]
and the active-prefix intrinsic dimension obeys
\[
  \sum_{h=1}^{J_m}\dim_\C\mcA_{m,h}\le C_0n_m .
\]

For each active \(h\le J_m\), tag the perturbation stream either \(\psi\) or
\(\mu\), choose \(\varepsilon_{m,h}\in(0,1]\), \(A_{m,h}\ge1\), and
\(s_{m,h}>0\), and define \(r_{m,h}\), \(V_{m,h}\), and \(q_{m,h}\) as in
Corollary~\ref{cor:mixed-dominated-operator-tail-active-libraries}, with
\(Y_{m,h,i}\) in place of \(X_{m,h,i}\).  Assume that, after adjoining
\(\mathcal H_m\) to the predictable sigma-fields, the corresponding
sub-Weibull or finite-moment dominated hypotheses and the matching truncated
tail-drift hypotheses hold outside events \(E_{m,h}\).

Let \(K_{\rm tail},\eta>0\) and set
\[
  \begin{aligned}
  \lambda_{m,h}&:=\max\{1,K_{\rm tail}+\eta\}\,q_{m,h},\\
  Z_{m,h,i}&:=\frac{Y_{m,h,i}}{\lambda_{m,h}},\\
  X_{m,h,i}&:=B_{m,h,i}+Z_{m,h,i}.
  \end{aligned}
\]
If \(\ell_m\le B\sqrt{n_m}\) eventually,
\[
  \sum_m\left[
    \mathbb P(F_m^{\rm alg})
    +\sum_{h=1}^{J_m}\bigl(\mathbb P(E_{m,h})+\varepsilon_{m,h}\bigr)
  \right]<\infty,
\]
and
\[
  \sum_mJ_m\Delta_m
  (e\Delta_m)^{-\gamma_{\rm tc}(K_{\rm tail},1,1,B)}<\infty,
\]
then, with probability one, for all sufficiently large \(m\), one
\(\mathcal H_m\)-measurable sign vector \(\eps_m^\star\) satisfies
\[
  \left\|\sum_{i=1}^{n_m}\eps_{m,i}^\star X_{m,h,i}\right\|_{\op}
  \le (K_{\rm alg}+1)\sqrt{n_m}
  \qquad(1\le h\le J_m).
\]
On the same event,
\[
  \max_{h\le J_m}\max_{i\le n_m}\|X_{m,h,i}\|_{\op}\le2 .
\]
Consequently the normalized contractions
\(\widehat X_{m,h,i}:=X_{m,h,i}/2\) have common active-prefix discrepancy at
most \((K_{\rm alg}+1)\sqrt{n_m}/2\), and every fixed menu element is
eventually covered whenever \(J_m\to\infty\).

In particular, if \(J_m\Delta_m\le(e\Delta_m)^\Gamma\),
\(J_m\le(e\Delta_m)^\Lambda\), \(n_m\ge a\ell_m^2\) eventually, and
\(\varepsilon_{m,h}=(e\Delta_m)^{-c}\), then after adding
\(\sum_m\mathbb P(F_m^{\rm alg})<\infty\) the cap-budget part is summable for
all integer dimensions when \(c>\Lambda+1\), and along every
stretched-exponential sparse dimension schedule when \(c>\Lambda\).  The
matrix-tail margins are
\[
  K_{\rm tail}>\Theta_{\rm tc}(\Gamma+1,a,1,1)
\]
and
\[
  K_{\rm tail}>\Theta_{\rm tc}(\Gamma,a,1,1),
\]
respectively.
\end{corollary}

\begin{proof}
Work outside \(F_m^{\rm alg}\).  Form the direct-sum algebra
\[
  \mcA_m^{(J)}:=\bigoplus_{h=1}^{J_m}\mcA_{m,h}
\]
and the direct-sum core tuple
\[
  \widetilde B_{m,i}:=\bigoplus_{h=1}^{J_m}B_{m,h,i}
  \qquad(1\le i\le n_m).
\]
Then \(\dim_\C\mcA_m^{(J)}\le C_0n_m\), each \(\widetilde B_{m,i}\) is a
self-adjoint contraction in \(\mcA_m^{(J)}\), and
Theorem~\ref{thm:akbassra} gives signs \(\eps_m^\star\) such that
\[
  \max_{h\le J_m}
  \left\|\sum_i\eps_{m,i}^\star B_{m,h,i}\right\|_{\op}
  =
  \left\|\sum_i\eps_{m,i}^\star\widetilde B_{m,i}\right\|_{\op}
  \le K_{\rm alg}\sqrt{n_m}.
\]
Choose the lexicographically first successful sign vector to make
\(\eps_m^\star\) \(\mathcal H_m\)-measurable.

Apply
Corollary~\ref{cor:signed-core-mixed-dominated-operator-tail-libraries} to the
present arrays \(B_{m,h,i}\) and \(Y_{m,h,i}\), with
\(K_{\rm core}=K_{\rm alg}\) and \(F_m^{\rm core}=F_m^{\rm alg}\).
The direct-sum Akbas--Sra signing verifies the signed-core hypothesis, and the
displayed mixed dominated assumptions are exactly the perturbation hypotheses.
The active-prefix conclusion, fixed-menu coverage, and all-integer/sparse
thresholds are the corresponding conclusions of the signed-core mixed
dominated corollary.
\end{proof}

\begin{corollary}[Structured core catalogue with mixed dominated noise]
\label{cor:structured-core-catalogue-mixed-dominated-operator-tail}
The mixed dominated operator-tail active-menu conclusion applies, with the
same normalizations and all-integer/sparse thresholds, to each of the following
structured random-matrix core catalogues.

\begin{enumerate}[label=\textup{(\alph*)}]
\item Finite-lift structured cores satisfying the active-prefix finite-lift
budget in
Corollary~\ref{cor:finite-lift-core-mixed-dominated-operator-tail-menu-borel}.
This includes scalar Toeplitz and Hankel cores, one-level fixed-block
Toeplitz/Hankel cores, fixed-level multilevel block Toeplitz/Hankel cores when
the active-prefix symbol dimension is \(O(n_m)\), and banded multilevel block
Toeplitz/Hankel cores when the active-prefix sum of anisotropic bandwidth
products is \(O(n_m)\).

\item Intrinsic algebraic cores satisfying the active-prefix dimension budget in
Corollary~\ref{cor:algebraic-core-mixed-dominated-operator-tail-menu-borel}.
This includes common finite-dimensional C*-algebra cores and finite group
convolution cores whenever the active-prefix sums of \(\dim_\C\mcA_{m,h}\), or
respectively \(|G_{m,h}|\), are \(O(n_m)\).
\end{enumerate}

In either case, each active perturbation stream may be tagged \(\psi\) or
\(\mu\) and scaled by its own
\(\lambda_{m,h}=\max\{1,K_{\rm tail}+\eta\}q_{m,h}\).  Under the corresponding
cap-budget, drift, and matrix-tail summability assumptions, the same
environment-measurable signs balance all active normalized sums
\[
  \left\|\sum_i\eps_{m,i}^\star
  \left(B_{m,h,i}+\frac{Y_{m,h,i}}{\lambda_{m,h}}\right)\right\|_{\op}
  =O(\sqrt{n_m})
  \qquad(1\le h\le J_m)
\]
eventually almost surely, and the normalized contractions have the same
active-prefix Matrix Spencer bound.
\end{corollary}

\begin{proof}
For the finite-lift catalogue, the cited Toeplitz/Hankel and banded multilevel
theorems give finite block lifts with active-prefix entropy equal, up to the
fixed structural constants, to the displayed symbol or bandwidth budget.
Corollary~\ref{cor:finite-lift-core-mixed-dominated-operator-tail-menu-borel}
therefore applies.

For the algebraic catalogue, Theorem~\ref{thm:akbassra} applies to the direct
sum of the active finite-dimensional C*-algebras, and
Corollary~\ref{cor:group-akbassra} identifies the group convolution case with
the dimension budget \(|G|\).  Hence
Corollary~\ref{cor:algebraic-core-mixed-dominated-operator-tail-menu-borel}
applies.  The displayed normalization and schedule assertions are exactly the
conclusions of those two mixed dominated corollaries.
\end{proof}

\begin{corollary}[Concrete structured size budgets with mixed dominated noise]
\label{cor:structured-size-budget-mixed-dominated-operator-tail}
Fix \(r_0,b_0,C_0\ge1\).  In
Corollary~\ref{cor:finite-lift-core-mixed-dominated-operator-tail-menu-borel},
the abstract active-prefix finite-lift assumption may be replaced by the
following concrete structural condition.  Outside a summable structural event,
every active core tuple is a self-adjoint contraction tuple of one of the
following types, with \(r\le r_0\) and \(b\le b_0\):
\begin{enumerate}[label=\textup{(\roman*)}]
\item scalar Toeplitz or Hankel on \(d_{m,h}\) coordinates, with
\(S_{m,h}=d_{m,h}\);
\item one-level \(b\times b\) block Toeplitz or Hankel on ambient dimension
\(D_{m,h}\), with \(S_{m,h}=D_{m,h}\);
\item full-bandwidth \(r\)-level \(b\times b\) block Toeplitz or Hankel on
ambient dimension \(D_{m,h}\), with \(S_{m,h}=D_{m,h}\);
\item banded \(r\)-level \(b\times b\) block Toeplitz or Hankel with
anisotropic half-bandwidths \(w_{m,h,1},\ldots,w_{m,h,r}\) and boxes satisfying
the hypotheses of Corollary~\ref{cor:anisotropic-bandml}, with
\[
  S_{m,h}:=\prod_{j=1}^r(2w_{m,h,j}+1).
\]
\end{enumerate}
If
\[
  \sum_{h=1}^{J_m}S_{m,h}\le C_0n_m
\]
eventually, then the full active-prefix conclusion, fixed-menu coverage, and
all-integer/sparse thresholds of
Corollary~\ref{cor:finite-lift-core-mixed-dominated-operator-tail-menu-borel}
hold, with constants depending only on \(r_0,b_0,C_0\).

Likewise, in
Corollary~\ref{cor:algebraic-core-mixed-dominated-operator-tail-menu-borel},
the active algebraic dimension condition is satisfied whenever, outside a
summable structural event, every active core tuple is either contained in a
finite-dimensional C*-algebra \(\mcA_{m,h}\) or is a finite group-convolution
core on a group \(G_{m,h}\), and
\[
  \sum_{h=1}^{J_m}s_{m,h}\le C_0n_m.
\]
Here
\[
  s_{m,h}:=
  \begin{cases}
  \dim_\C\mcA_{m,h},&\text{in the C*-algebra case},\\
  |G_{m,h}|,&\text{in the group-convolution case}.
  \end{cases}
\]
Thus those algebraic and group-convolution active menus have the same mixed
dominated perturbation conclusion and the same thresholds.
\end{corollary}

\begin{proof}
For the finite-lift alternatives, the proofs of
Theorem~\ref{thm:toep}, Corollary~\ref{cor:hankel},
Theorem~\ref{thm:blocktoep1}, Theorem~\ref{thm:blocktoepfull}, and
Corollary~\ref{cor:anisotropic-bandml} supply finite block lifts with
\(\Lambda K\) bounded in terms of \(r_0,b_0\) and active-prefix entropy bounded
by \(C(r_0,b_0)\sum_{h\le J_m}S_{m,h}\).  The displayed size budget therefore
implies the active-prefix finite-lift hypothesis with \(M_0\) and \(C_E\)
depending only on \(r_0,b_0,C_0\), and
Corollary~\ref{cor:finite-lift-core-mixed-dominated-operator-tail-menu-borel}
applies.

For the algebraic alternatives, finite-dimensional C*-algebra cores are already
in the hypothesis of
Corollary~\ref{cor:algebraic-core-mixed-dominated-operator-tail-menu-borel}.
For group convolution, the reduced convolution C*-algebra has complex dimension
\(\sum_{\pi\in\widehat G}d_\pi^2=|G|\), so the same size budget is precisely the
active-prefix intrinsic-dimension budget.  Applying the algebraic mixed
dominated corollary gives the stated conclusion.
\end{proof}

\begin{corollary}[Random orbit closures of structured mixed menus]
\label{cor:structured-orbit-closure-mixed-dominated-operator-tail}
Fix \(r_0,b_0,c_0,L_0,C_0\ge1\).  In the finite-lift part of
Corollary~\ref{cor:structured-size-budget-mixed-dominated-operator-tail}, the
concrete structural hypothesis may be replaced by the following closure
hypothesis.  Outside a summable structural event, for every active \(h\le J_m\)
there are finitely many components \(a\in I_{m,h}\), integers
\(1\le c_{m,h,a}\le c_0\), and size parameters \(S_{m,h,a}\).  For each
component, the base structured space is one of the four finite-lift spaces listed
in Corollary~\ref{cor:structured-size-budget-mixed-dominated-operator-tail}, with
levels and block sizes bounded by \(r_0,b_0\), and with base size
\(S_{m,h,a}\).  The component tuple
\[
  C_{m,h,a,1},\ldots,C_{m,h,a,n_m}
\]
is a self-adjoint contraction tuple in the \(c_{m,h,a}\)-matrix amplification of
that base space.  The active core tuple is obtained from the component direct sum
by an \(\mathcal H_m\)-measurable self-adjoint linear image
\(\Phi_{m,h}\), with \(\|\Phi_{m,h}\|\le L_0\):
\[
  B_{m,h,i}
  =
  \Phi_{m,h}\!\left(
    \bigoplus_{a\in I_{m,h}} C_{m,h,a,i}
  \right),
  \qquad
  \|B_{m,h,i}\|_{\op}\le1 .
\]
If the active-prefix charge satisfies
\[
  \sum_{h=1}^{J_m}\sum_{a\in I_{m,h}}
  c_{m,h,a}^2 S_{m,h,a}
  \le C_0 n_m
\]
eventually, then the active-prefix signing, fixed-menu coverage, and integer or
sparse schedule thresholds of
Corollary~\ref{cor:finite-lift-core-mixed-dominated-operator-tail-menu-borel}
hold.  The constants depend only on the five displayed structural parameters.

In particular, the structured finite-lift mixed-tail result is stable under
environment-measurable unitary or permutation conjugacy, orthogonal corners and
compressions, fixed-size matrix-valued tensor amplification, block direct sums,
and bounded self-adjoint images, provided the displayed active-prefix charge is
\(O(n_m)\).

Likewise, in the algebraic part of
Corollary~\ref{cor:structured-size-budget-mixed-dominated-operator-tail}, one
may replace each active algebraic or group-convolution component by an
environment-measurable unitary conjugate, by a finite direct sum of such
components, or by a \(c\)-matrix amplification with \(c\le c_0\).  The charge is
\[
  c^2\dim_\C\mcA
  \qquad\text{or}\qquad
  c^2|G|,
\]
respectively.  If the active-prefix sum of these algebraic charges is at most
\(C_0n_m\) eventually, then the conclusions and thresholds of
Corollary~\ref{cor:algebraic-core-mixed-dominated-operator-tail-menu-borel} hold
with constants depending only on \(c_0,C_0\).
\end{corollary}

\begin{proof}
For a base finite-lift structured component, the proof of
Corollary~\ref{cor:structured-size-budget-mixed-dominated-operator-tail} gives a
finite block lift with \(\Lambda K\) and maximum block size bounded in terms of
\(r_0,b_0\), and entropy bounded by \(C(r_0,b_0)S_{m,h,a}\).
Proposition~\ref{prop:cb} shows that \(c_{m,h,a}\)-matrix amplification multiplies the
entropy by \(c_{m,h,a}^2\) and keeps the lift constants bounded in terms of
\(r_0,b_0,c_0\).  Taking the direct sum over \(a\) adds the reduced dimensions
with the maximum block size still bounded in terms of \(r_0,b_0,c_0\).  Composing
the reconstruction map with \(\Phi_{m,h}\) only multiplies the lift map norm by
at most \(L_0\).  Hence the displayed charge bound gives the active-prefix
finite-lift hypothesis with \(M_0\) and \(C_E\) depending only on
\(r_0,b_0,c_0,L_0,C_0\).  The finite-lift mixed dominated corollary applies.
Unitary and permutation conjugacy, corners, and compressions are the special
case in which \(\Phi_{m,h}\) is a norm-one self-adjoint image.

For the algebraic closure statement, unitary conjugacy preserves the
finite-dimensional C*-algebra dimension, direct sums add dimensions, and tensor
amplification by \(\M_c(\C)\) multiplies dimension by \(c^2\).  A group
convolution component has reduced C*-algebra dimension \(|G|\), so its
\(c\)-amplification has dimension \(c^2|G|\).  The active-prefix charge bound is
therefore exactly the direct-sum intrinsic-dimension hypothesis of
Corollary~\ref{cor:algebraic-core-mixed-dominated-operator-tail-menu-borel}.
\end{proof}

\begin{corollary}[Randomly selected structured submenus with mixed dominated noise]
\label{cor:selected-structured-submenu-mixed-dominated-operator-tail}
Fix structural constants \(r_0,b_0,c_0,L_0,C_0\ge1\).  For each \(m\), let
\(\mathcal H_m\) be a row environment and let \(\mathscr S_m\) be an
\(\mathcal H_m\)-measurable finite set of labels chosen from a countable
structured catalogue.  Assume that there are deterministic \(J_m\ge1\) and
events \(F_m^{\rm sel}\in\mathcal H_m\) with
\(\sum_m\mathbb P(F_m^{\rm sel})<\infty\) such that, outside
\(F_m^{\rm sel}\),
\[
  |\mathscr S_m|\le J_m .
\]
Order \(\mathscr S_m\) by a fixed countable total order and write the selected
labels as \(\xi_{m,1},\ldots,\xi_{m,N_m}\), \(N_m:=|\mathscr S_m|\).
For each selected label \(h\le N_m\), let
\[
  B_{m,h,1},\ldots,B_{m,h,n_m},
  \qquad
  Y_{m,h,1},\ldots,Y_{m,h,n_m}
\]
be the selected self-adjoint core and perturbation streams in \(\M_{D_m}(\C)\).
Assume outside \(F_m^{\rm sel}\) that all \(B_{m,h,i}\) are contractions and that
each selected core has either a finite-lift orbit-closure certificate as in
Corollary~\ref{cor:structured-orbit-closure-mixed-dominated-operator-tail}, or
an algebraic/group certificate from the same corollary.  Let \(s_{m,h}\) denote
the corresponding charge: \(c^2S\) for a finite-lift component, \(c^2\dim_\C\mcA\)
for an algebraic component, and \(c^2|G|\) for a group-convolution component,
summed over the finitely many components of the selected label.  Suppose that
\[
  \sum_{h=1}^{N_m}s_{m,h}\le C_0 n_m
\]
eventually outside \(F_m^{\rm sel}\).

For every selected perturbation stream, choose a \(\psi\) or \(\mu\) tag and the
parameters \(\varepsilon_{m,h},A_{m,h},s'_{m,h}\) of
Corollary~\ref{cor:mixed-dominated-operator-tail-active-libraries}; define
\(q_{m,h}\) there with \(Y_{m,h,i}\) in place of \(X_{m,h,i}\).  Assume that the
corresponding conditional tail or moment hypotheses and truncated-tail drift
hypotheses hold after adjoining \(\mathcal H_m\), outside events
\(E^{\rm tail}_{m,h}\), and put
\[
  \lambda_{m,h}:=\max\{1,K_{\rm tail}+\eta\}q_{m,h}.
\]
If \(\ell_m:=\log(e(D_m+n_m+1))\le B\sqrt{n_m}\) eventually,
\[
  \sum_m\left[
    \mathbb P(F_m^{\rm sel})
    +\sum_{h=1}^{J_m}
      \bigl(\mathbb P(E^{\rm tail}_{m,h})+\varepsilon_{m,h}\bigr)
  \right]<\infty,
\]
with dummy zero terms for \(h>N_m\), and
\[
  \sum_m J_m(D_m+n_m+1)
  (e(D_m+n_m+1))^{-\gamma_{\rm tc}(K_{\rm tail},1,1,B)}<\infty,
\]
then, with probability one, for all sufficiently large \(m\), one
\(\mathcal H_m\)-measurable sign vector \(\eps_m^\star\) satisfies
\[
  \left\|
  \sum_{i=1}^{n_m}\eps_{m,i}^\star
  \left(B_{m,h,i}+\frac{Y_{m,h,i}}{\lambda_{m,h}}\right)
  \right\|_{\op}
  \le C\sqrt{n_m}
  \qquad(1\le h\le N_m),
\]
where \(C\) depends only on \(r_0,b_0,c_0,L_0,C_0,K_{\rm tail},\eta\).  On the
same event all selected normalized perturbations are contractions, all selected
core-plus-tail tuples have operator norm at most \(2\), and the normalized
tuples have the corresponding common Matrix Spencer bound.  The all-integer and
sparse threshold statements are the same as in
Corollary~\ref{cor:mixed-dominated-operator-tail-active-libraries}, with
\(J_m\) as the deterministic selected-cardinality envelope.
\end{corollary}

\begin{proof}
Work outside \(F_m^{\rm sel}\).  For every selected finite-lift orbit-closure
component, the proof of
Corollary~\ref{cor:structured-orbit-closure-mixed-dominated-operator-tail}
produces a finite-dimensional certificate algebra whose complex dimension is
bounded by a constant depending only on \(r_0,b_0,c_0\) times its charge, with
lift and image norms bounded by the structural constants.  Algebraic and group
components are already contained in certificate algebras of dimensions
\(c^2\dim_\C\mcA\) and \(c^2|G|\).  Direct-sum all selected certificates.  The
total certificate dimension is at most \(C(r_0,b_0,c_0,L_0,C_0)n_m\).

After normalizing the certificate representatives by the structural lift and
image constants, Theorem~\ref{thm:akbassra} gives one
\(\mathcal H_m\)-measurable sign vector whose pushed-forward discrepancy is at
most \(K_{\rm sel}\sqrt{n_m}\) for every selected core stream, with
\(K_{\rm sel}\) depending only on the displayed structural constants.  The
lexicographically first successful sign vector is measurable because the
selected catalogue and certificates are \(\mathcal H_m\)-measurable.

Pad the selected list by zero streams up to the deterministic length \(J_m\).
Apply
Corollary~\ref{cor:signed-core-mixed-dominated-operator-tail-libraries} to this
ordered active prefix, with \(K_{\rm core}=K_{\rm sel}\) and
\(F_m^{\rm core}=F_m^{\rm sel}\).  The displayed tail summability and matrix-tail
summability are exactly the hypotheses of that corollary after the padding, and
its conclusion gives the selected-submenu estimate.  The integer and sparse
thresholds are inherited from the mixed dominated active-library selector.
\end{proof}

\begin{corollary}[Bernstein-selected structured submenus with mixed dominated noise]
\label{cor:bernstein-selected-structured-submenu-mixed-dominated-operator-tail}
Fix structural constants \(r_0,b_0,c_0,L_0,C_0\ge1\).  For each \(m\), let
\(\mathcal G_m\) be a base environment and let
\[
  1,\ldots,M_m
\]
be a finite candidate list of structured labels.  The \(k\)-th candidate has a
charge \(\mathfrak s_{m,k}\in[0,b_m]\), core streams
\(B_{m,k,1},\ldots,B_{m,k,n_m}\), and perturbation streams
\(Y_{m,k,1},\ldots,Y_{m,k,n_m}\).  Let
\(\eta_{m,1},\ldots,\eta_{m,M_m}\in\{0,1\}\) be conditionally independent given
\(\mathcal G_m\), with
\[
  \mathbb P\{\eta_{m,k}=1\mid\mathcal G_m\}\le p_{m,k}
  \qquad(1\le k\le M_m),
\]
where the \(p_{m,k}\in[0,1]\) are deterministic.  Put
\[
  \mathscr S_m:=\{k:\eta_{m,k}=1\},
  \qquad
  \mathcal H_m:=\mathcal G_m\vee\sigma(\eta_{m,1},\ldots,\eta_{m,M_m}),
\]
and define
\[
  \mu_m^{\rm ch}:=\sum_{k=1}^{M_m}p_{m,k}\mathfrak s_{m,k},\qquad
  v_m^{\rm ch}:=\sum_{k=1}^{M_m}p_{m,k}\mathfrak s_{m,k}^2,\qquad
  \mu_m^{\#}:=\sum_{k=1}^{M_m}p_{m,k}.
\]

Assume that outside events \(F_m^{\rm cert}\in\mathcal G_m\), with
\(\sum_m\mathbb P(F_m^{\rm cert})<\infty\), every selected core is a contraction
and has one of the finite-lift orbit-closure or algebraic/group certificates from
Corollary~\ref{cor:selected-structured-submenu-mixed-dominated-operator-tail},
with charge bounded by the displayed \(\mathfrak s_{m,k}\).  Choose
\(t_m,u_m>0\) and a deterministic cardinality envelope \(J_m\ge1\) such that
\[
  \mu_m^{\rm ch}+t_m\le C_0n_m,\qquad
  \mu_m^{\#}+u_m\le J_m .
\]
If
\[
  \sum_m\left[
    \mathbb P(F_m^{\rm cert})
    +\exp\left\{
      -\frac{t_m^2}{2(v_m^{\rm ch}+b_mt_m/3)}
    \right\}
    +\exp\left\{
      -\frac{u_m^2}{2(\mu_m^{\#}+u_m/3)}
    \right\}
  \right]<\infty,
\]
then the selection event in
Corollary~\ref{cor:selected-structured-submenu-mixed-dominated-operator-tail} is
summable: almost surely, eventually
\[
  |\mathscr S_m|\le J_m,
  \qquad
  \sum_{k\in\mathscr S_m}\mathfrak s_{m,k}\le C_0n_m .
\]

Consequently, suppose each candidate perturbation stream is tagged \(\psi\) or
\(\mu\), with parameters \(\varepsilon_{m,k},A_{m,k},s'_{m,k}\), and satisfies
the corresponding mixed dominated tail and drift hypotheses after adjoining
\(\mathcal H_m\).  If
\[
  \ell_m:=\log(e(D_m+n_m+1))\le B\sqrt{n_m}
\]
eventually,
\[
  \sum_m\sum_{k=1}^{M_m}
  \bigl(\mathbb P(E^{\rm tail}_{m,k})+\varepsilon_{m,k}\bigr)<\infty,
\]
and
\[
  \sum_mJ_m(D_m+n_m+1)
  (e(D_m+n_m+1))^{-\gamma_{\rm tc}(K_{\rm tail},1,1,B)}<\infty,
\]
then the selected structured submenus have the same eventual common signing,
contraction normalization, fixed-menu coverage, and all-integer/sparse thresholds
as in
Corollary~\ref{cor:selected-structured-submenu-mixed-dominated-operator-tail}.
\end{corollary}

\begin{proof}
Condition on \(\mathcal G_m\).  For the selected charge,
\[
  Z_k:=\eta_{m,k}\mathfrak s_{m,k}
\]
are independent, satisfy \(0\le Z_k\le b_m\), and have conditional means at most
\(p_{m,k}\mathfrak s_{m,k}\).  Their conditional variance sum is at most
\(v_m^{\rm ch}\).  The scalar Bernstein inequality gives
\[
  \mathbb P\left\{
    \sum_k\eta_{m,k}\mathfrak s_{m,k}>C_0n_m
    \,\middle|\,\mathcal G_m
  \right\}
  \le
  \exp\left\{
    -\frac{t_m^2}{2(v_m^{\rm ch}+b_mt_m/3)}
  \right\}.
\]
The same argument with \(Z_k=\eta_{m,k}\), jump \(1\), and variance sum at most
\(\mu_m^{\#}\) gives
\[
  \mathbb P\left\{
    |\mathscr S_m|>J_m
    \,\middle|\,\mathcal G_m
  \right\}
  \le
  \exp\left\{
    -\frac{u_m^2}{2(\mu_m^{\#}+u_m/3)}
  \right\}.
\]
After averaging and adding \(F_m^{\rm cert}\), the displayed summability gives a
summable selection failure event.  On its complement, the selected set satisfies
the cardinality and charge hypotheses of
Corollary~\ref{cor:selected-structured-submenu-mixed-dominated-operator-tail}.
Ordering the selected labels and padding by zero streams converts the candidate
tail summability into the active-prefix tail summability required there.  The
final conclusions are exactly that corollary's conclusions.
\end{proof}

\begin{corollary}[Laplace-dominated selected structured submenus with mixed dominated noise]
\label{cor:laplace-dominated-selected-structured-submenu-mixed-dominated-operator-tail}
Fix structural constants \(r_0,b_0,c_0,L_0,C_0\ge1\).  For each \(m\), let
\(\mathcal G_m\) be a base environment and let \(1,\ldots,M_m\) be a finite
candidate list of structured labels.  The \(k\)-th candidate has charge
\(\mathfrak s_{m,k}\in[0,b_m]\), with \(b_m>0\), core streams
\(B_{m,k,1},\ldots,B_{m,k,n_m}\), and perturbation streams
\(Y_{m,k,1},\ldots,Y_{m,k,n_m}\).  Let
\(\eta_{m,1},\ldots,\eta_{m,M_m}\in\{0,1\}\) be
\(\mathcal H_m\)-measurable selection indicators, where
\[
  \mathcal H_m:=\mathcal G_m\vee\sigma(\eta_{m,1},\ldots,\eta_{m,M_m}),
  \qquad
  \mathscr S_m:=\{k:\eta_{m,k}=1\}.
\]
Assume there are deterministic \(p_{m,k}\in[0,1]\) such that, conditionally on
\(\mathcal G_m\), the selected cardinality and charge satisfy the Laplace
domination bounds
\[
  \mathbb E\left[
    \exp\left\{\theta\sum_{k=1}^{M_m}\eta_{m,k}\mathfrak s_{m,k}\right\}
    \,\middle|\,\mathcal G_m
  \right]
  \le
  \exp\left\{\sum_{k=1}^{M_m}p_{m,k}
    (e^{\theta\mathfrak s_{m,k}}-1)\right\}
\]
and
\[
  \mathbb E\left[
    \exp\left\{\theta\sum_{k=1}^{M_m}\eta_{m,k}\right\}
    \,\middle|\,\mathcal G_m
  \right]
  \le
  \exp\left\{\sum_{k=1}^{M_m}p_{m,k}(e^\theta-1)\right\}
\]
for every \(\theta\ge0\).  Put
\[
  \mu_m^{\rm ch}:=\sum_{k=1}^{M_m}p_{m,k}\mathfrak s_{m,k},\qquad
  v_m^{\rm ch}:=\sum_{k=1}^{M_m}p_{m,k}\mathfrak s_{m,k}^2,\qquad
  \mu_m^{\#}:=\sum_{k=1}^{M_m}p_{m,k}.
\]

Assume that outside events \(F_m^{\rm cert}\in\mathcal G_m\), with
\(\sum_m\mathbb P(F_m^{\rm cert})<\infty\), every selected core is a contraction
and has one of the finite-lift orbit-closure or algebraic/group certificates
from
Corollary~\ref{cor:selected-structured-submenu-mixed-dominated-operator-tail},
with charge bounded by the displayed \(\mathfrak s_{m,k}\).  Choose
\(t_m,u_m>0\) and a deterministic cardinality envelope \(J_m\ge1\) such that
\[
  \mu_m^{\rm ch}+t_m\le C_0n_m,\qquad
  \mu_m^{\#}+u_m\le J_m .
\]
If
\[
  \sum_m\left[
    \mathbb P(F_m^{\rm cert})
    +\exp\left\{
      -\frac{t_m^2}{2(v_m^{\rm ch}+b_mt_m/3)}
    \right\}
    +\exp\left\{
      -\frac{u_m^2}{2(\mu_m^{\#}+u_m/3)}
    \right\}
  \right]<\infty,
\]
then the selection event in
Corollary~\ref{cor:selected-structured-submenu-mixed-dominated-operator-tail} is
summable: almost surely, eventually
\[
  |\mathscr S_m|\le J_m,
  \qquad
  \sum_{k\in\mathscr S_m}\mathfrak s_{m,k}\le C_0n_m .
\]

Consequently, suppose each candidate perturbation stream is tagged \(\psi\) or
\(\mu\), with parameters \(\varepsilon_{m,k},A_{m,k},s'_{m,k}\), and satisfies
the corresponding mixed dominated tail and drift hypotheses after adjoining
\(\mathcal H_m\).  If
\[
  \ell_m:=\log(e(D_m+n_m+1))\le B\sqrt{n_m}
\]
eventually,
\[
  \sum_m\sum_{k=1}^{M_m}
  \bigl(\mathbb P(E^{\rm tail}_{m,k})+\varepsilon_{m,k}\bigr)<\infty,
\]
and
\[
  \sum_mJ_m(D_m+n_m+1)
  (e(D_m+n_m+1))^{-\gamma_{\rm tc}(K_{\rm tail},1,1,B)}<\infty,
\]
then the Laplace-dominated selected structured submenus have the same eventual
common signing, contraction normalization, fixed-menu coverage, and
all-integer/sparse thresholds as in
Corollary~\ref{cor:selected-structured-submenu-mixed-dominated-operator-tail}.

The displayed Laplace domination is automatic, for example, whenever
conditionally on \(\mathcal G_m\) the indicators are negatively associated and
\(\mathbb E[\eta_{m,k}\mid\mathcal G_m]\le p_{m,k}\) for every \(k\).
Here negative association is meant conditionally on \(\mathcal G_m\), in the
usual sense that increasing functions of disjoint coordinate sets have
nonpositive conditional covariance.
\end{corollary}

\begin{proof}
Condition on \(\mathcal G_m\).  The first Laplace domination hypothesis and the
standard Bennett bound
\[
  e^{\theta x}-1-\theta x
  \le \frac{x^2}{b_m^2}(e^{\theta b_m}-1-\theta b_m)
  \qquad(0\le x\le b_m)
\]
give, for \(S_m^{\rm ch}:=\sum_k\eta_{m,k}\mathfrak s_{m,k}\),
\[
  \log\mathbb E\left[
    e^{\theta(S_m^{\rm ch}-\mu_m^{\rm ch})}
    \,\middle|\,\mathcal G_m
  \right]
  \le
  \frac{v_m^{\rm ch}}{b_m^2}(e^{\theta b_m}-1-\theta b_m).
\]
Chernoff optimization gives
\[
  \mathbb P\{S_m^{\rm ch}>\mu_m^{\rm ch}+t_m\mid\mathcal G_m\}
  \le
  \exp\left\{-\frac{t_m^2}{2(v_m^{\rm ch}+b_mt_m/3)}\right\}.
\]
The second Laplace domination hypothesis gives the same estimate for
\(|\mathscr S_m|\), with \(b_m=1\), \(v_m^{\rm ch}\) replaced by \(\mu_m^{\#}\),
and \(t_m\) replaced by \(u_m\):
\[
  \mathbb P\{|\mathscr S_m|>\mu_m^{\#}+u_m\mid\mathcal G_m\}
  \le
  \exp\left\{-\frac{u_m^2}{2(\mu_m^{\#}+u_m/3)}\right\}.
\]
After averaging, adding \(F_m^{\rm cert}\), and applying Borel--Cantelli, the
selected set eventually satisfies the deterministic cardinality and charge
hypotheses of
Corollary~\ref{cor:selected-structured-submenu-mixed-dominated-operator-tail}.
Ordering the selected labels and padding by zero streams converts the candidate
tail summability into the active-prefix tail summability required there, and the
deterministic matrix-tail series is unchanged.

For the final sufficient condition, fix \(\theta\ge0\) and nonnegative weights
\(a_k\), and put
\[
  L_m(\theta,a)
  :=
  \mathbb E\left[
    \exp\left\{\theta\sum_k a_k\eta_{m,k}\right\}
    \,\middle|\,\mathcal G_m
  \right].
\]
Negative association gives
\[
\begin{aligned}
  L_m(\theta,a)
  &\le
  \prod_k\mathbb E[e^{\theta a_k\eta_{m,k}}\mid\mathcal G_m]  \\
  &\le
  \prod_k\left(1+p_{m,k}(e^{\theta a_k}-1)\right)  \\
  &\le
  \exp\left\{\sum_kp_{m,k}(e^{\theta a_k}-1)\right\}.
\end{aligned}
\]
Taking \(a_k=\mathfrak s_{m,k}\) and \(a_k=1\) gives the two Laplace domination
bounds.
\end{proof}

\begin{corollary}[Random-envelope Laplace-dominated selected structured submenus]
\label{cor:random-envelope-laplace-selected-structured-submenu-mixed-dominated-operator-tail}
Fix structural constants \(r_0,b_0,c_0,L_0,C_0\ge1\).  For each \(m\), let
\(\mathcal G_m\) be a base environment and let \(1,\ldots,M_m\) be a finite
candidate list of structured labels.  The \(k\)-th candidate has charge
\(\mathfrak s_{m,k}\in[0,b_m]\), with \(b_m>0\), core streams
\(B_{m,k,1},\ldots,B_{m,k,n_m}\), and perturbation streams
\(Y_{m,k,1},\ldots,Y_{m,k,n_m}\).  Let
\(\eta_{m,1},\ldots,\eta_{m,M_m}\in\{0,1\}\) be selection indicators, set
\[
  \mathcal H_m:=\mathcal G_m\vee\sigma(\eta_{m,1},\ldots,\eta_{m,M_m}),
  \qquad
  \mathscr S_m:=\{k:\eta_{m,k}=1\},
\]
and suppose there are \(\mathcal G_m\)-measurable envelopes
\(P_{m,k}\in[0,1]\) such that, conditionally on \(\mathcal G_m\), for every
\(\theta\ge0\),
\[
  \mathbb E\left[
    \exp\left\{\theta\sum_{k=1}^{M_m}\eta_{m,k}\mathfrak s_{m,k}\right\}
    \,\middle|\,\mathcal G_m
  \right]
  \le
  \exp\left\{\sum_{k=1}^{M_m}P_{m,k}
    (e^{\theta\mathfrak s_{m,k}}-1)\right\}
\]
and
\[
  \mathbb E\left[
    \exp\left\{\theta\sum_{k=1}^{M_m}\eta_{m,k}\right\}
    \,\middle|\,\mathcal G_m
  \right]
  \le
  \exp\left\{\sum_{k=1}^{M_m}P_{m,k}(e^\theta-1)\right\}.
\]
Put
\[
  M_m^{\rm ch}:=\sum_{k=1}^{M_m}P_{m,k}\mathfrak s_{m,k},\qquad
  V_m^{\rm ch}:=\sum_{k=1}^{M_m}P_{m,k}\mathfrak s_{m,k}^2,\qquad
  M_m^{\#}:=\sum_{k=1}^{M_m}P_{m,k}.
\]

Assume that there are events \(F_m^{\rm env}\in\mathcal G_m\), with
\(\sum_m\mathbb P(F_m^{\rm env})<\infty\), and deterministic bounds
\(\bar\mu_m^{\rm ch},\bar v_m^{\rm ch},\bar\mu_m^{\#}\ge0\) such that outside
\(F_m^{\rm env}\),
\[
  M_m^{\rm ch}\le\bar\mu_m^{\rm ch},\qquad
  V_m^{\rm ch}\le\bar v_m^{\rm ch},\qquad
  M_m^{\#}\le\bar\mu_m^{\#}.
\]
Assume also that outside events \(F_m^{\rm cert}\in\mathcal G_m\), with
\(\sum_m\mathbb P(F_m^{\rm cert})<\infty\), every selected core is a contraction
and has one of the finite-lift orbit-closure or algebraic/group certificates
from
Corollary~\ref{cor:selected-structured-submenu-mixed-dominated-operator-tail},
with charge bounded by the displayed \(\mathfrak s_{m,k}\).  Choose
\(t_m,u_m>0\) and a deterministic cardinality envelope \(J_m\ge1\) such that
\[
  \bar\mu_m^{\rm ch}+t_m\le C_0n_m,\qquad
  \bar\mu_m^{\#}+u_m\le J_m .
\]
Put
\[
\begin{aligned}
  \Lambda_m:={}&
    \mathbb P(F_m^{\rm env})+\mathbb P(F_m^{\rm cert}) \\
  &+\exp\left\{
      -\frac{t_m^2}{2(\bar v_m^{\rm ch}+b_mt_m/3)}
    \right\}
    +\exp\left\{
      -\frac{u_m^2}{2(\bar\mu_m^{\#}+u_m/3)}
    \right\}.
\end{aligned}
\]
If
\[
  \sum_m\Lambda_m<\infty,
\]
then the selection event in
Corollary~\ref{cor:selected-structured-submenu-mixed-dominated-operator-tail} is
summable: almost surely, eventually
\[
  |\mathscr S_m|\le J_m,
  \qquad
  \sum_{k\in\mathscr S_m}\mathfrak s_{m,k}\le C_0n_m .
\]

Consequently, suppose each candidate perturbation stream is tagged \(\psi\) or
\(\mu\), with parameters \(\varepsilon_{m,k},A_{m,k},s'_{m,k}\), and satisfies
the corresponding mixed dominated tail and drift hypotheses after adjoining
\(\mathcal H_m\).  If
\[
  \ell_m:=\log(e(D_m+n_m+1))\le B\sqrt{n_m}
\]
eventually,
\[
  \sum_m\sum_{k=1}^{M_m}
  \bigl(\mathbb P(E^{\rm tail}_{m,k})+\varepsilon_{m,k}\bigr)<\infty,
\]
and
\[
  \sum_mJ_m(D_m+n_m+1)
  (e(D_m+n_m+1))^{-\gamma_{\rm tc}(K_{\rm tail},1,1,B)}<\infty,
\]
then the random-envelope Laplace-dominated selected structured submenus have the
same eventual common signing, contraction normalization, fixed-menu coverage,
and all-integer/sparse thresholds as in
Corollary~\ref{cor:selected-structured-submenu-mixed-dominated-operator-tail}.

The hypotheses hold, for example, under conditional negative association with
\(\mathbb E[\eta_{m,k}\mid\mathcal G_m]\le P_{m,k}\), followed by any summable
high-probability control of the three aggregate profiles.  Thus the deterministic
envelopes in
Corollary~\ref{cor:laplace-dominated-selected-structured-submenu-mixed-dominated-operator-tail}
may be replaced by data-dependent or environment-random probabilities whenever
their total expected charge, quadratic charge, and cardinality are certified.
\end{corollary}

\begin{proof}
Work on \((F_m^{\rm env})^c\) and condition on \(\mathcal G_m\).  For
\(S_m^{\rm ch}:=\sum_k\eta_{m,k}\mathfrak s_{m,k}\), the Laplace domination
and the Bennett bound
\[
  e^{\theta x}-1-\theta x
  \le \frac{x^2}{b_m^2}(e^{\theta b_m}-1-\theta b_m)
  \qquad(0\le x\le b_m)
\]
give
\[
  \log\mathbb E\left[
    e^{\theta(S_m^{\rm ch}-M_m^{\rm ch})}
    \,\middle|\,\mathcal G_m
  \right]
  \le
  \frac{V_m^{\rm ch}}{b_m^2}(e^{\theta b_m}-1-\theta b_m).
\]
Since \(M_m^{\rm ch}\le\bar\mu_m^{\rm ch}\) and
\(V_m^{\rm ch}\le\bar v_m^{\rm ch}\) on the present event, Chernoff
optimization gives
\[
  \mathbb P\{S_m^{\rm ch}>\bar\mu_m^{\rm ch}+t_m\mid\mathcal G_m\}
  \le
  \exp\left\{-\frac{t_m^2}{2(\bar v_m^{\rm ch}+b_mt_m/3)}\right\}.
\]
The second Laplace domination bound, with \(b_m=1\), gives
\[
  \mathbb P\{|\mathscr S_m|>\bar\mu_m^{\#}+u_m\mid\mathcal G_m\}
  \le
  \exp\left\{-\frac{u_m^2}{2(\bar\mu_m^{\#}+u_m/3)}\right\}.
\]
After averaging, adding \(F_m^{\rm env}\cup F_m^{\rm cert}\), and applying
Borel--Cantelli, the selected catalogue eventually satisfies the deterministic
cardinality and charge hypotheses of
Corollary~\ref{cor:selected-structured-submenu-mixed-dominated-operator-tail}.
Ordering the selected labels and padding by zero streams converts the candidate
tail summability into the active-prefix tail summability required there.

The negative-association sufficient condition is the same product-mgf argument
as in
Corollary~\ref{cor:laplace-dominated-selected-structured-submenu-mixed-dominated-operator-tail},
with \(P_{m,k}\) in place of \(p_{m,k}\).
\end{proof}

\begin{corollary}[Random-envelope Poisson-count selected structured submenus]
\label{cor:random-envelope-poisson-count-selected-structured-submenu-mixed-dominated-operator-tail}
Fix structural constants \(r_0,b_0,c_0,L_0,C_0\ge1\).  For each \(m\), let
\(\mathcal G_m\) be a base environment and let \(1,\ldots,M_m\) be a finite
candidate list of structured labels.  The \(k\)-th candidate has charge
\(\mathfrak s_{m,k}\in[0,b_m]\), with \(b_m>0\), core streams
\(B_{m,k,1},\ldots,B_{m,k,n_m}\), and perturbation streams
\(Y_{m,k,1},\ldots,Y_{m,k,n_m}\).  Let
\(N_{m,k}\in\mathbb Z_{\ge0}\) be random proposal counts and select the support
of the proposal field,
\[
  \eta_{m,k}:={\bf 1}_{\{N_{m,k}\ge1\}},\qquad
  \mathscr S_m:=\{k:\eta_{m,k}=1\},
\]
with
\[
  \mathcal H_m:=\mathcal G_m\vee\sigma(N_{m,1},\ldots,N_{m,M_m}).
\]
Assume that there are \(\mathcal G_m\)-measurable intensity envelopes
\(\Lambda_{m,k}\ge0\) such that, conditionally on \(\mathcal G_m\), for every
\(\theta\ge0\),
\[
  \mathbb E\left[
    \exp\left\{\theta\sum_{k=1}^{M_m}N_{m,k}\mathfrak s_{m,k}\right\}
    \,\middle|\,\mathcal G_m
  \right]
  \le
  \exp\left\{\sum_{k=1}^{M_m}\Lambda_{m,k}
    (e^{\theta\mathfrak s_{m,k}}-1)\right\}
\]
and
\[
  \mathbb E\left[
    \exp\left\{\theta\sum_{k=1}^{M_m}N_{m,k}\right\}
    \,\middle|\,\mathcal G_m
  \right]
  \le
  \exp\left\{\sum_{k=1}^{M_m}\Lambda_{m,k}(e^\theta-1)\right\}.
\]
Define the random count profiles
\[
\begin{aligned}
  M_m^{\rm cnt,ch}&:=\sum_{k=1}^{M_m}\Lambda_{m,k}\mathfrak s_{m,k},\\
  V_m^{\rm cnt,ch}&:=\sum_{k=1}^{M_m}\Lambda_{m,k}\mathfrak s_{m,k}^2,\\
  M_m^{\rm cnt,\#}&:=\sum_{k=1}^{M_m}\Lambda_{m,k}.
\end{aligned}
\]

Assume that there are events \(F_m^{\rm env}\in\mathcal G_m\), with
\(\sum_m\mathbb P(F_m^{\rm env})<\infty\), and deterministic bounds
\(\bar\mu_m^{\rm cnt,ch},\bar v_m^{\rm cnt,ch},
\bar\mu_m^{\rm cnt,\#}\ge0\) such that outside \(F_m^{\rm env}\),
\[
  M_m^{\rm cnt,ch}\le\bar\mu_m^{\rm cnt,ch},\qquad
  V_m^{\rm cnt,ch}\le\bar v_m^{\rm cnt,ch},\qquad
  M_m^{\rm cnt,\#}\le\bar\mu_m^{\rm cnt,\#}.
\]
Assume also that outside events \(F_m^{\rm cert}\in\mathcal G_m\), with
\(\sum_m\mathbb P(F_m^{\rm cert})<\infty\), every selected core is a contraction
and has one of the finite-lift orbit-closure or algebraic/group certificates
from
Corollary~\ref{cor:selected-structured-submenu-mixed-dominated-operator-tail},
with charge bounded by the displayed \(\mathfrak s_{m,k}\).  Choose
\(t_m,u_m>0\) and a deterministic cardinality envelope \(J_m\ge1\) such that
\[
  \bar\mu_m^{\rm cnt,ch}+t_m\le C_0n_m,\qquad
  \bar\mu_m^{\rm cnt,\#}+u_m\le J_m .
\]
Put
\[
\begin{aligned}
  \Lambda_m^{\rm cnt}:={}&
    \mathbb P(F_m^{\rm env})+\mathbb P(F_m^{\rm cert}) \\
  &+\exp\left\{
      -\frac{t_m^2}{2(\bar v_m^{\rm cnt,ch}+b_mt_m/3)}
    \right\}
    &+\exp\left\{
      -\frac{u_m^2}{2(\bar\mu_m^{\rm cnt,\#}+u_m/3)}
    \right\}.
\end{aligned}
\]
If
\[
  \sum_m\Lambda_m^{\rm cnt}<\infty,
\]
then the selection event in
Corollary~\ref{cor:selected-structured-submenu-mixed-dominated-operator-tail} is
summable: almost surely, eventually
\[
  |\mathscr S_m|\le J_m,
  \qquad
  \sum_{k\in\mathscr S_m}\mathfrak s_{m,k}\le C_0n_m .
\]

Consequently, suppose each candidate perturbation stream is tagged \(\psi\) or
\(\mu\), with parameters \(\varepsilon_{m,k},A_{m,k},s'_{m,k}\), and satisfies
the corresponding mixed dominated tail and drift hypotheses after adjoining
\(\mathcal H_m\).  If
\[
  \ell_m:=\log(e(D_m+n_m+1))\le B\sqrt{n_m}
\]
eventually,
\[
  \sum_m\sum_{k=1}^{M_m}
  \bigl(\mathbb P(E^{\rm tail}_{m,k})+\varepsilon_{m,k}\bigr)<\infty,
\]
and
\[
  \sum_mJ_m(D_m+n_m+1)
  (e(D_m+n_m+1))^{-\gamma_{\rm tc}(K_{\rm tail},1,1,B)}<\infty,
\]
then the Poisson-count selected structured submenus have the same eventual
common signing, contraction normalization, fixed-menu coverage, and
all-integer/sparse thresholds as in
Corollary~\ref{cor:selected-structured-submenu-mixed-dominated-operator-tail}.

This applies to independent Poisson counts, Cox or environment-random Poisson
counts after summable certification of their aggregate intensity profiles, and
more generally to any proposal-count field whose cardinality and charge
transforms are dominated by the displayed compound-Poisson envelopes.  It
covers Poissonized model proposals, occupancy libraries, random hyperparameter
proposal clouds, and repeated-draw structured catalogues after passing from
counts to their occupied support.
\end{corollary}

\begin{proof}
Since \(N_{m,k}\) are nonnegative integers,
\[
  \eta_{m,k}\le N_{m,k},\qquad
  |\mathscr S_m|\le \sum_kN_{m,k},\qquad
  \sum_{k\in\mathscr S_m}\mathfrak s_{m,k}
  \le \sum_kN_{m,k}\mathfrak s_{m,k}.
\]
Work on \((F_m^{\rm env})^c\) and condition on \(\mathcal G_m\).  For
\[
  T_m^{\rm ch}:=\sum_kN_{m,k}\mathfrak s_{m,k},
\]
the first compound-Poisson domination hypothesis and the Bennett bound
\[
  e^{\theta x}-1-\theta x
  \le \frac{x^2}{b_m^2}(e^{\theta b_m}-1-\theta b_m)
  \qquad(0\le x\le b_m)
\]
give
\[
  \log\mathbb E\left[
    e^{\theta(T_m^{\rm ch}-M_m^{\rm cnt,ch})}
    \,\middle|\,\mathcal G_m
  \right]
  \le
  \frac{V_m^{\rm cnt,ch}}{b_m^2}
  (e^{\theta b_m}-1-\theta b_m).
\]
Using the deterministic bounds on the profiles and optimizing in \(\theta\),
\[
  \mathbb P\left\{
    (F_m^{\rm env})^c,\ T_m^{\rm ch}>
    \bar\mu_m^{\rm cnt,ch}+t_m
  \right\}
  \le
  \exp\left\{
    -\frac{t_m^2}{2(\bar v_m^{\rm cnt,ch}+b_mt_m/3)}
  \right\}.
\]
The second compound-Poisson domination bound, with atom size \(1\), gives
\[
  \mathbb P\left\{
    (F_m^{\rm env})^c,\ \sum_kN_{m,k}>
    \bar\mu_m^{\rm cnt,\#}+u_m
  \right\}
  \le
  \exp\left\{
    -\frac{u_m^2}{2(\bar\mu_m^{\rm cnt,\#}+u_m/3)}
  \right\}.
\]
Adding \(F_m^{\rm env}\cup F_m^{\rm cert}\) and summing over \(m\), the
Borel--Cantelli lemma gives the eventual selected-support cardinality and
charge bounds displayed above.

On the resulting good event, order the occupied labels by candidate order and
pad by zero streams up to \(J_m\).  The candidate tail summability dominates the
selected active-prefix tail summability required by
Corollary~\ref{cor:selected-structured-submenu-mixed-dominated-operator-tail},
and the deterministic matrix-tail series is unchanged.  Applying that corollary
gives the signing, normalization, coverage, and threshold conclusions.  The
listed Poisson and Cox examples have equality in the conditional transform
display after conditioning on their intensity field; the stated domination
variant is the same argument with the inequality in place of equality.
\end{proof}

\begin{corollary}[Cox-Poisson occupied-support selected structured submenus]
\label{cor:cox-poisson-occupied-support-selected-structured-submenu-mixed-dominated-operator-tail}
Fix structural constants \(r_0,b_0,c_0,L_0,C_0\ge1\).  In the setting of
Corollary~\ref{cor:random-envelope-poisson-count-selected-structured-submenu-mixed-dominated-operator-tail},
assume more specifically that, conditionally on \(\mathcal G_m\), the proposal
counts \(N_{m,1},\ldots,N_{m,M_m}\) are independent Poisson variables with
\(\mathcal G_m\)-measurable intensities \(\Lambda_{m,k}\ge0\).  Select the
occupied support
\[
  \eta_{m,k}:={\bf 1}_{\{N_{m,k}\ge1\}},
  \qquad
  \mathscr S_m:=\{k:\eta_{m,k}=1\},
\]
and put
\[
  Q_{m,k}:=1-e^{-\Lambda_{m,k}} .
\]
Define the occupied-support profiles
\[
\begin{aligned}
  M_m^{\rm occ,ch}&:=\sum_{k=1}^{M_m}Q_{m,k}\mathfrak s_{m,k},\\
  V_m^{\rm occ,ch}&:=\sum_{k=1}^{M_m}Q_{m,k}\mathfrak s_{m,k}^2,\\
  M_m^{\rm occ,\#}&:=\sum_{k=1}^{M_m}Q_{m,k}.
\end{aligned}
\]

Assume that there are events \(F_m^{\rm env}\in\mathcal G_m\), with
\(\sum_m\mathbb P(F_m^{\rm env})<\infty\), and deterministic bounds
\(\bar\mu_m^{\rm occ,ch},\bar v_m^{\rm occ,ch},
\bar\mu_m^{\rm occ,\#}\ge0\) such that outside \(F_m^{\rm env}\),
\[
  M_m^{\rm occ,ch}\le\bar\mu_m^{\rm occ,ch},\qquad
  V_m^{\rm occ,ch}\le\bar v_m^{\rm occ,ch},\qquad
  M_m^{\rm occ,\#}\le\bar\mu_m^{\rm occ,\#}.
\]
Assume the same candidate certificate condition as in
Corollary~\ref{cor:random-envelope-poisson-count-selected-structured-submenu-mixed-dominated-operator-tail},
outside summable events \(F_m^{\rm cert}\in\mathcal G_m\).  Choose
\(t_m,u_m>0\) and a deterministic cardinality envelope \(J_m\ge1\) so that
\[
  \bar\mu_m^{\rm occ,ch}+t_m\le C_0n_m,\qquad
  \bar\mu_m^{\rm occ,\#}+u_m\le J_m .
\]
If
\[
  \begin{aligned}
  \sum_m\Bigg[
    &\mathbb P(F_m^{\rm env})+\mathbb P(F_m^{\rm cert})
    \\
    &+\exp\left\{
      -\frac{t_m^2}{2(\bar v_m^{\rm occ,ch}+b_mt_m/3)}
    \right\}
    \\
    &+\exp\left\{
      -\frac{u_m^2}{2(\bar\mu_m^{\rm occ,\#}+u_m/3)}
    \right\}
  \Bigg]<\infty,
  \end{aligned}
\]
then the selection event in
Corollary~\ref{cor:selected-structured-submenu-mixed-dominated-operator-tail} is
summable: almost surely, eventually
\[
  |\mathscr S_m|\le J_m,
  \qquad
  \sum_{k\in\mathscr S_m}\mathfrak s_{m,k}\le C_0n_m .
\]

Consequently, under the same candidate perturbation, mixed dominated tail,
drift, cap-budget, and matrix-tail summability assumptions as in
Corollary~\ref{cor:random-envelope-poisson-count-selected-structured-submenu-mixed-dominated-operator-tail},
with this deterministic envelope \(J_m\), the Cox-Poisson occupied-support
submenus have the same eventual common signing, contraction normalization,
fixed-menu coverage, and all-integer/sparse thresholds as in
Corollary~\ref{cor:selected-structured-submenu-mixed-dominated-operator-tail}.

Thus Poissonized repeated-draw catalogues may be charged by occupied labels
rather than by total proposal multiplicity.  When some intensities are large,
the budget uses \(1-e^{-\Lambda_{m,k}}\le1\), improving the count profile
\(\Lambda_{m,k}\) from
Corollary~\ref{cor:random-envelope-poisson-count-selected-structured-submenu-mixed-dominated-operator-tail}.
\end{corollary}

\begin{proof}
Condition on \(\mathcal G_m\).  The occupied indicators are independent
Bernoulli variables and
\[
  \mathbb P\{\eta_{m,k}=1\mid\mathcal G_m\}
  =
  1-e^{-\Lambda_{m,k}}
  =
  Q_{m,k}.
\]
Therefore, for every \(\theta\ge0\),
\[
\begin{aligned}
  &\mathbb E\left[
    \exp\left\{\theta\sum_k\eta_{m,k}\mathfrak s_{m,k}\right\}
    \,\middle|\,\mathcal G_m
  \right]  \\
  &\qquad=
  \prod_k\left(1+Q_{m,k}(e^{\theta\mathfrak s_{m,k}}-1)\right)
  \le
  \exp\left\{\sum_kQ_{m,k}(e^{\theta\mathfrak s_{m,k}}-1)\right\},
\end{aligned}
\]
and the same calculation with \(\mathfrak s_{m,k}=1\) gives the cardinality
Laplace domination.  Hence
Corollary~\ref{cor:random-envelope-laplace-selected-structured-submenu-mixed-dominated-operator-tail}
applies with \(P_{m,k}=Q_{m,k}\).  Its conclusion is exactly the displayed
selection event and the inherited mixed-tail Matrix Spencer conclusion.
\end{proof}

\begin{corollary}[Product-categorical selected structured submenus with mixed dominated noise]
\label{cor:product-categorical-selected-structured-submenu-mixed-dominated-operator-tail}
Fix structural constants \(r_0,b_0,c_0,L_0,C_0\ge1\).  For each \(m\), let
\(\mathcal G_m\) be a base environment and let the finite candidate list be
partitioned into deterministic slots
\[
  \mathcal I_{m,1},\ldots,\mathcal I_{m,L_m}.
\]
The \(k\)-th candidate has charge \(\mathfrak s_{m,k}\in[0,b_m]\), with
\(b_m>0\), core streams \(B_{m,k,1},\ldots,B_{m,k,n_m}\), and perturbation
streams \(Y_{m,k,1},\ldots,Y_{m,k,n_m}\).  In slot \(g\), choose at most one
label by an \(\mathcal I_{m,g}\cup\{\varnothing\}\)-valued random variable
\(\xi_{m,g}\).  Conditionally on \(\mathcal G_m\), assume that
\(\xi_{m,1},\ldots,\xi_{m,L_m}\) are independent and that deterministic
probability envelopes \(\pi_{m,k}\in[0,1]\) satisfy
\[
  \mathbb P\{\xi_{m,g}=k\mid\mathcal G_m\}\le \pi_{m,k}
  \qquad(k\in\mathcal I_{m,g}).
\]
Let
\[
  \eta_{m,k}:=\mathbf 1_{\{\xi_{m,g}=k\}}\quad(k\in\mathcal I_{m,g}),\qquad
  \mathscr S_m:=\{k:\eta_{m,k}=1\},
\]
and put
\[
  \mathcal H_m:=
  \mathcal G_m\vee\sigma(\xi_{m,1},\ldots,\xi_{m,L_m}).
\]
Define
\[
  \mu_m^{\rm ch}:=\sum_k\pi_{m,k}\mathfrak s_{m,k},
  \qquad
  v_m^{\rm ch}:=\sum_k\pi_{m,k}\mathfrak s_{m,k}^2 .
\]

Assume that outside events \(F_m^{\rm cert}\in\mathcal G_m\), with
\(\sum_m\mathbb P(F_m^{\rm cert})<\infty\), every selected core is a contraction
and has one of the finite-lift orbit-closure or algebraic/group certificates
from
Corollary~\ref{cor:selected-structured-submenu-mixed-dominated-operator-tail},
with charge bounded by the displayed \(\mathfrak s_{m,k}\).  Choose \(t_m>0\)
and a deterministic cardinality envelope \(J_m\ge L_m\) such that
\[
  \mu_m^{\rm ch}+t_m\le C_0n_m .
\]
If
\[
  \sum_m\left[
    \mathbb P(F_m^{\rm cert})
    +\exp\left\{
      -\frac{t_m^2}{2(v_m^{\rm ch}+b_mt_m/3)}
    \right\}
  \right]<\infty,
\]
then the selection event in
Corollary~\ref{cor:selected-structured-submenu-mixed-dominated-operator-tail} is
summable: almost surely, eventually
\[
  |\mathscr S_m|\le J_m,
  \qquad
  \sum_{k\in\mathscr S_m}\mathfrak s_{m,k}\le C_0n_m .
\]

Consequently, suppose each candidate perturbation stream is tagged \(\psi\) or
\(\mu\), with parameters \(\varepsilon_{m,k},A_{m,k},s'_{m,k}\), and satisfies
the corresponding mixed dominated tail and drift hypotheses after adjoining
\(\mathcal H_m\).  If
\[
  \ell_m:=\log(e(D_m+n_m+1))\le B\sqrt{n_m}
\]
eventually,
\[
  \sum_m\sum_k
  \bigl(\mathbb P(E^{\rm tail}_{m,k})+\varepsilon_{m,k}\bigr)<\infty,
\]
and
\[
  \sum_mJ_m(D_m+n_m+1)
  (e(D_m+n_m+1))^{-\gamma_{\rm tc}(K_{\rm tail},1,1,B)}<\infty,
\]
then the product-categorical selected structured submenus have the same
eventual common signing, contraction normalization, fixed-menu coverage, and
all-integer/sparse thresholds as in
Corollary~\ref{cor:selected-structured-submenu-mixed-dominated-operator-tail}.

This covers independent randomized choice of one structured template per block,
product mixtures of finite-lift and algebraic core templates, randomized
hyperparameter menus, and sampling with replacement when each draw is treated as
a distinct candidate copy.
\end{corollary}

\begin{proof}
Condition on \(\mathcal G_m\) and put
\[
  X_{m,g}:=
  \begin{cases}
    \mathfrak s_{m,\xi_{m,g}},& \xi_{m,g}\ne\varnothing,\\
    0,& \xi_{m,g}=\varnothing .
  \end{cases}
\]
Then \(0\le X_{m,g}\le b_m\), the variables \(X_{m,g}\) are conditionally
independent, and
\[
  \sum_g\mathbb E[X_{m,g}\mid\mathcal G_m]\le\mu_m^{\rm ch},
  \qquad
  \sum_g\mathbb E[X_{m,g}^2\mid\mathcal G_m]\le v_m^{\rm ch}.
\]
Hence their conditional variance sum is at most \(v_m^{\rm ch}\).  Bernstein's
inequality gives
\[
  \mathbb P\left\{
    \sum_{k\in\mathscr S_m}\mathfrak s_{m,k}>\mu_m^{\rm ch}+t_m
    \,\middle|\,\mathcal G_m
  \right\}
  \le
  \exp\left\{-\frac{t_m^2}{2(v_m^{\rm ch}+b_mt_m/3)}\right\}.
\]
Also \(|\mathscr S_m|\le L_m\le J_m\) deterministically.  After averaging,
adding \(F_m^{\rm cert}\), and applying Borel--Cantelli, the selected catalogue
eventually satisfies the deterministic hypotheses of
Corollary~\ref{cor:selected-structured-submenu-mixed-dominated-operator-tail}.
Ordering the selected labels and padding by zero streams again converts the
candidate tail summability into the active-prefix summability required there.
The final assertions are exactly the conclusions of the selected-submenu
corollary.
\end{proof}

\begin{corollary}[Adaptive-categorical selected structured submenus with mixed dominated noise]
\label{cor:adaptive-categorical-selected-structured-submenu-mixed-dominated-operator-tail}
Fix structural constants \(r_0,b_0,c_0,L_0,C_0\ge1\).  For each \(m\), let
\[
  \mathcal G_m=\mathcal F_{m,0}\subset
  \mathcal F_{m,1}\subset\cdots\subset\mathcal F_{m,L_m}
\]
be a filtration and let the finite candidate list be partitioned into
deterministic slots \(\mathcal I_{m,1},\ldots,\mathcal I_{m,L_m}\).  The
\(k\)-th candidate has charge \(\mathfrak s_{m,k}\in[0,b_m]\), with \(b_m>0\),
core streams \(B_{m,k,1},\ldots,B_{m,k,n_m}\), and perturbation streams
\(Y_{m,k,1},\ldots,Y_{m,k,n_m}\).  In slot \(g\), choose at most one label by an
\(\mathcal I_{m,g}\cup\{\varnothing\}\)-valued random variable
\(\xi_{m,g}\) which is \(\mathcal F_{m,g}\)-measurable.  Assume deterministic
probability envelopes \(\pi_{m,k}\in[0,1]\) satisfy
\[
  \mathbb P\{\xi_{m,g}=k\mid\mathcal F_{m,g-1}\}\le \pi_{m,k}
  \qquad(k\in\mathcal I_{m,g}).
\]
Let
\[
\begin{aligned}
  \eta_{m,k}&:=\mathbf 1_{\{\xi_{m,g}=k\}}
    \quad(k\in\mathcal I_{m,g}),\\
  \mathscr S_m&:=\{k:\eta_{m,k}=1\},\\
  \mathcal H_m&:=\mathcal F_{m,L_m},
\end{aligned}
\]
and define
\[
  \mu_m^{\rm ch}:=\sum_k\pi_{m,k}\mathfrak s_{m,k},
  \qquad
  v_m^{\rm ch}:=\sum_k\pi_{m,k}\mathfrak s_{m,k}^2 .
\]

Assume that outside events \(F_m^{\rm cert}\in\mathcal G_m\), with
\(\sum_m\mathbb P(F_m^{\rm cert})<\infty\), every candidate core is a
contraction and has one of the finite-lift orbit-closure or algebraic/group
certificates from
Corollary~\ref{cor:selected-structured-submenu-mixed-dominated-operator-tail},
with charge bounded by the displayed \(\mathfrak s_{m,k}\).  Choose \(t_m>0\)
and a deterministic cardinality envelope \(J_m\ge L_m\) such that
\[
  \mu_m^{\rm ch}+t_m\le C_0n_m .
\]
If
\[
  \sum_m\left[
    \mathbb P(F_m^{\rm cert})
    +\exp\left\{
      -\frac{t_m^2}{2(v_m^{\rm ch}+b_mt_m/3)}
    \right\}
  \right]<\infty,
\]
then the selection event in
Corollary~\ref{cor:selected-structured-submenu-mixed-dominated-operator-tail} is
summable: almost surely, eventually
\[
  |\mathscr S_m|\le J_m,
  \qquad
  \sum_{k\in\mathscr S_m}\mathfrak s_{m,k}\le C_0n_m .
\]

Consequently, suppose each candidate perturbation stream is tagged \(\psi\) or
\(\mu\), with parameters \(\varepsilon_{m,k},A_{m,k},s'_{m,k}\), and satisfies
the corresponding mixed dominated tail and drift hypotheses after adjoining
\(\mathcal H_m\).  If
\[
  \ell_m:=\log(e(D_m+n_m+1))\le B\sqrt{n_m}
\]
eventually,
\[
  \sum_m\sum_k
  \bigl(\mathbb P(E^{\rm tail}_{m,k})+\varepsilon_{m,k}\bigr)<\infty,
\]
and
\[
  \sum_mJ_m(D_m+n_m+1)
  (e(D_m+n_m+1))^{-\gamma_{\rm tc}(K_{\rm tail},1,1,B)}<\infty,
\]
then the adaptively categorical selected structured submenus have the same
eventual common signing, contraction normalization, fixed-menu coverage, and
all-integer/sparse thresholds as in
Corollary~\ref{cor:selected-structured-submenu-mixed-dominated-operator-tail}.

This covers online or state-dependent randomized template choices, Markov
hyperparameter menus, adaptive blockwise model selection, and sequential
architecture/library choices whenever each stage chooses at most one structured
candidate and has deterministic predictable probability envelopes.
\end{corollary}

\begin{proof}
For each slot put
\[
  X_{m,g}:=
  \begin{cases}
    \mathfrak s_{m,\xi_{m,g}},& \xi_{m,g}\ne\varnothing,\\
    0,& \xi_{m,g}=\varnothing .
  \end{cases}
\]
Then \(0\le X_{m,g}\le b_m\) and
\[
  \sum_g\mathbb E[X_{m,g}\mid\mathcal F_{m,g-1}]\le\mu_m^{\rm ch},
  \qquad
  \sum_g\mathbb E[X_{m,g}^2\mid\mathcal F_{m,g-1}]\le v_m^{\rm ch}.
\]
Hence the martingale
\[
  M_j:=\sum_{g\le j}\left(X_{m,g}
  -\mathbb E[X_{m,g}\mid\mathcal F_{m,g-1}]\right)
\]
has increments bounded above by \(b_m\) and predictable quadratic variation at
most \(v_m^{\rm ch}\).  If
\(\sum_{k\in\mathscr S_m}\mathfrak s_{m,k}>\mu_m^{\rm ch}+t_m\), then
\(M_{L_m}>t_m\).  The scalar Freedman inequality gives
\[
  \mathbb P\left\{
    \sum_{k\in\mathscr S_m}\mathfrak s_{m,k}>\mu_m^{\rm ch}+t_m
  \right\}
  \le
  \exp\left\{-\frac{t_m^2}{2(v_m^{\rm ch}+b_mt_m/3)}\right\}.
\]
Also \(|\mathscr S_m|\le L_m\le J_m\) deterministically.  Adding
\(F_m^{\rm cert}\) and applying Borel--Cantelli gives the deterministic
selection hypotheses eventually.  Order the selected labels by slot and then by
candidate order inside each slot, pad by zero streams up to \(J_m\), and apply
Corollary~\ref{cor:selected-structured-submenu-mixed-dominated-operator-tail}.
\end{proof}

\begin{corollary}[Random-envelope adaptive-categorical selected structured submenus]
\label{cor:random-envelope-adaptive-categorical-selected-structured-submenu-mixed-dominated-operator-tail}
Fix structural constants \(r_0,b_0,c_0,L_0,C_0\ge1\).  For each \(m\), let
\[
  \mathcal G_m=\mathcal F_{m,0}\subset
  \mathcal F_{m,1}\subset\cdots\subset\mathcal F_{m,L_m}
\]
be a filtration and let the finite candidate list be partitioned into
deterministic slots \(\mathcal I_{m,1},\ldots,\mathcal I_{m,L_m}\).  The
\(k\)-th candidate has charge \(\mathfrak s_{m,k}\in[0,b_m]\), with \(b_m>0\),
core streams \(B_{m,k,1},\ldots,B_{m,k,n_m}\), and perturbation streams
\(Y_{m,k,1},\ldots,Y_{m,k,n_m}\).  In slot \(g\), choose at most one label by an
\(\mathcal I_{m,g}\cup\{\varnothing\}\)-valued random variable \(\xi_{m,g}\)
which is \(\mathcal F_{m,g}\)-measurable.  For each \(k\in\mathcal I_{m,g}\),
let \(\Pi_{m,k}\in[0,1]\) be \(\mathcal F_{m,g-1}\)-measurable and assume
\[
  \mathbb P\{\xi_{m,g}=k\mid\mathcal F_{m,g-1}\}\le \Pi_{m,k}.
\]
Let
\[
\begin{aligned}
  \eta_{m,k}&:=\mathbf 1_{\{\xi_{m,g}=k\}}
    \quad(k\in\mathcal I_{m,g}),\\
  \mathscr S_m&:=\{k:\eta_{m,k}=1\},\\
  \mathcal H_m&:=\mathcal F_{m,L_m},
\end{aligned}
\]
and define the predictable random charge profiles
\[
  M_m^{\rm ch}:=\sum_k\Pi_{m,k}\mathfrak s_{m,k},
  \qquad
  V_m^{\rm ch}:=\sum_k\Pi_{m,k}\mathfrak s_{m,k}^2 .
\]

Assume that there are events \(F_m^{\rm env}\in\mathcal F_{m,L_m}\), with
\(\sum_m\mathbb P(F_m^{\rm env})<\infty\), and deterministic bounds
\(\bar\mu_m^{\rm ch},\bar v_m^{\rm ch}\ge0\) such that outside
\(F_m^{\rm env}\),
\[
  M_m^{\rm ch}\le\bar\mu_m^{\rm ch},
  \qquad
  V_m^{\rm ch}\le\bar v_m^{\rm ch}.
\]
Assume also that outside events \(F_m^{\rm cert}\in\mathcal G_m\), with
\(\sum_m\mathbb P(F_m^{\rm cert})<\infty\), every candidate core is a
contraction and has one of the finite-lift orbit-closure or algebraic/group
certificates from
Corollary~\ref{cor:selected-structured-submenu-mixed-dominated-operator-tail},
with charge bounded by the displayed \(\mathfrak s_{m,k}\).  Choose \(t_m>0\)
and a deterministic cardinality envelope \(J_m\ge L_m\) such that
\[
  \bar\mu_m^{\rm ch}+t_m\le C_0n_m .
\]
If
\[
  \sum_m\left[
    \mathbb P(F_m^{\rm env})+\mathbb P(F_m^{\rm cert})
    +\exp\left\{
      -\frac{t_m^2}{2(\bar v_m^{\rm ch}+b_mt_m/3)}
    \right\}
  \right]<\infty,
\]
then the selection event in
Corollary~\ref{cor:selected-structured-submenu-mixed-dominated-operator-tail} is
summable: almost surely, eventually
\[
  |\mathscr S_m|\le J_m,
  \qquad
  \sum_{k\in\mathscr S_m}\mathfrak s_{m,k}\le C_0n_m .
\]

Consequently, suppose each candidate perturbation stream is tagged \(\psi\) or
\(\mu\), with parameters \(\varepsilon_{m,k},A_{m,k},s'_{m,k}\), and satisfies
the corresponding mixed dominated tail and drift hypotheses after adjoining
\(\mathcal H_m\).  If
\[
  \ell_m:=\log(e(D_m+n_m+1))\le B\sqrt{n_m}
\]
eventually,
\[
  \sum_m\sum_k
  \bigl(\mathbb P(E^{\rm tail}_{m,k})+\varepsilon_{m,k}\bigr)<\infty,
\]
and
\[
  \sum_mJ_m(D_m+n_m+1)
  (e(D_m+n_m+1))^{-\gamma_{\rm tc}(K_{\rm tail},1,1,B)}<\infty,
\]
then the random-envelope adaptive-categorical selected structured submenus have
the same eventual common signing, contraction normalization, fixed-menu
coverage, and all-integer/sparse thresholds as in
Corollary~\ref{cor:selected-structured-submenu-mixed-dominated-operator-tail}.

This covers state-dependent template probabilities, posterior or online
hyperparameter menus, adaptive blockwise architecture choices, and sequential
model-selection sublibraries whose predictable aggregate charge profiles are
certified with summable failure probability.
\end{corollary}

\begin{proof}
For each slot put
\[
  X_{m,g}:=
  \begin{cases}
    \mathfrak s_{m,\xi_{m,g}},& \xi_{m,g}\ne\varnothing,\\
    0,& \xi_{m,g}=\varnothing .
  \end{cases}
\]
Let \(q_{m,g}:=\mathbb E[X_{m,g}\mid\mathcal F_{m,g-1}]\).  Then
\[
  q_{m,g}\le\sum_{k\in\mathcal I_{m,g}}\Pi_{m,k}\mathfrak s_{m,k},
  \qquad
  \mathbb E[X_{m,g}^2\mid\mathcal F_{m,g-1}]
  \le\sum_{k\in\mathcal I_{m,g}}\Pi_{m,k}\mathfrak s_{m,k}^2 .
\]
Hence
\[
  M_j:=\sum_{g\le j}(X_{m,g}-q_{m,g})
\]
is a martingale with increments bounded above by \(b_m\), and its predictable
quadratic variation is at most \(V_m^{\rm ch}\).  On \((F_m^{\rm env})^c\), if
\[
  \sum_{k\in\mathscr S_m}\mathfrak s_{m,k}>\bar\mu_m^{\rm ch}+t_m,
\]
then \(M_{L_m}>t_m\) while the predictable quadratic variation is at most
\(\bar v_m^{\rm ch}\).  The scalar Freedman inequality therefore gives
\[
\begin{aligned}
  &\mathbb P\left\{
    (F_m^{\rm env})^c,
    \sum_{k\in\mathscr S_m}\mathfrak s_{m,k}>\bar\mu_m^{\rm ch}+t_m
  \right\} \\
  &\qquad\le
  \exp\left\{
    -\frac{t_m^2}{2(\bar v_m^{\rm ch}+b_mt_m/3)}
  \right\}.
\end{aligned}
\]
Together with \(\mathbb P(F_m^{\rm env})+\mathbb P(F_m^{\rm cert})\), the
displayed series is summable.  Since \(|\mathscr S_m|\le L_m\le J_m\)
deterministically, Borel--Cantelli gives the deterministic cardinality and
charge hypotheses eventually.

On that event, order selected labels by slot and candidate order, pad by zero
streams up to \(J_m\), and apply
Corollary~\ref{cor:selected-structured-submenu-mixed-dominated-operator-tail}.
The candidate tail summability dominates the selected active-prefix tail
summability, and the deterministic matrix-tail series is unchanged.
\end{proof}

\begin{corollary}[Freedman-selected adaptive structured submenus with mixed dominated noise]
\label{cor:freedman-selected-structured-submenu-mixed-dominated-operator-tail}
Fix structural constants \(r_0,b_0,c_0,L_0,C_0\ge1\).  For each \(m\), let
\[
  \mathcal G_m=\mathcal F_{m,0}\subset
  \mathcal F_{m,1}\subset\cdots\subset\mathcal F_{m,M_m}
\]
be a filtration exposing a finite candidate list \(1,\ldots,M_m\).  The
\(k\)-th candidate has charge \(\mathfrak s_{m,k}\in[0,b_m]\), with
\(b_m>0\), core streams \(B_{m,k,1},\ldots,B_{m,k,n_m}\), and perturbation
streams \(Y_{m,k,1},\ldots,Y_{m,k,n_m}\).  Let
\(\eta_{m,k}\in\{0,1\}\) be \(\mathcal F_{m,k}\)-measurable selection
indicators satisfying the predictable envelope
\[
  \mathbb E[\eta_{m,k}\mid\mathcal F_{m,k-1}]\le p_{m,k}
  \qquad(1\le k\le M_m),
\]
where the \(p_{m,k}\in[0,1]\) are deterministic.  No conditional independence
is assumed.  Put
\[
  \mathscr S_m:=\{k:\eta_{m,k}=1\},
  \qquad
  \mathcal H_m:=\mathcal F_{m,M_m},
\]
and define
\[
  \mu_m^{\rm ch}:=\sum_{k=1}^{M_m}p_{m,k}\mathfrak s_{m,k},\qquad
  v_m^{\rm ch}:=\sum_{k=1}^{M_m}p_{m,k}\mathfrak s_{m,k}^2,\qquad
  \mu_m^{\#}:=\sum_{k=1}^{M_m}p_{m,k}.
\]

Assume that outside events \(F_m^{\rm cert}\in\mathcal G_m\), with
\(\sum_m\mathbb P(F_m^{\rm cert})<\infty\), every selected core is a contraction
and has one of the finite-lift orbit-closure or algebraic/group certificates
from
Corollary~\ref{cor:selected-structured-submenu-mixed-dominated-operator-tail},
with charge bounded by the displayed \(\mathfrak s_{m,k}\).  Choose
\(t_m,u_m>0\) and a deterministic cardinality envelope \(J_m\ge1\) such that
\[
  \mu_m^{\rm ch}+t_m\le C_0n_m,\qquad
  \mu_m^{\#}+u_m\le J_m .
\]
If
\[
  \sum_m\left[
    \mathbb P(F_m^{\rm cert})
    +\exp\left\{
      -\frac{t_m^2}{2(v_m^{\rm ch}+b_mt_m/3)}
    \right\}
    +\exp\left\{
      -\frac{u_m^2}{2(\mu_m^{\#}+u_m/3)}
    \right\}
  \right]<\infty,
\]
then the selection event in
Corollary~\ref{cor:selected-structured-submenu-mixed-dominated-operator-tail} is
summable: almost surely, eventually
\[
  |\mathscr S_m|\le J_m,
  \qquad
  \sum_{k\in\mathscr S_m}\mathfrak s_{m,k}\le C_0n_m .
\]

Consequently, suppose each candidate perturbation stream is tagged \(\psi\) or
\(\mu\), with parameters \(\varepsilon_{m,k},A_{m,k},s'_{m,k}\), and satisfies
the corresponding mixed dominated tail and drift hypotheses after adjoining
\(\mathcal H_m\).  If
\[
  \ell_m:=\log(e(D_m+n_m+1))\le B\sqrt{n_m}
\]
eventually,
\[
  \sum_m\sum_{k=1}^{M_m}
  \bigl(\mathbb P(E^{\rm tail}_{m,k})+\varepsilon_{m,k}\bigr)<\infty,
\]
and
\[
  \sum_mJ_m(D_m+n_m+1)
  (e(D_m+n_m+1))^{-\gamma_{\rm tc}(K_{\rm tail},1,1,B)}<\infty,
\]
then the adaptively selected structured submenus have the same eventual common
signing, contraction normalization, fixed-menu coverage, and all-integer/sparse
thresholds as in
Corollary~\ref{cor:selected-structured-submenu-mixed-dominated-operator-tail}.
\end{corollary}

\begin{proof}
For each \(k\), set
\[
  \pi_{m,k}:=\mathbb E[\eta_{m,k}\mid\mathcal F_{m,k-1}].
\]
The process
\[
  M_j^{\rm ch}:=
  \sum_{k\le j}(\eta_{m,k}-\pi_{m,k})\mathfrak s_{m,k}
\]
is a martingale with increments bounded above by \(b_m\), and its predictable
quadratic variation is at most
\[
  \sum_k\pi_{m,k}\mathfrak s_{m,k}^2
  \le v_m^{\rm ch}.
\]
Since \(\sum_k(\pi_{m,k}-p_{m,k})\mathfrak s_{m,k}\le0\), the event
\[
  \sum_k\eta_{m,k}\mathfrak s_{m,k}>\mu_m^{\rm ch}+t_m
\]
implies \(M_{M_m}^{\rm ch}>t_m\).  The scalar Freedman inequality therefore
gives
\[
  \mathbb P\left\{
    \sum_k\eta_{m,k}\mathfrak s_{m,k}>C_0n_m
  \right\}
  \le
  \exp\left\{
    -\frac{t_m^2}{2(v_m^{\rm ch}+b_mt_m/3)}
  \right\}.
\]
The same argument with increments \(\eta_{m,k}-\pi_{m,k}\), jump bound \(1\),
and predictable variance at most \(\sum_kp_{m,k}=\mu_m^{\#}\) gives
\[
  \mathbb P\{|\mathscr S_m|>J_m\}
  \le
  \exp\left\{
    -\frac{u_m^2}{2(\mu_m^{\#}+u_m/3)}
  \right\}.
\]
Adding \(F_m^{\rm cert}\), the displayed series is summable, so Borel--Cantelli
gives the required eventual cardinality and charge bounds.

On the resulting good event, order the selected labels by the candidate order
and pad by zero streams up to the deterministic length \(J_m\).  The candidate
tail summability dominates the selected tail summability required by
Corollary~\ref{cor:selected-structured-submenu-mixed-dominated-operator-tail},
and the deterministic matrix-tail series is unchanged.  Applying that
corollary gives all stated conclusions.
\end{proof}

\begin{corollary}[Random-envelope Freedman-selected adaptive structured submenus]
\label{cor:random-envelope-freedman-selected-structured-submenu-mixed-dominated-operator-tail}
Fix structural constants \(r_0,b_0,c_0,L_0,C_0\ge1\).  For each \(m\), let
\[
  \mathcal G_m=\mathcal F_{m,0}\subset
  \mathcal F_{m,1}\subset\cdots\subset\mathcal F_{m,M_m}
\]
be a filtration exposing a finite candidate list \(1,\ldots,M_m\).  The
\(k\)-th candidate has charge \(\mathfrak s_{m,k}\in[0,b_m]\), with
\(b_m>0\), core streams \(B_{m,k,1},\ldots,B_{m,k,n_m}\), and perturbation
streams \(Y_{m,k,1},\ldots,Y_{m,k,n_m}\).  Let
\(\eta_{m,k}\in\{0,1\}\) be \(\mathcal F_{m,k}\)-measurable selection
indicators.  Assume there are predictable random envelopes
\(P_{m,k}\in[0,1]\), with \(P_{m,k}\) \(\mathcal F_{m,k-1}\)-measurable, such
that
\[
  \mathbb E[\eta_{m,k}\mid\mathcal F_{m,k-1}]\le P_{m,k}
  \qquad(1\le k\le M_m).
\]
Put
\[
  \mathscr S_m:=\{k:\eta_{m,k}=1\},
  \qquad
  \mathcal H_m:=\mathcal F_{m,M_m},
\]
and define the random predictable profiles
\[
  M_m^{\rm ch}:=\sum_{k=1}^{M_m}P_{m,k}\mathfrak s_{m,k},
  \qquad
  V_m^{\rm ch}:=\sum_{k=1}^{M_m}P_{m,k}\mathfrak s_{m,k}^2,
  \qquad
  M_m^{\#}:=\sum_{k=1}^{M_m}P_{m,k}.
\]

Assume that there are events \(F_m^{\rm env}\in\mathcal F_{m,M_m}\), with
\(\sum_m\mathbb P(F_m^{\rm env})<\infty\), and deterministic bounds
\(\bar\mu_m^{\rm ch},\bar v_m^{\rm ch},\bar\mu_m^{\#}\ge0\) such that outside
\(F_m^{\rm env}\),
\[
  M_m^{\rm ch}\le\bar\mu_m^{\rm ch},\qquad
  V_m^{\rm ch}\le\bar v_m^{\rm ch},\qquad
  M_m^{\#}\le\bar\mu_m^{\#}.
\]
Assume also that outside events \(F_m^{\rm cert}\in\mathcal G_m\), with
\(\sum_m\mathbb P(F_m^{\rm cert})<\infty\), every selected core is a contraction
and has one of the finite-lift orbit-closure or algebraic/group certificates
from
Corollary~\ref{cor:selected-structured-submenu-mixed-dominated-operator-tail},
with charge bounded by the displayed \(\mathfrak s_{m,k}\).  Choose
\(t_m,u_m>0\) and a deterministic cardinality envelope \(J_m\ge1\) such that
\[
  \bar\mu_m^{\rm ch}+t_m\le C_0n_m,\qquad
  \bar\mu_m^{\#}+u_m\le J_m .
\]
Put
\[
\begin{aligned}
  \Lambda_m:={}&
    \mathbb P(F_m^{\rm env})+\mathbb P(F_m^{\rm cert}) \\
  &+\exp\left\{
      -\frac{t_m^2}{2(\bar v_m^{\rm ch}+b_mt_m/3)}
    \right\}
    +\exp\left\{
      -\frac{u_m^2}{2(\bar\mu_m^{\#}+u_m/3)}
    \right\}.
\end{aligned}
\]
If
\[
  \sum_m\Lambda_m<\infty,
\]
then the selection event in
Corollary~\ref{cor:selected-structured-submenu-mixed-dominated-operator-tail} is
summable: almost surely, eventually
\[
  |\mathscr S_m|\le J_m,
  \qquad
  \sum_{k\in\mathscr S_m}\mathfrak s_{m,k}\le C_0n_m .
\]

Consequently, suppose each candidate perturbation stream is tagged \(\psi\) or
\(\mu\), with parameters \(\varepsilon_{m,k},A_{m,k},s'_{m,k}\), and satisfies
the corresponding mixed dominated tail and drift hypotheses after adjoining
\(\mathcal H_m\).  If
\[
  \ell_m:=\log(e(D_m+n_m+1))\le B\sqrt{n_m}
\]
eventually,
\[
  \sum_m\sum_{k=1}^{M_m}
  \bigl(\mathbb P(E^{\rm tail}_{m,k})+\varepsilon_{m,k}\bigr)<\infty,
\]
and
\[
  \sum_mJ_m(D_m+n_m+1)
  (e(D_m+n_m+1))^{-\gamma_{\rm tc}(K_{\rm tail},1,1,B)}<\infty,
\]
then the random-envelope Freedman-selected structured submenus have the same
eventual common signing, contraction normalization, fixed-menu coverage, and
all-integer/sparse thresholds as in
Corollary~\ref{cor:selected-structured-submenu-mixed-dominated-operator-tail}.

This covers adaptive thinning, stopping-rule catalogues, online selected
structured libraries, and general sequential model-selection submenus with
state-dependent predictable inclusion intensities, once their realized envelope
profiles are certified outside summable exceptional events.
\end{corollary}

\begin{proof}
Set
\[
  \pi_{m,k}:=\mathbb E[\eta_{m,k}\mid\mathcal F_{m,k-1}].
\]
For the charge process,
\[
  M_j^{\rm ch}:=
  \sum_{k\le j}(\eta_{m,k}-\pi_{m,k})\mathfrak s_{m,k}
\]
is a martingale with increments bounded above by \(b_m\), and predictable
quadratic variation at most
\[
  \sum_k\pi_{m,k}\mathfrak s_{m,k}^2\le V_m^{\rm ch}.
\]
On \((F_m^{\rm env})^c\), the event
\[
  \sum_k\eta_{m,k}\mathfrak s_{m,k}>\bar\mu_m^{\rm ch}+t_m
\]
implies \(M_{M_m}^{\rm ch}>t_m\), while the predictable quadratic variation is
at most \(\bar v_m^{\rm ch}\).  The scalar Freedman inequality gives
\[
\begin{aligned}
  &\mathbb P\left\{
    (F_m^{\rm env})^c,
    \sum_k\eta_{m,k}\mathfrak s_{m,k}>
      \bar\mu_m^{\rm ch}+t_m
  \right\} \\
  &\qquad\le
  \exp\left\{
    -\frac{t_m^2}{2(\bar v_m^{\rm ch}+b_mt_m/3)}
  \right\}.
\end{aligned}
\]
The same argument applied to
\[
  M_j^{\#}:=\sum_{k\le j}(\eta_{m,k}-\pi_{m,k})
\]
uses jump bound \(1\) and predictable quadratic variation at most
\(M_m^{\#}\).  Hence
\[
  \mathbb P\left\{
    (F_m^{\rm env})^c,\ |\mathscr S_m|>\bar\mu_m^{\#}+u_m
  \right\}
  \le
  \exp\left\{
    -\frac{u_m^2}{2(\bar\mu_m^{\#}+u_m/3)}
  \right\}.
\]
After adding \(F_m^{\rm env}\cup F_m^{\rm cert}\), the displayed series is
summable, so Borel--Cantelli gives the eventual deterministic cardinality and
charge bounds.

On the resulting good event, order the selected labels by candidate order and
pad by zero streams up to \(J_m\).  The candidate tail summability dominates the
selected tail summability required by
Corollary~\ref{cor:selected-structured-submenu-mixed-dominated-operator-tail},
and the deterministic matrix-tail series is unchanged.  Applying that corollary
gives all stated conclusions.
\end{proof}

\begin{corollary}[Serfling-selected fixed-size structured submenus with mixed dominated noise]
\label{cor:serfling-selected-structured-submenu-mixed-dominated-operator-tail}
Fix structural constants \(r_0,b_0,c_0,L_0,C_0\ge1\).  For each \(m\), let
\(\mathcal G_m\) be a base environment.  Let a finite candidate catalogue be
partitioned into deterministic strata
\[
  \mathcal I_{m,1},\ldots,\mathcal I_{m,A_m},
  \qquad |\mathcal I_{m,a}|=M_{m,a}.
\]
From stratum \(a\), choose exactly \(R_{m,a}\le M_{m,a}\) labels uniformly
without replacement, conditionally on \(\mathcal G_m\).  Write the chosen subset
as \(\mathscr S_{m,a}\), put
\[
  \mathscr S_m:=\bigcup_{a=1}^{A_m}\mathscr S_{m,a},
  \qquad
  J_m:=\sum_{a=1}^{A_m}R_{m,a},
\]
and let
\[
  \mathcal H_m:=\mathcal G_m\vee
  \sigma(\mathscr S_{m,1},\ldots,\mathscr S_{m,A_m}).
\]
The \(k\)-th candidate in stratum \(a\) has charge
\(\mathfrak s_{m,a,k}\in[0,b_{m,a}]\), core streams
\(B_{m,a,k,1},\ldots,B_{m,a,k,n_m}\), and perturbation streams
\(Y_{m,a,k,1},\ldots,Y_{m,a,k,n_m}\).  Define the finite-population mean charge
\[
  \bar{\mathfrak s}_{m,a}:=
  \frac{1}{M_{m,a}}\sum_{k\in\mathcal I_{m,a}}\mathfrak s_{m,a,k}.
\]

Assume that outside events \(F_m^{\rm cert}\in\mathcal G_m\), with
\(\sum_m\mathbb P(F_m^{\rm cert})<\infty\), every candidate core is a
contraction and has one of the finite-lift orbit-closure or algebraic/group
certificates from
Corollary~\ref{cor:selected-structured-submenu-mixed-dominated-operator-tail},
with charge bounded by the displayed \(\mathfrak s_{m,a,k}\).  Choose
numbers \(t_{m,a}>0\) for all strata with \(R_{m,a}>0\) so that
\[
  \sum_{a=1}^{A_m}
  \bigl(R_{m,a}\bar{\mathfrak s}_{m,a}+t_{m,a}\bigr)
  \le C_0n_m,
\]
where the \(a\)-term is read as \(0\) when \(R_{m,a}=0\).  If
\[
  \sum_m\left[
    \mathbb P(F_m^{\rm cert})
    +\sum_{\substack{1\le a\le A_m\\ R_{m,a}>0}}
      \exp\left\{-\frac{2t_{m,a}^2}{R_{m,a}b_{m,a}^2}\right\}
  \right]<\infty,
\]
then the selection event in
Corollary~\ref{cor:selected-structured-submenu-mixed-dominated-operator-tail} is
summable: almost surely, eventually
\[
  |\mathscr S_m|=J_m,
  \qquad
  \sum_{a=1}^{A_m}\sum_{k\in\mathscr S_{m,a}}
  \mathfrak s_{m,a,k}\le C_0n_m .
\]

Consequently, suppose each candidate perturbation stream is tagged \(\psi\) or
\(\mu\), with parameters \(\varepsilon_{m,a,k},A_{m,a,k},s'_{m,a,k}\), and
satisfies the corresponding mixed dominated tail and drift hypotheses after
adjoining \(\mathcal H_m\).  If
\[
  \ell_m:=\log(e(D_m+n_m+1))\le B\sqrt{n_m}
\]
eventually,
\[
  \sum_m\sum_{a=1}^{A_m}\sum_{k\in\mathcal I_{m,a}}
  \bigl(\mathbb P(E^{\rm tail}_{m,a,k})+\varepsilon_{m,a,k}\bigr)<\infty,
\]
and
\[
  \sum_mJ_m(D_m+n_m+1)
  (e(D_m+n_m+1))^{-\gamma_{\rm tc}(K_{\rm tail},1,1,B)}<\infty,
\]
then the fixed-size or stratified without-replacement selected submenus have the
same eventual common signing, contraction normalization, fixed-menu coverage,
and all-integer/sparse thresholds as in
Corollary~\ref{cor:selected-structured-submenu-mixed-dominated-operator-tail}.
\end{corollary}

\begin{proof}
The selected cardinality is deterministically \(J_m\).  Condition on
\(\mathcal G_m\).  In stratum \(a\), the selected charges are sampled without
replacement from the finite population
\(\{\mathfrak s_{m,a,k}:k\in\mathcal I_{m,a}\}\subset[0,b_{m,a}]\).
Hoeffding's comparison for sampling without replacement gives, for
\(R_{m,a}>0\),
\[
  \mathbb P\left\{
    \sum_{k\in\mathscr S_{m,a}}\mathfrak s_{m,a,k}
    >R_{m,a}\bar{\mathfrak s}_{m,a}+t_{m,a}
    \,\middle|\,\mathcal G_m
  \right\}
  \le
  \exp\left\{-\frac{2t_{m,a}^2}{R_{m,a}b_{m,a}^2}\right\}.
\]
For \(R_{m,a}=0\) the stratum contributes nothing.  A union bound over strata
shows that the total selected charge exceeds \(C_0n_m\) with conditional
probability bounded by the displayed stratum-tail sum.  After averaging and
adding \(F_m^{\rm cert}\), the assumed summability and Borel--Cantelli give the
summable selection event required by
Corollary~\ref{cor:selected-structured-submenu-mixed-dominated-operator-tail}.

On the eventual good event, order all selected labels by the fixed stratum order
and the catalogue order inside each stratum, then pad by zero streams up to
length \(J_m\).  The candidate tail summability dominates the selected tail
summability in the selected-submenu corollary, and the deterministic matrix-tail
series is the same one with cardinality envelope \(J_m\).  Applying that
corollary gives all stated signing, normalization, coverage, and threshold
conclusions.
\end{proof}

\begin{corollary}[Hybrid selected structured submenus with mixed dominated noise]
\label{cor:hybrid-selected-structured-submenu-mixed-dominated-operator-tail}
Fix structural constants \(r_0,b_0,c_0,L_0,C_0\ge1\).  For each \(m\), let
\(\mathcal G_m\) be a base environment and partition a finite candidate
catalogue into deterministic components
\[
  \mathcal I_{m,1},\ldots,\mathcal I_{m,R_m}.
\]
The candidate \(k\in\mathcal I_{m,h}\) has charge
\(\mathfrak s_{m,h,k}\ge0\), core stream
\((B_{m,h,k,i})_{i\le n_m}\), and perturbation stream
\((Y_{m,h,k,i})_{i\le n_m}\).  In component \(h\), select a random subset
\(\mathscr S_{m,h}\subset\mathcal I_{m,h}\), by any mechanism, and put
\[
  \mathscr S_m:=\bigcup_{h=1}^{R_m}\mathscr S_{m,h},
  \qquad
  \mathcal H_m:=
  \mathcal G_m\vee\sigma(\mathscr S_{m,1},\ldots,\mathscr S_{m,R_m}).
\]
Assume that for every component there are deterministic envelopes
\[
  J_{m,h}\in\mathbb Z_{\ge0},\qquad U_{m,h}\ge0,
\]
and selection-failure events \(F_{m,h}^{\rm sel}\), such that
\[
  \sum_m\sum_{h=1}^{R_m}\mathbb P(F_{m,h}^{\rm sel})<\infty
\]
and, outside \(F_{m,h}^{\rm sel}\),
\[
  |\mathscr S_{m,h}|\le J_{m,h},
  \qquad
  \sum_{k\in\mathscr S_{m,h}}\mathfrak s_{m,h,k}\le U_{m,h}.
\]
Let
\[
  J_m:=\max\left\{1,\sum_{h=1}^{R_m}J_{m,h}\right\},
  \qquad
  \sum_{h=1}^{R_m}U_{m,h}\le C_0n_m .
\]
Assume also that outside events \(F_m^{\rm cert}\in\mathcal G_m\), with
\(\sum_m\mathbb P(F_m^{\rm cert})<\infty\), every selected core is a
contraction and has one of the finite-lift orbit-closure or algebraic/group
certificates from
Corollary~\ref{cor:selected-structured-submenu-mixed-dominated-operator-tail},
with charge bounded by the displayed \(\mathfrak s_{m,h,k}\).

Then the selection event in
Corollary~\ref{cor:selected-structured-submenu-mixed-dominated-operator-tail} is
summable: almost surely, eventually
\[
  |\mathscr S_m|\le J_m,
  \qquad
  \sum_{h=1}^{R_m}\sum_{k\in\mathscr S_{m,h}}
  \mathfrak s_{m,h,k}\le C_0n_m .
\]
Consequently, suppose each candidate perturbation stream is tagged \(\psi\) or
\(\mu\), with parameters \(\varepsilon_{m,h,k},A_{m,h,k},s'_{m,h,k}\), and
satisfies the corresponding mixed dominated tail and drift hypotheses after
adjoining \(\mathcal H_m\).  If
\[
  \ell_m:=\log(e(D_m+n_m+1))\le B\sqrt{n_m}
\]
eventually,
\[
  \sum_m\sum_{h=1}^{R_m}\sum_{k\in\mathcal I_{m,h}}
  \bigl(\mathbb P(E^{\rm tail}_{m,h,k})+\varepsilon_{m,h,k}\bigr)<\infty,
\]
and
\[
  \sum_mJ_m(D_m+n_m+1)
  (e(D_m+n_m+1))^{-\gamma_{\rm tc}(K_{\rm tail},1,1,B)}<\infty,
\]
then the hybrid selected structured submenus have the same eventual common
signing, contraction normalization, fixed-menu coverage, and schedule thresholds
as in
Corollary~\ref{cor:selected-structured-submenu-mixed-dominated-operator-tail}.

Each component may be certified by any one of the preceding Bernstein,
Laplace-dominated, random-envelope, Poisson/Cox, categorical, Freedman, or
Serfling selected-submenu corollaries, or by a deterministic pathwise check.
No independence, negative association, or filtration relation between different
components is required beyond the displayed summability of the component
selection failures.
\end{corollary}

\begin{proof}
Let
\[
  F_m:=
  F_m^{\rm cert}\cup
  \bigcup_{h=1}^{R_m}F_{m,h}^{\rm sel}.
\]
By the assumed summability, \(\sum_m\mathbb P(F_m)<\infty\).  On \(F_m^c\),
the component bounds give
\[
  |\mathscr S_m|
  \le \sum_{h=1}^{R_m}|\mathscr S_{m,h}|
  \le J_m
\]
and
\[
  \sum_{h=1}^{R_m}\sum_{k\in\mathscr S_{m,h}}\mathfrak s_{m,h,k}
  \le
  \sum_{h=1}^{R_m}U_{m,h}
  \le C_0n_m .
\]
Thus Borel--Cantelli gives the eventual deterministic selected-submenu
cardinality and charge bounds.

On that eventual good event, order all selected labels first by component index
and then by the fixed catalogue order inside the component, padding by zero
streams up to \(J_m\).  The candidate tail summability dominates the selected
tail summability required by
Corollary~\ref{cor:selected-structured-submenu-mixed-dominated-operator-tail},
and the deterministic matrix-tail series is exactly the displayed one with
cardinality envelope \(J_m\).  Applying that corollary gives the stated signing,
normalization, coverage, and threshold conclusions.  The proof used only the
component good events, so cross-component dependence is irrelevant.
\end{proof}

\begin{corollary}[Random-budget hybrid selected structured submenus]
\label{cor:random-budget-hybrid-selected-structured-submenu-mixed-dominated-operator-tail}
In the setting and notation of
Corollary~\ref{cor:hybrid-selected-structured-submenu-mixed-dominated-operator-tail},
keep the candidate catalogue, component partition, selections
\(\mathscr S_{m,h}\), and filtration \(\mathcal H_m\).  Replace the
deterministic component envelopes by \(\mathcal G_m\)-measurable random
envelopes
\[
  \mathcal J_{m,h}\in\mathbb Z_{\ge0},
  \qquad
  \mathcal U_{m,h}\ge0 .
\]
Assume there are selection-failure events \(F_{m,h}^{\rm sel}\), aggregate
budget-failure events \(F_m^{\rm bud}\in\mathcal G_m\), and certificate-failure
events \(F_m^{\rm cert}\in\mathcal G_m\) such that
\[
  \sum_m\left[
    \mathbb P(F_m^{\rm bud})+\mathbb P(F_m^{\rm cert})
    +\sum_{h=1}^{R_m}\mathbb P(F_{m,h}^{\rm sel})
  \right]<\infty .
\]
Outside \(F_{m,h}^{\rm sel}\), the component selection obeys
\[
  |\mathscr S_{m,h}|\le\mathcal J_{m,h},
  \qquad
  \sum_{k\in\mathscr S_{m,h}}\mathfrak s_{m,h,k}
  \le \mathcal U_{m,h}.
\]
Outside \(F_m^{\rm bud}\), deterministic envelopes \(J_m\ge1\) and \(C_0n_m\)
certify the aggregate random budgets:
\[
  \sum_{h=1}^{R_m}\mathcal J_{m,h}\le J_m,
  \qquad
  \sum_{h=1}^{R_m}\mathcal U_{m,h}\le C_0n_m .
\]
Outside \(F_m^{\rm cert}\), every selected core is a contraction and has one of
the finite-lift orbit-closure or algebraic/group certificates from
Corollary~\ref{cor:selected-structured-submenu-mixed-dominated-operator-tail},
with charge bounded by \(\mathfrak s_{m,h,k}\).

Then the selection event in
Corollary~\ref{cor:selected-structured-submenu-mixed-dominated-operator-tail} is
summable: almost surely, eventually
\[
  |\mathscr S_m|\le J_m,
  \qquad
  \sum_{h=1}^{R_m}\sum_{k\in\mathscr S_{m,h}}
  \mathfrak s_{m,h,k}\le C_0n_m .
\]
Consequently, under the same mixed dominated tail, drift, candidate-tail
summability, and matrix-tail summability hypotheses as in
Corollary~\ref{cor:hybrid-selected-structured-submenu-mixed-dominated-operator-tail},
with this deterministic aggregate \(J_m\), the random-budget hybrid submenus
have the same eventual common signing, contraction normalization, fixed-menu
coverage, and schedule thresholds as in
Corollary~\ref{cor:selected-structured-submenu-mixed-dominated-operator-tail}.

Thus the mixture weights, posterior component masses, Cox intensities,
adaptive slot budgets, or stratum counts may be random and data-dependent.  It
is enough to certify their total realized cardinality and charge budgets
summably before applying the common Matrix Spencer signing.
\end{corollary}

\begin{proof}
Let
\[
  F_m:=
  F_m^{\rm bud}\cup F_m^{\rm cert}\cup
  \bigcup_{h=1}^{R_m}F_{m,h}^{\rm sel}.
\]
The displayed series gives \(\sum_m\mathbb P(F_m)<\infty\).  On \(F_m^c\),
\[
  |\mathscr S_m|
  \le\sum_{h=1}^{R_m}|\mathscr S_{m,h}|
  \le\sum_{h=1}^{R_m}\mathcal J_{m,h}
  \le J_m
\]
and
\[
  \sum_{h=1}^{R_m}\sum_{k\in\mathscr S_{m,h}}\mathfrak s_{m,h,k}
  \le
  \sum_{h=1}^{R_m}\mathcal U_{m,h}
  \le C_0n_m .
\]
Borel--Cantelli gives the eventual selected-submenu event.  Ordering selected
labels by component and catalogue order, then padding by zero streams up to
\(J_m\), reduces the conclusion to
Corollary~\ref{cor:selected-structured-submenu-mixed-dominated-operator-tail}
exactly as in
Corollary~\ref{cor:hybrid-selected-structured-submenu-mixed-dominated-operator-tail}.
\end{proof}

\begin{corollary}[Hierarchical random-budget selected structured submenus]
\label{cor:hierarchical-random-budget-selected-structured-submenu-mixed-dominated-operator-tail}
Fix structural constants \(r_0,b_0,c_0,L_0,C_0\ge1\).  For each \(m\), let
\(\mathfrak T_m\) be a countable rooted tree with a fixed depth-first order.
Each vertex \(v\in\mathfrak T_m\) carries a charge
\(\mathfrak s_{m,v}\ge0\), core stream
\((B_{m,v,i})_{i\le n_m}\), and perturbation stream
\((Y_{m,v,i})_{i\le n_m}\).  A random recursive procedure, possibly depending
on a base environment \(\mathcal G_m\), produces an
\(\mathcal H_m\)-measurable terminal antichain
\[
  \mathscr A_m\subset\mathfrak T_m,
  \qquad
  \mathcal H_m:=\mathcal G_m\vee\sigma(\mathscr A_m)
\]
such that no two selected vertices are in an ancestor--descendant relation.  The
fixed depth-first order is used only to list the selected vertices;
\(\mathscr A_m\) is the realized frontier of the random tree or adaptive search.

Assume that there are deterministic envelopes \(J_m\ge1\), tree-budget failure
events \(F_m^{\rm tree}\in\mathcal H_m\), and certificate-failure events
\(F_m^{\rm cert}\in\mathcal H_m\) such that
\[
  \sum_m\bigl(\mathbb P(F_m^{\rm tree})+\mathbb P(F_m^{\rm cert})\bigr)<\infty
\]
and, outside \(F_m^{\rm tree}\),
\[
  |\mathscr A_m|\le J_m,
  \qquad
  \sum_{v\in\mathscr A_m}\mathfrak s_{m,v}\le C_0n_m .
\]
Outside \(F_m^{\rm cert}\), every selected core is a contraction and has one of
the finite-lift orbit-closure or algebraic/group certificates from
Corollary~\ref{cor:selected-structured-submenu-mixed-dominated-operator-tail},
with charge bounded by \(\mathfrak s_{m,v}\).

For every candidate vertex \(v\), choose a tail type in \(\{\psi,\mu\}\).
Choose parameters
\[
  \varepsilon_{m,v},\qquad A_{m,v},\qquad s'_{m,v},
\]
and let \(q_{m,v}\) be the corresponding normalization from
Corollary~\ref{cor:mixed-dominated-operator-tail-active-libraries}.  Suppose
the conditional mixed dominated tail and truncated-drift hypotheses hold for
\((Y_{m,v,i})_{i\le n_m}\) after adjoining \(\mathcal H_m\), outside events
\(E^{\rm tail}_{m,v}\), and that
\[
  \ell_m:=\log(e(D_m+n_m+1))\le B\sqrt{n_m}
\]
eventually,
\[
  \sum_m\sum_{v\in\mathfrak T_m}
  \bigl(\mathbb P(E^{\rm tail}_{m,v})+\varepsilon_{m,v}\bigr)<\infty,
\]
and
\[
  \sum_mJ_m(D_m+n_m+1)
  (e(D_m+n_m+1))^{-\gamma_{\rm tc}(K_{\rm tail},1,1,B)}<\infty .
\]
Then the realized frontiers \(\mathscr A_m\) have the same eventual common
signing, contraction normalization, fixed-menu coverage, and all-integer and
sparse schedule thresholds as in
Corollary~\ref{cor:selected-structured-submenu-mixed-dominated-operator-tail},
with deterministic selected-cardinality envelope \(J_m\).

In particular, finite-depth or countably generated random trees, branching
Poisson/Cox catalogues, adaptive beam-search libraries, hierarchical mixtures,
random forests of structured cores, and recursive random basis/corner choices
are covered whenever their realized frontier size and charge have a summable
budget certificate.  No independence between generations, siblings, or
different branches is required beyond the displayed budget and candidate-tail
summability.
\end{corollary}

\begin{proof}
On \((F_m^{\rm tree}\cup F_m^{\rm cert})^c\), the terminal frontier
\(\mathscr A_m\) is an \(\mathcal H_m\)-measurable finite selected set from the
countable catalogue \(\mathfrak T_m\), ordered by the fixed depth-first order.
It satisfies the deterministic cardinality and charge bounds required by
Corollary~\ref{cor:selected-structured-submenu-mixed-dominated-operator-tail},
and its selected cores have the required certificates.  The displayed
summability gives \(\sum_m\mathbb P(F_m^{\rm tree}\cup F_m^{\rm cert})<\infty\),
so Borel--Cantelli gives the selected-submenu structural event eventually.

The candidate-tail series over the whole tree dominates the selected-tail
series after ordering the frontier and padding by zero streams up to \(J_m\).
The deterministic matrix-tail series is exactly the selected-submenu series with
cardinality envelope \(J_m\).  Applying
Corollary~\ref{cor:selected-structured-submenu-mixed-dominated-operator-tail}
therefore gives the claimed common signing and threshold conclusions.  The
antichain condition only records that the selected labels are terminal models;
the proof uses no independence between the recursive branching decisions.
\end{proof}

\begin{proposition}[Square-function Matrix Spencer certificate]
\label{prop:square-function-khintchine}
There is a universal constant \(C\) such that, for every self-adjoint tuple
\(A_1,\ldots,A_n\in\M_D(\C)\),
\[
  \disc(A_1,\ldots,A_n)
  \le
  C\sqrt{\log(eD)}
  \left\|\sum_{i=1}^n A_i^2\right\|_{\op}^{1/2}.
\]
More precisely, if
\[
  V:=\sum_{i=1}^nA_i^2,\qquad
  \rho(A):=\max\{1,\operatorname{rank}V\},
\]
then the ambient dimension may be replaced by the realized joint square-function
rank:
\[
  \disc(A_1,\ldots,A_n)
  \le
  C\sqrt{\log(e\rho(A))}
  \left\|\sum_{i=1}^n A_i^2\right\|_{\op}^{1/2}.
\]
Consequently, if \(A_1,\ldots,A_n\) are random self-adjoint contractions and
\[
  F_{C_0}:=
  \left\{
  \left\|\sum_{i=1}^n A_i^2\right\|_{\op}
  \le \frac{C_0 n}{\log(eD)}
  \right\},
\]
then
\[
  \mathbb P\left\{
  \disc(A_1,\ldots,A_n)>C\sqrt{C_0n}
  \right\}
  \le \mathbb P(F_{C_0}^{\,c}).
\]
Thus any random-matrix model whose realized noncommutative variance profile is
at most \(O(n/\log(eD))\) is in the Matrix Spencer regime, independently of the
ambient dimension \(D\).
\end{proposition}

\begin{proof}
Put \(V=\sum_iA_i^2\) and \(E=\operatorname{ran}V\).  If \(x\in\ker V\), then
\[
  0=\langle Vx,x\rangle=\sum_i\|A_ix\|^2,
\]
so every \(A_i\) kills \(\ker V\).  Since the \(A_i\)'s are self-adjoint, \(E\)
is invariant for every \(A_i\).  Restricting the tuple to \(E\) and applying the
finite-dimensional noncommutative Khintchine inequality
\cite{LustPisier,Pisier} in dimension \(\rho(A)\) gives, for independent
Rademacher signs \(\delta_1,\ldots,\delta_n\),
\[
  \mathbb E_\delta
  \left\|\sum_{i=1}^n\delta_i A_i\right\|_{\op}
  \le
  C\sqrt{\log(e\rho(A))}
  \left\|\left(\sum_{i=1}^nA_i^2\right)^{1/2}\right\|_{\op}.
\]
The minimum over sign vectors is bounded by this expectation, which proves the
rank-sensitive deterministic estimate.  Since \(\rho(A)\le D\), this also gives
the displayed ambient-dimensional form.  The probability statement is the same
bound applied pathwise on the event \(F_{C_0}\).
\end{proof}

\begin{corollary}[Random sign search under a square profile]
\label{cor:random-sign-search-square-profile}
There is a universal constant \(C\) with the following property.  Let
\(A_1,\ldots,A_n\in\M_D(\C)\) be self-adjoint, and let
\(\eta^{(1)},\ldots,\eta^{(m)}\) be independent uniform sign vectors in
\(\{\pm1\}^n\).  Then, for every \(m\ge1\),
\[
  \mathbb P_\eta\left\{
  \min_{1\le r\le m}
  \left\|\sum_{i=1}^n\eta_i^{(r)}A_i\right\|_{\op}
  >
  2C\sqrt{\log(e\rho(A))}
  \left\|\sum_{i=1}^nA_i^2\right\|_{\op}^{1/2}
  \right\}
  \le 2^{-m},
\]
where \(\rho(A)=\max\{1,\operatorname{rank}(\sum_iA_i^2)\}\).  In particular,
\[
  \mathbb P_\eta\left\{
  \min_{1\le r\le m}
  \left\|\sum_{i=1}^n\eta_i^{(r)}A_i\right\|_{\op}
  >
  2C\sqrt{\log(eD)}
  \left\|\sum_{i=1}^nA_i^2\right\|_{\op}^{1/2}
  \right\}
  \le 2^{-m}.
\]
Consequently, if \(A_1,\ldots,A_n\) are random self-adjoint contractions,
independent of the search signs, and \(F_{C_0}\) is the square-profile event
from Proposition~\ref{prop:square-function-khintchine}, then
\[
  \mathbb P\left\{
  \min_{1\le r\le m}
  \left\|\sum_{i=1}^n\eta_i^{(r)}A_i\right\|_{\op}
  >
  2C\sqrt{C_0n}
  \right\}
  \le
  \mathbb P(F_{C_0}^{\,c})+2^{-m}.
\]
Thus every random ensemble covered by the square-profile route also has a
simple randomized signing certificate: \(m\) independent trials lose only an
additional \(2^{-m}\) in failure probability.
\end{corollary}

\begin{proof}
For one uniform sign vector \(\eta\), the noncommutative Khintchine estimate
used in the rank-sensitive part of
Proposition~\ref{prop:square-function-khintchine} gives
\[
  \mathbb E_\eta
  \left\|\sum_i\eta_iA_i\right\|_{\op}
  \le
  C\sqrt{\log(e\rho(A))}
  \left\|\sum_iA_i^2\right\|_{\op}^{1/2}.
\]
Markov's inequality therefore makes one trial successful with probability at
least \(1/2\).  The \(m\) trials are independent, so all of them fail with
probability at most \(2^{-m}\).  Since \(\rho(A)\le D\), the displayed
ambient-dimensional search bound follows.  The random-tuple statement follows
by conditioning on \(A_1,\ldots,A_n\) and then restricting to \(F_{C_0}\).
\end{proof}

\begin{corollary}[Common signs for self-adjoint square-profile menus]
\label{cor:common-sign-selfadjoint-square-profile-menu}
Fix \(C_0\ge1\).  For each \(m\ge1\), let \(J_m\ge1\), \(D_m,n_m\ge1\), and
\[
  D_m^{\rm com}:=J_mD_m.
\]
For every \(1\le h\le J_m\), let
\[
  A_{m,h,1},\ldots,A_{m,h,n_m}\in\M_{D_m}(\C)
\]
be arbitrary random self-adjoint matrices.  Assume that there are events
\(E_{m,h}^{\rm sq}\) such that, outside \(E_{m,h}^{\rm sq}\),
\[
  \left\|\sum_{i=1}^{n_m}A_{m,h,i}^2\right\|_{\op}
  \le
  \frac{C_0n_m}{\log(eD_m^{\rm com})}.
\]
If
\[
  \sum_m\sum_{h=1}^{J_m}\mathbb P(E_{m,h}^{\rm sq})<\infty ,
\]
then, almost surely, for all sufficiently large \(m\), one sign vector
\(\eps^{(m)}\in\{\pm1\}^{n_m}\) balances the whole active self-adjoint menu:
\[
  \max_{1\le h\le J_m}
  \left\|\sum_{i=1}^{n_m}\eps_i^{(m)}A_{m,h,i}\right\|_{\op}
  \le C(C_0)\sqrt{n_m}.
\]
If \(M_m\) independent random sign vectors are sampled for the whole active
menu and
\[
  \sum_m2^{-M_m}<\infty ,
\]
then the sampled list almost surely eventually contains such a common
certificate.  If \(J_m\to\infty\), active prefixes may exhaust a countable
self-adjoint random-matrix library, and every fixed library element is
eventually covered by the common active-prefix signing.

Thus any self-adjoint active library whose realized noncommutative variance
profiles are summably certified at the common-sign scale has a single
simultaneous Matrix-Spencer signing.
\end{corollary}

\begin{proof}
Outside the union of the events \(E_{m,h}^{\rm sq}\), set
\[
  \widetilde A_{m,i}:=
  \bigoplus_{h=1}^{J_m}A_{m,h,i}\in\M_{D_m^{\rm com}}(\C).
\]
Then
\[
  \left\|\sum_i\widetilde A_{m,i}^2\right\|_{\op}
  =
  \max_{h\le J_m}
  \left\|\sum_iA_{m,h,i}^2\right\|_{\op}
  \le
  \frac{C_0n_m}{\log(eD_m^{\rm com})}.
\]
Proposition~\ref{prop:square-function-khintchine} applied in
\(\M_{D_m^{\rm com}}(\C)\) gives one signing of the block-diagonal tuple, which
is exactly one sign vector balancing every active stream.  Since
\[
  \sum_m\mathbb P\left(\bigcup_{h\le J_m}E_{m,h}^{\rm sq}\right)
  \le
  \sum_m\sum_{h\le J_m}\mathbb P(E_{m,h}^{\rm sq})<\infty ,
\]
Borel--Cantelli gives the eventual deterministic certificates.  The sampled
statement follows from Corollary~\ref{cor:random-sign-search-square-profile}
applied to the same block-diagonal tuple, adding only
\(\sum_m2^{-M_m}\).  If \(J_m\to\infty\), each fixed library stream belongs to
all sufficiently large active prefixes.
\end{proof}

\begin{corollary}[Common signs for rank-sensitive self-adjoint square-profile menus]
\label{cor:common-sign-rank-sensitive-selfadjoint-square-profile-menu}
Fix \(C_0\ge1\).  For each \(m\ge1\), let \(J_m\ge1\),
\(D_m,n_m\ge1\), and let \(R_m\ge1\).  For every \(1\le h\le J_m\), let
\[
  A_{m,h,1},\ldots,A_{m,h,n_m}\in\M_{D_m}(\C)
\]
be arbitrary random self-adjoint matrices, and put
\[
  V_{m,h}:=\sum_{i=1}^{n_m}A_{m,h,i}^2,\qquad
  \rho_{m,h}^{\rm sq}:=\max\{1,\operatorname{rank}V_{m,h}\}.
\]
Assume that there are events \(E_m^{\rm rsq}\) such that, outside
\(E_m^{\rm rsq}\),
\[
  \sum_{h=1}^{J_m}\rho_{m,h}^{\rm sq}\le R_m,
  \qquad
  \max_{1\le h\le J_m}\|V_{m,h}\|_{\op}
  \le
  \frac{C_0n_m}{\log(eR_m)}.
\]
If
\[
  \sum_m\mathbb P(E_m^{\rm rsq})<\infty ,
\]
then, almost surely, for all sufficiently large \(m\), one sign vector
\(\eps^{(m)}\in\{\pm1\}^{n_m}\) balances the whole active menu:
\[
  \max_{1\le h\le J_m}
  \left\|\sum_{i=1}^{n_m}\eps_i^{(m)}A_{m,h,i}\right\|_{\op}
  \le C(C_0)\sqrt{n_m}.
\]
If \(M_m\) independent random sign vectors are sampled for the whole active
menu and
\[
  \sum_m2^{-M_m}<\infty ,
\]
then the sampled list almost surely eventually contains such a common
certificate.

In particular, if each active stream is supported on a possibly random
\(\kappa_{m,h}\)-dimensional corner, or more generally
\(\operatorname{rank}V_{m,h}\le\kappa_{m,h}\), one may take
\[
  R_m=\sum_{h=1}^{J_m}\max\{1,\kappa_{m,h}\}.
\]
Thus common signs for low-rank active self-adjoint libraries pay the logarithm
of total active square-function rank, not the ambient block dimension
\(J_mD_m\).
\end{corollary}

\begin{proof}
On \((E_m^{\rm rsq})^c\), define the block-diagonal tuple
\[
  \widetilde A_{m,i}:=
  \bigoplus_{h=1}^{J_m}A_{m,h,i}\in\M_{J_mD_m}(\C).
\]
Its square function is
\[
  \widetilde V_m:=\sum_i\widetilde A_{m,i}^2
  =
  \bigoplus_{h=1}^{J_m}V_{m,h}.
\]
Hence
\[
  \|\widetilde V_m\|_{\op}
  =
  \max_{h\le J_m}\|V_{m,h}\|_{\op}
  \le
  \frac{C_0n_m}{\log(eR_m)}
\]
and, because \(R_m\ge1\),
\[
  \max\{1,\operatorname{rank}\widetilde V_m\}
  \le
  \sum_{h=1}^{J_m}\rho_{m,h}^{\rm sq}
  \le R_m .
\]
The rank-sensitive part of
Proposition~\ref{prop:square-function-khintchine} gives
\[
  \disc(\widetilde A_{m,1},\ldots,\widetilde A_{m,n_m})
  \le
  C\sqrt{\log(eR_m)}
  \|\widetilde V_m\|_{\op}^{1/2}
  \le C(C_0)\sqrt{n_m}.
\]
A signing of the block-diagonal tuple is exactly one sign vector balancing
every active stream.  The summability of \(\mathbb P(E_m^{\rm rsq})\) and
Borel--Cantelli give the eventual almost-sure statement.

For sampled signs, condition on the active menu on \((E_m^{\rm rsq})^c\) and
apply Corollary~\ref{cor:random-sign-search-square-profile} to the same
block-diagonal tuple in its rank-sensitive form.  The sampled failure for the
whole menu is \(2^{-M_m}\), so the displayed summability and Borel--Cantelli
complete the proof.  The final specialization follows from
\(\operatorname{rank}V_{m,h}\le\kappa_{m,h}\).
\end{proof}

\begin{corollary}[Rectangular square-profile certificate]
\label{cor:rectangular-square-profile-certificate}
Let \(X_1,\ldots,X_n\in\M_{d_1\times d_2}(\C)\).  Put
\[
  V_L:=\sum_{i=1}^nX_iX_i^*,
  \qquad
  V_R:=\sum_{i=1}^nX_i^*X_i,
  \qquad
  D_{\rm rec}:=d_1+d_2,
\]
and
\[
  \rho_{\rm rec}:=
  \max\{1,\operatorname{rank}V_L+\operatorname{rank}V_R\}.
\]
Then
\[
  \min_{\eps\in\{\pm1\}^n}
  \left\|\sum_{i=1}^n\eps_iX_i\right\|_{\op}
  \le
  C\sqrt{\log(e\rho_{\rm rec})}
  \max\{\|V_L\|_{\op},\|V_R\|_{\op}\}^{1/2}.
\]
In particular the rank factor may always be replaced by \(D_{\rm rec}\).

Moreover, if \(\eta^{(1)},\ldots,\eta^{(m)}\) are independent uniform sign
vectors, then
\[
  \mathbb P_\eta\left\{
  \min_{1\le r\le m}
  \left\|\sum_{i=1}^n\eta_i^{(r)}X_i\right\|_{\op}
  >
  2C\sqrt{\log(e\rho_{\rm rec})}
  \max\{\|V_L\|_{\op},\|V_R\|_{\op}\}^{1/2}
  \right\}
  \le 2^{-m}.
\]
Consequently, rectangular random contractions are in the Matrix-Spencer scale
whenever their realized left and right square profiles are both
\(O(n/\log(eD_{\rm rec}))\).
\end{corollary}

\begin{proof}
Apply the Hermitian dilation
\[
  \mathscr H(X_i)=
  \begin{pmatrix}0&X_i\\ X_i^*&0\end{pmatrix}
  \in\M_{D_{\rm rec}}(\C).
\]
Then
\[
  \mathscr H(X_i)^2=
  \begin{pmatrix}X_iX_i^*&0\\0&X_i^*X_i\end{pmatrix},
\]
so
\[
  \left\|\sum_i\mathscr H(X_i)^2\right\|_{\op}
  =
  \max\{\|V_L\|_{\op},\|V_R\|_{\op}\},
\]
and the rank of this square function is at most
\(\operatorname{rank}V_L+\operatorname{rank}V_R\).  Also
\[
  \left\|\sum_i\eps_i\mathscr H(X_i)\right\|_{\op}
  =
  \left\|\sum_i\eps_iX_i\right\|_{\op}.
\]
The deterministic and randomized claims are therefore exactly
Proposition~\ref{prop:square-function-khintchine} and
Corollary~\ref{cor:random-sign-search-square-profile} applied to the dilated
self-adjoint tuple.
\end{proof}

\begin{corollary}[Bernstein square-profile ensembles]
\label{cor:random-square-profile-bernstein}
Let \(D\ge2\), and let \(A_1,\ldots,A_n\in\M_D(\C)\) be independent random
self-adjoint contractions.  Put
\[
  \mu:=\left\|\sum_{i=1}^n \mathbb E A_i^2\right\|_{\op}.
\]
Then, for every \(C_0\ge1\), there is \(C(C_0)\) such that, for every
\(t>0\) satisfying
\[
  \mu+t\le \frac{C_0n}{\log(eD)},
\]
one has
\[
  \mathbb P\left\{
  \disc(A_1,\ldots,A_n)>C(C_0)\sqrt n
  \right\}
  \le
  D\exp\left\{-\frac{t^2}{2(\mu+t/3)}\right\}.
\]
In particular, if \(\mu\le C_0n/(2\log(eD))\), then
\[
  \mathbb P\left\{
  \disc(A_1,\ldots,A_n)>C(C_0)\sqrt n
  \right\}
  \le
  D\exp\left\{-\frac{3C_0n}{16\log(eD)}\right\}.
\]

Consequently, let \(T_1,\ldots,T_n\) be arbitrary, possibly random,
self-adjoint contractions, let \(U_1,\ldots,U_n\) be independent Haar unitary
matrices independent of the \(T_i\), and set \(A_i=U_iT_iU_i^*\).  If
\[
  G:=\left\{
  \frac1D\sum_{i=1}^n\operatorname{tr}(T_i^2)
  \le \frac{C_0n}{2\log(eD)}
  \right\},
\]
then
\[
  \mathbb P\left\{
  \disc(A_1,\ldots,A_n)>C(C_0)\sqrt n
  \right\}
  \le
  \mathbb P(G^c)+
  D\exp\left\{-\frac{3C_0n}{16\log(eD)}\right\}.
\]
Thus independent Haar-random rotations of fixed spectra, or of conditionally
random spectra, satisfy Matrix Spencer whenever their average Frobenius-square
energy is \(O(D/\log(eD))\).  This includes low-stable-rank and soft-spectrum
full-rank contractions, not only hard subspace-supported matrices.
\end{corollary}

\begin{proof}
Write \(B_i=A_i^2\).  Since \(A_i\) is a self-adjoint contraction,
\(0\le B_i\le I\).  Let \(X_i=B_i-\mathbb E B_i\).  Then
\(\mathbb E X_i=0\), \(\lambda_{\max}(X_i)\le1\), and
\[
  \sum_i\mathbb E X_i^2
  =
  \sum_i\left(\mathbb E B_i^2-(\mathbb E B_i)^2\right)
  \le
  \sum_i\mathbb E B_i^2
  \le
  \sum_i\mathbb E B_i .
\]
Hence the variance parameter in the one-sided matrix Bernstein inequality
\cite{Tropp} is at most \(\mu\).  Therefore
\[
\begin{aligned}
  \mathbb P\left\{
  \left\|\sum_i A_i^2\right\|_{\op}>\mu+t
  \right\}
  &\le
  \mathbb P\left\{
  \lambda_{\max}\left(\sum_iX_i\right)>t
  \right\} \\
  &\le
  D\exp\left\{-\frac{t^2}{2(\mu+t/3)}\right\}.
\end{aligned}
\]
On the complementary event, Proposition~\ref{prop:square-function-khintchine}
gives the discrepancy bound.  If \(\mu\le C_0n/(2\log(eD))\), take
\(t=C_0n/(2\log(eD))\), so that \(\mu\le t\) and
\[
  \frac{t^2}{2(\mu+t/3)}
  \ge
  \frac{t^2}{2(t+t/3)}
  =
  \frac{3C_0n}{16\log(eD)}.
\]

For the Haar-rotation specialization, condition on \(T_1,\ldots,T_n\).  Haar
invariance gives
\[
  \mathbb E(A_i^2\mid T_i)
  =
  \mathbb E(U_iT_i^2U_i^*\mid T_i)
  =
  \frac{\operatorname{tr}(T_i^2)}{D}I.
\]
Thus, on \(G\), the conditional value of \(\mu\) is at most
\(C_0n/(2\log(eD))\).  Applying the preceding bound conditionally and then
averaging proves the displayed estimate.
\end{proof}

\begin{corollary}[Conditional square-twirl ensembles]
\label{cor:conditional-square-twirl}
Let \(D\ge2\).  Let \(T_1,\ldots,T_n\in\M_D(\C)\) be arbitrary, possibly
dependent, random self-adjoint contractions.  Conditional on
\(T=(T_1,\ldots,T_n)\), let \(U_1,\ldots,U_n\) be independent random unitaries
and set
\[
  A_i=U_iT_iU_i^* .
\]
Define the conditional square profile
\[
  \mu(T):=
  \left\|
  \sum_{i=1}^n
  \mathbb E\left(U_iT_i^2U_i^*\,\middle|\,T\right)
  \right\|_{\op}.
\]
Then, for every \(C_0\ge1\) there is \(C(C_0)\) such that, whenever
\(v,u>0\) satisfy
\[
  v+u\le \frac{C_0n}{\log(eD)},
\]
one has
\[
  \mathbb P\left\{
  \disc(A_1,\ldots,A_n)>C(C_0)\sqrt n
  \right\}
  \le
  \mathbb P\{\mu(T)>v\}
  +
  D\exp\left\{-\frac{u^2}{2(v+u/3)}\right\}.
\]
Consequently, any independent random conjugation scheme whose conditional
square twirl has operator norm \(O(n/\log(eD))\) is Spencer-scale with the
displayed high probability.  The conjugating laws may be Haar, finite, or
sample-dependent; only the conditional independence and the square profile
enter the proof.
\end{corollary}

\begin{proof}
Condition on \(T\).  The matrices \(A_i\) are then independent self-adjoint
contractions and
\[
  \left\|\sum_i\mathbb E(A_i^2\mid T)\right\|_{\op}=\mu(T).
\]
On the event \(\{\mu(T)\le v\}\), the proof of
Corollary~\ref{cor:random-square-profile-bernstein}, applied conditionally
with \(t=u\), gives conditional failure probability at most
\[
  D\exp\left\{-\frac{u^2}{2(v+u/3)}\right\}.
\]
Averaging over \(T\) and adding the event \(\{\mu(T)>v\}\) proves the claim.
\end{proof}

\begin{corollary}[Finite-design and signed-permutation rotations]
\label{cor:finite-design-rotations}
Let \(D\ge2\), and let \(T_1,\ldots,T_n\in\M_D(\C)\) be arbitrary, possibly
dependent, random self-adjoint contractions.  Conditional on
\(T=(T_1,\ldots,T_n)\), let \(U_1,\ldots,U_n\) be independent random unitaries
such that, for every deterministic matrix \(X\in\M_D(\C)\),
\[
  \mathbb E\left(U_iXU_i^*\,\middle|\,T\right)
  =
  \frac{\operatorname{tr}X}{D}I .
\]
Set \(A_i=U_iT_iU_i^*\) and
\[
  \Phi(T):=\frac1D\sum_{i=1}^n\operatorname{tr}(T_i^2).
\]
Then, for every \(C_0\ge1\) there is \(C(C_0)\) such that, whenever
\(v,u>0\) satisfy \(v+u\le C_0n/\log(eD)\),
\[
  \mathbb P\left\{
  \disc(A_1,\ldots,A_n)>C(C_0)\sqrt n
  \right\}
  \le
  \mathbb P\{\Phi(T)>v\}
  +
  D\exp\left\{-\frac{u^2}{2(v+u/3)}\right\}.
\]
In particular, if \(\Phi(T)\le C_0n/(2\log(eD))\) outside an event \(E\), then
\[
  \mathbb P\left\{
  \disc(A_1,\ldots,A_n)>C(C_0)\sqrt n
  \right\}
  \le
  \mathbb P(E)+
  D\exp\left\{-\frac{3C_0n}{16\log(eD)}\right\}.
\]
This applies not only to Haar unitary or orthogonal rotations, but also to
uniform draws from any finite irreducible unitary group, to exact conjugation
unitary \(1\)-designs, and in particular to the signed-permutation/monomial
twirl.  Thus discrete random bases give the same Matrix Spencer conclusion as
Haar bases under the same average Frobenius-square profile.
\end{corollary}

\begin{proof}
The displayed twirl identity gives
\[
  \sum_i
  \mathbb E\left(U_iT_i^2U_i^*\,\middle|\,T\right)
  =
  \frac1D\sum_i\operatorname{tr}(T_i^2)I,
\]
so the conditional square profile in
Corollary~\ref{cor:conditional-square-twirl} is \(\Phi(T)\).  That corollary
gives the first estimate, and the second follows by taking
\(v=u=C_0n/(2\log(eD))\).

For a finite irreducible unitary group the twirl identity is Schur's lemma.
For the signed-permutation group it is also immediate: random signs kill the
off-diagonal entries of \(X\), and the random permutation averages the diagonal
entries to \((\operatorname{tr}X/D)I\).
\end{proof}

\begin{corollary}[Block-design twirl profile]
\label{cor:block-design-twirl}
Let \(D\ge2\).  Let \(\mcG\) be an environment, and let
\(T_1,\ldots,T_n\in\M_D(\C)\) be \(\mcG\)-measurable self-adjoint
contractions.  For each \(i\), let
\[
  I=\sum_{\alpha\in I_i}P_{i,\alpha}
\]
be a \(\mcG\)-measurable orthogonal decomposition, and put
\(d_{i,\alpha}:=\operatorname{rank}P_{i,\alpha}\).  Conditional on \(\mcG\),
let \(U_1,\ldots,U_n\) be independent random unitaries satisfying the block
twirl identity
\[
  \mathbb E\left(U_iXU_i^*\,\middle|\,\mcG\right)
  =
  \sum_{\alpha\in I_i}
  \frac{\operatorname{tr}(P_{i,\alpha}XP_{i,\alpha})}{d_{i,\alpha}}
  P_{i,\alpha}
  \qquad(X\in\M_D(\C)).
\]
Set \(A_i=U_iT_iU_i^*\), and define the block square profile
\[
  \Psi(\mcG):=
  \left\|
  \sum_{i=1}^n\sum_{\alpha\in I_i}
  \frac{\operatorname{tr}(P_{i,\alpha}T_i^2P_{i,\alpha})}{d_{i,\alpha}}
  P_{i,\alpha}
  \right\|_{\op}.
\]
Then, for every \(C_0\ge1\) there is \(C(C_0)\) such that, whenever
\(v,u>0\) satisfy \(v+u\le C_0n/\log(eD)\),
\[
  \mathbb P\left\{
  \disc(A_1,\ldots,A_n)>C(C_0)\sqrt n
  \right\}
  \le
  \mathbb P\{\Psi(\mcG)>v\}
  +
  D\exp\left\{-\frac{u^2}{2(v+u/3)}\right\}.
\]
In particular, if \(\Psi(\mcG)\le C_0n/(2\log(eD))\) outside an event \(E\),
then
\[
  \mathbb P\left\{
  \disc(A_1,\ldots,A_n)>C(C_0)\sqrt n
  \right\}
  \le
  \mathbb P(E)+
  D\exp\left\{-\frac{3C_0n}{16\log(eD)}\right\}.
\]
Thus independent block-Haar rotations, block finite-design rotations, and
block signed-permutation or monomial randomizations satisfy Matrix Spencer
whenever their block-energy profile is \(O(n/\log(eD))\).  The decomposition
may vary with \(i\), and may itself be random as part of the environment.
\end{corollary}

\begin{proof}
Apply the block twirl identity with \(X=T_i^2\).  Since
\(A_i^2=U_iT_i^2U_i^*\), the conditional square profile in
Corollary~\ref{cor:conditional-square-twirl} is exactly \(\Psi(\mcG)\).  That
corollary gives the first estimate; the second follows by taking
\(v=u=C_0n/(2\log(eD))\).

The listed examples satisfy the displayed identity block by block.  Haar
measure on each block and uniform draws from finite irreducible unitary groups
give the scalar block expectation by Schur's lemma; signed-permutation or
monomial randomization inside a block first kills off-diagonal entries and then
averages the diagonal entries over the block.
\end{proof}

\begin{corollary}[Dominated approximate twirl profiles]
\label{cor:dominated-approximate-twirl}
Let \(D\ge2\).  Let \(\mcG\) be an environment, and let
\(T_1,\ldots,T_n\in\M_D(\C)\) be \(\mcG\)-measurable self-adjoint
contractions.  Conditional on \(\mcG\), let \(U_1,\ldots,U_n\) be independent
random unitaries and set \(A_i=U_iT_iU_i^*\).  Suppose that there are
\(\mcG\)-measurable positive semidefinite matrices \(V_1,\ldots,V_n\) such
that
\[
  \mathbb E\left(U_iT_i^2U_i^*\,\middle|\,\mcG\right)\preceq V_i
  \qquad(1\le i\le n).
\]
Put
\[
  \Xi(\mcG):=\left\|\sum_{i=1}^nV_i\right\|_{\op}.
\]
Then, for every \(C_0\ge1\) there is \(C(C_0)\) such that, whenever
\(v,u>0\) satisfy \(v+u\le C_0n/\log(eD)\),
\[
  \mathbb P\left\{
  \disc(A_1,\ldots,A_n)>C(C_0)\sqrt n
  \right\}
  \le
  \mathbb P\{\Xi(\mcG)>v\}
  +
  D\exp\left\{-\frac{u^2}{2(v+u/3)}\right\}.
\]
Consequently, if \(\Xi(\mcG)\le C_0n/(2\log(eD))\) outside an event \(E\),
then
\[
  \mathbb P\left\{
  \disc(A_1,\ldots,A_n)>C(C_0)\sqrt n
  \right\}
  \le
  \mathbb P(E)+
  D\exp\left\{-\frac{3C_0n}{16\log(eD)}\right\}.
\]

Two useful specializations make the approximate-design cost explicit.
First, if \(\Gamma_i\succeq0\) are \(\mcG\)-measurable and, for every
\(0\preceq X\preceq I\),
\[
  \mathbb E\left(U_iXU_i^*\,\middle|\,\mcG\right)
  \preceq
  \frac{\operatorname{tr}X}{D}I+\Gamma_i ,
\]
then the theorem applies with
\[
  V_i=\frac{\operatorname{tr}(T_i^2)}{D}I+\Gamma_i .
\]
Thus approximate scalar unitary \(1\)-designs are Spencer-scale whenever
\[
  \left\|
  \frac1D\sum_i\operatorname{tr}(T_i^2)I+\sum_i\Gamma_i
  \right\|_{\op}
  =O\!\left(\frac{n}{\log(eD)}\right).
\]
Second, for \(\mcG\)-measurable decompositions
\(I=\sum_{\alpha\in I_i}P_{i,\alpha}\), if
\[
  \mathbb E\left(U_iXU_i^*\,\middle|\,\mcG\right)
  \preceq
  \sum_{\alpha\in I_i}
  \frac{\operatorname{tr}(P_{i,\alpha}XP_{i,\alpha})}{d_{i,\alpha}}
  P_{i,\alpha}
  +\Gamma_i
  \qquad(0\preceq X\preceq I),
\]
then the theorem applies with the block profile from
Corollary~\ref{cor:block-design-twirl} plus the aggregate error
\(\sum_i\Gamma_i\).  Hence approximate block designs and finite random-basis
samplers only pay their total positive semidefinite twirl error.
\end{corollary}

\begin{proof}
Since
\[
  \sum_i\mathbb E(A_i^2\mid\mcG)
  =
  \sum_i\mathbb E(U_iT_i^2U_i^*\mid\mcG)
  \preceq
  \sum_iV_i,
\]
the conditional square profile in
Corollary~\ref{cor:conditional-square-twirl} is at most \(\Xi(\mcG)\).  Applying
that corollary gives the finite-probability bounds.  The two displayed
specializations follow by taking \(X=T_i^2\), which satisfies
\(0\preceq T_i^2\preceq I\), in the corresponding approximate twirl
inequality.
\end{proof}

\begin{corollary}[Predictable square-profile ensembles]
\label{cor:random-square-profile-freedman}
Let \(D\ge2\), let \((\mcF_i)_{i=0}^n\) be a filtration, and let
\(A_1,\ldots,A_n\in\M_D(\C)\) be self-adjoint contractions such that \(A_i\) is
\(\mcF_i\)-measurable.  Put
\[
  \nu:=
  \left\|
  \sum_{i=1}^n \mathbb E(A_i^2\mid\mcF_{i-1})
  \right\|_{\op}.
\]
Then, for every \(C_0\ge1\), there is \(C(C_0)\) such that, for every \(v,t>0\)
satisfying
\[
  v+t\le \frac{C_0n}{\log(eD)},
\]
one has
\[
  \mathbb P\left\{
  \disc(A_1,\ldots,A_n)>C(C_0)\sqrt n
  \right\}
  \le
  \mathbb P\{\nu>v\}
  +
  D\exp\left\{-\frac{t^2}{2(v+t/3)}\right\}.
\]
Moreover, if \(\eta^{(1)},\ldots,\eta^{(m)}\) are independent uniform sign
vectors, independent of \(A_1,\ldots,A_n\), then, with possibly larger
\(C(C_0)\),
\[
  \begin{aligned}
  &\mathbb P\left\{
  \min_{1\le r\le m}
  \left\|\sum_{i=1}^n\eta_i^{(r)}A_i\right\|_{\op}
  >C(C_0)\sqrt n
  \right\} \\
  &\qquad\le
  \mathbb P\{\nu>v\}
  +
  D\exp\left\{-\frac{t^2}{2(v+t/3)}\right\}
  +2^{-m}.
  \end{aligned}
\]
In particular, if \(\nu\le C_0n/(2\log(eD))\) almost surely, then both
statements hold with failure at most
\[
  D\exp\left\{-\frac{3C_0n}{16\log(eD)}\right\}
\]
and, for the randomized signing search, with the additional \(2^{-m}\) term.
Thus the square-profile route covers adapted dependent random matrix streams,
not only independent ensembles, whenever their predictable noncommutative
variance profile is \(O(n/\log(eD))\).
\end{corollary}

\begin{proof}
Set \(B_i=A_i^2\) and
\(M_i=\mathbb E(B_i\mid\mcF_{i-1})\).  Then \(0\le B_i\le I\), and
\[
  X_i:=B_i-M_i
\]
is a self-adjoint matrix martingale difference with
\(\lambda_{\max}(X_i)\le1\).  Its predictable quadratic variation satisfies
\[
  \sum_i\mathbb E(X_i^2\mid\mcF_{i-1})
  =
  \sum_i\left(\mathbb E(B_i^2\mid\mcF_{i-1})-M_i^2\right)
  \le
  \sum_i\mathbb E(B_i^2\mid\mcF_{i-1})
  \le
  \sum_i M_i,
\]
because \(0\le B_i\le I\) implies \(B_i^2\le B_i\).  On the event
\(\{\nu\le v\}\), the variance process is therefore bounded by \(v\).  The
matrix Freedman inequality \cite{Tropp} gives
\[
  \mathbb P\left\{
  \lambda_{\max}\left(\sum_iX_i\right)>t
  \ \text{and}\ \nu\le v
  \right\}
  \le
  D\exp\left\{-\frac{t^2}{2(v+t/3)}\right\}.
\]
On the complementary high-probability event,
\[
  \left\|\sum_i A_i^2\right\|_{\op}
  =
  \lambda_{\max}\left(\sum_i B_i\right)
  \le v+t.
\]
The first claim follows from
Proposition~\ref{prop:square-function-khintchine}.  For the random-search
claim, condition on the realized tuple and apply
Corollary~\ref{cor:random-sign-search-square-profile} on the event
\(\|\sum_i A_i^2\|_{\op}\le v+t\).  The final displayed specialization is the
choice \(v=t=C_0n/(2\log(eD))\), as in
Corollary~\ref{cor:random-square-profile-bernstein}.
\end{proof}

\begin{corollary}[Tail-capped Hermitian entry-variance profiles]
\label{cor:tail-capped-hermitian-entry-variance-profile}
Fix \(C_0\ge1\), and put \(c_0:=3C_0/16\).  For each \(m\ge1\), let
\(D_m\ge2\), \(J_m\ge1\), and \(\ell_m=\log(eD_m)\).  For each active stream
\(1\le h\le J_m\), let
\[
  A_{m,h,i}\in\M_{D_m}(\C),\qquad 1\le i\le n_m,
\]
be self-adjoint and adapted to a filtration
\((\mcF_{m,h,i})_{i=0}^{n_m}\).  Write
\[
  A_{m,h,i}
  =
  \sum_{p=1}^{D_m}\xi_{m,h,i,p}E_{pp}
  +
  \sum_{1\le p<q\le D_m}
  \left(
  \zeta_{m,h,i,pq}E_{pq}
  +\overline{\zeta_{m,h,i,pq}}E_{qp}
  \right).
\]
Conditionally on \(\mcF_{m,h,i-1}\), assume that the upper-triangular
variables \(\{\xi_{m,h,i,p}\}_p\) and
\(\{\zeta_{m,h,i,pq}\}_{p<q}\) are independent and centered, with
\(\xi_{m,h,i,p}\) real-valued.  Let deterministic variance envelopes satisfy
\[
  \mathbb E(\xi_{m,h,i,p}^2\mid\mcF_{m,h,i-1})
  \le \alpha_{m,h,i,p}^2,\qquad
  \mathbb E(|\zeta_{m,h,i,pq}|^2\mid\mcF_{m,h,i-1})
  \le \beta_{m,h,i,pq}^2 .
\]
Let \(K_{m,h}>0\) be deterministic operator-norm caps and let
\(\eta_{m,h,i}\ge0\) satisfy
\[
  \mathbb P\{\|A_{m,h,i}\|_{\op}>K_{m,h}\}\le\eta_{m,h,i}.
\]
Define the cumulative Hermitian row-variance profile
\[
  W_{m,h}^{\rm herm}:=
  \max_{1\le p\le D_m}
  \sum_{i=1}^{n_m}
  \left(
  \alpha_{m,h,i,p}^2
  +\sum_{q<p}\beta_{m,h,i,qp}^2
  +\sum_{p<q}\beta_{m,h,i,pq}^2
  \right).
\]
Assume
\[
  W_{m,h}^{\rm herm}\le
  \frac{C_0K_{m,h}^2n_m}{2\ell_m}
  \qquad(1\le h\le J_m)
\]
and
\[
  \sum_m\sum_{h=1}^{J_m}
  \left[
  \sum_{i=1}^{n_m}\eta_{m,h,i}
  +
  D_m\exp\left\{-c_0\frac{n_m}{\ell_m}\right\}
  \right]<\infty .
\]
Then, almost surely, for all sufficiently large \(m\), every active Hermitian
entry-variance stream has a whole-matrix signing
\[
  \disc(A_{m,h,1},\ldots,A_{m,h,n_m})
  \le C(C_0)K_{m,h}\sqrt{n_m}
  \qquad(1\le h\le J_m).
\]
Independent or common sampled-sign batches are certified under the same
\(2^{-M}\)-summability hypotheses as in
Corollary~\ref{cor:finite-menu-predictable-square-profile-schedules}.

The bounded case is the specialization \(\eta_{m,h,i}=0\).  More generally,
put \(L_m=R\ell_m\), and let \(S_{m,h}>0\).  If, for fixed \(\alpha>0\) and
\(A_{\rm tail}\ge1\),
\[
  \mathbb P\{\|A_{m,h,i}\|_{\op}>S_{m,h}u\}
  \le A_{\rm tail}e^{-u^\alpha}
  \qquad(u\ge0),
\]
choose \(K_{m,h}=S_{m,h}L_m^{1/\alpha}\).  If instead, for fixed \(p_0>0\),
\[
  \mathbb E\|A_{m,h,i}\|_{\op}^{p_0}\le S_{m,h}^{p_0},
\]
choose \(K_{m,h}=S_{m,h}\exp\{L_m/p_0\}\).  In either case the cap-tail
contribution for stream \((m,h)\) is at most
\[
  A_*n_me^{-L_m},\qquad
  A_*=
  \begin{cases}
  A_{\rm tail},&\text{sub-Weibull case},\\
  1,&\text{moment case}.
  \end{cases}
\]
Consequently, if \(J_mn_m\le(eD_m)^\Gamma\) and
\(n_m\ge a\log^2(eD_m)\), then the full summability condition holds over all
integer dimensions \(D_m=m\) whenever
\[
  R>\Gamma+1,\qquad c_0a>\Gamma+2,
\]
and along stretched-exponential sparse dimension schedules whenever
\[
  R>\Gamma,\qquad c_0a>\Gamma+1.
\]
Thus Wigner, GOE/GUE, random band, sparse Hermitian, and variance-profile
upper-triangular ensembles are covered whenever their cumulative row variance
is \(O(K^2n/\log(eD))\), with one discrepancy sign assigned to each whole
Hermitian matrix.
\end{corollary}

\begin{proof}
For each stream set
\[
  Y_{m,h,i}:=
  A_{m,h,i}\mathbf 1_{\{\|A_{m,h,i}\|_{\op}\le K_{m,h}\}} .
\]
The cap-exception probability for stream \((m,h)\) is at most
\(\sum_i\eta_{m,h,i}\), and outside that event \(Y_{m,h,i}=A_{m,h,i}\) for all
\(i\).  The normalized matrices \(K_{m,h}^{-1}Y_{m,h,i}\) are self-adjoint
contractions, and
\[
  Y_{m,h,i}^2\preceq A_{m,h,i}^2 .
\]
By conditional independence and centering of the upper-triangular variables,
\(\mathbb E(A_{m,h,i}^2\mid\mcF_{m,h,i-1})\) is diagonal.  Its \(p\)-th
diagonal entry is bounded by
\[
  \alpha_{m,h,i,p}^2
  +\sum_{q<p}\beta_{m,h,i,qp}^2
  +\sum_{p<q}\beta_{m,h,i,pq}^2 .
\]
Therefore the predictable square profile of the capped normalized stream is at
most \(W_{m,h}^{\rm herm}/K_{m,h}^2\).  The half-budget hypothesis and
Corollary~\ref{cor:random-square-profile-freedman} give failure probability at
most
\[
  D_m\exp\left\{-c_0\frac{n_m}{\ell_m}\right\}
\]
for each active stream, plus the sampled-sign \(2^{-M}\) term when sampled
search is used.  Summability and Borel--Cantelli prove the eventual signing
statement.

The sub-Weibull cap gives
\(\eta_{m,h,i}\le A_{\rm tail}e^{-L_m}\), while Markov's inequality gives
\(\eta_{m,h,i}\le e^{-L_m}\) in the moment case.  The all-integer and sparse
thresholds are the same series tests for
\[
  J_mn_me^{-R\ell_m}
  \quad\text{and}\quad
  J_mD_me^{-c_0n_m/\ell_m}
\]
used in the rectangular tail-capped schedules below.
\end{proof}

\begin{corollary}[Predictable rectangular square-profile ensembles]
\label{cor:predictable-rectangular-square-profile}
Let \(d_1,d_2\ge1\), put \(D_{\rm rec}=d_1+d_2\), and let
\((\mcF_i)_{i=0}^n\) be a filtration.  Let
\[
  X_1,\ldots,X_n\in\M_{d_1\times d_2}(\C)
\]
be contractions such that \(X_i\) is \(\mcF_i\)-measurable.  Define the
predictable row and column square profiles
\[
  \nu_L:=
  \left\|
  \sum_{i=1}^n\mathbb E(X_iX_i^*\mid\mcF_{i-1})
  \right\|_{\op},
  \qquad
  \nu_R:=
  \left\|
  \sum_{i=1}^n\mathbb E(X_i^*X_i\mid\mcF_{i-1})
  \right\|_{\op},
\]
and put \(\nu_{\rm rec}:=\max\{\nu_L,\nu_R\}\).  Then, for every
\(C_0\ge1\), there is \(C(C_0)\) such that, for every \(v,t>0\) satisfying
\[
  v+t\le \frac{C_0n}{\log(eD_{\rm rec})},
\]
one has
\[
  \mathbb P\left\{
  \min_{\eps\in\{\pm1\}^n}
  \left\|\sum_{i=1}^n\eps_iX_i\right\|_{\op}
  >C(C_0)\sqrt n
  \right\}
  \le
  \mathbb P\{\nu_{\rm rec}>v\}
  +
  D_{\rm rec}\exp\left\{-\frac{t^2}{2(v+t/3)}\right\}.
\]
If \(\eta^{(1)},\ldots,\eta^{(m)}\) are independent uniform sign vectors,
independent of \(X_1,\ldots,X_n\), then, with possibly larger \(C(C_0)\),
\[
  \begin{aligned}
  &\mathbb P\left\{
  \min_{1\le r\le m}
  \left\|\sum_{i=1}^n\eta_i^{(r)}X_i\right\|_{\op}
  >C(C_0)\sqrt n
  \right\} \\
  &\qquad\le
  \mathbb P\{\nu_{\rm rec}>v\}
  +
  D_{\rm rec}\exp\left\{-\frac{t^2}{2(v+t/3)}\right\}
  +2^{-m}.
  \end{aligned}
\]
In particular, if \(\nu_{\rm rec}\le C_0n/(2\log(eD_{\rm rec}))\) almost
surely, then both statements hold with structural failure at most
\[
  D_{\rm rec}
  \exp\left\{-\frac{3C_0n}{16\log(eD_{\rm rec})}\right\},
\]
plus \(2^{-m}\) for sampled sign search.  For independent rectangular
contractions the same conclusion holds with
\[
  \nu_{\rm rec}
  =
  \max\left\{
  \left\|\sum_i\mathbb E X_iX_i^*\right\|_{\op},
  \left\|\sum_i\mathbb E X_i^*X_i\right\|_{\op}
  \right\}.
\]
Thus rectangular random operator models, bipartite support models, and
non-self-adjoint matrix-series models are Spencer-scale once both left and
right variance profiles are \(O(n/\log(e(d_1+d_2)))\).
\end{corollary}

\begin{proof}
For each \(i\), let
\[
  Y_i:=\mathscr H(X_i)=
  \begin{pmatrix}0&X_i\\ X_i^*&0\end{pmatrix}.
\]
Then \(Y_i\) is an adapted self-adjoint contraction in
\(\M_{D_{\rm rec}}(\C)\), and
\[
  \left\|\sum_i\mathbb E(Y_i^2\mid\mcF_{i-1})\right\|_{\op}
  =
  \max\{\nu_L,\nu_R\}
  =
  \nu_{\rm rec}.
\]
Also
\[
  \left\|\sum_i\eps_iY_i\right\|_{\op}
  =
  \left\|\sum_i\eps_iX_i\right\|_{\op}
\]
for every sign vector.  Applying
Corollary~\ref{cor:random-square-profile-freedman} to the self-adjoint tuple
\((Y_i)\) gives the two probability bounds.  The half-budget specialization is
the choice \(v=t=C_0n/(2\log(eD_{\rm rec}))\), and independence makes the
predictable profiles deterministic after conditioning on the past.
\end{proof}

\begin{corollary}[Centered independent-entry rectangular ensembles]
\label{cor:centered-independent-entry-rectangular}
Fix \(C_0\ge1\), and put \(c_0:=3C_0/16\).  For each \(m\ge1\), let
\(J_m\ge1\), \(d_{m,1},d_{m,2}\ge1\),
\(D_m^{\rm rec}=d_{m,1}+d_{m,2}\), and
\(\ell_m=\log(eD_m^{\rm rec})\).  For each active stream
\(1\le h\le J_m\), let
\[
  X_{m,h,i}=(x_{m,h,i,r,s})_{r,s}
  \in\M_{d_{m,1}\times d_{m,2}}(\C),
  \qquad 1\le i\le n_m,
\]
be adapted to a filtration \((\mcF_{m,h,i})_{i=0}^{n_m}\).  Assume that
\(\|X_{m,h,i}\|_{\op}\le K_{m,h}\) almost surely and that, conditionally on
\(\mcF_{m,h,i-1}\), the entries of \(X_{m,h,i}\) are independent and centered.
Suppose deterministic variance envelopes \(\sigma_{m,h,i,r,s}^2\ge0\) satisfy
\[
  \mathbb E\{|x_{m,h,i,r,s}|^2\mid\mcF_{m,h,i-1}\}
  \le \sigma_{m,h,i,r,s}^2
\]
almost surely, and define
\[
  V_{m,h}^{\rm ent}:=
  \max\left\{
  \begin{aligned}
  &\max_{1\le r\le d_{m,1}}
     \sum_{i=1}^{n_m}\sum_{s=1}^{d_{m,2}}
     \sigma_{m,h,i,r,s}^2,\\
  &\max_{1\le s\le d_{m,2}}
     \sum_{i=1}^{n_m}\sum_{r=1}^{d_{m,1}}
     \sigma_{m,h,i,r,s}^2
  \end{aligned}
  \right\}.
\]
If
\[
  V_{m,h}^{\rm ent}\le
  \frac{C_0K_{m,h}^2n_m}{2\ell_m}
  \qquad(1\le h\le J_m)
\]
and
\[
  \sum_m\sum_{h=1}^{J_m}
  D_m^{\rm rec}\exp\left\{-c_0\frac{n_m}{\ell_m}\right\}<\infty,
\]
then, almost surely, for all sufficiently large \(m\), every active
independent-entry rectangular stream has a per-matrix signing
\[
  \min_{\eps\in\{\pm1\}^{n_m}}
  \left\|\sum_{i=1}^{n_m}\eps_iX_{m,h,i}\right\|_{\op}
  \le C(C_0)K_{m,h}\sqrt{n_m}
  \qquad(1\le h\le J_m).
\]
Independent or common sampled-sign batches are certified under the same
\(2^{-M}\)-summability hypotheses as in
Corollary~\ref{cor:active-menu-rectangular-square-profile-schedules}.

In particular, if \(K_{m,h}\le K_*\), \(J_m\le(eD_m^{\rm rec})^\Gamma\), and
\[
  n_m\ge a\log^2(eD_m^{\rm rec}),
\]
then the structural series is summable over all integer dimensions
\(D_m^{\rm rec}=m\) whenever
\[
  c_0a>\Gamma+2,
\]
and along every stretched-exponential sparse dimension schedule
\(D_m^{\rm rec}\ge\exp(cm^\kappa)\) whenever
\[
  c_0a>\Gamma+1.
\]
Thus centered sparse or dense rectangular random matrices with independent
entries are Spencer-scale, with one sign per matrix, whenever their maximum
cumulative row and column variance is \(O(K^2n/\log(eD_{\rm rec}))\).
\end{corollary}

\begin{proof}
Apply Corollary~\ref{cor:predictable-rectangular-square-profile} to
\(\widetilde X_{m,h,i}:=K_{m,h}^{-1}X_{m,h,i}\).  Conditional independence and
centering make
\[
  \mathbb E(X_{m,h,i}X_{m,h,i}^*\mid\mcF_{m,h,i-1})
\]
diagonal, with \(r\)-th diagonal entry at most
\(\sum_s\sigma_{m,h,i,r,s}^2\).  Similarly,
\(\mathbb E(X_{m,h,i}^*X_{m,h,i}\mid\mcF_{m,h,i-1})\) is diagonal, with
\(s\)-th diagonal entry at most \(\sum_r\sigma_{m,h,i,r,s}^2\).  Hence the
predictable rectangular square profile of the normalized stream is at most
\(V_{m,h}^{\rm ent}/K_{m,h}^2\), which is bounded by
\(C_0n_m/(2\ell_m)\).  The half-budget form of
Corollary~\ref{cor:predictable-rectangular-square-profile} gives failure
probability at most
\[
  D_m^{\rm rec}\exp\left\{-c_0\frac{n_m}{\ell_m}\right\}
\]
for each active stream, plus \(2^{-M}\) for sampled sign search.  Summing over
\(m,h\) and applying Borel--Cantelli gives the eventual simultaneous statement.
The all-integer and sparse schedule thresholds are the same \(p\)-series and
stretched-exponential tests for
\[
  J_mD_m^{\rm rec}
  \exp\{-c_0n_m/\ell_m\}.
\]
\end{proof}

\begin{corollary}[Tail-capped independent-entry rectangular ensembles]
\label{cor:tail-capped-independent-entry-rectangular}
Fix \(C_0\ge1\), and put \(c_0:=3C_0/16\).  For each \(m\ge1\), let
\(J_m\ge1\), \(d_{m,1},d_{m,2}\ge1\),
\(D_m^{\rm rec}=d_{m,1}+d_{m,2}\), and
\(\ell_m=\log(eD_m^{\rm rec})\).  For each active stream
\(1\le h\le J_m\), let
\[
  X_{m,h,i}=(x_{m,h,i,r,s})_{r,s}
  \in\M_{d_{m,1}\times d_{m,2}}(\C),
  \qquad 1\le i\le n_m,
\]
be adapted to a filtration \((\mcF_{m,h,i})_{i=0}^{n_m}\).  Conditionally on
\(\mcF_{m,h,i-1}\), assume the entries of \(X_{m,h,i}\) are independent and
centered, and let deterministic variance envelopes satisfy
\[
  \mathbb E\{|x_{m,h,i,r,s}|^2\mid\mcF_{m,h,i-1}\}
  \le \sigma_{m,h,i,r,s}^2 .
\]
Let \(K_{m,h}>0\) be deterministic operator-norm caps and let
\(\eta_{m,h,i}\ge0\) satisfy
\[
  \mathbb P\{\|X_{m,h,i}\|_{\op}>K_{m,h}\}\le\eta_{m,h,i}.
\]
Define \(V_{m,h}^{\rm ent}\) as in
Corollary~\ref{cor:centered-independent-entry-rectangular}.  Assume
\[
  V_{m,h}^{\rm ent}\le
  \frac{C_0K_{m,h}^2n_m}{2\ell_m}
  \qquad(1\le h\le J_m)
\]
and
\[
  \sum_m\sum_{h=1}^{J_m}
  \left[
  \sum_{i=1}^{n_m}\eta_{m,h,i}
  +
  D_m^{\rm rec}\exp\left\{-c_0\frac{n_m}{\ell_m}\right\}
  \right]<\infty .
\]
Then, almost surely, for all sufficiently large \(m\), every active
independent-entry rectangular stream has a per-matrix signing
\[
  \min_{\eps\in\{\pm1\}^{n_m}}
  \left\|\sum_{i=1}^{n_m}\eps_iX_{m,h,i}\right\|_{\op}
  \le C(C_0)K_{m,h}\sqrt{n_m}
  \qquad(1\le h\le J_m).
\]
Independent or common sampled-sign batches are certified under the same
\(2^{-M}\)-summability hypotheses as in
Corollary~\ref{cor:active-menu-rectangular-square-profile-schedules}.

The cap may be chosen from standard whole-matrix tails.  Put
\(L_m=R\ell_m\), and let \(S_{m,h}>0\).  If, for fixed \(\alpha>0\) and
\(A_{\rm tail}\ge1\),
\[
  \mathbb P\{\|X_{m,h,i}\|_{\op}>S_{m,h}u\}
  \le A_{\rm tail}e^{-u^\alpha}
  \qquad(u\ge0),
\]
choose \(K_{m,h}=S_{m,h}L_m^{1/\alpha}\).  If instead, for fixed \(p>0\),
\[
  \mathbb E\|X_{m,h,i}\|_{\op}^p\le S_{m,h}^p,
\]
choose \(K_{m,h}=S_{m,h}\exp\{L_m/p\}\).  In the two cases the cap-tail
contribution for stream \((m,h)\) is at most
\[
  A_*n_m e^{-L_m},
  \qquad A_*=
  \begin{cases}
  A_{\rm tail},&\text{sub-Weibull case},\\
  1,&\text{moment case}.
  \end{cases}
\]
Consequently, if \(J_mn_m\le(eD_m^{\rm rec})^\Gamma\) and
\(n_m\ge a\log^2(eD_m^{\rm rec})\), then the full summability condition holds
over all integer dimensions \(D_m^{\rm rec}=m\) whenever
\[
  R>\Gamma+1,\qquad c_0a>\Gamma+2,
\]
and along stretched-exponential sparse dimension schedules whenever
\[
  R>\Gamma,\qquad c_0a>\Gamma+1.
\]
\end{corollary}

\begin{proof}
For each stream set
\[
  Y_{m,h,i}:=
  X_{m,h,i}\mathbf 1_{\{\|X_{m,h,i}\|_{\op}\le K_{m,h}\}} .
\]
The cap-exception probability for stream \((m,h)\) is at most
\(\sum_i\eta_{m,h,i}\), and outside that event \(Y_{m,h,i}=X_{m,h,i}\) for all
\(i\).  The normalized matrices \(K_{m,h}^{-1}Y_{m,h,i}\) are contractions.
Moreover,
\[
  Y_{m,h,i}Y_{m,h,i}^*
  \preceq X_{m,h,i}X_{m,h,i}^*,
  \qquad
  Y_{m,h,i}^*Y_{m,h,i}
  \preceq X_{m,h,i}^*X_{m,h,i}.
\]
By the same conditional independence and centering calculation as in
Corollary~\ref{cor:centered-independent-entry-rectangular}, the left and right
predictable square profiles of the capped normalized stream are therefore at
most \(V_{m,h}^{\rm ent}/K_{m,h}^2\).  The half-budget hypothesis and
Corollary~\ref{cor:predictable-rectangular-square-profile} give the displayed
structural failure term, plus the sampled-sign \(2^{-M}\) term when sampled
search is used.  Summability and Borel--Cantelli prove the eventual signing
statement.

The sub-Weibull cap gives
\(\eta_{m,h,i}\le A_{\rm tail}e^{-L_m}\), while Markov's inequality gives
\(\eta_{m,h,i}\le e^{-L_m}\) in the moment case.  The all-integer and sparse
thresholds are the same series tests for
\[
  J_mn_me^{-R\ell_m}
  \quad\text{and}\quad
  J_mD_m^{\rm rec}e^{-c_0n_m/\ell_m}
\]
used in Corollary~\ref{cor:closed-form-entry-intensity-cap-schedules}.
\end{proof}

\begin{corollary}[Active menus of rectangular square-profile streams]
\label{cor:active-menu-rectangular-square-profile-schedules}
For each \(m\ge1\), let \(J_m\ge1\), \(d_{m,1},d_{m,2}\ge1\),
\(D_m^{\rm rec}=d_{m,1}+d_{m,2}\), and \(n_m\ge1\).  For every
\(1\le h\le J_m\), let
\[
  X_{m,h,1},\ldots,X_{m,h,n_m}
  \in\M_{d_{m,1}\times d_{m,2}}(\C)
\]
be an adapted rectangular-contraction stream with predictable rectangular
profile \(\nu_{m,h}^{\rm rec}\) as in
Corollary~\ref{cor:predictable-rectangular-square-profile}.  Assume that,
outside events \(E_{m,h}\),
\[
  \nu_{m,h}^{\rm rec}
  \le
  \frac{C_0n_m}{2\log(eD_m^{\rm rec})}
  \qquad(1\le h\le J_m),
\]
and that
\[
  \sum_m\sum_{h=1}^{J_m}\mathbb P(E_{m,h})<\infty,
  \qquad
  \sum_m J_mD_m^{\rm rec}
  \exp\left\{
  -\frac{3C_0n_m}{16\log(eD_m^{\rm rec})}
  \right\}<\infty .
\]
Then, almost surely, for all sufficiently large \(m\), every active rectangular
stream has a Spencer-scale signing:
\[
  \min_{\eps\in\{\pm1\}^{n_m}}
  \left\|\sum_{i=1}^{n_m}\eps_iX_{m,h,i}\right\|_{\op}
  \le C(C_0)\sqrt{n_m}
  \qquad(1\le h\le J_m).
\]
Sampled certificates are also eventual almost surely if either
\[
  \sum_m\sum_{h=1}^{J_m}2^{-M_{m,h}}<\infty
\]
for independent per-stream batches, or
\[
  \sum_mJ_m2^{-M_m}<\infty
\]
for one common active-menu batch.  If \(J_m\to\infty\), the active prefixes may
exhaust a countable library of rectangular random-matrix models; every fixed
library element is then covered eventually.
\end{corollary}

\begin{proof}
Apply Corollary~\ref{cor:predictable-rectangular-square-profile} to each active
stream with \(v=t=C_0n_m/(2\log(eD_m^{\rm rec}))\), then sum the displayed
failure probabilities over \(m\) and \(h\).  The first Borel--Cantelli lemma
gives the simultaneous eventual signings.  The sampled-search statements add
the corresponding \(2^{-M_{m,h}}\) or common \(J_m2^{-M_m}\) losses before the
same Borel--Cantelli argument.  If \(J_m\to\infty\), each fixed library stream
is active for all sufficiently large \(m\).
\end{proof}

\begin{corollary}[Common signs for rectangular square-profile menus]
\label{cor:common-sign-rectangular-square-profile-menu}
Fix \(C_0\ge1\).  For each \(m\ge1\), let \(J_m\ge1\),
\(d_{m,1},d_{m,2}\ge1\), \(D_m^{\rm rec}=d_{m,1}+d_{m,2}\), and
\[
  D_m^{\rm com}:=J_mD_m^{\rm rec}.
\]
For every \(1\le h\le J_m\), let
\[
  X_{m,h,1},\ldots,X_{m,h,n_m}
  \in\M_{d_{m,1}\times d_{m,2}}(\C)
\]
be arbitrary random rectangular matrices.  Assume that there are events
\(E_{m,h}^{\rm sq}\) such that, outside \(E_{m,h}^{\rm sq}\),
\[
  \max\left\{
  \left\|\sum_{i=1}^{n_m}X_{m,h,i}X_{m,h,i}^*\right\|_{\op},
  \left\|\sum_{i=1}^{n_m}X_{m,h,i}^*X_{m,h,i}\right\|_{\op}
  \right\}
  \le
  \frac{C_0n_m}{\log(eD_m^{\rm com})}.
\]
If
\[
  \sum_m\sum_{h=1}^{J_m}\mathbb P(E_{m,h}^{\rm sq})<\infty ,
\]
then, almost surely, for all sufficiently large \(m\), one sign vector
\(\eps^{(m)}\in\{\pm1\}^{n_m}\) balances the whole active rectangular menu:
\[
  \max_{1\le h\le J_m}
  \left\|\sum_{i=1}^{n_m}\eps_i^{(m)}X_{m,h,i}\right\|_{\op}
  \le C(C_0)\sqrt{n_m}.
\]
If \(M_m\) independent random sign vectors are sampled for the whole active
menu and
\[
  \sum_m2^{-M_m}<\infty ,
\]
then the sampled list almost surely eventually contains such a common
certificate.  If \(J_m\to\infty\), active prefixes may exhaust a countable
library of rectangular or non-self-adjoint random-matrix models, and every
fixed library element is eventually covered by the common active-prefix
signing.

Thus any finite or countably exhausted rectangular model library whose realized
left and right square profiles are summably certified at the common-sign scale
has a single simultaneous Matrix-Spencer signing.
Its rank-sensitive form replaces \(J_m(d_{m,1}+d_{m,2})\) by the total active
left-plus-right square-profile rank, preserving the same shared sampled-sign
loss \(2^{-M_m}\).
\end{corollary}

\begin{proof}
Outside the union of the events \(E_{m,h}^{\rm sq}\), define
\[
  \widetilde X_{m,i}:=
  \bigoplus_{h=1}^{J_m}X_{m,h,i}
  \in\M_{J_md_{m,1}\times J_md_{m,2}}(\C).
\]
Then
\[
  \begin{aligned}
  \left\|\sum_i\widetilde X_{m,i}\widetilde X_{m,i}^*\right\|_{\op}
  &=
  \max_{h\le J_m}
  \left\|\sum_iX_{m,h,i}X_{m,h,i}^*\right\|_{\op},\\
  \left\|\sum_i\widetilde X_{m,i}^*\widetilde X_{m,i}\right\|_{\op}
  &=
  \max_{h\le J_m}
  \left\|\sum_iX_{m,h,i}^*X_{m,h,i}\right\|_{\op}.
  \end{aligned}
\]
Corollary~\ref{cor:rectangular-square-profile-certificate} in dimension
\(D_m^{\rm com}\) gives one signing of the block-direct-sum tuple, which is the
displayed common signing.  Since
\[
  \sum_m\mathbb P\left(\bigcup_{h\le J_m}E_{m,h}^{\rm sq}\right)
  \le
  \sum_m\sum_{h\le J_m}\mathbb P(E_{m,h}^{\rm sq})<\infty ,
\]
Borel--Cantelli gives the eventual deterministic certificates.  The sampled
statement applies the sampled part of
Corollary~\ref{cor:rectangular-square-profile-certificate} to the same
direct-sum tuple and adds only the failure series \(\sum_m2^{-M_m}\).  If
\(J_m\to\infty\), each fixed library stream belongs to all sufficiently large
active prefixes.
\end{proof}

\begin{corollary}[Common signs for rank-sensitive rectangular square-profile menus]
\label{cor:common-sign-rank-sensitive-rectangular-square-profile-menu}
Fix \(C_0\ge1\).  For each \(m\ge1\), let \(J_m\ge1\),
\(d_{m,1},d_{m,2},n_m\ge1\), and let \(R_m\ge1\).  For every
\(1\le h\le J_m\), let
\[
  X_{m,h,1},\ldots,X_{m,h,n_m}
  \in\M_{d_{m,1}\times d_{m,2}}(\C)
\]
be arbitrary random rectangular matrices.  Put
\[
  V_{m,h}^L:=\sum_{i=1}^{n_m}X_{m,h,i}X_{m,h,i}^*,
  \qquad
  V_{m,h}^R:=\sum_{i=1}^{n_m}X_{m,h,i}^*X_{m,h,i},
\]
and
\[
  \rho_{m,h}^{\rm rec}:=
  \max\{1,\operatorname{rank}V_{m,h}^L+\operatorname{rank}V_{m,h}^R\}.
\]
Assume that there are events \(E_m^{\rm rrec}\) such that, outside
\(E_m^{\rm rrec}\),
\[
  \sum_{h=1}^{J_m}\rho_{m,h}^{\rm rec}\le R_m,
  \qquad
  \max_{1\le h\le J_m}\max\{\|V_{m,h}^L\|_{\op},\|V_{m,h}^R\|_{\op}\}
  \le
  \frac{C_0n_m}{\log(eR_m)}.
\]
If
\[
  \sum_m\mathbb P(E_m^{\rm rrec})<\infty ,
\]
then, almost surely, for all sufficiently large \(m\), one sign vector
\(\eps^{(m)}\in\{\pm1\}^{n_m}\) balances the whole active rectangular menu:
\[
  \max_{1\le h\le J_m}
  \left\|\sum_{i=1}^{n_m}\eps_i^{(m)}X_{m,h,i}\right\|_{\op}
  \le C(C_0)\sqrt{n_m}.
\]
If \(M_m\) independent random sign vectors are sampled for the whole active
menu and
\[
  \sum_m2^{-M_m}<\infty ,
\]
then the sampled list almost surely eventually contains such a common
certificate.

In particular, if each active stream has certified realized profile ranks
\[
  \operatorname{rank}V_{m,h}^L\le\kappa_{m,h}^L,
  \qquad
  \operatorname{rank}V_{m,h}^R\le\kappa_{m,h}^R,
\]
one may take
\[
  R_m=\sum_{h=1}^{J_m}\max\{1,\kappa_{m,h}^L+\kappa_{m,h}^R\}.
\]
Thus common signs for low-rank rectangular or non-self-adjoint active
libraries pay the logarithm of total active left-plus-right square-function
rank, not the ambient block dimension \(J_m(d_{m,1}+d_{m,2})\).
\end{corollary}

\begin{proof}
On \((E_m^{\rm rrec})^c\), form the block-direct-sum rectangular tuple
\[
  \widetilde X_{m,i}:=
  \bigoplus_{h=1}^{J_m}X_{m,h,i}
  \in\M_{J_md_{m,1}\times J_md_{m,2}}(\C).
\]
Its left and right square profiles are
\[
  \sum_i\widetilde X_{m,i}\widetilde X_{m,i}^*
  =
  \bigoplus_{h=1}^{J_m}V_{m,h}^L,
  \qquad
  \sum_i\widetilde X_{m,i}^*\widetilde X_{m,i}
  =
  \bigoplus_{h=1}^{J_m}V_{m,h}^R.
\]
Therefore
\[
  \max\left\{
  \left\|\sum_i\widetilde X_{m,i}\widetilde X_{m,i}^*\right\|_{\op},
  \left\|\sum_i\widetilde X_{m,i}^*\widetilde X_{m,i}\right\|_{\op}
  \right\}
  \le
  \frac{C_0n_m}{\log(eR_m)}
\]
and the rectangular square-profile rank of the direct-sum tuple is at most
\[
  \sum_{h=1}^{J_m}
  \left(\operatorname{rank}V_{m,h}^L+\operatorname{rank}V_{m,h}^R\right)
  \le
  \sum_{h=1}^{J_m}\rho_{m,h}^{\rm rec}
  \le R_m .
\]
The rank-sensitive part of
Corollary~\ref{cor:rectangular-square-profile-certificate} gives one signing
of the direct-sum tuple with norm at most \(C(C_0)\sqrt{n_m}\), and this is
exactly one sign vector balancing every active rectangular stream.
Borel--Cantelli applied to \(\sum_m\mathbb P(E_m^{\rm rrec})\) gives the
eventual almost-sure statement.

For sampled signs, condition on the active menu on \((E_m^{\rm rrec})^c\) and
apply the sampled rank-sensitive part of
Corollary~\ref{cor:rectangular-square-profile-certificate} to the same
direct-sum tuple.  The sampled failure for the whole active menu is
\(2^{-M_m}\), so the displayed summability and Borel--Cantelli complete the
proof.  The final specialization is immediate from the rank bounds.
\end{proof}

\begin{corollary}[Rectangular low-overlap support ensembles]
\label{cor:rectangular-low-overlap-support}
For every \(C_0\ge1\) there is \(C(C_0)\) with the following property.  Let
\[
  X_1,\ldots,X_n\in\M_{d_1\times d_2}(\C)
\]
be random rectangular matrices and put \(D_{\rm rec}=d_1+d_2\).  Suppose that,
outside an event \(E\), there are orthogonal projections
\[
  P_i^L\in\M_{d_1}(\C),\qquad P_i^R\in\M_{d_2}(\C),
  \qquad 0\le\rho_i\le1,
\]
such that
\[
  X_i=P_i^LX_iP_i^R,\qquad \|X_i\|_{\op}\le\rho_i
  \qquad(1\le i\le n),
\]
and
\[
  \max\left\{
  \left\|\sum_{i=1}^n\rho_i^2P_i^L\right\|_{\op},
  \left\|\sum_{i=1}^n\rho_i^2P_i^R\right\|_{\op}
  \right\}
  \le
  \frac{C_0n}{\log(eD_{\rm rec})}.
\]
Then
\[
  \mathbb P\left\{
  \min_{\eps\in\{\pm1\}^n}
  \left\|\sum_{i=1}^n\eps_iX_i\right\|_{\op}
  > C(C_0)\sqrt n
  \right\}
  \le \mathbb P(E).
\]
If \(m\) independent random sign trials are used, the same conclusion holds
with an additional \(2^{-m}\) failure probability.  Thus rectangular sparse,
bipartite, row-column localized, and rectangular random-subspace models reduce
to the two weighted overlap operators on the range and domain sides.
\end{corollary}

\begin{proof}
On a realization outside \(E\), the support and norm assumptions give
\[
  0\preceq X_iX_i^*\preceq \rho_i^2P_i^L,
  \qquad
  0\preceq X_i^*X_i\preceq \rho_i^2P_i^R .
\]
Hence the two rectangular square profiles in
Corollary~\ref{cor:rectangular-square-profile-certificate} are at most
\(C_0n/\log(eD_{\rm rec})\).  That corollary gives the deterministic signing and
the sampled-search statement on this realization.
\end{proof}

\begin{corollary}[Active menus with dependent rectangular overlap budgets]
\label{cor:active-menu-dependent-rectangular-overlap-budget}
Fix \(C_0\ge1\).  For each \(m\ge1\), let \(J_m\ge1\),
\(d_{m,1},d_{m,2}\ge1\), \(D_m^{\rm rec}=d_{m,1}+d_{m,2}\), and \(n_m\ge1\).
For every \(1\le h\le J_m\), let
\[
  X_{m,h,1},\ldots,X_{m,h,n_m}\in\M_{d_{m,1}\times d_{m,2}}(\C)
\]
be arbitrary random rectangular matrices.  Assume that there are events
\(F_{m,h}^{\rm ov}\), random orthogonal projections
\[
  P_{m,h,i}^L\in\M_{d_{m,1}}(\C),\qquad
  P_{m,h,i}^R\in\M_{d_{m,2}}(\C),
\]
and random amplitudes \(0\le\rho_{m,h,i}\le1\) such that, outside
\(F_{m,h}^{\rm ov}\),
\[
  X_{m,h,i}=P_{m,h,i}^LX_{m,h,i}P_{m,h,i}^R,
  \qquad
  \|X_{m,h,i}\|_{\op}\le\rho_{m,h,i}
  \qquad(1\le i\le n_m),
\]
and
\[
  \max\left\{
  \left\|\sum_{i=1}^{n_m}\rho_{m,h,i}^2P_{m,h,i}^L\right\|_{\op},
  \left\|\sum_{i=1}^{n_m}\rho_{m,h,i}^2P_{m,h,i}^R\right\|_{\op}
  \right\}
  \le
  \frac{C_0n_m}{\log(eD_m^{\rm rec})}.
\]
If
\[
  \sum_m\sum_{h=1}^{J_m}\mathbb P(F_{m,h}^{\rm ov})<\infty ,
\]
then, almost surely, for all sufficiently large \(m\), every active stream has a
Spencer-scale rectangular signing:
\[
  \min_{\eps\in\{\pm1\}^{n_m}}
  \left\|\sum_{i=1}^{n_m}\eps_iX_{m,h,i}\right\|_{\op}
  \le C(C_0)\sqrt{n_m}
  \qquad(1\le h\le J_m).
\]
Sampled certificates are also eventual almost surely under the same
per-stream or common-batch \(2^{-M}\)-summability hypotheses as in
Corollary~\ref{cor:active-menu-rectangular-square-profile-schedules}.  If
\(J_m\to\infty\), active prefixes may exhaust a countable library.

Thus arbitrary dependent rectangular support laws, random regular or biregular
bipartite supports, matchings, without-replacement row/column choices,
determinantal or negatively dependent supports, and adaptive block/subspace
frontiers are covered once their realized left and right weighted overlap
budgets fail only summably.
\end{corollary}

\begin{proof}
Apply Corollary~\ref{cor:rectangular-low-overlap-support} to each active stream
on \((F_{m,h}^{\rm ov})^c\).  Summability of the overlap-budget failures and the
first Borel--Cantelli lemma give the simultaneous eventual signings.  The
sampled-certificate statements add the same per-stream
\(\sum_{m,h}2^{-M_{m,h}}\) or common-batch \(\sum_mJ_m2^{-M_m}\) failure series
used in Corollary~\ref{cor:active-menu-rectangular-square-profile-schedules}.
If \(J_m\to\infty\), each fixed library stream belongs to the active prefix for
all sufficiently large \(m\).
\end{proof}

\begin{corollary}[Tail-capped dependent rectangular overlap budgets]
\label{cor:active-menu-tail-capped-dependent-rectangular-overlap-budget}
Fix \(C_0\ge1\).  In the active-menu setting of
Corollary~\ref{cor:active-menu-dependent-rectangular-overlap-budget}, replace
the pointwise norm bounds by deterministic caps
\[
  0\le\rho_{m,h,i}\le1
  \qquad(1\le h\le J_m,\ 1\le i\le n_m).
\]
Assume that there are events \(F_{m,h}^{\rm ov}\), random orthogonal projections
\[
  P_{m,h,i}^L\in\M_{d_{m,1}}(\C),\qquad
  P_{m,h,i}^R\in\M_{d_{m,2}}(\C),
\]
and cap-tail budgets \(\delta_{m,h,i}\ge0\) such that
\[
  \mathbb P\{\|X_{m,h,i}\|_{\op}>\rho_{m,h,i}\}\le\delta_{m,h,i},
\]
and, outside \(F_{m,h}^{\rm ov}\),
\[
  X_{m,h,i}=P_{m,h,i}^LX_{m,h,i}P_{m,h,i}^R
  \qquad(1\le i\le n_m)
\]
and
\[
  \max\left\{
  \left\|\sum_{i=1}^{n_m}\rho_{m,h,i}^2P_{m,h,i}^L\right\|_{\op},
  \left\|\sum_{i=1}^{n_m}\rho_{m,h,i}^2P_{m,h,i}^R\right\|_{\op}
  \right\}
  \le
  \frac{C_0n_m}{\log(eD_m^{\rm rec})}.
\]
If
\[
  \sum_m\sum_{h=1}^{J_m}
  \left[
    \mathbb P(F_{m,h}^{\rm ov})
    +\sum_{i=1}^{n_m}\delta_{m,h,i}
  \right]<\infty ,
\]
then, almost surely, for all sufficiently large \(m\), every active stream has a
Spencer-scale rectangular signing:
\[
  \min_{\eps\in\{\pm1\}^{n_m}}
  \left\|\sum_{i=1}^{n_m}\eps_iX_{m,h,i}\right\|_{\op}
  \le C(C_0)\sqrt{n_m}
  \qquad(1\le h\le J_m).
\]
Sampled certificates and countable-library exhaustions follow under the same
additional summability hypotheses as in
Corollary~\ref{cor:active-menu-dependent-rectangular-overlap-budget}.

For example, if
\[
  \mathbb P\{\|X_{m,h,i}\|_{\op}>s_{m,h,i}u\}\le A_*e^{-u^\alpha}
  \qquad(u\ge0)
\]
with fixed \(\alpha>0\), choose
\(\rho_{m,h,i}=s_{m,h,i}\ell_{m,h,i}^{1/\alpha}\le1\) and
\(\delta_{m,h,i}=A_*e^{-\ell_{m,h,i}}\).  If instead
\(\mathbb E\|X_{m,h,i}\|_{\op}^p\le s_{m,h,i}^p\), choose
\(\rho_{m,h,i}=s_{m,h,i}e^{\ell_{m,h,i}/p}\le1\) and
\(\delta_{m,h,i}=e^{-\ell_{m,h,i}}\).  These deterministic caps feed the same
realized left/right overlap budget, so dependent heavy-tailed rectangular
supports are reduced to cap-tail summability plus a capped overlap certificate.
\end{corollary}

\begin{proof}
For every active stream put
\[
  \Omega_{m,h}^{\rm cap}
  :=\{\|X_{m,h,i}\|_{\op}\le\rho_{m,h,i}\text{ for all }i\le n_m\}.
\]
On \((F_{m,h}^{\rm ov})^c\cap\Omega_{m,h}^{\rm cap}\), the hypotheses of
Corollary~\ref{cor:active-menu-dependent-rectangular-overlap-budget} hold with
the deterministic amplitudes \(\rho_{m,h,i}\).  The exceptional probability for
stream \((m,h)\) is at most
\[
  \mathbb P(F_{m,h}^{\rm ov})
  +\sum_{i=1}^{n_m}\delta_{m,h,i}.
\]
The displayed summability and Borel--Cantelli give the eventual active-menu
signings; sampled certificates add the same sampled-sign failure series as in
the dependent-overlap corollary.  The sub-Weibull and moment cap choices follow
from the displayed tail bound and Markov's inequality, respectively.
\end{proof}

\begin{corollary}[Common signs for dependent rectangular overlap menus]
\label{cor:common-sign-dependent-rectangular-overlap-budget}
Fix \(C_0\ge1\).  For each \(m\ge1\), let \(J_m\ge1\),
\(d_{m,1},d_{m,2}\ge1\), \(D_m^{\rm rec}=d_{m,1}+d_{m,2}\), and
\[
  D_m^{\rm com}:=J_mD_m^{\rm rec}.
\]
For every \(1\le h\le J_m\), let
\[
  X_{m,h,1},\ldots,X_{m,h,n_m}\in\M_{d_{m,1}\times d_{m,2}}(\C)
\]
be arbitrary random rectangular matrices.  Assume that there are events
\(F_{m,h}^{\rm ov}\), random orthogonal projections
\[
  P_{m,h,i}^L\in\M_{d_{m,1}}(\C),\qquad
  P_{m,h,i}^R\in\M_{d_{m,2}}(\C),
\]
and random amplitudes \(0\le\rho_{m,h,i}\le1\) such that, outside
\(F_{m,h}^{\rm ov}\),
\[
  X_{m,h,i}=P_{m,h,i}^LX_{m,h,i}P_{m,h,i}^R,\qquad
  \|X_{m,h,i}\|_{\op}\le\rho_{m,h,i}
  \qquad(1\le i\le n_m),
\]
and
\[
  \max\left\{
  \left\|\sum_{i=1}^{n_m}\rho_{m,h,i}^2P_{m,h,i}^L\right\|_{\op},
  \left\|\sum_{i=1}^{n_m}\rho_{m,h,i}^2P_{m,h,i}^R\right\|_{\op}
  \right\}
  \le
  \frac{C_0n_m}{\log(eD_m^{\rm com})}.
\]
If
\[
  \sum_m\sum_{h=1}^{J_m}\mathbb P(F_{m,h}^{\rm ov})<\infty ,
\]
then, almost surely, for all sufficiently large \(m\), there is one sign vector
\(\eps^{(m)}\in\{\pm1\}^{n_m}\) such that
\[
  \max_{1\le h\le J_m}
  \left\|\sum_{i=1}^{n_m}\eps_i^{(m)}X_{m,h,i}\right\|_{\op}
  \le C(C_0)\sqrt{n_m}.
\]
If, for each \(m\), \(M_m\) independent random sign vectors are sampled and
tested against the whole active menu, and
\[
  \sum_m2^{-M_m}<\infty ,
\]
then the sampled list almost surely eventually contains such a common
certificate.  If \(J_m\to\infty\), active prefixes may exhaust a countable
library of dependent rectangular models; every fixed library element is then
covered eventually by the common active-prefix signing.

Thus arbitrary dependent bipartite supports, random regular or biregular
rectangular supports, matchings, without-replacement row/column choices, and
adaptive rectangular block or subspace frontiers have a single simultaneous
Matrix-Spencer signing whenever their realized left and right overlap budgets
are summably certified at the common-sign scale.
\end{corollary}

\begin{proof}
Outside the union of the events \(F_{m,h}^{\rm ov}\), define the rectangular
block-direct-sum tuple
\[
  \widetilde X_{m,i}:=
  \bigoplus_{h=1}^{J_m}X_{m,h,i}
  \in \M_{J_md_{m,1}\times J_md_{m,2}}(\C)
  \qquad(1\le i\le n_m).
\]
Then
\[
  \begin{aligned}
  \sum_i\widetilde X_{m,i}\widetilde X_{m,i}^*
  &=
  \bigoplus_{h=1}^{J_m}\sum_iX_{m,h,i}X_{m,h,i}^*,\\
  \sum_i\widetilde X_{m,i}^*\widetilde X_{m,i}
  &=
  \bigoplus_{h=1}^{J_m}\sum_iX_{m,h,i}^*X_{m,h,i}.
  \end{aligned}
\]
The support and norm bounds give
\[
  X_{m,h,i}X_{m,h,i}^*\preceq\rho_{m,h,i}^2P_{m,h,i}^L,
  \qquad
  X_{m,h,i}^*X_{m,h,i}\preceq\rho_{m,h,i}^2P_{m,h,i}^R,
\]
so the two square profiles of the direct-sum tuple are both bounded by
\(C_0n_m/\log(eD_m^{\rm com})\).  Corollary~\ref{cor:rectangular-square-profile-certificate} applied in rectangular
dimension \(D_m^{\rm com}\) gives one signing of the direct-sum tuple, which is
exactly one sign vector balancing every active stream.  The first
Borel--Cantelli lemma applies because
\[
  \sum_m\mathbb P\left(\bigcup_{h\le J_m}F_{m,h}^{\rm ov}\right)
  \le
  \sum_m\sum_{h\le J_m}\mathbb P(F_{m,h}^{\rm ov})<\infty .
\]
The sampled statement follows by applying the sampled part of
Corollary~\ref{cor:rectangular-square-profile-certificate} to the same direct-sum tuple,
adding only the failure series \(\sum_m2^{-M_m}\).  If \(J_m\to\infty\), each
fixed library element belongs to every sufficiently large active prefix.
\end{proof}

\begin{corollary}[Common signs for support-rank rectangular overlap menus]
\label{cor:common-sign-support-rank-rectangular-overlap-budget}
Fix \(C_0\ge1\).  For each \(m\ge1\), let \(J_m\ge1\),
\(d_{m,1},d_{m,2},n_m\ge1\), and let \(R_m\ge1\).  For every
\(1\le h\le J_m\), let
\[
  X_{m,h,1},\ldots,X_{m,h,n_m}
  \in\M_{d_{m,1}\times d_{m,2}}(\C)
\]
be arbitrary random rectangular matrices.  Assume that there are events
\(F_m^{\rm sr}\), random orthogonal projections
\[
  P_{m,h,i}^L\in\M_{d_{m,1}}(\C),\qquad
  P_{m,h,i}^R\in\M_{d_{m,2}}(\C),
\]
random support-envelope projections
\[
  Q_{m,h}^L\in\M_{d_{m,1}}(\C),\qquad
  Q_{m,h}^R\in\M_{d_{m,2}}(\C),
\]
and amplitudes \(0\le\rho_{m,h,i}\le1\) such that, outside \(F_m^{\rm sr}\),
\[
  X_{m,h,i}=P_{m,h,i}^LX_{m,h,i}P_{m,h,i}^R,\qquad
  \|X_{m,h,i}\|_{\op}\le\rho_{m,h,i},
\]
\[
  P_{m,h,i}^L\le Q_{m,h}^L,\qquad
  P_{m,h,i}^R\le Q_{m,h}^R
  \qquad(1\le h\le J_m,\ 1\le i\le n_m),
\]
and
\[
  \sum_{h=1}^{J_m}
  \max\{1,\operatorname{rank}Q_{m,h}^L+\operatorname{rank}Q_{m,h}^R\}
  \le R_m.
\]
Assume also that, outside \(F_m^{\rm sr}\),
\[
  \max_{1\le h\le J_m}
  \max\left\{
  \left\|\sum_{i=1}^{n_m}\rho_{m,h,i}^2P_{m,h,i}^L\right\|_{\op},
  \left\|\sum_{i=1}^{n_m}\rho_{m,h,i}^2P_{m,h,i}^R\right\|_{\op}
  \right\}
  \le
  \frac{C_0n_m}{\log(eR_m)}.
\]
If
\[
  \sum_m\mathbb P(F_m^{\rm sr})<\infty ,
\]
then, almost surely, for all sufficiently large \(m\), there is one sign vector
\(\eps^{(m)}\in\{\pm1\}^{n_m}\) such that
\[
  \max_{1\le h\le J_m}
  \left\|\sum_{i=1}^{n_m}\eps_i^{(m)}X_{m,h,i}\right\|_{\op}
  \le C(C_0)\sqrt{n_m}.
\]
If \(M_m\) independent random sign vectors are sampled and tested against the
whole active menu, and
\[
  \sum_m2^{-M_m}<\infty ,
\]
then the sampled list almost surely eventually contains such a common
certificate.  If \(J_m\to\infty\), active prefixes may exhaust a countable
library and every fixed library element is covered eventually.

Thus dependent rectangular support models with small active support-envelope
rank pay \(\log(eR_m)\), not \(\log(eJ_m(d_{m,1}+d_{m,2}))\).
\end{corollary}

\begin{proof}
On \((F_m^{\rm sr})^c\), put
\[
  V_{m,h}^L=\sum_iX_{m,h,i}X_{m,h,i}^*,
  \qquad
  V_{m,h}^R=\sum_iX_{m,h,i}^*X_{m,h,i}.
\]
The support and norm hypotheses give
\[
  0\preceq V_{m,h}^L
  \preceq
  \sum_i\rho_{m,h,i}^2P_{m,h,i}^L,
  \qquad
  0\preceq V_{m,h}^R
  \preceq
  \sum_i\rho_{m,h,i}^2P_{m,h,i}^R,
\]
and \(V_{m,h}^L=Q_{m,h}^LV_{m,h}^LQ_{m,h}^L\),
\(V_{m,h}^R=Q_{m,h}^RV_{m,h}^RQ_{m,h}^R\).  Hence
\[
  \operatorname{rank}V_{m,h}^L+\operatorname{rank}V_{m,h}^R
  \le
  \operatorname{rank}Q_{m,h}^L+\operatorname{rank}Q_{m,h}^R
\]
and the left/right square-profile norms are bounded by
\(C_0n_m/\log(eR_m)\).  Therefore the hypotheses of
Corollary~\ref{cor:common-sign-rank-sensitive-rectangular-square-profile-menu}
hold with exceptional event \(F_m^{\rm sr}\).  That corollary gives the
common deterministic signing and the sampled-sign conclusion.  The summability
of \(\mathbb P(F_m^{\rm sr})\) and Borel--Cantelli make the conclusions
eventual almost surely.
\end{proof}

\begin{corollary}[Tail-capped support-rank common signs for rectangular overlap menus]
\label{cor:common-sign-tail-capped-support-rank-rectangular-overlap-budget}
Fix \(C_0\ge1\).  For each \(m\ge1\), let \(J_m\ge1\),
\(d_{m,1},d_{m,2},n_m\ge1\), and let \(R_m\ge1\).  For every
\(1\le h\le J_m\), let
\[
  X_{m,h,1},\ldots,X_{m,h,n_m}
  \in\M_{d_{m,1}\times d_{m,2}}(\C)
\]
be arbitrary random rectangular matrices.  Assume that there are events
\(F_m^{\rm sr}\), random orthogonal projections
\[
  P_{m,h,i}^L\in\M_{d_{m,1}}(\C),\qquad
  P_{m,h,i}^R\in\M_{d_{m,2}}(\C),
\]
random support-envelope projections
\[
  Q_{m,h}^L\in\M_{d_{m,1}}(\C),\qquad
  Q_{m,h}^R\in\M_{d_{m,2}}(\C),
\]
deterministic caps \(0\le\rho_{m,h,i}\le1\), and cap-tail budgets
\(\delta_{m,h,i}\ge0\) such that
\[
  \mathbb P\{\|X_{m,h,i}\|_{\op}>\rho_{m,h,i}\}\le\delta_{m,h,i}.
\]
Outside \(F_m^{\rm sr}\), assume that
\[
  X_{m,h,i}=P_{m,h,i}^LX_{m,h,i}P_{m,h,i}^R,\qquad
  P_{m,h,i}^L\le Q_{m,h}^L,\qquad
  P_{m,h,i}^R\le Q_{m,h}^R
\]
for all \(h\le J_m\) and \(i\le n_m\), and
\[
  \sum_{h=1}^{J_m}
  \max\{1,\operatorname{rank}Q_{m,h}^L+\operatorname{rank}Q_{m,h}^R\}
  \le R_m.
\]
Assume also that, outside \(F_m^{\rm sr}\),
\[
  \max_{1\le h\le J_m}
  \max\left\{
  \left\|\sum_{i=1}^{n_m}\rho_{m,h,i}^2P_{m,h,i}^L\right\|_{\op},
  \left\|\sum_{i=1}^{n_m}\rho_{m,h,i}^2P_{m,h,i}^R\right\|_{\op}
  \right\}
  \le
  \frac{C_0n_m}{\log(eR_m)}.
\]
If
\[
  \sum_m\left[
    \mathbb P(F_m^{\rm sr})
    +\sum_{h=1}^{J_m}\sum_{i=1}^{n_m}\delta_{m,h,i}
  \right]<\infty ,
\]
then, almost surely, for all sufficiently large \(m\), there is one sign vector
\(\eps^{(m)}\in\{\pm1\}^{n_m}\) such that
\[
  \max_{1\le h\le J_m}
  \left\|\sum_{i=1}^{n_m}\eps_i^{(m)}X_{m,h,i}\right\|_{\op}
  \le C(C_0)\sqrt{n_m}.
\]
If \(M_m\) independent random sign vectors are sampled and tested against the
whole active menu, and
\[
  \sum_m2^{-M_m}<\infty ,
\]
then the sampled list almost surely eventually contains such a common
certificate.  If \(J_m\to\infty\), active prefixes may exhaust a countable
library and every fixed library element is covered eventually.

Thus unbounded localized rectangular libraries retain the support-envelope
logarithm \(\log(eR_m)\) after paying only summable cap-tail failures.
\end{corollary}

\begin{proof}
Let
\[
  \Omega_m^{\rm cap}
  :=
  \{\|X_{m,h,i}\|_{\op}\le\rho_{m,h,i}
  \text{ for all }h\le J_m,\ i\le n_m\}.
\]
The union bound gives
\[
  \sum_m\mathbb P\bigl((\Omega_m^{\rm cap})^c\bigr)
  \le
  \sum_m\sum_{h=1}^{J_m}\sum_{i=1}^{n_m}\delta_{m,h,i}<\infty .
\]
On the complement of
\[
  G_m:=F_m^{\rm sr}\cup(\Omega_m^{\rm cap})^c
\]
the support-rank rectangular corollary applies with amplitudes
\(\rho_{m,h,i}\).  Since the exceptional events \(G_m\) are summable,
Borel--Cantelli gives the eventual deterministic common signings and the
sampled-sign conclusion.  The countable-library conclusion is inherited from
the same corollary.
\end{proof}

\begin{corollary}[Common signs for rectangular coordinate-window Bernstein menus]
\label{cor:common-sign-rectangular-coordinate-window-bernstein}
Fix \(C_0\ge1\).  For each \(m\ge1\), let \(J_m\ge1\),
\(d_{m,1},d_{m,2},n_m\ge1\), and let \(R_m\ge1\).  For every
\(1\le h\le J_m\), let
\[
  E_{m,h}^L\subset\{1,\ldots,d_{m,1}\},\qquad
  E_{m,h}^R\subset\{1,\ldots,d_{m,2}\}
\]
be deterministic coordinate windows satisfying
\[
  \sum_{h=1}^{J_m}
  \max\{1,|E_{m,h}^L|+|E_{m,h}^R|\}\le R_m .
\]
For each \(h,i\), let
\[
  R_{m,h,i}\subset E_{m,h}^L,\qquad
  S_{m,h,i}\subset E_{m,h}^R
\]
be random row and column supports.  Assume that, for every
\(a\in E_{m,h}^L\), the indicators
\({\bf 1}_{\{a\in R_{m,h,i}\}}\), \(1\le i\le n_m\), are independent and
\[
  \mathbb P\{a\in R_{m,h,i}\}\le p^L_{m,h,i,a},
\]
and similarly, for every \(b\in E_{m,h}^R\),
\[
  \mathbb P\{b\in S_{m,h,i}\}\le p^R_{m,h,i,b}
\]
with independence in \(i\).  Conditional on the supports, let
\[
  X_{m,h,i}\in\M_{d_{m,1}\times d_{m,2}}(\C)
\]
be arbitrary random rectangular matrices satisfying
\[
  X_{m,h,i}=P_{R_{m,h,i}}X_{m,h,i}P_{S_{m,h,i}}.
\]
Let deterministic caps \(0\le\rho_{m,h,i}\le1\) and cap-tail budgets
\(\delta_{m,h,i}\ge0\) satisfy
\[
  \mathbb P\{\|X_{m,h,i}\|_{\op}>\rho_{m,h,i}\}\le\delta_{m,h,i}.
\]
Put
\[
  b_{m,h}:=\max_i\rho_{m,h,i}^2,
\]
\[
  \mu_{m,h}:=
  \max\left\{
  \max_{a\in E_{m,h}^L}\sum_i\rho_{m,h,i}^2p^L_{m,h,i,a},
  \max_{b\in E_{m,h}^R}\sum_i\rho_{m,h,i}^2p^R_{m,h,i,b}
  \right\},
\]
and
\[
  \sigma_{m,h}^2:=
  \max\left\{
  \max_{a\in E_{m,h}^L}\sum_i\rho_{m,h,i}^4p^L_{m,h,i,a},
  \max_{b\in E_{m,h}^R}\sum_i\rho_{m,h,i}^4p^R_{m,h,i,b}
  \right\},
\]
where a maximum over an empty window is interpreted as \(0\).  Suppose that
there are \(t_{m,h}>0\) such that
\[
  \mu_{m,h}+t_{m,h}
  \le
  \frac{C_0n_m}{\log(eR_m)}
  \qquad(1\le h\le J_m),
\]
and
\[
  \begin{aligned}
  \sum_m\Biggl[
  &\sum_{h=1}^{J_m}
  (|E_{m,h}^L|+|E_{m,h}^R|)
  \exp\left\{
  -\frac{t_{m,h}^2}
         {2(\sigma_{m,h}^2+b_{m,h}t_{m,h}/3)}
  \right\}\\
  &+
  \sum_{h=1}^{J_m}\sum_{i=1}^{n_m}\delta_{m,h,i}
  \Biggr]<\infty,
  \end{aligned}
\]
with zero denominator interpreted as zero contribution.  Then, almost surely,
for all sufficiently large \(m\), there is one sign vector
\(\eps^{(m)}\in\{\pm1\}^{n_m}\) such that
\[
  \max_{1\le h\le J_m}
  \left\|\sum_{i=1}^{n_m}\eps_i^{(m)}X_{m,h,i}\right\|_{\op}
  \le C(C_0)\sqrt{n_m}.
\]
The same sampled common-certificate and countable-library conclusions as in
Corollary~\ref{cor:common-sign-tail-capped-support-rank-rectangular-overlap-budget}
hold.

In particular, if
\[
  \mu_{m,h}\le \frac{C_0n_m}{2\log(eR_m)}
  \qquad(1\le h\le J_m),
\]
then one may take \(t_{m,h}=C_0n_m/(2\log(eR_m))\), and the displayed
Bernstein summand is bounded by
\[
  \sum_{h=1}^{J_m}
  (|E_{m,h}^L|+|E_{m,h}^R|)
  \exp\left\{-\frac{3C_0n_m}{16\log(eR_m)}\right\}.
\]
\end{corollary}

\begin{proof}
For each row coordinate define
\[
  N_{m,h,a}^L=\sum_i\rho_{m,h,i}^2{\bf 1}_{\{a\in R_{m,h,i}\}},
\]
and define \(N_{m,h,b}^R\) analogously on the column side.  Scalar Bernstein
and the union bound over the displayed windows show that the probability that
some active coordinate has load above \(\mu_{m,h}+t_{m,h}\) is bounded by the
displayed Bernstein summand.  Off this event, with
\(Q_{m,h}^L=P_{E_{m,h}^L}\), \(Q_{m,h}^R=P_{E_{m,h}^R}\),
\(P_{m,h,i}^L=P_{R_{m,h,i}}\), and \(P_{m,h,i}^R=P_{S_{m,h,i}}\), the
support-envelope rank and left/right overlap hypotheses of
Corollary~\ref{cor:common-sign-tail-capped-support-rank-rectangular-overlap-budget}
hold.  The cap-tail summability is exactly the remaining summand, so that
corollary gives the common signing, sampled certificates, and countable-library
conclusion.

For the final display, \(\sigma_{m,h}^2\le\mu_{m,h}\le t_{m,h}\) and
\(b_{m,h}\le1\), so
\[
  \frac{t_{m,h}^2}{2(\sigma_{m,h}^2+b_{m,h}t_{m,h}/3)}
  \ge \frac{3t_{m,h}}8
  =
  \frac{3C_0n_m}{16\log(eR_m)} .
\]
\end{proof}

\begin{corollary}[Common signs for rectangular predictable coordinate-window Freedman menus]
\label{cor:common-sign-rectangular-predictable-coordinate-window-freedman}
Fix \(C_0\ge1\).  For each \(m\ge1\), let \(J_m\ge1\),
\(d_{m,1},d_{m,2},n_m\ge1\), and let \(R_m\ge1\).  For every
\(1\le h\le J_m\), let
\[
  E_{m,h}^L\subset\{1,\ldots,d_{m,1}\},\qquad
  E_{m,h}^R\subset\{1,\ldots,d_{m,2}\}
\]
be deterministic coordinate windows satisfying
\[
  \sum_{h=1}^{J_m}
  \max\{1,|E_{m,h}^L|+|E_{m,h}^R|\}\le R_m .
\]
For each \(m,h\), let
\[
  \mcF_{m,h,0}\subset\mcF_{m,h,1}\subset\cdots\subset\mcF_{m,h,n_m}
\]
be a filtration.  For each \(h,i\), let
\[
  R_{m,h,i}\subset E_{m,h}^L,\qquad
  S_{m,h,i}\subset E_{m,h}^R
\]
be \(\mcF_{m,h,i}\)-measurable random row and column supports.  Assume that
there are deterministic envelopes \(p^L_{m,h,i,a},p^R_{m,h,i,b}\in[0,1]\)
such that, for every \(a\in E_{m,h}^L\) and \(b\in E_{m,h}^R\),
\[
  \mathbb P\{a\in R_{m,h,i}\mid\mcF_{m,h,i-1}\}\le p^L_{m,h,i,a},
  \qquad
  \mathbb P\{b\in S_{m,h,i}\mid\mcF_{m,h,i-1}\}\le p^R_{m,h,i,b}.
\]
Conditional on the supports, let
\[
  X_{m,h,i}\in\M_{d_{m,1}\times d_{m,2}}(\C)
\]
be arbitrary random rectangular matrices satisfying
\[
  X_{m,h,i}=P_{R_{m,h,i}}X_{m,h,i}P_{S_{m,h,i}}.
\]
Let deterministic caps \(0\le\rho_{m,h,i}\le1\) and cap-tail budgets
\(\delta_{m,h,i}\ge0\) satisfy
\[
  \mathbb P\{\|X_{m,h,i}\|_{\op}>\rho_{m,h,i}\}\le\delta_{m,h,i}.
\]
Put
\[
  b_{m,h}:=\max_i\rho_{m,h,i}^2,
\]
\[
  \mu_{m,h}:=
  \max\left\{
  \max_{a\in E_{m,h}^L}\sum_i\rho_{m,h,i}^2p^L_{m,h,i,a},
  \max_{b\in E_{m,h}^R}\sum_i\rho_{m,h,i}^2p^R_{m,h,i,b}
  \right\},
\]
and
\[
  \sigma_{m,h}^2:=
  \max\left\{
  \max_{a\in E_{m,h}^L}\sum_i\rho_{m,h,i}^4p^L_{m,h,i,a},
  \max_{b\in E_{m,h}^R}\sum_i\rho_{m,h,i}^4p^R_{m,h,i,b}
  \right\},
\]
where a maximum over an empty window is interpreted as \(0\).  Suppose that
there are \(t_{m,h}>0\) such that
\[
  \mu_{m,h}+t_{m,h}
  \le
  \frac{C_0n_m}{\log(eR_m)}
  \qquad(1\le h\le J_m),
\]
and
\[
  \begin{aligned}
  \sum_m\Biggl[
  &\sum_{h=1}^{J_m}
  (|E_{m,h}^L|+|E_{m,h}^R|)
  \exp\left\{
  -\frac{t_{m,h}^2}
         {2(\sigma_{m,h}^2+b_{m,h}t_{m,h}/3)}
  \right\}\\
  &+
  \sum_{h=1}^{J_m}\sum_{i=1}^{n_m}\delta_{m,h,i}
  \Biggr]<\infty,
  \end{aligned}
\]
with zero denominator interpreted as zero contribution.  Then, almost surely,
for all sufficiently large \(m\), there is one sign vector
\(\eps^{(m)}\in\{\pm1\}^{n_m}\) such that
\[
  \max_{1\le h\le J_m}
  \left\|\sum_{i=1}^{n_m}\eps_i^{(m)}X_{m,h,i}\right\|_{\op}
  \le C(C_0)\sqrt{n_m}.
\]
The same sampled common-certificate and countable-library conclusions as in
Corollary~\ref{cor:common-sign-tail-capped-support-rank-rectangular-overlap-budget}
hold.

In particular, if
\[
  \mu_{m,h}\le \frac{C_0n_m}{2\log(eR_m)}
  \qquad(1\le h\le J_m),
\]
then one may take \(t_{m,h}=C_0n_m/(2\log(eR_m))\), and the displayed Freedman
summand is bounded by
\[
  \sum_{h=1}^{J_m}
  (|E_{m,h}^L|+|E_{m,h}^R|)
  \exp\left\{-\frac{3C_0n_m}{16\log(eR_m)}\right\}.
\]
\end{corollary}

\begin{proof}
For each row coordinate set
\[
  N_{m,h,a}^L=\sum_i\rho_{m,h,i}^2{\bf 1}_{\{a\in R_{m,h,i}\}},
\]
and define \(N_{m,h,b}^R\) analogously on the column side.  For fixed
\(m,h,a\), the process
\[
  \sum_{i'\le i}\rho_{m,h,i'}^2
  \left(
  {\bf 1}_{\{a\in R_{m,h,i'}\}}
  -\mathbb E({\bf 1}_{\{a\in R_{m,h,i'}\}}\mid\mcF_{m,h,i'-1})
  \right)
\]
is a scalar martingale with increments bounded above by \(b_{m,h}\) and
predictable quadratic variation at most
\(\sum_i\rho_{m,h,i}^4p^L_{m,h,i,a}\).  Scalar Freedman and the union bound over
all row and column coordinates in the active windows show that the probability
that some active coordinate load exceeds \(\mu_{m,h}+t_{m,h}\) is bounded by the
displayed Freedman summand.  Off this event, with
\(Q_{m,h}^L=P_{E_{m,h}^L}\), \(Q_{m,h}^R=P_{E_{m,h}^R}\),
\(P_{m,h,i}^L=P_{R_{m,h,i}}\), and \(P_{m,h,i}^R=P_{S_{m,h,i}}\), the
support-envelope rank and left/right overlap hypotheses of
Corollary~\ref{cor:common-sign-tail-capped-support-rank-rectangular-overlap-budget}
hold.  The cap-tail summability is exactly the remaining summand, so that
corollary gives the common signing, sampled certificates, and countable-library
conclusion.  The half-budget estimate follows from
\(\sigma_{m,h}^2\le\mu_{m,h}\le t_{m,h}\) and \(b_{m,h}\le1\), giving exponent
at least \(3t_{m,h}/8\).
\end{proof}

\begin{corollary}[Common-sign rectangular predictable coordinate-window schedules]
\label{cor:common-sign-rectangular-predictable-coordinate-window-schedules}
In the setting of
Corollary~\ref{cor:common-sign-rectangular-predictable-coordinate-window-freedman},
suppose that
\[
  \mu_{m,h}\le \frac{C_0n_m}{2\log(eR_m)}
  \qquad(1\le h\le J_m)
\]
eventually, and that
\[
  \sum_m\sum_{h=1}^{J_m}\sum_{i=1}^{n_m}\delta_{m,h,i}<\infty .
\]
Then the same common-sign, sampled-certificate, and countable-library
conclusions hold whenever
\[
  \sum_m
  R_m\exp\left\{-\frac{3C_0n_m}{16\log(eR_m)}\right\}<\infty .
\]
In particular, after the displayed half-budget and cap-tail assumptions, the
structural tail is summable in each of the following support-rank schedules.
\begin{enumerate}[label=\textup{(\roman*)}]
\item Arbitrary growth.  If, for some \(\gamma>0\),
  \[
    n_m\ge \frac{16}{3C_0}\log(eR_m)
    \bigl(\log(eR_m)+(1+\gamma)\log(e+m)\bigr)
  \]
  eventually, then the almost-sure common-sign conclusion holds.  If a shared
  batch of \(M_m\) independent random sign vectors is used for the whole active
  menu, the sampled conclusion follows from
  \(M_m\ge(1+\delta)\log_2(e+m)\) eventually.
\item All support ranks.  Suppose the rows are indexed by
  \(R=2,3,\ldots\) and the \(R\)-th row has active coordinate-window rank
  parameter at most \(R\).  If
  \[
    n_R\ge a\log^2(eR)
  \]
  eventually for some \(a>32/(3C_0)\), then the common-sign conclusion holds
  along all support ranks.  The sampled conclusion follows if
  \(M_R\ge(1+\delta)\log_2(eR)\) eventually.
\item Sparse support ranks.  If \(R_m\ge\exp(c m^\kappa)\) eventually for some
  \(c,\kappa>0\), and
  \[
    n_m\ge a\log^2(eR_m)
  \]
  eventually for some \(a>16/(3C_0)\), then the common-sign conclusion holds.
  The sampled conclusion follows from
  \(M_m\ge(1+\delta)\log_2(e+m)\) eventually.
\end{enumerate}
\end{corollary}

\begin{proof}
Under the half-budget hypothesis, the final estimate in
Corollary~\ref{cor:common-sign-rectangular-predictable-coordinate-window-freedman}
bounds the structural Freedman tail by
\[
  \sum_{h=1}^{J_m}
  (|E_{m,h}^L|+|E_{m,h}^R|)
  \exp\left\{-\frac{3C_0n_m}{16\log(eR_m)}\right\}
  \le
  R_m\exp\left\{-\frac{3C_0n_m}{16\log(eR_m)}\right\}.
\]
The cap-tail series is summable by assumption, so the preceding corollary and
Borel--Cantelli give the common signing and the inherited sampled and
countable-library assertions.

For arbitrary growth, the lower bound on \(n_m\) makes the displayed structural
tail at most \(e^{-1}(e+m)^{-(1+\gamma)}\).  For all support ranks, the \(R\)-th
tail is bounded by
\[
  R\exp\left\{-\frac{3C_0a}{16}\log(eR)\right\},
\]
which is summable when \(3C_0a/16>2\).  For sparse support ranks the same
estimate is \(O(R_m^{1-3C_0a/16})\), summable under the stretched-exponential
growth of \(R_m\) as soon as \(3C_0a/16>1\).  The sampled thresholds are exactly
the summability of \(2^{-M_m}\) or \(2^{-M_R}\).
\end{proof}

\begin{corollary}[Common signs for rectangular subspace-envelope Bernstein menus]
\label{cor:common-sign-rectangular-subspace-envelope-bernstein}
Fix \(C_0\ge1\).  For each \(m\ge1\), let \(J_m\ge1\),
\(d_{m,1},d_{m,2},n_m\ge1\), and let \(R_m\ge1\).  For every
\(1\le h\le J_m\), let
\[
  Q_{m,h}^L\in\M_{d_{m,1}}(\C),\qquad
  Q_{m,h}^R\in\M_{d_{m,2}}(\C)
\]
be deterministic orthogonal projections satisfying
\[
  \sum_{h=1}^{J_m}
  \max\{1,\operatorname{rank}Q_{m,h}^L+\operatorname{rank}Q_{m,h}^R\}
  \le R_m .
\]
For each \(h,i\), let
\[
  P_{m,h,i}^L\le Q_{m,h}^L,\qquad
  P_{m,h,i}^R\le Q_{m,h}^R
\]
be random orthogonal projections, with the pairs
\((P_{m,h,i}^L,P_{m,h,i}^R)\), \(1\le i\le n_m\), independent for each fixed
\((m,h)\).  Conditional on these projections, let
\[
  X_{m,h,i}\in\M_{d_{m,1}\times d_{m,2}}(\C)
\]
be arbitrary random rectangular matrices satisfying
\[
  X_{m,h,i}=P_{m,h,i}^LX_{m,h,i}P_{m,h,i}^R.
\]
Let deterministic caps \(0\le\rho_{m,h,i}\le1\) and cap-tail budgets
\(\delta_{m,h,i}\ge0\) satisfy
\[
  \mathbb P\{\|X_{m,h,i}\|_{\op}>\rho_{m,h,i}\}\le\delta_{m,h,i}.
\]
Put
\[
  b_{m,h}:=\max_i\rho_{m,h,i}^2,
\]
and
\[
  \mu_{m,h}:=
  \max\left\{
  \left\|\sum_i\rho_{m,h,i}^2\mathbb EP_{m,h,i}^L\right\|_{\op},
  \left\|\sum_i\rho_{m,h,i}^2\mathbb EP_{m,h,i}^R\right\|_{\op}
  \right\}.
\]
Suppose that there are \(t_{m,h}>0\) such that
\[
  \mu_{m,h}+t_{m,h}
  \le
  \frac{C_0n_m}{\log(eR_m)}
  \qquad(1\le h\le J_m),
\]
and
\[
  \begin{aligned}
  \sum_m\Biggl[
  &\sum_{h=1}^{J_m}
  \bigl(\operatorname{rank}Q_{m,h}^L+\operatorname{rank}Q_{m,h}^R\bigr)
  \exp\left\{
  -\frac{t_{m,h}^2}{2(\mu_{m,h}+b_{m,h}t_{m,h}/3)}
  \right\}\\
  &+
  \sum_{h=1}^{J_m}\sum_{i=1}^{n_m}\delta_{m,h,i}
  \Biggr]<\infty,
  \end{aligned}
\]
with zero denominator interpreted as zero contribution.  Then, almost surely,
for all sufficiently large \(m\), there is one sign vector
\(\eps^{(m)}\in\{\pm1\}^{n_m}\) such that
\[
  \max_{1\le h\le J_m}
  \left\|\sum_{i=1}^{n_m}\eps_i^{(m)}X_{m,h,i}\right\|_{\op}
  \le C(C_0)\sqrt{n_m}.
\]
The same sampled common-certificate and countable-library conclusions as in
Corollary~\ref{cor:common-sign-tail-capped-support-rank-rectangular-overlap-budget}
hold.

In particular, if
\[
  \mu_{m,h}\le \frac{C_0n_m}{2\log(eR_m)}
  \qquad(1\le h\le J_m),
\]
then one may take \(t_{m,h}=C_0n_m/(2\log(eR_m))\), and the displayed
Bernstein summand is bounded by
\[
  \sum_{h=1}^{J_m}
  \bigl(\operatorname{rank}Q_{m,h}^L+\operatorname{rank}Q_{m,h}^R\bigr)
  \exp\left\{-\frac{3C_0n_m}{16\log(eR_m)}\right\}.
\]
For Haar-distributed rank-\(r_{m,h,i}^L\) and rank-\(r_{m,h,i}^R\) projections
inside the ranges of \(Q_{m,h}^L\) and \(Q_{m,h}^R\), the half-budget condition
is
\[
  \max\left\{
  \frac{\sum_i\rho_{m,h,i}^2r_{m,h,i}^L}{\operatorname{rank}Q_{m,h}^L},
  \frac{\sum_i\rho_{m,h,i}^2r_{m,h,i}^R}{\operatorname{rank}Q_{m,h}^R}
  \right\}
  \le
  \frac{C_0n_m}{2\log(eR_m)},
\]
omitting a ratio whose denominator is \(0\).
\end{corollary}

\begin{proof}
Apply matrix Bernstein separately in the row and column envelope subspaces.  On
the row side, set
\[
  Y_{m,h,i}^L=\rho_{m,h,i}^2
  (P_{m,h,i}^L-\mathbb EP_{m,h,i}^L)
\]
as an operator on the range of \(Q_{m,h}^L\).  Then
\(\mathbb EY_{m,h,i}^L=0\), \(\|Y_{m,h,i}^L\|_{\op}\le b_{m,h}\), and
\[
  \sum_i\mathbb E(Y_{m,h,i}^L)^2
  \preceq
  \sum_i\rho_{m,h,i}^2\mathbb EP_{m,h,i}^L,
\]
so the variance parameter is at most \(\mu_{m,h}\).  The same argument applies
on the column side.  Union over \(h\) and over the two envelope dimensions
gives the displayed Bernstein failure series for the event
\[
  \max_h\max\left\{
  \left\|\sum_i\rho_{m,h,i}^2P_{m,h,i}^L\right\|_{\op},
  \left\|\sum_i\rho_{m,h,i}^2P_{m,h,i}^R\right\|_{\op}
  \right\}
  \le \frac{C_0n_m}{\log(eR_m)}.
\]
Off this event, the support-envelope rank and overlap hypotheses of
Corollary~\ref{cor:common-sign-tail-capped-support-rank-rectangular-overlap-budget}
hold with envelopes \(Q_{m,h}^L,Q_{m,h}^R\).  The cap-tail summability is the
remaining summand, so that corollary gives the common signing, sampled
certificates, and countable-library conclusion.  The half-budget and Haar
specializations follow from the same computation as in
Corollary~\ref{cor:rectangular-random-subspace-bernstein}.
\end{proof}

\begin{corollary}[Common signs for rectangular predictable subspace-envelope Freedman menus]
\label{cor:common-sign-rectangular-predictable-subspace-envelope-freedman}
Fix \(C_0\ge1\).  For each \(m\ge1\), let \(J_m\ge1\),
\(d_{m,1},d_{m,2},n_m\ge1\), and let \(R_m\ge1\).  For every
\(1\le h\le J_m\), let
\[
  Q_{m,h}^L\in\M_{d_{m,1}}(\C),\qquad
  Q_{m,h}^R\in\M_{d_{m,2}}(\C)
\]
be deterministic orthogonal projections satisfying
\[
  \sum_{h=1}^{J_m}
  \max\{1,\operatorname{rank}Q_{m,h}^L+\operatorname{rank}Q_{m,h}^R\}
  \le R_m .
\]
For each \(m,h\), let
\[
  \mcF_{m,h,0}\subset\mcF_{m,h,1}\subset\cdots\subset\mcF_{m,h,n_m}
\]
be a filtration.  For each \(h,i\), let
\[
  P_{m,h,i}^L\le Q_{m,h}^L,\qquad
  P_{m,h,i}^R\le Q_{m,h}^R
\]
be \(\mcF_{m,h,i}\)-measurable random orthogonal projections.  Conditional on
these projections, let
\[
  X_{m,h,i}\in\M_{d_{m,1}\times d_{m,2}}(\C)
\]
be arbitrary random rectangular matrices satisfying
\[
  X_{m,h,i}=P_{m,h,i}^LX_{m,h,i}P_{m,h,i}^R .
\]
Let deterministic caps \(0\le\rho_{m,h,i}\le1\), cap-tail budgets
\(\delta_{m,h,i}\ge0\), deterministic predictable-overlap budgets
\(v_{m,h}\ge0\), and events \(F_m^{\rm pr}\) be such that
\[
  \mathbb P\{\|X_{m,h,i}\|_{\op}>\rho_{m,h,i}\}\le\delta_{m,h,i},
\]
and, outside \(F_m^{\rm pr}\), for every \(h\le J_m\),
\[
\begin{aligned}
  \left\|\sum_i\rho_{m,h,i}^2
  \mathbb E(P_{m,h,i}^L\mid\mcF_{m,h,i-1})\right\|_{\op}
  &\le v_{m,h},\\
  \left\|\sum_i\rho_{m,h,i}^2
  \mathbb E(P_{m,h,i}^R\mid\mcF_{m,h,i-1})\right\|_{\op}
  &\le v_{m,h}.
\end{aligned}
\]
Put
\[
  b_{m,h}:=\max_i\rho_{m,h,i}^2 .
\]
Suppose that there are \(t_{m,h}>0\) such that
\[
  v_{m,h}+t_{m,h}
  \le
  \frac{C_0n_m}{\log(eR_m)}
  \qquad(1\le h\le J_m),
\]
and
\[
  \begin{aligned}
  \sum_m\Biggl[
  &\mathbb P(F_m^{\rm pr})
  +\sum_{h=1}^{J_m}
  2\bigl(\operatorname{rank}Q_{m,h}^L+\operatorname{rank}Q_{m,h}^R\bigr)
  \exp\left\{
  -\frac{t_{m,h}^2}{2(v_{m,h}+b_{m,h}t_{m,h}/3)}
  \right\}\\
  &+
  \sum_{h=1}^{J_m}\sum_{i=1}^{n_m}\delta_{m,h,i}
  \Biggr]<\infty,
  \end{aligned}
\]
with zero denominator interpreted as zero contribution.  Then, almost surely,
for all sufficiently large \(m\), there is one sign vector
\(\eps^{(m)}\in\{\pm1\}^{n_m}\) such that
\[
  \max_{1\le h\le J_m}
  \left\|\sum_{i=1}^{n_m}\eps_i^{(m)}X_{m,h,i}\right\|_{\op}
  \le C(C_0)\sqrt{n_m}.
\]
The same sampled common-certificate and countable-library conclusions as in
Corollary~\ref{cor:common-sign-tail-capped-support-rank-rectangular-overlap-budget}
hold.

In particular, if
\[
  v_{m,h}\le \frac{C_0n_m}{2\log(eR_m)}
  \qquad(1\le h\le J_m),
\]
then one may take \(t_{m,h}=C_0n_m/(2\log(eR_m))\), and the Freedman summand is
bounded by
\[
  2\sum_{h=1}^{J_m}
  \bigl(\operatorname{rank}Q_{m,h}^L+\operatorname{rank}Q_{m,h}^R\bigr)
  \exp\left\{-\frac{3C_0n_m}{16\log(eR_m)}\right\}.
\]
If, conditionally on \(\mcF_{m,h,i-1}\), the row and column projections are
Haar-distributed inside the ranges of \(Q_{m,h}^L\) and \(Q_{m,h}^R\) with
predictable ranks \(r_{m,h,i}^L\) and \(r_{m,h,i}^R\), the half-budget
condition is
\[
  \max\left\{
  \frac{\sum_i\rho_{m,h,i}^2r_{m,h,i}^L}{\operatorname{rank}Q_{m,h}^L},
  \frac{\sum_i\rho_{m,h,i}^2r_{m,h,i}^R}{\operatorname{rank}Q_{m,h}^R}
  \right\}
  \le
  \frac{C_0n_m}{2\log(eR_m)},
\]
omitting a ratio whose denominator is \(0\).
\end{corollary}

\begin{proof}
Apply the two-sided matrix Freedman inequality separately in the row and column
envelope subspaces.  On the row side, set
\[
  Z_{m,h,i}^L=\rho_{m,h,i}^2
  \left(
  P_{m,h,i}^L-\mathbb E(P_{m,h,i}^L\mid\mcF_{m,h,i-1})
  \right)
\]
as an operator on the range of \(Q_{m,h}^L\).  Then
\((Z_{m,h,i}^L,\mcF_{m,h,i})\) is a self-adjoint martingale-difference
sequence, \(\|Z_{m,h,i}^L\|_{\op}\le b_{m,h}\), and
\[
  \sum_i
  \mathbb E((Z_{m,h,i}^L)^2\mid\mcF_{m,h,i-1})
  \preceq
  \sum_i\rho_{m,h,i}^2
  \mathbb E(P_{m,h,i}^L\mid\mcF_{m,h,i-1}).
\]
The same construction applies on the column side.  Hence, off
\(F_m^{\rm pr}\), the displayed Freedman series bounds the probability that
some active row or column realized overlap exceeds \(v_{m,h}+t_{m,h}\).  Off
that event and the cap-failure event, the support-envelope rank and overlap
hypotheses of
Corollary~\ref{cor:common-sign-tail-capped-support-rank-rectangular-overlap-budget}
hold with envelopes \(Q_{m,h}^L,Q_{m,h}^R\).  The summability assumptions and
Borel--Cantelli therefore give the common signing, sampled certificates, and
countable-library conclusion.  The half-budget estimate uses \(b_{m,h}\le1\)
and \(v_{m,h}\le t_{m,h}\).  Conditional Haar invariance gives
\[
  \mathbb E(P_{m,h,i}^L\mid\mcF_{m,h,i-1})
  =
  \frac{r_{m,h,i}^L}{\operatorname{rank}Q_{m,h}^L}Q_{m,h}^L
\]
and the analogous column identity, which gives the final specialization.
\end{proof}

\begin{corollary}[Common-sign rectangular predictable subspace-envelope schedules]
\label{cor:common-sign-rectangular-predictable-subspace-envelope-schedules}
In the setting of
Corollary~\ref{cor:common-sign-rectangular-predictable-subspace-envelope-freedman},
suppose that
\[
  v_{m,h}\le \frac{C_0n_m}{2\log(eR_m)}
  \qquad(1\le h\le J_m)
\]
eventually, and that
\[
  \sum_m\mathbb P(F_m^{\rm pr})<\infty,\qquad
  \sum_m\sum_{h=1}^{J_m}\sum_{i=1}^{n_m}\delta_{m,h,i}<\infty .
\]
Then the same common-sign, sampled-certificate, and countable-library
conclusions hold whenever
\[
  \sum_m
  2R_m\exp\left\{-\frac{3C_0n_m}{16\log(eR_m)}\right\}<\infty .
\]
In particular, after the displayed half-budget, predictable-exception, and
cap-tail assumptions, the structural tail is summable in each of the following
support-rank schedules.
\begin{enumerate}[label=\textup{(\roman*)}]
\item Arbitrary growth.  If, for some \(\gamma>0\),
  \[
    n_m\ge \frac{16}{3C_0}\log(eR_m)
    \bigl(\log(2eR_m)+(1+\gamma)\log(e+m)\bigr)
  \]
  eventually, then the almost-sure common-sign conclusion holds.  If a shared
  batch of \(M_m\) independent random sign vectors is used for the whole active
  menu, the sampled conclusion follows from
  \(M_m\ge(1+\delta)\log_2(e+m)\) eventually.
\item All support ranks.  Suppose the rows are indexed by
  \(R=2,3,\ldots\) and the \(R\)-th row has active subspace-envelope rank
  parameter at most \(R\).  If
  \[
    n_R\ge a\log^2(eR)
  \]
  eventually for some \(a>32/(3C_0)\), then the common-sign conclusion holds
  along all support ranks.  The sampled conclusion follows if
  \(M_R\ge(1+\delta)\log_2(eR)\) eventually.
\item Sparse support ranks.  If \(R_m\ge\exp(c m^\kappa)\) eventually for some
  \(c,\kappa>0\), and
  \[
    n_m\ge a\log^2(eR_m)
  \]
  eventually for some \(a>16/(3C_0)\), then the common-sign conclusion holds.
  The sampled conclusion follows from
  \(M_m\ge(1+\delta)\log_2(e+m)\) eventually.
\end{enumerate}
\end{corollary}

\begin{proof}
Under the half-budget hypothesis, the final estimate in
Corollary~\ref{cor:common-sign-rectangular-predictable-subspace-envelope-freedman}
bounds the structural Freedman tail by
\[
  \begin{aligned}
  &2\sum_{h=1}^{J_m}
  \bigl(\operatorname{rank}Q_{m,h}^L+\operatorname{rank}Q_{m,h}^R\bigr)
  \exp\left\{-\frac{3C_0n_m}{16\log(eR_m)}\right\}\\
  &\qquad\le
  2R_m\exp\left\{-\frac{3C_0n_m}{16\log(eR_m)}\right\}.
  \end{aligned}
\]
The predictable-exception and cap-tail series are summable by assumption, so the
preceding corollary and Borel--Cantelli give the common signing and the inherited
sampled and countable-library assertions.

For arbitrary growth, the lower bound on \(n_m\) makes the displayed structural
tail at most \(e^{-1}(e+m)^{-(1+\gamma)}\).  For all support ranks, the \(R\)-th
tail is bounded by
\[
  2R\exp\left\{-\frac{3C_0a}{16}\log(eR)\right\},
\]
which is summable when \(3C_0a/16>2\).  For sparse support ranks the same
estimate is \(O(R_m^{1-3C_0a/16})\), summable under the stretched-exponential
growth of \(R_m\) as soon as \(3C_0a/16>1\).  The sampled thresholds are exactly
the summability of \(2^{-M_m}\) or \(2^{-M_R}\).
\end{proof}

\begin{corollary}[Tail-capped common signs for dependent rectangular overlap menus]
\label{cor:common-sign-tail-capped-dependent-rectangular-overlap-budget}
In the setting of
Corollary~\ref{cor:common-sign-dependent-rectangular-overlap-budget}, replace
the pointwise norm bounds by deterministic caps
\[
  0\le\rho_{m,h,i}\le1
  \qquad(1\le h\le J_m,\ 1\le i\le n_m).
\]
Assume that there are events \(F_{m,h}^{\rm ov}\), random orthogonal projections
\[
  P_{m,h,i}^L\in\M_{d_{m,1}}(\C),\qquad
  P_{m,h,i}^R\in\M_{d_{m,2}}(\C),
\]
and cap-tail budgets \(\delta_{m,h,i}\ge0\) such that
\[
  \mathbb P\{\|X_{m,h,i}\|_{\op}>\rho_{m,h,i}\}\le\delta_{m,h,i},
\]
and, outside \(F_{m,h}^{\rm ov}\),
\[
  X_{m,h,i}=P_{m,h,i}^LX_{m,h,i}P_{m,h,i}^R
  \qquad(1\le i\le n_m)
\]
and
\[
  \max\left\{
  \left\|\sum_{i=1}^{n_m}\rho_{m,h,i}^2P_{m,h,i}^L\right\|_{\op},
  \left\|\sum_{i=1}^{n_m}\rho_{m,h,i}^2P_{m,h,i}^R\right\|_{\op}
  \right\}
  \le
  \frac{C_0n_m}{\log(eD_m^{\rm com})}.
\]
If
\[
  \sum_m\sum_{h=1}^{J_m}
  \left[
    \mathbb P(F_{m,h}^{\rm ov})
    +\sum_{i=1}^{n_m}\delta_{m,h,i}
  \right]<\infty ,
\]
then, almost surely, for all sufficiently large \(m\), one sign vector
\(\eps^{(m)}\) balances the whole active rectangular menu:
\[
  \max_{1\le h\le J_m}
  \left\|\sum_{i=1}^{n_m}\eps_i^{(m)}X_{m,h,i}\right\|_{\op}
  \le C(C_0)\sqrt{n_m}.
\]
The sampled common-certificate and countable-library conclusions of
Corollary~\ref{cor:common-sign-dependent-rectangular-overlap-budget} hold as
well.  The sub-Weibull and finite-moment cap choices in
Corollary~\ref{cor:active-menu-tail-capped-dependent-rectangular-overlap-budget}
apply verbatim, with the realized overlap budget checked at the common-sign
scale \(\log(eD_m^{\rm com})\).
\end{corollary}

\begin{proof}
Let
\[
  \Omega_{m,h}^{\rm cap}
  :=
  \{\|X_{m,h,i}\|_{\op}\le\rho_{m,h,i}\text{ for all }i\le n_m\}.
\]
The summability hypothesis implies
\[
  \sum_m\sum_{h=1}^{J_m}
  \mathbb P\bigl((\Omega_{m,h}^{\rm cap})^c\bigr)<\infty .
\]
Outside \(F_{m,h}^{\rm ov}\cup(\Omega_{m,h}^{\rm cap})^c\), the hypotheses of
Corollary~\ref{cor:common-sign-dependent-rectangular-overlap-budget} hold with
deterministic amplitudes \(\rho_{m,h,i}\).  Borel--Cantelli therefore reduces
the assertion to that corollary.  The displayed cap choices are the same
sub-Weibull tail estimate and Markov estimate used in
Corollary~\ref{cor:active-menu-tail-capped-dependent-rectangular-overlap-budget}.
\end{proof}

\begin{corollary}[Rectangular coordinate-support Bernstein]
\label{cor:rectangular-coordinate-support-bernstein}
Let \(d_1,d_2\ge1\), put \(D_{\rm rec}=d_1+d_2\), and let
\(\rho_1,\ldots,\rho_n\in[0,1]\) be deterministic.  For each \(i\), let
\[
  R_i\subset\{1,\ldots,d_1\},\qquad S_i\subset\{1,\ldots,d_2\}
\]
be random row and column supports.  Assume that, for every row coordinate \(a\),
the indicators \({\bf 1}_{\{a\in R_i\}}\), \(i=1,\ldots,n\), are independent and
\[
  \mathbb P\{a\in R_i\}\le p^L_{ia},
\]
and that, for every column coordinate \(b\), the indicators
\({\bf 1}_{\{b\in S_i\}}\), \(i=1,\ldots,n\), are independent and
\[
  \mathbb P\{b\in S_i\}\le p^R_{ib}.
\]
Conditional on the supports, let \(X_1,\ldots,X_n\) be arbitrary random
rectangular matrices satisfying
\[
  X_i=P_{R_i}X_iP_{S_i},\qquad \|X_i\|_{\op}\le\rho_i .
\]
Put
\[
  b:=\max_i\rho_i^2,
\]
\[
  \mu:=
  \max\left\{
  \max_{1\le a\le d_1}\sum_i\rho_i^2p^L_{ia},
  \max_{1\le b\le d_2}\sum_i\rho_i^2p^R_{ib}
  \right\},
\]
and
\[
  \sigma^2:=
  \max\left\{
  \max_{1\le a\le d_1}\sum_i\rho_i^4p^L_{ia},
  \max_{1\le b\le d_2}\sum_i\rho_i^4p^R_{ib}
  \right\}.
\]
Then, for every \(t>0\) and every \(C_0\ge1\) satisfying
\[
  \mu+t\le\frac{C_0n}{\log(eD_{\rm rec})},
\]
one has
\[
  \mathbb P\left\{
  \min_{\eps\in\{\pm1\}^n}
  \left\|\sum_{i=1}^n\eps_iX_i\right\|_{\op}
  >C(C_0)\sqrt n
  \right\}
  \le
  D_{\rm rec}\exp\left\{-\frac{t^2}{2(\sigma^2+bt/3)}\right\},
\]
with the convention that the exponential term is zero when
\(\sigma^2+bt/3=0\).

In particular, if \(p^L_{ia}\le p_L\), \(p^R_{ib}\le p_R\), and
\[
  \max\{p_L,p_R\}\sum_i\rho_i^2
  \le
  \frac{C_0n}{2\log(eD_{\rm rec})},
\]
then
\[
  \mathbb P\left\{
  \min_{\eps\in\{\pm1\}^n}
  \left\|\sum_i\eps_iX_i\right\|_{\op}
  >C(C_0)\sqrt n
  \right\}
  \le
  D_{\rm rec}
  \exp\left\{-\frac{3C_0n}{16\log(eD_{\rm rec})}\right\}.
\]
Uniformly chosen fixed-size row and column supports of sizes \(r_i\) and \(s_i\)
are covered by the half-budget condition
\[
  \max\left\{
  \frac1{d_1}\sum_i\rho_i^2r_i,
  \frac1{d_2}\sum_i\rho_i^2s_i
  \right\}
  \le
  \frac{C_0n}{2\log(eD_{\rm rec})}.
\]
\end{corollary}

\begin{proof}
For each row coordinate define
\[
  N_a^L:=\sum_i\rho_i^2{\bf 1}_{\{a\in R_i\}},
\]
and for each column coordinate define
\[
  N_b^R:=\sum_i\rho_i^2{\bf 1}_{\{b\in S_i\}}.
\]
By the one-coordinate independence hypotheses, scalar Bernstein gives for every
row \(a\) and every column \(b\)
\[
  \mathbb P\{N_a^L>\mu+t\},
  \mathbb P\{N_b^R>\mu+t\}
  \le
  \exp\left\{-\frac{t^2}{2(\sigma^2+bt/3)}\right\}.
\]
Union over the \(d_1+d_2\) coordinates.  On the complementary event,
\[
  \left\|\sum_i\rho_i^2P_{R_i}\right\|_{\op}\le\mu+t,
  \qquad
  \left\|\sum_i\rho_i^2P_{S_i}\right\|_{\op}\le\mu+t,
\]
so Corollary~\ref{cor:rectangular-low-overlap-support} gives the signing bound.
For the homogeneous display, take
\(t=C_0n/(2\log(eD_{\rm rec}))\).  Then \(\mu\le t\),
\(\sigma^2\le\mu\le t\), and \(b\le1\), so the exponent is at least
\(3C_0n/(16\log(eD_{\rm rec}))\).  Uniform fixed-size supports have
one-coordinate probabilities \(r_i/d_1\) and \(s_i/d_2\).
\end{proof}

\begin{corollary}[Tail-capped rectangular coordinate-support Bernstein]
\label{cor:rectangular-tail-capped-coordinate-support-bernstein}
Let \(d_1,d_2\ge1\), put \(D_{\rm rec}=d_1+d_2\), and let
\(\rho_1,\ldots,\rho_n\in[0,1]\) be deterministic caps.  For each \(i\), let
\[
  R_i\subset\{1,\ldots,d_1\},\qquad S_i\subset\{1,\ldots,d_2\}
\]
be random row and column supports satisfying the coordinatewise independence and
inclusion-probability hypotheses of
Corollary~\ref{cor:rectangular-coordinate-support-bernstein}.  Conditional on the
supports, let \(X_1,\ldots,X_n\) be arbitrary random rectangular matrices
satisfying
\[
  X_i=P_{R_i}X_iP_{S_i}.
\]
Assume that cap-tail budgets \(\delta_i\ge0\) satisfy
\[
  \mathbb P\{\|X_i\|_{\op}>\rho_i\}\le\delta_i .
\]
Define \(b\), \(\mu_{\rm cap}\), and \(\sigma_{\rm cap}^2\) from the displays in
Corollary~\ref{cor:rectangular-coordinate-support-bernstein}, with the caps
\(\rho_i\).  Then, for every \(t>0\) and every \(C_0\ge1\) satisfying
\[
  \mu_{\rm cap}+t\le\frac{C_0n}{\log(eD_{\rm rec})},
\]
one has
\[
  \begin{aligned}
  &\mathbb P\left\{
  \min_{\eps\in\{\pm1\}^n}
  \left\|\sum_i\eps_iX_i\right\|_{\op}
  >C(C_0)\sqrt n
  \right\} \\
  &\qquad\le
  \sum_{i=1}^n\delta_i
  +D_{\rm rec}
  \exp\left\{-\frac{t^2}{2(\sigma_{\rm cap}^2+bt/3)}\right\},
  \end{aligned}
\]
with the same zero-denominator convention as before.

In particular, if \(p^L_{ia}\le p_L\), \(p^R_{ib}\le p_R\), and
\[
  \max\{p_L,p_R\}\sum_i\rho_i^2
  \le
  \frac{C_0n}{2\log(eD_{\rm rec})},
\]
then the Bernstein term is at most
\[
  D_{\rm rec}
  \exp\left\{-\frac{3C_0n}{16\log(eD_{\rm rec})}\right\}.
\]
Uniformly chosen fixed-size supports have the same row/column half-budget
condition as in
Corollary~\ref{cor:rectangular-coordinate-support-bernstein}, now with the
deterministic caps \(\rho_i\).  Thus unbounded sparse rectangular support masks
are reduced to deterministic cap tails plus the capped row/column occupancy
budget.
\end{corollary}

\begin{proof}
Let
\[
  \Omega_{\rm cap}:=\{\|X_i\|_{\op}\le\rho_i\text{ for all }i\le n\}.
\]
On \(\Omega_{\rm cap}\), the hypotheses of
Corollary~\ref{cor:rectangular-coordinate-support-bernstein} hold with the caps
\(\rho_i\).  Add the probability of \(\Omega_{\rm cap}^c\), which is at most
\(\sum_i\delta_i\), and use the bounded coordinate-support Bernstein bound.  The
homogeneous and fixed-size support specializations are the corresponding
specializations of that corollary with \(t=C_0n/(2\log(eD_{\rm rec}))\).
\end{proof}

\begin{corollary}[Rectangular random-subspace Bernstein]
\label{cor:rectangular-random-subspace-bernstein}
Let \(d_1,d_2\ge1\), put \(D_{\rm rec}=d_1+d_2\), and let
\(\rho_1,\ldots,\rho_n\in[0,1]\) be deterministic.  For each \(i\), let
\[
  P_i^L\in\M_{d_1}(\C),\qquad P_i^R\in\M_{d_2}(\C)
\]
be random orthogonal projections, with the pairs \((P_i^L,P_i^R)\)
independent over \(i\).  Conditional on these projections, let
\[
  X_1,\ldots,X_n\in\M_{d_1\times d_2}(\C)
\]
be arbitrary random rectangular matrices satisfying
\[
  X_i=P_i^LX_iP_i^R,\qquad \|X_i\|_{\op}\le\rho_i .
\]
Put
\[
  b:=\max_i\rho_i^2,
\]
and
\[
  \mu:=
  \max\left\{
  \left\|\sum_i\rho_i^2\mathbb EP_i^L\right\|_{\op},
  \left\|\sum_i\rho_i^2\mathbb EP_i^R\right\|_{\op}
  \right\}.
\]
Then, for every \(t>0\) and every \(C_0\ge1\) satisfying
\[
  \mu+t\le \frac{C_0n}{\log(eD_{\rm rec})},
\]
one has
\[
  \mathbb P\left\{
  \min_{\eps\in\{\pm1\}^n}
  \left\|\sum_i\eps_iX_i\right\|_{\op}
  >C(C_0)\sqrt n
  \right\}
  \le
  D_{\rm rec}\exp\left\{-\frac{t^2}{2(\mu+bt/3)}\right\},
\]
with zero contribution when the denominator is zero.

In particular, if \(P_i^L\) and \(P_i^R\) are Haar-distributed projections of
ranks \(r_i^L\) and \(r_i^R\), and
\[
  \max\left\{
  \frac1{d_1}\sum_i\rho_i^2r_i^L,
  \frac1{d_2}\sum_i\rho_i^2r_i^R
  \right\}
  \le
  \frac{C_0n}{2\log(eD_{\rm rec})},
\]
then
\[
  \mathbb P\left\{
  \min_{\eps\in\{\pm1\}^n}
  \left\|\sum_i\eps_iX_i\right\|_{\op}
  >C(C_0)\sqrt n
  \right\}
  \le
  D_{\rm rec}
  \exp\left\{-\frac{3C_0n}{16\log(eD_{\rm rec})}\right\}.
\]
Thus rectangular random low-rank, row-column subspace-localized, and
rectangular GOE/GUE/Wigner-block models are Spencer-scale whenever the weighted
average row and column ranks are \(O(d_1/\log(eD_{\rm rec}))\) and
\(O(d_2/\log(eD_{\rm rec}))\), respectively.
\end{corollary}

\begin{proof}
Apply matrix Bernstein separately on the row and column sides.  For the row
side set
\[
  Y_i^L:=\rho_i^2(P_i^L-\mathbb EP_i^L).
\]
Then \(\mathbb EY_i^L=0\), \(\|Y_i^L\|_{\op}\le b\), and
\[
  \sum_i\mathbb E(Y_i^L)^2
  =
  \sum_i\rho_i^4\left(\mathbb EP_i^L-(\mathbb EP_i^L)^2\right)
  \preceq
  \sum_i\rho_i^2\mathbb EP_i^L,
\]
so the row variance parameter is at most \(\mu\).  The same argument applies to
the column projections.  A union bound over the row and column Bernstein tails
therefore gives
\[
  \mathbb P\left\{
  \max\left[
  \left\|\sum_i\rho_i^2P_i^L\right\|_{\op},
  \left\|\sum_i\rho_i^2P_i^R\right\|_{\op}
  \right]>\mu+t
  \right\}
  \le
  D_{\rm rec}\exp\left\{-\frac{t^2}{2(\mu+bt/3)}\right\}.
\]
On the complementary event, Corollary~\ref{cor:rectangular-low-overlap-support}
applies.  For Haar ranks,
\[
  \mathbb EP_i^L=\frac{r_i^L}{d_1}I_{d_1},\qquad
  \mathbb EP_i^R=\frac{r_i^R}{d_2}I_{d_2}.
\]
Taking \(t=C_0n/(2\log(eD_{\rm rec}))\) and using \(\mu\le t\), \(b\le1\),
gives the final displayed tail.
\end{proof}

\begin{corollary}[Tail-capped rectangular random subspaces]
\label{cor:rectangular-tail-capped-random-subspace-bernstein}
Let \(d_1,d_2\ge1\), put \(D_{\rm rec}=d_1+d_2\), and let
\(\rho_1,\ldots,\rho_n\in[0,1]\) be deterministic caps.  For each \(i\), let
\[
  P_i^L\in\M_{d_1}(\C),\qquad P_i^R\in\M_{d_2}(\C)
\]
be random orthogonal projections, with the pairs \((P_i^L,P_i^R)\)
independent over \(i\).  Conditional on these projections, let
\[
  X_1,\ldots,X_n\in\M_{d_1\times d_2}(\C)
\]
be arbitrary random rectangular matrices satisfying
\[
  X_i=P_i^LX_iP_i^R .
\]
Put
\[
  b:=\max_i\rho_i^2,
\]
and
\[
  \mu_{\rm cap}:=
  \max\left\{
  \left\|\sum_i\rho_i^2\mathbb EP_i^L\right\|_{\op},
  \left\|\sum_i\rho_i^2\mathbb EP_i^R\right\|_{\op}
  \right\}.
\]
Then, for every \(t>0\) and every \(C_0\ge1\) satisfying
\[
  \mu_{\rm cap}+t\le \frac{C_0n}{\log(eD_{\rm rec})},
\]
one has
\[
  \begin{aligned}
  &\mathbb P\left\{
  \min_{\eps\in\{\pm1\}^n}
  \left\|\sum_i\eps_iX_i\right\|_{\op}
  >C(C_0)\sqrt n
  \right\}  \\
  &\qquad\le
  \mathbb P\{\exists i:\ \|X_i\|_{\op}>\rho_i\}
  +
  D_{\rm rec}\exp\left\{-\frac{t^2}{2(\mu_{\rm cap}+bt/3)}\right\}.
  \end{aligned}
\]
In particular, the first term may be bounded by
\(\sum_i\mathbb P\{\|X_i\|_{\op}>\rho_i\}\), with no independence assumption on
the matrix amplitudes.  If \(P_i^L\) and \(P_i^R\) are Haar-distributed
projections of ranks \(r_i^L\) and \(r_i^R\), and
\[
  \max\left\{
  \frac1{d_1}\sum_i\rho_i^2r_i^L,
  \frac1{d_2}\sum_i\rho_i^2r_i^R
  \right\}
  \le
  \frac{C_0n}{2\log(eD_{\rm rec})},
\]
then
\[
  \mathbb P\left\{
  \min_{\eps\in\{\pm1\}^n}
  \left\|\sum_i\eps_iX_i\right\|_{\op}
  >C(C_0)\sqrt n
  \right\}
  \le
  \sum_i\mathbb P\{\|X_i\|_{\op}>\rho_i\}
  +
  D_{\rm rec}
  \exp\left\{-\frac{3C_0n}{16\log(eD_{\rm rec})}\right\}.
\]
Thus rectangular random low-rank, random row/column subspace-localized, and
random block-factor models remain Spencer-scale with unbounded internal
coefficients whenever the operator-norm cap tails and capped average row/column
rank profiles are controlled.
\end{corollary}

\begin{proof}
On the event
\[
  \Omega_{\rm cap}:=\{\|X_i\|_{\op}\le\rho_i\text{ for all }i\},
\]
Corollary~\ref{cor:rectangular-random-subspace-bernstein} applies with the same
projections and caps.  Add \(\mathbb P(\Omega_{\rm cap}^c)\), then use the
union bound for the displayed tail-sum variant.  The Haar-rank specialization is
the corresponding specialization in
Corollary~\ref{cor:rectangular-random-subspace-bernstein}.
\end{proof}

\begin{corollary}[Random rectangular dyads from frame laws]
\label{cor:rectangular-random-dyads-frame}
Let \(d_1,d_2\ge1\), put \(D_{\rm rec}=d_1+d_2\), and let
\(\rho_1,\ldots,\rho_n\in[0,1]\) be deterministic.  For each \(i\), let
\((u_i,v_i)\in\mathbb C^{d_1}\times\mathbb C^{d_2}\) be a random pair of unit
vectors, with the pairs independent over \(i\).  Conditional on these vectors,
let \(a_i\) be an arbitrary scalar random variable with \(|a_i|\le\rho_i\), and
set
\[
  X_i:=a_iu_iv_i^* \in\M_{d_1\times d_2}(\C).
\]
Put
\[
  b:=\max_i\rho_i^2,\qquad
  \mu_{\rm dyad}:=
  \max\left\{
  \left\|\sum_i\rho_i^2\mathbb E u_iu_i^*\right\|_{\op},
  \left\|\sum_i\rho_i^2\mathbb E v_iv_i^*\right\|_{\op}
  \right\}.
\]
Then, for every \(t>0\) and \(C_0\ge1\) satisfying
\[
  \mu_{\rm dyad}+t\le \frac{C_0n}{\log(eD_{\rm rec})},
\]
one has
\[
  \mathbb P\left\{
  \min_{\eps\in\{\pm1\}^n}
  \left\|\sum_i\eps_iX_i\right\|_{\op}
  >C(C_0)\sqrt n
  \right\}
  \le
  D_{\rm rec}\exp\left\{-\frac{t^2}{2(\mu_{\rm dyad}+bt/3)}\right\}.
\]
In the Haar-isotropic case,
\[
  \mathbb E u_iu_i^*=\frac1{d_1}I_{d_1},\qquad
  \mathbb E v_iv_i^*=\frac1{d_2}I_{d_2},
\]
so the half-budget condition becomes
\[
  \max\left\{\frac1{d_1}\sum_i\rho_i^2,\,
             \frac1{d_2}\sum_i\rho_i^2\right\}
  \le \frac{C_0n}{2\log(eD_{\rm rec})},
\]
and under this condition the failure probability is at most
\[
  D_{\rm rec}
  \exp\left\{-\frac{3C_0n}{16\log(eD_{\rm rec})}\right\}.
\]
Thus random rank-one measurement matrices, randomized feature dyads, and
nonuniformly sampled rectangular atom dictionaries are Spencer-scale whenever
their two weighted expected frame operators are \(O(n/\log(eD_{\rm rec}))\).
\end{corollary}

\begin{proof}
Let
\[
  P_i^L:=u_iu_i^*,\qquad P_i^R:=v_iv_i^* .
\]
These are rank-one orthogonal projections, independent as pairs over \(i\), and
\(X_i=P_i^LX_iP_i^R\) with \(\|X_i\|_{\op}=|a_i|\le\rho_i\).  The assertion is
therefore Corollary~\ref{cor:rectangular-random-subspace-bernstein}.  The
Haar-isotropic displays are the rank-one Haar specialization of that corollary.
\end{proof}

\begin{corollary}[Tail-capped rectangular dyads]
\label{cor:rectangular-tail-capped-random-dyads-frame}
Let \(d_1,d_2\ge1\), put \(D_{\rm rec}=d_1+d_2\), and let
\(\rho_1,\ldots,\rho_n\in[0,1]\) be deterministic caps.  For each \(i\), let
\((u_i,v_i)\in\mathbb C^{d_1}\times\mathbb C^{d_2}\) be a random pair of unit
vectors, with the pairs independent over \(i\).  Conditional on these vectors,
let \(a_i\) be an arbitrary scalar random variable, and put
\[
  X_i:=a_iu_iv_i^* .
\]
Set
\[
  b:=\max_i\rho_i^2,\qquad
  \mu_{\rm cap}:=
  \max\left\{
  \left\|\sum_i\rho_i^2\mathbb E u_iu_i^*\right\|_{\op},
  \left\|\sum_i\rho_i^2\mathbb E v_iv_i^*\right\|_{\op}
  \right\}.
\]
Then, for every \(t>0\) and \(C_0\ge1\) satisfying
\[
  \mu_{\rm cap}+t\le \frac{C_0n}{\log(eD_{\rm rec})},
\]
one has
\[
  \begin{aligned}
  &\mathbb P\left\{
  \min_{\eps\in\{\pm1\}^n}
  \left\|\sum_i\eps_iX_i\right\|_{\op}
  >C(C_0)\sqrt n
  \right\}  \\
  &\qquad\le
  \mathbb P\{\exists i:\ |a_i|>\rho_i\}
  +
  D_{\rm rec}\exp\left\{-\frac{t^2}{2(\mu_{\rm cap}+bt/3)}\right\}.
  \end{aligned}
\]
In particular, the first term may be bounded by
\(\sum_i\mathbb P\{|a_i|>\rho_i\}\), with no independence assumption on the
amplitudes.  If the dyad directions are Haar-isotropic and
\[
  \max\left\{\frac1{d_1}\sum_i\rho_i^2,\,
             \frac1{d_2}\sum_i\rho_i^2\right\}
  \le \frac{C_0n}{2\log(eD_{\rm rec})},
\]
then
\[
  \mathbb P\left\{
  \min_{\eps\in\{\pm1\}^n}
  \left\|\sum_i\eps_iX_i\right\|_{\op}
  >C(C_0)\sqrt n
  \right\}
  \le
  \sum_i\mathbb P\{|a_i|>\rho_i\}
  +
  D_{\rm rec}
  \exp\left\{-\frac{3C_0n}{16\log(eD_{\rm rec})}\right\}.
\]
Thus Gaussian, sub-Gaussian, sub-exponential, or heavy-tailed scalar weights
enter rectangular dyad and randomized-feature models through any deterministic
cap profile whose tail cost is acceptable and whose capped expected frame
operators have the Matrix Spencer scale.
\end{corollary}

\begin{proof}
On the event \(\{|a_i|\le\rho_i\ \text{for all }i\}\), apply
Corollary~\ref{cor:rectangular-random-dyads-frame} with the same dyad
directions and caps.  Add the probability of the complementary event.  The
union bound gives the displayed summable-tail variant, and the Haar-isotropic
case is the homogeneous specialization from
Corollary~\ref{cor:rectangular-random-dyads-frame}.
\end{proof}

\begin{corollary}[Active menus of rectangular coordinate-support streams]
\label{cor:active-menu-rectangular-coordinate-support-schedules}
Fix \(C_0\ge1\).  For each \(m\ge1\), let \(J_m\ge1\),
\(d_{m,1},d_{m,2}\ge1\), \(D_m^{\rm rec}=d_{m,1}+d_{m,2}\), and \(n_m\ge1\).
For every \(1\le h\le J_m\), let
\[
  X_{m,h,1},\ldots,X_{m,h,n_m}
  \in\M_{d_{m,1}\times d_{m,2}}(\C)
\]
be a rectangular coordinate-support model as in
Corollary~\ref{cor:rectangular-coordinate-support-bernstein}, with parameters
\(\rho_{m,h,i}\), \(b_{m,h}\), \(\mu_{m,h}\), and \(\sigma_{m,h}^2\).  Suppose
there are \(t_{m,h}>0\) such that
\[
  \mu_{m,h}+t_{m,h}
  \le
  \frac{C_0n_m}{\log(eD_m^{\rm rec})}
  \qquad(1\le h\le J_m),
\]
and
\[
  \sum_m\sum_{h=1}^{J_m}
  D_m^{\rm rec}
  \exp\left\{
  -\frac{t_{m,h}^2}
         {2(\sigma_{m,h}^2+b_{m,h}t_{m,h}/3)}
  \right\}<\infty,
\]
again with zero denominator interpreted as zero contribution.  Then, almost
surely, for all sufficiently large \(m\), every active rectangular
coordinate-support stream has a Spencer-scale signing:
\[
  \min_{\eps\in\{\pm1\}^{n_m}}
  \left\|\sum_{i=1}^{n_m}\eps_iX_{m,h,i}\right\|_{\op}
  \le C(C_0)\sqrt{n_m}
  \qquad(1\le h\le J_m).
\]
Sampled certificates are eventual almost surely under the same additional
summability conditions as in
Corollary~\ref{cor:active-menu-rectangular-square-profile-schedules}.  If
\(J_m\to\infty\), active prefixes exhaust countable libraries of rectangular
bipartite sparse-support models.
\end{corollary}

\begin{proof}
Apply Corollary~\ref{cor:rectangular-coordinate-support-bernstein} to each
active stream and sum the displayed failure probabilities.  The first
Borel--Cantelli lemma gives the eventual simultaneous signing.  The
sampled-search and countable-library statements are the same additions as in
Corollary~\ref{cor:active-menu-rectangular-square-profile-schedules}.
\end{proof}

\begin{corollary}[Active menus of tail-capped rectangular coordinate-support streams]
\label{cor:active-menu-rectangular-tail-capped-coordinate-support-schedules}
Fix \(C_0\ge1\).  For each \(m\ge1\), let \(J_m\ge1\),
\(d_{m,1},d_{m,2}\ge1\), \(D_m^{\rm rec}=d_{m,1}+d_{m,2}\), and \(n_m\ge1\).
For every \(1\le h\le J_m\), let
\[
  X_{m,h,1},\ldots,X_{m,h,n_m}
  \in\M_{d_{m,1}\times d_{m,2}}(\C)
\]
be a tail-capped rectangular coordinate-support model as in
Corollary~\ref{cor:rectangular-tail-capped-coordinate-support-bernstein}, with
deterministic caps \(\rho_{m,h,i}\), parameters \(b_{m,h}\),
\(\mu_{m,h}^{\rm cap}\), \(\sigma_{m,h,{\rm cap}}^2\), and tail budgets
\(\delta_{m,h,i}\ge0\).  Suppose there are \(t_{m,h}>0\) such that
\[
  \mu_{m,h}^{\rm cap}+t_{m,h}
  \le
  \frac{C_0n_m}{\log(eD_m^{\rm rec})}
  \qquad(1\le h\le J_m),
\]
and
\[
  \sum_m\sum_{h=1}^{J_m}
  \left[
  \sum_{i=1}^{n_m}\delta_{m,h,i}
  +
  D_m^{\rm rec}
  \exp\left\{
  -\frac{t_{m,h}^2}
         {2(\sigma_{m,h,{\rm cap}}^2+b_{m,h}t_{m,h}/3)}
  \right\}
  \right]<\infty,
\]
again with zero denominator interpreted as zero contribution.  Then, almost
surely, for all sufficiently large \(m\), every active tail-capped rectangular
coordinate-support stream has a Spencer-scale signing:
\[
  \min_{\eps\in\{\pm1\}^{n_m}}
  \left\|\sum_{i=1}^{n_m}\eps_iX_{m,h,i}\right\|_{\op}
  \le C(C_0)\sqrt{n_m}
  \qquad(1\le h\le J_m).
\]
Sampled certificates and countable-library exhaustions follow under the same
additional summability hypotheses as in
Corollary~\ref{cor:active-menu-rectangular-coordinate-support-schedules}.

In the half-budget case
\[
  \max\left\{
  \max_a\sum_i\rho_{m,h,i}^2p^L_{m,h,i,a},
  \max_b\sum_i\rho_{m,h,i}^2p^R_{m,h,i,b}
  \right\}
  \le \frac{C_0n_m}{2\log(eD_m^{\rm rec})},
\]
the matrix part of the active-menu failure is bounded by
\[
  J_mD_m^{\rm rec}
  \exp\left\{-\frac{3C_0n_m}{16\log(eD_m^{\rm rec})}\right\}.
\]
Consequently the arbitrary-growth, all-integer, sparse, and sampled-certificate
schedules of
Corollary~\ref{cor:active-menu-rectangular-weighted-support-explicit-schedules}
apply verbatim after adding the cap-tail summability
\[
  \sum_m\sum_{h=1}^{J_m}\sum_{i=1}^{n_m}\delta_{m,h,i}<\infty .
\]
\end{corollary}

\begin{proof}
Apply
Corollary~\ref{cor:rectangular-tail-capped-coordinate-support-bernstein} to each
active stream and sum its cap-tail and Bernstein failure probabilities.  The
first Borel--Cantelli lemma gives eventual simultaneous signings.  Sampled signs
add the same conditional search failures as in the bounded coordinate-support
case, because on the cap event the proof reduces to
Corollary~\ref{cor:rectangular-low-overlap-support}.  Under the half-budget
hypothesis, take \(t_{m,h}=C_0n_m/(2\log(eD_m^{\rm rec}))\); then
\(\sigma_{m,h,{\rm cap}}^2\le \mu_{m,h}^{\rm cap}\le t_{m,h}\) and
\(b_{m,h}\le1\), giving the displayed Bernstein exponent and the inherited
growth schedules.
\end{proof}

\begin{corollary}[Sub-Weibull and moment caps for coordinate-support menus]
\label{cor:active-menu-rectangular-tail-selected-coordinate-support-schedules}
Fix \(C_0\ge1\).  Work in the active-menu coordinate-support setting of
Corollary~\ref{cor:active-menu-rectangular-tail-capped-coordinate-support-schedules},
and write \(p^L_{m,h,i,a}\) and \(p^R_{m,h,i,b}\) for the row and column
inclusion-probability envelopes.  For each active triple \((m,h,i)\), choose a
level \(\ell_{m,h,i}>0\) and a scale \(s_{m,h,i}>0\).  Assume that the
deterministic cap \(\rho_{m,h,i}\le1\) is chosen by one of the following two
rules.

\begin{enumerate}[label=(\alph*)]
\item Sub-Weibull operator tails: for some fixed \(\alpha>0\) and
  \(A_{\rm tail}\ge1\),
  \[
    \mathbb P\{\|X_{m,h,i}\|_{\op}>s_{m,h,i}u\}
    \le A_{\rm tail}\exp\{-u^\alpha\}\qquad(u\ge0),
  \]
  and
  \[
    \rho_{m,h,i}:=s_{m,h,i}\ell_{m,h,i}^{1/\alpha}\le1.
  \]
\item Moment operator tails: for some fixed \(p>0\),
  \[
    \mathbb E\|X_{m,h,i}\|_{\op}^p\le s_{m,h,i}^p,
  \]
  and
  \[
    \rho_{m,h,i}:=s_{m,h,i}\exp\{\ell_{m,h,i}/p\}\le1.
  \]
\end{enumerate}

In case (a) put \(A_*=A_{\rm tail}\), and in case (b) put \(A_*=1\).  Let
\(b_{m,h}:=\max_i\rho_{m,h,i}^2\), and define the capped coordinate budgets
\[
  \mu_{m,h}^{\rm csel}:=
  \max\left\{
  \max_a\sum_i\rho_{m,h,i}^2p^L_{m,h,i,a},
  \max_b\sum_i\rho_{m,h,i}^2p^R_{m,h,i,b}
  \right\},
\]
and
\[
  (\sigma_{m,h}^{\rm csel})^2:=
  \max\left\{
  \max_a\sum_i\rho_{m,h,i}^4p^L_{m,h,i,a},
  \max_b\sum_i\rho_{m,h,i}^4p^R_{m,h,i,b}
  \right\}.
\]
If there are \(t_{m,h}>0\) such that
\[
  \mu_{m,h}^{\rm csel}+t_{m,h}
  \le
  \frac{C_0n_m}{\log(eD_m^{\rm rec})}
  \qquad(1\le h\le J_m),
\]
and
\[
  \sum_m\sum_{h=1}^{J_m}
  \left[
  A_*\sum_{i=1}^{n_m}\exp\{-\ell_{m,h,i}\}
  +
  D_m^{\rm rec}
  \exp\left\{
  -\frac{t_{m,h}^2}
         {2((\sigma_{m,h}^{\rm csel})^2+b_{m,h}t_{m,h}/3)}
  \right\}
  \right]<\infty,
\]
with zero denominator interpreted as zero contribution, then, almost surely,
for all sufficiently large \(m\), every active coordinate-support stream has a
Spencer-scale signing:
\[
  \min_{\eps\in\{\pm1\}^{n_m}}
  \left\|\sum_{i=1}^{n_m}\eps_iX_{m,h,i}\right\|_{\op}
  \le C(C_0)\sqrt{n_m}
  \qquad(1\le h\le J_m).
\]
Sampled certificates and countable-library exhaustions follow under the same
additional summability hypotheses as in
Corollary~\ref{cor:active-menu-rectangular-tail-capped-coordinate-support-schedules}.

For uniformly chosen fixed-size row and column supports, the two displayed
budgets are obtained by substituting
\[
  p^L_{m,h,i,a}=r^L_{m,h,i}/d_{m,1},
  \qquad
  p^R_{m,h,i,b}=r^R_{m,h,i}/d_{m,2}.
\]
Thus sub-Weibull sparse rectangular blocks pay \(\ell^{2/\alpha}\) in the
coordinate-load budget and \(\ell^{4/\alpha}\) in the Bernstein variance proxy,
while finite \(p\)-moment blocks pay \(\exp\{2\ell/p\}\) and
\(\exp\{4\ell/p\}\), respectively.

Consequently, if \(J_mn_m\le(eD_m^{\rm rec})^\Gamma\), the constants in the
tail bounds are uniform, and one takes the common level
\(\ell_{m,h,i}=R\log(eD_m^{\rm rec})\), then the cap-tail part of the series is
summable along the full integer dimension schedule \(D_m^{\rm rec}=m\) whenever
\(R>\Gamma+1\), and along every stretched-exponential sparse schedule whenever
\(R>\Gamma\).  The remaining condition is exactly the displayed capped
coordinate-occupancy Bernstein summability.
\end{corollary}

\begin{proof}
In the sub-Weibull case,
\[
  \mathbb P\{\|X_{m,h,i}\|_{\op}>\rho_{m,h,i}\}
  \le A_{\rm tail}\exp\{-\ell_{m,h,i}\}.
\]
In the moment case, Markov's inequality gives
\[
  \mathbb P\{\|X_{m,h,i}\|_{\op}>\rho_{m,h,i}\}
  \le
  \frac{s_{m,h,i}^p}{s_{m,h,i}^p\exp\{\ell_{m,h,i}\}}
  =
  \exp\{-\ell_{m,h,i}\}.
\]
Use these quantities as the tail budgets
\(\delta_{m,h,i}\) in
Corollary~\ref{cor:active-menu-rectangular-tail-capped-coordinate-support-schedules}.
The fixed-size support formulas are the coordinate-support probabilities
\(r^L_{m,h,i}/d_{m,1}\) and \(r^R_{m,h,i}/d_{m,2}\).  The final all-integer and
sparse cap-tail thresholds follow from summing
\((eD_m^{\rm rec})^{\Gamma-R}\).
\end{proof}

\begin{corollary}[Active menus of rectangular random-subspace streams]
\label{cor:active-menu-rectangular-random-subspace-schedules}
Fix \(C_0\ge1\).  For each \(m\ge1\), let \(J_m\ge1\),
\(d_{m,1},d_{m,2}\ge1\), \(D_m^{\rm rec}=d_{m,1}+d_{m,2}\), and \(n_m\ge1\).
For every \(1\le h\le J_m\), let
\[
  X_{m,h,1},\ldots,X_{m,h,n_m}
  \in\M_{d_{m,1}\times d_{m,2}}(\C)
\]
be a rectangular random-subspace model as in
Corollary~\ref{cor:rectangular-random-subspace-bernstein}, with parameters
\(b_{m,h}\) and \(\mu_{m,h}\).  Suppose there are \(t_{m,h}>0\) such that
\[
  \mu_{m,h}+t_{m,h}
  \le
  \frac{C_0n_m}{\log(eD_m^{\rm rec})}
  \qquad(1\le h\le J_m),
\]
and
\[
  \sum_m\sum_{h=1}^{J_m}
  D_m^{\rm rec}
  \exp\left\{
  -\frac{t_{m,h}^2}{2(\mu_{m,h}+b_{m,h}t_{m,h}/3)}
  \right\}<\infty,
\]
with zero denominator interpreted as zero contribution.  Then, almost surely,
for all sufficiently large \(m\), every active rectangular random-subspace stream
has a Spencer-scale signing:
\[
  \min_{\eps\in\{\pm1\}^{n_m}}
  \left\|\sum_{i=1}^{n_m}\eps_iX_{m,h,i}\right\|_{\op}
  \le C(C_0)\sqrt{n_m}
  \qquad(1\le h\le J_m).
\]
Sampled certificates and countable-library exhaustions follow under the same
additional summability hypotheses as in
Corollary~\ref{cor:active-menu-rectangular-square-profile-schedules}.
\end{corollary}

\begin{proof}
Apply Corollary~\ref{cor:rectangular-random-subspace-bernstein} to each active
stream and sum the displayed failure probabilities.  The first
Borel--Cantelli lemma gives the eventual simultaneous signing, and the sampled
and countable-library additions are the same as before.
\end{proof}

\begin{corollary}[Explicit rectangular weighted-support schedules]
\label{cor:active-menu-rectangular-weighted-support-explicit-schedules}
Fix \(C_0\ge1\), put \(c_0:=3C_0/16\), and, for each row \(m\), let
\[
  J_m\ge1,\qquad d_{m,1},d_{m,2}\ge1,\qquad n_m\ge1,
\]
and let
\[
  X_{m,h,1},\ldots,X_{m,h,n_m}
  \in \M_{d_{m,1}\times d_{m,2}}(\C)
  \qquad(1\le h\le J_m).
\]
Set
\[
  D_m^{\rm rec}:=d_{m,1}+d_{m,2},\qquad
  \ell_m:=\log(eD_m^{\rm rec}).
\]
For every active stream \(1\le h\le J_m\), assume that one of the following
half-budget models holds.
\begin{enumerate}[label=\textup{(\alph*)}]
\item Rectangular coordinate supports as in
  Corollary~\ref{cor:rectangular-coordinate-support-bernstein}, with
  deterministic amplitudes \(\rho_{m,h,i}\in[0,1]\), and
  \[
    \max\left\{
    \max_a\sum_i\rho_{m,h,i}^2p^L_{m,h,i,a},
    \max_b\sum_i\rho_{m,h,i}^2p^R_{m,h,i,b}
    \right\}
    \le \frac{C_0n_m}{2\ell_m}.
  \]
\item Rectangular random row/column subspaces as in
  Corollary~\ref{cor:rectangular-random-subspace-bernstein}, with independent
  projection pairs and deterministic amplitudes \(\rho_{m,h,i}\in[0,1]\), and
  \[
    \max\left\{
    \left\|\sum_i\rho_{m,h,i}^2\mathbb EP^L_{m,h,i}\right\|_{\op},
    \left\|\sum_i\rho_{m,h,i}^2\mathbb EP^R_{m,h,i}\right\|_{\op}
    \right\}
    \le \frac{C_0n_m}{2\ell_m}.
  \]
\end{enumerate}
Then, for all large \(m\),
\[
  \begin{aligned}
  &\mathbb P\left\{
  \exists h\le J_m:
  \min_{\eps\in\{\pm1\}^{n_m}}
  \left\|\sum_{i=1}^{n_m}\eps_iX_{m,h,i}\right\|_{\op}
  >C(C_0)\sqrt{n_m}
  \right\} \\
  &\qquad\le
  J_mD_m^{\rm rec}\exp\left\{-c_0\,\frac{n_m}{\ell_m}\right\}.
  \end{aligned}
\]
Consequently, if
\[
  \sum_m
  J_mD_m^{\rm rec}\exp\left\{-c_0\,\frac{n_m}{\ell_m}\right\}<\infty,
\]
then, almost surely, for all sufficiently large \(m\), every active rectangular
coordinate-support or random-subspace stream has a Spencer-scale signing.

The structural series is summable in each of the following explicit schedules.
\begin{enumerate}[label=\textup{(\roman*)}]
\item Arbitrary growth: for some \(\gamma>0\),
  \[
    n_m\ge
    \frac1{c_0}\,\ell_m\bigl(\log(eJ_mD_m^{\rm rec})
      +(1+\gamma)\log(e+m)\bigr)
  \]
  eventually.
\item All rectangular dimensions: if \(D_m^{\rm rec}=D=2,3,\ldots\),
  \(J_D\le(eD)^\Gamma\), and
  \[
    n_D\ge a\log^2(eD)
  \]
  eventually for some \(a>(\Gamma+2)/c_0\), then the conclusion holds along all
  integer \(D\).
\item Sparse rectangular dimensions: if
  \(D_m^{\rm rec}\ge\exp(c m^\kappa)\) eventually, \(J_m\le(eD_m^{\rm rec})^\Gamma\),
  and
  \[
    n_m\ge a\log^2(eD_m^{\rm rec})
  \]
  eventually for some \(a>(\Gamma+1)/c_0\), then the conclusion holds along the
  sparse schedule.
\end{enumerate}

The same assertions hold with sampled sign certificates.  More precisely, for
each active stream draw \(M_{m,h}\) independent uniform sign vectors, independent
of the matrix array.  If
\[
  \sum_m\left[
  J_mD_m^{\rm rec}\exp\left\{-c_0\,\frac{n_m}{\ell_m}\right\}
  +\sum_{h=1}^{J_m}2^{-M_{m,h}}
  \right]<\infty,
\]
then, almost surely for all sufficiently large \(m\), each active stream has one
sampled sign vector certifying the same \(C(C_0)\sqrt{n_m}\) bound.  In the
common batch-size case \(M_{m,h}=M_m\), it is enough to add the summability of
\(\sum_mJ_m2^{-M_m}\).  Thus the arbitrary-growth schedule is sampled-certified
whenever
\[
  M_m\ge \log_2(eJ_m)+(1+\delta)\log_2(e+m)
\]
eventually, and, under \(J_D\le(eD)^\Gamma\), the all-dimension schedule is
sampled-certified whenever \(M_D\ge(\Gamma+1+\delta)\log_2(eD)\) eventually.
Under the sparse schedule it is enough to take
\[
  M_m\ge \log_2(eJ_m)+(1+\delta)\log_2(e+m)
\]
eventually, or the displayed bound with
\(\log_2(eJ_m)\) replaced by \(\Gamma\log_2(eD_m^{\rm rec})\).

Uniform fixed-size row/column supports are included by the half-budget condition
\[
  \max\left\{
  \frac1{d_{m,1}}\sum_i\rho_{m,h,i}^2r^L_{m,h,i},
  \frac1{d_{m,2}}\sum_i\rho_{m,h,i}^2r^R_{m,h,i}
  \right\}
  \le \frac{C_0n_m}{2\ell_m},
\]
and Haar row/column subspaces use the identical formula with \(r^L,r^R\) the
projection ranks.
\end{corollary}

\begin{proof}
For a fixed active stream, take
\[
  t_m=\frac{C_0n_m}{2\ell_m}.
\]
In the coordinate-support model the half-budget hypothesis gives \(\mu_{m,h}\le
t_m\), while
\(\sigma_{m,h}^2\le\mu_{m,h}\le t_m\) and \(b_{m,h}\le1\).  Hence
Corollary~\ref{cor:rectangular-coordinate-support-bernstein} gives failure at
most
\[
  D_m^{\rm rec}\exp\left\{-c_0\,\frac{n_m}{\ell_m}\right\}.
\]
The random-subspace model is the same calculation with
Corollary~\ref{cor:rectangular-random-subspace-bernstein}.  Union over
\(h\le J_m\) gives the displayed active-menu tail, and Borel--Cantelli gives the
summable-array conclusion.

For arbitrary growth, the lower bound on \(n_m\) makes the active-menu tail at
most \(e^{-1}(e+m)^{-(1+\gamma)}\).  In the all-dimension schedule it is bounded
by
\[
  (eD)^{\Gamma+1-c_0a},
\]
which is summable over \(D\ge2\) when \(c_0a>\Gamma+2\).  In the sparse schedule
the same bound is stretched-exponentially summable when \(c_0a>\Gamma+1\).

For sampled certificates, condition on the complement of the coordinate or
projection overlap-failure event for a fixed active stream.  The proof of
Corollary~\ref{cor:rectangular-low-overlap-support} gives the deterministic
rectangular square-profile certificate, including sampled-search failure
\(2^{-M_{m,h}}\) after enlarging \(C(C_0)\).  Summing these failures over the
active streams gives the sampled Borel--Cantelli statement.  The displayed
choices of \(M_m\) make \(J_m2^{-M_m}\) summable in the arbitrary-growth,
all-dimension, and sparse cases.  The fixed-size support and Haar-rank formulas
are the corresponding one-coordinate probabilities and projection expectations.
\end{proof}

\begin{corollary}[Active menus of tail-capped rectangular random-subspace streams]
\label{cor:active-menu-rectangular-tail-capped-random-subspace-schedules}
Fix \(C_0\ge1\).  For each \(m\ge1\), let \(J_m\ge1\),
\(d_{m,1},d_{m,2}\ge1\), \(D_m^{\rm rec}=d_{m,1}+d_{m,2}\), and \(n_m\ge1\).
For every \(1\le h\le J_m\), let
\[
  X_{m,h,1},\ldots,X_{m,h,n_m}
  \in\M_{d_{m,1}\times d_{m,2}}(\C)
\]
be a rectangular random-subspace model as in
Corollary~\ref{cor:rectangular-tail-capped-random-subspace-bernstein}, with
deterministic caps \(\rho_{m,h,i}\), parameters \(b_{m,h}\) and
\(\mu_{m,h}^{\rm cap}\), and tail budgets \(\delta_{m,h,i}\ge0\) satisfying
\[
  \mathbb P\{\|X_{m,h,i}\|_{\op}>\rho_{m,h,i}\}\le\delta_{m,h,i}.
\]
Suppose there are \(t_{m,h}>0\) such that
\[
  \mu_{m,h}^{\rm cap}+t_{m,h}
  \le
  \frac{C_0n_m}{\log(eD_m^{\rm rec})}
  \qquad(1\le h\le J_m),
\]
and
\[
  \sum_m\sum_{h=1}^{J_m}
  \left[
  \sum_{i=1}^{n_m}\delta_{m,h,i}
  +
  D_m^{\rm rec}
  \exp\left\{
  -\frac{t_{m,h}^2}{2(\mu_{m,h}^{\rm cap}+b_{m,h}t_{m,h}/3)}
  \right\}
  \right]<\infty,
\]
with zero denominator interpreted as zero contribution.  Then, almost surely,
for all sufficiently large \(m\), every active tail-capped rectangular
random-subspace stream has a Spencer-scale signing:
\[
  \min_{\eps\in\{\pm1\}^{n_m}}
  \left\|\sum_{i=1}^{n_m}\eps_iX_{m,h,i}\right\|_{\op}
  \le C(C_0)\sqrt{n_m}
  \qquad(1\le h\le J_m).
\]
Sampled certificates and countable-library exhaustions follow under the same
additional summability hypotheses as in
Corollary~\ref{cor:active-menu-rectangular-square-profile-schedules}.
\end{corollary}

\begin{proof}
Apply Corollary~\ref{cor:rectangular-tail-capped-random-subspace-bernstein} to
each active stream, bound its cap-exception event by
\(\sum_i\delta_{m,h,i}\), and sum the displayed probabilities.  The first
Borel--Cantelli lemma gives eventual simultaneous signings.  The sampled and
countable-library additions are the same as in
Corollary~\ref{cor:active-menu-rectangular-random-subspace-schedules}.
\end{proof}

\begin{corollary}[Sub-Weibull and moment caps for random-subspace menus]
\label{cor:active-menu-rectangular-tail-selected-random-subspace-schedules}
Fix \(C_0\ge1\).  Work in the active-menu setting of
Corollary~\ref{cor:active-menu-rectangular-tail-capped-random-subspace-schedules}.
For each active triple \((m,h,i)\), choose a level
\(\ell_{m,h,i}>0\) and a scale \(s_{m,h,i}>0\).  Assume that the deterministic
cap \(\rho_{m,h,i}\le1\) is chosen by one of the following two rules.

\begin{enumerate}[label=(\alph*)]
\item Sub-Weibull operator tails: for some fixed \(\alpha>0\) and
  \(A_{\rm tail}\ge1\),
  \[
    \mathbb P\{\|X_{m,h,i}\|_{\op}>s_{m,h,i}u\}
    \le A_{\rm tail}\exp\{-u^\alpha\}\qquad(u\ge0),
  \]
  and
  \[
    \rho_{m,h,i}:=s_{m,h,i}\ell_{m,h,i}^{1/\alpha}\le1.
  \]
\item Moment operator tails: for some fixed \(p>0\),
  \[
    \mathbb E\|X_{m,h,i}\|_{\op}^p\le s_{m,h,i}^p,
  \]
  and
  \[
    \rho_{m,h,i}:=s_{m,h,i}\exp\{\ell_{m,h,i}/p\}\le1.
  \]
\end{enumerate}

In case (a) put \(A_*=A_{\rm tail}\), and in case (b) put \(A_*=1\).
Let \(b_{m,h}:=\max_i\rho_{m,h,i}^2\), and let
\(\mu_{m,h}^{\rm sel}\) be the capped row/column projection budget
\[
  \mu_{m,h}^{\rm sel}:=
  \max\left\{
  \left\|\sum_i\rho_{m,h,i}^2\mathbb EP_{m,h,i}^L\right\|_{\op},
  \left\|\sum_i\rho_{m,h,i}^2\mathbb EP_{m,h,i}^R\right\|_{\op}
  \right\}.
\]
If there are \(t_{m,h}>0\) such that
\[
  \mu_{m,h}^{\rm sel}+t_{m,h}
  \le
  \frac{C_0n_m}{\log(eD_m^{\rm rec})}
  \qquad(1\le h\le J_m),
\]
and
\[
  \sum_m\sum_{h=1}^{J_m}
  \left[
  A_*\sum_{i=1}^{n_m}\exp\{-\ell_{m,h,i}\}
  +
  D_m^{\rm rec}
  \exp\left\{
  -\frac{t_{m,h}^2}{2(\mu_{m,h}^{\rm sel}+b_{m,h}t_{m,h}/3)}
  \right\}
  \right]<\infty,
\]
with zero denominator interpreted as zero contribution, then, almost surely,
for all sufficiently large \(m\), every active random-subspace stream has a
Spencer-scale signing:
\[
  \min_{\eps\in\{\pm1\}^{n_m}}
  \left\|\sum_{i=1}^{n_m}\eps_iX_{m,h,i}\right\|_{\op}
  \le C(C_0)\sqrt{n_m}
  \qquad(1\le h\le J_m).
\]
Sampled certificates and countable-library exhaustions follow under the same
additional summability hypotheses as in
Corollary~\ref{cor:active-menu-rectangular-tail-capped-random-subspace-schedules}.

In the Haar row/column rank specialization, \(\mu_{m,h}^{\rm sel}\) is obtained
by replacing \(\rho_{m,h,i}^2\) with
\[
  s_{m,h,i}^2\ell_{m,h,i}^{2/\alpha}
  \quad\text{in case (a),}
  \qquad
  s_{m,h,i}^2\exp\{2\ell_{m,h,i}/p\}
  \quad\text{in case (b),}
\]
inside the two average-rank sums.  Thus sub-Weibull low-rank coefficients need
only lose the logarithmic cap factor \(\ell^{2/\alpha}\), while finite
\(p\)-moment coefficients are allowed whenever their normalization absorbs the
factor \(\exp\{2\ell/p\}\).

Consequently, if \(J_mn_m\le(eD_m^{\rm rec})^\Gamma\), the constants in the
tail bounds are uniform, and one takes the common level
\(\ell_{m,h,i}=R\log(eD_m^{\rm rec})\), then the cap-tail part of the series is
summable along the full integer dimension schedule \(D_m^{\rm rec}=m\) whenever
\(R>\Gamma+1\), and along every stretched-exponential sparse schedule whenever
\(R>\Gamma\).  The remaining condition is exactly the displayed capped
row/column projection Bernstein summability.
\end{corollary}

\begin{proof}
In the sub-Weibull case,
\[
  \mathbb P\{\|X_{m,h,i}\|_{\op}>\rho_{m,h,i}\}
  \le A_{\rm tail}\exp\{-\ell_{m,h,i}\}.
\]
In the moment case, Markov's inequality gives
\[
  \mathbb P\{\|X_{m,h,i}\|_{\op}>\rho_{m,h,i}\}
  \le
  \frac{s_{m,h,i}^p}{s_{m,h,i}^p\exp\{\ell_{m,h,i}\}}
  =
  \exp\{-\ell_{m,h,i}\}.
\]
Use these quantities as the tail budgets
\(\delta_{m,h,i}\) in
Corollary~\ref{cor:active-menu-rectangular-tail-capped-random-subspace-schedules}.
The Haar-rank formulas are the corresponding specialization of
\(\mathbb EP_{m,h,i}^L=(r_{m,h,i}^L/d_{m,1})I\) and
\(\mathbb EP_{m,h,i}^R=(r_{m,h,i}^R/d_{m,2})I\).  The final all-integer and
sparse cap-tail thresholds follow from summing
\((eD_m^{\rm rec})^{\Gamma-R}\).
\end{proof}

\begin{corollary}[Active menus of rectangular dyad/frame streams]
\label{cor:active-menu-rectangular-dyad-frame-schedules}
Fix \(C_0\ge1\).  For each \(m\ge1\), let \(J_m\ge1\),
\(d_{m,1},d_{m,2}\ge1\), \(D_m^{\rm rec}=d_{m,1}+d_{m,2}\), and \(n_m\ge1\).
For each active stream \(1\le h\le J_m\), let
\[
  X_{m,h,i}=a_{m,h,i}u_{m,h,i}v_{m,h,i}^*
  \in\M_{d_{m,1}\times d_{m,2}}(\C),
  \qquad 1\le i\le n_m,
\]
be as in Corollary~\ref{cor:rectangular-random-dyads-frame}: the unit-vector
pairs \((u_{m,h,i},v_{m,h,i})\) are independent over \(i\), and conditional on
them the amplitudes satisfy \(|a_{m,h,i}|\le\rho_{m,h,i}\le1\).  Put
\[
  b_{m,h}:=\max_i\rho_{m,h,i}^2
\]
and
\[
  \mu_{m,h}^{\rm dyad}:=
  \max\left\{
  \left\|\sum_i\rho_{m,h,i}^2\mathbb E u_{m,h,i}u_{m,h,i}^*\right\|_{\op},
  \left\|\sum_i\rho_{m,h,i}^2\mathbb E v_{m,h,i}v_{m,h,i}^*\right\|_{\op}
  \right\}.
\]
Suppose there are \(t_{m,h}>0\) such that
\[
  \mu_{m,h}^{\rm dyad}+t_{m,h}
  \le
  \frac{C_0n_m}{\log(eD_m^{\rm rec})}
  \qquad(1\le h\le J_m),
\]
and
\[
  \sum_m\sum_{h=1}^{J_m}
  D_m^{\rm rec}
  \exp\left\{
  -\frac{t_{m,h}^2}
         {2(\mu_{m,h}^{\rm dyad}+b_{m,h}t_{m,h}/3)}
  \right\}<\infty,
\]
with zero denominator interpreted as zero contribution.  Then, almost surely,
for all sufficiently large \(m\), every active rectangular dyad/frame stream has
a Spencer-scale signing:
\[
  \min_{\eps\in\{\pm1\}^{n_m}}
  \left\|\sum_{i=1}^{n_m}\eps_iX_{m,h,i}\right\|_{\op}
  \le C(C_0)\sqrt{n_m}
  \qquad(1\le h\le J_m).
\]
The same statement holds for active prefixes exhausting countable libraries of
randomized feature maps or sampled rectangular atom dictionaries.  Sampled
certificates follow under the same additional summability hypotheses as in
Corollary~\ref{cor:active-menu-rectangular-square-profile-schedules}.
\end{corollary}

\begin{proof}
Apply Corollary~\ref{cor:rectangular-random-dyads-frame} to each active stream
and sum the displayed failure probabilities.  The first Borel--Cantelli lemma
gives the eventual simultaneous signing.  The active-prefix and sampled-search
additions are the same as in
Corollary~\ref{cor:active-menu-rectangular-random-subspace-schedules}.
\end{proof}

\begin{corollary}[Active menus of tail-capped rectangular dyads]
\label{cor:active-menu-rectangular-tail-capped-dyad-frame-schedules}
Fix \(C_0\ge1\).  For each \(m\ge1\), let \(J_m\ge1\),
\(d_{m,1},d_{m,2}\ge1\), \(D_m^{\rm rec}=d_{m,1}+d_{m,2}\), and \(n_m\ge1\).
For each active stream \(1\le h\le J_m\), let
\[
  X_{m,h,i}=a_{m,h,i}u_{m,h,i}v_{m,h,i}^*,
  \qquad 1\le i\le n_m,
\]
where the unit-vector pairs \((u_{m,h,i},v_{m,h,i})\) are independent over
\(i\), and conditional on them the scalar amplitudes \(a_{m,h,i}\) are
arbitrary.  Let deterministic caps \(\rho_{m,h,i}\in[0,1]\) and tail budgets
\(\delta_{m,h,i}\ge0\) satisfy
\[
  \mathbb P\{|a_{m,h,i}|>\rho_{m,h,i}\}\le\delta_{m,h,i}.
\]
Put
\[
  b_{m,h}:=\max_i\rho_{m,h,i}^2
\]
and
\[
  \mu_{m,h}^{\rm cap}:=
  \max\left\{
  \left\|\sum_i\rho_{m,h,i}^2\mathbb E u_{m,h,i}u_{m,h,i}^*\right\|_{\op},
  \left\|\sum_i\rho_{m,h,i}^2\mathbb E v_{m,h,i}v_{m,h,i}^*\right\|_{\op}
  \right\}.
\]
Suppose there are \(t_{m,h}>0\) such that
\[
  \mu_{m,h}^{\rm cap}+t_{m,h}
  \le
  \frac{C_0n_m}{\log(eD_m^{\rm rec})}
  \qquad(1\le h\le J_m),
\]
and
\[
  \sum_m\sum_{h=1}^{J_m}
  \left[
  \sum_{i=1}^{n_m}\delta_{m,h,i}
  +
  D_m^{\rm rec}
  \exp\left\{
  -\frac{t_{m,h}^2}
         {2(\mu_{m,h}^{\rm cap}+b_{m,h}t_{m,h}/3)}
  \right\}
  \right]<\infty,
\]
with zero denominator interpreted as zero contribution.  Then, almost surely,
for all sufficiently large \(m\), every active tail-capped rectangular
dyad/frame stream has a Spencer-scale signing:
\[
  \min_{\eps\in\{\pm1\}^{n_m}}
  \left\|\sum_{i=1}^{n_m}\eps_iX_{m,h,i}\right\|_{\op}
  \le C(C_0)\sqrt{n_m}
  \qquad(1\le h\le J_m).
\]
Consequently, countable libraries of Gaussian, sub-Gaussian, sub-exponential,
or heavy-tailed randomized-feature and sampled-atom streams are covered whenever
their cap-tail costs and capped frame-Bernstein costs are summable.  Sampled
certificates follow under the same additional summability hypotheses as in
Corollary~\ref{cor:active-menu-rectangular-square-profile-schedules}.
\end{corollary}

\begin{proof}
Apply Corollary~\ref{cor:rectangular-tail-capped-random-dyads-frame} to each
active stream, bound the cap-exception event by
\(\sum_i\delta_{m,h,i}\), and sum over \(m,h\).  The first Borel--Cantelli
lemma gives eventual simultaneous signings.  The active-prefix and
sampled-search additions are the same as in
Corollary~\ref{cor:active-menu-rectangular-dyad-frame-schedules}.
\end{proof}

\begin{corollary}[Sub-Weibull and moment caps for dyad/frame menus]
\label{cor:active-menu-rectangular-tail-selected-dyad-frame-schedules}
Fix \(C_0\ge1\).  Work in the active-menu dyad setting of
Corollary~\ref{cor:active-menu-rectangular-tail-capped-dyad-frame-schedules}.
For each active triple \((m,h,i)\), choose a level
\(\ell_{m,h,i}>0\) and a scalar scale \(s_{m,h,i}>0\).  Assume that the
deterministic cap \(\rho_{m,h,i}\le1\) is chosen by one of the following two
rules.

\begin{enumerate}[label=(\alph*)]
\item Sub-Weibull scalar tails: for some fixed \(\alpha>0\) and
  \(A_{\rm tail}\ge1\),
  \[
    \mathbb P\{|a_{m,h,i}|>s_{m,h,i}u\}
    \le A_{\rm tail}\exp\{-u^\alpha\}\qquad(u\ge0),
  \]
  and
  \[
    \rho_{m,h,i}:=s_{m,h,i}\ell_{m,h,i}^{1/\alpha}\le1.
  \]
\item Moment scalar tails: for some fixed \(p>0\),
  \[
    \mathbb E|a_{m,h,i}|^p\le s_{m,h,i}^p,
  \]
  and
  \[
    \rho_{m,h,i}:=s_{m,h,i}\exp\{\ell_{m,h,i}/p\}\le1.
  \]
\end{enumerate}

In case (a) put \(A_*=A_{\rm tail}\), and in case (b) put \(A_*=1\).
Let \(b_{m,h}:=\max_i\rho_{m,h,i}^2\), and put
\[
  \mu_{m,h}^{\rm dsel}:=
  \max\left\{
  \left\|\sum_i\rho_{m,h,i}^2
     \mathbb E u_{m,h,i}u_{m,h,i}^*\right\|_{\op},
  \left\|\sum_i\rho_{m,h,i}^2
     \mathbb E v_{m,h,i}v_{m,h,i}^*\right\|_{\op}
  \right\}.
\]
If there are \(t_{m,h}>0\) such that
\[
  \mu_{m,h}^{\rm dsel}+t_{m,h}
  \le
  \frac{C_0n_m}{\log(eD_m^{\rm rec})}
  \qquad(1\le h\le J_m),
\]
and
\[
  \sum_m\sum_{h=1}^{J_m}
  \left[
  A_*\sum_{i=1}^{n_m}\exp\{-\ell_{m,h,i}\}
  +
  D_m^{\rm rec}
  \exp\left\{
  -\frac{t_{m,h}^2}{2(\mu_{m,h}^{\rm dsel}+b_{m,h}t_{m,h}/3)}
  \right\}
  \right]<\infty,
\]
with zero denominator interpreted as zero contribution, then, almost surely,
for all sufficiently large \(m\), every active dyad/frame stream has a
Spencer-scale signing:
\[
  \min_{\eps\in\{\pm1\}^{n_m}}
  \left\|\sum_{i=1}^{n_m}\eps_i
    a_{m,h,i}u_{m,h,i}v_{m,h,i}^*\right\|_{\op}
  \le C(C_0)\sqrt{n_m}
  \qquad(1\le h\le J_m).
\]
Sampled certificates and countable-library exhaustions follow under the same
additional summability hypotheses as in
Corollary~\ref{cor:active-menu-rectangular-tail-capped-dyad-frame-schedules}.

In the Haar-isotropic dyad case, \(\mu_{m,h}^{\rm dsel}\) is the maximum of
\[
  \frac1{d_{m,1}}\sum_i\rho_{m,h,i}^2,
  \qquad
  \frac1{d_{m,2}}\sum_i\rho_{m,h,i}^2.
\]
Thus sub-Weibull feature weights pay the logarithmic frame factor
\(\ell^{2/\alpha}\), while finite \(p\)-moment feature weights pay
\(\exp\{2\ell/p\}\), inside the two expected frame budgets.

Consequently, if \(J_mn_m\le(eD_m^{\rm rec})^\Gamma\), the constants in the
tail bounds are uniform, and one takes the common level
\(\ell_{m,h,i}=R\log(eD_m^{\rm rec})\), then the cap-tail part of the series is
summable along the full integer dimension schedule \(D_m^{\rm rec}=m\) whenever
\(R>\Gamma+1\), and along every stretched-exponential sparse schedule whenever
\(R>\Gamma\).  The remaining condition is the displayed capped frame-Bernstein
summability.
\end{corollary}

\begin{proof}
The sub-Weibull tail rule gives
\[
  \mathbb P\{|a_{m,h,i}|>\rho_{m,h,i}\}
  \le A_{\rm tail}\exp\{-\ell_{m,h,i}\}.
\]
The moment rule gives the same bound with \(A_{\rm tail}=1\) by Markov's
inequality:
\[
  \mathbb P\{|a_{m,h,i}|>\rho_{m,h,i}\}
  \le
  \frac{s_{m,h,i}^p}{s_{m,h,i}^p\exp\{\ell_{m,h,i}\}}
  =
  \exp\{-\ell_{m,h,i}\}.
\]
Use these as the \(\delta_{m,h,i}\) in
Corollary~\ref{cor:active-menu-rectangular-tail-capped-dyad-frame-schedules}.
The Haar-isotropic formula is the specialization
\(\mathbb Eu_{m,h,i}u_{m,h,i}^*=d_{m,1}^{-1}I\) and
\(\mathbb Ev_{m,h,i}v_{m,h,i}^*=d_{m,2}^{-1}I\).  The all-integer and sparse
cap-tail thresholds follow by summing \((eD_m^{\rm rec})^{\Gamma-R}\).
\end{proof}

\begin{corollary}[Homogeneous isotropic dyad cap schedules]
\label{cor:active-menu-homogeneous-isotropic-dyad-cap-schedules}
Fix \(C_0\ge1\), and put \(c_0:=3C_0/16\).  For each \(m\ge1\), let
\(J_m\ge1\), \(d_{m,1},d_{m,2}\ge1\), \(D_m^{\rm rec}=d_{m,1}+d_{m,2}\),
\(d_{m,*}:=\min\{d_{m,1},d_{m,2}\}\), and \(n_m\ge1\).  For each active
stream \(1\le h\le J_m\), let
\[
  X_{m,h,i}=a_{m,h,i}u_{m,h,i}v_{m,h,i}^*,
  \qquad 1\le i\le n_m,
\]
where the unit-vector pairs are independent over \(i\) and isotropic:
\[
  \mathbb E u_{m,h,i}u_{m,h,i}^*=\frac1{d_{m,1}}I_{d_{m,1}},
  \qquad
  \mathbb E v_{m,h,i}v_{m,h,i}^*=\frac1{d_{m,2}}I_{d_{m,2}}.
\]
Let \(L_m:=R\log(eD_m^{\rm rec})\).  Suppose that, for every active stream,
one of the following two homogeneous tail models holds with scale \(s_{m,h}>0\).

\begin{enumerate}[label=(\alph*)]
\item Sub-Weibull weights: for fixed \(\alpha>0\) and \(A_{\rm tail}\ge1\),
  \[
    \mathbb P\{|a_{m,h,i}|>s_{m,h}u\}
    \le A_{\rm tail}\exp\{-u^\alpha\}
    \qquad(u\ge0,\ 1\le i\le n_m),
  \]
  and
  \[
    q_{m,h}:=s_{m,h}L_m^{1/\alpha}.
  \]
\item Moment weights: for fixed \(p>0\),
  \[
    \mathbb E|a_{m,h,i}|^p\le s_{m,h}^p\qquad(1\le i\le n_m),
  \]
  and
  \[
    q_{m,h}:=s_{m,h}\exp\{L_m/p\}.
  \]
\end{enumerate}

In case (a) put \(A_*=A_{\rm tail}\), and in case (b) put \(A_*=1\).
Assume the homogeneous cap and frame half-budget bounds
\[
  q_{m,h}\le1,\qquad
  q_{m,h}^2\le \frac{C_0d_{m,*}}{2\log(eD_m^{\rm rec})}
  \qquad(1\le h\le J_m).
\]
If
\[
  \sum_m\sum_{h=1}^{J_m}
  \left[
  A_* n_m\exp\{-L_m\}
  +
  D_m^{\rm rec}\exp\left\{
  -\frac{c_0n_m}{\log(eD_m^{\rm rec})}
  \right\}
  \right]<\infty,
\]
then, almost surely, for all sufficiently large \(m\), every active
homogeneous isotropic dyad stream has a Spencer-scale signing:
\[
  \min_{\eps\in\{\pm1\}^{n_m}}
  \left\|\sum_{i=1}^{n_m}\eps_i
    a_{m,h,i}u_{m,h,i}v_{m,h,i}^*\right\|_{\op}
  \le C(C_0)\sqrt{n_m}
  \qquad(1\le h\le J_m).
\]

In particular, assume
\[
  D_m^{\rm rec}=m,\qquad
  J_mn_m\le(eD_m^{\rm rec})^\Gamma,\qquad
  n_m\ge a\log^2(eD_m^{\rm rec}).
\]
Then the displayed summability holds whenever
\[
  R>\Gamma+1,\qquad c_0a>\Gamma+2.
\]
Along any stretched-exponential sparse dimension schedule
\(D_m^{\rm rec}\ge \exp(cm^\kappa)\), it is enough that
\[
  R>\Gamma,\qquad c_0a>\Gamma+1.
\]
\end{corollary}

\begin{proof}
Apply
Corollary~\ref{cor:active-menu-rectangular-tail-selected-dyad-frame-schedules}
with the common levels \(\ell_{m,h,i}=L_m\) and common caps
\(\rho_{m,h,i}=q_{m,h}\).  The sub-Weibull or moment tail contribution is
bounded by \(A_*n_m\exp\{-L_m\}\) for stream \((m,h)\).  By isotropy,
\[
  \mu_{m,h}^{\rm dsel}
  =
  \max\left\{\frac{n_mq_{m,h}^2}{d_{m,1}},
             \frac{n_mq_{m,h}^2}{d_{m,2}}\right\}
  \le
  \frac{C_0n_m}{2\log(eD_m^{\rm rec})}.
\]
Taking \(t_{m,h}=C_0n_m/(2\log(eD_m^{\rm rec}))\) and using
\(b_{m,h}\le1\) gives
\[
  D_m^{\rm rec}
  \exp\left\{
  -\frac{t_{m,h}^2}{2(\mu_{m,h}^{\rm dsel}+b_{m,h}t_{m,h}/3)}
  \right\}
  \le
  D_m^{\rm rec}
  \exp\left\{-\frac{c_0n_m}{\log(eD_m^{\rm rec})}\right\}.
\]
Summability gives the almost-sure conclusion by Borel--Cantelli.  The
all-integer and sparse displays follow from
\[
  \sum_{h\le J_m}A_*n_m\exp\{-L_m\}
  \le A_*(eD_m^{\rm rec})^{\Gamma-R}
\]
and
\[
  \sum_{h\le J_m}D_m^{\rm rec}
  \exp\left\{-\frac{c_0n_m}{\log(eD_m^{\rm rec})}\right\}
  \le
  (eD_m^{\rm rec})^{\Gamma+1-c_0a}.
\]
\end{proof}

\begin{corollary}[Coordinate-sampled dyad cap schedules]
\label{cor:active-menu-coordinate-dyad-cap-schedules}
Fix \(C_0\ge1\), and put \(c_0:=3C_0/16\).  For each \(m\ge1\), let
\(J_m\ge1\), \(d_{m,1},d_{m,2}\ge1\), \(D_m^{\rm rec}=d_{m,1}+d_{m,2}\), and
\(n_m\ge1\).  For each active stream \(1\le h\le J_m\), let
\[
  X_{m,h,i}=a_{m,h,i}e_{I_{m,h,i}}e_{J_{m,h,i}}^*
  \in\M_{d_{m,1}\times d_{m,2}}(\C),
  \qquad 1\le i\le n_m,
\]
where the coordinate pairs \((I_{m,h,i},J_{m,h,i})\) are independent over
\(i\).  Assume that their row and column marginals have deterministic envelopes
\[
  \mathbb P\{I_{m,h,i}=r\}\le\lambda_{m,h}^L,\qquad
  \mathbb P\{J_{m,h,i}=s\}\le\lambda_{m,h}^R
\]
for every admissible \(i,r,s\), and put
\[
  \lambda_{m,h}:=\max\{\lambda_{m,h}^L,\lambda_{m,h}^R\}.
\]
Let \(L_m:=R\log(eD_m^{\rm rec})\).  Suppose that, for every active stream,
one of the following two homogeneous scalar-tail models holds with scale
\(s_{m,h}>0\).

\begin{enumerate}[label=(\alph*)]
\item Sub-Weibull weights: for fixed \(\alpha>0\) and \(A_{\rm tail}\ge1\),
  \[
    \mathbb P\{|a_{m,h,i}|>s_{m,h}u\}
    \le A_{\rm tail}\exp\{-u^\alpha\}
    \qquad(u\ge0,\ 1\le i\le n_m),
  \]
  and
  \[
    q_{m,h}:=s_{m,h}L_m^{1/\alpha}.
  \]
\item Moment weights: for fixed \(p>0\),
  \[
    \mathbb E|a_{m,h,i}|^p\le s_{m,h}^p\qquad(1\le i\le n_m),
  \]
  and
  \[
    q_{m,h}:=s_{m,h}\exp\{L_m/p\}.
  \]
\end{enumerate}

In case (a) put \(A_*=A_{\rm tail}\), and in case (b) put \(A_*=1\).
Assume the coordinate cap and marginal-load half-budget bounds
\[
  q_{m,h}\le1,\qquad
  q_{m,h}^2\lambda_{m,h}
  \le \frac{C_0}{2\log(eD_m^{\rm rec})}
  \qquad(1\le h\le J_m).
\]
If
\[
  \sum_m\sum_{h=1}^{J_m}
  \left[
  A_* n_m\exp\{-L_m\}
  +
  D_m^{\rm rec}\exp\left\{
  -\frac{c_0n_m}{\log(eD_m^{\rm rec})}
  \right\}
  \right]<\infty,
\]
then, almost surely, for all sufficiently large \(m\), every active
coordinate-sampled stream has a Spencer-scale signing:
\[
  \min_{\eps\in\{\pm1\}^{n_m}}
  \left\|\sum_{i=1}^{n_m}\eps_i
    a_{m,h,i}e_{I_{m,h,i}}e_{J_{m,h,i}}^*\right\|_{\op}
  \le C(C_0)\sqrt{n_m}
  \qquad(1\le h\le J_m).
\]

In particular, assume
\[
  D_m^{\rm rec}=m,\qquad
  J_mn_m\le(eD_m^{\rm rec})^\Gamma,\qquad
  n_m\ge a\log^2(eD_m^{\rm rec}).
\]
Then the displayed summability holds whenever
\[
  R>\Gamma+1,\qquad c_0a>\Gamma+2.
\]
Along any stretched-exponential sparse dimension schedule
\(D_m^{\rm rec}\ge \exp(cm^\kappa)\), it is enough that
\[
  R>\Gamma,\qquad c_0a>\Gamma+1.
\]
For uniform sampled entries, \(\lambda_{m,h}=1/d_{m,*}\), where
\(d_{m,*}=\min\{d_{m,1},d_{m,2}\}\), and the half-budget condition reduces to
\[
  q_{m,h}^2\le\frac{C_0d_{m,*}}{2\log(eD_m^{\rm rec})}.
\]
\end{corollary}

\begin{proof}
Apply
Corollary~\ref{cor:active-menu-rectangular-tail-selected-dyad-frame-schedules}
with
\[
  u_{m,h,i}=e_{I_{m,h,i}},\qquad
  v_{m,h,i}=e_{J_{m,h,i}},
\]
common levels \(\ell_{m,h,i}=L_m\), and common caps
\(\rho_{m,h,i}=q_{m,h}\).  The sub-Weibull or moment rule gives cap-tail cost
at most \(A_*n_m\exp\{-L_m\}\) for stream \((m,h)\).  The marginal envelopes
give
\[
  \left\|\mathbb E e_{I_{m,h,i}}e_{I_{m,h,i}}^*\right\|_{\op}
  \le \lambda_{m,h}^L,\qquad
  \left\|\mathbb E e_{J_{m,h,i}}e_{J_{m,h,i}}^*\right\|_{\op}
  \le \lambda_{m,h}^R,
\]
and hence
\[
  \mu_{m,h}^{\rm dsel}\le n_mq_{m,h}^2\lambda_{m,h}
  \le
  \frac{C_0n_m}{2\log(eD_m^{\rm rec})}.
\]
Taking \(t_{m,h}=C_0n_m/(2\log(eD_m^{\rm rec}))\) and using
\(q_{m,h}\le1\) gives the displayed Bernstein term with exponent
\(c_0n_m/\log(eD_m^{\rm rec})\), exactly as in
Corollary~\ref{cor:active-menu-homogeneous-isotropic-dyad-cap-schedules}.
Borel--Cantelli gives the almost-sure conclusion, and the all-integer and
sparse schedule thresholds are the same two summability computations from that
corollary.  In the uniform sampled-entry case the row and column marginal
loads are \(d_{m,1}^{-1}\) and \(d_{m,2}^{-1}\).
\end{proof}

\begin{corollary}[Predictable coordinate-dyad cap schedules]
\label{cor:active-menu-predictable-coordinate-dyad-cap-schedules}
Fix \(C_0\ge1\), and put \(c_0:=3C_0/16\).  For each \(m\ge1\), let
\(J_m\ge1\), \(d_{m,1},d_{m,2}\ge1\), \(D_m^{\rm rec}=d_{m,1}+d_{m,2}\), and
\(n_m\ge1\).  For each active stream \(1\le h\le J_m\), let
\((\mcF_{m,h,i})_{i=0}^{n_m}\) be a filtration and let
\[
  X_{m,h,i}=a_{m,h,i}e_{I_{m,h,i}}e_{J_{m,h,i}}^*
  \in\M_{d_{m,1}\times d_{m,2}}(\C),
  \qquad 1\le i\le n_m,
\]
be \(\mcF_{m,h,i}\)-measurable.  Assume that the one-step coordinate
marginals satisfy
\[
  \mathbb P\{I_{m,h,i}=r\mid\mcF_{m,h,i-1}\}\le\lambda_{m,h}^L,
  \qquad
  \mathbb P\{J_{m,h,i}=s\mid\mcF_{m,h,i-1}\}\le\lambda_{m,h}^R
\]
almost surely for every admissible \(i,r,s\), and put
\[
  \lambda_{m,h}:=\max\{\lambda_{m,h}^L,\lambda_{m,h}^R\}.
\]
Let \(L_m:=R\log(eD_m^{\rm rec})\).  Suppose that, for every active stream,
one of the following two scalar-tail models holds with scale \(s_{m,h}>0\).

\begin{enumerate}[label=(\alph*)]
\item Sub-Weibull weights: for fixed \(\alpha>0\) and \(A_{\rm tail}\ge1\),
  \[
    \mathbb P\{|a_{m,h,i}|>s_{m,h}u\}
    \le A_{\rm tail}\exp\{-u^\alpha\}
    \qquad(u\ge0,\ 1\le i\le n_m),
  \]
  and
  \[
    q_{m,h}:=s_{m,h}L_m^{1/\alpha}.
  \]
\item Moment weights: for fixed \(p>0\),
  \[
    \mathbb E|a_{m,h,i}|^p\le s_{m,h}^p\qquad(1\le i\le n_m),
  \]
  and
  \[
    q_{m,h}:=s_{m,h}\exp\{L_m/p\}.
  \]
\end{enumerate}

In case (a) put \(A_*=A_{\rm tail}\), and in case (b) put \(A_*=1\).
Assume
\[
  q_{m,h}\le1,\qquad
  q_{m,h}^2\lambda_{m,h}
  \le \frac{C_0}{2\log(eD_m^{\rm rec})}
  \qquad(1\le h\le J_m).
\]
If
\[
  \sum_m\sum_{h=1}^{J_m}
  \left[
  A_* n_m\exp\{-L_m\}
  +
  D_m^{\rm rec}\exp\left\{
  -\frac{c_0n_m}{\log(eD_m^{\rm rec})}
  \right\}
  \right]<\infty,
\]
then, almost surely, for all sufficiently large \(m\), every active
predictable coordinate-dyad stream has a Spencer-scale signing:
\[
  \min_{\eps\in\{\pm1\}^{n_m}}
  \left\|\sum_{i=1}^{n_m}\eps_i
    a_{m,h,i}e_{I_{m,h,i}}e_{J_{m,h,i}}^*\right\|_{\op}
  \le C(C_0)\sqrt{n_m}
  \qquad(1\le h\le J_m).
\]

In particular, if
\[
  D_m^{\rm rec}=m,\qquad
  J_mn_m\le(eD_m^{\rm rec})^\Gamma,\qquad
  n_m\ge a\log^2(eD_m^{\rm rec}),
\]
then the displayed summability holds whenever
\[
  R>\Gamma+1,\qquad c_0a>\Gamma+2.
\]
Along any stretched-exponential sparse dimension schedule
\(D_m^{\rm rec}\ge \exp(cm^\kappa)\), it is enough that
\[
  R>\Gamma,\qquad c_0a>\Gamma+1.
\]
Thus adaptive sampled-entry streams, without-replacement coordinate exposure,
and permutation or matching windows obey the same schedule whenever their
one-step row and column marginal envelopes satisfy the displayed half-budget.
\end{corollary}

\begin{proof}
For stream \((m,h)\), define the capped adapted contraction
\[
  Y_{m,h,i}:=
  a_{m,h,i}\mathbf 1_{\{|a_{m,h,i}|\le q_{m,h}\}}
  e_{I_{m,h,i}}e_{J_{m,h,i}}^* .
\]
By the sub-Weibull tail bound, or by Markov's inequality in the moment case,
\[
  \mathbb P\{\exists i:\ |a_{m,h,i}|>q_{m,h}\}
  \le A_*n_m\exp\{-L_m\}.
\]
On the complement of this cap-exception event, \(X_{m,h,i}=Y_{m,h,i}\) for all
\(i\).

For the capped stream,
\[
  Y_{m,h,i}Y_{m,h,i}^*
  \preceq q_{m,h}^2e_{I_{m,h,i}}e_{I_{m,h,i}}^*,\qquad
  Y_{m,h,i}^*Y_{m,h,i}
  \preceq q_{m,h}^2e_{J_{m,h,i}}e_{J_{m,h,i}}^* .
\]
Taking conditional expectations and using the marginal envelopes gives
\[
  \left\|\sum_i
  \mathbb E(Y_{m,h,i}Y_{m,h,i}^*\mid\mcF_{m,h,i-1})\right\|_{\op}
  \le n_mq_{m,h}^2\lambda_{m,h}^L,
\]
and similarly on the column side.  Hence the predictable rectangular profile
of the capped stream is at most
\[
  n_mq_{m,h}^2\lambda_{m,h}
  \le \frac{C_0n_m}{2\log(eD_m^{\rm rec})}.
\]
Apply Corollary~\ref{cor:predictable-rectangular-square-profile} to the capped
stream with \(v=t=C_0n_m/(2\log(eD_m^{\rm rec}))\).  The structural failure is
at most
\[
  D_m^{\rm rec}\exp\left\{
  -\frac{c_0n_m}{\log(eD_m^{\rm rec})}
  \right\}.
\]
Adding the cap-exception probability and summing over \(m,h\), the first
Borel--Cantelli lemma gives the eventual simultaneous conclusion for the
original streams.  The all-integer and sparse thresholds are the same two
summability computations as in
Corollary~\ref{cor:active-menu-coordinate-dyad-cap-schedules}.
\end{proof}

\begin{corollary}[Profile-capped predictable coordinate dyads]
\label{cor:active-menu-profile-capped-predictable-coordinate-dyads}
Fix \(C_0\ge1\).  For each \(m\ge1\), let \(J_m\ge1\),
\(d_{m,1},d_{m,2}\ge1\), \(D_m^{\rm rec}=d_{m,1}+d_{m,2}\), and \(n_m\ge1\).
For each active stream \(1\le h\le J_m\), let
\((\mcF_{m,h,i})_{i=0}^{n_m}\) be a filtration and let
\[
  X_{m,h,i}=a_{m,h,i}e_{I_{m,h,i}}e_{J_{m,h,i}}^*
  \in\M_{d_{m,1}\times d_{m,2}}(\C),
  \qquad 1\le i\le n_m,
\]
be \(\mcF_{m,h,i}\)-measurable.  Let deterministic caps
\(\rho_{m,h,i}\in[0,1]\) and tail budgets \(\delta_{m,h,i}\ge0\) satisfy
\[
  \mathbb P\{|a_{m,h,i}|>\rho_{m,h,i}\}\le\delta_{m,h,i}.
\]
Assume that deterministic one-step coordinate envelopes
\(\pi^L_{m,h,i,r}\) and \(\pi^R_{m,h,i,s}\) satisfy, almost surely,
\[
  \mathbb P\{I_{m,h,i}=r\mid\mcF_{m,h,i-1}\}
  \le \pi^L_{m,h,i,r},
  \qquad
  \mathbb P\{J_{m,h,i}=s\mid\mcF_{m,h,i-1}\}
  \le \pi^R_{m,h,i,s}.
\]
Define the capped row/column coordinate profile
\[
  \Lambda_{m,h}^{\rm coord}:=
  \max\left\{
  \max_{1\le r\le d_{m,1}}\sum_{i=1}^{n_m}
       \rho_{m,h,i}^2\pi^L_{m,h,i,r},
  \max_{1\le s\le d_{m,2}}\sum_{i=1}^{n_m}
       \rho_{m,h,i}^2\pi^R_{m,h,i,s}
  \right\}.
\]
Suppose there are \(t_{m,h}>0\) such that
\[
  \Lambda_{m,h}^{\rm coord}+t_{m,h}
  \le
  \frac{C_0n_m}{\log(eD_m^{\rm rec})}
  \qquad(1\le h\le J_m),
\]
and
\[
  \sum_m\sum_{h=1}^{J_m}
  \left[
  \sum_{i=1}^{n_m}\delta_{m,h,i}
  +
  D_m^{\rm rec}
  \exp\left\{
  -\frac{t_{m,h}^2}{2(\Lambda_{m,h}^{\rm coord}+t_{m,h}/3)}
  \right\}
  \right]<\infty.
\]
Then, almost surely, for all sufficiently large \(m\), every active
profile-capped predictable coordinate stream has a Spencer-scale signing:
\[
  \min_{\eps\in\{\pm1\}^{n_m}}
  \left\|\sum_{i=1}^{n_m}\eps_i
    a_{m,h,i}e_{I_{m,h,i}}e_{J_{m,h,i}}^*\right\|_{\op}
  \le C(C_0)\sqrt{n_m}
  \qquad(1\le h\le J_m).
\]

Consequently, sub-Weibull or finite-moment amplitudes may be inserted by
choosing caps exactly as in
Corollary~\ref{cor:active-menu-rectangular-tail-selected-dyad-frame-schedules}:
sub-Weibull tails use
\(\rho_{m,h,i}=s_{m,h,i}\ell_{m,h,i}^{1/\alpha}\le1\) and
\(\delta_{m,h,i}=A_{\rm tail}e^{-\ell_{m,h,i}}\), while finite \(p\)-moments
use
\(\rho_{m,h,i}=s_{m,h,i}\exp\{\ell_{m,h,i}/p\}\le1\) and
\(\delta_{m,h,i}=e^{-\ell_{m,h,i}}\).  Thus nonuniform sampled-entry,
adaptive sparse-observation, and weighted without-replacement coordinate
models are charged only by their maximum capped predictable row or column load.
\end{corollary}

\begin{proof}
For each active stream set
\[
  Y_{m,h,i}:=
  a_{m,h,i}\mathbf 1_{\{|a_{m,h,i}|\le\rho_{m,h,i}\}}
  e_{I_{m,h,i}}e_{J_{m,h,i}}^* .
\]
The cap-exception event has probability at most
\(\sum_i\delta_{m,h,i}\), and outside it \(X_{m,h,i}=Y_{m,h,i}\) for all
\(i\).  Also \(\|Y_{m,h,i}\|_{\op}\le\rho_{m,h,i}\le1\), and
\[
  Y_{m,h,i}Y_{m,h,i}^*
  \preceq
  \rho_{m,h,i}^2e_{I_{m,h,i}}e_{I_{m,h,i}}^*,
  \qquad
  Y_{m,h,i}^*Y_{m,h,i}
  \preceq
  \rho_{m,h,i}^2e_{J_{m,h,i}}e_{J_{m,h,i}}^* .
\]
Taking conditional expectations gives diagonal row and column profiles whose
\(r\)-th and \(s\)-th diagonal entries are bounded by
\[
  \sum_i\rho_{m,h,i}^2\pi^L_{m,h,i,r},
  \qquad
  \sum_i\rho_{m,h,i}^2\pi^R_{m,h,i,s}.
\]
Hence the predictable rectangular square profile of the capped stream is at
most \(\Lambda_{m,h}^{\rm coord}\).  Apply
Corollary~\ref{cor:predictable-rectangular-square-profile} with
\(v=\Lambda_{m,h}^{\rm coord}\) and \(t=t_{m,h}\), add the cap-exception
probability, and sum over \(m,h\).  Borel--Cantelli gives the almost-sure
simultaneous conclusion.  The displayed sub-Weibull and moment caps give the
claimed tail budgets by the same calculation as in
Corollary~\ref{cor:active-menu-rectangular-tail-selected-dyad-frame-schedules}.
\end{proof}

\begin{corollary}[Finite-menu profile-capped coordinate schedules]
\label{cor:finite-menu-profile-capped-predictable-coordinate-schedules}
Fix \(C_0\ge1\), and put \(\gamma_0:=3C_0/16\).  For each \(m\ge1\), let
\(J_m\ge1\), \(d_{m,1},d_{m,2}\ge1\),
\(D_m^{\rm rec}=d_{m,1}+d_{m,2}\), and
\[
  \ell_m:=\log(eD_m^{\rm rec}).
\]
For every active \(1\le h\le J_m\), let
\[
  X_{m,h,i}=a_{m,h,i}e_{I_{m,h,i}}e_{J_{m,h,i}}^*
  \in\M_{d_{m,1}\times d_{m,2}}(\C),
  \qquad 1\le i\le n_m,
\]
be an adapted coordinate-dyad stream satisfying the hypotheses of
Corollary~\ref{cor:active-menu-profile-capped-predictable-coordinate-dyads},
with deterministic caps \(\rho_{m,h,i}\), tail budgets
\(\delta_{m,h,i}\), one-step envelopes \(\pi^L_{m,h,i,r}\),
\(\pi^R_{m,h,i,s}\), and capped coordinate profile
\(\Lambda_{m,h}^{\rm coord}\).  Assume the half-budget profile bound
\[
  \Lambda_{m,h}^{\rm coord}
  \le \frac{C_0n_m}{2\ell_m}
  \qquad(1\le h\le J_m)
\]
and the two summability conditions
\[
  \sum_m\sum_{h=1}^{J_m}\sum_{i=1}^{n_m}
  \delta_{m,h,i}<\infty,
  \qquad
  \sum_m J_mD_m^{\rm rec}
  \exp\left\{-\gamma_0\frac{n_m}{\ell_m}\right\}<\infty .
\]
Then, almost surely, for all sufficiently large \(m\), every active
profile-capped predictable coordinate stream has a Spencer-scale signing:
\[
  \min_{\eps\in\{\pm1\}^{n_m}}
  \left\|\sum_{i=1}^{n_m}\eps_i
    a_{m,h,i}e_{I_{m,h,i}}e_{J_{m,h,i}}^*\right\|_{\op}
  \le C(C_0)\sqrt{n_m}
  \qquad(1\le h\le J_m).
\]

If each active stream is supplied with an independent batch of \(M_{m,h}\)
uniform random sign vectors, independent of the stream, then the same eventual
statement holds for sampled sign certificates provided
\[
  \sum_m\sum_{h=1}^{J_m}2^{-M_{m,h}}<\infty .
\]
If a common batch of \(M_m\) sign vectors is reused across the active menu and
is independent of all active streams, it is enough that
\[
  \sum_mJ_m2^{-M_m}<\infty .
\]

The structural summability follows from each of the following concrete
schedules.
\begin{enumerate}[label=\textup{(\alph*)}]
\item General schedules.  If, for some \(\delta>0\),
  \[
    n_m\ge
    \gamma_0^{-1}\ell_m
    \bigl(\log(eJ_mD_m^{\rm rec})+(1+\delta)\log(e+m)\bigr)
  \]
  eventually, then the structural term is summable.  A common sampled batch is
  summable if, for some \(\delta'>0\),
  \[
    M_m\ge \log_2(eJ_m)+(1+\delta')\log_2(e+m)
  \]
  eventually.
\item Full integer dimensions.  If the rows are indexed by
  \(D_m^{\rm rec}=D=2,3,\ldots\),
  \(J_DD\le(eD)^\Gamma\) eventually, and
  \[
    n_D\ge a\log^2(eD)
  \]
  eventually, then the structural term is summable whenever
  \[
    \gamma_0a>\Gamma+1 .
  \]
\item Sparse stretched-exponential dimensions.  If
  \(D_m^{\rm rec}\ge\exp(cm^\kappa)\) eventually for some \(c,\kappa>0\),
  \(J_mD_m^{\rm rec}\le(eD_m^{\rm rec})^\Gamma\) eventually, and
  \[
    n_m\ge a\log^2(eD_m^{\rm rec})
  \]
  eventually, then the structural term is summable whenever
  \[
    \gamma_0a>\Gamma .
  \]
\end{enumerate}

In particular, under the common sample-menu size control
\(J_mn_m\le(eD_m^{\rm rec})^\Gamma\), the structural term is summable on all
integer dimensions once \(\gamma_0a>\Gamma+2\), and on sparse
stretched-exponential dimension schedules once \(\gamma_0a>\Gamma+1\).  If,
in addition, the caps are selected with
\(\delta_{m,h,i}\le A_*e^{-R\ell_m}\), as in the sub-Weibull and moment
cap-selection rules above, then the cap-tail term is summable under the same
sample-menu control whenever \(R>\Gamma+1\) on all integer dimensions and
whenever \(R>\Gamma\) on sparse stretched-exponential dimension schedules.
If the streams are defined for all \(h\ge1\) and \(J_m\to\infty\), then every
fixed stream in the countable coordinate library is included for all
sufficiently large \(m\).  Thus finite active menus, deterministic exhaustions
of countable coordinate libraries, and sampled-sign certificate searches
inherit explicit profile-capped schedules.
\end{corollary}

\begin{proof}
For fixed \(m\) and \(h\le J_m\), apply
Corollary~\ref{cor:active-menu-profile-capped-predictable-coordinate-dyads} with
\[
  t_{m,h}=\frac{C_0n_m}{2\ell_m}.
\]
The half-budget assumption gives
\(\Lambda_{m,h}^{\rm coord}+t_{m,h}\le C_0n_m/\ell_m\).  Moreover
\[
  \frac{t_{m,h}^2}
       {2(\Lambda_{m,h}^{\rm coord}+t_{m,h}/3)}
  \ge
  \frac{3C_0n_m}{16\ell_m}
  =\gamma_0\frac{n_m}{\ell_m}.
\]
Hence the failure probability for the active stream is at most
\[
  \sum_{i=1}^{n_m}\delta_{m,h,i}
  +
  D_m^{\rm rec}
  \exp\left\{-\gamma_0\frac{n_m}{\ell_m}\right\}.
\]
Union over \(h\le J_m\) and sum over \(m\); the two displayed summability
hypotheses and Borel--Cantelli give the simultaneous eventual signing
conclusion.  The sampled-sign assertions use the sampled-search part of
Corollary~\ref{cor:predictable-rectangular-square-profile} in the proof of
Corollary~\ref{cor:active-menu-profile-capped-predictable-coordinate-dyads},
then add the displayed \(2^{-M_{m,h}}\) or \(J_m2^{-M_m}\) union cost.

For the general schedule,
\[
  J_mD_m^{\rm rec}
  \exp\left\{-\gamma_0\frac{n_m}{\ell_m}\right\}
  \le e^{-1}(e+m)^{-1-\delta}.
\]
The common sampled-batch estimate is identical.  In full integer dimensions,
\[
  J_DD\exp\left\{-\gamma_0\frac{n_D}{\log(eD)}\right\}
  \le (eD)^{\Gamma-\gamma_0a},
\]
which is summable when \(\gamma_0a>\Gamma+1\).  In the sparse schedule the
same bound is a stretched-exponential summable sequence when
\(\gamma_0a>\Gamma\).  If \(J_mn_m\le(eD_m^{\rm rec})^\Gamma\), then
\(J_mD_m^{\rm rec}\le(eD_m^{\rm rec})^{\Gamma+1}\), giving the stated
\(\Gamma+2\) and \(\Gamma+1\) structural margins.  Finally,
\[
  \sum_{h\le J_m}\sum_{i\le n_m}\delta_{m,h,i}
  \le A_*J_mn_m e^{-R\ell_m}
  \le A_*(eD_m^{\rm rec})^{\Gamma-R},
\]
so the cap-tail summability is the usual integer \(p\)-series or sparse
stretched-exponential series test.
\end{proof}

\begin{corollary}[Predictable inhomogeneous entry-intensity schedules]
\label{cor:predictable-inhomogeneous-entry-intensity-schedules}
Fix \(C_0\ge1\), and put \(\gamma_0:=3C_0/16\).  For each \(m\ge1\), let
\(J_m\ge1\), \(d_{m,1},d_{m,2}\ge1\),
\(D_m^{\rm rec}=d_{m,1}+d_{m,2}\), and
\[
  \ell_m:=\log(eD_m^{\rm rec}).
\]
For each active stream \(1\le h\le J_m\), let
\((\mcF_{m,h,i})_{i=0}^{n_m}\) be a filtration and let
\[
  X_{m,h,i}=a_{m,h,i}e_{I_{m,h,i}}e_{J_{m,h,i}}^*
  \in\M_{d_{m,1}\times d_{m,2}}(\C),
  \qquad 1\le i\le n_m,
\]
be \(\mcF_{m,h,i}\)-measurable.  Let deterministic caps
\(\rho_{m,h,i}\in[0,1]\), tail budgets \(\delta_{m,h,i}\ge0\), and
entry-intensity envelopes \(p_{m,h,i,r,s}\ge0\) satisfy
\[
  \begin{aligned}
  \mathbb P\{|a_{m,h,i}|>\rho_{m,h,i}\}
  &\le\delta_{m,h,i},\\
  \mathbb P\{I_{m,h,i}=r,\ J_{m,h,i}=s\mid\mcF_{m,h,i-1}\}
  &\le p_{m,h,i,r,s}
  \end{aligned}
\]
almost surely for every admissible \(i,r,s\).  Define the capped predictable
row/column entry-load profile
\[
  \Theta_{m,h}^{\rm ent}:=
  \max\left\{
  \begin{aligned}
  &\max_{1\le r\le d_{m,1}}\sum_{i=1}^{n_m}\rho_{m,h,i}^2
     \sum_{s=1}^{d_{m,2}}p_{m,h,i,r,s},\\
  &\max_{1\le s\le d_{m,2}}\sum_{i=1}^{n_m}\rho_{m,h,i}^2
     \sum_{r=1}^{d_{m,1}}p_{m,h,i,r,s}
  \end{aligned}
  \right\}.
\]
Assume the half-budget condition
\[
  \Theta_{m,h}^{\rm ent}\le \frac{C_0n_m}{2\ell_m}
  \qquad(1\le h\le J_m)
\]
and the two summability conditions
\[
  \sum_m\sum_{h=1}^{J_m}\sum_{i=1}^{n_m}
  \delta_{m,h,i}<\infty,
  \qquad
  \sum_mJ_mD_m^{\rm rec}
  \exp\left\{-\gamma_0\frac{n_m}{\ell_m}\right\}<\infty .
\]
Then, almost surely, for all sufficiently large \(m\), every active
inhomogeneous entry-intensity stream has a Spencer-scale signing:
\[
  \min_{\eps\in\{\pm1\}^{n_m}}
  \left\|\sum_{i=1}^{n_m}\eps_i
    a_{m,h,i}e_{I_{m,h,i}}e_{J_{m,h,i}}^*\right\|_{\op}
  \le C(C_0)\sqrt{n_m}
  \qquad(1\le h\le J_m).
\]
Independent or common sampled-sign batches are certified under the same
\(2^{-M}\)-summability hypotheses as in
Corollary~\ref{cor:finite-menu-profile-capped-predictable-coordinate-schedules}.
The general, full-integer, sparse stretched-exponential, countable-exhaustion,
and cap-selection schedules of that corollary apply verbatim with
\(\Theta_{m,h}^{\rm ent}\) in place of
\(\Lambda_{m,h}^{\rm coord}\).

In particular, adaptive inhomogeneous sampled-entry matrices, bipartite edge
streams, Poissonized sparse observations, weighted without-replacement edge
exposure, and permutation or matching windows are covered as soon as their
predictable weighted expected row and column degrees obey the displayed
half-budget and their cap tails are summable.
\end{corollary}

\begin{proof}
Set
\[
  \pi^L_{m,h,i,r}:=\sum_{s=1}^{d_{m,2}}p_{m,h,i,r,s},
  \qquad
  \pi^R_{m,h,i,s}:=\sum_{r=1}^{d_{m,1}}p_{m,h,i,r,s}.
\]
The pairwise envelope implies the one-step marginal bounds required in
Corollary~\ref{cor:active-menu-profile-capped-predictable-coordinate-dyads},
and the resulting capped coordinate profile is at most
\(\Theta_{m,h}^{\rm ent}\).  Therefore the half-budget and summability
hypotheses are exactly those of
Corollary~\ref{cor:finite-menu-profile-capped-predictable-coordinate-schedules}.
Applying that corollary gives the signing, sampled-certificate, and schedule
claims.
\end{proof}

\begin{corollary}[Closed-form entry-intensity cap schedules]
\label{cor:closed-form-entry-intensity-cap-schedules}
Fix \(C_0\ge1\), and put \(c_0:=3C_0/16\).  For each \(m\ge1\), let
\(J_m\ge1\), \(d_{m,1},d_{m,2}\ge1\),
\(D_m^{\rm rec}=d_{m,1}+d_{m,2}\), and \(n_m\ge1\).  For each active stream
\(1\le h\le J_m\), let
\[
  X_{m,h,i}=a_{m,h,i}e_{I_{m,h,i}}e_{J_{m,h,i}}^*
  \in\M_{d_{m,1}\times d_{m,2}}(\C),
  \qquad 1\le i\le n_m,
\]
be adapted to a filtration \((\mcF_{m,h,i})_{i=0}^{n_m}\).  Assume that
deterministic entry-intensity envelopes \(p_{m,h,i,r,s}\ge0\) satisfy
\[
  \mathbb P\{I_{m,h,i}=r,\ J_{m,h,i}=s\mid\mcF_{m,h,i-1}\}
  \le p_{m,h,i,r,s}
\]
almost surely for every admissible \(i,r,s\).  Put
\[
  \Delta_{m,h}^{\rm ent}:=
  \max\left\{
  \begin{aligned}
  &\max_{1\le r\le d_{m,1}}
     \sum_{i=1}^{n_m}\sum_{s=1}^{d_{m,2}}p_{m,h,i,r,s},\\
  &\max_{1\le s\le d_{m,2}}
     \sum_{i=1}^{n_m}\sum_{r=1}^{d_{m,1}}p_{m,h,i,r,s}
  \end{aligned}
  \right\}.
\]
Let \(L_m:=R\log(eD_m^{\rm rec})\).  Suppose that, for every active stream,
one of the following two scalar-tail models holds with scale \(s_{m,h}>0\).

\begin{enumerate}[label=(\alph*)]
\item Sub-Weibull weights: for fixed \(\alpha>0\) and \(A_{\rm tail}\ge1\),
  \[
    \mathbb P\{|a_{m,h,i}|>s_{m,h}u\}
    \le A_{\rm tail}\exp\{-u^\alpha\}
    \qquad(u\ge0,\ 1\le i\le n_m),
  \]
  and
  \[
    q_{m,h}:=s_{m,h}L_m^{1/\alpha}.
  \]
\item Moment weights: for fixed \(p>0\),
  \[
    \mathbb E|a_{m,h,i}|^p\le s_{m,h}^p
    \qquad(1\le i\le n_m),
  \]
  and
  \[
    q_{m,h}:=s_{m,h}\exp\{L_m/p\}.
  \]
\end{enumerate}

In case (a) put \(A_*:=A_{\rm tail}\), and in case (b) put \(A_*:=1\).
Assume
\[
  q_{m,h}\le1,\qquad
  q_{m,h}^2\Delta_{m,h}^{\rm ent}
  \le\frac{C_0n_m}{2\log(eD_m^{\rm rec})}
  \qquad(1\le h\le J_m).
\]
If
\[
  \sum_m\sum_{h=1}^{J_m}
  \left[
  A_*n_m\exp\{-L_m\}
  +
  D_m^{\rm rec}
  \exp\left\{-\frac{c_0n_m}{\log(eD_m^{\rm rec})}\right\}
  \right]<\infty,
\]
then, almost surely, for all sufficiently large \(m\), every active
entry-intensity stream has a Spencer-scale signing:
\[
  \min_{\eps\in\{\pm1\}^{n_m}}
  \left\|\sum_{i=1}^{n_m}\eps_i
    a_{m,h,i}e_{I_{m,h,i}}e_{J_{m,h,i}}^*\right\|_{\op}
  \le C(C_0)\sqrt{n_m}
  \qquad(1\le h\le J_m).
\]
Sampled certificates follow under the same additional \(2^{-M}\)-summability
hypotheses as in
Corollary~\ref{cor:finite-menu-profile-capped-predictable-coordinate-schedules}.

In particular, if
\[
  D_m^{\rm rec}=m,\qquad
  J_mn_m\le(eD_m^{\rm rec})^\Gamma,\qquad
  n_m\ge a\log^2(eD_m^{\rm rec}),
\]
then the displayed summability holds whenever
\[
  R>\Gamma+1,\qquad c_0a>\Gamma+2.
\]
Along any stretched-exponential sparse dimension schedule
\(D_m^{\rm rec}\ge\exp(cm^\kappa)\), it is enough that
\[
  R>\Gamma,\qquad c_0a>\Gamma+1.
\]
Thus inhomogeneous random bipartite-edge, sparse observation, and adaptive
sampled-entry models have a closed-form criterion: after scalar capping, the
maximum predictable expected row or column degree must be at most
\(C_0n_m/(2q_{m,h}^2\log(eD_m^{\rm rec}))\).
\end{corollary}

\begin{proof}
Use the deterministic cap \(\rho_{m,h,i}=q_{m,h}\).  The sub-Weibull tail
model gives
\[
  \mathbb P\{|a_{m,h,i}|>q_{m,h}\}
  \le A_{\rm tail}e^{-L_m},
\]
and the moment model gives the same bound with \(A_{\rm tail}=1\) by Markov's
inequality.  Hence the cap-exception contribution for stream \((m,h)\) is at
most \(A_*n_me^{-L_m}\).  With these caps, the entry-intensity profile in
Corollary~\ref{cor:predictable-inhomogeneous-entry-intensity-schedules} is
\[
  \Theta_{m,h}^{\rm ent}=q_{m,h}^2\Delta_{m,h}^{\rm ent}
  \le\frac{C_0n_m}{2\log(eD_m^{\rm rec})}.
\]
Applying that corollary with the half-budget choice gives the structural term
\[
  D_m^{\rm rec}
  \exp\left\{-\frac{c_0n_m}{\log(eD_m^{\rm rec})}\right\}.
\]
Summability and Borel--Cantelli give the eventual simultaneous signing
conclusion, and the sampled-certificate statement is inherited from the same
corollary.  Under \(J_mn_m\le(eD_m^{\rm rec})^\Gamma\), summing the cap tails
over \(h\le J_m\) gives at most
\(A_*(eD_m^{\rm rec})^{\Gamma-R}\).  The structural active-menu term is at
most \((eD_m^{\rm rec})^{\Gamma+1-c_0a}\).  These are summable over all integer
dimensions under \(R>\Gamma+1\) and \(c_0a>\Gamma+2\), and along
stretched-exponential sparse schedules under \(R>\Gamma\) and
\(c_0a>\Gamma+1\).
\end{proof}

\begin{corollary}[Position-dependent entry-intensity caps]
\label{cor:position-dependent-entry-intensity-caps}
Fix \(C_0\ge1\), and put \(c_0:=3C_0/16\).  For each \(m\ge1\), let
\(J_m\ge1\), \(d_{m,1},d_{m,2}\ge1\),
\(D_m^{\rm rec}=d_{m,1}+d_{m,2}\), and
\(\ell_m=\log(eD_m^{\rm rec})\).  For each active stream \(1\le h\le J_m\),
let
\[
  X_{m,h,i}=a_{m,h,i}e_{I_{m,h,i}}e_{J_{m,h,i}}^*
  \in\M_{d_{m,1}\times d_{m,2}}(\C),
  \qquad 1\le i\le n_m,
\]
be adapted to a filtration \((\mcF_{m,h,i})_{i=0}^{n_m}\).  Assume that
deterministic pairwise entry envelopes \(p_{m,h,i,r,s}\ge0\) satisfy
\[
  \mathbb P\{I_{m,h,i}=r,\ J_{m,h,i}=s\mid\mcF_{m,h,i-1}\}
  \le p_{m,h,i,r,s}
\]
almost surely.  Let deterministic site caps
\(q_{m,h,i,r,s}\in[0,1]\) and tail budgets
\(\delta_{m,h,i,r,s}\ge0\) satisfy
\[
  \mathbb P\{|a_{m,h,i}|>q_{m,h,i,r,s},\
  I_{m,h,i}=r,\ J_{m,h,i}=s\}\le \delta_{m,h,i,r,s}.
\]
Define the site-capped weighted entry profile
\[
  \Xi_{m,h}^{\rm ent}:=
  \max\left\{
  \begin{aligned}
  &\max_{1\le r\le d_{m,1}}
     \sum_{i=1}^{n_m}\sum_{s=1}^{d_{m,2}}
     q_{m,h,i,r,s}^2p_{m,h,i,r,s},\\
  &\max_{1\le s\le d_{m,2}}
     \sum_{i=1}^{n_m}\sum_{r=1}^{d_{m,1}}
     q_{m,h,i,r,s}^2p_{m,h,i,r,s}
  \end{aligned}
  \right\}.
\]
Assume the half-budget condition
\[
  \Xi_{m,h}^{\rm ent}\le \frac{C_0n_m}{2\ell_m}
  \qquad(1\le h\le J_m)
\]
and the summability condition
\[
  \sum_m\sum_{h=1}^{J_m}
  \left[
  \sum_{i=1}^{n_m}\sum_{r=1}^{d_{m,1}}\sum_{s=1}^{d_{m,2}}
  \delta_{m,h,i,r,s}
  +
  D_m^{\rm rec}\exp\left\{-c_0\frac{n_m}{\ell_m}\right\}
  \right]<\infty .
\]
Then, almost surely, for all sufficiently large \(m\), every active stream has
a Spencer-scale signing:
\[
  \min_{\eps\in\{\pm1\}^{n_m}}
  \left\|\sum_{i=1}^{n_m}\eps_i
    a_{m,h,i}e_{I_{m,h,i}}e_{J_{m,h,i}}^*\right\|_{\op}
  \le C(C_0)\sqrt{n_m}
  \qquad(1\le h\le J_m).
\]
Sampled certificates follow under the same additional \(2^{-M}\)-summability
hypotheses as in
Corollary~\ref{cor:finite-menu-profile-capped-predictable-coordinate-schedules}.

Moreover, the caps may be selected entry by entry.  Let
\(L_m=R\ell_m\).  If, for fixed \(\alpha>0\) and \(A_{\rm tail}\ge1\),
\[
  \mathbb P\{|a_{m,h,i}|>s_{m,h,i,r,s}u,\
  I_{m,h,i}=r,\ J_{m,h,i}=s\}
  \le A_{\rm tail}e^{-u^\alpha}p_{m,h,i,r,s},
\]
then choose
\[
  q_{m,h,i,r,s}=s_{m,h,i,r,s}L_m^{1/\alpha}.
\]
If instead, for a fixed \(p_0>0\),
\[
  \mathbb E\left[
  |a_{m,h,i}|^{p_0}
  \mathbf 1_{\{I_{m,h,i}=r,\ J_{m,h,i}=s\}}\right]
  \le s_{m,h,i,r,s}^{p_0}p_{m,h,i,r,s},
\]
then choose
\[
  q_{m,h,i,r,s}=s_{m,h,i,r,s}\exp\{L_m/p_0\}.
\]
In either case, assuming all chosen caps are at most \(1\), the cap-tail
contribution for stream \((m,h)\) is at most
\[
  A_*\exp\{-L_m\}B_{m,h}^{\rm pair},
  \qquad
  B_{m,h}^{\rm pair}:=
  \sum_{i=1}^{n_m}\sum_{r=1}^{d_{m,1}}\sum_{s=1}^{d_{m,2}}
  p_{m,h,i,r,s},
\]
where \(A_*=A_{\rm tail}\) in the sub-Weibull case and \(A_*=1\) in the
moment case.  Thus, whenever \(B_{m,h}^{\rm pair}\le n_m\), the same
all-integer and sparse stretched-exponential schedule thresholds as in
Corollary~\ref{cor:closed-form-entry-intensity-cap-schedules} apply, but with
the sharper site-weighted profile \(\Xi_{m,h}^{\rm ent}\) in place of the
coarser product \(q_{m,h}^2\Delta_{m,h}^{\rm ent}\).
\end{corollary}

\begin{proof}
For stream \((m,h)\), define the capped adapted contraction
\[
  Y_{m,h,i}:=
  a_{m,h,i}
  \mathbf 1_{\{|a_{m,h,i}|\le q_{m,h,i,I_{m,h,i},J_{m,h,i}}\}}
  e_{I_{m,h,i}}e_{J_{m,h,i}}^* .
\]
The cap-exception event has probability at most the displayed sum of
\(\delta_{m,h,i,r,s}\), and outside this event \(Y_{m,h,i}=X_{m,h,i}\) for all
\(i\).  Since all caps are at most \(1\), the \(Y_{m,h,i}\) are contractions.
For each row coordinate,
\[
  \mathbb E(Y_{m,h,i}Y_{m,h,i}^*\mid\mcF_{m,h,i-1})_{r,r}
  \le
  \sum_{s=1}^{d_{m,2}}q_{m,h,i,r,s}^2p_{m,h,i,r,s},
\]
and the off-diagonal entries vanish.  The column profile is bounded similarly
by the sums over \(r\).  Hence the predictable rectangular square profile of
the capped stream is at most \(\Xi_{m,h}^{\rm ent}\).  Apply
Corollary~\ref{cor:predictable-rectangular-square-profile} with
\[
  v=\Xi_{m,h}^{\rm ent},
  \qquad
  t=\frac{C_0n_m}{2\ell_m}.
\]
The half-budget gives \(v\le t\), so the structural exponent is at least
\[
  \frac{t^2}{2(v+t/3)}\ge \frac{3t}{8}
  =c_0\frac{n_m}{\ell_m}.
\]
Adding the cap-exception probability and summing over \(m,h\), Borel--Cantelli
gives the eventual simultaneous conclusion; the sampled-certificate statement
comes from the sampled version of the same rectangular square-profile
corollary.

For the sub-Weibull site-tail model, the chosen cap gives
\(\delta_{m,h,i,r,s}=A_{\rm tail}e^{-L_m}p_{m,h,i,r,s}\).  For the moment
site-tail model, Markov's inequality gives the same estimate with \(A_*=1\).
Summing these bounds over \(i,r,s\) gives the displayed cap-tail contribution.
When \(B_{m,h}^{\rm pair}\le n_m\), the cap-tail summability and structural
summability are exactly the two series tests computed in
Corollary~\ref{cor:closed-form-entry-intensity-cap-schedules}.
\end{proof}

\begin{corollary}[Rank-sensitive predictable square-profile ensembles]
\label{cor:rank-sensitive-predictable-square-profile}
Let \(D\ge2\), let \(1\le\kappa\le D\), and let
\((\mcF_i)_{i=0}^n\) be a filtration.  Let \(P\) be an
\(\mcF_0\)-measurable orthogonal projection with
\[
  \operatorname{rank}P\le\kappa .
\]
Let \(A_1,\ldots,A_n\in\M_D(\C)\) be self-adjoint contractions such that
\[
  A_i\text{ is }\mcF_i\text{-measurable},\qquad A_i=PA_iP .
\]
Put
\[
  \nu_\kappa:=
  \left\|
  \sum_{i=1}^n \mathbb E(A_i^2\mid\mcF_{i-1})
  \right\|_{\op}.
\]
Then, for every \(C_0\ge1\), there is \(C(C_0)\) such that, for every
\(v,t>0\) satisfying
\[
  v+t\le \frac{C_0n}{\log(e\kappa)},
\]
one has
\[
  \mathbb P\left\{
  \disc(A_1,\ldots,A_n)>C(C_0)\sqrt n
  \right\}
  \le
  \mathbb P\{\nu_\kappa>v\}
  +
  \kappa\exp\left\{-\frac{t^2}{2(v+t/3)}\right\}.
\]
Moreover, if \(\eta^{(1)},\ldots,\eta^{(m)}\) are independent uniform sign
vectors independent of \(\mcF_n\), then, after increasing \(C(C_0)\) if needed,
\[
  \begin{aligned}
  &\mathbb P\left\{
  \min_{1\le r\le m}
  \left\|\sum_{i=1}^n\eta_i^{(r)}A_i\right\|_{\op}
  >C(C_0)\sqrt n
  \right\} \\
  &\qquad\le
  \mathbb P\{\nu_\kappa>v\}
  +
  \kappa\exp\left\{-\frac{t^2}{2(v+t/3)}\right\}
  +2^{-m}.
  \end{aligned}
\]
In particular, if
\[
  \nu_\kappa\le \frac{C_0n}{2\log(e\kappa)}
\]
almost surely, then both bounds hold with structural failure at most
\[
  \kappa\exp\left\{-\frac{3C_0n}{16\log(e\kappa)}\right\},
\]
with the additional \(2^{-m}\) term for sampled sign search.  Thus adapted
random contractions supported in a common deterministic or random
\(\kappa\)-dimensional corner obey the square-profile Matrix Spencer theorem
with \(D\) replaced everywhere by the joint support rank.
\end{corollary}

\begin{proof}
Set \(B_i=A_i^2\) and \(M_i=\mathbb E(B_i\mid\mcF_{i-1})\).  Because
\(P\in\mcF_0\subseteq\mcF_{i-1}\) and \(B_i=PB_iP\), also \(M_i=PM_iP\).  Hence
the martingale differences
\[
  X_i:=B_i-M_i
\]
are supported on \(\operatorname{ran}P\).  Restrict the matrices to this space,
whose dimension is at most \(\kappa\).  The proof of
Corollary~\ref{cor:random-square-profile-freedman}, now in dimension
\(\operatorname{rank}P\), gives on \(\{\nu_\kappa\le v\}\)
\[
  \mathbb P\left\{
  \left\|\sum_i A_i^2\right\|_{\op}>v+t
  \ \text{and}\ \nu_\kappa\le v
  \right\}
  \le
  \kappa\exp\left\{-\frac{t^2}{2(v+t/3)}\right\}.
\]
On the complementary event the realized square function has rank at most
\(\kappa\) and norm at most \(v+t\).  The rank-sensitive part of
Proposition~\ref{prop:square-function-khintchine} therefore yields
\[
  \disc(A_1,\ldots,A_n)
  \le C\sqrt{\log(e\kappa)}\sqrt{v+t}
  \le C(C_0)\sqrt n .
\]
This proves the first estimate.  The sampled-search estimate follows by
conditioning on the realized tuple and applying
Corollary~\ref{cor:random-sign-search-square-profile} in its rank-sensitive
form.  The half-budget specialization is the choice
\(v=t=C_0n/(2\log(e\kappa))\).
\end{proof}

\begin{proposition}[Two-scale square-function certificate]
\label{prop:two-scale-square-function-certificate}
There is a universal constant \(C\) with the following property.  Let
\(D\ge2\), let \(1\le\kappa\le D\), and let \(P\) be an orthogonal projection
on \(\C^D\) with \(\operatorname{rank}P\le\kappa\).  Suppose
\[
  A_i=B_i+R_i,\qquad 1\le i\le n,
\]
where all \(A_i,B_i,R_i\) are self-adjoint and \(B_i=PB_iP\).  Put
\[
  \Theta(B,R):=
  \sqrt{\log(e\kappa)}
  \left\|\sum_{i=1}^nB_i^2\right\|_{\op}^{1/2}
  +
  \sqrt{\log(eD)}
  \left\|\sum_{i=1}^nR_i^2\right\|_{\op}^{1/2}.
\]
Then
\[
  \disc(A_1,\ldots,A_n)\le C\,\Theta(B,R).
\]
Moreover, if \(\eta^{(1)},\ldots,\eta^{(m)}\) are independent uniform sign
vectors in \(\{\pm1\}^n\), then, conditionally on the fixed decomposition,
\[
  \mathbb P_\eta\left\{
  \min_{1\le r\le m}
  \left\|\sum_{i=1}^n\eta_i^{(r)}A_i\right\|_{\op}
  >2C\,\Theta(B,R)
  \right\}
  \le 2^{-m}.
\]
Consequently, a tuple with a \(\kappa\)-dimensional square-function core and a
small ambient square-function residual has a single Spencer-scale signing:
if
\[
  \left\|\sum_iB_i^2\right\|_{\op}\le
  \frac{C_{\mathrm{lo}}n}{\log(e\kappa)},
  \qquad
  \left\|\sum_iR_i^2\right\|_{\op}\le
  \frac{C_{\mathrm{hi}}n}{\log(eD)},
\]
then
\[
  \disc(A_1,\ldots,A_n)
  \le C(C_{\mathrm{lo}},C_{\mathrm{hi}})\sqrt n,
\]
and \(m\) sampled sign vectors fail to find such a signing with probability at
most \(2^{-m}\).
\end{proposition}

\begin{proof}
For one uniform sign vector \(\eta\), the triangle inequality gives
\[
  \mathbb E_\eta
  \left\|\sum_i\eta_iA_i\right\|_{\op}
  \le
  \mathbb E_\eta
  \left\|\sum_i\eta_iB_i\right\|_{\op}
  +
  \mathbb E_\eta
  \left\|\sum_i\eta_iR_i\right\|_{\op}.
\]
The first tuple acts on \(\operatorname{ran}P\), whose dimension is at most
\(\kappa\), so the rank-sensitive noncommutative Khintchine estimate used in
Proposition~\ref{prop:square-function-khintchine} bounds the first term by
\[
  C\sqrt{\log(e\kappa)}
  \left\|\sum_iB_i^2\right\|_{\op}^{1/2}.
\]
The same inequality in the ambient dimension \(D\) bounds the residual term by
\[
  C\sqrt{\log(eD)}
  \left\|\sum_iR_i^2\right\|_{\op}^{1/2}.
\]
Hence the expected norm of one signed sum is at most \(C\Theta(B,R)\), and the
minimum over all signs is bounded by that expectation.  Markov's inequality
makes one sampled sign vector successful with probability at least \(1/2\);
independence gives the \(2^{-m}\) failure bound for \(m\) trials.  The final
display follows by substituting the two square-function budgets into
\(\Theta(B,R)\).
\end{proof}

\begin{corollary}[Two-scale predictable square-profile ensembles]
\label{cor:two-scale-predictable-square-profile}
Let \(D\ge2\), \(1\le\kappa\le D\), and let
\((\mcF_i)_{i=0}^n\) be a filtration.  Let \(P\) be an
\(\mcF_0\)-measurable orthogonal projection with
\(\operatorname{rank}P\le\kappa\).  Suppose
\[
  A_i=B_i+R_i,\qquad 1\le i\le n,
\]
where \(B_i\) and \(R_i\) are \(\mcF_i\)-measurable self-adjoint contractions
and \(B_i=PB_iP\).  Define the predictable core and residual square profiles by
\[
  \nu_{\mathrm{lo}}:=
  \left\|\sum_{i=1}^n\mathbb E(B_i^2\mid\mcF_{i-1})\right\|_{\op},
  \qquad
  \nu_{\mathrm{hi}}:=
  \left\|\sum_{i=1}^n\mathbb E(R_i^2\mid\mcF_{i-1})\right\|_{\op}.
\]
Then, for every \(C_{\mathrm{lo}},C_{\mathrm{hi}}\ge1\), there is
\(C(C_{\mathrm{lo}},C_{\mathrm{hi}})\) such that, whenever
\[
  v_{\mathrm{lo}}+t_{\mathrm{lo}}
  \le\frac{C_{\mathrm{lo}}n}{\log(e\kappa)},
  \qquad
  v_{\mathrm{hi}}+t_{\mathrm{hi}}
  \le\frac{C_{\mathrm{hi}}n}{\log(eD)},
\]
one has
\[
\begin{aligned}
  &\mathbb P\left\{
  \disc(A_1,\ldots,A_n)>
  C(C_{\mathrm{lo}},C_{\mathrm{hi}})\sqrt n
  \right\} \\
  &\quad\le
  \mathbb P\{\nu_{\mathrm{lo}}>v_{\mathrm{lo}}\}
  +
  \kappa
  \exp\left\{
  -\frac{t_{\mathrm{lo}}^2}
        {2(v_{\mathrm{lo}}+t_{\mathrm{lo}}/3)}
  \right\} \\
  &\qquad\quad+
  \mathbb P\{\nu_{\mathrm{hi}}>v_{\mathrm{hi}}\}
  +
  D
  \exp\left\{
  -\frac{t_{\mathrm{hi}}^2}
        {2(v_{\mathrm{hi}}+t_{\mathrm{hi}}/3)}
  \right\}.
\end{aligned}
\]
If \(\eta^{(1)},\ldots,\eta^{(m)}\) are independent uniform sign vectors
independent of \(\mcF_n\), then the same bound holds for sampled sign search
after adding \(2^{-m}\).  In particular, if the half-budget bounds
\[
  \nu_{\mathrm{lo}}\le
  \frac{C_{\mathrm{lo}}n}{2\log(e\kappa)},
  \qquad
  \nu_{\mathrm{hi}}\le
  \frac{C_{\mathrm{hi}}n}{2\log(eD)}
\]
hold outside events \(E_{\mathrm{lo}}\) and \(E_{\mathrm{hi}}\), respectively,
then the failure probability is at most
\[
  \mathbb P(E_{\mathrm{lo}}\cup E_{\mathrm{hi}})
  +
  \kappa
  \exp\left\{-\frac{3C_{\mathrm{lo}}n}{16\log(e\kappa)}\right\}
  +
  D
  \exp\left\{-\frac{3C_{\mathrm{hi}}n}{16\log(eD)}\right\},
\]
with an additional \(2^{-m}\) for randomized search.

Thus the exact common-corner theorem is stable under a full ambient residual,
provided that residual has its own logarithmic square-function budget.  This
covers random tuples which are not themselves low-rank but decompose into a
low-rank adaptive core plus diffuse martingale, twirl, scalar-multiplier, or
localized noise.
\end{corollary}

\begin{proof}
Apply the proof of
Corollary~\ref{cor:rank-sensitive-predictable-square-profile} to the adapted
core tuple \(B_i\).  On \(\{\nu_{\mathrm{lo}}\le v_{\mathrm{lo}}\}\), except
with probability at most
\[
  \kappa
  \exp\left\{
  -\frac{t_{\mathrm{lo}}^2}
        {2(v_{\mathrm{lo}}+t_{\mathrm{lo}}/3)}
  \right\},
\]
one has
\[
  \left\|\sum_iB_i^2\right\|_{\op}
  \le v_{\mathrm{lo}}+t_{\mathrm{lo}}.
\]
Similarly, the proof of
Corollary~\ref{cor:random-square-profile-freedman} applied to the residual
tuple \(R_i\) gives, on \(\{\nu_{\mathrm{hi}}\le v_{\mathrm{hi}}\}\), except
with probability at most
\[
  D
  \exp\left\{
  -\frac{t_{\mathrm{hi}}^2}
        {2(v_{\mathrm{hi}}+t_{\mathrm{hi}}/3)}
  \right\},
\]
the bound
\[
  \left\|\sum_iR_i^2\right\|_{\op}
  \le v_{\mathrm{hi}}+t_{\mathrm{hi}}.
\]
On the intersection of these two square-function events,
Proposition~\ref{prop:two-scale-square-function-certificate} gives
\[
  \disc(A_1,\ldots,A_n)
  \le
  C\sqrt{\log(e\kappa)}
  \sqrt{v_{\mathrm{lo}}+t_{\mathrm{lo}}}
  +
  C\sqrt{\log(eD)}
  \sqrt{v_{\mathrm{hi}}+t_{\mathrm{hi}}}
  \le
  C(C_{\mathrm{lo}},C_{\mathrm{hi}})\sqrt n.
\]
The sampled-search statement follows by conditioning on the realized
decomposition and applying the sampled part of
Proposition~\ref{prop:two-scale-square-function-certificate}.  The half-budget
form is the choice
\[
  v_{\mathrm{lo}}=t_{\mathrm{lo}}=
  \frac{C_{\mathrm{lo}}n}{2\log(e\kappa)},
  \qquad
  v_{\mathrm{hi}}=t_{\mathrm{hi}}=
  \frac{C_{\mathrm{hi}}n}{2\log(eD)}.
\]
\end{proof}

\begin{corollary}[Finite-menu two-scale predictable square-profile schedules]
\label{cor:finite-menu-two-scale-predictable-square-profile-schedules}
Fix \(C_{\mathrm{lo}},C_{\mathrm{hi}}\ge1\), and put
\[
  \gamma_{\mathrm{lo}}:=\frac{3C_{\mathrm{lo}}}{16},
  \qquad
  \gamma_{\mathrm{hi}}:=\frac{3C_{\mathrm{hi}}}{16}.
\]
For each \(m\ge1\), let \(D_m\ge2\), let \(J_m\ge1\), and, for
\(1\le h\le J_m\), let \(1\le\kappa_{m,h}\le D_m\).  Put
\[
  \ell_{m,h}^{\kappa}:=\log(e\kappa_{m,h}),
  \qquad
  \ell_m^D:=\log(eD_m).
\]
For every active pair \((m,h)\), suppose that
\[
  A_{m,h,i}=B_{m,h,i}+R_{m,h,i},
  \qquad 1\le i\le n_m,
\]
is a two-scale predictable stream as in
Corollary~\ref{cor:two-scale-predictable-square-profile}, with core rank
envelope \(\kappa_{m,h}\).  Let
\[
  \nu_{m,h}^{\mathrm{lo}}:=
  \left\|
  \sum_i\mathbb E(B_{m,h,i}^2\mid\mcF_{m,h,i-1})
  \right\|_{\op},
  \quad
  \nu_{m,h}^{\mathrm{hi}}:=
  \left\|
  \sum_i\mathbb E(R_{m,h,i}^2\mid\mcF_{m,h,i-1})
  \right\|_{\op}.
\]
Assume that
\[
  \nu_{m,h}^{\mathrm{lo}}
  \le\frac{C_{\mathrm{lo}}n_m}{2\ell_{m,h}^{\kappa}},
  \qquad
  \nu_{m,h}^{\mathrm{hi}}
  \le\frac{C_{\mathrm{hi}}n_m}{2\ell_m^D}
  \qquad(1\le h\le J_m)
\]
outside events \(E_{m,h}^{\mathrm{lo}}\) and \(E_{m,h}^{\mathrm{hi}}\), and
\[
  \sum_m\sum_{h=1}^{J_m}
  \mathbb P(E_{m,h}^{\mathrm{lo}}\cup E_{m,h}^{\mathrm{hi}})<\infty,
\]
\[
  \sum_m\sum_{h=1}^{J_m}
  \kappa_{m,h}
  \exp\left\{-\gamma_{\mathrm{lo}}
  \frac{n_m}{\ell_{m,h}^{\kappa}}\right\}<\infty,
  \qquad
  \sum_m J_mD_m
  \exp\left\{-\gamma_{\mathrm{hi}}
  \frac{n_m}{\ell_m^D}\right\}<\infty.
\]
Then, with probability one, for all sufficiently large \(m\), every active
two-scale stream has a Spencer-scale signing:
\[
  \disc(A_{m,h,1},\ldots,A_{m,h,n_m})
  \le C(C_{\mathrm{lo}},C_{\mathrm{hi}})\sqrt{n_m}
  \qquad(1\le h\le J_m).
\]
If the streams are supplied with independent sampled sign batches and
\[
  \sum_m\sum_{h=1}^{J_m}2^{-M_{m,h}}<\infty,
\]
then the same eventual statement holds for the sampled certificates.  If one
common batch of \(M_m\) sign vectors is reused across the active menu, it is
enough that
\[
  \sum_m J_m2^{-M_m}<\infty .
\]

A concrete sufficient growth condition is the following.  Define
\[
  \Lambda_m^\kappa:=\sum_{h=1}^{J_m}\kappa_{m,h},
  \qquad
  \ell_m^\kappa:=\max_{1\le h\le J_m}\ell_{m,h}^{\kappa}.
\]
If, for some \(\delta>0\),
\[
  n_m\ge
  \max\left\{
  \gamma_{\mathrm{lo}}^{-1}\ell_m^\kappa
  \bigl(\log(e\Lambda_m^\kappa)+(1+\delta)\log(e+m)\bigr),
  \gamma_{\mathrm{hi}}^{-1}\ell_m^D
  \bigl(\log(eJ_mD_m)+(1+\delta)\log(e+m)\bigr)
  \right\}
\]
eventually, then the two structural series above are summable.  With a common
sampled batch, it is enough in addition that
\[
  M_m\ge \log_2(eJ_m)+(1+\delta')\log_2(e+m)
\]
eventually for some \(\delta'>0\).  If the streams are defined for all
\(h\ge1\), \(J_m\to\infty\), and \(h\le J_m\) is a deterministic exhaustion,
then every fixed two-scale predictable random-matrix model in the countable
library is eventually covered.
\end{corollary}

\begin{proof}
For a fixed active pair \((m,h)\), apply
Corollary~\ref{cor:two-scale-predictable-square-profile} with half-budget
choices.  The failure probability is at most
\[
  \mathbb P(E_{m,h}^{\mathrm{lo}}\cup E_{m,h}^{\mathrm{hi}})
  +
  \kappa_{m,h}
  \exp\left\{-\gamma_{\mathrm{lo}}
  \frac{n_m}{\ell_{m,h}^{\kappa}}\right\}
  +
  D_m
  \exp\left\{-\gamma_{\mathrm{hi}}
  \frac{n_m}{\ell_m^D}\right\}.
\]
Union over \(1\le h\le J_m\), then apply the first Borel--Cantelli lemma.  The
sampled-sign assertions add either \(\sum_h2^{-M_{m,h}}\) or the common-batch
union bound \(J_m2^{-M_m}\).

For the concrete schedule, the low-rank structural cost is bounded by
\[
  \Lambda_m^\kappa
  \exp\left\{-\gamma_{\mathrm{lo}}
  \frac{n_m}{\ell_m^\kappa}\right\}
  \le e^{-1}(e+m)^{-(1+\delta)},
\]
and the ambient residual cost is bounded by
\[
  J_mD_m
  \exp\left\{-\gamma_{\mathrm{hi}}
  \frac{n_m}{\ell_m^D}\right\}
  \le e^{-1}(e+m)^{-(1+\delta)}.
\]
The sampled-batch condition is summable by the displayed lower bound on
\(M_m\).  If \(J_m\to\infty\), every fixed countable-menu element is active
eventually.
\end{proof}

\begin{proposition}[Finite multiscale square-function certificate]
\label{prop:finite-multiscale-square-function-certificate}
There is a universal constant \(C\) with the following property.  Let
\(D\ge2\), let \(Q\ge1\), let \(1\le\kappa_q\le D\), and, for
\(1\le q\le Q\), let \(P^{(q)}\) be an orthogonal projection on \(\C^D\) with
\(\operatorname{rank}P^{(q)}\le\kappa_q\).  Suppose
\[
  A_i=R_i+\sum_{q=1}^QB_i^{(q)},\qquad 1\le i\le n,
\]
where all matrices are self-adjoint and
\[
  B_i^{(q)}=P^{(q)}B_i^{(q)}P^{(q)}
  \qquad(1\le q\le Q,\ 1\le i\le n).
\]
Put
\[
  \Theta_Q(B,R):=
  \sum_{q=1}^Q
  \sqrt{\log(e\kappa_q)}
  \left\|\sum_{i=1}^n(B_i^{(q)})^2\right\|_{\op}^{1/2}
  +
  \sqrt{\log(eD)}
  \left\|\sum_{i=1}^nR_i^2\right\|_{\op}^{1/2}.
\]
Then
\[
  \disc(A_1,\ldots,A_n)\le C\,\Theta_Q(B,R).
\]
Moreover, if \(\eta^{(1)},\ldots,\eta^{(m)}\) are independent uniform sign
vectors in \(\{\pm1\}^n\), then, conditionally on the fixed decomposition,
\[
  \mathbb P_\eta\left\{
  \min_{1\le r\le m}
  \left\|\sum_{i=1}^n\eta_i^{(r)}A_i\right\|_{\op}
  >2C\,\Theta_Q(B,R)
  \right\}
  \le2^{-m}.
\]
In particular, if
\[
  \left\|\sum_i(B_i^{(q)})^2\right\|_{\op}
  \le \frac{a_qn}{\log(e\kappa_q)}
  \quad(1\le q\le Q),
  \qquad
  \left\|\sum_iR_i^2\right\|_{\op}
  \le\frac{a_\ast n}{\log(eD)},
\]
then
\[
  \disc(A_1,\ldots,A_n)
  \le
  C\left(\sum_{q=1}^Q\sqrt{a_q}+\sqrt{a_\ast}\right)\sqrt n,
\]
and the same threshold is found by \(m\) sampled sign vectors except with
probability at most \(2^{-m}\).
\end{proposition}

\begin{proof}
For one uniform sign vector \(\eta\), the triangle inequality gives
\[
  \mathbb E_\eta
  \left\|\sum_i\eta_iA_i\right\|_{\op}
  \le
  \sum_{q=1}^Q
  \mathbb E_\eta
  \left\|\sum_i\eta_iB_i^{(q)}\right\|_{\op}
  +
  \mathbb E_\eta
  \left\|\sum_i\eta_iR_i\right\|_{\op}.
\]
For each \(q\), the tuple \(B_i^{(q)}\) acts on
\(\operatorname{ran}P^{(q)}\), whose dimension is at most \(\kappa_q\).  The
rank-sensitive noncommutative Khintchine estimate from
Proposition~\ref{prop:square-function-khintchine} therefore bounds the
\(q\)-th term by
\[
  C\sqrt{\log(e\kappa_q)}
  \left\|\sum_i(B_i^{(q)})^2\right\|_{\op}^{1/2}.
\]
The residual term is bounded in the ambient dimension \(D\) by the same
inequality.  Thus the expected signed norm is at most \(C\Theta_Q(B,R)\), so
some signing has that norm.  Markov's inequality gives success probability at
least \(1/2\) for one sampled sign vector at threshold \(2C\Theta_Q(B,R)\), and
independence gives \(2^{-m}\) for \(m\) trials.  The final display follows by
substituting the square-function budgets into \(\Theta_Q(B,R)\).
\end{proof}

\begin{corollary}[Finite multiscale predictable square-profile ensembles]
\label{cor:finite-multiscale-predictable-square-profile}
Let \(D\ge2\), let \(Q\ge1\), let \(1\le\kappa_q\le D\), and let
\((\mcF_i)_{i=0}^n\) be a filtration.  For \(1\le q\le Q\), let
\(P^{(q)}\) be an \(\mcF_0\)-measurable orthogonal
projection with \(\operatorname{rank}P^{(q)}\le\kappa_q\).  Suppose
\[
  A_i=R_i+\sum_{q=1}^QB_i^{(q)},\qquad 1\le i\le n,
\]
where \(R_i\) and \(B_i^{(q)}\) are \(\mcF_i\)-measurable self-adjoint
contractions and \(B_i^{(q)}=P^{(q)}B_i^{(q)}P^{(q)}\).  Define
\[
  \nu_q:=
  \left\|\sum_{i=1}^n
  \mathbb E\bigl((B_i^{(q)})^2\mid\mcF_{i-1}\bigr)\right\|_{\op},
  \qquad
  \nu_\ast:=
  \left\|\sum_{i=1}^n
  \mathbb E(R_i^2\mid\mcF_{i-1})\right\|_{\op}.
\]
Let \(a_1,\ldots,a_Q,a_\ast>0\).  Whenever
\[
  v_q+t_q\le\frac{a_qn}{\log(e\kappa_q)}
  \quad(1\le q\le Q),
  \qquad
  v_\ast+t_\ast\le\frac{a_\ast n}{\log(eD)},
\]
one has
\[
\begin{aligned}
  &\mathbb P\left\{
  \disc(A_1,\ldots,A_n)>
  C\left(\sum_{q=1}^Q\sqrt{a_q}+\sqrt{a_\ast}\right)\sqrt n
  \right\} \\
  &\quad\le
  \sum_{q=1}^Q
  \left[
  \mathbb P\{\nu_q>v_q\}
  +
  \kappa_q
  \exp\left\{-\frac{t_q^2}{2(v_q+t_q/3)}\right\}
  \right] \\
  &\qquad\quad+
  \mathbb P\{\nu_\ast>v_\ast\}
  +
  D\exp\left\{-\frac{t_\ast^2}{2(v_\ast+t_\ast/3)}\right\}.
\end{aligned}
\]
If \(\eta^{(1)},\ldots,\eta^{(m)}\) are independent uniform sign vectors
independent of \(\mcF_n\), then the same bound holds for sampled sign search
after adding \(2^{-m}\).

In particular, if the half-budget inequalities
\[
  \nu_q\le\frac{a_qn}{2\log(e\kappa_q)}
  \quad(1\le q\le Q),
  \qquad
  \nu_\ast\le\frac{a_\ast n}{2\log(eD)}
\]
hold outside events \(E_q\) and \(E_\ast\), then the failure probability is at
most
\[
  \mathbb P\left(E_\ast\cup\bigcup_{q=1}^QE_q\right)
  +
  \sum_{q=1}^Q
  \kappa_q
  \exp\left\{-\frac{3a_qn}{16\log(e\kappa_q)}\right\}
  +
  D\exp\left\{-\frac{3a_\ast n}{16\log(eD)}\right\},
\]
with an additional \(2^{-m}\) for sampled sign search.

Thus finite-depth mixtures of adaptive low-rank cores and one full ambient
residual are Matrix-Spencer scale whenever the total multiscale budget
\(\sum_q\sqrt{a_q}+\sqrt{a_\ast}\) is bounded.
\end{corollary}

\begin{proof}
For each \(q\), apply the proof of
Corollary~\ref{cor:rank-sensitive-predictable-square-profile} to the adapted
tuple \(B_i^{(q)}\) on the \(\kappa_q\)-dimensional support
\(\operatorname{ran}P^{(q)}\).  On \(\{\nu_q\le v_q\}\), except with
probability at most
\[
  \kappa_q\exp\left\{-\frac{t_q^2}{2(v_q+t_q/3)}\right\},
\]
one has
\[
  \left\|\sum_i(B_i^{(q)})^2\right\|_{\op}\le v_q+t_q.
\]
Apply the proof of Corollary~\ref{cor:random-square-profile-freedman} to the
ambient residual tuple \(R_i\); on \(\{\nu_\ast\le v_\ast\}\), except with
probability at most
\[
  D\exp\left\{-\frac{t_\ast^2}{2(v_\ast+t_\ast/3)}\right\},
\]
one has \(\|\sum_iR_i^2\|_{\op}\le v_\ast+t_\ast\).  On the intersection of
all these square-function events,
Proposition~\ref{prop:finite-multiscale-square-function-certificate} gives the
displayed discrepancy bound.  The sampled-search statement follows by
conditioning on the realized multiscale decomposition and applying the sampled
part of that proposition.  The half-budget form is obtained by taking
\[
  v_q=t_q=\frac{a_qn}{2\log(e\kappa_q)}
  \quad(1\le q\le Q),
  \qquad
  v_\ast=t_\ast=\frac{a_\ast n}{2\log(eD)}.
\]
\end{proof}

\begin{corollary}[Finite-menu multiscale predictable schedules]
\label{cor:finite-menu-multiscale-predictable-square-profile-schedules}
For each \(m\ge1\), let \(D_m\ge2\) and \(J_m\ge1\).  For every active
\((m,h)\), let \(Q_{m,h}\ge1\), let
\[
  1\le\kappa_{m,h,q}\le D_m\qquad(1\le q\le Q_{m,h}),
\]
and suppose
\[
  A_{m,h,i}=R_{m,h,i}+\sum_{q=1}^{Q_{m,h}}B_{m,h,i}^{(q)},
  \qquad 1\le i\le n_m,
\]
is a finite multiscale predictable stream as in
Corollary~\ref{cor:finite-multiscale-predictable-square-profile}.  Let
\[
  \nu_{m,h,q}:=
  \left\|
  \sum_i
  \mathbb E\bigl((B_{m,h,i}^{(q)})^2\mid\mcF_{m,h,i-1}\bigr)
  \right\|_{\op},
  \qquad
  \nu_{m,h,\ast}:=
  \left\|
  \sum_i
  \mathbb E(R_{m,h,i}^2\mid\mcF_{m,h,i-1})
  \right\|_{\op}.
\]
Let
\[
  a_{m,h,q}>0,\qquad a_{m,h,\ast}>0,
\]
and assume the uniform multiscale size bound
\[
  S_\ast:=
  \sup_{m,h}
  \left(
  \sum_{q=1}^{Q_{m,h}}\sqrt{a_{m,h,q}}+\sqrt{a_{m,h,\ast}}
  \right)<\infty .
\]
Assume that the half-budget profile inequalities
\[
  \nu_{m,h,q}\le
  \frac{a_{m,h,q}n_m}{2\log(e\kappa_{m,h,q})}
  \qquad(1\le q\le Q_{m,h}),
  \qquad
  \nu_{m,h,\ast}\le
  \frac{a_{m,h,\ast}n_m}{2\log(eD_m)}
\]
hold outside events \(E_{m,h,q}\) and \(E_{m,h,\ast}\), and that
\[
  \sum_m\sum_{h=1}^{J_m}
  \mathbb P\left(E_{m,h,\ast}\cup
  \bigcup_{q=1}^{Q_{m,h}}E_{m,h,q}\right)<\infty,
\]
\[
  \sum_m\sum_{h=1}^{J_m}\sum_{q=1}^{Q_{m,h}}
  \kappa_{m,h,q}
  \exp\left\{-\frac{3a_{m,h,q}n_m}
  {16\log(e\kappa_{m,h,q})}\right\}<\infty,
\]
and
\[
  \sum_m\sum_{h=1}^{J_m}
  D_m
  \exp\left\{-\frac{3a_{m,h,\ast}n_m}{16\log(eD_m)}\right\}<\infty .
\]
Then, with probability one, for all sufficiently large \(m\), every active
multiscale stream has a Spencer-scale signing:
\[
  \disc(A_{m,h,1},\ldots,A_{m,h,n_m})
  \le C S_\ast\sqrt{n_m}
  \qquad(1\le h\le J_m).
\]
Sampled sign certificates are eventual almost surely if independent batches
satisfy
\[
  \sum_m\sum_{h=1}^{J_m}2^{-M_{m,h}}<\infty,
\]
or, for one common active-menu batch, if
\[
  \sum_mJ_m2^{-M_m}<\infty .
\]

For a concrete schedule, suppose \(a_{m,h,q}\ge a_{\mathrm{lo}}>0\) and
\(a_{m,h,\ast}\ge a_{\mathrm{hi}}>0\).  Put
\[
  \Lambda_m^{\mathrm{ms}}:=
  \sum_{h=1}^{J_m}\sum_{q=1}^{Q_{m,h}}\kappa_{m,h,q},
  \qquad
  \ell_m^{\mathrm{ms}}:=
  \max_{h,q}\log(e\kappa_{m,h,q}),
  \qquad
  \ell_m^D:=\log(eD_m).
\]
Then the two structural series are summable whenever, for some \(\delta>0\),
\[
  n_m\ge
  \max\left\{
  \frac{16}{3a_{\mathrm{lo}}}\,
  \ell_m^{\mathrm{ms}}
  \bigl(\log(e\Lambda_m^{\mathrm{ms}})+(1+\delta)\log(e+m)\bigr),
  \frac{16}{3a_{\mathrm{hi}}}\,
  \ell_m^D
  \bigl(\log(eJ_mD_m)+(1+\delta)\log(e+m)\bigr)
  \right\}
\]
eventually.  A common sampled batch is summable under
\[
  M_m\ge \log_2(eJ_m)+(1+\delta')\log_2(e+m)
\]
eventually.  If \(J_m\to\infty\), the active menus exhaust a countable library
of finite-depth multiscale random-matrix models.
\end{corollary}

\begin{proof}
Apply Corollary~\ref{cor:finite-multiscale-predictable-square-profile} to each
active pair \((m,h)\) with half-budget choices and union over
\(1\le h\le J_m\).  The three displayed summability hypotheses make the
resulting active-row failure probabilities summable, so Borel--Cantelli gives
the eventual discrepancy bound.  Sampled signs add the displayed search
failure series.

Under the concrete schedule,
\[
  \Lambda_m^{\mathrm{ms}}
  \exp\left\{-\frac{3a_{\mathrm{lo}}n_m}
  {16\ell_m^{\mathrm{ms}}}\right\}
  \le e^{-1}(e+m)^{-(1+\delta)}
\]
and
\[
  J_mD_m
  \exp\left\{-\frac{3a_{\mathrm{hi}}n_m}{16\ell_m^D}\right\}
  \le e^{-1}(e+m)^{-(1+\delta)}.
\]
The common sampled-batch estimate is identical to the two-scale schedule.
\end{proof}

\begin{proposition}[Summable multiscale square-function certificate]
\label{prop:summable-multiscale-square-function-certificate}
There is a universal constant \(C\) with the following property.  Let
\(D\ge2\).  For \(q\ge1\), let \(1\le\kappa_q\le D\), and let \(P^{(q)}\) be
an orthogonal projection on \(\C^D\) with
\(\operatorname{rank}P^{(q)}\le\kappa_q\).  Suppose
\[
  A_i=R_i+\sum_{q\ge1}B_i^{(q)},\qquad 1\le i\le n,
\]
where all matrices are self-adjoint, the series converges in operator norm for
each \(i\), and
\[
  B_i^{(q)}=P^{(q)}B_i^{(q)}P^{(q)}
  \qquad(q\ge1,\ 1\le i\le n).
\]
If
\[
  \Theta_\infty(B,R):=
  \sum_{q\ge1}
  \sqrt{\log(e\kappa_q)}
  \left\|\sum_{i=1}^n(B_i^{(q)})^2\right\|_{\op}^{1/2}
  +
  \sqrt{\log(eD)}
  \left\|\sum_{i=1}^nR_i^2\right\|_{\op}^{1/2}
  <\infty,
\]
then
\[
  \disc(A_1,\ldots,A_n)\le C\,\Theta_\infty(B,R).
\]
Moreover, if \(\eta^{(1)},\ldots,\eta^{(m)}\) are independent uniform sign
vectors in \(\{\pm1\}^n\), then, conditionally on the fixed decomposition,
\[
  \mathbb P_\eta\left\{
  \min_{1\le r\le m}
  \left\|\sum_{i=1}^n\eta_i^{(r)}A_i\right\|_{\op}
  >2C\,\Theta_\infty(B,R)
  \right\}
  \le2^{-m}.
\]
In particular, if
\[
  \left\|\sum_i(B_i^{(q)})^2\right\|_{\op}
  \le \frac{a_qn}{\log(e\kappa_q)}
  \quad(q\ge1),
  \qquad
  \left\|\sum_iR_i^2\right\|_{\op}
  \le\frac{a_\ast n}{\log(eD)}
\]
with \(\sum_q\sqrt{a_q}<\infty\), then
\[
  \disc(A_1,\ldots,A_n)
  \le
  C\left(\sum_{q\ge1}\sqrt{a_q}+\sqrt{a_\ast}\right)\sqrt n,
\]
and \(m\) sampled sign vectors fail to find that threshold with probability at
most \(2^{-m}\).
\end{proposition}

\begin{proof}
For \(Q\ge1\), put
\[
  A_i^{(Q)}=R_i+\sum_{q=1}^QB_i^{(q)}.
\]
By Proposition~\ref{prop:finite-multiscale-square-function-certificate},
\[
  \mathbb E_\eta
  \left\|\sum_i\eta_iA_i^{(Q)}\right\|_{\op}
  \le
  C\left[
  \sum_{q=1}^Q
  \sqrt{\log(e\kappa_q)}
  \left\|\sum_i(B_i^{(q)})^2\right\|_{\op}^{1/2}
  +
  \sqrt{\log(eD)}
  \left\|\sum_iR_i^2\right\|_{\op}^{1/2}
  \right].
\]
For every fixed sign vector, the operator-norm convergence of the layer series
implies
\[
  \sum_i\eta_iA_i^{(Q)}
  \longrightarrow
  \sum_i\eta_iA_i .
\]
Fatou's lemma therefore gives
\[
  \mathbb E_\eta
  \left\|\sum_i\eta_iA_i\right\|_{\op}
  \le C\,\Theta_\infty(B,R).
\]
The minimum over all signs is bounded by this expectation.  Markov's inequality
gives one-sample success probability at least \(1/2\) at threshold
\(2C\Theta_\infty(B,R)\), and independent trials give the \(2^{-m}\) failure
bound.  The final displayed estimate is the square-profile specialization of
\(\Theta_\infty(B,R)\).
\end{proof}

\begin{corollary}[Summable multiscale predictable square-profile ensembles]
\label{cor:summable-multiscale-predictable-square-profile}
Let \(D\ge2\), and let \((\mcF_i)_{i=0}^n\) be a filtration.  For
\(q\ge1\), let \(1\le\kappa_q\le D\), and let \(P^{(q)}\) be an
\(\mcF_0\)-measurable orthogonal projection with
\(\operatorname{rank}P^{(q)}\le\kappa_q\).  Suppose
\[
  A_i=R_i+\sum_{q\ge1}B_i^{(q)},\qquad 1\le i\le n,
\]
where \(R_i\) and \(B_i^{(q)}\) are \(\mcF_i\)-measurable self-adjoint
contractions, the series converges in operator norm for each \(i\), and
\[
  B_i^{(q)}=P^{(q)}B_i^{(q)}P^{(q)}.
\]
Define
\[
  \nu_q:=
  \left\|\sum_{i=1}^n
  \mathbb E\bigl((B_i^{(q)})^2\mid\mcF_{i-1}\bigr)\right\|_{\op},
  \qquad
  \nu_\ast:=
  \left\|\sum_{i=1}^n
  \mathbb E(R_i^2\mid\mcF_{i-1})\right\|_{\op}.
\]
Let \(a_\ast>0\) and \(a_q>0\), \(q\ge1\), with
\[
  A_{\mathrm{ms}}:=\sum_{q\ge1}\sqrt{a_q}+\sqrt{a_\ast}<\infty .
\]
Whenever
\[
  v_q+t_q\le\frac{a_qn}{\log(e\kappa_q)}
  \quad(q\ge1),
  \qquad
  v_\ast+t_\ast\le\frac{a_\ast n}{\log(eD)},
\]
one has
\[
\begin{aligned}
  &\mathbb P\left\{
  \disc(A_1,\ldots,A_n)>C A_{\mathrm{ms}}\sqrt n
  \right\} \\
  &\quad\le
  \sum_{q\ge1}
  \left[
  \mathbb P\{\nu_q>v_q\}
  +
  \kappa_q
  \exp\left\{-\frac{t_q^2}{2(v_q+t_q/3)}\right\}
  \right] \\
  &\qquad\quad+
  \mathbb P\{\nu_\ast>v_\ast\}
  +
  D\exp\left\{-\frac{t_\ast^2}{2(v_\ast+t_\ast/3)}\right\}.
\end{aligned}
\]
If \(\eta^{(1)},\ldots,\eta^{(m)}\) are independent uniform sign vectors
independent of \(\mcF_n\), then the same bound holds for sampled sign search
after adding \(2^{-m}\).

In particular, if the half-budget inequalities
\[
  \nu_q\le\frac{a_qn}{2\log(e\kappa_q)}
  \quad(q\ge1),
  \qquad
  \nu_\ast\le\frac{a_\ast n}{2\log(eD)}
\]
hold outside events \(E_q\) and \(E_\ast\), then the failure probability is at
most
\[
  \mathbb P\left(E_\ast\cup\bigcup_{q\ge1}E_q\right)
  +
  \sum_{q\ge1}
  \kappa_q
  \exp\left\{-\frac{3a_qn}{16\log(e\kappa_q)}\right\}
  +
  D\exp\left\{-\frac{3a_\ast n}{16\log(eD)}\right\},
\]
with an additional \(2^{-m}\) for sampled sign search.
\end{corollary}

\begin{proof}
For every fixed \(q\), the proof of
Corollary~\ref{cor:rank-sensitive-predictable-square-profile} gives, on
\(\{\nu_q\le v_q\}\), except with probability at most
\[
  \kappa_q\exp\left\{-\frac{t_q^2}{2(v_q+t_q/3)}\right\},
\]
the bound
\[
  \left\|\sum_i(B_i^{(q)})^2\right\|_{\op}\le v_q+t_q.
\]
The proof of Corollary~\ref{cor:random-square-profile-freedman} gives the
analogous residual event
\[
  \left\|\sum_iR_i^2\right\|_{\op}\le v_\ast+t_\ast
\]
outside the displayed ambient tail.  A countable union bound controls the
complement of the intersection of all layer and residual square-function
events.  On that intersection,
Proposition~\ref{prop:summable-multiscale-square-function-certificate} gives
the discrepancy bound.  The sampled-search statement follows by conditioning
on the realized decomposition.  The half-budget specialization is the choice
\[
  v_q=t_q=\frac{a_qn}{2\log(e\kappa_q)}
  \quad(q\ge1),
  \qquad
  v_\ast=t_\ast=\frac{a_\ast n}{2\log(eD)}.
\]
\end{proof}

\begin{corollary}[Active menus of summable multiscale streams]
\label{cor:active-menu-summable-multiscale-square-profile-schedules}
For each \(m\ge1\), let \(D_m\ge2\) and \(J_m\ge1\).  For each active
\((m,h)\), suppose
\[
  A_{m,h,i}=R_{m,h,i}+\sum_{q\ge1}B_{m,h,i}^{(q)},
  \qquad 1\le i\le n_m,
\]
is a summable multiscale predictable stream as in
Corollary~\ref{cor:summable-multiscale-predictable-square-profile}, with rank
envelopes \(\kappa_{m,h,q}\).  Let
\[
  \nu_{m,h,q}:=
  \left\|
  \sum_i
  \mathbb E\bigl((B_{m,h,i}^{(q)})^2\mid\mcF_{m,h,i-1}\bigr)
  \right\|_{\op},
  \qquad
  \nu_{m,h,\ast}:=
  \left\|
  \sum_i
  \mathbb E(R_{m,h,i}^2\mid\mcF_{m,h,i-1})
  \right\|_{\op}.
\]
Let
\[
  a_{m,h,q}>0,\qquad a_{m,h,\ast}>0,
\]
and assume
\[
  S_\infty:=
  \sup_{m,h}
  \left(
  \sum_{q\ge1}\sqrt{a_{m,h,q}}+\sqrt{a_{m,h,\ast}}
  \right)<\infty .
\]
Assume that the half-budget profile inequalities hold outside events
\(E_{m,h,q}\) and \(E_{m,h,\ast}\):
\[
  \nu_{m,h,q}\le
  \frac{a_{m,h,q}n_m}{2\log(e\kappa_{m,h,q})}
  \quad(q\ge1),
  \qquad
  \nu_{m,h,\ast}\le
  \frac{a_{m,h,\ast}n_m}{2\log(eD_m)}.
\]
If
\[
  \sum_m\sum_{h=1}^{J_m}
  \mathbb P\left(E_{m,h,\ast}\cup\bigcup_{q\ge1}E_{m,h,q}\right)<\infty,
\]
\[
  \sum_m\sum_{h=1}^{J_m}\sum_{q\ge1}
  \kappa_{m,h,q}
  \exp\left\{
  -\frac{3a_{m,h,q}n_m}{16\log(e\kappa_{m,h,q})}
  \right\}<\infty,
\]
and
\[
  \sum_m\sum_{h=1}^{J_m}
  D_m
  \exp\left\{
  -\frac{3a_{m,h,\ast}n_m}{16\log(eD_m)}
  \right\}<\infty,
\]
then, with probability one, for all sufficiently large \(m\), every active
summable-depth stream has a Spencer-scale signing:
\[
  \disc(A_{m,h,1},\ldots,A_{m,h,n_m})
  \le C S_\infty\sqrt{n_m}
  \qquad(1\le h\le J_m).
\]
Sampled certificates are eventual almost surely if
\[
  \sum_m\sum_{h=1}^{J_m}2^{-M_{m,h}}<\infty,
\]
or, for one common active-menu batch, if
\[
  \sum_mJ_m2^{-M_m}<\infty .
\]
If \(J_m\to\infty\), the active menus may exhaust a countable library of
summable-depth random-matrix models.
\end{corollary}

\begin{proof}
Apply
Corollary~\ref{cor:summable-multiscale-predictable-square-profile} to each
active pair \((m,h)\) with the half-budget choices, then union over
\(1\le h\le J_m\).  The three displayed summability assumptions make the
active-row failure probabilities summable.  Borel--Cantelli gives the eventual
simultaneous discrepancy bound.  The sampled-sign assertions add either the
independent-batch search series or the common-batch union bound
\(J_m2^{-M_m}\).
\end{proof}

\begin{corollary}[Flag-generated summable multiscale certificates]
\label{cor:flag-summable-multiscale-square-function-certificate}
There is a universal constant \(C\) with the following property.  Let
\(D\ge2\), let
\[
  0=P^{(0)}\preceq P^{(1)}\preceq P^{(2)}\preceq\cdots
\]
be an increasing sequence of orthogonal projections on \(\C^D\), and suppose
\[
  \operatorname{rank}P^{(q)}\le\kappa_q
  \qquad(q\ge1).
\]
Let \(P^{(\infty)}\) be the supremum projection.  For self-adjoint
\(A_1,\ldots,A_n\in\M_D(\C)\), define the flag shells and the residual by
\[
  \Delta_i^{(q)}
  :=
  P^{(q)}A_iP^{(q)}-P^{(q-1)}A_iP^{(q-1)},
  \qquad
  R_i:=A_i-P^{(\infty)}A_iP^{(\infty)} .
\]
If
\[
  \Theta_{\mathrm{flag}}:=
  \sum_{q\ge1}
  \sqrt{\log(e\kappa_q)}
  \left\|\sum_{i=1}^n(\Delta_i^{(q)})^2\right\|_{\op}^{1/2}
  +
  \sqrt{\log(eD)}
  \left\|\sum_{i=1}^nR_i^2\right\|_{\op}^{1/2}
  <\infty,
\]
then
\[
  \disc(A_1,\ldots,A_n)\le C\,\Theta_{\mathrm{flag}}.
\]
Moreover, if \(\eta^{(1)},\ldots,\eta^{(m)}\) are independent uniform sign
vectors, then
\[
  \mathbb P_\eta\left\{
  \min_{1\le r\le m}
  \left\|\sum_{i=1}^n\eta_i^{(r)}A_i\right\|_{\op}
  >2C\,\Theta_{\mathrm{flag}}
  \right\}
  \le2^{-m}.
\]
Consequently, whenever
\[
  \left\|\sum_i(\Delta_i^{(q)})^2\right\|_{\op}
  \le \frac{a_qn}{\log(e\kappa_q)}
  \quad(q\ge1),
  \qquad
  \left\|\sum_iR_i^2\right\|_{\op}
  \le\frac{a_\ast n}{\log(eD)}
\]
and \(\sum_q\sqrt{a_q}<\infty\), one has
\[
  \disc(A_1,\ldots,A_n)
  \le
  C\left(\sum_{q\ge1}\sqrt{a_q}+\sqrt{a_\ast}\right)\sqrt n,
\]
with sampled-search failure at most \(2^{-m}\) at twice that threshold.
\end{corollary}

\begin{proof}
For every \(Q\ge1\), the shells telescope:
\[
  P^{(Q)}A_iP^{(Q)}
  =
  \sum_{q=1}^Q\Delta_i^{(q)}.
\]
Since the projections are increasing in finite dimension, \(P^{(Q)}\) converges
in operator norm to \(P^{(\infty)}\).  Hence
\[
  A_i=R_i+\sum_{q\ge1}\Delta_i^{(q)}
\]
with operator-norm convergence.  Also
\(\Delta_i^{(q)}=P^{(q)}\Delta_i^{(q)}P^{(q)}\), so the \(q\)-th shell is
supported in a \(\kappa_q\)-dimensional corner.
Proposition~\ref{prop:summable-multiscale-square-function-certificate} applied to this
decomposition gives the discrepancy and sampled-search estimates.  The
square-profile specialization is immediate from the definition of
\(\Theta_{\mathrm{flag}}\).
\end{proof}

\begin{corollary}[Predictable flag-compressible square-profile streams]
\label{cor:predictable-flag-summable-square-profile}
Let \(D\ge2\), let \((\mcF_i)_{i=0}^n\) be a filtration, and let
\[
  0=P^{(0)}\preceq P^{(1)}\preceq P^{(2)}\preceq\cdots
\]
be \(\mcF_0\)-measurable orthogonal projections with
\(\operatorname{rank}P^{(q)}\le\kappa_q\).  Let
\(P^{(\infty)}=\sup_qP^{(q)}\).  Suppose
\(A_1,\ldots,A_n\in\M_D(\C)\) are self-adjoint contractions and
\(A_i\) is \(\mcF_i\)-measurable.  Define
\[
  \Delta_i^{(q)}
  :=
  P^{(q)}A_iP^{(q)}-P^{(q-1)}A_iP^{(q-1)},
  \qquad
  R_i:=A_i-P^{(\infty)}A_iP^{(\infty)} ,
\]
and the predictable shell profiles
\[
  \nu_q:=
  \left\|
  \sum_{i=1}^n
  \mathbb E\bigl((\Delta_i^{(q)})^2\mid\mcF_{i-1}\bigr)
  \right\|_{\op},
  \qquad
  \nu_\ast:=
  \left\|
  \sum_{i=1}^n
  \mathbb E(R_i^2\mid\mcF_{i-1})
  \right\|_{\op}.
\]
Let \(a_\ast>0\) and \(a_q>0\), \(q\ge1\), with
\[
  A_{\mathrm{flag}}:=\sum_{q\ge1}\sqrt{a_q}+\sqrt{a_\ast}<\infty .
\]
Whenever
\[
  v_q+t_q\le\frac{a_qn}{\log(e\kappa_q)}
  \quad(q\ge1),
  \qquad
  v_\ast+t_\ast\le\frac{a_\ast n}{\log(eD)},
\]
one has
\[
\begin{aligned}
  &\mathbb P\left\{
  \disc(A_1,\ldots,A_n)>C A_{\mathrm{flag}}\sqrt n
  \right\} \\
  &\quad\le
  \sum_{q\ge1}
  \left[
  \mathbb P\{\nu_q>v_q\}
  +
  \kappa_q
  \exp\left\{-\frac{t_q^2}{8(v_q+t_q/3)}\right\}
  \right] \\
  &\qquad\quad+
  \mathbb P\{\nu_\ast>v_\ast\}
  +
  D\exp\left\{-\frac{t_\ast^2}{8(v_\ast+t_\ast/3)}\right\}.
\end{aligned}
\]
If \(\eta^{(1)},\ldots,\eta^{(m)}\) are independent uniform sign vectors
independent of \(\mcF_n\), then the same estimate holds for sampled sign search
after adding \(2^{-m}\).

In particular, if the half-budget inequalities
\[
  \nu_q\le\frac{a_qn}{2\log(e\kappa_q)}
  \quad(q\ge1),
  \qquad
  \nu_\ast\le\frac{a_\ast n}{2\log(eD)}
\]
hold outside events \(E_q\) and \(E_\ast\), then the failure probability is at
most
\[
  \mathbb P\left(E_\ast\cup\bigcup_{q\ge1}E_q\right)
  +
  \sum_{q\ge1}
  \kappa_q
  \exp\left\{-\frac{3a_qn}{64\log(e\kappa_q)}\right\}
  +
  D\exp\left\{-\frac{3a_\ast n}{64\log(eD)}\right\},
\]
with an additional \(2^{-m}\) for sampled sign search.

Thus a predictable stream is Spencer-scale whenever its energy becomes
summable across a fixed or random initial flag: one pays \(\log(e\kappa_q)\)
for the \(q\)-th compressed shell and the ambient logarithm only for the part
outside the terminal flag subspace.
\end{corollary}

\begin{proof}
The shell and residual matrices are adapted, self-adjoint, and have operator
norm at most \(2\).  Fix a shell \(q\), and write \(Y_i=\Delta_i^{(q)}\),
\(B_i=Y_i^2\), and \(M_i=\mathbb E(B_i\mid\mcF_{i-1})\).  On the range of
\(P^{(q)}\), the martingale differences \(X_i=B_i-M_i\) satisfy
\[
  \lambda_{\max}(X_i)\le4,
  \qquad
  \sum_i\mathbb E(X_i^2\mid\mcF_{i-1})
  \preceq
  4\sum_iM_i,
\]
because \(0\preceq B_i\preceq4I\) implies \(B_i^2\preceq4B_i\).  Matrix
Freedman in dimension at most \(\kappa_q\) gives
\[
  \mathbb P\left\{
  \left\|\sum_i(\Delta_i^{(q)})^2\right\|_{\op}>v_q+t_q
  \ \text{and}\ \nu_q\le v_q
  \right\}
  \le
  \kappa_q
  \exp\left\{-\frac{t_q^2}{8(v_q+t_q/3)}\right\}.
\]
The same argument in the ambient dimension \(D\), with \(Y_i=R_i\), gives the
residual tail.  A countable union bound controls the complement of the event on
which all shell and residual square functions are bounded by \(v_q+t_q\) and
\(v_\ast+t_\ast\).  On that event,
Corollary~\ref{cor:flag-summable-multiscale-square-function-certificate} gives
the discrepancy bound.  Conditioning on the realized tuple and applying the
sampled-search part of that corollary adds \(2^{-m}\).  The half-budget form is
obtained by taking
\[
  v_q=t_q=\frac{a_qn}{2\log(e\kappa_q)}
  \quad(q\ge1),
  \qquad
  v_\ast=t_\ast=\frac{a_\ast n}{2\log(eD)}.
\]
\end{proof}

\begin{corollary}[Active menus of predictable flag-compressible streams]
\label{cor:active-menu-predictable-flag-summable-square-profile-schedules}
For each \(m\ge1\), let \(D_m\ge2\) and \(J_m\ge1\).  For every active
\((m,h)\), let
\[
  0=P_{m,h}^{(0)}\preceq P_{m,h}^{(1)}
  \preceq P_{m,h}^{(2)}\preceq\cdots
\]
be an \(\mcF_{m,h,0}\)-measurable flag in \(\M_{D_m}(\C)\), with
\[
  \operatorname{rank}P_{m,h}^{(q)}\le\kappa_{m,h,q}.
\]
Let \(A_{m,h,1},\ldots,A_{m,h,n_m}\) be adapted self-adjoint contractions, and
let \(\nu_{m,h,q}\) and \(\nu_{m,h,\ast}\) be the flag-shell and residual
predictable profiles from
Corollary~\ref{cor:predictable-flag-summable-square-profile}.  Let
\[
  a_{m,h,q}>0,\qquad a_{m,h,\ast}>0,
\]
and assume
\[
  S_{\mathrm{flag}}:=
  \sup_{m,h}
  \left(
  \sum_{q\ge1}\sqrt{a_{m,h,q}}+\sqrt{a_{m,h,\ast}}
  \right)<\infty .
\]
Assume that the half-budget profile inequalities hold outside events
\(E_{m,h,q}\) and \(E_{m,h,\ast}\):
\[
  \nu_{m,h,q}\le
  \frac{a_{m,h,q}n_m}{2\log(e\kappa_{m,h,q})}
  \quad(q\ge1),
  \qquad
  \nu_{m,h,\ast}\le
  \frac{a_{m,h,\ast}n_m}{2\log(eD_m)}.
\]
If
\[
  \sum_m\sum_{h=1}^{J_m}
  \mathbb P\left(E_{m,h,\ast}\cup\bigcup_{q\ge1}E_{m,h,q}\right)<\infty,
\]
\[
  \sum_m\sum_{h=1}^{J_m}\sum_{q\ge1}
  \kappa_{m,h,q}
  \exp\left\{
  -\frac{3a_{m,h,q}n_m}{64\log(e\kappa_{m,h,q})}
  \right\}<\infty,
\]
and
\[
  \sum_m\sum_{h=1}^{J_m}
  D_m
  \exp\left\{
  -\frac{3a_{m,h,\ast}n_m}{64\log(eD_m)}
  \right\}<\infty,
\]
then, with probability one, for all sufficiently large \(m\), every active
flag-compressible stream has a Spencer-scale signing:
\[
  \disc(A_{m,h,1},\ldots,A_{m,h,n_m})
  \le C S_{\mathrm{flag}}\sqrt{n_m}
  \qquad(1\le h\le J_m).
\]
Sampled certificates are eventual almost surely if
\[
  \sum_m\sum_{h=1}^{J_m}2^{-M_{m,h}}<\infty,
\]
or, for one common active-menu batch, if
\[
  \sum_mJ_m2^{-M_m}<\infty .
\]
If \(J_m\to\infty\), the active menus may exhaust a countable library of
flag-compressible random-matrix models.
\end{corollary}

\begin{proof}
Apply
Corollary~\ref{cor:predictable-flag-summable-square-profile} to each active
pair \((m,h)\) with the half-budget choices, then union over
\(1\le h\le J_m\).  The displayed summability assumptions make the active-row
failure probabilities summable, so Borel--Cantelli gives the eventual
simultaneous discrepancy bound.  The sampled-sign assertions add either the
independent-batch series or the common-batch union bound \(J_m2^{-M_m}\).
\end{proof}

\begin{corollary}[Compression-error flag certificates]
\label{cor:compression-error-flag-square-profile}
There is a universal constant \(C\) with the following property.  Let
\(D\ge2\), and let
\[
  0=P^{(0)}\preceq P^{(1)}\preceq P^{(2)}\preceq\cdots
\]
be an increasing sequence of orthogonal projections on \(\C^D\), with
\[
  \operatorname{rank}P^{(q)}\le\kappa_q
  \qquad(q\ge1).
\]
Let \(P^{(\infty)}=\sup_qP^{(q)}\).  Let
\(A_1,\ldots,A_n\in\M_D(\C)\) be self-adjoint contractions.  Suppose that
nonnegative numbers \(\tau_{i,q}\) and \(\tau_{i,\ast}\) satisfy
\[
  \left\|A_i-P^{(q)}A_iP^{(q)}\right\|_{\op}\le\tau_{i,q}
  \quad(q\ge0),
  \qquad
  \left\|A_i-P^{(\infty)}A_iP^{(\infty)}\right\|_{\op}\le\tau_{i,\ast}.
\]
Put
\[
  \Theta_{\mathrm{ce}}:=
  \sum_{q\ge1}\sqrt{\log(e\kappa_q)}
  \left(\sum_{i=1}^n(\tau_{i,q-1}+\tau_{i,q})^2\right)^{1/2}
  +
  \sqrt{\log(eD)}
  \left(\sum_{i=1}^n\tau_{i,\ast}^2\right)^{1/2}.
\]
If \(\Theta_{\mathrm{ce}}<\infty\), then
\[
  \disc(A_1,\ldots,A_n)\le C\,\Theta_{\mathrm{ce}},
\]
and \(m\) independent sampled sign vectors fail to find the threshold
\(2C\Theta_{\mathrm{ce}}\) with probability at most \(2^{-m}\).

In particular, if \(P^{(\infty)}=I\) and there are deterministic shell-error
levels \(b_q\) such that \(\tau_{i,q}\le b_q\) for every \(i,q\), then
\[
  \disc(A_1,\ldots,A_n)
  \le
  C\sqrt n
  \sum_{q\ge1}\sqrt{\log(e\kappa_q)}(b_{q-1}+b_q).
\]
Thus any tuple that is uniformly compressible along a common flag is
Spencer-scale whenever its weighted approximation-error increments are
summable.
\end{corollary}

\begin{proof}
Let
\[
  \Delta_i^{(q)}
  =
  P^{(q)}A_iP^{(q)}-P^{(q-1)}A_iP^{(q-1)},\qquad
  R_i=A_i-P^{(\infty)}A_iP^{(\infty)}.
\]
Then
\[
  \|\Delta_i^{(q)}\|_{\op}
  \le
  \left\|A_i-P^{(q-1)}A_iP^{(q-1)}\right\|_{\op}
  +
  \left\|A_i-P^{(q)}A_iP^{(q)}\right\|_{\op}
  \le \tau_{i,q-1}+\tau_{i,q}.
\]
Since \((\Delta_i^{(q)})^2\preceq\|\Delta_i^{(q)}\|_{\op}^2P^{(q)}\),
\[
  \left\|\sum_i(\Delta_i^{(q)})^2\right\|_{\op}^{1/2}
  \le
  \left(\sum_i(\tau_{i,q-1}+\tau_{i,q})^2\right)^{1/2}.
\]
Similarly,
\(\|\sum_iR_i^2\|_{\op}^{1/2}\le(\sum_i\tau_{i,\ast}^2)^{1/2}\).
Corollary~\ref{cor:flag-summable-multiscale-square-function-certificate}
therefore gives the discrepancy and sampled-search bounds.  The uniform
compressibility specialization substitutes \(\tau_{i,q}\le b_q\) and
\(\tau_{i,\ast}=0\).
\end{proof}

\begin{corollary}[Predictable compression-error flag streams]
\label{cor:predictable-compression-error-flag-square-profile}
Let \(D\ge2\), let \((\mcF_i)_{i=0}^n\) be a filtration, and let
\[
  0=P^{(0)}\preceq P^{(1)}\preceq P^{(2)}\preceq\cdots
\]
be an \(\mcF_0\)-measurable flag with
\(\operatorname{rank}P^{(q)}\le\kappa_q\).  Let
\(P^{(\infty)}=\sup_qP^{(q)}\).  Let
\(A_1,\ldots,A_n\) be adapted self-adjoint contractions.  Suppose that
nonnegative \(\mcF_{i-1}\)-measurable variables \(\tau_{i,q}\) and
\(\tau_{i,\ast}\) satisfy
\[
  \left\|A_i-P^{(q)}A_iP^{(q)}\right\|_{\op}\le\tau_{i,q}
  \quad(q\ge0),
  \qquad
  \left\|A_i-P^{(\infty)}A_iP^{(\infty)}\right\|_{\op}\le\tau_{i,\ast}.
\]
Define the predictable compression-error budgets
\[
  U_q:=\sum_{i=1}^n(\tau_{i,q-1}+\tau_{i,q})^2,
  \qquad
  U_\ast:=\sum_{i=1}^n\tau_{i,\ast}^2 .
\]
Let \(a_\ast>0\) and \(a_q>0\), \(q\ge1\), with
\[
  A_{\mathrm{ce}}:=\sum_{q\ge1}\sqrt{a_q}+\sqrt{a_\ast}<\infty .
\]
Assume that the half-budget inequalities
\[
  U_q\le\frac{a_qn}{2\log(e\kappa_q)}
  \quad(q\ge1),
  \qquad
  U_\ast\le\frac{a_\ast n}{2\log(eD)}
\]
hold outside events \(E_q\) and \(E_\ast\).  Then
\[
\begin{aligned}
  &\mathbb P\left\{
  \disc(A_1,\ldots,A_n)>C A_{\mathrm{ce}}\sqrt n
  \right\} \\
  &\quad\le
  \mathbb P\left(E_\ast\cup\bigcup_{q\ge1}E_q\right)
  +
  \sum_{q\ge1}
  \kappa_q
  \exp\left\{-\frac{3a_qn}{64\log(e\kappa_q)}\right\}
  +
  D\exp\left\{-\frac{3a_\ast n}{64\log(eD)}\right\}.
\end{aligned}
\]
If \(\eta^{(1)},\ldots,\eta^{(m)}\) are independent uniform sign vectors
independent of \(\mcF_n\), then the same bound holds for sampled sign search
after adding \(2^{-m}\).
\end{corollary}

\begin{proof}
Use the notation \(\Delta_i^{(q)}\) and \(R_i\) from the previous proof.  The
same triangle inequality gives
\(\|\Delta_i^{(q)}\|_{\op}\le\tau_{i,q-1}+\tau_{i,q}\), hence
\[
  \mathbb E\bigl((\Delta_i^{(q)})^2\mid\mcF_{i-1}\bigr)
  \preceq
  (\tau_{i,q-1}+\tau_{i,q})^2P^{(q)}.
\]
Therefore the shell profile \(\nu_q\) in
Corollary~\ref{cor:predictable-flag-summable-square-profile} is at most
\(U_q\), and the residual profile \(\nu_\ast\) is at most \(U_\ast\).  Applying
that corollary with
\[
  v_q=t_q=\frac{a_qn}{2\log(e\kappa_q)},\qquad
  v_\ast=t_\ast=\frac{a_\ast n}{2\log(eD)}
\]
gives the displayed estimate and the sampled-search variant.
\end{proof}

\begin{corollary}[Active menus of predictable compression-error flag streams]
\label{cor:active-menu-predictable-compression-error-flag-schedules}
For each \(m\ge1\), let \(D_m\ge2\) and \(J_m\ge1\).  For every active
\((m,h)\), let \(A_{m,h,1},\ldots,A_{m,h,n_m}\) be an adapted self-adjoint
contraction stream with an initial flag
\[
  0=P_{m,h}^{(0)}\preceq P_{m,h}^{(1)}
  \preceq P_{m,h}^{(2)}\preceq\cdots ,
  \qquad
  \operatorname{rank}P_{m,h}^{(q)}\le\kappa_{m,h,q}.
\]
Let \(U_{m,h,q}\) and \(U_{m,h,\ast}\) be predictable compression-error budgets
as in Corollary~\ref{cor:predictable-compression-error-flag-square-profile}.
Let \(a_{m,h,q}>0\) and \(a_{m,h,\ast}>0\), and assume
\[
  S_{\mathrm{ce}}:=
  \sup_{m,h}
  \left(
  \sum_{q\ge1}\sqrt{a_{m,h,q}}+\sqrt{a_{m,h,\ast}}
  \right)<\infty .
\]
Suppose the half-budget inequalities
\[
  U_{m,h,q}\le
  \frac{a_{m,h,q}n_m}{2\log(e\kappa_{m,h,q})}
  \quad(q\ge1),
  \qquad
  U_{m,h,\ast}\le
  \frac{a_{m,h,\ast}n_m}{2\log(eD_m)}
\]
hold outside events \(E_{m,h,q}\) and \(E_{m,h,\ast}\).  If
\[
  \sum_m\sum_{h=1}^{J_m}
  \mathbb P\left(E_{m,h,\ast}\cup\bigcup_{q\ge1}E_{m,h,q}\right)<\infty,
\]
\[
  \sum_m\sum_{h=1}^{J_m}\sum_{q\ge1}
  \kappa_{m,h,q}
  \exp\left\{
  -\frac{3a_{m,h,q}n_m}{64\log(e\kappa_{m,h,q})}
  \right\}<\infty,
\]
and
\[
  \sum_m\sum_{h=1}^{J_m}
  D_m
  \exp\left\{
  -\frac{3a_{m,h,\ast}n_m}{64\log(eD_m)}
  \right\}<\infty,
\]
then, with probability one, for all sufficiently large \(m\), every active
compression-error flag stream has a Spencer-scale signing:
\[
  \disc(A_{m,h,1},\ldots,A_{m,h,n_m})
  \le C S_{\mathrm{ce}}\sqrt{n_m}
  \qquad(1\le h\le J_m).
\]
Sampled certificates are eventual almost surely if
\[
  \sum_m\sum_{h=1}^{J_m}2^{-M_{m,h}}<\infty,
\]
or, for one common active-menu batch, if
\[
  \sum_mJ_m2^{-M_m}<\infty .
\]
\end{corollary}

\begin{proof}
Apply
Corollary~\ref{cor:predictable-compression-error-flag-square-profile} to each
active pair \((m,h)\) and union over \(1\le h\le J_m\).  The displayed
summability assumptions make the active-row failure probabilities summable, so
Borel--Cantelli gives the simultaneous eventual bound.  The sampled-sign
assertions add the independent-batch or common-batch search series.
\end{proof}

\begin{corollary}[Common-flag tail envelopes]
\label{cor:common-flag-tail-envelope-square-profile}
There is a universal constant \(C\) with the following property.  Let
\(D\ge2\), let
\[
  0=P^{(0)}\preceq P^{(1)}\preceq P^{(2)}\preceq\cdots
\]
be an increasing flag on \(\C^D\), and suppose
\(\operatorname{rank}P^{(q)}\le\kappa_q\).  Let
\(P^{(\infty)}=\sup_qP^{(q)}\).  Let \(b_q\ge0\), and put
\[
  B_{\mathrm{flag}}(b,\kappa):=
  \sum_{q\ge1}
  \sqrt{\log(e\kappa_q)}(b_{q-1}+b_q).
\]
Let \(A_1,\ldots,A_n\in\M_D(\C)\) be self-adjoint contractions.  Suppose that
there are nonnegative amplitudes \(L_i\) and residual amplitudes \(r_i\) such
that
\[
  \left\|A_i-P^{(q)}A_iP^{(q)}\right\|_{\op}\le L_i b_q
  \quad(q\ge0),
  \qquad
  \left\|A_i-P^{(\infty)}A_iP^{(\infty)}\right\|_{\op}\le r_i.
\]
If \(B_{\mathrm{flag}}(b,\kappa)<\infty\), then
\[
  \disc(A_1,\ldots,A_n)
  \le
  C\left[
  B_{\mathrm{flag}}(b,\kappa)
  \left(\sum_iL_i^2\right)^{1/2}
  +
  \sqrt{\log(eD)}
  \left(\sum_ir_i^2\right)^{1/2}
  \right].
\]
The same threshold, enlarged by a factor \(2\), is found by \(m\) independent
sampled sign vectors with failure probability at most \(2^{-m}\).

Consequently, if
\[
  \sum_iL_i^2\le L^2 n,
  \qquad
  \sum_ir_i^2\le \frac{R^2n}{\log(eD)},
\]
then
\[
  \disc(A_1,\ldots,A_n)
  \le
  C\bigl(LB_{\mathrm{flag}}(b,\kappa)+R\bigr)\sqrt n.
\]
In particular the flag sum is finite in each of the following common regimes:
\[
  \kappa_q\le\exp(\gamma q^s),\qquad b_q\le M(1+q)^{-p},
  \qquad p>1+\frac{s}{2},
\]
and
\[
  \kappa_q\le K_0\rho^q,\qquad b_q\le M\rho^{-\alpha q},
  \qquad \rho>1,\ \alpha>0.
\]
Thus polynomial spectral-window tails beat stretched-exponential flag growth
when \(p>1+s/2\), and any positive power tail in geometrically growing rank
windows gives a Spencer-scale common-flag model.
\end{corollary}

\begin{proof}
Apply Corollary~\ref{cor:compression-error-flag-square-profile} with
\[
  \tau_{i,q}=L_i b_q,\qquad \tau_{i,\ast}=r_i .
\]
Then
\[
  \left(\sum_i(\tau_{i,q-1}+\tau_{i,q})^2\right)^{1/2}
  =
  (b_{q-1}+b_q)\left(\sum_iL_i^2\right)^{1/2},
\]
and the displayed bound follows.  The sampled-search statement is inherited
from the same corollary.

For the polynomial example,
\(\sqrt{\log(e\kappa_q)}\le C_{\gamma,s}(1+q)^{s/2}\), so the flag sum is
dominated by a constant multiple of
\(\sum_q(1+q)^{s/2-p}\), which converges exactly under \(p>1+s/2\).  For the
geometric rank/power-tail example,
\(\sqrt{\log(e\kappa_q)}\le C(1+q)^{1/2}\), while
\(b_{q-1}+b_q\le C_{\rho,\alpha,M}\rho^{-\alpha q}\); the resulting series is
geometric up to a polynomial factor and therefore converges.
\end{proof}

\begin{corollary}[Active menus of common-flag tail-envelope streams]
\label{cor:active-menu-common-flag-tail-envelope-schedules}
For each \(m\ge1\), let \(D_m\ge2\) and \(J_m\ge1\).  For every active
\((m,h)\), let
\[
  0=P_{m,h}^{(0)}\preceq P_{m,h}^{(1)}
  \preceq P_{m,h}^{(2)}\preceq\cdots
\]
be a flag in \(\M_{D_m}(\C)\), with
\(\operatorname{rank}P_{m,h}^{(q)}\le\kappa_{m,h,q}\).  Let
\(P_{m,h}^{(\infty)}=\sup_qP_{m,h}^{(q)}\), and let
\(b_{m,h,q}\ge0\).  Put
\[
  B_{m,h}^{\mathrm{flag}}:=
  \sum_{q\ge1}
  \sqrt{\log(e\kappa_{m,h,q})}
  (b_{m,h,q-1}+b_{m,h,q}).
\]
Suppose that outside an event \(E_{m,h}\), the active stream
\(A_{m,h,1},\ldots,A_{m,h,n_m}\) consists of self-adjoint contractions and
there are amplitudes \(L_{m,h,i}\) and \(r_{m,h,i}\) such that
\[
  \left\|
  A_{m,h,i}-P_{m,h}^{(q)}A_{m,h,i}P_{m,h}^{(q)}
  \right\|_{\op}
  \le L_{m,h,i}b_{m,h,q}
  \quad(q\ge0),
\]
\[
  \left\|
  A_{m,h,i}
  -P_{m,h}^{(\infty)}A_{m,h,i}P_{m,h}^{(\infty)}
  \right\|_{\op}
  \le r_{m,h,i},
\]
and
\[
  B_{m,h}^{\mathrm{flag}}
  \left(n_m^{-1}\sum_iL_{m,h,i}^2\right)^{1/2}
  +
  \sqrt{\log(eD_m)}
  \left(n_m^{-1}\sum_ir_{m,h,i}^2\right)^{1/2}
  \le S .
\]
If
\[
  \sum_m\sum_{h=1}^{J_m}\mathbb P(E_{m,h})<\infty,
\]
then, with probability one, for all sufficiently large \(m\), every active
common-flag tail-envelope stream has a Spencer-scale signing:
\[
  \disc(A_{m,h,1},\ldots,A_{m,h,n_m})
  \le C S\sqrt{n_m}
  \qquad(1\le h\le J_m).
\]
Sampled certificates are eventual almost surely if
\[
  \sum_m\sum_{h=1}^{J_m}2^{-M_{m,h}}<\infty,
\]
or, for one common active-menu batch, if
\[
  \sum_mJ_m2^{-M_m}<\infty .
\]
If \(J_m\to\infty\), the active prefixes may exhaust a countable library of
common-flag spectral-truncation, wavelet-tail, or low-rank-approximation
random-matrix models.  The polynomial and geometric tail regimes in
Corollary~\ref{cor:common-flag-tail-envelope-square-profile} give checkable
uniform choices of \(S\).
\end{corollary}

\begin{proof}
On \(E_{m,h}^c\), Corollary~\ref{cor:common-flag-tail-envelope-square-profile}
gives
\[
  \disc(A_{m,h,1},\ldots,A_{m,h,n_m})\le CS\sqrt{n_m}
\]
and the corresponding sampled-search statement.  The assumed summability of
the active bad-event probabilities, together with Borel--Cantelli, gives the
eventual simultaneous deterministic-signing conclusion.  The sampled-sign
assertions add the independent-batch or common-batch search series.
\end{proof}

\begin{corollary}[Rank-sensitive dominated square-twirl ensembles]
\label{cor:rank-sensitive-dominated-square-twirl}
Let \(D\ge2\), let \(1\le\kappa\le D\), and let \(\mcG\) be an
environment.  Let \(P\) be a \(\mcG\)-measurable orthogonal projection with
\[
  1\le \operatorname{rank}P\le\kappa .
\]
Let \(A_1,\ldots,A_n\in\M_D(\C)\) be conditionally independent self-adjoint
contractions over \(\mcG\) such that
\[
  A_i=PA_iP .
\]
Suppose that there are \(\mcG\)-measurable positive semidefinite matrices
\(V_1,\ldots,V_n\) satisfying
\[
  \mathbb E(A_i^2\mid\mcG)\preceq V_i \qquad(1\le i\le n),
\]
and put
\[
  \Xi_\kappa(\mcG):=
  \left\|\sum_{i=1}^nV_i\right\|_{\op}.
\]
Then, for every \(C_0\ge1\), there is \(C(C_0)\) such that, whenever
\[
  v+t\le \frac{C_0n}{\log(e\kappa)},
\]
one has
\[
  \mathbb P\left\{
  \disc(A_1,\ldots,A_n)>C(C_0)\sqrt n
  \right\}
  \le
  \mathbb P\{\Xi_\kappa(\mcG)>v\}
  +
  \kappa\exp\left\{-\frac{t^2}{2(v+t/3)}\right\}.
\]
If \(\eta^{(1)},\ldots,\eta^{(m)}\) are independent uniform sign vectors,
independent of the tuple, then the same bound holds for sampled sign search
after adding \(2^{-m}\).  In particular, if
\[
  \Xi_\kappa(\mcG)\le \frac{C_0n}{2\log(e\kappa)}
\]
outside an event \(E\), then
\[
  \mathbb P\left\{
  \disc(A_1,\ldots,A_n)>C(C_0)\sqrt n
  \right\}
  \le
  \mathbb P(E)+
  \kappa\exp\left\{-\frac{3C_0n}{16\log(e\kappa)}\right\},
\]
with the additional \(2^{-m}\) term for sampled sign search.

Equivalently, the ambient dimension is irrelevant for random-basis ensembles
that live inside a common random \(\kappa\)-dimensional corner.  For example,
let \(T_i=PT_iP\) be \(\mcG\)-measurable self-adjoint contractions, and,
conditional on \(\mcG\), let \(U_i\) be independent Haar unitaries, exact
finite unitary \(1\)-designs, or signed-permutation/monomial designs on
\(\operatorname{ran}P\), extended by the identity on \(P^\perp\).  With
\(A_i=U_iT_iU_i^*\) and \(r=\operatorname{rank}P\),
\[
  \mathbb E(A_i^2\mid\mcG)=\frac{\operatorname{tr}(T_i^2)}{r}P,
\]
so the theorem applies with
\[
  \Xi_\kappa(\mcG)=\frac1r\sum_{i=1}^n\operatorname{tr}(T_i^2).
\]
More generally, approximate scalar designs in the corner satisfying, for
\(0\preceq X\preceq P\),
\[
  \mathbb E(U_iXU_i^*\mid\mcG)
  \preceq
  \frac{\operatorname{tr}X}{r}P+\Gamma_i
\]
are covered with
\[
  V_i=\frac{\operatorname{tr}(T_i^2)}{r}P+\Gamma_i ,
\]
and block designs inside \(P=\sum_{\alpha\in I_i}P_{i,\alpha}\) are covered by
replacing the scalar term by
\[
  \sum_{\alpha\in I_i}
  \frac{\operatorname{tr}(P_{i,\alpha}T_i^2P_{i,\alpha})}
       {\operatorname{rank}P_{i,\alpha}}
  P_{i,\alpha}
\]
plus any positive semidefinite block-twirl error.  Thus random low-rank Haar
subspace rotations, finite-design rotations, signed-permutation rotations, and
their dominated approximate versions are Spencer-scale under the
\(\kappa\)-dimensional square profile \(O(n/\log(e\kappa))\), rather than the
ambient \(D\)-dimensional profile.
\end{corollary}

\begin{proof}
Let
\[
  \mcF_i:=\sigma(\mcG,A_1,\ldots,A_i)\qquad(0\le i\le n),
\]
with \(\mcF_0=\mcG\).  Conditional independence over \(\mcG\) gives
\[
  \mathbb E(A_i^2\mid\mcF_{i-1})=\mathbb E(A_i^2\mid\mcG)\preceq V_i .
\]
Therefore the predictable square profile \(\nu_\kappa\) in
Corollary~\ref{cor:rank-sensitive-predictable-square-profile} is at most
\(\Xi_\kappa(\mcG)\).  Applying that corollary proves the discrepancy and
sampled-search estimates, and the half-budget form follows by taking
\(v=t=C_0n/(2\log(e\kappa))\).

For the scalar twirl specializations, the Haar, finite \(1\)-design, and
signed-permutation laws on \(\operatorname{ran}P\) satisfy
\[
  \mathbb E(U_iXU_i^*\mid\mcG)=\frac{\operatorname{tr}X}{r}P
  \qquad (X=PXP),
\]
so substituting \(X=T_i^2\) gives the displayed profile.  The approximate
scalar and block forms follow in the same way from the assumed Loewner
domination, again with \(X=T_i^2\).
\end{proof}

\begin{corollary}[Finite-menu rank-sensitive dominated square-twirl schedules]
\label{cor:finite-menu-rank-sensitive-dominated-square-twirl-schedules}
Fix \(C_0\ge1\) and put \(\gamma_0:=3C_0/16\).  For each \(m\ge1\), let
\(D_m\ge2\), let \(J_m\ge1\), and, for every \(1\le h\le J_m\), let
\[
  1\le\kappa_{m,h}\le D_m,\qquad
  \ell_{m,h}:=\log(e\kappa_{m,h}).
\]
Let \(\mcG_{m,h}\) be an environment, and let \(P_{m,h}\) be a
\(\mcG_{m,h}\)-measurable orthogonal projection with
\[
  1\le \operatorname{rank}P_{m,h}\le\kappa_{m,h}.
\]
Let
\[
  A_{m,h,1},\ldots,A_{m,h,n_m}\in\M_{D_m}(\C)
\]
be conditionally independent self-adjoint contractions over \(\mcG_{m,h}\)
such that
\[
  A_{m,h,i}=P_{m,h}A_{m,h,i}P_{m,h}.
\]
Suppose that \(\mcG_{m,h}\)-measurable positive semidefinite matrices
\(V_{m,h,1},\ldots,V_{m,h,n_m}\) satisfy
\[
  \mathbb E(A_{m,h,i}^2\mid\mcG_{m,h})\preceq V_{m,h,i},
\]
and put
\[
  \Xi_{m,h}^{\kappa}:=
  \left\|\sum_{i=1}^{n_m}V_{m,h,i}\right\|_{\op}.
\]
Assume that
\[
  \Xi_{m,h}^{\kappa}\le\frac{C_0n_m}{2\ell_{m,h}}
  \qquad(1\le h\le J_m)
\]
outside events \(E_{m,h}\), and that
\[
  \sum_m\sum_{h=1}^{J_m}\mathbb P(E_{m,h})<\infty,
  \qquad
  \sum_m\sum_{h=1}^{J_m}
  \kappa_{m,h}\exp\left\{-\gamma_0\frac{n_m}{\ell_{m,h}}\right\}<\infty .
\]
Then, with probability one, for all sufficiently large \(m\), every active
rank-sensitive dominated-twirl stream has a Spencer-scale signing:
\[
  \disc(A_{m,h,1},\ldots,A_{m,h,n_m})
  \le C(C_0)\sqrt{n_m}
  \qquad(1\le h\le J_m).
\]
If the streams are supplied with independent sampled sign batches and
\[
  \sum_m\sum_{h=1}^{J_m}2^{-M_{m,h}}<\infty,
\]
then, after enlarging \(C(C_0)\), almost surely every active stream has a
sampled signing certificate for all sufficiently large \(m\).  If one common
batch of \(M_m\) signs is reused across the active menu and is independent of
all active streams, it is enough that
\[
  \sum_mJ_m2^{-M_m}<\infty .
\]

The concrete rank-sensitive schedule from
Corollary~\ref{cor:finite-menu-rank-sensitive-square-profile-schedules}
applies verbatim.  Namely, with
\[
  \Lambda_m^\kappa:=\sum_{h=1}^{J_m}\kappa_{m,h},
  \qquad
  \ell_m^\kappa:=\max_{1\le h\le J_m}\ell_{m,h},
\]
the structural tail is summable whenever, for some \(\delta>0\),
\[
  n_m\ge
  \gamma_0^{-1}\ell_m^\kappa
  \left(\log(e\Lambda_m^\kappa)+(1+\delta)\log(e+m)\right)
\]
eventually.  With a common sign batch, the sampled conclusion follows if
\[
  M_m\ge \log_2(eJ_m)+(1+\delta')\log_2(e+m)
\]
eventually for some \(\delta'>0\).  If the streams are defined for all
\(h\ge1\), \(J_m\to\infty\), and \(h\le J_m\) is a deterministic exhaustion,
then every fixed low-rank dominated-twirl model in the countable library is
covered for all sufficiently large \(m\).

In particular, the statement applies to active libraries of low-rank Haar,
finite-design, signed-permutation, block-design, and approximate random-basis
ensembles.  In the scalar corner-twirl case \(A_{m,h,i}=U_{m,h,i}
T_{m,h,i}U_{m,h,i}^*\), with \(T_{m,h,i}=P_{m,h}T_{m,h,i}P_{m,h}\) and
\(r_{m,h}=\operatorname{rank}P_{m,h}\), the half-budget exception may be taken
as
\[
  \frac1{r_{m,h}}\sum_{i=1}^{n_m}\operatorname{tr}(T_{m,h,i}^2)
  >
  \frac{C_0n_m}{2\ell_{m,h}},
\]
with the corresponding block-energy or aggregate PSD-error dominated profile
for block and approximate designs.
\end{corollary}

\begin{proof}
For a fixed active pair \((m,h)\), apply
Corollary~\ref{cor:rank-sensitive-dominated-square-twirl} with
\[
  v=t=\frac{C_0n_m}{2\ell_{m,h}} .
\]
This bounds its discrepancy failure by
\[
  \mathbb P(E_{m,h})+
  \kappa_{m,h}\exp\left\{-\gamma_0\frac{n_m}{\ell_{m,h}}\right\},
\]
and adds \(2^{-M_{m,h}}\), or \(J_m2^{-M_m}\) after union over the active menu,
for sampled sign search.  The assumed summability and Borel--Cantelli prove
the almost-sure simultaneous conclusions.  The displayed concrete schedule is
the same estimate as in
Corollary~\ref{cor:finite-menu-rank-sensitive-square-profile-schedules}:
\[
  \sum_{h=1}^{J_m}
  \kappa_{m,h}\exp\left\{-\gamma_0\frac{n_m}{\ell_{m,h}}\right\}
  \le
  \Lambda_m^\kappa
  \exp\left\{-\gamma_0\frac{n_m}{\ell_m^\kappa}\right\}.
\]
The low-rank random-basis examples are exactly the specialization recorded in
Corollary~\ref{cor:rank-sensitive-dominated-square-twirl}.
\end{proof}

\begin{corollary}[Finite-menu rank-sensitive square-profile schedules]
\label{cor:finite-menu-rank-sensitive-square-profile-schedules}
Fix \(C_0\ge1\) and put \(\gamma_0:=3C_0/16\).  For each \(m\ge1\), let
\(D_m\ge2\), let \(J_m\ge1\), and, for every \(1\le h\le J_m\), let
\[
  1\le\kappa_{m,h}\le D_m,\qquad
  \ell_{m,h}:=\log(e\kappa_{m,h}).
\]
Let \((\mcF_{m,h,i})_{i=0}^{n_m}\) be a filtration, let
\(P_{m,h}\) be an \(\mcF_{m,h,0}\)-measurable orthogonal projection with
\[
  \operatorname{rank}P_{m,h}\le\kappa_{m,h},
\]
and let
\[
  A_{m,h,1},\ldots,A_{m,h,n_m}\in\M_{D_m}(\C)
\]
be self-adjoint contractions with \(A_{m,h,i}\) \(\mcF_{m,h,i}\)-measurable and
\[
  A_{m,h,i}=P_{m,h}A_{m,h,i}P_{m,h}.
\]
Put
\[
  \nu_{m,h}^{\kappa}:=
  \left\|
  \sum_{i=1}^{n_m}
  \mathbb E(A_{m,h,i}^2\mid\mcF_{m,h,i-1})
  \right\|_{\op}.
\]
Assume that
\[
  \nu_{m,h}^{\kappa}\le \frac{C_0n_m}{2\ell_{m,h}}
  \qquad(1\le h\le J_m)
\]
outside events \(E_{m,h}\), and that
\[
  \sum_m\sum_{h=1}^{J_m}\mathbb P(E_{m,h})<\infty,
  \qquad
  \sum_m\sum_{h=1}^{J_m}
  \kappa_{m,h}\exp\left\{-\gamma_0\frac{n_m}{\ell_{m,h}}\right\}<\infty .
\]
Then, with probability one, for all sufficiently large \(m\), every active
rank-sensitive stream has a Spencer-scale signing:
\[
  \disc(A_{m,h,1},\ldots,A_{m,h,n_m})
  \le C(C_0)\sqrt{n_m}
  \qquad(1\le h\le J_m).
\]
If the streams are supplied with independent sampled sign batches and
\[
  \sum_m\sum_{h=1}^{J_m}2^{-M_{m,h}}<\infty,
\]
then, after enlarging \(C(C_0)\), almost surely every active stream has a
sampled certificate for all sufficiently large \(m\).  If one common batch of
\(M_m\) signs is reused across the active menu and is independent of all active
streams, it is enough that
\[
  \sum_mJ_m2^{-M_m}<\infty .
\]

In particular, define
\[
  \Lambda_m^\kappa:=\sum_{h=1}^{J_m}\kappa_{m,h},
  \qquad
  \ell_m^\kappa:=\max_{1\le h\le J_m}\ell_{m,h}.
\]
The structural tail is summable whenever, for some \(\delta>0\),
\[
  n_m\ge
  \gamma_0^{-1}\ell_m^\kappa
  \left(\log(e\Lambda_m^\kappa)+(1+\delta)\log(e+m)\right)
\]
eventually, and the exceptional events are summable.  With a common sign batch,
the sampled conclusion follows if
\[
  M_m\ge \log_2(eJ_m)+(1+\delta')\log_2(e+m)
\]
eventually for some \(\delta'>0\).  If the streams are defined for all
\(h\ge1\), \(J_m\to\infty\), and the active menu is a deterministic exhaustion,
then every fixed rank-sensitive common-corner stream is covered for all
sufficiently large \(m\).
\end{corollary}

\begin{proof}
For a fixed active pair \((m,h)\), apply
Corollary~\ref{cor:rank-sensitive-predictable-square-profile} with
\[
  v=t=\frac{C_0n_m}{2\ell_{m,h}} .
\]
This gives failure probability at most
\[
  \mathbb P(E_{m,h})+
  \kappa_{m,h}\exp\left\{-\gamma_0\frac{n_m}{\ell_{m,h}}\right\}.
\]
Summing over active \(h\) and over \(m\), then applying the first
Borel--Cantelli lemma, proves the simultaneous almost-sure signing conclusion.
The sampled-sign assertions add the displayed \(2^{-M_{m,h}}\) failures, or the
common-batch union bound \(J_m2^{-M_m}\).

For the concrete schedule, use \(\ell_{m,h}\le\ell_m^\kappa\) and
\[
  \sum_{h=1}^{J_m}
  \kappa_{m,h}\exp\left\{-\gamma_0\frac{n_m}{\ell_{m,h}}\right\}
  \le
  \Lambda_m^\kappa
  \exp\left\{-\gamma_0\frac{n_m}{\ell_m^\kappa}\right\}
  \le e^{-1}(e+m)^{-1-\delta}.
\]
The common sampled-batch estimate is identical.  If \(J_m\to\infty\), each fixed
countable-menu element belongs to the active prefix eventually.
\end{proof}

\begin{corollary}[Finite-menu predictable square-profile schedules]
\label{cor:finite-menu-predictable-square-profile-schedules}
Fix \(C_0\ge1\) and put
\[
  \gamma_0:=\frac{3C_0}{16}.
\]
For each \(m\ge1\), let \(D_m\ge2\), put
\[
  \ell_m:=\log(eD_m),
\]
and let \(J_m\ge1\) be a deterministic active-menu size.  For every
\(1\le h\le J_m\), let \((\mcF_{m,h,i})_{i=0}^{n_m}\) be a filtration and let
\[
  A_{m,h,1},\ldots,A_{m,h,n_m}\in\M_{D_m}(\C)
\]
be self-adjoint contractions with \(A_{m,h,i}\) \(\mcF_{m,h,i}\)-measurable.
Put
\[
  \nu_{m,h}:=
  \left\|
  \sum_{i=1}^{n_m}
  \mathbb E(A_{m,h,i}^2\mid\mcF_{m,h,i-1})
  \right\|_{\op},
  \qquad
  L_m^{\rm sq}:=J_mD_m .
\]
Assume that
\[
  \nu_{m,h}\le \frac{C_0n_m}{2\ell_m}
  \qquad(1\le h\le J_m)
\]
outside events \(E_{m,h}\), and that
\[
  \sum_m\sum_{h=1}^{J_m}\mathbb P(E_{m,h})<\infty,
  \qquad
  \sum_m L_m^{\rm sq}
  \exp\left\{-\gamma_0\frac{n_m}{\ell_m}\right\}<\infty .
\]
Then, with probability one, for all sufficiently large \(m\), every active
adapted stream has a Spencer-scale signing:
\[
  \disc(A_{m,h,1},\ldots,A_{m,h,n_m})
  \le C(C_0)\sqrt{n_m}
  \qquad(1\le h\le J_m).
\]

If the menu elements are supplied with independent sampled sign batches and
\[
  \sum_m\sum_{h=1}^{J_m}2^{-M_{m,h}}<\infty,
\]
then, after enlarging \(C(C_0)\), almost surely every active menu element has a
sampled certificate for all sufficiently large \(m\).  If instead one common
batch of \(M_m\) sign vectors is reused for the entire active menu and is
independent of all active streams, it is enough that
\[
  \sum_mJ_m2^{-M_m}<\infty .
\]

In particular, under the displayed half-budget exception hypothesis, the
structural tail is summable in each of the following schedules.
\begin{enumerate}[label=\textup{(\alph*)}]
\item General schedules.  If, for some \(\delta>0\),
  \[
    n_m\ge
    \gamma_0^{-1}\ell_m
    \bigl(\log(eL_m^{\rm sq})+(1+\delta)\log(e+m)\bigr)
  \]
  eventually, then the simultaneous active-menu conclusion holds.  With a
  common sign batch, the sampled conclusion follows if
  \[
    M_m\ge \log_2(eJ_m)+(1+\delta')\log_2(e+m)
  \]
  eventually for some \(\delta'>0\).
\item Full integer dimensions.  Suppose the rows are indexed by
  \(D=2,3,\ldots\), with \(D_m=D\), and
  \(L_D^{\rm sq}\le(eD)^\Gamma\) eventually for some \(\Gamma\ge0\).  If
  \[
    n_D\ge a\log^2(eD)
  \]
  eventually for some
  \[
    a>\frac{16(\Gamma+1)}{3C_0},
  \]
  then the simultaneous active-menu conclusion holds along all integer
  dimensions.
\item Sparse stretched-exponential dimensions.  If there are \(c,\kappa>0\)
  such that \(D_m\ge\exp(c m^\kappa)\) eventually,
  \(L_m^{\rm sq}\le(eD_m)^\Gamma\) eventually, and
  \[
    n_m\ge a\log^2(eD_m)
  \]
  eventually for some
  \[
    a>\frac{16\Gamma}{3C_0},
  \]
  then the simultaneous active-menu conclusion holds along the sparse schedule.
\end{enumerate}

If the streams are defined for all \(h\ge1\), \(J_m\to\infty\), and the active
menu is a deterministic exhaustion of a countable library, then every fixed
adapted square-profile model is covered for all sufficiently large \(m\).
This applies directly to finite or countable libraries of predictable twirl,
predictable scalar-multiplier, predictable localized, frame-localized, and
other adapted random-matrix streams once their predictable square profiles obey
the displayed half-budget condition.
\end{corollary}

\begin{proof}
For fixed \(m\) and \(h\le J_m\), the half-budget specialization of
Corollary~\ref{cor:random-square-profile-freedman} gives failure probability at
most
\[
  \mathbb P(E_{m,h})
  +
  D_m\exp\left\{-\gamma_0\frac{n_m}{\ell_m}\right\}.
\]
Union over \(1\le h\le J_m\), then apply the first Borel--Cantelli lemma.  The
sampled-sign assertions add the corresponding \(2^{-M_{m,h}}\) failures, or the
common-batch union bound \(J_m2^{-M_m}\).

For the general schedule,
\[
  L_m^{\rm sq}
  \exp\left\{-\gamma_0\frac{n_m}{\ell_m}\right\}
  \le e^{-1}(e+m)^{-(1+\delta)}.
\]
The common sampled-batch condition gives
\(J_m2^{-M_m}\le e^{-1}(e+m)^{-(1+\delta')}\).  In the all-integer schedule,
\[
  L_D^{\rm sq}
  \exp\left\{-\gamma_0\frac{n_D}{\log(eD)}\right\}
  \le
  (eD)^{\Gamma-\gamma_0a},
\]
which is summable over \(D\) when \(\gamma_0a>\Gamma+1\).  In the sparse
schedule the same estimate is bounded by
\(\exp\{-(\gamma_0a-\Gamma)\log(eD_m)\}\), which is summable under the
stretched-exponential lower bound on \(D_m\) when \(\gamma_0a>\Gamma\).  If
\(J_m\to\infty\), then every fixed countable-menu element lies in the active
menu eventually.
\end{proof}

\begin{corollary}[Predictable dominated twirl streams]
\label{cor:predictable-dominated-twirl-freedman}
Let \(D\ge2\), let \((\mcF_i)_{i=0}^n\) be a filtration, and let
\(T_1,\ldots,T_n\in\M_D(\C)\) be self-adjoint contractions such that \(T_i\) is
\(\mcF_{i-1}\)-measurable.  Let \(U_i\) be an \(\mcF_i\)-measurable random
unitary, set
\[
  A_i=U_iT_iU_i^* ,
\]
and assume that there are \(\mcF_{i-1}\)-measurable positive semidefinite
matrices \(V_i\) such that
\[
  \mathbb E\left(U_iT_i^2U_i^*\,\middle|\,\mcF_{i-1}\right)\preceq V_i
  \qquad(1\le i\le n).
\]
Put
\[
  \Xi=\left\|\sum_{i=1}^nV_i\right\|_{\op}.
\]
Then, for every \(C_0\ge1\), there is \(C(C_0)\) such that, for every
\(v,t>0\) satisfying
\[
  v+t\le \frac{C_0n}{\log(eD)},
\]
one has
\[
  \mathbb P\left\{
  \disc(A_1,\ldots,A_n)>C(C_0)\sqrt n
  \right\}
  \le
  \mathbb P\{\Xi>v\}
  +
  D\exp\left\{-\frac{t^2}{2(v+t/3)}\right\}.
\]
Moreover, if \(\eta^{(1)},\ldots,\eta^{(m)}\) are independent uniform sign
vectors independent of \(\mcF_n\), then, with possibly larger \(C(C_0)\),
\[
  \begin{aligned}
  &\mathbb P\left\{
  \min_{1\le r\le m}
  \left\|\sum_{i=1}^n\eta_i^{(r)}A_i\right\|_{\op}
  >C(C_0)\sqrt n
  \right\} \\
  &\qquad\le
  \mathbb P\{\Xi>v\}
  +
  D\exp\left\{-\frac{t^2}{2(v+t/3)}\right\}
  +2^{-m}.
  \end{aligned}
\]
In particular, if \(\Xi\le C_0n/(2\log(eD))\) outside an event \(E\), then the
discrepancy failure is at most
\[
  \mathbb P(E)+
  D\exp\left\{-\frac{3C_0n}{16\log(eD)}\right\},
\]
with an additional \(2^{-m}\) for the randomized search.

This covers online adaptive random-basis samplers.  For example, suppose
\(\Gamma_i\succeq0\) are \(\mcF_{i-1}\)-measurable and the one-step conditional
law satisfies, for every \(\mcF_{i-1}\)-measurable positive contraction \(X\),
\[
  \mathbb E\left(U_iXU_i^*\,\middle|\,\mcF_{i-1}\right)
  \preceq
  \frac{\operatorname{tr}X}{D}I+\Gamma_i .
\]
Then the theorem applies with
\[
  V_i=\frac{\operatorname{tr}(T_i^2)}{D}I+\Gamma_i .
\]
Likewise, if \(I=\sum_{\alpha\in I_i}P_{i,\alpha}\) is an
\(\mcF_{i-1}\)-measurable orthogonal decomposition,
\(d_{i,\alpha}=\operatorname{rank}P_{i,\alpha}\), and, for every
\(\mcF_{i-1}\)-measurable positive contraction \(X\),
\[
  \mathbb E\left(U_iXU_i^*\,\middle|\,\mcF_{i-1}\right)
  \preceq
  \sum_{\alpha\in I_i}
  \frac{\operatorname{tr}(P_{i,\alpha}XP_{i,\alpha})}{d_{i,\alpha}}
  P_{i,\alpha}
  +\Gamma_i
\]
then the theorem applies with the predictable block profile
\[
  V_i=
  \sum_{\alpha\in I_i}
  \frac{\operatorname{tr}(P_{i,\alpha}T_i^2P_{i,\alpha})}{d_{i,\alpha}}
  P_{i,\alpha}
  +\Gamma_i .
\]
Thus exact or approximate scalar designs, block designs, signed-permutation
samplers, and finite random-basis samplers may be selected adaptively from the
past; conditional independence between different times is not required.
\end{corollary}

\begin{proof}
The matrices \(A_i\) are \(\mcF_i\)-measurable self-adjoint contractions and
\[
  A_i^2=U_iT_i^2U_i^* .
\]
Therefore the predictable square profile \(\nu\) in
Corollary~\ref{cor:random-square-profile-freedman} satisfies
\[
  \nu
  =
  \left\|\sum_i\mathbb E(A_i^2\mid\mcF_{i-1})\right\|_{\op}
  \le
  \left\|\sum_iV_i\right\|_{\op}
  =
  \Xi .
\]
Applying Corollary~\ref{cor:random-square-profile-freedman} gives the
discrepancy and randomized-search estimates.  The final specialization follows
by taking \(v=t=C_0n/(2\log(eD))\).  The scalar and block forms are obtained by
applying the corresponding one-step domination inequality with \(X=T_i^2\).
\end{proof}

\begin{corollary}[Rank-sensitive predictable dominated twirl streams]
\label{cor:rank-sensitive-predictable-dominated-twirl}
Let \(D\ge2\), let \(1\le\kappa\le D\), and let
\((\mcF_i)_{i=0}^n\) be a filtration.  Let \(P\) be an
\(\mcF_0\)-measurable orthogonal projection with
\[
  1\le \operatorname{rank}P\le\kappa .
\]
Let \(T_1,\ldots,T_n\in\M_D(\C)\) be self-adjoint contractions such that
\[
  T_i\text{ is }\mcF_{i-1}\text{-measurable},
  \qquad
  T_i=PT_iP .
\]
Let \(U_i\) be an \(\mcF_i\)-measurable random unitary satisfying
\[
  U_iP=PU_i,
\]
and set
\[
  A_i=U_iT_iU_i^* .
\]
Assume that there are \(\mcF_{i-1}\)-measurable positive semidefinite matrices
\(V_i\) such that
\[
  \mathbb E\left(U_iT_i^2U_i^*\,\middle|\,\mcF_{i-1}\right)\preceq V_i
  \qquad(1\le i\le n),
\]
and put
\[
  \Xi_\kappa:=
  \left\|\sum_{i=1}^nV_i\right\|_{\op}.
\]
Then, for every \(C_0\ge1\), there is \(C(C_0)\) such that, whenever
\[
  v+t\le\frac{C_0n}{\log(e\kappa)},
\]
one has
\[
  \mathbb P\left\{
  \disc(A_1,\ldots,A_n)>C(C_0)\sqrt n
  \right\}
  \le
  \mathbb P\{\Xi_\kappa>v\}
  +
  \kappa\exp\left\{-\frac{t^2}{2(v+t/3)}\right\}.
\]
Moreover, if \(\eta^{(1)},\ldots,\eta^{(m)}\) are independent uniform sign
vectors independent of \(\mcF_n\), then the same bound holds for sampled sign
search after adding \(2^{-m}\).  In particular, if
\[
  \Xi_\kappa\le \frac{C_0n}{2\log(e\kappa)}
\]
outside an event \(E\), then the discrepancy failure is at most
\[
  \mathbb P(E)+
  \kappa\exp\left\{-\frac{3C_0n}{16\log(e\kappa)}\right\},
\]
with an additional \(2^{-m}\) for randomized search.

This covers online adaptive low-rank random-basis samplers.  If
\(r=\operatorname{rank}P\), \(\Gamma_i\succeq0\) are
\(\mcF_{i-1}\)-measurable, and the one-step conditional law on
\(\operatorname{ran}P\) satisfies, for every \(\mcF_{i-1}\)-measurable
positive contraction \(X=PXP\),
\[
  \mathbb E(U_iXU_i^*\mid\mcF_{i-1})
  \preceq
  \frac{\operatorname{tr}X}{r}P+\Gamma_i ,
\]
then the theorem applies with
\[
  V_i=\frac{\operatorname{tr}(T_i^2)}{r}P+\Gamma_i .
\]
Likewise, if \(P=\sum_{\alpha\in I_i}P_{i,\alpha}\) is an
\(\mcF_{i-1}\)-measurable orthogonal decomposition and the one-step law
satisfies
\[
  \mathbb E(U_iXU_i^*\mid\mcF_{i-1})
  \preceq
  \sum_{\alpha\in I_i}
  \frac{\operatorname{tr}(P_{i,\alpha}XP_{i,\alpha})}
       {\operatorname{rank}P_{i,\alpha}}
  P_{i,\alpha}
  +\Gamma_i
  \qquad(0\preceq X=PXP\preceq P),
\]
then the theorem applies with the corresponding predictable block profile plus
\(\Gamma_i\).  Thus exact or approximate scalar designs, block designs,
signed-permutation samplers, and finite random-basis samplers may be selected
adaptively from the past inside a common random low-rank corner; conditional
independence between different times is not required.
\end{corollary}

\begin{proof}
Because \(U_i\) preserves \(\operatorname{ran}P\) and \(T_i=PT_iP\), the
matrices \(A_i\) are \(\mcF_i\)-measurable self-adjoint contractions satisfying
\[
  A_i=PA_iP .
\]
Also \(A_i^2=U_iT_i^2U_i^*\), so the predictable square profile in
Corollary~\ref{cor:rank-sensitive-predictable-square-profile} satisfies
\[
  \left\|
  \sum_i\mathbb E(A_i^2\mid\mcF_{i-1})
  \right\|_{\op}
  \le
  \left\|\sum_iV_i\right\|_{\op}
  =
  \Xi_\kappa .
\]
Applying Corollary~\ref{cor:rank-sensitive-predictable-square-profile} gives
the discrepancy and sampled-search estimates.  The half-budget form follows
from \(v=t=C_0n/(2\log(e\kappa))\).  The scalar and block specializations are
obtained by substituting \(X=T_i^2\) in the displayed one-step Loewner
domination assumptions.
\end{proof}

\begin{corollary}[Finite-menu rank-sensitive predictable twirl schedules]
\label{cor:finite-menu-rank-sensitive-predictable-dominated-twirl-schedules}
Fix \(C_0\ge1\) and put \(\gamma_0:=3C_0/16\).  For each \(m\ge1\), let
\(D_m\ge2\), let \(J_m\ge1\), and, for \(1\le h\le J_m\), let
\[
  1\le\kappa_{m,h}\le D_m,\qquad
  \ell_{m,h}:=\log(e\kappa_{m,h}).
\]
For every active pair \((m,h)\), let
\[
  A_{m,h,i}=U_{m,h,i}T_{m,h,i}U_{m,h,i}^*,
  \qquad 1\le i\le n_m,
\]
be a rank-sensitive predictable dominated-twirl stream as in
Corollary~\ref{cor:rank-sensitive-predictable-dominated-twirl}, with
filtration \((\mcF_{m,h,i})_{i=0}^{n_m}\), common support projection
\(P_{m,h}\), and rank envelope \(\kappa_{m,h}\).  Thus \(P_{m,h}\) is
\(\mcF_{m,h,0}\)-measurable, \(\operatorname{rank}P_{m,h}\le\kappa_{m,h}\),
\(T_{m,h,i}=P_{m,h}T_{m,h,i}P_{m,h}\), \(U_{m,h,i}P_{m,h}=P_{m,h}U_{m,h,i}\),
and there are predictable positive semidefinite
profiles \(V_{m,h,i}\) satisfying
\[
  \mathbb E\left(
  U_{m,h,i}T_{m,h,i}^2U_{m,h,i}^*
  \,\middle|\,\mcF_{m,h,i-1}
  \right)
  \preceq V_{m,h,i}.
\]
Put
\[
  \Xi_{m,h}^{\kappa}:=
  \left\|\sum_{i=1}^{n_m}V_{m,h,i}\right\|_{\op}.
\]
Assume that
\[
  \Xi_{m,h}^{\kappa}\le\frac{C_0n_m}{2\ell_{m,h}}
  \qquad(1\le h\le J_m)
\]
outside events \(E_{m,h}\), and that
\[
  \sum_m\sum_{h=1}^{J_m}\mathbb P(E_{m,h})<\infty,
  \qquad
  \sum_m\sum_{h=1}^{J_m}
  \kappa_{m,h}\exp\left\{-\gamma_0\frac{n_m}{\ell_{m,h}}\right\}<\infty .
\]
Then, with probability one, for all sufficiently large \(m\), every active
online low-rank twirl stream has a Spencer-scale signing:
\[
  \disc(A_{m,h,1},\ldots,A_{m,h,n_m})
  \le C(C_0)\sqrt{n_m}
  \qquad(1\le h\le J_m).
\]
If the streams are supplied with independent sampled sign batches and
\[
  \sum_m\sum_{h=1}^{J_m}2^{-M_{m,h}}<\infty,
\]
then, after enlarging \(C(C_0)\), almost surely every active stream has a
sampled certificate for all sufficiently large \(m\).  If one common batch of
\(M_m\) signs is reused across the active menu and is independent of all active
streams, it is enough that
\[
  \sum_mJ_m2^{-M_m}<\infty .
\]

Equivalently, with
\[
  \Lambda_m^\kappa:=\sum_{h=1}^{J_m}\kappa_{m,h},
  \qquad
  \ell_m^\kappa:=\max_{1\le h\le J_m}\ell_{m,h},
\]
the structural tail is summable whenever, for some \(\delta>0\),
\[
  n_m\ge
  \gamma_0^{-1}\ell_m^\kappa
  \left(\log(e\Lambda_m^\kappa)+(1+\delta)\log(e+m)\right)
\]
eventually.  With a common sign batch, the sampled conclusion follows if
\[
  M_m\ge \log_2(eJ_m)+(1+\delta')\log_2(e+m)
\]
eventually for some \(\delta'>0\).  If the streams are defined for all
\(h\ge1\), \(J_m\to\infty\), and \(h\le J_m\) is a deterministic exhaustion,
then every fixed adaptive low-rank random-basis stream in the countable library
is covered for all sufficiently large \(m\).
\end{corollary}

\begin{proof}
For a fixed active pair \((m,h)\), apply
Corollary~\ref{cor:rank-sensitive-predictable-dominated-twirl} with
\[
  v=t=\frac{C_0n_m}{2\ell_{m,h}} .
\]
Its failure probability is at most
\[
  \mathbb P(E_{m,h})+
  \kappa_{m,h}\exp\left\{-\gamma_0\frac{n_m}{\ell_{m,h}}\right\},
\]
with the additional sampled-search terms displayed in the statement.  The
Borel--Cantelli lemma gives the almost-sure simultaneous conclusions.  The
explicit schedule follows from
\[
  \sum_{h=1}^{J_m}
  \kappa_{m,h}\exp\left\{-\gamma_0\frac{n_m}{\ell_{m,h}}\right\}
  \le
  \Lambda_m^\kappa
  \exp\left\{-\gamma_0\frac{n_m}{\ell_m^\kappa}\right\}
  \le e^{-1}(e+m)^{-1-\delta}.
\]
The common sampled-batch estimate is identical.
\end{proof}

\begin{corollary}[Almost-sure predictable twirl arrays]
\label{cor:predictable-dominated-twirl-borel}
Fix \(C_0\ge1\).  For each \(m\ge1\), let
\((\mcF_{m,i})_{i=0}^{n_m}\) be a filtration and let
\[
  A_{m,i}=U_{m,i}T_{m,i}U_{m,i}^*\in\M_{D_m}(\C),
  \qquad 1\le i\le n_m,
\]
be a predictable dominated-twirl stream as in
Corollary~\ref{cor:predictable-dominated-twirl-freedman}.  Thus \(T_{m,i}\) is
\(\mcF_{m,i-1}\)-measurable, \(U_{m,i}\) is \(\mcF_{m,i}\)-measurable, and
there are \(\mcF_{m,i-1}\)-measurable positive semidefinite matrices \(V_{m,i}\)
such that
\[
  \mathbb E\left(
  U_{m,i}T_{m,i}^2U_{m,i}^*\,\middle|\,\mcF_{m,i-1}
  \right)
  \preceq V_{m,i}.
\]
Put
\[
  \Xi_m=\left\|\sum_{i=1}^{n_m}V_{m,i}\right\|_{\op}.
\]
Assume that there are deterministic \(v_m,t_m>0\) satisfying
\[
  v_m+t_m\le \frac{C_0n_m}{\log(eD_m)}
\]
and
\[
  \sum_m
  \left[
  \mathbb P\{\Xi_m>v_m\}
  +
  D_m\exp\left\{-\frac{t_m^2}{2(v_m+t_m/3)}\right\}
  \right]<\infty .
\]
Then, with probability one, for all sufficiently large \(m\),
\[
  \disc(A_{m,1},\ldots,A_{m,n_m})\le C(C_0)\sqrt{n_m}.
\]

If, in addition, for each \(m\) we sample \(M_m\) independent uniform sign
vectors \(\eta_m^{(1)},\ldots,\eta_m^{(M_m)}\in\{\pm1\}^{n_m}\), independent of
the \(m\)-th matrix row, and
\[
  \sum_m2^{-M_m}<\infty,
\]
then, with probability one, for all sufficiently large \(m\) at least one sampled
sign vector satisfies
\[
  \left\|
  \sum_{i=1}^{n_m}\eta_{m,i}^{(r)}A_{m,i}
  \right\|_{\op}
  \le C(C_0)\sqrt{n_m}.
\]

In particular, if
\[
  \Xi_m\le \frac{C_0n_m}{2\log(eD_m)}
\]
outside an event \(E_m\), and
\[
  \sum_m
  \left[
  \mathbb P(E_m)
  +
  D_m\exp\left\{-\frac{3C_0n_m}{16\log(eD_m)}\right\}
  \right]<\infty,
\]
then the same almost-sure conclusions hold, with the sampled-certificate
conclusion again requiring \(\sum_m2^{-M_m}<\infty\).  This applies in
particular to predictable exact scalar designs, predictable approximate scalar
designs, predictable block designs, and predictable approximate block designs
by using the one-step profiles in
Corollary~\ref{cor:predictable-dominated-twirl-freedman}.
\end{corollary}

\begin{proof}
For the \(m\)-th row, Corollary~\ref{cor:predictable-dominated-twirl-freedman}
gives failure probability at most
\[
  \mathbb P\{\Xi_m>v_m\}
  +
  D_m\exp\left\{-\frac{t_m^2}{2(v_m+t_m/3)}\right\}.
\]
The displayed summability and the first Borel--Cantelli lemma imply that only
finitely many row failures occur.  The sampled-sign statement follows from the
same finite-row corollary with the additional \(2^{-M_m}\) term and another
application of Borel--Cantelli.  The final specialization is obtained by taking
\(v_m=t_m=C_0n_m/(2\log(eD_m))\).
\end{proof}

\begin{corollary}[Explicit predictable twirl-array schedules]
\label{cor:predictable-dominated-twirl-schedules}
Fix \(C_0\ge1\) and put
\[
  \gamma_0:=\frac{3C_0}{16}.
\]
In the setting of Corollary~\ref{cor:predictable-dominated-twirl-borel}, assume
that
\[
  \Xi_m\le \frac{C_0n_m}{2\log(eD_m)}
\]
outside events \(E_m\) with
\[
  \sum_m\mathbb P(E_m)<\infty.
\]
Then the following deterministic schedule conditions imply that, almost surely,
\[
  \disc(A_{m,1},\ldots,A_{m,n_m})\le C(C_0)\sqrt{n_m}
\]
for all sufficiently large \(m\).  The same conclusions hold for sampled sign
certificates whenever \(\sum_m2^{-M_m}<\infty\).

\begin{enumerate}[label=\textup{(\alph*)}]
\item General schedules.  If, for some \(\delta>0\),
\[
  n_m\ge
  \gamma_0^{-1}\log(eD_m)
  \bigl(\log(eD_m)+(1+\delta)\log(e+m)\bigr)
\]
eventually, then the conclusion holds.

\item Full integer dimensions.  Suppose the rows are indexed by
\(D=2,3,\ldots\), with ambient dimension \(D\).  If
\[
  n_D\ge a\,\log^2(eD)
\]
eventually for some
\[
  a>\frac{32}{3C_0},
\]
then the conclusion holds along all integer dimensions.

\item Sparse stretched-exponential dimensions.  If there are \(c,\kappa>0\) such
that
\[
  D_m\ge \exp(c m^\kappa)
\]
eventually, and
\[
  n_m\ge a\,\log^2(eD_m)
\]
eventually for some
\[
  a>\frac{16}{3C_0},
\]
then the conclusion holds along the sparse schedule.
\end{enumerate}

In each case, the same statement applies to predictable exact scalar designs,
predictable approximate scalar designs, predictable block designs, and
predictable approximate block designs by inserting the one-step profiles from
Corollary~\ref{cor:predictable-dominated-twirl-freedman}.
\end{corollary}

\begin{proof}
By Corollary~\ref{cor:predictable-dominated-twirl-borel}, it remains only to
check
\[
  \sum_m
  D_m\exp\left\{-\gamma_0\frac{n_m}{\log(eD_m)}\right\}<\infty .
\]
For part (a), write \(\ell_m=\log(eD_m)\).  Since \(D_m\le e^{\ell_m}\), the
summand is bounded by
\[
  \exp\{\ell_m-\ell_m-(1+\delta)\log(e+m)\}
  =(e+m)^{-1-\delta},
\]
which is summable.

For part (b), choose \(\delta>0\) so small that
\(a>\gamma_0^{-1}(2+\delta)\).  With \(m=D\), part (a) applies because
\(\log(e+m)=\log(eD)\).

For part (c), choose \(\theta>0\) with \(\gamma_0a-1>\theta\).  Then, eventually,
\[
  D_m\exp\left\{-\gamma_0\frac{n_m}{\log(eD_m)}\right\}
  \le
  \exp\{-\theta\log(eD_m)\}
  \le
  \exp\{-\theta c m^\kappa\},
\]
which is summable.
\end{proof}

\begin{corollary}[Finite-menu predictable twirl-array schedules]
\label{cor:finite-menu-predictable-dominated-twirl-schedules}
Fix \(C_0\ge1\) and put
\[
  \gamma_0:=\frac{3C_0}{16}.
\]
For each \(m\ge1\), let \(D_m\ge2\), put
\[
  \ell_m:=\log(eD_m),
\]
and let \(J_m\ge1\) be a deterministic active-menu size.  For every
\(1\le h\le J_m\), let
\[
  A_{m,i}^{(h)}
  =
  U_{m,h,i}T_{m,h,i}U_{m,h,i}^*
  \in\M_{D_m}(\C),
  \qquad 1\le i\le n_m,
\]
be a predictable dominated-twirl stream as in
Corollary~\ref{cor:predictable-dominated-twirl-freedman}, with its own
filtration \((\mcF_{m,h,i})_{i=0}^{n_m}\) allowed.  Thus there are predictable
positive semidefinite profiles
\(V_{m,h,i}\) such that
\[
  \mathbb E\left(
  U_{m,h,i}T_{m,h,i}^2U_{m,h,i}^*
  \,\middle|\,\mcF_{m,h,i-1}
  \right)
  \preceq V_{m,h,i}.
\]
Put
\[
  \Xi_{m,h}
  :=
  \left\|\sum_{i=1}^{n_m}V_{m,h,i}\right\|_{\op},
  \qquad
  L_m^{\rm tw}:=J_mD_m .
\]
Assume that
\[
  \Xi_{m,h}\le \frac{C_0n_m}{2\ell_m}
  \qquad(1\le h\le J_m)
\]
outside events \(E_{m,h}\), and that
\[
  \sum_m\sum_{h=1}^{J_m}\mathbb P(E_{m,h})<\infty,
  \qquad
  \sum_m L_m^{\rm tw}
  \exp\left\{-\gamma_0\frac{n_m}{\ell_m}\right\}<\infty .
\]
Then, with probability one, for all sufficiently large \(m\), every active menu
element has a Spencer-scale signing:
\[
  \disc(A_{m,1}^{(h)},\ldots,A_{m,n_m}^{(h)})
  \le C(C_0)\sqrt{n_m}
  \qquad(1\le h\le J_m).
\]

If the menu elements are supplied with independent sampled sign batches and
\[
  \sum_m\sum_{h=1}^{J_m}2^{-M_{m,h}}<\infty,
\]
then, after enlarging \(C(C_0)\), almost surely every active menu element has a
sampled certificate for all sufficiently large \(m\).  If instead one common
batch of \(M_m\) sign vectors is reused for the entire active menu and is
independent of all active streams, it is enough that
\[
  \sum_m J_m2^{-M_m}<\infty .
\]

In particular, under the displayed half-budget exception hypothesis, the
structural tail is summable in each of the following schedules.
\begin{enumerate}[label=\textup{(\alph*)}]
\item General schedules.  If, for some \(\delta>0\),
  \[
    n_m\ge
    \gamma_0^{-1}\ell_m
    \bigl(\log(eL_m^{\rm tw})+(1+\delta)\log(e+m)\bigr)
  \]
  eventually, then the simultaneous active-menu conclusion holds.  With a
  common sign batch, the sampled conclusion follows if
  \[
    M_m\ge \log_2(eJ_m)+(1+\delta')\log_2(e+m)
  \]
  eventually for some \(\delta'>0\).
\item Full integer dimensions.  Suppose the rows are indexed by
  \(D=2,3,\ldots\), with \(D_m=D\), and
  \(L_D^{\rm tw}\le(eD)^\Gamma\) eventually for some \(\Gamma\ge0\).  If
  \[
    n_D\ge a\log^2(eD)
  \]
  eventually for some
  \[
    a>\frac{16(\Gamma+1)}{3C_0},
  \]
  then the simultaneous active-menu conclusion holds along all integer
  dimensions.
\item Sparse stretched-exponential dimensions.  If there are \(c,\kappa>0\)
  such that \(D_m\ge\exp(c m^\kappa)\) eventually,
  \(L_m^{\rm tw}\le(eD_m)^\Gamma\) eventually, and
  \[
    n_m\ge a\log^2(eD_m)
  \]
  eventually for some
  \[
    a>\frac{16\Gamma}{3C_0},
  \]
  then the simultaneous active-menu conclusion holds along the sparse schedule.
\end{enumerate}

The same statements apply to finite active menus of predictable exact scalar
designs, predictable approximate scalar designs, predictable block designs, and
predictable approximate block designs by inserting the one-step profiles from
Corollary~\ref{cor:predictable-dominated-twirl-freedman}.  If the streams are
defined for all \(h\ge1\), \(J_m\to\infty\), and the active menu is a
deterministic exhaustion of a countable library, then every fixed library
element is covered for all sufficiently large \(m\).
\end{corollary}

\begin{proof}
For fixed \(m\) and \(h\le J_m\), the half-budget specialization of
Corollary~\ref{cor:predictable-dominated-twirl-borel} gives failure probability
at most
\[
  \mathbb P(E_{m,h})
  +
  D_m\exp\left\{-\gamma_0\frac{n_m}{\ell_m}\right\}.
\]
Union over \(1\le h\le J_m\), then apply the first Borel--Cantelli lemma.  The
sampled-sign assertions add the corresponding \(2^{-M_{m,h}}\) failures, or the
common-batch union bound \(J_m2^{-M_m}\).

For the general schedule,
\[
  L_m^{\rm tw}
  \exp\left\{-\gamma_0\frac{n_m}{\ell_m}\right\}
  \le e^{-1}(e+m)^{-(1+\delta)}.
\]
The common sampled-batch condition gives
\(J_m2^{-M_m}\le e^{-1}(e+m)^{-(1+\delta')}\).  In the all-integer schedule,
\[
  L_D^{\rm tw}
  \exp\left\{-\gamma_0\frac{n_D}{\log(eD)}\right\}
  \le
  (eD)^{\Gamma-\gamma_0a},
\]
which is summable over \(D\) when \(\gamma_0a>\Gamma+1\).  In the sparse
schedule the same estimate is bounded by
\(\exp\{-(\gamma_0a-\Gamma)\log(eD_m)\}\), which is summable under the
stretched-exponential lower bound on \(D_m\) when \(\gamma_0a>\Gamma\).  The
design specializations are exactly the dominated-twirl specializations already
recorded in Corollary~\ref{cor:predictable-dominated-twirl-freedman}.  If
\(J_m\to\infty\), then every fixed countable-menu element lies in the active
menu eventually.
\end{proof}

\begin{corollary}[Predictable subspace-localized ensembles]
\label{cor:predictable-subspace-localized-freedman}
Let \(D\ge2\), let \((\mcF_i)_{i=0}^n\) be a filtration, and let
\(P_1,\ldots,P_n\in\M_D(\C)\) be random orthogonal projections such that \(P_i\)
is \(\mcF_i\)-measurable.  Let \(A_i\) be \(\mcF_i\)-measurable self-adjoint
contractions satisfying
\[
  A_i=P_iA_iP_i,
\]
and put
\[
  \mu:=
  \left\|
  \sum_{i=1}^n\mathbb E(P_i\mid\mcF_{i-1})
  \right\|_{\op}.
\]
Then, for every \(C_0\ge1\), there is \(C(C_0)\) such that, for every \(v,t>0\)
satisfying
\[
  v+t\le \frac{C_0n}{\log(eD)},
\]
one has
\[
  \mathbb P\left\{
  \disc(A_1,\ldots,A_n)>C(C_0)\sqrt n
  \right\}
  \le
  \mathbb P\{\mu>v\}
  +
  D\exp\left\{-\frac{t^2}{2(v+t/3)}\right\}.
\]
Moreover, if \(\eta^{(1)},\ldots,\eta^{(m)}\) are independent uniform sign
vectors independent of \(\mcF_n\), then, with possibly larger \(C(C_0)\),
\[
  \begin{aligned}
  &\mathbb P\left\{
  \min_{1\le r\le m}
  \left\|\sum_{i=1}^n\eta_i^{(r)}A_i\right\|_{\op}
  >C(C_0)\sqrt n
  \right\} \\
  &\qquad\le
  \mathbb P\{\mu>v\}
  +
  D\exp\left\{-\frac{t^2}{2(v+t/3)}\right\}
  +2^{-m}.
  \end{aligned}
\]
In particular, if
\[
  \sum_{i=1}^n\mathbb E(P_i\mid\mcF_{i-1})\preceq pn\,I
  \qquad\text{almost surely}
\]
and \(p\le C_0/(2\log(eD))\), then the discrepancy failure is at most
\[
  D\exp\left\{-\frac{3C_0n}{16\log(eD)}\right\},
\]
with the additional \(2^{-m}\) term for randomized search.  This covers adaptive
coordinate supports, online random subspaces, and conditionally sampled
subspace-localized Wigner/GOE/GUE blocks whenever their predictable overlap
operator is \(O(n/\log(eD))\).
\end{corollary}

\begin{proof}
Since \(A_i=P_iA_iP_i\) and \(\|A_i\|_{\op}\le1\), we have
\[
  0\le A_i^2\preceq P_i.
\]
Indeed, for every vector \(x\),
\[
  \langle A_i^2x,x\rangle=\|A_ix\|^2=\|A_iP_ix\|^2\le \|P_ix\|^2
  =\langle P_ix,x\rangle .
\]
Therefore
\[
  \sum_i\mathbb E(A_i^2\mid\mcF_{i-1})
  \preceq
  \sum_i\mathbb E(P_i\mid\mcF_{i-1}),
\]
so the predictable square-profile \(\nu\) from
Corollary~\ref{cor:random-square-profile-freedman} satisfies \(\nu\le\mu\).
The first two estimates follow by applying that corollary.  If the displayed
predictable overlap bound holds, then
\(\mu\le pn\le C_0n/(2\log(eD))\); choose
\(v=t=C_0n/(2\log(eD))\), as in
Corollary~\ref{cor:random-square-profile-freedman}.
\end{proof}

\begin{corollary}[Finite-menu predictable subspace-localized schedules]
\label{cor:finite-menu-predictable-subspace-localized-schedules}
Fix \(C_0\ge1\) and put
\[
  \gamma_0:=\frac{3C_0}{16}.
\]
For each \(m\ge1\), let \(D_m\ge2\), put
\[
  \ell_m:=\log(eD_m),
\]
and let \(J_m\ge1\) be a deterministic active-menu size.  For every
\(1\le h\le J_m\), let \((\mcF_{m,h,i})_{i=0}^{n_m}\) be a filtration, let
\[
  P_{m,h,1},\ldots,P_{m,h,n_m}\in\M_{D_m}(\C)
\]
be orthogonal projections with \(P_{m,h,i}\) \(\mcF_{m,h,i}\)-measurable, and
let \(A_{m,h,i}\) be \(\mcF_{m,h,i}\)-measurable self-adjoint contractions
satisfying
\[
  A_{m,h,i}=P_{m,h,i}A_{m,h,i}P_{m,h,i}.
\]
Put
\[
  \mu_{m,h}:=
  \left\|
  \sum_{i=1}^{n_m}
  \mathbb E(P_{m,h,i}\mid\mcF_{m,h,i-1})
  \right\|_{\op},
  \qquad
  L_m^{\rm loc}:=J_mD_m .
\]
Assume that
\[
  \mu_{m,h}\le \frac{C_0n_m}{2\ell_m}
  \qquad(1\le h\le J_m)
\]
outside events \(E_{m,h}\), and that
\[
  \sum_m\sum_{h=1}^{J_m}\mathbb P(E_{m,h})<\infty,
  \qquad
  \sum_m L_m^{\rm loc}
  \exp\left\{-\gamma_0\frac{n_m}{\ell_m}\right\}<\infty .
\]
Then, with probability one, for all sufficiently large \(m\), every active
localized stream has a Spencer-scale signing:
\[
  \disc(A_{m,h,1},\ldots,A_{m,h,n_m})
  \le C(C_0)\sqrt{n_m}
  \qquad(1\le h\le J_m).
\]

If the menu elements are supplied with independent sampled sign batches and
\[
  \sum_m\sum_{h=1}^{J_m}2^{-M_{m,h}}<\infty,
\]
then, after enlarging \(C(C_0)\), almost surely every active menu element has a
sampled certificate for all sufficiently large \(m\).  If instead one common
batch of \(M_m\) sign vectors is reused for the entire active menu and is
independent of all active streams, it is enough that
\[
  \sum_mJ_m2^{-M_m}<\infty .
\]

In particular, under the displayed half-budget exception hypothesis, the
structural tail is summable in each of the following schedules.
\begin{enumerate}[label=\textup{(\alph*)}]
\item General schedules.  If, for some \(\delta>0\),
  \[
    n_m\ge
    \gamma_0^{-1}\ell_m
    \bigl(\log(eL_m^{\rm loc})+(1+\delta)\log(e+m)\bigr)
  \]
  eventually, then the simultaneous active-menu conclusion holds.  With a
  common sign batch, the sampled conclusion follows if
  \[
    M_m\ge \log_2(eJ_m)+(1+\delta')\log_2(e+m)
  \]
  eventually for some \(\delta'>0\).
\item Full integer dimensions.  Suppose the rows are indexed by
  \(D=2,3,\ldots\), with \(D_m=D\), and
  \(L_D^{\rm loc}\le(eD)^\Gamma\) eventually for some \(\Gamma\ge0\).  If
  \[
    n_D\ge a\log^2(eD)
  \]
  eventually for some
  \[
    a>\frac{16(\Gamma+1)}{3C_0},
  \]
  then the simultaneous active-menu conclusion holds along all integer
  dimensions.
\item Sparse stretched-exponential dimensions.  If there are \(c,\kappa>0\)
  such that \(D_m\ge\exp(c m^\kappa)\) eventually,
  \(L_m^{\rm loc}\le(eD_m)^\Gamma\) eventually, and
  \[
    n_m\ge a\log^2(eD_m)
  \]
  eventually for some
  \[
    a>\frac{16\Gamma}{3C_0},
  \]
  then the simultaneous active-menu conclusion holds along the sparse schedule.
\end{enumerate}

The half-budget hypothesis is automatic, with no \(E_{m,h}\), whenever
\[
  \sum_{i=1}^{n_m}
  \mathbb E(P_{m,h,i}\mid\mcF_{m,h,i-1})
  \preceq
  p_mn_mI,
  \qquad
  p_m\le \frac{C_0}{2\ell_m},
\]
for all active \(h\).  In the adaptive coordinate-support case this is the
condition
\[
  \max_j\sum_i
  \mathbb P\{j\in S_{m,h,i}\mid\mcF_{m,h,i-1}\}
  \le \frac{C_0n_m}{2\ell_m},
\]
and for conditional Haar rank-\(r_{m,h,i}\) subspaces it is the average-rank
condition
\[
  \sum_i \frac{r_{m,h,i}}{D_m}
  \le \frac{C_0n_m}{2\ell_m}.
\]
If the streams are defined for all \(h\ge1\), \(J_m\to\infty\), and the active
menu is a deterministic exhaustion of a countable library, then every fixed
adaptive localized model is covered for all sufficiently large \(m\).
\end{corollary}

\begin{proof}
For fixed \(m\) and \(h\le J_m\), the half-budget specialization of
Corollary~\ref{cor:predictable-subspace-localized-freedman} gives failure
probability at most
\[
  \mathbb P(E_{m,h})
  +
  D_m\exp\left\{-\gamma_0\frac{n_m}{\ell_m}\right\}.
\]
Union over \(1\le h\le J_m\), then apply the first Borel--Cantelli lemma.  The
sampled-sign assertions add the corresponding \(2^{-M_{m,h}}\) failures, or the
common-batch union bound \(J_m2^{-M_m}\).

The schedule estimates are identical to those in
Corollary~\ref{cor:finite-menu-predictable-dominated-twirl-schedules}, with
\(L_m^{\rm tw}\) replaced by \(L_m^{\rm loc}\).  The coordinate and Haar
specializations are the displayed forms of the predictable overlap operator:
coordinate projections are diagonal with conditional diagonal entries given by
the conditional inclusion probabilities, while conditional Haar rank-\(r\)
projections satisfy
\(\mathbb E(P_{m,h,i}\mid\mcF_{m,h,i-1})=(r_{m,h,i}/D_m)I\).  If
\(J_m\to\infty\), then every fixed countable-menu element lies in the active
menu eventually.
\end{proof}

\begin{corollary}[Random scalar-multiplier ensembles]
\label{cor:random-scalar-multiplier-bernstein}
Let \(D\ge2\).  Let \(B_1,\ldots,B_n\in\M_D(\C)\) be arbitrary, possibly
dependent, random self-adjoint contractions.  Conditional on
\(B=(B_1,\ldots,B_n)\), let \(\theta_1,\ldots,\theta_n\) be independent real
random variables satisfying \(|\theta_i|\le1\), and set
\[
  A_i=\theta_iB_i,\qquad
  q_i=\mathbb E(\theta_i^2\mid B),\qquad
  \nu=\left\|\sum_{i=1}^n q_iB_i^2\right\|_{\op}.
\]
Then, for every \(C_0\ge1\), there is \(C(C_0)\) such that, for all deterministic
\(v\ge0\) and \(t>0\) satisfying
\[
  v+t\le \frac{C_0n}{\log(eD)},
\]
one has
\[
  \mathbb P\left\{
  \disc(A_1,\ldots,A_n)>C(C_0)\sqrt n
  \right\}
  \le
  \mathbb P\{\nu>v\}
  +
  D\exp\left\{-\frac{t^2}{2(v+t/3)}\right\}.
\]
In particular, if \(q_i\le p\) almost surely for every \(i\) and
\[
  p\le \frac{C_0}{2\log(eD)},
\]
then
\[
  \mathbb P\left\{
  \disc(A_1,\ldots,A_n)>C(C_0)\sqrt n
  \right\}
  \le
  D\exp\left\{-\frac{3C_0n}{16\log(eD)}\right\}.
\]
Thus arbitrary deterministic or random self-adjoint contraction tuples become
Spencer-scale after independent Bernoulli thinning of density
\(O(1/\log(eD))\), or after any independent bounded scalar amplitudes with the
same conditional second-moment bound.  No support, commutation, or low-rank
assumption is imposed on the underlying directions \(B_i\).
\end{corollary}

\begin{proof}
Condition on the environment \(B\).  The matrices \(A_i=\theta_iB_i\) are then
independent self-adjoint contractions and
\[
  \left\|\sum_i\mathbb E(A_i^2\mid B)\right\|_{\op}
  =
  \left\|\sum_i q_iB_i^2\right\|_{\op}
  =
  \nu.
\]
On the event \(\{\nu\le v\}\), Corollary~\ref{cor:random-square-profile-bernstein}
applied conditionally gives
\[
  \mathbb P\left(
  \disc(A_1,\ldots,A_n)>C(C_0)\sqrt n\,\middle|\,B
  \right)
  \le
  D\exp\left\{-\frac{t^2}{2(v+t/3)}\right\},
\]
because \(\nu+t\le v+t\le C_0n/\log(eD)\).  Averaging over \(B\) and adding the
exceptional event \(\{\nu>v\}\) proves the first display.

If \(q_i\le p\), then \(0\le B_i^2\le I\), so
\[
  \nu
  \le
  p\left\|\sum_iB_i^2\right\|_{\op}
  \le pn.
\]
Take \(v=pn\) and \(t=C_0n/(2\log(eD))\).  The assumption gives \(v\le t\), and
the same elementary calculation as above gives
\[
  \frac{t^2}{2(v+t/3)}
  \ge
  \frac{3C_0n}{16\log(eD)}.
\]
\end{proof}

\begin{corollary}[Predictable scalar-multiplier ensembles]
\label{cor:predictable-scalar-multiplier-freedman}
Let \(D\ge2\), let \((\mcF_i)_{i=0}^n\) be a filtration, and let
\(H_1,\ldots,H_n\in\M_D(\C)\) be self-adjoint contractions such that \(H_i\) is
\(\mcF_{i-1}\)-measurable.  Let \(\theta_i\) be real \(\mcF_i\)-measurable
random variables with \(|\theta_i|\le1\), and put
\[
  A_i=\theta_iH_i,\qquad
  q_i=\mathbb E(\theta_i^2\mid\mcF_{i-1}),\qquad
  \nu=\left\|\sum_{i=1}^n q_iH_i^2\right\|_{\op}.
\]
Then, for every \(C_0\ge1\), there is \(C(C_0)\) such that, for all
deterministic \(v\ge0\) and \(t>0\) satisfying
\[
  v+t\le \frac{C_0n}{\log(eD)},
\]
one has
\[
  \mathbb P\left\{
  \disc(A_1,\ldots,A_n)>C(C_0)\sqrt n
  \right\}
  \le
  \mathbb P\{\nu>v\}
  +
  D\exp\left\{-\frac{t^2}{2(v+t/3)}\right\}.
\]
If, in addition, \(\eta^{(1)},\ldots,\eta^{(m)}\) are independent uniform sign
vectors independent of \(\mcF_n\), then, with possibly larger \(C(C_0)\),
\[
  \begin{aligned}
  &\mathbb P\left\{
  \min_{1\le r\le m}
  \left\|\sum_{i=1}^n\eta_i^{(r)}A_i\right\|_{\op}
  >C(C_0)\sqrt n
  \right\} \\
  &\qquad\le
  \mathbb P\{\nu>v\}
  +
  D\exp\left\{-\frac{t^2}{2(v+t/3)}\right\}
  +2^{-m}.
  \end{aligned}
\]
In particular, if \(q_i\le p\) almost surely for every \(i\) and
\[
  p\le \frac{C_0}{2\log(eD)},
\]
then the discrepancy failure is at most
\[
  D\exp\left\{-\frac{3C_0n}{16\log(eD)}\right\},
\]
with the additional \(2^{-m}\) term for the randomized search.  Thus the
scalar-multiplier theorem remains valid for online or adaptive matrix
directions and bounded scalar amplitudes; conditional independence of the
amplitudes is not needed once the one-step conditional second moments are
controlled.
\end{corollary}

\begin{proof}
The matrices \(A_i\) are self-adjoint contractions and are
\(\mcF_i\)-measurable.  Since \(H_i\) is \(\mcF_{i-1}\)-measurable and
\(\theta_i\) is scalar,
\[
  \mathbb E(A_i^2\mid\mcF_{i-1})
  =
  \mathbb E(\theta_i^2\mid\mcF_{i-1})H_i^2
  =
  q_iH_i^2.
\]
The first two bounds are therefore exactly
Corollary~\ref{cor:random-square-profile-freedman}.  If \(q_i\le p\), then
\[
  \nu
  \le
  p\left\|\sum_iH_i^2\right\|_{\op}
  \le pn,
\]
and the final estimate follows by taking \(v=pn\) and
\(t=C_0n/(2\log(eD))\), as in
Corollary~\ref{cor:random-scalar-multiplier-bernstein}.
\end{proof}

\begin{corollary}[Almost-sure predictable scalar-multiplier arrays]
\label{cor:predictable-scalar-multiplier-borel}
Fix \(C_0\ge1\).  For each \(m\ge1\), let \(D_m\ge2\), let
\((\mcF_{m,i})_{i=0}^{n_m}\) be a filtration, and let
\[
  A_{m,i}=\theta_{m,i}H_{m,i}\in\M_{D_m}(\C),
  \qquad 1\le i\le n_m,
\]
be a predictable scalar-multiplier stream as in
Corollary~\ref{cor:predictable-scalar-multiplier-freedman}.  Thus
\(H_{m,i}\) is an \(\mcF_{m,i-1}\)-measurable self-adjoint contraction,
\(\theta_{m,i}\) is real, \(\mcF_{m,i}\)-measurable, and
\(|\theta_{m,i}|\le1\).  Put
\[
  q_{m,i}:=\mathbb E(\theta_{m,i}^2\mid\mcF_{m,i-1}),
  \qquad
  \nu_m:=
  \left\|\sum_{i=1}^{n_m}q_{m,i}H_{m,i}^2\right\|_{\op}.
\]
Assume that there are deterministic \(v_m,t_m>0\) satisfying
\[
  v_m+t_m\le \frac{C_0n_m}{\log(eD_m)}
\]
and
\[
  \sum_m
  \left[
  \mathbb P\{\nu_m>v_m\}
  +
  D_m\exp\left\{-\frac{t_m^2}{2(v_m+t_m/3)}\right\}
  \right]<\infty .
\]
Then, with probability one, for all sufficiently large \(m\),
\[
  \disc(A_{m,1},\ldots,A_{m,n_m})\le C(C_0)\sqrt{n_m}.
\]
If, in addition, for each \(m\) we sample \(M_m\) independent uniform sign
vectors \(\eta_m^{(1)},\ldots,\eta_m^{(M_m)}\in\{\pm1\}^{n_m}\), independent of
the \(m\)-th matrix row, and
\[
  \sum_m2^{-M_m}<\infty,
\]
then, with probability one, for all sufficiently large \(m\) at least one sampled
sign vector satisfies
\[
  \left\|
  \sum_{i=1}^{n_m}\eta_{m,i}^{(r)}A_{m,i}
  \right\|_{\op}
  \le C(C_0)\sqrt{n_m}.
\]

In particular, if
\[
  \nu_m\le \frac{C_0n_m}{2\log(eD_m)}
\]
outside an event \(E_m\), and
\[
  \sum_m
  \left[
  \mathbb P(E_m)
  +
  D_m\exp\left\{-\frac{3C_0n_m}{16\log(eD_m)}\right\}
  \right]<\infty,
\]
then the same almost-sure conclusions hold.  This half-budget hypothesis is
implied by the pointwise conditional second-moment bound
\[
  q_{m,i}\le p_m\le\frac{C_0}{2\log(eD_m)}
  \qquad(1\le i\le n_m)
\]
outside \(E_m\), because \(0\le H_{m,i}^2\le I\).
\end{corollary}

\begin{proof}
For the \(m\)-th row, Corollary~\ref{cor:predictable-scalar-multiplier-freedman}
gives failure probability at most
\[
  \mathbb P\{\nu_m>v_m\}
  +
  D_m\exp\left\{-\frac{t_m^2}{2(v_m+t_m/3)}\right\}.
\]
The summability assumption and the first Borel--Cantelli lemma imply that only
finitely many discrepancy failures occur.  The sampled-sign statement follows
from the same finite-row corollary with the additional \(2^{-M_m}\) term.
The half-budget specialization takes
\(v_m=t_m=C_0n_m/(2\log(eD_m))\).  Finally, on the event
\(q_{m,i}\le p_m\) for all \(i\),
\[
  \nu_m
  \le
  p_m\left\|\sum_{i=1}^{n_m}H_{m,i}^2\right\|_{\op}
  \le p_m n_m
  \le\frac{C_0n_m}{2\log(eD_m)} .
\]
\end{proof}

\begin{corollary}[Explicit predictable scalar-multiplier schedules]
\label{cor:predictable-scalar-multiplier-schedules}
Fix \(C_0\ge1\) and put \(\gamma_0:=3C_0/16\).  In the setting of
Corollary~\ref{cor:predictable-scalar-multiplier-borel}, assume that the
half-budget event
\[
  \nu_m\le \frac{C_0n_m}{2\log(eD_m)}
\]
holds outside events \(E_m\) with \(\sum_m\mathbb P(E_m)<\infty\).  Then the
following deterministic schedule conditions imply that, almost surely,
\[
  \disc(A_{m,1},\ldots,A_{m,n_m})\le C(C_0)\sqrt{n_m}
\]
for all sufficiently large \(m\).  The same conclusions hold for sampled sign
certificates whenever \(\sum_m2^{-M_m}<\infty\).

\begin{enumerate}[label=\textup{(\alph*)}]
\item General schedules.  If, for some \(\delta>0\),
\[
  n_m\ge
  \gamma_0^{-1}\log(eD_m)
  \bigl(\log(eD_m)+(1+\delta)\log(e+m)\bigr)
\]
eventually, then the conclusion holds.

\item Full integer dimensions.  Suppose the rows are indexed by
\(D=2,3,\ldots\), with ambient dimension \(D\).  If
\[
  n_D\ge a\,\log^2(eD)
\]
eventually for some \(a>32/(3C_0)\), then the conclusion holds along all integer
dimensions.

\item Sparse stretched-exponential dimensions.  If there are \(c,\kappa>0\) such
that \(D_m\ge\exp(c m^\kappa)\) eventually, and
\[
  n_m\ge a\,\log^2(eD_m)
\]
eventually for some \(a>16/(3C_0)\), then the conclusion holds along the sparse
schedule.
\end{enumerate}

In particular, the same explicit schedules apply whenever the adaptive
amplitudes satisfy \(q_{m,i}\le p_m\le C_0/(2\log(eD_m))\) outside summable
events.  Thus arbitrary adaptive matrix directions become eventually
Spencer-scale after bounded online thinning or scalar amplitudes at the
logarithmic conditional second-moment density.
\end{corollary}

\begin{proof}
By Corollary~\ref{cor:predictable-scalar-multiplier-borel}, it remains only to
check
\[
  \sum_m
  D_m\exp\left\{-\gamma_0\frac{n_m}{\log(eD_m)}\right\}<\infty .
\]
The verification is identical to the proof of
Corollary~\ref{cor:predictable-dominated-twirl-schedules}.  For part (a), with
\(\ell_m=\log(eD_m)\), the summand is at most
\[
  \exp\{\ell_m-\ell_m-(1+\delta)\log(e+m)\}
  =(e+m)^{-1-\delta}.
\]
For part (b), choose \(\delta>0\) so small that
\(a>\gamma_0^{-1}(2+\delta)\); then part (a) applies because along the full
integer sequence \(\log(e+D)\le\log(eD)\).  For part (c), if
\(a>\gamma_0^{-1}\), then
\[
  D_m\exp\left\{-\gamma_0\frac{n_m}{\log(eD_m)}\right\}
  \le
  \exp\{-(\gamma_0a-1)\log(eD_m)\},
\]
which is summable along every stretched-exponential dimension schedule.
\end{proof}

\begin{corollary}[Finite menus of predictable scalar-multiplier schedules]
\label{cor:predictable-scalar-multiplier-menu-schedules}
Fix \(C_0\ge1\).  For each \(m\ge1\), let \(D_m\ge2\), put
\(\ell_m=\log(eD_m)\), and let \(\mathcal H_m\) be a nonempty finite menu.  For
each \(h\in\mathcal H_m\), let
\[
  A_{m,i}^{(h)}=\theta_{m,i}^{(h)}H_{m,i}^{(h)}\in\M_{D_m}(\C),
  \qquad 1\le i\le n_m,
\]
be a predictable scalar-multiplier stream as in
Corollary~\ref{cor:predictable-scalar-multiplier-freedman}, with its own
filtration \((\mcF_{m,i}^{(h)})_{i=0}^{n_m}\) allowed.  Put
\[
  q_{m,i}^{(h)}
  =\mathbb E\bigl((\theta_{m,i}^{(h)})^2\mid \mcF_{m,i-1}^{(h)}\bigr),
  \qquad
  \nu_{m,h}
  =
  \left\|\sum_i q_{m,i}^{(h)}(H_{m,i}^{(h)})^2\right\|_{\op},
\]
and set
\[
  L_m:=|\mathcal H_m|D_m.
\]
Assume that the half-budget condition
\[
  \nu_{m,h}\le \frac{C_0n_m}{2\ell_m}
  \qquad(h\in\mathcal H_m)
\]
holds outside events \(E_{m,h}\), and that
\[
  \sum_m\sum_{h\in\mathcal H_m}\mathbb P(E_{m,h})<\infty,
  \qquad
  \sum_m L_m\exp\left\{-\frac{3C_0n_m}{16\ell_m}\right\}<\infty .
\]
Then, with probability one, for all sufficiently large \(m\), every
\(h\in\mathcal H_m\) has a Spencer-scale signing:
\[
  \disc(A_{m,1}^{(h)},\ldots,A_{m,n_m}^{(h)})
  \le C(C_0)\sqrt{n_m}.
\]
If each menu element is also supplied with \(M_{m,h}\) independent uniform sign
vectors in \(\{\pm1\}^{n_m}\), independent of the corresponding row, and
\[
  \sum_m\sum_{h\in\mathcal H_m}2^{-M_{m,h}}<\infty,
\]
then after enlarging \(C(C_0)\), almost surely for all sufficiently large \(m\)
every menu element has one sampled sign certifying the same bound.  If a common
batch of \(M_m\) signs is reused for all \(h\in\mathcal H_m\), independent of
the whole \(m\)-th menu, it is enough that
\[
  \sum_m|\mathcal H_m|\,2^{-M_m}<\infty .
\]

In particular, under the displayed half-budget exception hypothesis, the
structural menu tail is summable in each of the following schedules.
\begin{enumerate}[label=\textup{(\alph*)}]
\item General schedules.  If, for some \(\delta>0\),
  \[
    n_m\ge
    \frac{16}{3C_0}\ell_m
    \bigl(\log(eL_m)+(1+\delta)\log(e+m)\bigr)
  \]
  eventually, then the simultaneous almost-sure signing conclusion holds.  With
  a common sign batch, the simultaneous sampled conclusion follows if
  \[
    M_m\ge \log_2(e|\mathcal H_m|)+(1+\delta')\log_2(e+m)
  \]
  eventually for some \(\delta'>0\).
\item Full integer dimensions.  Suppose the rows are indexed by
  \(D=2,3,\ldots\), with \(D_m=D\), and \(L_D\le(eD)^\Gamma\) eventually for
  some \(\Gamma\ge0\).  If
  \[
    n_D\ge a\log^2(eD)
  \]
  eventually for some \(a>16(\Gamma+1)/(3C_0)\), then the simultaneous
  almost-sure signing conclusion holds.
\item Sparse stretched-exponential dimensions.  If there are \(c,\kappa>0\)
  such that \(D_m\ge\exp(c m^\kappa)\) eventually, \(L_m\le(eD_m)^\Gamma\)
  eventually, and
  \[
    n_m\ge a\log^2(eD_m)
  \]
  eventually for some \(a>16\Gamma/(3C_0)\), then the simultaneous
  almost-sure signing conclusion holds.
\end{enumerate}
The half-budget hypothesis is automatic, with no \(E_{m,h}\), whenever
\[
  q_{m,i}^{(h)}\le p_m\le \frac{C_0}{2\ell_m}
  \qquad(h\in\mathcal H_m,\ 1\le i\le n_m),
\]
because \(0\le (H_{m,i}^{(h)})^2\le I\).
\end{corollary}

\begin{proof}
For each fixed \(m\) and \(h\), Corollary~\ref{cor:predictable-scalar-multiplier-borel}
with \(v_m=t_m=C_0n_m/(2\ell_m)\) gives failure probability at most
\[
  \mathbb P(E_{m,h})
  +
  D_m\exp\left\{-\frac{3C_0n_m}{16\ell_m}\right\}.
\]
Union over \(h\in\mathcal H_m\), and apply the first Borel--Cantelli lemma.
The sampled-sign statements follow from
Corollary~\ref{cor:predictable-scalar-multiplier-freedman} by adding the
corresponding \(2^{-M_{m,h}}\) failures, or by union bounding the reused common
batch as \(|\mathcal H_m|2^{-M_m}\).

For the general schedule, the lower bound on \(n_m\) gives
\[
  L_m\exp\left\{-\frac{3C_0n_m}{16\ell_m}\right\}
  \le e^{-1}(e+m)^{-(1+\delta)}.
\]
The common sampled-batch condition gives
\(|\mathcal H_m|2^{-M_m}\le e^{-1}(e+m)^{-(1+\delta')}\).  The all-integer and
sparse cases are the same estimates with \(L_m\le(eD_m)^\Gamma\).  Finally,
the pointwise \(q\)-bound implies
\[
  \nu_{m,h}
  \le p_m\left\|\sum_i(H_{m,i}^{(h)})^2\right\|_{\op}
  \le p_m n_m
  \le \frac{C_0n_m}{2\ell_m}.
\]
\end{proof}

\begin{corollary}[Exhaustive countable menus of predictable scalar-multiplier schedules]
\label{cor:predictable-scalar-multiplier-countable-exhaustion-schedules}
Fix \(C_0\ge1\).  For each \(m\ge1\), let \(D_m\ge2\), put
\(\ell_m=\log(eD_m)\), and let
\[
  J_1\le J_2\le\cdots,\qquad J_m\to\infty,
\]
be a deterministic exhaustion of the countable menu
\(\mathbb N=\{1,2,\ldots\}\).  For each \(m\) and \(h\ge1\), let
\[
  A_{m,i}^{(h)}=\theta_{m,i}^{(h)}H_{m,i}^{(h)}\in\M_{D_m}(\C),
  \qquad 1\le i\le n_m,
\]
be a predictable scalar-multiplier stream as in
Corollary~\ref{cor:predictable-scalar-multiplier-freedman}, with its own
filtration \((\mcF_{m,i}^{(h)})_{i=0}^{n_m}\) allowed.  Put
\[
  q_{m,i}^{(h)}
  =\mathbb E\bigl((\theta_{m,i}^{(h)})^2\mid \mcF_{m,i-1}^{(h)}\bigr),
  \qquad
  \nu_{m,h}
  =
  \left\|\sum_i q_{m,i}^{(h)}(H_{m,i}^{(h)})^2\right\|_{\op},
\]
and set
\[
  L_m^{(J)}:=J_mD_m .
\]
Assume that the half-budget condition
\[
  \nu_{m,h}\le \frac{C_0n_m}{2\ell_m}
  \qquad(1\le h\le J_m)
\]
holds outside events \(E_{m,h}\), and that
\[
  \sum_m\sum_{1\le h\le J_m}\mathbb P(E_{m,h})<\infty,
  \qquad
  \sum_m L_m^{(J)}\exp\left\{-\frac{3C_0n_m}{16\ell_m}\right\}<\infty .
\]
Then, with probability one, for all sufficiently large \(m\), every active
prefix stream \(1\le h\le J_m\) has a Spencer-scale signing:
\[
  \disc(A_{m,1}^{(h)},\ldots,A_{m,n_m}^{(h)})
  \le C(C_0)\sqrt{n_m}\qquad(1\le h\le J_m).
\]
Consequently, for each fixed \(h\), the same bound holds for all sufficiently
large \(m\).  If the active-prefix streams are supplied with independent sampled
signs and
\[
  \sum_m\sum_{1\le h\le J_m}2^{-M_{m,h}}<\infty,
\]
then after enlarging \(C(C_0)\), almost surely every active-prefix stream has a
sampled certificate for all sufficiently large \(m\).  If one common batch of
\(M_m\) signs is reused for all \(1\le h\le J_m\), independent of the whole
active prefix, it is enough that
\[
  \sum_mJ_m\,2^{-M_m}<\infty .
\]

In particular, under the displayed half-budget exception hypothesis, the
structural prefix tail is summable in each of the following schedules.
\begin{enumerate}[label=\textup{(\alph*)}]
\item General schedules.  If, for some \(\delta>0\),
  \[
    n_m\ge
    \frac{16}{3C_0}\ell_m
    \bigl(\log(eL_m^{(J)})+(1+\delta)\log(e+m)\bigr)
  \]
  eventually, then the active-prefix and fixed-stream almost-sure signing
  conclusions hold.  With a common sign batch, the sampled conclusion follows
  if
  \[
    M_m\ge \log_2(eJ_m)+(1+\delta')\log_2(e+m)
  \]
  eventually for some \(\delta'>0\).
\item Full integer dimensions.  Suppose the rows are indexed by
  \(D=2,3,\ldots\), with \(D_m=D\), \(J_D\to\infty\), and
  \(L_D^{(J)}\le(eD)^\Gamma\) eventually for some \(\Gamma\ge0\).  If
  \[
    n_D\ge a\log^2(eD)
  \]
  eventually for some \(a>16(\Gamma+1)/(3C_0)\), then the active-prefix and
  fixed-stream almost-sure signing conclusions hold.
\item Sparse stretched-exponential dimensions.  If there are \(c,\kappa>0\)
  such that \(D_m\ge\exp(c m^\kappa)\) eventually, \(J_m\to\infty\),
  \(L_m^{(J)}\le(eD_m)^\Gamma\) eventually, and
  \[
    n_m\ge a\log^2(eD_m)
  \]
  eventually for some \(a>16\Gamma/(3C_0)\), then the active-prefix and
  fixed-stream almost-sure signing conclusions hold.
\end{enumerate}
The half-budget hypothesis is automatic, with no \(E_{m,h}\), whenever
\[
  q_{m,i}^{(h)}\le p_m\le \frac{C_0}{2\ell_m}
  \qquad(1\le h\le J_m,\ 1\le i\le n_m),
\]
because \(0\le (H_{m,i}^{(h)})^2\le I\).
\end{corollary}

\begin{proof}
Apply Corollary~\ref{cor:predictable-scalar-multiplier-menu-schedules} to the
finite active menu \(\{1,\ldots,J_m\}\) in the \(m\)-th row.  The displayed
exception and structural summability hypotheses are exactly the finite-menu
hypotheses for those active prefixes, so Borel--Cantelli gives the simultaneous
active-prefix conclusion.  Since \(J_m\to\infty\), each fixed \(h\) lies in the
active prefix for all sufficiently large \(m\).  The sampled-sign claims are
the finite-menu sampled claims summed over \(1\le h\le J_m\), or the common
batch union bound \(J_m2^{-M_m}\).

For the general schedule,
\[
  L_m^{(J)}\exp\left\{-\frac{3C_0n_m}{16\ell_m}\right\}
  \le e^{-1}(e+m)^{-(1+\delta)}.
\]
The common sampled-batch condition gives
\(J_m2^{-M_m}\le e^{-1}(e+m)^{-(1+\delta')}\).  The all-integer and sparse
dimension cases are the same estimates with
\(L_m^{(J)}\le(eD_m)^\Gamma\).  Finally, the pointwise \(q\)-bound implies
\[
  \nu_{m,h}
  \le p_m\left\|\sum_i(H_{m,i}^{(h)})^2\right\|_{\op}
  \le p_m n_m
  \le \frac{C_0n_m}{2\ell_m}
\]
for every active-prefix stream.
\end{proof}

\begin{corollary}[Core signings with centered noise]
\label{cor:core-plus-centered-scalar-noise}
Let \(D\ge2\).  Let \(B_i,H_i\in\M_D(\C)\), \(i=1,\ldots,n\), be arbitrary,
possibly dependent, random self-adjoint contractions.  Write \(B,H\) for the two
tuples.  Conditional on \((B,H)\), let
\(\theta_1,\ldots,\theta_n\) be independent real random variables satisfying
\[
  \mathbb E(\theta_i\mid B,H)=0,\qquad |\theta_i|\le1,
\]
and set
\[
  X_i=B_i+\theta_iH_i,\qquad
  q_i=\mathbb E(\theta_i^2\mid B,H),\qquad
  \nu=\left\|\sum_{i=1}^nq_iH_i^2\right\|_{\op}.
\]
Assume that, outside an event \(E\), there are signs
\(\eps_i^\star=\eps_i^\star(B,H)\in\{-1,1\}\) such that
\[
  \left\|\sum_{i=1}^n\eps_i^\star B_i\right\|_{\op}\le K\sqrt n .
\]
Then for every deterministic \(v\ge0\) and \(s>0\),
\[
\begin{aligned}
  \mathbb P\left\{
  \disc(X_1,\ldots,X_n)>K\sqrt n+s
  \right\}
  &\le
  \mathbb P(E)+\mathbb P\{\nu>v\} \\
  &\quad+
  2D\exp\left\{-\frac{s^2}{2(v+s/3)}\right\}.
\end{aligned}
\]
Since \(\|X_i\|_{\op}\le2\), the normalized contractions \(X_i/2\) satisfy the
same estimate with the right-hand discrepancy threshold divided by \(2\).

In particular, if \(q_i\le p\) almost surely for every \(i\), then
\[
  \mathbb P\left\{
  \disc(X_1,\ldots,X_n)>K\sqrt n+s
  \right\}
  \le
  \mathbb P(E)
  +
  2D\exp\left\{-\frac{s^2}{2(pn+s/3)}\right\}.
\]
Thus algebraic cores, finite-lift cores, Toeplitz/Hankel cores, or any other
class already carrying a \(K\sqrt n\) signing remain Spencer-scale after adding
centered independent Bernoulli or bounded scalar-amplitude noise in arbitrary
self-adjoint contraction directions, whenever the displayed Bernstein tail is
small at \(s=O(\sqrt n)\).
\end{corollary}

\begin{proof}
Work on \(E^c\) and condition on the environment \((B,H)\).  For the
environment-measurable signing \(\eps^\star\), set
\[
  Z_i=\eps_i^\star\theta_iH_i.
\]
Then \(Z_i\) are conditionally independent centered self-adjoint matrices with
\(\|Z_i\|_{\op}\le1\), and
\[
  \left\|\sum_i\mathbb E(Z_i^2\mid B,H)\right\|_{\op}
  =
  \left\|\sum_i q_iH_i^2\right\|_{\op}
  =
  \nu.
\]
The two-sided matrix Bernstein inequality \cite{Tropp} gives, on the event
\(\{\nu\le v\}\),
\[
  \mathbb P\left(
  \left\|\sum_iZ_i\right\|_{\op}>s\,\middle|\,B,H
  \right)
  \le
  2D\exp\left\{-\frac{s^2}{2(v+s/3)}\right\}.
\]
On the complementary good event,
\[
  \left\|\sum_i\eps_i^\star X_i\right\|_{\op}
  \le
  \left\|\sum_i\eps_i^\star B_i\right\|_{\op}
  +
  \left\|\sum_iZ_i\right\|_{\op}
  \le K\sqrt n+s.
\]
Taking expectations over the environment and adding \(E\cup\{\nu>v\}\) proves
the first display.  If \(q_i\le p\), then \(0\le H_i^2\le I\), so
\(\nu\le p\|\sum_iH_i^2\|_{\op}\le pn\), which gives the final form.
\end{proof}

\begin{corollary}[Core signings with matrix martingale noise]
\label{cor:core-plus-matrix-martingale-noise}
Let \(D\ge2\).  Let \(\mathcal G\) be a sigma-field, and let
\(B_1,\ldots,B_n\in\M_D(\C)\) be \(\mathcal G\)-measurable random self-adjoint
contractions.  Let \((\mcF_i)_{i=0}^n\) be a filtration with
\(\mathcal G\subset\mcF_0\), and let \(Y_i\in\M_D(\C)\) be self-adjoint
\(\mcF_i\)-measurable random matrices satisfying
\[
  \mathbb E(Y_i\mid\mcF_{i-1})=0,\qquad \|Y_i\|_{\op}\le1.
\]
Set
\[
  X_i=B_i+Y_i,\qquad
  \nu=\left\|\sum_{i=1}^n\mathbb E(Y_i^2\mid\mcF_{i-1})\right\|_{\op}.
\]
Assume that, outside an event \(E\in\mathcal G\), there are
\(\mathcal G\)-measurable signs \(\eps_i^\star\in\{-1,1\}\) such that
\[
  \left\|\sum_{i=1}^n\eps_i^\star B_i\right\|_{\op}\le K\sqrt n .
\]
Then, for every deterministic \(v\ge0\) and \(s>0\),
\[
\begin{aligned}
  \mathbb P\left\{
  \disc(X_1,\ldots,X_n)>K\sqrt n+s
  \right\}
  &\le
  \mathbb P(E)+\mathbb P\{\nu>v\} \\
  &\quad+
  2D\exp\left\{-\frac{s^2}{2(v+s/3)}\right\}.
\end{aligned}
\]
Since \(\|X_i\|_{\op}\le2\), the normalized contractions \(X_i/2\) satisfy the
same estimate with the discrepancy threshold divided by \(2\).  Thus any random
core class with an environment-measurable \(K\sqrt n\) signing admits the
displayed tail after arbitrary bounded self-adjoint matrix-valued
martingale-difference noise.  Equivalently, the random noise has the usual
Freedman signing scale
\(s\asymp \sqrt{v\log(eD)}+\log(eD)\); this is \(O(\sqrt n)\) when
\(v=O(n/\log(eD))\) and \(\log(eD)=O(\sqrt n)\).
\end{corollary}

\begin{proof}
Work on \(E^c\).  For the \(\mathcal G\)-measurable signing \(\eps^\star\), set
\[
  Z_i=\eps_i^\star Y_i.
\]
Since \(\mathcal G\subset\mcF_{i-1}\), the matrices \(Z_i\) are self-adjoint
martingale differences with \(\|Z_i\|_{\op}\le1\), and
\[
  \mathbb E(Z_i^2\mid\mcF_{i-1})
  =
  \mathbb E(Y_i^2\mid\mcF_{i-1}).
\]
On \(\{\nu\le v\}\), the predictable quadratic variation of the martingale
\(\sum_iZ_i\) has operator norm at most \(v\).  Applying the matrix Freedman
inequality \cite{Tropp} to \(Z_i\) and to \(-Z_i\), and taking a union bound,
gives
\[
  \mathbb P\left\{
  \left\|\sum_iZ_i\right\|_{\op}>s\ \text{and}\ \nu\le v
  \right\}
  \le
  2D\exp\left\{-\frac{s^2}{2(v+s/3)}\right\}.
\]
On the complementary good event,
\[
  \left\|\sum_i\eps_i^\star X_i\right\|_{\op}
  \le
  \left\|\sum_i\eps_i^\star B_i\right\|_{\op}
  +
  \left\|\sum_iZ_i\right\|_{\op}
  \le K\sqrt n+s.
\]
Adding the exceptional events \(E\) and \(\{\nu>v\}\) proves the estimate.
\end{proof}

\begin{corollary}[Almost-sure signed-core matrix martingale arrays]
\label{cor:core-plus-matrix-martingale-noise-borel}
For each \(m\ge1\), let \(D_m\ge2\), let \(\mathcal G_m\) be a sigma-field, and
let \((\mcF_{m,i})_{i=0}^{n_m}\) be a filtration with
\(\mathcal G_m\subset\mcF_{m,0}\).  Let
\[
  X_{m,i}=B_{m,i}+Y_{m,i}\in\M_{D_m}(\C),
  \qquad 1\le i\le n_m,
\]
where \(B_{m,i}\) is a \(\mathcal G_m\)-measurable random self-adjoint
contraction and \(Y_{m,i}\) is a self-adjoint \(\mcF_{m,i}\)-measurable
martingale difference satisfying
\[
  \mathbb E(Y_{m,i}\mid\mcF_{m,i-1})=0,
  \qquad \|Y_{m,i}\|_{\op}\le1.
\]
Put
\[
  \nu_m:=
  \left\|
  \sum_{i=1}^{n_m}\mathbb E(Y_{m,i}^2\mid\mcF_{m,i-1})
  \right\|_{\op}.
\]
Assume that, outside an event \(E_m\in\mathcal G_m\), there are
\(\mathcal G_m\)-measurable signs \(\eps_{m,i}^\star\) such that
\[
  \left\|\sum_{i=1}^{n_m}\eps_{m,i}^\star B_{m,i}\right\|_{\op}
  \le K_m\sqrt{n_m}.
\]
Let \(v_m\ge0\) and \(s_m>0\) be deterministic.  If
\[
  \sum_m
  \left[
  \mathbb P(E_m)+\mathbb P\{\nu_m>v_m\}
  +
  2D_m\exp\left\{-\frac{s_m^2}{2(v_m+s_m/3)}\right\}
  \right]<\infty,
\]
then, with probability one, for all sufficiently large \(m\),
\[
  \disc(X_{m,1},\ldots,X_{m,n_m})
  \le K_m\sqrt{n_m}+s_m .
\]

In particular, put \(\ell_m=\log(eD_m)\).  Suppose that \(K_m\le K\), that
\[
  \nu_m\le Q\frac{n_m}{\ell_m}
\]
outside events \(F_m\), that
\[
  \sum_m[\mathbb P(E_m)+\mathbb P(F_m)]<\infty,
\]
and that, for some \(A,B>0\),
\[
  \sum_m(eD_m)^{-A}<\infty,
  \qquad
  \ell_m\le B\sqrt{n_m}
\]
eventually.  Choose \(L>0\) with
\[
  \frac{L^2}{2(Q+LB/3)}\ge A+1 .
\]
Then, almost surely, for all sufficiently large \(m\),
\[
  \disc(X_{m,1},\ldots,X_{m,n_m})
  \le (K+L)\sqrt{n_m}.
\]

Consequently:
\begin{enumerate}[label=\textup{(\alph*)}]
\item Along all integer dimensions \(D=2,3,\ldots\), if \(K_D\le K\),
\(\nu_D\le Qn_D/\log(eD)\) outside summable events, and
\[
  n_D\ge a\,\log^2(eD)
\]
eventually for some \(a>0\), then for every \(A>1\) and every \(L\) satisfying
\[
  \frac{L^2}{2(Q+L/(3\sqrt a))}\ge A+1
\]
one has
\[
  \disc(X_{D,1},\ldots,X_{D,n_D})\le (K+L)\sqrt{n_D}
\]
eventually almost surely.

\item If \(D_m\ge\exp(c m^\kappa)\) eventually for some \(c,\kappa>0\), if
\(K_m\le K\), if \(\nu_m\le Qn_m/\log(eD_m)\) outside summable events, and if
\[
  n_m\ge a\,\log^2(eD_m)
\]
eventually for some \(a>0\), then the same conclusion holds along the sparse
schedule for every \(A>0\) and every \(L\) satisfying
\[
  \frac{L^2}{2(Q+L/(3\sqrt a))}\ge A+1 .
\]
\end{enumerate}
The normalized contractions \(X_{m,i}/2\) satisfy the same statements with the
discrepancy thresholds divided by \(2\).
\end{corollary}

\begin{proof}
For the \(m\)-th row, write
\(\Delta_m=\disc(X_{m,1},\ldots,X_{m,n_m})\).
Corollary~\ref{cor:core-plus-matrix-martingale-noise} gives
\[
\begin{aligned}
  \mathbb P\{\Delta_m>K_m\sqrt{n_m}+s_m\}
  &\le
  \mathbb P(E_m)+\mathbb P\{\nu_m>v_m\}
  \\
  &\quad+
  2D_m\exp\left\{-\frac{s_m^2}{2(v_m+s_m/3)}\right\}.
\end{aligned}
\]
The displayed summability and the first Borel--Cantelli lemma prove the first
claim.

For the concrete profile, take \(s_m=L\sqrt{n_m}\) and
\(v_m=Qn_m/\ell_m\).  On \(F_m^c\),
\[
  \frac{s_m^2}{2(v_m+s_m/3)}
  =
  \frac{L^2n_m}{2(Qn_m/\ell_m+L\sqrt{n_m}/3)}
  \ge
  \frac{L^2}{2(Q+LB/3)}\,\ell_m
  \ge (A+1)\ell_m .
\]
Therefore
\[
  2D_m\exp\left\{-\frac{s_m^2}{2(v_m+s_m/3)}\right\}
  \le 2(eD_m)^{-A},
\]
which is summable by assumption.

For part (a), take \(B=1/\sqrt a\) and any \(A>1\), since
\(\sum_{D\ge2}(eD)^{-A}<\infty\).  For part (b), again take
\(B=1/\sqrt a\); the stretched-exponential lower bound on \(D_m\) makes
\(\sum_m(eD_m)^{-A}\) finite for every \(A>0\).  The final normalized statement
is obtained by homogeneity.
\end{proof}

\begin{corollary}[Exhaustive countable menus of signed-core matrix martingale-noise arrays]
\label{cor:core-plus-matrix-martingale-noise-countable-exhaustion}
For each \(m\ge1\), let \(D_m\ge2\), put
\(\ell_m=\log(eD_m)\), and let
\[
  J_1\le J_2\le\cdots,\qquad J_m\to\infty,
\]
be a deterministic exhaustion of the countable menu
\(\mathbb N=\{1,2,\ldots\}\).  For each \(m\) and \(h\ge1\), let
\(\mathcal G_m^{(h)}\) be a sigma-field and let
\((\mcF_{m,i}^{(h)})_{i=0}^{n_m}\) be a filtration with
\(\mathcal G_m^{(h)}\subset\mcF_{m,0}^{(h)}\).  Let
\[
  X_{m,i}^{(h)}=B_{m,i}^{(h)}+Y_{m,i}^{(h)}\in\M_{D_m}(\C),
  \qquad 1\le i\le n_m,
\]
where \(B_{m,i}^{(h)}\) is a
\(\mathcal G_m^{(h)}\)-measurable self-adjoint contraction and
\(Y_{m,i}^{(h)}\) is a self-adjoint
\(\mcF_{m,i}^{(h)}\)-measurable martingale difference satisfying
\[
  \mathbb E(Y_{m,i}^{(h)}\mid\mcF_{m,i-1}^{(h)})=0,
  \qquad \|Y_{m,i}^{(h)}\|_{\op}\le1 .
\]
Put
\[
  \nu_{m,h}:=
  \left\|
  \sum_{i=1}^{n_m}
  \mathbb E\bigl((Y_{m,i}^{(h)})^2\mid\mcF_{m,i-1}^{(h)}\bigr)
  \right\|_{\op}.
\]
Assume that there is \(K<\infty\) such that, outside events
\(E_{m,h}\in\mathcal G_m^{(h)}\), there are
\(\mathcal G_m^{(h)}\)-measurable signs \(\eps_{m,i}^{(h),\star}\) with
\[
  \left\|\sum_{i=1}^{n_m}\eps_{m,i}^{(h),\star}B_{m,i}^{(h)}\right\|_{\op}
  \le K\sqrt{n_m}
  \qquad(1\le h\le J_m).
\]
Assume also that, for some \(Q,A,B_0>0\), the variance budgets obey
\[
  \nu_{m,h}\le Q\frac{n_m}{\ell_m}\qquad(1\le h\le J_m)
\]
outside events \(F_{m,h}\), and that
\[
  \sum_m\sum_{1\le h\le J_m}
  \bigl[\mathbb P(E_{m,h})+\mathbb P(F_{m,h})\bigr]<\infty,
  \qquad
  \sum_mJ_m(eD_m)^{-A}<\infty,
\]
while
\[
  \ell_m\le B_0\sqrt{n_m}
\]
eventually.  If \(L>0\) satisfies
\[
  \frac{L^2}{2(Q+LB_0/3)}\ge A+1,
\]
then, with probability one, for all sufficiently large \(m\), every active
prefix member \(1\le h\le J_m\) satisfies
\[
  \disc(X_{m,1}^{(h)},\ldots,X_{m,n_m}^{(h)})
  \le (K+L)\sqrt{n_m}.
\]
Consequently, for every fixed \(h\), the same bound holds for all sufficiently
large \(m\).  The normalized contractions \(X_{m,i}^{(h)}/2\) satisfy the same
statements with the discrepancy threshold divided by \(2\).  The active-prefix
conclusion also applies to arbitrary row-dependent finite menus with \(J_m\)
members; only the final fixed-\(h\) conclusion uses the countable exhaustion.

In particular, if \(n_m\ge a\log^2(eD_m)\) eventually, then one may take
\(B_0=a^{-1/2}\).  Along all integer dimensions \(D=2,3,\ldots\), if
\(J_D\le(eD)^\Gamma\) eventually and \(n_D\ge a\log^2(eD)\), then the conclusion
holds for every \(A>\Gamma+1\) and every \(L\) satisfying
\[
  \frac{L^2}{2(Q+L/(3\sqrt a))}\ge A+1 .
\]
Along a sparse schedule with \(D_m\ge\exp(c m^\kappa)\) eventually,
\(J_m\le(eD_m)^\Gamma\) eventually, and
\(n_m\ge a\log^2(eD_m)\), the same conclusion holds for every \(A>\Gamma\) and
every \(L\) satisfying the same display.
\end{corollary}

\begin{proof}
For fixed \(m\) and \(h\le J_m\), apply
Corollary~\ref{cor:core-plus-matrix-martingale-noise} with
\[
  v_m=Q\frac{n_m}{\ell_m},
  \qquad
  s_m=L\sqrt{n_m}.
\]
Outside \(F_{m,h}\), the variance parameter is at most \(v_m\).  Therefore
\[
\begin{aligned}
  &\mathbb P\left\{
  \disc(X_{m,1}^{(h)},\ldots,X_{m,n_m}^{(h)})
  >(K+L)\sqrt{n_m}\right\} \\
  &\qquad\le
  \mathbb P(E_{m,h})+\mathbb P(F_{m,h})
  +2D_m\exp\left\{
  -\frac{L^2n_m}{2(Qn_m/\ell_m+L\sqrt{n_m}/3)}
  \right\}.
\end{aligned}
\]
For all sufficiently large \(m\), \(\ell_m\le B_0\sqrt{n_m}\), so the
exponent is at least \((A+1)\ell_m\).  Hence the last term is at most
\[
  2D_m(eD_m)^{-(A+1)}\le 2(eD_m)^{-A}.
\]
Summing over \(1\le h\le J_m\) and then over \(m\), the displayed hypotheses
make the total failure probability finite.  The first Borel--Cantelli lemma
gives the simultaneous active-prefix conclusion.  Since \(J_m\to\infty\), each
fixed \(h\) belongs to the active prefix eventually.  The normalized statement
is homogeneity, and the finite-menu statement is the same union bound over the
row's finite menu.

If \(n_m\ge a\log^2(eD_m)\), then
\(\ell_m\le a^{-1/2}\sqrt{n_m}\).  For all integer dimensions, the prefix
series is bounded by a constant multiple of
\(\sum_D(eD)^{\Gamma-A}\), which converges when \(A>\Gamma+1\).  For the sparse
case it is bounded by
\(\sum_m\exp\{-(A-\Gamma)\log(eD_m)\}\), which converges under the
stretched-exponential lower bound whenever \(A>\Gamma\).
\end{proof}

\begin{corollary}[Scalar-dominated matrix martingale noise with random amplitudes]
\label{cor:core-plus-scalar-variance-random-amplitude-noise}
Let \(D\ge2\).  Let \(\mathcal G\) be a sigma-field, and let
\(B_1,\ldots,B_n\in\M_D(\C)\) be \(\mathcal G\)-measurable random self-adjoint
contractions.  Assume that, outside an event \(E\in\mathcal G\), there are
\(\mathcal G\)-measurable signs \(\eps_i^\star\in\{-1,1\}\) such that
\[
  \left\|\sum_{i=1}^n\eps_i^\star B_i\right\|_{\op}\le K\sqrt n .
\]
Let \((\theta_i,\rho_i)\), \(i=1,\ldots,n\), be independent pairs, independent
of \(\mathcal G\), with \(\theta_i,\rho_i\in[0,1]\).  Let
\(\theta_*\in[0,1]\) be deterministic with \(\theta_i\le\theta_*\) almost
surely for every \(i\).  Suppose that \((\mcF_i)_{i=0}^n\) is a filtration with
\(\mathcal G,(\theta_i,\rho_i)_{i=1}^n\subset\mcF_0\).  Let \(Y_i\) be
self-adjoint and \(\mcF_i\)-measurable with
\[
  \mathbb E(Y_i\mid\mcF_{i-1})=0,\qquad
  \|Y_i\|_{\op}\le\rho_i,\qquad
  \mathbb E(Y_i^2\mid\mcF_{i-1})\preceq \theta_i\rho_i^2 I ,
\]
and set \(X_i=B_i+Y_i\).  Define
\[
  \mu:=\sum_{i=1}^n\mathbb E(\theta_i\rho_i^2),
  \qquad
  \sigma^2:=\sum_{i=1}^n\mathbb E(\theta_i^2\rho_i^2).
\]
Then, for every \(u\ge0\) and \(s>0\),
\[
\begin{aligned}
  \mathbb P\left\{
  \disc(X_1,\ldots,X_n)>K\sqrt n+s
  \right\}
  &\le
  \mathbb P(E)
  +\exp\left\{-\frac{u^2}{2(\sigma^2+\theta_*u/3)}\right\}\\
  &\quad+
  2D\exp\left\{-\frac{s^2}{2(\mu+u+s/3)}\right\},
\end{aligned}
\]
with the convention that the first exponential term is omitted when
\(\sigma^2+\theta_*u/3=0\).  In particular, if \(\theta_i\le\theta\) almost
surely for all \(i\) and \(A=\sum_i\mathbb E\rho_i^2\), then
\[
\begin{aligned}
  \mathbb P\left\{
  \disc(X_1,\ldots,X_n)>K\sqrt n+s
  \right\}
  &\le
  \mathbb P(E)
  +\exp\left\{-\frac{u^2}{2(\theta^2A+\theta u/3)}\right\}\\
  &\quad+
  2D\exp\left\{-\frac{s^2}{2(\theta A+u+s/3)}\right\}.
\end{aligned}
\]
This covers bounded matrix martingale noise whose conditional second moment is
scalar-dominated, without requiring a coordinate, block, or projection support
model.
\end{corollary}

\begin{proof}
The variance-domination assumption gives
\[
  \sum_i\mathbb E(Y_i^2\mid\mcF_{i-1})
  \preceq
  \left(\sum_i\theta_i\rho_i^2\right)I .
\]
Thus the variance parameter \(\nu\) in
Corollary~\ref{cor:core-plus-matrix-martingale-noise} is bounded by
\(W:=\sum_i\theta_i\rho_i^2\).  Applying that corollary with \(v=\mu+u\)
reduces the exceptional term to \(\mathbb P\{W>\mu+u\}\).  For
\(Z_i:=\theta_i\rho_i^2\), independence gives
\[
\begin{aligned}
  0\le Z_i&\le\theta_*,
  \qquad
  \sum_i\mathbb E Z_i=\mu,\\
  \sum_i\operatorname{Var}(Z_i)
  &\le\sum_i\mathbb E Z_i^2
  \le\sum_i\mathbb E(\theta_i^2\rho_i^2)=\sigma^2 .
\end{aligned}
\]
Scalar Bernstein gives the budget-exception bound.  The homogeneous display
follows from \(\mu\le\theta A\), \(\sigma^2\le\theta^2A\), and
\(\theta_*\le\theta\).
\end{proof}

\begin{corollary}[Matrix-profile dominated martingale noise with random amplitudes]
\label{cor:core-plus-matrix-profile-variance-random-amplitude-noise}
Let \(D\ge2\).  Let \(\mathcal G\), \(B_i\), and the event \(E\in\mathcal G\)
be as in Corollary~\ref{cor:core-plus-matrix-martingale-noise}, and assume that
outside \(E\) there are \(\mathcal G\)-measurable signs
\(\eps_i^\star\in\{-1,1\}\) such that
\[
  \left\|\sum_{i=1}^n\eps_i^\star B_i\right\|_{\op}\le K\sqrt n .
\]
Let \((V_i,\rho_i)\), \(i=1,\ldots,n\), be independent pairs, independent of
\(\mathcal G\), where \(V_i\in\M_D(\C)\) is positive semidefinite,
\(\rho_i\in[0,1]\), and
\[
  0\preceq V_i\preceq \tau_* I
\]
almost surely for a deterministic \(\tau_*\in[0,1]\).  Suppose that
\((\mcF_i)_{i=0}^n\) is a filtration with
\(\mathcal G,(V_i,\rho_i)_{i=1}^n\subset\mcF_0\).  Let \(Y_i\) be self-adjoint
and \(\mcF_i\)-measurable with
\[
  \mathbb E(Y_i\mid\mcF_{i-1})=0,\qquad
  \|Y_i\|_{\op}\le\rho_i,\qquad
  \mathbb E(Y_i^2\mid\mcF_{i-1})\preceq \rho_i^2 V_i ,
\]
and set \(X_i=B_i+Y_i\).  Define
\[
  \mu:=\left\|\sum_{i=1}^n\mathbb E(\rho_i^2V_i)\right\|_{\op},
  \qquad
  \omega^2:=\left\|\sum_{i=1}^n\mathbb E(\rho_i^4V_i^2)\right\|_{\op}.
\]
Then, for every \(u\ge0\) and \(s>0\),
\[
\begin{aligned}
  \mathbb P\left\{
  \disc(X_1,\ldots,X_n)>K\sqrt n+s
  \right\}
  &\le
  \mathbb P(E)
  +D\exp\left\{-\frac{u^2}{2(\omega^2+\tau_*u/3)}\right\}\\
  &\quad+
  2D\exp\left\{-\frac{s^2}{2(\mu+u+s/3)}\right\},
\end{aligned}
\]
with the convention that the first exponential term is omitted when
\(\omega^2+\tau_*u/3=0\).  In particular, if \(V_i\preceq\tau I\) almost
surely for all \(i\) and \(A=\sum_i\mathbb E\rho_i^2\), then
\[
\begin{aligned}
  \mathbb P\left\{
  \disc(X_1,\ldots,X_n)>K\sqrt n+s
  \right\}
  &\le
  \mathbb P(E)
  +D\exp\left\{-\frac{u^2}{2(\tau^2A+\tau u/3)}\right\}\\
  &\quad+
  2D\exp\left\{-\frac{s^2}{2(\tau A+u+s/3)}\right\}.
\end{aligned}
\]
This covers anisotropic random conditional-variance profiles, random covariance
mixtures, and leverage-profile matrix martingale noise, without asking the
variance profile to be scalar.
\end{corollary}

\begin{proof}
The variance-domination assumption gives
\[
  \sum_i\mathbb E(Y_i^2\mid\mcF_{i-1})
  \preceq
  \sum_i\rho_i^2V_i .
\]
Hence the variation parameter \(\nu\) in
Corollary~\ref{cor:core-plus-matrix-martingale-noise} is bounded by
\[
  W:=\left\|\sum_i\rho_i^2V_i\right\|_{\op}.
\]
Applying that corollary with \(v=\mu+u\) leaves only the budget exception
\(\mathbb P\{W>\mu+u\}\).  Put \(Z_i:=\rho_i^2V_i\).  Then the \(Z_i\) are
independent positive semidefinite matrices, \(0\preceq Z_i\preceq\tau_*I\), and
\[
  \left\|\sum_i\mathbb E Z_i\right\|_{\op}=\mu .
\]
Moreover, the one-sided matrix Bernstein variance proxy for
\(Z_i-\mathbb E Z_i\) is at most
\[
  \left\|\sum_i\mathbb E(Z_i-\mathbb E Z_i)^2\right\|_{\op}
  \le
  \left\|\sum_i\mathbb E Z_i^2\right\|_{\op}
  =
  \omega^2 .
\]
Tropp's matrix Bernstein inequality gives
\[
  \mathbb P\{W>\mu+u\}
  \le
  D\exp\left\{-\frac{u^2}{2(\omega^2+\tau_*u/3)}\right\}.
\]
For the homogeneous display, use
\[
  \mu\le\tau A,\qquad
  \omega^2\le\tau^2A,
\]
and take the deterministic jump bound to be \(\tau_*=\tau\).
\end{proof}

\begin{corollary}[Trace-budget dominated martingale noise with random amplitudes]
\label{cor:core-plus-trace-budget-variance-random-amplitude-noise}
Let \(D\ge2\).  Let \(\mathcal G\), \(B_i\), and the event \(E\in\mathcal G\)
be as in Corollary~\ref{cor:core-plus-matrix-martingale-noise}, and assume that
outside \(E\) there are \(\mathcal G\)-measurable signs
\(\eps_i^\star\in\{-1,1\}\) such that
\[
  \left\|\sum_{i=1}^n\eps_i^\star B_i\right\|_{\op}\le K\sqrt n .
\]
Let \((V_i,\rho_i)\), \(i=1,\ldots,n\), be independent pairs, independent of
\(\mathcal G\), where \(V_i\in\M_D(\C)\) is positive semidefinite and
\(\rho_i\in[0,1]\).  Assume that
\[
  \operatorname{tr}V_i\le \gamma_*
\]
almost surely for a deterministic \(\gamma_*\ge0\).  Suppose that
\((\mcF_i)_{i=0}^n\) is a filtration with
\(\mathcal G,(V_i,\rho_i)_{i=1}^n\subset\mcF_0\).  Let \(Y_i\) be self-adjoint
and \(\mcF_i\)-measurable with
\[
  \mathbb E(Y_i\mid\mcF_{i-1})=0,\qquad
  \|Y_i\|_{\op}\le\rho_i,\qquad
  \mathbb E(Y_i^2\mid\mcF_{i-1})\preceq \rho_i^2 V_i ,
\]
and set \(X_i=B_i+Y_i\).  Define
\[
  \mu_{\operatorname{tr}}:=
  \sum_{i=1}^n\mathbb E(\rho_i^2\operatorname{tr}V_i),
  \qquad
  \sigma_{\operatorname{tr}}^2:=
  \sum_{i=1}^n\mathbb E(\rho_i^4(\operatorname{tr}V_i)^2).
\]
Then, for every \(u\ge0\) and \(s>0\),
\[
\begin{aligned}
  \mathbb P\left\{
  \disc(X_1,\ldots,X_n)>K\sqrt n+s
  \right\}
  &\le
  \mathbb P(E)
  +\exp\left\{-\frac{u^2}{2(\sigma_{\operatorname{tr}}^2+\gamma_*u/3)}\right\}\\
  &\quad+
  2D\exp\left\{-\frac{s^2}{2(\mu_{\operatorname{tr}}+u+s/3)}\right\},
\end{aligned}
\]
with the convention that the first exponential term is omitted when
\(\sigma_{\operatorname{tr}}^2+\gamma_*u/3=0\).  In particular, if
\(\operatorname{tr}V_i\le\gamma\) almost surely for all \(i\) and
\(A=\sum_i\mathbb E\rho_i^2\), then
\[
\begin{aligned}
  \mathbb P\left\{
  \disc(X_1,\ldots,X_n)>K\sqrt n+s
  \right\}
  &\le
  \mathbb P(E)
  +\exp\left\{-\frac{u^2}{2(\gamma^2A+\gamma u/3)}\right\}\\
  &\quad+
  2D\exp\left\{-\frac{s^2}{2(\gamma A+u+s/3)}\right\}.
\end{aligned}
\]
Thus low-total-trace random variance profiles give a dimension-free budget
exception; the price is that the Freedman variance parameter is controlled by
the trace budget rather than the sharper operator-norm budget.
\end{corollary}

\begin{proof}
As in the matrix-profile corollary,
\[
  \sum_i\mathbb E(Y_i^2\mid\mcF_{i-1})
  \preceq
  \sum_i\rho_i^2V_i .
\]
Because the right-hand side is positive semidefinite,
\[
  \left\|\sum_i\rho_i^2V_i\right\|_{\op}
  \le
  \operatorname{tr}\left(\sum_i\rho_i^2V_i\right)
  =
  \sum_i\rho_i^2\operatorname{tr}V_i .
\]
Hence Corollary~\ref{cor:core-plus-matrix-martingale-noise} applies with the
random scalar budget
\[
  T:=\sum_i\rho_i^2\operatorname{tr}V_i .
\]
Taking \(v=\mu_{\operatorname{tr}}+u\), it remains to bound
\(\mathbb P\{T>\mu_{\operatorname{tr}}+u\}\).  For
\(Z_i:=\rho_i^2\operatorname{tr}V_i\), the variables are independent and satisfy
\[
\begin{aligned}
  0\le Z_i&\le\gamma_*,
  \qquad
  \sum_i\mathbb E Z_i=\mu_{\operatorname{tr}},\\
  \sum_i\operatorname{Var}(Z_i)
  &\le\sum_i\mathbb E Z_i^2
  =\sum_i\mathbb E(\rho_i^4(\operatorname{tr}V_i)^2)
  =\sigma_{\operatorname{tr}}^2 .
\end{aligned}
\]
Scalar Bernstein gives the first display.  The homogeneous display follows from
\(\mu_{\operatorname{tr}}\le\gamma A\) and
\(\sigma_{\operatorname{tr}}^2\le\gamma^2A\).
\end{proof}

\begin{corollary}[Trace-envelope and low-rank variance profiles]
\label{cor:core-plus-trace-envelope-low-rank-random-amplitude-noise}
Let \(D\ge2\).  Let \(\mathcal G\), \(B_i\), and the event \(E\in\mathcal G\)
be as in Corollary~\ref{cor:core-plus-matrix-martingale-noise}, and assume that
outside \(E\) there are \(\mathcal G\)-measurable signs
\(\eps_i^\star\in\{-1,1\}\) such that
\[
  \left\|\sum_{i=1}^n\eps_i^\star B_i\right\|_{\op}\le K\sqrt n .
\]
Let \((V_i,\Gamma_i,\rho_i)\), \(i=1,\ldots,n\), be independent triples,
independent of \(\mathcal G\), where \(V_i\in\M_D(\C)\) is positive
semidefinite, \(\rho_i\in[0,1]\), and
\[
  0\le \Gamma_i\le \gamma_*,
  \qquad
  \operatorname{tr}V_i\le \Gamma_i
\]
almost surely for a deterministic \(\gamma_*\ge0\).  Suppose that
\((\mcF_i)_{i=0}^n\) is a filtration with
\(\mathcal G,(V_i,\Gamma_i,\rho_i)_{i=1}^n\subset\mcF_0\).  Let \(Y_i\) be
self-adjoint and \(\mcF_i\)-measurable with
\[
  \mathbb E(Y_i\mid\mcF_{i-1})=0,\qquad
  \|Y_i\|_{\op}\le\rho_i,\qquad
  \mathbb E(Y_i^2\mid\mcF_{i-1})\preceq \rho_i^2 V_i ,
\]
and set \(X_i=B_i+Y_i\).  Define
\[
  \mu_{\Gamma}:=\sum_{i=1}^n\mathbb E(\rho_i^2\Gamma_i),
  \qquad
  \sigma_{\Gamma}^2:=\sum_{i=1}^n\mathbb E(\rho_i^4\Gamma_i^2).
\]
Then, for every \(u\ge0\) and \(s>0\),
\[
\begin{aligned}
  \mathbb P\left\{
  \disc(X_1,\ldots,X_n)>K\sqrt n+s
  \right\}
  &\le
  \mathbb P(E)
  +\exp\left\{-\frac{u^2}{2(\sigma_{\Gamma}^2+\gamma_*u/3)}\right\}\\
  &\quad+
  2D\exp\left\{-\frac{s^2}{2(\mu_{\Gamma}+u+s/3)}\right\},
\end{aligned}
\]
with the convention that the first exponential term is omitted when
\(\sigma_{\Gamma}^2+\gamma_*u/3=0\).

In particular, suppose that the envelope is supplied by random low-rank data:
there are independent tuples \((V_i,R_i,\tau_i,\rho_i)\), independent of
\(\mathcal G\), such that \(V_i\succeq0\),
\[
  \operatorname{rank}V_i\le R_i,\qquad
  \|V_i\|_{\op}\le \tau_i,\qquad
  R_i\tau_i\le\gamma_*
\]
almost surely.  Then the same conclusion holds with
\[
  \mu_{\operatorname{lr}}
  :=\sum_{i=1}^n\mathbb E(R_i\tau_i\rho_i^2),
  \qquad
  \sigma_{\operatorname{lr}}^2
  :=\sum_{i=1}^n\mathbb E(R_i^2\tau_i^2\rho_i^4)
\]
in place of \(\mu_\Gamma,\sigma_\Gamma^2\).  If moreover
\(R_i\le r\) and \(\tau_i\le\tau\) almost surely for every \(i\), and
\(A=\sum_i\mathbb E\rho_i^2\), then
\[
\begin{aligned}
  \mathbb P\left\{
  \disc(X_1,\ldots,X_n)>K\sqrt n+s
  \right\}
  &\le
  \mathbb P(E)
  +\exp\left\{-\frac{u^2}{2(r^2\tau^2A+r\tau u/3)}\right\}\\
  &\quad+
  2D\exp\left\{-\frac{s^2}{2(r\tau A+u+s/3)}\right\}.
\end{aligned}
\]
Thus trace-enveloped low-rank and spiked variance profiles are controlled by
their scalar envelope rather than by the ambient dimension.
\end{corollary}

\begin{proof}
As in the proof of
Corollary~\ref{cor:core-plus-trace-budget-variance-random-amplitude-noise}, the
variation in Corollary~\ref{cor:core-plus-matrix-martingale-noise} is bounded by
the random scalar trace budget.  Since \(\operatorname{tr}V_i\le\Gamma_i\), that
budget satisfies
\[
  T=\sum_i\rho_i^2\operatorname{tr}V_i
  \le
  \sum_i\rho_i^2\Gamma_i .
\]
For \(Z_i:=\rho_i^2\Gamma_i\), the variables are independent,
\(0\le Z_i\le\gamma_*\), and
\[
  \sum_i\mathbb E Z_i=\mu_\Gamma,\qquad
  \sum_i\operatorname{Var}(Z_i)\le
  \sum_i\mathbb E Z_i^2=\sigma_\Gamma^2 .
\]
Scalar Bernstein therefore gives
\[
  \mathbb P\{T>\mu_\Gamma+u\}
  \le
  \exp\left\{-\frac{u^2}{2(\sigma_\Gamma^2+\gamma_*u/3)}\right\}.
\]
Combining this with Corollary~\ref{cor:core-plus-matrix-martingale-noise} at
\(v=\mu_\Gamma+u\) proves the trace-envelope display.

For the low-rank specialization, take \(\Gamma_i=R_i\tau_i\).  The trace bound
follows from
\[
  \operatorname{tr}V_i
  \le
  \operatorname{rank}(V_i)\|V_i\|_{\op}
  \le R_i\tau_i .
\]
The last display uses
\[
  \mu_{\operatorname{lr}}\le r\tau A,
  \qquad
  \sigma_{\operatorname{lr}}^2\le r^2\tau^2A,
  \qquad
  \gamma_*=r\tau,
\]
where \(\rho_i^4\le\rho_i^2\).
\end{proof}

\begin{corollary}[Schatten-envelope variance profiles with random amplitudes]
\label{cor:core-plus-schatten-envelope-random-amplitude-noise}
Let \(1\le p\le\infty\), and put \(\beta=1-1/p\), with the convention
\(\beta=1\) when \(p=\infty\).  Keep the signed-core, filtration, and
martingale-noise hypotheses of
Corollary~\ref{cor:core-plus-trace-envelope-low-rank-random-amplitude-noise},
but replace the scalar trace-envelope assumption by the following Schatten
envelope.  Let \((V_i,R_i,L_i,\rho_i)\), \(i=1,\ldots,n\), be independent
tuples, independent of \(\mathcal G\), where \(V_i\succeq0\),
\(\rho_i\in[0,1]\), \(1\le R_i\le D\), and
\[
  \operatorname{rank}V_i\le R_i,\qquad
  \|V_i\|_{S_p}\le L_i,\qquad
  R_i^\beta L_i\le\gamma_*
\]
almost surely.  Define
\[
  \mu_p:=\sum_{i=1}^n\mathbb E(R_i^\beta L_i\rho_i^2),
  \qquad
  \sigma_p^2:=\sum_{i=1}^n\mathbb E(R_i^{2\beta}L_i^2\rho_i^4).
\]
Then, for every \(u\ge0\) and \(s>0\),
\[
\begin{aligned}
  \mathbb P\left\{
  \disc(X_1,\ldots,X_n)>K\sqrt n+s
  \right\}
  &\le
  \mathbb P(E)
  +\exp\left\{-\frac{u^2}{2(\sigma_p^2+\gamma_*u/3)}\right\}\\
  &\quad+
  2D\exp\left\{-\frac{s^2}{2(\mu_p+u+s/3)}\right\}.
\end{aligned}
\]
In particular, if \(R_i\le r\), \(L_i\le L\) almost surely, and
\(A=\sum_i\mathbb E\rho_i^2\), then
\[
\begin{aligned}
  \mathbb P\left\{
  \disc(X_1,\ldots,X_n)>K\sqrt n+s
  \right\}
  &\le
  \mathbb P(E)
  +\exp\left\{-\frac{u^2}{2(r^{2\beta}L^2A+r^\beta Lu/3)}\right\}\\
  &\quad+
  2D\exp\left\{-\frac{s^2}{2(r^\beta LA+u+s/3)}\right\}.
\end{aligned}
\]
Thus trace control (\(p=1\)), Frobenius-plus-rank control (\(p=2\)), and
operator-norm-plus-rank control (\(p=\infty\)) are all instances of the same
signed-core random variance-profile theorem.
\end{corollary}

\begin{proof}
For a positive semidefinite matrix \(V_i\) with at most \(R_i\) nonzero
singular values, Holder's inequality for the singular-value vector gives
\[
  \operatorname{tr}V_i=\|V_i\|_{S_1}
  \le R_i^{1-1/p}\|V_i\|_{S_p}
  \le R_i^\beta L_i .
\]
Apply
Corollary~\ref{cor:core-plus-trace-envelope-low-rank-random-amplitude-noise}
with \(\Gamma_i=R_i^\beta L_i\).  The homogeneous display follows from
\(\rho_i^4\le\rho_i^2\).
\end{proof}

\begin{corollary}[Quantile trace-envelope variance profiles]
\label{cor:core-plus-quantile-trace-envelope-random-amplitude-noise}
Let \(D\ge2\).  Keep the signed-core, filtration, and martingale-noise
hypotheses of
Corollary~\ref{cor:core-plus-trace-envelope-low-rank-random-amplitude-noise},
but replace the boundedness and independence assumptions on the trace
envelopes by the following confidence-envelope assumption.  The triples
\((V_i,\Gamma_i,\rho_i)\) are \(\mcF_0\)-measurable, \(V_i\succeq0\),
\(\rho_i\in[0,1]\), \(\Gamma_i\ge0\), and
\[
  \operatorname{tr}V_i\le\Gamma_i
  \qquad\text{almost surely}.
\]
Let \(S\) be any nonnegative random variable such that
\[
  \sum_{i=1}^n\rho_i^2\Gamma_i\le S
  \qquad\text{almost surely}.
\]
Then, for every deterministic \(v\ge0\) and \(s>0\),
\[
\begin{aligned}
  \mathbb P\left\{
  \disc(X_1,\ldots,X_n)>K\sqrt n+s
  \right\}
  &\le
  \mathbb P(E)+\mathbb P\{S>v\}\\
  &\quad+
  2D\exp\left\{-\frac{s^2}{2(v+s/3)}\right\}.
\end{aligned}
\]
Consequently, if \(q_\delta\) is any deterministic upper
\(\delta\)-quantile of \(S\), meaning
\(\mathbb P\{S>q_\delta\}\le\delta\), then
\[
\begin{aligned}
  \mathbb P\left\{
  \disc(X_1,\ldots,X_n)>K\sqrt n+s
  \right\}
  &\le
  \mathbb P(E)+\delta\\
  &\quad+
  2D\exp\left\{-\frac{s^2}{2(q_\delta+s/3)}\right\}.
\end{aligned}
\]
If, more concretely, \(\mathbb E S^r\le M_r<\infty\) for some \(r>0\), then
for every \(0<\delta\le1\),
\[
\begin{aligned}
  \mathbb P\left\{
  \disc(X_1,\ldots,X_n)>K\sqrt n+s
  \right\}
  &\le
  \mathbb P(E)+\delta\\
  &\quad+
  2D\exp\left\{-\frac{s^2}
  {2((M_r/\delta)^{1/r}+s/3)}\right\}.
\end{aligned}
\]
Thus heavy-tailed trace-class random covariance, spike, low-rank, and Schatten
variance profiles still admit the same signed-core Matrix Spencer reduction at
any confidence level for which the total weighted trace envelope has a usable
upper quantile.  In the Schatten setting one may take
\(S=\sum_i\rho_i^2R_i^{1-1/p}L_i\) whenever
\(\operatorname{rank}V_i\le R_i\) and \(\|V_i\|_{S_p}\le L_i\); no independence
between these envelope data is needed for this quantile form.
\end{corollary}

\begin{proof}
The proof of
Corollary~\ref{cor:core-plus-trace-envelope-low-rank-random-amplitude-noise}
uses only the trace domination to bound the predictable variation in
Corollary~\ref{cor:core-plus-matrix-martingale-noise}.  Here the same argument
gives
\[
  \nu
  =
  \left\|\sum_i\mathbb E(Y_i^2\mid\mcF_{i-1})\right\|_{\op}
  \le
  \sum_i\rho_i^2\operatorname{tr}V_i
  \le
  \sum_i\rho_i^2\Gamma_i
  \le S .
\]
Applying Corollary~\ref{cor:core-plus-matrix-martingale-noise} with the
deterministic parameter \(v\) and using
\(\{\nu>v\}\subseteq\{S>v\}\) proves the first display.  The quantile form is
the same estimate with \(v=q_\delta\).  The moment form follows from Markov's
inequality,
\[
  \mathbb P\{S>(M_r/\delta)^{1/r}\}\le\delta,
\]
with the evident interpretation when \(M_r=0\).  The Schatten sentence follows
from Holder's inequality for the singular-value vector, as in
Corollary~\ref{cor:core-plus-schatten-envelope-random-amplitude-noise}.
\end{proof}

\begin{corollary}[Almost-sure quantile trace-envelope families]
\label{cor:core-plus-quantile-trace-envelope-borel-random-amplitude-noise}
For each \(m\ge1\), let
\[
  X_{m,1},\ldots,X_{m,n_m}\in\M_{D_m}(\C),\qquad D_m\ge2,
\]
be a random instance satisfying the hypotheses of
Corollary~\ref{cor:core-plus-quantile-trace-envelope-random-amplitude-noise},
with parameters \(K_m\), exceptional event \(E_m\), and dominating trace budget
\(S_m\).  Let \(q_m,\delta_m,s_m\) be deterministic numbers with
\(q_m,s_m\ge0\), \(0\le\delta_m\le1\), and
\[
  \mathbb P\{S_m>q_m\}\le\delta_m .
\]
If
\[
  \sum_{m=1}^\infty
  \left[
    \mathbb P(E_m)+\delta_m
    +2D_m\exp\left\{-\frac{s_m^2}{2(q_m+s_m/3)}\right\}
  \right]<\infty,
\]
then, almost surely, for all sufficiently large \(m\),
\[
  \disc(X_{m,1},\ldots,X_{m,n_m})
  \le K_m\sqrt{n_m}+s_m .
\]

In particular, put \(\ell_m=\log(eD_m)\).  Suppose that \(K_m\le K\), that
\[
  q_m\le Q\frac{n_m}{\ell_m},
  \qquad
  \ell_m\le B\sqrt{n_m}
\]
for fixed \(Q,B<\infty\), and that for some \(A>0\),
\[
  \sum_{m=1}^\infty
  \left[\mathbb P(E_m)+\delta_m+(eD_m)^{-A}\right]<\infty .
\]
Choose \(L>0\) so that
\[
  \frac{L^2}{2(Q+LB/3)}\ge A+1 .
\]
Then almost surely, for all sufficiently large \(m\),
\[
  \disc(X_{m,1},\ldots,X_{m,n_m})
  \le (K+L)\sqrt{n_m}.
\]
The same criterion may be checked from moments: if
\(\mathbb E S_m^r\le M_{m,r}\), one may take
\(q_m=(M_{m,r}/\delta_m)^{1/r}\).  Thus summable quantile or finite-moment
control of heavy-tailed trace budgets at scale \(n_m/\log(eD_m)\) gives
eventual Matrix Spencer bounds for triangular arrays, with no independence
assumption across the dimensions \(m\).
\end{corollary}

\begin{proof}
Corollary~\ref{cor:core-plus-quantile-trace-envelope-random-amplitude-noise}
gives, for each \(m\), with
\(\Delta_m=\disc(X_{m,1},\ldots,X_{m,n_m})\),
\[
\begin{aligned}
  \mathbb P\{\Delta_m>K_m\sqrt{n_m}+s_m\}
  &\le
  \mathbb P(E_m)+\delta_m\\
  &\quad+
  2D_m\exp\left\{-\frac{s_m^2}{2(q_m+s_m/3)}\right\}.
\end{aligned}
\]
The displayed summability hypothesis and the first Borel--Cantelli lemma give
the almost-sure eventual estimate.

For the concrete profile, take \(s_m=L\sqrt{n_m}\).  Since
\(\ell_m\le B\sqrt{n_m}\), we have
\[
  L\sqrt{n_m}/3\le (LB/3)\,n_m/\ell_m .
\]
Together with \(q_m\le Qn_m/\ell_m\), this gives
\[
  \frac{s_m^2}{2(q_m+s_m/3)}
  \ge
  \frac{L^2}{2(Q+LB/3)}\,\ell_m
  \ge (A+1)\ell_m .
\]
Therefore
\[
  2D_m\exp\left\{-\frac{s_m^2}{2(q_m+s_m/3)}\right\}
  \le 2D_m(eD_m)^{-(A+1)}
  \le 2(eD_m)^{-A}.
\]
The assumed summability and the first part apply.  The moment assertion is
Markov's inequality, exactly as in
Corollary~\ref{cor:core-plus-quantile-trace-envelope-random-amplitude-noise}.
\end{proof}

\begin{corollary}[Moment-summable heavy-tail trace-envelope families]
\label{cor:core-plus-moment-trace-envelope-borel-random-amplitude-noise}
In the setting of
Corollary~\ref{cor:core-plus-quantile-trace-envelope-borel-random-amplitude-noise},
put \(\ell_m=\log(eD_m)\).  Suppose that \(K_m\le K\) and
\(\ell_m\le B\sqrt{n_m}\) for fixed \(K,B<\infty\).  Let \(r>0\), \(Q>0\),
and \(A>0\).  If
\[
  \sum_{m=1}^\infty
  \left[
    \mathbb P(E_m)
    +\mathbb E\left(\frac{\ell_m S_m}{Qn_m}\right)^r
    +(eD_m)^{-A}
  \right]<\infty,
\]
then, for every \(L>0\) satisfying
\[
  \frac{L^2}{2(Q+LB/3)}\ge A+1,
\]
almost surely, for all sufficiently large \(m\),
\[
  \disc(X_{m,1},\ldots,X_{m,n_m})
  \le (K+L)\sqrt{n_m}.
\]
Equivalently, it is enough that
\[
  \mathbb E S_m^r
  \le
  Q^r\left(\frac{n_m}{\ell_m}\right)^r a_m
  \qquad\text{with}\qquad
  \sum_m a_m<\infty,
\]
along with the summability of \(\mathbb P(E_m)\) and \((eD_m)^{-A}\).  Thus
heavy-tailed trace envelopes whose normalized \(r\)-th moments are summable
still give eventual \(O(\sqrt{n_m})\) Matrix Spencer bounds; no exponential tail
or independence across the dimensions is required.
\end{corollary}

\begin{proof}
Set
\[
  q_m=Q\frac{n_m}{\ell_m},
  \qquad
  \delta_m=\mathbb E\left(\frac{\ell_m S_m}{Qn_m}\right)^r .
\]
Markov's inequality gives
\[
  \mathbb P\{S_m>q_m\}\le\delta_m .
\]
The displayed moment summability is exactly the summability hypothesis in the
concrete profile of
Corollary~\ref{cor:core-plus-quantile-trace-envelope-borel-random-amplitude-noise}.
The stated almost-sure conclusion follows.  The equivalent form is just the
same criterion after bounding the normalized moment by \(a_m\).
\end{proof}

\begin{corollary}[Termwise-moment heavy-tail trace-envelope families]
\label{cor:core-plus-termwise-moment-trace-envelope-borel-random-amplitude-noise}
In the setting of
Corollary~\ref{cor:core-plus-moment-trace-envelope-borel-random-amplitude-noise},
suppose that, for each \(m\), the dominating trace budget has an almost-sure
termwise bound
\[
  S_m\le \sum_{i=1}^{n_m} W_{m,i},
  \qquad W_{m,i}\ge0 .
\]
This includes the basic trace-envelope choice
\(W_{m,i}=\rho_{m,i}^2\Gamma_{m,i}\), and in the Schatten-envelope case the
choice \(W_{m,i}=\rho_{m,i}^2R_{m,i}^{1-1/p}L_{m,i}\).  Assume that
\(K_m\le K\) and \(\ell_m=\log(eD_m)\le B\sqrt{n_m}\), and fix \(Q,A>0\).
If, for some \(0<r\le1\),
\[
  \sum_{m=1}^\infty
  \left[
    \mathbb P(E_m)
    +\left(\frac{\ell_m}{Qn_m}\right)^r
      \sum_{i=1}^{n_m}\mathbb E W_{m,i}^r
    +(eD_m)^{-A}
  \right]<\infty,
\]
then, for every \(L>0\) satisfying
\[
  \frac{L^2}{2(Q+LB/3)}\ge A+1,
\]
almost surely, for all sufficiently large \(m\),
\[
  \disc(X_{m,1},\ldots,X_{m,n_m})
  \le (K+L)\sqrt{n_m}.
\]
The same conclusion holds for \(r\ge1\) under the alternative summability
condition
\[
  \sum_{m=1}^\infty
  \left[
    \mathbb P(E_m)
    +\left(\frac{\ell_m}{Qn_m}\right)^r
      \left(\sum_{i=1}^{n_m}(\mathbb E W_{m,i}^r)^{1/r}\right)^r
    +(eD_m)^{-A}
  \right]<\infty .
\]
Thus random spike, low-rank, trace-class, and Schatten variance profiles with
only termwise heavy-tail moments still give almost-sure eventual Matrix
Spencer bounds, without any independence among the terms \(W_{m,i}\).
\end{corollary}

\begin{proof}
Put \(\widetilde S_m=\sum_iW_{m,i}\).  Since \(S_m\le\widetilde S_m\), it is
enough to verify the normalized-moment hypothesis of
Corollary~\ref{cor:core-plus-moment-trace-envelope-borel-random-amplitude-noise}
with \(\widetilde S_m\) in place of \(S_m\).  If \(0<r\le1\), then
\((\sum_i w_i)^r\le\sum_i w_i^r\) for nonnegative \(w_i\), so
\[
  \mathbb E\widetilde S_m^r
  \le \sum_{i=1}^{n_m}\mathbb E W_{m,i}^r .
\]
If \(r\ge1\), Minkowski's inequality in \(L^r\) gives
\[
  \left(\mathbb E\widetilde S_m^r\right)^{1/r}
  \le
  \sum_{i=1}^{n_m}\left(\mathbb E W_{m,i}^r\right)^{1/r}.
\]
Each displayed summability condition therefore implies the corresponding
normalized-moment summability for \(\widetilde S_m\), and the preceding
corollary applies.
\end{proof}

\begin{corollary}[Termwise-tail heavy-tail trace-envelope families]
\label{cor:core-plus-termwise-tail-trace-envelope-borel-random-amplitude-noise}
In the setting of
Corollary~\ref{cor:core-plus-termwise-moment-trace-envelope-borel-random-amplitude-noise},
suppose again that
\[
  S_m\le \sum_{i=1}^{n_m}W_{m,i},
  \qquad W_{m,i}\ge0 .
\]
Let \(K_m\le K\), \(\ell_m=\log(eD_m)\le B\sqrt{n_m}\), and fix \(Q,A>0\).
For each \(m\), let \(b_{m,1},\ldots,b_{m,n_m}\ge0\) and
\(\delta_{m,1},\ldots,\delta_{m,n_m}\ge0\) be deterministic numbers such that
\[
  \sum_{i=1}^{n_m} b_{m,i}\le Q\frac{n_m}{\ell_m},
  \qquad
  \mathbb P\{W_{m,i}>b_{m,i}\}\le\delta_{m,i}
  \quad(1\le i\le n_m).
\]
If
\[
  \sum_{m=1}^\infty
  \left[
    \mathbb P(E_m)
    +\sum_{i=1}^{n_m}\delta_{m,i}
    +(eD_m)^{-A}
  \right]<\infty,
\]
then, for every \(L>0\) satisfying
\[
  \frac{L^2}{2(Q+LB/3)}\ge A+1,
\]
almost surely, for all sufficiently large \(m\),
\[
  \disc(X_{m,1},\ldots,X_{m,n_m})
  \le (K+L)\sqrt{n_m}.
\]
Thus it is enough to truncate each heavy-tailed spike, low-rank, trace-class,
or Schatten atom at deterministic levels whose total mass is
\(O(n_m/\log(eD_m))\), with summable total truncation failure; no independence
among the atoms is required.
\end{corollary}

\begin{proof}
Set
\[
  q_m=Q\frac{n_m}{\ell_m},
  \qquad
  \delta_m=\sum_{i=1}^{n_m}\delta_{m,i}.
\]
If \(W_{m,i}\le b_{m,i}\) for every \(i\), then
\[
  S_m\le\sum_i W_{m,i}\le\sum_i b_{m,i}\le q_m .
\]
Therefore, by the union bound,
\[
  \mathbb P\{S_m>q_m\}
  \le
  \sum_{i=1}^{n_m}\mathbb P\{W_{m,i}>b_{m,i}\}
  \le\delta_m .
\]
The concrete profile in
Corollary~\ref{cor:core-plus-quantile-trace-envelope-borel-random-amplitude-noise}
now applies with these \(q_m\) and \(\delta_m\).
\end{proof}

\begin{corollary}[Scale-mixture termwise-tail trace-envelope families]
\label{cor:core-plus-scale-mixture-tail-trace-envelope-borel-random-amplitude-noise}
In the setting of
Corollary~\ref{cor:core-plus-termwise-tail-trace-envelope-borel-random-amplitude-noise},
suppose that
\[
  S_m\le \sum_{i=1}^{n_m}W_{m,i},
  \qquad
  W_{m,i}\le a_{m,i}Z_{m,i}
\]
almost surely, where \(a_{m,i}\ge0\) are deterministic and
\(Z_{m,i}\ge0\) are arbitrary random multipliers.  No independence is assumed
among the \(Z_{m,i}\).  Let \(K_m\le K\),
\(\ell_m=\log(eD_m)\le B\sqrt{n_m}\), and fix \(Q,A>0\).  If there are
deterministic thresholds \(\tau_m\ge0\) such that
\[
  \tau_m\sum_{i=1}^{n_m}a_{m,i}\le Q\frac{n_m}{\ell_m}
\]
and
\[
  \sum_{m=1}^\infty
  \left[
    \mathbb P(E_m)
    +\sum_{i=1}^{n_m}\mathbb P\{Z_{m,i}>\tau_m\}
    +(eD_m)^{-A}
  \right]<\infty,
\]
then, for every \(L>0\) satisfying
\[
  \frac{L^2}{2(Q+LB/3)}\ge A+1,
\]
almost surely, for all sufficiently large \(m\),
\[
  \disc(X_{m,1},\ldots,X_{m,n_m})
  \le (K+L)\sqrt{n_m}.
\]

In particular, put \(A_m=\sum_i a_{m,i}\).  If, for some \(\alpha>0\), the
multipliers satisfy the uniform polynomial tail
\[
  \mathbb P\{Z_{m,i}>t\}\le C_m t^{-\alpha}
  \qquad(t>0,\ 1\le i\le n_m),
\]
then the conclusion follows from
\[
  \sum_{m=1}^\infty\left[\mathbb P(E_m)+(eD_m)^{-A}\right]
  +
  \sum_{m:A_m>0}
  n_mC_m\left(\frac{\ell_m A_m}{Qn_m}\right)^\alpha
  <\infty,
\]
with the zero-scale rows \(A_m=0\) carrying no multiplier-tail term.  Thus
deterministic trace, rank, or Schatten
atom weights may carry arbitrary dependent heavy-tailed multipliers as long as
their aggregate upper tail is summable at the \(n_m/\log(eD_m)\) budget scale.
\end{corollary}

\begin{proof}
Apply
Corollary~\ref{cor:core-plus-termwise-tail-trace-envelope-borel-random-amplitude-noise}
with \(b_{m,i}=\tau_m a_{m,i}\) and
\(\delta_{m,i}=\mathbb P\{Z_{m,i}>\tau_m\}\).  Indeed, if
\(Z_{m,i}\le\tau_m\), then \(W_{m,i}\le b_{m,i}\), while if \(a_{m,i}=0\) the
bound \(W_{m,i}\le0=b_{m,i}\) is automatic.  Hence
\[
  \mathbb P\{W_{m,i}>b_{m,i}\}
  \le \mathbb P\{Z_{m,i}>\tau_m\}.
\]
The mass condition gives \(\sum_i b_{m,i}\le Qn_m/\ell_m\), and the displayed
summability hypothesis is precisely the termwise-tail hypothesis.

For the polynomial-tail specialization, take
\(\tau_m=Qn_m/(\ell_m A_m)\) when \(A_m>0\).  Then
\(\tau_mA_m=Qn_m/\ell_m\), and
\[
  \sum_{i=1}^{n_m}\mathbb P\{Z_{m,i}>\tau_m\}
  \le
  n_mC_m\left(\frac{\ell_m A_m}{Qn_m}\right)^\alpha .
\]
When \(A_m=0\), all \(W_{m,i}\) vanish almost surely, so no multiplier tail
term is needed.
\end{proof}

\begin{corollary}[Common-scale heavy-tail trace-envelope families]
\label{cor:core-plus-common-scale-tail-trace-envelope-borel-random-amplitude-noise}
In the setting of
Corollary~\ref{cor:core-plus-quantile-trace-envelope-borel-random-amplitude-noise},
put \(\ell_m=\log(eD_m)\).  Suppose that \(K_m\le K\) and
\(\ell_m\le B\sqrt{n_m}\) for fixed \(K,B<\infty\).  Let \(Q,A>0\).
Assume that there are deterministic weights \(a_{m,i}\ge0\) and a nonnegative
random multiplier \(Z_m\) such that, with \(A_m=\sum_i a_{m,i}\),
\[
  S_m\le Z_m A_m
  \qquad\text{almost surely}.
\]
This includes the atomwise situation \(S_m\le\sum_iW_{m,i}\) with
\(W_{m,i}\le a_{m,i}Z_m\), where the same random scale multiplies every atom.
If
\[
  \sum_{m=1}^\infty\left[\mathbb P(E_m)+(eD_m)^{-A}\right]
  +
  \sum_{m:A_m>0}
  \mathbb P\left\{Z_m>Q\frac{n_m}{\ell_m A_m}\right\}
  <\infty,
\]
then, for every \(L>0\) satisfying
\[
  \frac{L^2}{2(Q+LB/3)}\ge A+1,
\]
almost surely, for all sufficiently large \(m\),
\[
  \disc(X_{m,1},\ldots,X_{m,n_m})
  \le (K+L)\sqrt{n_m}.
\]

In particular, if for some \(\alpha>0\)
\[
  \mathbb P\{Z_m>t\}\le C_m t^{-\alpha}\qquad(t>0),
\]
then the conclusion follows from
\[
  \sum_{m=1}^\infty\left[\mathbb P(E_m)+(eD_m)^{-A}\right]
  +
  \sum_{m:A_m>0}
  C_m\left(\frac{\ell_m A_m}{Qn_m}\right)^\alpha
  <\infty .
\]
Thus global random volatility, elliptical covariance scales, and other
common-multiplier trace-envelope models pay only the tail of the shared scale,
not a union bound over all atoms.
\end{corollary}

\begin{proof}
Put \(q_m=Qn_m/\ell_m\).  If \(A_m=0\), then \(S_m=0\) almost surely.  If
\(A_m>0\), then
\[
  \mathbb P\{S_m>q_m\}
  \le
  \mathbb P\{Z_mA_m>q_m\}
  =
  \mathbb P\left\{Z_m>Q\frac{n_m}{\ell_m A_m}\right\}.
\]
Thus the concrete profile in
Corollary~\ref{cor:core-plus-quantile-trace-envelope-borel-random-amplitude-noise}
applies with these deterministic quantiles.  The polynomial tail
specialization is the same estimate after inserting
\(t=Qn_m/(\ell_m A_m)\) for rows with \(A_m>0\).
\end{proof}

\begin{corollary}[Grouped-scale heavy-tail trace-envelope families]
\label{cor:core-plus-grouped-scale-tail-trace-envelope-borel-random-amplitude-noise}
In the setting of
Corollary~\ref{cor:core-plus-quantile-trace-envelope-borel-random-amplitude-noise},
put \(\ell_m=\log(eD_m)\).  Suppose that \(K_m\le K\) and
\(\ell_m\le B\sqrt{n_m}\) for fixed \(K,B<\infty\).  Let \(Q,A>0\).
For each \(m\), let \(G_m\) be a finite set, let
\(A_{m,g}\ge0\) be deterministic group masses, and let
\(Z_{m,g}\ge0\) be arbitrary random group multipliers.  Assume that
\[
  S_m\le\sum_{g\in G_m} A_{m,g}Z_{m,g}
  \qquad\text{almost surely}.
\]
If there are deterministic thresholds \(\tau_{m,g}\ge0\) such that
\[
  \sum_{g\in G_m}\tau_{m,g}A_{m,g}
  \le Q\frac{n_m}{\ell_m}
\]
and
\[
  \sum_{m=1}^\infty
  \left[
    \mathbb P(E_m)
    +\sum_{g\in G_m}\mathbb P\{Z_{m,g}>\tau_{m,g}\}
    +(eD_m)^{-A}
  \right]<\infty,
\]
then, for every \(L>0\) satisfying
\[
  \frac{L^2}{2(Q+LB/3)}\ge A+1,
\]
almost surely, for all sufficiently large \(m\),
\[
  \disc(X_{m,1},\ldots,X_{m,n_m})
  \le (K+L)\sqrt{n_m}.
\]

In particular, suppose that, for some \(\alpha>0\),
\[
  \mathbb P\{Z_{m,g}>t\}\le C_{m,g}t^{-\alpha}
  \qquad(t>0,\ g\in G_m),
\]
and that there are deterministic budget shares
\(\lambda_{m,g}>0\) for every \(g\) with \(A_{m,g}>0\), satisfying
\(\sum_{g:A_{m,g}>0}\lambda_{m,g}\le1\).  Then the conclusion follows from
\[
  \sum_{m=1}^\infty\left[\mathbb P(E_m)+(eD_m)^{-A}\right]
  +
  \sum_m\sum_{\substack{g\in G_m\\ A_{m,g}>0}}
  C_{m,g}
  \left(\frac{\ell_m A_{m,g}}{Qn_m\lambda_{m,g}}\right)^\alpha
  <\infty .
\]
Thus block-volatility, sector-volatility, and multiscale elliptical
trace-envelope models pay one heavy-tail cost per group, interpolating between
the atomwise scale-mixture and common-scale criteria.
\end{corollary}

\begin{proof}
Put \(q_m=Qn_m/\ell_m\).  If \(Z_{m,g}\le\tau_{m,g}\) for every \(g\in G_m\),
then
\[
  S_m\le
  \sum_{g\in G_m} A_{m,g}Z_{m,g}
  \le
  \sum_{g\in G_m}\tau_{m,g}A_{m,g}
  \le q_m .
\]
The union bound therefore gives
\[
  \mathbb P\{S_m>q_m\}
  \le
  \sum_{g\in G_m}\mathbb P\{Z_{m,g}>\tau_{m,g}\}.
\]
The concrete Borel profile in
Corollary~\ref{cor:core-plus-quantile-trace-envelope-borel-random-amplitude-noise}
applies.

For the polynomial-tail specialization, take
\[
  \tau_{m,g}
  =
  Q\frac{n_m}{\ell_m}\frac{\lambda_{m,g}}{A_{m,g}}
  \qquad(A_{m,g}>0).
\]
Rows or groups with \(A_{m,g}=0\) contribute no mass.  The budget-share
condition gives \(\sum_g\tau_{m,g}A_{m,g}\le Qn_m/\ell_m\), and the displayed
tail sum follows after inserting this value of \(\tau_{m,g}\).
\end{proof}

\begin{corollary}[Countable grouped-scale heavy-tail trace-envelope families]
\label{cor:core-plus-countable-grouped-scale-tail-trace-envelope-borel-random-amplitude-noise}
In the setting of
Corollary~\ref{cor:core-plus-quantile-trace-envelope-borel-random-amplitude-noise},
put \(\ell_m=\log(eD_m)\).  Suppose that \(K_m\le K\) and
\(\ell_m\le B\sqrt{n_m}\) for fixed \(K,B<\infty\).  Let \(Q,A>0\).
For each \(m\), let \(G_m\) be finite or countable, let
\(A_{m,g}\ge0\) be deterministic group masses, and let
\(Z_{m,g}\ge0\) be arbitrary random group multipliers.  Assume that
\[
  S_m\le\sum_{g\in G_m} A_{m,g}Z_{m,g}
  \qquad\text{almost surely}.
\]
If there are deterministic thresholds \(\tau_{m,g}\ge0\) such that
\[
  \sum_{g\in G_m}\tau_{m,g}A_{m,g}
  \le Q\frac{n_m}{\ell_m}
\]
and
\[
  \sum_{m=1}^\infty
  \left[
    \mathbb P(E_m)
    +\sum_{g\in G_m}\mathbb P\{Z_{m,g}>\tau_{m,g}\}
    +(eD_m)^{-A}
  \right]<\infty,
\]
then, for every \(L>0\) satisfying
\[
  \frac{L^2}{2(Q+LB/3)}\ge A+1,
\]
almost surely, for all sufficiently large \(m\),
\[
  \disc(X_{m,1},\ldots,X_{m,n_m})
  \le (K+L)\sqrt{n_m}.
\]

In particular, suppose that, for some \(\alpha>0\),
\[
  \mathbb P\{Z_{m,g}>t\}\le C_{m,g}t^{-\alpha}
  \qquad(t>0,\ g\in G_m),
\]
and that there are deterministic budget shares
\(\lambda_{m,g}>0\) for every \(g\) with \(A_{m,g}>0\), satisfying
\(\sum_{g:A_{m,g}>0}\lambda_{m,g}\le1\).  Then the conclusion follows from
\[
  \sum_{m=1}^\infty\left[\mathbb P(E_m)+(eD_m)^{-A}\right]
  +
  \sum_m\sum_{\substack{g\in G_m\\ A_{m,g}>0}}
  C_{m,g}
  \left(\frac{\ell_m A_{m,g}}{Qn_m\lambda_{m,g}}\right)^\alpha
  <\infty .
\]
Thus countable sector, frequency, scale, and trace-class volatility
decompositions can be handled with one truncation failure per active group and
no independence assumptions.
\end{corollary}

\begin{proof}
Put \(q_m=Qn_m/\ell_m\).  If \(Z_{m,g}\le\tau_{m,g}\) for every \(g\in G_m\),
then monotone convergence for nonnegative series gives
\[
  S_m\le
  \sum_{g\in G_m} A_{m,g}Z_{m,g}
  \le
  \sum_{g\in G_m}\tau_{m,g}A_{m,g}
  \le q_m .
\]
The countable union bound therefore gives
\[
  \mathbb P\{S_m>q_m\}
  \le
  \sum_{g\in G_m}\mathbb P\{Z_{m,g}>\tau_{m,g}\}.
\]
The concrete Borel profile in
Corollary~\ref{cor:core-plus-quantile-trace-envelope-borel-random-amplitude-noise}
applies.

For the polynomial-tail specialization, take
\[
  \tau_{m,g}
  =
  Q\frac{n_m}{\ell_m}\frac{\lambda_{m,g}}{A_{m,g}}
  \qquad(A_{m,g}>0),
\]
and ignore the zero-mass groups.  The share condition gives
\(\sum_g\tau_{m,g}A_{m,g}\le Qn_m/\ell_m\), and substituting this
\(\tau_{m,g}\) into the countable tail sum gives the displayed criterion.
\end{proof}

\begin{corollary}[Optimized tail-profile grouped trace envelopes]
\label{cor:core-plus-optimized-tail-profile-grouped-trace-envelope-borel-random-amplitude-noise}
In the setting of
Corollary~\ref{cor:core-plus-countable-grouped-scale-tail-trace-envelope-borel-random-amplitude-noise},
let \(G_m\) be finite or countable and write \(\ell_m=\log(eD_m)\).  Suppose
that \(K_m\le K\), \(\ell_m\le B\sqrt{n_m}\), and
\[
  S_m\le\sum_{g\in G_m}A_{m,g}Z_{m,g}
  \qquad\text{almost surely},
\]
where \(A_{m,g}\ge0\) are deterministic and \(Z_{m,g}\ge0\) are arbitrary
random multipliers.  For each positive-mass group let
\(\pi_{m,g}:(0,\infty)\to[0,\infty]\) be a deterministic tail profile such
that
\[
  \mathbb P\{Z_{m,g}>t\}\le \pi_{m,g}(t)
  \qquad(t>0).
\]
For fixed \(Q>0\), define the optimized profile cost
\[
\begin{aligned}
  \Theta_m(Q)
  =
  \inf_{\lambda}
  \sum_{\substack{g\in G_m\\ A_{m,g}>0}}
  \pi_{m,g}\left(
    Q\frac{n_m}{\ell_m}\frac{\lambda_g}{A_{m,g}}
  \right),
\end{aligned}
\]
where the infimum is over all deterministic \(\lambda_g>0\) on the
positive-mass groups with
\(\sum_{g:A_{m,g}>0}\lambda_g\le1\); if there are no positive-mass groups, the
empty infimum is \(0\).  If, for some \(A>0\),
\[
  \sum_{m=1}^\infty
  \left[
    \mathbb P(E_m)
    +(eD_m)^{-A}
    +\Theta_m(Q)
  \right]<\infty,
\]
then, for every \(L>0\) satisfying
\[
  \frac{L^2}{2(Q+LB/3)}\ge A+1,
\]
almost surely, for all sufficiently large \(m\),
\[
  \disc(X_{m,1},\ldots,X_{m,n_m})
  \le (K+L)\sqrt{n_m}.
\]

In particular, it is enough to exhibit deterministic shares
\(\lambda_{m,g}>0\), \(\sum_{g:A_{m,g}>0}\lambda_{m,g}\le1\), for which
\[
\begin{aligned}
  &\sum_{m=1}^\infty\left[\mathbb P(E_m)+(eD_m)^{-A}\right]\\
  &\quad+
  \sum_m\sum_{\substack{g\in G_m\\ A_{m,g}>0}}
  \pi_{m,g}\left(
    Q\frac{n_m}{\ell_m}\frac{\lambda_{m,g}}{A_{m,g}}
  \right)
  <\infty .
\end{aligned}
\]
Thus one may use exact survival bounds, distribution-specific quantile
estimates, polynomial tails, stretched-exponential tails, or log-Weibull tails
such as \(C\exp\{-c(\log(e+t))^\beta\}\) in the same grouped trace-budget
criterion.  In particular, lognormal and multiplicative volatility envelopes
fit the same almost-sure Matrix Spencer route whenever this optimized profile
cost is summable.
\end{corollary}

\begin{proof}
For a fixed family of shares set, for \(A_{m,g}>0\),
\[
  \tau_{m,g}
  =
  Q\frac{n_m}{\ell_m}\frac{\lambda_{m,g}}{A_{m,g}} .
\]
Then
\[
  \sum_g \tau_{m,g}A_{m,g}
  =
  Q\frac{n_m}{\ell_m}
  \sum_{g:A_{m,g}>0}\lambda_{m,g}
  \le Q\frac{n_m}{\ell_m},
\]
and
\[
  \mathbb P\{Z_{m,g}>\tau_{m,g}\}
  \le \pi_{m,g}(\tau_{m,g}).
\]
The displayed fixed-share condition is therefore exactly the countable grouped
threshold criterion of
Corollary~\ref{cor:core-plus-countable-grouped-scale-tail-trace-envelope-borel-random-amplitude-noise}.

For the optimized criterion, the summability hypothesis implies
\(\Theta_m(Q)<\infty\) for every \(m\).  If positive-mass groups are present,
choose deterministic shares \(\lambda_{m,g}>0\) with
\(\sum_g\lambda_{m,g}\le1\) and
\[
\begin{aligned}
  \sum_{\substack{g\in G_m\\ A_{m,g}>0}}
  \pi_{m,g}\left(
    Q\frac{n_m}{\ell_m}\frac{\lambda_{m,g}}{A_{m,g}}
  \right)
  \le \Theta_m(Q)+2^{-m}.
\end{aligned}
\]
This is possible by the definition of the infimum; when there are no
positive-mass groups there is no multiplier-tail cost.  The resulting fixed
share series is summable, and the fixed-share criterion gives the conclusion.
\end{proof}

\begin{corollary}[Profile-menu grouped trace envelopes]
\label{cor:core-plus-profile-menu-grouped-trace-envelope-borel-random-amplitude-noise}
In the setting of
Corollary~\ref{cor:core-plus-optimized-tail-profile-grouped-trace-envelope-borel-random-amplitude-noise},
suppose that, for each \(m\), the positive-mass groups are partitioned into a
finite or countable family of profile classes \(G_{m,h}\), \(h\in H_m\).  For
each class there is a deterministic shape
\(\phi_{m,h}:(0,\infty)\to[0,\infty]\) and constants \(C_{m,g}\ge0\) such that
\[
  \mathbb P\{Z_{m,g}>t\}\le C_{m,g}\phi_{m,h}(t)
  \qquad(t>0,\ g\in G_{m,h},\ A_{m,g}>0).
\]
Put
\[
  \mathcal R_{m,h}
  =
  \sum_{\substack{g\in G_{m,h}\\ A_{m,g}>0}}A_{m,g},
  \qquad
  \mathcal C_{m,h}
  =
  \sum_{\substack{g\in G_{m,h}\\ A_{m,g}>0}}C_{m,g},
\]
and assume that \(\mathcal R_{m,h}<\infty\) whenever
\(\mathcal R_{m,h}>0\).  Suppose there are deterministic class shares
\(\eta_{m,h}>0\), defined for active classes \(\mathcal R_{m,h}>0\), with
\[
  \sum_{h:\mathcal R_{m,h}>0}\eta_{m,h}\le1.
\]
If, for some \(Q,A>0\),
\[
\begin{aligned}
  &\sum_{m=1}^\infty\left[\mathbb P(E_m)+(eD_m)^{-A}\right]\\
  &\quad+
  \sum_m\sum_{h:\mathcal R_{m,h}>0}
  \mathcal C_{m,h}\,
  \phi_{m,h}\left(
    Q\frac{n_m}{\ell_m}\frac{\eta_{m,h}}{\mathcal R_{m,h}}
  \right)
  <\infty,
\end{aligned}
\]
then, for every \(L>0\) satisfying
\[
  \frac{L^2}{2(Q+LB/3)}\ge A+1,
\]
almost surely, for all sufficiently large \(m\),
\[
  \disc(X_{m,1},\ldots,X_{m,n_m})
  \le (K+L)\sqrt{n_m}.
\]
Thus mixed volatility menus, for instance a finite or countable collection of
exponential, log-Weibull, regularly varying, and distribution-specific tail
classes, can be checked using only class-level total trace masses
\(\mathcal R_{m,h}\), aggregate coefficients \(\mathcal C_{m,h}\), and a
class-level budget allocation.  No independence between the group multipliers is
needed.
\end{corollary}

\begin{proof}
For each active class and each positive-mass group \(g\in G_{m,h}\), set
\[
  \lambda_{m,g}
  =
  \eta_{m,h}\frac{A_{m,g}}{\mathcal R_{m,h}}.
\]
Then
\[
  \sum_{\substack{g\in G_m\\ A_{m,g}>0}}\lambda_{m,g}
  =
  \sum_{h:\mathcal R_{m,h}>0}\eta_{m,h}
  \le1.
\]
Use the deterministic profile
\(\pi_{m,g}(t)=C_{m,g}\phi_{m,h}(t)\) for \(g\in G_{m,h}\).  The corresponding
fixed-share threshold is constant inside each profile class:
\[
  Q\frac{n_m}{\ell_m}\frac{\lambda_{m,g}}{A_{m,g}}
  =
  Q\frac{n_m}{\ell_m}\frac{\eta_{m,h}}{\mathcal R_{m,h}}.
\]
Therefore the fixed-share cost in
Corollary~\ref{cor:core-plus-optimized-tail-profile-grouped-trace-envelope-borel-random-amplitude-noise}
is bounded by the displayed class-level aggregate series.  The fixed-share
criterion gives the conclusion.
\end{proof}

\begin{corollary}[Optimized profile-menu grouped trace envelopes]
\label{cor:core-plus-optimized-profile-menu-grouped-trace-envelope-borel-random-amplitude-noise}
In the setting of
Corollary~\ref{cor:core-plus-profile-menu-grouped-trace-envelope-borel-random-amplitude-noise},
define, for \(Q>0\), the optimized class-profile cost
\[
\begin{aligned}
  \Xi_m(Q)
  =
  \inf_{\eta}
  \sum_{h:\mathcal R_{m,h}>0}
  \mathcal C_{m,h}\,
  \phi_{m,h}\left(
    Q\frac{n_m}{\ell_m}\frac{\eta_h}{\mathcal R_{m,h}}
  \right),
\end{aligned}
\]
where the infimum is over all deterministic \(\eta_h>0\) on the active profile
classes with \(\sum_{h:\mathcal R_{m,h}>0}\eta_h\le1\); if there are no active
classes, the empty infimum is \(0\).  If, for some \(A>0\),
\[
  \sum_{m=1}^\infty
  \left[
    \mathbb P(E_m)
    +(eD_m)^{-A}
    +\Xi_m(Q)
  \right]<\infty,
\]
then, for every \(L>0\) satisfying
\[
  \frac{L^2}{2(Q+LB/3)}\ge A+1,
\]
almost surely, for all sufficiently large \(m\),
\[
  \disc(X_{m,1},\ldots,X_{m,n_m})
  \le (K+L)\sqrt{n_m}.
\]
Thus mixed tail menus can be checked at the best deterministic allocation of
the trace budget among tail classes, while still avoiding optimization over
individual groups inside a class.
\end{corollary}

\begin{proof}
The summability hypothesis implies \(\Xi_m(Q)<\infty\) for every \(m\).  When
active profile classes are present, choose deterministic class shares
\(\eta_{m,h}>0\) with \(\sum_h\eta_{m,h}\le1\) and
\[
\begin{aligned}
  \sum_{h:\mathcal R_{m,h}>0}
  \mathcal C_{m,h}\,
  \phi_{m,h}\left(
    Q\frac{n_m}{\ell_m}\frac{\eta_{m,h}}{\mathcal R_{m,h}}
  \right)
  \le
  \Xi_m(Q)+2^{-m}.
\end{aligned}
\]
If there are no active profile classes, there is no multiplier-tail cost.  The
resulting class-level series is summable, so
Corollary~\ref{cor:core-plus-profile-menu-grouped-trace-envelope-borel-random-amplitude-noise}
applies.
\end{proof}

\begin{corollary}[Finite-menu heterogeneous polynomial profile schedules]
\label{cor:core-plus-finite-menu-heterogeneous-polynomial-profile-menu-schedule-trace-envelope-borel-random-amplitude-noise}
In the setting of
Corollary~\ref{cor:core-plus-optimized-profile-menu-grouped-trace-envelope-borel-random-amplitude-noise},
assume that, for each \(m\), the active profile classes are finite and every
active class has a polynomial tail shape
\[
  \phi_{m,h}(t)=t^{-\alpha_{m,h}},
  \qquad
  \alpha_{m,h}\ge\alpha_0>0.
\]
Let
\[
  \mathcal N_m
  =
  \#\{h:\mathcal R_{m,h}>0\},
  \qquad
  \Lambda_m=\log(e+m).
\]
Assume that the base exceptional series is finite and that there are constants
\(N_0,C_0,R_0<\infty\), \(w,p,u\ge0\), and \(z,q,v\in\mathbb R\) such that,
for all sufficiently large \(m\),
\[
  \mathcal N_m\le N_0m^w\Lambda_m^z,
\]
and, for every active class \(h\),
\[
  \mathcal C_{m,h}\le C_0m^p\Lambda_m^q,
  \qquad
  \mathcal R_{m,h}\le R_0m^u\Lambda_m^v .
\]
If there are \(a>0\) and \(b,c\in\mathbb R\) such that, for all sufficiently
large \(m\),
\[
  Q\frac{n_m}{\ell_m}
  \ge
  a m^b\Lambda_m^c,
\]
then the same almost-sure conclusion holds whenever either
\[
  \alpha_0 b>p+\alpha_0 u+(\alpha_0+1)w+1,
\]
or
\[
  \alpha_0 b=p+\alpha_0 u+(\alpha_0+1)w+1
  \qquad\text{and}\qquad
  q+\alpha_0 v+(\alpha_0+1)z-\alpha_0 c<-1.
\]
Thus finitely many polynomial profile classes with different Pareto indices are
controlled by the worst tail exponent \(\alpha_0\), the active class count, and
the worst aggregate coefficient and trace-mass growth.
\end{corollary}

\begin{proof}
If there are no active classes, there is no profile-menu cost.  Otherwise set
\[
  T_m=Q\frac{n_m}{\ell_m},
  \qquad
  \eta_{m,h}=\frac1{\mathcal N_m}
\]
on the active classes.  The growth hypotheses and either displayed exponent
condition imply \(b>w+u\), so for every active class and all sufficiently large
\(m\),
\[
  \frac{\mathcal N_m\mathcal R_{m,h}}{T_m}
  \le
  \frac{N_0R_0}{a}\,
  m^{w+u-b}\Lambda_m^{z+v-c}
  \le 1.
\]
Since \(\alpha_{m,h}\ge\alpha_0\), this gives
\[
\begin{aligned}
  \mathcal C_{m,h}
  \left(
    T_m\frac{\eta_{m,h}}{\mathcal R_{m,h}}
  \right)^{-\alpha_{m,h}}
  &=
  \mathcal C_{m,h}
  \left(
    \frac{\mathcal N_m\mathcal R_{m,h}}{T_m}
  \right)^{\alpha_{m,h}}
  \\
  &\le
  \mathcal C_{m,h}
  \left(
    \frac{\mathcal N_m\mathcal R_{m,h}}{T_m}
  \right)^{\alpha_0}.
\end{aligned}
\]
Summing over the at most \(\mathcal N_m\) active classes,
\[
\begin{aligned}
  \Xi_m(Q)
  &\le
  \sum_{h:\mathcal R_{m,h}>0}
  \mathcal C_{m,h}
  \left(
    \frac{\mathcal N_m\mathcal R_{m,h}}{T_m}
  \right)^{\alpha_0}
  \\
  &\le
  C\,
  m^{p+\alpha_0 u+(\alpha_0+1)w-\alpha_0 b}
  \Lambda_m^{q+\alpha_0 v+(\alpha_0+1)z-\alpha_0 c},
\end{aligned}
\]
with \(C=C_0N_0^{\alpha_0+1}R_0^{\alpha_0}a^{-\alpha_0}\).  This upper bound is
summable under the stated strict inequality, and at the endpoint precisely under
the displayed logarithmic inequality.  The optimized profile-menu corollary now
applies.
\end{proof}

\begin{corollary}[Countable weighted heterogeneous polynomial profile schedules]
\label{cor:core-plus-countable-weighted-heterogeneous-polynomial-profile-menu-schedule-trace-envelope-borel-random-amplitude-noise}
In the setting of
Corollary~\ref{cor:core-plus-optimized-profile-menu-grouped-trace-envelope-borel-random-amplitude-noise},
suppose the active profile classes have been indexed by positive integers and
every active class has a polynomial tail shape
\[
  \phi_{m,h}(t)=t^{-\alpha_{m,h}},
  \qquad
  \alpha_{m,h}\ge\alpha_0>0 .
\]
Let \((\omega_h)_{h\ge1}\) be positive weights with
\[
  W=\sum_{h=1}^\infty \omega_h<\infty .
\]
Assume that the base exceptional series
\[
  \sum_{m=1}^\infty\left[\mathbb P(E_m)+(eD_m)^{-A}\right]
\]
is finite and that, for all sufficiently large \(m\) and every active class
\(h\),
\[
  W\,\frac{\mathcal R_{m,h}}{\omega_h}
  \le
  Q\frac{n_m}{\ell_m}.
\]
If
\[
  \sum_{m=1}^\infty
  \left(\frac{\ell_m W}{Qn_m}\right)^{\alpha_0}
  \sum_{h:\mathcal R_{m,h}>0}
  \mathcal C_{m,h}
  \left(\frac{\mathcal R_{m,h}}{\omega_h}\right)^{\alpha_0}
  <\infty,
\]
then, for every \(L>0\) satisfying
\[
  \frac{L^2}{2(Q+LB/3)}\ge A+1,
\]
almost surely, for all sufficiently large \(m\),
\[
  \disc(X_{m,1},\ldots,X_{m,n_m})
  \le (K+L)\sqrt{n_m}.
\]

In particular, put \(\Lambda_m=\log(e+m)\).  Suppose that there are constants
\(C_0,R_0<\infty\), \(p,u\ge0\), \(q,v\in\mathbb R\), \(r,s\in\mathbb R\),
and \(\tau>1\) such that \(\tau<s\),
\[
  r+\alpha_0(s-\tau)>1,
\]
and, for all sufficiently large \(m\) and every active class \(h\),
\[
  \mathcal C_{m,h}\le C_0m^p\Lambda_m^q h^{-r},
  \qquad
  \mathcal R_{m,h}\le R_0m^u\Lambda_m^v h^{-s}.
\]
If there are \(a>0\) and \(b,c\in\mathbb R\) such that, for all sufficiently
large \(m\),
\[
  Q\frac{n_m}{\ell_m}
  \ge
  a m^b\Lambda_m^c,
\]
then the same almost-sure conclusion holds whenever either
\[
  \alpha_0 b>p+\alpha_0 u+1,
\]
or
\[
  \alpha_0 b=p+\alpha_0 u+1
  \qquad\text{and}\qquad
  q+\alpha_0 v-\alpha_0 c<-1.
\]
Equivalently, the displayed power-decay condition can be arranged whenever
\(s>1\) and \(r+\alpha_0(s-1)>1\).  Thus countably many mixed-Pareto profile
classes are covered by a fixed summable class-share distribution whenever the
trace masses decay slightly faster than that distribution and the weighted
coefficient--trace tail is summable.
\end{corollary}

\begin{proof}
Set
\[
  T_m=Q\frac{n_m}{\ell_m}.
\]
If there are no active profile classes, there is no profile-menu cost.
Otherwise choose deterministic class shares
\[
  \eta_{m,h}=\frac{\omega_h}{W}
\]
on the active classes; their total is at most one.  By the size hypothesis,
\[
  x_{m,h}
  :=
  \frac{W\mathcal R_{m,h}}{T_m\omega_h}
  \le 1
\]
for all sufficiently large \(m\) and every active class \(h\).  Since
\(\alpha_{m,h}\ge\alpha_0\),
\[
  x_{m,h}^{\alpha_{m,h}}\le x_{m,h}^{\alpha_0}.
\]
Hence the optimized profile-menu cost satisfies, for all large \(m\),
\[
\begin{aligned}
  \Xi_m(Q)
  &\le
  \sum_{h:\mathcal R_{m,h}>0}
  \mathcal C_{m,h}
  \left(
    \frac{W\mathcal R_{m,h}}{T_m\omega_h}
  \right)^{\alpha_{m,h}}
  \\
  &\le
  \left(\frac{\ell_m W}{Qn_m}\right)^{\alpha_0}
  \sum_{h:\mathcal R_{m,h}>0}
  \mathcal C_{m,h}
  \left(\frac{\mathcal R_{m,h}}{\omega_h}\right)^{\alpha_0}.
\end{aligned}
\]
The displayed summability hypothesis therefore implies the summability of
\(\sum_m\Xi_m(Q)\), and the optimized profile-menu corollary applies.

For the power-decaying specialization, take \(\omega_h=h^{-\tau}\) and
\(W=\sum_hh^{-\tau}\).  The strict or endpoint exponent condition implies
\(b>u\).  Since \(\tau<s\),
\[
  \frac{W\mathcal R_{m,h}}{T_m\omega_h}
  \le
  \frac{WR_0}{a}\,
  m^{u-b}\Lambda_m^{v-c}h^{\tau-s}
  \le
  \frac{WR_0}{a}\,
  m^{u-b}\Lambda_m^{v-c},
\]
which is at most one for all sufficiently large \(m\).  Also
\[
\begin{aligned}
  \sum_{h:\mathcal R_{m,h}>0}
  \mathcal C_{m,h}
  \left(\frac{\mathcal R_{m,h}}{\omega_h}\right)^{\alpha_0}
  &\le
  C_0R_0^{\alpha_0}
  m^{p+\alpha_0 u}
  \Lambda_m^{q+\alpha_0 v}
  \sum_{h=1}^\infty
  h^{-r-\alpha_0(s-\tau)} .
\end{aligned}
\]
The last series is finite by \(r+\alpha_0(s-\tau)>1\).  Therefore
\[
  \Xi_m(Q)
  \le
  C\,
  m^{p+\alpha_0 u-\alpha_0 b}
  \Lambda_m^{q+\alpha_0 v-\alpha_0 c}
\]
for a finite constant \(C\).  This is summable under the strict displayed
inequality, and at equality precisely under the displayed logarithmic endpoint
condition.  Finally, if \(s>1\) and \(r+\alpha_0(s-1)>1\), choose
\(\tau>1\) close enough to \(1\) so that \(\tau<s\) and
\(r+\alpha_0(s-\tau)>1\).
\end{proof}

\begin{corollary}[Finite-menu heterogeneous stretched-exponential profile schedules]
\label{cor:core-plus-finite-menu-heterogeneous-stretched-exponential-profile-menu-schedule-trace-envelope-borel-random-amplitude-noise}
In the setting of
Corollary~\ref{cor:core-plus-optimized-profile-menu-grouped-trace-envelope-borel-random-amplitude-noise},
assume that, for each \(m\), the active profile classes are finite and every
active class has a stretched-exponential tail shape
\[
  \phi_{m,h}(t)=\exp\{-c_{m,h}t^{\beta_{m,h}}\},
  \qquad
  c_{m,h}\ge c_0>0,
  \qquad
  \beta_{m,h}\ge\beta_0>0 .
\]
Let
\[
  \mathcal N_m=\#\{h:\mathcal R_{m,h}>0\},
  \qquad
  \mathcal R_m^*=\max_{h:\mathcal R_{m,h}>0}\mathcal R_{m,h},
\]
with \(\mathcal R_m^*=0\) if there are no active classes.  Put
\[
  \Lambda_m=\log(e+m),
  \qquad
  M_m=\log(e+\Lambda_m).
\]
Assume that the base exceptional series
\[
  \sum_{m=1}^\infty\left[\mathbb P(E_m)+(eD_m)^{-A}\right]
\]
is finite and that there are constants
\(N_0,C_0<\infty\), \(w,p\ge0\), and \(z,q\in\mathbb R\) such that, for all
sufficiently large \(m\),
\[
  \mathcal N_m\le N_0m^w\Lambda_m^z,
\]
and, for every active class \(h\),
\[
  \mathcal C_{m,h}\le C_0m^p\Lambda_m^q .
\]
Suppose that there are \(a>0\) and \(b\in\mathbb R\) such that, for all
sufficiently large \(m\) with \(\mathcal R_m^*>0\),
\[
  Q\frac{n_m}{\ell_m}\frac1{\mathcal N_m\mathcal R_m^*}
  \ge
  a(\Lambda_m+bM_m)^{1/\beta_0}.
\]
Let \(\delta=c_0a^{\beta_0}\).  Then the same almost-sure conclusion holds
whenever either
\[
  \delta>p+w+1,
\]
or
\[
  \delta=p+w+1
  \qquad\text{and}\qquad
  q+z-\delta b<-1.
\]

Equivalently, if
\(\mathcal R_m^*\le R_0m^u\Lambda_m^v\) eventually and
\[
  Q\frac{n_m}{\ell_m}
  \ge
  a_0m^{u+w}\Lambda_m^{v+z}
  (\Lambda_m+bM_m)^{1/\beta_0},
\]
then the same conclusion follows with
\(\delta=c_0(a_0/(N_0R_0))^{\beta_0}\).  Thus finitely many mixed exponential,
Weibull, and sub-Weibull profile classes are absorbed by a logarithmic
trace-budget margin whose polynomial exponent only pays for the class count and
aggregate coefficient growth.
\end{corollary}

\begin{proof}
If there are no active classes, there is no profile-menu cost.  Otherwise set
\[
  T_m=Q\frac{n_m}{\ell_m},
  \qquad
  \eta_{m,h}=\frac1{\mathcal N_m}
\]
on the active classes.  The budget hypothesis gives, for every active class and
all large \(m\),
\[
  T_m\frac{\eta_{m,h}}{\mathcal R_{m,h}}
  \ge
  T_m\frac1{\mathcal N_m\mathcal R_m^*}
  \ge
  a(\Lambda_m+bM_m)^{1/\beta_0}.
\]
The right-hand side tends to infinity, so it is at least \(1\) for all
sufficiently large \(m\).  For \(t\ge1\), the lower bounds
\(c_{m,h}\ge c_0\) and \(\beta_{m,h}\ge\beta_0\) imply
\[
  \exp\{-c_{m,h}t^{\beta_{m,h}}\}
  \le
  \exp\{-c_0t^{\beta_0}\}.
\]
Hence, with \(\delta=c_0a^{\beta_0}\),
\[
\begin{aligned}
  \Xi_m(Q)
  &\le
  \sum_{h:\mathcal R_{m,h}>0}
  \mathcal C_{m,h}
  \exp\left\{
    -c_{m,h}
    \left(
      T_m\frac{\eta_{m,h}}{\mathcal R_{m,h}}
    \right)^{\beta_{m,h}}
  \right\}
  \\
  &\le
  \mathcal N_m C_0m^p\Lambda_m^q
  \exp\{-\delta(\Lambda_m+bM_m)\}
  \\
  &\le
  C\,m^{p+w-\delta}\Lambda_m^{q+z-\delta b}
\end{aligned}
\]
for a finite constant \(C\).  This is summable under the strict inequality, and
at equality precisely when \(q+z-\delta b<-1\).  The optimized profile-menu
corollary gives the asserted almost-sure Matrix Spencer bound.

The absolute-budget variant follows because
\[
  \frac{T_m}{\mathcal N_m\mathcal R_m^*}
  \ge
  \frac{a_0}{N_0R_0}
  (\Lambda_m+bM_m)^{1/\beta_0}
\]
for all sufficiently large active \(m\).
\end{proof}

\begin{corollary}[Countable weighted heterogeneous stretched-exponential profile schedules]
\label{cor:core-plus-countable-weighted-heterogeneous-stretched-exponential-profile-menu-schedule-trace-envelope-borel-random-amplitude-noise}
In the setting of
Corollary~\ref{cor:core-plus-optimized-profile-menu-grouped-trace-envelope-borel-random-amplitude-noise},
suppose the active profile classes have been indexed by positive integers and
every active class has a stretched-exponential tail shape
\[
  \phi_{m,h}(t)=\exp\{-c_{m,h}t^{\beta_{m,h}}\},
  \qquad
  c_{m,h}\ge c_0>0,
  \qquad
  \beta_{m,h}\ge\beta_0>0 .
\]
Let \((\omega_h)_{h\ge1}\) be positive weights with
\[
  W=\sum_{h=1}^\infty\omega_h<\infty .
\]
Assume that the base exceptional series
\[
  \sum_{m=1}^\infty\left[\mathbb P(E_m)+(eD_m)^{-A}\right]
\]
is finite and that, for all sufficiently large \(m\) and every active class
\(h\),
\[
  Q\frac{n_m}{\ell_m}\frac{\omega_h}{W\mathcal R_{m,h}}\ge1.
\]
If
\[
  \sum_{m=1}^\infty
  \sum_{h:\mathcal R_{m,h}>0}
  \mathcal C_{m,h}
  \exp\left\{
    -c_0
    \left(
      Q\frac{n_m}{\ell_m}\frac{\omega_h}{W\mathcal R_{m,h}}
    \right)^{\beta_0}
  \right\}
  <\infty,
\]
then, for every \(L>0\) satisfying
\[
  \frac{L^2}{2(Q+LB/3)}\ge A+1,
\]
almost surely, for all sufficiently large \(m\),
\[
  \disc(X_{m,1},\ldots,X_{m,n_m})
  \le (K+L)\sqrt{n_m}.
\]

In particular, put \(\Lambda_m=\log(e+m)\) and
\(M_m=\log(e+\Lambda_m)\).  Suppose that there are constants
\(C_0,R_0<\infty\), \(p,u\ge0\), \(q,v,r\in\mathbb R\), and
\(s,\tau\in\mathbb R\) such that \(\tau>1\), \(s>\tau\), and, for all
sufficiently large \(m\) and every active class \(h\),
\[
  \mathcal C_{m,h}\le C_0m^p\Lambda_m^q h^{-r},
  \qquad
  \mathcal R_{m,h}\le R_0m^u\Lambda_m^v h^{-s}.
\]
If there are \(a>0\) and \(b\in\mathbb R\) such that, for all sufficiently
large \(m\),
\[
  Q\frac{n_m}{\ell_m}
  \ge
  a m^u\Lambda_m^v(\Lambda_m+bM_m)^{1/\beta_0},
\]
then the same almost-sure conclusion holds with
\[
  \delta=c_0\left(\frac{a}{WR_0}\right)^{\beta_0},
  \qquad
  W=\sum_{h=1}^\infty h^{-\tau},
\]
whenever either
\[
  \delta>p+1,
\]
or
\[
  \delta=p+1
  \qquad\text{and}\qquad
  q-\delta b<-1.
\]
Thus countably many mixed exponential, Weibull, and sub-Weibull profile classes
are allowed when the deterministic trace masses decay faster than a fixed
summable class-share distribution; no class-count factor appears in the final
triangular-array budget rule.
\end{corollary}

\begin{proof}
If there are no active classes, there is no profile-menu cost.  Otherwise choose
deterministic class shares
\[
  \eta_{m,h}=\frac{\omega_h}{W}.
\]
Their total over active classes is at most one.  For all sufficiently large
\(m\), the displayed size hypothesis makes
\[
  t_{m,h}
  =
  Q\frac{n_m}{\ell_m}\frac{\eta_{m,h}}{\mathcal R_{m,h}}
  =
  Q\frac{n_m}{\ell_m}\frac{\omega_h}{W\mathcal R_{m,h}}
  \ge1
\]
on every active class.  Hence
\[
  \phi_{m,h}(t_{m,h})
  =
  \exp\{-c_{m,h}t_{m,h}^{\beta_{m,h}}\}
  \le
  \exp\{-c_0t_{m,h}^{\beta_0}\}.
\]
The fixed class-share criterion in
Corollary~\ref{cor:core-plus-profile-menu-grouped-trace-envelope-borel-random-amplitude-noise}
therefore gives the first assertion.

For the power-decaying specialization, take \(\omega_h=h^{-\tau}\).  Since
\(s>\tau\), the budget hypothesis gives, for every active class and all large
\(m\),
\[
  t_{m,h}
  \ge
  \frac{a}{WR_0}
  (\Lambda_m+bM_m)^{1/\beta_0}
  h^{s-\tau}.
\]
The right-hand side is eventually at least \(1\).  With
\[
  \delta=c_0\left(\frac{a}{WR_0}\right)^{\beta_0},
  \qquad
  \gamma=\beta_0(s-\tau)>0,
\]
the class-profile cost is bounded by
\[
\begin{aligned}
  \sum_{h:\mathcal R_{m,h}>0}
  \mathcal C_{m,h}\phi_{m,h}(t_{m,h})
  &\le
  C_0m^p\Lambda_m^q
  \sum_{h=1}^\infty
  h^{-r}\exp\{-\delta(\Lambda_m+bM_m)h^\gamma\}.
\end{aligned}
\]
For all sufficiently large \(m\), \(\Lambda_m+bM_m\ge1\), and
\[
  \sum_{h=1}^\infty
  h^{-r}\exp\{-\delta(\Lambda_m+bM_m)h^\gamma\}
  \le
  C_{r,\gamma,\delta}\exp\{-\delta(\Lambda_m+bM_m)\}.
\]
Indeed, after factoring the \(h=1\) exponential term, the remaining summand is
bounded by
\[
  h^{-r}\exp\{-\delta(h^\gamma-1)\},
\]
whose series is finite.  Thus
\[
  \Xi_m(Q)
  \le
  C\,m^{p-\delta}\Lambda_m^{q-\delta b}.
\]
This series is summable if \(\delta>p+1\), and at equality precisely when
\(q-\delta b<-1\).  The optimized profile-menu corollary now gives the
asserted conclusion.
\end{proof}

\begin{corollary}[Finite-menu heterogeneous log-Weibull profile schedules]
\label{cor:core-plus-finite-menu-heterogeneous-logweibull-profile-menu-schedule-trace-envelope-borel-random-amplitude-noise}
In the setting of
Corollary~\ref{cor:core-plus-optimized-profile-menu-grouped-trace-envelope-borel-random-amplitude-noise},
assume that, for each \(m\), the active profile classes are finite and every
active class has a log-Weibull tail shape
\[
\begin{aligned}
  \phi_{m,h}(t)
  &=
  \exp\{-c_{m,h}(\log(e+t))^{\beta_{m,h}}\},\\
  c_{m,h}&\ge c_0>0,
  \qquad
  \beta_{m,h}\ge\beta_0>0 .
\end{aligned}
\]
Let
\[
  \mathcal N_m=\#\{h:\mathcal R_{m,h}>0\},
  \qquad
  \mathcal R_m^*=\max_{h:\mathcal R_{m,h}>0}\mathcal R_{m,h},
\]
with \(\mathcal R_m^*=0\) if there are no active classes.  Put
\[
  \Lambda_m=\log(e+m),
  \qquad
  M_m=\log(e+\Lambda_m).
\]
Assume that the base exceptional series
\[
  \sum_{m=1}^\infty\left[\mathbb P(E_m)+(eD_m)^{-A}\right]
\]
is finite and that there are constants \(N_0,C_0<\infty\), \(w,p\ge0\), and
\(z,q\in\mathbb R\) such that, for all sufficiently large \(m\),
\[
  \mathcal N_m\le N_0m^w\Lambda_m^z,
\]
and, for every active class \(h\),
\[
  \mathcal C_{m,h}\le C_0m^p\Lambda_m^q .
\]
Suppose that there are \(a>0\) and \(b\in\mathbb R\) such that, for all
sufficiently large \(m\) with \(\mathcal R_m^*>0\),
\[
  Q\frac{n_m}{\ell_m}\frac1{\mathcal N_m\mathcal R_m^*}
  \ge
  \exp\{a(\Lambda_m+bM_m)^{1/\beta_0}\}.
\]
Let \(\delta=c_0a^{\beta_0}\).  Then the same almost-sure conclusion holds
whenever either
\[
  \delta>p+w+1,
\]
or
\[
  \delta=p+w+1
  \qquad\text{and}\qquad
  q+z-\delta b<-1.
\]

Equivalently, if \(\mathcal R_m^*\le R_0m^u\Lambda_m^v\) eventually and
\[
  Q\frac{n_m}{\ell_m}
  \ge
  N_0R_0m^{u+w}\Lambda_m^{v+z}
  \exp\{a(\Lambda_m+bM_m)^{1/\beta_0}\},
\]
then the same conclusion follows.  Thus finitely many mixed log-Weibull,
including lognormal-type, profile classes are absorbed by a subpolynomial
trace-budget margin; the class-count and aggregate coefficient growth enter
only through the polynomial summability exponent \(p+w\).
\end{corollary}

\begin{proof}
If there are no active classes, there is no profile-menu cost.  Otherwise set
\[
  T_m=Q\frac{n_m}{\ell_m},
  \qquad
  \eta_{m,h}=\frac1{\mathcal N_m}
\]
on the active classes.  For all sufficiently large \(m\),
\(\Lambda_m+bM_m\ge1\).  The budget hypothesis gives, for every active class
and all large \(m\),
\[
  T_m\frac{\eta_{m,h}}{\mathcal R_{m,h}}
  \ge
  T_m\frac1{\mathcal N_m\mathcal R_m^*}
  \ge
  \exp\{a(\Lambda_m+bM_m)^{1/\beta_0}\}.
\]
Therefore
\[
  \log\left(
    e+T_m\frac{\eta_{m,h}}{\mathcal R_{m,h}}
  \right)
  \ge
  a(\Lambda_m+bM_m)^{1/\beta_0}
\]
for all large \(m\).  Since \(\log(e+t)\ge1\), the lower bounds
\(c_{m,h}\ge c_0\) and \(\beta_{m,h}\ge\beta_0\) imply
\[
  \exp\{-c_{m,h}(\log(e+t))^{\beta_{m,h}}\}
  \le
  \exp\{-c_0(\log(e+t))^{\beta_0}\}.
\]
Hence, with \(\delta=c_0a^{\beta_0}\),
\[
\begin{aligned}
  \Xi_m(Q)
  &\le
  \sum_{h:\mathcal R_{m,h}>0}
  \mathcal C_{m,h}
  \exp\left\{
    -c_{m,h}
    \left(
      \log\left(e+
        T_m\frac{\eta_{m,h}}{\mathcal R_{m,h}}
      \right)
    \right)^{\beta_{m,h}}
  \right\}
  \\
  &\le
  \mathcal N_m C_0m^p\Lambda_m^q
  \exp\{-\delta(\Lambda_m+bM_m)\}
  \\
  &\le
  C\,m^{p+w-\delta}\Lambda_m^{q+z-\delta b}
\end{aligned}
\]
for a finite constant \(C\).  This is summable under the strict inequality, and
at equality precisely when \(q+z-\delta b<-1\).  The optimized profile-menu
corollary gives the asserted almost-sure Matrix Spencer bound.

The absolute-budget variant follows immediately from the displayed bounds on
\(\mathcal N_m\) and \(\mathcal R_m^*\).
\end{proof}

\begin{corollary}[Countable weighted heterogeneous log-Weibull profile schedules]
\label{cor:core-plus-countable-weighted-heterogeneous-logweibull-profile-menu-schedule-trace-envelope-borel-random-amplitude-noise}
In the setting of
Corollary~\ref{cor:core-plus-optimized-profile-menu-grouped-trace-envelope-borel-random-amplitude-noise},
suppose the active profile classes have been indexed by positive integers and
every active class has a log-Weibull tail shape
\[
\begin{aligned}
  \phi_{m,h}(t)
  &=
  \exp\{-c_{m,h}(\log(e+t))^{\beta_{m,h}}\},\\
  c_{m,h}&\ge c_0>0,
  \qquad
  \beta_{m,h}\ge\beta_0>0 .
\end{aligned}
\]
Let \((\omega_h)_{h\ge1}\) be positive weights with
\[
  W=\sum_{h=1}^\infty\omega_h<\infty .
\]
Assume that the base exceptional series
\[
  \sum_{m=1}^\infty\left[\mathbb P(E_m)+(eD_m)^{-A}\right]
\]
is finite.  If
\[
  \sum_{m=1}^\infty
  \sum_{h:\mathcal R_{m,h}>0}
  \mathcal C_{m,h}
  \exp\left\{
    -c_0
    \left(
      \log\left(e+
        Q\frac{n_m}{\ell_m}\frac{\omega_h}{W\mathcal R_{m,h}}
      \right)
    \right)^{\beta_0}
  \right\}
  <\infty,
\]
then, for every \(L>0\) satisfying
\[
  \frac{L^2}{2(Q+LB/3)}\ge A+1,
\]
almost surely, for all sufficiently large \(m\),
\[
  \disc(X_{m,1},\ldots,X_{m,n_m})
  \le (K+L)\sqrt{n_m}.
\]

In particular, put \(\Lambda_m=\log(e+m)\) and
\(M_m=\log(e+\Lambda_m)\).  Suppose that there are constants
\(C_0,R_0<\infty\), \(p,u\ge0\), \(q,v,r\in\mathbb R\), and
\(s,\tau\in\mathbb R\) such that \(\tau>1\), \(s>\tau\), and, for all
sufficiently large \(m\) and every active class \(h\),
\[
  \mathcal C_{m,h}\le C_0m^p\Lambda_m^q h^{-r},
  \qquad
  \mathcal R_{m,h}\le R_0m^u\Lambda_m^v h^{-s}.
\]
Let
\[
  \kappa_{\beta_0}=\min\{1,2^{\beta_0-1}\},
  \qquad
  d=s-\tau,
\]
and assume
\[
  H_{\beta_0,d,r}
  =
  \sum_{h=1}^\infty
  h^{-r}
  \exp\{-c_0\kappa_{\beta_0}d^{\beta_0}(\log h)^{\beta_0}\}
  <\infty .
\]
If there are \(a>0\) and \(b\in\mathbb R\) such that, for all sufficiently
large \(m\),
\[
  Q\frac{n_m}{\ell_m}
  \ge
  WR_0m^u\Lambda_m^v
  \exp\{a(\Lambda_m+bM_m)^{1/\beta_0}\},
\]
then the same almost-sure conclusion holds with
\[
  \delta=c_0\kappa_{\beta_0}a^{\beta_0}
\]
whenever either
\[
  \delta>p+1,
\]
or
\[
  \delta=p+1
  \qquad\text{and}\qquad
  q-\delta b<-1.
\]
The class-index summability hypothesis is automatic for every \(r\) when
\(\beta_0>1\), becomes \(r+c_0d>1\) when \(\beta_0=1\), and reduces to
\(r>1\) when \(0<\beta_0<1\).  Thus countable mixed log-Weibull and
lognormal-type profile classes are controlled by a subpolynomial trace-budget
margin whenever their trace masses decay faster than a summable class-share
distribution and the residual class-index tail is summable.
\end{corollary}

\begin{proof}
If there are no active classes, there is no profile-menu cost.  Otherwise choose
deterministic class shares
\[
  \eta_{m,h}=\frac{\omega_h}{W}.
\]
Their total over active classes is at most one.  With
\[
  t_{m,h}
  =
  Q\frac{n_m}{\ell_m}\frac{\eta_{m,h}}{\mathcal R_{m,h}}
  =
  Q\frac{n_m}{\ell_m}\frac{\omega_h}{W\mathcal R_{m,h}},
\]
the lower bounds on \(c_{m,h}\) and \(\beta_{m,h}\), together with
\(\log(e+t_{m,h})\ge1\), give
\[
  \phi_{m,h}(t_{m,h})
  \le
  \exp\{-c_0(\log(e+t_{m,h}))^{\beta_0}\}.
\]
The fixed class-share criterion in
Corollary~\ref{cor:core-plus-profile-menu-grouped-trace-envelope-borel-random-amplitude-noise}
therefore gives the first assertion.

For the power-decaying specialization, take \(\omega_h=h^{-\tau}\).  For all
sufficiently large \(m\), \(\Lambda_m+bM_m\ge1\).  The budget and trace-mass
hypotheses give, for every active class,
\[
  t_{m,h}
  \ge
  \exp\{a(\Lambda_m+bM_m)^{1/\beta_0}\}h^{s-\tau}.
\]
Hence
\[
  \log(e+t_{m,h})
  \ge
  a(\Lambda_m+bM_m)^{1/\beta_0}+d\log h.
\]
For \(x,y\ge0\),
\[
  (x+y)^{\beta_0}
  \ge
  \kappa_{\beta_0}(x^{\beta_0}+y^{\beta_0}),
\]
where \(\kappa_{\beta_0}=\min\{1,2^{\beta_0-1}\}\).  Therefore the class-profile
cost is bounded by
\[
\begin{aligned}
  \sum_{h:\mathcal R_{m,h}>0}
  \mathcal C_{m,h}\phi_{m,h}(t_{m,h})
  &\le
  C_0m^p\Lambda_m^q
  \exp\{-\delta(\Lambda_m+bM_m)\}
  H_{\beta_0,d,r}
  \\
  &\le
  C\,m^{p-\delta}\Lambda_m^{q-\delta b}.
\end{aligned}
\]
This series is summable if \(\delta>p+1\), and at equality precisely when
\(q-\delta b<-1\).  The optimized profile-menu corollary gives the asserted
almost-sure Matrix Spencer bound.  The displayed descriptions of when
\(H_{\beta_0,d,r}\) is finite follow by comparing
\(\exp\{-C(\log h)^{\beta_0}\}\) with powers of \(h\).
\end{proof}

\begin{corollary}[Common-index polynomial profile-menu trace envelopes]
\label{cor:core-plus-common-index-polynomial-profile-menu-trace-envelope-borel-random-amplitude-noise}
In the setting of
Corollary~\ref{cor:core-plus-optimized-profile-menu-grouped-trace-envelope-borel-random-amplitude-noise},
assume that, for a fixed \(\alpha>0\), every active profile class has the
polynomial shape
\[
  \phi_{m,h}(t)=t^{-\alpha}.
\]
Define the class-level effective polynomial mass
\[
  \mathcal M_m
  =
  \sum_{h:\mathcal R_{m,h}>0}
  \mathcal C_{m,h}^{1/(\alpha+1)}
  \mathcal R_{m,h}^{\alpha/(\alpha+1)}
\]
with the convention that \(\mathcal M_m=0\) if there are no active classes.
Then
\[
  \Xi_m(Q)
  =
  \left(\frac{\ell_m}{Qn_m}\right)^\alpha
  \mathcal M_m^{\alpha+1}.
\]
Consequently, if, for some \(A,Q>0\),
\[
  \sum_{m=1}^\infty
  \left[
    \mathbb P(E_m)
    +(eD_m)^{-A}
    +
    \left(\frac{\ell_m}{Qn_m}\right)^\alpha
    \mathcal M_m^{\alpha+1}
  \right]<\infty,
\]
then, for every \(L>0\) satisfying
\[
  \frac{L^2}{2(Q+LB/3)}\ge A+1,
\]
almost surely, for all sufficiently large \(m\),
\[
  \disc(X_{m,1},\ldots,X_{m,n_m})
  \le (K+L)\sqrt{n_m}.
\]
This is the class-level analogue of the optimized grouped polynomial-tail
criterion: it needs only aggregate class coefficients and total deterministic
trace mass inside each profile class, not individual group-tail data.
\end{corollary}

\begin{proof}
Fix \(m\), suppress the index \(m\), and put
\[
  a_h=\mathcal C_h\mathcal R_h^\alpha,
  \qquad
  w_h=a_h^{1/(\alpha+1)}
  =
  \mathcal C_h^{1/(\alpha+1)}\mathcal R_h^{\alpha/(\alpha+1)}.
\]
For every admissible class-share family \((\eta_h)\),
\[
\begin{aligned}
  \sum_h w_h
  &=
  \sum_h (a_h\eta_h^{-\alpha})^{1/(\alpha+1)}
         \eta_h^{\alpha/(\alpha+1)}
  \\
  &\le
  \left(\sum_h a_h\eta_h^{-\alpha}\right)^{1/(\alpha+1)}
  \left(\sum_h\eta_h\right)^{\alpha/(\alpha+1)}
  \\
  &\le
  \left(\sum_h a_h\eta_h^{-\alpha}\right)^{1/(\alpha+1)}.
\end{aligned}
\]
Thus
\[
  \sum_h a_h\eta_h^{-\alpha}\ge \mathcal M_m^{\alpha+1}.
\]
Conversely, if \(0<\mathcal M_m<\infty\), the proportional shares
\[
  \eta_h=\frac{w_h}{\mathcal M_m}
\]
on the positive-weight classes have total one and give equality.  If zero-weight
active classes are also present, assign them total share \(\varepsilon>0\) and
scale the positive-weight shares by \(1-\varepsilon\); their cost increases by
the factor \((1-\varepsilon)^{-\alpha}\), and letting \(\varepsilon\downarrow0\)
gives the same infimum.  If \(\mathcal M_m=0\), all terms have zero cost.  If
\(\mathcal M_m=\infty\), the lower bound applied to finite subfamilies forces
the infimum to be infinite.  Multiplying by \((\ell_m/(Qn_m))^\alpha\) gives
the displayed formula for \(\Xi_m(Q)\).  The summability criterion is now
Corollary~\ref{cor:core-plus-optimized-profile-menu-grouped-trace-envelope-borel-random-amplitude-noise}.
\end{proof}

\begin{corollary}[Common-index polynomial profile-menu schedules]
\label{cor:core-plus-common-index-polynomial-profile-menu-schedule-trace-envelope-borel-random-amplitude-noise}
In the setting of
Corollary~\ref{cor:core-plus-common-index-polynomial-profile-menu-trace-envelope-borel-random-amplitude-noise},
assume that the base exceptional series
\[
  \sum_{m=1}^\infty\left[\mathbb P(E_m)+(eD_m)^{-A}\right]
\]
is finite.  Put \(\Lambda_m=\log(e+m)\).  Suppose that there are constants
\(M_0<\infty\), \(p\ge0\), \(q\in\mathbb R\), \(a>0\), and
\(b,c\in\mathbb R\) such that, for all sufficiently large \(m\),
\[
  \mathcal M_m\le M_0m^p\Lambda_m^q,
  \qquad
  Q\frac{n_m}{\ell_m}
  \ge
  a m^b\Lambda_m^c .
\]
If either
\[
  \alpha b>(\alpha+1)p+1,
\]
or
\[
  \alpha b=(\alpha+1)p+1
  \qquad\text{and}\qquad
  (\alpha+1)q-\alpha c<-1,
\]
then, for every \(L>0\) satisfying
\[
  \frac{L^2}{2(Q+LB/3)}\ge A+1,
\]
almost surely, for all sufficiently large \(m\),
\[
  \disc(X_{m,1},\ldots,X_{m,n_m})
  \le (K+L)\sqrt{n_m}.
\]
Thus a triangular array with polynomially many effective profile classes, or
with polynomially growing aggregate class coefficients and trace masses, is
covered as soon as the total trace budget \(Qn_m/\ell_m\) outruns the effective
class mass at the exponent dictated by the optimized polynomial allocation.
\end{corollary}

\begin{proof}
By
Corollary~\ref{cor:core-plus-common-index-polynomial-profile-menu-trace-envelope-borel-random-amplitude-noise},
it is enough to prove the summability of
\[
  \sum_m
  \left(\frac{\ell_m}{Qn_m}\right)^\alpha
  \mathcal M_m^{\alpha+1}.
\]
For all sufficiently large \(m\), the growth hypotheses give
\[
  \left(\frac{\ell_m}{Qn_m}\right)^\alpha
  \mathcal M_m^{\alpha+1}
  \le
  a^{-\alpha}M_0^{\alpha+1}
  m^{(\alpha+1)p-\alpha b}
  \Lambda_m^{(\alpha+1)q-\alpha c}.
\]
The latter series converges if
\(\alpha b>(\alpha+1)p+1\), and at equality precisely under
\((\alpha+1)q-\alpha c<-1\).  The preceding corollary gives the asserted
almost-sure Matrix Spencer bound.
\end{proof}

\begin{corollary}[Finite-menu common-index polynomial profile schedules]
\label{cor:core-plus-finite-menu-common-index-polynomial-profile-menu-schedule-trace-envelope-borel-random-amplitude-noise}
In the setting of
Corollary~\ref{cor:core-plus-common-index-polynomial-profile-menu-schedule-trace-envelope-borel-random-amplitude-noise},
assume that the active profile classes are finite for each \(m\), and let
\[
  \mathcal N_m
  =
  \#\{h:\mathcal R_{m,h}>0\}.
\]
Let \(\Lambda_m=\log(e+m)\).  Assume that the base exceptional series is finite
and that there are constants \(N_0,C_0,R_0<\infty\), \(w,p,u\ge0\), and
\(z,q,v\in\mathbb R\) such that, for all sufficiently large \(m\),
\[
  \mathcal N_m\le N_0m^w\Lambda_m^z,
\]
and, for every active class \(h\),
\[
  \mathcal C_{m,h}\le C_0m^p\Lambda_m^q,
  \qquad
  \mathcal R_{m,h}\le R_0m^u\Lambda_m^v .
\]
If there are \(a>0\) and \(b,c\in\mathbb R\) such that, for all sufficiently
large \(m\),
\[
  Q\frac{n_m}{\ell_m}
  \ge
  a m^b\Lambda_m^c,
\]
then the same almost-sure conclusion holds whenever either
\[
  \alpha b>p+\alpha u+(\alpha+1)w+1,
\]
or
\[
  \alpha b=p+\alpha u+(\alpha+1)w+1
  \qquad\text{and}\qquad
  q+\alpha v+(\alpha+1)z-\alpha c<-1.
\]
Thus polynomially many common-index polynomial tail classes can be checked from
only worst-class aggregate coefficient and trace-mass bounds, plus the class
count.
\end{corollary}

\begin{proof}
For all sufficiently large \(m\),
\[
\begin{aligned}
  \mathcal M_m
  &=
  \sum_{h:\mathcal R_{m,h}>0}
  \mathcal C_{m,h}^{1/(\alpha+1)}
  \mathcal R_{m,h}^{\alpha/(\alpha+1)}
  \\
  &\le
  N_0C_0^{1/(\alpha+1)}R_0^{\alpha/(\alpha+1)}
  m^{w+(p+\alpha u)/(\alpha+1)}
  \Lambda_m^{z+(q+\alpha v)/(\alpha+1)} .
\end{aligned}
\]
Corollary~\ref{cor:core-plus-common-index-polynomial-profile-menu-schedule-trace-envelope-borel-random-amplitude-noise}
therefore applies with
\[
  p'=w+\frac{p+\alpha u}{\alpha+1},
  \qquad
  q'=z+\frac{q+\alpha v}{\alpha+1}.
\]
The conditions \(\alpha b>(\alpha+1)p'+1\) and
\((\alpha+1)q'-\alpha c<-1\) are exactly the displayed inequalities.
\end{proof}

\begin{corollary}[Power-decaying common-index polynomial profile menus]
\label{cor:core-plus-power-decaying-common-index-polynomial-profile-menu-schedule-trace-envelope-borel-random-amplitude-noise}
In the setting of
Corollary~\ref{cor:core-plus-common-index-polynomial-profile-menu-schedule-trace-envelope-borel-random-amplitude-noise},
suppose the active profile classes have been indexed by positive integers.  Let
\(\Lambda_m=\log(e+m)\).  Assume that the base exceptional series is finite and
that there are constants \(C_0,R_0<\infty\), \(p,u\ge0\),
\(q,v\in\mathbb R\), and \(r,s\in\mathbb R\) such that, for all sufficiently
large \(m\) and every active class \(h\),
\[
  \mathcal C_{m,h}\le C_0m^p\Lambda_m^q h^{-r},
  \qquad
  \mathcal R_{m,h}\le R_0m^u\Lambda_m^v h^{-s}.
\]
If
\[
  r+\alpha s>\alpha+1,
\]
and there are \(a>0\) and \(b,c\in\mathbb R\) such that, for all sufficiently
large \(m\),
\[
  Q\frac{n_m}{\ell_m}
  \ge
  a m^b\Lambda_m^c,
\]
then the same almost-sure conclusion holds whenever either
\[
  \alpha b>p+\alpha u+1,
\]
or
\[
  \alpha b=p+\alpha u+1
  \qquad\text{and}\qquad
  q+\alpha v-\alpha c<-1.
\]
Thus countably many polynomial-tail profile classes are allowed when their
aggregate coefficients and deterministic trace masses have enough class-index
decay; only the aggregate growth exponents \(p,u,q,v\) enter the final
triangular-array budget rule.
\end{corollary}

\begin{proof}
The effective class mass from
Corollary~\ref{cor:core-plus-common-index-polynomial-profile-menu-trace-envelope-borel-random-amplitude-noise}
satisfies, for all large \(m\),
\[
\begin{aligned}
  \mathcal M_m
  &=
  \sum_{h:\mathcal R_{m,h}>0}
  \mathcal C_{m,h}^{1/(\alpha+1)}
  \mathcal R_{m,h}^{\alpha/(\alpha+1)}
  \\
  &\le
  C_0^{1/(\alpha+1)}R_0^{\alpha/(\alpha+1)}
  m^{(p+\alpha u)/(\alpha+1)}
  \Lambda_m^{(q+\alpha v)/(\alpha+1)}
  \sum_{h=1}^\infty
  h^{-(r+\alpha s)/(\alpha+1)} .
\end{aligned}
\]
The last series is finite exactly under \(r+\alpha s>\alpha+1\).  Hence
Corollary~\ref{cor:core-plus-common-index-polynomial-profile-menu-schedule-trace-envelope-borel-random-amplitude-noise}
applies with effective growth exponents
\[
  p'=\frac{p+\alpha u}{\alpha+1},
  \qquad
  q'=\frac{q+\alpha v}{\alpha+1}.
\]
Its conditions \(\alpha b>(\alpha+1)p'+1\) and
\((\alpha+1)q'-\alpha c<-1\) are precisely the displayed inequalities.
\end{proof}

\begin{corollary}[Common-profile grouped trace envelopes]
\label{cor:core-plus-common-profile-grouped-trace-envelope-borel-random-amplitude-noise}
In the setting of
Corollary~\ref{cor:core-plus-optimized-tail-profile-grouped-trace-envelope-borel-random-amplitude-noise},
suppose that the group tails have a common deterministic shape: for each \(m\)
there is \(\phi_m:(0,\infty)\to[0,\infty]\) and constants \(C_{m,g}\ge0\) such
that, for every positive-mass group,
\[
  \mathbb P\{Z_{m,g}>t\}\le C_{m,g}\phi_m(t)
  \qquad(t>0).
\]
Put
\[
  \mathcal R_m
  =
  \sum_{\substack{g\in G_m\\ A_{m,g}>0}} A_{m,g},
\]
and assume that \(\mathcal R_m<\infty\) whenever positive-mass groups are
present.  If, for some \(Q,A>0\),
\[
\begin{aligned}
  &\sum_{m=1}^\infty\left[\mathbb P(E_m)+(eD_m)^{-A}\right]\\
  &\quad+
  \sum_{m:\mathcal R_m>0}
  \sum_{\substack{g\in G_m\\ A_{m,g}>0}}
  C_{m,g}\,
  \phi_m\left(Q\frac{n_m}{\ell_m\mathcal R_m}\right)
  <\infty,
\end{aligned}
\]
then, for every \(L>0\) satisfying
\[
  \frac{L^2}{2(Q+LB/3)}\ge A+1,
\]
almost surely, for all sufficiently large \(m\),
\[
  \disc(X_{m,1},\ldots,X_{m,n_m})
  \le (K+L)\sqrt{n_m}.
\]

In particular, if
\[
  \mathbb P\{Z_{m,g}>t\}
  \le C_{m,g}\exp\{-c_m(\log(e+t))^{\beta_m}\}
  \qquad(t>0),
\]
then it is enough that
\[
\begin{aligned}
  &\sum_{m=1}^\infty\left[\mathbb P(E_m)+(eD_m)^{-A}\right]\\
  &\quad+
  \sum_{m:\mathcal R_m>0}
  \sum_{\substack{g\in G_m\\ A_{m,g}>0}}
  C_{m,g}
  \exp\left\{
    -c_m
    \left(\log\left(e+Q\frac{n_m}{\ell_m\mathcal R_m}\right)\right)^{\beta_m}
  \right\}
  <\infty .
\end{aligned}
\]
Thus countable grouped volatility decompositions with a common log-Weibull,
including lognormal-type, multiplier tail are controlled by one deterministic
total trace mass \(\mathcal R_m\) and one aggregate group-tail coefficient,
without solving the optimized share problem.
\end{corollary}

\begin{proof}
If there are no positive-mass groups, there is no multiplier-tail cost.  When
\(\mathcal R_m>0\), choose
\[
  \lambda_{m,g}
  =
  \frac{A_{m,g}}{\mathcal R_m}
  \qquad(A_{m,g}>0).
\]
These shares are positive and sum to one.  With
\(\pi_{m,g}(t)=C_{m,g}\phi_m(t)\), the corresponding threshold in
Corollary~\ref{cor:core-plus-optimized-tail-profile-grouped-trace-envelope-borel-random-amplitude-noise}
is
\[
  Q\frac{n_m}{\ell_m}\frac{\lambda_{m,g}}{A_{m,g}}
  =
  Q\frac{n_m}{\ell_m\mathcal R_m},
\]
independent of \(g\).  The fixed-share criterion in that corollary is exactly
the displayed aggregate-tail summability.  The log-Weibull specialization is the
same estimate with the displayed choice of \(\phi_m\).
\end{proof}

\begin{corollary}[Log-Weibull common-profile schedules]
\label{cor:core-plus-logweibull-common-profile-schedule-trace-envelope-borel-random-amplitude-noise}
In the setting of
Corollary~\ref{cor:core-plus-common-profile-grouped-trace-envelope-borel-random-amplitude-noise},
assume that, for fixed \(c,\beta>0\),
\[
  \mathbb P\{Z_{m,g}>t\}
  \le C_{m,g}\exp\{-c(\log(e+t))^\beta\}
  \qquad(t>0)
\]
for every positive-mass group, and put
\[
  \mathcal C_m
  =
  \sum_{\substack{g\in G_m\\ A_{m,g}>0}} C_{m,g}.
\]
Assume also that the base exceptional series
\[
  \sum_{m=1}^\infty\left[\mathbb P(E_m)+(eD_m)^{-A}\right]
\]
is finite.  If there are constants \(C_0<\infty\), \(p\ge0\), and \(a>0\)
such that, for all sufficiently large \(m\) with \(\mathcal R_m>0\),
\[
  \mathcal C_m\le C_0m^p,
  \qquad
  Q\frac{n_m}{\ell_m\mathcal R_m}
  \ge
  \exp\{a(\log m)^{1/\beta}\},
\]
and
\[
  ca^\beta>p+1,
\]
then, for every \(L>0\) satisfying
\[
  \frac{L^2}{2(Q+LB/3)}\ge A+1,
\]
almost surely, for all sufficiently large \(m\),
\[
  \disc(X_{m,1},\ldots,X_{m,n_m})
  \le (K+L)\sqrt{n_m}.
\]
In particular, lognormal-type grouped volatility, corresponding to
\(\beta=2\), is covered whenever the trace-budget ratio
\(Qn_m/(\ell_m\mathcal R_m)\) is at least
\(\exp\{a\sqrt{\log m}\}\) and \(ca^2>p+1\).  Thus a subpolynomial safety
margin in the total deterministic trace mass is already enough to absorb
polynomially many common lognormal-type volatility groups.
\end{corollary}

\begin{proof}
By Corollary~\ref{cor:core-plus-common-profile-grouped-trace-envelope-borel-random-amplitude-noise}
it is enough to show
\[
  \sum_{m:\mathcal R_m>0}
  \mathcal C_m
  \exp\left\{
    -c
    \left(\log\left(e+Q\frac{n_m}{\ell_m\mathcal R_m}\right)\right)^\beta
  \right\}
  <\infty .
\]
For all sufficiently large \(m\) with \(\mathcal R_m>0\), the growth
hypothesis gives
\[
  \log\left(e+Q\frac{n_m}{\ell_m\mathcal R_m}\right)
  \ge a(\log m)^{1/\beta}.
\]
Hence the \(m\)-th summand is at most
\[
  C_0m^p\exp\{-ca^\beta\log m\}
  =
  C_0m^{p-ca^\beta}.
\]
This is summable because \(ca^\beta>p+1\).  The common-profile corollary now
gives the asserted almost-sure Matrix Spencer bound.
\end{proof}

\begin{corollary}[Stretched-exponential common-profile schedules]
\label{cor:core-plus-stretched-exponential-common-profile-schedule-trace-envelope-borel-random-amplitude-noise}
In the setting of
Corollary~\ref{cor:core-plus-common-profile-grouped-trace-envelope-borel-random-amplitude-noise},
assume that, for fixed \(c,\beta>0\),
\[
  \mathbb P\{Z_{m,g}>t\}
  \le C_{m,g}\exp\{-ct^\beta\}
  \qquad(t>0)
\]
for every positive-mass group, and put
\[
  \mathcal C_m
  =
  \sum_{\substack{g\in G_m\\ A_{m,g}>0}} C_{m,g}.
\]
Assume also that
\[
  \sum_{m=1}^\infty\left[\mathbb P(E_m)+(eD_m)^{-A}\right]<\infty .
\]
Let \(\Lambda_m=\log(e+m)\) and \(M_m=\log(e+\Lambda_m)\).  Suppose that
there are constants \(C_0<\infty\), \(p\ge0\), \(q\in\mathbb R\), \(a>0\),
and \(b\in\mathbb R\) such that, for all sufficiently large \(m\) with
\(\mathcal R_m>0\),
\[
  \mathcal C_m\le C_0m^p\Lambda_m^q,
  \qquad
  Q\frac{n_m}{\ell_m\mathcal R_m}
  \ge
  a(\Lambda_m+bM_m)^{1/\beta}.
\]
If either
\[
  ca^\beta>p+1,
\]
or
\[
  ca^\beta=p+1
  \qquad\text{and}\qquad
  q-ca^\beta b<-1,
\]
then, for every \(L>0\) satisfying
\[
  \frac{L^2}{2(Q+LB/3)}\ge A+1,
\]
almost surely, for all sufficiently large \(m\),
\[
  \disc(X_{m,1},\ldots,X_{m,n_m})
  \le (K+L)\sqrt{n_m}.
\]
Thus common exponential, Weibull, and sub-Weibull volatility groups require
only logarithmic growth of the trace-budget ratio against polynomially many
aggregate tail coefficients.  At the critical balance \(ca^\beta=p+1\), the
log-log budget correction \(b\) and the aggregate logarithmic coefficient \(q\)
decide summability through \(q-ca^\beta b<-1\).
\end{corollary}

\begin{proof}
By Corollary~\ref{cor:core-plus-common-profile-grouped-trace-envelope-borel-random-amplitude-noise}
it is enough to show that
\[
  \sum_{m:\mathcal R_m>0}
  \mathcal C_m
  \exp\left\{
    -c\left(Q\frac{n_m}{\ell_m\mathcal R_m}\right)^\beta
  \right\}
  <\infty .
\]
For all sufficiently large \(m\), the quantity
\(\Lambda_m+bM_m\) is positive.  With \(\delta=ca^\beta\), the growth
hypotheses give
\[
  \mathcal C_m
  \exp\left\{
    -c\left(Q\frac{n_m}{\ell_m\mathcal R_m}\right)^\beta
  \right\}
  \le
  C_0m^p\Lambda_m^q
  \exp\{-\delta\Lambda_m-\delta bM_m\}.
\]
Since \(e^{-\delta\Lambda_m}=(e+m)^{-\delta}\) and
\(e^{-\delta bM_m}=(e+\Lambda_m)^{-\delta b}\), the \(m\)-th summand is at
most
\[
  C m^{p-\delta}\Lambda_m^{q-\delta b}.
\]
This is summable if \(\delta>p+1\), and in the endpoint case
\(\delta=p+1\) precisely when \(q-\delta b<-1\).  The common-profile corollary
therefore gives the asserted almost-sure Matrix Spencer bound.
\end{proof}

\begin{corollary}[Regularly varying common-profile schedules]
\label{cor:core-plus-regularly-varying-common-profile-schedule-trace-envelope-borel-random-amplitude-noise}
In the setting of
Corollary~\ref{cor:core-plus-common-profile-grouped-trace-envelope-borel-random-amplitude-noise},
assume that, for fixed \(\alpha>0\) and \(\gamma\in\mathbb R\),
\[
  \mathbb P\{Z_{m,g}>t\}
  \le C_{m,g}t^{-\alpha}(\log(e+t))^\gamma
  \qquad(t>0)
\]
for every positive-mass group, and put
\[
  \mathcal C_m
  =
  \sum_{\substack{g\in G_m\\ A_{m,g}>0}} C_{m,g}.
\]
Assume also that
\[
  \sum_{m=1}^\infty\left[\mathbb P(E_m)+(eD_m)^{-A}\right]<\infty .
\]
Let \(\Lambda_m=\log(e+m)\).  Suppose that there are constants
\(C_0<\infty\), \(p\ge0\), \(q\in\mathbb R\), \(a>0\), and
\(b\in\mathbb R\) such that, for all sufficiently large \(m\) with
\(\mathcal R_m>0\),
\[
  \mathcal C_m\le C_0m^p\Lambda_m^q,
  \qquad
  Q\frac{n_m}{\ell_m\mathcal R_m}
  \ge
  m^a\Lambda_m^b .
\]
If either
\[
  \alpha a>p+1,
\]
or
\[
  \alpha a=p+1
  \qquad\text{and}\qquad
  q+\gamma-\alpha b<-1,
\]
then, for every \(L>0\) satisfying
\[
  \frac{L^2}{2(Q+LB/3)}\ge A+1,
\]
almost surely, for all sufficiently large \(m\),
\[
  \disc(X_{m,1},\ldots,X_{m,n_m})
  \le (K+L)\sqrt{n_m}.
\]
Thus common regularly varying volatility tails with logarithmic corrections are
controlled by the polynomial growth of the aggregate tail coefficient and by
the polynomial/logarithmic safety margin in the deterministic trace budget.  At
the critical polynomial balance \(\alpha a=p+1\), the logarithmic inequality
\(q+\gamma-\alpha b<-1\) is the remaining summability condition.
\end{corollary}

\begin{proof}
By Corollary~\ref{cor:core-plus-common-profile-grouped-trace-envelope-borel-random-amplitude-noise}
it is enough to prove the summability of
\[
  \sum_{m:\mathcal R_m>0}
  \mathcal C_m
  T_m^{-\alpha}(\log(e+T_m))^\gamma,
  \qquad
  T_m=Q\frac{n_m}{\ell_m\mathcal R_m}.
\]
The function \(t\mapsto t^{-\alpha}(\log(e+t))^\gamma\) is decreasing for all
sufficiently large \(t\).  Since \(T_m\ge m^a\Lambda_m^b\) and \(a>0\), this
monotonicity applies for all sufficiently large active \(m\).  Hence, after
enlarging the implicit constant,
\[
  T_m^{-\alpha}(\log(e+T_m))^\gamma
  \le
  C m^{-\alpha a}\Lambda_m^{\gamma-\alpha b}.
\]
The \(m\)-th summand is therefore bounded by
\[
  C m^{p-\alpha a}\Lambda_m^{q+\gamma-\alpha b}.
\]
This series converges if \(\alpha a>p+1\), and in the endpoint case
\(\alpha a=p+1\) precisely under \(q+\gamma-\alpha b<-1\).  The common-profile
corollary gives the asserted almost-sure Matrix Spencer bound.
\end{proof}

\begin{corollary}[Heterogeneous polynomial grouped trace envelopes]
\label{cor:core-plus-heterogeneous-polynomial-grouped-tail-trace-envelope-borel-random-amplitude-noise}
In the setting of
Corollary~\ref{cor:core-plus-countable-grouped-scale-tail-trace-envelope-borel-random-amplitude-noise},
allow \(G_m\) to be finite or countable.  Suppose that \(K_m\le K\) and
\(\ell_m=\log(eD_m)\le B\sqrt{n_m}\) for fixed \(K,B<\infty\), and let
\(Q,A>0\).  Assume that
\[
  S_m\le\sum_{g\in G_m} A_{m,g}Z_{m,g}
  \qquad\text{almost surely},
\]
with deterministic \(A_{m,g}\ge0\) and arbitrary random \(Z_{m,g}\ge0\).
For every \(g\) with \(A_{m,g}>0\), suppose that there are constants
\(\alpha_{m,g}>0\) and \(C_{m,g}>0\) such that
\[
  \mathbb P\{Z_{m,g}>t\}\le C_{m,g}t^{-\alpha_{m,g}}
  \qquad(t>0),
\]
and deterministic budget shares \(\lambda_{m,g}>0\) satisfying
\[
  \sum_{g:A_{m,g}>0}\lambda_{m,g}\le1 .
\]
If
\[
  \sum_{m=1}^\infty\left[\mathbb P(E_m)+(eD_m)^{-A}\right]
  +
  \sum_m\sum_{\substack{g\in G_m\\ A_{m,g}>0}}
  C_{m,g}
  \left(\frac{\ell_m A_{m,g}}{Qn_m\lambda_{m,g}}\right)^{\alpha_{m,g}}
  <\infty,
\]
then, for every \(L>0\) satisfying
\[
  \frac{L^2}{2(Q+LB/3)}\ge A+1,
\]
almost surely, for all sufficiently large \(m\),
\[
  \disc(X_{m,1},\ldots,X_{m,n_m})
  \le (K+L)\sqrt{n_m}.
\]
Thus countable grouped trace-envelope models may mix tail indices across
sectors, scales, frequencies, or trace-class atoms, rather than using a single
polynomial exponent.
\end{corollary}

\begin{proof}
For each positive-mass group set
\[
  \tau_{m,g}
  =
  Q\frac{n_m}{\ell_m}\frac{\lambda_{m,g}}{A_{m,g}} .
\]
Then
\[
  \sum_{g\in G_m}\tau_{m,g}A_{m,g}
  =
  Q\frac{n_m}{\ell_m}
  \sum_{g:A_{m,g}>0}\lambda_{m,g}
  \le Q\frac{n_m}{\ell_m}.
\]
Moreover
\[
  \mathbb P\{Z_{m,g}>\tau_{m,g}\}
  \le
  C_{m,g}
  \left(\frac{\ell_m A_{m,g}}{Qn_m\lambda_{m,g}}\right)^{\alpha_{m,g}} .
\]
The displayed summability condition is therefore exactly the countable
threshold criterion of
Corollary~\ref{cor:core-plus-countable-grouped-scale-tail-trace-envelope-borel-random-amplitude-noise}
with these thresholds.
\end{proof}

\begin{corollary}[Optimized heterogeneous polynomial grouped trace envelopes]
\label{cor:core-plus-optimized-heterogeneous-polynomial-grouped-tail-trace-envelope-borel-random-amplitude-noise}
In the setting of
Corollary~\ref{cor:core-plus-heterogeneous-polynomial-grouped-tail-trace-envelope-borel-random-amplitude-noise},
let \(G_m\) be finite or countable and write \(\ell_m=\log(eD_m)\).  For fixed
\(Q>0\), define the optimized share cost
\[
\begin{aligned}
  H_m(Q)
  =
  \inf_{\lambda}
  \sum_{\substack{g\in G_m\\ A_{m,g}>0}}
  C_{m,g}
  \left(\frac{\ell_m A_{m,g}}{Qn_m\lambda_g}\right)^{\alpha_{m,g}},
\end{aligned}
\]
where the infimum is over all deterministic \(\lambda_g>0\) on the
positive-mass groups with
\(\sum_{g:A_{m,g}>0}\lambda_g\le1\); if there are no positive-mass groups, the
empty infimum is \(0\).  Suppose that \(K_m\le K\),
\(\ell_m\le B\sqrt{n_m}\), and
\[
  \sum_{m=1}^\infty
  \left[
    \mathbb P(E_m)
    +(eD_m)^{-A}
    +H_m(Q)
  \right]<\infty .
\]
Then, for every \(L>0\) satisfying
\[
  \frac{L^2}{2(Q+LB/3)}\ge A+1,
\]
almost surely, for all sufficiently large \(m\),
\[
  \disc(X_{m,1},\ldots,X_{m,n_m})
  \le (K+L)\sqrt{n_m}.
\]
Thus heterogeneous grouped polynomial tails can be checked at the best
available deterministic allocation of the \(n_m/\log(eD_m)\) trace budget.
\end{corollary}

\begin{proof}
The summability hypothesis implies \(H_m(Q)<\infty\) for every \(m\).  If there
are no positive-mass groups, there is no multiplier-tail contribution.  When
positive-mass groups are present, choose deterministic shares
\(\lambda_{m,g}>0\) with \(\sum_g\lambda_{m,g}\le1\) and
\[
  \sum_{\substack{g\in G_m\\ A_{m,g}>0}}
  C_{m,g}
  \left(\frac{\ell_m A_{m,g}}{Qn_m\lambda_{m,g}}\right)^{\alpha_{m,g}}
  \le H_m(Q)+2^{-m}.
\]
This is possible by the definition of the infimum.  Summing over \(m\), the
right-hand side is finite.  Applying
Corollary~\ref{cor:core-plus-heterogeneous-polynomial-grouped-tail-trace-envelope-borel-random-amplitude-noise}
with these shares gives the conclusion.
\end{proof}

\begin{corollary}[Optimized sub-Weibull grouped trace envelopes]
\label{cor:core-plus-optimized-subweibull-grouped-tail-trace-envelope-borel-random-amplitude-noise}
In the setting of
Corollary~\ref{cor:core-plus-countable-grouped-scale-tail-trace-envelope-borel-random-amplitude-noise},
let \(G_m\) be finite or countable and write \(\ell_m=\log(eD_m)\).  Suppose
that \(K_m\le K\), \(\ell_m\le B\sqrt{n_m}\), and
\[
  S_m\le\sum_{g\in G_m}A_{m,g}Z_{m,g}
  \qquad\text{almost surely},
\]
where \(A_{m,g}\ge0\) are deterministic and \(Z_{m,g}\ge0\) are arbitrary
random multipliers.  For each positive-mass group assume that there are
constants \(C_{m,g},c_{m,g},\beta_{m,g}>0\) such that
\[
  \mathbb P\{Z_{m,g}>t\}
  \le C_{m,g}\exp\{-c_{m,g}t^{\beta_{m,g}}\}
  \qquad(t>0).
\]
For fixed \(Q>0\), define the optimized sub-Weibull share cost
\[
\begin{aligned}
  \Psi_m(Q)
  =
  \inf_{\lambda}
  \sum_{\substack{g\in G_m\\ A_{m,g}>0}}
  C_{m,g}
  \exp\left\{
    -c_{m,g}
    \left(\frac{Qn_m\lambda_g}{\ell_m A_{m,g}}\right)^{\beta_{m,g}}
  \right\},
\end{aligned}
\]
where the infimum is over all deterministic \(\lambda_g>0\) on the
positive-mass groups with
\(\sum_{g:A_{m,g}>0}\lambda_g\le1\); if there are no positive-mass groups, the
empty infimum is \(0\).  If, for some \(A>0\),
\[
  \sum_{m=1}^\infty
  \left[
    \mathbb P(E_m)
    +(eD_m)^{-A}
    +\Psi_m(Q)
  \right]<\infty,
\]
then, for every \(L>0\) satisfying
\[
  \frac{L^2}{2(Q+LB/3)}\ge A+1,
\]
almost surely, for all sufficiently large \(m\),
\[
  \disc(X_{m,1},\ldots,X_{m,n_m})
  \le (K+L)\sqrt{n_m}.
\]

In particular, it is enough to exhibit deterministic shares
\(\lambda_{m,g}>0\), \(\sum_{g:A_{m,g}>0}\lambda_{m,g}\le1\), for which
\[
\begin{aligned}
  &\sum_{m=1}^\infty\left[\mathbb P(E_m)+(eD_m)^{-A}\right]\\
  &\quad+
  \sum_m\sum_{\substack{g\in G_m\\ A_{m,g}>0}}
  C_{m,g}
  \exp\left\{
    -c_{m,g}
    \left(\frac{Qn_m\lambda_{m,g}}{\ell_m A_{m,g}}\right)^{\beta_{m,g}}
  \right\}
  <\infty .
\end{aligned}
\]
Thus sector, scale, frequency, and trace-class volatility decompositions with
mixed stretched-exponential tail indices still give almost-sure eventual
Spencer-scale discrepancy at the trace-envelope threshold.
\end{corollary}

\begin{proof}
The displayed fixed-share criterion is the threshold criterion of
Corollary~\ref{cor:core-plus-countable-grouped-scale-tail-trace-envelope-borel-random-amplitude-noise}
with
\[
  \tau_{m,g}
  =
  Q\frac{n_m}{\ell_m}\frac{\lambda_{m,g}}{A_{m,g}}
  \qquad(A_{m,g}>0),
\]
zero-mass groups being ignored.

For the optimized criterion, the summability hypothesis implies
\(\Psi_m(Q)<\infty\) for every \(m\).  If positive-mass groups are present,
choose deterministic shares \(\lambda_{m,g}>0\) with
\(\sum_g\lambda_{m,g}\le1\) and
\[
\begin{aligned}
  \sum_{\substack{g\in G_m\\ A_{m,g}>0}}
  C_{m,g}
  \exp\left\{
    -c_{m,g}
    \left(\frac{Qn_m\lambda_{m,g}}{\ell_m A_{m,g}}\right)^{\beta_{m,g}}
  \right\}
  \le \Psi_m(Q)+2^{-m}.
\end{aligned}
\]
This is possible by the definition of the infimum; when there are no
positive-mass groups there is no multiplier-tail cost.  The resulting fixed
share series is summable, and the fixed-share criterion gives the conclusion.
\end{proof}

\begin{corollary}[Optimized grouped polynomial-tail trace envelopes]
\label{cor:core-plus-optimized-grouped-polynomial-tail-trace-envelope-borel-random-amplitude-noise}
In the setting of
Corollary~\ref{cor:core-plus-quantile-trace-envelope-borel-random-amplitude-noise},
put \(\ell_m=\log(eD_m)\).  Suppose that \(K_m\le K\) and
\(\ell_m\le B\sqrt{n_m}\) for fixed \(K,B<\infty\).  Let
\(\alpha,Q,A>0\).  For each \(m\), let \(G_m\) be finite, let
\(A_{m,g}\ge0\) be deterministic, and let \(Z_{m,g}\ge0\) be arbitrary random
multipliers such that
\[
  S_m\le \sum_{g\in G_m} A_{m,g}Z_{m,g}
  \qquad\text{almost surely}.
\]
Assume that, for every \(g\) with \(A_{m,g}>0\), there is a constant
\(C_{m,g}>0\) such that
\[
  \mathbb P\{Z_{m,g}>t\}\le C_{m,g}t^{-\alpha}
  \qquad(t>0).
\]
Set
\[
  R_m
  =
  \sum_{\substack{g\in G_m\\ A_{m,g}>0}}
  C_{m,g}^{1/(\alpha+1)}A_{m,g}^{\alpha/(\alpha+1)} .
\]
If
\[
  \sum_{m=1}^\infty
  \left[
    \mathbb P(E_m)
    +(eD_m)^{-A}
    +\left(\frac{\ell_m}{Qn_m}\right)^\alpha R_m^{\alpha+1}
  \right]<\infty,
\]
then, for every \(L>0\) satisfying
\[
  \frac{L^2}{2(Q+LB/3)}\ge A+1,
\]
almost surely, for all sufficiently large \(m\),
\[
  \disc(X_{m,1},\ldots,X_{m,n_m})
  \le (K+L)\sqrt{n_m}.
\]
Thus the polynomial grouped-scale condition can be checked by a single
effective group-tail mass, without choosing budget shares by hand.
\end{corollary}

\begin{proof}
If \(R_m=0\), then all \(A_{m,g}\) vanish and the grouped-scale corollary has no
multiplier-tail contribution.  Otherwise set
\[
  \lambda_{m,g}
  =
  \frac{
    C_{m,g}^{1/(\alpha+1)}A_{m,g}^{\alpha/(\alpha+1)}
  }{R_m}
  \qquad(A_{m,g}>0).
\]
These shares are positive and sum to one.  The polynomial-tail specialization
of
Corollary~\ref{cor:core-plus-grouped-scale-tail-trace-envelope-borel-random-amplitude-noise}
then requires the summability of
\[
  \sum_m\sum_{\substack{g\in G_m\\ A_{m,g}>0}}
  C_{m,g}
  \left(\frac{\ell_m A_{m,g}}{Qn_m\lambda_{m,g}}\right)^\alpha ,
\]
besides the base exceptional series.  For the displayed choice of
\(\lambda_{m,g}\), this double sum is
\[
  \sum_m
  \left(\frac{\ell_m}{Qn_m}\right)^\alpha R_m^\alpha
  \sum_{\substack{g\in G_m\\ A_{m,g}>0}}
  C_{m,g}^{1/(\alpha+1)}A_{m,g}^{\alpha/(\alpha+1)}
  =
  \sum_m
  \left(\frac{\ell_m}{Qn_m}\right)^\alpha R_m^{\alpha+1}.
\]
The grouped-scale corollary gives the asserted almost-sure Spencer-scale
bound.
\end{proof}

\begin{corollary}[Countable optimized grouped polynomial-tail trace envelopes]
\label{cor:core-plus-countable-optimized-grouped-polynomial-tail-trace-envelope-borel-random-amplitude-noise}
In the setting of
Corollary~\ref{cor:core-plus-quantile-trace-envelope-borel-random-amplitude-noise},
put \(\ell_m=\log(eD_m)\).  Suppose that \(K_m\le K\) and
\(\ell_m\le B\sqrt{n_m}\) for fixed \(K,B<\infty\).  Let
\(\alpha,Q,A>0\).  For each \(m\), let \(G_m\) be finite or countable, let
\(A_{m,g}\ge0\) be deterministic, and let \(Z_{m,g}\ge0\) be arbitrary random
multipliers such that
\[
  S_m\le \sum_{g\in G_m} A_{m,g}Z_{m,g}
  \qquad\text{almost surely}.
\]
Assume that, for every \(g\) with \(A_{m,g}>0\), there is a constant
\(C_{m,g}>0\) such that
\[
  \mathbb P\{Z_{m,g}>t\}\le C_{m,g}t^{-\alpha}
  \qquad(t>0).
\]
Define
\[
  R_m
  =
  \sum_{\substack{g\in G_m\\ A_{m,g}>0}}
  C_{m,g}^{1/(\alpha+1)}A_{m,g}^{\alpha/(\alpha+1)}
  \in[0,\infty].
\]
If
\[
  \sum_{m=1}^\infty
  \left[
    \mathbb P(E_m)
    +(eD_m)^{-A}
    +\left(\frac{\ell_m}{Qn_m}\right)^\alpha R_m^{\alpha+1}
  \right]<\infty,
\]
then, for every \(L>0\) satisfying
\[
  \frac{L^2}{2(Q+LB/3)}\ge A+1,
\]
almost surely, for all sufficiently large \(m\),
\[
  \disc(X_{m,1},\ldots,X_{m,n_m})
  \le (K+L)\sqrt{n_m}.
\]
Consequently, countable sector, frequency, scale, or trace-class decompositions
with polynomial group scales are controlled by the same effective group-tail
mass as the finite grouped model.
\end{corollary}

\begin{proof}
Put \(q_m=Qn_m/\ell_m\).  If \(R_m=0\), then there are no positive-mass groups
and \(S_m=0\) almost surely.  Otherwise the summability hypothesis forces
\(R_m<\infty\) for the relevant \(m\).  Set
\[
  \lambda_{m,g}
  =
  \frac{
    C_{m,g}^{1/(\alpha+1)}A_{m,g}^{\alpha/(\alpha+1)}
  }{R_m}
  \qquad(A_{m,g}>0),
\]
so that \(\sum_{g:A_{m,g}>0}\lambda_{m,g}=1\).  Define
\[
  \tau_{m,g}
  =
  Q\frac{n_m}{\ell_m}\frac{\lambda_{m,g}}{A_{m,g}}
  \qquad(A_{m,g}>0).
\]
If \(Z_{m,g}\le\tau_{m,g}\) for every positive-mass group, monotone convergence
for the nonnegative series gives
\[
  S_m
  \le \sum_{g\in G_m} A_{m,g}Z_{m,g}
  \le \sum_{A_{m,g}>0} A_{m,g}\tau_{m,g}
  = q_m .
\]
Hence, by the countable union bound and the polynomial tails,
\[
\begin{aligned}
  \mathbb P\{S_m>q_m\}
  &\le
  \sum_{A_{m,g}>0}\mathbb P\{Z_{m,g}>\tau_{m,g}\}  \\
  &\le
  \sum_{A_{m,g}>0}
  C_{m,g}\left(\frac{\ell_m A_{m,g}}{Qn_m\lambda_{m,g}}\right)^\alpha \\
  &=
  \left(\frac{\ell_m}{Qn_m}\right)^\alpha R_m^{\alpha+1}.
\end{aligned}
\]
The concrete Borel profile in
Corollary~\ref{cor:core-plus-quantile-trace-envelope-borel-random-amplitude-noise}
applies with this deterministic quantile-failure bound.
\end{proof}

\begin{corollary}[Sub-gamma trace-envelope variance profiles]
\label{cor:core-plus-subgamma-trace-envelope-random-amplitude-noise}
Let \(D\ge2\).  Keep the signed-core, filtration, and martingale-noise
hypotheses of
Corollary~\ref{cor:core-plus-trace-envelope-low-rank-random-amplitude-noise},
but do not assume that the trace envelopes are uniformly bounded.  Thus
\((V_i,\Gamma_i,\rho_i)\), \(i=1,\ldots,n\), are independent triples,
independent of \(\mathcal G\), with \(V_i\succeq0\), \(\rho_i\in[0,1]\), and
\[
  \operatorname{tr}V_i\le\Gamma_i
  \qquad\text{almost surely}.
\]
Put
\[
  Z_i=\rho_i^2\Gamma_i,
  \qquad
  \mu=\sum_{i=1}^n\mathbb E Z_i .
\]
Assume that the centered envelope sum is sub-gamma: for some
\(\sigma^2,c\ge0\),
\[
  \log\mathbb E
  \exp\left\{\lambda\sum_{i=1}^n(Z_i-\mathbb E Z_i)\right\}
  \le
  \frac{\lambda^2\sigma^2}{2(1-c\lambda)}
  \qquad(0\le\lambda<1/c),
\]
with the usual interpretation \(1/c=\infty\) when \(c=0\).  Then, for every
\(u\ge0\) and \(s>0\),
\[
\begin{aligned}
  \mathbb P\left\{
  \disc(X_1,\ldots,X_n)>K\sqrt n+s
  \right\}
  &\le
  \mathbb P(E)
  +\exp\left\{-\frac{u^2}{2(\sigma^2+cu)}\right\}\\
  &\quad+
  2D\exp\left\{-\frac{s^2}{2(\mu+u+s/3)}\right\},
\end{aligned}
\]
with the convention that the first exponential term is omitted when
\(\sigma^2+cu=0\).  Thus the signed-core theorem also covers unbounded random
trace envelopes whenever their total weighted envelope has a Bernstein-type
upper tail.
\end{corollary}

\begin{proof}
As before,
\[
  \left\|\sum_i\mathbb E(Y_i^2\mid\mcF_{i-1})\right\|_{\op}
  \le
  T:=\sum_i\rho_i^2\operatorname{tr}V_i
  \le
  \sum_i Z_i .
\]
The sub-gamma assumption and Chernoff's bound imply
\[
  \mathbb P\left\{\sum_iZ_i>\mu+u\right\}
  \le
  \exp\left\{-\frac{u^2}{2(\sigma^2+cu)}\right\}.
\]
Indeed, when \(\sigma^2+cu>0\), take
\(\lambda=u/(\sigma^2+cu)\) if \(\sigma^2>0\).  If \(\sigma^2=0<c\), use
\(\lambda=(1-\eta)/c\) and let \(\eta\downarrow0\).  The displayed
moment-generating-function bound gives the stated exponent in either case.
Combining this budget-exception estimate with
Corollary~\ref{cor:core-plus-matrix-martingale-noise} at \(v=\mu+u\) proves the
claim.
\end{proof}

\begin{corollary}[Termwise sub-gamma trace envelopes]
\label{cor:core-plus-termwise-subgamma-trace-envelope-random-amplitude-noise}
Keep the hypotheses and notation of
Corollary~\ref{cor:core-plus-subgamma-trace-envelope-random-amplitude-noise},
but replace the global sub-gamma assumption by the following termwise one.
For each \(i\), assume that there are \(\sigma_i^2,c_i\ge0\) such that, for
\(Z_i=\rho_i^2\Gamma_i\),
\[
  \log\mathbb E\exp\{\lambda(Z_i-\mathbb E Z_i)\}
  \le
  \frac{\lambda^2\sigma_i^2}{2(1-c_i\lambda)}
  \qquad(0\le\lambda<1/c_i).
\]
Put
\[
  \mu=\sum_i\mathbb E Z_i,\qquad
  \sigma^2=\sum_i\sigma_i^2,\qquad
  c=\max_{1\le i\le n}c_i .
\]
Then, for every \(u\ge0\) and \(s>0\),
\[
\begin{aligned}
  \mathbb P\left\{
  \disc(X_1,\ldots,X_n)>K\sqrt n+s
  \right\}
  &\le
  \mathbb P(E)
  +\exp\left\{-\frac{u^2}{2(\sigma^2+cu)}\right\}\\
  &\quad+
  2D\exp\left\{-\frac{s^2}{2(\mu+u+s/3)}\right\}.
\end{aligned}
\]
The same conclusion therefore applies whenever the weighted trace envelopes,
or the low-rank/Schatten envelopes \(R_i\tau_i\rho_i^2\) and
\(R_i^\beta L_i\rho_i^2\), are independently sub-gamma with summable variance
proxies and bounded scale.
\end{corollary}

\begin{proof}
Independence gives, for \(0\le\lambda<1/c\),
\[
\begin{aligned}
  \log\mathbb E\exp\left\{
  \lambda\sum_i(Z_i-\mathbb E Z_i)
  \right\}
  &=
  \sum_i\log\mathbb E\exp\{\lambda(Z_i-\mathbb E Z_i)\}\\
  &\le
  \sum_i\frac{\lambda^2\sigma_i^2}{2(1-c_i\lambda)}
  \le
  \frac{\lambda^2\sigma^2}{2(1-c\lambda)} .
\end{aligned}
\]
This is exactly the sub-gamma hypothesis of
Corollary~\ref{cor:core-plus-subgamma-trace-envelope-random-amplitude-noise}.
\end{proof}

\begin{corollary}[Gamma-dominated trace-envelope variance profiles]
\label{cor:core-plus-gamma-trace-envelope-random-amplitude-noise}
Keep the signed-core, filtration, and martingale-noise hypotheses of
Corollary~\ref{cor:core-plus-subgamma-trace-envelope-random-amplitude-noise}.
Assume that there are independent random variables \(W_i\), independent of
\(\mathcal G\), such that
\[
  \rho_i^2\Gamma_i\le W_i
  \qquad\text{almost surely}
\]
and \(W_i\) has the Gamma law with shape \(\kappa_i\ge0\) and scale
\(a_i\ge0\), with the convention that \(W_i=0\) when \(a_i\kappa_i=0\).
Define
\[
  \mu_\gamma=\sum_i\kappa_i a_i,\qquad
  \sigma_\gamma^2=\sum_i\kappa_i a_i^2,\qquad
  a_*=\max_i a_i .
\]
Then, for every \(u\ge0\) and \(s>0\),
\[
\begin{aligned}
  \mathbb P\left\{
  \disc(X_1,\ldots,X_n)>K\sqrt n+s
  \right\}
  &\le
  \mathbb P(E)
  +\exp\left\{-\frac{u^2}{2(\sigma_\gamma^2+a_*u)}\right\}\\
  &\quad+
  2D\exp\left\{-\frac{s^2}{2(\mu_\gamma+u+s/3)}\right\}.
\end{aligned}
\]
with the convention that the first exponential term is omitted when
\(\sigma_\gamma^2+a_*u=0\).
Equivalently, if \(W_i=(a_i/2)\chi_{\nu_i}^2\), then
\(\kappa_i=\nu_i/2\),
\[
  \mu_\gamma=\frac12\sum_i\nu_i a_i,
  \qquad
  \sigma_\gamma^2=\frac12\sum_i\nu_i a_i^2 .
\]
Thus chi-square and Wishart-type random trace envelopes give an explicit
signed-core Matrix Spencer tail.
\end{corollary}

\begin{proof}
The weighted trace budget satisfies
\[
  T=\sum_i\rho_i^2\operatorname{tr}V_i
  \le
  \sum_i\rho_i^2\Gamma_i
  \le
  \sum_i W_i .
\]
For \(W_i\sim\operatorname{Gamma}(\kappa_i,a_i)\),
\[
  \log\mathbb E\exp\{\lambda(W_i-\kappa_i a_i)\}
  =
  \kappa_i\bigl[-\log(1-a_i\lambda)-a_i\lambda\bigr]
  \le
  \frac{\lambda^2\kappa_i a_i^2}{2(1-a_i\lambda)}
\]
for \(0\le\lambda<1/a_i\); terms with \(a_i=0\) are deterministic and are
omitted.  Here we used the elementary bound
\(-\log(1-x)-x\le x^2/(2(1-x))\).  Summing over \(i\) gives the global
sub-gamma bound with parameters \(\sigma_\gamma^2\) and \(a_*\).  Hence
\[
  \mathbb P\{T>\mu_\gamma+u\}
  \le
  \exp\left\{-\frac{u^2}{2(\sigma_\gamma^2+a_*u)}\right\}.
\]
Combining this with Corollary~\ref{cor:core-plus-matrix-martingale-noise} at
\(v=\mu_\gamma+u\) proves the displayed tail.  The chi-square specialization is
the identity \(\chi_\nu^2\sim\operatorname{Gamma}(\nu/2,2)\).
\end{proof}

\begin{corollary}[Gamma-atom trace-envelope variance profiles]
\label{cor:core-plus-gamma-atom-trace-envelope-random-amplitude-noise}
Keep the signed-core, filtration, and martingale-noise hypotheses of
Corollary~\ref{cor:core-plus-subgamma-trace-envelope-random-amplitude-noise}.
Assume that there are independent nonnegative random variables
\(H_{i\ell}\), independent of \(\mathcal G\), such that
\[
  \rho_i^2\Gamma_i\le W_i:=\sum_{\ell=1}^{m_i}H_{i\ell}
  \qquad\text{almost surely}
\]
and \(H_{i\ell}\) has the Gamma law with shape
\(\kappa_{i\ell}\ge0\) and scale \(a_{i\ell}\ge0\), with the convention that
\(H_{i\ell}=0\) when \(a_{i\ell}\kappa_{i\ell}=0\).  Define
\[
  \mu_{\Gamma\text{-}\mathrm{at}}
  =\sum_{i,\ell}\kappa_{i\ell}a_{i\ell},\qquad
  \sigma_{\Gamma\text{-}\mathrm{at}}^2
  =\sum_{i,\ell}\kappa_{i\ell}a_{i\ell}^2,\qquad
  a_*=\max_{i,\ell}a_{i\ell}.
\]
Then, for every \(u\ge0\) and \(s>0\),
\[
\begin{aligned}
  \mathbb P\left\{
  \disc(X_1,\ldots,X_n)>K\sqrt n+s
  \right\}
  &\le
  \mathbb P(E)
  +\exp\left\{-\frac{u^2}{2(\sigma_{\Gamma\text{-}\mathrm{at}}^2+a_*u)}\right\}\\
  &\quad+
  2D\exp\left\{-\frac{s^2}{2(\mu_{\Gamma\text{-}\mathrm{at}}+u+s/3)}\right\},
\end{aligned}
\]
with the convention that the first exponential term is omitted when
\(\sigma_{\Gamma\text{-}\mathrm{at}}^2+a_*u=0\).
In particular, independent exponential trace atoms have
\(\kappa_{i\ell}=1\), while independent chi-square atoms
\((a_{i\ell}/2)\chi^2_{\nu_{i\ell}}\) have
\(\kappa_{i\ell}=\nu_{i\ell}/2\).  Thus heterogeneous gamma, exponential, and
chi-square atom decompositions give explicit signed-core Matrix Spencer tails
with the largest atom scale as the jump parameter.
\end{corollary}

\begin{proof}
The weighted trace budget satisfies
\[
  T=\sum_i\rho_i^2\operatorname{tr}V_i
  \le
  \sum_i\rho_i^2\Gamma_i
  \le
  \sum_{i,\ell}H_{i\ell}.
\]
For \(H_{i\ell}\sim\operatorname{Gamma}(\kappa_{i\ell},a_{i\ell})\),
\[
  \log\mathbb E\exp\{\lambda(H_{i\ell}-\kappa_{i\ell}a_{i\ell})\}
  \le
  \frac{\lambda^2\kappa_{i\ell}a_{i\ell}^2}{2(1-a_{i\ell}\lambda)}
\]
for \(0\le\lambda<1/a_{i\ell}\); deterministic terms with \(a_{i\ell}=0\)
are omitted.  This is the same elementary Gamma mgf bound used in
Corollary~\ref{cor:core-plus-gamma-trace-envelope-random-amplitude-noise}.
Summing over all atoms gives the sub-gamma bound with parameters
\(\sigma_{\Gamma\text{-}\mathrm{at}}^2\) and \(a_*\).  Hence
\[
  \mathbb P\{T>\mu_{\Gamma\text{-}\mathrm{at}}+u\}
  \le
  \exp\left\{-\frac{u^2}{2(\sigma_{\Gamma\text{-}\mathrm{at}}^2+a_*u)}\right\}.
\]
Combining this with Corollary~\ref{cor:core-plus-matrix-martingale-noise} at
\(v=\mu_{\Gamma\text{-}\mathrm{at}}+u\) proves the displayed tail.
\end{proof}

\begin{corollary}[Countable gamma-atom trace-envelope variance profiles]
\label{cor:core-plus-countable-gamma-atom-trace-envelope-random-amplitude-noise}
Keep the signed-core, filtration, and martingale-noise hypotheses of
Corollary~\ref{cor:core-plus-subgamma-trace-envelope-random-amplitude-noise}.
Assume that there are independent nonnegative random variables
\(H_{i\ell}\), independent of \(\mathcal G\), such that
\[
  \rho_i^2\Gamma_i\le W_i:=\sum_{\ell\ge1}H_{i\ell}
  \qquad\text{almost surely}
\]
and \(H_{i\ell}\) has the Gamma law with shape
\(\kappa_{i\ell}\ge0\) and scale \(a_{i\ell}\ge0\), with the convention that
\(H_{i\ell}=0\) when \(a_{i\ell}\kappa_{i\ell}=0\).  Suppose
\[
  \mu_{\Gamma,\infty}
  :=\sum_{i,\ell}\kappa_{i\ell}a_{i\ell}<\infty,\qquad
  \sigma_{\Gamma,\infty}^2
  :=\sum_{i,\ell}\kappa_{i\ell}a_{i\ell}^2<\infty,\qquad
  a_*:=\sup_{i,\ell}a_{i\ell}<\infty .
\]
Then, for every \(u\ge0\) and \(s>0\),
\[
\begin{aligned}
  \mathbb P\left\{
  \disc(X_1,\ldots,X_n)>K\sqrt n+s
  \right\}
  &\le
  \mathbb P(E)
  +\exp\left\{-\frac{u^2}{2(\sigma_{\Gamma,\infty}^2+a_*u)}\right\}\\
  &\quad+
  2D\exp\left\{-\frac{s^2}{2(\mu_{\Gamma,\infty}+u+s/3)}\right\},
\end{aligned}
\]
with the convention that the first exponential term is omitted when
\(\sigma_{\Gamma,\infty}^2+a_*u=0\).
Thus countable exponential and chi-square trace-atom expansions, including
trace-class real and complex Gaussian quadratic-form/Wishart envelopes, inherit
the same signed-core Matrix Spencer tail whenever the total mean trace budget
and the largest atom scale are finite.
\end{corollary}

\begin{proof}
The assumption \(\mu_{\Gamma,\infty}<\infty\) implies
\(\sum_{i,\ell}H_{i\ell}<\infty\) almost surely, so the weighted trace budget
satisfies
\[
  T=\sum_i\rho_i^2\operatorname{tr}V_i
  \le
  \sum_i\rho_i^2\Gamma_i
  \le
  \sum_{i,\ell}H_{i\ell}.
\]
If \(a_*=0\), then all atoms vanish almost surely and the claim follows
directly from Corollary~\ref{cor:core-plus-matrix-martingale-noise}.  Assume
\(a_*>0\).
For a finite set \(F\) of index pairs put
\[
  S_F=\sum_{(i,\ell)\in F}H_{i\ell},\qquad
  \mu_F=\sum_{(i,\ell)\in F}\kappa_{i\ell}a_{i\ell},\qquad
  \sigma_F^2=\sum_{(i,\ell)\in F}\kappa_{i\ell}a_{i\ell}^2 .
\]
By the finite gamma-atom bound in
Corollary~\ref{cor:core-plus-gamma-atom-trace-envelope-random-amplitude-noise},
for \(0\le\lambda<1/a_*\),
\[
  \mathbb E\exp\{\lambda(S_F-\mu_F)\}
  \le
  \exp\left\{
    \frac{\lambda^2\sigma_F^2}{2(1-a_*\lambda)}
  \right\}
  \le
  \exp\left\{
    \frac{\lambda^2\sigma_{\Gamma,\infty}^2}{2(1-a_*\lambda)}
  \right\}.
\]
Since \(\mu_F\le\mu_{\Gamma,\infty}\), the same bound holds with
\(S_F-\mu_{\Gamma,\infty}\) in place of \(S_F-\mu_F\).  Letting \(F\)
increase to all index pairs and using monotone convergence gives the same
bound for \(S=\sum_{i,\ell}H_{i\ell}\) centered at \(\mu_{\Gamma,\infty}\).
Hence
\[
  \mathbb P\{T>\mu_{\Gamma,\infty}+u\}
  \le
  \exp\left\{-\frac{u^2}{2(\sigma_{\Gamma,\infty}^2+a_*u)}\right\}.
\]
Combining this with Corollary~\ref{cor:core-plus-matrix-martingale-noise} at
\(v=\mu_{\Gamma,\infty}+u\) proves the displayed tail.
\end{proof}

\begin{corollary}[Compound-Poisson dominated trace-envelope variance profiles]
\label{cor:core-plus-compound-poisson-trace-envelope-random-amplitude-noise}
Keep the signed-core, filtration, and martingale-noise hypotheses of
Corollary~\ref{cor:core-plus-subgamma-trace-envelope-random-amplitude-noise}.
Assume that there are independent random variables \(W_i\), independent of
\(\mathcal G\), such that
\[
  \rho_i^2\Gamma_i\le W_i
  \qquad\text{almost surely}
\]
and each \(W_i\) has a compound-Poisson representation
\[
  W_i=\sum_{\ell=1}^{N_i} Z_{i\ell}.
\]
Here \(N_i\sim\operatorname{Poisson}(\lambda_i)\), the variables
\(Z_{i\ell}\) are conditionally iid given \(N_i\), independent of \(N_i\),
nonnegative, and satisfy \(Z_{i\ell}\le a_i\) almost surely.  Put
\[
  \mu_{\mathrm P}=\sum_i\lambda_i\mathbb E Z_{i1},\qquad
  \sigma_{\mathrm P}^2=\sum_i\lambda_i\mathbb E Z_{i1}^2,\qquad
  a_*=\max_i a_i .
\]
Then, for every \(u\ge0\) and \(s>0\),
\[
\begin{aligned}
  \mathbb P\left\{
  \disc(X_1,\ldots,X_n)>K\sqrt n+s
  \right\}
  &\le
  \mathbb P(E)
  +\exp\left\{-\frac{u^2}{2(\sigma_{\mathrm P}^2+a_*u/3)}\right\}\\
  &\quad+
  2D\exp\left\{-\frac{s^2}{2(\mu_{\mathrm P}+u+s/3)}\right\},
\end{aligned}
\]
with the convention that the first exponential term is omitted when
\(\sigma_{\mathrm P}^2+a_*u/3=0\).
In particular, if \(Z_{i\ell}=a_i\) almost surely, then
\(W_i=a_iN_i\),
\[
  \mu_{\mathrm P}=\sum_i\lambda_i a_i,\qquad
  \sigma_{\mathrm P}^2=\sum_i\lambda_i a_i^2 .
\]
Thus Poisson random counts of bounded trace-envelope atoms, including random
bounded spike-count variance profiles, give an explicit signed-core Matrix
Spencer tail.
\end{corollary}

\begin{proof}
The weighted trace budget is bounded by
\[
  T=\sum_i\rho_i^2\operatorname{tr}V_i
  \le
  \sum_i\rho_i^2\Gamma_i
  \le
  \sum_i W_i .
\]
For \(t\ge0\),
\[
  \log\mathbb E\exp\{t(W_i-\mathbb E W_i)\}
  =
  \lambda_i\mathbb E\bigl(e^{tZ_{i1}}-1-tZ_{i1}\bigr).
\]
The elementary bound \(e^x-1-x\le x^2/(2(1-x/3))\), valid for
\(0\le x<3\), gives
\[
  \log\mathbb E\exp\{t(W_i-\mathbb E W_i)\}
  \le
  \frac{t^2\lambda_i\mathbb E Z_{i1}^2}{2(1-a_it/3)}
\]
for \(0\le t<3/a_i\); terms with \(a_i=0\) are deterministic and are omitted.
Summing over \(i\) gives the sub-gamma bound with parameters
\(\sigma_{\mathrm P}^2\) and \(a_*/3\).  Hence
\[
  \mathbb P\{T>\mu_{\mathrm P}+u\}
  \le
  \exp\left\{-\frac{u^2}{2(\sigma_{\mathrm P}^2+a_*u/3)}\right\}.
\]
Combining this with Corollary~\ref{cor:core-plus-matrix-martingale-noise} at
\(v=\mu_{\mathrm P}+u\) proves the displayed tail.
\end{proof}

\begin{corollary}[Atomized bounded trace-envelope variance profiles]
\label{cor:core-plus-atomized-bounded-trace-envelope-random-amplitude-noise}
Keep the signed-core, filtration, and martingale-noise hypotheses of
Corollary~\ref{cor:core-plus-subgamma-trace-envelope-random-amplitude-noise}.
Assume that there are independent nonnegative random variables
\(Z_{i\ell}\), independent of \(\mathcal G\), such that
\[
  0\le Z_{i\ell}\le a_{i\ell}\qquad\text{almost surely}
\]
and
\[
  \rho_i^2\Gamma_i\le W_i:=\sum_{\ell=1}^{m_i} Z_{i\ell}
  \qquad\text{almost surely}.
\]
Define
\[
  \mu_{\mathrm A}=\sum_{i,\ell}\mathbb E Z_{i\ell},\qquad
  \sigma_{\mathrm A}^2=\sum_{i,\ell}\mathbb E Z_{i\ell}^2,\qquad
  a_*=\max_{i,\ell} a_{i\ell}.
\]
Then, for every \(u\ge0\) and \(s>0\),
\[
\begin{aligned}
  \mathbb P\left\{
  \disc(X_1,\ldots,X_n)>K\sqrt n+s
  \right\}
  &\le
  \mathbb P(E)
  +\exp\left\{-\frac{u^2}{2(\sigma_{\mathrm A}^2+a_*u/3)}\right\}\\
  &\quad+
  2D\exp\left\{-\frac{s^2}{2(\mu_{\mathrm A}+u+s/3)}\right\},
\end{aligned}
\]
with the convention that the first exponential term is omitted when
\(\sigma_{\mathrm A}^2+a_*u/3=0\).
In particular, if \(Z_{i\ell}=a_{i\ell}B_{i\ell}\) with independent
Bernoulli variables \(B_{i\ell}\) of means \(p_{i\ell}\), then
\[
  \mu_{\mathrm A}=\sum_{i,\ell}a_{i\ell}p_{i\ell},\qquad
  \sigma_{\mathrm A}^2=\sum_{i,\ell}a_{i\ell}^2p_{i\ell}.
\]
Thus finite Bernoulli/binomial random sums of bounded trace-envelope atoms,
including bounded random spike and low-rank atom variance profiles, give an
explicit signed-core Matrix Spencer tail with atom-level jump parameter.
\end{corollary}

\begin{proof}
The weighted trace budget satisfies
\[
  T=\sum_i\rho_i^2\operatorname{tr}V_i
  \le
  \sum_i\rho_i^2\Gamma_i
  \le
  \sum_{i,\ell} Z_{i\ell}.
\]
For a nonnegative \(Z\le a\),
\[
  \log\mathbb E e^{t(Z-\mathbb E Z)}
  \le
  \mathbb E(e^{tZ}-1-tZ)
  \le
  \frac{t^2\mathbb E Z^2}{2(1-at/3)}
\]
for \(0\le t<3/a\), using \(\log x\le x-1\) and
\(e^x-1-x\le x^2/(2(1-x/3))\); terms with \(a=0\) are deterministic and are
omitted.  Summing over all atoms gives the sub-gamma bound with parameters
\(\sigma_{\mathrm A}^2\) and \(a_*/3\).  Therefore
\[
  \mathbb P\{T>\mu_{\mathrm A}+u\}
  \le
  \exp\left\{-\frac{u^2}{2(\sigma_{\mathrm A}^2+a_*u/3)}\right\}.
\]
Combining this with Corollary~\ref{cor:core-plus-matrix-martingale-noise} at
\(v=\mu_{\mathrm A}+u\) proves the displayed tail.
\end{proof}

\begin{corollary}[Countable bounded-atom trace-envelope variance profiles]
\label{cor:core-plus-countable-atomized-bounded-trace-envelope-random-amplitude-noise}
Keep the signed-core, filtration, and martingale-noise hypotheses of
Corollary~\ref{cor:core-plus-subgamma-trace-envelope-random-amplitude-noise}.
Assume that there are independent nonnegative random variables
\(Z_{i\ell}\), independent of \(\mathcal G\), such that
\[
  0\le Z_{i\ell}\le a_{i\ell}\qquad\text{almost surely},
  \qquad
  \rho_i^2\Gamma_i\le W_i:=\sum_{\ell\ge1} Z_{i\ell}
  \qquad\text{almost surely}.
\]
Suppose
\[
  \mu_{\mathrm A,\infty}:=\sum_{i,\ell}\mathbb E Z_{i\ell}<\infty,\qquad
  \sigma_{\mathrm A,\infty}^2:=\sum_{i,\ell}\mathbb E Z_{i\ell}^2<\infty,
  \qquad
  a_*:=\sup_{i,\ell}a_{i\ell}<\infty .
\]
Then, for every \(u\ge0\) and \(s>0\),
\[
\begin{aligned}
  \mathbb P\left\{
  \disc(X_1,\ldots,X_n)>K\sqrt n+s
  \right\}
  &\le
  \mathbb P(E)
  +\exp\left\{-\frac{u^2}{2(\sigma_{\mathrm A,\infty}^2+a_*u/3)}\right\}\\
  &\quad+
  2D\exp\left\{-\frac{s^2}{2(\mu_{\mathrm A,\infty}+u+s/3)}\right\},
\end{aligned}
\]
with the convention that the first exponential term is omitted when
\(\sigma_{\mathrm A,\infty}^2+a_*u/3=0\).
In particular, if \(Z_{i\ell}=a_{i\ell}B_{i\ell}\) with independent
Bernoulli variables \(B_{i\ell}\) of means \(p_{i\ell}\), then
\[
  \mu_{\mathrm A,\infty}=\sum_{i,\ell}a_{i\ell}p_{i\ell},\qquad
  \sigma_{\mathrm A,\infty}^2=\sum_{i,\ell}a_{i\ell}^2p_{i\ell}.
\]
Thus countable Bernoulli and bounded-atom trace-envelope expansions, including
trace-class random spike and low-rank atom variance profiles, give explicit
signed-core Matrix Spencer tails with the supremal atom size as the jump
parameter.
\end{corollary}

\begin{proof}
The assumption \(\mu_{\mathrm A,\infty}<\infty\) implies
\(\sum_{i,\ell}Z_{i\ell}<\infty\) almost surely.  Hence
\[
  T=\sum_i\rho_i^2\operatorname{tr}V_i
  \le
  \sum_i\rho_i^2\Gamma_i
  \le
  \sum_{i,\ell}Z_{i\ell}.
\]
If \(a_*=0\), then all atoms vanish almost surely and the claim follows
directly from Corollary~\ref{cor:core-plus-matrix-martingale-noise}.  Assume
\(a_*>0\).  For a finite set \(F\) of index pairs put
\[
  S_F=\sum_{(i,\ell)\in F}Z_{i\ell},\qquad
  \mu_F=\sum_{(i,\ell)\in F}\mathbb E Z_{i\ell},\qquad
  \sigma_F^2=\sum_{(i,\ell)\in F}\mathbb E Z_{i\ell}^2 .
\]
The finite-atom mgf estimate in the proof of
Corollary~\ref{cor:core-plus-atomized-bounded-trace-envelope-random-amplitude-noise}
gives, for \(0\le t<3/a_*\),
\[
  \mathbb E\exp\{t(S_F-\mu_F)\}
  \le
  \exp\left\{
    \frac{t^2\sigma_F^2}{2(1-a_*t/3)}
  \right\}
  \le
  \exp\left\{
    \frac{t^2\sigma_{\mathrm A,\infty}^2}{2(1-a_*t/3)}
  \right\}.
\]
Since \(\mu_F\le\mu_{\mathrm A,\infty}\), the same bound holds with
\(S_F-\mu_{\mathrm A,\infty}\) in place of \(S_F-\mu_F\).  Letting \(F\)
increase to all index pairs and using monotone convergence gives the same
sub-gamma mgf bound for \(S=\sum_{i,\ell}Z_{i\ell}\) centered at
\(\mu_{\mathrm A,\infty}\).  Chernoff's bound yields
\[
  \mathbb P\{T>\mu_{\mathrm A,\infty}+u\}
  \le
  \exp\left\{-\frac{u^2}{2(\sigma_{\mathrm A,\infty}^2+a_*u/3)}\right\}.
\]
Combining this with Corollary~\ref{cor:core-plus-matrix-martingale-noise} at
\(v=\mu_{\mathrm A,\infty}+u\) proves the displayed tail.
\end{proof}

\begin{corollary}[Mixed atom trace-envelope variance profiles]
\label{cor:core-plus-mixed-atom-trace-envelope-random-amplitude-noise}
Keep the signed-core, filtration, and martingale-noise hypotheses of
Corollary~\ref{cor:core-plus-subgamma-trace-envelope-random-amplitude-noise}.
Let \(\mathcal I_\Gamma\), \(\mathcal I_{\mathrm A}\), and
\(\mathcal I_{\mathrm P}\) be finite or countable index sets.  Assume that all
random variables below are mutually independent and independent of
\(\mathcal G\), and that
\[
  \sum_i\rho_i^2\Gamma_i
  \le
  \sum_{\alpha\in\mathcal I_\Gamma}H_\alpha
  +\sum_{\beta\in\mathcal I_{\mathrm A}}Z_\beta
  +\sum_{j\in\mathcal I_{\mathrm P}}C_j
  \qquad\text{almost surely}.
\]
Here \(H_\alpha\sim\operatorname{Gamma}(\kappa_\alpha,a_\alpha)\) with
\(\kappa_\alpha,a_\alpha\ge0\),
\(0\le Z_\beta\le b_\beta\) almost surely, and
\[
  C_j=\sum_{\ell=1}^{N_j}U_{j\ell},
  \qquad
  N_j\sim\operatorname{Poisson}(\lambda_j),\qquad
  0\le U_{j\ell}\le d_j\quad\text{almost surely},
\]
with the \(U_{j\ell}\) iid for fixed \(j\) and independent of \(N_j\).  Suppose
\[
\begin{aligned}
  \mu_{\mathrm{mix}}
  &:=\sum_\alpha\kappa_\alpha a_\alpha
    +\sum_\beta\mathbb E Z_\beta
    +\sum_j\lambda_j\mathbb E U_{j1}<\infty,\\
  \sigma_{\mathrm{mix}}^2
  &:=\sum_\alpha\kappa_\alpha a_\alpha^2
    +\sum_\beta\mathbb E Z_\beta^2
    +\sum_j\lambda_j\mathbb E U_{j1}^2<\infty,
\end{aligned}
\]
and put
\[
  c_{\mathrm{mix}}
  :=\max\left\{
    \sup_\alpha a_\alpha,\,
    \frac13\sup_\beta b_\beta,\,
    \frac13\sup_j d_j
  \right\}<\infty,
\]
where the supremum over an empty family is \(0\).  Then, for every \(u\ge0\)
and \(s>0\),
\[
\begin{aligned}
  \mathbb P\left\{
  \disc(X_1,\ldots,X_n)>K\sqrt n+s
  \right\}
  &\le
  \mathbb P(E)
  +\exp\left\{-\frac{u^2}{2(\sigma_{\mathrm{mix}}^2+c_{\mathrm{mix}}u)}\right\}\\
  &\quad+
  2D\exp\left\{-\frac{s^2}{2(\mu_{\mathrm{mix}}+u+s/3)}\right\},
\end{aligned}
\]
with the convention that the first exponential term is omitted when
\(\sigma_{\mathrm{mix}}^2+c_{\mathrm{mix}}u=0\).  Thus a single signed-core
Matrix Spencer tail covers mixed trace envelopes containing, at the same time,
bounded Bernoulli/spike atoms, gamma or chi-square/Wishart atoms, and
compound-Poisson random counts of bounded atoms.
\end{corollary}

\begin{proof}
Let
\[
  S=
  \sum_{\alpha}H_\alpha+\sum_\beta Z_\beta+\sum_j C_j .
\]
The assumption \(\mu_{\mathrm{mix}}<\infty\) implies \(S<\infty\) almost
surely, and the weighted trace budget satisfies
\[
  T=\sum_i\rho_i^2\operatorname{tr}V_i
  \le \sum_i\rho_i^2\Gamma_i
  \le S .
\]
If \(c_{\mathrm{mix}}=0\), then all components vanish almost surely and the
claim follows directly from
Corollary~\ref{cor:core-plus-matrix-martingale-noise}.  Assume
\(c_{\mathrm{mix}}>0\).  For finite subfamilies
\(F_\Gamma,F_{\mathrm A},F_{\mathrm P}\), write \(S_F\), \(\mu_F\), and
\(\sigma_F^2\) for the corresponding truncated sum, mean, and variance proxy.
The gamma, bounded-atom, and compound-Poisson estimates proved above give, for
\(0\le t<1/c_{\mathrm{mix}}\),
\[
  \mathbb E\exp\{t(S_F-\mu_F)\}
  \le
  \exp\left\{
    \frac{t^2\sigma_F^2}{2(1-c_{\mathrm{mix}}t)}
  \right\}
  \le
  \exp\left\{
    \frac{t^2\sigma_{\mathrm{mix}}^2}{2(1-c_{\mathrm{mix}}t)}
  \right\}.
\]
Indeed the denominator \(1-c_{\mathrm{mix}}t\) is no larger than the
denominator coming from any individual gamma, bounded, or compound-Poisson
component.  Since \(\mu_F\le\mu_{\mathrm{mix}}\), the same bound holds with
\(S_F-\mu_{\mathrm{mix}}\) in place of \(S_F-\mu_F\).  Letting the finite
subfamilies increase to all indices and using monotone convergence gives the
same mgf bound for \(S-\mu_{\mathrm{mix}}\).  Chernoff's bound yields
\[
  \mathbb P\{T>\mu_{\mathrm{mix}}+u\}
  \le
  \exp\left\{-\frac{u^2}{2(\sigma_{\mathrm{mix}}^2+c_{\mathrm{mix}}u)}\right\}.
\]
Combining this budget-exception estimate with
Corollary~\ref{cor:core-plus-matrix-martingale-noise} at
\(v=\mu_{\mathrm{mix}}+u\) proves the displayed tail.
\end{proof}

\begin{corollary}[Weighted chi-square trace-envelope variance profiles]
\label{cor:core-plus-weighted-chisquare-trace-envelope-random-amplitude-noise}
Keep the signed-core, filtration, and martingale-noise hypotheses of
Corollary~\ref{cor:core-plus-subgamma-trace-envelope-random-amplitude-noise}.
Assume that there are independent random variables \(W_i\), independent of
\(\mathcal G\), such that
\[
  \rho_i^2\Gamma_i\le W_i
  \qquad\text{almost surely}
\]
and
\[
  W_i=\sum_{\ell=1}^{r_i} a_{i\ell}g_{i\ell}^2,
\]
where \(a_{i\ell}\ge0\) and the \(g_{i\ell}\) are independent standard real
Gaussians.  Define
\[
  \mu_\chi=\sum_{i,\ell}a_{i\ell},\qquad
  \sigma_\chi^2=2\sum_{i,\ell}a_{i\ell}^2,\qquad
  a_*=\max_{i,\ell}a_{i\ell}.
\]
Then, for every \(u\ge0\) and \(s>0\),
\[
\begin{aligned}
  \mathbb P\left\{
  \disc(X_1,\ldots,X_n)>K\sqrt n+s
  \right\}
  &\le
  \mathbb P(E)
  +\exp\left\{-\frac{u^2}{2(\sigma_\chi^2+2a_*u)}\right\}\\
  &\quad+
  2D\exp\left\{-\frac{s^2}{2(\mu_\chi+u+s/3)}\right\},
\end{aligned}
\]
with the convention that the first exponential term is omitted when
\(\sigma_\chi^2+2a_*u=0\).
In particular, if \(G_i\) is a \(d_i\times m_i\) matrix with independent
standard real Gaussian entries and \(A_i\succeq0\), then
\[
  W_i=\operatorname{tr}(G_i^*A_iG_i)
\]
is covered by taking the eigenvalues of \(A_i\), each repeated \(m_i\) times.
The corresponding parameters are
\[
  \mu_\chi=\sum_i m_i\operatorname{tr}A_i,\qquad
  \sigma_\chi^2=2\sum_i m_i\operatorname{tr}(A_i^2),\qquad
  a_*=\max_i\|A_i\|_{\op}.
\]
Thus weighted chi-square and real Wishart trace envelopes give explicit
signed-core Matrix Spencer tails.
\end{corollary}

\begin{proof}
The weighted trace budget satisfies
\[
  T=\sum_i\rho_i^2\operatorname{tr}V_i
  \le
  \sum_i\rho_i^2\Gamma_i
  \le
  \sum_{i,\ell} a_{i\ell}g_{i\ell}^2 .
\]
For a standard real Gaussian \(g\) and \(a\ge0\),
\[
  \log\mathbb E\exp\{t(a g^2-a)\}
  =
  -\frac12\log(1-2at)-at
  \le
  \frac{t^2a^2}{1-2at}
  =
  \frac{t^2(2a^2)}{2(1-2at)}
\]
for \(0\le t<1/(2a)\); terms with \(a=0\) are deterministic and are omitted.
The inequality is the bound
\(-\log(1-x)-x\le x^2/(2(1-x))\) with \(x=2at\).  Summing over all Gaussian
coordinates gives the sub-gamma bound with parameters \(\sigma_\chi^2\) and
\(2a_*\).  Therefore
\[
  \mathbb P\{T>\mu_\chi+u\}
  \le
  \exp\left\{-\frac{u^2}{2(\sigma_\chi^2+2a_*u)}\right\}.
\]
Combining this with Corollary~\ref{cor:core-plus-matrix-martingale-noise} at
\(v=\mu_\chi+u\) proves the displayed tail.  The Wishart specialization follows
by diagonalizing each \(A_i\) and writing
\(\operatorname{tr}(G_i^*A_iG_i)\) as the corresponding sum of independent
weighted squares.
\end{proof}

\begin{corollary}[Weighted exponential and complex-Wishart trace envelopes]
\label{cor:core-plus-complex-wishart-trace-envelope-random-amplitude-noise}
Keep the signed-core, filtration, and martingale-noise hypotheses of
Corollary~\ref{cor:core-plus-subgamma-trace-envelope-random-amplitude-noise}.
Assume that there are independent random variables \(W_i\), independent of
\(\mathcal G\), such that
\[
  \rho_i^2\Gamma_i\le W_i
  \qquad\text{almost surely}
\]
and
\[
  W_i=\sum_{\ell=1}^{r_i} a_{i\ell}|z_{i\ell}|^2,
\]
where \(a_{i\ell}\ge0\) and the \(z_{i\ell}\) are independent standard complex
Gaussians, normalized so that \(\mathbb E|z_{i\ell}|^2=1\).  Define
\[
  \mu_{\mathbb C}=\sum_{i,\ell}a_{i\ell},\qquad
  \sigma_{\mathbb C}^2=\sum_{i,\ell}a_{i\ell}^2,\qquad
  a_*=\max_{i,\ell}a_{i\ell}.
\]
Then, for every \(u\ge0\) and \(s>0\),
\[
\begin{aligned}
  \mathbb P\left\{
  \disc(X_1,\ldots,X_n)>K\sqrt n+s
  \right\}
  &\le
  \mathbb P(E)
  +\exp\left\{-\frac{u^2}{2(\sigma_{\mathbb C}^2+a_*u)}\right\}\\
  &\quad+
  2D\exp\left\{-\frac{s^2}{2(\mu_{\mathbb C}+u+s/3)}\right\},
\end{aligned}
\]
with the convention that the first exponential term is omitted when
\(\sigma_{\mathbb C}^2+a_*u=0\).
In particular, if \(G_i\) is a \(d_i\times m_i\) matrix with independent
standard complex Gaussian entries and \(A_i\succeq0\), then
\[
  W_i=\operatorname{tr}(G_i^*A_iG_i)
\]
is covered by taking the eigenvalues of \(A_i\), each repeated \(m_i\) times,
and the parameters are
\[
  \mu_{\mathbb C}=\sum_i m_i\operatorname{tr}A_i,\qquad
  \sigma_{\mathbb C}^2=\sum_i m_i\operatorname{tr}(A_i^2),\qquad
  a_*=\max_i\|A_i\|_{\op}.
\]
Thus complex Wishart and weighted exponential trace envelopes give explicit
signed-core Matrix Spencer tails.
\end{corollary}

\begin{proof}
The weighted trace budget satisfies
\[
  T=\sum_i\rho_i^2\operatorname{tr}V_i
  \le
  \sum_i\rho_i^2\Gamma_i
  \le
  \sum_{i,\ell} a_{i\ell}|z_{i\ell}|^2 .
\]
Since \(|z|^2\) is exponential with mean one,
\[
  \log\mathbb E\exp\{t(a|z|^2-a)\}
  =
  -\log(1-at)-at
  \le
  \frac{t^2a^2}{2(1-at)}
\]
for \(0\le t<1/a\); terms with \(a=0\) are deterministic and are omitted.
Summing over all complex Gaussian coordinates gives the sub-gamma bound with
parameters \(\sigma_{\mathbb C}^2\) and \(a_*\).  Therefore
\[
  \mathbb P\{T>\mu_{\mathbb C}+u\}
  \le
  \exp\left\{-\frac{u^2}{2(\sigma_{\mathbb C}^2+a_*u)}\right\}.
\]
Combining this with Corollary~\ref{cor:core-plus-matrix-martingale-noise} at
\(v=\mu_{\mathbb C}+u\) proves the displayed tail.  The complex-Wishart
specialization follows by diagonalizing each \(A_i\).
\end{proof}

\begin{corollary}[Noncentral Gaussian and Wishart trace envelopes]
\label{cor:core-plus-noncentral-gaussian-wishart-trace-envelope-random-amplitude-noise}
Keep the signed-core, filtration, and martingale-noise hypotheses of
Corollary~\ref{cor:core-plus-subgamma-trace-envelope-random-amplitude-noise}.
Assume that there are independent random variables \(W_i\), independent of
\(\mathcal G\), such that
\[
  \rho_i^2\Gamma_i\le W_i
  \qquad\text{almost surely}.
\]
Suppose either that
\[
  W_i=\delta_i+\sum_{\ell=1}^{r_i}a_{i\ell}(g_{i\ell}+b_{i\ell})^2,
\]
where \(\delta_i\ge0\), \(a_{i\ell}\ge0\), \(b_{i\ell}\in\R\) are
deterministic, and the \(g_{i\ell}\) are independent standard real Gaussians,
or that
\[
  W_i=\delta_i+\sum_{\ell=1}^{r_i}a_{i\ell}|z_{i\ell}+b_{i\ell}|^2,
\]
where \(\delta_i\ge0\), \(a_{i\ell}\ge0\), \(b_{i\ell}\in\C\) are
deterministic, and the \(z_{i\ell}\) are independent standard complex Gaussians
normalized by \(\mathbb E|z_{i\ell}|^2=1\).  In the real case put
\[
  \mu_{\mathrm R,\mathrm{nc}}
  =\sum_i\delta_i+\sum_{i,\ell}a_{i\ell}(1+b_{i\ell}^2),
  \qquad
  \sigma_{\mathrm R,\mathrm{nc}}^2
  =2\sum_{i,\ell}a_{i\ell}^2(1+2b_{i\ell}^2),
  \qquad
  c_{\mathrm R}=2a_*,
\]
and in the complex case put
\[
  \mu_{\mathbb C,\mathrm{nc}}
  =\sum_i\delta_i+\sum_{i,\ell}a_{i\ell}(1+|b_{i\ell}|^2),
  \qquad
  \sigma_{\mathbb C,\mathrm{nc}}^2
  =\sum_{i,\ell}a_{i\ell}^2(1+2|b_{i\ell}|^2),
  \qquad
  c_{\mathbb C}=a_*,
\]
where \(a_*=\max_{i,\ell}a_{i\ell}\), with \(a_*=0\) if all weights vanish.
Let \((\mu,\sigma^2,c)\) denote the corresponding real or complex triple.  Then,
for every \(u\ge0\) and \(s>0\),
\[
\begin{aligned}
  \mathbb P\left\{
  \disc(X_1,\ldots,X_n)>K\sqrt n+s
  \right\}
  &\le
  \mathbb P(E)
  +\exp\left\{-\frac{u^2}{2(\sigma^2+cu)}\right\}\\
  &\quad+
  2D\exp\left\{-\frac{s^2}{2(\mu+u+s/3)}\right\},
\end{aligned}
\]
with the convention that the first exponential term is omitted when
\(\sigma^2+cu=0\).

In particular, if \(G_i=M_i+Z_i\) is a real \(d_i\times m_i\) Gaussian matrix
with deterministic mean \(M_i\), independent standard real Gaussian noise
\(Z_i\), and \(A_i\succeq0\), then
\[
  W_i=\operatorname{tr}(G_i^*A_iG_i)
\]
is covered with
\[
  \mu_{\mathrm R,\mathrm{nc}}
  =\sum_i\left(m_i\operatorname{tr}A_i
      +\operatorname{tr}(M_i^*A_iM_i)\right),
\]
\[
  \sigma_{\mathrm R,\mathrm{nc}}^2
  =2\sum_i\left(m_i\operatorname{tr}(A_i^2)
      +2\operatorname{tr}(M_i^*A_i^2M_i)\right),
  \qquad
  c_{\mathrm R}=2\max_i\|A_i\|_{\op}.
\]
If \(G_i=M_i+Z_i\) is instead complex Gaussian with the same normalization as
above, then the same trace envelope is covered with
\[
  \mu_{\mathbb C,\mathrm{nc}}
  =\sum_i\left(m_i\operatorname{tr}A_i
      +\operatorname{tr}(M_i^*A_iM_i)\right),
\]
\[
  \sigma_{\mathbb C,\mathrm{nc}}^2
  =\sum_i\left(m_i\operatorname{tr}(A_i^2)
      +2\operatorname{tr}(M_i^*A_i^2M_i)\right),
  \qquad
  c_{\mathbb C}=\max_i\|A_i\|_{\op}.
\]
Thus nonzero-mean Gaussian directions and noncentral real or complex Wishart
trace envelopes give explicit signed-core Matrix Spencer tails.

More generally, if
\[
  W_i=\|C_i\xi_i+h_i\|_2^2
\]
with deterministic real or complex \(C_i,h_i\) and a standard real or complex
Gaussian vector \(\xi_i\), then the same conclusion holds after taking a
singular-value decomposition of \(C_i\).  The weights are the nonzero squared
singular values, the shifts are the coordinates of the projection of \(h_i\) to
\(\operatorname{Ran}C_i\) divided by the corresponding singular values, and
\(\delta_i\) is the squared norm of the orthogonal residual.  Hence correlated
Gaussian vector, matrix-normal, and correlated noncentral Wishart trace
envelopes are included after vectorizing the quadratic form.
\end{corollary}

\begin{proof}
Let \(T=\sum_i\rho_i^2\operatorname{tr}V_i\).  By the trace-envelope
hypothesis,
\[
  T\le \sum_i\rho_i^2\Gamma_i\le \sum_i W_i .
\]
For a standard real Gaussian \(g\), \(b\in\R\), \(a\ge0\), and
\(0\le t<1/(2a)\),
\[
\begin{aligned}
  &\log\mathbb E\exp\{t(a(g+b)^2-a(1+b^2))\}\\
  &\qquad=
  -\frac12\log(1-2at)-at
  +\frac{2a^2b^2t^2}{1-2at}\\
  &\qquad\le
  \frac{a^2(1+2b^2)t^2}{1-2at}
  =
  \frac{2a^2(1+2b^2)t^2}{2(1-2at)} .
\end{aligned}
\]
For a standard complex Gaussian \(z\) with \(\mathbb E|z|^2=1\),
\(b\in\C\), \(a\ge0\), and \(0\le t<1/a\),
\[
\begin{aligned}
  &\log\mathbb E\exp\{t(a|z+b|^2-a(1+|b|^2))\}\\
  &\qquad=
  -\log(1-at)-at
  +\frac{a^2|b|^2t^2}{1-at}\\
  &\qquad\le
  \frac{a^2(1+2|b|^2)t^2}{2(1-at)} .
\end{aligned}
\]
Summing over the independent coordinates gives a sub-gamma bound for
\(\sum_iW_i-\mu\) with variance proxy \(\sigma^2\) and scale \(c\).  Therefore
\[
  \mathbb P\{T>\mu+u\}
  \le
  \exp\left\{-\frac{u^2}{2(\sigma^2+cu)}\right\}.
\]
Combining this with Corollary~\ref{cor:core-plus-matrix-martingale-noise} at
\(v=\mu+u\) proves the displayed tail.

The Wishart specializations follow by diagonalizing each \(A_i\).  Orthogonal
or unitary invariance preserves the law of the centered Gaussian noise, while
the deterministic mean is transformed to the same eigenbasis.  The displayed
parameters are exactly the sums of \(a(1+|b|^2)\) and of
\(a^2(1+2|b|^2)\), with the extra factor \(2\) in the real variance proxy.
The final correlated Gaussian statement is the singular-value decomposition of
\(C_i\), followed by the orthogonal decomposition of \(h_i\) into the range of
\(C_i\) and its orthogonal complement; the latter contributes exactly the
deterministic residual \(\delta_i\).
\end{proof}

\begin{corollary}[Core signings with localized matrix martingale noise]
\label{cor:core-plus-localized-matrix-martingale-noise}
Let \(D\ge2\).  Let \(\mathcal G\) be a sigma-field, and let
\(B_1,\ldots,B_n\in\M_D(\C)\) be \(\mathcal G\)-measurable random self-adjoint
contractions.  Let \((\mcF_i)_{i=0}^n\) be a filtration with
\(\mathcal G\subset\mcF_0\).  For each \(i\), let \(P_i\) be an
\(\mcF_{i-1}\)-measurable orthogonal projection, and let \(Y_i\) be a
self-adjoint \(\mcF_i\)-measurable matrix satisfying
\[
  \mathbb E(Y_i\mid\mcF_{i-1})=0,\qquad
  \|Y_i\|_{\op}\le1,\qquad
  Y_i=P_iY_iP_i .
\]
Set \(X_i=B_i+Y_i\) and
\[
  \mu=\left\|\sum_{i=1}^n P_i\right\|_{\op}.
\]
Assume that, outside an event \(E\in\mathcal G\), there are
\(\mathcal G\)-measurable signs \(\eps_i^\star\in\{-1,1\}\) such that
\[
  \left\|\sum_{i=1}^n\eps_i^\star B_i\right\|_{\op}\le K\sqrt n .
\]
Then, for every deterministic \(v\ge0\) and \(s>0\),
\[
\begin{aligned}
  \mathbb P\left\{
  \disc(X_1,\ldots,X_n)>K\sqrt n+s
  \right\}
  &\le
  \mathbb P(E)+\mathbb P\{\mu>v\} \\
  &\quad+
  2D\exp\left\{-\frac{s^2}{2(v+s/3)}\right\}.
\end{aligned}
\]
Thus algebraic, finite-lift, Toeplitz/Hankel, and other already-signed random
cores obey the same Freedman tail after online coordinate-supported,
subspace-localized, or frame-localized matrix martingale-difference noise, with
variation controlled by the predictable overlap \(\|\sum_iP_i\|_{\op}\).
\end{corollary}

\begin{proof}
Since \(Y_i=P_iY_iP_i\) and \(\|Y_i\|_{\op}\le1\), we have \(0\le Y_i^2\preceq
P_i\).  As \(P_i\) is \(\mcF_{i-1}\)-measurable,
\[
  \sum_i\mathbb E(Y_i^2\mid\mcF_{i-1})
  \preceq
  \sum_iP_i.
\]
Therefore the variation parameter \(\nu\) in
Corollary~\ref{cor:core-plus-matrix-martingale-noise} satisfies \(\nu\le\mu\),
and the displayed bound follows from that corollary.
\end{proof}

\begin{corollary}[Core signings with conditionally sampled localized matrix noise]
\label{cor:core-plus-sampled-localized-matrix-martingale-noise}
Let \(D\ge2\).  Let \(\mathcal G\) be a sigma-field, and let
\(B_1,\ldots,B_n\in\M_D(\C)\) be \(\mathcal G\)-measurable random self-adjoint
contractions.  Let \((\mcF_i)_{i=0}^n\) be a filtration with
\(\mathcal G\subset\mcF_0\).  For each \(i\), let \(P_i\) be an
\(\mcF_i\)-measurable orthogonal projection, and let \(Y_i\) be a self-adjoint
\(\mcF_i\)-measurable matrix satisfying
\[
  \mathbb E(Y_i\mid\mcF_{i-1})=0,\qquad
  \|Y_i\|_{\op}\le1,\qquad
  Y_i=P_iY_iP_i .
\]
Set \(X_i=B_i+Y_i\) and
\[
  \mu=
  \left\|
  \sum_{i=1}^n\mathbb E(P_i\mid\mcF_{i-1})
  \right\|_{\op}.
\]
Assume that, outside an event \(E\in\mathcal G\), there are
\(\mathcal G\)-measurable signs \(\eps_i^\star\in\{-1,1\}\) such that
\[
  \left\|\sum_{i=1}^n\eps_i^\star B_i\right\|_{\op}\le K\sqrt n .
\]
Then, for every deterministic \(v\ge0\) and \(s>0\),
\[
\begin{aligned}
  \mathbb P\left\{
  \disc(X_1,\ldots,X_n)>K\sqrt n+s
  \right\}
  &\le
  \mathbb P(E)+\mathbb P\{\mu>v\} \\
  &\quad+
  2D\exp\left\{-\frac{s^2}{2(v+s/3)}\right\}.
\end{aligned}
\]
In particular, if
\[
  \sum_{i=1}^n\mathbb E(P_i\mid\mcF_{i-1})\preceq pn\,I
  \qquad\text{almost surely},
\]
then
\[
  \mathbb P\left\{
  \disc(X_1,\ldots,X_n)>K\sqrt n+s
  \right\}
  \le
  \mathbb P(E)+
  2D\exp\left\{-\frac{s^2}{2(pn+s/3)}\right\}.
\]
Thus a signed core remains balanced when the matrix martingale noise is
supported on coordinate sets, random subspaces, or any other sampled projection
model at time \(i\), as long as the predictable conditional overlap is small.
\end{corollary}

\begin{proof}
Since \(Y_i=P_iY_iP_i\) and \(\|Y_i\|_{\op}\le1\), the same compression estimate
as above gives \(0\le Y_i^2\preceq P_i\).  Hence
\[
  \sum_i\mathbb E(Y_i^2\mid\mcF_{i-1})
  \preceq
  \sum_i\mathbb E(P_i\mid\mcF_{i-1}),
\]
so the variation parameter \(\nu\) in
Corollary~\ref{cor:core-plus-matrix-martingale-noise} satisfies \(\nu\le\mu\).
The first bound follows from that corollary.  The final display is the same
estimate with \(v=pn\).
\end{proof}

\begin{corollary}[Core signings with weighted sampled localized matrix noise]
\label{cor:core-plus-weighted-sampled-localized-matrix-noise}
Let \(D\ge2\).  Let \(\mathcal G\) be a sigma-field, and let
\(B_1,\ldots,B_n\in\M_D(\C)\) be \(\mathcal G\)-measurable random self-adjoint
contractions.  Let \((\mcF_i)_{i=0}^n\) be a filtration with
\(\mathcal G\subset\mcF_0\).  For each \(i\), let \(P_i\) be an
\(\mcF_i\)-measurable orthogonal projection and let
\(\rho_i\in[0,1]\) be \(\mcF_{i-1}\)-measurable.  Let \(Y_i\) be a self-adjoint
\(\mcF_i\)-measurable matrix satisfying
\[
  \mathbb E(Y_i\mid\mcF_{i-1})=0,\qquad
  \|Y_i\|_{\op}\le\rho_i,\qquad
  Y_i=P_iY_iP_i .
\]
Set \(X_i=B_i+Y_i\) and
\[
  \mu_\rho=
  \left\|
  \sum_{i=1}^n\rho_i^2\,\mathbb E(P_i\mid\mcF_{i-1})
  \right\|_{\op}.
\]
Assume that, outside an event \(E\in\mathcal G\), there are
\(\mathcal G\)-measurable signs \(\eps_i^\star\in\{-1,1\}\) such that
\[
  \left\|\sum_{i=1}^n\eps_i^\star B_i\right\|_{\op}\le K\sqrt n .
\]
Then, for every deterministic \(v\ge0\) and \(s>0\),
\[
\begin{aligned}
  \mathbb P\left\{
  \disc(X_1,\ldots,X_n)>K\sqrt n+s
  \right\}
  &\le
  \mathbb P(E)+\mathbb P\{\mu_\rho>v\} \\
  &\quad+
  2D\exp\left\{-\frac{s^2}{2(v+s/3)}\right\}.
\end{aligned}
\]
In particular, if
\[
  \sum_{i=1}^n\rho_i^2\,\mathbb E(P_i\mid\mcF_{i-1})\preceq V I
  \qquad\text{almost surely}
\]
for deterministic \(V\ge0\), then
\[
  \mathbb P\left\{
  \disc(X_1,\ldots,X_n)>K\sqrt n+s
  \right\}
  \le
  \mathbb P(E)+
  2D\exp\left\{-\frac{s^2}{2(V+s/3)}\right\}.
\]
\end{corollary}

\begin{proof}
Since \(Y_i=P_iY_iP_i\) and \(\|Y_i\|_{\op}\le\rho_i\), we have
\[
  0\le Y_i^2\preceq \rho_i^2 P_i.
\]
Because \(\rho_i\) is \(\mcF_{i-1}\)-measurable,
\[
  \sum_i\mathbb E(Y_i^2\mid\mcF_{i-1})
  \preceq
  \sum_i\rho_i^2\mathbb E(P_i\mid\mcF_{i-1}).
\]
Thus the variation parameter \(\nu\) in
Corollary~\ref{cor:core-plus-matrix-martingale-noise} satisfies
\(\nu\le\mu_\rho\), and that corollary gives the first display.  The deterministic
weighted-overlap specialization follows by taking \(v=V\).
\end{proof}

\begin{corollary}[Signed cores with weighted sampled coordinate or Haar-subspace matrix noise]
\label{cor:core-plus-sampled-coordinate-haar-matrix-noise}
Let \(D\ge2\).  Let \(\mathcal G\), \(B_i\), \((\mcF_i)\), \(P_i\), \(\rho_i\),
\(Y_i\), and \(X_i=B_i+Y_i\) be as in
Corollary~\ref{cor:core-plus-weighted-sampled-localized-matrix-noise}.  Assume
that, outside an event \(E\in\mathcal G\), there are \(\mathcal G\)-measurable
signs \(\eps_i^\star\in\{-1,1\}\) such that
\[
  \left\|\sum_{i=1}^n\eps_i^\star B_i\right\|_{\op}\le K\sqrt n .
\]

Suppose that one of the following two sampled-projection hypotheses holds.
\begin{enumerate}[label=\textup{(\alph*)}]
\item Weighted coordinate supports.  \(P_i=P_{S_i}\) is the coordinate projection
  onto an \(\mcF_i\)-measurable random set \(S_i\subset\{1,\ldots,D\}\).  There
  are \(\mcF_{i-1}\)-measurable random variables \(q_{ij}\in[0,1]\) such that
  \[
    \mathbb P\{j\in S_i\mid\mcF_{i-1}\}\le q_{ij}
    \qquad\text{for every }i,j .
  \]
  In this case put
  \[
    W=\max_{1\le j\le D}\sum_{i=1}^n\rho_i^2q_{ij}.
  \]
\item Weighted Haar subspaces.  There are \(\mcF_{i-1}\)-measurable ranks
  \(r_i\in\{0,\ldots,D\}\) such that, conditionally on \(\mcF_{i-1}\), \(P_i\)
  is Haar-distributed on the Grassmannian of rank-\(r_i\) orthogonal
  projections.  In this case put
  \[
    W=\sum_{i=1}^n\rho_i^2\frac{r_i}{D}.
  \]
\end{enumerate}
Then, for every deterministic \(v\ge0\) and \(s>0\),
\[
\begin{aligned}
  \mathbb P\left\{
  \disc(X_1,\ldots,X_n)>K\sqrt n+s
  \right\}
  &\le
  \mathbb P(E)+\mathbb P\{W>v\}\\
  &\quad+
  2D\exp\left\{-\frac{s^2}{2(v+s/3)}\right\}.
\end{aligned}
\]
In particular, if \(W\le V\) almost surely for a deterministic \(V\ge0\), then
for every \(s>0\),
\[
  \mathbb P\left\{
  \disc(X_1,\ldots,X_n)>K\sqrt n+s
  \right\}
  \le
  \mathbb P(E)+
  2D\exp\left\{-\frac{s^2}{2(V+s/3)}\right\}.
\]
Consequently, adaptive coordinate-supported or conditional Haar-subspace
matrix martingale noise preserves the core signing at the usual weighted
Freedman scale.  The previous unweighted form is recovered by taking
\(\rho_i=1\) and \(q_{ij}=p_i\).  If \(q_{ij}\le p_i\) deterministically and
\(\sum_i\rho_i^2p_i\le V\) almost surely, or if
\(\sum_i\rho_i^2r_i/D\le V\) in the Haar-subspace case, then the same tail holds.
\end{corollary}

\begin{proof}
It is enough to prove that the weighted predictable conditional overlap in
Corollary~\ref{cor:core-plus-weighted-sampled-localized-matrix-noise} is at most
\(W\).  In the coordinate case,
\(\mathbb E(P_i\mid\mcF_{i-1})\) is diagonal with \(j\)-th diagonal entry
\(\mathbb P\{j\in S_i\mid\mcF_{i-1}\}\), hence
\[
  \sum_i\rho_i^2\mathbb E(P_i\mid\mcF_{i-1})
  \preceq
  \operatorname{diag}\left(\sum_i\rho_i^2q_{ij}\right)_{j=1}^D
  \preceq W I .
\]
In the Haar-subspace case, Haar invariance gives
\[
  \mathbb E(P_i\mid\mcF_{i-1})=\frac{r_i}{D}I,
\]
and therefore the same bound is exactly the definition of \(W\).  Thus
\(\mu_\rho\le W\).  Apply
Corollary~\ref{cor:core-plus-weighted-sampled-localized-matrix-noise}; since
\(\{W\le v\}\subseteq\{\mu_\rho\le v\}\), its exceptional term
\(\mathbb P\{\mu_\rho>v\}\) is bounded by \(\mathbb P\{W>v\}\).  The deterministic
budget case follows by taking \(v=V\).
\end{proof}

\begin{corollary}[Almost-sure signed-core sampled localized matrix-noise arrays]
\label{cor:core-plus-sampled-localized-matrix-noise-borel}
For each \(m\ge1\), let \(D_m\ge2\), let \(\mathcal G_m\) be a sigma-field,
and let \((\mcF_{m,i})_{i=0}^{n_m}\) be a filtration with
\(\mathcal G_m\subset\mcF_{m,0}\).  Let
\[
  X_{m,i}=B_{m,i}+Y_{m,i}\in\M_{D_m}(\C)
  \qquad(1\le i\le n_m)
\]
be a signed-core weighted sampled localized matrix-noise row as in
Corollary~\ref{cor:core-plus-weighted-sampled-localized-matrix-noise}: the
\(B_{m,i}\) are \(\mathcal G_m\)-measurable self-adjoint contractions,
\(\rho_{m,i}\in[0,1]\) is \(\mcF_{m,i-1}\)-measurable,
\(P_{m,i}\) is an \(\mcF_{m,i}\)-measurable orthogonal projection, and
\[
  \mathbb E(Y_{m,i}\mid\mcF_{m,i-1})=0,\qquad
  \|Y_{m,i}\|_{\op}\le\rho_{m,i},\qquad
  Y_{m,i}=P_{m,i}Y_{m,i}P_{m,i}.
\]
Put
\[
  \mu_m=
  \left\|
  \sum_{i=1}^{n_m}\rho_{m,i}^2\,
  \mathbb E(P_{m,i}\mid\mcF_{m,i-1})
  \right\|_{\op}.
\]
Assume that, outside an event \(E_m\in\mathcal G_m\), there are
\(\mathcal G_m\)-measurable signs \(\eps_{m,i}^\star\in\{-1,1\}\) such that
\[
  \left\|\sum_{i=1}^{n_m}\eps_{m,i}^\star B_{m,i}\right\|_{\op}
  \le K_m\sqrt{n_m}.
\]
Let \(v_m\ge0\) and \(s_m>0\) be deterministic.  If
\[
  \sum_m
  \left[
  \mathbb P(E_m)+\mathbb P\{\mu_m>v_m\}
  +
  2D_m\exp\left\{-\frac{s_m^2}{2(v_m+s_m/3)}\right\}
  \right]<\infty,
\]
then, with probability one, for all sufficiently large \(m\),
\[
  \disc(X_{m,1},\ldots,X_{m,n_m})
  \le K_m\sqrt{n_m}+s_m .
\]

In particular, put \(\ell_m=\log(eD_m)\).  Suppose that \(K_m\le K\), that
\[
  \mu_m\le Q\frac{n_m}{\ell_m}
\]
outside events \(F_m\), that
\[
  \sum_m[\mathbb P(E_m)+\mathbb P(F_m)]<\infty,
\]
and that, for some \(A,B>0\),
\[
  \sum_m(eD_m)^{-A}<\infty,
  \qquad
  \ell_m\le B\sqrt{n_m}
\]
eventually.  Choose \(L>0\) with
\[
  \frac{L^2}{2(Q+LB/3)}\ge A+1 .
\]
Then, almost surely, for all sufficiently large \(m\),
\[
  \disc(X_{m,1},\ldots,X_{m,n_m})
  \le (K+L)\sqrt{n_m}.
\]

The same conclusions hold with \(\mu_m\) replaced by any random budget \(W_m\)
that dominates the predictable weighted overlap.  In particular, this applies
to the two model classes in
Corollary~\ref{cor:core-plus-sampled-coordinate-haar-matrix-noise}.  For
coordinate supports \(P_{m,i}=P_{S_{m,i}}\), it is enough to have predictable
bounds
\[
  \mathbb P\{j\in S_{m,i}\mid\mcF_{m,i-1}\}\le q_{m,ij},
  \qquad
  W_m=\max_{1\le j\le D_m}\sum_{i=1}^{n_m}\rho_{m,i}^2q_{m,ij}.
\]
For conditional Haar subspaces of predictable rank \(r_{m,i}\), it is enough to
take
\[
  W_m=\sum_{i=1}^{n_m}\rho_{m,i}^2\frac{r_{m,i}}{D_m}.
\]
Thus, along all integer dimensions \(D=2,3,\ldots\), if \(K_D\le K\), if either
\(\mu_D\le Qn_D/\log(eD)\) or the corresponding coordinate/Haar budget
\(W_D\le Qn_D/\log(eD)\) outside summable events, and if
\[
  n_D\ge a\,\log^2(eD)
\]
eventually for some \(a>0\), then for every \(A>1\) and every \(L\) satisfying
\[
  \frac{L^2}{2(Q+L/(3\sqrt a))}\ge A+1
\]
one has
\[
  \disc(X_{D,1},\ldots,X_{D,n_D})\le (K+L)\sqrt{n_D}
\]
eventually almost surely.  If \(D_m\ge\exp(c m^\kappa)\) eventually for some
\(c,\kappa>0\), the same conclusion holds along the sparse schedule for every
\(A>0\) under the same logarithmic-square sample lower bound.  The normalized
contractions \(X_{m,i}/2\) inherit the same almost-sure conclusions with the
displayed thresholds divided by \(2\).
\end{corollary}

\begin{proof}
For the \(m\)-th row, write
\(\Delta_m=\disc(X_{m,1},\ldots,X_{m,n_m})\).
Corollary~\ref{cor:core-plus-weighted-sampled-localized-matrix-noise} gives
\[
\begin{aligned}
  \mathbb P\{\Delta_m>K_m\sqrt{n_m}+s_m\}
  &\le
  \mathbb P(E_m)+\mathbb P\{\mu_m>v_m\}\\
  &\quad+
  2D_m\exp\left\{-\frac{s_m^2}{2(v_m+s_m/3)}\right\}.
\end{aligned}
\]
The first assertion is the first Borel--Cantelli lemma.

For the concrete profile, take \(s_m=L\sqrt{n_m}\) and
\(v_m=Qn_m/\ell_m\).  On \(F_m^c\),
\[
  \frac{s_m^2}{2(v_m+s_m/3)}
  =
  \frac{L^2n_m}{2(Qn_m/\ell_m+L\sqrt{n_m}/3)}
  \ge
  \frac{L^2}{2(Q+LB/3)}\,\ell_m
  \ge (A+1)\ell_m .
\]
Therefore the exponential term is at most \(2(eD_m)^{-A}\), and the assumed
summability again gives Borel--Cantelli.

The coordinate and Haar-subspace budgets dominate \(\mu_m\) by the proof of
Corollary~\ref{cor:core-plus-sampled-coordinate-haar-matrix-noise}.  The
all-integer case is the preceding profile with \(m=D\),
\(B=1/\sqrt a\), and \(\sum_D(eD)^{-A}<\infty\) for \(A>1\).  The sparse case
uses \(\sum_m\exp(-Acm^\kappa)<\infty\), valid for every \(A>0\).  Finally
\(\|X_{m,i}\|_{\op}\le2\), and discrepancy is homogeneous.
\end{proof}

\begin{corollary}[Exhaustive countable menus of signed-core sampled localized matrix-noise arrays]
\label{cor:core-plus-sampled-localized-matrix-noise-countable-exhaustion}
For each \(m\ge1\), let \(D_m\ge2\), put
\(\ell_m=\log(eD_m)\), and let
\[
  J_1\le J_2\le\cdots,\qquad J_m\to\infty,
\]
be a deterministic exhaustion of the countable menu
\(\mathbb N=\{1,2,\ldots\}\).  For each \(m\) and \(h\ge1\), let
\(\mathcal G_m^{(h)}\) be a sigma-field and let
\((\mcF_{m,i}^{(h)})_{i=0}^{n_m}\) be a filtration with
\(\mathcal G_m^{(h)}\subset\mcF_{m,0}^{(h)}\).  Let
\[
  X_{m,i}^{(h)}=B_{m,i}^{(h)}+Y_{m,i}^{(h)}\in\M_{D_m}(\C)
  \qquad(1\le i\le n_m)
\]
be a signed-core weighted sampled localized matrix-noise row as in
Corollary~\ref{cor:core-plus-weighted-sampled-localized-matrix-noise}, with its
own filtration allowed.  Thus \(B_{m,i}^{(h)}\) is
\(\mathcal G_m^{(h)}\)-measurable, \(\rho_{m,i}^{(h)}\in[0,1]\) is
\(\mcF_{m,i-1}^{(h)}\)-measurable, \(P_{m,i}^{(h)}\) is
\(\mcF_{m,i}^{(h)}\)-measurable, and
\[
  \mathbb E(Y_{m,i}^{(h)}\mid\mcF_{m,i-1}^{(h)})=0,\qquad
  \|Y_{m,i}^{(h)}\|_{\op}\le\rho_{m,i}^{(h)},\qquad
  Y_{m,i}^{(h)}=P_{m,i}^{(h)}Y_{m,i}^{(h)}P_{m,i}^{(h)} .
\]
Put
\[
  \mu_{m,h}:=
  \left\|
  \sum_{i=1}^{n_m}(\rho_{m,i}^{(h)})^2\,
  \mathbb E(P_{m,i}^{(h)}\mid\mcF_{m,i-1}^{(h)})
  \right\|_{\op}.
\]
Assume that there is \(K<\infty\) such that, outside events
\(E_{m,h}\in\mathcal G_m^{(h)}\), there are
\(\mathcal G_m^{(h)}\)-measurable signs \(\eps_{m,i}^{(h),\star}\) with
\[
  \left\|\sum_{i=1}^{n_m}\eps_{m,i}^{(h),\star}B_{m,i}^{(h)}\right\|_{\op}
  \le K\sqrt{n_m}
  \qquad(1\le h\le J_m).
\]
Assume also that, for some \(Q,A,B_0>0\),
\[
  \mu_{m,h}\le Q\frac{n_m}{\ell_m}
  \qquad(1\le h\le J_m)
\]
outside events \(F_{m,h}\), and that
\[
  \sum_m\sum_{1\le h\le J_m}
  \bigl[\mathbb P(E_{m,h})+\mathbb P(F_{m,h})\bigr]<\infty,
  \qquad
  \sum_mJ_m(eD_m)^{-A}<\infty,
\]
while
\[
  \ell_m\le B_0\sqrt{n_m}
\]
eventually.  If \(L>0\) satisfies
\[
  \frac{L^2}{2(Q+LB_0/3)}\ge A+1,
\]
then, with probability one, for all sufficiently large \(m\), every active
prefix member \(1\le h\le J_m\) satisfies
\[
  \disc(X_{m,1}^{(h)},\ldots,X_{m,n_m}^{(h)})
  \le (K+L)\sqrt{n_m}.
\]
Consequently, for every fixed \(h\), the same bound holds for all sufficiently
large \(m\).  The normalized contractions \(X_{m,i}^{(h)}/2\) satisfy the same
statements with the discrepancy threshold divided by \(2\).

The same conclusions hold with \(\mu_{m,h}\) replaced by any random budget
\(W_{m,h}\) that dominates the predictable weighted overlap.  In particular,
for active-prefix coordinate supports it is enough to have predictable bounds
\[
  \mathbb P\{j\in S_{m,i}^{(h)}\mid\mcF_{m,i-1}^{(h)}\}\le q_{m,ij}^{(h)},
  \qquad
  W_{m,h}=\max_{1\le j\le D_m}
  \sum_{i=1}^{n_m}(\rho_{m,i}^{(h)})^2q_{m,ij}^{(h)}.
\]
For active-prefix conditional Haar subspaces of predictable ranks
\(r_{m,i}^{(h)}\), it is enough to take
\[
  W_{m,h}=
  \sum_{i=1}^{n_m}(\rho_{m,i}^{(h)})^2\frac{r_{m,i}^{(h)}}{D_m}.
\]

In particular, if \(n_m\ge a\log^2(eD_m)\) eventually, then one may take
\(B_0=a^{-1/2}\).  Along all integer dimensions \(D=2,3,\ldots\), if
\(J_D\le(eD)^\Gamma\) eventually and \(n_D\ge a\log^2(eD)\), then the conclusion
holds for every \(A>\Gamma+1\) and every \(L\) satisfying
\[
  \frac{L^2}{2(Q+L/(3\sqrt a))}\ge A+1 .
\]
Along a sparse schedule with \(D_m\ge\exp(c m^\kappa)\) eventually,
\(J_m\le(eD_m)^\Gamma\) eventually, and
\(n_m\ge a\log^2(eD_m)\), the same conclusion holds for every \(A>\Gamma\) and
every \(L\) satisfying the same display.
\end{corollary}

\begin{proof}
For fixed \(m\) and \(h\le J_m\), apply
Corollary~\ref{cor:core-plus-weighted-sampled-localized-matrix-noise} with
\[
  v_m=Q\frac{n_m}{\ell_m},
  \qquad
  s_m=L\sqrt{n_m}.
\]
Outside \(F_{m,h}\), the weighted-overlap parameter is at most \(v_m\).  Hence
\[
\begin{aligned}
  &\mathbb P\left\{
  \disc(X_{m,1}^{(h)},\ldots,X_{m,n_m}^{(h)})
  >(K+L)\sqrt{n_m}\right\} \\
  &\qquad\le
  \mathbb P(E_{m,h})+\mathbb P(F_{m,h})
  +2D_m\exp\left\{
  -\frac{L^2n_m}{2(Qn_m/\ell_m+L\sqrt{n_m}/3)}
  \right\}.
\end{aligned}
\]
For all sufficiently large \(m\), the exponent is at least
\((A+1)\ell_m\), so the last term is bounded by \(2(eD_m)^{-A}\).  Summing over
\(1\le h\le J_m\) and then over \(m\), the hypotheses make the total failure
probability finite.  Borel--Cantelli gives the active-prefix conclusion, and
\(J_m\to\infty\) gives the fixed-\(h\) statement.

The coordinate and Haar budgets dominate \(\mu_{m,h}\) by the proof of
Corollary~\ref{cor:core-plus-sampled-coordinate-haar-matrix-noise}.  The
normalized statement is homogeneity.  If \(n_m\ge a\log^2(eD_m)\), then
\(\ell_m\le a^{-1/2}\sqrt{n_m}\).  The all-integer and sparse schedule
summability checks are the same as in
Corollary~\ref{cor:core-plus-matrix-martingale-noise-countable-exhaustion}.
\end{proof}

\begin{corollary}[Signed cores with coordinate matrix noise and random amplitudes]
\label{cor:core-plus-bernoulli-coordinate-random-amplitude-noise}
Let \(D\ge2\).  Let \(\mathcal G\), \(B_i\), and the event \(E\in\mathcal G\)
be as in Corollary~\ref{cor:core-plus-weighted-sampled-localized-matrix-noise},
and assume that outside \(E\) there are \(\mathcal G\)-measurable signs
\(\eps_i^\star\in\{-1,1\}\) such that
\[
  \left\|\sum_{i=1}^n\eps_i^\star B_i\right\|_{\op}\le K\sqrt n .
\]
Let \(\rho_1,\ldots,\rho_n\) be independent random variables in \([0,1]\),
independent of \(\mathcal G\), revealed in \(\mcF_0\), and put
\[
  a_i:=\mathbb E\rho_i^2 .
\]
Let \(p_{ij}\in[0,1]\) be deterministic numbers, and suppose that
\((\mcF_i)_{i=0}^n\) is a filtration with
\(\mathcal G,\rho_1,\ldots,\rho_n\subset\mcF_0\).  For each \(i\), let
\(S_i\subset\{1,\ldots,D\}\) be \(\mcF_i\)-measurable and set
\(P_i=P_{S_i}\).  Assume
\[
  \mathbb P\{j\in S_i\mid\mcF_{i-1}\}\le p_{ij}
  \qquad\text{for every }i,j .
\]
Let \(Y_i\) be self-adjoint and \(\mcF_i\)-measurable with
\[
  \mathbb E(Y_i\mid\mcF_{i-1})=0,\qquad
  \|Y_i\|_{\op}\le\rho_i,\qquad
  Y_i=P_iY_iP_i ,
\]
and set \(X_i=B_i+Y_i\).  Define
\[
  \mu:=\max_{1\le j\le D}\sum_{i=1}^n p_{ij}a_i,\qquad
  \sigma^2:=\max_{1\le j\le D}\sum_{i=1}^n p_{ij}^2a_i,\qquad
  p_*:=\max_{i,j}p_{ij}.
\]
Then, for every \(u\ge0\) and \(s>0\),
\[
\begin{aligned}
  \mathbb P\left\{
  \disc(X_1,\ldots,X_n)>K\sqrt n+s
  \right\}
  &\le
  \mathbb P(E)
  +D\exp\left\{-\frac{u^2}{2(\sigma^2+p_*u/3)}\right\}\\
  &\quad+
  2D\exp\left\{-\frac{s^2}{2(\mu+u+s/3)}\right\},
\end{aligned}
\]
with the convention that the first exponential term is omitted when
\(\sigma^2+p_*u/3=0\).  In particular, if the amplitudes are deterministic, or more
generally if \(\max_j\sum_i p_{ij}\rho_i^2\le V\) almost surely, the first
exceptional term may be omitted and one recovers the deterministic budget
\(V\).
\end{corollary}

\begin{proof}
Condition on the amplitudes \(\rho=(\rho_1,\ldots,\rho_n)\).  Since \(\rho\) is
independent of \(\mathcal G\), the core exception contributes \(\mathbb P(E)\)
after averaging.  The preceding corollary applies in the coordinate-support case
with
\[
  W(\rho)=\max_{1\le j\le D}\sum_{i=1}^n p_{ij}\rho_i^2.
\]
Thus, with \(v=\mu+u\),
\[
\begin{aligned}
  \mathbb P\left\{
  \disc(X_1,\ldots,X_n)>K\sqrt n+s
  \right\}
  &\le
  \mathbb P(E)+\mathbb P\{W>\mu+u\}\\
  &\quad+
  2D\exp\left\{-\frac{s^2}{2(\mu+u+s/3)}\right\}.
\end{aligned}
\]
It remains to bound the budget tail.  For each fixed \(j\), the variables
\(Z_{ij}:=p_{ij}\rho_i^2\) are independent, satisfy
\[
  0\le Z_{ij}\le p_*,
  \qquad
  \sum_i\mathbb E Z_{ij}\le\mu,
\]
and
\[
\begin{aligned}
  \sum_i\operatorname{Var}(Z_{ij})
  &\le
  \sum_i\mathbb E Z_{ij}^2
  \le
  \sum_i p_{ij}^2a_i
  \le\sigma^2 .
\end{aligned}
\]
Scalar Bernstein gives
\[
  \mathbb P\left\{\sum_i Z_{ij}>\mu+u\right\}
  \le
  \exp\left\{-\frac{u^2}{2(\sigma^2+p_*u/3)}\right\}.
\]
Taking a union bound over \(j=1,\ldots,D\) gives the displayed estimate.  The
deterministic-budget specialization is just the \(W\le V\) case of
Corollary~\ref{cor:core-plus-sampled-coordinate-haar-matrix-noise}.
\end{proof}

\begin{corollary}[Coordinate matrix noise with random profiles and amplitudes]
\label{cor:core-plus-random-coordinate-profile-random-amplitude-noise}
Let \(D\ge2\).  Let \(\mathcal G\), \(B_i\), and the event \(E\in\mathcal G\)
be as in Corollary~\ref{cor:core-plus-weighted-sampled-localized-matrix-noise},
and assume that outside \(E\) there are \(\mathcal G\)-measurable signs
\(\eps_i^\star\in\{-1,1\}\) such that
\[
  \left\|\sum_{i=1}^n\eps_i^\star B_i\right\|_{\op}\le K\sqrt n .
\]
Let \((Q_i,\rho_i)\), \(i=1,\ldots,n\), be independent pairs, independent of
\(\mathcal G\), where \(Q_i=(Q_{i1},\ldots,Q_{iD})\in[0,1]^D\) and
\(\rho_i\in[0,1]\).  Let \(q_*\in[0,1]\) be deterministic with
\(Q_{ij}\le q_*\) almost surely for every \(i,j\).  Suppose that
\((\mcF_i)_{i=0}^n\) is a filtration with
\(\mathcal G,(Q_i,\rho_i)_{i=1}^n\subset\mcF_0\).  For each \(i\), let
\(S_i\subset\{1,\ldots,D\}\) be \(\mcF_i\)-measurable, set \(P_i=P_{S_i}\), and
assume
\[
  \mathbb P\{j\in S_i\mid\mcF_{i-1}\}\le Q_{ij}
  \qquad\text{for every }i,j .
\]
Let \(Y_i\) be self-adjoint and \(\mcF_i\)-measurable with
\[
  \mathbb E(Y_i\mid\mcF_{i-1})=0,\qquad
  \|Y_i\|_{\op}\le\rho_i,\qquad
  Y_i=P_iY_iP_i ,
\]
and set \(X_i=B_i+Y_i\).  Define
\[
  \mu:=\max_{1\le j\le D}\sum_{i=1}^n\mathbb E(Q_{ij}\rho_i^2),
  \qquad
  \sigma^2:=\max_{1\le j\le D}\sum_{i=1}^n
    \mathbb E(Q_{ij}^2\rho_i^2).
\]
Then, for every \(u\ge0\) and \(s>0\),
\[
\begin{aligned}
  \mathbb P\left\{
  \disc(X_1,\ldots,X_n)>K\sqrt n+s
  \right\}
  &\le
  \mathbb P(E)
  +D\exp\left\{-\frac{u^2}{2(\sigma^2+q_*u/3)}\right\}\\
  &\quad+
  2D\exp\left\{-\frac{s^2}{2(\mu+u+s/3)}\right\},
\end{aligned}
\]
with the convention that the first exponential term is omitted when
\(\sigma^2+q_*u/3=0\).  In particular, if \(Q_{ij}\le p\) almost surely for all
\(i,j\) and \(A=\sum_i\mathbb E\rho_i^2\), then
\[
\begin{aligned}
  \mathbb P\left\{
  \disc(X_1,\ldots,X_n)>K\sqrt n+s
  \right\}
  &\le
  \mathbb P(E)
  +D\exp\left\{-\frac{u^2}{2(p^2A+pu/3)}\right\}\\
  &\quad+
  2D\exp\left\{-\frac{s^2}{2(pA+u+s/3)}\right\}.
\end{aligned}
\]
\end{corollary}

\begin{proof}
Condition on the random profiles and amplitudes.  Since they are independent of
\(\mathcal G\), the core exception contributes \(\mathbb P(E)\) after averaging.
The coordinate case of
Corollary~\ref{cor:core-plus-sampled-coordinate-haar-matrix-noise} applies with
\[
  W=\max_{1\le j\le D}\sum_{i=1}^n Q_{ij}\rho_i^2 .
\]
Thus, with \(v=\mu+u\),
\[
\begin{aligned}
  \mathbb P\left\{
  \disc(X_1,\ldots,X_n)>K\sqrt n+s
  \right\}
  &\le
  \mathbb P(E)+\mathbb P\{W>\mu+u\}\\
  &\quad+
  2D\exp\left\{-\frac{s^2}{2(\mu+u+s/3)}\right\}.
\end{aligned}
\]
For each fixed coordinate \(j\), the variables \(Z_{ij}:=Q_{ij}\rho_i^2\) are
independent and satisfy
\[
\begin{aligned}
  0\le Z_{ij}&\le q_*,
  \qquad
  \sum_i\mathbb E Z_{ij}\le\mu,\\
  \sum_i\operatorname{Var}(Z_{ij})
  &\le\sum_i\mathbb E Z_{ij}^2
  \le\sum_i\mathbb E(Q_{ij}^2\rho_i^2)
  \le\sigma^2 .
\end{aligned}
\]
Scalar Bernstein gives
\[
  \mathbb P\left\{\sum_i Z_{ij}>\mu+u\right\}
  \le
  \exp\left\{-\frac{u^2}{2(\sigma^2+q_*u/3)}\right\}.
\]
Union over \(j=1,\ldots,D\).  The homogeneous-profile display follows by
choosing the jump bound \(q_*=p\), and using
\(\mu\le pA\) and \(\sigma^2\le p^2A\).
\end{proof}

\begin{corollary}[Block-coordinate matrix noise with random profiles and amplitudes]
\label{cor:core-plus-random-block-profile-random-amplitude-noise}
Let \(D\ge2\).  Let \(\Pi_1,\ldots,\Pi_m\) be nonzero pairwise orthogonal
projections on \(\C^D\) with \(\sum_{\ell=1}^m\Pi_\ell=I\).  Let
\(\mathcal G\), \(B_i\), and the event \(E\in\mathcal G\) be as in
Corollary~\ref{cor:core-plus-weighted-sampled-localized-matrix-noise}, and
assume that outside \(E\) there are \(\mathcal G\)-measurable signs
\(\eps_i^\star\in\{-1,1\}\) such that
\[
  \left\|\sum_{i=1}^n\eps_i^\star B_i\right\|_{\op}\le K\sqrt n .
\]
Let \((Q_i,\rho_i)\), \(i=1,\ldots,n\), be independent pairs, independent of
\(\mathcal G\), where \(Q_i=(Q_{i1},\ldots,Q_{im})\in[0,1]^m\) and
\(\rho_i\in[0,1]\).  Let \(q_*\in[0,1]\) be deterministic with
\(Q_{i\ell}\le q_*\) almost surely for every \(i,\ell\).  Suppose that
\((\mcF_i)_{i=0}^n\) is a filtration with
\(\mathcal G,(Q_i,\rho_i)_{i=1}^n\subset\mcF_0\).  For each \(i\), let
\(T_i\subset\{1,\ldots,m\}\) be \(\mcF_i\)-measurable and put
\[
  P_i=\sum_{\ell\in T_i}\Pi_\ell .
\]
Assume
\[
  \mathbb P\{\ell\in T_i\mid\mcF_{i-1}\}\le Q_{i\ell}
  \qquad\text{for every }i,\ell .
\]
Let \(Y_i\) be self-adjoint and \(\mcF_i\)-measurable with
\[
  \mathbb E(Y_i\mid\mcF_{i-1})=0,\qquad
  \|Y_i\|_{\op}\le\rho_i,\qquad
  Y_i=P_iY_iP_i ,
\]
and set \(X_i=B_i+Y_i\).  Define
\[
  \mu:=\max_{1\le \ell\le m}\sum_{i=1}^n\mathbb E(Q_{i\ell}\rho_i^2),
  \qquad
  \sigma^2:=\max_{1\le \ell\le m}\sum_{i=1}^n
    \mathbb E(Q_{i\ell}^2\rho_i^2).
\]
Then, for every \(u\ge0\) and \(s>0\),
\[
\begin{aligned}
  \mathbb P\left\{
  \disc(X_1,\ldots,X_n)>K\sqrt n+s
  \right\}
  &\le
  \mathbb P(E)
  +m\exp\left\{-\frac{u^2}{2(\sigma^2+q_*u/3)}\right\}\\
  &\quad+
  2D\exp\left\{-\frac{s^2}{2(\mu+u+s/3)}\right\},
\end{aligned}
\]
with the convention that the first exponential term is omitted when
\(\sigma^2+q_*u/3=0\).  In particular, if \(Q_{i\ell}\le p\) almost surely for
all \(i,\ell\) and \(A=\sum_i\mathbb E\rho_i^2\), then
\[
\begin{aligned}
  \mathbb P\left\{
  \disc(X_1,\ldots,X_n)>K\sqrt n+s
  \right\}
  &\le
  \mathbb P(E)
  +m\exp\left\{-\frac{u^2}{2(p^2A+pu/3)}\right\}\\
  &\quad+
  2D\exp\left\{-\frac{s^2}{2(pA+u+s/3)}\right\}.
\end{aligned}
\]
\end{corollary}

\begin{proof}
Set
\[
  W=\max_{1\le\ell\le m}\sum_{i=1}^n Q_{i\ell}\rho_i^2 .
\]
Since the \(\Pi_\ell\) are orthogonal projections,
\[
  \mathbb E(P_i\mid\mcF_{i-1})
  =
  \sum_{\ell=1}^m
  \mathbb P\{\ell\in T_i\mid\mcF_{i-1}\}\Pi_\ell
  \preceq
  \sum_{\ell=1}^m Q_{i\ell}\Pi_\ell .
\]
Therefore
\[
  \sum_i\rho_i^2\mathbb E(P_i\mid\mcF_{i-1})
  \preceq
  \sum_{\ell=1}^m\left(\sum_i Q_{i\ell}\rho_i^2\right)\Pi_\ell
  \preceq W I .
\]
Corollary~\ref{cor:core-plus-weighted-sampled-localized-matrix-noise} gives,
with \(v=\mu+u\),
\[
\begin{aligned}
  \mathbb P\left\{
  \disc(X_1,\ldots,X_n)>K\sqrt n+s
  \right\}
  &\le
  \mathbb P(E)+\mathbb P\{W>\mu+u\}\\
  &\quad+
  2D\exp\left\{-\frac{s^2}{2(\mu+u+s/3)}\right\}.
\end{aligned}
\]
For fixed \(\ell\), the variables \(Z_{i\ell}:=Q_{i\ell}\rho_i^2\) are
independent and satisfy
\[
\begin{aligned}
  0\le Z_{i\ell}&\le q_*,
  \qquad
  \sum_i\mathbb E Z_{i\ell}\le\mu,\\
  \sum_i\operatorname{Var}(Z_{i\ell})
  &\le\sum_i\mathbb E Z_{i\ell}^2
  \le\sum_i\mathbb E(Q_{i\ell}^2\rho_i^2)
  \le\sigma^2 .
\end{aligned}
\]
Scalar Bernstein and a union bound over the \(m\) blocks give the displayed
estimate.  The homogeneous display follows from
\(\mu\le pA\), \(\sigma^2\le p^2A\), and \(q_*\le p\).
\end{proof}

\begin{corollary}[Homogeneous coordinate matrix noise with random amplitudes]
\label{cor:core-plus-homogeneous-coordinate-random-amplitude-noise}
In the setting of
Corollary~\ref{cor:core-plus-bernoulli-coordinate-random-amplitude-noise},
suppose in addition that \(p_{ij}\le p\) for all \(i,j\), and put
\[
  A:=\sum_{i=1}^n\mathbb E\rho_i^2 .
\]
Then, for every \(u\ge0\) and \(s>0\),
\[
\begin{aligned}
  \mathbb P\left\{
  \disc(X_1,\ldots,X_n)>K\sqrt n+s
  \right\}
  &\le
  \mathbb P(E)
  +D\exp\left\{-\frac{u^2}{2(p^2A+pu/3)}\right\}\\
  &\quad+
  2D\exp\left\{-\frac{s^2}{2(pA+u+s/3)}\right\}.
\end{aligned}
\]
with the same convention that the first exponential term is omitted when
\(p^2A+pu/3=0\).  Thus, if the average amplitude energy is small, then the coordinate
sampling budget is governed by \(pA\) rather than by \(pn\).  In particular,
taking \(u=pA\) gives the simpler bound
\[
\begin{aligned}
  \mathbb P\left\{
  \disc(X_1,\ldots,X_n)>K\sqrt n+s
  \right\}
  &\le
  \mathbb P(E)+D\exp\{-3A/8\}\\
  &\quad+
  2D\exp\left\{-\frac{s^2}{2(2pA+s/3)}\right\}.
\end{aligned}
\]
\end{corollary}

\begin{proof}
In Corollary~\ref{cor:core-plus-bernoulli-coordinate-random-amplitude-noise} we
have
\[
  \mu\le pA,\qquad \sigma^2\le p^2A,\qquad p_*\le p.
\]
Substitute these bounds in the displayed estimate there.  The last display is
the specialization \(u=pA\); when \(pA=0\), the first budget-exception term can
be omitted, and the displayed \(D\exp\{-3A/8\}\) term is only a harmless weaker
upper bound.
\end{proof}

\begin{corollary}[Haar-subspace matrix noise with random amplitudes]
\label{cor:core-plus-haar-rank-random-amplitude-noise}
Let \(D\ge2\).  Let \(\mathcal G\), \(B_i\), and the event \(E\in\mathcal G\)
be as in Corollary~\ref{cor:core-plus-weighted-sampled-localized-matrix-noise},
and assume that outside \(E\) there are \(\mathcal G\)-measurable signs
\(\eps_i^\star\in\{-1,1\}\) such that
\[
  \left\|\sum_{i=1}^n\eps_i^\star B_i\right\|_{\op}\le K\sqrt n .
\]
Let \(\rho_1,\ldots,\rho_n\) be independent random variables in \([0,1]\),
independent of \(\mathcal G\), revealed in \(\mcF_0\), and put
\[
  a_i:=\mathbb E\rho_i^2 .
\]
Let \(r_i\in\{0,\ldots,D\}\) be deterministic ranks and write
\(\alpha_i=r_i/D\).  Suppose that \((\mcF_i)_{i=0}^n\) is a filtration with
\(\mathcal G,\rho_1,\ldots,\rho_n\subset\mcF_0\).  For each \(i\), let \(P_i\)
be \(\mcF_i\)-measurable and, conditionally on \(\mcF_{i-1}\),
Haar-distributed on the Grassmannian of rank-\(r_i\) orthogonal projections.
Let \(Y_i\) be self-adjoint and \(\mcF_i\)-measurable with
\[
  \mathbb E(Y_i\mid\mcF_{i-1})=0,\qquad
  \|Y_i\|_{\op}\le\rho_i,\qquad
  Y_i=P_iY_iP_i ,
\]
and set \(X_i=B_i+Y_i\).  Define
\[
  \mu:=\sum_{i=1}^n \alpha_i a_i,\qquad
  \sigma^2:=\sum_{i=1}^n \alpha_i^2 a_i,\qquad
  \alpha_*:=\max_i\alpha_i .
\]
Then, for every \(u\ge0\) and \(s>0\),
\[
\begin{aligned}
  \mathbb P\left\{
  \disc(X_1,\ldots,X_n)>K\sqrt n+s
  \right\}
  &\le
  \mathbb P(E)
  +\exp\left\{-\frac{u^2}{2(\sigma^2+\alpha_*u/3)}\right\}\\
  &\quad+
  2D\exp\left\{-\frac{s^2}{2(\mu+u+s/3)}\right\},
\end{aligned}
\]
with the convention that the first exponential term is omitted when
\(\sigma^2+\alpha_*u/3=0\).  In particular, if \(\alpha_i\le\alpha\) for all
\(i\) and \(A=\sum_i\mathbb E\rho_i^2\), then
\[
\begin{aligned}
  \mathbb P\left\{
  \disc(X_1,\ldots,X_n)>K\sqrt n+s
  \right\}
  &\le
  \mathbb P(E)
  +\exp\left\{-\frac{u^2}{2(\alpha^2A+\alpha u/3)}\right\}\\
  &\quad+
  2D\exp\left\{-\frac{s^2}{2(\alpha A+u+s/3)}\right\}.
\end{aligned}
\]
Taking \(u=\alpha A\) gives the simpler bound
\[
\begin{aligned}
  \mathbb P\left\{
  \disc(X_1,\ldots,X_n)>K\sqrt n+s
  \right\}
  &\le
  \mathbb P(E)+\exp\{-3A/8\}\\
  &\quad+
  2D\exp\left\{-\frac{s^2}{2(2\alpha A+s/3)}\right\}.
\end{aligned}
\]
\end{corollary}

\begin{proof}
Condition on the amplitudes \(\rho=(\rho_1,\ldots,\rho_n)\).  The
Haar-subspace case of
Corollary~\ref{cor:core-plus-sampled-coordinate-haar-matrix-noise} applies with
\[
  W(\rho)=\sum_{i=1}^n\alpha_i\rho_i^2 .
\]
Thus, with \(v=\mu+u\),
\[
\begin{aligned}
  \mathbb P\left\{
  \disc(X_1,\ldots,X_n)>K\sqrt n+s
  \right\}
  &\le
  \mathbb P(E)+\mathbb P\{W>\mu+u\}\\
  &\quad+
  2D\exp\left\{-\frac{s^2}{2(\mu+u+s/3)}\right\}.
\end{aligned}
\]
The variables \(Z_i:=\alpha_i\rho_i^2\) are independent and satisfy
\[
  0\le Z_i\le\alpha_*,
  \qquad
  \sum_i\mathbb E Z_i=\mu,
  \qquad
  \sum_i\operatorname{Var}(Z_i)\le\sum_i\mathbb E Z_i^2\le\sigma^2 .
\]
Scalar Bernstein gives
\[
  \mathbb P\{W>\mu+u\}
  \le
  \exp\left\{-\frac{u^2}{2(\sigma^2+\alpha_*u/3)}\right\}.
\]
This proves the first display.  The homogeneous-rank display follows from
\(\mu\le\alpha A\), \(\sigma^2\le\alpha^2A\), and \(\alpha_*\le\alpha\); the
last display is the specialization \(u=\alpha A\), with the same harmless
zero-denominator convention as above.
\end{proof}

\begin{corollary}[Random-rank Haar-subspace matrix noise with random amplitudes]
\label{cor:core-plus-random-rank-haar-random-amplitude-noise}
Let \(D\ge2\).  Let \(\mathcal G\), \(B_i\), and the event \(E\in\mathcal G\)
be as in Corollary~\ref{cor:core-plus-weighted-sampled-localized-matrix-noise},
and assume that outside \(E\) there are \(\mathcal G\)-measurable signs
\(\eps_i^\star\in\{-1,1\}\) such that
\[
  \left\|\sum_{i=1}^n\eps_i^\star B_i\right\|_{\op}\le K\sqrt n .
\]
Let \((R_i,\rho_i)\), \(i=1,\ldots,n\), be independent pairs, independent of
\(\mathcal G\), with \(R_i\in\{0,\ldots,D\}\) and \(\rho_i\in[0,1]\).  Put
\(\alpha_i=R_i/D\), and let \(\alpha_*\in[0,1]\) be deterministic with
\(\alpha_i\le\alpha_*\) almost surely for every \(i\).  Suppose that
\((\mcF_i)_{i=0}^n\) is a filtration with
\(\mathcal G,(R_i,\rho_i)_{i=1}^n\subset\mcF_0\).  For each \(i\), let \(P_i\)
be \(\mcF_i\)-measurable and, conditionally on \(\mcF_{i-1}\),
Haar-distributed on the Grassmannian of rank-\(R_i\) orthogonal projections.
Let \(Y_i\) be self-adjoint and \(\mcF_i\)-measurable with
\[
  \mathbb E(Y_i\mid\mcF_{i-1})=0,\qquad
  \|Y_i\|_{\op}\le\rho_i,\qquad
  Y_i=P_iY_iP_i ,
\]
and set \(X_i=B_i+Y_i\).  Define
\[
  \mu:=\sum_{i=1}^n \mathbb E(\alpha_i\rho_i^2),
  \qquad
  \sigma^2:=\sum_{i=1}^n \mathbb E(\alpha_i^2\rho_i^2).
\]
Then, for every \(u\ge0\) and \(s>0\),
\[
\begin{aligned}
  \mathbb P\left\{
  \disc(X_1,\ldots,X_n)>K\sqrt n+s
  \right\}
  &\le
  \mathbb P(E)
  +\exp\left\{-\frac{u^2}{2(\sigma^2+\alpha_*u/3)}\right\}\\
  &\quad+
  2D\exp\left\{-\frac{s^2}{2(\mu+u+s/3)}\right\},
\end{aligned}
\]
with the convention that the first exponential term is omitted when
\(\sigma^2+\alpha_*u/3=0\).  In particular, if \(\alpha_i\le\alpha\) almost
surely for all \(i\) and \(A=\sum_i\mathbb E\rho_i^2\), then
\[
\begin{aligned}
  \mathbb P\left\{
  \disc(X_1,\ldots,X_n)>K\sqrt n+s
  \right\}
  &\le
  \mathbb P(E)
  +\exp\left\{-\frac{u^2}{2(\alpha^2A+\alpha u/3)}\right\}\\
  &\quad+
  2D\exp\left\{-\frac{s^2}{2(\alpha A+u+s/3)}\right\}.
\end{aligned}
\]
\end{corollary}

\begin{proof}
Condition on the random ranks and amplitudes.  The Haar-subspace case of
Corollary~\ref{cor:core-plus-sampled-coordinate-haar-matrix-noise} applies with
\[
  W=\sum_{i=1}^n \alpha_i\rho_i^2 .
\]
Therefore, with \(v=\mu+u\),
\[
\begin{aligned}
  \mathbb P\left\{
  \disc(X_1,\ldots,X_n)>K\sqrt n+s
  \right\}
  &\le
  \mathbb P(E)+\mathbb P\{W>\mu+u\}\\
  &\quad+
  2D\exp\left\{-\frac{s^2}{2(\mu+u+s/3)}\right\}.
\end{aligned}
\]
The variables \(Z_i:=\alpha_i\rho_i^2\) are independent and satisfy
\[
\begin{aligned}
  0\le Z_i&\le\alpha_*,
  \qquad
  \sum_i\mathbb E Z_i=\mu,\\
  \sum_i\operatorname{Var}(Z_i)
  &\le\sum_i\mathbb E Z_i^2
  \le\sum_i\mathbb E(\alpha_i^2\rho_i^2)=\sigma^2 .
\end{aligned}
\]
Scalar Bernstein gives the displayed budget-exception bound.  The final display
follows from \(\mu\le\alpha A\), \(\sigma^2\le\alpha^2A\), and
\(\alpha_*\le\alpha\).
\end{proof}

\begin{corollary}[Isotropic projection matrix noise with random amplitudes]
\label{cor:core-plus-isotropic-projection-random-amplitude-noise}
Let \(D\ge2\).  Let \(\mathcal G\), \(B_i\), and the event \(E\in\mathcal G\)
be as in Corollary~\ref{cor:core-plus-weighted-sampled-localized-matrix-noise},
and assume that outside \(E\) there are \(\mathcal G\)-measurable signs
\(\eps_i^\star\in\{-1,1\}\) such that
\[
  \left\|\sum_{i=1}^n\eps_i^\star B_i\right\|_{\op}\le K\sqrt n .
\]
Let \((\alpha_i,\rho_i)\), \(i=1,\ldots,n\), be independent pairs,
independent of \(\mathcal G\), with \(\alpha_i,\rho_i\in[0,1]\).  Let
\(\alpha_*\in[0,1]\) be deterministic with \(\alpha_i\le\alpha_*\) almost
surely for every \(i\).  Suppose that \((\mcF_i)_{i=0}^n\) is a filtration with
\(\mathcal G,(\alpha_i,\rho_i)_{i=1}^n\subset\mcF_0\).  For each \(i\), let
\(P_i\) be an \(\mcF_i\)-measurable orthogonal projection such that
\[
  \mathbb E(P_i\mid\mcF_{i-1})=\alpha_i I .
\]
Let \(Y_i\) be self-adjoint and \(\mcF_i\)-measurable with
\[
  \mathbb E(Y_i\mid\mcF_{i-1})=0,\qquad
  \|Y_i\|_{\op}\le\rho_i,\qquad
  Y_i=P_iY_iP_i ,
\]
and set \(X_i=B_i+Y_i\).  Define
\[
  \mu:=\sum_{i=1}^n\mathbb E(\alpha_i\rho_i^2),
  \qquad
  \sigma^2:=\sum_{i=1}^n\mathbb E(\alpha_i^2\rho_i^2).
\]
Then, for every \(u\ge0\) and \(s>0\),
\[
\begin{aligned}
  \mathbb P\left\{
  \disc(X_1,\ldots,X_n)>K\sqrt n+s
  \right\}
  &\le
  \mathbb P(E)
  +\exp\left\{-\frac{u^2}{2(\sigma^2+\alpha_*u/3)}\right\}\\
  &\quad+
  2D\exp\left\{-\frac{s^2}{2(\mu+u+s/3)}\right\},
\end{aligned}
\]
with the convention that the first exponential term is omitted when
\(\sigma^2+\alpha_*u/3=0\).  In particular, if \(\alpha_i\le\alpha\) almost
surely for all \(i\) and \(A=\sum_i\mathbb E\rho_i^2\), then
\[
\begin{aligned}
  \mathbb P\left\{
  \disc(X_1,\ldots,X_n)>K\sqrt n+s
  \right\}
  &\le
  \mathbb P(E)
  +\exp\left\{-\frac{u^2}{2(\alpha^2A+\alpha u/3)}\right\}\\
  &\quad+
  2D\exp\left\{-\frac{s^2}{2(\alpha A+u+s/3)}\right\}.
\end{aligned}
\]
\end{corollary}

\begin{proof}
The conditional isotropy assumption gives
\[
  \sum_i\rho_i^2\mathbb E(P_i\mid\mcF_{i-1})
  =
  \left(\sum_i\alpha_i\rho_i^2\right)I .
\]
Therefore Corollary~\ref{cor:core-plus-weighted-sampled-localized-matrix-noise}
applies with
\[
  W=\sum_{i=1}^n\alpha_i\rho_i^2 .
\]
Taking \(v=\mu+u\) reduces the matter to a scalar budget tail for \(W\).  The
variables \(Z_i:=\alpha_i\rho_i^2\) are independent and satisfy
\[
\begin{aligned}
  0\le Z_i&\le\alpha_*,
  \qquad
  \sum_i\mathbb E Z_i=\mu,\\
  \sum_i\operatorname{Var}(Z_i)
  &\le\sum_i\mathbb E Z_i^2
  \le\sum_i\mathbb E(\alpha_i^2\rho_i^2)=\sigma^2 .
\end{aligned}
\]
Scalar Bernstein gives the first display.  The homogeneous display follows from
\(\mu\le\alpha A\), \(\sigma^2\le\alpha^2A\), and \(\alpha_*\le\alpha\).
\end{proof}

\begin{corollary}[Spectrally dominated projection matrix noise with random amplitudes]
\label{cor:core-plus-dominated-projection-random-amplitude-noise}
Let \(D\ge2\).  Let \(\mathcal G\), \(B_i\), and the event \(E\in\mathcal G\)
be as in Corollary~\ref{cor:core-plus-weighted-sampled-localized-matrix-noise},
and assume that outside \(E\) there are \(\mathcal G\)-measurable signs
\(\eps_i^\star\in\{-1,1\}\) such that
\[
  \left\|\sum_{i=1}^n\eps_i^\star B_i\right\|_{\op}\le K\sqrt n .
\]
Let \((\alpha_i,\rho_i)\), \(i=1,\ldots,n\), be independent pairs,
independent of \(\mathcal G\), with \(\alpha_i,\rho_i\in[0,1]\).  Let
\(\alpha_*\in[0,1]\) be deterministic with \(\alpha_i\le\alpha_*\) almost
surely for every \(i\).  Suppose that \((\mcF_i)_{i=0}^n\) is a filtration with
\(\mathcal G,(\alpha_i,\rho_i)_{i=1}^n\subset\mcF_0\).  For each \(i\), let
\(P_i\) be an \(\mcF_i\)-measurable orthogonal projection such that
\[
  \mathbb E(P_i\mid\mcF_{i-1})\preceq \alpha_i I .
\]
Let \(Y_i\) be self-adjoint and \(\mcF_i\)-measurable with
\[
  \mathbb E(Y_i\mid\mcF_{i-1})=0,\qquad
  \|Y_i\|_{\op}\le\rho_i,\qquad
  Y_i=P_iY_iP_i ,
\]
and set \(X_i=B_i+Y_i\).  Define
\[
  \mu:=\sum_{i=1}^n\mathbb E(\alpha_i\rho_i^2),
  \qquad
  \sigma^2:=\sum_{i=1}^n\mathbb E(\alpha_i^2\rho_i^2).
\]
Then, for every \(u\ge0\) and \(s>0\),
\[
\begin{aligned}
  \mathbb P\left\{
  \disc(X_1,\ldots,X_n)>K\sqrt n+s
  \right\}
  &\le
  \mathbb P(E)
  +\exp\left\{-\frac{u^2}{2(\sigma^2+\alpha_*u/3)}\right\}\\
  &\quad+
  2D\exp\left\{-\frac{s^2}{2(\mu+u+s/3)}\right\},
\end{aligned}
\]
with the convention that the first exponential term is omitted when
\(\sigma^2+\alpha_*u/3=0\).  In particular, if \(\alpha_i\le\alpha\) almost
surely for all \(i\) and \(A=\sum_i\mathbb E\rho_i^2\), then
\[
\begin{aligned}
  \mathbb P\left\{
  \disc(X_1,\ldots,X_n)>K\sqrt n+s
  \right\}
  &\le
  \mathbb P(E)
  +\exp\left\{-\frac{u^2}{2(\alpha^2A+\alpha u/3)}\right\}\\
  &\quad+
  2D\exp\left\{-\frac{s^2}{2(\alpha A+u+s/3)}\right\}.
\end{aligned}
\]
\end{corollary}

\begin{proof}
The domination assumption gives
\[
  \sum_i\rho_i^2\mathbb E(P_i\mid\mcF_{i-1})
  \preceq
  \left(\sum_i\alpha_i\rho_i^2\right)I .
\]
Thus Corollary~\ref{cor:core-plus-weighted-sampled-localized-matrix-noise}
applies with \(W=\sum_i\alpha_i\rho_i^2\).  Taking \(v=\mu+u\), it remains only
to bound \(\mathbb P\{W>\mu+u\}\).  The variables
\(Z_i:=\alpha_i\rho_i^2\) are independent and satisfy
\[
\begin{aligned}
  0\le Z_i&\le\alpha_*,
  \qquad
  \sum_i\mathbb E Z_i=\mu,\\
  \sum_i\operatorname{Var}(Z_i)
  &\le\sum_i\mathbb E Z_i^2
  \le\sum_i\mathbb E(\alpha_i^2\rho_i^2)=\sigma^2 .
\end{aligned}
\]
Scalar Bernstein gives the first display.  The homogeneous display follows from
\(\mu\le\alpha A\), \(\sigma^2\le\alpha^2A\), and \(\alpha_*\le\alpha\).
\end{proof}

\begin{corollary}[Mixed coordinate or Haar-subspace matrix noise with random amplitudes]
\label{cor:core-plus-mixed-coordinate-haar-random-amplitude-noise}
Let \(D\ge2\).  Let \(\mathcal G\), \(B_i\), and the event \(E\in\mathcal G\)
be as in Corollary~\ref{cor:core-plus-weighted-sampled-localized-matrix-noise},
and assume that outside \(E\) there are \(\mathcal G\)-measurable signs
\(\eps_i^\star\in\{-1,1\}\) such that
\[
  \left\|\sum_{i=1}^n\eps_i^\star B_i\right\|_{\op}\le K\sqrt n .
\]
Let \((M_i,Q_i,R_i,\rho_i)\), \(i=1,\ldots,n\), be independent tuples,
independent of \(\mathcal G\), where
\[
  \begin{gathered}
  M_i\in\{0,1\},\qquad R_i\in\{0,\ldots,D\},\qquad \rho_i\in[0,1],\\
  Q_i=(Q_{i1},\ldots,Q_{iD})\in[0,1]^D .
  \end{gathered}
\]
Here \(M_i=1\) denotes the coordinate-support branch and \(M_i=0\) denotes the
Haar-subspace branch.  Put
\[
  \beta_{ij}:=M_iQ_{ij}+(1-M_i)\frac{R_i}{D},
\]
and let \(b_*\in[0,1]\) be deterministic with \(\beta_{ij}\le b_*\) almost
surely for every \(i,j\).  Suppose that \((\mcF_i)_{i=0}^n\) is a filtration
with \(\mathcal G,(M_i,Q_i,R_i,\rho_i)_{i=1}^n\subset\mcF_0\).  For each \(i\),
let \(P_i\) be an \(\mcF_i\)-measurable orthogonal projection.  On the event
\(\{M_i=1\}\), assume that there is an \(\mcF_i\)-measurable set
\(S_i\subset\{1,\ldots,D\}\) such that \(P_i=P_{S_i}\) and
\[
  \mathbb P\{j\in S_i\mid\mcF_{i-1}\}\le Q_{ij}
  \qquad\text{for every }j .
\]
On the event \(\{M_i=0\}\), assume that, conditionally on \(\mcF_{i-1}\), \(P_i\)
is Haar-distributed on the Grassmannian of rank-\(R_i\) orthogonal projections.
Let \(Y_i\) be self-adjoint and \(\mcF_i\)-measurable with
\[
  \mathbb E(Y_i\mid\mcF_{i-1})=0,\qquad
  \|Y_i\|_{\op}\le\rho_i,\qquad
  Y_i=P_iY_iP_i ,
\]
and set \(X_i=B_i+Y_i\).  Define
\[
  \mu:=\max_{1\le j\le D}\sum_{i=1}^n\mathbb E(\beta_{ij}\rho_i^2),
  \qquad
  \sigma^2:=\max_{1\le j\le D}\sum_{i=1}^n
    \mathbb E(\beta_{ij}^2\rho_i^2).
\]
Then, for every \(u\ge0\) and \(s>0\),
\[
\begin{aligned}
  \mathbb P\left\{
  \disc(X_1,\ldots,X_n)>K\sqrt n+s
  \right\}
  &\le
  \mathbb P(E)
  +D\exp\left\{-\frac{u^2}{2(\sigma^2+b_*u/3)}\right\}\\
  &\quad+
  2D\exp\left\{-\frac{s^2}{2(\mu+u+s/3)}\right\},
\end{aligned}
\]
with the convention that the first exponential term is omitted when
\(\sigma^2+b_*u/3=0\).  In particular, if \(\beta_{ij}\le b\) almost surely for
all \(i,j\) and \(A=\sum_i\mathbb E\rho_i^2\), then
\[
\begin{aligned}
  \mathbb P\left\{
  \disc(X_1,\ldots,X_n)>K\sqrt n+s
  \right\}
  &\le
  \mathbb P(E)
  +D\exp\left\{-\frac{u^2}{2(b^2A+bu/3)}\right\}\\
  &\quad+
  2D\exp\left\{-\frac{s^2}{2(bA+u+s/3)}\right\}.
\end{aligned}
\]
\end{corollary}

\begin{proof}
Set
\[
  W=\max_{1\le j\le D}\sum_{i=1}^n\beta_{ij}\rho_i^2 .
\]
We first show that \(W\) controls the weighted conditional overlap in
Corollary~\ref{cor:core-plus-weighted-sampled-localized-matrix-noise}.  Because
\((M_i,Q_i,R_i,\rho_i)\) is \(\mcF_{i-1}\)-measurable, the coordinate branch
gives
\[
  \mathbb E(P_i\mid\mcF_{i-1})
  \preceq \operatorname{diag}(Q_{ij})_{j=1}^D
  \qquad\text{on }\{M_i=1\},
\]
while Haar invariance gives
\[
  \mathbb E(P_i\mid\mcF_{i-1})=\frac{R_i}{D}I
  \qquad\text{on }\{M_i=0\}.
\]
Thus, pointwise,
\[
  \mathbb E(P_i\mid\mcF_{i-1})
  \preceq \operatorname{diag}(\beta_{ij})_{j=1}^D .
\]
Consequently
\[
  \sum_i\rho_i^2\mathbb E(P_i\mid\mcF_{i-1})
  \preceq
  \operatorname{diag}\left(\sum_i\beta_{ij}\rho_i^2\right)_{j=1}^D
  \preceq W I .
\]
Corollary~\ref{cor:core-plus-weighted-sampled-localized-matrix-noise} therefore
applies with \(v=\mu+u\) and gives
\[
\begin{aligned}
  \mathbb P\left\{
  \disc(X_1,\ldots,X_n)>K\sqrt n+s
  \right\}
  &\le
  \mathbb P(E)+\mathbb P\{W>\mu+u\}\\
  &\quad+
  2D\exp\left\{-\frac{s^2}{2(\mu+u+s/3)}\right\}.
\end{aligned}
\]
It remains to control the random budget.  For fixed \(j\), the variables
\(Z_{ij}:=\beta_{ij}\rho_i^2\) are independent and satisfy
\[
\begin{aligned}
  0\le Z_{ij}&\le b_*,
  \qquad
  \sum_i\mathbb E Z_{ij}\le\mu,\\
  \sum_i\operatorname{Var}(Z_{ij})
  &\le\sum_i\mathbb E Z_{ij}^2
  \le\sum_i\mathbb E(\beta_{ij}^2\rho_i^2)
  \le\sigma^2 .
\end{aligned}
\]
Scalar Bernstein gives
\[
  \mathbb P\left\{\sum_i Z_{ij}>\mu+u\right\}
  \le
  \exp\left\{-\frac{u^2}{2(\sigma^2+b_*u/3)}\right\}.
\]
Union over \(j=1,\ldots,D\).  The homogeneous display follows from
\(\mu\le bA\), \(\sigma^2\le b^2A\), and \(b_*\le b\).
\end{proof}

\begin{corollary}[Almost-sure random-amplitude localized matrix-noise menus]
\label{cor:core-plus-localized-random-amplitude-menu-borel}
For each \(m\ge1\), let \(D_m\ge2\), and let
\[
  X_{m,i}=B_{m,i}+Y_{m,i}\in\M_{D_m}(\C),
  \qquad 1\le i\le n_m,
\]
be a signed-core localized matrix-noise row with independent random profile or
amplitude data as in one of
Corollaries~\ref{cor:core-plus-bernoulli-coordinate-random-amplitude-noise},
\ref{cor:core-plus-random-coordinate-profile-random-amplitude-noise},
\ref{cor:core-plus-random-block-profile-random-amplitude-noise},
\ref{cor:core-plus-haar-rank-random-amplitude-noise},
\ref{cor:core-plus-random-rank-haar-random-amplitude-noise},
\ref{cor:core-plus-isotropic-projection-random-amplitude-noise},
\ref{cor:core-plus-dominated-projection-random-amplitude-noise}, or
\ref{cor:core-plus-mixed-coordinate-haar-random-amplitude-noise}.  Thus,
outside an event \(E_m\), the core matrices have
\(\mathcal G_m\)-measurable signs with
\[
  \left\|\sum_{i=1}^{n_m}\eps_{m,i}^\star B_{m,i}\right\|_{\op}
  \le K_m\sqrt{n_m}.
\]
Let \(N_m\ge1\), \(\mu_m\ge0\), \(\sigma_m^2\ge0\), and \(b_m\in[0,1]\) be
the budget-union count, mean budget, Bernstein variance proxy, and jump
parameter in the corresponding one-row corollary.  Concretely, \(N_m=D_m\) for
coordinate-profile and mixed coordinate/Haar menus, \(N_m\) is the number of
structural blocks in the block-profile menu, and \(N_m=1\) for Haar-rank,
isotropic, or spectrally dominated projection menus.

Let \(u_m\ge0\) and \(s_m>0\) be deterministic.  If
\[
  \sum_m
  \left[
  \mathbb P(E_m)
  +N_m\exp\left\{-\frac{u_m^2}{2(\sigma_m^2+b_mu_m/3)}\right\}
  +2D_m\exp\left\{-\frac{s_m^2}{2(\mu_m+u_m+s_m/3)}\right\}
  \right]<\infty,
\]
with the usual convention that the Bernstein term is omitted when its
denominator is zero, then, with probability one, for all sufficiently large
\(m\),
\[
  \disc(X_{m,1},\ldots,X_{m,n_m})
  \le K_m\sqrt{n_m}+s_m .
\]

In particular, put \(\ell_m=\log(eD_m)\).  Suppose that \(K_m\le K\), that
\[
  \mu_m+u_m\le Q\frac{n_m}{\ell_m}
\]
eventually, and that, for some \(A,B>0\),
\[
  \sum_m\left[
  \mathbb P(E_m)
  +N_m\exp\left\{-\frac{u_m^2}{2(\sigma_m^2+b_mu_m/3)}\right\}
  +(eD_m)^{-A}
  \right]<\infty,
  \qquad
  \ell_m\le B\sqrt{n_m}
\]
eventually.  Choose \(L>0\) with
\[
  \frac{L^2}{2(Q+LB/3)}\ge A+1 .
\]
Then, almost surely, for all sufficiently large \(m\),
\[
  \disc(X_{m,1},\ldots,X_{m,n_m})
  \le (K+L)\sqrt{n_m}.
\]

A fully explicit half-budget schedule is obtained as follows.  Assume
\(\sum_m\mathbb P(E_m)<\infty\) and
\[
  \mu_m\le \frac{Qn_m}{2\ell_m},\qquad
  \sigma_m^2\le S\frac{n_m}{\ell_m},\qquad
  b_m\le b,\qquad
  N_m\le(eD_m)^\Gamma
\]
eventually, and set
\[
  u_m=\frac{Qn_m}{2\ell_m}.
\]
If \(n_m\ge a\ell_m^2\) eventually, then the random-budget Bernstein exception
is bounded by
\[
  (eD_m)^{\Gamma-\delta},
  \qquad
  \delta=\frac{Q^2a}{8S+\frac43 bQ}.
\]
Consequently, along all integer dimensions \(D=2,3,\ldots\), if
\(\delta-\Gamma>1\), then for every \(A>1\) and every \(L\) satisfying
\[
  \frac{L^2}{2(Q+L/(3\sqrt a))}\ge A+1
\]
one has
\[
  \disc(X_{D,1},\ldots,X_{D,n_D})\le (K+L)\sqrt{n_D}
\]
eventually almost surely.  Along any sparse schedule
\(D_m\ge\exp(c m^\kappa)\), the same conclusion holds whenever
\(\delta-\Gamma>0\), for every \(A>0\) and every \(L\) satisfying the same
display.
These criteria apply directly to homogeneous coordinate noise
\((\mu,\sigma^2,b)=(pA,p^2A,p)\), Haar or dominated projection noise
\((\mu,\sigma^2,b)=(\alpha A,\alpha^2A,\alpha)\), block-profile menus with
the block count in \(N_m\), and the mixed coordinate/Haar menu with
\((\mu,\sigma^2,b)=(bA,b^2A,b)\) after the corresponding homogeneous bounds.
\end{corollary}

\begin{proof}
Each listed one-row corollary has the common form
\[
\begin{aligned}
  \mathbb P\left\{
  \disc(X_{m,1},\ldots,X_{m,n_m})>K_m\sqrt{n_m}+s_m
  \right\}
  &\le
  \mathbb P(E_m)
  +N_m\exp\left\{-\frac{u_m^2}{2(\sigma_m^2+b_mu_m/3)}\right\}\\
  &\quad+
  2D_m\exp\left\{-\frac{s_m^2}{2(\mu_m+u_m+s_m/3)}\right\}.
\end{aligned}
\]
The first assertion is the first Borel--Cantelli lemma.

For the profile form, take \(s_m=L\sqrt{n_m}\).  On the event
\(\mu_m+u_m\le Qn_m/\ell_m\),
\[
  \frac{s_m^2}{2(\mu_m+u_m+s_m/3)}
  \ge
  \frac{L^2n_m}{2(Qn_m/\ell_m+L\sqrt{n_m}/3)}
  \ge
  \frac{L^2}{2(Q+LB/3)}\,\ell_m
  \ge (A+1)\ell_m .
\]
Thus the matrix-Freedman term is at most \(2(eD_m)^{-A}\), and the stated
summability gives the almost-sure conclusion.

For the explicit half-budget schedule, the choice
\(u_m=Qn_m/(2\ell_m)\) gives
\[
  \frac{u_m^2}{2(\sigma_m^2+b_mu_m/3)}
  \ge
  \frac{Q^2 n_m}{(8S+\frac43 bQ)\ell_m}
  \ge \delta\,\ell_m .
\]
Hence the budget exception is at most
\((eD_m)^\Gamma e^{-\delta\ell_m}=(eD_m)^{\Gamma-\delta}\).  The all-integer
and sparse summability conditions follow from the usual \(p\)-series and
stretched-exponential series tests.  The final examples are obtained by
substituting the homogeneous parameter bounds from the corresponding one-row
corollaries.
\end{proof}

\begin{corollary}[Exhaustive countable random-amplitude localized matrix-noise menus]
\label{cor:core-plus-countable-localized-random-amplitude-menu-borel}
For each \(m\ge1\), let \(D_m\ge2\), put \(\ell_m=\log(eD_m)\), and let
\[
  J_1\le J_2\le\cdots,\qquad J_m\to\infty,
\]
be a deterministic exhaustion of the countable menu
\(\mathbb N=\{1,2,\ldots\}\).  For each \(m\) and \(h\ge1\), let
\[
  X_{m,i}^{(h)}=B_{m,i}^{(h)}+Y_{m,i}^{(h)}\in\M_{D_m}(\C),
  \qquad 1\le i\le n_m,
\]
be a signed-core localized matrix-noise row with independent random profile or
amplitude data as in one of the one-row corollaries listed in
Corollary~\ref{cor:core-plus-localized-random-amplitude-menu-borel}, with its
own filtration and random environment allowed.  Thus, outside an event
\(E_{m,h}\), the \(h\)-th core has environment-measurable signs satisfying
\[
  \left\|\sum_{i=1}^{n_m}\eps_{m,i}^{(h),\star}B_{m,i}^{(h)}\right\|_{\op}
  \le K_{m,h}\sqrt{n_m}.
\]
Let \(N_{m,h}\ge1\), \(\mu_{m,h}\ge0\), \(\sigma_{m,h}^2\ge0\), and
\(b_{m,h}\in[0,1]\) be the budget-union count, mean budget, Bernstein variance
proxy, and jump parameter supplied by the corresponding one-row corollary.
Let \(u_{m,h}\ge0\) and \(s_m>0\) be deterministic.  If
\[
  \sum_m\sum_{1\le h\le J_m}
  \left[
  \mathbb P(E_{m,h})
  +N_{m,h}\exp\left\{
    -\frac{u_{m,h}^2}{2(\sigma_{m,h}^2+b_{m,h}u_{m,h}/3)}
  \right\}
  +2D_m\exp\left\{
    -\frac{s_m^2}{2(\mu_{m,h}+u_{m,h}+s_m/3)}
  \right\}
  \right]<\infty,
\]
with the usual convention that the Bernstein term is omitted when its
denominator is zero, then, with probability one, for all sufficiently large
\(m\), every active-prefix element \(1\le h\le J_m\) satisfies
\[
  \disc(X_{m,1}^{(h)},\ldots,X_{m,n_m}^{(h)})
  \le K_{m,h}\sqrt{n_m}+s_m .
\]
Consequently, each fixed \(h\) is signed at that scale for all sufficiently
large \(m\).

In particular, suppose that \(K_{m,h}\le K\) for \(1\le h\le J_m\), that
\[
  \mu_{m,h}+u_{m,h}\le Q\frac{n_m}{\ell_m}
  \qquad(1\le h\le J_m)
\]
eventually, and that, for some \(A,B>0\),
\[
  \sum_m\sum_{1\le h\le J_m}
  \left[
  \mathbb P(E_{m,h})
  +N_{m,h}\exp\left\{
    -\frac{u_{m,h}^2}{2(\sigma_{m,h}^2+b_{m,h}u_{m,h}/3)}
  \right\}
  \right]
  +
  \sum_mJ_m(eD_m)^{-A}<\infty,
\]
while \(\ell_m\le B\sqrt{n_m}\) eventually.  If \(L>0\) satisfies
\[
  \frac{L^2}{2(Q+LB/3)}\ge A+1,
\]
then, almost surely, for all sufficiently large \(m\),
\[
  \disc(X_{m,1}^{(h)},\ldots,X_{m,n_m}^{(h)})
  \le (K+L)\sqrt{n_m}\qquad(1\le h\le J_m),
\]
and every fixed menu element is eventually covered.

For an explicit half-budget schedule, set
\[
  \Lambda_m^{(J)}:=\sum_{1\le h\le J_m}N_{m,h}.
\]
Assume that, eventually for every \(1\le h\le J_m\),
\[
  \mu_{m,h}\le \frac{Qn_m}{2\ell_m},\qquad
  \sigma_{m,h}^2\le S\frac{n_m}{\ell_m},\qquad
  b_{m,h}\le b,
\]
that \(K_{m,h}\le K\), \(\sum_m\sum_{h\le J_m}\mathbb P(E_{m,h})<\infty\),
\(\Lambda_m^{(J)}\le(eD_m)^\Gamma\), and \(n_m\ge a\ell_m^2\).  Put
\[
  u_{m,h}=\frac{Qn_m}{2\ell_m},
  \qquad
  \delta=\frac{Q^2a}{8S+\frac43bQ}.
\]
Along all integer dimensions \(D=2,3,\ldots\), if \(\delta>\Gamma+1\), then
the active-prefix and fixed-menu almost-sure conclusions hold for every
\(A>\Gamma+1\) and every \(L\) satisfying
\[
  \frac{L^2}{2(Q+L/(3\sqrt a))}\ge A+1 .
\]
Along any sparse schedule \(D_m\ge\exp(c m^\kappa)\), if
\(\delta>\Gamma\), the same conclusions hold for every \(A>\Gamma\) and every
\(L\) satisfying the same display.
\end{corollary}

\begin{proof}
For each fixed \(m\) and \(h\), the corresponding one-row corollary has the
common tail bound displayed in the proof of
Corollary~\ref{cor:core-plus-localized-random-amplitude-menu-borel}.  Summing
that bound over \(1\le h\le J_m\) and over \(m\) gives a finite total failure
probability by hypothesis, so the first Borel--Cantelli lemma proves the
active-prefix conclusion.  Since \(J_m\to\infty\), every fixed \(h\) belongs to
the active prefix eventually.

For the profile form, take \(s_m=L\sqrt{n_m}\).  On the active prefix,
\[
  \frac{s_m^2}{2(\mu_{m,h}+u_{m,h}+s_m/3)}
  \ge
  \frac{L^2n_m}{2(Qn_m/\ell_m+L\sqrt{n_m}/3)}
  \ge
  \frac{L^2}{2(Q+LB/3)}\,\ell_m
  \ge (A+1)\ell_m,
\]
so the matrix-Freedman term is at most \(2(eD_m)^{-A}\) for each active
\(h\).  The displayed \(J_m(eD_m)^{-A}\) summability controls these matrix
tails, while the assumed double series controls the core and random-budget
exceptions.

For the half-budget schedule, the same computation as in
Corollary~\ref{cor:core-plus-localized-random-amplitude-menu-borel} gives
\[
  N_{m,h}\exp\left\{
    -\frac{u_{m,h}^2}{2(\sigma_{m,h}^2+b_{m,h}u_{m,h}/3)}
  \right\}
  \le N_{m,h}(eD_m)^{-\delta}.
\]
Summing over \(h\le J_m\) bounds the active-prefix budget exception by
\(\Lambda_m^{(J)}(eD_m)^{-\delta}\le(eD_m)^{\Gamma-\delta}\).  Also
\(J_m\le\Lambda_m^{(J)}\), since \(N_{m,h}\ge1\), so the matrix tails are
summable under the stated conditions on \(A\).  The all-integer and sparse
criteria are the usual \(p\)-series and stretched-exponential summability
checks.
\end{proof}

\begin{corollary}[Homogeneous random-amplitude localized schedules]
\label{cor:core-plus-homogeneous-localized-random-amplitude-schedules}
In the setting of
Corollary~\ref{cor:core-plus-localized-random-amplitude-menu-borel}, suppose
that each row belongs to one of the homogeneous localized models below.  Let
\[
  H_m=\sum_{i=1}^{n_m}\mathbb E\rho_{m,i}^2
\]
be the amplitude energy, let \(\beta_m\in[0,1]\) be a deterministic
homogeneous intensity bound, and let \(N_m\ge1\) be the corresponding budget
union count:
\begin{enumerate}[label=\textup{(\alph*)}]
\item coordinate or random-coordinate-profile supports, where
  \(p_{m,ij}\le\beta_m\) or \(Q_{m,ij}\le\beta_m\) and \(N_m=D_m\);
\item block-profile supports over \(N_m\) fixed orthogonal blocks, where
  \(Q_{m,i\ell}\le\beta_m\);
\item fixed-rank Haar, random-rank Haar, isotropic, or spectrally dominated
  projection laws, where the rank fraction or intensity satisfies
  \(\alpha_{m,i}\le\beta_m\) and \(N_m=1\);
\item mixed coordinate/Haar menus, where
  \(M_{m,i}Q_{m,ij}+(1-M_{m,i})R_{m,i}/D_m\le\beta_m\) and \(N_m=D_m\).
\end{enumerate}
Assume the same signed-core event \(E_m\) and core constant \(K_m\) as in
Corollary~\ref{cor:core-plus-localized-random-amplitude-menu-borel}.  Then, for
every \(s_m>0\),
\[
\begin{aligned}
  \mathbb P\left\{
  \disc(X_{m,1},\ldots,X_{m,n_m})>K_m\sqrt{n_m}+s_m
  \right\}
  &\le
  \mathbb P(E_m)+N_m\exp\{-3H_m/8\}\\
  &\quad+
  2D_m\exp\left\{-\frac{s_m^2}{2(2\beta_mH_m+s_m/3)}\right\}.
\end{aligned}
\]
Consequently, if
\[
  \sum_m\left[
  \mathbb P(E_m)+N_m\exp\{-3H_m/8\}
  +2D_m\exp\left\{-\frac{s_m^2}{2(2\beta_mH_m+s_m/3)}\right\}
  \right]<\infty,
\]
then, with probability one, for all sufficiently large \(m\),
\[
  \disc(X_{m,1},\ldots,X_{m,n_m})
  \le K_m\sqrt{n_m}+s_m .
\]

For a concrete \(n/\log(eD)\) schedule, put \(\ell_m=\log(eD_m)\).  Suppose
that \(K_m\le K\), that
\[
  2\beta_mH_m\le Q\frac{n_m}{\ell_m},
  \qquad
  H_m\ge R\ell_m,
  \qquad
  N_m\le(eD_m)^\Gamma
\]
eventually, and that \(\sum_m\mathbb P(E_m)<\infty\).  If
\(\ell_m\le B\sqrt{n_m}\) eventually and
\[
  \sum_m(eD_m)^{\Gamma-3R/8}<\infty,
\]
then, for every \(A>0\) with \(\sum_m(eD_m)^{-A}<\infty\) and every \(L>0\)
satisfying
\[
  \frac{L^2}{2(Q+LB/3)}\ge A+1,
\]
one has, almost surely for all sufficiently large \(m\),
\[
  \disc(X_{m,1},\ldots,X_{m,n_m})\le (K+L)\sqrt{n_m}.
\]
In particular, along all integer dimensions \(D=2,3,\ldots\), it is enough that
\[
  n_D\ge a\log^2(eD),\qquad
  2\beta_DH_D\le Q\frac{n_D}{\log(eD)},\qquad
  H_D\ge R\log(eD),
\]
\(N_D\le(eD)^\Gamma\), and \(R>8(\Gamma+1)/3\).  Then the same conclusion
holds for every \(A>1\) and every \(L\) with
\[
  \frac{L^2}{2(Q+L/(3\sqrt a))}\ge A+1.
\]
Along any sparse schedule \(D_m\ge\exp(c m^\kappa)\), the same conclusion holds
under the same \(n_m\), budget, and energy hypotheses whenever
\(R>8\Gamma/3\), for every \(A>0\) and the same displayed condition on \(L\).
\end{corollary}

\begin{proof}
In each listed homogeneous case, the corresponding one-row corollary gives
\[
  \mu_m\le\beta_mH_m,\qquad
  \sigma_m^2\le\beta_m^2H_m,\qquad
  b_m\le\beta_m,
\]
with budget-union count \(N_m\).  Taking \(u_m=\beta_mH_m\) in
Corollary~\ref{cor:core-plus-localized-random-amplitude-menu-borel}, or
directly in the one-row tail, gives
when \(\beta_mH_m>0\),
\[
  N_m
  \exp\left\{
  -\frac{\beta_m^2H_m^2}{2(\beta_m^2H_m+\beta_m^2H_m/3)}
  \right\}
  \le N_m e^{-3H_m/8},
\]
while the zero-intensity case has the Bernstein term omitted.
Also \(\mu_m+u_m\le2\beta_mH_m\), which proves the first displayed tail and the
Borel--Cantelli consequence.

For the schedule form, take \(s_m=L\sqrt{n_m}\).  The budget hypothesis gives
\[
  \frac{s_m^2}{2(2\beta_mH_m+s_m/3)}
  \ge
  \frac{L^2n_m}{2(Qn_m/\ell_m+L\sqrt{n_m}/3)}
  \ge
  \frac{L^2}{2(Q+LB/3)}\,\ell_m
  \ge (A+1)\ell_m .
\]
Thus the matrix tail is at most \(2(eD_m)^{-A}\).  The scalar-budget exception
is bounded by \(N_m e^{-3H_m/8}\le(eD_m)^{\Gamma-3R/8}\), which is summable by
assumption.  The all-integer and sparse criteria are the corresponding
\(p\)-series and stretched-exponential summability tests, with
\(\ell_D/\sqrt{n_D}\le1/\sqrt a\).
\end{proof}

\begin{corollary}[Common-sign countable random-amplitude localized menus]
\label{cor:core-plus-common-sign-countable-localized-random-amplitude-menu-borel}
For each \(m\ge1\), let \(D_m\ge2\), put \(\ell_m=\log(eD_m)\), and let
\[
  J_1\le J_2\le\cdots,\qquad J_m\to\infty,
\]
be a deterministic exhaustion of the countable menu
\(\mathbb N=\{1,2,\ldots\}\).  Let \(\mathcal G_m\) be a row environment.  For
each \(m\) and \(h\ge1\), let
\[
  X_{m,i}^{(h)}=B_{m,i}^{(h)}+Y_{m,i}^{(h)}\in\M_{D_m}(\C),
  \qquad 1\le i\le n_m,
\]
be a signed-core localized matrix-noise row with independent random profile or
amplitude data as in one of the one-row corollaries listed in
Corollary~\ref{cor:core-plus-localized-random-amplitude-menu-borel}.  The menu
elements may be dependent on one another, but for each fixed \(h\) the
corresponding one-row hypotheses hold relative to a filtration containing
\(\mathcal G_m\) at time zero.  Suppose that there is an event
\(E_m^{\rm com}\in\mathcal G_m\) such that, outside \(E_m^{\rm com}\), one
\(\mathcal G_m\)-measurable sign vector
\(\eps_m^\star\in\{-1,1\}^{n_m}\) satisfies
\[
  \left\|\sum_{i=1}^{n_m}\eps_{m,i}^\star B_{m,i}^{(h)}\right\|_{\op}
  \le K_{m,h}\sqrt{n_m}
  \qquad(1\le h\le J_m).
\]
Let \(N_{m,h}\ge1\), \(\mu_{m,h}\ge0\), \(\sigma_{m,h}^2\ge0\), and
\(b_{m,h}\in[0,1]\) be the budget-union count, mean budget, Bernstein variance
proxy, and jump parameter supplied by the corresponding one-row corollary for
the \(h\)-th noise model.  Let \(u_{m,h}\ge0\) and \(s_m>0\) be deterministic.
If
\[
  \sum_m
  \left[
  \mathbb P(E_m^{\rm com})
  +\sum_{1\le h\le J_m}
  \left(
  N_{m,h}\exp\left\{
    -\frac{u_{m,h}^2}{2(\sigma_{m,h}^2+b_{m,h}u_{m,h}/3)}
  \right\}
  +2D_m\exp\left\{
    -\frac{s_m^2}{2(\mu_{m,h}+u_{m,h}+s_m/3)}
  \right\}
  \right)
  \right]<\infty,
\]
with the usual convention that the Bernstein term is omitted when its
denominator is zero, then, with probability one, for all sufficiently large
\(m\), the same sign vector \(\eps_m^\star\) satisfies
\[
  \left\|\sum_{i=1}^{n_m}\eps_{m,i}^\star X_{m,i}^{(h)}\right\|_{\op}
  \le K_{m,h}\sqrt{n_m}+s_m
  \qquad(1\le h\le J_m).
\]
Consequently, every fixed \(h\) is eventually covered by the row's common
active-prefix signing.

In particular, suppose that \(K_{m,h}\le K\) for \(1\le h\le J_m\), that
\[
  \mu_{m,h}+u_{m,h}\le Q\frac{n_m}{\ell_m}
  \qquad(1\le h\le J_m)
\]
eventually, and that, for some \(A,B>0\),
\[
  \sum_m
  \left[
  \mathbb P(E_m^{\rm com})
  +\sum_{1\le h\le J_m}
  N_{m,h}\exp\left\{
    -\frac{u_{m,h}^2}{2(\sigma_{m,h}^2+b_{m,h}u_{m,h}/3)}
  \right\}
  \right]
  +
  \sum_mJ_m(eD_m)^{-A}<\infty,
\]
while \(\ell_m\le B\sqrt{n_m}\) eventually.  If \(L>0\) satisfies
\[
  \frac{L^2}{2(Q+LB/3)}\ge A+1,
\]
then, almost surely, for all sufficiently large \(m\), one common sign vector
\(\eps_m^\star\) obeys
\[
  \left\|\sum_{i=1}^{n_m}\eps_{m,i}^\star X_{m,i}^{(h)}\right\|_{\op}
  \le (K+L)\sqrt{n_m}\qquad(1\le h\le J_m),
\]
and every fixed menu element is eventually covered by these common row signs.

For an explicit half-budget schedule, set
\[
  \Lambda_m^{(J)}:=\sum_{1\le h\le J_m}N_{m,h}.
\]
Assume that, eventually for every \(1\le h\le J_m\),
\[
  \mu_{m,h}\le \frac{Qn_m}{2\ell_m},\qquad
  \sigma_{m,h}^2\le S\frac{n_m}{\ell_m},\qquad
  b_{m,h}\le b,
\]
that \(K_{m,h}\le K\), \(\sum_m\mathbb P(E_m^{\rm com})<\infty\),
\(\Lambda_m^{(J)}\le(eD_m)^\Gamma\), and \(n_m\ge a\ell_m^2\).  Put
\[
  u_{m,h}=\frac{Qn_m}{2\ell_m},
  \qquad
  \delta=\frac{Q^2a}{8S+\frac43bQ}.
\]
Along all integer dimensions \(D=2,3,\ldots\), if \(\delta>\Gamma+1\), then
the active-prefix common-sign and fixed-menu conclusions hold for every
\(A>\Gamma+1\) and every \(L\) satisfying
\[
  \frac{L^2}{2(Q+L/(3\sqrt a))}\ge A+1 .
\]
Along any sparse schedule \(D_m\ge\exp(c m^\kappa)\), if
\(\delta>\Gamma\), the same conclusions hold for every \(A>\Gamma\) and every
\(L\) satisfying the same display.
\end{corollary}

\begin{proof}
Work on the common core event \((E_m^{\rm com})^c\), and fix the common signs
\(\eps_m^\star\).  For each \(h\le J_m\), the proof of the corresponding
one-row random-amplitude localized corollary applies with these prescribed
signs: after the random budget is at most \(\mu_{m,h}+u_{m,h}\), matrix
Freedman bounds
\[
  \left\|\sum_i\eps_{m,i}^\star Y_{m,i}^{(h)}\right\|_{\op}
\]
by \(s_m\) except with probability at most the displayed \(2D_m\)-term, while
scalar Bernstein gives the displayed budget-exception term.  Hence
\[
\begin{aligned}
  &\mathbb P\left\{
  \exists h\le J_m:
  \left\|\sum_i\eps_{m,i}^\star X_{m,i}^{(h)}\right\|_{\op}
  >K_{m,h}\sqrt{n_m}+s_m
  \right\}\\
  &\quad\le
  \mathbb P(E_m^{\rm com})
  +\sum_{h\le J_m}
  \left[
  N_{m,h}\exp\left\{
    -\frac{u_{m,h}^2}{2(\sigma_{m,h}^2+b_{m,h}u_{m,h}/3)}
  \right\}
  +2D_m\exp\left\{
    -\frac{s_m^2}{2(\mu_{m,h}+u_{m,h}+s_m/3)}
  \right\}
  \right].
\end{aligned}
\]
The first assertion follows from the first Borel--Cantelli lemma.  Since
\(J_m\to\infty\), every fixed \(h\) is in the active prefix for all large \(m\).

The profile and half-budget schedule proofs are the same summability estimates
as in
Corollary~\ref{cor:core-plus-countable-localized-random-amplitude-menu-borel}.
Indeed, with \(s_m=L\sqrt{n_m}\), the active-prefix matrix tail is bounded by
\(2(eD_m)^{-A}\) under the displayed condition on \(L\), so the
\(J_m(eD_m)^{-A}\) series controls all matrix tails.  In the half-budget
setting,
\[
  N_{m,h}\exp\left\{
    -\frac{u_{m,h}^2}{2(\sigma_{m,h}^2+b_{m,h}u_{m,h}/3)}
  \right\}
  \le N_{m,h}(eD_m)^{-\delta},
\]
and summing over \(h\le J_m\) gives
\(\Lambda_m^{(J)}(eD_m)^{-\delta}\le(eD_m)^{\Gamma-\delta}\).  Since
\(J_m\le\Lambda_m^{(J)}\), the stated all-integer and sparse thresholds are the
usual \(p\)-series and stretched-exponential summability tests.
\end{proof}

\begin{corollary}[Common-sign homogeneous random-amplitude localized schedules]
\label{cor:core-plus-common-sign-homogeneous-localized-random-amplitude-schedules}
For each \(m\ge1\), let \(D_m\ge2\), put \(\ell_m=\log(eD_m)\), and let
\[
  J_1\le J_2\le\cdots,\qquad J_m\to\infty,
\]
be a deterministic exhaustion of the countable menu
\(\mathbb N=\{1,2,\ldots\}\).  Let \(\mathcal G_m\) be a row environment.  For
each \(m\) and \(h\ge1\), let
\[
  X_{m,i}^{(h)}=B_{m,i}^{(h)}+Y_{m,i}^{(h)}\in\M_{D_m}(\C),
  \qquad 1\le i\le n_m,
\]
be one of the homogeneous localized random-amplitude models listed in
Corollary~\ref{cor:core-plus-homogeneous-localized-random-amplitude-schedules},
with energy
\[
  H_{m,h}=\sum_{i=1}^{n_m}\mathbb E(\rho_{m,i}^{(h)})^2,
\]
homogeneous intensity \(\beta_{m,h}\in[0,1]\), and budget-union count
\(N_{m,h}\ge1\).  Suppose that there is an event
\(E_m^{\rm com}\in\mathcal G_m\) such that, outside \(E_m^{\rm com}\), one
\(\mathcal G_m\)-measurable sign vector
\(\eps_m^\star\in\{-1,1\}^{n_m}\) satisfies
\[
  \left\|\sum_{i=1}^{n_m}\eps_{m,i}^\star B_{m,i}^{(h)}\right\|_{\op}
  \le K_{m,h}\sqrt{n_m}
  \qquad(1\le h\le J_m).
\]
Then, for every deterministic \(s_m>0\), if
\[
  \sum_m
  \left[
  \mathbb P(E_m^{\rm com})
  +\sum_{1\le h\le J_m}
  \left(
  N_{m,h}e^{-3H_{m,h}/8}
  +2D_m\exp\left\{
    -\frac{s_m^2}{2(2\beta_{m,h}H_{m,h}+s_m/3)}
  \right\}
  \right)
  \right]<\infty,
\]
then, with probability one, for all sufficiently large \(m\), the same sign
vector \(\eps_m^\star\) satisfies
\[
  \left\|\sum_{i=1}^{n_m}\eps_{m,i}^\star X_{m,i}^{(h)}\right\|_{\op}
  \le K_{m,h}\sqrt{n_m}+s_m
  \qquad(1\le h\le J_m).
\]
Consequently, every fixed homogeneous localized random-amplitude model is
eventually covered by the row's common active-prefix signing.

Set
\[
  \Lambda_m^{(J)}:=\sum_{1\le h\le J_m}N_{m,h}.
\]
Suppose that \(K_{m,h}\le K\) for \(1\le h\le J_m\), and that there are
constants \(Q,R,\Gamma\ge0\) such that, eventually on the active prefix,
\[
  2\beta_{m,h}H_{m,h}\le Q\frac{n_m}{\ell_m},
  \qquad
  H_{m,h}\ge R\ell_m,
  \qquad
  \Lambda_m^{(J)}\le(eD_m)^\Gamma,
\]
while \(\sum_m\mathbb P(E_m^{\rm com})<\infty\).  If
\(\ell_m\le B\sqrt{n_m}\) eventually and, for some \(A>0\),
\[
  \sum_m(eD_m)^{\Gamma-3R/8}<\infty,
  \qquad
  \sum_mJ_m(eD_m)^{-A}<\infty,
\]
then for every \(L>0\) satisfying
\[
  \frac{L^2}{2(Q+LB/3)}\ge A+1,
\]
one has, almost surely for all sufficiently large \(m\), one common sign vector
\(\eps_m^\star\) obeys
\[
  \left\|\sum_{i=1}^{n_m}\eps_{m,i}^\star X_{m,i}^{(h)}\right\|_{\op}
  \le (K+L)\sqrt{n_m}\qquad(1\le h\le J_m),
\]
and each fixed menu element is eventually covered by these common row signs.

Along all integer dimensions \(D=2,3,\ldots\), it is enough that
\[
  n_D\ge a\log^2(eD),\qquad
  2\beta_{D,h}H_{D,h}\le Q\frac{n_D}{\log(eD)},\qquad
  H_{D,h}\ge R\log(eD)
\]
for \(h\le J_D\), that \(\Lambda_D^{(J)}\le(eD)^\Gamma\), and that
\[
  R>\frac{8(\Gamma+1)}{3}.
\]
Then the active-prefix common-sign and fixed-model conclusions hold for every
\(A>\Gamma+1\) and every \(L\) satisfying
\[
  \frac{L^2}{2(Q+L/(3\sqrt a))}\ge A+1 .
\]
Along any sparse schedule \(D_m\ge\exp(c m^\kappa)\), the same conclusions hold
under the same active-prefix budget and energy hypotheses whenever
\[
  R>\frac{8\Gamma}{3},
\]
for every \(A>\Gamma\) and the same displayed condition on \(L\).
\end{corollary}

\begin{proof}
For each active-prefix model, the homogeneous one-row bounds give
\[
  \mu_{m,h}\le\beta_{m,h}H_{m,h},\qquad
  \sigma_{m,h}^2\le\beta_{m,h}^2H_{m,h},\qquad
  b_{m,h}\le\beta_{m,h}.
\]
Apply
Corollary~\ref{cor:core-plus-common-sign-countable-localized-random-amplitude-menu-borel}
with \(u_{m,h}=\beta_{m,h}H_{m,h}\).  When \(\beta_{m,h}H_{m,h}>0\),
\[
  N_{m,h}\exp\left\{
    -\frac{u_{m,h}^2}{2(\sigma_{m,h}^2+b_{m,h}u_{m,h}/3)}
  \right\}
  \le
  N_{m,h}e^{-3H_{m,h}/8},
\]
and in the zero-intensity case the Bernstein term is omitted.  Also
\(\mu_{m,h}+u_{m,h}\le2\beta_{m,h}H_{m,h}\), which proves the first assertion.

For the profile form, take \(s_m=L\sqrt{n_m}\).  On the active prefix,
\[
  \frac{s_m^2}{2(2\beta_{m,h}H_{m,h}+s_m/3)}
  \ge
  \frac{L^2n_m}{2(Qn_m/\ell_m+L\sqrt{n_m}/3)}
  \ge
  \frac{L^2}{2(Q+LB/3)}\,\ell_m
  \ge (A+1)\ell_m .
\]
Thus each matrix tail is at most \(2(eD_m)^{-A}\), and the
\(J_m(eD_m)^{-A}\) series controls all active-prefix matrix tails.  The
homogeneous scalar-budget exceptions satisfy
\[
  \sum_{h\le J_m}N_{m,h}e^{-3H_{m,h}/8}
  \le
  \Lambda_m^{(J)}e^{-3R\ell_m/8}
  \le
  (eD_m)^{\Gamma-3R/8},
\]
which is summable by assumption.  The common core exceptions are summable by
hypothesis.

The all-integer and sparse schedule checks are the same \(p\)-series and
stretched-exponential estimates: in integer dimensions
\((eD)^{\Gamma-3R/8}\) is summable when \(3R/8-\Gamma>1\), while
\(J_D(eD)^{-A}\le(eD)^{\Gamma-A}\) is summable when \(A>\Gamma+1\); in sparse
dimensions the corresponding stretched-exponential series only require
\(3R/8-\Gamma>0\) and \(A>\Gamma\).  Finally
\(\ell_D/\sqrt{n_D}\le1/\sqrt a\) under \(n_D\ge a\log^2(eD)\).
\end{proof}

\begin{corollary}[Finite-lift core libraries with random-amplitude localized noise]
\label{cor:finite-lift-core-random-amplitude-libraries}
In the triangular-array setting of
Corollary~\ref{cor:core-plus-common-sign-countable-localized-random-amplitude-menu-borel},
replace the common-core signing hypothesis by the following finite-lift menu
hypothesis.  There are events \(F_m^{\rm lift}\in\mathcal G_m\) with
\(\sum_m\mathbb P(F_m^{\rm lift})<\infty\) such that, outside
\(F_m^{\rm lift}\), the active-prefix core tuples
\[
  B_{m,1}^{(h)},\ldots,B_{m,n_m}^{(h)}\in\M_{D_m}(\C),
  \qquad 1\le h\le J_m,
\]
are contractions and admit finite block lifts
\[
  \mcB_{m,h}=\bigoplus_{\alpha=1}^{q_{m,h}}\M_{r_{m,h,\alpha}}(\C)
\]
with lift constants \(K_{m,h}^{\rm lift},\Lambda_{m,h}^{\rm lift}\).  Put
\[
\begin{aligned}
M_m^{(J)}
&=\max_{1\le h\le J_m}
  \Lambda_{m,h}^{\rm lift}K_{m,h}^{\rm lift},\\
h_m^{(J)}
&=\max_{\substack{1\le h\le J_m\\1\le\alpha\le q_{m,h}}}
  r_{m,h,\alpha},
\qquad
R_m^{(J)}
=\sum_{1\le h\le J_m}\sum_{\alpha=1}^{q_{m,h}}r_{m,h,\alpha},
\end{aligned}
\]
and \(E_m^{(J)}=h_m^{(J)}R_m^{(J)}\).  If
\[
  M_m^{(J)}\le M_0,\qquad E_m^{(J)}\le C_E n_m
\]
eventually outside \(F_m^{\rm lift}\), then every conclusion of
Corollary~\ref{cor:core-plus-common-sign-countable-localized-random-amplitude-menu-borel}
holds with
\[
  E_m^{\rm com}=F_m^{\rm lift},\qquad
  K_{m,h}=K_{\rm fl}:=C M_0\sqrt{\log(e+C_E)} .
\]
In particular, if the lift condition is deterministic then there is no core
exception.  In the homogeneous localized subcases, every conclusion and explicit
all-integer or sparse schedule of
Corollary~\ref{cor:core-plus-common-sign-homogeneous-localized-random-amplitude-schedules}
holds with the same substitution \(K=K_{\rm fl}\) and
\(E_m^{\rm com}=F_m^{\rm lift}\).
\end{corollary}

\begin{proof}
On \((F_m^{\rm lift})^c\), apply
Corollary~\ref{cor:finite-lift-common-menu} to the finite menu
\(\{B_{m}^{(h)}:1\le h\le J_m\}\).  The displayed bounds on \(M_m^{(J)}\) and
\(E_m^{(J)}\) give one sign vector \(\eps_m^\star\) such that
\[
  \left\|\sum_{i=1}^{n_m}\eps_{m,i}^\star B_{m,i}^{(h)}\right\|_{\op}
  \le K_{\rm fl}\sqrt{n_m}
  \qquad(1\le h\le J_m).
\]
Since the set of sign vectors is finite, choosing the lexicographically first
successful sign vector makes \(\eps_m^\star\) \(\mathcal G_m\)-measurable.  This
is exactly the common-core hypothesis in the two common-sign random-amplitude
corollaries, with \(E_m^{\rm com}=F_m^{\rm lift}\).  Applying those corollaries
gives the stated conclusions.
\end{proof}

\begin{corollary}[Exhaustive homogeneous random-amplitude localized schedules]
\label{cor:core-plus-countable-homogeneous-localized-random-amplitude-schedules}
For each \(m\ge1\), let \(D_m\ge2\), put \(\ell_m=\log(eD_m)\), and let
\[
  J_1\le J_2\le\cdots,\qquad J_m\to\infty,
\]
be a deterministic exhaustion of the countable menu
\(\mathbb N=\{1,2,\ldots\}\).  For each \(m\) and \(h\ge1\), let
\[
  X_{m,i}^{(h)}=B_{m,i}^{(h)}+Y_{m,i}^{(h)}\in\M_{D_m}(\C),
  \qquad 1\le i\le n_m,
\]
be one of the homogeneous localized random-amplitude models listed in
Corollary~\ref{cor:core-plus-homogeneous-localized-random-amplitude-schedules},
with its own signed-core event \(E_{m,h}\), core constant \(K_{m,h}\), energy
\[
  H_{m,h}=\sum_{i=1}^{n_m}\mathbb E(\rho_{m,i}^{(h)})^2,
\]
homogeneous intensity \(\beta_{m,h}\in[0,1]\), and budget-union count
\(N_{m,h}\ge1\).  Set
\[
  \Lambda_m^{(J)}:=\sum_{1\le h\le J_m}N_{m,h}.
\]
Then, for every deterministic \(s_m>0\), if
\[
  \sum_m\sum_{1\le h\le J_m}
  \left[
  \mathbb P(E_{m,h})+N_{m,h}e^{-3H_{m,h}/8}
  +2D_m\exp\left\{
    -\frac{s_m^2}{2(2\beta_{m,h}H_{m,h}+s_m/3)}
  \right\}
  \right]<\infty,
\]
then, with probability one, for all sufficiently large \(m\), every active
prefix element \(1\le h\le J_m\) satisfies
\[
  \disc(X_{m,1}^{(h)},\ldots,X_{m,n_m}^{(h)})
  \le K_{m,h}\sqrt{n_m}+s_m .
\]
Consequently, every fixed homogeneous localized random-amplitude model is
covered for all sufficiently large \(m\).

In particular, suppose that \(K_{m,h}\le K\) for \(1\le h\le J_m\), and that
there are constants \(Q,R,\Gamma\ge0\) such that, eventually on the active
prefix,
\[
  2\beta_{m,h}H_{m,h}\le Q\frac{n_m}{\ell_m},
  \qquad
  H_{m,h}\ge R\ell_m,
  \qquad
  \Lambda_m^{(J)}\le(eD_m)^\Gamma,
\]
while
\[
  \sum_m\sum_{1\le h\le J_m}\mathbb P(E_{m,h})<\infty .
\]
If \(\ell_m\le B\sqrt{n_m}\) eventually and, for some \(A>0\),
\[
  \sum_m(eD_m)^{\Gamma-3R/8}<\infty,
  \qquad
  \sum_mJ_m(eD_m)^{-A}<\infty,
\]
then for every \(L>0\) satisfying
\[
  \frac{L^2}{2(Q+LB/3)}\ge A+1,
\]
one has, almost surely for all sufficiently large \(m\),
\[
  \disc(X_{m,1}^{(h)},\ldots,X_{m,n_m}^{(h)})
  \le (K+L)\sqrt{n_m}\qquad(1\le h\le J_m),
\]
and each fixed menu element is eventually covered.

Along all integer dimensions \(D=2,3,\ldots\), it is enough that
\[
  n_D\ge a\log^2(eD),\qquad
  2\beta_{D,h}H_{D,h}\le Q\frac{n_D}{\log(eD)},\qquad
  H_{D,h}\ge R\log(eD)
\]
for \(h\le J_D\), that \(\Lambda_D^{(J)}\le(eD)^\Gamma\), and that
\[
  R>\frac{8(\Gamma+1)}{3}.
\]
Then the active-prefix and fixed-model conclusions hold for every
\(A>\Gamma+1\) and every \(L\) satisfying
\[
  \frac{L^2}{2(Q+L/(3\sqrt a))}\ge A+1 .
\]
Along any sparse schedule \(D_m\ge\exp(c m^\kappa)\), the same conclusions hold
under the same active-prefix budget and energy hypotheses whenever
\[
  R>\frac{8\Gamma}{3},
\]
for every \(A>\Gamma\) and the same displayed condition on \(L\).
\end{corollary}

\begin{proof}
For each fixed \(m\) and \(h\),
Corollary~\ref{cor:core-plus-homogeneous-localized-random-amplitude-schedules}
gives the one-row failure bound with parameters
\(E_{m,h}\), \(K_{m,h}\), \(H_{m,h}\), \(\beta_{m,h}\), and \(N_{m,h}\).
Summing this bound over \(1\le h\le J_m\) and then over \(m\) gives a finite
total failure probability, so the first Borel--Cantelli lemma proves the
active-prefix statement.  Since \(J_m\to\infty\), each fixed \(h\) lies in the
active prefix for all large \(m\).

For the profile form, take \(s_m=L\sqrt{n_m}\).  On the active prefix,
\[
  \frac{s_m^2}{2(2\beta_{m,h}H_{m,h}+s_m/3)}
  \ge
  \frac{L^2n_m}{2(Qn_m/\ell_m+L\sqrt{n_m}/3)}
  \ge
  \frac{L^2}{2(Q+LB/3)}\,\ell_m
  \ge (A+1)\ell_m .
\]
Thus each matrix tail is at most \(2(eD_m)^{-A}\), and the active-prefix
matrix tails are summable by the displayed \(J_m(eD_m)^{-A}\) hypothesis.  The
homogeneous scalar-budget exceptions satisfy
\[
  \sum_{h\le J_m}N_{m,h}e^{-3H_{m,h}/8}
  \le
  \Lambda_m^{(J)}e^{-3R\ell_m/8}
  \le
  (eD_m)^{\Gamma-3R/8},
\]
which is summable by assumption.  The core exceptions are summable by
hypothesis.

The all-integer case follows because
\((eD)^{\Gamma-3R/8}\) is summable when
\(3R/8-\Gamma>1\), while \(J_D(eD)^{-A}\le(eD)^{\Gamma-A}\) is summable when
\(A>\Gamma+1\).  Also \(\ell_D/\sqrt{n_D}\le1/\sqrt a\).  The sparse case is
the same calculation with stretched-exponential summability, requiring only
\(3R/8-\Gamma>0\) and \(A>\Gamma\).
\end{proof}

\begin{corollary}[Core signings with predictable martingale scalar noise]
\label{cor:core-plus-martingale-scalar-noise}
Let \(D\ge2\).  Let \(\mathcal G\) be a sigma-field, and let
\(B_i\in\M_D(\C)\), \(i=1,\ldots,n\), be \(\mathcal G\)-measurable random
self-adjoint contractions.  Let \((\mcF_i)_{i=0}^n\) be a filtration with
\(\mathcal G\subset\mcF_0\), and let \(H_i\in\M_D(\C)\) be self-adjoint
contractions such that \(H_i\) is \(\mcF_{i-1}\)-measurable.  Let \(\theta_i\)
be real \(\mcF_i\)-measurable random variables satisfying
\[
  \mathbb E(\theta_i\mid\mcF_{i-1})=0,\qquad |\theta_i|\le1.
\]
Set
\[
  X_i=B_i+\theta_iH_i,\qquad
  q_i=\mathbb E(\theta_i^2\mid\mcF_{i-1}),\qquad
  \nu=\left\|\sum_{i=1}^n q_iH_i^2\right\|_{\op}.
\]
Assume that, outside an event \(E\in\mathcal G\), there are
\(\mathcal G\)-measurable signs
\(\eps_i^\star\in\{-1,1\}\) such that
\[
  \left\|\sum_{i=1}^n\eps_i^\star B_i\right\|_{\op}\le K\sqrt n .
\]
Then, for every deterministic \(v\ge0\) and \(s>0\),
\[
\begin{aligned}
  \mathbb P\left\{
  \disc(X_1,\ldots,X_n)>K\sqrt n+s
  \right\}
  &\le
  \mathbb P(E)+\mathbb P\{\nu>v\} \\
  &\quad+
  2D\exp\left\{-\frac{s^2}{2(v+s/3)}\right\}.
\end{aligned}
\]
Since \(\|X_i\|_{\op}\le2\), the normalized contractions \(X_i/2\) satisfy the
same estimate with the discrepancy threshold divided by \(2\).  In particular,
if \(q_i\le p\) almost surely for every \(i\), then
\[
  \mathbb P\left\{
  \disc(X_1,\ldots,X_n)>K\sqrt n+s
  \right\}
  \le
  \mathbb P(E)
  +
  2D\exp\left\{-\frac{s^2}{2(pn+s/3)}\right\}.
\]
Thus the signed-core noise theorem does not require conditional independence of
the scalar amplitudes, nor does it require the noise directions to be fixed at
time zero: after the core signing is fixed, predictable matrix directions with
bounded martingale-difference scalar amplitudes are absorbed by the same core
signing.
\end{corollary}

\begin{proof}
Work on \(E^c\).  For the \(\mathcal G\)-measurable signing \(\eps^\star\), define
\[
  Z_i=\eps_i^\star\theta_iH_i.
\]
Since \(\mathcal G\subset\mcF_{i-1}\) and \(H_i\) is
\(\mcF_{i-1}\)-measurable, the matrices \(\eps_i^\star H_i\) are
\(\mcF_{i-1}\)-measurable.  Hence \(Z_i\) is a self-adjoint matrix martingale
difference and \(\|Z_i\|_{\op}\le1\).  Moreover,
\[
  \mathbb E(Z_i^2\mid\mcF_{i-1})
  =
  q_iH_i^2.
\]
On \(\{\nu\le v\}\), the predictable quadratic variation therefore has
operator norm at most \(v\).  Applying the matrix Freedman inequality
\cite{Tropp} to \(Z_i\) and to \(-Z_i\), and taking a union bound, gives
\[
  \mathbb P\left\{
  \left\|\sum_iZ_i\right\|_{\op}>s\ \text{and}\ \nu\le v
  \right\}
  \le
  2D\exp\left\{-\frac{s^2}{2(v+s/3)}\right\}.
\]
On the complementary good event,
\[
  \left\|\sum_i\eps_i^\star X_i\right\|_{\op}
  \le
  \left\|\sum_i\eps_i^\star B_i\right\|_{\op}
  +
  \left\|\sum_iZ_i\right\|_{\op}
  \le K\sqrt n+s.
\]
Adding the exceptional events \(E\) and \(\{\nu>v\}\) proves the first estimate.
If \(q_i\le p\), then \(0\le H_i^2\le I\), so
\(\nu\le p\|\sum_iH_i^2\|_{\op}\le pn\), which gives the final form.
\end{proof}

\begin{corollary}[Core signings with predictable localized martingale noise]
\label{cor:core-plus-predictable-localized-martingale-noise}
Let \(D\ge2\).  Let \(\mathcal G\) be a sigma-field, and let
\(B_i\in\M_D(\C)\), \(i=1,\ldots,n\), be \(\mathcal G\)-measurable random
self-adjoint contractions.  Let \((\mcF_i)_{i=0}^n\) be a filtration with
\(\mathcal G\subset\mcF_0\).  For each \(i\), let \(P_i\) be an
\(\mcF_{i-1}\)-measurable orthogonal projection and let \(H_i\) be an
\(\mcF_{i-1}\)-measurable self-adjoint contraction satisfying
\[
  H_i=P_iH_iP_i.
\]
Let \(\theta_i\) be real \(\mcF_i\)-measurable random variables satisfying
\[
  \mathbb E(\theta_i\mid\mcF_{i-1})=0,\qquad |\theta_i|\le1,
\]
and set
\[
  X_i=B_i+\theta_iH_i,\qquad
  q_i=\mathbb E(\theta_i^2\mid\mcF_{i-1}),\qquad
  \mu=\left\|\sum_{i=1}^n q_iP_i\right\|_{\op}.
\]
Assume that, outside an event \(E\in\mathcal G\), there are
\(\mathcal G\)-measurable signs \(\eps_i^\star\in\{-1,1\}\) with
\[
  \left\|\sum_{i=1}^n\eps_i^\star B_i\right\|_{\op}\le K\sqrt n .
\]
Then, for every deterministic \(v\ge0\) and \(s>0\),
\[
\begin{aligned}
  \mathbb P\left\{
  \disc(X_1,\ldots,X_n)>K\sqrt n+s
  \right\}
  &\le
  \mathbb P(E)+\mathbb P\{\mu>v\} \\
  &\quad+
  2D\exp\left\{-\frac{s^2}{2(v+s/3)}\right\}.
\end{aligned}
\]
In particular, if \(q_i\le p\) almost surely for every \(i\), then
\[
  \begin{aligned}
  \mathbb P\left\{
  \disc(X_1,\ldots,X_n)>K\sqrt n+s
  \right\}
  &\le
  \mathbb P(E)
  +
  \mathbb P\left\{\left\|\sum_iP_i\right\|_{\op}>M\right\} \\
  &\quad+
  2D\exp\left\{-\frac{s^2}{2(pM+s/3)}\right\}
  \end{aligned}
\]
for every deterministic \(M\ge0\).  Thus already-signed random cores admit an
explicit tail bound after online coordinate-supported, subspace-localized, or
frame-localized martingale scalar noise, and they remain Spencer-scale at
\(s=O(\sqrt n)\) when the predictable weighted overlap
\(\|\sum_i q_iP_i\|_{\op}\) is \(O(n/\log(eD))\).
\end{corollary}

\begin{proof}
Since \(H_i=P_iH_iP_i\) and \(\|H_i\|_{\op}\le1\), we have \(0\le H_i^2\preceq
P_i\).  Hence
\[
  \left\|\sum_i q_iH_i^2\right\|_{\op}
  \le
  \left\|\sum_i q_iP_i\right\|_{\op}
  =
  \mu.
\]
The first estimate follows from
Corollary~\ref{cor:core-plus-martingale-scalar-noise}.  If \(q_i\le p\), then
\(\mu\le p\|\sum_iP_i\|_{\op}\); applying the first estimate on the event
\(\|\sum_iP_i\|_{\op}\le M\) gives the displayed specialization.
\end{proof}

\begin{corollary}[Almost-sure signed-core predictable scalar-noise arrays]
\label{cor:core-plus-martingale-scalar-noise-borel}
For each \(m\ge1\), let \(D_m\ge2\), let \(\mathcal G_m\) be a sigma-field, let
\((\mcF_{m,i})_{i=0}^{n_m}\) be a filtration with
\(\mathcal G_m\subset\mcF_{m,0}\), and let
\[
  X_{m,i}=B_{m,i}+\theta_{m,i}H_{m,i}\in\M_{D_m}(\C)
  \qquad(1\le i\le n_m)
\]
be a signed-core predictable scalar-noise row as in
Corollary~\ref{cor:core-plus-martingale-scalar-noise}.  Thus \(B_{m,i}\) is
\(\mathcal G_m\)-measurable, \(H_{m,i}\) is
\(\mcF_{m,i-1}\)-measurable, \(\mathbb E(\theta_{m,i}\mid\mcF_{m,i-1})=0\), and
\(|\theta_{m,i}|\le1\).  Put
\[
  q_{m,i}:=\mathbb E(\theta_{m,i}^2\mid\mcF_{m,i-1}),
  \qquad
  \nu_m:=
  \left\|\sum_{i=1}^{n_m}q_{m,i}H_{m,i}^2\right\|_{\op}.
\]
Assume that, outside an event \(E_m\in\mathcal G_m\), there are
\(\mathcal G_m\)-measurable signs \(\eps_{m,i}^\star\) such that
\[
  \left\|\sum_{i=1}^{n_m}\eps_{m,i}^\star B_{m,i}\right\|_{\op}
  \le K_m\sqrt{n_m}.
\]
Let \(v_m\ge0\) and \(s_m>0\) be deterministic.  If
\[
  \sum_m
  \left[
  \mathbb P(E_m)+\mathbb P\{\nu_m>v_m\}
  +
  2D_m\exp\left\{-\frac{s_m^2}{2(v_m+s_m/3)}\right\}
  \right]<\infty,
\]
then, with probability one, for all sufficiently large \(m\),
\[
  \disc(X_{m,1},\ldots,X_{m,n_m})
  \le K_m\sqrt{n_m}+s_m .
\]

In particular, put \(\ell_m=\log(eD_m)\).  Suppose that \(K_m\le K\), that
\[
  \nu_m\le Q\frac{n_m}{\ell_m}
\]
outside events \(F_m\), that
\[
  \sum_m[\mathbb P(E_m)+\mathbb P(F_m)]<\infty,
\]
and that, for some \(A,B>0\),
\[
  \sum_m(eD_m)^{-A}<\infty,
  \qquad
  \ell_m\le B\sqrt{n_m}
\]
eventually.  Choose \(L>0\) with
\[
  \frac{L^2}{2(Q+LB/3)}\ge A+1 .
\]
Then, almost surely, for all sufficiently large \(m\),
\[
  \disc(X_{m,1},\ldots,X_{m,n_m})
  \le (K+L)\sqrt{n_m}.
\]
The variance-budget hypothesis is implied by the pointwise condition
\(q_{m,i}\le Q/\ell_m\) outside \(F_m\).

Consequently:
\begin{enumerate}[label=\textup{(\alph*)}]
\item Along all integer dimensions \(D=2,3,\ldots\), if \(K_D\le K\),
\(\nu_D\le Qn_D/\log(eD)\) outside summable events, and
\[
  n_D\ge a\,\log^2(eD)
\]
eventually for some \(a>0\), then for every \(A>1\) and every \(L\) satisfying
\[
  \frac{L^2}{2(Q+L/(3\sqrt a))}\ge A+1
\]
one has
\[
  \disc(X_{D,1},\ldots,X_{D,n_D})\le (K+L)\sqrt{n_D}
\]
eventually almost surely.

\item If \(D_m\ge\exp(c m^\kappa)\) eventually for some \(c,\kappa>0\), if
\(K_m\le K\), if \(\nu_m\le Qn_m/\log(eD_m)\) outside summable events, and if
\[
  n_m\ge a\,\log^2(eD_m)
\]
eventually for some \(a>0\), then the same conclusion holds along the sparse
schedule for every \(A>0\) and every \(L\) satisfying
\[
  \frac{L^2}{2(Q+L/(3\sqrt a))}\ge A+1 .
\]
\end{enumerate}
\end{corollary}

\begin{proof}
For the \(m\)-th row, write
\(\Delta_m=\disc(X_{m,1},\ldots,X_{m,n_m})\).
Corollary~\ref{cor:core-plus-martingale-scalar-noise} gives
\[
\begin{aligned}
  \mathbb P\{\Delta_m>K_m\sqrt{n_m}+s_m\}
  &\le
  \mathbb P(E_m)+\mathbb P\{\nu_m>v_m\}
  \\
  &\quad+
  2D_m\exp\left\{-\frac{s_m^2}{2(v_m+s_m/3)}\right\}.
\end{aligned}
\]
The displayed summability and the first Borel--Cantelli lemma prove the first
claim.

For the concrete profile, take \(s_m=L\sqrt{n_m}\) and
\(v_m=Qn_m/\ell_m\).  On \(F_m^c\),
\[
  \frac{s_m^2}{2(v_m+s_m/3)}
  =
  \frac{L^2n_m}{2(Qn_m/\ell_m+L\sqrt{n_m}/3)}
  \ge
  \frac{L^2}{2(Q+LB/3)}\,\ell_m
  \ge (A+1)\ell_m .
\]
Therefore
\[
  2D_m\exp\left\{-\frac{s_m^2}{2(v_m+s_m/3)}\right\}
  \le 2(eD_m)^{-A},
\]
which is summable by assumption.  If \(q_{m,i}\le Q/\ell_m\), then
\[
  \nu_m
  \le
  \frac{Q}{\ell_m}
  \left\|\sum_{i=1}^{n_m}H_{m,i}^2\right\|_{\op}
  \le Q\frac{n_m}{\ell_m}.
\]

For part (a), take \(B=1/\sqrt a\) and any \(A>1\), since
\(\sum_{D\ge2}(eD)^{-A}<\infty\).  For part (b), again take
\(B=1/\sqrt a\); the stretched-exponential lower bound on \(D_m\) makes
\(\sum_m(eD_m)^{-A}\) finite for every \(A>0\).
\end{proof}

\begin{corollary}[Weighted random low-overlap support ensembles]
\label{cor:random-low-overlap-support}
For every \(C_0\ge1\) there is \(C(C_0)\) with the following property.  Let
\(A_1,\ldots,A_n\) be random self-adjoint matrices in \(\M_D(\C)\).  Suppose
that, outside an event \(E\), there are orthogonal projections
\(P_1,\ldots,P_n\) and numbers \(\rho_1,\ldots,\rho_n\in[0,1]\) such that
\[
  A_i=P_iA_iP_i,\qquad \|A_i\|_{\op}\le\rho_i,\qquad
  \left\|\sum_{i=1}^n\rho_i^2P_i\right\|_{\op}\le
  \frac{C_0n}{\log(eD)}
\]
for every \(i\).  Then
\[
  \mathbb P\left\{
  \disc(A_1,\ldots,A_n)>C(C_0)\sqrt n
  \right\}
  \le \mathbb P(E).
\]
In particular, the Matrix Spencer bound holds with the same probability as the
displayed weighted-overlap event for coordinate-supported random matrices,
random localized Wigner/GOE/GUE blocks, and random projection or frame
ensembles.  The unweighted form is recovered by taking \(\rho_i=1\), in which
case the coordinate condition is maximum occupancy \(O(n/\log(eD))\).
\end{corollary}

\begin{proof}
On a realization outside \(E\), each \(A_i\) is self-adjoint and supported on
the range of \(P_i\), with \(\|A_i\|_{\op}\le\rho_i\).  Hence
\(0\le A_i^2\le \rho_i^2P_i\), and therefore
\[
  \left\|\sum_iA_i^2\right\|_{\op}
  \le
  \left\|\sum_i\rho_i^2P_i\right\|_{\op}
  \le \frac{C_0n}{\log(eD)}.
\]
Proposition~\ref{prop:square-function-khintchine} gives the claimed
discrepancy bound on that realization.  For coordinate supports, \(P_i\) is the
coordinate projection and \(\|\sum_i\rho_i^2P_i\|_{\op}\) is exactly the
maximum weighted coordinate occupancy; for subspace or frame models it is the
weighted empirical frame-operator norm.
\end{proof}

\begin{corollary}[Active menus with dependent self-adjoint overlap budgets]
\label{cor:active-menu-dependent-selfadjoint-overlap-budget}
Fix \(C_0\ge1\).  For each \(m\ge1\), let \(D_m\ge2\), \(n_m\ge1\), and
\(J_m\ge1\).  For each active index \(1\le h\le J_m\), let
\[
  A_{m,h,1},\ldots,A_{m,h,n_m}\in\M_{D_m}(\C)
\]
be random self-adjoint matrices.  Suppose that there are events
\(F_{m,h}^{\rm ov}\), random orthogonal projections \(P_{m,h,i}\), and random
amplitudes \(0\le\rho_{m,h,i}\le1\) such that, outside \(F_{m,h}^{\rm ov}\),
\[
  A_{m,h,i}=P_{m,h,i}A_{m,h,i}P_{m,h,i},
  \qquad
  \|A_{m,h,i}\|_{\op}\le\rho_{m,h,i}
  \qquad(1\le i\le n_m),
\]
and
\[
  \left\|\sum_{i=1}^{n_m}\rho_{m,h,i}^2P_{m,h,i}\right\|_{\op}
  \le
  \frac{C_0n_m}{\log(eD_m)}.
\]
If
\[
  \sum_m\sum_{h=1}^{J_m}\mathbb P(F_{m,h}^{\rm ov})<\infty ,
\]
then, almost surely, for all sufficiently large \(m\), every active stream has a
Matrix Spencer signing:
\[
  \disc(A_{m,h,1},\ldots,A_{m,h,n_m})
  \le C(C_0)\sqrt{n_m}
  \qquad(1\le h\le J_m).
\]
Sampled sign certificates are eventual almost surely if the per-stream candidate
lists also satisfy \(\sum_{m,h}2^{-M_{m,h}}<\infty\); a common list length
\(M_m\) for all active streams only requires
\(\sum_mJ_m2^{-M_m}<\infty\).  If \(J_m\to\infty\), active prefixes may exhaust
a countable library.

Thus arbitrary dependent self-adjoint support laws, random regular graphs,
matchings, without-replacement coordinate choices, determinantal or negatively
dependent supports, and adaptive projection/frontier models are covered once
their realized weighted overlap budgets fail only summably.
\end{corollary}

\begin{proof}
On \((F_{m,h}^{\rm ov})^c\), the pathwise hypotheses of
Corollary~\ref{cor:random-low-overlap-support} hold for stream \((m,h)\).
Summability of the overlap-budget failures and the first Borel--Cantelli lemma
give the simultaneous eventual signings.  The sampled-certificate assertion is
the same conditional random-sign argument as in
Corollary~\ref{cor:random-weighted-support-search}, with the displayed
per-stream or common-list failure series.  If \(J_m\to\infty\), each fixed
library stream belongs to the active prefix for all sufficiently large \(m\).
\end{proof}

\begin{corollary}[Tail-capped dependent self-adjoint overlap budgets]
\label{cor:active-menu-tail-capped-dependent-selfadjoint-overlap-budget}
Fix \(C_0\ge1\).  In the active-menu setting of
Corollary~\ref{cor:active-menu-dependent-selfadjoint-overlap-budget}, replace
the pointwise norm bounds by deterministic caps
\[
  0\le\rho_{m,h,i}\le1
  \qquad(1\le h\le J_m,\ 1\le i\le n_m).
\]
Assume that there are events \(F_{m,h}^{\rm ov}\), random orthogonal projections
\(P_{m,h,i}\), and cap-tail budgets \(\delta_{m,h,i}\ge0\) such that
\[
  \mathbb P\{\|A_{m,h,i}\|_{\op}>\rho_{m,h,i}\}\le\delta_{m,h,i},
\]
and, outside \(F_{m,h}^{\rm ov}\),
\[
  A_{m,h,i}=P_{m,h,i}A_{m,h,i}P_{m,h,i}
  \qquad(1\le i\le n_m)
\]
and
\[
  \left\|\sum_{i=1}^{n_m}\rho_{m,h,i}^2P_{m,h,i}\right\|_{\op}
  \le
  \frac{C_0n_m}{\log(eD_m)}.
\]
If
\[
  \sum_m\sum_{h=1}^{J_m}
  \left[
    \mathbb P(F_{m,h}^{\rm ov})
    +\sum_{i=1}^{n_m}\delta_{m,h,i}
  \right]<\infty ,
\]
then, almost surely, for all sufficiently large \(m\), every active stream has
\[
  \disc(A_{m,h,1},\ldots,A_{m,h,n_m})
  \le C(C_0)\sqrt{n_m}
  \qquad(1\le h\le J_m).
\]
Sampled certificates and countable-library exhaustions follow under the same
additional summability hypotheses as in
Corollary~\ref{cor:active-menu-dependent-selfadjoint-overlap-budget}.

For example, if
\[
  \mathbb P\{\|A_{m,h,i}\|_{\op}>s_{m,h,i}u\}\le A_*e^{-u^\alpha}
  \qquad(u\ge0)
\]
with fixed \(\alpha>0\), choose
\(\rho_{m,h,i}=s_{m,h,i}\ell_{m,h,i}^{1/\alpha}\le1\) and
\(\delta_{m,h,i}=A_*e^{-\ell_{m,h,i}}\).  If instead
\(\mathbb E\|A_{m,h,i}\|_{\op}^p\le s_{m,h,i}^p\), choose
\(\rho_{m,h,i}=s_{m,h,i}e^{\ell_{m,h,i}/p}\le1\) and
\(\delta_{m,h,i}=e^{-\ell_{m,h,i}}\).  These deterministic caps feed the same
realized self-adjoint overlap budget, so dependent heavy-tailed localized
models reduce to cap-tail summability plus a capped overlap certificate.
\end{corollary}

\begin{proof}
For each active stream set
\[
  \Omega_{m,h}^{\rm cap}
  :=\{\|A_{m,h,i}\|_{\op}\le\rho_{m,h,i}\text{ for all }i\le n_m\}.
\]
On \((F_{m,h}^{\rm ov})^c\cap\Omega_{m,h}^{\rm cap}\), the hypotheses of
Corollary~\ref{cor:active-menu-dependent-selfadjoint-overlap-budget} hold with
the deterministic caps \(\rho_{m,h,i}\).  The exceptional probability for
stream \((m,h)\) is at most
\[
  \mathbb P(F_{m,h}^{\rm ov})
  +\sum_{i=1}^{n_m}\delta_{m,h,i}.
\]
The displayed summability and Borel--Cantelli give the eventual active-menu
signings; sampled certificates add the same sampled-sign failure series as in
the dependent-overlap corollary.  The sub-Weibull and moment cap choices follow
from the displayed tail bound and Markov's inequality, respectively.
\end{proof}

\begin{corollary}[Common signs for dependent self-adjoint overlap menus]
\label{cor:common-sign-dependent-selfadjoint-overlap-budget}
Fix \(C_0\ge1\).  For each \(m\ge1\), let \(D_m\ge2\), \(n_m\ge1\),
\(J_m\ge1\), and put
\[
  D_m^{\rm com}:=J_mD_m .
\]
For \(1\le h\le J_m\), let
\[
  A_{m,h,1},\ldots,A_{m,h,n_m}\in\M_{D_m}(\C)
\]
be random self-adjoint matrices.  Suppose that there are events
\(F_{m,h}^{\rm ov}\), random orthogonal projections \(P_{m,h,i}\), and random
amplitudes \(0\le\rho_{m,h,i}\le1\) such that, outside \(F_{m,h}^{\rm ov}\),
\[
  A_{m,h,i}=P_{m,h,i}A_{m,h,i}P_{m,h,i},
  \qquad
  \|A_{m,h,i}\|_{\op}\le\rho_{m,h,i}
  \qquad(1\le i\le n_m),
\]
and
\[
  \left\|\sum_{i=1}^{n_m}\rho_{m,h,i}^2P_{m,h,i}\right\|_{\op}
  \le
  \frac{C_0n_m}{\log(eD_m^{\rm com})}.
\]
If
\[
  \sum_m\sum_{h=1}^{J_m}\mathbb P(F_{m,h}^{\rm ov})<\infty ,
\]
then, almost surely, for all sufficiently large \(m\), there is one sign vector
\(\eps^{(m)}\in\{\pm1\}^{n_m}\) such that
\[
  \max_{1\le h\le J_m}
  \left\|\sum_{i=1}^{n_m}\eps_i^{(m)}A_{m,h,i}\right\|_{\op}
  \le C(C_0)\sqrt{n_m}.
\]
If \(M_m\) independent random sign vectors are sampled for the whole active
menu and \(\sum_m2^{-M_m}<\infty\), then the sampled list contains such a common
certificate eventually almost surely.  If \(J_m\to\infty\), the common signs
cover each active prefix in a countable library eventually.
\end{corollary}

\begin{proof}
Outside the union of the events \(F_{m,h}^{\rm ov}\), define block-diagonal
self-adjoint matrices
\[
  \widetilde A_{m,i}:=
  \bigoplus_{h=1}^{J_m}A_{m,h,i}\in\M_{D_m^{\rm com}}(\C).
\]
Then
\[
  \left\|\sum_i\widetilde A_{m,i}^2\right\|_{\op}
  =
  \max_{h\le J_m}\left\|\sum_i A_{m,h,i}^2\right\|_{\op}
  \le
  \frac{C_0n_m}{\log(eD_m^{\rm com})}.
\]
Proposition~\ref{prop:square-function-khintchine} applied in
\(\M_{D_m^{\rm com}}(\C)\) gives one signing of the block-diagonal tuple, which
is exactly the displayed common signing.  The failure event is summable by
assumption, so Borel--Cantelli gives the eventual deterministic common signing.
The sampled version follows from
Corollary~\ref{cor:random-sign-search-square-profile} applied to the same
block-diagonal tuple, adding only the failure series \(\sum_m2^{-M_m}\).
\end{proof}

\begin{corollary}[Common signs for support-rank self-adjoint overlap menus]
\label{cor:common-sign-support-rank-selfadjoint-overlap-budget}
Fix \(C_0\ge1\).  For each \(m\ge1\), let \(D_m\ge2\), \(n_m\ge1\),
\(J_m\ge1\), and let \(R_m\ge1\).  For \(1\le h\le J_m\), let
\[
  A_{m,h,1},\ldots,A_{m,h,n_m}\in\M_{D_m}(\C)
\]
be random self-adjoint matrices.  Suppose that there are events \(F_m^{\rm sr}\),
random orthogonal projections \(P_{m,h,i}\), random support-envelope
projections \(Q_{m,h}\), and random amplitudes \(0\le\rho_{m,h,i}\le1\) such
that, outside \(F_m^{\rm sr}\),
\[
  A_{m,h,i}=P_{m,h,i}A_{m,h,i}P_{m,h,i},\qquad
  \|A_{m,h,i}\|_{\op}\le\rho_{m,h,i},\qquad
  P_{m,h,i}\le Q_{m,h},
\]
for all \(h\le J_m\) and \(i\le n_m\), and
\[
  \sum_{h=1}^{J_m}\max\{1,\operatorname{rank}Q_{m,h}\}\le R_m.
\]
Assume also that, outside \(F_m^{\rm sr}\),
\[
  \max_{1\le h\le J_m}
  \left\|\sum_{i=1}^{n_m}\rho_{m,h,i}^2P_{m,h,i}\right\|_{\op}
  \le
  \frac{C_0n_m}{\log(eR_m)}.
\]
If
\[
  \sum_m\mathbb P(F_m^{\rm sr})<\infty ,
\]
then, almost surely, for all sufficiently large \(m\), there is one sign vector
\(\eps^{(m)}\in\{\pm1\}^{n_m}\) such that
\[
  \max_{1\le h\le J_m}
  \left\|\sum_{i=1}^{n_m}\eps_i^{(m)}A_{m,h,i}\right\|_{\op}
  \le C(C_0)\sqrt{n_m}.
\]
If \(M_m\) independent random sign vectors are sampled for the whole active
menu and
\[
  \sum_m2^{-M_m}<\infty ,
\]
then the sampled list contains such a common certificate eventually almost
surely.  If \(J_m\to\infty\), the common signs cover each active prefix in a
countable library eventually.

Thus dependent localized self-adjoint models with small active support-envelope
rank pay \(\log(eR_m)\), not \(\log(eJ_mD_m)\).
\end{corollary}

\begin{proof}
On \((F_m^{\rm sr})^c\), put
\[
  V_{m,h}:=\sum_iA_{m,h,i}^2 .
\]
The support and norm hypotheses give
\[
  0\preceq V_{m,h}
  \preceq
  \sum_i\rho_{m,h,i}^2P_{m,h,i},
  \qquad
  V_{m,h}=Q_{m,h}V_{m,h}Q_{m,h}.
\]
Thus
\[
  \operatorname{rank}V_{m,h}\le\operatorname{rank}Q_{m,h}
\]
and the realized square-profile norms are bounded by
\(C_0n_m/\log(eR_m)\).  The hypotheses of
Corollary~\ref{cor:common-sign-rank-sensitive-selfadjoint-square-profile-menu}
therefore hold with exceptional event \(F_m^{\rm sr}\).  That corollary gives
the common deterministic signing and the sampled-sign conclusion.  The
summability of \(\mathbb P(F_m^{\rm sr})\) and Borel--Cantelli make the
conclusions eventual almost surely.
\end{proof}

\begin{corollary}[Tail-capped support-rank common signs for self-adjoint overlap menus]
\label{cor:common-sign-tail-capped-support-rank-selfadjoint-overlap-budget}
Fix \(C_0\ge1\).  For each \(m\ge1\), let \(D_m\ge2\), \(n_m\ge1\),
\(J_m\ge1\), and let \(R_m\ge1\).  For \(1\le h\le J_m\), let
\[
  A_{m,h,1},\ldots,A_{m,h,n_m}\in\M_{D_m}(\C)
\]
be random self-adjoint matrices.  Suppose that there are events \(F_m^{\rm sr}\),
random orthogonal projections \(P_{m,h,i}\), random support-envelope
projections \(Q_{m,h}\), deterministic caps \(0\le\rho_{m,h,i}\le1\), and
cap-tail budgets \(\delta_{m,h,i}\ge0\) such that
\[
  \mathbb P\{\|A_{m,h,i}\|_{\op}>\rho_{m,h,i}\}\le\delta_{m,h,i}.
\]
Outside \(F_m^{\rm sr}\), assume that
\[
  A_{m,h,i}=P_{m,h,i}A_{m,h,i}P_{m,h,i},\qquad
  P_{m,h,i}\le Q_{m,h}
\]
for all \(h\le J_m\) and \(i\le n_m\), and
\[
  \sum_{h=1}^{J_m}\max\{1,\operatorname{rank}Q_{m,h}\}\le R_m.
\]
Assume also that, outside \(F_m^{\rm sr}\),
\[
  \max_{1\le h\le J_m}
  \left\|\sum_{i=1}^{n_m}\rho_{m,h,i}^2P_{m,h,i}\right\|_{\op}
  \le
  \frac{C_0n_m}{\log(eR_m)}.
\]
If
\[
  \sum_m\left[
    \mathbb P(F_m^{\rm sr})
    +\sum_{h=1}^{J_m}\sum_{i=1}^{n_m}\delta_{m,h,i}
  \right]<\infty ,
\]
then, almost surely, for all sufficiently large \(m\), there is one sign vector
\(\eps^{(m)}\in\{\pm1\}^{n_m}\) such that
\[
  \max_{1\le h\le J_m}
  \left\|\sum_{i=1}^{n_m}\eps_i^{(m)}A_{m,h,i}\right\|_{\op}
  \le C(C_0)\sqrt{n_m}.
\]
If \(M_m\) independent random sign vectors are sampled for the whole active
menu and
\[
  \sum_m2^{-M_m}<\infty ,
\]
then the sampled list contains such a common certificate eventually almost
surely.  If \(J_m\to\infty\), the common signs cover each active prefix in a
countable library eventually.

Thus unbounded localized self-adjoint libraries retain the support-envelope
logarithm \(\log(eR_m)\) after paying only summable cap-tail failures.
\end{corollary}

\begin{proof}
For each \(m\), set
\[
  \Omega_m^{\rm cap}
  :=
  \{\|A_{m,h,i}\|_{\op}\le\rho_{m,h,i}
  \text{ for all }h\le J_m,\ i\le n_m\}.
\]
Then
\[
  \sum_m\mathbb P\bigl((\Omega_m^{\rm cap})^c\bigr)
  \le
  \sum_m\sum_{h=1}^{J_m}\sum_{i=1}^{n_m}\delta_{m,h,i}<\infty .
\]
On the complement of
\[
  G_m:=F_m^{\rm sr}\cup(\Omega_m^{\rm cap})^c
\]
the hypotheses of
Corollary~\ref{cor:common-sign-support-rank-selfadjoint-overlap-budget} hold
with amplitudes \(\rho_{m,h,i}\).  The summability of \(\mathbb P(G_m)\) lets
that corollary and Borel--Cantelli supply the eventual deterministic common
signings, sampled common certificates, and countable-library conclusion.
\end{proof}

\begin{corollary}[Common signs for self-adjoint coordinate-window Bernstein menus]
\label{cor:common-sign-selfadjoint-coordinate-window-bernstein}
Fix \(C_0\ge1\).  For each \(m\ge1\), let \(D_m\ge2\), \(J_m,n_m\ge1\),
and let \(R_m\ge1\).  For every \(1\le h\le J_m\), let
\[
  E_{m,h}\subset\{1,\ldots,D_m\}
\]
be a deterministic coordinate window satisfying
\[
  \sum_{h=1}^{J_m}\max\{1,|E_{m,h}|\}\le R_m .
\]
For each \(h,i\), let \(S_{m,h,i}\subset E_{m,h}\) be a random support.
Assume that, for every coordinate \(j\in E_{m,h}\), the indicators
\({\bf 1}_{\{j\in S_{m,h,i}\}}\), \(1\le i\le n_m\), are independent and
\[
  \mathbb P\{j\in S_{m,h,i}\}\le p_{m,h,i,j}.
\]
Conditional on the supports, let
\[
  A_{m,h,i}\in\M_{D_m}(\C)
\]
be arbitrary random self-adjoint matrices satisfying
\[
  A_{m,h,i}=P_{S_{m,h,i}}A_{m,h,i}P_{S_{m,h,i}}.
\]
Let deterministic caps \(0\le\rho_{m,h,i}\le1\) and cap-tail budgets
\(\delta_{m,h,i}\ge0\) satisfy
\[
  \mathbb P\{\|A_{m,h,i}\|_{\op}>\rho_{m,h,i}\}\le\delta_{m,h,i}.
\]
Put
\[
  b_{m,h}:=\max_i\rho_{m,h,i}^2,\qquad
  \mu_{m,h}:=\max_{j\in E_{m,h}}\sum_i\rho_{m,h,i}^2p_{m,h,i,j},
\]
and
\[
  \sigma_{m,h}^2:=
  \max_{j\in E_{m,h}}\sum_i\rho_{m,h,i}^4p_{m,h,i,j},
\]
with maxima over empty windows interpreted as \(0\).  Suppose that there are
\(t_{m,h}>0\) such that
\[
  \mu_{m,h}+t_{m,h}
  \le
  \frac{C_0n_m}{\log(eR_m)}
  \qquad(1\le h\le J_m),
\]
and
\[
  \begin{aligned}
  \sum_m\Biggl[
  &\sum_{h=1}^{J_m}
  |E_{m,h}|
  \exp\left\{
  -\frac{t_{m,h}^2}
         {2(\sigma_{m,h}^2+b_{m,h}t_{m,h}/3)}
  \right\}\\
  &+
  \sum_{h=1}^{J_m}\sum_{i=1}^{n_m}\delta_{m,h,i}
  \Biggr]<\infty,
  \end{aligned}
\]
with zero denominator interpreted as zero contribution.  Then, almost surely,
for all sufficiently large \(m\), there is one sign vector
\(\eps^{(m)}\in\{\pm1\}^{n_m}\) such that
\[
  \max_{1\le h\le J_m}
  \left\|\sum_{i=1}^{n_m}\eps_i^{(m)}A_{m,h,i}\right\|_{\op}
  \le C(C_0)\sqrt{n_m}.
\]
The same sampled common-certificate and countable-library conclusions as in
Corollary~\ref{cor:common-sign-tail-capped-support-rank-selfadjoint-overlap-budget}
hold.

In particular, if
\[
  \mu_{m,h}\le \frac{C_0n_m}{2\log(eR_m)}
  \qquad(1\le h\le J_m),
\]
then one may take \(t_{m,h}=C_0n_m/(2\log(eR_m))\), and the displayed
Bernstein summand is bounded by
\[
  \sum_{h=1}^{J_m}
  |E_{m,h}|
  \exp\left\{-\frac{3C_0n_m}{16\log(eR_m)}\right\}.
\]
\end{corollary}

\begin{proof}
For each active coordinate define
\[
  N_{m,h,j}:=\sum_i\rho_{m,h,i}^2{\bf 1}_{\{j\in S_{m,h,i}\}}.
\]
Scalar Bernstein and the union bound over the windows \(E_{m,h}\) show that the
probability that some \(N_{m,h,j}\) exceeds \(\mu_{m,h}+t_{m,h}\) is bounded by
the displayed Bernstein summand.  Off this event, with
\(Q_{m,h}=P_{E_{m,h}}\) and \(P_{m,h,i}=P_{S_{m,h,i}}\), the support-envelope
rank and overlap hypotheses of
Corollary~\ref{cor:common-sign-tail-capped-support-rank-selfadjoint-overlap-budget}
hold.  The cap-tail summability is the remaining summand, so that corollary
gives the common signing, sampled certificates, and countable-library
conclusion.

For the final display, \(\sigma_{m,h}^2\le\mu_{m,h}\le t_{m,h}\) and
\(b_{m,h}\le1\), hence the Bernstein exponent is at least \(3t_{m,h}/8\), as
claimed.
\end{proof}

\begin{corollary}[Common signs for self-adjoint predictable coordinate-window Freedman menus]
\label{cor:common-sign-selfadjoint-predictable-coordinate-window-freedman}
Fix \(C_0\ge1\).  For each \(m\ge1\), let \(D_m\ge2\), \(J_m,n_m\ge1\),
and let \(R_m\ge1\).  For every \(1\le h\le J_m\), let
\[
  E_{m,h}\subset\{1,\ldots,D_m\}
\]
be a deterministic coordinate window satisfying
\[
  \sum_{h=1}^{J_m}\max\{1,|E_{m,h}|\}\le R_m .
\]
For each \(m,h\), let
\[
  \mcF_{m,h,0}\subset\mcF_{m,h,1}\subset\cdots\subset\mcF_{m,h,n_m}
\]
be a filtration.  For each \(h,i\), let
\[
  S_{m,h,i}\subset E_{m,h}
\]
be an \(\mcF_{m,h,i}\)-measurable random support.  Assume that there are
deterministic envelopes \(p_{m,h,i,j}\in[0,1]\) such that, for every
\(j\in E_{m,h}\),
\[
  \mathbb P\{j\in S_{m,h,i}\mid\mcF_{m,h,i-1}\}\le p_{m,h,i,j}.
\]
Conditional on the supports, let
\[
  A_{m,h,i}\in\M_{D_m}(\C)
\]
be arbitrary random self-adjoint matrices satisfying
\[
  A_{m,h,i}=P_{S_{m,h,i}}A_{m,h,i}P_{S_{m,h,i}}.
\]
Let deterministic caps \(0\le\rho_{m,h,i}\le1\) and cap-tail budgets
\(\delta_{m,h,i}\ge0\) satisfy
\[
  \mathbb P\{\|A_{m,h,i}\|_{\op}>\rho_{m,h,i}\}\le\delta_{m,h,i}.
\]
Put
\[
  b_{m,h}:=\max_i\rho_{m,h,i}^2,\qquad
  \mu_{m,h}:=\max_{j\in E_{m,h}}\sum_i\rho_{m,h,i}^2p_{m,h,i,j},
\]
and
\[
  \sigma_{m,h}^2:=
  \max_{j\in E_{m,h}}\sum_i\rho_{m,h,i}^4p_{m,h,i,j},
\]
with maxima over empty windows interpreted as \(0\).  Suppose that there are
\(t_{m,h}>0\) such that
\[
  \mu_{m,h}+t_{m,h}
  \le
  \frac{C_0n_m}{\log(eR_m)}
  \qquad(1\le h\le J_m),
\]
and
\[
  \begin{aligned}
  \sum_m\Biggl[
  &\sum_{h=1}^{J_m}
  |E_{m,h}|
  \exp\left\{
  -\frac{t_{m,h}^2}
         {2(\sigma_{m,h}^2+b_{m,h}t_{m,h}/3)}
  \right\}\\
  &+
  \sum_{h=1}^{J_m}\sum_{i=1}^{n_m}\delta_{m,h,i}
  \Biggr]<\infty,
  \end{aligned}
\]
with zero denominator interpreted as zero contribution.  Then, almost surely,
for all sufficiently large \(m\), there is one sign vector
\(\eps^{(m)}\in\{\pm1\}^{n_m}\) such that
\[
  \max_{1\le h\le J_m}
  \left\|\sum_{i=1}^{n_m}\eps_i^{(m)}A_{m,h,i}\right\|_{\op}
  \le C(C_0)\sqrt{n_m}.
\]
The same sampled common-certificate and countable-library conclusions as in
Corollary~\ref{cor:common-sign-tail-capped-support-rank-selfadjoint-overlap-budget}
hold.

In particular, if
\[
  \mu_{m,h}\le \frac{C_0n_m}{2\log(eR_m)}
  \qquad(1\le h\le J_m),
\]
then one may take \(t_{m,h}=C_0n_m/(2\log(eR_m))\), and the displayed Freedman
summand is bounded by
\[
  \sum_{h=1}^{J_m}
  |E_{m,h}|
  \exp\left\{-\frac{3C_0n_m}{16\log(eR_m)}\right\}.
\]
\end{corollary}

\begin{proof}
For each active coordinate define
\[
  N_{m,h,j}:=\sum_i\rho_{m,h,i}^2{\bf 1}_{\{j\in S_{m,h,i}\}}.
\]
For fixed \(m,h,j\), the centered partial sums
\[
  \sum_{i'\le i}\rho_{m,h,i'}^2
  \left(
  {\bf 1}_{\{j\in S_{m,h,i'}\}}
  -\mathbb E({\bf 1}_{\{j\in S_{m,h,i'}\}}\mid\mcF_{m,h,i'-1})
  \right)
\]
form a scalar martingale with increments bounded above by \(b_{m,h}\) and
predictable quadratic variation at most
\(\sum_i\rho_{m,h,i}^4p_{m,h,i,j}\).  Scalar Freedman and the union bound over
the windows \(E_{m,h}\) show that the probability that some \(N_{m,h,j}\)
exceeds \(\mu_{m,h}+t_{m,h}\) is bounded by the displayed Freedman summand.
Off this event, with \(Q_{m,h}=P_{E_{m,h}}\) and
\(P_{m,h,i}=P_{S_{m,h,i}}\), the support-envelope rank and overlap hypotheses of
Corollary~\ref{cor:common-sign-tail-capped-support-rank-selfadjoint-overlap-budget}
hold.  The cap-tail summability is the remaining summand, so that corollary
gives the common signing, sampled certificates, and countable-library
conclusion.  The half-budget estimate follows from
\(\sigma_{m,h}^2\le\mu_{m,h}\le t_{m,h}\) and \(b_{m,h}\le1\), giving exponent
at least \(3t_{m,h}/8\).
\end{proof}

\begin{corollary}[Common-sign self-adjoint predictable coordinate-window schedules]
\label{cor:common-sign-selfadjoint-predictable-coordinate-window-schedules}
In the setting of
Corollary~\ref{cor:common-sign-selfadjoint-predictable-coordinate-window-freedman},
suppose that
\[
  \mu_{m,h}\le \frac{C_0n_m}{2\log(eR_m)}
  \qquad(1\le h\le J_m)
\]
eventually, and that
\[
  \sum_m\sum_{h=1}^{J_m}\sum_{i=1}^{n_m}\delta_{m,h,i}<\infty .
\]
Then the same common-sign, sampled-certificate, and countable-library
conclusions hold whenever
\[
  \sum_m
  R_m\exp\left\{-\frac{3C_0n_m}{16\log(eR_m)}\right\}<\infty .
\]
In particular, after the displayed half-budget and cap-tail assumptions, the
structural tail is summable in each of the following support-rank schedules.
\begin{enumerate}[label=\textup{(\roman*)}]
\item Arbitrary growth.  If, for some \(\gamma>0\),
  \[
    n_m\ge \frac{16}{3C_0}\log(eR_m)
    \bigl(\log(eR_m)+(1+\gamma)\log(e+m)\bigr)
  \]
  eventually, then the almost-sure common-sign conclusion holds.  If a shared
  batch of \(M_m\) independent random sign vectors is used for the whole active
  menu, the sampled conclusion follows from
  \(M_m\ge(1+\delta)\log_2(e+m)\) eventually.
\item All support ranks.  Suppose the rows are indexed by
  \(R=2,3,\ldots\) and the \(R\)-th row has active coordinate-window rank
  parameter at most \(R\).  If
  \[
    n_R\ge a\log^2(eR)
  \]
  eventually for some \(a>32/(3C_0)\), then the common-sign conclusion holds
  along all support ranks.  The sampled conclusion follows if
  \(M_R\ge(1+\delta)\log_2(eR)\) eventually.
\item Sparse support ranks.  If \(R_m\ge\exp(c m^\kappa)\) eventually for some
  \(c,\kappa>0\), and
  \[
    n_m\ge a\log^2(eR_m)
  \]
  eventually for some \(a>16/(3C_0)\), then the common-sign conclusion holds.
  The sampled conclusion follows from
  \(M_m\ge(1+\delta)\log_2(e+m)\) eventually.
\end{enumerate}
\end{corollary}

\begin{proof}
Under the half-budget hypothesis, the final estimate in
Corollary~\ref{cor:common-sign-selfadjoint-predictable-coordinate-window-freedman}
bounds the structural Freedman tail by
\[
  \sum_{h=1}^{J_m}
  |E_{m,h}|
  \exp\left\{-\frac{3C_0n_m}{16\log(eR_m)}\right\}
  \le
  R_m\exp\left\{-\frac{3C_0n_m}{16\log(eR_m)}\right\}.
\]
The cap-tail series is summable by assumption, so the preceding corollary and
Borel--Cantelli give the common signing and the inherited sampled and
countable-library assertions.

For arbitrary growth, the lower bound on \(n_m\) makes the displayed structural
tail at most \(e^{-1}(e+m)^{-(1+\gamma)}\).  For all support ranks, the \(R\)-th
tail is bounded by
\[
  R\exp\left\{-\frac{3C_0a}{16}\log(eR)\right\},
\]
which is summable when \(3C_0a/16>2\).  For sparse support ranks the same
estimate is \(O(R_m^{1-3C_0a/16})\), summable under the stretched-exponential
growth of \(R_m\) as soon as \(3C_0a/16>1\).  The sampled thresholds are exactly
the summability of \(2^{-M_m}\) or \(2^{-M_R}\).
\end{proof}

\begin{corollary}[Common signs for self-adjoint subspace-envelope Bernstein menus]
\label{cor:common-sign-selfadjoint-subspace-envelope-bernstein}
Fix \(C_0\ge1\).  For each \(m\ge1\), let \(D_m\ge2\), \(J_m,n_m\ge1\),
and let \(R_m\ge1\).  For every \(1\le h\le J_m\), let
\[
  Q_{m,h}\in\M_{D_m}(\C)
\]
be a deterministic orthogonal projection satisfying
\[
  \sum_{h=1}^{J_m}\max\{1,\operatorname{rank}Q_{m,h}\}\le R_m .
\]
For each \(h,i\), let \(P_{m,h,i}\le Q_{m,h}\) be random orthogonal
projections, with \(P_{m,h,1},\ldots,P_{m,h,n_m}\) independent for each fixed
\((m,h)\).  Conditional on these projections, let
\[
  A_{m,h,i}\in\M_{D_m}(\C)
\]
be arbitrary random self-adjoint matrices satisfying
\[
  A_{m,h,i}=P_{m,h,i}A_{m,h,i}P_{m,h,i}.
\]
Let deterministic caps \(0\le\rho_{m,h,i}\le1\) and cap-tail budgets
\(\delta_{m,h,i}\ge0\) satisfy
\[
  \mathbb P\{\|A_{m,h,i}\|_{\op}>\rho_{m,h,i}\}\le\delta_{m,h,i}.
\]
Put
\[
  b_{m,h}:=\max_i\rho_{m,h,i}^2,\qquad
  \mu_{m,h}:=
  \left\|\sum_i\rho_{m,h,i}^2\mathbb EP_{m,h,i}\right\|_{\op}.
\]
Suppose that there are \(t_{m,h}>0\) such that
\[
  \mu_{m,h}+t_{m,h}
  \le
  \frac{C_0n_m}{\log(eR_m)}
  \qquad(1\le h\le J_m),
\]
and
\[
  \begin{aligned}
  \sum_m\Biggl[
  &\sum_{h=1}^{J_m}
  \operatorname{rank}Q_{m,h}
  \exp\left\{
  -\frac{t_{m,h}^2}{2(\mu_{m,h}+b_{m,h}t_{m,h}/3)}
  \right\}\\
  &+
  \sum_{h=1}^{J_m}\sum_{i=1}^{n_m}\delta_{m,h,i}
  \Biggr]<\infty,
  \end{aligned}
\]
with zero denominator interpreted as zero contribution.  Then, almost surely,
for all sufficiently large \(m\), there is one sign vector
\(\eps^{(m)}\in\{\pm1\}^{n_m}\) such that
\[
  \max_{1\le h\le J_m}
  \left\|\sum_{i=1}^{n_m}\eps_i^{(m)}A_{m,h,i}\right\|_{\op}
  \le C(C_0)\sqrt{n_m}.
\]
The same sampled common-certificate and countable-library conclusions as in
Corollary~\ref{cor:common-sign-tail-capped-support-rank-selfadjoint-overlap-budget}
hold.

In particular, if
\[
  \mu_{m,h}\le \frac{C_0n_m}{2\log(eR_m)}
  \qquad(1\le h\le J_m),
\]
then one may take \(t_{m,h}=C_0n_m/(2\log(eR_m))\), and the displayed
Bernstein summand is bounded by
\[
  \sum_{h=1}^{J_m}
  \operatorname{rank}Q_{m,h}
  \exp\left\{-\frac{3C_0n_m}{16\log(eR_m)}\right\}.
\]
For Haar-distributed rank-\(r_{m,h,i}\) projections inside the range of
\(Q_{m,h}\), the half-budget condition is
\[
  \frac1{\operatorname{rank}Q_{m,h}}
  \sum_i\rho_{m,h,i}^2r_{m,h,i}
  \le
  \frac{C_0n_m}{2\log(eR_m)},
\]
with the zero-rank case omitted.
\end{corollary}

\begin{proof}
For fixed \(m,h\), apply matrix Bernstein on the range of \(Q_{m,h}\) to
\[
  Y_{m,h,i}:=\rho_{m,h,i}^2(P_{m,h,i}-\mathbb EP_{m,h,i}).
\]
Then \(\mathbb EY_{m,h,i}=0\), \(\|Y_{m,h,i}\|_{\op}\le b_{m,h}\), and
\[
  \sum_i\mathbb EY_{m,h,i}^2
  \preceq
  \sum_i\rho_{m,h,i}^2\mathbb EP_{m,h,i},
\]
so the variance parameter is at most \(\mu_{m,h}\).  The displayed Bernstein
series therefore bounds the probability that some active \(h\) violates
\[
  \left\|\sum_i\rho_{m,h,i}^2P_{m,h,i}\right\|_{\op}
  \le
  \frac{C_0n_m}{\log(eR_m)}.
\]
Off this event, the support-envelope rank and overlap hypotheses of
Corollary~\ref{cor:common-sign-tail-capped-support-rank-selfadjoint-overlap-budget}
hold with envelopes \(Q_{m,h}\).  The cap-tail summability is the remaining
summand, so that corollary gives the common signing, sampled certificates, and
countable-library conclusion.  The half-budget and Haar specializations follow
from the same computation as in
Corollary~\ref{cor:random-subspace-matrix-bernstein}.
\end{proof}

\begin{corollary}[Common signs for self-adjoint predictable subspace-envelope Freedman menus]
\label{cor:common-sign-selfadjoint-predictable-subspace-envelope-freedman}
Fix \(C_0\ge1\).  For each \(m\ge1\), let \(D_m\ge2\), \(J_m,n_m\ge1\),
and let \(R_m\ge1\).  For every \(1\le h\le J_m\), let
\[
  Q_{m,h}\in\M_{D_m}(\C)
\]
be a deterministic orthogonal projection satisfying
\[
  \sum_{h=1}^{J_m}\max\{1,\operatorname{rank}Q_{m,h}\}\le R_m .
\]
For each \(m,h\), let
\[
  \mcF_{m,h,0}\subset\mcF_{m,h,1}\subset\cdots\subset\mcF_{m,h,n_m}
\]
be a filtration.  For each \(h,i\), let \(P_{m,h,i}\le Q_{m,h}\) be an
\(\mcF_{m,h,i}\)-measurable random orthogonal projection.  Conditional on
these projections, let
\[
  A_{m,h,i}\in\M_{D_m}(\C)
\]
be arbitrary random self-adjoint matrices satisfying
\[
  A_{m,h,i}=P_{m,h,i}A_{m,h,i}P_{m,h,i}.
\]
Let deterministic caps \(0\le\rho_{m,h,i}\le1\), cap-tail budgets
\(\delta_{m,h,i}\ge0\), deterministic predictable-overlap budgets
\(v_{m,h}\ge0\), and events \(F_m^{\rm pr}\) be such that
\[
  \mathbb P\{\|A_{m,h,i}\|_{\op}>\rho_{m,h,i}\}\le\delta_{m,h,i},
\]
and, outside \(F_m^{\rm pr}\), for every \(h\le J_m\),
\[
  \left\|\sum_i\rho_{m,h,i}^2
  \mathbb E(P_{m,h,i}\mid\mcF_{m,h,i-1})\right\|_{\op}
  \le v_{m,h}.
\]
Put
\[
  b_{m,h}:=\max_i\rho_{m,h,i}^2 .
\]
Suppose that there are \(t_{m,h}>0\) such that
\[
  v_{m,h}+t_{m,h}
  \le
  \frac{C_0n_m}{\log(eR_m)}
  \qquad(1\le h\le J_m),
\]
and
\[
  \begin{aligned}
  \sum_m\Biggl[
  &\mathbb P(F_m^{\rm pr})
  +\sum_{h=1}^{J_m}
  2\operatorname{rank}Q_{m,h}
  \exp\left\{
  -\frac{t_{m,h}^2}{2(v_{m,h}+b_{m,h}t_{m,h}/3)}
  \right\}\\
  &+
  \sum_{h=1}^{J_m}\sum_{i=1}^{n_m}\delta_{m,h,i}
  \Biggr]<\infty,
  \end{aligned}
\]
with zero denominator interpreted as zero contribution.  Then, almost surely,
for all sufficiently large \(m\), there is one sign vector
\(\eps^{(m)}\in\{\pm1\}^{n_m}\) such that
\[
  \max_{1\le h\le J_m}
  \left\|\sum_{i=1}^{n_m}\eps_i^{(m)}A_{m,h,i}\right\|_{\op}
  \le C(C_0)\sqrt{n_m}.
\]
The same sampled common-certificate and countable-library conclusions as in
Corollary~\ref{cor:common-sign-tail-capped-support-rank-selfadjoint-overlap-budget}
hold.

In particular, if
\[
  v_{m,h}\le \frac{C_0n_m}{2\log(eR_m)}
  \qquad(1\le h\le J_m),
\]
then one may take \(t_{m,h}=C_0n_m/(2\log(eR_m))\), and the Freedman summand is
bounded by
\[
  2\sum_{h=1}^{J_m}
  \operatorname{rank}Q_{m,h}
  \exp\left\{-\frac{3C_0n_m}{16\log(eR_m)}\right\}.
\]
If, conditionally on \(\mcF_{m,h,i-1}\), the projections are Haar-distributed
inside the range of \(Q_{m,h}\) with predictable ranks \(r_{m,h,i}\), the
half-budget condition is
\[
  \frac1{\operatorname{rank}Q_{m,h}}
  \sum_i\rho_{m,h,i}^2r_{m,h,i}
  \le
  \frac{C_0n_m}{2\log(eR_m)},
\]
with the zero-rank case omitted.
\end{corollary}

\begin{proof}
For fixed \(m,h\), apply the two-sided matrix Freedman inequality on the range
of \(Q_{m,h}\) to
\[
  Z_{m,h,i}:=\rho_{m,h,i}^2
  \left(P_{m,h,i}-\mathbb E(P_{m,h,i}\mid\mcF_{m,h,i-1})\right).
\]
Then \((Z_{m,h,i},\mcF_{m,h,i})\) is a self-adjoint
martingale-difference sequence.  Also
\(\|Z_{m,h,i}\|_{\op}\le b_{m,h}\), and
\[
  \sum_i\mathbb E(Z_{m,h,i}^2\mid\mcF_{m,h,i-1})
  \preceq
  \sum_i\rho_{m,h,i}^2
  \mathbb E(P_{m,h,i}\mid\mcF_{m,h,i-1}).
\]
Thus, off \(F_m^{\rm pr}\), the displayed Freedman series bounds the
probability that some active realized overlap exceeds \(v_{m,h}+t_{m,h}\).
Off that event and the cap-failure event, the support-envelope rank and
overlap hypotheses of
Corollary~\ref{cor:common-sign-tail-capped-support-rank-selfadjoint-overlap-budget}
hold with envelopes \(Q_{m,h}\).  The summability assumptions and
Borel--Cantelli therefore give the common signing, sampled certificates, and
countable-library conclusion.  The half-budget estimate uses \(b_{m,h}\le1\)
and \(v_{m,h}\le t_{m,h}\).  Conditional Haar invariance gives
\[
  \mathbb E(P_{m,h,i}\mid\mcF_{m,h,i-1})
  =
  \frac{r_{m,h,i}}{\operatorname{rank}Q_{m,h}}Q_{m,h},
\]
which is the displayed specialization.
\end{proof}

\begin{corollary}[Common-sign self-adjoint predictable subspace-envelope schedules]
\label{cor:common-sign-selfadjoint-predictable-subspace-envelope-schedules}
In the setting of
Corollary~\ref{cor:common-sign-selfadjoint-predictable-subspace-envelope-freedman},
suppose that
\[
  v_{m,h}\le \frac{C_0n_m}{2\log(eR_m)}
  \qquad(1\le h\le J_m)
\]
eventually, and that
\[
  \sum_m\mathbb P(F_m^{\rm pr})<\infty,\qquad
  \sum_m\sum_{h=1}^{J_m}\sum_{i=1}^{n_m}\delta_{m,h,i}<\infty .
\]
Then the same common-sign, sampled-certificate, and countable-library
conclusions hold whenever
\[
  \sum_m
  2R_m\exp\left\{-\frac{3C_0n_m}{16\log(eR_m)}\right\}<\infty .
\]
In particular, after the displayed half-budget, predictable-exception, and
cap-tail assumptions, the structural tail is summable in each of the following
support-rank schedules.
\begin{enumerate}[label=\textup{(\roman*)}]
\item Arbitrary growth.  If, for some \(\gamma>0\),
  \[
    n_m\ge \frac{16}{3C_0}\log(eR_m)
    \bigl(\log(2eR_m)+(1+\gamma)\log(e+m)\bigr)
  \]
  eventually, then the almost-sure common-sign conclusion holds.  If a shared
  batch of \(M_m\) independent random sign vectors is used for the whole active
  menu, the sampled conclusion follows from
  \(M_m\ge(1+\delta)\log_2(e+m)\) eventually.
\item All support ranks.  Suppose the rows are indexed by
  \(R=2,3,\ldots\) and the \(R\)-th row has active subspace-envelope rank
  parameter at most \(R\).  If
  \[
    n_R\ge a\log^2(eR)
  \]
  eventually for some \(a>32/(3C_0)\), then the common-sign conclusion holds
  along all support ranks.  The sampled conclusion follows if
  \(M_R\ge(1+\delta)\log_2(eR)\) eventually.
\item Sparse support ranks.  If \(R_m\ge\exp(c m^\kappa)\) eventually for some
  \(c,\kappa>0\), and
  \[
    n_m\ge a\log^2(eR_m)
  \]
  eventually for some \(a>16/(3C_0)\), then the common-sign conclusion holds.
  The sampled conclusion follows from
  \(M_m\ge(1+\delta)\log_2(e+m)\) eventually.
\end{enumerate}
\end{corollary}

\begin{proof}
Under the half-budget hypothesis, the final estimate in
Corollary~\ref{cor:common-sign-selfadjoint-predictable-subspace-envelope-freedman}
bounds the structural Freedman tail by
\[
  2\sum_{h=1}^{J_m}
  \operatorname{rank}Q_{m,h}
  \exp\left\{-\frac{3C_0n_m}{16\log(eR_m)}\right\}
  \le
  2R_m\exp\left\{-\frac{3C_0n_m}{16\log(eR_m)}\right\}.
\]
The predictable-exception and cap-tail series are summable by assumption, so the
preceding corollary and Borel--Cantelli give the common signing and the inherited
sampled and countable-library assertions.

For arbitrary growth, the lower bound on \(n_m\) makes the displayed structural
tail at most \(e^{-1}(e+m)^{-(1+\gamma)}\).  For all support ranks, the \(R\)-th
tail is bounded by
\[
  2R\exp\left\{-\frac{3C_0a}{16}\log(eR)\right\},
\]
which is summable when \(3C_0a/16>2\).  For sparse support ranks the same
estimate is \(O(R_m^{1-3C_0a/16})\), summable under the stretched-exponential
growth of \(R_m\) as soon as \(3C_0a/16>1\).  The sampled thresholds are exactly
the summability of \(2^{-M_m}\) or \(2^{-M_R}\).
\end{proof}

\begin{corollary}[Tail-capped common signs for dependent self-adjoint overlap menus]
\label{cor:common-sign-tail-capped-dependent-selfadjoint-overlap-budget}
Fix \(C_0\ge1\).  In the setting of
Corollary~\ref{cor:common-sign-dependent-selfadjoint-overlap-budget}, replace
the pointwise norm bounds by deterministic caps
\[
  0\le\rho_{m,h,i}\le1.
\]
Assume that there are cap-tail budgets \(\delta_{m,h,i}\ge0\) such that
\[
  \mathbb P\{\|A_{m,h,i}\|_{\op}>\rho_{m,h,i}\}\le\delta_{m,h,i},
\]
and, outside \(F_{m,h}^{\rm ov}\),
\[
  A_{m,h,i}=P_{m,h,i}A_{m,h,i}P_{m,h,i}
  \qquad(1\le i\le n_m)
\]
and
\[
  \left\|\sum_{i=1}^{n_m}\rho_{m,h,i}^2P_{m,h,i}\right\|_{\op}
  \le
  \frac{C_0n_m}{\log(eD_m^{\rm com})}.
\]
If
\[
  \sum_m\sum_{h=1}^{J_m}
  \left[
    \mathbb P(F_{m,h}^{\rm ov})
    +\sum_{i=1}^{n_m}\delta_{m,h,i}
  \right]<\infty ,
\]
then, almost surely, for all sufficiently large \(m\), one sign vector balances
the whole active menu:
\[
  \max_{1\le h\le J_m}
  \left\|\sum_{i=1}^{n_m}\eps_i^{(m)}A_{m,h,i}\right\|_{\op}
  \le C(C_0)\sqrt{n_m}.
\]
The same sampled common-certificate and countable-library conclusions as in
Corollary~\ref{cor:common-sign-dependent-selfadjoint-overlap-budget} hold.

The sub-Weibull and finite-moment cap choices in
Corollary~\ref{cor:active-menu-tail-capped-dependent-selfadjoint-overlap-budget}
apply verbatim, with the overlap budget measured at the common-sign logarithmic
scale \(\log(eD_m^{\rm com})\).
\end{corollary}

\begin{proof}
For each active stream set
\[
  \Omega_{m,h}^{\rm cap}
  :=\{\|A_{m,h,i}\|_{\op}\le\rho_{m,h,i}\text{ for all }i\le n_m\}.
\]
On the complement of \(F_{m,h}^{\rm ov}\) and on \(\Omega_{m,h}^{\rm cap}\), the
hypotheses of
Corollary~\ref{cor:common-sign-dependent-selfadjoint-overlap-budget} hold for
stream \(h\) at the common dimension \(D_m^{\rm com}\).  The exceptional
probabilities are bounded by the displayed cap-tail sum, so Borel--Cantelli and
the preceding corollary give the eventual common signings.  The sampled and cap
specializations are inherited from the preceding corollary and the one-stream
tail-capped overlap corollary.
\end{proof}

\begin{corollary}[Weighted inhomogeneous Bernstein coordinate-support ensembles]
\label{cor:random-coordinate-support-bernstein}
Let \(D\ge2\), let \(\rho_1,\ldots,\rho_n\in[0,1]\) be deterministic, and let
\(S_1,\ldots,S_n\) be random subsets of \(\{1,\ldots,D\}\).  Suppose that for
every coordinate \(j\) the indicators \(\mathbf 1_{\{j\in S_i\}}\),
\(i=1,\ldots,n\), are independent and satisfy
\[
  \mathbb P\{j\in S_i\}\le p_{ij}
  \qquad(1\le i\le n,\ 1\le j\le D)
\]
for deterministic \(p_{ij}\in[0,1]\).  Conditional on these supports, let
\(A_1,\ldots,A_n\in\M_D(\C)\) be arbitrary random self-adjoint matrices such
that \(A_i=P_{S_i}A_iP_{S_i}\) and \(\|A_i\|_{\op}\le\rho_i\), where
\(P_{S_i}\) is the coordinate projection onto \(S_i\).  Put
\[
  b:=\max_i\rho_i^2,\qquad
  \mu:=\max_{1\le j\le D}\sum_{i=1}^n\rho_i^2p_{ij},\qquad
  \sigma^2:=\max_{1\le j\le D}\sum_{i=1}^n\rho_i^4p_{ij}.
\]
Then for every \(t>0\) and every \(C_0\ge1\) satisfying
\[
  \mu+t\le \frac{C_0n}{\log(eD)},
\]
one has
\[
  \mathbb P\left\{
  \disc(A_1,\ldots,A_n)>C(C_0)\sqrt n
  \right\}
  \le
  D\exp\left\{-\frac{t^2}{2(\sigma^2+bt/3)}\right\},
\]
with the convention that the exponential term is zero when
\(\sigma^2+bt/3=0\).
In particular, if \(p_{ij}\le p\) for all \(i,j\) and
\[
  p\sum_{i=1}^n\rho_i^2\le \frac{C_0n}{2\log(eD)},
\]
then
\[
  \mathbb P\left\{
  \disc(A_1,\ldots,A_n)>C(C_0)\sqrt n
  \right\}
  \le
  D\exp\left\{-\frac{3C_0n}{16\log(eD)}\right\}.
\]
Thus inhomogeneous random coordinate-supported self-adjoint matrices with
weighted coordinate load \(O(n/\log(eD))\) satisfy the Matrix Spencer bound with
high probability as soon as the displayed Bernstein tail is small.  The model
includes Bernoulli supports with nonidentical probabilities, uniformly chosen
fixed-size supports of sizes \(r_i\) satisfying
\(D^{-1}\sum_i\rho_i^2r_i\le C_0n/(2\log(eD))\), and the older homogeneous
case \(p_{ij}=p\), \(\rho_i=1\).
\end{corollary}

\begin{proof}
Let
\[
  N_j=\sum_i\rho_i^2\mathbf 1_{\{j\in S_i\}}.
\]
For each fixed \(j\), the summands in \(N_j\) are independent, lie in
\([0,b]\), have total mean at most \(\mu\), and have total variance at most
\(\sigma^2\).  Bernstein's inequality gives
\[
  \mathbb P\{N_j>\mu+t\}
  \le
  \exp\left\{-\frac{t^2}{2(\sigma^2+bt/3)}\right\}.
\]
A union bound over \(j=1,\ldots,D\) shows that, outside the displayed
exceptional event,
\[
  \left\|\sum_i\rho_i^2P_{S_i}\right\|_{\op}
  =\max_{1\le j\le D}N_j
  \le \mu+t
  \le \frac{C_0n}{\log(eD)}.
\]
Corollary~\ref{cor:random-low-overlap-support} then gives the discrepancy
bound.  For the final display take \(t=C_0n/(2\log(eD))\).  The assumptions give
\(\mu\le t\), while \(\sigma^2\le p\sum_i\rho_i^4\le p\sum_i\rho_i^2\le t\) and
\(b\le1\).  Hence
\[
  \frac{t^2}{2(\sigma^2+bt/3)}
  \ge
  \frac{t^2}{2(t+t/3)}
  =
  \frac{3t}{8}
  =
  \frac{3C_0n}{16\log(eD)}.
\]
For uniformly chosen fixed-size supports, each coordinate is present in each
support with probability \(r_i/D\), and independence holds across \(i\).
\end{proof}

\begin{corollary}[Weighted Matrix-Bernstein random subspace ensembles]
\label{cor:random-subspace-matrix-bernstein}
Let \(P_1,\ldots,P_n\) be independent random orthogonal projections in
\(\M_D(\C)\), let \(\rho_1,\ldots,\rho_n\in[0,1]\) be deterministic, and put
\[
  b:=\max_i\rho_i^2,\qquad
  \mu:=\left\|\sum_{i=1}^n\rho_i^2\mathbb E P_i\right\|_{\op}.
\]
Conditional on these projections, let \(A_1,\ldots,A_n\in\M_D(\C)\) be
arbitrary random self-adjoint matrices satisfying \(A_i=P_iA_iP_i\) and
\(\|A_i\|_{\op}\le\rho_i\).  Then for every \(t>0\) and every \(C_0\ge1\)
satisfying
\[
  \mu+t\le \frac{C_0n}{\log(eD)},
\]
one has
\[
  \mathbb P\left\{
  \disc(A_1,\ldots,A_n)>C(C_0)\sqrt n
  \right\}
  \le
  D\exp\left\{-\frac{t^2}{2(\mu+bt/3)}\right\}.
\]
In particular, if \(P_i\) are independent Haar-distributed projections of ranks
\(r_i\) and
\[
  \frac1D\sum_{i=1}^n\rho_i^2r_i\le \frac{C_0n}{2\log(eD)},
\]
then
\[
  \mathbb P\left\{
  \disc(A_1,\ldots,A_n)>C(C_0)\sqrt n
  \right\}
  \le
  D\exp\left\{-\frac{3C_0n}{16\log(eD)}\right\}.
\]
Thus random subspace-localized Wigner, GOE/GUE, or arbitrary self-adjoint
blocks are Spencer-scale with high probability whenever the weighted average
random subspace rank is \(O(D/\log(eD))\), uniformly in the ambient dimension
apart from the displayed tail.  The old unweighted rank condition is the
special case \(\rho_i=1\).
\end{corollary}

\begin{proof}
Set \(Y_i=\rho_i^2(P_i-\mathbb EP_i)\).  Then \(\mathbb EY_i=0\),
\(\|Y_i\|_{\op}\le b\), and
\[
  \sum_i\mathbb EY_i^2
  =
  \sum_i\rho_i^4\left(\mathbb EP_i-(\mathbb EP_i)^2\right)
  \preceq
  \sum_i\rho_i^2\mathbb EP_i,
\]
so the matrix variance parameter is at most \(\mu\).  The matrix Bernstein
inequality \cite{Tropp} gives
\[
  \mathbb P\left\{
  \left\|\sum_i\rho_i^2P_i\right\|_{\op}>\mu+t
  \right\}
  \le
  D\exp\left\{-\frac{t^2}{2(\mu+bt/3)}\right\}.
\]
On the complementary event,
\[
  \left\|\sum_i\rho_i^2P_i\right\|_{\op}
  \le \mu+t
  \le \frac{C_0n}{\log(eD)},
\]
and Corollary~\ref{cor:random-low-overlap-support} gives the discrepancy
bound.  For Haar rank-\(r_i\) projections,
\(\mathbb EP_i=(r_i/D)I\), hence
\(\mu=D^{-1}\sum_i\rho_i^2r_i\).  Taking
\(t=C_0n/(2\log(eD))\) and using \(\mu\le t\) gives the final displayed
failure bound, since \(b\le1\).
\end{proof}

\begin{corollary}[Random-amplitude weighted support ensembles]
\label{cor:random-amplitude-weighted-support-bernstein}
Fix \(C_0\ge1\).  Let \(D\ge2\).

\begin{enumerate}[label=\textup{(\alph*)}]
\item Coordinate supports.  Let \((\rho_i,S_i)\), \(i=1,\ldots,n\), be
  independent pairs, where \(0\le\rho_i\le1\) and
  \(S_i\subset\{1,\ldots,D\}\).  Assume that for deterministic
  \(p_{ij}\in[0,1]\),
  \[
    \mathbb P\{j\in S_i\mid \rho_i\}\le p_{ij}
    \qquad\text{a.s.}
  \]
  Conditional on the amplitudes and supports, let
  \(A_1,\ldots,A_n\in\M_D(\C)\) be arbitrary random self-adjoint matrices such
  that \(A_i=P_{S_i}A_iP_{S_i}\) and \(\|A_i\|_{\op}\le\rho_i\).  Put
  \[
    \mu_c:=\max_j\sum_i p_{ij}\mathbb E\rho_i^2,
    \qquad
    \sigma_c^2:=\max_j\sum_i p_{ij}\mathbb E\rho_i^4 .
  \]
  Then, for every \(t>0\) with
  \[
    \mu_c+t\le\frac{C_0n}{\log(eD)},
  \]
  one has
  \[
    \mathbb P\left\{
    \disc(A_1,\ldots,A_n)>C(C_0)\sqrt n
    \right\}
    \le
    D\exp\left\{-\frac{t^2}{2(\sigma_c^2+t/3)}\right\}.
  \]
\item Random subspaces.  Let \((\rho_i,P_i)\), \(i=1,\ldots,n\), be independent
  pairs, where \(0\le\rho_i\le1\) and \(P_i\in\M_D(\C)\) is an orthogonal
  projection.  Conditional on the amplitudes and projections, let
  \(A_1,\ldots,A_n\in\M_D(\C)\) be arbitrary random self-adjoint matrices such
  that \(A_i=P_iA_iP_i\) and \(\|A_i\|_{\op}\le\rho_i\).  Put
  \[
    \mu_s:=\left\|\sum_i\mathbb E(\rho_i^2P_i)\right\|_{\op}.
  \]
  Then, for every \(t>0\) with
  \[
    \mu_s+t\le\frac{C_0n}{\log(eD)},
  \]
  one has
  \[
    \mathbb P\left\{
    \disc(A_1,\ldots,A_n)>C(C_0)\sqrt n
    \right\}
    \le
    D\exp\left\{-\frac{t^2}{2(\mu_s+t/3)}\right\}.
  \]
\end{enumerate}

In either case, the half-budget condition
\[
  \mu_c\le\frac{C_0n}{2\log(eD)}
  \qquad\text{or}\qquad
  \mu_s\le\frac{C_0n}{2\log(eD)}
\]
gives the common bound
\[
  \mathbb P\left\{
  \disc(A_1,\ldots,A_n)>C(C_0)\sqrt n
  \right\}
  \le
  D\exp\left\{-\frac{3C_0n}{16\log(eD)}\right\}.
\]
Consequently, the Borel--Cantelli, all-dimension, growth-rate, and sampled
certificate conclusions of
Corollaries~\ref{cor:random-weighted-support-borel},
\ref{cor:random-weighted-support-all-dim-schedules},
\ref{cor:random-weighted-support-search-all-dim-schedules}, and
\ref{cor:random-weighted-support-growth-schedules} remain valid with the
deterministic overlap budgets replaced by \(\mu_c\) or \(\mu_s\).
For uniformly chosen random supports of random sizes \(R_i\), or Haar
projections of random ranks \(R_i\), the homogeneous budget is
\[
  \frac1D\sum_i\mathbb E(\rho_i^2R_i).
\]
\end{corollary}

\begin{proof}
In the coordinate case, set
\[
  W_j=\sum_i\rho_i^2\mathbf 1_{\{j\in S_i\}}.
\]
For fixed \(j\), the summands are independent, lie in \([0,1]\), have total
mean at most \(\sum_i p_{ij}\mathbb E\rho_i^2\), and have total variance at
most
\[
  \sum_i\mathbb E\bigl[\rho_i^4\mathbf 1_{\{j\in S_i\}}\bigr]
  \le \sum_i p_{ij}\mathbb E\rho_i^4 .
\]
Bernstein's inequality and a union bound over \(j\) give
\[
  \mathbb P\left\{\max_j W_j>\mu_c+t\right\}
  \le
  D\exp\left\{-\frac{t^2}{2(\sigma_c^2+t/3)}\right\}.
\]
On the complementary event,
\[
  \left\|\sum_i\rho_i^2P_{S_i}\right\|_{\op}=\max_jW_j
  \le \frac{C_0n}{\log(eD)},
\]
so Corollary~\ref{cor:random-low-overlap-support} gives the discrepancy bound.

In the subspace case, put \(X_i=\rho_i^2P_i\).  Then \(0\preceq X_i\preceq I\)
and
\[
  \left\|\sum_i\mathbb EX_i\right\|_{\op}=\mu_s.
\]
For \(Y_i=X_i-\mathbb EX_i\), the summands are independent, self-adjoint, and
\(\|Y_i\|_{\op}\le1\).  Moreover,
\[
  \sum_i\mathbb EY_i^2
  =
  \sum_i\bigl(\mathbb EX_i^2-(\mathbb EX_i)^2\bigr)
  \preceq
  \sum_i\mathbb EX_i^2
  \preceq
  \sum_i\mathbb EX_i,
\]
because \(X_i^2=\rho_i^4P_i\preceq\rho_i^2P_i=X_i\).  Matrix Bernstein gives
\[
  \mathbb P\left\{\left\|\sum_iX_i\right\|_{\op}>\mu_s+t\right\}
  \le D\exp\left\{-\frac{t^2}{2(\mu_s+t/3)}\right\}.
\]
On the complementary event, Corollary~\ref{cor:random-low-overlap-support}
again applies.

For the half-budget display, take \(t=C_0n/(2\log(eD))\).  In the coordinate
case \(\sigma_c^2\le\mu_c\le t\), and in the subspace case \(\mu_s\le t\), so
both exponents are at least \(3t/8=3C_0n/(16\log(eD))\).  The array and
sampled-certificate assertions are the same Borel--Cantelli applications as in
the cited corollaries, using this identical row-failure bound.
\end{proof}

\begin{corollary}[Random-amplitude block-support ensembles]
\label{cor:random-amplitude-block-support-bernstein}
Fix \(C_0\ge1\).  Let \(D\ge2\), and let
\(\Pi_1,\ldots,\Pi_B\) be pairwise orthogonal projections in \(\M_D(\C)\) with
\(\sum_{\alpha=1}^B\Pi_\alpha=I\).  Let
\((\rho_i,T_i)\), \(i=1,\ldots,n\), be independent pairs, where
\(0\le\rho_i\le1\) and \(T_i\subset\{1,\ldots,B\}\).  Assume that for
deterministic \(q_{i\alpha}\in[0,1]\),
\[
  \mathbb P\{\alpha\in T_i\mid \rho_i\}\le q_{i\alpha}
  \qquad\text{a.s.}
\]
Put \(P_{T_i}=\sum_{\alpha\in T_i}\Pi_\alpha\).  Conditional on the amplitudes
and selected blocks, let \(A_1,\ldots,A_n\in\M_D(\C)\) be arbitrary random
self-adjoint matrices such that
\[
  A_i=P_{T_i}A_iP_{T_i},
  \qquad
  \|A_i\|_{\op}\le\rho_i .
\]
Define
\[
  \mu_B:=\max_{\alpha\le B}\sum_i q_{i\alpha}\mathbb E\rho_i^2,
  \qquad
  \sigma_B^2:=\max_{\alpha\le B}\sum_i q_{i\alpha}\mathbb E\rho_i^4 .
\]
Then, for every \(t>0\) with
\[
  \mu_B+t\le \frac{C_0n}{\log(eD)},
\]
one has
\[
  \mathbb P\left\{
  \disc(A_1,\ldots,A_n)>C(C_0)\sqrt n
  \right\}
  \le
  B\exp\left\{-\frac{t^2}{2(\sigma_B^2+t/3)}\right\}.
\]
In particular, if
\[
  \mu_B\le \frac{C_0n}{2\log(eD)},
\]
then
\[
  \mathbb P\left\{
  \disc(A_1,\ldots,A_n)>C(C_0)\sqrt n
  \right\}
  \le
  B\exp\left\{-\frac{3C_0n}{16\log(eD)}\right\}.
\]
For uniformly chosen block sets of random sizes \(R_i\), the homogeneous
budget is
\[
  \frac1B\sum_i\mathbb E(\rho_i^2R_i).
\]
\end{corollary}

\begin{proof}
Set
\[
  W_\alpha=\sum_i\rho_i^2\mathbf 1_{\{\alpha\in T_i\}}.
\]
For fixed \(\alpha\), the summands are independent, lie in \([0,1]\), have
total mean at most \(\sum_iq_{i\alpha}\mathbb E\rho_i^2\), and have total
variance at most
\[
  \sum_i\mathbb E\bigl[\rho_i^4\mathbf 1_{\{\alpha\in T_i\}}\bigr]
  \le
  \sum_iq_{i\alpha}\mathbb E\rho_i^4 .
\]
Scalar Bernstein and a union bound over the \(B\) blocks give
\[
  \mathbb P\left\{\max_{\alpha\le B}W_\alpha>\mu_B+t\right\}
  \le
  B\exp\left\{-\frac{t^2}{2(\sigma_B^2+t/3)}\right\}.
\]
On the complementary event,
\[
  \left\|\sum_i\rho_i^2P_{T_i}\right\|_{\op}
  =
  \left\|\sum_{\alpha=1}^B W_\alpha\Pi_\alpha\right\|_{\op}
  =
  \max_{\alpha\le B}W_\alpha
  \le \frac{C_0n}{\log(eD)}.
\]
Corollary~\ref{cor:random-low-overlap-support} gives the discrepancy bound.
The half-budget display follows with \(t=C_0n/(2\log(eD))\), since
\(\sigma_B^2\le\mu_B\le t\).  If, conditional on \(R_i\) and \(\rho_i\), the
set \(T_i\) is uniformly chosen among block sets of size \(R_i\), then
\(\mathbb P\{\alpha\in T_i\mid R_i,\rho_i\}=R_i/B\), giving the final budget.
\end{proof}

\begin{corollary}[Random-profile random-amplitude support ensembles]
\label{cor:random-profile-weighted-support-bernstein}
Fix \(C_0\ge1\) and \(D\ge2\).  The following two random-profile variants of
the preceding support theorems hold.
\begin{enumerate}[label=\textup{(\alph*)}]
\item Coordinate profiles.  Let
  \((\rho_i,Q_i,S_i)\), \(i=1,\ldots,n\), be independent triples, where
  \(0\le\rho_i\le1\), \(Q_i=(Q_{ij})_{j=1}^D\in[0,1]^D\), and
  \(S_i\subset\{1,\ldots,D\}\).  Assume that
  \[
    \mathbb P\{j\in S_i\mid \rho_i,Q_i\}\le Q_{ij}
    \qquad\text{a.s.}
  \]
  Conditional on the amplitudes and supports, let
  \(A_1,\ldots,A_n\in\M_D(\C)\) be arbitrary random self-adjoint matrices such
  that \(A_i=P_{S_i}A_iP_{S_i}\) and \(\|A_i\|_{\op}\le\rho_i\).  Put
  \[
    \mu_c^{\rm prof}:=\max_j\sum_i\mathbb E(\rho_i^2Q_{ij}),
    \qquad
    (\sigma_c^{\rm prof})^2:=\max_j\sum_i\mathbb E(\rho_i^4Q_{ij}).
  \]
  If \(t>0\) and
  \[
    \mu_c^{\rm prof}+t\le\frac{C_0n}{\log(eD)},
  \]
  then
  \[
    \mathbb P\left\{
    \disc(A_1,\ldots,A_n)>C(C_0)\sqrt n
    \right\}
    \le
    D\exp\left\{-\frac{t^2}
    {2((\sigma_c^{\rm prof})^2+t/3)}\right\}.
  \]
\item Block profiles.  Let
  \(\Pi_1,\ldots,\Pi_B\) be pairwise orthogonal projections in \(\M_D(\C)\)
  with \(\sum_{\alpha=1}^B\Pi_\alpha=I\).  Let
  \((\rho_i,Q_i,T_i)\), \(i=1,\ldots,n\), be independent triples, where
  \(0\le\rho_i\le1\), \(Q_i=(Q_{i\alpha})_{\alpha=1}^B\in[0,1]^B\), and
  \(T_i\subset\{1,\ldots,B\}\).  Assume that
  \[
    \mathbb P\{\alpha\in T_i\mid \rho_i,Q_i\}\le Q_{i\alpha}
    \qquad\text{a.s.}
  \]
  Put \(P_{T_i}=\sum_{\alpha\in T_i}\Pi_\alpha\).  Conditional on the
  amplitudes and selected blocks, let
  \(A_1,\ldots,A_n\in\M_D(\C)\) be arbitrary random self-adjoint matrices such
  that \(A_i=P_{T_i}A_iP_{T_i}\) and \(\|A_i\|_{\op}\le\rho_i\).  Put
  \[
    \mu_B^{\rm prof}:=\max_{\alpha\le B}\sum_i
      \mathbb E(\rho_i^2Q_{i\alpha}),
    \qquad
    (\sigma_B^{\rm prof})^2:=\max_{\alpha\le B}\sum_i
      \mathbb E(\rho_i^4Q_{i\alpha}).
  \]
  If \(t>0\) and
  \[
    \mu_B^{\rm prof}+t\le\frac{C_0n}{\log(eD)},
  \]
  then
  \[
    \mathbb P\left\{
    \disc(A_1,\ldots,A_n)>C(C_0)\sqrt n
    \right\}
    \le
    B\exp\left\{-\frac{t^2}
    {2((\sigma_B^{\rm prof})^2+t/3)}\right\}.
  \]
\end{enumerate}
In either part, the half-budget condition
\[
  \mu_c^{\rm prof}\le\frac{C_0n}{2\log(eD)}
  \qquad\text{or}\qquad
  \mu_B^{\rm prof}\le\frac{C_0n}{2\log(eD)}
\]
gives the corresponding row failure
\[
  N\exp\left\{-\frac{3C_0n}{16\log(eD)}\right\},
\]
where \(N=D\) in the coordinate case and \(N=B\) in the block case.  Therefore
the all-dimension, sparse, arbitrary-growth, and sampled-certificate schedules
in Corollaries~\ref{cor:random-amplitude-weighted-support-schedules} and
\ref{cor:random-amplitude-block-support-schedules} remain valid after replacing
the deterministic profile budgets by \(\mu_c^{\rm prof}\) or
\(\mu_B^{\rm prof}\).  In particular, uniformly chosen coordinate supports of
random sizes \(R_i\), or uniformly chosen block sets of random sizes \(R_i\),
use the effective budgets
\[
  \frac1D\sum_i\mathbb E(\rho_i^2R_i)
  \qquad\text{and}\qquad
  \frac1B\sum_i\mathbb E(\rho_i^2R_i),
\]
respectively, with no independence required between \(R_i\) and \(\rho_i\)
inside a row variable.
\end{corollary}

\begin{proof}
In the coordinate case, set
\[
  W_j=\sum_i\rho_i^2\mathbf 1_{\{j\in S_i\}}.
\]
For fixed \(j\), the summands are independent, lie in \([0,1]\), and satisfy
\[
  \mathbb E[\rho_i^2\mathbf 1_{\{j\in S_i\}}]
  \le \mathbb E(\rho_i^2Q_{ij}),
  \qquad
  \mathbb E[\rho_i^4\mathbf 1_{\{j\in S_i\}}]
  \le \mathbb E(\rho_i^4Q_{ij}).
\]
Scalar Bernstein and a union bound over \(j\) give
\[
  \mathbb P\left\{\max_jW_j>\mu_c^{\rm prof}+t\right\}
  \le
  D\exp\left\{-\frac{t^2}
  {2((\sigma_c^{\rm prof})^2+t/3)}\right\}.
\]
On the complementary event,
\[
  \left\|\sum_i\rho_i^2P_{S_i}\right\|_{\op}
  =\max_j W_j
  \le \frac{C_0n}{\log(eD)},
\]
and Corollary~\ref{cor:random-low-overlap-support} gives the discrepancy bound.

The block case is identical with
\[
  W_\alpha=\sum_i\rho_i^2\mathbf 1_{\{\alpha\in T_i\}}
\]
and the union bound over \(\alpha=1,\ldots,B\); on the good event
\[
  \left\|\sum_i\rho_i^2P_{T_i}\right\|_{\op}
  =
  \left\|\sum_{\alpha=1}^B W_\alpha\Pi_\alpha\right\|_{\op}
  =
  \max_{\alpha\le B}W_\alpha.
\]
For the half-budget statement take \(t=C_0n/(2\log(eD))\) and use
\(\rho_i^4\le\rho_i^2\).  The schedule and sampled-certificate assertions are
the same Borel--Cantelli and conditional sign-search arguments as in the cited
corollaries, using these identical row tails.  For uniform random sizes, take
\(Q_{ij}=R_i/D\) or \(Q_{i\alpha}=R_i/B\).
\end{proof}

\begin{corollary}[Predictable random-profile support ensembles]
\label{cor:predictable-profile-weighted-support-freedman}
Fix \(C_0\ge1\), \(D\ge2\), and a filtration \((\mcF_i)_{i=0}^n\).  The
following two predictable support models have the same Spencer-scale conclusion.
\begin{enumerate}[label=\textup{(\alph*)}]
\item Coordinate profiles.  Let \(\rho_i\) and
  \(Q_i=(Q_{ij})_{j=1}^D\) be \(\mcF_{i-1}\)-measurable, with
  \(0\le\rho_i\le1\) and \(Q_i\in[0,1]^D\).  Let
  \(S_i\subset\{1,\ldots,D\}\) be \(\mcF_i\)-measurable and assume that
  \[
    \mathbb P\{j\in S_i\mid\mcF_{i-1}\}\le Q_{ij}
    \qquad\text{a.s.}
  \]
  Let \(A_1,\ldots,A_n\in\M_D(\C)\) be arbitrary random self-adjoint matrices
  such that
  \[
    A_i=P_{S_i}A_iP_{S_i},
    \qquad
    \|A_i\|_{\op}\le\rho_i .
  \]
  Put
  \[
    U=\max_{j\le D}\sum_i\rho_i^2Q_{ij},
    \qquad
    V=\max_{j\le D}\sum_i\rho_i^4Q_{ij},
    \qquad
    N=D .
  \]
\item Block profiles.  Let \(\Pi_1,\ldots,\Pi_B\) be pairwise orthogonal
  projections in \(\M_D(\C)\) with \(\sum_{\alpha=1}^B\Pi_\alpha=I\).  Let
  \(\rho_i\) and \(Q_i=(Q_{i\alpha})_{\alpha=1}^B\) be
  \(\mcF_{i-1}\)-measurable, with \(0\le\rho_i\le1\) and
  \(Q_i\in[0,1]^B\).  Let \(T_i\subset\{1,\ldots,B\}\) be
  \(\mcF_i\)-measurable and assume that
  \[
    \mathbb P\{\alpha\in T_i\mid\mcF_{i-1}\}\le Q_{i\alpha}
    \qquad\text{a.s.}
  \]
  Put \(P_{T_i}=\sum_{\alpha\in T_i}\Pi_\alpha\), and let
  \(A_1,\ldots,A_n\in\M_D(\C)\) be arbitrary random self-adjoint matrices such
  that
  \[
    A_i=P_{T_i}A_iP_{T_i},
    \qquad
    \|A_i\|_{\op}\le\rho_i .
  \]
  Put
  \[
    U=\max_{\alpha\le B}\sum_i\rho_i^2Q_{i\alpha},
    \qquad
    V=\max_{\alpha\le B}\sum_i\rho_i^4Q_{i\alpha},
    \qquad
    N=B .
  \]
\end{enumerate}
In either part, if deterministic \(\mu,\sigma^2,t\) satisfy
\[
  \mu+t\le\frac{C_0n}{\log(eD)},
\]
then
\[
\begin{aligned}
  \mathbb P\left\{
  \disc(A_1,\ldots,A_n)>C(C_0)\sqrt n
  \right\}
  &\le
  \mathbb P\{U>\mu\ \text{or}\ V>\sigma^2\}\\
  &\quad+
  N\exp\left\{-\frac{t^2}{2(\sigma^2+t/3)}\right\}.
\end{aligned}
\]
In particular, if
\[
  U\le\frac{C_0n}{2\log(eD)}
  \qquad\text{a.s.},
\]
then
\[
  \mathbb P\left\{
  \disc(A_1,\ldots,A_n)>C(C_0)\sqrt n
  \right\}
  \le
  N\exp\left\{-\frac{3C_0n}{16\log(eD)}\right\}.
\]
If \(\eta^{(1)},\ldots,\eta^{(M)}\) are independent uniform sign vectors in
\(\{\pm1\}^n\), independent of the matrix array, then after enlarging \(C(C_0)\)
if necessary, the same right-hand sides gain only \(+2^{-M}\) for the event that
no sampled sign is a \(C(C_0)\sqrt n\) certificate.  For predictably chosen
uniform coordinate supports, or predictably chosen uniform block sets, of sizes
\(R_i\), the realized half-budget is respectively
\[
  \frac1D\sum_i\rho_i^2R_i
  \qquad\text{or}\qquad
  \frac1B\sum_i\rho_i^2R_i .
\]
Consequently, triangular arrays of either predictable model have eventual
almost-sure signings, and eventual sampled certificates, whenever the displayed
row-failure bounds and the sampled-search terms are summable.
\end{corollary}

\begin{proof}
We prove the coordinate case; the block case is identical after replacing
coordinates by blocks.  For fixed \(j\), set
\[
  X_i=\rho_i^2\mathbf 1_{\{j\in S_i\}},
  \qquad
  m_i=\mathbb E(X_i\mid\mcF_{i-1}).
\]
Then \(0\le X_i\le1\), \(m_i\le\rho_i^2Q_{ij}\), and
\[
  \sum_i\mathbb E\bigl[(X_i-m_i)^2\mid\mcF_{i-1}\bigr]
  \le
  \sum_i\mathbb E(X_i^2\mid\mcF_{i-1})
  \le
  \sum_i\rho_i^4Q_{ij}.
\]
On \(\{U\le\mu,\ V\le\sigma^2\}\), if
\(\sum_iX_i>\mu+t\), then
\(\sum_i(X_i-m_i)>t\), while the scalar predictable quadratic variation is at
most \(\sigma^2\).  The scalar Freedman inequality, as the one-dimensional case
of \cite{Tropp}, gives
\[
  \mathbb P\left\{
  \sum_iX_i>\mu+t,\ U\le\mu,\ V\le\sigma^2
  \right\}
  \le
  \exp\left\{-\frac{t^2}{2(\sigma^2+t/3)}\right\}.
\]
Union over \(j\).  On the complementary event to the resulting occupancy
failure and the budget exception,
\[
  \left\|\sum_i\rho_i^2P_{S_i}\right\|_{\op}
  =
  \max_j\sum_i\rho_i^2\mathbf 1_{\{j\in S_i\}}
  \le \frac{C_0n}{\log(eD)}.
\]
Corollary~\ref{cor:random-low-overlap-support} gives the discrepancy bound.

In the block case the same argument controls
\[
  W_\alpha=\sum_i\rho_i^2\mathbf 1_{\{\alpha\in T_i\}},
\]
and the orthogonality of the blocks gives
\[
  \left\|\sum_i\rho_i^2P_{T_i}\right\|_{\op}
  =
  \left\|\sum_{\alpha=1}^B W_\alpha\Pi_\alpha\right\|_{\op}
  =
  \max_{\alpha\le B}W_\alpha .
\]
For the half-budget estimate take \(\mu=t=C_0n/(2\log(eD))\) and use
\(V\le U\).  The sampled-sign statement follows by conditioning on the
weighted-overlap event and applying
Corollary~\ref{cor:random-sign-search-square-profile}.  The triangular-array
assertion is the first Borel--Cantelli lemma applied to the row failures.  The
uniform-size budgets follow from \(Q_{ij}=R_i/D\) and
\(Q_{i\alpha}=R_i/B\).
\end{proof}

\begin{corollary}[Predictable-profile support schedules]
\label{cor:predictable-profile-weighted-support-schedules}
Fix \(C_0\ge1\).  For each \(m\ge1\), let \(D_m\ge2\), put
\(\ell_m=\log(eD_m)\), and let
\(A_{m,1},\ldots,A_{m,n_m}\in\M_{D_m}(\C)\) be a row satisfying one of the two
predictable support models of
Corollary~\ref{cor:predictable-profile-weighted-support-freedman}.  In the
coordinate case put \(N_m=D_m\) and let \(U_m\) be the coordinate predictable
overlap budget; in the block case put \(N_m=B_m\) and let \(U_m\) be the block
predictable overlap budget.  Assume that
\[
  \sum_m
  \mathbb P\left\{U_m>\frac{C_0n_m}{2\ell_m}\right\}<\infty
\]
and
\[
  \sum_m
  N_m\exp\left\{-\frac{3C_0n_m}{16\ell_m}\right\}<\infty .
\]
Then, with probability one, for all sufficiently large \(m\),
\[
  \disc(A_{m,1},\ldots,A_{m,n_m})\le C(C_0)\sqrt{n_m}.
\]
If each row is also supplied with \(M_m\) independent uniform sign vectors in
\(\{\pm1\}^{n_m}\), independent of the matrix array, and
\(\sum_m2^{-M_m}<\infty\), then after enlarging \(C(C_0)\), almost surely for
all sufficiently large \(m\) one sampled sign certifies the same bound.

In particular, the structural tail is summable in each of the following
schedules, provided the displayed predictable-budget exception is summable.
\begin{enumerate}[label=\textup{(\roman*)}]
\item Arbitrary growth.  If, for some \(\gamma>0\),
  \[
    n_m\ge \frac{16}{3C_0}\,\ell_m
    \bigl(\log(eN_m)+(1+\gamma)\log(e+m)\bigr)
  \]
  eventually, then the almost-sure signing conclusion holds.  The sampled
  conclusion follows if \(M_m\ge(1+\delta)\log_2(e+m)\) eventually for some
  \(\delta>0\).
\item All integer dimensions.  Suppose the rows are indexed by
  \(D=2,3,\ldots\), with \(D_m=D\), and \(N_D\le(eD)^\Gamma\) eventually for
  some \(\Gamma\ge0\).  If
  \[
    n_D\ge a\log^2(eD)
  \]
  eventually for some \(a>16(\Gamma+1)/(3C_0)\), then the almost-sure signing
  conclusion holds.  The sampled conclusion follows if
  \(M_D\ge(1+\delta)\log_2(eD)\) eventually.
\item Sparse dimensions.  If \(D_m\ge\exp(c m^\kappa)\) eventually for some
  \(c,\kappa>0\), \(N_m\le(eD_m)^\Gamma\) eventually, and
  \[
    n_m\ge a\log^2(eD_m)
  \]
  eventually for some \(a>16\Gamma/(3C_0)\), then the almost-sure signing
  conclusion holds.  The sampled conclusion follows if
  \(M_m\ge(1+\delta)\log_2(e+m)\) eventually.
\end{enumerate}
For predictably chosen uniform coordinate supports or uniform block sets of
sizes \(R_{m,i}\), the predictable-budget exception above is respectively
\[
  \left\{\frac1{D_m}\sum_i\rho_{m,i}^2R_{m,i}>
  \frac{C_0n_m}{2\ell_m}\right\}
  \qquad\text{or}\qquad
  \left\{\frac1{B_m}\sum_i\rho_{m,i}^2R_{m,i}>
  \frac{C_0n_m}{2\ell_m}\right\}.
\]
\end{corollary}

\begin{proof}
On the event \(U_m\le C_0n_m/(2\ell_m)\), the corresponding variance budget
satisfies \(V_m\le U_m\).  Applying
Corollary~\ref{cor:predictable-profile-weighted-support-freedman} to the
\(m\)-th row with
\[
  \mu=t=\frac{C_0n_m}{2\ell_m},\qquad
  \sigma^2=\frac{C_0n_m}{2\ell_m},
\]
gives the row failure
\[
  \mathbb P\left\{U_m>\frac{C_0n_m}{2\ell_m}\right\}
  +N_m\exp\left\{-\frac{3C_0n_m}{16\ell_m}\right\}.
\]
The first Borel--Cantelli lemma proves the almost-sure signing conclusion, and
the sampled version follows by adding the summable \(2^{-M_m}\) row failures
from the same corollary.

For arbitrary growth, the displayed lower bound gives
\[
  N_m\exp\left\{-\frac{3C_0n_m}{16\ell_m}\right\}
  \le e^{-1}(e+m)^{-(1+\gamma)}.
\]
For all integer dimensions,
\[
  N_D\exp\left\{-\frac{3C_0n_D}{16\log(eD)}\right\}
  \le (eD)^{\Gamma-3C_0a/16},
\]
which is summable when \(3C_0a/16>\Gamma+1\).  Along sparse schedules the same
bound is at most
\[
  \exp\left\{-(3C_0a/16-\Gamma)\log(eD_m)\right\},
\]
which is summable under the stated stretched-exponential lower bound on \(D_m\).
The sampled-search thresholds are immediate from the summability of
\((e+m)^{-(1+\delta)}\) or \((eD)^{-(1+\delta)}\).  The final display is just the
uniform-size specialization from
Corollary~\ref{cor:predictable-profile-weighted-support-freedman}.
\end{proof}

\begin{corollary}[Finite menus of predictable-profile support schedules]
\label{cor:predictable-profile-menu-schedules}
Fix \(C_0\ge1\).  For each \(m\ge1\), let \(D_m\ge2\), put
\(\ell_m=\log(eD_m)\), and let \(\mathcal H_m\) be a finite menu.  For each
\(h\in\mathcal H_m\), let
\[
  A_{m,1}^{(h)},\ldots,A_{m,n_m}^{(h)}\in\M_{D_m}(\C)
\]
be a row satisfying one of the two predictable support models of
Corollary~\ref{cor:predictable-profile-weighted-support-freedman}, with its own
filtration allowed.  In the coordinate case put \(N_{m,h}=D_m\) and let
\(U_{m,h}\) be the coordinate predictable overlap budget.  In the block case
put \(N_{m,h}=B_{m,h}\), the number of structural blocks for that menu element,
and let \(U_{m,h}\) be the block predictable overlap budget.  Set
\[
  L_m:=\sum_{h\in\mathcal H_m}N_{m,h}.
\]
Assume that
\[
  \sum_m\sum_{h\in\mathcal H_m}
  \mathbb P\left\{U_{m,h}>\frac{C_0n_m}{2\ell_m}\right\}<\infty
\]
and
\[
  \sum_m
  L_m\exp\left\{-\frac{3C_0n_m}{16\ell_m}\right\}<\infty .
\]
Then, with probability one, for all sufficiently large \(m\), every
\(h\in\mathcal H_m\) has a Spencer-scale signing:
\[
  \disc(A_{m,1}^{(h)},\ldots,A_{m,n_m}^{(h)})
  \le C(C_0)\sqrt{n_m}.
\]
If each menu element is also supplied with \(M_{m,h}\) independent uniform sign
vectors in \(\{\pm1\}^{n_m}\), independent of the corresponding row, and
\[
  \sum_m\sum_{h\in\mathcal H_m}2^{-M_{m,h}}<\infty,
\]
then after enlarging \(C(C_0)\), almost surely for all sufficiently large \(m\)
every menu element has one sampled sign certifying the same bound.  If a common
batch of \(M_m\) signs is reused for all \(h\in\mathcal H_m\), independent of
the whole \(m\)-th menu, it is enough that
\[
  \sum_m|\mathcal H_m|\,2^{-M_m}<\infty .
\]

In particular, once the displayed predictable-budget exceptions are summable,
the structural menu tail is summable in each of the following schedules.
\begin{enumerate}[label=\textup{(\roman*)}]
\item Arbitrary growth.  If, for some \(\gamma>0\),
  \[
    n_m\ge \frac{16}{3C_0}\,\ell_m
    \bigl(\log(eL_m)+(1+\gamma)\log(e+m)\bigr)
  \]
  eventually, then the simultaneous almost-sure signing conclusion holds.  With
  a common sign batch, the simultaneous sampled conclusion follows if
  \[
    M_m\ge \log_2(e|\mathcal H_m|)+(1+\delta)\log_2(e+m)
  \]
  eventually for some \(\delta>0\).
\item All integer dimensions.  Suppose the rows are indexed by
  \(D=2,3,\ldots\), with \(D_m=D\), and \(L_D\le(eD)^\Gamma\) eventually for
  some \(\Gamma\ge0\).  If
  \[
    n_D\ge a\log^2(eD)
  \]
  eventually for some \(a>16(\Gamma+1)/(3C_0)\), then the simultaneous
  almost-sure signing conclusion holds.
\item Sparse dimensions.  If \(D_m\ge\exp(c m^\kappa)\) eventually for some
  \(c,\kappa>0\), \(L_m\le(eD_m)^\Gamma\) eventually, and
  \[
    n_m\ge a\log^2(eD_m)
  \]
  eventually for some \(a>16\Gamma/(3C_0)\), then the simultaneous almost-sure
  signing conclusion holds.
\end{enumerate}
For predictably chosen uniform coordinate supports or uniform block sets of
sizes \(R_{m,i}^{(h)}\), the corresponding menu budget exception is respectively
\[
  \left\{\frac1{D_m}\sum_i(\rho_{m,i}^{(h)})^2R_{m,i}^{(h)}>
  \frac{C_0n_m}{2\ell_m}\right\}
  \qquad\text{or}\qquad
  \left\{\frac1{B_{m,h}}\sum_i(\rho_{m,i}^{(h)})^2R_{m,i}^{(h)}>
  \frac{C_0n_m}{2\ell_m}\right\}.
\]
\end{corollary}

\begin{proof}
For fixed \(m\) and \(h\), Corollary~\ref{cor:predictable-profile-weighted-support-freedman}
with
\[
  \mu=t=\sigma^2=\frac{C_0n_m}{2\ell_m}
\]
gives row failure at most
\[
  \mathbb P\left\{U_{m,h}>\frac{C_0n_m}{2\ell_m}\right\}
  +N_{m,h}\exp\left\{-\frac{3C_0n_m}{16\ell_m}\right\}.
\]
Union over \(h\in\mathcal H_m\), then apply the first Borel--Cantelli lemma.
The sampled-sign statements follow by adding the conditional search failure
\(2^{-M_{m,h}}\) for each menu element, or by union bounding a reused common
batch as \(|\mathcal H_m|2^{-M_m}\).

For arbitrary growth, the displayed lower bound gives
\[
  L_m\exp\left\{-\frac{3C_0n_m}{16\ell_m}\right\}
  \le e^{-1}(e+m)^{-(1+\gamma)}.
\]
The common sampled-batch condition gives
\(|\mathcal H_m|2^{-M_m}\le e^{-1}(e+m)^{-(1+\delta)}\).  The all-integer and
sparse cases are the same estimates with \(L_m\le(eD_m)^\Gamma\).  The final
display is the uniform-size specialization from
Corollary~\ref{cor:predictable-profile-weighted-support-freedman}.
\end{proof}

\begin{corollary}[Exhaustive countable menus of predictable-profile support schedules]
\label{cor:predictable-profile-countable-exhaustion-schedules}
Fix \(C_0\ge1\).  For each \(m\ge1\), let \(D_m\ge2\), put
\(\ell_m=\log(eD_m)\), and let
\[
  J_1\le J_2\le\cdots,\qquad J_m\to\infty,
\]
be a deterministic exhaustion of the countable menu
\(\mathbb N=\{1,2,\ldots\}\).  For each \(m\) and \(h\ge1\), let
\[
  A_{m,1}^{(h)},\ldots,A_{m,n_m}^{(h)}\in\M_{D_m}(\C)
\]
be a row satisfying one of the two predictable support models of
Corollary~\ref{cor:predictable-profile-weighted-support-freedman}, with its own
filtration allowed.  In the coordinate case put \(N_{m,h}=D_m\) and let
\(U_{m,h}\) be the coordinate predictable overlap budget.  In the block case
put \(N_{m,h}=B_{m,h}\), the number of structural blocks for that menu element,
and let \(U_{m,h}\) be the block predictable overlap budget.  Set
\[
  L_m^{(J)}:=\sum_{1\le h\le J_m}N_{m,h}.
\]
Assume that
\[
  \sum_m\sum_{1\le h\le J_m}
  \mathbb P\left\{U_{m,h}>\frac{C_0n_m}{2\ell_m}\right\}<\infty
\]
and
\[
  \sum_m
  L_m^{(J)}\exp\left\{-\frac{3C_0n_m}{16\ell_m}\right\}<\infty .
\]
Then, with probability one, for all sufficiently large \(m\), every model in
the active prefix \(1\le h\le J_m\) has a Spencer-scale signing:
\[
  \disc(A_{m,1}^{(h)},\ldots,A_{m,n_m}^{(h)})
  \le C(C_0)\sqrt{n_m}\qquad(1\le h\le J_m).
\]
Consequently, for each fixed \(h\), the same bound holds for all sufficiently
large \(m\).  If the \(h\)-th menu element is also supplied with \(M_{m,h}\)
independent uniform sign vectors in \(\{\pm1\}^{n_m}\), independent of the
corresponding row, and
\[
  \sum_m\sum_{1\le h\le J_m}2^{-M_{m,h}}<\infty,
\]
then after enlarging \(C(C_0)\), almost surely every active prefix element has
one sampled sign certifying the same bound for all sufficiently large \(m\).
If a common batch of \(M_m\) signs is reused for all \(1\le h\le J_m\),
independent of the whole active prefix, it is enough that
\[
  \sum_m J_m\,2^{-M_m}<\infty .
\]

In particular, once the displayed predictable-budget exceptions over the active
prefixes are summable, the structural prefix tail is summable in each of the
following schedules.
\begin{enumerate}[label=\textup{(\roman*)}]
\item Arbitrary growth.  If, for some \(\gamma>0\),
  \[
    n_m\ge \frac{16}{3C_0}\,\ell_m
    \bigl(\log(eL_m^{(J)})+(1+\gamma)\log(e+m)\bigr)
  \]
  eventually, then the active-prefix and fixed-model almost-sure signing
  conclusions hold.  With a common sign batch, the sampled conclusion follows
  if
  \[
    M_m\ge \log_2(eJ_m)+(1+\delta)\log_2(e+m)
  \]
  eventually for some \(\delta>0\).
\item All integer dimensions.  Suppose the rows are indexed by
  \(D=2,3,\ldots\), with \(D_m=D\), \(J_D\to\infty\), and
  \(L_D^{(J)}\le(eD)^\Gamma\) eventually for some \(\Gamma\ge0\).  If
  \[
    n_D\ge a\log^2(eD)
  \]
  eventually for some \(a>16(\Gamma+1)/(3C_0)\), then the active-prefix and
  fixed-model almost-sure signing conclusions hold.
\item Sparse dimensions.  If \(D_m\ge\exp(c m^\kappa)\) eventually for some
  \(c,\kappa>0\), \(J_m\to\infty\), \(L_m^{(J)}\le(eD_m)^\Gamma\) eventually,
  and
  \[
    n_m\ge a\log^2(eD_m)
  \]
  eventually for some \(a>16\Gamma/(3C_0)\), then the active-prefix and
  fixed-model almost-sure signing conclusions hold.
\end{enumerate}
For predictably chosen uniform coordinate supports or uniform block sets of
sizes \(R_{m,i}^{(h)}\), the prefix budget exceptions are the corresponding
events
\[
  \left\{\frac1{D_m}\sum_i(\rho_{m,i}^{(h)})^2R_{m,i}^{(h)}>
  \frac{C_0n_m}{2\ell_m}\right\}
  \qquad\text{or}\qquad
  \left\{\frac1{B_{m,h}}\sum_i(\rho_{m,i}^{(h)})^2R_{m,i}^{(h)}>
  \frac{C_0n_m}{2\ell_m}\right\},
\]
summed over \(m\) and \(1\le h\le J_m\).
\end{corollary}

\begin{proof}
Apply Corollary~\ref{cor:predictable-profile-menu-schedules} to the finite
active menu \(\{1,\ldots,J_m\}\) in the \(m\)-th row.  The two displayed
summability assumptions are exactly the finite-menu budget-exception series and
the finite-menu structural series for these active prefixes, so
Borel--Cantelli gives the simultaneous active-prefix conclusion.  Since
\(J_m\to\infty\), each fixed \(h\) belongs to the active prefix for all
sufficiently large \(m\), giving the fixed-model statement.  The sampled-sign
claims are the finite-menu sampled claims summed over \(1\le h\le J_m\), or the
common-batch union bound \(J_m2^{-M_m}\).

The arbitrary-growth estimate is
\[
  L_m^{(J)}\exp\left\{-\frac{3C_0n_m}{16\ell_m}\right\}
  \le e^{-1}(e+m)^{-(1+\gamma)}.
\]
The common sampled-batch condition gives
\(J_m2^{-M_m}\le e^{-1}(e+m)^{-(1+\delta)}\).  The all-integer and sparse
dimension cases are the same estimates with
\(L_m^{(J)}\le(eD_m)^\Gamma\).  The uniform-size events are inherited from
Corollary~\ref{cor:predictable-profile-weighted-support-freedman}.
\end{proof}

\begin{corollary}[Adaptive deterministic-envelope support ensembles]
\label{cor:adaptive-envelope-weighted-support-freedman}
Fix \(C_0\ge1\), \(D\ge2\), and a filtration \((\mcF_i)_{i=0}^n\).  The
following adaptive deterministic-envelope models satisfy Spencer-scale tails
with no budget-exception term.
\begin{enumerate}[label=\textup{(\alph*)}]
\item Coordinate supports.  Let \(a_i\in[0,1]\) and
  \(p_{ij}\in[0,1]\) be deterministic.  Let \(\rho_i\) be
  \(\mcF_{i-1}\)-measurable with \(0\le\rho_i\le a_i\), and let
  \(S_i\subset\{1,\ldots,D\}\) be \(\mcF_i\)-measurable with
  \[
    \mathbb P\{j\in S_i\mid\mcF_{i-1}\}\le p_{ij}
    \qquad\text{a.s.}
  \]
  Conditional on the amplitudes and supports, let
  \(A_1,\ldots,A_n\in\M_D(\C)\) be arbitrary random self-adjoint matrices such
  that \(A_i=P_{S_i}A_iP_{S_i}\) and \(\|A_i\|_{\op}\le\rho_i\).  Put
  \[
    \mu_c^{\rm env}:=\max_{j\le D}\sum_i a_i^2p_{ij}.
  \]
  If \(t>0\) and
  \[
    \mu_c^{\rm env}+t\le\frac{C_0n}{\log(eD)},
  \]
  then
  \[
    \mathbb P\left\{
    \disc(A_1,\ldots,A_n)>C(C_0)\sqrt n
    \right\}
    \le
    D\exp\left\{-\frac{t^2}{2(\mu_c^{\rm env}+t/3)}\right\}.
  \]
\item Block supports.  Let \(\Pi_1,\ldots,\Pi_B\) be pairwise orthogonal
  projections in \(\M_D(\C)\) with \(\sum_{\alpha=1}^B\Pi_\alpha=I\).  Let
  \(a_i\in[0,1]\) and \(q_{i\alpha}\in[0,1]\) be deterministic.  Let \(\rho_i\)
  be \(\mcF_{i-1}\)-measurable with \(0\le\rho_i\le a_i\), and let
  \(T_i\subset\{1,\ldots,B\}\) be \(\mcF_i\)-measurable with
  \[
    \mathbb P\{\alpha\in T_i\mid\mcF_{i-1}\}\le q_{i\alpha}
    \qquad\text{a.s.}
  \]
  Put \(P_{T_i}=\sum_{\alpha\in T_i}\Pi_\alpha\).  Conditional on the
  amplitudes and selected blocks, let
  \(A_1,\ldots,A_n\in\M_D(\C)\) be arbitrary random self-adjoint matrices such
  that \(A_i=P_{T_i}A_iP_{T_i}\) and \(\|A_i\|_{\op}\le\rho_i\).  Put
  \[
    \mu_B^{\rm env}:=\max_{\alpha\le B}\sum_i a_i^2q_{i\alpha}.
  \]
  If \(t>0\) and
  \[
    \mu_B^{\rm env}+t\le\frac{C_0n}{\log(eD)},
  \]
  then
  \[
    \mathbb P\left\{
    \disc(A_1,\ldots,A_n)>C(C_0)\sqrt n
    \right\}
    \le
    B\exp\left\{-\frac{t^2}{2(\mu_B^{\rm env}+t/3)}\right\}.
  \]
\end{enumerate}
In either part, the half-budget condition
\[
  \mu_c^{\rm env}\le\frac{C_0n}{2\log(eD)}
  \qquad\text{or}\qquad
  \mu_B^{\rm env}\le\frac{C_0n}{2\log(eD)}
\]
gives the corresponding row failure
\[
  N\exp\left\{-\frac{3C_0n}{16\log(eD)}\right\},
\]
where \(N=D\) in the coordinate case and \(N=B\) in the block case.  Therefore
the schedule and sampled-certificate conclusions of
Corollary~\ref{cor:predictable-profile-weighted-support-schedules} apply with
zero predictable-budget exception.  In particular, if the row supports are
predictably uniform of sizes \(R_i\le r_i\), then the deterministic envelope
budgets are
\[
  \frac1D\sum_i a_i^2r_i
  \qquad\text{and}\qquad
  \frac1B\sum_i a_i^2r_i
\]
for coordinate supports and block supports, respectively.
\end{corollary}

\begin{proof}
In the coordinate case apply
Corollary~\ref{cor:predictable-profile-weighted-support-freedman} with the
predictable profile \(Q_{ij}=p_{ij}\).  Then
\[
  U=\max_j\sum_i\rho_i^2p_{ij}\le \mu_c^{\rm env},
  \qquad
  V=\max_j\sum_i\rho_i^4p_{ij}\le \mu_c^{\rm env},
\]
so the finite tail follows with \(\mu=\sigma^2=\mu_c^{\rm env}\).  The block
case is identical with \(Q_{i\alpha}=q_{i\alpha}\).  The half-budget estimate
uses \(t=C_0n/(2\log(eD))\), and the schedule and sampled-certificate assertions
are the corresponding zero-exception specialization of
Corollary~\ref{cor:predictable-profile-weighted-support-schedules}.  For
predictably uniform support sizes, take \(p_{ij}=r_i/D\) or
\(q_{i\alpha}=r_i/B\).
\end{proof}

\begin{corollary}[Adaptive deterministic-envelope support schedules]
\label{cor:adaptive-envelope-weighted-support-schedules}
Fix \(C_0\ge1\).  For each \(m\ge1\), let \(D_m\ge2\), put
\(\ell_m=\log(eD_m)\), and let
\(A_{m,1},\ldots,A_{m,n_m}\in\M_{D_m}(\C)\) be a row satisfying one of the two
deterministic-envelope models of
Corollary~\ref{cor:adaptive-envelope-weighted-support-freedman}.  In the
coordinate case set \(N_m=D_m\) and
\[
  \mu_m^{\rm env}=\max_{j\le D_m}
  \sum_i a_{m,i}^2p_{m,ij}.
\]
In the block case, with \(B_m\) structural blocks, set \(N_m=B_m\) and
\[
  \mu_m^{\rm env}=\max_{\alpha\le B_m}
  \sum_i a_{m,i}^2q_{m,i\alpha}.
\]
Assume that the deterministic half-budget condition
\[
  \mu_m^{\rm env}\le \frac{C_0n_m}{2\ell_m}
\]
holds eventually, and that
\[
  \sum_m N_m\exp\left\{-\frac{3C_0n_m}{16\ell_m}\right\}<\infty .
\]
Then, with probability one, for all sufficiently large \(m\),
\[
  \disc(A_{m,1},\ldots,A_{m,n_m})\le C(C_0)\sqrt{n_m}.
\]
If each row is also supplied with \(M_m\) independent uniform sign vectors in
\(\{\pm1\}^{n_m}\), independent of the matrix array, and
\(\sum_m2^{-M_m}<\infty\), then after enlarging \(C(C_0)\), almost surely for
all sufficiently large \(m\) one sampled sign certifies the same bound.

In particular, under the deterministic half-budget condition, the following
growth schedules imply the corresponding almost-sure conclusions.
\begin{enumerate}[label=\textup{(\roman*)}]
\item Arbitrary growth.  If, for some \(\gamma>0\),
  \[
    n_m\ge \frac{16}{3C_0}\,\ell_m
    \bigl(\log(eN_m)+(1+\gamma)\log(e+m)\bigr)
  \]
  eventually, then the almost-sure signing conclusion holds.  The sampled
  conclusion follows if \(M_m\ge(1+\delta)\log_2(e+m)\) eventually for some
  \(\delta>0\).
\item All integer dimensions.  Suppose the rows are indexed by
  \(D=2,3,\ldots\), with \(D_m=D\), and \(N_D\le(eD)^\Gamma\) eventually for
  some \(\Gamma\ge0\).  If
  \[
    n_D\ge a\log^2(eD)
  \]
  eventually for some \(a>16(\Gamma+1)/(3C_0)\), then the almost-sure signing
  conclusion holds.  The sampled conclusion follows if
  \(M_D\ge(1+\delta)\log_2(eD)\) eventually.
\item Sparse dimensions.  If \(D_m\ge\exp(c m^\kappa)\) eventually for some
  \(c,\kappa>0\), \(N_m\le(eD_m)^\Gamma\) eventually, and
  \[
    n_m\ge a\log^2(eD_m)
  \]
  eventually for some \(a>16\Gamma/(3C_0)\), then the almost-sure signing
  conclusion holds.  The sampled conclusion follows if
  \(M_m\ge(1+\delta)\log_2(e+m)\) eventually.
\end{enumerate}
For predictably uniform coordinate supports or uniform block sets with
deterministic size caps \(R_{m,i}\le r_{m,i}\), the half-budget condition is
respectively
\[
  \frac1{D_m}\sum_i a_{m,i}^2r_{m,i}
  \le\frac{C_0n_m}{2\ell_m}
  \qquad\text{or}\qquad
  \frac1{B_m}\sum_i a_{m,i}^2r_{m,i}
  \le\frac{C_0n_m}{2\ell_m}.
\]
\end{corollary}

\begin{proof}
Corollary~\ref{cor:adaptive-envelope-weighted-support-freedman}, with
\(t=C_0n_m/(2\ell_m)\), gives the row failure
\[
  N_m\exp\left\{-\frac{3C_0n_m}{16\ell_m}\right\}
\]
whenever the deterministic half-budget holds.  The first Borel--Cantelli lemma
proves the almost-sure signing conclusion, and the sampled version follows by
adding the summable \(2^{-M_m}\) row failures.  The three growth schedules are
exactly the summability estimates in
Corollary~\ref{cor:predictable-profile-weighted-support-schedules}, now with no
predictable-budget exception term.  The final display is the uniform-size
specialization from
Corollary~\ref{cor:adaptive-envelope-weighted-support-freedman}.
\end{proof}

\begin{corollary}[Finite menus of adaptive deterministic-envelope schedules]
\label{cor:adaptive-envelope-menu-schedules}
Fix \(C_0\ge1\).  For each \(m\ge1\), let \(D_m\ge2\), put
\(\ell_m=\log(eD_m)\), and let \(\mathcal H_m\) be a finite menu.  For each
\(h\in\mathcal H_m\), let
\[
  A_{m,1}^{(h)},\ldots,A_{m,n_m}^{(h)}\in\M_{D_m}(\C)
\]
be a random self-adjoint row satisfying one of the two deterministic-envelope
models of Corollary~\ref{cor:adaptive-envelope-weighted-support-freedman}, with
its own filtration allowed.  If \(h\) is a coordinate-support model, put
\[
  N_{m,h}=D_m,\qquad
  \mu_{m,h}^{\rm env}=\max_{j\le D_m}\sum_i a_{m,i,h}^2p_{m,ij,h}.
\]
If \(h\) is a block-support model with \(B_{m,h}\) structural blocks, put
\[
  N_{m,h}=B_{m,h},\qquad
  \mu_{m,h}^{\rm env}
  =\max_{\alpha\le B_{m,h}}\sum_i a_{m,i,h}^2q_{m,i\alpha,h}.
\]
Finally set
\[
  L_m=\sum_{h\in\mathcal H_m}N_{m,h}.
\]
Assume that the deterministic half-budget condition
\[
  \max_{h\in\mathcal H_m}\mu_{m,h}^{\rm env}
  \le \frac{C_0n_m}{2\ell_m}
\]
holds eventually, and that
\[
  \sum_m L_m\exp\left\{-\frac{3C_0n_m}{16\ell_m}\right\}<\infty .
\]
Then, with probability one, for all sufficiently large \(m\), every
\(h\in\mathcal H_m\) has a Spencer-scale signing:
\[
  \disc(A_{m,1}^{(h)},\ldots,A_{m,n_m}^{(h)})
  \le C(C_0)\sqrt{n_m}.
\]
If each menu element is also supplied with \(M_{m,h}\) independent uniform sign
vectors in \(\{\pm1\}^{n_m}\), independent of the corresponding row, and
\[
  \sum_m\sum_{h\in\mathcal H_m}2^{-M_{m,h}}<\infty,
\]
then after enlarging \(C(C_0)\), almost surely for all sufficiently large \(m\)
every \(h\in\mathcal H_m\) has one sampled sign certifying the same bound.  If a
common batch of \(M_m\) signs is reused for all \(h\in\mathcal H_m\), the same
conclusion follows from
\[
  \sum_m |\mathcal H_m|\,2^{-M_m}<\infty .
\]

In particular, under the deterministic half-budget condition, the structural
menu tail is summable in each of the following schedules.
\begin{enumerate}[label=\textup{(\roman*)}]
\item Arbitrary growth.  If, for some \(\gamma>0\),
  \[
    n_m\ge \frac{16}{3C_0}\,\ell_m
    \bigl(\log(eL_m)+(1+\gamma)\log(e+m)\bigr)
  \]
  eventually, then the simultaneous almost-sure signing conclusion holds.  With
  a common sign batch, the simultaneous sampled conclusion follows if
  \[
    M_m\ge \log_2(e|\mathcal H_m|)+(1+\delta)\log_2(e+m)
  \]
  eventually for some \(\delta>0\).
\item All integer dimensions.  Suppose the rows are indexed by
  \(D=2,3,\ldots\), with \(D_m=D\), and \(L_D\le(eD)^\Gamma\) eventually for
  some \(\Gamma\ge0\).  If
  \[
    n_D\ge a\log^2(eD)
  \]
  eventually for some \(a>16(\Gamma+1)/(3C_0)\), then the simultaneous
  almost-sure signing conclusion holds.
\item Sparse dimensions.  If \(D_m\ge\exp(c m^\kappa)\) eventually for some
  \(c,\kappa>0\), \(L_m\le(eD_m)^\Gamma\) eventually, and
  \[
    n_m\ge a\log^2(eD_m)
  \]
  eventually for some \(a>16\Gamma/(3C_0)\), then the simultaneous almost-sure
  signing conclusion holds.
\end{enumerate}
\end{corollary}

\begin{proof}
For each fixed \(m\) and \(h\), Corollary~\ref{cor:adaptive-envelope-weighted-support-schedules}
gives the row failure
\[
  N_{m,h}\exp\left\{-\frac{3C_0n_m}{16\ell_m}\right\}
\]
under the deterministic half-budget condition.  Union over \(h\in\mathcal H_m\);
the structural failure is at most
\[
  L_m\exp\left\{-\frac{3C_0n_m}{16\ell_m}\right\}.
\]
The first Borel--Cantelli lemma gives the simultaneous eventual signing
statement.  Conditional sign search gives an additional \(2^{-M_{m,h}}\) failure
for each menu element, and union over the menu gives either the displayed
double sum or, for a reused common batch, the bound \(|\mathcal H_m|2^{-M_m}\).

The arbitrary-growth lower bound gives
\[
  L_m\exp\left\{-\frac{3C_0n_m}{16\ell_m}\right\}
  \le e^{-1}(e+m)^{-(1+\gamma)}.
\]
The common sampled-batch condition similarly gives
\(|\mathcal H_m|2^{-M_m}\le e^{-1}(e+m)^{-(1+\delta)}\).  The all-integer and sparse
dimension estimates are the same summability computations with \(L_m\) in
place of \(N_m\).
\end{proof}

\begin{corollary}[Exhaustive countable menus of adaptive deterministic-envelope schedules]
\label{cor:adaptive-envelope-countable-exhaustion-schedules}
Fix \(C_0\ge1\).  For each \(m\ge1\), let \(D_m\ge2\), put
\(\ell_m=\log(eD_m)\), and let
\[
  J_1\le J_2\le\cdots,\qquad J_m\to\infty,
\]
be a deterministic exhaustion of the countable menu
\(\mathbb N=\{1,2,\ldots\}\).  For each \(m\) and \(h\ge1\), let
\[
  A_{m,1}^{(h)},\ldots,A_{m,n_m}^{(h)}\in\M_{D_m}(\C)
\]
be a random self-adjoint row satisfying one of the two deterministic-envelope
models of Corollary~\ref{cor:adaptive-envelope-weighted-support-freedman}, with
its own filtration allowed.  If \(h\) is a coordinate-support model, put
\[
  N_{m,h}=D_m,\qquad
  \mu_{m,h}^{\rm env}=\max_{j\le D_m}\sum_i a_{m,i,h}^2p_{m,ij,h}.
\]
If \(h\) is a block-support model with \(B_{m,h}\) structural blocks, put
\[
  N_{m,h}=B_{m,h},\qquad
  \mu_{m,h}^{\rm env}
  =\max_{\alpha\le B_{m,h}}\sum_i a_{m,i,h}^2q_{m,i\alpha,h}.
\]
Set
\[
  L_m^{(J)}:=\sum_{1\le h\le J_m}N_{m,h}.
\]
Assume that the deterministic prefix half-budget condition
\[
  \max_{1\le h\le J_m}\mu_{m,h}^{\rm env}
  \le \frac{C_0n_m}{2\ell_m}
\]
holds eventually, and that
\[
  \sum_m L_m^{(J)}
  \exp\left\{-\frac{3C_0n_m}{16\ell_m}\right\}<\infty .
\]
Then, with probability one, for all sufficiently large \(m\), every active
prefix model \(1\le h\le J_m\) has a Spencer-scale signing:
\[
  \disc(A_{m,1}^{(h)},\ldots,A_{m,n_m}^{(h)})
  \le C(C_0)\sqrt{n_m}\qquad(1\le h\le J_m).
\]
Consequently, for each fixed \(h\), the same bound holds for all sufficiently
large \(m\).  If the active-prefix models are supplied with independent sampled
signs and
\[
  \sum_m\sum_{1\le h\le J_m}2^{-M_{m,h}}<\infty,
\]
then after enlarging \(C(C_0)\), almost surely every active-prefix model has a
sampled certificate for all sufficiently large \(m\).  If one common batch of
\(M_m\) signs is reused for all \(1\le h\le J_m\), independent of the whole
active prefix, it is enough that
\[
  \sum_m J_m\,2^{-M_m}<\infty .
\]

In particular, under the deterministic prefix half-budget condition, the
structural prefix tail is summable in each of the following schedules.
\begin{enumerate}[label=\textup{(\roman*)}]
\item Arbitrary growth.  If, for some \(\gamma>0\),
  \[
    n_m\ge \frac{16}{3C_0}\,\ell_m
    \bigl(\log(eL_m^{(J)})+(1+\gamma)\log(e+m)\bigr)
  \]
  eventually, then the active-prefix and fixed-model almost-sure signing
  conclusions hold.  With a common sign batch, the sampled conclusion follows
  if
  \[
    M_m\ge \log_2(eJ_m)+(1+\delta)\log_2(e+m)
  \]
  eventually for some \(\delta>0\).
\item All integer dimensions.  Suppose the rows are indexed by
  \(D=2,3,\ldots\), with \(D_m=D\), \(J_D\to\infty\), and
  \(L_D^{(J)}\le(eD)^\Gamma\) eventually for some \(\Gamma\ge0\).  If
  \[
    n_D\ge a\log^2(eD)
  \]
  eventually for some \(a>16(\Gamma+1)/(3C_0)\), then the active-prefix and
  fixed-model almost-sure signing conclusions hold.
\item Sparse dimensions.  If \(D_m\ge\exp(c m^\kappa)\) eventually for some
  \(c,\kappa>0\), \(J_m\to\infty\), \(L_m^{(J)}\le(eD_m)^\Gamma\) eventually,
  and
  \[
    n_m\ge a\log^2(eD_m)
  \]
  eventually for some \(a>16\Gamma/(3C_0)\), then the active-prefix and
  fixed-model almost-sure signing conclusions hold.
\end{enumerate}
For predictably uniform coordinate supports or uniform block sets with
deterministic size caps \(R_{m,i}^{(h)}\le r_{m,i}^{(h)}\), the active-prefix
half-budget condition is implied by
\[
  \max_{1\le h\le J_m}
  \frac1{D_m}\sum_i(a_{m,i,h})^2r_{m,i}^{(h)}
  \le\frac{C_0n_m}{2\ell_m}
\]
in the coordinate case, or by
\[
  \max_{1\le h\le J_m}
  \frac1{B_{m,h}}\sum_i(a_{m,i,h})^2r_{m,i}^{(h)}
  \le\frac{C_0n_m}{2\ell_m}
\]
in the block case.
\end{corollary}

\begin{proof}
Apply Corollary~\ref{cor:adaptive-envelope-menu-schedules} to the finite active
menu \(\{1,\ldots,J_m\}\) in the \(m\)-th row.  The displayed prefix
half-budget and structural summability hypotheses are exactly the finite-menu
hypotheses for those active prefixes, so Borel--Cantelli gives the simultaneous
active-prefix conclusion.  Since \(J_m\to\infty\), each fixed \(h\) is in the
active prefix for all sufficiently large \(m\).  The sampled-sign statements
are the finite-menu sampled statements summed over \(1\le h\le J_m\), or the
common-batch union bound \(J_m2^{-M_m}\).

The arbitrary-growth estimate is
\[
  L_m^{(J)}\exp\left\{-\frac{3C_0n_m}{16\ell_m}\right\}
  \le e^{-1}(e+m)^{-(1+\gamma)}.
\]
The common sampled-batch condition gives
\(J_m2^{-M_m}\le e^{-1}(e+m)^{-(1+\delta)}\).  The all-integer and sparse
dimension cases are the same estimates with
\(L_m^{(J)}\le(eD_m)^\Gamma\).  The uniform-size budgets are the deterministic
uniform-support specialization in
Corollary~\ref{cor:adaptive-envelope-weighted-support-freedman}.
\end{proof}

\begin{corollary}[Almost-sure random-profile support arrays]
\label{cor:random-profile-weighted-support-borel}
Fix \(C_0\ge1\).  For each \(m\ge1\), let \(D_m\ge2\), put
\(\ell_m=\log(eD_m)\), and let
\(A_{m,1},\ldots,A_{m,n_m}\) be random self-adjoint matrices in
\(\M_{D_m}(\C)\).  Suppose that, for each row \(m\), one of the following two
random-profile models holds.
\begin{enumerate}[label=\textup{(\alph*)}]
\item Coordinate profiles.  The triples
  \((\rho_{m,i},Q_{m,i},S_{m,i})\), \(i=1,\ldots,n_m\), are independent,
  \(0\le\rho_{m,i}\le1\),
  \(Q_{m,i}=(Q_{m,ij})_{j=1}^{D_m}\in[0,1]^{D_m}\), and
  \(S_{m,i}\subset\{1,\ldots,D_m\}\), with
  \[
    \mathbb P\{j\in S_{m,i}\mid \rho_{m,i},Q_{m,i}\}\le Q_{m,ij}
    \qquad\text{a.s.}
  \]
  Conditional on the amplitudes, profiles, and supports,
  \[
    A_{m,i}=P_{S_{m,i}}A_{m,i}P_{S_{m,i}},
    \qquad
    \|A_{m,i}\|_{\op}\le\rho_{m,i}.
  \]
  Put
  \[
    N_m=D_m,\qquad
    \mu_m=\max_{j\le D_m}\sum_i\mathbb E(\rho_{m,i}^2Q_{m,ij}),
  \]
  and
  \[
    \sigma_m^2=\max_{j\le D_m}\sum_i
      \mathbb E(\rho_{m,i}^4Q_{m,ij}).
  \]
\item Block profiles.  There are pairwise orthogonal
  \(\Pi_{m,\alpha}\in\M_{D_m}(\C)\), \(1\le\alpha\le B_m\), with
  \(\sum_{\alpha=1}^{B_m}\Pi_{m,\alpha}=I\).  The triples
  \((\rho_{m,i},Q_{m,i},T_{m,i})\), \(i=1,\ldots,n_m\), are independent,
  \(0\le\rho_{m,i}\le1\),
  \(Q_{m,i}=(Q_{m,i\alpha})_{\alpha=1}^{B_m}\in[0,1]^{B_m}\), and
  \(T_{m,i}\subset\{1,\ldots,B_m\}\), with
  \[
    \mathbb P\{\alpha\in T_{m,i}\mid \rho_{m,i},Q_{m,i}\}\le Q_{m,i\alpha}
    \qquad\text{a.s.}
  \]
  Put \(P_{m,T_{m,i}}=\sum_{\alpha\in T_{m,i}}\Pi_{m,\alpha}\).  Conditional
  on the amplitudes, profiles, and selected blocks,
  \[
    A_{m,i}=P_{m,T_{m,i}}A_{m,i}P_{m,T_{m,i}},
    \qquad
    \|A_{m,i}\|_{\op}\le\rho_{m,i}.
  \]
  Put
  \[
    N_m=B_m,\qquad
    \mu_m=\max_{\alpha\le B_m}\sum_i
      \mathbb E(\rho_{m,i}^2Q_{m,i\alpha}),
  \]
  and
  \[
    \sigma_m^2=\max_{\alpha\le B_m}\sum_i
      \mathbb E(\rho_{m,i}^4Q_{m,i\alpha}).
  \]
\end{enumerate}
Assume that there are \(t_m>0\) such that
\[
  \mu_m+t_m\le \frac{C_0n_m}{\ell_m}
\]
and
\[
  \sum_m
  N_m\exp\left\{-\frac{t_m^2}{2(\sigma_m^2+t_m/3)}\right\}<\infty .
\]
Then, with probability one, for all sufficiently large \(m\),
\[
  \disc(A_{m,1},\ldots,A_{m,n_m})\le C(C_0)\sqrt{n_m}.
\]
In particular, under the half-budget condition
\[
  \mu_m\le \frac{C_0n_m}{2\ell_m},
\]
the same conclusion holds whenever
\[
  \sum_m
  N_m\exp\left\{-\frac{3C_0n_m}{16\ell_m}\right\}<\infty .
\]
This summability is implied, for any \(\gamma>0\), by the direct growth rule
\[
  n_m\ge \frac{16}{3C_0}\,\ell_m
  \bigl(\log(eN_m)+(1+\gamma)\log(e+m)\bigr)
\]
eventually.  If, in addition, each row is supplied with \(M_m\) independent
uniform sign vectors in \(\{\pm1\}^{n_m}\), independent of the matrix array, and
\(\sum_m2^{-M_m}<\infty\), then after enlarging \(C(C_0)\) if necessary, almost
surely for all sufficiently large \(m\) one sampled sign certifies the same
\(C(C_0)\sqrt{n_m}\) bound.
\end{corollary}

\begin{proof}
Corollary~\ref{cor:random-profile-weighted-support-bernstein} applied to the
\(m\)-th row gives the displayed one-row failure probability in either model.
The summability assumption and the first Borel--Cantelli lemma give the
almost-sure discrepancy conclusion.  Under the half-budget assumption take
\(t_m=C_0n_m/(2\ell_m)\); since \(\rho_{m,i}^4\le\rho_{m,i}^2\), one has
\(\sigma_m^2\le\mu_m\le t_m\), and the exponent is at least
\(3C_0n_m/(16\ell_m)\).

The direct growth rule gives
\[
  N_m\exp\left\{-\frac{3C_0n_m}{16\ell_m}\right\}
  \le e^{-1}(e+m)^{-(1+\gamma)},
\]
which is summable.  For sampled signs, condition on the realized row and on the
weighted-overlap event used in the proof of
Corollary~\ref{cor:random-profile-weighted-support-bernstein}.  On that event,
\(\|\sum_iA_{m,i}^2\|_{\op}\le C_0n_m/\ell_m\), so
Corollary~\ref{cor:random-sign-search-square-profile} gives conditional
sampled-search failure at most \(2^{-M_m}\), after enlarging \(C(C_0)\).  Another
application of Borel--Cantelli proves the sampled-certificate statement.
\end{proof}

\begin{corollary}[Almost-sure weighted support arrays]
\label{cor:random-weighted-support-borel}
Fix \(C_0\ge1\).  For each \(m\ge1\), let
\(A_{m,1},\ldots,A_{m,n_m}\) be random self-adjoint matrices in
\(\M_{D_m}(\C)\), with deterministic amplitudes
\(\rho_{m,i}\in[0,1]\).  Suppose that, for each \(m\), one of the following two
models holds.
\begin{enumerate}[label=\textup{(\alph*)}]
\item Coordinate supports.  There are random sets
  \(S_{m,i}\subset\{1,\ldots,D_m\}\) such that, for every coordinate \(j\), the
  indicators \(\mathbf 1_{\{j\in S_{m,i}\}}\), \(i=1,\ldots,n_m\), are
  independent and satisfy
  \[
    \mathbb P\{j\in S_{m,i}\}\le p_{m,ij}.
  \]
  Conditional on the supports, \(A_{m,i}=P_{S_{m,i}}A_{m,i}P_{S_{m,i}}\) and
  \(\|A_{m,i}\|_{\op}\le\rho_{m,i}\).  Put
  \[
    b_m=\max_i\rho_{m,i}^2,\qquad
    \mu_m=\max_j\sum_i\rho_{m,i}^2p_{m,ij},\qquad
    \sigma_m^2=\max_j\sum_i\rho_{m,i}^4p_{m,ij}.
  \]
  Assume that there are \(t_m>0\) with
  \[
    \mu_m+t_m\le \frac{C_0n_m}{\log(eD_m)}
  \]
  and
  \[
    \sum_m
    D_m\exp\left\{-\frac{t_m^2}{2(\sigma_m^2+b_mt_m/3)}\right\}<\infty,
  \]
  with the zero-denominator convention of
  Corollary~\ref{cor:random-coordinate-support-bernstein}.
\item Random subspaces.  There are independent random orthogonal projections
  \(P_{m,1},\ldots,P_{m,n_m}\) such that, conditional on the projections,
  \(A_{m,i}=P_{m,i}A_{m,i}P_{m,i}\) and
  \(\|A_{m,i}\|_{\op}\le\rho_{m,i}\).  Put
  \[
    b_m=\max_i\rho_{m,i}^2,\qquad
    \mu_m=\left\|\sum_i\rho_{m,i}^2\mathbb E P_{m,i}\right\|_{\op}.
  \]
  Assume that there are \(t_m>0\) with
  \[
    \mu_m+t_m\le \frac{C_0n_m}{\log(eD_m)}
  \]
  and
  \[
    \sum_m
    D_m\exp\left\{-\frac{t_m^2}{2(\mu_m+b_mt_m/3)}\right\}<\infty.
  \]
\end{enumerate}
Then, with probability one, for all sufficiently large \(m\),
\[
  \disc(A_{m,1},\ldots,A_{m,n_m})\le C(C_0)\sqrt{n_m}.
\]
In particular, the almost-sure conclusion holds in the homogeneous coordinate
case \(p_{m,ij}\le p_m\) whenever
\[
  p_m\sum_i\rho_{m,i}^2\le \frac{C_0n_m}{2\log(eD_m)}
  \qquad\text{and}\qquad
  \sum_m D_m\exp\left\{-\frac{3C_0n_m}{16\log(eD_m)}\right\}<\infty,
\]
and it holds for Haar rank-\(r_{m,i}\) projections under the same summability
condition whenever
\[
  \frac1{D_m}\sum_i\rho_{m,i}^2r_{m,i}
  \le \frac{C_0n_m}{2\log(eD_m)}.
\]
\end{corollary}

\begin{proof}
In the coordinate case, apply
Corollary~\ref{cor:random-coordinate-support-bernstein} to the \(m\)-th row.
In the random-subspace case, apply
Corollary~\ref{cor:random-subspace-matrix-bernstein} to the \(m\)-th row.  The
displayed summability assumptions make the resulting failure probabilities
summable, so the first Borel--Cantelli lemma gives the eventual almost-sure
bound.  The homogeneous coordinate and Haar specializations are exactly the
``in particular'' estimates in those two finite-dimensional corollaries, with
\(t_m=C_0n_m/(2\log(eD_m))\).
\end{proof}

\begin{corollary}[All-dimension weighted random-support schedules]
\label{cor:random-weighted-support-all-dim-schedules}
Fix \(C_0\ge1\).  For each row \(m\), let
\(\ell_m=\log(eD_m)\), let \(D_m\ge2\), and let
\(A_{m,1},\ldots,A_{m,n_m}\) be random self-adjoint matrices in
\(\M_{D_m}(\C)\) with deterministic amplitudes
\(\rho_{m,i}\in[0,1]\).  Assume that one of the following two half-budget
models holds for all large \(m\).
\begin{enumerate}[label=\textup{(\alph*)}]
\item Coordinate supports.  There are random supports
  \(S_{m,i}\subset\{1,\ldots,D_m\}\) such that, for every coordinate \(j\), the
  indicators \(\mathbf 1_{\{j\in S_{m,i}\}}\), \(i=1,\ldots,n_m\), are
  independent and satisfy
  \[
    \mathbb P\{j\in S_{m,i}\}\le p_{m,ij},
    \qquad
    \max_j\sum_{i=1}^{n_m}\rho_{m,i}^2p_{m,ij}
    \le \frac{C_0n_m}{2\ell_m}.
  \]
  Conditional on the supports,
  \(A_{m,i}=P_{S_{m,i}}A_{m,i}P_{S_{m,i}}\) and
  \(\|A_{m,i}\|_{\op}\le\rho_{m,i}\).
\item Random subspaces.  There are independent random orthogonal projections
  \(P_{m,1},\ldots,P_{m,n_m}\) such that, conditional on the projections,
  \(A_{m,i}=P_{m,i}A_{m,i}P_{m,i}\),
  \(\|A_{m,i}\|_{\op}\le\rho_{m,i}\), and
  \[
    \left\|\sum_{i=1}^{n_m}\rho_{m,i}^2\mathbb EP_{m,i}\right\|_{\op}
    \le \frac{C_0n_m}{2\ell_m}.
  \]
\end{enumerate}
Then
\[
  \mathbb P\left\{
  \disc(A_{m,1},\ldots,A_{m,n_m})>C(C_0)\sqrt{n_m}
  \right\}
  \le
  D_m\exp\left\{-\frac{3C_0n_m}{16\ell_m}\right\}
\]
for all large \(m\).  Consequently, if
\[
  \sum_mD_m\exp\left\{-\frac{3C_0n_m}{16\log(eD_m)}\right\}<\infty,
\]
then, almost surely for all sufficiently large \(m\),
\[
  \disc(A_{m,1},\ldots,A_{m,n_m})\le C(C_0)\sqrt{n_m}.
\]

In particular, along all integer dimensions \(D=2,3,\ldots\), the almost-sure
conclusion holds whenever the half-budget coordinate or subspace hypothesis
above holds with \(D_m=D\), \(n_m=n_D\), and
\[
  n_D\ge a\log^2(eD)
\]
for some \(a>32/(3C_0)\).  Along any sparse schedule
\(D_m\ge\exp(c m^\kappa)\), it is enough that
\[
  n_m\ge a\log^2(eD_m)
\]
eventually for some \(a>16/(3C_0)\).

For uniformly chosen fixed-size coordinate supports of sizes \(r_{m,i}\), or
for Haar-distributed projections of ranks \(r_{m,i}\), the half-budget
condition is
\[
  \frac1{D_m}\sum_i\rho_{m,i}^2r_{m,i}
  \le \frac{C_0n_m}{2\log(eD_m)}.
\]
Thus the unweighted homogeneous coordinate case is covered by
\[
  p_m\le \frac{C_0}{2\log(eD_m)},
\]
and the unweighted Haar-rank case is covered by
\[
  \frac1{n_m}\sum_i r_{m,i}\le \frac{C_0D_m}{2\log(eD_m)},
\]
with the same all-integer and sparse almost-sure schedules.
\end{corollary}

\begin{proof}
For the coordinate model, apply
Corollary~\ref{cor:random-coordinate-support-bernstein} with
\[
  t_m=\frac{C_0n_m}{2\ell_m}.
\]
The half-budget hypothesis gives \(\mu_m\le t_m\), and because
\(\rho_{m,i}\le1\),
\[
  \sigma_m^2
  =\max_j\sum_i\rho_{m,i}^4p_{m,ij}
  \le \max_j\sum_i\rho_{m,i}^2p_{m,ij}
  \le t_m,
  \qquad b_m\le1.
\]
Hence
\[
  \frac{t_m^2}{2(\sigma_m^2+b_mt_m/3)}
  \ge
  \frac{t_m^2}{2(t_m+t_m/3)}
  =
  \frac{3C_0n_m}{16\ell_m}.
\]
The random-subspace model is identical, using
Corollary~\ref{cor:random-subspace-matrix-bernstein}; its variance parameter is
\(\mu_m\le t_m\), and again \(b_m\le1\).  This proves the displayed one-row
tail, and the first Borel--Cantelli lemma gives the summable-array conclusion.

If \(n_D\ge a\log^2(eD)\), then
\[
  D\exp\left\{-\frac{3C_0n_D}{16\log(eD)}\right\}
  \le
  (eD)^{1-3C_0a/16}.
\]
The series over all integer \(D\ge2\) converges when
\(3C_0a/16>2\).  Along \(D_m\ge\exp(c m^\kappa)\), the same bound is
stretched-exponentially summable as soon as \(3C_0a/16>1\).  The fixed-size and
Haar-rank specializations are the corresponding expectations
\(\mathbb P\{j\in S_{m,i}\}=r_{m,i}/D_m\) and
\(\mathbb EP_{m,i}=(r_{m,i}/D_m)I\).
\end{proof}

\begin{corollary}[Sampled sign certificates for weighted support ensembles]
\label{cor:random-weighted-support-search}
Fix \(C_0\ge1\), and let \(C(C_0)\) be sufficiently large.  Let
\(\eta^{(1)},\ldots,\eta^{(M)}\) be independent uniform sign vectors in
\(\{\pm1\}^n\), independent of the random matrices below.

\begin{enumerate}[label=\textup{(\alph*)}]
\item If \(A_1,\ldots,A_n\) satisfy the hypotheses of
  Corollary~\ref{cor:random-low-overlap-support} with exceptional event \(E\),
  then
  \[
    \mathbb P\left\{
    \min_{1\le r\le M}
    \left\|\sum_{i=1}^n\eta_i^{(r)}A_i\right\|_{\op}
    >C(C_0)\sqrt n
    \right\}
    \le \mathbb P(E)+2^{-M}.
  \]
\item Under the hypotheses of
  Corollary~\ref{cor:random-coordinate-support-bernstein}, the same probability
  is bounded by
  \[
    D\exp\left\{-\frac{t^2}{2(\sigma^2+bt/3)}\right\}+2^{-M},
  \]
  with the same zero-denominator convention; under the hypotheses of
  Corollary~\ref{cor:random-subspace-matrix-bernstein}, it is bounded by
  \[
    D\exp\left\{-\frac{t^2}{2(\mu+bt/3)}\right\}+2^{-M}.
  \]
\item In the triangular-array setting of
  Corollary~\ref{cor:random-weighted-support-borel}, let
  \(\eta_m^{(1)},\ldots,\eta_m^{(M_m)}\) be independent uniform sign vectors in
  \(\{\pm1\}^{n_m}\), independent across \(m\) and independent of the matrix
  array.  If, in addition to the summability hypotheses in that corollary,
  \[
    \sum_m2^{-M_m}<\infty,
  \]
  then, with probability one, for all sufficiently large \(m\) there is
  \(1\le r\le M_m\) such that
  \[
    \left\|\sum_{i=1}^{n_m}\eta_{m,i}^{(r)}A_{m,i}\right\|_{\op}
    \le C(C_0)\sqrt{n_m}.
  \]
\end{enumerate}
Thus every weighted low-overlap coordinate, projection, frame, or random
subspace ensemble covered above has an explicit randomized certificate: a
constant number of independent sign trials gives a constant success probability,
and any choice of \(M_m\) with \(\sum_m2^{-M_m}<\infty\) makes the algorithmic
failure summable along triangular arrays.
\end{corollary}

\begin{proof}
On the weighted-overlap event in
Corollary~\ref{cor:random-low-overlap-support}, its proof shows
\[
  \left\|\sum_iA_i^2\right\|_{\op}\le \frac{C_0n}{\log(eD)}.
\]
Conditioning on the realized tuple and applying
Corollary~\ref{cor:random-sign-search-square-profile} gives conditional
failure at most \(2^{-M}\) for the sampled signs, after enlarging \(C(C_0)\).
This proves part (a).  Parts (b) and (c) are obtained by combining the same
conditional \(2^{-M}\) (or \(2^{-M_m}\)) sign-search failure with the matrix
failure bounds in Corollaries~\ref{cor:random-coordinate-support-bernstein},
\ref{cor:random-subspace-matrix-bernstein}, and
\ref{cor:random-weighted-support-borel}.  The array statement follows from the
first Borel--Cantelli lemma applied to the sum of the matrix-row failures and
the sign-search failures.
\end{proof}

\begin{corollary}[All-dimension sampled certificates for weighted random supports]
\label{cor:random-weighted-support-search-all-dim-schedules}
In the triangular-array setting of
Corollary~\ref{cor:random-weighted-support-all-dim-schedules}, let
\(\eta_m^{(1)},\ldots,\eta_m^{(M_m)}\) be independent uniform sign vectors in
\(\{\pm1\}^{n_m}\), independent across \(m\) and independent of the matrix
array.  Under either half-budget coordinate-support or random-subspace model in
that corollary, for all large \(m\),
\[
\begin{aligned}
  \mathbb P\left\{
  \min_{1\le r\le M_m}
  \left\|\sum_{i=1}^{n_m}\eta_{m,i}^{(r)}A_{m,i}\right\|_{\op}
  >C(C_0)\sqrt{n_m}
  \right\}
  &\le
  D_m\exp\left\{-\frac{3C_0n_m}{16\log(eD_m)}\right\}\\
  &\quad+2^{-M_m}.
\end{aligned}
\]
Consequently, if
\[
  \sum_m\left[
  D_m\exp\left\{-\frac{3C_0n_m}{16\log(eD_m)}\right\}
  +2^{-M_m}
  \right]<\infty,
\]
then, with probability one, for all sufficiently large \(m\), one of the
sampled sign vectors certifies
\[
  \left\|\sum_{i=1}^{n_m}\eta_{m,i}^{(r_m)}A_{m,i}\right\|_{\op}
  \le C(C_0)\sqrt{n_m}.
\]

In particular:
\begin{enumerate}[label=\textup{(\alph*)}]
\item Along all integer dimensions \(D=2,3,\ldots\), the sampled-certificate
  conclusion holds under the half-budget hypotheses of
  Corollary~\ref{cor:random-weighted-support-all-dim-schedules} if
  \(n_D\ge a\log^2(eD)\) for some \(a>32/(3C_0)\) and, for some
  \(\delta>0\),
  \[
    M_D\ge (1+\delta)\log_2(eD)
  \]
  eventually.
\item Along any sparse schedule \(D_m\ge\exp(c m^\kappa)\), the same conclusion
  holds if \(n_m\ge a\log^2(eD_m)\) eventually for some \(a>16/(3C_0)\) and,
  for some \(\delta>0\),
  \[
    M_m\ge (1+\delta)\log_2(e+m)
  \]
  eventually.
\end{enumerate}
Thus the existence statement for weighted coordinate and random-subspace
ensembles can be made into an eventual randomized certificate with only
logarithmically many independent sign trials per row.
\end{corollary}

\begin{proof}
Apply Corollary~\ref{cor:random-weighted-support-search} to the \(m\)-th row
with \(t_m=C_0n_m/(2\log(eD_m))\).  The half-budget hypotheses and the proof of
Corollary~\ref{cor:random-weighted-support-all-dim-schedules} give the matrix
failure term
\[
  D_m\exp\left\{-\frac{3C_0n_m}{16\log(eD_m)}\right\},
\]
and Corollary~\ref{cor:random-weighted-support-search} adds the independent
sign-search failure \(2^{-M_m}\).  The first Borel--Cantelli lemma gives the
summable-array conclusion.

For integer dimensions, the matrix tail is summable by
Corollary~\ref{cor:random-weighted-support-all-dim-schedules}.  The sign-search
tail satisfies
\[
  2^{-M_D}\le(eD)^{-(1+\delta)},
\]
which is summable over \(D\ge2\).  In the sparse case, the matrix tail is again
summable by Corollary~\ref{cor:random-weighted-support-all-dim-schedules}, and
\[
  2^{-M_m}\le(e+m)^{-(1+\delta)}
\]
is summable over \(m\).
\end{proof}

\begin{corollary}[Growth-rate weighted random-support schedules]
\label{cor:random-weighted-support-growth-schedules}
In the half-budget coordinate-support or random-subspace setting of
Corollary~\ref{cor:random-weighted-support-all-dim-schedules}, put
\[
  \ell_m=\log(eD_m),\qquad L_m=\log(e+m).
\]
Assume that, for some \(\gamma>0\),
\[
  n_m\ge \frac{16}{3C_0}\,\ell_m\bigl(\ell_m+(1+\gamma)L_m\bigr)
\]
eventually.  Then, almost surely for all sufficiently large \(m\),
\[
  \disc(A_{m,1},\ldots,A_{m,n_m})\le C(C_0)\sqrt{n_m}.
\]
Moreover, if
\(\eta_m^{(1)},\ldots,\eta_m^{(M_m)}\) are independent uniform sign vectors as
in Corollary~\ref{cor:random-weighted-support-search-all-dim-schedules} and,
for some \(\delta>0\),
\[
  M_m\ge(1+\delta)\log_2(e+m)
\]
eventually, then, almost surely for all sufficiently large \(m\), one sampled
sign vector certifies
\[
  \left\|\sum_{i=1}^{n_m}\eta_{m,i}^{(r_m)}A_{m,i}\right\|_{\op}
  \le C(C_0)\sqrt{n_m}.
\]
Thus the weighted coordinate-support and random-subspace models have eventual
Spencer signings, and eventual sampled certificates, for every dimension growth
regime after paying one logarithmic-in-row-index summability margin.
\end{corollary}

\begin{proof}
By Corollary~\ref{cor:random-weighted-support-all-dim-schedules}, the row
failure probability is at most
\[
  D_m\exp\left\{-\frac{3C_0n_m}{16\ell_m}\right\}.
\]
The growth hypothesis gives
\[
  \frac{3C_0n_m}{16\ell_m}\ge \ell_m+(1+\gamma)L_m,
\]
and hence
\[
  D_m\exp\left\{-\frac{3C_0n_m}{16\ell_m}\right\}
  \le e^{-1}(e+m)^{-(1+\gamma)}.
\]
This is summable, so the first Borel--Cantelli lemma proves the discrepancy
conclusion.  With sampled signs,
Corollary~\ref{cor:random-weighted-support-search-all-dim-schedules} adds the
failure probability \(2^{-M_m}\), and
\[
  2^{-M_m}\le(e+m)^{-(1+\delta)}
\]
is summable.  Another application of Borel--Cantelli gives the sampled
certificate conclusion.
\end{proof}

\begin{corollary}[Random-amplitude weighted-support schedules]
\label{cor:random-amplitude-weighted-support-schedules}
Fix \(C_0\ge1\).  For each row \(m\), put
\(\ell_m=\log(eD_m)\), let \(D_m\ge2\), and let
\(A_{m,1},\ldots,A_{m,n_m}\) be random self-adjoint matrices in
\(\M_{D_m}(\C)\).  Assume that, for all large \(m\), one of the following
random-amplitude half-budget models holds.
\begin{enumerate}[label=\textup{(\alph*)}]
\item Coordinate supports.  The pairs
  \((\rho_{m,i},S_{m,i})\), \(1\le i\le n_m\), are independent,
  \(0\le\rho_{m,i}\le1\),
  \(S_{m,i}\subset\{1,\ldots,D_m\}\), and for deterministic
  \(p_{m,ij}\in[0,1]\),
  \[
    \mathbb P\{j\in S_{m,i}\mid \rho_{m,i}\}\le p_{m,ij}
    \qquad\text{a.s.}
  \]
  Conditional on the amplitudes and supports,
  \(A_{m,i}=P_{S_{m,i}}A_{m,i}P_{S_{m,i}}\) and
  \(\|A_{m,i}\|_{\op}\le\rho_{m,i}\).  Suppose that
  \[
    \mu_{m,c}:=\max_j\sum_{i=1}^{n_m}
      p_{m,ij}\mathbb E\rho_{m,i}^2
    \le \frac{C_0n_m}{2\ell_m}.
  \]
\item Random subspaces.  The pairs
  \((\rho_{m,i},P_{m,i})\), \(1\le i\le n_m\), are independent,
  \(0\le\rho_{m,i}\le1\), and \(P_{m,i}\) is an orthogonal projection in
  \(\M_{D_m}(\C)\).  Conditional on the amplitudes and projections,
  \(A_{m,i}=P_{m,i}A_{m,i}P_{m,i}\) and
  \(\|A_{m,i}\|_{\op}\le\rho_{m,i}\).  Suppose that
  \[
    \mu_{m,s}:=
    \left\|\sum_{i=1}^{n_m}\mathbb E(\rho_{m,i}^2P_{m,i})\right\|_{\op}
    \le \frac{C_0n_m}{2\ell_m}.
  \]
\end{enumerate}
Then, for all large \(m\),
\[
  \mathbb P\left\{
  \disc(A_{m,1},\ldots,A_{m,n_m})>C(C_0)\sqrt{n_m}
  \right\}
  \le
  D_m\exp\left\{-\frac{3C_0n_m}{16\ell_m}\right\}.
\]
Consequently, if
\[
  \sum_mD_m\exp\left\{-\frac{3C_0n_m}{16\log(eD_m)}\right\}<\infty,
\]
then, almost surely for all sufficiently large \(m\),
\[
  \disc(A_{m,1},\ldots,A_{m,n_m})\le C(C_0)\sqrt{n_m}.
\]

Let \(\eta_m^{(1)},\ldots,\eta_m^{(M_m)}\) be independent uniform sign vectors
in \(\{\pm1\}^{n_m}\), independent across \(m\) and independent of the matrix
array.  After enlarging \(C(C_0)\) if necessary,
\[
\begin{aligned}
  \mathbb P\left\{
  \min_{1\le r\le M_m}
  \left\|\sum_{i=1}^{n_m}\eta_{m,i}^{(r)}A_{m,i}\right\|_{\op}
  >C(C_0)\sqrt{n_m}
  \right\}
  &\le
  D_m\exp\left\{-\frac{3C_0n_m}{16\ell_m}\right\}\\
  &\quad+2^{-M_m}.
\end{aligned}
\]
If the right-hand side is summable in \(m\), then almost surely for all
sufficiently large \(m\) one of the sampled signs certifies the same
\(C(C_0)\sqrt{n_m}\) bound.

In particular:
\begin{enumerate}[label=\textup{(\roman*)}]
\item Along all integer dimensions \(D=2,3,\ldots\), the almost-sure signing
  conclusion holds whenever the random-amplitude half-budget hypothesis holds
  with \(D_m=D\), \(n_m=n_D\), and
  \[
    n_D\ge a\log^2(eD)
  \]
  eventually for some \(a>32/(3C_0)\).  The sampled-certificate conclusion
  holds as well if, for some \(\delta>0\),
  \[
    M_D\ge(1+\delta)\log_2(eD)
  \]
  eventually.
\item Along any sparse schedule \(D_m\ge\exp(c m^\kappa)\), the same
  almost-sure signing conclusion holds whenever
  \[
    n_m\ge a\log^2(eD_m)
  \]
  eventually for some \(a>16/(3C_0)\).  The sampled-certificate conclusion
  holds as well if, for some \(\delta>0\),
  \[
    M_m\ge(1+\delta)\log_2(e+m)
  \]
  eventually.
\item For an arbitrary dimension schedule, if for some \(\gamma>0\)
  \[
    n_m\ge
    \frac{16}{3C_0}\,\ell_m\bigl(\ell_m+(1+\gamma)\log(e+m)\bigr)
  \]
  eventually, then the almost-sure signing conclusion holds; the same
  \(M_m\ge(1+\delta)\log_2(e+m)\) condition gives eventual sampled
  certificates.
\end{enumerate}
For uniformly chosen coordinate supports of random sizes \(R_{m,i}\), or
Haar-distributed projections of random ranks \(R_{m,i}\), the half-budget
condition in either model is
\[
  \frac1{D_m}\sum_i\mathbb E(\rho_{m,i}^2R_{m,i})
  \le \frac{C_0n_m}{2\log(eD_m)}.
\]
\end{corollary}

\begin{proof}
Apply Corollary~\ref{cor:random-amplitude-weighted-support-bernstein} to the
\(m\)-th row with \(t_m=C_0n_m/(2\ell_m)\).  In the coordinate model,
\[
  \max_j\sum_i p_{m,ij}\mathbb E\rho_{m,i}^4
  \le
  \max_j\sum_i p_{m,ij}\mathbb E\rho_{m,i}^2
  \le t_m,
\]
and in the subspace model \(\mu_{m,s}\le t_m\).  The half-budget estimate in
that corollary gives the displayed one-row matrix tail.  The first
Borel--Cantelli lemma gives the almost-sure signing conclusion under
summability.

For sampled signs, on the overlap event used in the proof of
Corollary~\ref{cor:random-amplitude-weighted-support-bernstein},
Corollary~\ref{cor:random-low-overlap-support} gives
\(\|\sum_iA_{m,i}^2\|_{\op}\le C_0n_m/\ell_m\).  Conditioning on the realized
tuple and applying Corollary~\ref{cor:random-sign-search-square-profile} gives
conditional sampled-sign failure at most \(2^{-M_m}\), after enlarging
\(C(C_0)\).  Combining this with the matrix-row failure proves the sampled
tail and the sampled Borel--Cantelli statement.

The integer and sparse schedule thresholds are exactly the summability
calculations in
Corollary~\ref{cor:random-weighted-support-all-dim-schedules}.  The growth-rate
condition gives
\[
  D_m\exp\left\{-\frac{3C_0n_m}{16\ell_m}\right\}
  \le e^{-1}(e+m)^{-(1+\gamma)},
\]
and \(2^{-M_m}\le(e+m)^{-(1+\delta)}\).  For uniform random support sizes and
Haar random ranks, the budget identity follows from
\(\mathbb P\{j\in S_{m,i}\mid R_{m,i},\rho_{m,i}\}=R_{m,i}/D_m\) and
\(\mathbb E(P_{m,i}\mid R_{m,i},\rho_{m,i})=(R_{m,i}/D_m)I\).
\end{proof}

\begin{corollary}[Block-count random-amplitude support schedules]
\label{cor:random-amplitude-block-support-schedules}
Fix \(C_0\ge1\).  For each row \(m\), put
\(\ell_m=\log(eD_m)\).  Let
\(\Pi_{m,1},\ldots,\Pi_{m,B_m}\) be pairwise orthogonal projections in
\(\M_{D_m}(\C)\) summing to \(I\).  Let
\((\rho_{m,i},T_{m,i})\), \(1\le i\le n_m\), be independent pairs, where
\(0\le\rho_{m,i}\le1\) and \(T_{m,i}\subset\{1,\ldots,B_m\}\).  Assume that
for deterministic \(q_{m,i\alpha}\in[0,1]\),
\[
  \mathbb P\{\alpha\in T_{m,i}\mid \rho_{m,i}\}\le q_{m,i\alpha}
  \qquad\text{a.s.}
\]
Put \(P_{m,T_{m,i}}=\sum_{\alpha\in T_{m,i}}\Pi_{m,\alpha}\), and suppose that,
conditional on the amplitudes and selected blocks,
\[
  A_{m,i}=P_{m,T_{m,i}}A_{m,i}P_{m,T_{m,i}},
  \qquad
  \|A_{m,i}\|_{\op}\le\rho_{m,i}.
\]
If the block half-budget condition
\[
  \mu_{m,B}:=\max_{\alpha\le B_m}\sum_{i=1}^{n_m}
  q_{m,i\alpha}\mathbb E\rho_{m,i}^2
  \le \frac{C_0n_m}{2\ell_m}
\]
holds for all large \(m\), then
\[
  \mathbb P\left\{
  \disc(A_{m,1},\ldots,A_{m,n_m})>C(C_0)\sqrt{n_m}
  \right\}
  \le
  B_m\exp\left\{-\frac{3C_0n_m}{16\ell_m}\right\}
\]
for all large \(m\).  Hence the Spencer-scale conclusion holds almost surely
eventually whenever
\[
  \sum_m B_m\exp\left\{-\frac{3C_0n_m}{16\log(eD_m)}\right\}<\infty.
\]

Moreover, if
\(\eta_m^{(1)},\ldots,\eta_m^{(M_m)}\) are independent uniform sign vectors in
\(\{\pm1\}^{n_m}\), independent across \(m\) and independent of the matrix
array, then after enlarging \(C(C_0)\) if necessary,
\[
\begin{aligned}
  \mathbb P\left\{
  \min_{1\le r\le M_m}
  \left\|\sum_{i=1}^{n_m}\eta_{m,i}^{(r)}A_{m,i}\right\|_{\op}
  >C(C_0)\sqrt{n_m}
  \right\}
  &\le
  B_m\exp\left\{-\frac{3C_0n_m}{16\ell_m}\right\}\\
  &\quad+2^{-M_m}.
\end{aligned}
\]
Thus sampled certificates are eventual almost surely whenever the displayed
right-hand side is summable.

Concrete summable forms are as follows.
\begin{enumerate}[label=\textup{(\alph*)}]
\item If the rows are indexed by all integer dimensions \(D=2,3,\ldots\), if
  \(B_D\le(eD)^\Gamma\) eventually, and if
  \[
    n_D\ge a\log^2(eD)
  \]
  eventually for some \(a>16(\Gamma+1)/(3C_0)\), then the almost-sure signing
  conclusion holds.  The sampled-certificate conclusion holds as well if
  \(M_D\ge(1+\delta)\log_2(eD)\) eventually for some \(\delta>0\).
\item If \(D_m\ge\exp(c m^\kappa)\) eventually, if
  \(B_m\le(eD_m)^\Gamma\) eventually, and if
  \[
    n_m\ge a\log^2(eD_m)
  \]
  eventually for some \(a>16\Gamma/(3C_0)\), then the almost-sure signing
  conclusion holds.  The sampled-certificate conclusion holds as well if
  \(M_m\ge(1+\delta)\log_2(e+m)\) eventually for some \(\delta>0\).
\item For an arbitrary dimension and block-count schedule, if for some
  \(\gamma>0\),
  \[
    n_m\ge
    \frac{16}{3C_0}\,\ell_m
    \bigl(\log(eB_m)+(1+\gamma)\log(e+m)\bigr)
  \]
  eventually, then the almost-sure signing conclusion holds.  The same
  \(M_m\ge(1+\delta)\log_2(e+m)\) condition gives eventual sampled
  certificates.
\end{enumerate}
For uniformly chosen block sets of random sizes \(R_{m,i}\), the half-budget
condition is
\[
  \frac1{B_m}\sum_i\mathbb E(\rho_{m,i}^2R_{m,i})
  \le \frac{C_0n_m}{2\log(eD_m)}.
\]
\end{corollary}

\begin{proof}
Apply Corollary~\ref{cor:random-amplitude-block-support-bernstein} to the
\(m\)-th row with \(t_m=C_0n_m/(2\ell_m)\).  Since
\(\mathbb E\rho_{m,i}^4\le\mathbb E\rho_{m,i}^2\), its half-budget estimate
gives the displayed block-count row failure.  The first Borel--Cantelli lemma
proves the summable-array conclusion.

For sampled signs, condition on the realized matrices and on the block-overlap
event in the proof of
Corollary~\ref{cor:random-amplitude-block-support-bernstein}.  There,
\(\|\sum_iA_{m,i}^2\|_{\op}\le C_0n_m/\ell_m\).  Corollary~\ref{cor:random-sign-search-square-profile}
therefore gives conditional sampled-sign failure at most \(2^{-M_m}\), after
enlarging \(C(C_0)\).  This proves the sampled tail and its Borel--Cantelli
consequence.

For all integer dimensions,
\[
  B_D\exp\left\{-\frac{3C_0n_D}{16\log(eD)}\right\}
  \le
  (eD)^{\Gamma-3C_0a/16},
\]
which is summable when \(3C_0a/16>\Gamma+1\).  Along
stretched-exponential sparse schedules the same bound is summable when
\(3C_0a/16>\Gamma\).  The arbitrary-growth condition gives
\[
  B_m\exp\left\{-\frac{3C_0n_m}{16\ell_m}\right\}
  \le e^{-1}(e+m)^{-(1+\gamma)}.
\]
The sampled-sign summability follows from the displayed lower bounds on
\(M_D\) or \(M_m\).  The uniform block-size identity follows from
\(\mathbb P\{\alpha\in T_{m,i}\mid R_{m,i},\rho_{m,i}\}=R_{m,i}/B_m\).
\end{proof}

\begin{remark}
These pathwise statements are not the full Matrix Spencer conjecture for
generic \(D\times D\) random matrix ensembles when \(D^2\gg n\), unless the
realization has one of the structural decompositions above or one uses an
average-case input such as Corollary~\ref{cor:random-gaussian-tail-full}, or a
variance-profile input such as
Proposition~\ref{prop:square-function-khintchine}. The entropy
obstruction below applies to full-unit-ball or full-span faithful finite-block
reductions; it does not rule out tuple-dependent lifts, variance-sensitive
arguments, or genuinely average-case discrepancy methods for random inputs.
Average-case matrix discrepancy for such ensembles is a separate line of work;
see, for example, Kunisky--Zhang \cite{KuniskyZhang}, Maillard
\cite{Maillard}, and the later combined work of Kunisky, Oertel, Wengiel and
Zhang \cite{KOWZ}.
\end{remark}

\section{Boundary of the finite-lift method}

The framework gives a precise limitation for exact finite block lifts.

\begin{proposition}[Entropy obstruction for the full matrix algebra]\label{prop:obstruct}
Suppose the full algebra $\M_D(\C)$ admits an exact finite block lift through
$\mcB=\bigoplus_\alpha\M_{r_\alpha}(\C)$ in the sense that there is a linear map
from $\mcB$ onto $\M_D(\C)$. Then
\[
\Ent(\mcB)
=\left(\max_\alpha r_\alpha\right)\sum_\alpha r_\alpha
\ge D^2 .
\]
Consequently, Theorem \ref{thm:main} cannot yield better than
$C\sqrt{n\log(e+D^2/n)}$ when applied through an exact finite block lift that is
onto the full algebra \(\M_D(\C)\).
\end{proposition}

\begin{proof}
Surjectivity onto the $D^2$-dimensional vector space $\M_D(\C)$ forces
\[
\dim\mcB=\sum_\alpha r_\alpha^2\ge D^2 .
\]
Since
\[
\sum_\alpha r_\alpha^2
\le \left(\max_\alpha r_\alpha\right)\sum_\alpha r_\alpha
=\Ent(\mcB),
\]
the claim follows.
\end{proof}

\begin{lemma}[Proper subspaces miss a unit vector at distance one]
\label{lem:proper-subspace-distance}
Let \(X\) be a finite-dimensional real or complex normed space and let
\(L\subsetneq X\) be a proper linear subspace. Then there is \(x\in X\) with
\(\norm{x}=1\) and
\[
  \operatorname{dist}(x,L):=\inf_{\ell\in L}\norm{x-\ell}=1.
\]
\end{lemma}

\begin{proof}
By Hahn--Banach, choose a nonzero continuous linear functional
\(\varphi\in X^*\) that vanishes on \(L\). Normalize it so that
\(\norm{\varphi}_{X^*}=1\). Since \(X\) is finite-dimensional, the compact unit
ball contains a point \(x\) with \(\norm{x}=1\) and \(\abs{\varphi(x)}=1\). For
every \(\ell\in L\),
\[
  \norm{x-\ell}\ge \abs{\varphi(x-\ell)}=\abs{\varphi(x)}=1.
\]
The opposite inequality follows from \(\ell=0\), so
\(\operatorname{dist}(x,L)=1\).
\end{proof}

\begin{proposition}[Robust entropy obstruction for approximate faithful lifts]
\label{prop:robustobstruct}
Let \(\mcB=\bigoplus_\alpha\M_{r_\alpha}(\C)\).

\begin{enumerate}[leftmargin=2em]
\item Suppose \(W\subseteq\mcB\) is a complex linear subspace and
\(\Pi:W\to\M_D(\C)\) is complex-linear. If for some fixed \(\delta<1\), every
complex contraction \(A\in\M_D(\C)\), \(\norm{A}_{\op}\le1\), is within
operator-norm distance \(\delta\) of \(\Pi(W)\), then
\[
  \Ent(\mcB)\ge D^2 .
\]
\item Suppose \(W_{\rm sa}\subseteq\mcB_{\rm sa}\) is a real linear subspace and
\(\Pi:W_{\rm sa}\to\M_D(\C)_{\rm sa}\) is real-linear. If for some fixed
\(\delta<1\), every self-adjoint contraction \(A=A^*\), \(\norm{A}_{\op}\le1\),
is within operator-norm distance \(\delta\) of \(\Pi(W_{\rm sa})\), then again
\[
  \Ent(\mcB)\ge D^2 .
\]
\end{enumerate}
Thus the obstruction in Proposition~\ref{prop:obstruct} is not merely an exact
surjectivity obstruction: any fixed-error faithful reduction of the full matrix
unit ball through finite block lifts already pays quadratic entropy.
\end{proposition}

\begin{proof}
For (1), set \(L=\Pi(W)\subseteq\M_D(\C)\). If \(L\) were a proper complex
subspace, Lemma~\ref{lem:proper-subspace-distance}, applied to
\(X=\M_D(\C)\) with the operator norm, would produce a contraction \(A\) whose
distance to \(L\) is \(1\), contradicting the assumed \(\delta<1\)
approximation. Hence \(L=\M_D(\C)\), so
\[
  \sum_\alpha r_\alpha^2=\dim_\C\mcB\ge \dim_\C W\ge \dim_\C L=D^2.
\]
Since \(\sum_\alpha r_\alpha^2\le\Ent(\mcB)\), the entropy bound follows.

The self-adjoint case is identical over the real vector spaces
\(X=\M_D(\C)_{\rm sa}\) and \(L=\Pi(W_{\rm sa})\). If \(L\) is proper, the lemma
gives a self-adjoint contraction at distance \(1\) from \(L\), contradicting
\(\delta<1\). Therefore \(L=\M_D(\C)_{\rm sa}\), and
\[
  \sum_\alpha r_\alpha^2=\dim_\R\mcB_{\rm sa}
  \ge \dim_\R W_{\rm sa}\ge \dim_\R\M_D(\C)_{\rm sa}=D^2.
\]
Again \(\Ent(\mcB)\ge\sum_\alpha r_\alpha^2\).
\end{proof}

\begin{corollary}[No cheap universal finite-lift reduction]
\label{cor:nouniversal}
Suppose a proposed reduction of the full \(D\times D\) Matrix Spencer problem
assigns, to every tuple of contractions \(A_1,\ldots,A_n\in\M_D(\C)\), a finite
block lift \(Y_i\in\mcB=\bigoplus_\alpha\M_{r_\alpha}(\C)\) and a linear map
\(\Pi\) as in Theorem~\ref{thm:main}, with constants independent of \(D\).
Then, applied to any tuple whose span is all of \(\M_D(\C)\), the reduction must
use
\[
\Ent(\mcB)\ge D^2 .
\]
The same conclusion holds for any fixed-error version in which, for the lifted
subspace \(W\subseteq\mcB\), the image \(\Pi(W)\) \(\delta\)-approximates the
operator-norm unit ball of \(\M_D(\C)\), or of its self-adjoint part, for some
fixed \(\delta<1\). Consequently, within this exact or fixed-error faithful
finite-lift route, no structured target family whose relevant finite-block
realizations all have entropy \(o(D^2)\) can linearly realize all \(D\times D\)
matrix contractions. In particular, same-scale diagonal, circulant, Toeplitz,
and bounded-block models are not universal for the full matrix unit ball. They
could serve as universal targets in this sense only after enlarging the target
until the associated finite-block entropy is at least quadratic in \(D\); thus
their existing structured Matrix Spencer bounds do not by themselves settle the
full square conjecture at the class/full-unit-ball level.
\end{corollary}

\begin{proof}
Let \(A_1,\ldots,A_n\) be any tuple whose complex span is \(\M_D(\C)\), and set
\(W=\Span\{Y_1,\ldots,Y_n\}\subseteq\mcB\). Since \(\Pi(Y_i)=A_i\), we have
\(\Pi(W)=\M_D(\C)\). Extending \(\Pi\) linearly from \(W\) to all of \(\mcB\)
gives a surjection \(\mcB\to\M_D(\C)\), so Proposition~\ref{prop:obstruct}
gives \(\Ent(\mcB)\ge D^2\). The same argument applies over real spans in the
self-adjoint exact setting. The fixed-error claim is precisely
Proposition~\ref{prop:robustobstruct}, applied to \(W\) or \(W_{\rm sa}\). The
final assertions are only about structured target families used through this
same exact or fixed-error faithful finite-lift mechanism.
\end{proof}

The preceding corollary suggests isolating the complexity parameter actually
seen by Theorem~\ref{thm:main}.

\begin{definition}[Finite-lift entropy of a tuple]\label{def:liftentropy}
Let \(A_1,\ldots,A_n\in\M_D(\C)\) and let \(K,\Lambda\ge1\). Define
\[
\LiftEnt_{K,\Lambda}(A_1,\ldots,A_n)
=\inf \Ent(\mcB),
\]
where the infimum runs over all finite block algebras
\(\mcB=\bigoplus_\alpha\M_{r_\alpha}(\C)\) for which there are elements
\(Y_i\in\mcB\) and a fixed linear map
\[
\Pi:\Span\{Y_1,\ldots,Y_n\}\longrightarrow\M_D(\C)
\]
satisfying
\[
\Pi(Y_i)=A_i,\qquad
\norm{Y_i}_{\mcB}\le K,\qquad
\norm{\Pi(Y)}_{\op}\le\Lambda\norm{Y}_{\mcB}
\quad(Y\in\Span\{Y_i\}).
\]
If no such lift exists, set
\(\LiftEnt_{K,\Lambda}(A_1,\ldots,A_n)=+\infty\).
\end{definition}

\begin{proposition}[Discrepancy in terms of lift entropy]
\label{prop:liftentropy}
For every tuple of operator-norm contractions and every \(K,\Lambda\ge1\),
\[
\discop(A_1,\ldots,A_n)
\le C\Lambda K
\sqrt{n\log\!\left(e+
\frac{\LiftEnt_{K,\Lambda}(A_1,\ldots,A_n)}{n}\right)}
\]
whenever the right-hand side is finite. Moreover:
\begin{enumerate}[leftmargin=2em]
\item \(\LiftEnt_{1,1}(A_1,\ldots,A_n)\le D^2\) for every tuple, by the
identity lift through \(\M_D(\C)\).
\item If \(\Span_\C\{A_1,\ldots,A_n\}=\M_D(\C)\), then
\(\LiftEnt_{K,\Lambda}(A_1,\ldots,A_n)\ge D^2\) for all fixed \(K,\Lambda\).
The same lower bound holds for self-adjoint tuples whose real span is
\(\M_D(\C)_{\rm sa}\) when one restricts to self-adjoint lifts and real-linear
maps.
\item If the tuple lies in a common finite-dimensional C*-algebra with reduced
entropy \(hR\), then \(\LiftEnt_{1,1}(A_1,\ldots,A_n)\le hR\). If all matrices
are supported on a common \(\rho\)-dimensional reducing subspace, then
\(\LiftEnt_{1,1}\le\rho^2\). For scalar Toeplitz or Hankel contractions,
\(\LiftEnt_{C,1}\le C D\), and for fixed-level banded multilevel block Toeplitz
or Hankel tuples it is bounded by the corresponding sampled-symbol entropy
\(O_{r,b}(\prod_j(2w_j+1))\).
\end{enumerate}
Thus \(\LiftEnt\) is small on precisely the structured classes handled by the
finite-lift theorem, while full-span generic matrix tuples have maximal
finite-lift entropy \(D^2\).
\end{proposition}

\begin{proof}
For any \(E>\LiftEnt_{K,\Lambda}(A_1,\ldots,A_n)\), choose a lift with
\(\Ent(\mcB)\le E\) and apply Theorem~\ref{thm:main}; then let
\(E\downarrow\LiftEnt_{K,\Lambda}\). Item (1) uses
\(\mcB=\M_D(\C)\), \(Y_i=A_i\), and \(\Pi=\operatorname{id}\). Item (2) is
Proposition~\ref{prop:obstruct} applied to the span of the lifted tuple; in the
self-adjoint formulation it is the real self-adjoint part of
Proposition~\ref{prop:robustobstruct} with exact image. Item (3) records the
lifts already proved in Theorem~\ref{thm:cstar}, the identity lift on the
common \(\rho\times\rho\) corner, Theorem~\ref{thm:toep} and
Corollary~\ref{cor:hankel}, and the sampled multilevel estimates of
Theorem~\ref{thm:blocktoepml} and Corollary~\ref{cor:anisotropic-bandml}.
\end{proof}

This explains why Toeplitz and Hankel classes are tractable by this method:
the relevant finite Fourier data has only $O(D)$ scalar degrees of freedom and
admits a bounded $L^\infty$ extension. It also explains why common block
diagonalization with bounded blocks is tractable: the reduced noncommutative
entropy is linear in $D$. Conversely, the full matrix algebra has entropy
$D^2$, and Proposition~\ref{prop:robustobstruct} shows that allowing a fixed
operator-norm approximation error below \(1\) for the full unit ball does not
lower that entropy floor.

\begin{question}[Bounded finite-entropy lift problem]
Classify structured operator systems $S_D\subset\M_D(\C)$ for which every finite
tuple in \(S_D\) has a common uniformly bounded linear lift on its span either
to a finite block algebra $\bigoplus_\alpha\M_{r_\alpha}(\C)$ with
\[
\left(\max_\alpha r_\alpha\right)\sum_\alpha r_\alpha=O(D),
\]
or to a matrix-valued multiplier model with an $O(D)$-size norming set and
bounded block size. Every such class satisfies Matrix Spencer in the regime
$D\le n$.
\end{question}

This remains a useful classification problem for nonalgebraic structured
operator systems, multiplier models, and compression/extension mechanisms.  For
the special case where the tuple is already contained in a finite-dimensional
C*-algebra, Theorem~\ref{thm:akbassra} is the sharper current input: it asks for
intrinsic dimension \(O(n)\), not finite-lift entropy \(O(n)\).

The Toeplitz, Hankel, one-level block Toeplitz, fixed-level multilevel block
Toeplitz, banded multilevel Toeplitz, finite abelian convolution,
bounded-irrep convolution, simultaneous diagonalization, bounded-block
diagonalization, and commutative association-scheme examples are instances of
this program. The full-unit-ball/full-span faithful finite-lift route would require either a
way to beat the entropy obstruction for generic matrix algebras or a method not
reducible to exact or fixed-error approximate faithful lifts of the full matrix
unit ball.

\section{Extension to AF and von Neumann algebras}
\label{sec:af}

The finite-lift principle is a statement about a single finite-dimensional
C*-algebra: Theorem~\ref{thm:cstar} signs contractions that lie in a common
$\mcA\simeq\bigoplus_\alpha(\M_{r_\alpha}\otimes I_{m_\alpha})$ at the rate set
by the reduced entropy $\Ent=hR$, $h=\max_\alpha r_\alpha$,
$R=\sum_\alpha r_\alpha$. It is natural to ask whether anything survives in the
inductive limit, that is, for contractions in an infinite-dimensional
approximately finite-dimensional (AF) C*-algebra, or more generally in a von
Neumann algebra. The answer is a clean reduction together with an equally clean
obstruction, and the two together delimit what the finite-stage/full-span lift
route proves and where it runs into the stated obstruction.

Throughout, $A$ is a separable unital AF C*-algebra: the norm closure
\[
A=\overline{\bigcup_{k\ge1}A_k}^{\,\|\cdot\|},
\qquad
A_1\subseteq A_2\subseteq\cdots,
\qquad
A_k=\bigoplus_{\alpha=1}^{m_k}\M_{r_{k,\alpha}}(\C),
\]
of an increasing tower of finite-dimensional C*-subalgebras with unital
inclusions, as encoded by a Bratteli diagram \cite{Bratteli,Davidson}. We write
\[
h_k=\max_{\alpha}r_{k,\alpha},
\qquad
R_k=\sum_{\alpha}r_{k,\alpha},
\qquad
\Ent(A_k)=h_kR_k,
\]
and
\[
d_k=\dim_\C A_k=\sum_\alpha r_{k,\alpha}^2
\]
for the entropy of the $k$th stage, computed by Theorem~\ref{thm:cstar} from the
reduced algebra $\bigoplus_\alpha\M_{r_{k,\alpha}}(\C)$; multiplicities in the
embedding $A_k\hookrightarrow A$ do not enter, exactly as in
Theorem~\ref{thm:cstar}. The UHF algebra $M_{2^\infty}=\varinjlim\M_{2^k}(\C)$,
with $A_k\simeq\M_{2^k}(\C)$ a single block of size $2^k$, is the running
example; there $h_k=R_k=2^k$ and $\Ent(A_k)=d_k=4^k$.  Always
\(d_k\le\Ent(A_k)\), and the inequality can be strict for heterogeneous direct
sums.  The elementary finite-lift theorem sees \(\Ent(A_k)\); the Akbas--Sra
theorem sees \(d_k\).

The point of this section is sixfold.
\begin{enumerate}[leftmargin=2em]
\item For contractions that already lie in a fixed stage $A_k$, the bound is
immediate from Theorem~\ref{thm:cstar}: the supporting stage controls the
discrepancy (Lemma~\ref{lem:fixedstage}).
\item The bound is \emph{stable} under the inductive limit: contractions only
$\eps$-approximable by a common finite stage obey the same bound up to an
$O(n\eps)$ additive error, which vanishes as $\eps\to0$
(Lemma~\ref{lem:trunc}, Theorem~\ref{thm:af}).
\item Whether this is a $\sqrt n$ bound or a weaker bound is decided entirely by
the entropy of the \emph{supporting} stage, which is small only for special
classes of elements (\S\ref{subsec:linear}); non-additivity of entropy along the
tower (Remark~\ref{rem:nonadd}) prevents a stage-uniform bound from this entropy
accounting without restrictions on the supporting stage.
\item Using the external theorem of Akbas--Sra, exact or approximate support on a
stage with \(d_k=O(n)\) already gives the Spencer-scale bound, even if the
finite-lift entropy \(h_kR_k\) is larger.
\item The perturbative theorem of Akbas--Sra also survives the AF reduction in
finite-dimensional representations: low-Frobenius self-adjoint tails may be
added to a dimension-controlled AF-stage core (Corollary~\ref{cor:afdim-pert}).
\item For diffuse elements with no finite-dimensional supporting subalgebra, the
exact finite-stage route gives no exact-stage statement
(Proposition~\ref{prop:vnobstruct}); moreover any approximate finite-stage route
must pay at least the spectral packing number of the element
(Proposition~\ref{prop:spectralfloor}).
\end{enumerate}

\subsection{Reduction to the supporting stage}

The first statement requires no new idea; it is recorded to fix the form of the
bound.

\begin{lemma}[Fixed-stage bound]\label{lem:fixedstage}
Let $A_1,\ldots,A_n$ be contractions, $\|A_i\|_{\op}\le1$, all lying in a common
finite stage $A_k$ of the tower. Then
\[
\discop(A_1,\ldots,A_n)
\le C\sqrt{n\log\!\left(e+\frac{\Ent(A_k)}{n}\right)},
\]
with $C$ the universal constant of Theorem~\ref{thm:cstar}; the self-adjoint
version holds for self-adjoint $A_i$. The discrepancy is computed in $A$, and the
signs achieving it are the same as those achieving it inside $A_k$.
\end{lemma}

\begin{proof}
$A_k$ is a finite-dimensional C*-algebra; apply Theorem~\ref{thm:cstar} to its
reduced form $\bigoplus_\alpha\M_{r_{k,\alpha}}(\C)$, of entropy
$h_kR_k=\Ent(A_k)$. The inclusion $A_k\hookrightarrow A$ is an isometric
$*$-homomorphism, so the operator norm of any fixed signed sum
$\sum_i\eps_iA_i$ is the same whether computed in $A_k$ or in $A$. Hence the
minimum over signs agrees as well.
\end{proof}

The content is the word ``common.'' A signing is a single choice of
$\eps\in\{-1,1\}^n$ that must work for all $i$ simultaneously, so the relevant
object is the smallest stage containing the \emph{whole tuple}, and its entropy
is the cost.

\subsection{A truncation lemma}

To pass from the algebraic union to its norm closure we need a finite-stage
approximant that is itself a contraction. Norm-density gives an approximant in
some $A_k$, but it need not be a contraction; the following clamp fixes this
with no loss, using only the spectral mapping theorem. We avoid any
operator-Lipschitz inequality: the naive bound ``$f$ is $1$-Lipschitz on $\R$,
hence $\|f(x)-f(y)\|\le\|x-y\|$'' is \emph{false} for non-commuting
self-adjoints (the absolute value is the standard counterexample), so we use a
triangle-inequality argument that costs only a harmless factor.

\begin{lemma}[Contraction approximant]\label{lem:trunc}
Let $a=a^*\in A$ with $\|a\|_{\op}\le1$, and suppose $b_0=b_0^*\in A_k$ satisfies
$\|a-b_0\|_{\op}<\eps$. Let $f(t)=\max(-1,\min(1,t))$ and put $b=f(b_0)$, by
continuous functional calculus. Then
\[
b=b^*\in A_k,
\qquad
\|b\|_{\op}\le1,
\qquad
\|a-b\|_{\op}<2\eps .
\]
\end{lemma}

\begin{proof}
Since $b_0=b_0^*\in A_k$ and $A_k$ is a C*-subalgebra, the continuous functional
calculus of $b_0$ stays in $A_k$, so $b=f(b_0)\in A_k$ and $b=b^*$. By the
spectral mapping theorem $\mathrm{sp}(b)=f(\mathrm{sp}(b_0))\subseteq[-1,1]$,
hence $\|b\|_{\op}\le1$. For the error, set $g(t)=t-f(t)$, so that
$g\equiv0$ on $[-1,1]$ and $|g(t)|=\operatorname{dist}(t,[-1,1])$ for all $t$.
Then $b_0-b=g(b_0)$ and
\[
\|b_0-b\|_{\op}
=\max_{t\in\mathrm{sp}(b_0)}\operatorname{dist}\bigl(t,[-1,1]\bigr).
\]
For each $t\in\mathrm{sp}(b_0)$, self-adjoint spectral continuity (Weyl) gives
$\operatorname{dist}(t,\mathrm{sp}(a))\le\|a-b_0\|<\eps$, and
$\mathrm{sp}(a)\subseteq[-1,1]$ because $a$ is a self-adjoint contraction, so
$\operatorname{dist}(t,[-1,1])\le\operatorname{dist}(t,\mathrm{sp}(a))<\eps$.
Thus $\|b_0-b\|<\eps$, and
$\|a-b\|\le\|a-b_0\|+\|b_0-b\|<2\eps$.
\end{proof}

\begin{remark}
The factor $2$ is immaterial: one applies Lemma~\ref{lem:trunc} with $\eps/2$ in
place of $\eps$, or simply absorbs it, since in
Theorem~\ref{thm:af} the approximation parameter is sent to $0$. For non-self-adjoint
$a$ one truncates the self-adjoint and anti-self-adjoint parts
$\operatorname{Re}a=\tfrac12(a+a^*)$ and $\operatorname{Im}a=\tfrac1{2i}(a-a^*)$
separately and recombines, losing only a further universal constant; this is the
same realification already used in the proof of Theorem~\ref{thm:main}. We state
the self-adjoint case and note the general case in Theorem~\ref{thm:af}.
\end{remark}

\subsection{The AF stability theorem}

We now combine Lemma~\ref{lem:fixedstage} with the truncation. The hypothesis is
that the tuple is jointly $\eps$-close to a single stage. This is exactly the
local characterization of AF algebras: for a separable AF algebra, every finite
subset and every $\eps$ are captured by some finite-dimensional C*-subalgebra,
and that subalgebra may be taken inside the tower
\cite{Bratteli,Davidson}, so the hypothesis can always be met for some
$k=k(\eps)$.

\begin{theorem}[AF stability]\label{thm:af}
Let $A=\overline{\bigcup_kA_k}$ be a separable unital AF C*-algebra and let
$A_1,\ldots,A_n\in A$ be contractions. Fix $\eps>0$ and a stage index
$k$ for which there are contractions $B_i\in A_k$ with
$\|A_i-B_i\|_{\op}<\eps$ for all $i$ (in the self-adjoint case the $B_i$ may be
taken self-adjoint, by Lemma~\ref{lem:trunc}). Then
\[
\discop(A_1,\ldots,A_n)
\le C\sqrt{n\log\!\left(e+\frac{\Ent(A_k)}{n}\right)}
+2n\eps .
\]
Consequently
\[
\discop(A_1,\ldots,A_n)
\le \inf_{k\ge1}\left\{
C\sqrt{n\log\!\left(e+\frac{\Ent(A_k)}{n}\right)}
+2n\,\delta_k(A_1,\ldots,A_n)
\right\},
\]
where
$\delta_k(A_1,\ldots,A_n)=\max_i\operatorname{dist}_{\op}(A_i,A_k^{\le1})$ is the
joint distance of the tuple to the contractions of the $k$th stage, and
$\delta_k\to0$ as $k\to\infty$. The same statements hold for $\disc$ in the
self-adjoint case.
\end{theorem}

\begin{proof}
By Lemma~\ref{lem:fixedstage} there are signs $\eps^{\star}\in\{-1,1\}^n$ with
\[
\Bigl\|\sum_i\eps^{\star}_iB_i\Bigr\|_{\op}
\le C\sqrt{n\log\!\left(e+\frac{\Ent(A_k)}{n}\right)} .
\]
For these same signs, the triangle inequality and $\|A_i-B_i\|<\eps$ give
\[
\Bigl\|\sum_i\eps^{\star}_iA_i\Bigr\|_{\op}
\le\Bigl\|\sum_i\eps^{\star}_iB_i\Bigr\|_{\op}
+\sum_{i=1}^n\|A_i-B_i\|_{\op}
< C\sqrt{n\log\!\left(e+\frac{\Ent(A_k)}{n}\right)}+n\eps .
\]
Since the $B_i$ produced by Lemma~\ref{lem:trunc} from an
$(\eps/2)$-approximant satisfy $\|A_i-B_i\|<\eps$, the displayed bound with the
stated $2n\eps$ (from running Lemma~\ref{lem:trunc} at parameter $\eps$, giving
$\|A_i-B_i\|<2\eps$) follows; either normalization is the claim after relabelling
$\eps$. Taking the infimum over $k$, and over the choice of $B_i$ realizing
$\delta_k$ up to an arbitrarily small slack, gives the second display. Finally
$\delta_k\to0$ is the local characterization of AF algebras: the union
$\bigcup_kA_k$ is norm-dense, so each $A_i$ is approximated to within any $\eps$
by some stage, and after Lemma~\ref{lem:trunc} the approximant is a contraction;
taking the largest stage among the $n$ approximants makes the approximation
joint. The general (non-self-adjoint) case truncates real and imaginary parts
separately and absorbs the resulting constant into $C$.
\end{proof}

The estimate is genuinely an infimum over a \emph{trade-off}. Lowering the
approximation error $\eps=\delta_k$ forces a higher stage $k$, hence a larger
$\Ent(A_k)$; the optimal $k$ balances the additive sampling-type error
$2n\delta_k$ against the entropy term. Two regimes make this concrete.

\begin{corollary}[Elements of finite type]\label{cor:finitetype}
If the tuple lies exactly in some stage, $\{A_1,\ldots,A_n\}\subseteq A_k$ with
$\delta_k=0$, then the additive error vanishes and
\[
\discop(A_1,\ldots,A_n)
\le C\sqrt{n\log\!\left(e+\frac{\Ent(A_k)}{n}\right)} .
\]
In particular, if moreover $\Ent(A_k)=O(n)$, the tuple satisfies the conjectured
Matrix Spencer bound $\discop=O(\sqrt n)$, with a constant independent of every
later stage and of the ambient (infinite) algebra $A$.
\end{corollary}

\begin{corollary}[Quantitative limit elements]\label{cor:limit}
Suppose the tower has stage entropies $\Ent(A_k)$ and the tuple has joint
approximation rate $\delta_k$. If there is a stage $k^{\star}=k^{\star}(n)$ with
$\Ent(A_{k^{\star}})=O(n)$ and $\delta_{k^{\star}}=O(n^{-1/2})$, then
\[
\discop(A_1,\ldots,A_n)=O(\sqrt n) .
\]
\end{corollary}

\begin{proof}
Take $k=k^{\star}$ in Theorem~\ref{thm:af}. The entropy term is
$C\sqrt{n\log(e+O(1))}=O(\sqrt n)$ and the additive term is
$2n\cdot O(n^{-1/2})=O(\sqrt n)$.
\end{proof}

Corollary~\ref{cor:limit} is the precise sense in which the bound is ``stable
under the inductive limit'': one never needs the element to sit in a stage, only
to be approximable by a controlled-entropy stage at the scale $n^{-1/2}$ that the
$\sqrt n$ target already tolerates.

\begin{corollary}[Intrinsic-dimension AF stage bound]\label{cor:afdim}
For every \(C_0\ge1\) there is \(C(C_0)\) such that the following holds.  Let
\(A=\overline{\bigcup_kA_k}\) be a separable unital AF algebra and let
\(A_1,\ldots,A_n\in A\) be contractions.  Suppose that for some stage \(k\) with
\[
d_k=\dim_\C A_k\le C_0n
\]
there are contractions \(B_i\in A_k\) satisfying
\(\|A_i-B_i\|_{\op}<\eps\) for all \(i\).  Then
\[
\discop(A_1,\ldots,A_n)\le C(C_0)\sqrt n+2n\eps .
\]
Consequently, if the tuple is jointly approximable at scale \(O(n^{-1/2})\) by
stages with intrinsic dimension \(O(n)\), then it satisfies the
Spencer-scale bound \(O(\sqrt n)\).  The self-adjoint version holds for
self-adjoint tuples.
\end{corollary}

\begin{proof}
Apply Theorem~\ref{thm:akbassra} to the stage tuple
\(B_1,\ldots,B_n\in A_k\), and then use the same triangle-inequality
perturbation as in Theorem~\ref{thm:af}.  The factor \(2\) is kept to match the
normalization of Theorem~\ref{thm:af}, where contraction approximants may be
obtained by the truncation lemma.
\end{proof}

\begin{corollary}[Low-Frobenius perturbations of AF-stage cores]
\label{cor:afdim-pert}
For every \(C_0\ge1\) there is \(C(C_0)\) such that the following holds.  Let
\(A=\overline{\bigcup_kA_k}\) be a separable unital AF algebra and let
\(\pi:A\to\M_D(\C)\) be a unital *-representation.  Let
\[
X_i=H_i+L_i,\qquad i=1,\ldots,n,
\]
be self-adjoint matrices, where \(H_i\in\pi(A)\) are self-adjoint contractions
and the perturbations \(L_i=L_i^*\) satisfy
\[
\|L_i\|_{\op}\le1,\qquad
\|L_i\|_{\mathrm F}^2\le \frac{n}{\log^3(en)}
\]
for all \(i\).  Suppose that for some stage \(k\) with
\(\dim_\C A_k\le C_0n\) there are self-adjoint contractions
\(B_i\in\pi(A_k)\) satisfying \(\|H_i-B_i\|_{\op}<\eps\) for all \(i\).  Then
\[
\disc(X_1,\ldots,X_n)\le C(C_0)\sqrt n+2n\eps .
\]
Consequently, an AF algebraic core that is jointly captured at scale
\(O(n^{-1/2})\) by stages of intrinsic dimension \(O(n)\) remains in the
Spencer regime after such low-Frobenius self-adjoint perturbations.
In particular, the same conclusion holds if the Frobenius hypothesis on the
tails is replaced by
\[
\operatorname{rank}(L_i)\le \frac{n}{\log^3(en)}
\]
together with \(\|L_i\|_{\op}\le1\).
\end{corollary}

\begin{proof}
Apply Theorem~\ref{thm:akbassra-pert} to the self-adjoint matrices
\(B_i+L_i\).  For the resulting signs,
\[
\left\|\sum_i\eps_iX_i\right\|_{\op}
\le
\left\|\sum_i\eps_i(B_i+L_i)\right\|_{\op}
\;+\;
\sum_i\|H_i-B_i\|_{\op}
\le C(C_0)\sqrt n+n\eps,
\]
and the displayed \(2n\eps\) form again matches the truncation normalization
used above.
\end{proof}

\subsection{Which AF elements have a controlled supporting stage}
\label{subsec:linear}

Corollary~\ref{cor:finitetype} is only as useful as the supply of tuples whose
supporting stage has entropy $O(n)$. Here the structure of the Bratteli diagram
is decisive, and the obstruction of \S\ref{subsec:obstruct} is the reason a
\emph{uniform} statement from supporting-stage entropy alone is unavailable.

\paragraph{The favourable case: bounded-block towers.}
Suppose the tower has uniformly bounded block size,
$h_k=\max_\alpha r_{k,\alpha}\le b$ for all $k$ and some fixed $b$. Then
$\Ent(A_k)=h_kR_k\le bR_k$ is linear in the reduced dimension $R_k$, and any
tuple supported on a stage with $R_k=O(n)$ is in the Spencer regime by
Corollary~\ref{cor:finitetype}. Such AF algebras are exactly the ones whose
Bratteli diagram has bounded vertex labels; the commutative AF algebras
(all $r_{k,\alpha}=1$, $b=1$) are the extreme case, where $\Ent(A_k)=R_k$ is the
number of minimal projections of the stage. A self-adjoint contraction supported
on a commutative stage with $s=O(n)$ atoms is signed at the rate
$\sqrt{n\log(e+s/n)}$, recovering the association-scheme and
simultaneous-diagonalization corollaries of the main text in the limit.

\paragraph{The current larger favourable case: dimension-controlled stages.}
The bounded-block hypothesis is no longer the limit of what is known for exact
finite-dimensional algebraic support.  By Corollary~\ref{cor:afdim}, any tuple
supported, or \(O(n^{-1/2})\)-approximated, on a stage with
\[
d_k=\sum_\alpha r_{k,\alpha}^2=O(n)
\]
has discrepancy \(O(\sqrt n)\).  This can be strictly stronger than the
finite-lift criterion \(\Ent(A_k)=h_kR_k=O(n)\) when a stage is a heterogeneous
direct sum: a few larger blocks may make \(h_kR_k\) large by multiplying the
whole reduced dimension, while their actual square-sum contribution remains
small.  For a single full matrix block, however, \(d_k=\Ent(A_k)=D^2\), so the
new theorem does not change the full-block threshold.

\paragraph{The unfavourable case: UHF and growing blocks.}
For $M_{2^\infty}$ the block size doubles at every stage,
$h_k=R_k=2^k$, so $\Ent(A_k)=4^k$ grows quadratically in the stage dimension
$2^k$; here the Akbas--Sra parameter \(d_k\) is the same \(4^k\). An arbitrary
tuple controlled only through the full-stage bound lies in
$\M_{2^k}(\C)$ and gets the rate
$\sqrt{n\log(e+4^k/n)}$; to keep this at $\sqrt n$ one needs $4^k=O(n)$, i.e.
$2^k=O(\sqrt n)$, so the full-stage entropy criterion reaches arbitrary
full-span tuples only at \emph{logarithmically shallow} stages
$k\le\tfrac12\log_2 n+O(1)$. The full unit ball of
\(\M_{2^k}(\C)\), and tuples whose compressed span is the whole block, are the
full-matrix instances against which Proposition~\ref{prop:obstruct} warns: a
single block of size \(2^k\) carries \(4^k\) scalar degrees of freedom, and no
full-span faithful finite block lift of it can have entropy below \(4^k\).

\paragraph{Entropy is not additive along the tower.}
The two cases are not reconciled by summing. If a tuple decomposes as
$A_i=\sum_{\ell}A_i^{(\ell)}$ with the $\ell$th piece supported on a sub-block
$B_\ell\subseteq A_k$ of entropy $\Ent(B_\ell)=h(B_\ell)R(B_\ell)$, then by
Proposition~\ref{prop:sums} and Remark~\ref{rem:nonadd} the entropy that governs
the signed sum is
\[
\Ent\Bigl(\bigoplus_\ell B_\ell\Bigr)
=\Bigl(\max_\ell h(B_\ell)\Bigr)\sum_\ell R(B_\ell),
\]
\emph{not} $\sum_\ell\Ent(B_\ell)$. A single large block anywhere in the support
raises $\max_\ell h(B_\ell)$ and multiplies the entire reduced dimension by it.
Along a UHF tower this is the quantitative obstruction to a stage-uniform bound:
the maximal block size $h_k=2^k$ is a prefactor on \emph{all} of $R_k$, so the
entropy cannot be kept \(O(n)\) for the full supporting stage once the support
reaches depth $k>\tfrac12\log_2 n$, however the tuple is decomposed.

\paragraph{Commutative decompositions do not evade the full-block floor.}
There is one useful special case of the direct-sum accounting. Let
\(S_j\subset\M_D(\C)_{\rm sa}\) be the self-adjoint part of a maximal abelian
subalgebra. If
\[
S_1+\cdots+S_s=\M_D(\C)_{\rm sa},
\]
then \(s\ge D+1\). Indeed, each \(S_j\) contains the common scalar direction
\(\R I\), and \(S_j\cap\{X:\operatorname{tr}X=0\}\) has real dimension \(D-1\).
Projecting away \(\R I\), the traceless space has dimension \(D^2-1\), so
\[
D^2-1\le s(D-1).
\]
Thus \(s\ge D+1\). Consequently a faithful decomposition of a full self-adjoint
block into commutative \(D\)-dimensional pieces, followed by the direct-sum lift
of Proposition~\ref{prop:sums}, has combined entropy
\[
1\cdot sD\ge D(D+1)>D^2.
\]
Splitting a full block across many MASAs therefore preserves, rather than
removes, the quadratic entropy floor.

\subsection{The obstruction: von Neumann algebras and generic AF elements}
\label{subsec:obstruct}

Theorem~\ref{thm:af} requires a \emph{finite}-dimensional approximating stage.
Two distinct failures of this hypothesis bound the method from outside.

The first is internal to the finite-lift mechanism and is
Proposition~\ref{prop:robustobstruct} read along the tower: if the compressed
span of the tuple on a full matrix block \(\M_D(\C)\) is all of that block, or
even if a proposed finite-block image approximates the full unit ball in that
block to a fixed operator-norm error \(<1\), then the lift algebra has entropy
at least \(D^2\). This is a ceiling on full-span faithful finite-block reductions
along an AF tower.

\begin{proposition}[Robust entropy ceiling along an AF tower]\label{prop:afceiling}
Let $A=\varinjlim A_k$ be an AF algebra and suppose a supporting stage contains
a full matrix block $\M_D(\C)$, with central projection \(P\). Suppose a finite
block lift through \(\mcB=\bigoplus_\alpha\M_{r_\alpha}(\C)\) is used for the
compressed tuple \(P A_iP\). If the image of the lifted span is all of
\(\M_D(\C)\), or more generally if for some fixed \(\delta<1\) it is within
operator-norm distance \(\delta\) of every contraction in \(\M_D(\C)\), then
\[
\Ent(\mcB)\ge D^2 .
\]
The same conclusion holds in the self-adjoint setting when the real image
\(\delta\)-approximates every self-adjoint contraction. Consequently the entropy
term in Theorem~\ref{thm:main}, applied through any such full-span faithful lift,
cannot be better than the scale
\[
\sqrt{n\log\!\left(e+\frac{D^2}{n}\right)} .
\]
In the regime $D=\Theta(\sqrt n)$ this is the conjectured $O(\sqrt n)$; for
$D\gg\sqrt n$ it is strictly weaker, and improving it within this route would
require leaving the full-span faithful finite-block setup or avoiding the entropy
dependence.
\end{proposition}

\begin{proof}
Apply Proposition~\ref{prop:robustobstruct} to the image of the proposed lift on
the compressed block \(P A_kP\simeq\M_D(\C)\), using the complex version in the
non-self-adjoint case and the real self-adjoint version for self-adjoint tuples.
It gives
\[
\sum_\alpha r_\alpha^2\ge D^2 .
\]
Since \(\sum_\alpha r_\alpha^2\le\Ent(\mcB)\), the entropy bound follows, and the
stated scale is the entropy contribution in Theorem~\ref{thm:main}.
\end{proof}

The second failure is external to the exact finite-stage route: it removes \emph{any} finite
supporting stage. It applies to von Neumann algebras that are not AF, and to AF
algebras for elements whose spectral data is genuinely infinite.

\begin{proposition}[No finite stage for diffuse elements]\label{prop:vnobstruct}
Let $a=a^*$ be a self-adjoint element of a C*-algebra $M$ whose spectrum
$\mathrm{sp}(a)$ is an infinite set. Then $a$ lies in no finite-dimensional
C*-subalgebra of $M$. Consequently:
\begin{enumerate}[leftmargin=2em]
\item In any von Neumann algebra without minimal projections (a diffuse algebra;
in particular any $\mathrm{II}_1$ factor), there exist self-adjoint contractions
$a$ -- for instance the real part $\tfrac12(u+u^*)$ of a Haar unitary, with
atomless arcsine spectral distribution on $[-1,1]$ -- that lie in no
finite-dimensional subalgebra. For such an element neither
Theorem~\ref{thm:cstar} nor Theorem~\ref{thm:af} applies, and the finite-lift
method gives no exact-stage signing theorem; this does not rule out
tuple-dependent lifts or the trivial one-element signing bound.
\item In an infinite-dimensional AF algebra $A=\overline{\bigcup_kA_k}$ every self-adjoint element has
a finite-spectrum approximant in some stage (Lemma~\ref{lem:trunc}), but a
generic self-adjoint element of $A$ has infinite spectrum and therefore lies in
no single stage; only the \emph{approximate} statement of Theorem~\ref{thm:af} is
available, at the entropy of the stage that realizes the approximation.
\end{enumerate}
\end{proposition}

\begin{proof}
By the continuous functional calculus for a single self-adjoint element,
$C^*(a,1)\cong C(\mathrm{sp}(a))$. If $a$ were contained in a finite-dimensional
C*-subalgebra $C\subseteq M$, then $C^*(a,1)\subseteq C$ would be
finite-dimensional, forcing $C(\mathrm{sp}(a))$ finite-dimensional, i.e.\
$\mathrm{sp}(a)$ finite. This proves the displayed claim. For (1): a diffuse von
Neumann algebra contains a copy of $L^\infty[0,1]$, hence a self-adjoint element
with Lebesgue (atomless) spectral measure and continuous spectrum; equivalently
any $\mathrm{II}_1$ factor contains a Haar unitary $u$, and
$\operatorname{Re}u=\tfrac12(u+u^*)$ has the arcsine law on $[-1,1]$, an infinite
(indeed uncountable, atomless) spectrum. Such $a$ lies in no finite-dimensional
subalgebra by the displayed claim, so it is contained in no stage of any tower
and is contained in no finite-dimensional C*-subalgebra/stage; the finite-lift
method has no finite-dimensional object to feed Theorem~\ref{thm:cstar}. For
(2): in an infinite-dimensional AF algebra, finite-spectrum self-adjoints are
dense, but they are meagre: for each fixed \(m\), the set of self-adjoint
elements with at most \(m\) spectral values is closed with empty interior. Hence
a Baire-generic self-adjoint element has infinite spectrum, while the
approximate bound is exactly Theorem~\ref{thm:af}.
\end{proof}

\begin{proposition}[Spectral floor for finite-stage approximation]
\label{prop:spectralfloor}
Let $a=a^*$ be a self-adjoint contraction in a C*-algebra $M$, and let
\[
  B\simeq\bigoplus_{\alpha=1}^m\M_{r_\alpha}(\C)
\]
be a unital finite-dimensional C*-subalgebra of $M$. Put
$R(B)=\sum_\alpha r_\alpha$ and $\Ent(B)=(\max_\alpha r_\alpha)R(B)$. If
$b=b^*\in B$ and $\|a-b\|_{\op}<\eta$, then
\[
  R(B)\ge P_\eta(\mathrm{sp}(a)),
  \qquad
  \Ent(B)\ge P_\eta(\mathrm{sp}(a)),
\]
where $P_\eta(E)$ denotes the largest cardinality of a subset of
$E\subset\R$ whose distinct points are separated by more than $2\eta$.
Consequently, if $\mathrm{sp}(a)$ is infinite, then the entropy of any
finite-dimensional subalgebra containing self-adjoint approximants
$b_\eta\to a$ must tend to infinity. If $\mathrm{sp}(a)$ contains a nontrivial
interval, then this lower bound is $\Omega(\eta^{-1})$ as $\eta\downarrow0$.
\end{proposition}

\begin{proof}
For self-adjoint elements, spectra move by at most the operator-norm
perturbation: the Hausdorff distance between $\mathrm{sp}(a)$ and
$\mathrm{sp}(b)$ is at most $\|a-b\|_{\op}<\eta$. Let
$\lambda_1,\ldots,\lambda_p\in\mathrm{sp}(a)$ be pairwise separated by more than
$2\eta$, with $p=P_\eta(\mathrm{sp}(a))$. For each $j$ choose
$\mu_j\in\mathrm{sp}(b)$ with $|\lambda_j-\mu_j|<\eta$. If $\mu_i=\mu_j$ for
$i\ne j$, then
$|\lambda_i-\lambda_j|<2\eta$, a contradiction. Hence
$\mathrm{sp}(b)$ has at least $p$ points.

On the other hand, a self-adjoint element of
$\bigoplus_\alpha\M_{r_\alpha}(\C)$ has at most
$\sum_\alpha r_\alpha=R(B)$ distinct spectral values, since each block
contributes at most its matrix size. Thus $R(B)\ge p$, and
$\Ent(B)\ge R(B)\ge p$. If $\mathrm{sp}(a)$ is infinite, its packing numbers
$P_\eta(\mathrm{sp}(a))$ are unbounded as $\eta\downarrow0$; if it contains an
interval of positive length, the elementary packing bound on an interval gives
$P_\eta(\mathrm{sp}(a))\ge c\,\eta^{-1}$ for small $\eta$.
\end{proof}

\begin{remark}[What the spectral floor does and does not prove]
Proposition~\ref{prop:spectralfloor} is a lower bound on the finite-stage
entropy required to approximate a diffuse self-adjoint element; it is not a
Matrix Spencer lower bound. For example, if
$\mathrm{sp}(a)=[-1,1]$ and one wants the approximation term in
Theorem~\ref{thm:af} to be $O(\sqrt n)$, then one takes
$\eta=O(n^{-1/2})$, and Proposition~\ref{prop:spectralfloor} only forces
$\Ent(B)\gtrsim\sqrt n$. This is still compatible with the Spencer-scale upper
bound. The proposition's role is narrower: it rules out a dimension-free
finite-stage model for diffuse spectral data and makes the approximation/entropy
trade-off in Theorem~\ref{thm:af} unavoidable.
\end{remark}

\begin{corollary}[Commutative AF partition trade-off]
\label{cor:comm-af-partition}
Let \(A=C(X)\) be a commutative AF algebra, so finite stages are algebras of
functions constant on finite clopen partitions.  Let
\(f_1,\ldots,f_n\in C(X)\) be real-valued contractions, and let
\(\mathcal P\) be a finite clopen partition of \(X\) with
\(M=\#\mathcal P\).  Suppose
\[
  \operatorname{osc}_{\mathcal P}(f_i)
  :=\max_{P\in\mathcal P}\sup_{x,y\in P}|f_i(x)-f_i(y)|
  \le \eta
  \qquad(1\le i\le n).
\]
Then
\[
  \disc(f_1,\ldots,f_n)
  \le C\sqrt{n\log\!\left(e+\frac{M}{n}\right)}+2n\eta .
\]
Consequently, if for each \(n\) the tuple is captured by a clopen partition with
\(M=O(n)\) and \(\eta=O(n^{-1/2})\), then
\(\disc(f_1,\ldots,f_n)=O(\sqrt n)\).

Conversely, if one \(f_i(X)\) contains an interval of positive length, then any
finite clopen partition that supports an \(\eta\)-approximant to that \(f_i\)
must have \(M=\Omega(\eta^{-1})\) atoms as \(\eta\downarrow0\).  In particular,
any commutative finite-stage certificate for the \(O(n^{-1/2})\) approximation
scale must use \(M=\Omega(\sqrt n)\) atoms.  This is a lower bound on the stage
complexity required by the finite-stage certificate, not a lower bound on the
discrepancy itself.
\end{corollary}

\begin{proof}
For each \(P\in\mathcal P\) and each \(i\), choose a point \(x_P\in P\) and let
\(g_i\) be the function that is identically \(f_i(x_P)\) on \(P\).  Then
\(g_i\) lies in the finite-dimensional algebra \(B_{\mathcal P}\simeq\C^M\),
\(\|f_i-g_i\|_\infty\le\eta\), and
\(\Ent(B_{\mathcal P})=M\).  The displayed upper bound is Theorem~\ref{thm:af}
in the self-adjoint/commutative case.

For the lower statement, let \(f=f_i\) and suppose \(f(X)\) contains an interval
\(I\) of length \(\ell>0\).  If \(g\in B_{\mathcal P}\) and
\(\|f-g\|_\infty<\eta\), then Proposition~\ref{prop:spectralfloor} gives
\[
  M=\Ent(B_{\mathcal P})\ge P_\eta(f(X)).
\]
The interval \(I\) contains at least \(c_\ell\eta^{-1}\) points separated by
more than \(2\eta\), for all sufficiently small \(\eta\).  Hence
\(M\ge c_\ell\eta^{-1}\).  Taking \(\eta\asymp n^{-1/2}\) gives the
\(\Omega(\sqrt n)\) certificate-size floor.
\end{proof}

\begin{corollary}[Metric clopen partitions]
\label{cor:comm-af-metric}
Let \(X\) be a compact zero-dimensional metric space and, for \(r>0\), let
\(M_X(r)\) be the least cardinality of a finite clopen partition of \(X\) whose
atoms all have diameter at most \(r\).  Let \(f_1,\ldots,f_n\in C(X)\) be
real-valued contractions, and put
\[
  \omega(r):=\max_{1\le i\le n}\sup_{d(x,y)\le r}|f_i(x)-f_i(y)|.
\]
Then, for every \(r>0\) with \(M_X(r)<\infty\),
\[
  \disc(f_1,\ldots,f_n)
  \le C\sqrt{n\log\!\left(e+\frac{M_X(r)}{n}\right)}
      +2n\omega(r) .
\]
In particular, if \(M_X(r)\le A r^{-d}\) for \(0<r\le1\) and the tuple is
uniformly \((L,\alpha)\)-Hölder, meaning \(\omega(r)\le Lr^\alpha\), then
\[
  \disc(f_1,\ldots,f_n)
  \le C_{A,d,L,\alpha}
     \sqrt{n\log\!\left(e+n^{d/(2\alpha)-1}\right)} .
\]
Thus polynomial clopen partition growth gives an \(O(\sqrt{n\log n})\) bound
for fixed \(d\) and \(\alpha\), and the dimension-regularity threshold
\(d\le2\alpha\) gives the Spencer-scale \(O(\sqrt n)\) bound.
\end{corollary}

\begin{proof}
If every atom of \(\mathcal P\) has diameter at most \(r\), then
\(\operatorname{osc}_{\mathcal P}(f_i)\le \omega(r)\) for every \(i\).  The
first display is Corollary~\ref{cor:comm-af-partition}.  Under the polynomial
partition-growth and Hölder hypotheses, take
\(r=\min(1,(L\sqrt n)^{-1/\alpha})\).  The additive term is
\(O_{L,\alpha}(\sqrt n)\), and
\[
  \frac{M_X(r)}{n}\le C_{A,d,L,\alpha} n^{d/(2\alpha)-1}
\]
after increasing the constant to absorb the finitely many small \(n\).  This
gives the displayed bound and its two stated consequences.
\end{proof}

\begin{lemma}[Clopen shrinkings of finite covers]
\label{lem:clopen-shrinking}
Let \(X\) be a compact zero-dimensional metric space.  If
\(U_1,\ldots,U_m\) is a finite open cover of \(X\), then, after allowing empty
sets, there are pairwise disjoint clopen sets \(V_1,\ldots,V_m\) such that
\[
  V_j\subset U_j\qquad(1\le j\le m),
  \qquad
  \bigcup_{j=1}^m V_j=X.
\]
Consequently, if \(N_X(\rho)\) is the least number of open \(\rho\)-balls
needed to cover \(X\), then
\[
  M_X(2\rho)\le N_X(\rho)\qquad(\rho>0).
\]
\end{lemma}

\begin{proof}
We first recall the elementary separation property used below.  If
\(F\subset U\subset X\), \(F\) is closed and \(U\) is open, then zero
dimensionality gives, for every \(x\in F\), a clopen neighborhood
\(C_x\) with \(x\in C_x\subset U\).  Compactness of \(F\) gives a finite
subcover, and the union of the selected \(C_x\)'s is clopen; hence there is a
clopen \(C\) with \(F\subset C\subset U\).

The cover refinement follows by induction on \(m\).  For \(m=1\) there is
nothing to prove.  For \(m>1\), set
\[
  F_1=X\setminus \bigcup_{j=2}^m U_j\subset U_1.
\]
Choose a clopen \(V_1\) with \(F_1\subset V_1\subset U_1\).  Then
\(X\setminus V_1\) is compact, zero-dimensional, and covered by
\(U_2,\ldots,U_m\).  Applying the induction hypothesis in this clopen subspace
gives \(V_2,\ldots,V_m\), which are also clopen in \(X\).

For the last assertion, cover \(X\) by \(N_X(\rho)\) open \(\rho\)-balls and
apply the first part.  Each resulting clopen atom lies in one of those balls
and therefore has diameter at most \(2\rho\).
\end{proof}

\begin{corollary}[Ordinary covering entropy]
\label{cor:comm-af-covering}
Let \(X\) be a compact zero-dimensional metric space, let \(N_X(\rho)\) be the
least number of open \(\rho\)-balls needed to cover \(X\), and let
\(f_1,\ldots,f_n\in C(X)\) be real-valued contractions with common modulus
\[
  \omega(r):=\max_{1\le i\le n}\sup_{d(x,y)\le r}|f_i(x)-f_i(y)|.
\]
Then, for every \(\rho>0\),
\[
  \disc(f_1,\ldots,f_n)
  \le C\sqrt{n\log\!\left(e+\frac{N_X(\rho)}{n}\right)}
      +2n\omega(2\rho).
\]
In particular, if
\[
  N_X(\rho)\le A\rho^{-d}\qquad(0<\rho\le1)
\]
and the tuple is uniformly \((L,\alpha)\)-Holder, then
\[
  \disc(f_1,\ldots,f_n)
  \le C_{A,d,L,\alpha}
     \sqrt{n\log\!\left(e+n^{d/(2\alpha)-1}\right)}.
\]
Thus ordinary polynomial metric entropy gives the same
\(O(\sqrt{n\log n})\) bound as clopen partition growth, and the threshold
\(d\le2\alpha\) gives the Spencer-scale bound \(O(\sqrt n)\).  More generally,
if
\[
  N_X(\rho)\le A\rho^{-2\alpha}\bigl(\log(e/\rho)\bigr)^\beta
  \qquad(0<\rho\le1),
\]
then
\[
  \disc(f_1,\ldots,f_n)
  \le
  \begin{cases}
  C_{A,L,\alpha,\beta}\sqrt n,
    & \beta\le0,\\[3pt]
  C_{A,L,\alpha,\beta}\sqrt{n\log\log(e^e+n)},
    & \beta>0.
  \end{cases}
\]
\end{corollary}

\begin{proof}
Lemma~\ref{lem:clopen-shrinking} gives \(M_X(2\rho)\le N_X(\rho)\), so the
first display is Corollary~\ref{cor:comm-af-metric} applied at scale \(2\rho\).
If \(N_X(\rho)\le A\rho^{-d}\) and \(\omega(r)\le Lr^\alpha\), take
\[
  \rho=\min\left(1,\frac12(L\sqrt n)^{-1/\alpha}\right).
\]
Then \(2n\omega(2\rho)\le2nL(2\rho)^\alpha=O_{L,\alpha}(\sqrt n)\), while
\[
  \frac{N_X(\rho)}{n}
  \le C_{A,d,L,\alpha}n^{d/(2\alpha)-1}
\]
after enlarging the constant for the finitely many small \(n\).  The
logarithmic statement is identical, with
\(N_X(\rho)/n\le C_{A,L,\alpha}(\log(e+n))^\beta\).
\end{proof}

\begin{definition}[Ordinary metric Minkowski data]
\label{def:metric-minkowski}
For a compact metric space \(X\), let
\[
  \overline{\dim}_{\mathrm M}(X)
  :=
  \limsup_{\rho\downarrow0}
  \frac{\log N_X(\rho)}{\log(1/\rho)}
\]
be its upper metric Minkowski dimension, where \(N_X(\rho)\) denotes the least
number of open \(\rho\)-balls needed to cover \(X\).  For \(s\ge0\), let
\[
  \overline{\mathcal M}^{\,s}_{\mathrm M}(X)
  :=
  \limsup_{\rho\downarrow0}\rho^s N_X(\rho)
\]
be the corresponding upper metric \(s\)-dimensional Minkowski content.
\end{definition}

\begin{corollary}[Ordinary Minkowski threshold]
\label{cor:comm-af-metric-minkowski}
Let \(X\) be a compact zero-dimensional metric space and let
\(f_1,\ldots,f_n\in C(X)\) be real-valued contractions with common
\((L,\alpha)\)-H\"older modulus.  If
\(\overline{\dim}_{\mathrm M}(X)<\infty\), then
\[
  \disc(f_1,\ldots,f_n)\le C_{X,L,\alpha}\sqrt{n\log(e+n)} .
\]
If, in addition,
\[
  \overline{\dim}_{\mathrm M}(X)<2\alpha,
\]
then
\[
  \disc(f_1,\ldots,f_n)\le C_{X,L,\alpha}\sqrt n .
\]
At the critical exponent, the same Spencer-scale conclusion holds whenever
\[
  \overline{\mathcal M}^{\,2\alpha}_{\mathrm M}(X)<\infty .
\]
Thus the metric AF bound can be read from the ordinary covering dimension and
ordinary critical content of \(X\), without separately estimating clopen
partition growth.
\end{corollary}

\begin{proof}
If \(D>\overline{\dim}_{\mathrm M}(X)\), then there are constants
\(C_D<\infty\) and \(\rho_D>0\) such that
\[
  N_X(\rho)\le C_D\rho^{-D}\qquad(0<\rho\le\rho_D).
\]
After enlarging \(C_D\) to cover the finitely many larger scales, and after the
same harmless metric rescaling used below, Corollary~\ref{cor:comm-af-covering}
gives
\[
  \disc(f_1,\ldots,f_n)
  \le C_{D,L,\alpha,X}
     \sqrt{n\log\!\left(e+n^{D/(2\alpha)-1}\right)}.
\]
This is \(O_{X,L,\alpha}(\sqrt{n\log(e+n)})\) for finite \(D\).  If
\(\overline{\dim}_{\mathrm M}(X)<2\alpha\), choose \(D\) strictly between
\(\overline{\dim}_{\mathrm M}(X)\) and \(2\alpha\), making the logarithm
bounded.

Finally, finite upper metric \(2\alpha\)-content gives
\[
  N_X(\rho)\le C_0\rho^{-2\alpha}\qquad(0<\rho\le\rho_0)
\]
for all sufficiently small \(\rho\).  Enlarging the constant for larger scales
puts the hypotheses of Corollary~\ref{cor:comm-af-covering} in force with
\(d=2\alpha\), hence the logarithmic factor is bounded and the Spencer-scale
estimate follows.
\end{proof}

\begin{definition}[Upper Assouad growth]
\label{def:upper-assouad-growth}
Let \(X\) be a compact metric space with \(\operatorname{diam}X\le1\).  We say
that \(X\) has upper Assouad growth \((A,s)\) if every closed ball in \(X\) of
radius \(R\le1\) can be covered by at most
\[
  A\left(\frac R\rho\right)^s
\]
open balls of radius \(\rho\), for all \(0<\rho\le R\le1\).
\end{definition}

\begin{definition}[Logarithmic upper Assouad growth]
\label{def:upper-log-assouad-growth}
Let \(X\) be a compact metric space with \(\operatorname{diam}X\le1\).  We say
that \(X\) has logarithmic upper Assouad growth \((A,s,\gamma)\) if every
closed ball in \(X\) of radius \(R\le1\) can be covered by at most
\[
  A\left(\frac R\rho\right)^s\bigl(\log(eR/\rho)\bigr)^\gamma
\]
open balls of radius \(\rho\), for all \(0<\rho\le R\le1\).
\end{definition}

\begin{corollary}[Doubling and Assouad entropy]
\label{cor:comm-af-assouad}
Let \(X\) be a compact zero-dimensional metric space with
\(\operatorname{diam}X\le1\), and let \(f_1,\ldots,f_n\in C(X)\) be
real-valued contractions with common \((L,\alpha)\)-H\"older modulus.
If \(X\) has upper Assouad growth \((A,s)\), then
\[
  \disc(f_1,\ldots,f_n)
  \le C_{A,s,L,\alpha}
     \sqrt{n\log\!\left(e+n^{s/(2\alpha)-1}\right)}.
\]
In particular, \(s\le2\alpha\) gives the Spencer-scale bound
\[
  \disc(f_1,\ldots,f_n)\le C_{A,s,L,\alpha}\sqrt n .
\]
Consequently, if \(X\) is \(D\)-doubling, meaning that every closed ball of
radius \(R\le1\) is covered by at most \(D\) open balls of radius \(R/2\), then
the same display holds with \(s=\log_2 D\) and constants depending on
\(D,L,\alpha\).  Thus a \(D\)-doubling zero-dimensional model is Spencer-scale
for
\((L,\alpha)\)-H\"older tuples whenever \(\log_2D\le2\alpha\).
\end{corollary}

\begin{proof}
Fix \(x_0\in X\).  Since \(\operatorname{diam}X\le1\), \(X\) is contained in
the closed ball \(\overline B(x_0,1)\).  The upper Assouad growth assumption
with \(R=1\) gives
\[
  N_X(\rho)\le A\rho^{-s}\qquad(0<\rho\le1).
\]
The first two assertions are therefore Corollary~\ref{cor:comm-af-covering}.

If \(X\) is \(D\)-doubling, iterate the doubling cover.  For each integer
\(k\ge0\), \(X\) is covered by at most \(D^k\) balls of radius \(2^{-k}\).
Given \(0<\rho\le1\), choose \(k\) with \(2^{-k}\le\rho<2^{-k+1}\).  Then
\[
  N_X(\rho)\le D^k
  \le D\,\rho^{-\log_2D}.
\]
Thus \(X\) has the global polynomial covering bound with exponent
\(\log_2D\), and another application of
Corollary~\ref{cor:comm-af-covering} gives the stated doubling consequence.
\end{proof}

\begin{corollary}[Finite products add ordinary entropy]
\label{cor:comm-af-ordinary-product}
Let \(X_j\), \(1\le j\le m\), be compact zero-dimensional metric spaces and
let \(X=\prod_{j=1}^m X_j\) with the max product metric.  Suppose
\[
  N_{X_j}(\rho)\le A_j\rho^{-d_j}\qquad(0<\rho\le1),
\]
and put \(d_*=\sum_{j=1}^m d_j\).  If
\(f_1,\ldots,f_n\in C(X)\) are real-valued contractions with common
\((L,\alpha)\)-H\"older modulus, then
\[
  \disc(f_1,\ldots,f_n)
  \le C_{\{A_j,d_j\},L,\alpha,m}
     \sqrt{n\log\!\left(e+n^{d_*/(2\alpha)-1}\right)}.
\]
In particular, \(d_*\le2\alpha\) gives the Spencer-scale bound
\[
  \disc(f_1,\ldots,f_n)\le C_{\{A_j,d_j\},L,\alpha,m}\sqrt n .
\]
Equivalently, finite ordinary upper metric dimensions add under finite
max-metric products: if
\[
  \sum_{j=1}^m \overline{\dim}_{\mathrm M}(X_j)<2\alpha,
\]
then every common \((L,\alpha)\)-H\"older tuple on \(X\) is Spencer-scale.  At
the critical endpoint, the same conclusion holds if
\[
  \sum_{j=1}^m s_j=2\alpha
  \quad\text{and}\quad
  \overline{\mathcal M}^{\,s_j}_{\mathrm M}(X_j)<\infty
  \quad(1\le j\le m).
\]
Likewise, if each \(X_j\) has upper Assouad growth \((A_j,s_j)\) and
\(\operatorname{diam}X_j\le1\), then \(X\) has upper Assouad growth
\((\prod_j A_j,\sum_j s_j)\), so \(\sum_j s_j\le2\alpha\) gives Spencer scale.
In particular, a finite max product of \(D_j\)-doubling factors with
\(\operatorname{diam}X_j\le1\) is
Spencer-scale whenever
\[
  \sum_{j=1}^m \log_2D_j\le2\alpha .
\]
\end{corollary}

\begin{proof}
Finite products of compact zero-dimensional metric spaces are again compact
and zero-dimensional.  A product of open \(\rho\)-balls in the factors is an
open \(\rho\)-ball for the max product metric.  Hence
\[
  N_X(\rho)\le\prod_{j=1}^m N_{X_j}(\rho)
  \le \left(\prod_{j=1}^m A_j\right)\rho^{-d_*},
  \qquad 0<\rho\le1.
\]
The displayed discrepancy estimate and the \(d_*\le2\alpha\) consequence follow
from Corollary~\ref{cor:comm-af-covering}.

For the dimension statement, choose \(D_j>\overline{\dim}_{\mathrm M}(X_j)\)
with \(\sum_jD_j<2\alpha\) and apply the preceding paragraph.  For the critical
content endpoint, finite upper metric \(s_j\)-content gives
\(N_{X_j}(\rho)\le C_j\rho^{-s_j}\) at small scales, and enlarging constants
handles the larger scales; the product bound then has exponent
\(\sum_j s_j=2\alpha\).

For the Assouad statement, a closed \(R\)-ball in the max product metric is the
product of the closed \(R\)-balls in the factors.  Cover each factor ball by at
most \(A_j(R/\rho)^{s_j}\) open \(\rho\)-balls and multiply the covers.  The
doubling consequence is the special case \(s_j=\log_2D_j\) from
Corollary~\ref{cor:comm-af-assouad}.
\end{proof}

\begin{definition}[Zero-dimensional H\"older resolution]
\label{def:zero-dim-holder-resolution}
Let \(Y\) be a compact metric space.  An
\((A,s,H,\beta)\)-zero-dimensional H\"older resolution of \(Y\) consists of a
compact zero-dimensional metric space \(X\) and a continuous surjection
\(\pi:X\to Y\) such that
\[
  N_X(\rho)\le A\rho^{-s}\qquad(0<\rho\le1)
\]
and
\[
  d_Y(\pi x,\pi y)\le H\,d_X(x,y)^\beta\qquad(x,y\in X).
\]
\end{definition}

\begin{corollary}[H\"older resolutions transfer the entropy bound]
\label{cor:comm-af-holder-resolution}
Let \(Y\) admit an \((A,s,H,\beta)\)-zero-dimensional H\"older resolution.
If \(g_1,\ldots,g_n\in C(Y)\) are real-valued contractions with common
\((L,\alpha)\)-H\"older modulus, then
\[
  \disc(g_1,\ldots,g_n)
  \le C_{A,s,H,L,\alpha,\beta}
     \sqrt{n\log\!\left(e+n^{s/(2\alpha\beta)-1}\right)}.
\]
In particular, \(s\le2\alpha\beta\) gives the Spencer-scale bound
\[
  \disc(g_1,\ldots,g_n)\le C_{A,s,H,L,\alpha,\beta}\sqrt n .
\]
\end{corollary}

\begin{proof}
Let \(f_i=g_i\circ\pi\).  Since \(\pi\) is onto, pullback is isometric:
\[
  \left\|\sum_i\eps_i g_i\right\|_{C(Y)}
  =
  \left\|\sum_i\eps_i f_i\right\|_{C(X)}
  \qquad(\eps_i\in\{-1,1\}).
\]
Thus the discrepancies of the two tuples are equal.  Moreover
\[
  |f_i(x)-f_i(y)|
  \le L\,d_Y(\pi x,\pi y)^\alpha
  \le L H^\alpha d_X(x,y)^{\alpha\beta}.
\]
Corollary~\ref{cor:comm-af-covering}, applied on \(X\) with exponent
\(\alpha\beta\), gives the result.
\end{proof}

\begin{proposition}[Entropy-profile zero-dimensional resolution]
\label{prop:compact-metric-resolution}
Let \(Y\) be a compact metric space with \(\operatorname{diam}Y\le1\), and let
\(0<\theta<1\).  Then there are a compact zero-dimensional metric space \(X\)
and a continuous surjection \(\pi:X\to Y\) such that
\[
  d_Y(\pi x,\pi x')\le\frac{2}{1-\theta}d_X(x,x')
  \qquad(x,x'\in X)
\]
and
\[
  N_X(\rho)\le N_Y(\theta\rho)\qquad(0<\rho\le1).
\]
In particular, if \(N_Y(\rho)\le A\rho^{-d}\) for \(0<\rho\le1\), then, for
every \(0<\theta<1\), \(Y\) admits a
\((\theta^{-d}A,d,2/(1-\theta),1)\)-zero-dimensional H\"older resolution.
\end{proposition}

\begin{proof}
Put \(r_k=\theta^k\).  For each \(k\ge0\), choose a set
\(C_k\subset Y\) with
\[
  |C_k|=N_Y(r_k)
  \qquad\text{and}\qquad
  Y\subset \bigcup_{c\in C_k}B(c,r_k).
\]
For each \(c\in C_{k+1}\), choose a parent \(p(c)\in C_k\) with
\(d_Y(c,p(c))<r_k\).  Let \(X\) be the space of all infinite parent chains
\[
  x=(c_0,c_1,\ldots),\qquad c_k\in C_k,\qquad p(c_{k+1})=c_k .
\]
As a closed subspace of the product of finite discrete spaces, \(X\) is compact
and zero-dimensional.

For \(x=(c_k)\in X\), the estimate
\[
  d_Y(c_\ell,c_k)
  \le \sum_{j=k}^{\ell-1}r_j < \frac{r_k}{1-\theta}
  \qquad(\ell>k)
\]
shows that \((c_k)\) converges in \(Y\).  Define
\(\pi(x)=\lim_k c_k\).  This map is continuous.

To see that \(\pi\) is onto, fix \(y\in Y\).  For every \(m\), choose
\(a_m\in C_m\) with \(d_Y(y,a_m)<r_m\), and take the ancestor chain of
\(a_m\) down to level \(0\).  A diagonal subsequence stabilizes at every fixed
level \(k\), producing a chain \(x=(c_k)\in X\).  Along the stabilizing
subsequence,
\[
  d_Y(y,c_k)
  \le r_m+\sum_{j=k}^{m-1}r_j
  < r_m+\frac{r_k}{1-\theta},
\]
so \(\pi(x)=y\).

Define an ultrametric on \(X\) by setting \(d_X(x,x')=0\) if \(x=x'\), and
otherwise
\[
  d_X(x,x')=
  \begin{cases}
  1, & c_0\ne c'_0,\\
  \theta^k, & k=\max\{j:c_j=c'_j\}.
  \end{cases}
\]
If \(x\) and \(x'\) agree through level \(k\), then both limits are within
\(r_k/(1-\theta)\) of the common center \(c_k\).  Hence
\[
  d_Y(\pi x,\pi x')\le\frac{2r_k}{1-\theta}
  =\frac{2}{1-\theta}d_X(x,x')
\]
whenever \(d_X(x,x')=\theta^k\), and the same bound is trivial when
\(d_X(x,x')=1\) because \(\operatorname{diam}Y\le1\).

It remains to check the covering entropy of \(X\).  Given \(0<\rho\le1\), let
\(k\) be minimal with \(\theta^k<\rho\).  The level-\(k\) cylinders cover \(X\),
are contained in open \(\rho\)-balls for \(d_X\), and their number is at most
\(|C_k|\).  Since \(\theta^k\ge\theta\rho\) and covering numbers decrease with the
radius,
\[
  N_X(\rho)\le |C_k|=N_Y(\theta^k)\le N_Y(\theta\rho).
\]
The polynomial consequence follows by applying the assumed bound at scale
\(\theta\rho\).
\end{proof}

\begin{corollary}[Compact metric entropy without zero-dimensionality]
\label{cor:comm-af-compact-metric}
Let \(Y\) be a compact metric space with \(\operatorname{diam}Y\le1\), let
\(N_Y(\rho)\) be its ordinary covering number, and let
\(g_1,\ldots,g_n\in C(Y)\) be real-valued contractions with common modulus
\[
  \omega(r):=\max_{1\le i\le n}\sup_{d_Y(y,z)\le r}|g_i(y)-g_i(z)|.
\]
Then, for every \(0<\theta<1\) and \(0<s\le1\),
\[
  \disc(g_1,\ldots,g_n)
  \le C\sqrt{n\log\!\left(e+\frac{N_Y(\theta s)}{n}\right)}
      +2n\omega\!\left(\frac{4s}{1-\theta}\right).
\]
In particular, for every \(0<\rho\le1/2\),
\[
  \disc(g_1,\ldots,g_n)
  \le C\sqrt{n\log\!\left(e+\frac{N_Y(\rho)}{n}\right)}
      +2n\omega(16\rho).
\]
In particular, if
\[
  N_Y(\rho)\le A\rho^{-d}\qquad(0<\rho\le1),
\]
and the tuple is uniformly \((L,\alpha)\)-H\"older, then
\[
  \disc(g_1,\ldots,g_n)
  \le C_{A,d,L,\alpha}
     \sqrt{n\log\!\left(e+n^{d/(2\alpha)-1}\right)}.
\]
In particular, \(d\le2\alpha\) gives
\[
  \disc(g_1,\ldots,g_n)\le C_{A,d,L,\alpha}\sqrt n .
\]
More generally, if
\[
  N_Y(\rho)\le A\rho^{-2\alpha}\bigl(\log(e/\rho)\bigr)^\beta
  \qquad(0<\rho\le1),
\]
then
\[
  \disc(g_1,\ldots,g_n)
  \le
  \begin{cases}
  C_{A,L,\alpha,\beta}\sqrt n,
    & \beta\le0,\\[3pt]
  C_{A,L,\alpha,\beta}\sqrt{n\log\log(e^e+n)},
    & \beta>0.
  \end{cases}
\]
Consequently, the following hold for an arbitrary compact metric \(Y\).  Finite
upper metric Minkowski dimension gives
\[
  \disc(g_1,\ldots,g_n)\le C_{Y,L,\alpha}\sqrt{n\log(e+n)} .
\]
Upper metric dimension \(<2\alpha\) gives Spencer scale, and finite upper
metric \(2\alpha\)-content gives the critical Spencer-scale endpoint.
\end{corollary}

\begin{proof}
Fix \(0<\theta<1\).  Let \(X\) and \(\pi:X\to Y\) be the resolution from
Proposition~\ref{prop:compact-metric-resolution} with this value of \(\theta\),
and put \(f_i=g_i\circ\pi\).  As in
Corollary~\ref{cor:comm-af-holder-resolution}, the discrepancies of \((g_i)\)
and \((f_i)\) are equal.  The common modulus \(\omega_X\) of the pulled-back
tuple satisfies \(\omega_X(r)\le\omega(2r/(1-\theta))\).  Applying
Corollary~\ref{cor:comm-af-covering} on \(X\) at scale \(s\) gives
\[
  \disc(g_1,\ldots,g_n)
  \le C\sqrt{n\log\!\left(e+\frac{N_Y(\theta s)}{n}\right)}
      +2n\omega\!\left(\frac{4s}{1-\theta}\right).
\]
Taking \(\theta=1/2\) and \(s=2\rho\) gives the displayed direct profile bound.

If \(N_Y(\rho)\le A\rho^{-d}\) and \(\omega(r)\le Lr^\alpha\), take
\[
  \rho=\min\left(\frac12,(16L\sqrt n)^{-1/\alpha}\right).
\]
Then the approximation term is \(O_{L,\alpha}(\sqrt n)\), while
\[
  \frac{N_Y(\rho)}{n}
  \le C_{A,d,L,\alpha}n^{d/(2\alpha)-1}.
\]
This proves the polynomial estimate.  The logarithmic critical-window estimate
is identical, using
\[
  \frac{N_Y(\rho)}{n}
  \le C_{A,L,\alpha}\bigl(\log(e+n)\bigr)^\beta
\]
at the same choice of \(\rho\).

For the dimension and content consequences, first rescale the metric by a
constant depending only on \(Y\) so that \(\operatorname{diam}Y\le1\).
The upper metric dimension and critical-content consequences are exactly the
same enlarging-of-constants argument as in
Corollary~\ref{cor:comm-af-metric-minkowski}: choose any exponent strictly
above \(\overline{\dim}_{\mathrm M}(Y)\), or use the finite
\(2\alpha\)-content bound at small scales and enlarge constants for the
remaining scales.
\end{proof}

\begin{corollary}[Qualitative compact equicontinuity]
\label{cor:comm-af-compact-equicontinuity}
Let \(Y\) be a compact metric space with \(\operatorname{diam}Y\le1\), and let
\(\omega\) be a modulus with \(\omega(r)\to0\) as \(r\downarrow0\).  For each
\(n\ge1\), let \(g_{n,1},\ldots,g_{n,n}\in C(Y)\) be real-valued contractions
such that
\[
  |g_{n,i}(y)-g_{n,i}(z)|\le \omega(d_Y(y,z))
  \qquad(y,z\in Y,\ 1\le i\le n).
\]
Then
\[
  \disc(g_{n,1},\ldots,g_{n,n})=o(n).
\]
More quantitatively, if \(\rho_n\in(0,1/2]\) satisfies
\[
  \log N_Y(\rho_n)=o(n)
  \qquad\text{and}\qquad
  \omega(16\rho_n)=o(1),
\]
then the same conclusion holds.
\end{corollary}

\begin{proof}
The quantitative statement follows directly from
Corollary~\ref{cor:comm-af-compact-metric}:
\[
  \disc(g_{n,1},\ldots,g_{n,n})
  \le C\sqrt{n\log\!\left(e+\frac{N_Y(\rho_n)}{n}\right)}
       +2n\omega(16\rho_n).
\]
If \(E_n=\log N_Y(\rho_n)\), then
\[
  \frac1n
  \sqrt{n\log\!\left(e+\frac{N_Y(\rho_n)}{n}\right)}
  \le
  \sqrt{\frac{1+E_n}{n}}=o(1),
\]
and the approximation term is \(o(n)\) by assumption.

It remains only to note that such a scale sequence exists for every compact
\(Y\) and every fixed modulus \(\omega\).  Choose
\(\rho_k\in(0,1/16]\) with \(\rho_k\downarrow0\) and
\(\omega(16\rho_k)\le1/k\).  Put \(E_k=\log N_Y(\rho_k)\) and
\[
  m_k=\left\lceil k\max\{1,E_k\}\right\rceil .
\]
For each sufficiently large \(n\), let \(k(n)\) be the largest \(k\) with
\(m_k\le n\), and set \(\rho_n=\rho_{k(n)}\).  Since \(m_k\ge k\) and each
\(m_k\) is finite, \(k(n)\to\infty\).  Moreover
\[
  \frac{\log N_Y(\rho_n)}{n}
  \le \frac1{k(n)}
  \qquad\text{and}\qquad
  \omega(16\rho_n)\le\frac1{k(n)} .
\]
The quantitative statement applies.
\end{proof}

\begin{corollary}[Precompact function classes]
\label{cor:comm-af-compact-function-class}
Let \(Y\) be a compact metric space with \(\operatorname{diam}Y\le1\), and let
\(\mathcal F\subset C(Y)\) be a relatively compact family of real-valued
contractions.  Define
\[
  \Delta_n(\mathcal F)
  :=
  \sup_{f_1,\ldots,f_n\in\mathcal F}\disc(f_1,\ldots,f_n).
\]
Then
\[
  \Delta_n(\mathcal F)=o(n).
\]
\end{corollary}

\begin{proof}
Let \(K=\overline{\mathcal F}\) in \(C(Y)\).  Compactness of \(K\) gives a
common modulus of continuity.  Indeed, given \(\eps>0\), choose an
\(\eps/3\)-net \(h_1,\ldots,h_m\) for \(K\) in the uniform norm, then choose
\(\delta>0\) so that
\[
  |h_j(y)-h_j(z)|<\eps/3
  \qquad(d_Y(y,z)<\delta,\ 1\le j\le m).
\]
If \(f\in K\) and \(\|f-h_j\|_\infty<\eps/3\), then
\(|f(y)-f(z)|<\eps\) whenever \(d_Y(y,z)<\delta\).  Thus
\[
  \omega(r):=\sup_{f\in K}\sup_{d_Y(y,z)\le r}|f(y)-f(z)|
\]
satisfies \(\omega(r)\to0\) as \(r\downarrow0\).

Apply Corollary~\ref{cor:comm-af-compact-equicontinuity} with this common
modulus.  Its proof gives scales \(\rho_n\), depending only on \(Y\) and \(K\),
for which
\[
  \disc(f_1,\ldots,f_n)
  \le C\sqrt{n\log\!\left(e+\frac{N_Y(\rho_n)}{n}\right)}
       +2n\omega(16\rho_n)
  =o(n)
\]
uniformly for all \(f_1,\ldots,f_n\in K\).  Taking the supremum over
\(\mathcal F^n\) gives the claim.
\end{proof}

\begin{corollary}[Precompact Banach classes]
\label{cor:precompact-banach-class}
Let \(E\) be a normed space and let \(\mathcal K\) be a totally bounded subset
of the unit ball of \(E\).  Set
\[
  \Delta_n^E(\mathcal K)
  :=
  \sup_{x_1,\ldots,x_n\in\mathcal K}
  \min_{\eps_i\in\{-1,1\}}
  \left\|\sum_{i=1}^n\eps_i x_i\right\|_E .
\]
Let \(N_{\mathcal K}(\eta)\) denote the least size of an \(\eta\)-net for
\(\mathcal K\) whose centers lie in \(\mathcal K\).  Then
\[
  \Delta_n^E(\mathcal K)\le N_{\mathcal K}(\eta)+n\eta
  \qquad(\eta>0).
\]
Consequently
\[
  \Delta_n^E(\mathcal K)=o(n).
\]
\end{corollary}

\begin{proof}
Choose an \(\eta\)-net \(u_1,\ldots,u_M\) for \(\mathcal K\), with
\(M=N_{\mathcal K}(\eta)\), and take \(x_1,\ldots,x_n\in\mathcal K\).  For each
\(i\), choose \(a(i)\) with \(\|x_i-u_{a(i)}\|_E\le\eta\).  For every
\(1\le j\le M\), choose signs on the fiber
\[
  I_j=\{i:a(i)=j\}
\]
so that
\[
  \left|\sum_{i\in I_j}\eps_i\right|\le1 .
\]
Then
\[
  \left\|\sum_i\eps_i u_{a(i)}\right\|_E
  \le \sum_{j=1}^M
       \left|\sum_{i\in I_j}\eps_i\right|\|u_j\|_E
  \le M,
\]
while
\[
  \left\|\sum_i\eps_i(x_i-u_{a(i)})\right\|_E\le n\eta .
\]
This proves the quantitative estimate.

For the qualitative consequence, choose \(\eta_k\downarrow0\) with
\(\eta_k\le1/k\), put \(M_k=N_{\mathcal K}(\eta_k)\), and set
\[
  m_k=\left\lceil k\max\{1,M_k\}\right\rceil .
\]
For sufficiently large \(n\), let \(k(n)\) be the largest \(k\) with
\(m_k\le n\).  Then \(k(n)\to\infty\), and the quantitative estimate with
\(\eta=\eta_{k(n)}\) gives
\[
  \frac{\Delta_n^E(\mathcal K)}{n}
  \le \frac1{k(n)}+\eta_{k(n)}
  \le \frac2{k(n)}\to0 .
\]
\end{proof}

\begin{corollary}[Banach entropy rates]
\label{cor:banach-entropy-rates}
Let \(E\), \(\mathcal K\), and \(N_{\mathcal K}\) be as in
Corollary~\ref{cor:precompact-banach-class}.  Then
\[
  \Delta_n^E(\mathcal K)
  \le \inf_{0<\eta\le1}\bigl\{N_{\mathcal K}(\eta)+n\eta\bigr\}.
\]
Consequently:
\begin{enumerate}[label=\textup{(\roman*)}]
\item If \(N_{\mathcal K}(\eta)\le A\eta^{-d}\) for
\(0<\eta\le1\), with \(A\ge1\) and \(d>0\), then
\[
  \Delta_n^E(\mathcal K)
  \le 2A^{1/(d+1)}\,n^{d/(d+1)}
  \qquad(n\ge1).
\]
\item If \(N_{\mathcal K}(\eta)\le \exp(A\eta^{-q})\) for
\(0<\eta\le1\), with \(A,q>0\), then
\[
  \Delta_n^E(\mathcal K)
  \le C_{A,q}\,\frac{n}{(\log(e+n))^{1/q}}
  \qquad(n\ge1).
\]
\item More generally, if
\[
  N_{\mathcal K}(\eta)
  \le
  \exp\!\left(A\eta^{-q}\bigl(\log(e/\eta)\bigr)^\gamma\right)
  \qquad(0<\eta\le1/2),
\]
with \(A,q>0\) and \(\gamma\in\mathbb R\), then
\[
  \Delta_n^E(\mathcal K)
  \le C_{A,q,\gamma}\,
  n\,(\log(e+n))^{-1/q}
  \bigl(\log\log(e^e+n)\bigr)^{\gamma/q}
  \qquad(n\ge2).
\]
\end{enumerate}
\end{corollary}

\begin{proof}
The infimum estimate is Corollary~\ref{cor:precompact-banach-class}, restricted
to \(0<\eta\le1\).

For the polynomial bound, use the trivial estimate
\(\Delta_n^E(\mathcal K)\le n\) when \(n<A\).  Since
\[
  n\le A^{1/(d+1)}n^{d/(d+1)}
  \qquad(n<A),
\]
this is enough in the small-\(n\) range.  When \(n\ge A\), choose
\(\eta=(A/n)^{1/(d+1)}\).  Then \(\eta\le1\), and the infimum estimate gives
\[
  \Delta_n^E(\mathcal K)
  \le A\eta^{-d}+n\eta
  =2A^{1/(d+1)}n^{d/(d+1)} .
\]

For the stretched-exponential bound, enlarge \(C_{A,q}\) to cover the bounded
range where \((2A/\log(e+n))^{1/q}>1\).  In the remaining range take
\[
  \eta=\left(\frac{2A}{\log(e+n)}\right)^{1/q}.
\]
Then \(0<\eta\le1\) and
\[
  N_{\mathcal K}(\eta)\le \exp(A\eta^{-q})
  =\exp\!\left(\frac12\log(e+n)\right)
  \le C_q\,\frac{n}{(\log(e+n))^{1/q}},
\]
after increasing \(C_q\) for small \(n\).  The term \(n\eta\) has the displayed
order with constant depending on \(A\) and \(q\), so the claim follows.

For the signed-log refinement, write
\[
  L_n=\log(e+n),\qquad M_n=\log\log(e^e+n).
\]
After enlarging constants, assume
\[
  \eta=B L_n^{-1/q}M_n^{\gamma/q}\le\frac12
\]
with \(B\ge1\) to be chosen.  Then
\[
  c_{q,\gamma}M_n\le\log(e/\eta)\le C_{q,\gamma}M_n
\]
for all sufficiently large \(n\), and hence
\[
  A\eta^{-q}\bigl(\log(e/\eta)\bigr)^\gamma
  \le C_{A,q,\gamma}B^{-q}L_n .
\]
Choosing \(B\) large enough makes the last quantity at most \(L_n/2\).  Thus
\[
  N_{\mathcal K}(\eta)\le e^{L_n/2}\le C\sqrt n.
\]
Corollary~\ref{cor:precompact-banach-class} gives
\[
  \Delta_n^E(\mathcal K)
  \le C\sqrt n+nB L_n^{-1/q}M_n^{\gamma/q}.
\]
The second term dominates \(\sqrt n\) for large \(n\), and bounded \(n\) are
absorbed into the constant.
\end{proof}

\begin{corollary}[Compact operator images]
\label{cor:compact-operator-image}
Let \(X\) and \(E\) be normed spaces, and let \(T:X\to E\) be a compact linear
operator with \(\|T\|\le1\).  Set
\[
  \Delta_n(T):=
  \sup_{\|x_1\|_X,\ldots,\|x_n\|_X\le1}
  \min_{\eps_i\in\{-1,1\}}
  \left\|\sum_{i=1}^n\eps_iTx_i\right\|_E .
\]
Let \(N_T(\eta)\) be the least size of an \(\eta\)-net for \(T(B_X)\), with
centers in \(T(B_X)\), where \(B_X=\{x:\|x\|_X\le1\}\).  Then
\[
  \Delta_n(T)\le N_T(\eta)+n\eta
  \qquad(\eta>0),
\]
and in particular
\[
  \Delta_n(T)=o(n).
\]
Moreover, all rate estimates in Corollary~\ref{cor:banach-entropy-rates} apply
with \(N_{\mathcal K}\) replaced by \(N_T\).
\end{corollary}

\begin{proof}
The set \(\mathcal K=T(B_X)\) is totally bounded because \(T\) is compact, and
it is contained in the unit ball of \(E\) because \(\|T\|\le1\).  Applying
Corollaries~\ref{cor:precompact-banach-class} and
\ref{cor:banach-entropy-rates} to this \(\mathcal K\) gives all assertions.
\end{proof}

\begin{corollary}[Dyadic entropy-number rates for compact images]
\label{cor:compact-operator-entropy-number-rates}
Let \(E\) be a real normed space and let \(\mathcal K\subset E\) be totally
bounded.  Put
\[
  R_{\mathcal K}:=\sup_{x\in\mathcal K}\|x\|_E<\infty
\]
and define \(\Delta_n^E(\mathcal K)\) as in
Corollary~\ref{cor:precompact-banach-class}.  Define the internal dyadic
entropy numbers
\[
  e_k^{\mathrm{int}}(\mathcal K)
  :=
  \inf\{\eta>0:N_{\mathcal K}(\eta)\le2^{k-1}\},
  \qquad k\ge1 .
\]
Then, for every \(k\ge1\),
\[
  \Delta_n^E(\mathcal K)
  \le R_{\mathcal K}2^{k-1}
      +n\,e_k^{\mathrm{int}}(\mathcal K).
\]
Consequently:
\begin{enumerate}[label=\textup{(\roman*)}]
\item If \(R_{\mathcal K}\le R\) and
\(e_k^{\mathrm{int}}(\mathcal K)\le A k^{-s}\) for all \(k\ge1\),
with \(A,R,s>0\), then
\[
  \Delta_n^E(\mathcal K)
  \le C_{A,R,s}\,n(\log(e+n))^{-s}
  \qquad(n\ge2).
\]
\item More generally, if
\(R_{\mathcal K}\le R\) and
\[
  e_k^{\mathrm{int}}(\mathcal K)
  \le A k^{-s}\bigl(\log(e+k)\bigr)^\beta
  \qquad(k\ge1),
\]
with \(A,R,s>0\) and \(\beta\in\mathbb R\), then
\[
  \Delta_n^E(\mathcal K)
  \le C_{A,R,s,\beta}\,
  n(\log(e+n))^{-s}
  \bigl(\log\log(e^e+n)\bigr)^\beta
  \qquad(n\ge2).
\]
\item If \(R_{\mathcal K}\le R\) and
\(e_k^{\mathrm{int}}(\mathcal K)\le A e^{-bk}\) for all \(k\ge1\),
with \(A,R,b>0\), then, with
\(\alpha_b=\log2/(b+\log2)\in(0,1)\),
\[
  \Delta_n^E(\mathcal K)
  \le C_{A,R,b}\,n^{\alpha_b}
  \qquad(n\ge2).
\]
\end{enumerate}
For a compact operator \(T:X\to E\), the same first estimate holds with
\(\mathcal K=T(B_X)\), \(R_{\mathcal K}\le\|T\|\), and
\(e_k^{\mathrm{int}}(T):=e_k^{\mathrm{int}}(T(B_X))\); in particular the three
rates hold uniformly for compact operators with \(\|T\|\le R\).
\end{corollary}

\begin{proof}
Choose any \(\eta>e_k^{\mathrm{int}}(\mathcal K)\) and an internal
\(\eta\)-net \(u_1,\ldots,u_M\in\mathcal K\) with \(M\le2^{k-1}\).  Given
\(x_1,\ldots,x_n\in\mathcal K\), assign \(x_i\) to a net point \(u_{a(i)}\)
within distance \(\eta\).  On each fiber choose signs with imbalance at most
one.  Then
\[
  \left\|\sum_i\eps_i u_{a(i)}\right\|_E
  \le \sum_{j=1}^M
       \left|\sum_{i:a(i)=j}\eps_i\right|\|u_j\|_E
  \le R_{\mathcal K}M,
\]
and the approximation error contributes at most \(n\eta\).  Letting
\(\eta\downarrow e_k^{\mathrm{int}}(\mathcal K)\) proves the first display.

For the polynomial and log-refined bounds, put
\[
  k_n=\left\lceil\frac{\log(e+n)}{2\log2}\right\rceil .
\]
Then \(R_{\mathcal K}2^{k_n-1}\le C_R\sqrt{e+n}\) and
\(k_n\asymp\log(e+n)\).  Substituting \(k=k_n\) in the first display gives the
two asserted logarithmic rates; the \(R_{\mathcal K}2^{k_n-1}\) term is absorbed because
\[
  \sqrt n=o\!\left(
    n(\log(e+n))^{-s}
    \bigl(\log\log(e^e+n)\bigr)^\beta
  \right)
\]
for every fixed \(s>0\) and \(\beta\in\mathbb R\).

For the exponential entropy-number bound, take, for large \(n\),
\[
  k_n=\left\lfloor\frac{\log(e+n)}{b+\log2}\right\rfloor ,
\]
and enlarge the constant to cover the finite range where \(k_n<1\).  Then
\[
  R_{\mathcal K}2^{k_n}\le C_{R,b} n^{\alpha_b},
  \qquad
  n e^{-bk_n}\le C_b n^{\alpha_b},
  \qquad
  \alpha_b=\frac{\log2}{b+\log2}.
\]
The first display gives the claimed \(n^{\alpha_b}\) rate.  The operator
version is the specialization \(\mathcal K=T(B_X)\).
\end{proof}

\begin{corollary}[External entropy numbers, unions, and sums]
\label{cor:banach-external-entropy-union-sum}
Let \(E\) be a real normed space and let \(\mathcal K\subset E\) be totally
bounded.  Define the usual external dyadic entropy numbers
\[
  e_k(\mathcal K;E)
  :=
  \inf\Bigl\{\eta>0:
    \mathcal K\subset\bigcup_{j=1}^{2^{k-1}}(y_j+\eta B_E)
    \text{ for some }y_j\in E\Bigr\}.
\]
Then
\[
  e_k^{\mathrm{int}}(\mathcal K)\le 2e_k(\mathcal K;E)
  \qquad(k\ge1),
\]
so all estimates of
Corollary~\ref{cor:compact-operator-entropy-number-rates} remain valid, up to
constant changes, with the standard external entropy numbers \(e_k(\mathcal K;E)\)
in place of the internal ones.

Moreover, if \(\mathcal K=\bigcup_{\ell=1}^r\mathcal K_\ell\), then
\[
  e_{k+\lceil\log_2 r\rceil}(\mathcal K;E)
  \le \max_{1\le\ell\le r}e_k(\mathcal K_\ell;E)
  \qquad(k\ge1).
\]
If instead
\[
  \mathcal K\subset \mathcal K_1+\cdots+\mathcal K_r
  :=\{x_1+\cdots+x_r:x_\ell\in\mathcal K_\ell\},
\]
then, for every \(k_1,\ldots,k_r\ge1\),
\[
  e_{1+\sum_{\ell=1}^r(k_\ell-1)}(\mathcal K;E)
  \le \sum_{\ell=1}^r e_{k_\ell}(\mathcal K_\ell;E).
\]
Consequently finite unions and finite Minkowski sums of bounded compact classes
preserve the polynomial and log-refined dyadic entropy decay classes, hence the
corresponding discrepancy rates above; exponential dyadic entropy decay is
also preserved, with the exponent for a finite Minkowski sum allowed to depend
on the number of summands.
\end{corollary}

\begin{proof}
For the internal--external comparison, fix \(\eta>e_k(\mathcal K;E)\) and cover
\(\mathcal K\) by at most \(2^{k-1}\) balls \(y_j+\eta B_E\) with arbitrary
centers \(y_j\in E\).  Discard empty balls.  For every remaining ball choose
one point \(u_j\in\mathcal K\cap(y_j+\eta B_E)\).  If
\(x\in\mathcal K\cap(y_j+\eta B_E)\), then
\(\|x-u_j\|_E\le2\eta\).  Thus these chosen points form an internal
\(2\eta\)-net of the same cardinality.  Letting
\(\eta\downarrow e_k(\mathcal K;E)\) proves
\(e_k^{\mathrm{int}}(\mathcal K)\le2e_k(\mathcal K;E)\).

For a finite union, choose \(e_k(\mathcal K_\ell;E)\)-near covers of each
\(\mathcal K_\ell\).  Combining the \(r\) covers uses at most
\[
  r\,2^{k-1}\le 2^{k+\lceil\log_2r\rceil-1}
\]
balls and radius \(\max_\ell e_k(\mathcal K_\ell;E)\), up to an arbitrary
positive enlargement.  Letting that enlargement tend to zero proves the union
bound.

For a finite Minkowski sum, choose covers of \(\mathcal K_\ell\) with radii
arbitrarily close to \(e_{k_\ell}(\mathcal K_\ell;E)\).  Sums of the cover
centers cover \(\mathcal K_1+\cdots+\mathcal K_r\) with radius equal to the sum
of the radii and with at most
\[
  \prod_{\ell=1}^r2^{k_\ell-1}
  =2^{\sum_\ell(k_\ell-1)}
  =2^{(1+\sum_\ell(k_\ell-1))-1}
\]
balls.  This proves the displayed sum bound.  The final sentence follows by
taking the same dyadic scale in each summand, up to fixed integer shifts.
\end{proof}

\begin{corollary}[Summable Minkowski series]
\label{cor:banach-summable-minkowski}
Let \(E\) be a real Banach space.  For \(j\ge1\), let
\(\mathcal K_j\subset E\) be totally bounded and put
\[
  R_j:=\sup_{x\in\mathcal K_j}\|x\|_E<\infty .
\]
Let \(a_j\ge0\) satisfy \(\sum_{j\ge1}a_jR_j<\infty\), and set
\[
  \mathcal K
  :=
  \left\{\sum_{j=1}^\infty a_jx_j:
    x_j\in\mathcal K_j\right\}.
\]
Then the displayed series converge absolutely, \(\mathcal K\) is totally
bounded, and for every \(r\ge1\) and \(k_1,\ldots,k_r\ge1\),
\[
  e_{1+\sum_{j=1}^r(k_j-1)}(\mathcal K;E)
  \le
  \sum_{j=1}^r a_j e_{k_j}(\mathcal K_j;E)
  +\sum_{j>r}a_jR_j .
\]
Consequently, suppose in addition that \(R_j\le R\) and
\(a_j\le A_0e^{-cj}\) for all \(j\), with \(A_0,R,c>0\).
\begin{enumerate}[label=\textup{(\roman*)}]
\item If, uniformly in \(j\),
\[
  e_k(\mathcal K_j;E)
  \le B k^{-s}\bigl(\log(e+k)\bigr)^\beta
  \qquad(k\ge1),
\]
where \(B,s>0\) and \(\beta\in\R\), and
\(\beta_+=\max(\beta,0)\), then
\[
  e_m(\mathcal K;E)
  \le C m^{-s}\bigl(\log(e+m)\bigr)^{s+\beta_+}
  \qquad(m\ge2),
\]
and
\[
  \Delta_n^E(\mathcal K)
  \le C'\,
  n(\log(e+n))^{-s}
  \bigl(\log\log(e^e+n)\bigr)^{s+\beta_+}
  \qquad(n\ge2).
\]
\item If, uniformly in \(j\),
\[
  e_k(\mathcal K_j;E)\le B e^{-bk}\qquad(k\ge1),
\]
with \(B,b>0\), then there are \(\gamma,\gamma'>0\) such that
\[
  e_m(\mathcal K;E)\le C e^{-\gamma\sqrt m}\qquad(m\ge2),
\]
and
\[
  \Delta_n^E(\mathcal K)
  \le C'\,n\exp\!\bigl(-\gamma'\sqrt{\log(e+n)}\bigr)
  \qquad(n\ge2).
\]
\end{enumerate}
All constants depend only on the displayed parameters.
\end{corollary}

\begin{proof}
The absolute convergence follows from
\[
  \sum_{j\ge1}\|a_jx_j\|_E\le\sum_{j\ge1}a_jR_j<\infty,
\]
and completeness of \(E\).  Fix \(r\) and \(k_1,\ldots,k_r\).  For each
\(1\le j\le r\), choose an external cover of \(\mathcal K_j\) by at most
\(2^{k_j-1}\) balls with radius arbitrarily close to
\(e_{k_j}(\mathcal K_j;E)\).  Sums of the first \(r\) chosen centers cover all
partial sums \(\sum_{j=1}^r a_jx_j\) with radius arbitrarily close to
\(\sum_{j=1}^r a_j e_{k_j}(\mathcal K_j;E)\), and the number of such centers is
\[
  \prod_{j=1}^r2^{k_j-1}
  =
  2^{(1+\sum_{j=1}^r(k_j-1))-1}.
\]
The remaining tail has norm at most \(\sum_{j>r}a_jR_j\).  Letting the cover
enlargements tend to zero gives the entropy-number estimate.  Taking \(r\) so
that the tail is small and then taking the finitely many \(k_j\) large proves
total boundedness of \(\mathcal K\).

Assume now \(R_j\le R\) and \(a_j\le A_0e^{-cj}\).  For the polynomial/log
case, choose
\[
  r=\left\lceil C_0\log(e+m)\right\rceil,
  \qquad
  k=1+\left\lfloor\frac{m-1}{r}\right\rfloor ,
\]
with \(C_0\) large enough that \(\sum_{j>r}a_jR_j\le C m^{-s}\).  For bounded
\(m\) the claim is absorbed into the constant, so assume \(m\) is large enough
that \(k\ge1\) and \(k\asymp m/r\).  The entropy estimate with
\(k_1=\cdots=k_r=k\) gives an index at most \(m\), hence
\[
  e_m(\mathcal K;E)
  \le C k^{-s}\bigl(\log(e+k)\bigr)^\beta
      +C m^{-s}.
\]
If \(\beta\ge0\), then
\[
  k^{-s}\bigl(\log(e+k)\bigr)^\beta
  \le
  C m^{-s}\bigl(\log(e+m)\bigr)^{s+\beta}.
\]
If \(\beta<0\), the logarithmic factor is at most a constant and
\(k^{-s}\le C m^{-s}(\log(e+m))^s\).  This proves the displayed entropy bound
with exponent \(s+\beta_+\).  The discrepancy estimate follows from
Corollaries~\ref{cor:compact-operator-entropy-number-rates} and
\ref{cor:banach-external-entropy-union-sum}.

For the exponential component-entropy case, choose
\[
  r=\left\lceil C_1\sqrt m\right\rceil,
  \qquad
  k=1+\left\lfloor\frac{m-1}{r}\right\rfloor .
\]
For large \(m\), \(k\asymp\sqrt m\), while the tail is \(O(e^{-c r})\).  The
same entropy estimate therefore gives
\[
  e_m(\mathcal K;E)\le C e^{-\gamma\sqrt m}
\]
after decreasing \(\gamma>0\) and enlarging \(C\) to cover small \(m\).  Also
\[
  R_{\mathcal K}:=\sup_{x\in\mathcal K}\|x\|_E
  \le\sum_{j\ge1}a_jR_j<\infty .
\]
The external-to-internal comparison gives
\(e_q^{\mathrm{int}}(\mathcal K)\le C e^{-\gamma\sqrt q}\).  Insert
\[
  q=\left\lfloor\frac{\log(e+n)}{2\log2}\right\rfloor
\]
in the first display of
Corollary~\ref{cor:compact-operator-entropy-number-rates}; after changing
constants and covering bounded \(n\),
\[
  \Delta_n^E(\mathcal K)
  \le C\sqrt n+C n e^{-\gamma\sqrt q}
  \le C' n\exp\!\bigl(-\gamma'\sqrt{\log(e+n)}\bigr).
\]
\end{proof}

\begin{corollary}[Profile form for weighted Minkowski series]
\label{cor:banach-weighted-profile}
In the setting of Corollary~\ref{cor:banach-summable-minkowski}, suppose that
there are numbers \(b_j\ge0\) and a function
\(\varphi:\mathbb N\to[0,\infty)\) such that
\[
  e_k(\mathcal K_j;E)\le b_j\varphi(k)
  \qquad(j,k\ge1).
\]
Put
\[
  S_r:=\sum_{j=1}^r a_jb_j,
  \qquad
  T_r:=\sum_{j>r}a_jR_j,
  \qquad
  R_{\mathcal K}:=\sum_{j\ge1}a_jR_j .
\]
Then, for every \(m\ge2\),
\[
  e_m(\mathcal K;E)
  \le
  \inf_{1\le r\le m}
  \left\{
    S_r\,
    \varphi\!\left(1+\left\lfloor\frac{m-1}{r}\right\rfloor\right)
    +T_r
  \right\}.
\]
Consequently, for every \(q\ge1\),
\[
  \Delta_n^E(\mathcal K)
  \le
  R_{\mathcal K}2^{q-1}
  +2n\inf_{1\le r\le q}
  \left\{
    S_r\,
    \varphi\!\left(1+\left\lfloor\frac{q-1}{r}\right\rfloor\right)
    +T_r
  \right\}.
\]
In particular, with
\[
  q_n=\left\lfloor\frac{\log(e+n)}{2\log2}\right\rfloor
  \qquad(n\ge2),
\]
one has
\[
  \Delta_n^E(\mathcal K)
  \le
  C R_{\mathcal K}\sqrt n
  +2n\inf_{1\le r\le q_n}
  \left\{
    S_r\,
    \varphi\!\left(1+\left\lfloor\frac{q_n-1}{r}\right\rfloor\right)
    +T_r
  \right\}.
\]
\end{corollary}

\begin{proof}
Fix \(m\ge2\) and \(1\le r\le m\), and set
\[
  k=1+\left\lfloor\frac{m-1}{r}\right\rfloor .
\]
Then \(1+r(k-1)\le m\).  Applying
Corollary~\ref{cor:banach-summable-minkowski} with
\(k_1=\cdots=k_r=k\), and using monotonicity of entropy numbers in the index,
gives
\[
  e_m(\mathcal K;E)
  \le e_{1+r(k-1)}(\mathcal K;E)
  \le
  \sum_{j=1}^r a_j e_k(\mathcal K_j;E)+T_r
  \le S_r\varphi(k)+T_r .
\]
Taking the infimum over \(r\) proves the entropy profile bound.

The external-to-internal comparison from
Corollary~\ref{cor:banach-external-entropy-union-sum} gives
\[
  e_q^{\mathrm{int}}(\mathcal K)\le2e_q(\mathcal K;E).
\]
Substitute the profile estimate for \(e_q(\mathcal K;E)\) in the first display
of Corollary~\ref{cor:compact-operator-entropy-number-rates}.  Since
\(\sup_{x\in\mathcal K}\|x\|_E\le R_{\mathcal K}\), the displayed discrepancy
bound follows.  The final estimate is the same bound with \(q=q_n\), using
\(2^{q_n-1}\le C\sqrt n\).
\end{proof}

\begin{corollary}[Polynomial profile optimization for weighted series]
\label{cor:banach-weighted-polynomial-profile}
In the setting of Corollary~\ref{cor:banach-weighted-profile}, with \(S_r\) and
\(T_r\) as there, assume that, for some \(A>0\), \(\gamma>0\),
\(\delta\ge0\), and \(\eta,\tau\in\R\),
\[
  S_r\le A r^\delta\bigl(\log(e+r)\bigr)^\eta,
  \qquad
  T_r\le A r^{-\gamma}\bigl(\log(e+r)\bigr)^{-\tau}
  \qquad(r\ge1).
\]
Then the following hold.
\begin{enumerate}[label=\textup{(\roman*)}]
\item If
\[
  \varphi(k)\le B k^{-s}\bigl(\log(e+k)\bigr)^\beta
  \qquad(k\ge1),
\]
with \(B,s>0\) and \(\beta\in\R\), put
\[
  D=s+\delta+\gamma,\qquad
  \sigma=\frac{s\gamma}{D},\qquad
  \lambda=\frac{\gamma(\beta+\eta)-\tau(s+\delta)}{D}.
\]
Then
\[
  e_m(\mathcal K;E)
  \le C m^{-\sigma}\bigl(\log(e+m)\bigr)^\lambda
  \qquad(m\ge2),
\]
and
\[
  \Delta_n^E(\mathcal K)
  \le C'\,
  n(\log(e+n))^{-\sigma}
  \bigl(\log\log(e^e+n)\bigr)^\lambda
  \qquad(n\ge2).
\]
\item If
\[
  \varphi(k)\le B e^{-bk}\qquad(k\ge1)
\]
with \(B,b>0\), then
\[
  e_m(\mathcal K;E)
  \le C m^{-\gamma}\bigl(\log(e+m)\bigr)^{\gamma-\tau}
  \qquad(m\ge2),
\]
and
\[
  \Delta_n^E(\mathcal K)
  \le C'\,
  n(\log(e+n))^{-\gamma}
  \bigl(\log\log(e^e+n)\bigr)^{\gamma-\tau}
  \qquad(n\ge2).
\]
\end{enumerate}
The constants depend only on the displayed parameters and on
\(R_{\mathcal K}\) in the discrepancy bounds.
\end{corollary}

\begin{proof}
For the polynomial/log profile, write \(L_m=\log(e+m)\) and
\(D=s+\delta+\gamma\).  For large \(m\), choose
\[
  r=\left\lceil
      m^{s/D}L_m^{-(\beta+\eta+\tau)/D}
    \right\rceil,
  \qquad
  k=1+\left\lfloor\frac{m-1}{r}\right\rfloor .
\]
The finite range of \(m\) is absorbed into the constant.  Since
\(0<s/D<1\), for large \(m\) one has \(1\le r\le m\), \(k\asymp m/r\), and
\[
  \log(e+r)\asymp L_m,\qquad \log(e+k)\asymp L_m .
\]
Corollary~\ref{cor:banach-weighted-profile} gives
\[
  e_m(\mathcal K;E)
  \le
  C m^{-s}r^{s+\delta}L_m^{\beta+\eta}
  +C r^{-\gamma}L_m^{-\tau}.
\]
With the displayed choice of \(r\), both terms are bounded by
\[
  C m^{-s\gamma/D}
    L_m^{(\gamma(\beta+\eta)-\tau(s+\delta))/D},
\]
which proves the entropy estimate.  The discrepancy estimate follows from
Corollaries~\ref{cor:compact-operator-entropy-number-rates}
and~\ref{cor:banach-external-entropy-union-sum}; the term
\(R_{\mathcal K}2^{q-1}\) is absorbed by taking
\(q\asymp\log(e+n)\), since \(\sigma>0\).

For the exponential profile, choose \(D_0>2(\gamma+\delta)/b\) and set
\[
  r=\left\lceil\frac{m}{D_0\log(e+m)}\right\rceil,
  \qquad
  k=1+\left\lfloor\frac{m-1}{r}\right\rfloor .
\]
For large \(m\), \(k\ge(D_0/2)\log(e+m)\), and the profile bound gives
\[
  e_m(\mathcal K;E)
  \le
  C r^\delta\bigl(\log(e+r)\bigr)^\eta e^{-bk}
  +C r^{-\gamma}\bigl(\log(e+r)\bigr)^{-\tau}.
\]
The second term is at most
\[
  C m^{-\gamma}\bigl(\log(e+m)\bigr)^{\gamma-\tau}.
\]
The first is
\[
  \le C
  m^{\delta-bD_0/2}\bigl(\log(e+m)\bigr)^{\eta-\delta},
\]
and is absorbed by the same tail scale because \(bD_0/2>\gamma+\delta\).
The entropy estimate follows, and the discrepancy estimate is again the
dyadic entropy-number bound and the external-to-internal comparison.
\end{proof}

\begin{corollary}[Polynomial weights with growing component amplitudes]
\label{cor:banach-polynomial-amplitude-minkowski}
Let \(E\) be a real Banach space.  For \(j\ge1\), let
\(\mathcal K_j\subset E\) be totally bounded, put
\[
  R_j:=\sup_{x\in\mathcal K_j}\|x\|_E,
\]
and assume \(R_j\le R\).  Let \(a_j\ge0\) satisfy
\[
  a_j\le A_0 j^{-u}\bigl(\log(e+j)\bigr)^{-v}\qquad(j\ge1)
\]
for some \(A_0,R>0\), \(u>1\), and \(v\in\R\).  Define
\[
  \mathcal K
  :=
  \left\{\sum_{j=1}^\infty a_jx_j:
    x_j\in\mathcal K_j\right\}.
\]
Let \(\rho>u-1\) and \(\eta\in\R\).  Then \(\mathcal K\) is totally bounded.
Moreover:
\begin{enumerate}[label=\textup{(\roman*)}]
\item If
\[
  e_k(\mathcal K_j;E)
  \le B j^\rho\bigl(\log(e+j)\bigr)^\eta
       k^{-s}\bigl(\log(e+k)\bigr)^\beta
  \qquad(j,k\ge1),
\]
where \(B,s>0\) and \(\beta\in\R\), then, with
\[
  \sigma=\frac{s(u-1)}{s+\rho},
  \qquad
  \lambda=\frac{(u-1)(\beta+\eta)-v(s+\rho)}{s+\rho},
\]
one has
\[
  e_m(\mathcal K;E)
  \le C m^{-\sigma}\bigl(\log(e+m)\bigr)^\lambda
  \qquad(m\ge2),
\]
and therefore
\[
  \Delta_n^E(\mathcal K)
  \le C'\,
  n(\log(e+n))^{-\sigma}
  \bigl(\log\log(e^e+n)\bigr)^\lambda
  \qquad(n\ge2).
\]
\item If
\[
  e_k(\mathcal K_j;E)
  \le B j^\rho\bigl(\log(e+j)\bigr)^\eta e^{-bk}
  \qquad(j,k\ge1),
\]
with \(B,b>0\), then
\[
  e_m(\mathcal K;E)
  \le C m^{-(u-1)}\bigl(\log(e+m)\bigr)^{u-1-v}
  \qquad(m\ge2),
\]
and hence
\[
  \Delta_n^E(\mathcal K)
  \le C'\,
  n(\log(e+n))^{-(u-1)}
  \bigl(\log\log(e^e+n)\bigr)^{u-1-v}
  \qquad(n\ge2).
\]
\end{enumerate}
All constants depend only on the displayed parameters.
\end{corollary}

\begin{proof}
Since \(u>1\) and the radii are uniformly bounded,
\[
  \sum_{j\ge1}a_jR_j<\infty,
\]
so Corollary~\ref{cor:banach-summable-minkowski} gives total boundedness.  The
regularly varying tail bound gives
\[
  T_r=\sum_{j>r}a_jR_j
  \le C r^{-(u-1)}\bigl(\log(e+r)\bigr)^{-v}.
\]
For the profile theorem take
\[
  b_j=Bj^\rho\bigl(\log(e+j)\bigr)^\eta .
\]
Since \(\rho>u-1\),
\[
  S_r=\sum_{j=1}^r a_jb_j
  \le C r^{\rho-u+1}\bigl(\log(e+r)\bigr)^{\eta-v}.
\]
Apply Corollary~\ref{cor:banach-weighted-polynomial-profile} with
\[
  \gamma=u-1,\qquad
  \delta=\rho-u+1,\qquad
  \eta_0=\eta-v,\qquad
  \tau=v .
\]
For the polynomial/log component profile this gives
\[
  D=s+\delta+\gamma=s+\rho,
  \qquad
  \sigma=\frac{s(u-1)}{s+\rho},
\]
and logarithmic power
\[
  \frac{(u-1)(\beta+\eta-v)-v(s+\rho-u+1)}{s+\rho}
  =
  \frac{(u-1)(\beta+\eta)-v(s+\rho)}{s+\rho}.
\]
This proves the first pair of estimates.  The exponential component profile is
the second part of
Corollary~\ref{cor:banach-weighted-polynomial-profile}, which gives the
tail-limited entropy and discrepancy rates displayed above.
\end{proof}

\begin{corollary}[Stretched-exponentially weighted Minkowski series]
\label{cor:banach-stretched-minkowski}
Let \(E\) be a real Banach space.  For \(j\ge1\), let
\(\mathcal K_j\subset E\) be totally bounded, put
\[
  R_j:=\sup_{x\in\mathcal K_j}\|x\|_E,
\]
and assume \(R_j\le R\).  Let \(a_j\ge0\) satisfy
\[
  a_j\le A_0\exp(-c j^\theta)\qquad(j\ge1)
\]
for some \(A_0,R,c,\theta>0\), and define
\[
  \mathcal K
  :=
  \left\{\sum_{j=1}^\infty a_jx_j:
    x_j\in\mathcal K_j\right\}.
\]
Then \(\mathcal K\) is totally bounded.  Moreover:
\begin{enumerate}[label=\textup{(\roman*)}]
\item If, uniformly in \(j\),
\[
  e_k(\mathcal K_j;E)
  \le B k^{-s}\bigl(\log(e+k)\bigr)^\beta
  \qquad(k\ge1),
\]
where \(B,s>0\) and \(\beta\in\R\), and
\(\beta_+=\max(\beta,0)\), then
\[
  e_m(\mathcal K;E)
  \le C m^{-s}\bigl(\log(e+m)\bigr)^{s/\theta+\beta_+}
  \qquad(m\ge2),
\]
and
\[
  \Delta_n^E(\mathcal K)
  \le C'\,
  n(\log(e+n))^{-s}
  \bigl(\log\log(e^e+n)\bigr)^{s/\theta+\beta_+}
  \qquad(n\ge2).
\]
\item If, uniformly in \(j\),
\[
  e_k(\mathcal K_j;E)\le B e^{-bk}\qquad(k\ge1),
\]
with \(B,b>0\), then there are \(\gamma,\gamma'>0\) such that
\[
  e_m(\mathcal K;E)
  \le C\exp\!\bigl(-\gamma m^{\theta/(1+\theta)}\bigr)
  \qquad(m\ge2),
\]
and
\[
  \Delta_n^E(\mathcal K)
  \le C'\,n
  \exp\!\bigl(-\gamma'(\log(e+n))^{\theta/(1+\theta)}\bigr)
  \qquad(n\ge2).
\]
\end{enumerate}
All constants depend only on the displayed parameters.
\end{corollary}

\begin{proof}
Since \(\sum_j e^{-c j^\theta}<\infty\), the series defining
\(\mathcal K\) is covered by
Corollary~\ref{cor:banach-summable-minkowski}.  We shall also use the tail
bound
\[
  \sum_{j>r}a_jR_j\le C\exp(-c_1r^\theta)\qquad(r\ge1),
\]
for some \(c_1>0\), which follows from the integral estimate for
\(\sum_{j>r}e^{-c j^\theta}\) after decreasing \(c\) and increasing \(C\).

For the polynomial/log component entropy, choose \(C_0\) large and set
\[
  r=\left\lceil C_0(\log(e+m))^{1/\theta}\right\rceil,
  \qquad
  k=1+\left\lfloor\frac{m-1}{r}\right\rfloor .
\]
For bounded \(m\) the claim is absorbed into the constant; for large \(m\),
\(k\asymp m/r\) and \(1+r(k-1)\le m\).  Choosing \(C_0\) large enough makes
\(\sum_{j>r}a_jR_j\le C m^{-s}\).  The finite part is bounded by
\[
  C k^{-s}\bigl(\log(e+k)\bigr)^\beta
  \le
  C m^{-s}\bigl(\log(e+m)\bigr)^{s/\theta+\beta_+},
\]
with the same convention as above for negative \(\beta\).  The entropy
estimate follows, and the discrepancy estimate is
Corollaries~\ref{cor:compact-operator-entropy-number-rates}
and~\ref{cor:banach-external-entropy-union-sum}.

For exponential component entropy, choose
\[
  r=\left\lceil C_1 m^{1/(1+\theta)}\right\rceil,
  \qquad
  k=1+\left\lfloor\frac{m-1}{r}\right\rfloor .
\]
Then \(k\asymp m^{\theta/(1+\theta)}\) and
\(r^\theta\asymp m^{\theta/(1+\theta)}\).  The finite part
\(C e^{-bk}\) and the tail \(C e^{-c_1r^\theta}\) are both bounded by
\(C\exp(-\gamma m^{\theta/(1+\theta)})\), after decreasing \(\gamma>0\) and
covering small \(m\).  Finally, the first display of
Corollary~\ref{cor:compact-operator-entropy-number-rates}, applied with
\[
  q=\left\lfloor\frac{\log(e+n)}{2\log2}\right\rfloor ,
\]
gives
\[
  \Delta_n^E(\mathcal K)
  \le
  C\sqrt n+
  Cn\exp\!\bigl(-\gamma q^{\theta/(1+\theta)}\bigr),
\]
and the \(\sqrt n\) term is absorbed by the displayed bound after decreasing
\(\gamma'\) and enlarging \(C'\).
\end{proof}

\begin{corollary}[Polynomially weighted Minkowski series]
\label{cor:banach-polynomial-minkowski}
Let \(E\) be a real Banach space.  For \(j\ge1\), let
\(\mathcal K_j\subset E\) be totally bounded, let
\[
  R_j:=\sup_{x\in\mathcal K_j}\|x\|_E,
\]
and assume \(R_j\le R\).  Let \(a_j\ge0\) satisfy
\[
  a_j\le A_0 j^{-u}\qquad(j\ge1)
\]
for some \(A_0,R>0\) and \(u>1\), and define
\[
  \mathcal K
  :=
  \left\{\sum_{j=1}^\infty a_jx_j:
    x_j\in\mathcal K_j\right\}.
\]
Then \(\mathcal K\) is totally bounded.  Moreover:
\begin{enumerate}[label=\textup{(\roman*)}]
\item If, uniformly in \(j\),
\[
  e_k(\mathcal K_j;E)
  \le B k^{-s}\bigl(\log(e+k)\bigr)^\beta
  \qquad(k\ge1),
\]
where \(B,s>0\) and \(\beta\in\R\), then, with
\[
  \sigma=\frac{s(u-1)}{s+u-1},
  \qquad
  \beta_+=\max(\beta,0),
\]
one has
\[
  e_m(\mathcal K;E)
  \le C m^{-\sigma}\bigl(\log(e+m)\bigr)^{\beta_+}
  \qquad(m\ge2),
\]
and therefore
\[
  \Delta_n^E(\mathcal K)
  \le C'\,
  n(\log(e+n))^{-\sigma}
  \bigl(\log\log(e^e+n)\bigr)^{\beta_+}
  \qquad(n\ge2).
\]
\item If, uniformly in \(j\),
\[
  e_k(\mathcal K_j;E)\le B e^{-bk}\qquad(k\ge1),
\]
with \(B,b>0\), then
\[
  e_m(\mathcal K;E)
  \le C m^{-(u-1)}\bigl(\log(e+m)\bigr)^{u-1}
  \qquad(m\ge2),
\]
and hence
\[
  \Delta_n^E(\mathcal K)
  \le C'\,
  n(\log(e+n))^{-(u-1)}
  \bigl(\log\log(e^e+n)\bigr)^{u-1}
  \qquad(n\ge2).
\]
\end{enumerate}
All constants depend only on the displayed parameters.
\end{corollary}

\begin{proof}
Since \(u>1\),
\[
  \sum_{j\ge1}a_jR_j\le A_0R\sum_{j\ge1}j^{-u}<\infty .
\]
Thus Corollary~\ref{cor:banach-summable-minkowski} gives total boundedness and
the truncation estimate.  We shall use
\[
  \sum_{j=1}^r a_j\le C_{A_0,u},
  \qquad
  \sum_{j>r}a_jR_j\le C_{A_0,R,u}r^{1-u}.
\]

For the polynomial/log component entropy, put
\[
  r=\left\lceil m^{s/(s+u-1)}\right\rceil,
  \qquad
  k=1+\left\lfloor\frac{m-1}{r}\right\rfloor .
\]
After enlarging constants to cover bounded \(m\), \(1\le r\le m\) and
\(k\asymp m/r\).  The truncation estimate with \(k_1=\cdots=k_r=k\) has
entropy index \(1+r(k-1)\le m\).  Therefore
\[
  e_m(\mathcal K;E)
  \le C k^{-s}\bigl(\log(e+k)\bigr)^\beta
      +C r^{1-u}.
\]
If \(\beta\ge0\), then
\((\log(e+k))^\beta\le(\log(e+m))^\beta\); if \(\beta<0\), this logarithmic
factor is at most \(1\).  Hence
\[
  k^{-s}\bigl(\log(e+k)\bigr)^\beta
  \le
  C m^{-s}r^s\bigl(\log(e+m)\bigr)^{\beta_+}.
\]
The chosen \(r\) makes \(m^{-s}r^s\) and \(r^{1-u}\) both bounded by
\(C m^{-\sigma}\), which proves the entropy estimate.  The discrepancy bound
then follows from Corollaries~\ref{cor:compact-operator-entropy-number-rates}
and~\ref{cor:banach-external-entropy-union-sum}.

For exponential component entropy, choose \(D>2(u-1)/b\) and set
\[
  r=\left\lceil\frac{m}{D\log(e+m)}\right\rceil,
  \qquad
  k=1+\left\lfloor\frac{m-1}{r}\right\rfloor .
\]
For large \(m\), \(k\ge (D/2)\log(e+m)\), while
\[
  r^{1-u}\le C m^{-(u-1)}\bigl(\log(e+m)\bigr)^{u-1}.
\]
The finite part is
\[
  C e^{-bk}
  \le C m^{-bD/2},
\]
which is absorbed by the displayed tail bound because \(bD/2>u-1\).  The
entropy estimate follows, and the discrepancy bound is again
Corollaries~\ref{cor:compact-operator-entropy-number-rates}
and~\ref{cor:banach-external-entropy-union-sum}, now applied with polynomial
exponent \(u-1\) and logarithmic power \(u-1\).
\end{proof}

\begin{corollary}[Log-balanced polynomially weighted Minkowski series]
\label{cor:banach-log-polynomial-minkowski}
Let \(E\) be a real Banach space.  For \(j\ge1\), let
\(\mathcal K_j\subset E\) be totally bounded, put
\[
  R_j:=\sup_{x\in\mathcal K_j}\|x\|_E,
\]
and assume \(R_j\le R\).  Let \(a_j\ge0\) satisfy
\[
  a_j\le A_0 j^{-u}\bigl(\log(e+j)\bigr)^{-v}\qquad(j\ge1)
\]
for some \(A_0,R>0\), \(u>1\), and \(v\in\R\), and define
\[
  \mathcal K
  :=
  \left\{\sum_{j=1}^\infty a_jx_j:
    x_j\in\mathcal K_j\right\}.
\]
Put \(\gamma=u-1\).  Then \(\mathcal K\) is totally bounded.  Moreover:
\begin{enumerate}[label=\textup{(\roman*)}]
\item If, uniformly in \(j\),
\[
  e_k(\mathcal K_j;E)
  \le B k^{-s}\bigl(\log(e+k)\bigr)^\beta
  \qquad(k\ge1),
\]
where \(B,s>0\) and \(\beta\in\R\), then, with
\[
  \sigma=\frac{s\gamma}{s+\gamma},
  \qquad
  \lambda=\frac{\gamma\beta-sv}{s+\gamma},
\]
one has
\[
  e_m(\mathcal K;E)
  \le C m^{-\sigma}\bigl(\log(e+m)\bigr)^\lambda
  \qquad(m\ge2),
\]
and therefore
\[
  \Delta_n^E(\mathcal K)
  \le C'\,
  n(\log(e+n))^{-\sigma}
  \bigl(\log\log(e^e+n)\bigr)^\lambda
  \qquad(n\ge2).
\]
\item If, uniformly in \(j\),
\[
  e_k(\mathcal K_j;E)\le B e^{-bk}\qquad(k\ge1),
\]
with \(B,b>0\), then
\[
  e_m(\mathcal K;E)
  \le C m^{-\gamma}\bigl(\log(e+m)\bigr)^{\gamma-v}
  \qquad(m\ge2),
\]
and hence
\[
  \Delta_n^E(\mathcal K)
  \le C'\,
  n(\log(e+n))^{-\gamma}
  \bigl(\log\log(e^e+n)\bigr)^{\gamma-v}
  \qquad(n\ge2).
\]
\end{enumerate}
All constants depend only on the displayed parameters.
\end{corollary}

\begin{proof}
The dyadic, or integral, estimate for regularly varying tails gives, for
\(\gamma=u-1>0\),
\[
  \sum_{j=1}^r a_j\le C,
  \qquad
  \sum_{j>r}a_jR_j
  \le C r^{-\gamma}\bigl(\log(e+r)\bigr)^{-v}
  \qquad(r\ge1).
\]
In particular \(\sum_j a_jR_j<\infty\), so
Corollary~\ref{cor:banach-summable-minkowski} gives total boundedness.  We use
the profile form in Corollary~\ref{cor:banach-weighted-profile}.

For the polynomial/log component entropy, let \(L_m=\log(e+m)\) and, for large
\(m\), choose
\[
  r=\left\lceil
      m^{s/(s+\gamma)}L_m^{-(\beta+v)/(s+\gamma)}
    \right\rceil,
  \qquad
  k=1+\left\lfloor\frac{m-1}{r}\right\rfloor .
\]
The finite range of \(m\) is absorbed into the constant.  For large \(m\),
\(1\le r\le m\), \(k\asymp m/r\), and both
\(\log(e+r)\) and \(\log(e+k)\) are comparable to \(L_m\).  Hence the profile
bound gives
\[
  e_m(\mathcal K;E)
  \le
  C m^{-s}r^s L_m^\beta
  +C r^{-\gamma}L_m^{-v}.
\]
With the displayed choice of \(r\), both terms are bounded by
\[
  C m^{-s\gamma/(s+\gamma)}
    L_m^{(\gamma\beta-sv)/(s+\gamma)}.
\]
This proves the first entropy estimate.  The discrepancy bound follows from
Corollaries~\ref{cor:compact-operator-entropy-number-rates}
and~\ref{cor:banach-external-entropy-union-sum}, applied with polynomial
exponent \(\sigma\) and logarithmic power \(\lambda\).

For exponential component entropy, choose \(D>2\gamma/b\) and set
\[
  r=\left\lceil\frac{m}{D\log(e+m)}\right\rceil,
  \qquad
  k=1+\left\lfloor\frac{m-1}{r}\right\rfloor .
\]
For large \(m\), \(k\ge(D/2)\log(e+m)\), and therefore
\[
  e_m(\mathcal K;E)
  \le C e^{-bk}
      +C r^{-\gamma}\bigl(\log(e+r)\bigr)^{-v}
  \le C m^{-bD/2}
      +C m^{-\gamma}\bigl(\log(e+m)\bigr)^{\gamma-v}.
\]
Since \(bD/2>\gamma\), the first term is absorbed by the second after enlarging
the constant.  This proves the second entropy estimate, and the corresponding
discrepancy estimate follows from the same dyadic entropy-number corollaries
with exponent \(\gamma\) and logarithmic power \(\gamma-v\).
\end{proof}

\begin{corollary}[H\"older images of Banach compact classes]
\label{cor:banach-holder-image-entropy}
Let \(E,F\) be real normed spaces, let \(\mathcal K\subset E\) be totally
bounded, and let \(\Phi:\mathcal K\to F\) be \((L,\alpha)\)-H\"older,
\[
  \|\Phi(x)-\Phi(y)\|_F\le L\|x-y\|_E^\alpha
  \qquad(x,y\in\mathcal K),
\]
with \(0<\alpha\le1\).  Then
\[
  e_k(\Phi(\mathcal K);F)
  \le L\,2^\alpha e_k(\mathcal K;E)^\alpha
  \qquad(k\ge1).
\]
Consequently, if \(R_{\Phi}:=\sup_{x\in\mathcal K}\|\Phi(x)\|_F<\infty\), the
following discrepancy bounds hold for
\(\Delta_n^F(\Phi(\mathcal K))\).
\begin{enumerate}[label=\textup{(\roman*)}]
\item If \(e_k(\mathcal K;E)\le A k^{-s}\) for all \(k\ge1\), with
\(A,s>0\), then
\[
  \Delta_n^F(\Phi(\mathcal K))
  \le C_{A,L,R_\Phi,s,\alpha}\,
     n(\log(e+n))^{-\alpha s}
  \qquad(n\ge2).
\]
\item If
\[
  e_k(\mathcal K;E)
  \le A k^{-s}\bigl(\log(e+k)\bigr)^\beta
  \qquad(k\ge1),
\]
with \(A,s>0\) and \(\beta\in\mathbb R\), then
\[
  \Delta_n^F(\Phi(\mathcal K))
  \le C\,
     n(\log(e+n))^{-\alpha s}
     \bigl(\log\log(e^e+n)\bigr)^{\alpha\beta}
  \qquad(n\ge2),
\]
where \(C\) depends on \(A,L,R_\Phi,s,\alpha,\beta\).
\item If \(e_k(\mathcal K;E)\le A e^{-bk}\) for all \(k\ge1\), with
\(A,b>0\), then
\[
  \Delta_n^F(\Phi(\mathcal K))
  \le C_{A,L,R_\Phi,b,\alpha}\,
     n^{\log2/(\alpha b+\log2)}
  \qquad(n\ge2).
\]
\end{enumerate}
\end{corollary}

\begin{proof}
Fix \(\eta>e_k(\mathcal K;E)\) and cover \(\mathcal K\) by at most \(2^{k-1}\)
external balls \(y_j+\eta B_E\).  As in the proof of
Corollary~\ref{cor:banach-external-entropy-union-sum}, choose one point
\(u_j\in\mathcal K\) from each nonempty ball.  These points form an internal
\(2\eta\)-net for \(\mathcal K\).  Hence the points \(\Phi(u_j)\) cover
\(\Phi(\mathcal K)\) in \(F\) with radius \(L(2\eta)^\alpha\).  Letting
\(\eta\downarrow e_k(\mathcal K;E)\) proves the displayed entropy-number
transfer.

The three rate estimates are now
Corollary~\ref{cor:compact-operator-entropy-number-rates}, together with the
external-to-internal comparison in
Corollary~\ref{cor:banach-external-entropy-union-sum}, applied to the compact
class \(\Phi(\mathcal K)\).  The polynomial exponent \(s\), logarithmic power
\(\beta\), and exponential rate \(b\) are replaced respectively by
\(\alpha s\), \(\alpha\beta\), and \(\alpha b\).
\end{proof}

\begin{corollary}[Finite-rank approximation envelope]
\label{cor:banach-finite-rank-envelope}
Let \(E\) be a real normed space and let \(\mathcal K\) be a totally bounded
subset of the unit ball of \(E\).  For \(m\ge0\), put
\[
  a_m(\mathcal K):=
  \inf_{\dim F\le m}\ \sup_{x\in\mathcal K}\inf_{y\in F}\|x-y\|_E,
\]
where the infimum is over real linear subspaces \(F\subset E\) of dimension at
most \(m\).  Then, for every \(m\ge0\) and \(0<\eta\le1\),
\[
  \Delta_n^E(\mathcal K)
  \le 3\left(1+\frac6\eta\right)^m+n\bigl(\eta+a_m(\mathcal K)\bigr).
\]
Consequently
\[
  \Delta_n^E(\mathcal K)
  \le
  \inf_{\substack{m\ge0\\0<\eta\le1}}
  \left\{3\left(1+\frac6\eta\right)^m
       +n\bigl(\eta+a_m(\mathcal K)\bigr)\right\}.
\]
If \(X,E\) are real normed spaces and \(T:X\to E\) is a compact operator with
\(\|T\|\le1\), the same estimate holds for \(\Delta_n(T)\) with
\[
  a_m(T):=\inf_{\operatorname{rank}S\le m}\|T-S\|_{X\to E}
\]
in place of \(a_m(\mathcal K)\).
\end{corollary}

\begin{proof}
Fix \(0<\delta\le1\), and choose \(F\) with \(\dim F\le m\) and
\[
  \sup_{x\in\mathcal K}\inf_{y\in F}\|x-y\|_E\le a_m(\mathcal K)+\delta .
\]
For each \(x\in\mathcal K\) choose \(y_x\in F\) with
\(\|x-y_x\|_E\le a_m(\mathcal K)+\delta\).  Since \(\mathcal K\) is contained in
the unit ball, \(a_m(\mathcal K)\le1\), and \(\delta\le1\), all such \(y_x\) lie in
\(3B_E\cap F\).  The standard volumetric covering bound in an \(m\)-dimensional
normed space gives an \(\eta\)-net of \(3B_E\cap F\) with at most
\((1+6/\eta)^m\) points.  Thus
\[
  \mathcal K\subset
  \bigcup_{j=1}^M
  \bigl(z_j+(a_m(\mathcal K)+\delta+\eta)B_E\bigr),
  \qquad
  M\le\left(1+\frac6\eta\right)^m,
\]
with centers \(z_j\) lying in \(3B_E\cap F\).  Given
\(x_1,\ldots,x_n\in\mathcal K\), assign each \(x_i\) to a center \(z_{a(i)}\)
within distance \(a_m(\mathcal K)+\delta+\eta\).  On every fiber
\(\{i:a(i)=j\}\), choose signs with imbalance at most one.  Then
\[
  \left\|\sum_i\eps_i z_{a(i)}\right\|_E
  \le 3\left(1+\frac6\eta\right)^m
\]
and the approximation error contributes at most
\(n(a_m(\mathcal K)+\delta+\eta)\).  Letting \(\delta\downarrow0\) proves the
first two displays.

For the operator version, take \(\mathcal K=T(B_X)\).  If \(S\) has rank at most
\(m\), then \(T(B_X)\) lies within \(\|T-S\|\) of the \(m\)-dimensional space
\(\operatorname{range}S\).  Hence \(a_m(\mathcal K)\le a_m(T)\), and the
preceding estimate gives the claim.
\end{proof}

\begin{corollary}[Approximation-number rates for compact images]
\label{cor:compact-operator-approx-rates}
Let \(X,E\) be real normed spaces and let \(T:X\to E\) be a compact operator
with \(\|T\|\le1\).  With \(a_m(T)\) as in
Corollary~\ref{cor:banach-finite-rank-envelope}, the following hold.
\begin{enumerate}[label=\textup{(\roman*)}]
\item If \(a_m(T)\le A m^{-s}\) for all \(m\ge1\), with \(A,s>0\), then
\[
  \Delta_n(T)\le C_{A,s}\,
  n\left(\frac{\log\log(e^e+n)}{\log(e+n)}\right)^s
  \qquad(n\ge2).
\]
\item More generally, if
\[
  a_m(T)\le A m^{-s}\bigl(\log(e+m)\bigr)^\beta
  \qquad(m\ge1),
\]
with \(A,s>0\) and \(\beta\in\mathbb R\), then
\[
  \Delta_n(T)\le C_{A,s,\beta}\,
  n\left(\frac{\log\log(e^e+n)}{\log(e+n)}\right)^s
  \bigl(\log\log(e^e+n)\bigr)^\beta
  \qquad(n\ge2).
\]
\item If \(a_m(T)\le A\exp(-bm^\theta)\) for all \(m\ge1\), with
\(A,b,\theta>0\), then
\[
  \Delta_n(T)\le C_{A,b,\theta}\,
  n\exp\!\left(-c_{b,\theta}(\log(e+n))^{\theta/(1+\theta)}\right)
  \qquad(n\ge2).
\]
For \(\theta=1\) this is
\(\Delta_n(T)\le C_{A,b}n\exp(-c_b\sqrt{\log(e+n)})\).
\end{enumerate}
The same conclusions hold for any totally bounded \(\mathcal K\subset B_E\)
whose widths \(a_m(\mathcal K)\) satisfy the corresponding decay.
\end{corollary}

\begin{proof}
The final assertion follows from the same proof with
\(a_m(\mathcal K)\) in place of \(a_m(T)\), so we prove the operator statements.

Assume first that \(a_m(T)\le A m^{-s}\).  For bounded \(n\) the claim follows
from the trivial estimate \(\Delta_n(T)\le n\).  For large \(n\), set
\[
  m=\left\lfloor
     \frac{\log(e+n)}
          {4s\,\log\log(e^e+n)}
     \right\rfloor,
  \qquad
  \eta=A m^{-s},
\]
enlarging constants to cover the range where \(m<1\) or \(\eta>1\).  Then
\[
  m\log\left(1+\frac6\eta\right)
  \le m\bigl(C_{A,s}+s\log m\bigr)
  \le \frac12\log(e+n)
\]
for all sufficiently large \(n\).  Hence the finite-rank net term in
Corollary~\ref{cor:banach-finite-rank-envelope} is at most \(C\sqrt n\), while
\[
  n(\eta+a_m(T))\le 2A n m^{-s}
  \le C_{A,s}n
     \left(\frac{\log\log(e^e+n)}{\log(e+n)}\right)^s .
\]
This term dominates \(C\sqrt n\) after increasing the constant, giving (i).

The log-refined polynomial case is identical.  Put
\[
  L_n=\log(e+n),\qquad M_n=\log\log(e^e+n),
\]
and, for large \(n\),
\[
  m=\left\lfloor
     c_{s,\beta}\frac{L_n}{M_n}
     \right\rfloor,
  \qquad
  \eta=A m^{-s}\bigl(\log(e+m)\bigr)^\beta,
\]
where \(c_{s,\beta}>0\) is small.  Enlarging constants covers the finite range
where \(m<1\) or \(\eta>1\).  Since
\[
  \log(1/\eta)
  \le C_{A,s,\beta}+s\log m+|\beta|\log\log(e+m),
\]
the choice of \(c_{s,\beta}\) ensures
\[
  m\log\left(1+\frac6\eta\right)\le\frac12L_n
\]
for all large \(n\).  The net term is therefore \(O(\sqrt n)\), while
\[
  n(\eta+a_m(T))
  \le C_{A,s,\beta} n m^{-s}\bigl(\log(e+m)\bigr)^\beta
  \le C_{A,s,\beta} n
  \left(\frac{M_n}{L_n}\right)^sM_n^\beta .
\]
This proves the second item, after again absorbing the \(O(\sqrt n)\) net term.

If \(a_m(T)\le A\exp(-bm^\theta)\), take for large \(n\)
\[
  m=\left\lfloor c_{b,\theta}(\log(e+n))^{1/(1+\theta)}\right\rfloor,
  \qquad
  \eta=A\exp(-bm^\theta),
\]
where \(c_{b,\theta}>0\) is small enough that
\[
  m\log\left(1+\frac6\eta\right)\le \frac12\log(e+n)
\]
for all large \(n\).  Corollary~\ref{cor:banach-finite-rank-envelope} then gives
a net term \(O(\sqrt n)\) and a tail term
\[
  n(\eta+a_m(T))\le 2A n e^{-bm^\theta}
  \le C_{A,b,\theta}n
       \exp\!\left(-c_{b,\theta}(\log(e+n))^{\theta/(1+\theta)}\right).
\]
Again the latter term dominates after enlarging the constant for small \(n\).
\end{proof}

\begin{corollary}[Hilbert singular-value rates]
\label{cor:hilbert-singular-rates}
Let \(H_1,H_2\) be real Hilbert spaces and let \(T:H_1\to H_2\) be a compact
operator with \(\|T\|\le1\).  Let
\[
  s_1(T)\ge s_2(T)\ge\cdots
\]
be its singular values, repeated with multiplicity and continued by zeros if
\(T\) has finite rank.  Then the following hold.
\begin{enumerate}[label=\textup{(\roman*)}]
\item If \(s_j(T)\le A j^{-s}\) for all \(j\ge1\), with \(A,s>0\), then
\[
  \Delta_n(T)\le C_{A,s}\,
  n\left(\frac{\log\log(e^e+n)}{\log(e+n)}\right)^s
  \qquad(n\ge2).
\]
\item More generally, if
\[
  s_j(T)\le A j^{-s}\bigl(\log(e+j)\bigr)^\beta
  \qquad(j\ge1),
\]
with \(A,s>0\) and \(\beta\in\mathbb R\), then
\[
  \Delta_n(T)\le C_{A,s,\beta}\,
  n\left(\frac{\log\log(e^e+n)}{\log(e+n)}\right)^s
  \bigl(\log\log(e^e+n)\bigr)^\beta
  \qquad(n\ge2).
\]
\item If \(s_j(T)\le A\exp(-bj^\theta)\) for all \(j\ge1\), with
\(A,b,\theta>0\), then
\[
  \Delta_n(T)\le C_{A,b,\theta}\,
  n\exp\!\left(-c_{b,\theta}(\log(e+n))^{\theta/(1+\theta)}\right)
  \qquad(n\ge2).
\]
For \(\theta=1\) this gives
\(\Delta_n(T)\le C_{A,b}n\exp(-c_b\sqrt{\log(e+n)})\).
\end{enumerate}
\end{corollary}

\begin{proof}
For compact operators between Hilbert spaces, the approximation numbers in
Corollary~\ref{cor:banach-finite-rank-envelope} satisfy
\[
  a_m(T)=s_{m+1}(T)\qquad(m\ge0),
\]
by the standard Schmidt approximation theorem.  The displayed rates are
therefore exactly the three cases of
Corollary~\ref{cor:compact-operator-approx-rates}, after shifting the index by
one and absorbing constants.
\end{proof}

\begin{corollary}[Schatten-class Hilbert images]
\label{cor:hilbert-schatten-rates}
Let \(H_1,H_2\) be real Hilbert spaces and let \(T:H_1\to H_2\) be a compact
operator with \(\|T\|\le1\).  If for some \(p,S>0\) its singular values satisfy
\[
  \left(\sum_{j\ge1}s_j(T)^p\right)^{1/p}\le S,
\]
then
\[
  \Delta_n(T)\le C_{S,p}\,
  n\left(\frac{\log\log(e^e+n)}{\log(e+n)}\right)^{1/p}
  \qquad(n\ge2).
\]
The same conclusion holds under the weak Schatten bound
\(\sup_{j\ge1} j^{1/p}s_j(T)\le S\).
\end{corollary}

\begin{proof}
Since the singular values are nonincreasing,
\[
  j\,s_j(T)^p\le \sum_{\ell=1}^j s_\ell(T)^p\le S^p,
\]
and hence \(s_j(T)\le S j^{-1/p}\).  The weak Schatten hypothesis is exactly
the same pointwise estimate.  Apply Corollary~\ref{cor:hilbert-singular-rates}
with \(s=1/p\).
\end{proof}

\begin{corollary}[Compact Assouad and logarithmic Assouad targets]
\label{cor:comm-af-compact-assouad}
Let \(Y\) be a compact metric space with \(\operatorname{diam}Y\le1\), and let
\(g_1,\ldots,g_n\in C(Y)\) be real-valued contractions with common
\((L,\alpha)\)-H\"older modulus.  If \(Y\) has upper Assouad growth \((A,s)\),
then
\[
  \disc(g_1,\ldots,g_n)
  \le C_{A,s,L,\alpha}
     \sqrt{n\log\!\left(e+n^{s/(2\alpha)-1}\right)}.
\]
In particular, \(s\le2\alpha\) gives Spencer scale.  Thus every compact
\(D\)-doubling target is Spencer-scale whenever \(\log_2D\le2\alpha\).

More generally, if \(Y\) has logarithmic upper Assouad growth
\((A,s,\gamma)\), then
\[
  \disc(g_1,\ldots,g_n)
  \le C\sqrt{
    n\log\!\left(
      e+n^{s/(2\alpha)-1}\bigl(\log(e+n)\bigr)^\gamma
    \right)} .
\]
Consequently \(s<2\alpha\) is Spencer-scale.  At \(s=2\alpha\), the target is
still Spencer-scale if \(\gamma\le0\), and otherwise has only the
critical-window loss
\[
  \disc(g_1,\ldots,g_n)\le C\sqrt{n\log\log(e^e+n)} .
\]
\end{corollary}

\begin{proof}
Fix \(y_0\in Y\).  Since \(\operatorname{diam}Y\le1\), the whole space lies in
\(\overline B(y_0,1)\).  Upper Assouad growth with \(R=1\) gives
\[
  N_Y(\rho)\le A\rho^{-s}\qquad(0<\rho\le1),
\]
and Corollary~\ref{cor:comm-af-compact-metric} gives the first estimate.
The doubling assertion follows from the standard iteration of the doubling
cover, which gives global covering exponent \(\log_2D\).

For logarithmic upper Assouad growth, the same \(R=1\) specialization gives
\[
  N_Y(\rho)\le A\rho^{-s}\bigl(\log(e/\rho)\bigr)^\gamma .
\]
Apply the direct compact-metric profile bound from
Corollary~\ref{cor:comm-af-compact-metric} with
\[
  \rho=\min\left(\frac12,(16L\sqrt n)^{-1/\alpha}\right).
\]
The approximation term is \(O_{L,\alpha}(\sqrt n)\), while
\[
  \frac{N_Y(\rho)}{n}
  \le C n^{s/(2\alpha)-1}\bigl(\log(e+n)\bigr)^\gamma .
\]
This proves the displayed logarithmic estimate, and the endpoint conclusions
follow from the sign of the polynomial and logarithmic powers.
\end{proof}

\begin{corollary}[Stretched-exponential compact entropy]
\label{cor:comm-af-stretched-entropy}
Let \(Y\) be a compact metric space with \(\operatorname{diam}Y\le1\), and let
\(g_1,\ldots,g_n\in C(Y)\) be real-valued contractions with common
\((L,\alpha)\)-H\"older modulus.  Suppose that, for some \(A<\infty\), \(q>0\),
and \(\gamma\in\mathbb R\),
\[
  \log N_Y(\rho)\le
  A\rho^{-q}\bigl(\log(e/\rho)\bigr)^\gamma
  \qquad(0<\rho\le1).
\]
Then
\[
  \disc(g_1,\ldots,g_n)
  \le
  C\,n^{(\alpha+q)/(2\alpha+q)}
  \bigl(\log(e+n)\bigr)^{\alpha\gamma/(2\alpha+q)} .
\]
In particular, if \(\log N_Y(\rho)=O(\rho^{-q})\) for some \(q>0\), then
\(\disc(g_1,\ldots,g_n)=o(n)\) for every \(\alpha>0\).
\end{corollary}

\begin{proof}
It is enough to treat large \(n\); enlarging \(C\) handles the remaining
finite range.  Put
\[
  \rho=n^{-1/(2\alpha+q)}
       \bigl(\log(e+n)\bigr)^{\gamma/(2\alpha+q)} .
\]
For large \(n\), \(0<\rho\le1/2\), and
\(\log(e/\rho)\asymp_{\alpha,q,\gamma}\log(e+n)\).  The direct compact profile
bound in Corollary~\ref{cor:comm-af-compact-metric} gives
\[
  \disc(g_1,\ldots,g_n)
  \le C\sqrt{n\log\!\left(e+\frac{N_Y(\rho)}{n}\right)}
       +C_{L,\alpha}n\rho^\alpha .
\]
Writing
\[
  T=A\rho^{-q}\bigl(\log(e/\rho)\bigr)^\gamma ,
\]
we have
\[
  \log\!\left(e+\frac{N_Y(\rho)}{n}\right)
  \le \log(e+\exp T)\le C(1+T).
\]
The choice of \(\rho\) gives
\[
  \sqrt{nT}
  \le
  C n^{(\alpha+q)/(2\alpha+q)}
    \bigl(\log(e+n)\bigr)^{\alpha\gamma/(2\alpha+q)}
\]
and the same bound for \(n\rho^\alpha\).  This proves the estimate.
\end{proof}

\begin{corollary}[H\"older quotients of compact metric models]
\label{cor:comm-af-holder-quotient}
Let \(X\) and \(Y\) be compact metric spaces with diameter at most \(1\), and
let \(\pi:X\to Y\) be a continuous surjection such that
\[
  d_Y(\pi x,\pi x')\le H\,d_X(x,x')^\beta
  \qquad(x,x'\in X),
  \qquad 0<\beta\le1 .
\]
Assume that
\[
  N_X(\rho)\le A\rho^{-s}\qquad(0<\rho\le1).
\]
If \(g_1,\ldots,g_n\in C(Y)\) are real-valued contractions with common
\((L,\alpha)\)-H\"older modulus on \(Y\), then
\[
  \disc(g_1,\ldots,g_n)
  \le C_{A,s,H,L,\alpha,\beta}
     \sqrt{n\log\!\left(e+n^{s/(2\alpha\beta)-1}\right)}.
\]
In particular, \(s\le2\alpha\beta\) gives the Spencer-scale estimate
\[
  \disc(g_1,\ldots,g_n)\le C_{A,s,H,L,\alpha,\beta}\sqrt n .
\]
Equivalently, if \(Y\) is a \(\beta\)-H\"older quotient of a compact model
\(X\) with upper metric dimension \(<2\alpha\beta\), then every common
\((L,\alpha)\)-H\"older tuple on \(Y\) is Spencer-scale.  The endpoint is also
Spencer-scale when the model has finite upper metric \(2\alpha\beta\)-content.
Likewise, if \(X\) has upper Assouad growth \((A,s)\), then the same estimate
holds with that exponent \(s\).  In particular, a \(\beta\)-H\"older quotient of
a compact \(D\)-doubling model is Spencer-scale whenever
\[
  \log_2D\le2\alpha\beta .
\]
\end{corollary}

\begin{proof}
Let \(f_i=g_i\circ\pi\).  Since \(\pi\) is onto, pullback is an isometry for
signed sums:
\[
  \left\|\sum_i\eps_i g_i\right\|_{C(Y)}
  =
  \left\|\sum_i\eps_i f_i\right\|_{C(X)}
  \qquad(\eps_i\in\{-1,1\}).
\]
Thus the two tuples have the same discrepancy.  Moreover
\[
  |f_i(x)-f_i(x')|
  \le L\,d_Y(\pi x,\pi x')^\alpha
  \le L H^\alpha d_X(x,x')^{\alpha\beta}.
\]
Apply Corollary~\ref{cor:comm-af-compact-metric} to \(X\), with H\"older
exponent \(\alpha\beta\).  The dimension and endpoint statements follow by
choosing an exponent above the model's upper metric dimension, respectively by
using finite upper metric \(2\alpha\beta\)-content at small scales and enlarging
constants for the remaining scales.
If \(X\) has upper Assouad growth \((A,s)\), then, after fixing one point in
\(X\), the whole space is contained in a closed unit ball and
\(N_X(\rho)\le A\rho^{-s}\); the first part applies.  The doubling conclusion
follows from the standard iteration of the doubling cover, which gives upper
Assouad exponent \(\log_2D\).
\end{proof}

\begin{corollary}[Logarithmic windows for H\"older quotients]
\label{cor:comm-af-holder-quotient-log-window}
In the setting of Corollary~\ref{cor:comm-af-holder-quotient}, replace the
covering hypothesis on \(X\) by
\[
  N_X(\rho)\le A\rho^{-s}\bigl(\log(e/\rho)\bigr)^\gamma
  \qquad(0<\rho\le1),
\]
where \(\gamma\in\mathbb R\).  If \(g_1,\ldots,g_n\in C(Y)\) are real-valued
contractions with common \((L,\alpha)\)-H\"older modulus on \(Y\), then
\[
  \disc(g_1,\ldots,g_n)
  \le C\sqrt{
    n\log\!\left(
      e+n^{s/(2\alpha\beta)-1}
       \bigl(\log(e+n)\bigr)^\gamma
    \right)} .
\]
Consequently \(s<2\alpha\beta\) is Spencer-scale.  At
\(s=2\alpha\beta\), the quotient is still Spencer-scale if \(\gamma\le0\), and
otherwise has only the critical-window loss
\[
  \disc(g_1,\ldots,g_n)
  \le C\sqrt{n\log\log(e^e+n)} .
\]
\end{corollary}

\begin{proof}
As in Corollary~\ref{cor:comm-af-holder-quotient}, pull back to
\(f_i=g_i\circ\pi\).  The discrepancy is unchanged and the pulled-back tuple is
common \((LH^\alpha,\alpha\beta)\)-H\"older on \(X\).  Apply the direct compact
metric profile bound from Corollary~\ref{cor:comm-af-compact-metric} at
\[
  \rho=\min\left(\frac12,(16\max(1,LH^\alpha)\sqrt n)^{-1/(\alpha\beta)}\right).
\]
The approximation term is \(O\sqrt n\), while
\[
  \frac{N_X(\rho)}{n}
  \le
  C n^{s/(2\alpha\beta)-1}\bigl(\log(e+n)\bigr)^\gamma .
\]
This proves the displayed estimate.  The subcritical and critical conclusions
are immediate from the sign of the polynomial exponent and, at equality, from
the sign of the logarithmic power.
\end{proof}

\begin{corollary}[Snowflaked compact targets]
\label{cor:comm-af-snowflake}
Let \((Y,d)\) be a compact metric space with \(\operatorname{diam}Y\le1\), and
let \(0<\vartheta\le1\).  Write
\[
  d_\vartheta(y,z)=d(y,z)^\vartheta .
\]
Assume that
\[
  N_{(Y,d)}(\rho)\le A\rho^{-s}\qquad(0<\rho\le1).
\]
If \(g_1,\ldots,g_n\in C(Y)\) are real-valued contractions with common
\((L,\alpha)\)-H\"older modulus with respect to \(d_\vartheta\), then
\[
  \disc(g_1,\ldots,g_n)
  \le C_{A,s,L,\alpha,\vartheta}
     \sqrt{n\log\!\left(e+n^{s/(2\alpha\vartheta)-1}\right)}.
\]
In particular, \(s\le2\alpha\vartheta\) gives the Spencer-scale estimate
\[
  \disc(g_1,\ldots,g_n)\le C_{A,s,L,\alpha,\vartheta}\sqrt n .
\]
Equivalently, if \(Y\) has finite upper metric dimension \(<2\alpha\vartheta\)
in the original metric, then every common \((L,\alpha)\)-H\"older tuple for the
snowflaked metric \(d^\vartheta\) is Spencer-scale.  The endpoint
\(s=2\alpha\vartheta\) is also Spencer-scale under finite upper metric
\(s\)-content in the original metric.
\end{corollary}

\begin{proof}
For \(0<r\le1\), an \(r\)-ball for the snowflaked metric \(d_\vartheta\) is the
same set as an \(r^{1/\vartheta}\)-ball for \(d\).  Hence
\[
  N_{(Y,d_\vartheta)}(r)
  =
  N_{(Y,d)}(r^{1/\vartheta})
  \le A r^{-s/\vartheta}.
\]
Apply Corollary~\ref{cor:comm-af-compact-metric} to the compact metric space
\((Y,d_\vartheta)\).  The displayed estimate is exactly the polynomial-covering
estimate there with covering exponent \(s/\vartheta\).  The dimension and
content statements follow by the same enlarging-of-constants argument used in
Corollary~\ref{cor:comm-af-compact-metric}.
\end{proof}

\begin{corollary}[Logarithmic windows for snowflaked compact targets]
\label{cor:comm-af-snowflake-log-window}
In the setting of Corollary~\ref{cor:comm-af-snowflake}, replace the covering
hypothesis on the original metric by
\[
  N_{(Y,d)}(\rho)\le A\rho^{-s}\bigl(\log(e/\rho)\bigr)^\gamma
  \qquad(0<\rho\le1),
\]
where \(\gamma\in\mathbb R\).  If \(g_1,\ldots,g_n\in C(Y)\) are real-valued
contractions with common \((L,\alpha)\)-H\"older modulus with respect to
\(d_\vartheta=d^\vartheta\), then
\[
  \disc(g_1,\ldots,g_n)
  \le C\sqrt{
    n\log\!\left(
      e+n^{s/(2\alpha\vartheta)-1}
       \bigl(\log(e+n)\bigr)^\gamma
    \right)} .
\]
Consequently \(s<2\alpha\vartheta\) is Spencer-scale.  At
\(s=2\alpha\vartheta\), the snowflaked target is still Spencer-scale if
\(\gamma\le0\), and otherwise has only the critical-window loss
\[
  \disc(g_1,\ldots,g_n)\le C\sqrt{n\log\log(e^e+n)} .
\]
\end{corollary}

\begin{proof}
For \(0<r\le1\),
\[
  N_{(Y,d_\vartheta)}(r)
  =
  N_{(Y,d)}(r^{1/\vartheta})
  \le
  A r^{-s/\vartheta}
  \bigl(\log(e/r^{1/\vartheta})\bigr)^\gamma .
\]
Since
\[
  \log(e/r)\le\log(e/r^{1/\vartheta})
  \le \vartheta^{-1}\log(e/r)
  \qquad(0<r\le1),
\]
the signed logarithmic power is bounded by a constant times
\((\log(e/r))^\gamma\).  Thus
\[
  N_{(Y,d_\vartheta)}(r)
  \le C r^{-s/\vartheta}\bigl(\log(e/r)\bigr)^\gamma .
\]
Apply the direct compact-metric profile bound from
Corollary~\ref{cor:comm-af-compact-metric} to \((Y,d_\vartheta)\), using
\[
  r=\min\left(\frac12,(16\max(1,L)\sqrt n)^{-1/\alpha}\right).
\]
The approximation term is \(O\sqrt n\), while
\[
  \frac{N_{(Y,d_\vartheta)}(r)}{n}
  \le
  C n^{s/(2\alpha\vartheta)-1}\bigl(\log(e+n)\bigr)^\gamma .
\]
This proves the displayed estimate.  The subcritical and critical conclusions
follow from the polynomial exponent and, at equality, the sign of the
logarithmic power.
\end{proof}

\begin{corollary}[Finite compact-metric unions keep the largest exponent]
\label{cor:comm-af-compact-union}
Let \(Y\) be a compact metric space with \(\operatorname{diam}Y\le1\), and
suppose \(Y=\bigcup_{\ell=1}^mY_\ell\) is a finite union of compact subsets.
Assume that, for each \(\ell\),
\[
  N_{Y_\ell}(\rho)\le A_\ell\rho^{-d_\ell}\qquad(0<\rho\le1),
\]
where covering numbers are computed in the inherited metric.  Put
\(d_*=\max_\ell d_\ell\).  If
\(g_1,\ldots,g_n\in C(Y)\) are real-valued contractions with common
\((L,\alpha)\)-H\"older modulus, then
\[
  \disc(g_1,\ldots,g_n)
  \le C_{\{A_\ell,d_\ell\},L,\alpha,m}
     \sqrt{n\log\!\left(e+n^{d_*/(2\alpha)-1}\right)}.
\]
In particular, \(d_*\le2\alpha\) gives the Spencer-scale bound
\[
  \disc(g_1,\ldots,g_n)
  \le C_{\{A_\ell,d_\ell\},L,\alpha,m}\sqrt n .
\]
Equivalently, finite upper metric dimensions take the maximum under finite
unions.  Thus if
\[
  \max_{\ell}\overline{\dim}_{\mathrm M}(Y_\ell)<2\alpha,
\]
then every common \((L,\alpha)\)-H\"older tuple on \(Y\) is Spencer-scale.  At
the critical endpoint, the same conclusion holds if each \(Y_\ell\) has finite
upper metric \(2\alpha\)-content.
Likewise, if each \(Y_\ell\) has upper Assouad growth \((A_\ell,s_\ell)\), then
the same conclusion holds with \(d_*=\max_\ell s_\ell\).  In particular, if
each \(Y_\ell\) is \(D_\ell\)-doubling, then the finite union is Spencer-scale
whenever
\[
  \max_\ell\log_2D_\ell\le2\alpha .
\]
\end{corollary}

\begin{proof}
For \(0<\rho\le1\), choose an inherited-metric \(\rho\)-cover of each
\(Y_\ell\).  Replacing each subspace ball by the ambient ball with the same
center gives a cover of \(Y\), so
\[
  N_Y(\rho)\le\sum_{\ell=1}^mN_{Y_\ell}(\rho)
  \le \sum_{\ell=1}^m A_\ell\rho^{-d_\ell}
  \le A_*\rho^{-d_*},
  \qquad A_*:=\sum_{\ell=1}^m A_\ell ,
\]
because \(\rho\le1\).  Corollary~\ref{cor:comm-af-compact-metric} gives the
displayed estimate and the Spencer-scale consequence.

For the dimension statement, write
\(d_\ell^{\rm M}:=\overline{\dim}_{\mathrm M}(Y_\ell)\).  Choose
\(D_\ell>d_\ell^{\rm M}\) with \(\max_\ell D_\ell<2\alpha\) and apply the
preceding paragraph.  If every
\(Y_\ell\) has finite upper metric \(2\alpha\)-content, then
\(N_{Y_\ell}(\rho)\le C_\ell\rho^{-2\alpha}\) at small scales; enlarging
constants for the remaining scales gives the critical endpoint.

If \(Y_\ell\) has upper Assouad growth \((A_\ell,s_\ell)\), then, after fixing
one point in \(Y_\ell\), the whole piece is contained in its closed
unit-radius ball, so \(N_{Y_\ell}(\rho)\le A_\ell\rho^{-s_\ell}\).  The first
paragraph applies.  The doubling assertion follows from the standard iteration
of the doubling cover, giving upper Assouad exponent \(\log_2D_\ell\) for each
piece.
\end{proof}

\begin{corollary}[Logarithmic windows for finite compact-metric unions]
\label{cor:comm-af-compact-union-log-window}
In the setting of Corollary~\ref{cor:comm-af-compact-union}, replace the
covering hypothesis on each piece by
\[
  N_{Y_\ell}(\rho)\le
  A_\ell\rho^{-d_\ell}\bigl(\log(e/\rho)\bigr)^{\gamma_\ell}
  \qquad(0<\rho\le1),
\]
where \(\gamma_\ell\in\mathbb R\).  Put
\[
  d_*=\max_\ell d_\ell,\qquad
  \gamma_*=\max_{\ell:\,d_\ell=d_*}\gamma_\ell .
\]
If \(g_1,\ldots,g_n\in C(Y)\) are real-valued contractions with common
\((L,\alpha)\)-H\"older modulus, then
\[
  \disc(g_1,\ldots,g_n)
  \le C\sqrt{
    n\log\!\left(
      e+n^{d_*/(2\alpha)-1}\bigl(\log(e+n)\bigr)^{\gamma_*}
    \right)} .
\]
Consequently \(d_*<2\alpha\) is Spencer-scale.  At \(d_*=2\alpha\), the finite
union is still Spencer-scale if \(\gamma_*\le0\), and otherwise has only the
critical-window loss
\[
  \disc(g_1,\ldots,g_n)
  \le C\sqrt{n\log\log(e^e+n)} .
\]
\end{corollary}

\begin{proof}
The inherited covers of the pieces give
\[
  N_Y(\rho)\le
  \sum_{\ell=1}^m
  A_\ell\rho^{-d_\ell}\bigl(\log(e/\rho)\bigr)^{\gamma_\ell}.
\]
For \(d_\ell=d_*\), the definition of \(\gamma_*\) dominates the logarithmic
factor.  For \(d_\ell<d_*\), the function
\[
  \rho^{d_*-d_\ell}\bigl(\log(e/\rho)\bigr)^{\gamma_\ell-\gamma_*}
\]
is bounded on \(0<\rho\le1\).  Hence
\[
  N_Y(\rho)\le C'\rho^{-d_*}\bigl(\log(e/\rho)\bigr)^{\gamma_*}.
\]
Apply Corollary~\ref{cor:comm-af-compact-metric} with
\[
  \rho=\min\left(\frac12,(16L\sqrt n)^{-1/\alpha}\right),
\]
exactly as in Corollary~\ref{cor:comm-af-compact-holder-product-log-window}.
\end{proof}

\begin{corollary}[Compact metric products add ordinary entropy]
\label{cor:comm-af-compact-product}
Let \(Y_j\), \(1\le j\le m\), be compact metric spaces with
\(\operatorname{diam}Y_j\le1\), and let \(Y=\prod_{j=1}^mY_j\) with the max
product metric.  Suppose
\[
  N_{Y_j}(\rho)\le A_j\rho^{-d_j}\qquad(0<\rho\le1),
\]
and put \(d_*=\sum_{j=1}^m d_j\).  If
\(g_1,\ldots,g_n\in C(Y)\) are real-valued contractions with common
\((L,\alpha)\)-H\"older modulus, then
\[
  \disc(g_1,\ldots,g_n)
  \le C_{\{A_j,d_j\},L,\alpha,m}
     \sqrt{n\log\!\left(e+n^{d_*/(2\alpha)-1}\right)}.
\]
In particular, \(d_*\le2\alpha\) gives the Spencer-scale bound
\[
  \disc(g_1,\ldots,g_n)\le C_{\{A_j,d_j\},L,\alpha,m}\sqrt n .
\]
Equivalently, finite ordinary upper metric dimensions add for arbitrary
compact metric targets: if
\[
  \sum_{j=1}^m \overline{\dim}_{\mathrm M}(Y_j)<2\alpha,
\]
then every common \((L,\alpha)\)-H\"older tuple on \(Y\) is Spencer-scale.  At
the critical endpoint, the same conclusion holds if
\[
  \sum_{j=1}^m s_j=2\alpha
  \quad\text{and}\quad
  \overline{\mathcal M}^{\,s_j}_{\mathrm M}(Y_j)<\infty
  \quad(1\le j\le m).
\]
Likewise, if each \(Y_j\) has upper Assouad growth \((A_j,s_j)\), then \(Y\)
has upper Assouad growth \((\prod_jA_j,\sum_js_j)\), so
\(\sum_js_j\le2\alpha\) gives Spencer scale.  In particular, a finite max
product of compact \(D_j\)-doubling metric spaces is Spencer-scale whenever
\[
  \sum_{j=1}^m\log_2D_j\le2\alpha .
\]
\end{corollary}

\begin{proof}
A product of open \(\rho\)-balls in the factors is an open \(\rho\)-ball in
the max product metric.  Hence
\[
  N_Y(\rho)\le\prod_{j=1}^mN_{Y_j}(\rho)
  \le\left(\prod_{j=1}^mA_j\right)\rho^{-d_*},
  \qquad 0<\rho\le1.
\]
The displayed discrepancy estimate and the \(d_*\le2\alpha\) consequence follow
from Corollary~\ref{cor:comm-af-compact-metric}.

The dimension statement follows by choosing
\(D_j>\overline{\dim}_{\mathrm M}(Y_j)\) with \(\sum_jD_j<2\alpha\) and
applying the preceding paragraph.  Finite upper metric \(s_j\)-content gives
\(N_{Y_j}(\rho)\le C_j\rho^{-s_j}\) at small scales; enlarging constants to
cover the remaining scales gives the critical endpoint.

For the Assouad statement, a closed \(R\)-ball in the max product metric is the
product of the closed \(R\)-balls in the factors.  Cover each factor ball by at
most \(A_j(R/\rho)^{s_j}\) open \(\rho\)-balls and multiply the covers.  The
doubling consequence is the special case \(s_j=\log_2D_j\), after the standard
iteration of the doubling cover.
\end{proof}

\begin{corollary}[Logarithmic windows for compact metric products]
\label{cor:comm-af-compact-product-log-window}
In the setting of Corollary~\ref{cor:comm-af-compact-product}, replace the
factor covering hypotheses by
\[
  N_{Y_j}(\rho)\le
  A_j\rho^{-d_j}\bigl(\log(e/\rho)\bigr)^{\gamma_j}
  \qquad(0<\rho\le1),
\]
where \(\gamma_j\in\mathbb R\).  Put
\[
  d_*=\sum_{j=1}^m d_j,\qquad
  \Gamma_*=\sum_{j=1}^m \gamma_j .
\]
If \(g_1,\ldots,g_n\in C(Y)\) are real-valued contractions with common
\((L,\alpha)\)-H\"older modulus, then
\[
  \disc(g_1,\ldots,g_n)
  \le C\sqrt{
    n\log\!\left(
      e+n^{d_*/(2\alpha)-1}\bigl(\log(e+n)\bigr)^{\Gamma_*}
    \right)} .
\]
Thus \(d_*<2\alpha\) gives Spencer scale.  At \(d_*=2\alpha\), products with
\(\Gamma_*\le0\) remain Spencer-scale, while \(\Gamma_*>0\) gives the
\(\sqrt{n\log\log n}\) critical-window loss.
\end{corollary}

\begin{proof}
A product of factor \(\rho\)-covers gives
\[
  N_Y(\rho)\le
  \left(\prod_{j=1}^mA_j\right)\rho^{-d_*}
  \bigl(\log(e/\rho)\bigr)^{\Gamma_*}.
\]
The discrepancy estimate and endpoint consequences are the same direct
application of Corollary~\ref{cor:comm-af-compact-metric} used in
Corollary~\ref{cor:comm-af-compact-union-log-window}.
\end{proof}

\begin{corollary}[Finite compact product strata keep the largest stratum exponent]
\label{cor:comm-af-compact-product-strata}
Let \(Y\) be a compact metric space with \(\operatorname{diam}Y\le1\), and
suppose
\[
  Y=\bigcup_{\ell=1}^M Z_\ell
\]
is a finite union of compact subsets.  For each \(\ell\), let
\[
  P_\ell=\prod_{j=1}^{m_\ell}Y_{\ell,j}
\]
with the max product metric, where each \(Y_{\ell,j}\) is compact and
\(\operatorname{diam}Y_{\ell,j}\le1\).  Assume that there is a
\(B_\ell\)-Lipschitz surjection \(\pi_\ell:P_\ell\to Z_\ell\), where
\(Z_\ell\) carries the metric inherited from \(Y\) and \(B_\ell\ge1\).
Suppose
\[
  N_{Y_{\ell,j}}(\rho)\le A_{\ell,j}\rho^{-d_{\ell,j}}
  \qquad(0<\rho\le1),
\]
and put
\[
  s_\ell=\sum_{j=1}^{m_\ell}d_{\ell,j},
  \qquad
  s_*=\max_{1\le\ell\le M}s_\ell .
\]
If \(g_1,\ldots,g_n\in C(Y)\) are real-valued contractions with common
\((L,\alpha)\)-H\"older modulus, then
\[
  \disc(g_1,\ldots,g_n)
  \le C\sqrt{n\log\!\left(e+n^{s_*/(2\alpha)-1}\right)},
\]
where \(C\) depends only on \(L,\alpha\) and the finite stratum data
\(\{A_{\ell,j},d_{\ell,j},B_\ell,m_\ell\}_{\ell,j}\).  In particular,
\(s_*\le2\alpha\) gives the Spencer-scale bound
\[
  \disc(g_1,\ldots,g_n)\le C\sqrt n .
\]
Equivalently, a finite union of product-parametrized compact strata is governed
by the largest stratum sum: if
\[
  \max_{\ell}\sum_{j=1}^{m_\ell}
  \overline{\dim}_{\mathrm M}(Y_{\ell,j})<2\alpha,
\]
then every common \((L,\alpha)\)-H\"older tuple on \(Y\) is Spencer-scale.  At
the critical endpoint, the same conclusion holds if, for every \(\ell\), there
are \(s_{\ell,j}\ge0\) with
\[
  \sum_{j=1}^{m_\ell}s_{\ell,j}\le2\alpha
  \quad\text{and}\quad
  \overline{\mathcal M}^{\,s_{\ell,j}}_{\mathrm M}(Y_{\ell,j})<\infty
  \quad(1\le j\le m_\ell).
\]
Likewise, if each \(Y_{\ell,j}\) has upper Assouad growth
\((A_{\ell,j},a_{\ell,j})\), then the same conclusion holds with
\[
  s_*=\max_\ell\sum_{j=1}^{m_\ell}a_{\ell,j}.
\]
In particular, if each \(Y_{\ell,j}\) is \(D_{\ell,j}\)-doubling, then \(Y\)
is Spencer-scale whenever
\[
  \max_\ell\sum_{j=1}^{m_\ell}\log_2D_{\ell,j}\le2\alpha .
\]
\end{corollary}

\begin{proof}
Fix a stratum \(\ell\).  A product of open \(\rho/B_\ell\)-balls in the
factors is an open \(\rho/B_\ell\)-ball in \(P_\ell\), and its image under
\(\pi_\ell\) is contained in an inherited-metric \(\rho\)-ball in \(Z_\ell\).
Therefore, for \(0<\rho\le1\),
\[
  N_{Z_\ell}(\rho)
  \le N_{P_\ell}(\rho/B_\ell)
  \le \prod_{j=1}^{m_\ell}N_{Y_{\ell,j}}(\rho/B_\ell)
  \le A'_\ell\rho^{-s_\ell},
\]
where
\[
  A'_\ell=\prod_{j=1}^{m_\ell} A_{\ell,j}B_\ell^{d_{\ell,j}} .
\]
Corollary~\ref{cor:comm-af-compact-union} applied to the pieces \(Z_\ell\)
gives the displayed estimate and the \(s_*\le2\alpha\) consequence.

For the dimension assertion, choose
\(D_{\ell,j}>\overline{\dim}_{\mathrm M}(Y_{\ell,j})\) so that
\[
  \max_\ell\sum_{j=1}^{m_\ell}D_{\ell,j}<2\alpha,
\]
and use the preceding paragraph with \(d_{\ell,j}=D_{\ell,j}\).  The
critical-content statement is the same
argument with
\[
  N_{Y_{\ell,j}}(\rho)\le C_{\ell,j}\rho^{-s_{\ell,j}}
\]
at small scales, after enlarging constants for the remaining scales.  The
Assouad and doubling assertions follow by the same global-covering reduction
used in Corollaries~\ref{cor:comm-af-compact-union} and
\ref{cor:comm-af-compact-product}.
\end{proof}

\begin{corollary}[H\"older product strata divide the parametrizing exponent]
\label{cor:comm-af-compact-holder-product-strata}
Let \(Y\) be a compact metric space with \(\operatorname{diam}Y\le1\), and
suppose
\[
  Y=\bigcup_{\ell=1}^M Z_\ell
\]
is a finite union of compact subsets.  For each \(\ell\), let
\[
  P_\ell=\prod_{j=1}^{m_\ell}Y_{\ell,j}
\]
with the max product metric, where each \(Y_{\ell,j}\) is compact and
\(\operatorname{diam}Y_{\ell,j}\le1\).  Assume that there is a
\((B_\ell,\beta_\ell)\)-H\"older surjection \(\pi_\ell:P_\ell\to Z_\ell\),
where \(0<\beta_\ell\le1\), \(B_\ell\ge1\), and
\[
  d_Y(\pi_\ell x,\pi_\ell y)\le B_\ell d_{P_\ell}(x,y)^{\beta_\ell}.
\]
Suppose
\[
  N_{Y_{\ell,j}}(\rho)\le A_{\ell,j}\rho^{-d_{\ell,j}}
  \qquad(0<\rho\le1),
\]
and put
\[
  s_\ell=\sum_{j=1}^{m_\ell}d_{\ell,j},
  \qquad
  t_\ell=\frac{s_\ell}{\beta_\ell},
  \qquad
  t_*=\max_{1\le\ell\le M}t_\ell .
\]
If \(g_1,\ldots,g_n\in C(Y)\) are real-valued contractions with common
\((L,\alpha)\)-H\"older modulus, then
\[
  \disc(g_1,\ldots,g_n)
  \le C\sqrt{n\log\!\left(e+n^{t_*/(2\alpha)-1}\right)},
\]
where \(C\) depends only on \(L,\alpha\) and on the finite collection of
parameters
\[
  A_{\ell,j},\ d_{\ell,j},\ B_\ell,\ \beta_\ell,\ m_\ell
  \qquad(1\le\ell\le M,\ 1\le j\le m_\ell).
\]
In
particular, \(t_*\le2\alpha\) gives the Spencer-scale bound
\[
  \disc(g_1,\ldots,g_n)\le C\sqrt n .
\]
Equivalently, a finite union of H\"older product-parametrized compact strata
is governed by
\[
  \max_\ell\frac{1}{\beta_\ell}
  \sum_{j=1}^{m_\ell}\overline{\dim}_{\mathrm M}(Y_{\ell,j}).
\]
If this quantity is \(<2\alpha\), then every common \((L,\alpha)\)-H\"older
tuple on \(Y\) is Spencer-scale.  At the critical endpoint, the same
conclusion holds if, for every \(\ell\), there are \(s_{\ell,j}\ge0\) with
\[
  \frac{1}{\beta_\ell}\sum_{j=1}^{m_\ell}s_{\ell,j}\le2\alpha
  \quad\text{and}\quad
  \overline{\mathcal M}^{\,s_{\ell,j}}_{\mathrm M}(Y_{\ell,j})<\infty
  \quad(1\le j\le m_\ell).
\]
Likewise, if each \(Y_{\ell,j}\) has upper Assouad growth
\((A_{\ell,j},a_{\ell,j})\), then the same conclusion holds with
\[
  t_*=\max_\ell\frac{1}{\beta_\ell}
       \sum_{j=1}^{m_\ell}a_{\ell,j}.
\]
In particular, if each \(Y_{\ell,j}\) is \(D_{\ell,j}\)-doubling, then \(Y\)
is Spencer-scale whenever
\[
  \max_\ell\frac{1}{\beta_\ell}
       \sum_{j=1}^{m_\ell}\log_2D_{\ell,j}\le2\alpha .
\]
\end{corollary}

\begin{proof}
Fix a stratum \(\ell\).  Put
\[
  r_\ell(\rho)=\left(\frac{\rho}{B_\ell}\right)^{1/\beta_\ell}.
\]
Since \(0<\rho\le1\) and \(B_\ell\ge1\), one has \(0<r_\ell(\rho)\le1\).
A product of open \(r_\ell(\rho)\)-balls in the factors is an open
\(r_\ell(\rho)\)-ball in \(P_\ell\), and its image under \(\pi_\ell\) is
contained in an inherited-metric \(\rho\)-ball in \(Z_\ell\).  Hence
\[
  N_{Z_\ell}(\rho)
  \le N_{P_\ell}(r_\ell(\rho))
  \le \prod_{j=1}^{m_\ell}N_{Y_{\ell,j}}(r_\ell(\rho))
  \le A''_\ell\rho^{-s_\ell/\beta_\ell},
\]
where
\[
  A''_\ell=\prod_{j=1}^{m_\ell}
  A_{\ell,j}B_\ell^{d_{\ell,j}/\beta_\ell}.
\]
Corollary~\ref{cor:comm-af-compact-union} applied to the pieces \(Z_\ell\)
gives the displayed estimate and the \(t_*\le2\alpha\) consequence.

The ordinary dimension statement follows by choosing
\(D_{\ell,j}>\overline{\dim}_{\mathrm M}(Y_{\ell,j})\) so that
\[
  \max_\ell\frac{1}{\beta_\ell}\sum_{j=1}^{m_\ell}D_{\ell,j}<2\alpha
\]
and applying the preceding paragraph.  The finite-content endpoint, Assouad
statement, and doubling statement are the same covering reduction with
\(d_{\ell,j}=s_{\ell,j}\), \(d_{\ell,j}=a_{\ell,j}\), and
\(a_{\ell,j}=\log_2D_{\ell,j}\), respectively.
\end{proof}

\begin{corollary}[Logarithmic windows for H\"older product strata]
\label{cor:comm-af-compact-holder-product-log-window}
In the setting of
Corollary~\ref{cor:comm-af-compact-holder-product-strata}, replace the factor
covering hypothesis by
\[
  N_{Y_{\ell,j}}(\rho)
  \le A_{\ell,j}\rho^{-d_{\ell,j}}
       \bigl(\log(e/\rho)\bigr)^{\gamma_{\ell,j}}
  \qquad(0<\rho\le1),
\]
with \(\gamma_{\ell,j}\in\mathbb R\).  Put
\[
  t_\ell=\frac{1}{\beta_\ell}\sum_{j=1}^{m_\ell}d_{\ell,j},\qquad
  \Gamma_\ell=\sum_{j=1}^{m_\ell}\gamma_{\ell,j},\qquad
  t_*=\max_\ell t_\ell,\qquad
  \Gamma_*=\max_{\ell:\,t_\ell=t_*}\Gamma_\ell .
\]
If \(g_1,\ldots,g_n\in C(Y)\) are real-valued contractions with common
\((L,\alpha)\)-H\"older modulus, then
\[
  \disc(g_1,\ldots,g_n)
  \le C\sqrt{
      n\log\!\left(
        e+n^{t_*/(2\alpha)-1}\bigl(\log(e+n)\bigr)^{\Gamma_*}
      \right)} ,
\]
where \(C\) depends only on \(L,\alpha\) and the finite logarithmic stratum
data.  Consequently \(t_*<2\alpha\) is still Spencer-scale.  At the critical
exponent \(t_*=2\alpha\), one has the endpoint bound
\[
  \disc(g_1,\ldots,g_n)
  \le
  \begin{cases}
  C\sqrt n, & \Gamma_*\le0,\\[3pt]
  C\sqrt{n\log\log(e^e+n)}, & \Gamma_*>0 .
  \end{cases}
\]
\end{corollary}

\begin{proof}
The proof of
Corollary~\ref{cor:comm-af-compact-holder-product-strata} gives, with
\[
  r_\ell(\rho)=\left(\frac{\rho}{B_\ell}\right)^{1/\beta_\ell},
\]
the estimate
\[
  N_{Z_\ell}(\rho)
  \le C_\ell\rho^{-t_\ell}
       \bigl(\log(e/r_\ell(\rho))\bigr)^{\Gamma_\ell}.
\]
Since \(B_\ell\ge1\) and \(0<\beta_\ell\le1\), the last logarithm is comparable
to \(\log(e/\rho)\), with constants depending only on \(B_\ell,\beta_\ell\).
Thus the same bound holds with \(\log(e/\rho)^{\Gamma_\ell}\).  Because there
are only finitely many strata and polynomial gaps dominate all logarithmic
powers, every term with \(t_\ell<t_*\) is also bounded by a constant multiple of
\(\rho^{-t_*}(\log(e/\rho))^{\Gamma_*}\).  For \(t_\ell=t_*\), the definition of
\(\Gamma_*\) gives the same domination.  Hence
\[
  N_Y(\rho)\le C'\rho^{-t_*}\bigl(\log(e/\rho)\bigr)^{\Gamma_*}
  \qquad(0<\rho\le1).
\]
Apply the direct compact-metric profile bound from
Corollary~\ref{cor:comm-af-compact-metric} with
\[
  \rho=\min\left(\frac12,(16L\sqrt n)^{-1/\alpha}\right).
\]
The approximation term is \(O_{L,\alpha}(\sqrt n)\), and
\[
  \frac{N_Y(\rho)}{n}
  \le C'' n^{t_*/(2\alpha)-1}\bigl(\log(e+n)\bigr)^{\Gamma_*}.
\]
This is the displayed estimate.  If \(t_*<2\alpha\), the factor after \(e+\) is
bounded uniformly in \(n\).  If \(t_*=2\alpha\), it is bounded for
\(\Gamma_*\le0\), and for \(\Gamma_*>0\) contributes only
\(\log\log(e^e+n)\).
\end{proof}

\begin{corollary}[Euclidean cubes via Cantor resolutions]
\label{cor:comm-af-euclidean-cube}
Let \(D\ge1\), and give \([0,1]^D\) the max metric.  If
\(g_1,\ldots,g_n\in C([0,1]^D)\) are real-valued contractions with common
\((L,\alpha)\)-H\"older modulus, then
\[
  \disc(g_1,\ldots,g_n)
  \le C_{D,L,\alpha}
     \sqrt{n\log\!\left(e+n^{D/(2\alpha)-1}\right)}.
\]
In particular, \(D\le2\alpha\) gives
\[
  \disc(g_1,\ldots,g_n)\le C_{D,L,\alpha}\sqrt n .
\]
More generally, if \(K\subset[0,1]^D\) is compact and \(0<\alpha\le1\), the
same estimates hold for common \((L,\alpha)\)-H\"older contractions on \(K\).
\end{corollary}

\begin{proof}
Let
\[
  \Sigma_D=(\{0,1\}^D)^{\mathbb N}.
\]
For distinct \(\sigma,\tau\in\Sigma_D\), put
\[
  d_\Sigma(\sigma,\tau)
  =2^{-\min\{k\ge0:\sigma_k\ne\tau_k\}},
\]
and put \(d_\Sigma(\sigma,\sigma)=0\).  The length-\(k\) cylinder partition
has \(2^{Dk}\) atoms of diameter at most \(2^{-k}\).  Given \(0<\rho\le1\),
choose \(k\) minimal with \(2^{-k}<\rho\).  These cylinders are contained in
open \(\rho\)-balls and \(2^k\le2\rho^{-1}\), so
\[
  N_{\Sigma_D}(\rho)\le 2^D\rho^{-D}\qquad(0<\rho\le1).
\]
Define the binary coding map \(\pi:\Sigma_D\to[0,1]^D\) by
\[
  \pi(\sigma)_j=\sum_{k=0}^\infty \sigma_{k,j}2^{-(k+1)},
  \qquad 1\le j\le D .
\]
It is onto, and if the first differing symbol occurs at level \(k\), then
\[
  \|\pi(\sigma)-\pi(\tau)\|_\infty\le 2^{-k}=d_\Sigma(\sigma,\tau).
\]
Thus \([0,1]^D\) has a \((2^D,D,1,1)\)-zero-dimensional H\"older resolution,
and Corollary~\ref{cor:comm-af-holder-resolution} proves the cube statement.

For a compact \(K\subset[0,1]^D\) and \(0<\alpha\le1\), each real-valued
\((L,\alpha)\)-H\"older contraction \(g_i\) on \(K\) has the McShane extension
\[
  \widetilde g_i(x)
  =\inf_{y\in K}\bigl(g_i(y)+L\|x-y\|_\infty^\alpha\bigr),
  \qquad x\in[0,1]^D,
\]
because \(\|x-y\|_\infty^\alpha\) is a metric.  Clipping
\(\widetilde g_i\) to the interval \([-1,1]\) preserves both the restriction to
\(K\) and the \(L\)-H\"older constant.  A signing for the clipped extensions on
the cube restricts to the same signing on \(K\), so the cube estimate gives the
compact-subset estimate.
\end{proof}

\begin{corollary}[Finite Lipschitz atlases]
\label{cor:comm-af-lipschitz-atlas}
Let \(Y\) be a compact metric space with \(\operatorname{diam}Y\le1\).  Suppose
that
\[
  Y=\bigcup_{\ell=1}^m\phi_\ell(K_\ell),
\]
where \(K_\ell\subset[0,1]^{D_\ell}\) is compact, the cube has the max metric,
and \(\phi_\ell:K_\ell\to Y\) is \(B_\ell\)-Lipschitz.  Put
\[
  D_*=\max_{1\le\ell\le m}D_\ell .
\]
If \(g_1,\ldots,g_n\in C(Y)\) are real-valued contractions with common
\((L,\alpha)\)-H\"older modulus, then
\[
  \disc(g_1,\ldots,g_n)
  \le C_{\{D_\ell,B_\ell\},m,L,\alpha}
     \sqrt{n\log\!\left(e+n^{D_*/(2\alpha)-1}\right)}.
\]
In particular, every compact target admitting a finite \(D\)-dimensional
Lipschitz atlas is Spencer-scale for common \((L,\alpha)\)-H\"older tuples
whenever \(D\le2\alpha\):
\[
  \disc(g_1,\ldots,g_n)
  \le C\sqrt n .
\]
Thus compact Lipschitz \(D\)-manifolds, finite \(D\)-polyhedra, and finite
unions of Lipschitz images of compact subsets of \(\R^D\) obey the same
Euclidean threshold \(D\le2\alpha\).
\end{corollary}

\begin{proof}
It is enough to bound the covering number of each chart image.  If
\(B_\ell^+=\max(1,B_\ell)\), then a max-metric grid cover of
\([0,1]^{D_\ell}\) gives
\[
  N_{K_\ell}(r)\le C_{D_\ell}r^{-D_\ell}\qquad(0<r\le1).
\]
Since \(\phi_\ell\) is \(B_\ell\)-Lipschitz,
\[
  N_{\phi_\ell(K_\ell)}(\rho)
  \le N_{K_\ell}(\rho/B_\ell^+)
  \le C_{D_\ell}(B_\ell^+)^{D_\ell}\rho^{-D_\ell}
  \qquad(0<\rho\le1).
\]
Corollary~\ref{cor:comm-af-compact-union}, applied to the finite union of the
chart images, gives the displayed estimate with exponent \(D_*\).  The
Spencer-scale consequence follows from \(D_*\le2\alpha\).  The manifold and
polyhedral examples have finite Lipschitz atlases of the stated dimension after
choosing finitely many compact coordinate patches.
\end{proof}

\begin{corollary}[Finite H\"older atlases]
\label{cor:comm-af-holder-atlas}
Let \(Y\) be a compact metric space with \(\operatorname{diam}Y\le1\).  Suppose
that
\[
  Y=\bigcup_{\ell=1}^m\phi_\ell(K_\ell),
\]
where \(K_\ell\subset[0,1]^{D_\ell}\) is compact, the cube has the max metric,
and \(\phi_\ell:K_\ell\to Y\) is \((B_\ell,\beta_\ell)\)-H\"older:
\[
  d_Y(\phi_\ell(x),\phi_\ell(y))
  \le B_\ell\|x-y\|_\infty^{\beta_\ell},
  \qquad 0<\beta_\ell\le1 .
\]
Put
\[
  s_*=\max_{1\le\ell\le m}\frac{D_\ell}{\beta_\ell}.
\]
If \(g_1,\ldots,g_n\in C(Y)\) are real-valued contractions with common
\((L,\alpha)\)-H\"older modulus, then
\[
  \disc(g_1,\ldots,g_n)
  \le C_{\{D_\ell,B_\ell,\beta_\ell\},m,L,\alpha}
     \sqrt{n\log\!\left(e+n^{s_*/(2\alpha)-1}\right)}.
\]
In particular, \(s_*\le2\alpha\) gives the Spencer-scale estimate
\[
  \disc(g_1,\ldots,g_n)\le C\sqrt n .
\]
Equivalently, a finite \(D\)-dimensional \(\beta\)-H\"older atlas is
Spencer-scale whenever \(D\le2\alpha\beta\).  The Lipschitz-atlas corollary is
the special case \(\beta_\ell=1\).
\end{corollary}

\begin{proof}
Let \(B_\ell^+=\max(1,B_\ell)\).  The grid-cover estimate used in the proof of
Corollary~\ref{cor:comm-af-lipschitz-atlas} gives
\[
  N_{K_\ell}(r)\le C_{D_\ell}r^{-D_\ell}\qquad(0<r\le1).
\]
A set in \(K_\ell\) of diameter at most
\((\rho/B_\ell^+)^{1/\beta_\ell}\) has image of diameter at most \(\rho\).
Hence, for \(0<\rho\le1\),
\[
  N_{\phi_\ell(K_\ell)}(\rho)
  \le
  N_{K_\ell}\!\left((\rho/B_\ell^+)^{1/\beta_\ell}\right)
  \le
  C_{D_\ell}(B_\ell^+)^{D_\ell/\beta_\ell}
  \rho^{-D_\ell/\beta_\ell}.
\]
Apply Corollary~\ref{cor:comm-af-compact-union} to the finite union of the
chart images.  The exponent is
\(\max_\ell D_\ell/\beta_\ell=s_*\), giving the displayed estimate.  If
\(s_*\le2\alpha\), Corollary~\ref{cor:comm-af-compact-union} gives the
Spencer-scale conclusion.
\end{proof}

\begin{corollary}[Logarithmic H\"older atlases]
\label{cor:comm-af-holder-atlas-log-window}
Let \(Y\) be a compact metric space with \(\operatorname{diam}Y\le1\).  Suppose
that
\[
  Y=\bigcup_{\ell=1}^m\phi_\ell(K_\ell),
\]
where each \(K_\ell\) is a compact metric space with \(\operatorname{diam}K_\ell\le1\), and
\(\phi_\ell:K_\ell\to Y\) is \((B_\ell,\beta_\ell)\)-H\"older:
\[
  d_Y(\phi_\ell(x),\phi_\ell(y))
  \le B_\ell d_{K_\ell}(x,y)^{\beta_\ell},
  \qquad 0<\beta_\ell\le1 .
\]
Assume that, for each \(\ell\),
\[
  N_{K_\ell}(\rho)\le
  A_\ell\rho^{-d_\ell}\bigl(\log(e/\rho)\bigr)^{\gamma_\ell}
  \qquad(0<\rho\le1),
\]
where \(\gamma_\ell\in\mathbb R\).  Put
\[
  t_\ell=\frac{d_\ell}{\beta_\ell},\qquad
  t_*=\max_{1\le\ell\le m}t_\ell,\qquad
  \Gamma_*=\max_{\ell:\,t_\ell=t_*}\gamma_\ell .
\]
If \(g_1,\ldots,g_n\in C(Y)\) are real-valued contractions with common
\((L,\alpha)\)-H\"older modulus, then
\[
  \disc(g_1,\ldots,g_n)
  \le C\sqrt{
    n\log\!\left(
      e+n^{t_*/(2\alpha)-1}\bigl(\log(e+n)\bigr)^{\Gamma_*}
    \right)} .
\]
Consequently \(t_*<2\alpha\) is Spencer-scale.  At \(t_*=2\alpha\), the atlas
is still Spencer-scale if \(\Gamma_*\le0\), and otherwise has only the
critical-window loss
\[
  \disc(g_1,\ldots,g_n)\le C\sqrt{n\log\log(e^e+n)} .
\]
\end{corollary}

\begin{proof}
Let \(B_\ell^+=\max(1,B_\ell)\).  A cover of \(K_\ell\) by balls of radius
\((\rho/B_\ell^+)^{1/\beta_\ell}\) maps to a cover of
\(\phi_\ell(K_\ell)\) by \(Y\)-balls of radius \(\rho\).  Hence
\[
  N_{\phi_\ell(K_\ell)}(\rho)
  \le
  A_\ell (B_\ell^+)^{d_\ell/\beta_\ell}\rho^{-d_\ell/\beta_\ell}
  \bigl(\log(e(B_\ell^+/\rho)^{1/\beta_\ell})\bigr)^{\gamma_\ell}.
\]
For \(0<\rho\le1\), the last logarithm is comparable to \(\log(e/\rho)\),
with constants depending on \(B_\ell\) and \(\beta_\ell\).  This remains true
for signed powers \(\gamma_\ell\).  Therefore
\[
  N_{\phi_\ell(K_\ell)}(\rho)
  \le
  C_\ell\rho^{-t_\ell}\bigl(\log(e/\rho)\bigr)^{\gamma_\ell}.
\]
Now apply Corollary~\ref{cor:comm-af-compact-union-log-window} to the finite
union of chart images.  The top polynomial exponent is \(t_*\), and among
charts with this exponent the top logarithmic power is \(\Gamma_*\), giving
the displayed estimate and the endpoint consequences.
\end{proof}

\begin{corollary}[Exhaustive finite-scale H\"older atlases]
\label{cor:comm-af-exhaustive-finite-scale-atlas}
Let \(Y\) be a compact metric space with \(\operatorname{diam}Y\le1\).  Let
\(Y_r\subset Y\), \(r\ge1\), be compact sets whose one-sided Hausdorff errors
\[
  \delta_r:=\sup_{y\in Y}\operatorname{dist}(y,Y_r)
\]
tend to \(0\).  Assume that there are parameters \(s\ge0\), \(\gamma\in\mathbb R\),
and \(A_r\ge1\) such that
\[
  N_{Y_r}(u)\le A_ru^{-s}\bigl(\log(e/u)\bigr)^\gamma
  \qquad(0<u\le1,\ r\ge1).
\]
Let \(g_1,\ldots,g_n\in C(Y)\) be real-valued contractions with common
\((L,\alpha)\)-H\"older modulus, and put
\[
  L_+:=\max(1,L),\qquad
  \rho_n:=\min\left(\frac12,\frac{1}{16L_+^{1/\alpha}n^{1/(2\alpha)}}\right).
\]
Choose any \(r_n\) with \(\delta_{r_n}\le\rho_n/2\).  Then
\[
  \disc(g_1,\ldots,g_n)
  \le
  C_{s,\gamma,L,\alpha}\sqrt{
    n\log\!\left(
      e+A_{r_n}n^{s/(2\alpha)-1}\bigl(\log(e+n)\bigr)^\gamma
    \right)} .
\]
Consequently, exhaustive stages may grow up to the unused entropy slack: if
\[
  A_{r_n}\le C_0n^{1-s/(2\alpha)}\bigl(\log(e+n)\bigr)^{-\gamma}
  \qquad(n\ge1),
\]
then the tuple is Spencer-scale, \(\disc(g_1,\ldots,g_n)\le C\sqrt n\).  At
the critical exponent \(s=2\alpha\), if
\[
  A_{r_n}\le C_0\bigl(\log(e+n)\bigr)^\kappa,
\]
then the same Spencer-scale conclusion holds when \(\kappa+\gamma\le0\), while
\(\kappa+\gamma>0\) gives the critical-window bound
\[
  \disc(g_1,\ldots,g_n)\le C\sqrt{n\log\log(e^e+n)} .
\]
The stage covering hypothesis is automatic when the finite-atlas covering
calculation on each \(Y_r\) gives the same \(A_r\)-budgeted profile.
\end{corollary}

\begin{proof}
Fix \(0<\rho\le1/2\) and choose \(r\) with \(\delta_r\le\rho/2\).  If
\(\{B(y_j,\rho/2)\}_{j=1}^M\) covers \(Y_r\), then every \(y\in Y\) has
\(y'\in Y_r\) with \(d_Y(y,y')\le\rho/2\), and \(y'\) lies in some
\(B(y_j,\rho/2)\).  Hence \(Y\subset\bigcup_jB(y_j,\rho)\), so
\[
  N_Y(\rho)\le N_{Y_r}(\rho/2)
  \le C_{s,\gamma}A_r\rho^{-s}\bigl(\log(e/\rho)\bigr)^\gamma .
\]
The signed logarithmic power is harmless because
\(\log(2e/\rho)\) and \(\log(e/\rho)\) are comparable on \(0<\rho\le1/2\).

Apply the direct profile bound in
Corollary~\ref{cor:comm-af-compact-metric} with \(\rho=\rho_n\) and
\(r=r_n\).  The H\"older approximation term satisfies
\[
  2nL(16\rho_n)^\alpha\le2\sqrt n ,
\]
and the entropy term is bounded by
\[
  \frac{N_Y(\rho_n)}{n}
  \le C_{s,\gamma,L,\alpha}
     A_{r_n}n^{s/(2\alpha)-1}\bigl(\log(e+n)\bigr)^\gamma .
\]
This proves the displayed estimate.  The Spencer-scale and critical-window
consequences are immediate from the size of the quantity inside the logarithm.
The final atlas assertion is the finite-stage covering calculation from
Corollary~\ref{cor:comm-af-holder-atlas-log-window}, applied separately to
each \(Y_r\).
\end{proof}

\begin{corollary}[Polynomial exhaustive atlas schedules]
\label{cor:comm-af-exhaustive-polynomial-atlas}
In the setting of
Corollary~\ref{cor:comm-af-exhaustive-finite-scale-atlas}, assume in addition
that, for some \(H,A,\eta>0\), \(\lambda\ge0\), and \(\kappa\in\mathbb R\),
\[
  \delta_r\le Hr^{-\eta},\qquad
  A_r\le Ar^\lambda\bigl(\log(e+r)\bigr)^\kappa
  \qquad(r\ge1).
\]
Put
\[
  s_{\mathrm{eff}}:=s+\frac{\lambda}{\eta},
  \qquad
  \Gamma_{\mathrm{eff}}:=\gamma+\kappa .
\]
Then every real-valued contraction tuple on \(Y\) with common
\((L,\alpha)\)-H\"older modulus satisfies
\[
  \disc(g_1,\ldots,g_n)
  \le C\sqrt{
    n\log\!\left(
      e+n^{s_{\mathrm{eff}}/(2\alpha)-1}
        \bigl(\log(e+n)\bigr)^{\Gamma_{\mathrm{eff}}}
    \right)} .
\]
Consequently \(s_{\mathrm{eff}}<2\alpha\) is Spencer-scale.  At
\(s_{\mathrm{eff}}=2\alpha\), the same Spencer-scale conclusion holds when
\(\Gamma_{\mathrm{eff}}\le0\), while \(\Gamma_{\mathrm{eff}}>0\) gives only the
critical-window loss
\[
  \disc(g_1,\ldots,g_n)\le C\sqrt{n\log\log(e^e+n)} .
\]
Thus a polynomially improving exhaustion trades approximation rate against
stage complexity by adding \(\lambda/\eta\) to the intrinsic stage exponent.
\end{corollary}

\begin{proof}
Use the signing scale \(\rho_n\) from
Corollary~\ref{cor:comm-af-exhaustive-finite-scale-atlas}, and choose
\[
  r_n=\left\lceil\max\left(1,\left(\frac{2H}{\rho_n}\right)^{1/\eta}\right)\right\rceil .
\]
Then \(\delta_{r_n}\le\rho_n/2\).  Since
\(\rho_n\asymp_{L,\alpha} n^{-1/(2\alpha)}\), the choice of \(r_n\) gives
\[
  r_n^\lambda\le C n^{\lambda/(2\alpha\eta)}
  \quad\text{and}\quad
  \bigl(\log(e+r_n)\bigr)^\kappa
  \le C\bigl(\log(e+n)\bigr)^\kappa ,
\]
with the usual comparison for signed logarithmic powers.  Therefore
\[
  A_{r_n}n^{s/(2\alpha)-1}\bigl(\log(e+n)\bigr)^\gamma
  \le
  C n^{s_{\mathrm{eff}}/(2\alpha)-1}
    \bigl(\log(e+n)\bigr)^{\Gamma_{\mathrm{eff}}}.
\]
Substituting this in
Corollary~\ref{cor:comm-af-exhaustive-finite-scale-atlas} proves the displayed
estimate.  The Spencer-scale and critical-window conclusions follow by
inspecting the term inside the outer logarithm.
\end{proof}

\begin{corollary}[Stretched-exponential exhaustive atlas schedules]
\label{cor:comm-af-exhaustive-stretched-atlas}
In the setting of
Corollary~\ref{cor:comm-af-exhaustive-finite-scale-atlas}, assume in addition
that, for some \(H,A,c,\theta>0\), \(\lambda\ge0\), and
\(\kappa\in\mathbb R\),
\[
  \delta_r\le H\exp(-cr^\theta),\qquad
  A_r\le Ar^\lambda\bigl(\log(e+r)\bigr)^\kappa
  \qquad(r\ge1).
\]
Put
\[
  \Gamma_{\mathrm{fast}}:=\gamma+\frac{\lambda}{\theta}.
\]
Then every real-valued contraction tuple on \(Y\) with common
\((L,\alpha)\)-H\"older modulus satisfies
\[
  \disc(g_1,\ldots,g_n)
  \le C\sqrt{
    n\log\!\left(
      e+n^{s/(2\alpha)-1}
        \bigl(\log(e+n)\bigr)^{\Gamma_{\mathrm{fast}}}
        \bigl(\log\log(e^e+n)\bigr)^\kappa
    \right)} .
\]
Consequently \(s<2\alpha\) is Spencer-scale.  At \(s=2\alpha\), the same
Spencer-scale conclusion holds if either \(\Gamma_{\mathrm{fast}}<0\), or
\(\Gamma_{\mathrm{fast}}=0\) and \(\kappa\le0\).  If
\(s=2\alpha\) and \(\Gamma_{\mathrm{fast}}>0\), then
\[
  \disc(g_1,\ldots,g_n)\le C\sqrt{n\log\log(e^e+n)} ,
\]
while the borderline case \(\Gamma_{\mathrm{fast}}=0<\kappa\) has the sharper
triple-log loss
\[
  \disc(g_1,\ldots,g_n)\le C\sqrt{n\log\log\log(e^{e^e}+n)} .
\]
\end{corollary}

\begin{proof}
Use the signing scale \(\rho_n\) from
Corollary~\ref{cor:comm-af-exhaustive-finite-scale-atlas}.  Choose
\[
  r_n=\left\lceil C_0\bigl(\log(e+n)\bigr)^{1/\theta}\right\rceil ,
\]
with \(C_0\) large enough, depending on \(H,c,L,\alpha,\theta\), so that
\(\delta_{r_n}\le\rho_n/2\) for every \(n\ge1\).  Then
\[
  r_n^\lambda\le C\bigl(\log(e+n)\bigr)^{\lambda/\theta}
\]
and, for signed \(\kappa\),
\[
  \bigl(\log(e+r_n)\bigr)^\kappa
  \le C\bigl(\log\log(e^e+n)\bigr)^\kappa .
\]
Thus
\[
  A_{r_n}n^{s/(2\alpha)-1}\bigl(\log(e+n)\bigr)^\gamma
  \le
  C n^{s/(2\alpha)-1}
    \bigl(\log(e+n)\bigr)^{\Gamma_{\mathrm{fast}}}
    \bigl(\log\log(e^e+n)\bigr)^\kappa .
\]
Substitution in
Corollary~\ref{cor:comm-af-exhaustive-finite-scale-atlas} gives the displayed
estimate.  The endpoint cases follow by inspecting the logarithm: polynomial
decay in \(\log n\) absorbs any fixed power of \(\log\log n\), positive powers
of \(\log n\) cost \(\log\log n\), and a pure positive power of
\(\log\log n\) costs \(\log\log\log n\).
\end{proof}

\begin{definition}[Clopen Minkowski dimension]
\label{def:clopen-minkowski}
For a compact zero-dimensional metric space \(X\), define the upper clopen
Minkowski dimension by
\[
  \overline{\dim}_{\mathrm{cl}}(X)
  :=
  \limsup_{r\downarrow0}
  \frac{\log M_X(r)}{\log(1/r)},
\]
where \(M_X(r)\) is the clopen partition growth from
Corollary~\ref{cor:comm-af-metric}.
\end{definition}

\begin{definition}[Upper clopen Minkowski content]
\label{def:clopen-content}
For \(s\ge0\), define the upper clopen \(s\)-dimensional Minkowski content by
\[
  \overline{\mathcal M}^{\,s}_{\mathrm{cl}}(X)
  :=
  \limsup_{r\downarrow0} r^s M_X(r),
\]
again with \(M_X(r)\) as in Corollary~\ref{cor:comm-af-metric}.
\end{definition}

\begin{corollary}[Clopen-dimension threshold]
\label{cor:comm-af-dim-threshold}
Let \(X\) be a compact zero-dimensional metric space with
\(\overline{\dim}_{\mathrm{cl}}(X)<\infty\), and let
\(f_1,\ldots,f_n\in C(X)\) be real-valued contractions with common
\((L,\alpha)\)-Hölder modulus.  Then
\[
  \disc(f_1,\ldots,f_n)\le C_{X,L,\alpha}\sqrt{n\log(e+n)} .
\]
If, in addition,
\[
  \overline{\dim}_{\mathrm{cl}}(X)<2\alpha,
\]
then the Spencer-scale bound holds:
\[
  \disc(f_1,\ldots,f_n)\le C_{X,L,\alpha}\sqrt n .
\]
In particular, every compact zero-dimensional metric space of upper clopen
Minkowski dimension \(0\) satisfies the Spencer-scale estimate for every
positive Hölder exponent \(\alpha\).
\end{corollary}

\begin{proof}
If \(D>\overline{\dim}_{\mathrm{cl}}(X)\), then by the definition of the
limsup there are \(C_D<\infty\) and \(r_D>0\) such that
\[
  M_X(r)\le C_D r^{-D}\qquad(0<r\le r_D).
\]
For the finitely many larger scales \(r_D\le r\le\operatorname{diam}X\), the
constant may be enlarged, so the same polynomial-growth hypothesis holds on
\((0,1]\) after rescaling the metric by an absolute factor.  Applying
Corollary~\ref{cor:comm-af-metric} with this exponent \(D\) gives
\[
  \disc(f_1,\ldots,f_n)
  \le C_{D,L,\alpha,X}
     \sqrt{n\log\!\left(e+n^{D/(2\alpha)-1}\right)}.
\]
If \(D\) is finite this is \(O_{X,L,\alpha}(\sqrt{n\log(e+n)})\).  If
\(\overline{\dim}_{\mathrm{cl}}(X)<2\alpha\), choose \(D\) strictly between
these two numbers; then \(D/(2\alpha)-1<0\), so the logarithm is bounded and the
estimate is \(O_{X,L,\alpha}(\sqrt n)\).  The dimension-zero case follows for
every \(\alpha>0\).
\end{proof}

\begin{corollary}[Critical clopen-content endpoint]
\label{cor:comm-af-critical-content}
Let \(X\) be a compact zero-dimensional metric space and let
\(f_1,\ldots,f_n\in C(X)\) be real-valued contractions with common
\((L,\alpha)\)-Hölder modulus.  If
\[
  \overline{\mathcal M}^{\,2\alpha}_{\mathrm{cl}}(X)<\infty,
\]
then the Spencer-scale estimate holds:
\[
  \disc(f_1,\ldots,f_n)\le C_{X,L,\alpha}\sqrt n .
\]
Thus the strict inequality
\(\overline{\dim}_{\mathrm{cl}}(X)<2\alpha\) in
Corollary~\ref{cor:comm-af-dim-threshold} can be replaced at the critical
dimension by finite upper clopen \(2\alpha\)-content.
\end{corollary}

\begin{proof}
Finite upper content gives constants \(C_0<\infty\) and \(r_0>0\) such that
\[
  M_X(r)\le C_0 r^{-2\alpha}\qquad(0<r\le r_0).
\]
Enlarging \(C_0\) to cover the finitely many larger scales, after the same
harmless metric rescaling used in Corollary~\ref{cor:comm-af-dim-threshold},
puts \(X\) under the polynomial clopen-growth hypothesis of
Corollary~\ref{cor:comm-af-metric} with \(d=2\alpha\).  The logarithmic factor
there is \(\log(e+n^{0})=O(1)\), giving the stated \(O_{X,L,\alpha}(\sqrt n)\)
bound.
\end{proof}

\begin{corollary}[Logarithmic critical window]
\label{cor:comm-af-log-window}
Let \(X\) be a compact zero-dimensional metric space, and let
\(f_1,\ldots,f_n\in C(X)\) be real-valued contractions with common
\((L,\alpha)\)-Hölder modulus.  Suppose that for some \(A<\infty\) and
\(\beta\in\mathbb R\),
\[
  M_X(r)\le A r^{-2\alpha}\bigl(\log(e/r)\bigr)^\beta
  \qquad(0<r\le1).
\]
Then
\[
  \disc(f_1,\ldots,f_n)
  \le
  \begin{cases}
  C_{A,L,\alpha,\beta}\sqrt n,
    & \beta\le0,\\[3pt]
  C_{A,L,\alpha,\beta}\sqrt{n\log\log(e^e+n)},
    & \beta>0.
  \end{cases}
\]
Thus a polylogarithmic excess over the critical growth \(r^{-2\alpha}\) costs
only a \(\sqrt{\log\log n}\) factor, while finite or subcritical logarithmic
content remains at Spencer scale.
\end{corollary}

\begin{proof}
Take \(r=\min(1,(L\sqrt n)^{-1/\alpha})\), as in
Corollary~\ref{cor:comm-af-metric}.  Then \(2nLr^\alpha=O_{L,\alpha}(\sqrt n)\)
and
\[
  \frac{M_X(r)}{n}
  \le C_{A,L,\alpha}\bigl(\log(e+n)\bigr)^\beta
\]
after increasing the constant to absorb the finitely many small \(n\).  Insert
this into Corollary~\ref{cor:comm-af-metric}.  If \(\beta\le0\), the logarithm
there is \(O_{A,L,\alpha,\beta}(1)\).  If \(\beta>0\), it is
\[
  O_{A,L,\alpha,\beta}\bigl(\log\log(e^e+n)\bigr),
\]
which gives the displayed estimates.
\end{proof}

\begin{corollary}[Ultrametric and symbolic models]
\label{cor:comm-af-ultrametric}
Let \(X\) be a compact ultrametric space, and let \(N_X(r)\) be the least
number of closed \(r\)-balls needed to cover \(X\).  Suppose
\[
  N_X(r)\le A r^{-d}\qquad(0<r\le1).
\]
If \(f_1,\ldots,f_n\in C(X)\) are real-valued contractions with common
\((L,\alpha)\)-Hölder modulus, then
\[
  \disc(f_1,\ldots,f_n)
  \le C_{A,d,L,\alpha}
     \sqrt{n\log\!\left(e+n^{d/(2\alpha)-1}\right)}.
\]
In particular, \(d\le2\alpha\) gives
\(\disc(f_1,\ldots,f_n)\le C_{A,d,L,\alpha}\sqrt n\).

For the one-sided \(q\)-symbol Cantor shift
\(\Sigma_q=\{1,\ldots,q\}^{\mathbb N}\) with metric
\[
  d_\theta(x,y)=\theta^{\min\{k\ge0:x_k\ne y_k\}}
  \qquad(0<\theta<1),
\]
with \(d_\theta(x,x)=0\), put
\[
  s=\frac{\log q}{\log(1/\theta)}.
\]
Then every common \((L,\alpha)\)-Hölder tuple on \(\Sigma_q\) satisfies
\[
  \disc(f_1,\ldots,f_n)
  \le C_{q,\theta,L,\alpha}
     \sqrt{n\log\!\left(e+n^{s/(2\alpha)-1}\right)},
\]
and the Spencer-scale bound holds whenever \(\alpha\ge s/2\).
\end{corollary}

\begin{proof}
For a fixed \(r>0\), the relation \(x\sim_r y\) iff \(d(x,y)\le r\) partitions
an ultrametric space into closed \(r\)-balls.  Each class has diameter at most
\(r\), and it is clopen: if \(y\) lies in the class, then the open ball
\(\{z:d(y,z)<r\}\) is contained in the same class.  Thus \(M_X(r)\le N_X(r)\),
and the first assertion is Corollary~\ref{cor:comm-af-metric}.

For \(\Sigma_q\), the length-\(k\) cylinder partition has \(q^k\) clopen atoms
of diameter at most \(\theta^k\).  If \(\theta^k\le r<\theta^{k-1}\), then
\[
  M_{\Sigma_q}(r)\le q^k\le q\,r^{-s}.
\]
The previous paragraph gives the displayed bound with \(d=s\).  The endpoint
\(\alpha=s/2\) is included because this is finite upper clopen \(s\)-content,
equivalently by Corollary~\ref{cor:comm-af-critical-content}.
\end{proof}

\begin{corollary}[Finite products add clopen dimension]
\label{cor:comm-af-product}
Let \(X_j\), \(1\le j\le m\), be compact zero-dimensional metric spaces with
clopen partition growth
\[
  M_{X_j}(r)\le A_j r^{-d_j}\qquad(0<r\le1),
\]
and let \(X=\prod_{j=1}^m X_j\) with the max product metric.  Put
\[
  d_*=\sum_{j=1}^m d_j .
\]
If \(f_1,\ldots,f_n\in C(X)\) are real-valued contractions with common
\((L,\alpha)\)-Hölder modulus, then
\[
  \disc(f_1,\ldots,f_n)
  \le C_{\{A_j,d_j\},L,\alpha,m}
     \sqrt{n\log\!\left(e+n^{d_*/(2\alpha)-1}\right)}.
\]
In particular, \(d_*\le2\alpha\) gives the Spencer-scale bound
\[
  \disc(f_1,\ldots,f_n)\le C_{\{A_j,d_j\},L,\alpha,m}\sqrt n .
\]

Consequently, for a finite product of symbolic Cantor shifts
\[
  \prod_{j=1}^m \bigl(\{1,\ldots,q_j\}^{\mathbb N},d_{\theta_j}\bigr)
\]
with the max metric, every common \((L,\alpha)\)-Hölder tuple has
Spencer-scale discrepancy whenever
\[
  \alpha\ge \frac12\sum_{j=1}^m
  \frac{\log q_j}{\log(1/\theta_j)} .
\]
\end{corollary}

\begin{proof}
For each \(j\), choose a clopen partition \(\mathcal P_j\) of \(X_j\) with at
most \(M_{X_j}(r)\) atoms and atom diameters at most \(r\).  Then
\[
  \mathcal P_1\times\cdots\times\mathcal P_m
\]
is a clopen partition of \(X\) whose atoms have max-metric diameter at most
\(r\).  Therefore
\[
  M_X(r)\le\prod_{j=1}^m M_{X_j}(r)
  \le \left(\prod_{j=1}^m A_j\right)r^{-d_*}.
\]
Corollary~\ref{cor:comm-af-metric} gives the first display and the
Spencer-scale consequence.  The symbolic product statement follows by applying
Corollary~\ref{cor:comm-af-ultrametric} to each factor, where
\(d_j=\log q_j/\log(1/\theta_j)\).
\end{proof}

\begin{remark}[Why there is no ``$\Ent(A)$'' for the limit]
Proposition~\ref{prop:vnobstruct} explains the absence of any single entropy
attached to the AF limit $A$ itself. An infinite-dimensional C*-algebra has no
finite $hR$; the only entropies in play are the stage entropies $\Ent(A_k)$, and
Theorem~\ref{thm:af} is forced to be a statement \emph{relative to a chosen
stage}. The supporting-stage entropy, not any intrinsic invariant of $A$, is what
controls the discrepancy. A bounded-block tower has $\Ent(A_k)=O(R_k)$ and is the
favourable case; a UHF tower has $\Ent(A_k)=R_k^2$ and is the boundary case where
the entropy obstruction of Proposition~\ref{prop:obstruct} reappears at every
depth.
\end{remark}

\subsection{Summary statement}

Collecting the favourable reduction and the two obstructions gives the cleanest
theorem the method supports in the limit. It is, honestly, ``reduce to the
supporting finite stage,'' made precise and made stable.

\begin{theorem}[Matrix Spencer for AF algebras, summary form]\label{thm:afsummary}
Let $A=\overline{\bigcup_kA_k}$ be a separable unital AF C*-algebra with stage
entropies $\Ent(A_k)=h_kR_k$, and let $A_1,\ldots,A_n\in A$ be contractions.
\begin{enumerate}[leftmargin=2em]
\item \textup{(Reduction.)} For every stage $k$ and every joint contraction
approximant of error $\delta_k=\delta_k(A_1,\ldots,A_n)$,
\[
\discop(A_1,\ldots,A_n)
\le C\sqrt{n\log\!\left(e+\frac{\Ent(A_k)}{n}\right)}+2n\delta_k,
\]
and one may take the infimum over $k$.
\item \textup{(Spencer regime.)} If the tuple is supported, exactly or to within
$O(n^{-1/2})$, on a stage of entropy $\Ent(A_k)=O(n)$, then
$\discop(A_1,\ldots,A_n)=O(\sqrt n)$. This holds for every tuple in a
bounded-block tower whose support has reduced dimension $O(n)$, and in particular
for every commutative AF algebra with $O(n)$ active minimal projections.
\item \textup{(Intrinsic-dimension upgrade.)} By the Akbas--Sra algebraic
Matrix Spencer theorem, the same \(O(\sqrt n)\) conclusion holds whenever the
tuple is supported, exactly or to within \(O(n^{-1/2})\), on a stage with
\(\dim_\C A_k=O(n)\).  This improves the preceding entropy criterion for
heterogeneous direct-sum stages, but coincides with it on a single full matrix
block.
\item \textup{(Low-Frobenius perturbative upgrade.)} In a finite-dimensional
matrix representation, the preceding intrinsic-dimension AF-stage conclusion is
stable under self-adjoint perturbations \(L_i\) with
\(\|L_i\|_{\op}\le1\) and
\(\|L_i\|_{\mathrm F}^2\le n/\log^3(en)\)
(in particular, under rank-\(n/\log^3(en)\) tails; see
Corollaries~\ref{cor:akbassra-lowrank} and~\ref{cor:afdim-pert}).
\item \textup{(Full-span finite-lift ceiling.)} If the supporting stage contains a full
\(\M_D(\C)\) block on which the compressed tuple spans the whole block, then the
full-span faithful application of Theorem~\ref{thm:main} cannot certify better
than \(\sqrt{n\log(e+D^2/n)}\)
(Proposition~\ref{prop:afceiling}); for full-span tuples in the \(M_{2^k}\)
stage of a UHF tower this becomes the full-stage rate once
\(k>\tfrac12\log_2n\), and improving it would require leaving the full-span
faithful finite-block route or avoiding the entropy dependence.
\item \textup{(Hard obstruction.)} If $A$ is replaced by a non-AF von Neumann
algebra, the finite-stage route can fail completely. For a generic AF element
with infinite spectrum, the exact-stage statement is unavailable and only the
approximate theorem applies, at the entropy of the approximating stages
(Proposition~\ref{prop:vnobstruct}). Quantitatively, those stage entropies must
at least dominate the spectral packing numbers of the element being
approximated (Proposition~\ref{prop:spectralfloor}). In a $\mathrm{II}_1$ factor
a single diffuse self-adjoint contraction already escapes the exact finite-stage
method.
\end{enumerate}
The self-adjoint versions hold for self-adjoint tuples. The constant $C$ is that
of Theorem~\ref{thm:cstar} and is independent of $k$, of $D$, and of the ambient
algebra.
\end{theorem}

\begin{proof}
Parts (1) and (2) are Theorem~\ref{thm:af} with Corollaries~\ref{cor:finitetype}
and~\ref{cor:limit}, together with the bounded-block discussion of
\S\ref{subsec:linear}.  The intrinsic-dimension upgrade is
Corollary~\ref{cor:afdim}, and the low-Frobenius perturbative upgrade is
Corollary~\ref{cor:afdim-pert}.  The full-span finite-lift ceiling is
Proposition~\ref{prop:afceiling}.  The hard obstruction is
Propositions~\ref{prop:vnobstruct} and~\ref{prop:spectralfloor}.
\end{proof}

\begin{remark}[The honest takeaway]
The finite-lift theorem proved here is not a bound attached to an
infinite-dimensional algebra as a single object. What survives the inductive
limit is precisely the finite-stage theorem, made
uniform in the ambient algebra and stable under norm approximation: the
discrepancy of an AF tuple is controlled by the entropy of the finite stage that
supports it, at the same rate $\sqrt{n\log(e+\Ent/n)}$ as in the finite case,
with an approximation error that vanishes faster than $\sqrt n$ whenever the
support is captured at scale $n^{-1/2}$. This stage-entropy bound is
\(O(\sqrt n)\) for bounded-block towers and shallow UHF elements, where the
supporting-stage entropy is $O(n)$.  For exact finite-dimensional algebraic
support, the stronger Akbas--Sra theorem replaces the sufficient entropy
condition by the intrinsic-dimension condition \(\dim_\C A_k=O(n)\).  The
Akbas--Sra perturbative theorem also allows low-Frobenius self-adjoint tails on
top of such a represented AF-stage core; low-rank tails with
\(\operatorname{rank}L_i\le n/\log^3(en)\) are an immediate special case.  The
finite-lift stage-entropy bound still degrades to the generic
\(\sqrt{n\log(e+D^2/n)}\) rate for
full-span tuples/full-unit-ball models, by the same $D^2$ entropy obstruction
(Proposition~\ref{prop:obstruct}) that bars full-unit-ball faithful finite lifts
from the square case of the conjecture in finite dimensions; and the exact
finite-stage theorem disappears for diffuse elements of a von Neumann algebra,
which have no finite supporting stage at all.
\end{remark}

\section*{Acknowledgements and disclosure}
The author is responsible for all mathematical claims and bibliographic
information herein, and has verified the arguments and references.

\begin{thebibliography}{99}

\bibitem{BBvH}
A.~S. Bandeira, M.~T. Boedihardjo, and R. van Handel,
\emph{Matrix concentration inequalities and free probability},
Invent. Math. \textbf{234} (2023), 419--487.

\bibitem{AkbasSra}
E. Akbas and S. Sra,
\emph{An algebraic Matrix Spencer theorem},
arXiv:2606.16005, 2026.

\bibitem{BandeiraBolcskei}
A.~S. Bandeira and H. B\"olcskei,
\emph{Matrix discrepancy for representations of finite groups},
arXiv:2606.12181, 2026.

\bibitem{BJM}
N. Bansal, H. Jiang, and R. Meka,
\emph{Resolving Matrix Spencer conjecture up to poly-logarithmic rank},
SIAM J. Comput., Special Section STOC 2023,
DOI: 10.1137/23M1592201.

\bibitem{BottcherGrudsky}
A. B\"ottcher and S.~M. Grudsky,
\emph{Spectral Properties of Banded Toeplitz Matrices},
SIAM, Philadelphia, 2005.

\bibitem{Bratteli}
O. Bratteli,
\emph{Inductive limits of finite dimensional C*-algebras},
Trans. Amer. Math. Soc. \textbf{171} (1972), 195--234.

\bibitem{CarlssonPerfekt}
M. Carlsson and K.-M. Perfekt,
\emph{Nehari's Theorem for Convex Domain Hankel and Toeplitz Operators in Several Variables},
Int. Math. Res. Not. IMRN \textbf{2021}, no.~5, 3331--3361.

\bibitem{Chazelle}
B. Chazelle,
\emph{The Discrepancy Method: Randomness and Complexity},
Cambridge University Press, 2000.

\bibitem{CourtneySarason}
D. Courtney and D. Sarason,
\emph{A mini-max problem for self-adjoint Toeplitz matrices},
Math. Scand. \textbf{110} (2012), 82--98.

\bibitem{Davidson}
K.~R. Davidson,
\emph{C*-Algebras by Example},
Fields Institute Monographs 6, Amer. Math. Soc., 1996.

\bibitem{DJR}
D. Dadush, H. Jiang, and V. Reis,
\emph{A new framework for matrix discrepancy: partial coloring bounds via mirror descent},
Proceedings of the 54th Annual ACM SIGACT Symposium on Theory of Computing
(STOC), 2022, 649--658; arXiv:2111.03171.

\bibitem{FergusonLacey}
S.~H. Ferguson and M.~T. Lacey,
\emph{A characterization of product BMO by commutators},
Acta Math. \textbf{189} (2002), no.~2, 143--160; arXiv:math/0104144.

\bibitem{LaceyTerwilleger}
M.~T. Lacey and E. Terwilleger,
\emph{Hankel operators in several complex variables and product BMO},
Houston J. Math. \textbf{35} (2009), no.~1, 159--183.

\bibitem{LRR}
A. Levy, H. Ramadas, and T. Rothvoss,
\emph{Deterministic discrepancy minimization via the multiplicative weight update method},
Integer Programming and Combinatorial Optimization (IPCO 2017), LNCS 10328,
380--391; arXiv:1611.08752.

\bibitem{Gluskin}
E.~D. Gluskin,
\emph{Extremal properties of orthogonal parallelepipeds and their applications to the geometry of Banach spaces},
Math. USSR-Sb. \textbf{64} (1989), 85--96.

\bibitem{HRS}
S.~B. Hopkins, P. Raghavendra, and A. Shetty,
\emph{Matrix discrepancy from quantum communication},
Proceedings of the 54th Annual ACM SIGACT Symposium on Theory of Computing
(STOC), 2022, 637--648.

\bibitem{KuniskyZhang}
D. Kunisky and P. Zhang,
\emph{Average-case matrix discrepancy: asymptotics and online algorithms},
arXiv:2307.10055, 2024.

\bibitem{LustPisier}
F. Lust-Piquard and G. Pisier,
\emph{Non-commutative Khintchine and Paley inequalities},
Ark. Mat. \textbf{29} (1991), 241--260.

\bibitem{Maillard}
A. Maillard,
\emph{Average-case matrix discrepancy: satisfiability bounds},
arXiv:2410.17887, 2024.

\bibitem{Matousek}
J. Matou\v{s}ek,
\emph{Geometric Discrepancy: An Illustrated Guide},
Springer, 1999.

\bibitem{NikolskayaFarforovskaya}
L.~N. Nikolskaya and Yu.~B. Farforovskaya,
\emph{Toeplitz and Hankel matrices as Hadamard--Schur multipliers},
Algebra i Analiz \textbf{15} (2003), no.~6, 141--160.

\bibitem{OrtegaCerdaSeip}
J. Ortega-Cerd\`a and K. Seip,
\emph{A lower bound in Nehari's theorem on the polydisc},
J. Anal. Math. \textbf{118} (2012), 339--342; arXiv:1107.0175.

\bibitem{Pisier}
G. Pisier,
\emph{Introduction to Operator Space Theory},
Cambridge University Press, 2003.

\bibitem{Spencer}
J. Spencer,
\emph{Six standard deviations suffice},
Trans. Amer. Math. Soc. \textbf{289} (1985), 679--706.

\bibitem{Tropp}
J.~A. Tropp,
\emph{User-friendly tail bounds for sums of random matrices},
Found. Comput. Math. \textbf{12} (2012), no.~4, 389--434.

\bibitem{KOWZ}
D. Kunisky, T. Oertel, N. Wengiel, and P. Zhang,
\emph{Asymptotic bounds and online algorithms for average-case matrix discrepancy},
arXiv:2410.23915v2, 2025.

\bibitem{Shmalo-paper5}
Y. Shmalo,
\emph{A parity obstruction to the stability of random banded Toeplitz matrices},
2026.

\end{thebibliography}

\end{document}
